% Auto-generated from data/segments.json + layout.json (+ transforms) — do not edit by hand.
% volume: 05Vin05
% mode: printing
% segments: 2822
% transform_rules: 8
% toc_marks: 508

% Volume layout defaults (layout.json ``layout``).
\csromanlayoutapply{1}{1.5}{21.6pt}{8pt}{8pt}{65pt}{2.5em}%

\pitaka{\textbf{วินยปิฏก}}
\csromanmatikaanchor{mtk.0}
\csromantocmarknopage{chapter}{ปริวารปาฬิ}{mtk.0}
\bookmark[dest={mtk.0},level=0]{ปริวารปาฬิ}
\gambhira{\textbf{ปริวารปาฬิ}}
\namakkaram{นโม ตสฺส ภควโต อรหโต สมฺมา\-สมฺพุทฺธสฺส.}
\csromanmatikaanchor{mtk.1}
\csromantocmarknopage{chapter}{ภิกฺขุวิภงฺค}{mtk.1}
\bookmark[dest={mtk.1},level=0]{ภิกฺขุวิภงฺค}
\chapterhead{ภิกฺขุ\-วิภงฺค}
\csromanmatikaanchor{mtk.2}
\csromanmatikaanchor{mtk.3}
\csromantocmarknopage{chapter}{1. กตฺถปญฺญตฺติวาร}{mtk.2}
\bookmark[dest={mtk.2},level=0]{1. กตฺถปญฺญตฺติวาร}
\csromantocmarkref{chapter}{1. ปาราชิกกณฺฑ}{mtk.3}
\bookmark[dest={mtk.3},level=0]{1. ปาราชิกกณฺฑ}
\chapterhead{1. กตฺถ\-ปญฺญตฺติวาร 1. ปาราชิกกณฺฑ}
\proseitem{1}{\csromanfolio{1}ยํ เตน ภควตา ชานตา ปสฺสตา อรหตา สมฺมา\-สมฺพุทฺเธน ปฐมํ ปาราชิกํ กตฺถ ปญฺญตฺตํ, กํ อารพฺภ, กิสฺมิํ วตฺถุสฺมิํ, อตฺถิ ตตฺถ ปญฺญตฺติ อนุปญฺญตฺติ อนุปฺปนฺนปญฺญตฺติ, สพฺพตฺถ\-ปญฺญตฺติ ปเทสปญฺญตฺติ, สาธารณ\-ปญฺญตฺติ อสาธารณ\-ปญฺญตฺติ, เอกโตปญฺญตฺติ อุภโตปญฺญตฺติ, ปญฺจนฺนํ ปาติโมกฺขุ\-ทฺเทสานํ กตฺโถคธํ กตฺถ ปริยาปนฺนํ, กตเมน อุทฺเทเสน อุทฺเทสํ อาคจฺฉติ, จตุนฺนํ วิปตฺตีนํ กตมา วิปตฺติ, สตฺตนฺนํ อาปตฺติ\-กฺขนฺธานํ กตโม อาปตฺติกฺขนฺโธ, ฉนฺนํ อาปตฺติ\-สมุฏฺฐานานํ กติหิ สมุฏฺฐาเนหิ สมุฏฺฐาติ, จตุนฺนํ อธิกรณานํ กตมํ อธิกรณํ, สตฺตนฺนํ สมถานํ กติหิ สมเถหิ สมฺมติ, โก ตตฺถ วินโย, โก ตตฺถ อภิวินโย, กิํ ตตฺถ ปาติโมกฺขํ, กิํ ตตฺถ อธิปาติโมกฺขํ, กา วิปตฺติ, กา สมฺปตฺติ กา ปฏิปตฺติ, กติ อตฺถวเส ปฏิจฺจ ภควตา ปฐมํ ปาราชิกํ ปญฺญตฺตํ, เก สิกฺขนฺติ, เก สิกฺขิตสิกฺขา, กตฺถ ฐิตํ, เก ธาเรนฺติ, กสฺส วจนํ, เกนาภตนฺติ?}

\proseitem{2}{ยํ เตน ภควตา ชานตา ปสฺสตา อรหตา สมฺมา\-สมฺพุทฺเธน ปฐมํ ปาราชิกํ กตฺถ ปญฺญตฺตนฺติ เวสาลิยํ ปญฺญตฺตํ. กํ อารพฺภาติ สุทินฺนํ กลนฺทปุตฺตํ อารพฺภ. กิสฺมิํ วตฺถุสฺมินฺติ สุทินฺโน กลนฺทปุตฺโต ปุราณ\-ทุติยิกาย เมถุนํ ธมฺมํ ปฏิเสวิ, ตสฺมิํ วตฺถุสฺมิํ. อตฺถิ ตตฺถ ปญฺญตฺติ อนุปญฺญตฺติ อนุปฺปนฺนปญฺญตฺตีติ เอกา ปญฺญตฺติ, เทฺว อนุปญฺญตฺติโย, \csromanfolio{2}อนุปฺปนฺนปญฺญตฺติ ตสฺมิํ นตฺถิ. สพฺพตฺถ\-ปญฺญตฺติ ปเทสปญฺญตฺตีติ สพฺพตฺถ\-ปญฺญตฺติ. สาธารณ\-ปญฺญตฺติ อสาธารณปญฺญตฺตีติ สาธารณ\-ปญฺญตฺติ. เอกโตปญฺญตฺติ อุภโตปญฺญตฺตีติ อุภโตปญฺญตฺติ. ปญฺจนฺนํ ปาติโมกฺขุ\-ทฺเทสานํ กตฺโถคธํ กตฺถ ปริยาปนฺนนฺติ นิทาโนคธํ นิทาน\-ปริยาปนฺนํ. กตเมน อุทฺเทเสน อุทฺเทสํ อาคจฺฉตีติ ทุติเยน อุทฺเทเสน อุทฺเทสํ อาคจฺฉติ. จตุนฺนํ วิปตฺตีนํ กตมา วิปตฺตีติ สีลวิปตฺติ. สตฺตนฺนํ อาปตฺติ\-กฺขนฺธานํ กตโม อาปตฺติ\-กฺขนฺโธติ ปาราชิกาปตฺติ\-กฺขนฺโธ. ฉนฺนํ อาปตฺติ\-สมุฏฺฐานานํ กติหิ สมุฏฺฐาเนหิ สมุฏฺฐาตีติ เอเกน สมุฏฺฐาเนน สมุฏฺฐาติ, กายโต จ จิตฺตโต จ สมุฏฺฐาติ, น วาจโต. จตุนฺนํ อธิกรณานํ กตมํ อธิกรณนฺติ อาปตฺตา\-ธิกรณํ. สตฺตนฺนํ สมถานํ กติหิ สมเถหิ สมฺมตีติ ทฺวีหิ สมเถหิ สมฺมติ สมฺมุขา\-วินเยนจ ปฏิญฺญาต\-กรเณนจ. โก ตตฺถ วินโย โก ตตฺถ อภิวินโยติ ปญฺญตฺติ วินโย, วิภตฺติ อภิวินโย. กิํ ตตฺถ ปาติโมกฺขํ กิํ ตตฺถ อธิปาติ\-โมกฺขนฺติ ปญฺญตฺติ ปาติโมกฺขํ, วิภตฺติ อธิปาติโมกฺขํ. กาวิปตฺตีติ อสํวโร วิปตฺติ. กา สมฺปตฺตีติ สํวโร สมฺปตฺติ. กา ปฏิปตฺตีติ “น เอวรูปํ กริสฺสามี”ติ ยาวชีวํ อาปาณโกฏิกํ สมาทาย สิกฺขติ สิกฺขาปเทสุ. * กติ อตฺถวเส ปฏิจฺจ ภควตา ปฐมํ ปาราชิกํ ปญฺญตฺตนฺติ ทส อตฺถวเส ปฏิจฺจ ภควตา ปฐมํ ปาราชิกํ ปญฺญตฺตํ สํฆสุฏฺฐุตาย สํฆผาสุตาย ทุมฺมงฺกูนํ ปุคฺคลานํ นิคฺคหาย เปสลานํ ภิกฺขูนํ ผาสุวิหาราย ทิฏฺฐ\-ธมฺมิกานํ อาสวานํ สํวราย สมฺปรายิกานํ อาสวานํ ปฏิฆาตาย อปฺปสนฺนานํ ปสาทาย ปสนฺนานํ ภิยฺโยภาวาย สทฺธมฺม\-ฏฺฐิติยา วินยา\-นุคฺคหาย. เก สิกฺขนฺตีติ เสกฺขา จ ปุถุชฺชน\-กลฺยาณกา จ สิกฺขนฺติ. เก สิกฺขิต\-สิกฺขาติ อรหนฺโต สิกฺขิตสิกฺขา. กตฺถ ฐิตนฺติ สิกฺขากาเมสุ ฐิตํ. เก ธาเรนฺตีติ เยสํ วตฺตติ, เต ธาเรนฺติ. กสฺส วจนนฺติ ภควโต วจนํ อรหโต สมฺมา\-สมฺพุทฺธสฺส. เกนาภตนฺติ ปรมฺปราภตํ.}

\proseitem{3}{อุปาลิ ทาสโก เจว, โสณโก สิคฺคโวตถา.}

\prose{โมคฺคลิปุตฺเตน\csromansharedfootnote{da111f94614764c2}{โมคฺคลีปุตฺเตน (สฺยา)} ปญฺจมา, เอเต ชมฺพุสิริวฺหเย\csromansharedfootnote{ae3f8709600149a1}{ชมฺพุสริวฺหเย (สารตฺถ)}.}

\prose{ตโต มหินฺโท อิฏฺฏิโย\csromansharedfootnote{93435dd17a03a9c4}{อิฏฺฐิโย (สี)}, อุตฺติโย สมฺพโล\csromansharedfootnote{ac705b8ecad1c4d7}{อุฏฺฏิโย สมฺปโล (อิติปิ)} ตถา.}

\prose{\csromanfolio{3}วินยํ เต วาจยิํสุ, ปิฏกํ ตมฺพปณฺณิยา.}

\prose{นิกาเย ปญฺจ วาเจสุํ, สตฺต เจว ปกรเณ.}

\prose{ตโต อริฏฺโฐ เมธาวี, ติสฺสทตฺโต จปณฺฑิโต.}

\prose{วิสารโท กาฬสุมโน, เถโร จ ทีฆนามโก.}

\prose{ปุนเทว\csromansharedfootnote{72ffd47748444be5}{ปุนเรว (สฺยา)} กาฬสุมโน, นาคตฺเถโร จ พุทฺธรกฺขิโต.}

\prose{ติสฺสตฺเถโร จ เมธาวี, เทวตฺเถโร\csromansharedfootnote{b06395dcef8a432d}{เรวตฺเถโร (อิติปิ)} จ ปณฺฑิโต.}

\prose{ปุนเทว สุมโน เมธาวี, วินเย จ วิสารโท.}

\prose{พหุสฺสุโต จูฬนาโค, คโชว ทุปฺปธํสิโย.}

\prose{ธมฺมปาลิต\-นาโม จ, โรหเณ\csromansharedfootnote{068dead288b081c6}{โรหโณ (สี)} สาธุปูชิโต.}

\prose{ตสฺส สิสฺโส มหาปญฺโญ, เขมนาโม ติเปฏโก.}

\prose{ทีเป ตารกราชาว, ปญฺญาย อติโรจถ.}

\prose{อุปติสฺโส จ เมธาวี, ผุสฺสเทโว มหากถี.}

\prose{ปุนเทว สุมโน เมธาวี, ปุปฺผนาโม พหุสฺสุโต.}

\prose{มหากถี มหาสิโว, ปิฏเก สพฺพตฺถ โกวิโท.}

\prose{ปุนเทว อุปาลิ เมธาวี, วินเย จ วิสารโท.}

\prose{มหานาโค มหาปญฺโญ, สทฺธมฺม\-วํส\-โกวิโท.}

\prose{ปุนเทว อภโย เมธาวี, ปิฏเก สพฺพตฺถ โกวิโท.}

\prose{ติสฺสตฺเถโร จ เมธาวี, วินเย จ วิสารโท.}

\prose{ตสฺส สิสฺโส มหาปญฺโญ, ปุปฺผนาโม พหุสฺสุโต.}

\prose{สาสนํ อนุรกฺขนฺโต, ชมฺพุทีเป ปติฏฺฐิโต.}

\prose{จูฬาภโย จ เมธาวี, วินเย จ วิสารโท.}

\prose{ติสฺสตฺเถโร จ เมธาวี, สทฺธมฺม\-วํส\-โกวิโท.}

\prose{จูฬเทโว\csromansharedfootnote{71492dce967fb530}{ผุสฺสเทโว (สี)} จ เมธาวี, วินเย จ วิสารโท,}

\prose{สิวตฺเถโร จ เมธาวี, วินเย สพฺพตฺถ โกวิโท.}

\prose{\csromanfolio{4}เอเต นาคา มหาปญฺญา, วินยญฺญู มคฺคโกวิทา.}

\prose{วินยํ ทีเป ปกาเสสุํ, ปิฏกํ ตมฺพปณฺณิยาติ.}

\proseitem{4}{ยํ เตน ภควตา ชานตา ปสฺสตา อรหตา สมฺมา\-สมฺพุทฺเธน ทุติยํ ปาราชิกํ กตฺถ ปญฺญตฺตนฺติ ราชคเห ปญฺญตฺตํ. กํ อารพฺภาติ ธนิยํ กุมฺภการ\-ปุตฺตํ อารพฺภ. กิสฺมิํ วตฺถุสฺมินฺติ ธนิโย กุมฺภการ\-ปุตฺโต รญฺโญ ทารูนิ อทินฺนํ อาทิยิ, ตสฺมิํ วตฺถุสฺมิํ. เอกา ปญฺญตฺติ, เอกา อนุปญฺญตฺติ. ฉนฺนํ อาปตฺติ\-สมุฏฺฐานานํ กติหิ สมุฏฺฐาเนหิ สมุฏฺฐาตีติ ตีหิ สมุฏฺฐาเนหิ สมุฏฺฐาติ, สิยา กายโต จ จิตฺตโต จ สมุฏฺฐาติ, น วาจโต, สิยา วาจโต จ จิตฺตโต จ สมุฏฺฐาติ, น กายโต, สิยา กายโต จ วาจโต จ จิตฺตโต จ สมุฏฺฐาติ ฯเปฯ.}

\proseitem{5}{ตติยํ ปาราชิกํ กตฺถ ปญฺญตฺตนฺติ เวสาลิยํ ปญฺญตฺตํ. กํ อารพฺภาติ สมฺพหุเล ภิกฺขู อารพฺภ. กิสฺมิํ วตฺถุสฺมินฺติ สมฺพหุลา ภิกฺขู อญฺญมญฺญํ ชีวิตา โวโรเปสุํ, ตสฺมิํ วตฺถุสฺมิํ. เอกา ปญฺญตฺติ, เอกา อนุปญฺญตฺติ. ฉนฺนํ อาปตฺติ\-สมุฏฺฐานานํ กติหิ สมุฏฺฐาเนหิ สมุฏฺฐาตีติ ตีหิ สมุฏฺฐาเนหิ สมุฏฺฐาติ, สิยา กายโต จ จิตฺตโต จ สมุฏฺฐาติ, น วาจโต, สิยา วาจโต จ จิตฺตโต จ สมุฏฺฐาติ, น กายโต, สิยา กายโต จ วาจโต จ จิตฺตโต จ สมุฏฺฐาติ ฯเปฯ.}

\proseitemwithcloser{6}{จตุตฺถํ ปาราชิกํ กตฺถ ปญฺญตฺตนฺติ เวสาลิยํ ปญฺญตฺตํ. กํ อารพฺภาติ วคฺคุมุทา\-ตีริเย ภิกฺขู อารพฺภ. กิสฺมิํ วตฺถุสฺมินฺติ วคฺคุมุทา\-ตีริยา ภิกฺขู คิหีนํ อญฺญมญฺญสฺส อุตฺตริ\-มนุสฺสธมฺมสฺส วณฺณํ ภาสิํสุ, ตสฺมิํ วตฺถุสฺมิํ. เอกา ปญฺญตฺติ, เอกา อนุปญฺญตฺติ. ฉนฺนํ อาปตฺติ\-สมุฏฺฐานานํ กติหิ สมุฏฺฐาเนหิ สมุฏฺฐาตีติ ตีหิ สมุฏฺฐาเนหิ สมุฏฺฐาติ, สิยา กายโต จ จิตฺตโต จ สมุฏฺฐาติ, น วาจโต, สิยา วาจโต จ จิตฺตโต จ สมุฏฺฐาติ, น กายโต, สิยา กายโต จ วาจโต จ จิตฺตโต จ}{จตฺตาโร ปาราชิกา นิฏฺฐิตา.}
\csromancenter{ตสฺสุทฺทานํ.}
\setlength{\csromangathapagewidth}{0pt}
\setlength{\csromangathawakwidth}{0pt}
\csromangathameasuregroup{เมถุนาทินฺนาทานญฺจ,}{\csromangathabat{เมถุนาทินฺนาทานญฺจ,}{มนุสฺสวิคฺคหุตฺตริ.}}
\csromangathasetleft{เมถุนาทินฺนาทานญฺจ,}
\csromangathagroup{\csromangathastanza{\csromanfolio{5}\csromangathabat{เมถุนาทินฺนาทานญฺจ,}{มนุสฺสวิคฺคหุตฺตริ.}}}

\prose{ปาราชิกานิ จตฺตาริ, เฉชฺชวตฺถู อสํสยาติ.}
\csromansectionrule

\csromanmatikaanchor{mtk.4}
\csromantocmarkref{chapter}{2. สํฆาทิเสสกณฺฑ}{mtk.4}
\bookmark[dest={mtk.4},level=0]{2. สํฆาทิเสสกณฺฑ}
\chapterhead{2. สํฆาทิเสส\-กณฺฑ}
\proseitem{7}{ยํ เตน ภควตา ชานตา ปสฺสตา อรหตา สมฺมา\-สมฺพุทฺเธน อุปกฺกมิตฺวา อสุจิํ โมเจนฺตสฺส สํฆาทิเสโส กตฺถ ปญฺญตฺโต, กํ อารพฺภ, กิสฺมิํ วตฺถุสฺมิํ, อตฺถิ ตตฺถ ปญฺญตฺติ อนุปญฺญตฺติ อนุปฺปนฺนปญฺญตฺติ, สพฺพตฺถ\-ปญฺญตฺติ ปเทสปญฺญตฺติ, สาธารณ\-ปญฺญตฺติ อสาธารณ\-ปญฺญตฺติ, เอกโตปญฺญตฺติ อุภโตปญฺญตฺติ, ปญฺจนฺนํ ปาติโมกฺขุ\-ทฺเทสานํ กตฺโถคธํ กตฺถ ปริยาปนฺนํ, กตเมน อุทฺเทเสน อุทฺเทสํ อาคจฺฉติ, จตุนฺนํ วิปตฺตีนํ กตมา วิปตฺติ, สตฺตนฺนํ อาปตฺติ\-กฺขนฺธานํ กตโม อาปตฺติกฺขนฺโธ, ฉนฺนํ อาปตฺติ\-สมุฏฺฐานานํ กติหิ สมุฏฺฐาเนหิ สมุฏฺฐาติ, จตุนฺนํ อธิกรณานํ กตมํ อธิกรณํ, สตฺตนฺนํ สมถานํ กติหิ สมเถหิ สมฺมติ, โก ตตฺถ วินโย, โก ตตฺถ อภิวินโย, กิํ ตตฺถ ปาติโมกฺขํ, กิํ ตตฺถ อธิปาติโมกฺขํ, กา วิปตฺติ, กา สมฺปตฺติ, กา ปฏิปตฺติ, กติ อตฺถวเส ปฏิจฺจ ภควตา อุปกฺกมิตฺวา อสุจิํ โมเจนฺตสฺส สํฆาทิเสโส ปญฺญตฺโต, เก สิกฺขนฺติ, เก สิกฺขิตสิกฺขา, กตฺถ ฐิตํ, เก ธาเรนฺติ, กสฺส วจนํ, เกนาภตนฺติ.}

\proseitem{8}{ยํ เตน ภควตา ชานตา ปสฺสตา อรหตา สมฺมา\-สมฺพุทฺเธน อุปกฺกมิตฺวา อสุจิํ โมเจนฺตสฺส สํฆาทิเสโส กตฺถ ปญฺญตฺโตติ สาวตฺถิยํ ปญฺญตฺโต. กํ อารพฺภาติ อายสฺมนฺตํ เสยฺยสกํ อารพฺภ. กิสฺมิํ วตฺถุสฺมินฺติ อายสฺมา เสยฺยสโก หตฺเถน อุปกฺกมิตฺวา อสุจิํ โมเจสิ, ตสฺมิํ วตฺถุสฺมิํ. อตฺถิ ตตฺถ ปญฺญตฺติ อนุปญฺญตฺติ อนุปฺปนฺนปญฺญตฺตีติ เอกา ปญฺญตฺติ, เอกา อนุปญฺญตฺติ, อนุปฺปนฺนปญฺญตฺติ ตสฺมิํ นตฺถิ. สพฺพตฺถ\-ปญฺญตฺติ ปเทสปญฺญตฺตีติ สพฺพตฺถ\-ปญฺญตฺติ. สาธารณ\-ปญฺญตฺติ อสาธารณปญฺญตฺตีติ อสาธารณ\-ปญฺญตฺติ. เอกโตปญฺญตฺติ อุภโตปญฺญตฺตีติ เอกโตปญฺญตฺติ. ปญฺจนฺนํ ปาติโมกฺขุ\-ทฺเทสานํ กตฺโถคธํ กตฺถ ปริยาปนฺนนฺติ นิทาโนคธํ นิทาน\-ปริยาปนฺนํ. กตเมน อุทฺเทเสน อุทฺเทสํ อาคจฺฉตีติ \csromanfolio{6}ตติเยน อุทฺเทเสน อุทฺเทสํ อาคจฺฉติ. จตุนฺนํ วิปตฺตีนํ กตมา วิปตฺตีติ สีลวิปตฺติ. สตฺตนฺนํ อาปตฺติ\-กฺขนฺธานํ กตโม อาปตฺติ\-กฺขนฺโธติ สํฆาทิเสโส อาปตฺติกฺขนฺโธ. ฉนฺนํ อาปตฺติ\-สมุฏฺฐานานํ กติหิ สมุฏฺฐาเนหิ สมุฏฺฐาตีติ เอเกน สมุฏฺฐาเนน สมุฏฺฐาติ, กายโต จ จิตฺตโต จ สมุฏฺฐาติ, น วาจโต. จตุนฺนํ อธิกรณานํ กตมํ อธิกรณนฺติ อาปตฺตา\-ธิกรณํ. สตฺตนฺนํ สมถานํ กติหิ สมเถหิ สมฺมตีติ ทฺวีหิ สมเถหิ สมฺมติ สมฺมุขา\-วินเยน จ ปฏิญฺญาต\-กรเณน จ. โก ตตฺถ วินโย โก ตตฺถ อภิวินโยติ ปญฺญตฺติ วินโย, วิภตฺติ อภิวินโย. กิํ ตตฺถ ปาติโมกฺขํ กิํ ตตฺถ อธิปาติ\-โมกฺขนฺติ ปญฺญตฺติ ปาติโมกฺขํ, วิภตฺติ อธิปาติโมกฺขํ. กา วิปตฺตีติ อสํวโร วิปตฺติ. กา สมฺปตฺตีติ สํวโร สมฺปตฺติ. กา ปฏิปตฺตีติ “น เอวรูปํ กริสฺสามี”ติ ยาวชีวํ อาปาณโกฏิกํ สมาทาย สิกฺขติ สิกฺขาปเทสุ. * กติ อตฺถวเส ปฏิจฺจ ภควตา อุปกฺกมิตฺวา อสุจิํ โมเจนฺตสฺส สํฆาทิเสโส ปญฺญตฺโตติ ทส อตฺถวเส ปฏิจฺจ ภควตา อุปกฺกมิตฺวา อสุจิํ โมเจนฺตสฺส สํฆาทิเสโส ปญฺญตฺโต สํฆสุฏฺฐุตาย สํฆผาสุตาย ทุมฺมงฺกูนํ ปุคฺคลานํ นิคฺคหาย เปสลานํ ภิกฺขูนํ ผาสุวิหาราย ทิฏฺฐ\-ธมฺมิกานํ อาสวานํ สํวราย สมฺปรายิกานํอาสวานํ ปฏิฆาตาย อปฺปสนฺนานํ ปสาทาย ปสนฺนานํ ภิยฺโยภาวาย สทฺธมฺม\-ฏฺฐิติยา วินยา\-นุคฺคหาย. เก สิกฺขนฺตีติ เสกฺขา จ ปุถุชฺชน\-กลฺยาณกาจ สิกฺขนฺติ. เก สิกฺขิต สิกฺขาติ อรหนฺโต สิกฺขิตสิกฺขา.กตฺถ ฐิตนฺติ สิกฺขากาเมสุ ฐิตํ. เก ธาเรนฺตีติ เยสํ วตฺตติ, เต ธาเรนฺติ. กสฺส วจนนฺติ ภควโต วจนํ อรหโต สมฺมา\-สมฺพุทฺธสฺส. เกนาภตนฺติ ปรมฺปราภตํ.}

\prose{อุปาลิ ทาสโก เจว, โสณโก สิคฺคโว ตถา.}

\prose{โมคฺคลิปุตฺเตน ปญฺจมา, เอเต ชมฺพุสิริวฺหเย.}

\prose{ตโต มหินฺโท อิฏฺฏิโย, อุตฺติโย สมฺพโล ตถา.}

\prose{วินยํ เต วาจยิํสุ, ปิฏกํ ตมฺพปณฺณิยา.}

\prose{นิกาเย ปญฺจ วาเจสุํ, สตฺต เจว ปกรเณ.}

\prose{ตโต อริฏฺโฐ เมธาวี, ติสฺสทตฺโต จ ปณฺฑิโต.}

\prose{\csromanfolio{7}วิสารโท กาฬสุมโน, เถโร จ ทีฆนามโก.}

\prose{ปุนเทว กาฬสุมโน, นาคตฺเถโร จ พุทฺธรกฺขิโต.}

\prose{ติสฺสตฺเถโร จ เมธาวี, เทวตฺเถโร จ ปณฺฑิโต.}

\prose{ปุนเทว สุมโน เมธาวี, วินเย จ วิสารโท.}

\prose{พหุสฺสุโต จูฬนาโค, คโชว ทุปฺปธํสิโย.}

\prose{ธมฺมปาลิต\-นาโม จ, โรหเณ สาธุปูชิโต.}

\prose{ตสฺส สิสฺโส มหาปญฺโญ, เขมนาโม ติเปฏโก.}

\prose{ทีเป ตารกราชาว, ปญฺญาย อติโรจถ.}

\prose{อุปติสฺโส จ เมธาวี, ผุสฺสเทโว มหากถี.}

\prose{ปุนเทว สุมโน เมธาวี, ปุปฺผนาโม พหุสฺสุโต.}

\prose{มหากถี มหาสิโว, ปิฏเก สพฺพตฺถ โกวิโท.}

\prose{ปุนเทว อุปาลิ เมธาวี, วินเย จ วิสารโท.}

\prose{มหานาโค มหาปญฺโญ, สทฺธมฺม\-วํส\-โกวิโท.}

\prose{ปุนเทว อภโย เมธาวี, ปิฏเก สพฺพตฺถ โกวิโท.}

\prose{ติสฺสตฺเถโร จ เมธาวี, วินเย จ วิสารโท.}

\prose{ตสฺส สิสฺโส มหาปญฺโญ, ปุปฺผนาโม พหุสฺสุโต.}

\prose{สาสนํ อนุรกฺขนฺโต, ชมฺพุทีเป ปติฏฺฐิโต.}

\prose{จูฬาภโย จ เมธาวี, วินเย จ วิสารโท.}

\prose{ติสฺสตฺเถโร จ เมธาวี, สทฺธมฺม\-วํส\-โกวิโท.}

\prose{จูฬเทโว จ เมธาวี, วินเย จ วิสารโท.}

\prose{สิวตฺเถโร จ เมธาวี, วินเย สพฺพตฺถ โกวิโท.}

\prose{เอเต นาคา มหาปญฺญา, วินยญฺญู มคฺคโกวิทา.}

\prose{วินยํ ทีเป ปกาเสสุํ, ปิฏกํ ตมฺพปณฺณิยาติ.}

\proseitem{9}{ยํ เตน ภควตา ชานตา ปสฺสตา อรหตา สมฺมา\-สมฺพุทฺเธน มาตุคาเมน สทฺธิํ กายสํสคฺคํ สมา\-ป\-ชฺช\-นฺตสฺส สํฆาทิเสโส กตฺถ ปญฺญตฺโตติ สาวตฺถิยํ ปญฺญตฺโต. กํ อารพฺภาติ อายสฺมนฺตํ อุทายิํ \csromanfolio{8}อารพฺภ. กิสฺมิํ วตฺถุสฺมินฺติ อายสฺมา อุทายี มาตุคาเมน สทฺธิํ กายสํสคฺคํ สมาปชฺชิ, ตสฺมิํ วตฺถุสฺมิํ. เอกา ปญฺญตฺติ. ฉนฺนํ อาปตฺติ\-สมุฏฺฐานานํ เอเกน สมุฏฺฐาเนน สมุฏฺฐาติ, กายโต จ จิตฺตโต จ สมุฏฺฐาติ, น วาจโต ฯเปฯ.}

\proseitem{10}{มาตุคามํ ทุฏฺฐุลฺลาหิ วาจาหิ โอภาสนฺตสฺส สํฆาทิเสโส กตฺถ ปญฺญตฺโตติ สาวตฺถิยํ ปญฺญตฺโต. กํ อารพฺภาติ อายสฺมนฺตํ อุทายิํ อารพฺภ. กิสฺมิํ วตฺถุสฺมินฺติ อายสฺมา อุทายี มาตุคามํ ทุฏฺฐุลฺลาหิ วาจาหิ โอภาสิ, ตสฺมิํ วตฺถุสฺมิํ. เอกา ปญฺญตฺติ. ฉนฺนํ อาปตฺติ\-สมุฏฺฐานานํ ตีหิ สมุฏฺฐาเนหิ สมุฏฺฐาติ, สิยา กายโต จ จิตฺตโต จ สมุฏฺฐาติ, น วาจโต, สิยา วาจโต จ จิตฺตโต จ สมุฏฺฐาติ, น กายโต, สิยา กายโต จ วาจโต จจิตฺตโต จ สมุฏฺฐาติ ฯเปฯ.}

\proseitem{11}{มาตุคามสฺส สนฺติเก อตฺต\-กาม\-ปาริจริยาย วณฺณํ ภาสนฺตสฺส สํฆาทิเสโส กตฺถ ปญฺญตฺโตติ สาวตฺถิยํ ปญฺญตฺโต. กํ อารพฺภาติ อายสฺมนฺตํ อุทายิํ อารพฺภ. กิสฺมิํ วตฺถุสฺมินฺติ อายสฺมา อุทายี มาตุคามสฺส สนฺติเก อตฺต\-กาม\-ปาริจริยาย วณฺณํ อภาสิ, ตสฺมิํ วตฺถุสฺมิํ. เอกา ปญฺญตฺติ. ฉนฺนํ อาปตฺติ\-สมุฏฺฐานานํ ตีหิ สมุฏฺฐาเนหิ สมุฏฺฐาติ, สิยา กายโต จ จิตฺตโต จ สมุฏฺฐาติ, น วาจโต, สิยา วาจโต จ จิตฺตโต จ สมุฏฺฐาติ, น กายโต, สิยา กายโต จ วาจโต จ จิตฺตโต จ}

\proseitem{12}{สญฺจริตฺตํ สมา\-ป\-ชฺช\-นฺตสฺส สํฆาทิเสโส กตฺถ ปญฺญตฺโตติ สาวตฺถิยํ ปญฺญตฺโต. กํ อารพฺภาติ อายสฺมนฺตํ อุทายิํ อารพฺภ. กิสฺมิํ วตฺถุสฺมินฺติ อายสฺมา อุทายี สญฺจริตฺตํ สมาปชฺชิ, ตสฺมิํ วตฺถุสฺมิํ. เอกา ปญฺญตฺติ, เอกา อนุปญฺญตฺติ. ฉนฺนํ อาปตฺติ\-สมุฏฺฐานานํ ฉหิ สมุฏฺฐาเนหิ สมุฏฺฐาติ, สิยา กายโต สมุฏฺฐาติ, น วาจโต, น จิตฺตโต,สิยา วาจโต สมุฏฺฐาติ, น กายโต, น จิตฺตโต, สิยา กายโต จ วาจโต จ สมุฏฺฐาติ, น จิตฺตโต, สิยา กายโต จ จิตฺตโต จ สมุฏฺฐาติ, น วาจโต, สิยา วาจโต จ จิตฺตโต จ สมุฏฺฐาติ, น กายโต, สิยา กายโต จ วาจโต จ จิตฺตโต จ}

\proseitem{13}{\csromanfolio{9}สญฺญาจิกาย กุฏิํ การาเปนฺตสฺส สํฆาทิเสโส กตฺถ ปญฺญตฺโตติ อาฬวิยํ ปญฺญตฺโต. กํ อารพฺภาติ อาฬวเก ภิกฺขู อารพฺภ. กิสฺมิํ วตฺถุสฺมินฺติ อาฬวกา ภิกฺขู สญฺญาจิกาย กุฏิโย การาเปสุํ, ตสฺมิํ วตฺถุสฺมิํ. เอกา ปญฺญตฺติ. ฉนฺนํ อาปตฺติ\-สมุฏฺฐานานํ ฉหิ สมุฏฺฐาเนหิ สมุฏฺฐาติ ฯเปฯ.}

\proseitem{14}{มหลฺลกํ วิหารํ การาเปนฺตสฺส สํฆาทิเสโส กตฺถ ปญฺญตฺโตติ โกสมฺพิยํ ปญฺญตฺโต. กํ อารพฺภาติ อายสฺมนฺตํ ฉนฺนํ อารพฺภ. กิสฺมิํ วตฺถุสฺมินฺติ อายสฺมา ฉนฺโน วิหารวตฺถุํ โสเธนฺโต อญฺญตรํ เจติยรุกฺขํ เฉทาเปสิ, ตสฺมิํ วตฺถุสฺมิํ. เอกา ปญฺญตฺติ. ฉนฺนํ อาปตฺติ\-สมุฏฺฐานานํ ฉหิ สมุฏฺฐาเนหิ สมุฏฺฐาติ ฯเปฯ.}

\proseitem{15}{ภิกฺขุํ อมูลเกน ปาราชิเกน ธมฺเมน อนุทฺธํเส\-นฺตสฺส สํฆาทิเสโส กตฺถ ปญฺญตฺโตติ ราชคเห ปญฺญตฺโต. กํ อารพฺภาติ เมตฺติย\-ภูมชเก ภิกฺขู อารพฺภ. กิสฺมิํ วตฺถุสฺมินฺติ เมตฺติย\-ภูมชกา ภิกฺขู อายสฺมนฺตํ ทพฺพํ มลฺลปุตฺตํ อมูลเกน ปาราชิเกน ธมฺเมน อนุทฺธํเสสุํ, ตสฺมิํ วตฺถุสฺมิํ. เอกา ปญฺญตฺติ. ฉนฺนํ อาปตฺติ\-สมุฏฺฐานานํ ตีหิ สมุฏฺฐาเนหิ สมุฏฺฐาติ ฯเปฯ.}

\proseitem{16}{ภิกฺขุํ อญฺญภาคิยสฺส อธิกรณสฺส กิญฺจิ เทสํ เลสมตฺตํ อุปาทาย ปาราชิเกน ธมฺเมน อนุทฺธํเส\-นฺตสฺส สํฆาทิเสโส กตฺถ ปญฺญตฺโตติ ราชคเห ปญฺญตฺโต. กํ อารพฺภาติ เมตฺติย\-ภูมชเก ภิกฺขู อารพฺภ. กิสฺมิํ วตฺถุสฺมินฺติ เมตฺติย\-ภูมชกา ภิกฺขู อายสฺมนฺตํ ทพฺพํ มลฺลปุตฺตํ อญฺญภาคิยสฺส อธิกรณสฺส กิญฺจิ เทสํ เลสมตฺตํ อุปาทาย ปาราชิเกน ธมฺเมน อนุทฺธํเสสุํ, ตสฺมิํ วตฺถุสฺมิํ. เอกา ปญฺญตฺติ. ฉนฺนํ อาปตฺติ\-สมุฏฺฐานานํ ตีหิ สมุฏฺฐาเนหิ}

\proseitem{17}{สํฆเภทกสฺส ภิกฺขุโน ยาวตติยํ สมนุภาสนาย น ปฏินิสฺสชฺชนฺตสฺส สํฆาทิเสโส กตฺถ ปญฺญตฺโตติ ราชคเห ปญฺญตฺโต. กํ อารพฺภาติ เทวทตฺตํ อารพฺภ. กิสฺมิํ วตฺถุสฺมินฺติ เทวทตฺโต สมคฺคสฺส สํฆสฺส เภทาย ปรกฺกมิ, ตสฺมิํ วตฺถุสฺมิํ. เอกา ปญฺญตฺติ. ฉนฺนํ อาปตฺติ\-สมุฏฺฐานานํ เอเกน สมุฏฺฐาเนน สมุฏฺฐาติ, กายโต จ วาจโต จ}

\proseitem{18}{\csromanfolio{10}เภทา\-นุวตฺตกานํ ภิกฺขูนํ ยาวตติยํ สมนุภาสนาย น ปฏินิสฺสชฺชนฺตานํ สํฆาทิเสโส กตฺถ ปญฺญตฺโตติ ราชคเห ปญฺญตฺโต. กํ อารพฺภาติ สมฺพหุเล ภิกฺขู อารพฺภ. กิํสฺมิํ วตฺถุสฺมินฺติ สมฺพหุลา ภิกฺขู เทวทตฺตสฺส สํฆเภทาย ปรกฺกม\-นฺตสฺส อนุวตฺตกา อเหสุํ วคฺควาทกา, ตสฺมิํ วตฺถุสฺมิํ. เอกา ปญฺญตฺติ. ฉนฺนํ อาปตฺติ\-สมุฏฺฐานานํ เอเกน สมุฏฺฐาเนน สมุฏฺฐาติ, กายโต จ วาจโต จ}

\proseitem{19}{ทุพฺพจสฺส ภิกฺขุโน ยาวตติยํ สมนุภาสนาย น ปฏินิสฺสชฺชนฺตสฺส สํฆาทิเสโส กตฺถ ปญฺญตฺโตติ โกสมฺพิยํ ปญฺญตฺโต. กํ อารพฺภาติ อายสฺมนฺตํ ฉนฺนํ อารพฺภ. กิสฺมิํ วตฺถุสฺมินฺติ อายสฺมา ฉนฺโน ภิกฺขูหิ สหธมฺมิกํ วุจฺจมาโน อตฺตานํ อวจนียํ อกาสิ, ตสฺมิํ วตฺถุสฺมิํ. เอกา ปญฺญตฺติ. ฉนฺนํ อาปตฺติ\-สมุฏฺฐานานํ เอเกน สมุฏฺฐาเนน สมุฏฺฐาติ, กายโต จ วาจโต จ}

\proseitemwithcloser{20}{กุลทูสกสฺส ภิกฺขุโน ยาวตติยํ สมนุภาสนาย น ปฏินิสฺสชฺชนฺตสฺส สํฆาทิเสโส กตฺถ ปญฺญตฺโตติ สาวตฺถิยํ ปญฺญตฺโต. กํ อารพฺภาติ อสฺสชิปุนพฺพสุเก ภิกฺขู อารพฺภ. กิสฺมิํ วตฺถุสฺมินฺติ อสฺสชิ\-ปุนพฺพสุกา ภิกฺขู สํเฆน ปพฺพาชนีย\-กมฺมกตา ภิกฺขู ฉนฺทคามิตา โทสคามิตา โมหคามิตา ภยคามิตา ปาเปสุํ, ตสฺมิํ วตฺถุสฺมิํ. เอกา ปญฺญตฺติ. ฉนฺนํ อาปตฺติ\-สมุฏฺฐานานํ เอเกน สมุฏฺฐาเนน สมุฏฺฐาติ, กายโต จ วาจโต จ จิตฺตโต จ สมุฏฺฐาติ ฯเปฯ.}{เตรส สํฆาทิเสสา นิฏฺฐิตา.}
\csromancenter{ตสฺสุทฺทานํ.}
\setlength{\csromangathapagewidth}{0pt}
\setlength{\csromangathawakwidth}{0pt}
\csromangathameasuregroup{วิสฺสฏฺฐิ กายสํสคฺคํ, \\ สญฺจริตฺตํ กุฏี เจว, \\ กิญฺจิเทสญฺจ เภโท จ, \\ ทุพฺพจํ กุลทูสญฺจ,}{\csromangathabat{วิสฺสฏฺฐิ กายสํสคฺคํ,}{ทุฏฺฐุลฺลํ อตฺตกามญฺจ.} \\ \csromangathabat{สญฺจริตฺตํ กุฏี เจว,}{วิหาโร จ อมูลกํ.} \\ \csromangathabat{กิญฺจิเทสญฺจ เภโท จ,}{ตสฺเสว\textsuperscript{1} อนุวตฺตกา.} \\ \csromangathabat{ทุพฺพจํ กุลทูสญฺจ,}{สํฆาทิเสสา เตรสาติ.}}
\csromangathasetleft{วิสฺสฏฺฐิ กายสํสคฺคํ, \\ สญฺจริตฺตํ กุฏี เจว, \\ กิญฺจิเทสญฺจ เภโท จ, \\ ทุพฺพจํ กุลทูสญฺจ,}
\csromangathagroup{\csromangathastanza{\csromangathabat{วิสฺสฏฺฐิ กายสํสคฺคํ,}{ทุฏฺฐุลฺลํ อตฺตกามญฺจ.} \\ \csromangathabat{สญฺจริตฺตํ กุฏี เจว,}{วิหาโร จ อมูลกํ.}}\csromangathastanza{\csromangathabat{กิญฺจิเทสญฺจ เภโท จ,}{ตสฺเสว\csromansharedfootnote{f7e0a71730974247}{ตเถว (ก)} อนุวตฺตกา.} \\ \csromangathabat{ทุพฺพจํ กุลทูสญฺจ,}{สํฆาทิเสสา เตรสาติ.}}}

\csromanmatikaanchor{mtk.5}
\csromantocmarkref{chapter}{3. อนิยตกณฺฑ}{mtk.5}
\bookmark[dest={mtk.5},level=0]{3. อนิยตกณฺฑ}
\chapterheadpageread{3. \textbf{อนิยตกณฺฑ}}
\proseitem{21}{\csromanfolio{11}ยํ เตน ภควตา ชานตา ปสฺสตา อรหตา สมฺมา\-สมฺพุทฺเธน ปฐโม อนิยโต กตฺถ ปญฺญตฺโต,กํ อารพฺภ, กิสฺมิํ วตฺถุสฺมิํ, อตฺถิ ตตฺถ ปญฺญตฺติ อนุปญฺญตฺติ อนุปฺปนฺนปญฺญตฺติ, สพฺพตฺถ\-ปญฺญตฺติ ปเทสปญฺญตฺติ, สาธารณ\-ปญฺญตฺติ อสาธารณ\-ปญฺญตฺติ, เอกโต ปญฺญตฺติ อุภโตปญฺญตฺติ, ปญฺจนฺนํ ปาติโมกฺขุ\-ทฺเทสานํ กตฺโถคธํ กตฺถ ปริยาปนฺนํ, กตเมน อุทฺเทเสน อุทฺเทสํ อาคจฺฉติ, จตุนฺนํ วิปตฺตีนํ กตมา วิปตฺติ, สตฺตนฺนํ อาปตฺติ\-กฺขนฺธานํ กตโม อาปตฺติกฺขนฺโธ, ฉนฺนํ อาปตฺติ\-สมุฏฺฐานานํ กติหิ สมุฏฺฐาเนหิ สมุฏฺฐาติ, จตุนฺนํ อธิกรณานํ กตมํ อธิกรณํ, สตฺตนฺนํ สมถานํ กติหิ สมเถหิ สมฺมติ, โก ตตฺถ วินโย, โก ตตฺถ อภิวินโย, กิํ ต ตฺถ ปาติโมกฺขํ, กิํ ตตฺถ อธิปาติโมกฺขํ, กาวิปตฺติ, กา สมฺปตฺติ, กา ปฏิปตฺติ, กติ อตฺถวเส ปฏิจฺจ ภควตา ปฐโม อนิยโต ปญฺญตฺโต, เก สิกฺขนฺติ, เก สิกฺขิตสิกฺขา, กตฺถ ฐิตํ, เก ธาเรนฺติ, กสฺส วจนํ, เกนาภตนฺติ.}

\proseitem{22}{ยํ เตน ภควตา ชานตา ปสฺสตา อรหตา สมฺมา\-สมฺพุทฺเธน ปฐโม อนิยโต กตฺถ ปญฺญตฺโตติ สาวตฺถิยํ ปญฺญตฺโต. กํ อารพฺภาติ อายสฺมนฺตํ อุทายิํ อารพฺภ. กิสฺมิํ วตฺถุสฺมินฺติ อายสฺมา อุทายี มาตุคาเมน สทฺธิํ เอโก เอกาย รโห ปฏิจฺฉนฺเน อาสเน อลํกมฺมนิเย นิสชฺชํ กปฺเปสิ,ตสฺมิํ วตฺถุสฺมิํ. อตฺถิ ตตฺถ ปญฺญตฺติ อนุปญฺญตฺติ อนุปฺปนฺนปญฺญตฺตีติ เอกา ปญฺญตฺติ, อนุปญฺญตฺติ อนุปฺปนฺนปญฺญตฺติ ตสฺมิํ นตฺถิ. สพฺพตฺถ\-ปญฺญตฺติ ปเทสปญฺญตฺตีติ สพฺพตฺถ\-ปญฺญตฺติ. สาธารณ\-ปญฺญตฺติ อสาธารณปญฺญตฺตีติ อสาธารณ\-ปญฺญตฺติ. เอกโตปญฺญตฺติ อุภโตปญฺญตฺตีติ เอกโตปญฺญตฺติ. ปญฺจนฺนํ ปาติโมกฺขุ\-ทฺเทสานํ กตฺโถคธํ กตฺถ ปริยาปนฺนนฺติ นิทาโนคธํ นิทาน\-ปริยาปนฺนํ. กตเมน อุทฺเทเสน อุทฺเทสํ อาคจฺฉตีติ จตุตฺเถน อุทฺเทเสน อุทฺเทสํ อาคจฺฉติ. จตุนฺนํ วิปตฺตีนํ กตมา วิปตฺตีติ สิยา สีลวิปตฺติ, สิยา อาจารวิปตฺติ. สตฺตนฺนํ อาปตฺติ\-กฺขนฺธานํ กตโม อาปตฺติ\-กฺขนฺโธติ สิยา ปาราชิกาปตฺติ\-กฺขนฺโธ, สิยา สํฆาทิเสสา\-ปตฺติกฺขนฺโธ, สิยา ปาจิตฺติยาปตฺติ\-กฺขนฺโธ. ฉนฺนํ อาปตฺติ\-สมุฏฺฐานานํ กติหิ สมุฏฺฐาเนหิ สมุฏฺฐาตีติ \csromanfolio{12}เอเกน สมุฏฺฐาเนน สมุฏฺฐาติ, กายโต จ จิตฺตโต จ สมุฏฺฐาติ, น วาจโต. จตุนฺนํ อธิกรณานํ กตมํ อธิกรณนฺติ อาปตฺตา\-ธิกรณํ. สตฺตนฺนํ สมถานํ กติหิ สมเถหิ สมฺมตีติ ตีหิ สมเถหิ สมฺมติ, สิยา สมฺมุขา\-วินเยน จ ปฏิญฺญาต\-กรเณน จ, สิยา สมฺมุขา\-วินเยน จ ติณวตฺถารเกน จ. โก ตตฺถ วินโย โก ตตฺถ อภิวินโยติ ปญฺญตฺติ วินโย, วิภตฺติ อภิวินโย. กิํ ตตฺถ ปาติโมกฺขํ กิํ ตตฺถ อธิปาติ\-โมกฺขนฺติ ปญฺญตฺติ ปาติโมกฺขํ, วิภตฺติ อธิปาติโมกฺขํ. กา วิปตฺตีติ อสํวโร วิปตฺติ. กา สมฺปตฺตีติ สํวโร สมฺปตฺติ. กา ปฏิปตฺตีติ “น เอวรูปํ กริสฺสามี”ติ ยาวชีวํ อาปาณโกฏิกํ สมาทาย สิกฺขติ สิกฺขาปเทสุ. * กติ อตฺถวเส ปฏิจฺจ ภควตา ปฐโม อนิยโต ปญฺญตฺโตติ ทส อตฺถวเส ปฏิจฺจ ภควตา ปฐโม อนิยโต ปญฺญตฺโต สํฆสุฏฺฐุตาย สํฆผาสุตาย ทุมฺมงฺกูนํ ปุคฺคลานํ นิคฺคหาย เปสลานํ ภิกฺขูนํ ผาสุวิหาราย ทิฏฺฐ\-ธมฺมิกานํ อาสวานํ สํวราย สมฺปรายิกานํ อาสวานํ ปฏิฆาตาย อปฺปสนฺนานํ ปสาทาย ปสนฺนานํ ภิยฺโยภาวาย สทฺธมฺม\-ฏฺฐิติยา วินยา\-นุคฺคหาย. เก สิกฺขนฺตีติ เสกฺขา จ ปุถุชฺชน\-กลฺยาณกา จ สิกฺขนฺติ. เก สิกฺขิต\-สิกฺขาติ อรหนฺโต สิกฺขิตสิกฺขา. กตฺถ ฐิตนฺติ สิกฺขากาเมสุ ฐิตํ. เก ธาเรนฺตีติ เยสํ วตฺตติ, เต ธาเรนฺติ. กสฺส วจนนฺติ ภควโต วจนํ อรหโต สมฺมา\-สมฺพุทฺธสฺส. เกนาภตนฺติ ปรมฺปราภตํ.}

\prose{อุปาลิ ทาสโก เจว, โสณโก สิคฺคโว ตถา.}

\prose{โมคฺคลิปุตฺเตน ปญฺจมา, เอเต ชมฺพุสิริวฺหเย.}

\prose{ตโต มหินฺโท อิฏฺฏิโย, อุตฺติโย สมฺพโล ตถา.}

\prose{วินยํ เต วาจยิํสุ, ปิฏกํ ตมฺพปณฺณิยา.}

\prose{นิกาเย ปญฺจ วาเจสุํ, สตฺต เจว ปกรเณ.}

\prose{ตโต อริฏฺโฐ เมธาวี, ติสฺสทตฺโต จ ปณฺฑิโต.}

\prose{วิสารโท กาฬสุมโน, เถโร จ ทีฆนามโก.}

\prose{\csromanfolio{13}ปุนเทว กาฬสุมโน, นาคตฺเถโร จ พุทฺธรกฺขิโต.}

\prose{ติสฺสตฺเถโร จ เมธาวี, เทวตฺเถโร จ ปณฺฑิโต.}

\prose{ปุนเทว สุมโน เมธาวี, วินเย จ วิสารโท.}

\prose{พหุสฺสุโต จูฬนาโค, คโชว ทุปฺปธํสิโย.}

\prose{ธมฺมปาลิต\-นาโม จ, โรหเณ สาธุปูชิโต.}

\prose{ตสฺส สิสฺโส มหาปญฺโญ, เขมนาโม ติเปฏโก.}

\prose{ทีเป ตารกราชาว, ปญฺญาย อติโรจถ.}

\prose{อุปติสฺโส จ เมธาวี, ผุสฺสเทโว มหากถี.}

\prose{ปุนเทว สุมโน เมธาวี, ปุปฺผนาโม พหุสฺสุโต.}

\prose{มหากถี มหาสิโว, ปิฏเก สพฺพตฺถ โกวิโท.}

\prose{ปุนเทว อุปาลิ เมธาวี, วินเย จ วิสารโท.}

\prose{มหานาโค มหาปญฺโญ, สทฺธมฺม\-วํส\-โกวิโท.}

\prose{ปุนเทว อภโย เมธาวี, ปิฏเก สพฺพตฺถ โกวิโท.}

\prose{ติสฺสตฺเถโร จ เมธาวี, วินเย จ วิสารโท.}

\prose{ตสฺส สิสฺโส มหาปญฺโญ, ปุปฺผนาโม พหุสฺสุโต.}

\prose{สาสนํ อนุรกฺขนฺโต, ชมฺพุทีเป ปติฏฺฐิโต.}

\prose{จูฬาภโย จ เมธาวี, วินเย จ วิสารโท.}

\prose{ติสฺสตฺเถโร จ เมธาวี, สทฺธมฺม\-วํส\-โกวิโท.}

\prose{จูฬเทโว จ เมธาวี,วินเย จ วิสารโท.}

\prose{สิวตฺเถโร จ เมธาวี,วินเย สพฺพตฺถ โกวิโท.}

\prose{เอเต นาคา มหาปญฺญา, วินยญฺญู มคฺคโกวิทา.}

\prose{วินยํ ทีเป ปกาเสสุํ, ปิฏกํ ตมฺพปณฺณิยาติ.}

\proseitemwithcloser{23}{ยํ เตน ภควตา ชานตา ปสฺสตา อรหตา สมฺมา\-สมฺพุทฺเธน ทุติโย อนิยโต กตฺถ ปญฺญตฺโตติ สาวตฺถิยํ ปญฺญตฺโต. กํ อารพฺภาติ อายสฺมนฺตํ อุทายิํ อารพฺภ. กิสฺมิํ วตฺถุสฺมินฺติ อายสฺมา อุทายี มาตุคาเมน สทฺธิํ เอโก เอกาย รโห นิสชฺชํ กปฺเปสิ, ตสฺมิํ วตฺถุสฺมิํ. อตฺถิ ตตฺถ ปญฺญตฺติ อนุปญฺญตฺติ อนุปฺปนฺนปญฺญตฺตีติ เอกา ปญฺญตฺติ, \csromanfolio{14}อนุปญฺญตฺติ อนุปฺปนฺนปญฺญตฺติ ตสฺมิํ นตฺถิ. สพฺพตฺถ\-ปญฺญตฺติ ปเทสปญฺญตฺตีติ สพฺพตฺถ\-ปญฺญตฺติ. สาธารณ\-ปญฺญตฺติ อสาธารณปญฺญตฺตีติ อสาธารณ\-ปญฺญตฺติ. เอกโตปญฺญตฺติ อุภโตปญฺญตฺตีติ เอกโตปญฺญตฺติ. ปญฺจนฺนํ ปาติโมกฺขุ\-ทฺเทสานํ กตฺโถคธํ กตฺถ ปริยาปนฺนนฺติ นิทาโนคธํ นิทาน\-ปริยาปนฺนํ. กตเมน อุทฺเทเสน อุทฺเทสํ อาคจฺฉตีติ จตุตฺเถน อุทฺเทเสน อุทฺเทสํ อาคจฺฉติ. จตุนฺนํ วิปตฺตีนํ กตมา วิปตฺตีติ สิยา สีลวิปตฺติ, สิยา อาจารวิปตฺติ. สตฺตนฺนํ อาปตฺติ\-กฺขนฺธานํ กตโม อาปตฺติ\-กฺขนฺโธติ สิยา สํฆาทิเสสา\-ปตฺติกฺขนฺโธ, สิยา ปาจิตฺติยาปตฺติ\-กฺขนฺโธ. ฉนฺนํ อาปตฺติ\-สมุฏฺฐานานํ กติหิ สมุฏฺฐาเนหิ สมุฏฺฐาตีติ ตีหิ สมุฏฺฐาเนหิ สมุฏฺฐาติ, สิยา กายโต จ จิตฺตโต จ สมุฏฺฐาติ, น วาจโต, สิยา วาจโต จ จิตฺตโต จ สมุฏฺฐาติ, น กายโต, สิยา กายโต จ วาจโต จ จิตฺตโต จ สมุฏฺฐาติ. จตุนฺนํ อธิกรณานํ กตมํ อธิกรณนฺติ อาปตฺตา\-ธิกรณํ. สตฺตนฺนํ สมถานํ กติหิ สมเถหิ สมฺมตีติ ตีหิ สมเถหิ สมฺมติ, สิยา สมฺมุขา\-วินเยน จ ปฏิญฺญาต\-กรเณน จ, สิยา สมฺมุขา\-วินเยน จ ติณวตฺถารเกน จ ฯเปฯ.}{เทฺว อนิยตา นิฏฺฐิตา.}
\csromancenter{ตสฺสุทฺทานํ.}
\setlength{\csromangathapagewidth}{0pt}
\setlength{\csromangathawakwidth}{0pt}
\csromangathameasuregroup{อลํกมฺมนิยญฺเจว,}{\csromangathabat{อลํกมฺมนิยญฺเจว,}{ตเถว จ น เหว โข.}}
\csromangathasetleft{อลํกมฺมนิยญฺเจว,}
\csromangathagroup{\csromangathastanza{\csromangathabat{อลํกมฺมนิยญฺเจว,}{ตเถว จ น เหว โข.}}}

\vspace{\tipitakaparskip}
\csromancenter{อนิยตา สุปญฺญตฺตา, พุทฺธเสฏฺเฐน ตาทินาติ.}
\csromansectionrule
\csromanmatikaanchor{mtk.6}
\csromanmatikaanchor{mtk.7}
\csromantocmarkref{chapter}{4. นิสฺสคฺคิยกณฺฑ}{mtk.6}
\bookmark[dest={mtk.6},level=0]{4. นิสฺสคฺคิยกณฺฑ}
\csromantocmarkref{section}{1. กถินวคฺค}{mtk.7}
\bookmark[dest={mtk.7},level=1]{1. กถินวคฺค}
\csromanmarkh{4. นิสฺสคฺคิย\-กณฺฑ}\csromanmarkhii{1. กถินวคฺค}\csromanheadernomark{1}{4. นิสฺสคฺคิย\-กณฺฑ 1. กถินวคฺค}
\proseitem{24}{ยํ เตน ภควตา ชานตา ปสฺสตา อรหตา สมฺมา\-สมฺพุทฺเธน อติเรกจีวรํ ทสาหํ อติกฺกาเมนฺตสฺส นิสฺสคฺคิยํ ปาจิตฺติยํ กตฺถ ปญฺญตฺตนฺติ เวสาลิยํ ปญฺญตฺตํ. กํ อารพฺภาติ ฉพฺพคฺคิเย ภิกฺขู อารพฺภ. กิสฺมิํ วตฺถุสฺมินฺติ ฉพฺพคฺคิยา ภิกฺขู อติเรกจีวรํ ธาเรสุํ, ตสฺมิํ วตฺถุสฺมิํ. เอกา ปญฺญตฺติ, เอกา อนุปญฺญตฺติ. ฉนฺนํ อาปตฺติ\-สมุฏฺฐานานํ ทฺวีหิ \csromanfolio{15}สมุฏฺฐาเนหิ สมุฏฺฐาติ, สิยา กายโต จ วาจโต จ สมุฏฺฐาติ, น จิตฺตโต, สิยา กายโต จ วาจโต จ จิตฺตโต จ สมุฏฺฐาติ ฯเปฯ.}

\proseitem{25}{เอกรตฺตํ ติจีวเรน วิปฺป\-วส\-นฺตสฺส นิสฺสคฺคิยํ ปาจิตฺติยํ กตฺถ ปญฺญตฺตนฺติ สาวตฺถิยํ ปญฺญตฺตํ. กํ อารพฺภาติ สมฺพหุเล ภิกฺขู อารพฺภ. กิสฺมิํ วตฺถุสฺมินฺติ สมฺพหุลา ภิกฺขู ภิกฺขูนํ หตฺเถ จีวรํ นิกฺขิปิตฺวา สนฺตรุตฺตเรน ชนปท\-จาริกํ ปกฺกมิํสุ, ตสฺมิํ วตฺถุสฺมิํ. เอกา ปญฺญตฺติ, เอกา อนุปญฺญตฺติ. ฉนฺนํ อาปตฺติ\-สมุฏฺฐานานํ ทฺวีหิ สมุฏฺฐาเนหิ สมุฏฺฐาติ, สิยา กายโต จ วาจโต จ สมุฏฺฐาติ, น จิตฺตโต, สิยา กายโต จ วาจโต จ จิตฺตโต จ สมุฏฺฐาติ ฯเปฯ.}

\proseitem{26}{อกาลจีวรํ ปฏิคฺคเหตฺวา มาสํ อติกฺกาเมนฺตสฺส นิสฺสคฺคิยํ ปาจิตฺติยํ กตฺถ ปญฺญตฺตนฺติ สาวตฺถิยํ ปญฺญตฺตํ. กํ อารพฺภาติ สมฺพหุเล ภิกฺขู อารพฺภ. กิสฺมิํ วตฺถุสฺมินฺติ สมฺพหุลา ภิกฺขู อกาลจีวรํ ปฏิคฺคเหตฺวา มาสํ อติกฺกาเมสุํ, ตสฺมิํ วตฺถุสฺมิํ. เอกา ปญฺญตฺติ, เอกา อนุปญฺญตฺติ. ฉนฺนํ อาปตฺติ\-สมุฏฺฐานานํ ทฺวีหิ สมุฏฺฐาเนหิ สมุฏฺฐาติ, สิยา กายโต จ วาจโต จ สมุฏฺฐาติ, น จิตฺตโต, สิยา กายโต จ วาจโต จ จิตฺตโต จ สมุฏฺฐาติ ฯเปฯ.}

\proseitem{27}{อญฺญาติกาย ภิกฺขุนิยา ปุราณจีวรํ โธวาเปนฺตสฺส นิสฺสคฺคิยํ ปาจิตฺติยํ กตฺถ ปญฺญตฺตนฺติ สาวตฺถิยํ ปญฺญตฺตํ. กํ อารพฺภาติ อายสฺมนฺตํ อุทายิํ อารพฺภ. กิสฺมิํ วตฺถุสฺมินฺติ อายสฺมา อุทายี อญฺญาติกาย ภิกฺขุนิยา ปุราณจีวรํ โธวาเปสิ, ตสฺมิํ วตฺถุสฺมิํ. เอกา ปญฺญตฺติ. ฉนฺนํ อาปตฺติ\-สมุฏฺฐานานํ ฉหิ สมุฏฺฐาเนหิ สมุฏฺฐาติ ฯเปฯ.}

\proseitem{28}{อญฺญาติกาย ภิกฺขุนิยา หตฺถโต จีวรํ ปฏิคฺคณฺหนฺตสฺส นิสฺสคฺคิยํ ปาจิตฺติยํ กตฺถ ปญฺญตฺตนฺติ ราชคเห ปญฺญตฺตํ. กํ อารพฺภาติ อายสฺมนฺตํ อุทายิํ อารพฺภ. กิสฺมิํ วตฺถุสฺมินฺติ อายสฺมา อุทายี อญฺญาติกาย ภิกฺขุนิยา หตฺถโต จีวรํ ปฏิคฺคเหสิ, ตสฺมิํ วตฺถุสฺมิํ. เอกา ปญฺญตฺติ, เอกา อนุปญฺญตฺติ. ฉนฺนํ อาปตฺติ\-สมุฏฺฐานานํ ฉหิ สมุฏฺฐาเนหิ สมุฏฺฐาติ ฯเปฯ.}

\proseitem{29}{\csromanfolio{16}อญฺญาตกํ คหปติํ วา คหปตานิํ วา จีวรํ วิญฺญาเปนฺตสฺส นิสฺสคฺคิยํ ปาจิตฺติยํ กตฺถ ปญฺญตฺตนฺติ สาวตฺถิยํ ปญฺญตฺตํ. กํ อารพฺภาติ อายสฺมนฺตํ อุปนนฺทํ สกฺยปุตฺตํ อารพฺภ. กิสฺมิํ วตฺถุสฺมินฺติ อายสฺมา อุปนนฺโท สกฺยปุตฺโต อญฺญาตกํ เสฏฺฐิปุตฺตํ จีวรํ วิญฺญาเปสิ, ตสฺมิํ วตฺถุสฺมิํ. เอกา ปญฺญตฺติ, เอกา อนุปญฺญตฺติ. ฉนฺนํ อาปตฺติ\-สมุฏฺฐานานํ ฉหิ สมุฏฺฐาเนหิ สมุฏฺฐาติ ฯเปฯ.}

\proseitem{30}{อญฺญาตกํ คหปติํ วา คหปตานิํ วา ตตุตฺตริ จีวรํ วิญฺญาเปนฺตสฺส นิสฺสคฺคิยํ ปาจิตฺติยํ กตฺถ ปญฺญตฺตนฺติ สาวตฺถิยํ ปญฺญตฺตํ. กํ อารพฺภาติ ฉพฺพคฺคิเย ภิกฺขู อารพฺภ. กิสฺมิํ วตฺถุสฺมินฺติ ฉพฺพคฺคิยา ภิกฺขู น มตฺตํ ชานิตฺวา พหุํ จีวรํ วิญฺญาเปสุํ, ตสฺมิํ วตฺถุสฺมิํ. เอกา ปญฺญตฺติ. ฉนฺนํ อาปตฺติ\-สมุฏฺฐานานํ ฉหิ สมุฏฺฐาเนหิ สมุฏฺฐาติ ฯเปฯ.}

\proseitem{31}{ปุพฺเพ อปฺปวาริตสฺส อญฺญาตกํ คหปติกํ อุปสงฺกมิตฺวา จีวเร วิกปฺปํ อาปชฺชนฺตสฺส นิสฺสคฺคิยํ ปาจิตฺติยํ กตฺถ ปญฺญตฺตนฺติ สาวตฺถิยํ ปญฺญตฺตํ. กํ อารพฺภาติ อายสฺมนฺตํ อุปนนฺทํ สกฺยปุตฺตํ อารพฺภ. กิสฺมิํ วตฺถุสฺมินฺติ อายสฺมา อุปนนฺโท สกฺยปุตฺโต ปุพฺเพ อปฺปวาริโต อญฺญาตกํ คหปติกํ อุปสงฺกมิตฺวา จีวเร วิกปฺปํ อาปชฺชิ, ตสฺมิํ วตฺถุสฺมิํ. เอกา ปญฺญตฺติ. ฉนฺนํ อาปตฺติ\-สมุฏฺฐานานํ ฉหิ สมุฏฺฐาเนหิ สมุฏฺฐาติ ฯเปฯ.}

\proseitem{32}{ปุพฺเพ อปฺปวาริตสฺส อญฺญาตเก คหปติเก อุปสงฺกมิตฺวา จีวเร วิกปฺปํ อาปชฺชนฺตสฺส นิสฺสคฺคิยํ ปาจิตฺติยํ กตฺถ ปญฺญตฺตนฺติ สาวตฺถิยํ ปญฺญตฺตํ. กํ อารพฺภาติ อายสฺมนฺตํ อุปนนฺทํ สกฺยปุตฺตํ อารพฺภ. กิสฺมิํ วตฺถุสฺมินฺติ อายสฺมา อุปนนฺโท สกฺยปุตฺโต ปุพฺเพ อปฺปวาริโต อญฺญาตเก คหปติเก อุปสงฺกมิตฺวา จีวเร วิกปฺปํ อาปชฺชิ, ตสฺมิํ วตฺถุสฺมิํ. เอกา ปญฺญตฺติ. ฉนฺนํ อาปตฺติ\-สมุฏฺฐานานํ ฉหิ สมุฏฺฐาเนหิ สมุฏฺฐาติ ฯเปฯ.}

\proseitemwithclosermidruled{33}{อติเรก\-ติ\-กฺขตฺตุํ โจทนาย อติเรก\-ฉ\-กฺขตฺตุํ ฐาเนน จีวรํ อภินิ\-ปฺผาเท\-นฺตสฺส นิสฺสคฺคิยํ ปาจิตฺติยํ กตฺถ ปญฺญตฺตนฺติ สาวตฺถิยํ ปญฺญตฺตํ. กํ อารพฺภาติ อายสฺมนฺตํ อุปนนฺทํ สกฺยปุตฺตํ อารพฺภ. กิสฺมิํ วตฺถุสฺมินฺติ อายสฺมา อุปนนฺโท สกฺยปุตฺโต อุปาสเกน “อชฺชโณฺห ภนฺเต \csromanfolio{17}อาคเมหี”ติ วุจฺจมาโน นาคเมสิ, ตสฺมิํ วตฺถุสฺมิํ. เอกา ปญฺญตฺติ. ฉนฺนํ อาปตฺติ\-สมุฏฺฐานานํ ฉหิ สมุฏฺฐาเนหิ สมุฏฺฐาติ ฯเปฯ.}{กถินวคฺโค ปฐโม.}
\csromanmatikaanchor{mtk.8}
\csromantocmarkref{section}{2. โกสิยวคฺค}{mtk.8}
\bookmark[dest={mtk.8},level=1]{2. โกสิยวคฺค}
\csromanheader{1}{2. \textbf{โกสิยวคฺค}}
\proseitem{34}{โกสิยมิสฺสกํ สนฺถตํ การาเปนฺตสฺส นิสฺสคฺคิยํ ปาจิตฺติยํ กตฺถ ปญฺญตฺตนฺติ อาฬวิยํ ปญฺญตฺตํ. กํ อารพฺภาติ ฉพฺพคฺคิเย ภิกฺขู อารพฺภ. กิสฺมิํ วตฺถุสฺมินฺติ ฉพฺพคฺคิยา ภิกฺขู โกสิยการเก อุปสงฺกมิตฺวา เอวมาหํสุ “พหู อาวุโส โกสการเก ปจถ, อมฺหากมฺปิ ทสฺสถ, มยมฺปิ อิจฺฉาม โกสิยมิสฺสกํ สนฺถตํ กาตุนฺ”ติ, ตสฺมิํ วตฺถุสฺมิํ. เอกา ปญฺญาตฺติ. ฉนฺนํ อาปตฺติ\-สมุฏฺฐานานํ ฉหิ สมุฏฺฐาเนหิ สมุฏฺฐาติ ฯเปฯ.}

\proseitem{35}{สุทฺธ\-กาฬกานํ เอฬกโลมานํ สนฺถตํ การาเปนฺตสฺส นิสฺสคฺคิยํ ปาจิตฺติยํ กตฺถ ปญฺญตฺตนฺติ เวสาลิยํ ปญฺญตฺตํ. กํ อารพฺภาติ ฉพฺพคฺคิเย ภิกฺขู อารพฺภ. กิสฺมิํ วตฺถุสฺมินฺติ ฉพฺพคฺคิยา ภิกฺขู สุทฺธ\-กาฬกานํ เอฬกโลมานํ สนฺถตํ การาเปสุํ, ตสฺมิํ วตฺถุสฺมิํ. เอกา ปญฺญตฺติ. ฉนฺนํ อาปตฺติ\-สมุฏฺฐานานํ ฉหิ สมุฏฺฐาเนหิ สมุฏฺฐาติ ฯเปฯ.}

\proseitem{36}{อนาทิยิตฺวา ตุลํ โอทาตานํ ตุลํ โคจริยานํ นวํ สนฺถตํ การาเปนฺตสฺส นิสฺสคฺคิยํ ปาจิตฺติยํ กตฺถ ปญฺญตฺตนฺติ สาวตฺถิยํ ปญฺญตฺตํ. กํ อารพฺภาติ ฉพฺพคฺคิเย ภิกฺขู อารพฺภ. กิสฺมิํ วตฺถุสฺมินฺติ ฉพฺพคฺคิยา ภิกฺขู โถกญฺเญว โอทาตํ อนฺเต\csromansharedfootnote{a93e55e7bfee4816}{โอทาตานํ อนฺเต อนฺเต (สี), โอทาตานํ อนฺเต (สฺยา)} อาทิยิตฺวา ตเถว สุทฺธ\-กาฬกานํ เอฬกโลมานํ สนฺถตํ การาเปสุํ, ตสฺมิํ วตฺถุสฺมิํ. เอกา ปญฺญตฺติ. ฉนฺนํ อาปตฺติ\-สมุฏฺฐานานํ ฉหิ สมุฏฺฐาเนหิ สมุฏฺฐาติ ฯเปฯ.}

\proseitem{37}{อนุวสฺสํ สนฺถตํ การาเปนฺตสฺส นิสฺสคฺคิยํ ปาจิตฺติยํ กตฺถ ปญฺญตฺตนฺติ สาวตฺถิยํ ปญฺญตฺตํ. กํ อารพฺภาติ สมฺพหุเล ภิกฺขู อารพฺภ. กิสฺมิํ วตฺถุสฺมินฺติ สมฺพหุลา ภิกฺขู อนุวสฺสํ สนฺถตํ การาเปสุํ, ตสฺมิํ \csromanfolio{18}วตฺถุสฺมิํ. เอกา ปญฺญตฺติ, เอกา อนุปญฺญตฺติ. ฉนฺนํ อาปตฺติ\-สมุฏฺฐานานํ ฉหิ สมุฏฺฐาเนหิ สมุฏฺฐาติ ฯเปฯ.}

\proseitem{38}{อนาทิยิตฺวา ปุราณ\-สนฺถตสฺส สามนฺตา สุคต\-วิทตฺถิํ นวํ นิสีทน\-สนฺถตํ การาเปนฺตสฺส นิสฺสคฺคิยํ ปาจิตฺติยํ กตฺถ ปญฺญตฺตนฺติ สาวตฺถิยํ ปญฺญตฺตํ. กํ อารพฺภาติ สมฺพหุเล ภิกฺขู อารพฺภ. กิสฺมิํ วตฺถุสฺมินฺติ สมฺพหุลา ภิกฺขู สนฺถตานิ อุชฺฌิตฺวา อารญฺญิกงฺคํ ปิณฺฑปาติ\-กงฺคํ ปํสุกูลิ\-กงฺคํ สมาทิยิํสุ, ตสฺมิํ วตฺถุสฺมิํ. เอกา ปญฺญตฺติ. ฉนฺนํ อาปตฺติ\-สมุฏฺฐานานํ ฉหิ สมุฏฺฐาเนหิ สมุฏฺฐาติ ฯเปฯ.}

\proseitem{39}{เอฬกโลมานิ ปฏิคฺคเหตฺวา ติโยชนํ อติกฺกาเมนฺตสฺส นิสฺสคฺคิยํ ปาจิตฺติยํ กตฺถ ปญฺญตฺตนฺติ สาวตฺถิยํ ปญฺญตฺตํ. กํ อารพฺภาติ อญฺญตรํ ภิกฺขุํ อารพฺภ. กิสฺมิํ วตฺถุสฺมินฺติ อญฺญตโร ภิกฺขุ เอฬกโลมานิ ปฏิคฺคเหตฺวา ติโยชนํ อติกฺกาเมสิ, ตสฺมิํ วตฺถุสฺมิํ. เอกา ปญฺญตฺติ. ฉนฺนํ อาปตฺติ\-สมุฏฺฐานานํ ทฺวีหิ สมุฏฺฐาเนหิ สมุฏฺฐาติ, สิยา กายโต สมุฏฺฐาติ, น วาจโต, น จิตฺตโต, สิยา กายโต จ จิตฺตโต จ สมุฏฺฐาติ, น วาจโต ฯเปฯ.}

\proseitem{40}{อญฺญาติกาย ภิกฺขุนิยา เอฬกโลมานิ โธวาเปนฺตสฺส นิสฺสคฺคิยํ ปาจิตฺติยํ กตฺถ ปญฺญตฺตนฺติ สกฺเกสุ ปญฺญตฺตํ. กํ อารพฺภาติ ฉพฺพคฺคิเย ภิกฺขู อารพฺภ. กิสฺมิํ วตฺถุสฺมินฺติ ฉพฺพคฺคิยา ภิกฺขู อญฺญาติกาหิ ภิกฺขุนีหิ เอฬกโลมานิ โธวาเปสุํ, ตสฺมิํ วตฺถุสฺมิํ. เอกา ปญฺญตฺติ. ฉนฺนํ อาปตฺติ\-สมุฏฺฐานานํ ฉหิ สมุฏฺฐาเนหิ สมุฏฺฐาติ ฯเปฯ.}

\proseitem{41}{รูปิยํ ปฏิคฺคณฺหนฺตสฺส นิสฺสคฺคิยํ ปาจิตฺติยํ กตฺถ ปญฺญตฺตนฺติ ราชคเห ปญฺญตฺตํ. กํ อารพฺภาติ อายสฺมนฺตํ อุปนนฺทํ สกฺยปุตฺตํ อารพฺภ. กิสฺมิํ วตฺถุสฺมินฺติ อายสฺมา อุปนนฺโท สกฺยปุตฺโต รูปิยํ ปฏิคฺคเหสิ, ตสฺมิํ วตฺถุสฺมิํ. เอกา ปญฺญตฺติ. ฉนฺนํ อาปตฺติ\-สมุฏฺฐานานํ ฉหิ สมุฏฺฐาเนหิ สมุฏฺฐาติ ฯเปฯ.}

\proseitem{42}{นานปฺปการกํ รูปิย\-สํโวหารํ สมา\-ป\-ชฺช\-นฺตสฺส นิสฺสคฺคิยํ ปาจิตฺติยํ กตฺถ ปญฺญตฺตนฺติ สาวตฺถิยํ ปญฺญตฺตํ. กํ อารพฺภาติ ฉพฺพคฺคิเย \csromanfolio{19}ภิกฺขู อารพฺภ. กิสฺมิํ วตฺถุสฺมินฺติ ฉพฺพคฺคิยา ภิกฺขู นานปฺปการกํ รูปิย\-สํโวหารํ สมาปชฺชิํสุ, ตสฺมิํ วตฺถุสฺมิํ. เอกา ปญฺญตฺติ. ฉนฺนํ อาปตฺติ\-สมุฏฺฐานานํ ฉหิ สมุฏฺฐาเนหิ สมุฏฺฐาติ ฯเปฯ.}

\proseitemwithclosermidruled{43}{นานปฺปการกํ กยวิกฺกยํ สมา\-ป\-ชฺช\-นฺตสฺส นิสฺสคฺคิยํ ปาจิตฺติยํ กตฺถ ปญฺญตฺตนฺติ สาวตฺถิยํ ปญฺญตฺตํ. กํ อารพฺภาติ อายสฺมนฺตํ อุปนนฺทํ สกฺยปุตฺตํ อารพฺภ. กิสฺมิํ วตฺถุสฺมินฺติ อายสฺมา อุปนนฺโท สกฺยปุตฺโต ปริพฺพาชเกน สทฺธิํ กยวิกฺกยํ สมาปชฺชิ, ตสฺมิํ วตฺถุสฺมิํ. เอกา ปญฺญตฺติ. ฉนฺนํ อาปตฺติ\-สมุฏฺฐานานํ ฉหิ สมุฏฺฐาเนหิ สมุฏฺฐาติ ฯเปฯ.}{โกสิยวคฺโค ทุติโย.}
\csromanmatikaanchor{mtk.9}
\csromantocmarkref{section}{3. ปตฺตวคฺค}{mtk.9}
\bookmark[dest={mtk.9},level=1]{3. ปตฺตวคฺค}
\csromanheader{1}{3. \textbf{ปตฺตวคฺค}}
\proseitem{44}{อติเรกปตฺตํ ทสาหํ อติกฺกาเมนฺตสฺส นิสฺสคฺคิยํ ปาจิตฺติยํ กตฺถ ปญฺญตฺตนฺติ สาวตฺถิยํ ปญฺญตฺตํ. กํ อารพฺภาติ ฉพฺพคฺคิเย ภิกฺขู อารพฺภ. กิสฺมิํ วตฺถุสฺมินฺติ ฉพฺพคฺคิยา ภิกฺขู อติเรกปตฺตํ ธาเรสุํ, ตสฺมิํ วตฺถุสฺมิํ. เอกา ปญฺญตฺติ, เอกา อนุปญฺญตฺติ. ฉนฺนํ อาปตฺติ\-สมุฏฺฐานานํ ทฺวีหิ สมุฏฺฐาเนหิ สมุฏฺฐาติ, สิยา กายโต จ วาจโต จ สมุฏฺฐาติ, น จิตฺตโต, สิยา กายโต จ วาจโต จ จิตฺตโต จ สมุฏฺฐาติ ฯเปฯ.}

\proseitem{45}{อูนปญฺจ\-พนฺธเนน ปตฺเตน อญฺญํ นวํ ปตฺตํ เจตาเปนฺตสฺส นิสฺสคฺคิยํ ปาจิตฺติยํ กตฺถ ปญฺญตฺตนฺติ สกฺเกสุ ปญฺญตฺตํ. กํ อารพฺภาติ ฉพฺพคฺคิเย ภิกฺขู อารพฺภ. กิสฺมิํ วตฺถุสฺมินฺติ ฉพฺพคฺคิยา ภิกฺขู อปฺป\-มตฺต\-เกนปิ ภินฺเนน อปฺป\-มตฺต\-เกนปิ ขณฺเฑน วิลิขิตมตฺเตนปิ พหู ปตฺเต วิญฺญาเปสุํ, ตสฺมิํ วตฺถุสฺมิํ. เอกา ปญฺญตฺติ. ฉนฺนํ อาปตฺติ\-สมุฏฺฐานานํ ฉหิ สมุฏฺฐาเนหิ สมุฏฺฐาติ ฯเปฯ.}

\proseitem{46}{เภสชฺชานิ ปฏิคฺคเหตฺวา สตฺตาหํ อติกฺกาเมนฺตสฺส นิสฺสคฺคิยํ ปาจิตฺติย กตฺถ ปญฺญตฺตนฺติ สาวตฺถิยํ ปญฺญตฺตํ. กํ อารพฺภาติ สมฺพหุเล ภิกฺขู อารพฺภ. กิสฺมิํ วตฺถุสฺมินฺติ สมฺพหุลา ภิกฺขู เภสชฺชานิ ปฏิคฺคเหตฺวา สตฺตาหํ อติกฺกาเมสุํ, ตสฺมิํ วตฺถุสฺมิํ. เอกา ปญฺญตฺติ. ฉนฺนํ อาปตฺติ\-สมุฏฺฐานานํ ทฺวีหิ สมุฏฺฐาเนหิ สมุฏฺฐาติ, กถินเก ฯเปฯ.}

\proseitem{47}{\csromanfolio{20}อติเรกมาเส เสเส คิมฺหาเน วสฺสิก\-สาฏิก\-จีวรํ ปริเยสนฺตสฺส นิสฺสคฺคิยํ ปาจิตฺติยํ กตฺถ ปญฺญตฺตนฺติ สาวตฺถิยํ ปญฺญตฺตํ. กํ อารพฺภาติ ฉพฺพคฺคิเย ภิกฺขู อารพฺภ. กิสฺมิํ วตฺถุสฺมินฺติ ฉพฺพคฺคิยา ภิกฺขู อติเรกมาเส เสเส คิมฺหาเน วสฺสิก\-สาฏิก\-จีวรํ ปริเยสิํสุ, ตสฺมิํ วตฺถุสฺมิํ. เอกา ปญฺญตฺติ. ฉนฺนํ อาปตฺติ\-สมุฏฺฐานานํ ฉหิ สมุฏฺฐาเนหิ สมุฏฺฐาติ ฯเปฯ.}

\proseitem{48}{ภิกฺขุสฺส สามํ จีวรํ ทตฺวา กุปิเตน อนตฺตมเนน อจฺฉินฺทนฺตสฺส นิสฺสคฺคิยํ ปาจิตฺติยํ กตฺถ ปญฺญตฺตนฺติ สาวตฺถิยํ ปญฺญตฺตํ. กํ อารพฺภาติ อายสฺมนฺตํ อุปนนฺทํ สกฺยปุตฺตํ อารพฺภ. กิสฺมิํ วตฺถุสฺมินฺติ อายสฺมา อุปนนฺโท สกฺยปุตฺโต ภิกฺขุสฺส สามํ จีวรํ ทตฺวา กุปิโต อนตฺตมโน อจฺฉินฺทิ, ตสฺมิํ วตฺถุสฺมิํ. เอกา ปญฺญตฺติ. ฉนฺนํ อาปตฺติ\-สมุฏฺฐานานํ ตีหิ สมุฏฺฐาเนหิ สมุฏฺฐาติ ฯเปฯ.}

\proseitem{49}{สามํ สุตฺตํ วิญฺญาเปตฺวา ตนฺตวาเยหิ จีวรํ วายาเปนฺตสฺส นิสฺสคฺคิยํ ปาจิตฺติยํ กตฺถ ปญฺญตฺตนฺติ ราชคเห ปญฺญตฺตํ. กํ อารพฺภาติ ฉพฺพคฺคิเย ภิกฺขู อารพฺภ. กิสฺมิํ วตฺถุสฺมินฺติ ฉพฺพคฺคิยา ภิกฺขู สามํ สุตฺตํ วิญฺญาเปตฺวา ตนฺตวาเยหิ จีวรํ วายาเปสุํ, ตสฺมิํ วตฺถุสฺมิํ. เอกา ปญฺญตฺติ. ฉนฺนํ อาปตฺติ\-สมุฏฺฐานานํ ฉหิ สมุฏฺฐาเนหิ สมุฏฺฐาติ ฯเปฯ.}

\proseitem{50}{ปุพฺเพ อปฺปวาริตสฺส อญฺญาตกสฺส คหปติกสฺส ตนฺตวาเย อุปสงฺกมิตฺวา จีวเร วิกปฺปํ อาปชฺชนฺตสฺส นิสฺสคฺคิยํ ปาจิตฺติยํ กตฺถ ปญฺญตฺตนฺติ สาวตฺถิยํ ปญฺญตฺตํ. กํ อารพฺภาติ อายสฺมนฺตํ อุปนนฺทํ สกฺยปุตฺตํ อารพฺภ. กิสฺมิํ วตฺถุสฺมินฺติ อายสฺมา อุปนนฺโท สกฺยปุตฺโต ปุพฺเพ อปฺปวาริโต อญฺญาตกสฺส คหปติกสฺส ตนฺตวาเย อุปสงฺกมิตฺวา จีวเร วิกปฺปํ อาปชฺชิ, ตสฺมิํ วตฺถุสฺมิํ. เอกา ปญฺญตฺติ. ฉนฺนํ อาปตฺติ\-สมุฏฺฐานานํ ฉหิ สมุฏฺฐาเนหิ สมุฏฺฐาติ ฯเปฯ.}

\proseitem{51}{อจฺเจกจีวรํ ปฏิคฺคเหตฺวา จีวร\-กาล\-สมยํ อติกฺกาเมนฺตสฺส นิสฺสคฺคิยํ ปาจิตฺติยํ กตฺถ ปญฺญตฺตนฺติ สาวตฺถิยํ ปญฺญตฺตํ. กํ อารพฺภาติ สมฺพหุเล ภิกฺขู อารพฺภ. กิสฺมิํ วตฺถุสฺมินฺติ สมฺพหุลา ภิกฺขู อจฺเจกจีวรํ ปฏิคฺคเหตฺวา จีวร\-กาล\-สมยํ อติกฺกาเมสุํ, ตสฺมิํ วตฺถุสฺมิํ. เอกา ปญฺญตฺติ. ฉนฺนํ อาปตฺติ\-สมุฏฺฐานานํ ทฺวีหิ สมุฏฺฐาเนหิ สมุฏฺฐาติ, กถินเก ฯเปฯ.}

\proseitem{52}{\csromanfolio{21}ติณฺณํ จีวรานํ อญฺญตรํ จีวรํ อนฺตรฆเร นิกฺขิปิตฺวา อติเรกฉารตฺตํ วิปฺป\-วส\-นฺตสฺส นิสฺสคฺคิยํ ปาจิตฺติยํ กตฺถ ปญฺญตฺตนฺติ สาวตฺถิยํ ปญฺญตฺตํ. กํ อารพฺภาติ สมฺพหุเล ภิกฺขู อารพฺภ. กิสฺมิํ วตฺถุสฺมินฺติ สมฺพหุลา ภิกฺขู ติณฺณํ จีวรานํ อญฺญตรํ จีวรํ อนฺตรฆเร นิกฺขิปิตฺวา อติเรกฉารตฺตํ วิปฺปวสิํสุ, ตสฺมิํ วตฺถุสฺมิํ. เอกา ปญฺญตฺติ. ฉนฺนํ อาปตฺติ\-สมุฏฺฐานานํ ทฺวีหิ สมุฏฺฐาเนหิ สมุฏฺฐาติ, กถินเก ฯเปฯ.}

\proseitemwithcloser{53}{ชานํ สํฆิกํ ลาภํ ปริณตํ อตฺตโน ปริณาเมนฺตสฺส นิสฺสคฺคิยํ ปาจิตฺติยํ กตฺถ ปญฺญตฺตนฺติ สาวตฺถิยํ ปญฺญตฺตํ. กํ อารพฺภาติ ฉพฺพคฺคิเย ภิกฺขู อารพฺภ. กิสฺมิํ วตฺถุสฺมินฺติ ฉพฺพคฺคิยา ภิกฺขู ชานํ สํฆิกํ ลาภํ ปริณตํ อตฺตโน ปริณาเมสุํ, ตสฺมิํ วตฺถุสฺมิํ. เอกา ปญฺญตฺติ. ฉนฺนํ อาปตฺติ\-สมุฏฺฐานานํ ตีหิ สมุฏฺฐาเนหิ}{ปตฺตวคฺโค ตติโย. ติํส นิสฺสคฺคิยา ปาจิตฺติยา นิฏฺฐิตา.}
\csromancenter{\textbf{ตสฺสุทฺทานํ.}}
\setlength{\csromangathapagewidth}{0pt}
\setlength{\csromangathawakwidth}{0pt}
\csromangathameasuregroup{ทเสกรตฺติมาโส จ, \\ อญฺญาตํ ตญฺจ\textsuperscript{1} อุทฺทิสฺส, \\ โกสิยา สุทฺธเทฺวภาคา, \\ เทฺว จ โลมานิ อุคฺคเณฺห, \\ เทฺว จ ปตฺตานิ เภสชฺชํ,}{\csromangathabat{ทเสกรตฺติมาโส จ,}{โธวนญฺจ ปฏิคฺคโห.} \\ \csromangathabat{อญฺญาตํ ตญฺจ\textsuperscript{1} อุทฺทิสฺส,}{อุภินฺนํ ทูตเกน จ.} \\ \csromangathabat{โกสิยา สุทฺธเทฺวภาคา,}{ฉพฺพสฺสานิ นิสีทนํ.} \\ \csromangathabat{เทฺว จ โลมานิ อุคฺคเณฺห,}{อุโภ นานปฺปการกา.} \\ \csromangathabat{เทฺว จ ปตฺตานิ เภสชฺชํ,}{วสฺสิกา ทานปญฺจมํ.}}
\csromangathasetleft{ทเสกรตฺติมาโส จ, \\ อญฺญาตํ ตญฺจ\textsuperscript{1} อุทฺทิสฺส, \\ โกสิยา สุทฺธเทฺวภาคา, \\ เทฺว จ โลมานิ อุคฺคเณฺห, \\ เทฺว จ ปตฺตานิ เภสชฺชํ,}
\csromangathagroup{\csromangathastanza{\csromangathabat{ทเสกรตฺติมาโส จ,}{โธวนญฺจ ปฏิคฺคโห.} \\ \csromangathabat{อญฺญาตํ ตญฺจ\csromansharedfootnote{e63ba6f4612c2c15}{อญฺญาตกญฺจ (ก)} อุทฺทิสฺส,}{อุภินฺนํ ทูตเกน จ.}}\csromangathastanza{\csromangathabat{โกสิยา สุทฺธเทฺวภาคา,}{ฉพฺพสฺสานิ นิสีทนํ.} \\ \csromangathabat{เทฺว จ โลมานิ อุคฺคเณฺห,}{อุโภ นานปฺปการกา.} \\ \csromangathabat{เทฺว จ ปตฺตานิ เภสชฺชํ,}{วสฺสิกา ทานปญฺจมํ.}}}

\vspace{\tipitakaparskip}
\csromancenter{สามํ วายาปนจฺเจโก, สาสงฺกํ สํฆิเกน จาติ.}
\csromansectionrule
\csromanmatikaanchor{mtk.10}
\csromanmatikaanchor{mtk.11}
\csromantocmarkref{chapter}{5. ปาจิตฺติยกณฺฑ}{mtk.10}
\bookmark[dest={mtk.10},level=0]{5. ปาจิตฺติยกณฺฑ}
\csromantocmarkref{section}{1. มุสาวาทวคฺค}{mtk.11}
\bookmark[dest={mtk.11},level=1]{1. มุสาวาทวคฺค}
\csromanmarkh{5. ปาจิตฺติยกณฺฑ}\csromanmarkhii{1. มุสาวาทวคฺค}\csromanheadernomark{1}{5. \textbf{ปาจิตฺติยกณฺฑ} \textbf{1. มุสาวาทวคฺค}}
\proseitem{54}{ยํ เตน ภควตา ชานตา ปสฺสตา อรหตา สมฺมา\-สมฺพุทฺเธน สมฺปชาน\-มุสาวาเท ปาจิตฺติยํ กตฺถ ปญฺญตฺตนฺติ สาวตฺถิยํ ปญฺญตฺตํ. กํ \csromanfolio{22}อารพฺภาติ หตฺถกํ สกฺยปุตฺตํ อารพฺภ. กิสฺมิํ วตฺถุสฺมินฺติ อายสฺมา หตฺถโก สกฺยปุตฺโต ติตฺถิเยหิ สทฺธิํ สลฺลปนฺโต อวชานิตฺวา ปฏิชานิ, ปฏิชานิตฺวา อวชานิ, ตสฺมิํ วตฺถุสฺมิํ. เอกา ปญฺญตฺติ. ฉนฺนํ อาปตฺติ\-สมุฏฺฐานานํ ตีหิ สมุฏฺฐาเนหิ สมุฏฺฐาติ, สิยา กายโต จ จิตฺตโต จ สมุฏฺฐาติ, น วาจโต, สิยา วาจโต จ จิตฺตโต จ สมุฏฺฐาติ, น กายโต, สิยา กายโต จ วาจโต จ จิตฺตโต จ สมุฏฺฐาติ ฯเปฯ.}

\proseitem{55}{โอมสวาเท ปาจิตฺติยํ กตฺถ ปญฺญตฺตนฺติ สาวตฺถิยํ ปญฺญตฺตํ. กํ อารพฺภาติ ฉพฺพคฺคิเย ภิกฺขู อารพฺภ. กิสฺมิํ วตฺถุสฺมินฺติ ฉพฺพคฺคิยา ภิกฺขู เปสเลหิ ภิกฺขูหิ สทฺธิํ ภณฺฑนฺตา\csromansharedfootnote{8777e97785fb71da}{ภณฺเฑนฺตา (ก)} เปสเล ภิกฺขู โอมสิํสุ, ตสฺมิํ วตฺถุสฺมิํ. เอกา ปญฺญตฺติ. ฉนฺนํ อาปตฺติ\-สมุฏฺฐานานํ ตีหิ สมุฏฺฐาเนหิ สมุฏฺฐาติ ฯเปฯ.}

\proseitem{56}{ภิกฺขุเปสุญฺเญ ปาจิตฺติยํ กตฺถ ปญฺญตฺตนฺติ สาวตฺถิยํ ปญฺญตฺตํ. กํ อารพฺภาติ ฉพฺพคฺคิเย ภิกฺขู อารพฺภ. กิสฺมิํ วตฺถุสฺมินฺติ ฉพฺพคฺคิยา ภิกฺขู ภิกฺขูนํ ภณฺฑน\-ชาตานํ กลหชาตานํ วิวาทาปนฺนานํ เปสุญฺญํ อุปสํหริํสุ, ตสฺมิํ วตฺถุสฺมิํ. เอกา ปญฺญตฺติ. ฉนฺนํ อาปตฺติ\-สมุฏฺฐานานํ ตีหิ สมุฏฺฐาเนหิ}

\proseitem{57}{อนุปสมฺปนฺนํ ปทโส ธมฺมํ วาเจนฺตสฺส ปาจิตฺติยํ กตฺถ ปญฺญตฺตนฺติ สาวตฺถิยํ ปญฺญตฺตํ. กํ อารพฺภาติ ฉพฺพคฺคิเย ภิกฺขู อารพฺภ. กิสฺมิํ วตฺถุสฺมินฺติ ฉพฺพคฺคิยา ภิกฺขู อุปาสเก ปทโส ธมฺมํ วาเจสุํ, ตสฺมิํ วตฺถุสฺมิํ. เอกา ปญฺญตฺติ. ฉนฺนํ อาปตฺติ\-สมุฏฺฐานานํ ทฺวีหิ สมุฏฺฐาเนหิ สมุฏฺฐาติ, สิยา วาจโต สมุฏฺฐาติ, น กายโต, น จิตฺตโต, สิยา วาจโต จ จิตฺตโต จ สมุฏฺฐาติ, น กายโต ฯเปฯ.}

\proseitem{58}{อนุปสมฺปนฺเนน อุตฺตริทิรตฺตติรตฺตํ สหเสยฺยํ กปฺเปนฺตสฺส ปาจิตฺติยํ กตฺถ ปญฺญตฺตนฺติ อาฬวิยํ ปญฺญตฺตํ. กํ อารพฺภาติ สมฺพหุเล ภิกฺขู อารพฺภ. กิสฺมิํ วตฺถุสฺมินฺติ สมฺพหุลา ภิกฺขู อนุปสมฺปนฺเนน สหเสยฺยํ กปฺเปสุํ, ตสฺมิํ วตฺถุสฺมิํ. เอกา ปญฺญตฺติ, เอกา อนุปญฺญตฺติ. \csromanfolio{23}ฉนฺนํ อาปตฺติ\-สมุฏฺฐานานํ ทฺวีหิ สมุฏฺฐาเนหิ สมุฏฺฐาติ, สิยา กายโต สมุฏฺฐาติ, น วาจโต, น จิตฺตโต, สิยา กายโต จ จิตฺตโต จ สมุฏฺฐาติ, น วาจโต ฯเปฯ.}

\proseitem{59}{มาตุคาเมน สหเสยฺยํ กปฺเปนฺตสฺส ปาจิตฺติยํ กตฺถ ปญฺญตฺตนฺติ สาวตฺถิยํ ปญฺญตฺตํ. กํ อารพฺภาติ อายสฺมนฺตํ อนุรุทฺธํ อารพฺภ. กิสฺมิํ วตฺถุสฺมินฺติ อายสฺมา อนุรุทฺโธ มาตุคาเมน สหเสยฺยํ กปฺเปสิ, ตสฺมิํ วตฺถุสฺมิํ. เอกา ปญฺญตฺติ. ฉนฺนํ อาปตฺติ\-สมุฏฺฐานานํ ทฺวีหิ สมุฏฺฐาเนหิ สมุฏฺฐาติ, เอฬกโลมเก ฯเปฯ.}

\proseitem{60}{มาตุคามสฺส อุตฺตริ\-ฉปฺปญฺจ\-วาจาหิ ธมฺมํ เทเสนฺตสฺส ปาจิตฺติยํ กตฺถ ปญฺญตฺตนฺติ สาวตฺถิยํ ปญฺญตฺตํ. กํ อารพฺภาติ อายสฺมนฺตํ อุทายิํ อารพฺภ. กิสฺมิํ วตฺถุสฺมินฺติ อายสฺมา อุทายี มาตุคามสฺส ธมฺมํ เทเสสิ, ตสฺมิํ วตฺถุสฺมิํ. เอกา ปญฺญตฺติ, เทฺว อนุปญฺญตฺติโย. ฉนฺนํ อาปตฺติ\-สมุฏฺฐานานํ ทฺวีหิ สมุฏฺฐาเนหิ สมุฏฺฐาติ, ปทโสธมฺเม ฯเปฯ.}

\proseitem{61}{อนุปสมฺปนฺนสฺส อุตฺตริ\-มนุสฺสธมฺมํ ภูตํ อาโรเจนฺตสฺส ปาจิตฺติยํ กตฺถ ปญฺญตฺตนฺติ เวสาลิยํ ปญฺญตฺตํ. กํ อารพฺภาติ วคฺคุมุทา\-ตีริเย ภิกฺขู อารพฺภ. กิสฺมิํ วตฺถุสฺมินฺติ วคฺคุมุทา\-ตีริยา ภิกฺขู คิหีนํ อญฺญมญฺญสฺส อุตฺตริ\-มนุสฺสธมฺมสฺส วณฺณํ ภาสิํสุ, ตสฺมิํ วตฺถุสฺมิํ. เอกา ปญฺญตฺติ. ฉนฺนํ อาปตฺติ\-สมุฏฺฐานานํ ตีหิ สมุฏฺฐาเนหิ สมุฏฺฐาติ, สิยา กายโต สมุฏฺฐาติ, น วาจโต, น จิตฺตโต, สิยา วาจโต สมุฏฺฐาติ, น กายโต, น จิตฺตโต, สิยา กายโต จ วาจโต จ สมุฏฺฐาติ, น จิตฺตโต ฯเปฯ.}

\proseitem{62}{ภิกฺขุสฺส ทุฏฺฐุลฺลํ อาปตฺติํ อนุปสมฺปนฺนสฺส อาโรเจนฺตสฺส ปาจิตฺติยํ กตฺถ ปญฺญตฺตนฺติ สาวตฺถิยํ ปญฺญตฺตํ. กํ อารพฺภาติ ฉพฺพคฺคิเย ภิกฺขู อารพฺภ. กิสฺมิํ วตฺถุสฺมินฺติ ฉพฺพคฺคิยา ภิกฺขู ภิกฺขุสฺส ทุฏฺฐุลฺลา\-ปตฺติํ อนุปสมฺปนฺนสฺส อาโรเจสุํ, ตสฺมิํ วตฺถุสฺมิํ. เอกา ปญฺญตฺติ. ฉนฺนํ อาปตฺติ\-สมุฏฺฐานานํ ตีหิ สมุฏฺฐาเนหิ}

\proseitemwithclosermidruled{63}{ปถวิํ ขณนฺตสฺส ปาจิตฺติยํ กตฺถ ปญฺญตฺตนฺติ อาฬวิยํ ปญฺญตฺตํ. กํ อารพฺภาติ อาฬวเก ภิกฺขู อารพฺภ. กิสฺมิํ วตฺถุสฺมินฺติ \csromanfolio{24}อาฬวกา ภิกฺขู ปถวิํ ขณิํสุ, ตสฺมิํ วตฺถุสฺมิํ. เอกา ปญฺญตฺติ. ฉนฺนํ อาปตฺติ\-สมุฏฺฐานานํ ตีหิ สมุฏฺฐาเนหิ สมุฏฺฐาติ ฯเปฯ.}{มุสาวาทวคฺโค ปฐโม.}
\csromanmatikaanchor{mtk.12}
\csromantocmarkref{section}{2. ภูตคามวคฺค}{mtk.12}
\bookmark[dest={mtk.12},level=1]{2. ภูตคามวคฺค}
\csromanheader{1}{2. \textbf{ภูตคามวคฺค}}
\proseitem{64}{ภูตคาม\-ปาตพฺยตาย ปาจิตฺติยํ กตฺถ ปญฺญตฺตนฺติ อาฬวิยํ ปญฺญตฺตํ. กํ อารพฺภาติ อาฬวเก ภิกฺขู อารพฺภ. กิสฺมิํ วตฺถุสฺมินฺติ อาฬวกา ภิกฺขู รุกฺขํ ฉินฺทิํสุ, ตสฺมิํ วตฺถุสฺมิํ. เอกา ปญฺญตฺติ. ฉนฺนํ อาปตฺติ\-สมุฏฺฐานานํ ตีหิ สมุฏฺฐาเนหิ สมุฏฺฐาติ ฯเปฯ.}

\proseitem{65}{อญฺญวาทเก วิเหสเก ปาจิตฺติยํ กตฺถ ปญฺญตฺตนฺติ โกสมฺพิยํ ปญฺญตฺตํ. กํ อารพฺภาติ อายสฺมนฺตํ ฉนฺนํ อารพฺภ. กิสฺมิํ วตฺถุสฺมินฺติ อายสฺมา ฉนฺโน สํฆมชฺเฌ อาปตฺติยา อนุยุญฺชิยมาโน อญฺเญนญฺญํ ปฏิจริ, ตสฺมิํ วตฺถุสฺมิํ. เอกา ปญฺญตฺติ, เอกา อนุปญฺญตฺติ. ฉนฺนํ อาปตฺติ\-สมุฏฺฐานานํ ตีหิ สมุฏฺฐาเนหิ สมุฏฺฐาติ ฯเปฯ.}

\proseitem{66}{อุชฺฌาปนเก ขิยฺยนเก ปาจิตฺติยํ กตฺถ ปญฺญตฺตนฺติ ราชคเห ปญฺญตฺตํ. กํ อารพฺภาติ เมตฺติย\-ภูมชเก ภิกฺขู อารพฺภ. กิสฺมิํ วตฺถุสฺมินฺติ เมตฺติย\-ภูมชกา ภิกฺขู อายสฺมนฺตํ ทพฺพํ มลฺลปุตฺตํ ภิกฺขู อุชฺฌาเปสุํ, ตสฺมิํ วตฺถุสฺมิํ. เอกา ปญฺญตฺติ, เอกา อนุปญฺญตฺติ. ฉนฺนํ อาปตฺติ\-สมุฏฺฐานานํ ตีหิ สมุฏฺฐาเนหิ สมุฏฺฐาติ ฯเปฯ.}

\proseitem{67}{สํฆิกํ มญฺจํ วา ปีฐํ วา ภิสิํ วา โกจฺฉํ วา อชฺโฌกาเส สนฺถริตฺวา อนุทฺธริตฺวา อนาปุจฺฉา ปกฺกมนฺตสฺส ปาจิตฺติยํ กตฺถ ปญฺญตฺตนฺติ สาวตฺถิยํ ปญฺญตฺตํ. กํ อารพฺภาติ สมฺพหุเล ภิกฺขู อารพฺภ. กิสฺมิํ วตฺถุสฺมินฺติ สมฺพหุลา ภิกฺขู สํฆิกํ เสนาสนํ อชฺโฌกาเส สนฺถริตฺวา อนุทฺธริตฺวา อนาปุจฺฉา ปกฺกมิํสุ, ตสฺมิํ วตฺถุสฺมิํ. เอกา ปญฺญตฺติ, เอกา อนุปญฺญตฺติ. ฉนฺนํ อาปตฺติ\-สมุฏฺฐานานํ ทฺวีหิ สมุฏฺฐาเนหิ สมุฏฺฐาติ, กถินเก ฯเปฯ.}

\proseitem{68}{\csromanfolio{25}สํฆิเก วิหาเร เสยฺยํ สนฺถริตฺวา อนุทฺธริตฺวา อนาปุจฺฉา ปกฺกมนฺตสฺส ปาจิตฺติยํ กตฺถ ปญฺญตฺตนฺติ สาวตฺถิยํ ปญฺญตฺตํ. กํ อารพฺภาติ สตฺตรส\-วคฺคิเย ภิกฺขู อารพฺภ. กิสฺมิํ วตฺถุสฺมินฺติ สตฺตรส\-วคฺคิยา ภิกฺขู สํฆิเก วิหาเร เสยฺยํ สนฺถริตฺวา อนุทฺธริตฺวา อนาปุจฺฉา ปกฺกมิํสุ, ตสฺมิํ วตฺถุสฺมิํ. เอกา ปญฺญตฺติ. ฉนฺนํ อาปตฺติ\-สมุฏฺฐานานํ ทฺวีหิ สมุฏฺฐาเนหิ สมุฏฺฐาติ, กถินเก ฯเปฯ.}

\proseitem{69}{สํฆิเก วิหาเร ชานํ ปุพฺพุปคตํ ภิกฺขุํ อนุปขชฺช เสยฺยํ กปฺเปนฺตสฺส ปาจิตฺติยํ กตฺถ ปญฺญตฺตนฺติ สาวตฺถิยํ ปญฺญตฺตํ. กํ อารพฺภาติ ฉพฺพคฺคิเย ภิกฺขู อารพฺภ. กิสฺมิํ วตฺถุสฺมินฺติ ฉพฺพคฺคิยา ภิกฺขู เถเร ภิกฺขู อนุปขชฺช เสยฺยํ กปฺเปสุํ, ตสฺมิํ วตฺถุสฺมิํ. เอกา ปญฺญตฺติ. ฉนฺนํ อาปตฺติ\-สมุฏฺฐานานํ เอเกน สมุฏฺฐาเนน สมุฏฺฐาติ, กายโต จ จิตฺตโต จ สมุฏฺฐาติ, น วาจโต ฯเปฯ.}

\proseitem{70}{ภิกฺขุํ กุปิเตน อนตฺตมเนน สํฆิกา วิหารา นิกฺกฑฺฒ\-นฺตสฺส ปาจิตฺติยํ กตฺถ ปญฺญตฺตนฺติ สาวตฺถิยํ ปญฺญตฺตํ. กํ อารพฺภาติ ฉพฺพคฺคิเย ภิกฺขู อารพฺภ. กิสฺมิํ วตฺถุสฺมินฺติ ฉพฺพคฺคิยา ภิกฺขู กุปิตา อนตฺตมนา ภิกฺขู สํฆิกา วิหารา นิกฺกฑฺฒิํสุ, ตสฺมิํ วตฺถุสฺมิํ. เอกา ปญฺญตฺติ. ฉนฺนํ อาปตฺติ\-สมุฏฺฐานานํ ตีหิ สมุฏฺฐาเนหิ}

\proseitem{71}{สํฆิเก วิหาเร อุปริ\-เวหาส\-กุฏิยา อาหจฺจปาทกํ มญฺจํ วา ปีฐํ วา อภินิ\-สีท\-นฺตสฺส ปาจิตฺติยํ กตฺถ ปญฺญตฺตนฺติ สาวตฺถิยํ ปญฺญตฺตํ. กํ อารพฺภาติ อญฺญตรํ ภิกฺขุํ อารพฺภ. กิสฺมิํ วตฺถุสฺมินฺติ อญฺญตโร ภิกฺขุ สํฆิเก วิหาเร อุปริ\-เวหาส\-กุฏิยา อาหจฺจปาทกํ มญฺจํ สหสา อภินิสีทิ, ตสฺมิํ วตฺถุสฺมิํ. เอกา ปญฺญตฺติ. ฉนฺนํ อาปตฺติ\-สมุฏฺฐานานํ ทฺวีหิ สมุฏฺฐาเนหิ สมุฏฺฐาติ, สิยา กายโต สมุฏฺฐาติ, น วาจโต, นจิตฺตโต, สิยา กายโต จ จิตฺตโต จ สมุฏฺฐาติ, น วาจโต ฯเปฯ.}

\proseitem{72}{ทฺวตฺติปริยาเย อธิฏฺฐหิตฺวา ตตุตฺตริ อธิฏฺฐห\-นฺตสฺส ปาจิตฺติยํ กตฺถ ปญฺญตฺตนฺติ โกสมฺพิยํ ปญฺญตฺตํ. กํ อารพฺภาติ อายสฺมนฺตํ ฉนฺนํ อารพฺภ. กิสฺมิํ วตฺถุสฺมินฺติ อายสฺมา ฉนฺโน กตปริโยสิตํ วิหารํ ปุนปฺปุนํ ฉาทาเปสิ ปุนปฺปุนํ ลิมฺปาเปสิ, อติภาริโก วิหาโร \csromanfolio{26}ปริปติ, ตสฺมิํ วตฺถุสฺมิํ. เอกา ปญฺญตฺติ. ฉนฺนํ อาปตฺติ\-สมุฏฺฐานานํ ฉหิ สมุฏฺฐาเนหิ สมุฏฺฐาติ ฯเปฯ.}

\proseitemwithclosermidruled{73}{ชานํ สปฺปาณกํ อุทกํ ติณํ วา มตฺติกํ วา สิญฺจนฺตสฺส ปาจิตฺติยํ กตฺถ ปญฺญตฺตนฺติ อาฬวิยํ ปญฺญตฺตํ. กํ อารพฺภาติ อาฬวเก ภิกฺขู อารพฺภ. กิสฺมิํ วตฺถุสฺมินฺติ อาฬวกา ภิกฺขู ชานํ สปฺปาณกํ อุทกํ ติณมฺปิ มตฺติกมฺปิ สิญฺจิํสุ, ตสฺมิํ วตฺถุสฺมิํ. เอกา ปญฺญตฺติ. ฉนฺนํ อาปตฺติ\-สมุฏฺฐานานํ ตีหิ สมุฏฺฐาเนหิ สมุฏฺฐาติ ฯเปฯ.}{ภูตคามวคฺโค ทุติโย.}
\csromanmatikaanchor{mtk.13}
\csromantocmarkref{section}{3. โอวาทวคฺค}{mtk.13}
\bookmark[dest={mtk.13},level=1]{3. โอวาทวคฺค}
\csromanheader{1}{3. \textbf{โอวาทวคฺค}}
\proseitem{74}{อสมฺมเตน ภิกฺขุนิโย โอวทนฺตสฺส ปาจิตฺติยํ กตฺถ ปญฺญตฺตนฺติ สาวตฺถิยํ ปญฺญตฺตํ. กํ อารพฺภาติ ฉพฺพคฺคิเย ภิกฺขู อารพฺภ. กิสฺมิํ วตฺถุสฺมินฺติ ฉพฺพคฺคิยา ภิกฺขู อสมฺมตา ภิกฺขุนิโย โอวทิํสุ, ตสฺมิํ วตฺถุสฺมิํ. อตฺถิ ตตฺถ ปญฺญตฺติ อนุปญฺญตฺติ อนุปฺปนฺนปญฺญตฺตีติ เอกา ปญฺญตฺติ, เอกา อนุปญฺญตฺติ, อนุปฺปนฺนปญฺญตฺติ ตสฺมิํ นตฺถิ. ฉนฺนํ อาปตฺติ\-สมุฏฺฐานานํ ทฺวีหิ สมุฏฺฐาเนหิ สมุฏฺฐาติ, สิยา วาจโต สมุฏฺฐาติ, น กายโต, น จิตฺตโต, สิยา วาจโต จ จิตฺตโต จ สมุฏฺฐาติ, น กายโต ฯเปฯ.}

\proseitem{75}{อตฺถงฺคเต สูริเย ภิกฺขุนิโย โอวทนฺตสฺส ปาจิตฺติยํ กตฺถ ปญฺญตฺตนฺติ สาวตฺถิยํ ปญฺญตฺตํ. กํ อารพฺภาติ อายสฺมนฺตํ จูฬปนฺถกํ อารพฺภ. กิสฺมิํ วตฺถุสฺมินฺติ อายสฺมา จูฬปนฺถโก อตฺถงฺคเต สูริเย ภิกฺขุนิโย โอวทิ, ตสฺมิํ วตฺถุสฺมิํ. เอกา ปญฺญตฺติ. ฉนฺนํ อาปตฺติ\-สมุฏฺฐานานํ ทฺวีหิ สมุฏฺฐาเนหิ สมุฏฺฐาติ, ปทโสธมฺเม ฯเปฯ.}

\proseitem{76}{ภิกฺขุนุ\-ปสฺสยํ อุปสงฺกมิตฺวา ภิกฺขุนิโย โอวทนฺตสฺส ปาจิตฺติยํ กตฺถ ปญฺญตฺตนฺติ สกฺเกสุ ปญฺญตฺตํ. กํ อารพฺภาติ ฉพฺพคฺคิเย ภิกฺขู อารพฺภ. กิสฺมิํ วตฺถุสฺมินฺติ ฉพฺพคฺคิยา ภิกฺขู ภิกฺขุนุ\-ปสฺสยํ อุปสงฺกมิตฺวา ภิกฺขุนิโย โอวทิํสุ, ตสฺมิํ วตฺถุสฺมิํ. เอกา ปญตฺติ, เอกา อนุปญฺญตฺติ. ฉนฺนํ อาปตฺติ\-สมุฏฺฐานานํ ทฺวีหิ สมุฏฺฐาเนหิ สมุฏฺฐาติ, กถินเก ฯเปฯ.}

\proseitem{77}{\csromanfolio{27}“อามิสเหตุ ภิกฺขู ภิกฺขุนิโย โอวทนฺตี”ติ ภณนฺตสฺส ปาจิตฺติยํ กตฺถ ปญฺญตฺตนฺติ สาวตฺถิยํ ปญฺญตฺตํ. กํ อารพฺภาติ ฉพฺพคฺคิเย ภิกฺขู อารพฺภ. กิสฺมิํ วตฺถุสฺมินฺติ ฉพฺพคฺคิยา ภิกฺขู “อามิสเหตุ ภิกฺขู ภิกฺขุนิโย โอวทนฺตี”ติ ภณิํสุ, ตสฺมิํ วตฺถุสฺมิํ. เอกา ปญฺญตฺติ. ฉนฺนํ อาปตฺติ\-สมุฏฺฐานานํ ตีหิ สมุฏฺฐาเนหิ สมุฏฺฐาติ ฯเปฯ.}

\proseitem{78}{อญฺญาติกาย ภิกฺขุนิยา จีวรํ เทนฺตสฺส ปาจิตฺติยํ กตฺถ ปญฺญตฺตนฺติ สาวตฺถิยํ ปญฺญตฺตํ. กํ อารพฺภาติ อญฺญตรํ ภิกฺขุํ อารพฺภ. กิสฺมิํ วตฺถุสฺมินฺติ อญฺญตโร ภิกฺขุ อญฺญาติกาย ภิกฺขุนิยา จีวรํ อทาสิ, ตสฺมิํ วตฺถุสฺมิํ. เอกา ปญฺญตฺติ, เอกา อนุปญฺญตฺติ. ฉนฺนํ อาปตฺติ\-สมุฏฺฐานานํ ฉหิ สมุฏฺฐาเนหิ สมุฏฺฐาติ ฯเปฯ.}

\proseitem{79}{อญฺญาติกาย ภิกฺขุนิยา จีวรํ สิพฺเพนฺตสฺส ปาจิตฺติยํ กตฺถ ปญฺญตฺตนฺติ สาวตฺถิยํ ปญฺญตฺตํ. กํ อารพฺภาติ อายสฺมนฺตํ อุทายิํ อารพฺภ. กิสฺมิํ วตฺถุสฺมินฺติ อายสฺมา อุทายี อญฺญาติกาย ภิกฺขุนิยา จีวรํ สิพฺเพสิ, ตสฺมิํ วตฺถุสฺมิํ. เอกา ปญฺญตฺติ. ฉนฺนํ อาปตฺติ\-สมุฏฺฐานานํ ฉหิ สมุฏฺฐาเนหิ สมุฏฺฐาติ ฯเปฯ.}

\proseitem{80}{ภิกฺขุนิยา สทฺธิํ สํวิธาย เอก\-ทฺธาน\-มคฺคํ ปฏิปชฺชนฺตสฺส ปาจิตฺติยํ กตฺถ ปญฺญตฺตนฺติ สาวตฺถิยํ ปญฺญตฺตํ. กํ อารพฺภาติ ฉพฺพคฺคิเย ภิกฺขู อารพฺภ. กิสฺมิํ วตฺถุสฺมินฺติ ฉพฺพคฺคิยา ภิกฺขู ภิกฺขุนีหิ สทฺธิํ สํวิธาย เอก\-ทฺธาน\-มคฺคํ ปฏิปชฺชิํสุ, ตสฺมิํ วตฺถุสฺมิํ. เอกา ปญฺญตฺติ, เอกา อนุปญฺญตฺติ. ฉนฺนํ อาปตฺติ\-สมุฏฺฐานานํ จตูหิ สมุฏฺฐาเนหิ สมุฏฺฐาติ, สิยา กายโต สมุฏฺฐาติ, น วาจโต, น จิตฺตโต, สิยา กายโต จ วาจโต จ สมุฏฺฐาติ, น จิตฺตโต, สิยา กายโต จ จิตฺตโต จ สมุฏฺฐาติ, น วาจโต, สิยา กายโต จ วาจโต จ จิตฺตโต จ}

\proseitem{81}{ภิกฺขุนิยา สทฺธิํ สํวิธาย เอกํ นาวํ อภิรุหนฺตสฺส ปาจิตฺติยํ กตฺถ ปญฺญตฺตนฺติ สาวตฺถิยํ ปญฺญตฺตํ. กํ อารพฺภาติ ฉพฺพคฺคิเย ภิกฺขู อารพฺภ. กิสฺมิํ วตฺถุสฺมินฺติ ฉพฺพคฺคิยา ภิกฺขู ภิกฺขุนีหิ สทฺธิํ สํวิธาย เอกํ นาวํ อภิรุหิํสุ, ตสฺมิํ วตฺถุสฺมิํ. เอกา ปญฺญตฺติ, เอกา อนุปญฺญตฺติ. ฉนฺนํ อาปตฺติ\-สมุฏฺฐานานํ จตูหิ สมุฏฺฐาเนหิ สมุฏฺฐาติ ฯเปฯ.}

\proseitem{82}{\csromanfolio{28}ชานํ ภิกฺขุนิ\-ปริปาจิตํ ปิณฺฑปาตํ ภุญฺชนฺตสฺส ปาจิตฺติยํ กตฺถ ปญฺญตฺตนฺติ ราชคเห ปญฺญตฺตํ. กํ อารพฺภาติ เทวทตฺตํ อารพฺภ. กิสฺมิํ วตฺถุสฺมินฺติ เทวทตฺโต ชานํ ภิกฺขุนิ\-ปริปาจิตํ ปิณฺฑปาตํ ภุญฺชิ, ตสฺมิํ วตฺถุสฺมิํ. เอกา ปญฺญตฺติ, เอกา อนุปญฺญตฺติ. ฉนฺนํ อาปตฺติ\-สมุฏฺฐานานํ เอเกน สมุฏฺฐาเนน สมุฏฺฐาติ, กายโต จ จิตฺตโต จ สมุฏฺฐาติ, น วาจโต ฯเปฯ.}

\proseitemwithclosermidruled{83}{ภิกฺขุนิยา สทฺธิํ เอโก เอกาย รโห นิสชฺชํ กปฺเปนฺตสฺส ปาจิตฺติยํ กตฺถ ปญฺญตฺตนฺติ สาวตฺถิยํ ปญฺญตฺตํ. กํ อารพฺภาติ อายสฺมนฺตํ อุทายิํ อารพฺภ. กิสฺมิํ วตฺถุสฺมินฺติ อายสฺมา อุทายี ภิกฺขุนิยา สทฺธิํ เอโก เอกาย รโห นิสชฺชํ กปฺเปสิ, ตสฺมิํ วตฺถุสฺมิํ. เอกา ปญฺญตฺติ. ฉนฺนํ อาปตฺติ\-สมุฏฺฐานานํ เอเกน สมุฏฺฐาเนน สมุฏฺฐาติ, กายโต จ จิตฺตโต จ สมุฏฺฐาติ, น วาจโต ฯเปฯ.}{โอวาทวคฺโค ตติโย.}
\csromanmatikaanchor{mtk.14}
\csromantocmarkref{section}{4. โภชนวคฺค}{mtk.14}
\bookmark[dest={mtk.14},level=1]{4. โภชนวคฺค}
\csromanheader{1}{4. \textbf{โภชนวคฺค}}
\proseitem{84}{ตตุตฺตริ อาวสถปิณฺฑํ ภุญฺชนฺตสฺส ปาจิตฺติยํ กตฺถ ปญฺญตฺตนฺติ สาวตฺถิยํ ปญฺญตฺตํ. กํ อารพฺภาติ ฉพฺพคฺคิเย ภิกฺขู อารพฺภ. กิสฺมิํ วตฺถุสฺมินฺติ ฉพฺพคฺคิยา ภิกฺขู อนุวสิตฺวา อนุวสิตฺวา อาวสถปิณฺฑํ ภุญฺชิํสุ, ตสฺมิํ วตฺถุสฺมิํ. เอกา ปญฺญตฺติ, เอกา อนุปญฺญตฺติ. ฉนฺนํ อาปตฺติ\-สมุฏฺฐานานํ ทฺวีหิ สมุฏฺฐาเนหิ สมุฏฺฐาติ, เอฬกโลมเก ฯเปฯ.}

\proseitem{85}{คณโภชเน ปาจิตฺติยํ กตฺถ ปญฺญตฺตนฺติ ราชคเห ปญฺญตฺตํ. กํ อารพฺภาติ เทวทตฺตํ อารพฺภ. กิสฺมิํ วตฺถุสฺมินฺติ เทวทตฺโต สปริโส กุเลสุ วิญฺญาเปตฺวา วิญฺญาเปตฺวา ภุญฺชิ, ตสฺมิํ วตฺถุสฺมิํ. เอกา ปญฺญตฺติ, สตฺต อนุปญฺญตฺติโย. ฉนฺนํ อาปตฺติ\-สมุฏฺฐานานํ ทฺวีหิ สมุฏฺฐาเนหิ สมุฏฺฐาติ, เอฬกโลมเก ฯเปฯ.}

\proseitem{86}{ปรมฺปร\-โภชเน ปาจิตฺติยํ กตฺถ ปญฺญตฺตนฺติ เวสาลิยํ ปญฺญตฺตํ. กํ อารพฺภาติ สมฺพหุเล ภิกฺขู อารพฺภ. กิสฺมิํ วตฺถุสฺมินฺติ สมฺพหุลา \csromanfolio{29}ภิกฺขู อญฺญตฺร นิมนฺติตา อญฺญตฺร ภุญฺชิํสุ, ตสฺมิํ วตฺถุสฺมิํ. เอกา ปญฺญตฺติ, จตสฺโส อนุปญฺญตฺติโย. ฉนฺนํ อาปตฺติ\-สมุฏฺฐานานํ ทฺวีหิ สมุฏฺฐาเนหิ สมุฏฺฐาติ, กถินเก ฯเปฯ.}

\proseitem{87}{ทฺวตฺติ\-ปตฺตปูเร ปูเว ปฏิคฺคเหตฺวา ตตุตฺตริ ปฏิคฺคณฺหนฺตสฺส ปาจิตฺติยํ กตฺถ ปญฺญตฺตนฺติ สาวตฺถิยํ ปญฺญตฺตํ. กํ อารพฺภาติ สมฺพหุเล ภิกฺขู อารพฺภ. กิสฺมิํ วตฺถุสฺมินฺติ สมฺพหุลา ภิกฺขู น มตฺตํ ชานิตฺวา ปฏิคฺคเหสุํ, ตสฺมิํ วตฺถุสฺมิํ. เอกา ปญฺญตฺติ. ฉนฺนํ อาปตฺติ\-สมุฏฺฐานานํ ฉหิ สมุฏฺฐาเนหิ สมุฏฺฐาติ ฯเปฯ.}

\proseitem{88}{ภุตฺตาวินา ปวาริเตน อนติริตฺตํ ขาทนียํ วา โภชนียํ วา ภุญฺชนฺตสฺส ปาจิตฺติยํ กตฺถ ปญฺญตฺตนฺติ สาวตฺถิยํ ปญฺญตฺตํ. กํ อารพฺภาติ สมฺพหุเล ภิกฺขู อารพฺภ. กิสฺมิํ วตฺถุสฺมินฺติ สมฺพหุลา ภิกฺขู ภุตฺตาวี ปวาริตา อญฺญตฺร ภุญฺชิํสุ, ตสฺมิํ วตฺถุสฺมิํ. เอกา ปญฺญตฺติ, เอกา อนุปญฺญตฺติ. ฉนฺนํ อาปตฺติ\-สมุฏฺฐานานํ ทฺวีหิ สมุฏฺฐาเนหิ สมุฏฺฐาติ, กถินเก ฯเปฯ.}

\proseitem{89}{ภิกฺขุํ ภุตฺตาวิํ ปวาริตํ อนติริตฺเตน ขาทนีเยน วา โภชนีเยน วา อภิหฏฺฐุํ ปวาเรนฺตสฺส ปาจิตฺติยํ กตฺต ปญฺญตฺตนฺติ สาวตฺถิยํ ปญฺญตฺตํ. กํ อารพฺภติ อญฺญตรํ ภิกฺขุํ อารพฺภ. กิสฺมิํ วตฺถุสฺมินฺติ อญฺญตโร ภิกฺขุ ภิกฺขุํ ภุตฺตาวิํ ปวาริตํ อนติริตฺเตน โภชนีเยน อภิหฏฺฐุํ ปวาเรสิ, ตสฺมิํ วตฺถุสฺมิํ. เอกา ปญฺญตฺติ. ฉนฺนํ อาปตฺติ\-สมุฏฺฐานานํ ตีหิ สมุฏฺฐาเนหิ สมุฏฺฐาติ ฯเปฯ.}

\proseitem{90}{วิกาเล ขาทนียํ วา โภชนียํ วา ภุญฺชนฺตสฺส ปาจิตฺติยํ กตฺถ ปญฺญตฺตนฺติ ราชคเห ปญฺญตฺตํ. กํ อารพฺภาติ สตฺตรส\-วคฺคิเย ภิกฺขู อารพฺภ. กิสฺมิํ วตฺถุสฺมินฺติ สตฺตรส\-วคฺคิยา ภิกฺขู วิกาเล โภชนํ ภุญฺชิํสุ, ตสฺมิํ วตฺถุสฺมิํ. เอกา ปญฺญตฺติ. ฉนฺนํ อาปตฺติ\-สมุฏฺฐานานํ ทฺวีหิ สมุฏฺฐาเนหิ สมุฏฺฐาติ, เอฬกโลมเก ฯเปฯ.}

\proseitem{91}{สนฺนิธิการกํ ขาทนียํ วา โภชนียํ วาภุญฺชนฺตสฺส ปาจิตฺติยํ กตฺถ ปญฺญตฺตนฺติ สาวตฺถิยํ ปญฺญตฺตํ. กํ อารพฺภาติ อายสฺมนฺตํ เพลฏฺฐสีสํ อารพฺภ. กิสฺมิํ วตฺถุสฺมินฺติ อายสฺมา เพลฏฺฐสีโส สนฺนิธิการกํ \csromanfolio{30}โภชนํ ภุญฺชิ, ตสฺมิํ วตฺถุสฺมิํ. เอกา ปญฺญตฺติ. ฉนฺนํ อาปตฺติ\-สมุฏฺฐานานํ ทฺวีหิ สมุฏฺฐาเนหิ สมุฏฺฐาติ, เอฬกโลมเก ฯเปฯ.}

\proseitem{92}{ปณีต\-โภชนานิ อตฺตโน อตฺถาย วิญฺญาเปตฺวา ภุญฺชนฺตสฺส ปาจิตฺติยํ กตฺถ ปญฺญตฺตนฺติ สาวตฺถิยํ ปญฺญตฺตํ. กํ อารพฺภาติ ฉพฺพคฺคิเย ภิกฺขู อารพฺภ. กิสฺมิํ วตฺถุสฺมินฺติ ฉพฺพคฺคิยา ภิกฺขู ปณีต\-โภชนานิ อตฺตโน อตฺถาย วิญฺญาเปตฺวา ภุญฺชิํสุ, ตสฺมิํ วตฺถุสฺมิํ. เอกา ปญฺญตฺติ, เอกา อนุปญฺญตฺติ. ฉนฺนํ อาปตฺติ\-สมุฏฺฐานานํ จตูหิ สมุฏฺฐาเนหิ สมุฏฺฐาติ ฯเปฯ.}

\proseitemwithclosermidruled{93}{อทินฺนํ มุขทฺวารํ อาหารํ อาหรนฺตสฺส ปาจิตฺติยํ กตฺถ ปญฺญตฺตนฺติ เวสาลิยํ ปญฺญตฺตํ. กํ อารพฺภาติ อญฺญตรํ ภิกฺขุํ อารพฺภ. กิสฺมิํ วตฺถุสฺมินฺติ อญฺญตโร ภิกฺขุ อทินฺนํ มุขทฺวารํ อาหารํ อาหริ, ตสฺมิํ วตฺถุสฺมิํ. เอกา ปญฺญตฺติ, เอกา อนุปญฺญตฺติ. ฉนฺนํ อาปตฺติ\-สมุฏฺฐานานํ ทฺวีหิ สมุฏฺฐาเนหิ สมุฏฺฐาติ, เอฬกโลมเก ฯเปฯ.}{โภชนวคฺโค จตุตฺโถ.}
\csromanmatikaanchor{mtk.15}
\csromantocmarkref{section}{5. อเจลกวคฺค}{mtk.15}
\bookmark[dest={mtk.15},level=1]{5. อเจลกวคฺค}
\csromanheader{1}{5. \textbf{อเจลกวคฺค}}
\proseitem{94}{อเจลกสฺส วา ปริพฺพาชกสฺส วา ปริพฺพาชิกาย วา สหตฺถา ขาทนียํ วา โภชนียํ วา เทนฺตสฺส ปาจิตฺติยํ กตฺถ ปญฺญตฺตนฺติ เวสาลิยํ ปญฺญตฺตํ. กํ อารพฺภาติ อายสฺมนฺตํ อานนฺทํ อารพฺภ. กิสฺมิํ วตฺถุสฺมินฺติ อายสฺมา อานนฺโท อญฺญตริสฺสา ปริพฺพาชิกาย เอกํ มญฺญมาโน เทฺว ปูเว อทาสิ, ตสฺมิํ วตฺถุสฺมิํ. เอกา ปญฺญตฺติ. ฉนฺนํ อาปตฺติ\-สมุฏฺฐานานํ ทฺวีหิ สมุฏฺฐาเนหิ สมุฏฺฐาติ, เอฬกโลมเก ฯเปฯ.}

\proseitem{95}{ภิกฺขุํ “เอหาวุโส, คามํ วา นิคมํ วา ปิณฺฑาย ปวิสิสฺสามา”ติ ตสฺส ทาเปตฺวา วา อทาเปตฺวา วา อุยฺโยเชนฺตสฺส ปาจิตฺติยํ กตฺถ ปญฺญตฺตนฺติ สาวตฺถิยํ ปญฺญตฺตํ. กํ อารพฺภาติ อายสฺมนฺตํ อุปนนฺทํ สกฺยปุตฺตํ อารพฺภ. กิสฺมิํ วตฺถุสฺมินฺติ อายสฺมา อุปนนฺโท สกฺยปุตฺโต ภิกฺขุํ “เอหาวุโส, คามํ ปิณฺฑาย ปวิสิสฺสามา”ติ ตสฺส อทาเปตฺวา อุยฺโยเชสิ, ตสฺมิํ วตฺถุสฺมิํ. เอกา ปญฺญตฺติ. ฉนฺนํ อาปตฺติ\-สมุฏฺฐานานํ ตีหิ สมุฏฺฐาเนหิ สมุฏฺฐาติ ฯเปฯ.}

\proseitem{96}{\csromanfolio{31}สโภชเน กุเล อนุปขชฺช นิสชฺชํ กปฺเปนฺตสฺส ปาจิตฺติยํ กตฺถ ปญฺญตฺตนฺติ สาวตฺถิยํ ปญฺญตฺตํ. กํ อารพฺภาติ อายสฺมนฺตํ อุปนนฺทํ สกฺยปุตฺตํ อารพฺภ. กิสฺมิํ วตฺถุสฺมินฺติ อายสฺมา อุปนนฺโท สกฺยปุตฺโต สโภชเน กุเล อนุปขชฺช นิสชฺชํ กปฺเปสิ, ตสฺมิํ วตฺถุสฺมิํ. เอกา ปญฺญตฺติ. ฉนฺนํ อาปตฺติ\-สมุฏฺฐานานํ เอเกน สมุฏฺฐาเนน สมุฏฺฐาติ, กายโต จ จิตฺตโต จ สมุฏฺฐาติ, น วาจโต ฯเปฯ.}

\proseitem{97}{มาตุคาเมน สทฺธิํ รโห ปฏิจฺฉนฺเน อาสเน นิสชฺชํ กปฺเปนฺตสฺส ปาจิตฺติยํ กตฺถ ปญฺญตฺตนฺติ สาวตฺถิยํ ปญฺญตฺตํ. กํ อารพฺภาติ อายสฺมนฺตํ อุปนนฺทํ สกฺยปุตฺตํ อารพฺภ. กิสฺมิํ วตฺถุสฺมินฺติ อายสฺมา อุปนนฺโท สกฺยปุตฺโต มาตุคาเมน สทฺธิํ รโห ปฏิจฺฉนฺเน อาสเน นิสชฺชํ กปฺเปสิ, ตสฺมิํ วตฺถุสฺมิํ. เอกา ปญฺญตฺติ. ฉนฺนํ อาปตฺติ\-สมุฏฺฐานานํ เอเกน สมุฏฺฐาเนน สมุฏฺฐาติ, กายโต จ จิตฺตโต จ สมุฏฺฐาติ, น วาจโต ฯเปฯ.}

\proseitem{98}{มาตุคาเมน สทฺธิํ เอโก เอกาย รโห นิสชฺชํ กปฺเปนฺตสฺส ปาจิตฺติยํ กตฺถ ปญฺญตฺตนฺติ สาวตฺถิยํ ปญฺญตฺตํ. กํ อารพฺภาติ อายสฺมนฺตํ อุปนนฺทํ สกฺยปุตฺตํ อารพฺภ. กิสฺมิํ วตฺถุสฺมินฺติ อายสฺมา อุปนนฺโท สกฺยปุตฺโต มาตุคาเมน สทฺธิํ เอโก เอกาย รโห นิสชฺชํ กปฺเปสิ, ตสฺมิํ วตฺถุสฺมิํ. เอกา ปญฺญตฺติ. ฉนฺนํ อาปตฺติ\-สมุฏฺฐานานํ เอเกน สมุฏฺฐาเนน สมุฏฺฐาติ, กายโต จ จิตฺตโต จ สมุฏฺฐาติ, น วาจโต ฯเปฯ.}

\proseitem{99}{นิมนฺติเตน สภตฺเตน สนฺตํ ภิกฺขุํ อนาปุจฺฉา ปุเรภตฺตํ ปจฺฉาภตฺตํ กุเลสุ จาริตฺตํ อาปชฺชนฺตสฺส ปาจิตฺติยํ กตฺถ ปญฺญตฺตนฺติ ราชคเห ปญฺญตฺตํ. กํ อารพฺภาติ อายสฺมนฺตํ อุปนนฺทํ สกฺยปุตฺตํ อารพฺภ. กิสฺมิํ วตฺถุสฺมินฺติ อายสฺมา อุปนนฺโท สกฺยปุตฺโต นิมนฺติโต สภตฺโต สมาโน ปุเรภตฺตํ ปจฺฉาภตฺตํ กุเลสุ จาริตฺตํ อาปชฺชิ, ตสฺมิํ วตฺถุสฺมิํ. เอกา ปญฺญตฺติ, จตสฺโส อนุปญฺญตฺติโย. ฉนฺนํ อาปตฺติ\-สมุฏฺฐานานํ ทฺวีหิ สมุฏฺฐาเนหิ สมุฏฺฐาติ, กถินเก ฯเปฯ.}

\proseitem{100}{ตตุตฺตริ เภสชฺชํ วิญฺญาเปนฺตสฺส ปาจิตฺติยํ กตฺถ ปญฺญตฺตนฺติ สกฺเกสุ ปญฺญตฺตํ. กํ อารพฺภาติ ฉพฺพคฺคิเย ภิกฺขู อารพฺภ. กิสฺมิํ วตฺถุสฺมินฺติ ฉพฺพคฺคิยา ภิกฺขู มหานาเมน สกฺเกน “อชฺชโณฺห ภนฺเต \csromanfolio{32}อาคเมถา”ติ วุจฺจมานา นาคเมสุํ, ตสฺมิํ วตฺถุสฺมิํ. เอกา ปญฺญตฺติ. ฉนฺนํ อาปตฺติ\-สมุฏฺฐานานํ ฉหิ สมุฏฺฐาเนหิ สมุฏฺฐาติ ฯเปฯ.}

\proseitem{101}{อุยฺยุตฺตํ เสนํ ทสฺสนาย คจฺฉนฺตสฺส ปาจิตฺติยํ กตฺถ ปญฺญตฺตนฺติ สาวตฺติยํ ปญฺญตฺตํ. กํ อารพฺภาติ ฉพฺพคฺคิเย ภิกฺขู อารพฺภ. กิสฺมิํ วตฺถุสฺมินฺติ ฉพฺพคฺคิยา ภิกฺขู อุยฺยุตฺตํ เสนํ ทสฺสนาย อคมํสุ, ตสฺมิํ วตฺถุสฺมิํ. เอกา ปญฺญตฺติ, เอกา อนุปญฺญตฺติ. ฉนฺนํ อาปตฺติ\-สมุฏฺฐานานํ ทฺวีหิ สมุฏฺฐาเนหิ สมุฏฺฐาติ, เอฬกโลมเก ฯเปฯ.}

\proseitem{102}{อติเรก\-ติรตฺตํ เสนาย วสนฺตสฺส ปาจิตฺติยํ กตฺถ ปญฺญตฺตนฺติ สาวตฺถิยํ ปญฺญตฺตํ. กํ อารพฺภาติ ฉพฺพคฺคิเย ภิกฺขู อารพฺภ. กิสฺมิํ วตฺถุสฺมินฺติ ฉพฺพคฺคิยา ภิกฺขู อติเรก\-ติรตฺตํ เสนาย วสิํสุ, ตสฺมิํ วตฺถุสฺมิํ. เอกา ปญฺญตฺติ. ฉนฺนํ อาปตฺติ\-สมุฏฺฐานานํ ทฺวีหิ สมุฏฺฐาเนหิ สมุฏฺฐาติ, เอฬกโลมเก ฯเปฯ.}

\proseitemwithclosermidruled{103}{อุยฺโยธิกํ คจฺฉนฺตสฺส ปาจิตฺติยํ กตฺถ ปญฺญตฺตนฺติ สาวตฺถิยํ ปญฺญตฺตํ. กํ อารพฺภาติ ฉพฺพคิเย ภิกฺขู อารพฺภ. กิสฺมิํ วตฺถุสฺมินฺติ ฉพฺพคฺคิยา ภิกฺขู อุยฺโยธิกํ อคมํสุ, ตสฺมิํ วตฺถุสฺมิํ. เอกา ปญฺญตฺติ. ฉนฺนํ อาปตฺติ\-สมุฏฺฐานานํ ทฺวีหิ สมุฏฺฐาเนหิ สมุฏฺฐาติ, เอฬกโลมเก ฯเปฯ.}{อเจลกวคฺโค ปญฺจโม.}
\csromanmatikaanchor{mtk.16}
\csromantocmarkref{section}{6. สุราปานวคฺค}{mtk.16}
\bookmark[dest={mtk.16},level=1]{6. สุราปานวคฺค}
\csromanheader{1}{6. \textbf{สุราปานวคฺค}}
\proseitem{104}{สุราเมรยปาเน ปาจิตฺติยํ กตฺถ ปญฺญตฺตนฺติ โกสมฺพิยํ ปญฺญตฺตํ. กํ อารพฺภาติ อายสฺมนฺตํ สาคตํ อารพฺภ. กิสฺมิํ วตฺถุสฺมินฺติ อายสฺมา สาคโต มชฺชํ ปิวิ, ตสฺมิํ วตฺถุสฺมิํ. เอกา ปญฺญตฺติ. ฉนฺนํ อาปตฺติ\-สมุฏฺฐานานํ ทฺวีหิ สมุฏฺฐาเนหิ สมุฏฺฐาติ, สิยา กายโต สมุฏฺฐาติ, น วาจโต, น จิตฺตโต, สิยา กายโต จ จิตฺตโต จ สมุฏฺฐาติ, น วาจโต ฯเปฯ.}

\proseitem{105}{องฺคุลิปโตทเก ปาจิตฺติยํ กตฺถ ปญฺญตฺตนฺติ สาวตฺถิยํ ปญฺญตฺตํ. กํ อารพฺภาติ ฉพฺพคฺคิเย ภิกฺขู อารพฺภ. กิสฺมิํ วตฺถุสฺมินฺติ \csromanfolio{33}ฉพฺพคฺคิยา ภิกฺขู ภิกฺขุํ องฺคุลิ\-ปโตทเกน หาเสสุํ, ตสฺมิํ วตฺถุสฺมิํ. เอกา ปญฺญตฺติ. ฉนฺนํ อาปตฺติ\-สมุฏฺฐานานํ เอเกน สมุฏฺฐาเนน สมุฏฺฐาติ, กายโต จ จิตฺตโต จ สมุฏฺฐาติ, น วาจโต ฯเปฯ.}

\proseitem{106}{อุทเก หสธมฺเม\csromansharedfootnote{823fa0c09a95b1da}{หสฺสธมฺเม (สี, สฺยา)} ปาจิตฺติยํ กตฺถ ปญฺญตฺตนฺติ สาวตฺถิยํ ปญฺญตฺตํ. กํ อารพฺภาติ สตฺตรส\-วคฺคิเย ภิกฺขู อารพฺภ. กิสฺมิํ วตฺถุสฺมินฺติ สตฺตรส\-วคฺคิยา ภิกฺขู อจิรวติยา นทิยา อุทเก กีฬิํสุ, ตสฺมิํ วตฺถุสฺมิํ. เอกา ปญฺญตฺติ. ฉนฺนํ อาปตฺติ\-สมุฏฺฐานานํ เอเกน สมุฏฺฐาเนน สมุฏฺฐาติ, กายโต จ จิตฺตโต จ สมุฏฺฐาติ, น วาจโต ฯเปฯ.}

\proseitem{107}{อนาทริเย ปาจิตฺติยํ กตฺถ ปญฺญตฺตนฺติ โกสมฺพิยํ ปญฺญตฺตํ. กํ อารพฺภาติ อายสฺมนฺตํ ฉนฺนํ อารพฺภ. กิสฺมิํ วตฺถุสฺมินฺติ อายสฺมา ฉนฺโน อนาทริยํ อกาสิ, ตสฺมิํ วตฺถุสฺมิํ. เอกา ปญฺญตฺติ. ฉนฺนํ อาปตฺติ\-สมุฏฺฐานานํ ตีหิ สมุฏฺฐาเนหิ สมุฏฺฐาติ ฯเปฯ.}

\proseitem{108}{ภิกฺขุํ ภิํสาเปนฺตสฺส ปาจิตฺติยํ กตฺถ ปญฺญตฺตนฺติ สาวตฺถิยํ ปญฺญตฺตํ. กํ อารพฺภาติ ฉพฺพคฺคิเย ภิกฺขู อารพฺภ. กิสฺมิํ วตฺถุสฺมินฺติ ฉพฺพคฺคิยา ภิกฺขู ภิกฺขุํ ภิสาเปสุํ, ตสฺมิํ วตฺถุสฺมิํ. เอกา ปญฺญตฺติ. ฉนฺนํ อาปตฺติ\-สมุฏฺฐานานํ ตีหิ สมุฏฺฐาเนหิ}

\proseitem{109}{โชติํ สมาทหิตฺวา วิสิพฺเพนฺตสฺส ปาจิตฺติยํ กตฺถ ปญฺญตฺถนฺติ ภคฺเคสุ ปญฺญตฺตํ. กํ อารพฺภาติ สมฺพหุเล ภิกฺขู อารพฺภ. กิสฺมิํ วตฺถุสฺมินฺติ สมฺพหุลา ภิกฺขู โชติํ สมาทหิตฺวา วิสิพฺเพสุํ, ตสฺมิํ วตฺถุสฺมิํ. เอกา ปญฺญตฺติ, เทฺว อนุปญฺญตฺติโย. ฉนฺนํ อาปตฺติ\-สมุฏฺฐานานํ ฉหิ สมุฏฺฐาเนหิ สมุฏฺฐาติ ฯเปฯ.}

\proseitem{110}{โอเรนทฺธมาสํ นหายนฺตสฺส ปาจิตฺติยํ กตฺถ ปญฺญตฺตนฺติ ราชคเห ปญฺญตฺตํ. กํ อารพฺภาติ สมฺพหุเล ภิกฺขู อารพฺภ. กิสฺมิํ วตฺถุสฺมินฺติ สมฺพหุลา ภิกฺขู ราชานมฺปิ ปสฺสิตฺวา น มตฺตํ ชานิตฺวา นหายิํสุ, ตสฺมิํ วตฺถุสฺมิํ. เอกา ปญฺญตฺติ, ฉ อนุปญฺญตฺติโย. สพฺพตฺถ\-ปญฺญตฺติ ปเทสปญฺญตฺตีติ ปเทสปญฺญตฺติ. ฉนฺนํ อาปตฺติ\-สมุฏฺฐานานํ ทฺวีหิ สมุฏฺฐาเนหิ สมุฏฺฐาติ, เอฬกโลมเก ฯเปฯ.}

\proseitem{111}{\csromanfolio{34}อนาทิยิตฺวา ติณฺณํ ทุพฺพณฺณ\-กรณานํ อญฺญตรํ ทุพฺพณฺณ\-กรณํ นวํ จีวรํ ปริภุญฺช\-นฺตสฺส ปาจิตฺติยํ กตฺถ ปญฺญตฺตนฺติ สาวตฺถิยํ ปญฺญตฺตํ. กํ อารพฺภาติ สมฺพหุเล ภิกฺขู อารพฺภ. กิสฺมิํ วตฺถุสฺมินฺติ สมฺพหุลา ภิกฺขู อตฺตโน จีวรํ น สญฺชานิํสุ, ตสฺมิํ วตฺถุสฺมิํ. เอกา ปญฺญตฺติ. ฉนฺนํ อาปตฺติ\-สมุฏฺฐานานํ ทฺวีหิ สมุฏฺฐาเนหิ สมุฏฺฐาติ, เอฬกโลมเก ฯเปฯ.}

\proseitem{112}{ภิกฺขุสฺส วา ภิกฺขุนิยา วา สิกฺขมานาย วา สามเณรสฺส วา สามเณริยา วา สามํ จีวรํ วิกปฺเปตฺวา อปฺปจฺจุทฺธารณํ ปริภุญฺช\-นฺตสฺส ปาจิตฺติยํ กตฺถ ปญฺญตฺตนฺติ สาวตฺถิยํ ปญฺญตฺตํ. กํ อารพฺภาติ อายสฺมนฺตํ อุปนนฺทํ สกฺยปุตฺตํ อารพฺภ. กิสฺมิํ วตฺถุสฺมินฺติ อายสฺมา อุปนนฺโท สกฺยปุตฺโต ภิกฺขุสฺส สามํ จีวรํ วิกปฺเปตฺวา อปฺปจฺจุทฺธารณํ ปริภุญฺชิ, ตสฺมิํ วตฺถุสฺมิํ. เอกา ปญฺญตฺติ. ฉนฺนํ อาปตฺติ\-สมุฏฺฐานานํ ทฺวีหิ สมุฏฺฐาเนหิ สมุฏฺฐาติ, กถินเก ฯเปฯ.}

\proseitemwithclosermidruled{113}{ภิกฺขุสฺส ปตฺตํ วา จีวรํ วา นิสีทนํ วา สูจิฆรํ วา กายพนฺธนํ วา อปนิเธนฺตสฺส ปาจิตฺติยํ กตฺถ ปญฺญตฺตนฺติ สาวตฺถิยํ ปญฺญตฺตํ. กํ อารพฺภาติ ฉพฺพคฺคิเย ภิกฺขู อารพฺภ. กิสฺมิํ วตฺถุสฺมินฺติ ฉพฺพคฺคิยา ภิกฺขู ภิกฺขูนํ ปตฺตมฺปิ จีวรมฺปิ อปนิเธสุํ, ตสฺมิํ วตฺถุสฺมิํ. เอกา ปญฺญตฺติ. ฉนฺนํ อาปตฺติ\-สมุฏฺฐานานํ ตีหิ สมุฏฺฐาเนหิ สมุฏฺฐาติ ฯเปฯ.}{สุราเมรย\-วคฺโค ฉฏฺโฐ.}
\csromanmatikaanchor{mtk.17}
\csromantocmarkref{section}{7. สปฺปาณกวคฺค}{mtk.17}
\bookmark[dest={mtk.17},level=1]{7. สปฺปาณกวคฺค}
\csromanheader{1}{7. \textbf{สปฺปาณกวคฺค}}
\proseitem{114}{สญฺจิจฺจ ปาณํ ชีวิตา โวโรเปนฺตสฺส ปาจิตฺติยํ กตฺถ ปญฺญตฺตนฺติ สาวตฺถิยํ ปญฺญตฺตํ. กํ อารพฺภาติ อายสฺมนฺตํ อุทายิํ อารพฺภ. กิสฺมิํ วตฺถุสฺมินฺติ อายสฺมา อุทายี สญฺจิจฺจ ปาณํ ชีวิตา โวโรเปสิ, ตสฺมิํ วตฺถุสฺมิํ. เอกา ปญฺญตฺติ. ฉนฺนํ อาปตฺติ\-สมุฏฺฐานานํ ตีหิ สมุฏฺฐาเนหิ}

\proseitem{115}{ชานํ สปฺปาณกํ อุทกํ ปริภุญฺช\-นฺตสฺส ปาจิตฺติยํ กตฺถ ปญฺญตฺตนฺติ สาวตฺถิยํ ปญฺญตฺตํ. กํ อารพฺภาติ ฉพฺพคฺคิเย ภิกฺขู อารพฺภ. กิสฺมิํ \csromanfolio{35}วตฺถุสฺมินฺติ ฉพฺพคฺคิยา ภิกฺขู ชานํ สปฺปาณกํ อุทกํ ปริภุญฺจิํสุ, ตสฺมิํ วตฺถุสฺมิํ. เอกา ปญฺญตฺติ. ฉนฺนํ อาปตฺติ\-สมุฏฺฐานานํ ตีหิ สมุฏฺฐาเนหิ สมุฏฺฐาติ ฯเปฯ.}

\proseitem{116}{ชานํ ยถาธมฺมํ นิหตา\-ธิกรณํ ปุน กมฺมาย อุกฺโกเฏนฺตสฺส ปาจิตฺติยํ กตฺถ ปญฺญตฺตนฺติ สาวตฺถิยํ ปญฺญตฺตํ. กํ อารพฺภาติ ฉพฺพคฺคิเย ภิกฺขู อารพฺภ. กิสฺมิํ วตฺถุสฺมินฺติ ฉพฺพคฺคิยา ภิกฺขู ชานํ ยถาธมฺมํ นิหตา\-ธิกรณํ ปุน กมฺมาย อุกฺโกเฏสุํ, ตสฺมิํ วตฺถุสฺมิํ. เอกา ปญฺญตฺติ. ฉนฺนํ อาปตฺติ\-สมุฏฺฐานานํ ตีหิ สมุฏฺฐาเนหิ สมุฏฺฐาติ ฯเปฯ.}

\proseitem{117}{ภิกฺขุสฺส ชานํ ทุฏฺฐุลฺลํ อาปตฺติํ ปฏิจฺฉาเทนฺตสฺส ปาจิตฺติยํ กตฺถ ปญฺญตฺตนฺติ สาวตฺถิยํ ปญฺญตฺตํ. กํ อารพฺภาติ อญฺญตรํ ภิกฺขุํ อารพฺภ, กิสฺมิํ วตฺถุสฺมินฺติ อญฺญตโร ภิกฺขุ ภิกฺขุสฺส ชานํ ทุฏฺฐุลฺลํ อาปตฺติํ ปฏิจฺฉาเทสิ, ตสฺมิํ วตฺถุสฺมิํ. เอกา ปญฺญตฺติ. ฉนฺนํ อาปตฺติ\-สมุฏฺฐานานํ เอเกน สมุฏฺฐาเนน สมุฏฺฐาติ, กายโต จ วาจโต จ จิตฺตโต จ สมุฏฺฐาติ ฯเปฯ.}

\proseitem{118}{ชานํ อูนวีสติวสฺสํ ปุคฺคลํ อุปสมฺปาเทนฺตสฺส ปาจิตฺติยํ กตฺถ ปญฺญตฺตนฺติ ราชคเห ปญฺญตฺตํ. กํ อารพฺภาติ สมฺพหุเล ภิกฺขู อารพฺภ. กิสฺมิํ วตฺถุสฺมินฺติ สมฺพหุลา ภิกฺขู ชานํ อูนวีสติวสฺสํ ปุคฺคลํ อุปสมฺปาเทสุํ, ตสฺมิํ วตฺถุสฺมิํ. เอกา ปญฺญตฺติ. ฉนฺนํ อาปตฺติ\-สมุฏฺฐานานํ ตีหิ สมุฏฺฐาเนหิ สมุฏฺฐาติ ฯเปฯ.}

\proseitem{119}{ชานํ เถยฺยสตฺเถน สทฺธิํ สํวิธาย เอก\-ทฺธาน\-มคฺคํ ปฏิปชฺชนฺตสฺส ปาจิตฺติยํ กตฺถ ปญฺญตฺตนฺติ สาวตฺถิยํ ปญฺญตฺตํ. กํ อารพฺภาติ อญฺญตรํ ภิกฺขุํ อารพฺภ. กิสฺมิํ วตฺถุสฺมินฺติ อญฺญตโร ภิกฺขุ ชานํ เถยฺยสตฺเถน สทฺธิํ สํวิธาย เอก\-ทฺธาน\-มคฺคํ ปฏิปชฺชิ, ตสฺมิํ วตฺถุสฺมิํ. เอกา ปญฺญตฺติ. ฉนฺนํ อาปตฺติ\-สมุฏฺฐานานํ ทฺวีหิ สมุฏฺฐาเนหิ สมุฏฺฐาติ, สิยา กายโต จ จิตฺตโต จ สมุฏฺฐาติ, น วาจโต, สิยา กายโต จ วาจโต จ จิตฺตโต จ สมุฏฺฐาติ ฯเปฯ.}

\proseitem{120}{มาตุคาเมน สทฺธิํ สํวิธาย เอก\-ทฺธาน\-มคฺคํ ปฏิปชฺชนฺตสฺส ปาจิตฺติยํ กตฺถ ปญฺญตฺตนฺติ สาวตฺถิยํ ปญฺญตฺตํ. กํ อารพฺภาติ อญฺญตรํ ภิกฺขุํ อารพฺภ. \csromanfolio{36}กิสฺมิํ วตฺถุสฺมินฺติ อญฺญตโร ภิกฺขุ มาตุคาเมน สทฺธิํ สํวิธาย เอก\-ทฺธาน\-มคฺคํ ปฏิปชฺชิ, ตสฺมิํ วตฺถุสฺมิํ. เอกา ปญฺญตฺติ. ฉนฺนํ อาปตฺติ\-สมุฏฺฐานานํ จตูหิ สมุฏฺฐาเนหิ สมุฏฺฐาติ ฯเปฯ.}

\proseitem{121}{ปาปิกาย ทิฏฺฐิยา ยาวตติยํ สมนุภาสนาย น ปฏินิสฺสชฺชนฺตสฺส ปาจิตฺติยํ กตฺถ ปญฺญตฺตนฺติ สาวตฺถิยํ ปญฺญตฺตํ. กํ อารพฺภาติ อริฏฺฐํ ภิกฺขุํ คทฺธพาธิ\-ปุพฺพํ อารพฺภ. กิสฺมิํ วตฺถุสฺมินฺติ อริฏฺโฐ ภิกฺขุ คทฺธพาธิ\-ปุพฺโพ ปาปิกาย ทิฏฺฐิยา ยาวตติยํ สมนุภาสนาย น ปฏินิสฺสชฺชิ, ตสฺมิํ วตฺถุสฺมิํ. เอกา ปญฺญตฺติ. ฉนฺนํ อาปตฺติ\-สมุฏฺฐานานํ เอเกน สมุฏฺฐาเนน สมุฏฺฐาติ, กายโต จ วาจโต จ จิตฺตโต จ สมุฏฺฐาติ ฯเปฯ.}

\proseitem{122}{ชานํ ตถาวาทินา ภิกฺขุนา\csromansharedfootnote{d007ee3ea711319e}{อริฏฺเฐน ภิกฺขุนา (ก)} อกฏา\-นุธมฺเมน ตํ ทิฏฺฐิํ อปฺปฏินิสฺสฏฺเฐน สทฺธิํ สมฺภุญฺช\-นฺตสฺส ปาจิตฺติยํ กตฺถ ปญฺญตฺตนฺติ สาวตฺถิยํ ปญฺญตฺตํ. กํ อารพฺภาติ ฉพฺพคฺคิเย ภิกฺขู อารพฺภ. กิสฺมิํ วตฺถุสฺมินฺติ ฉพฺภคฺคิยา ภิกฺขู ชานํ ตถาวาทินา อริฏฺเฐน ภิกฺขุนา อกฏา\-นุธมฺเมน ตํ ทิฏฺฐิํ อปฺปฏินิสฺสฏฺเฐน สทฺธิํ สมฺภุญฺชิํสุ, ตสฺมิํ วตฺถุสฺมิํ. เอกา ปญฺญตฺติ. ฉนฺนํ อาปตฺติ\-สมุฏฺฐานานํ ตีหิ สมุฏฺฐาเนหิ สมุฏฺฐาติ ฯเปฯ.}

\proseitemwithclosermidruled{123}{ชานํ ตถานาสิตํ สมณุทฺเทสํ อุปลาเปนฺตสฺส ปาจิตฺติยํ กตฺถ ปญฺญตฺตนฺติ สาวตฺถิยํ ปญฺญตฺตํ. กํ อารพฺภาติ ฉพฺพคฺคิเย ภิกฺขู อารพฺภ. กิสฺมิํ วตฺถุสฺมินฺติ ฉพฺพคฺคิยา ภิกฺขู ชานํ ตถานาสิตํ กณฺฏกํ สมณุทฺเทสํ อุปลาเปสุํ, ตสฺมิํ วตฺถุสฺมิํ. เอกา ปญฺญตฺติ. ฉนฺนํ อาปตฺติ\-สมุฏฺฐานานํ ตีหิ สมุฏฺฐาเนหิ สมุฏฺฐาติ ฯเปฯ.}{สปฺปาณกวคฺโค สตฺตโม.}
\csromanmatikaanchor{mtk.18}
\csromantocmarkref{section}{8. สหธมฺมิกวคฺค}{mtk.18}
\bookmark[dest={mtk.18},level=1]{8. สหธมฺมิกวคฺค}
\csromanheader{1}{8. สหธมฺมิก\-วคฺค}
\proseitem{124}{ภิกฺขูหิ สหธมฺมิกํ วุจฺจมาเนน “น ตาวาหํ อาวุโส เอตสฺมิํ สิกฺขาปเท สิกฺขิสฺสามิ, ยาว น อญฺญํ ภิกฺขุํ พฺยตฺตํ วินยธรํ \csromanfolio{37}ปริปุจฺฉิสฺสามี”ติ ภณนฺตสฺส ปาจิตฺติยํ กตฺถ ปญฺญตฺตนฺติ โกสมฺพิยํ ปญฺญตฺตํ. กํ อารพฺภาติ อายสฺมนฺตํ ฉนฺนํ อารพฺภ. กิสฺมิํ วตฺถุสฺมินฺติ อายสฺมา ฉนฺโน ภิกฺขูหิ สหธมฺมิกํ วุจฺจมาโน “น ตาวาหํ อาวุโส เอตสฺมิํ สิกฺขาปเท สิกฺขิสฺสามิ, ยาว น อญฺญํ ภิกฺขุํ พฺยตฺตํ วินยธรํ ปริปุจฺฉิสฺสามี”ติ ภณิ, ตสฺมิํ วตฺถุสฺมิํ. เอกา ปญฺญตฺติ. ฉนฺนํ อาปตฺติ\-สมุฏฺฐานานํ ตีหิ สมุฏฺฐาเนหิ สมุฏฺฐาติ ฯเปฯ.}

\proseitem{125}{วินยํ วิวณฺเณนฺตสฺส ปาจิตฺติยํ กตฺถ ปญฺญตฺตนฺติ สาวตฺถิยํ ปญฺญตฺตํ. กํ อารพฺภาติ ฉพฺพคฺคิเย ภิกฺขู อารพฺภ. กิสฺมิํ วตฺถุสฺมินฺติ ฉพฺพคฺคิยา ภิกฺขู วินยํ วิวณฺเณสุํ, ตสฺมิํ วตฺถุสฺมิํ. เอกา ปญฺญตฺติ. ฉนฺนํ อาปตฺติ\-สมุฏฺฐานานํ ตีหิ สมุฏฺฐาเนหิ}

\proseitem{126}{โมหนเก ปาจิตฺติยํ กตฺถ ปญฺญตฺตนฺติ สาวตฺถิยํ ปญฺญตฺตํ. กํ อารพฺภาติ ฉพฺพคฺคิเย ภิกฺขู อารพฺภ. กิสฺมิํ วตฺถุสฺมินฺติ ฉพฺพคฺคิยา ภิกฺขู โมเหสุํ, ตสฺมิํ วตฺถุสฺมิํ. เอกา ปญฺญตฺติ. ฉนฺนํ อาปตฺติ\-สมุฏฺฐานานํ ตีหิ สมุฏฺฐาเนหิ สมุฏฺฐาติ ฯเปฯ.}

\proseitem{127}{ภิกฺขุสฺส กุปิเตน อนตฺตมเนน ปหารํ เทนฺตสฺส ปาจิตฺติยํ กตฺถ ปญฺญตฺตนฺติ สาวตฺถิยํ ปญฺญตฺตํ. กํ อารพฺภาติ ฉพฺพคฺคิเย ภิกฺขู อารพฺภ. กิสฺมิํ วตฺถุสฺมินฺติ ฉพฺพคฺคิยา ภิกฺขู กุปิตา อนตฺตมนา ภิกฺขูนํ ปหารํ อทํสุ, ตสฺมิํ วตฺถุสฺมิํ. เอกา ปญฺญตฺติ. ฉนฺนํ อาปตฺติ\-สมุฏฺฐานานํ เอเกน สมุฏฺฐาเนน สมุฏฺฐาติ, กายโต จ จิตฺตโต จ สมุฏฺฐาติ, น วาจโต ฯเปฯ.}

\proseitem{128}{ภิกฺขุสฺส กุปิเตน อนตฺตมเนน ตลสตฺติกํ อุคฺคิรนฺตสฺส ปาจิตฺติยํ กตฺถ ปญฺญตฺตนฺติ สาวตฺถิยํ ปญฺญตฺตํ. กํ อารพฺภาติ ฉพฺพคฺคิเย ภิกฺขู อารพฺภ. กิสฺมิํ วตฺถุสฺมินฺติ ฉพฺพคฺคิยา ภิกฺขู กุปิตา อนตฺตมนา ภิกฺขูนํ ตลสตฺติกํ อุคฺคิริํสุ, ตสฺมิํ วตฺถุสฺมิํ. เอกา ปญฺญตฺติ. ฉนฺนํ อาปตฺติ\-สมุฏฺฐานานํ เอเกน สมุฏฺฐาเนน สมุฏฺฐาติ, กายโต จ จิตฺตโต จ สมุฏฺฐาติ, น วาจโต ฯเปฯ.}

\proseitem{129}{ภิกฺขุํ อมูลเกน สํฆาทิเสเสน อนุทฺธํเส\-นฺตสฺส ปาจิตฺติยํ กตฺถ ปญฺญตฺตนฺติ สาวตฺถิยํ ปญฺญตฺตํ. กํ อารพฺภาติ ฉพฺพคฺคิเย ภิกฺขู \csromanfolio{38}อารพฺภ. กิสฺมิํ วตฺถุสฺมินฺติ ฉพฺพคฺคิยา ภิกฺขู ภิกฺขุํ อมูลเกน สํฆาทิเสเสน อนุทฺธํเสสุํ, ตสฺมิํ วตฺถุสฺมิํ. เอกา ปญฺญตฺติ. ฉนฺนํ อาปตฺติ\-สมุฏฺฐานานํ ตีหิ สมุฏฺฐาเนหิ สมุฏฺฐาติ ฯเปฯ.}

\proseitem{130}{ภิกฺขุสฺส สญฺจิจฺจ กุกฺกุจฺจํ อุปทหนฺตสฺส1 ปาจิตฺติยํ กตฺถ ปญฺญตฺตนฺติ สาวตฺถิยํ ปญฺญตฺตํ. กํ อารพฺภาติ ฉพฺพคฺคิเย ภิกฺขู อารพฺภ. กิสฺมิํ วตฺถุสฺมินฺติ ฉพฺพคฺคิยา ภิกฺขู ภิกฺขูนํ สญฺจิจฺจ กุกฺกุจฺจํ อุปทหิํสุ, ตสฺมิํ วตฺถุสฺมิํ. เอกา ปญฺญตฺติ. ฉนฺนํ อาปตฺติ\-สมุฏฺฐานานํ ตีหิ สมุฏฺฐาเนหิ สมุฏฺฐาติ ฯเปฯ.}

\proseitem{131}{ภิกฺขูนํ ภณฺฑน\-ชาตานํ กลหชาตานํ วิวาทาปนฺนานํ อุปสฺสุติํ2 ติฏฺฐนฺตสฺส ปาจิตฺติยํ กตฺถ ปญฺญตฺตนฺติ สาวตฺถิยํ ปญฺญตฺตํ. กํ อารพฺภาติ ฉพฺพคฺคิเย ภิกฺขู อารพฺภ. กิสฺมิํ วตฺถุสฺมินฺติ ฉพฺพคฺคิยา ภิกฺขู ภิกฺขูนํ ภณฺฑน\-ชาตานํ กลหชาตานํ วิวาทาปนฺนานํ อุปสฺสุติํ ติฏฺฐหิํสุ, ตสฺมิํ วตฺถุสฺมิํ. เอกา ปญฺญตฺติ. ฉนฺนํ อาปตฺติ\-สมุฏฺฐานานํ ทฺวีหิ สมุฏฺฐาเนหิ สมุฏฺฐาติ, สิยา กายโต จ จิตฺตโต จ สมุฏฺฐาติ, น วาจโต, สิยา กายโต จ วาจโต จ}

\proseitem{132}{ธมฺมิกานํ กมฺมานํ ฉนฺทํ ทตฺวา ปจฺฉา ขียนธมฺมํ3 อาปชฺชนฺตสฺส ปาจิตฺติยํ กตฺถ ปญฺญตฺตนฺติ สาวตฺถิยํ ปญฺญตฺตํ. กํ อารพฺภาติ ฉพฺพคฺคิเย ภิกฺขู อารพฺภ. กิสฺมิํ วตฺถุสฺมินฺติ ฉพฺพคฺคิยา ภิกฺขู ธมฺมิกานํ กมฺมานํ ฉนฺทํ ทตฺวา ปจฺฉา ขียนธมฺมํ อาปชฺชิํสุ, ตสฺมิํ วตฺถุสฺมิํ. เอกา ปญฺญตฺติ. ฉนฺนํ อาปตฺติ\-สมุฏฺฐานานํ ตีหิ สมุฏฺฐาเนหิ สมุฏฺฐาติ ฯเปฯ.}

\proseitem{133}{สํเฆ วินิจฺฉย\-กถาย วตฺตมานาย ฉนฺทํ อทตฺวา อุฏฺฐายาสนา ปกฺกมนฺตสฺส ปาจิตฺติยํ กตฺถ ปญฺญตฺตนฺติ สาวตฺถิยํ ปญฺญตฺตํ. กํ อารพฺภาติ อญฺญตรํ ภิกฺขุํ อารพฺภ. กิสฺมิํ วตฺถุสฺมินฺติ อญฺญตโร ภิกฺขุ สํเฆ ปินิจฺฉยกถาย วตฺตมานาย ฉนฺทํ อทตฺวา อุฏฺฐายาสนา ปกฺกามิ, ตสฺมิํ วตฺถุสฺมิํ. เอกา ปญฺญตฺติ. ฉนฺนํ อาปตฺติ\-สมุฏฺฐานานํ เอเกน สมุฏฺฐาเนน สมุฏฺฐาติ, กายโต จ วาจโต จ จิตฺตโต จ สมุฏฺฐาติ ฯเปฯ.}

\proseitem{1}{อุปฺปาเทนฺตสฺส(สฺยา) 2. อุปสฺสุติ (?) 3. ขิยฺยธมฺมํ (ก)}

\proseitem{134}{\csromanfolio{39}สมคฺเคน สํเฆน จีวรํ ทตฺวา ปจฺฉา ขียนธมฺมํ อาปชฺชนฺตสฺส ปาจิตฺติยํ กตฺถ ปญฺญตฺตนฺติ ราชคเห ปญฺญตฺตํ. กํ อารพฺภาติ ฉพฺพคฺคิเย ภิกฺขู อารพฺภ. กิสฺมิํ วตฺถุสฺมินฺติ ฉพฺพคฺคิยา ภิกฺขู สมคฺเคน สํเฆน จีวรํ ทตฺวา ปจฺฉา ขียนธมฺมํ อาปชฺชิํสุ, ตสฺมิํ วตฺถุสฺมิํ. เอกา ปญฺญตฺติ. ฉนฺนํ อาปตฺติ\-สมุฏฺฐานานํ ตีหิ สมุฏฺฐาเนหิ สมุฏฺฐาติ ฯเปฯ.}

\proseitemwithclosermidruled{135}{ชานํ สํฆิกํ ลาภํ ปริณตํ ปุคฺคลสฺส ปริณาเมนฺตสฺส ปาจิตฺติยํ กตฺถ ปญฺญตฺตนฺติ สาวตฺถิยํ ปญฺญตฺตํ. กํ อารพฺภาติ ฉพฺพคฺคิเย ภิกฺขู อารพฺภ. กิสฺมิํ วตฺถุสฺมินฺติ ฉพฺพคฺคิยา ภิกฺขู ชานํ สํฆิกํ ลาภํ ปริณตํ ปุคฺคลสฺส ปริณาเมสุํ, ตสฺมิํ วตฺถุสฺมิํ. เอกา ปญฺญตฺติ. ฉนฺนํ อาปตฺติ\-สมุฏฺฐานานํ ตีหิ สมุฏฺฐาเนหิ}{สหธมฺมิก\-วคฺโค อฏฺฐโม.}
\csromanmatikaanchor{mtk.19}
\csromantocmarkref{section}{9. ราชวคฺค}{mtk.19}
\bookmark[dest={mtk.19},level=1]{9. ราชวคฺค}
\csromanheader{1}{9. \textbf{ราชวคฺค}}
\proseitem{136}{ปุพฺเพ อปฺปฏิสํวิทิเตน รญฺโญ อนฺเตปุรํ ปวิสนฺตสฺส ปาจิตฺติยํ กตฺถ ปญฺญตฺตนฺติ สาวตฺถิยํ ปญฺญตฺตํ. กํ อารพฺภาติ อายสฺมนฺตํ อานนฺทํ อารพฺภ. กิสฺมิํ วตฺถุสฺมินฺติ อายสฺมา อานนฺโท ปุพฺเพ อปฺปฏิสํวิทิโต รญฺโญ อนฺเตปุรํ ปาวิสิ, ตสฺมิํ วตฺถุสฺมิํ. เอกา ปญฺญตฺติ. ฉนฺนํ อาปตฺติ\-สมุฏฺฐานานํ ทฺวีหิ สมุฏฺฐาเนหิ สมุฏฺฐาติ, กถินเก ฯเปฯ.}

\proseitem{137}{รตนํ อุคฺคณฺหนฺตสฺส ปาจิตฺติยํ กตฺถ ปญฺญตฺตนฺติ สาวตฺถิยํ ปญฺญตฺตํ. กํ อารพฺภาติ อญฺญตรํ ภิกฺขุํ อารพฺภ. กิสฺมิํ วตฺถุสฺมินฺติ อญฺญตโร ภิกฺขุ รตนํ อุคฺคเหสิ, ตสฺมิํ วตฺถุสฺมิํ. เอกา ปญฺญตฺติ, เทฺว อนุปญฺญตฺติโย. ฉนฺนํ อาปตฺติ\-สมุฏฺฐานานํ ฉหิ สมุฏฺฐาเนหิ สมุฏฺฐาติ ฯเปฯ.}

\proseitem{138}{สนฺตํ ภิกฺขุํ อนาปุจฺฉา วิกาเล คามํ ปวิสนฺตสฺส ปาจิตฺติยํ กตฺถ ปญฺญตฺตนฺติ สาวตฺถิยํ ปญฺญตฺตํ. กํ อารพฺภาติ ฉพฺพคฺคิเย ภิกฺขู อารพฺภ. กิสฺมิํ วตฺถุสฺมินฺติ ฉพฺพคฺคิยา ภิกฺขู วิกาเล คามํ ปวิสิํสุ, ตสฺมิํ วตฺถุสฺมิํ. เอกา ปญฺญตฺติ, ติสฺโส อนุปญฺญตฺติโย. ฉนฺนํ อาปตฺติ\-สมุฏฺฐานานํ ทฺวีหิ สมุฏฺฐาเนหิ สมุฏฺฐาติ, กถินเก ฯเปฯ.}

\proseitem{139}{\csromanfolio{40}อฏฺฐิมยํ วา ทนฺตมยํ วา วิสาณมยํ วา สูจิฆรํ การาเปนฺตสฺส ปาจิตฺติยํ กตฺถ ปญฺญตฺตนฺติ สกฺเกสุ ปญฺญตฺตํ. กํ อารพฺภาติ สมฺพหุเล ภิกฺขู อารพฺภ. กิสฺมิํ วตฺถุสฺมินฺติ สมฺพหุลา ภิกฺขู น มตฺตํ ชานิตฺวา พหู สูจิฆเร วิญฺญาเปสุํ, ตสฺมิํ วตฺถุสฺมิํ. เอกา ปญฺญตฺติ. ฉนฺนํ อาปตฺติ\-สมุฏฺฐานานํ ฉหิ สมุฏฺฐาเนหิ สมุฏฺฐาติ ฯเปฯ.}

\proseitem{140}{ปมาณา\-ติกฺกนฺตํ มญฺจํ วา ปีฐํ วา การาเปนฺตสฺส ปาจิตฺติยํ กตฺถ ปญฺญตฺตนฺติ สาวตฺถิยํ ปญฺญตฺตํ. กํ อารพฺภาติ อายสฺมนฺตํ อุปนนฺทํ สกฺยปุตฺตํ อารพฺภ. กิสฺมิํ วตฺถุสฺมินฺติ อายสฺมา อุปนนฺโท สกฺยปุตฺโต อุจฺเจ มญฺเจ สยิ, ตสฺมิํ วตฺถุสฺมิํ. เอกา ปญฺญตฺติ. ฉนฺนํ อาปตฺติ\-สมุฏฺฐานานํ ฉหิ สมุฏฺฐาเนหิ สมุฏฺฐาติ ฯเปฯ.}

\proseitem{141}{มญฺจํ วา ปีฐํ วา ตูโลนทฺธํ การาเปนฺตสฺส ปาจิตฺติยํ กตฺถ ปญฺญตฺตนฺติ สาวตฺถิยํ ปญฺญตฺตํ. กํ อารพฺภาติ ฉพฺพคฺคิเย ภิกฺขู อารพฺภ. กิสฺมิํ วตฺถุสฺมินฺติ ฉพฺพคฺคิยา ภิกฺขู มญฺจํ วา ปีฐํ วา ตูโลนทฺธํ การาเปสุํ, ตสฺมิํ วตฺถุสฺมิํ. เอกา ปญฺญตฺติ. ฉนฺนํ อาปตฺติ\-สมุฏฺฐานานํ ฉหิ สมุฏฺฐาเนหิ สมุฏฺฐาติ ฯเปฯ.}

\proseitem{142}{ปมาณา\-ติกฺกนฺตํ นิสีทนํ การาเปนฺตสฺส ปาจิตฺติยํ กตฺถ ปญฺญตฺตนฺติ สาวตฺถิยํ ปญฺญตฺตํ. กํ อารพฺภาติ ฉพฺพคฺคิเย ภิกฺขู อารพฺภ. กิสฺมิํ วตฺถุสฺมินฺติ ฉพฺพคฺคิยา ภิกฺขู อปฺปมาณิกานิ นิสีทนานิ ธาเรสุํ, ตสฺมิํ วตฺถุสฺมิํ. เอกา ปญฺญตฺติ, เอกา อนุปญฺญตฺติ. ฉนฺนํ อาปตฺติ\-สมุฏฺฐานานํ ฉหิ สมุฏฺฐาเนหิ สมุฏฺฐาติ ฯเปฯ.}

\proseitem{143}{ปมาณา\-ติกฺกนฺตํ กณฺฑุ\-ปฺปฏิจฺฉาทิํ การาเปนฺตสฺส ปาจิตฺติยํ กตฺถ ปญฺญตฺตนฺติ สาวตฺถิยํ ปญฺญตฺตํ. กํ อารพฺภาติ ฉพฺพคฺคิเย ภิกฺขู อารพฺภ. กิสฺมิํ วตฺถุสฺมินฺติ ฉพฺพคฺคิยา ภิกฺขู อปฺปมาณิกาโย กณฺฑุ\-ปฺปฏิจฺฉาทิโย ธาเรสุํ, ตสฺมิํ วตฺถุสฺมิํ. เอกา ปญฺญตฺติ. ฉนฺนํ อาปตฺติ\-สมุฏฺฐานานํ ฉหิ สมุฏฺฐาเนหิ สมุฏฺฐาติ ฯเปฯ.}

\proseitem{144}{ปมาณา\-ติกฺกนฺตํ วสฺสิกสาฏิกํ การาเปนฺตสฺส ปาจิตฺติยํ กตฺถ ปญฺญตฺตนฺติ สาวตฺถิยํ ปญฺญตฺตํ. กํ อารพฺภาติ ฉพฺพคฺคิเย ภิกฺขู \csromanfolio{41}อารพฺภ. กิสฺมิํ วตฺถุสฺมินฺติ ฉพฺพคฺคิยา ภิกฺขู อปฺปมาณิกาโย วสฺสิก\-สาฏิกาโย ธาเรสุํ, ตสฺมิํ วตฺถุสฺมิํ. เอกา ปญฺญตฺติ. ฉนฺนํ อาปตฺติ\-สมุฏฺฐานานํ ฉหิ สมุฏฺฐาเนหิ สมุฏฺฐาติ ฯเปฯ.}

\proseitemwithclosermidruled{145}{สุคต\-จีวร\-ปฺปมาณํ จีวรํ การาเปนฺตสฺส ปาจิตฺติยํ กตฺถ ปญฺญตฺตนฺติ สาวตฺถิยํ ปญฺญตฺตํ. กํ อารพฺภาติ อายสฺมนฺตํ นนฺทํ อารพฺภ. กิสฺมิํ วตฺถุสฺมินฺติ อายสฺมา นนฺโท สุคต\-จีวร\-ปฺปมาณํ จีวรํ ธาเรสิ, ตสฺมิํ วตฺถุสฺมิํ. เอกา ปญฺญตฺติ. ฉนฺนํ อาปตฺติ\-สมุฏฺฐานานํ ฉหิ สมุฏฺฐาเนหิ สมุฏฺฐาติ ฯเปฯ.}{ราชวคฺโค นวโม.}
\nitthitam{เทฺวนวุติ ปาจิตฺติยา นิฏฺฐิตา.}
\nitthitam{ขุทฺทกํ สมตฺตํ.}
\csromancenter{ตสฺสุทฺทานํ.}
\setlength{\csromangathapagewidth}{0pt}
\setlength{\csromangathawakwidth}{0pt}
\csromangathameasuregroup{มุสา โอมสเปสุญฺญํ, \\ อญฺญตฺร วิญฺญุนา ภูตา1, \\ ภูตํ อญฺญาย อุชฺฌายิ2, \\ ปุพฺเพ นิกฺกฑฺฒนาหจฺจ, \\ อสมฺมตา อตฺถงฺคเต, \\ ทเท สิพฺเพ วิธาเนน, \\ ปิณฺฑํ คณํ ปรํ ปูวํ, \\ วิกาลํ สนฺนิธิ ขีรํ, \\ อเจลกํ อุยฺโยขชฺช3, \\ นิมนฺติโต ปจฺจเยหิ, \\ สุรา องฺคุลิ หาโส จ, \\ โชติ นหาน ทุพฺพณฺณํ,}{\csromangathabat{มุสา โอมสเปสุญฺญํ,}{ปทเสยฺยา จ อิตฺถิยา.} \\ \csromangathabat{อญฺญตฺร วิญฺญุนา ภูตา1,}{ทุฏฺฐุลฺลาปตฺติ ขณนา.} \\ \csromangathabat{ภูตํ อญฺญาย อุชฺฌายิ2,}{มญฺโจ เสยฺโย จ วุจฺจติ.} \\ \csromangathabat{ปุพฺเพ นิกฺกฑฺฒนาหจฺจ,}{ทฺวารํ สปฺปาณเกน จ.} \\ \csromangathabat{อสมฺมตา อตฺถงฺคเต,}{อุปสฺสยามิเสน จ.} \\ \csromangathabat{ทเท สิพฺเพ วิธาเนน,}{นาวา ภุญฺเชยฺย เอกโต.} \\ \csromangathabat{ปิณฺฑํ คณํ ปรํ ปูวํ,}{ปวาริโต ปวาริตํ.} \\ \csromangathabat{วิกาลํ สนฺนิธิ ขีรํ,}{ทนฺตโปเนน เต ทส.} \\ \csromangathabat{อเจลกํ อุยฺโยขชฺช3,}{ปฏิจฺฉนฺนํ รเหน จ.} \\ \csromangathabat{นิมนฺติโต ปจฺจเยหิ,}{เสนาวสนุยฺโยธิกํ.} \\ \csromangathabat{สุรา องฺคุลิ หาโส จ,}{อนาทริยญฺจ ภิํสนํ.} \\ \csromangathabat{โชติ นหาน ทุพฺพณฺณํ,}{สามํ อปนิเธน จ.}}
\csromangathasetleft{มุสา โอมสเปสุญฺญํ, \\ อญฺญตฺร วิญฺญุนา ภูตา1, \\ ภูตํ อญฺญาย อุชฺฌายิ2, \\ ปุพฺเพ นิกฺกฑฺฒนาหจฺจ, \\ อสมฺมตา อตฺถงฺคเต, \\ ทเท สิพฺเพ วิธาเนน, \\ ปิณฺฑํ คณํ ปรํ ปูวํ, \\ วิกาลํ สนฺนิธิ ขีรํ, \\ อเจลกํ อุยฺโยขชฺช3, \\ นิมนฺติโต ปจฺจเยหิ, \\ สุรา องฺคุลิ หาโส จ, \\ โชติ นหาน ทุพฺพณฺณํ,}
\csromangathagroup{\csromangathastanza{\csromangathabat{มุสา โอมสเปสุญฺญํ,}{ปทเสยฺยา จ อิตฺถิยา.} \\ \csromangathabat{อญฺญตฺร วิญฺญุนา ภูตา1,}{ทุฏฺฐุลฺลาปตฺติ ขณนา.}}\csromangathastanza{\csromangathabat{ภูตํ อญฺญาย อุชฺฌายิ2,}{มญฺโจ เสยฺโย จ วุจฺจติ.} \\ \csromangathabat{ปุพฺเพ นิกฺกฑฺฒนาหจฺจ,}{ทฺวารํ สปฺปาณเกน จ.}}\csromangathastanza{\csromangathabat{อสมฺมตา อตฺถงฺคเต,}{อุปสฺสยามิเสน จ.} \\ \csromangathabat{ทเท สิพฺเพ วิธาเนน,}{นาวา ภุญฺเชยฺย เอกโต.}}\csromangathastanza{\csromangathabat{ปิณฺฑํ คณํ ปรํ ปูวํ,}{ปวาริโต ปวาริตํ.} \\ \csromangathabat{วิกาลํ สนฺนิธิ ขีรํ,}{ทนฺตโปเนน เต ทส.}}\csromangathastanza{\csromangathabat{อเจลกํ อุยฺโยขชฺช3,}{ปฏิจฺฉนฺนํ รเหน จ.} \\ \csromangathabat{นิมนฺติโต ปจฺจเยหิ,}{เสนาวสนุยฺโยธิกํ.}}\csromangathastanza{\csromangathabat{สุรา องฺคุลิ หาโส จ,}{อนาทริยญฺจ ภิํสนํ.} \\ \csromangathabat{โชติ นหาน ทุพฺพณฺณํ,}{สามํ อปนิเธน จ.}}}

\proseitem{1}{เทสนาโรจนา เจว (สี, สฺยา) 2. ภูตญฺญวาทอุชฺฌายิ (สี)}

\proseitem{3}{อเจลกา\-นุปขชฺช (ก)}

\setlength{\csromangathapagewidth}{0pt}
\setlength{\csromangathawakwidth}{0pt}
\csromangathameasuregroup{สญฺจิจฺจุทกกมฺมา จ, \\ เถยฺยอิตฺถิอวเทสํ\textsuperscript{1}, \\ สหธมฺมิกวิเลขา, \\ อมูลกญฺจ สญฺจิจฺจ, \\ สํเฆน จีวรํ ทตฺวา, \\ รญฺญญฺจ รตนํ สนฺตํ, \\ นิสีทนํ กณฺฑุจฺฉาทิ,}{\csromangathabat{สญฺจิจฺจุทกกมฺมา จ,}{ทุฏฺฐุลฺลํ อูนวีสติ.} \\ \csromangathabat{เถยฺยอิตฺถิอวเทสํ\textsuperscript{1},}{สํวาเส นาสิเตน จ.} \\ \csromangathabat{สหธมฺมิกวิเลขา,}{โมโห ปหาเรนุคฺคิเร.} \\ \csromangathabat{อมูลกญฺจ สญฺจิจฺจ,}{โสสฺสามิ ขิยฺยปกฺกเม.} \\ \csromangathabat{สํเฆน จีวรํ ทตฺวา,}{ปริณาเมยฺย ปุคฺคเล.} \\ \csromangathabat{รญฺญญฺจ รตนํ สนฺตํ,}{สูจิ มญฺโจ จ ตูลิกา.} \\ \csromangathabat{นิสีทนํ กณฺฑุจฺฉาทิ,}{วสฺสิกา สุคเตน จาติ.} \\ เตสํ วคฺคานํ อุทฺทานํ. \\ มุสา ภูตา จ โอวาโท, \\ โภชนาเจลเกน จ.}
\csromangathasetleft{สญฺจิจฺจุทกกมฺมา จ, \\ เถยฺยอิตฺถิอวเทสํ\textsuperscript{1}, \\ สหธมฺมิกวิเลขา, \\ อมูลกญฺจ สญฺจิจฺจ, \\ สํเฆน จีวรํ ทตฺวา, \\ รญฺญญฺจ รตนํ สนฺตํ, \\ นิสีทนํ กณฺฑุจฺฉาทิ,}
\csromangathagroup{\csromangathastanza{\csromanfolio{42}\csromangathabat{สญฺจิจฺจุทกกมฺมา จ,}{ทุฏฺฐุลฺลํ อูนวีสติ.} \\ \csromangathabat{เถยฺยอิตฺถิอวเทสํ\csromansharedfootnote{adecc06dc80c299b}{อริฏฺฐกํ (วิภงฺเค)},}{สํวาเส นาสิเตน จ.}}\csromangathastanza{\csromangathabat{สหธมฺมิกวิเลขา,}{โมโห ปหาเรนุคฺคิเร.} \\ \csromangathabat{อมูลกญฺจ สญฺจิจฺจ,}{โสสฺสามิ ขิยฺยปกฺกเม.}}\csromangathastanza{\csromangathabat{สํเฆน จีวรํ ทตฺวา,}{ปริณาเมยฺย ปุคฺคเล.} \\ \csromangathabat{รญฺญญฺจ รตนํ สนฺตํ,}{สูจิ มญฺโจ จ ตูลิกา.}}\csromangathastanza{\csromangathabat{นิสีทนํ กณฺฑุจฺฉาทิ,}{วสฺสิกา สุคเตน จาติ.} \\ เตสํ วคฺคานํ อุทฺทานํ. \\ มุสา ภูตา จ โอวาโท,}\csromangathastanza{โภชนาเจลเกน จ.}}

\vspace{\tipitakaparskip}
\csromancenter{สุรา สปฺปาณกา ธมฺโม, ราชวคฺเคน เต นวาติ.}
\csromansectionrule
\csromanmatikaanchor{mtk.20}
\csromantocmarkref{chapter}{6. ปาฏิเทสนียกณฺฑ}{mtk.20}
\bookmark[dest={mtk.20},level=0]{6. ปาฏิเทสนียกณฺฑ}
\chapterhead{6. ปาฏิเทสนีย\-กณฺฑ}
\proseitem{146}{ยํ เตน ภควตา ชานตา ปสฺสตา อรหตา สมฺมา\-สมฺพุทฺเธน อญฺญาติกาย ภิกฺขุนิยา อนฺตรฆรํ ปวิฏฺฐาย หตฺถโต ขาทนียํ วา โภชนียํ วา สหตฺถา ปฏิคฺคเหตฺวา ภุญฺชนฺตสฺส ปาฏิเทสนียํ กตฺถ ปญฺญตฺตนฺติ สาวตฺถิยํ ปญฺญตฺตํ. กํ อารพฺภาติ อญฺญตรํ ภิกฺขุํ อารพฺภ. กิสฺมิํ วตฺถุสฺมินฺติ อญฺญตโร ภิกฺขุ อญฺญาติกาย ภิกฺขุนิยา อนฺตรฆรํ ปวิฏฺฐาย หตฺถโต อามิสํ ปฏิคฺคเหสิ, ตสฺมิํ วตฺถุสฺมิํ. เอกา ปญฺญตฺติ. ฉนฺนํ อาปตฺติ\-สมุฏฺฐานานํ ทฺวีหิ สมุฏฺฐาเนหิ สมุฏฺฐาติ, สิยา กายโต สมุฏฺฐาติ, น วาจโต, น จิตฺตโต, สิยา กายโต จ จิตฺตโต จ สมุฏฺฐาติ, น วาจโต ฯเปฯ.}

\proseitem{147}{ภิกฺขุนิยา โวสาสนฺติยา น นิวาเรตฺวา ภุญฺชนฺตสฺส ปาฏิเทสนียํ กตฺถ ปญฺญตฺตนฺติ ราชคเห ปญฺญตฺตํ. กํ อารพฺภาติ ฉพฺพคฺคิเย ภิกฺขู อารพฺภ. กิสฺมิํ วตฺถุสฺมินฺติ ฉพฺพคฺคิยา ภิกฺขู ภิกฺขุนิโย โวสาสนฺติโย น นิวาเรสุํ, ตสฺมิํ วตฺถุสฺมิํ. เอกา ปญฺญตฺติ. ฉนฺนํ อาปตฺติ\-สมุฏฺฐานานํ ทฺวีหิ สมุฏฺฐาเนหิ สมุฏฺฐาติ, สิยา กายโต จ วาจโต จ สมุฏฺฐาติ, น จิตฺตโต, สิยา กายโต จ วาจโต จ จิตฺตโต จ สมุฏฺฐาติ ฯเปฯ.}

\proseitem{148}{\csromanfolio{43}เสกฺข\-สมฺมเตสุ กุเลสุ ขาทนียํ วา โภชนียํ วา สหตฺถา ปฏิคฺคเหตฺวา ภุญฺชนฺตสฺส ปาฏิเทสนียํ กตฺถ ปญฺญตฺตนฺติ สาวตฺถิยํ ปญฺญตฺตํ. กํ อารพฺภาติ สมฺพหุเล ภิกฺขู อารพฺภ. กิสฺมิํ วตฺถุสฺมินฺติ สมฺพหุลา ภิกฺขู น มตฺตํ ชานิตฺวา ปฏิคฺคเหสุํ, ตสฺมิํ วตฺถุสฺมิํ. เอกา ปญฺญตฺติ, เทฺว อนุปญฺญตฺติโย. ฉนฺนํ อาปตฺติ สมุฏฺฐานานํ ทฺวีหิ สมุฏฺฐาเนหิ สมุฏฺฐาติ, สิยา กายโต สมุฏฺฐาติ, น วาจโต, น จิตฺตโต, สิยา กายโต จ จิตฺตโต จ สมุฏฺฐาติ, น วาจโต ฯเปฯ.}

\proseitemwithcloser{149}{อารญฺญเกสุ เสนาสเนสุ ปุพฺเพ อปฺปฏิสํวิทิตํ ขาทนียํ วา โภชนียํ วา อชฺฌาราเม สหตฺถา ปฏิคฺคเหตฺวา ภุญฺชนฺตสฺส ปาฏิเทสนียํ กตฺถ ปญฺญตฺตนฺติ สกฺเกสุ ปญฺญตฺตํ. กํ อารพฺภาติ สมฺพหุเล ภิกฺขู อารพฺภ. กิสฺมิํ วตฺถุสฺมินฺติ สมฺพหุลา ภิกฺขู อาราเม โจเร ปฏิวสนฺเต นาโรเจสุํ, ตสฺมิํ วตฺถุสฺมิํ. เอกา ปญฺญตฺติ, เอกา อนุปญฺญตฺติ. ฉนฺนํ อาปตฺติ\-สมุฏฺฐานานํ ทฺวีหิ สมุฏฺฐาเนหิ สมุฏฺฐาติ, สิยา กายโต จ วาจโต จ สมุฏฺฐาติ, น จิตฺตโต, สิยา กายโต จ วาจโต จ จิตฺตโต จ สมุฏฺฐาติ ฯเปฯ.}{จตฺตาโร ปาฏิเทสนียา นิฏฺฐิตา.}
\csromancenter{ตสฺสุทฺทานํ.}
\setlength{\csromangathapagewidth}{0pt}
\setlength{\csromangathawakwidth}{0pt}
\csromangathameasuregroup{อญฺญาติกาย โวสาสํ,}{\csromangathabat{อญฺญาติกาย โวสาสํ,}{เสกฺขอารญฺญเกน จ.}}
\csromangathasetleft{อญฺญาติกาย โวสาสํ,}
\csromangathagroup{\csromangathastanza{\csromangathabat{อญฺญาติกาย โวสาสํ,}{เสกฺขอารญฺญเกน จ.}}}

\vspace{\tipitakaparskip}
\csromancenter{ปาฏิเทสนียา จตฺตาโร, สมฺพุทฺเธน ปกาสิตาติ.}
\csromansectionrule
\csromanmatikaanchor{mtk.21}
\csromantocmarkref{chapter}{7. เสขิยกณฺฑ}{mtk.21}
\bookmark[dest={mtk.21},level=0]{7. เสขิยกณฺฑ}
\chapterhead{7. \textbf{เสขิยกณฺฑ}}
\csromanmatikaanchor{mtk.22}
\csromantocmarkref{section}{1. ปริมณฺฑลวคฺค}{mtk.22}
\bookmark[dest={mtk.22},level=1]{1. ปริมณฺฑลวคฺค}
\csromanheader{1}{1. ปริมณฺฑล\-วคฺค}
\proseitem{150}{ยํ เตน ภควตา ชานตา ปสฺสตา อรหตา สมฺมา\-สมฺพุทฺเธน อนาทริยํ ปฏิจฺจ ปุรโต วา ปจฺฉโต วา โอลมฺเพนฺเตน นิวาเสนฺตสฺส ทุกฺกฏํ กตฺถ ปญฺญตฺตนฺติ สาวตฺถิยํ ปญฺญตฺตํ. กํ อารพฺภาติ ฉพฺพคฺคิเย ภิกฺขู อารพฺภ. กิสฺมิํ วตฺถุสฺมินฺติ ฉพฺพคฺคิยา ภิกฺขู ปุรโตปิ ปจฺฉโตปิ โอลมฺเพนฺตา นิวาเสสุํ, ตสฺมิํ วตฺถุสฺมิํ. เอกา ปญฺญตฺติ. ฉนฺนํ \csromanfolio{44}อาปตฺติ\-สมุฏฺฐานานํ เอเกน สมุฏฺฐาเนน สมุฏฺฐาติ, กายโต จ จิตฺตโต จ สมุฏฺฐาติ, น วาจโต ฯเปฯ.}

\prose{อนาทริยํ ปฏิจฺจ ปุรโต วา ปจฺฉโต วา โอลมฺเพนฺเตน ปารุปนฺตสฺส ทุกฺกฏํ กตฺถ ปญฺญตฺตนฺติ สาวตฺถิยํ ปญฺญตฺตํ. กํ อารพฺภาติ ฉพฺพคฺคิเย ภิกฺขู อารพฺภ. กิสฺมิํ วตฺถุสฺมินฺติ ฉพฺพคฺคิยา ภิกฺขู ปุรโตปิ ปจฺฉโตปิ โอลมฺเพนฺตา ปารุปิํสุ, ตสฺมิํ วตฺถุสฺมิํ. เอกา ปญฺญตฺติ. ฉนฺนํ อาปตฺติ\-สมุฏฺฐานานํ เอเกน สมุฏฺฐาเนน สมุฏฺฐาติ, กายโต จ จิตฺตโต จ สมุฏฺฐาติ, น วาจโต ฯเปฯ.}

\prose{อนาทริยํ ปฏิจฺจ กายํ วิวริตฺวา อนฺตรฆเร คจฺฉนฺตสฺส ทุกฺกฏํ ฯเปฯ. เอกา ปญฺญตฺติ. เอเกน สมุฏฺฐาเนน สมุฏฺฐาติ, กายโต จ จิตฺตโต จ สมุฏฺฐาติ, น วาจโต ฯเปฯ.}

\prose{อนาทริยํ ปฏิจฺจ กายํ วิวริตฺวา อนฺตรฆเร นิสีทนฺตสฺส ทุกฺกฏํ ฯเปฯ. เอกา ปญฺญตฺติ. เอเกน สมุฏฺฐาเนน สมุฏฺฐาติ, กายโต จ จิตฺตโต จ สมุฏฺฐาติ, น วาจโต ฯเปฯ.}

\prose{อนาทริยํ ปฏิจฺจ หตฺถํ วา ปาทํ วา กีฬาเปนฺเตน อนฺตรฆเร คจฺฉนฺตสฺส ทุกฺกฏํ ฯเปฯ. เอกา ปญฺญตฺติ. เอเกน สมุฏฺฐาเนน สมุฏฺฐาติ, กายโต จ จิตฺตโต จ สมุฏฺฐาติ, น วาจโต ฯเปฯ.}

\prose{อนาทริยํ ปฏิจฺจ หตฺถํ วา ปาทํ วา กีฬาเปนฺเตน อนฺตรฆเร นิสีทนฺตสฺส ทุกฺกฏํ ฯเปฯ. เอกา ปญฺญตฺติ. เอเกน สมุฏฺฐาเนน สมุฏฺฐาติ, กายโต จ จิตฺตโต จ สมุฏฺฐาติ, น วาจโต ฯเปฯ.}

\prose{อนาทริยํ ปฏิจฺจ ตหํ ตหํ โอโลเกนฺเตน อนฺตรฆเร คจฺฉนฺตสฺส ทุกฺกฏํ ฯเปฯ. เอกา ปญฺญตฺติ. เอเกน สมุฏฺฐาเนน สมุฏฺฐาติ, กายโต จ จิตฺตโต จ สมุฏฺฐาติ, น วาจโต ฯเปฯ.}

\prose{อนาทริยํ ปฏิจฺจ ตหํ ตหํ โอโลเกนฺเตน อนฺตรฆเร นิสีทนฺตสฺส ทุกฺกฏํ ฯเปฯ. เอกา ปญฺญตฺติ. เอเกน สมุฏฺฐาเนน สมุฏฺฐาติ, กายโต จ จิตฺตโต จ สมุฏฺฐาติ, น วาจโต ฯเปฯ.}

\prose{อนาทริยํ ปฏิจฺจ อุกฺขิตฺตกาย อนฺตรฆเร คจฺฉนฺตสฺส ทุกฺกฏํ ฯเปฯ. เอกา ปญฺญตฺติ. เอเกน สมุฏฺฐาเนน สมุฏฺฐาติ, กายโต จ จิตฺตโต จ สมุฏฺฐาติ, น วาจโต ฯเปฯ.}

\prosewithclosermidruled{\csromanfolio{45}อนาทริยํ ปฏิจฺจ อุกฺขิตฺตกาย อนฺตรฆเร นิสีทนฺตสฺส ทุกฺกฏํ ฯเปฯ. เอกา ปญฺญตฺติ. เอเกน สมุฏฺฐาเนน สมุฏฺฐาติ, กายโต จ จิตฺตโต จ สมุฏฺฐาติ, น วาจโต ฯเปฯ.}{ปริมณฺฑล\-วคฺโค ปฐโม.}
\csromanmatikaanchor{mtk.23}
\csromantocmarkref{section}{2. อุชฺชคฺฆิกวคฺค}{mtk.23}
\bookmark[dest={mtk.23},level=1]{2. อุชฺชคฺฆิกวคฺค}
\csromanheader{1}{2. อุชฺชคฺฆิก\-วคฺค}
\proseitem{151}{อนาทริยํ ปฏิจฺจ อุชฺชคฺฆิกาย อนฺตรฆเร คจฺฉนฺตสฺส ทุกฺกฏํ กตฺถ ปญฺญตฺตนฺติ สาวตฺถิยํ ปญฺญตฺตํ. กํ อารพฺภาติ ฉพฺพคฺคิเย ภิกฺขู อารพฺภ. กิสฺมิํ วตฺถุสฺมินฺติ ฉพฺพคฺคิยา ภิกฺขู มหาหสิตํ หสนฺตา อนฺตรฆเร คจฺฉิํสุ, ตสฺมิํ วตฺถุสฺมิํ. เอกา ปญฺญตฺติ. ฉนฺนํ อาปตฺติ\-สมุฏฺฐานานํ เอเกน สมุฏฺฐาเนน สมุฏฺฐาติ, กายโต จ วาจโต จ จิตฺตโต จ สมุฏฺฐาติ ฯเปฯ.}

\prose{อนาทริยํ ปฏิจฺจ อุชฺชคฺฆิกาย อนฺตรฆเร นิสีทนฺตสฺส ทุกฺกฏํ กตฺถ ปญฺญตฺตนฺติ สาวตฺถิยํ ปญฺญตฺตํ. กํ อารพฺภาติ ฉพฺพคฺคิเย ภิกฺขู อารพฺภ. กิสฺมิํ วตฺถุสฺมินฺติ ฉพฺพคฺคิยา ภิกฺขู มหาหสิตํ หสนฺตา อนฺตรฆเร นิสีทิํสุ, ตสฺมิํ วตฺถิสฺมิํ. เอกา ปญฺญตฺติ. ฉนฺนํ อาปตฺติ\-สมุฏฺฐานานํ เอเกน สมุฏฺฐาเนน สมุฏฺฐาติ, กายโต จ วาจโต จ}

\prose{อนาทริยํ ปฏิจฺจ อุจฺจาสทฺทํ มหาสทฺทํ กโรนฺเตน อนฺตรฆเร คจฺฉนฺตสฺส ทุกฺกฏํ กตฺถ ปญฺญตฺตนฺติ สาวตฺถิยํ ปญฺญตฺตํ. กํ อารพฺภาติ ฉพฺพคฺคิเย ภิกฺขู อารพฺภ. กิสฺมิํ วตฺถุสฺมินฺติ ฉพฺพคฺคิยา ภิกฺขู อุจฺจาสทฺทํ มหาสทฺทํ กโรนฺตา อนฺตรฆเร คจฺฉิํสุ, ตสฺมิํ วตฺถุสฺมิํ. เอกา ปญฺญตฺติ. ฉนฺนํ อาปตฺติ\-สมุฏฺฐานานํ เอเกน สมุฏฺฐาเนน สมุฏฺฐาติ, กายโต จ วาจโต จ จิตฺตโต จ สมุฏฺฐาติ ฯเปฯ.}

\prose{อนาทริยํ ปฏิจฺจ อุจฺจาสทฺทํ มหาสทฺทํ กโรนฺเตน อนฺตรฆเร นิสีทนฺตสฺส ทุกฺกฏํ กตฺถ ปญฺญตฺตนฺติ สาวตฺถิยํ ปญฺญตฺตํ. กํ อารพฺภาติ ฉพฺพคฺคิเย ภิกฺขู อารพฺภ. กิสฺมิํ วตฺถุสฺมินฺติ ฉพฺพคฺคิยา ภิกฺขู อุจฺจาสทฺทํ มหาสทฺทํ กโรนฺตา อนฺตรฆเร นิสีทิํสุ, ตสฺมิํ วตฺถุสฺมิํ. เอกา \csromanfolio{46}ปญฺญตฺติ. ฉนฺนํ อาปตฺติ\-สมุฏฺฐานานํ เอเกน สมุฏฺฐาเนน สุมุฏฺฐาติ, กายโต จ วาจโต จ จิตฺตโต จ สมุฏฺฐาติ ฯเปฯ.}

\prose{อนาทริยํ ปฏิจฺจ กายปฺปจาลกํ อนฺตรฆเร คจฺฉนฺตสฺส ทุกฺกฏํ ฯเปฯ. เอกา ปญฺญตฺติ. เอเกน สมุฏฺฐาเนน สมุฏฺฐาติ, กายโต จ จิตฺตโต จ สมุฏฺฐาติ, น วาจโต ฯเปฯ.}

\prose{อนาทริยํ ปฏิจฺจ กายปฺปจาลกํ อนฺตรฆเร นิสีทนฺตสฺส ทุกฺกฏํ ฯเปฯ. เอกา ปญฺญตฺติ. เอเกน สมุฏฺฐาเนน สมุฏฺฐาติ, กายโต จ จิตฺตโต จ สมุฏฺฐาติ, น วาจโต ฯเปฯ.}

\prose{อนาทริยํ ปฏิจฺจ พาหุปฺปจาลกํ อนฺตรฆเร คจฺฉนฺตสฺส ทุกฺกฏํ ฯเปฯ. เอกา ปญฺญตฺติ. เอเกน สมุฏฺฐาเนน สมุฏฺฐาติ, กายโต จ จิตฺตโต จ สมุฏฺฐาติ, น วาจโต ฯเปฯ.}

\prose{อนาทริยํ ปฏิจฺจ พาหุปฺปจาลกํ อนฺตรฆเร นิสีทนฺตสฺส ทุกฺกฏํ ฯเปฯ. เอกา ปญฺญตฺติ. เอเกน สมุฏฺฐาเนน สมุฏฺฐาติ, กายโต จ จิตฺตโต จ สมุฏฺฐาติ, น วาจโต ฯเปฯ.}

\prose{อนาทริยํ ปฏิจฺจ สีสปฺปจาลกํ อนฺตรฆเร คจฺฉนฺตสฺส ทุกฺกฏํ ฯเปฯ. เอกา ปญฺญตฺติ. เอเกน สมุฏฺฐาเนน สมุฏฺฐาติ, กายโต จ จิตฺตโต จ สมุฏฺฐาติ, น วาจโต ฯเปฯ.}

\prosewithclosermidruled{อนาทริยํ ปฏิจฺจ สีสปฺปจาลกํ อนฺตรฆเร นิสีทนฺตสฺส ทุกฺกฏํ ฯเปฯ. เอกา ปญฺญตฺติ. เอเกน สมุฏฺฐาเนน สมุฏฺฐาติ, กายโต จ จิตฺตโต จ สมุฏฺฐาติ, น วาจโต ฯเปฯ.}{อุชฺชคฺฆิก\-วคฺโค ทุติโย.}
\csromanmatikaanchor{mtk.24}
\csromantocmarkref{section}{3. ขมฺภกตวคฺค}{mtk.24}
\bookmark[dest={mtk.24},level=1]{3. ขมฺภกตวคฺค}
\csromanheader{1}{3. ขมฺภกต\-วคฺค}
\proseitem{152}{อนาทริยํ ปฏิจฺจ ขมฺภกเตน อนฺตรฆเร คจฺฉนฺตสฺส ทุกฺกฏํ ฯเปฯ. เอกา ปญฺญตฺติ. เอเกน สมุฏฺฐาเนน สมุฏฺฐาติ, กายโต จ จิตฺตโต จ สมุฏฺฐาติ, น วาจโต ฯเปฯ.}

\prose{\csromanfolio{47}อนาทริยํ ปฏิจฺจ ขมฺภกเตน อนฺตรฆเร นิสีทนฺตสฺส ทุกฺกฏํ ฯเปฯ. เอกา ปญฺญตฺติ. เอเกน สมุฏฺฐาเนน สมุฏฺฐาติ, กายโต จ จิตฺตโต จ สมุฏฺฐาติ, น วาจโต ฯเปฯ.}

\prose{อนาทริยํ ปฏิจฺจ โอคุณฺฐิเตน อนฺตรฆเร คจฺฉนฺตสฺส ทุกฺกฏํ กตฺถ ปญฺญตฺตนฺติ สาวตฺถิยํ ปญฺญตฺตํ. กํ อารพฺภาติ ฉพฺพคฺคิเย ภิกฺขู อารพฺภ. กิสฺมิํ วตฺถุสฺมินฺติ ฉพฺพคฺคิยา ภิกฺขู สสีสํ ปารุปิตฺวา อนฺตรฆเร คจฺฉิํสุ, ตสฺมิํ วตฺถุสฺมิํ. เอกา ปญฺญตฺติ. ฉนฺนํ อาปตฺติ\-สมุฏฺฐานานํ เอเกน สมุฏฺฐาเนน สมุฏฺฐาติ, กายโต จ จิตฺตโต จ สมุฏฺฐาติ, น วาจโต ฯเปฯ.}

\prose{อนาทริยํ ปฏิจฺจ โอคุณฺฐิเตน อนฺตรฆเร นิสีทนฺตสฺส ทุกฺกฏํ กตฺถ ปญฺญตฺตนฺติ สาวตฺถิยํ ปญฺญตฺตํ. กํ อารพฺภาติ ฉพฺพคฺคิเย ภิกฺขู อารพฺภ. กิสฺมิํ วตฺถุสฺมินฺติ ฉพฺพคฺคิยา ภิกฺขู สสีสํ ปารุปิตฺวา อนฺตรฆเร นิสีทิํสุ, ตสฺมิํ วตฺถุสฺมิํ. เอกา ปญฺญตฺติ. ฉนฺนํ อาปตฺติ\-สมุฏฺฐานานํ เอเกน สมุฏฺฐาเนน สมุฏฺฐาติ, กายโต จ จิตฺตโต จ สมุฏฺฐาติ, น วาจโต ฯเปฯ.}

\prose{อนาทริยํ ปฏิจฺจ อุกฺกุฏิกาย อนฺตรฆเร คจฺฉนฺตสฺส ทุกฺกฏํ ฯเปฯ. เอกา ปญฺญตฺติ. เอเกน สมุฏฺฐาเนน สมุฏฺฐาติ, กายโต จ จิตฺตโต จ สมุฏฺฐาติ, น วาจโต ฯเปฯ.}

\prose{อนาทริยํ ปฏิจฺจ ปลฺลตฺถิกาย\csromansharedfootnote{d5ec3658f0e08d1a}{ปลฺลตฺติกาย (ก)} อนฺตรฆเร นิสีทนฺตสฺส ทุกฺกฏํ ฯเปฯ. เอกา ปญฺญตฺติ. เอเกน สมุฏฺฐาเนน สมุฏฺฐาติ, กายโต จ จิตฺตโต จ สมุฏฺฐาติ, น วาจโต ฯเปฯ.}

\prose{อนาทริยํ ปฏิจฺจ อสกฺกจฺจํ ปิณฺฑปาตํ ปฏิคฺคณฺหนฺตสฺส ทุกฺกฏํ ฯเปฯ. เอกา ปญฺญตฺติ. เอเกน สมุฏฺฐาเนน สมุฏฺฐาติ, กายโต จ จิตฺตโต จ สมุฏฺฐาติ, น วาจโต ฯเปฯ.}

\prose{อนาทริยํ ปฏิจฺจ ตหํ ตหํ โอโลเกนฺเตน ปิณฺฑปาตํ ปฏิคฺคณฺหนฺตสฺส ทุกฺกฏํ ฯเปฯ. เอกา ปญฺญตฺติ. เอเกน สมุฏฺฐาเนน สมุฏฺฐาติ, กายโต จ จิตฺตโต จ สมุฏฺฐาติ, น วาจโต ฯเปฯ.}

\prose{อนาทริยํ ปฏิจฺจ สูปญฺเญว พหุํ ปฏิคฺคณฺหนฺตสฺส ทุกฺกฏํ ฯเปฯ. เอกา ปญฺญตฺติ. เอเกน สมุฏฺฐาเนน สมุฏฺฐาติ, กายโต จ จิตฺตโต จ สมุฏฺฐาติ, น วาจโต ฯเปฯ.}

\prosewithclosermidruled{\csromanfolio{48}อนาทริยํ ปฏิจฺจ ถูปีกตํ ปิณฺฑปาตํ ปฏิคฺคณฺหนฺตสฺส ทุกฺกฏํ ฯเปฯ. เอกา ปญฺญตฺติ. เอเกน สมุฏฺฐาเนน สมุฏฺฐาติ, กายโต จ จิตฺตโต จ สมุฏฺฐาติ, น วาจโต ฯเปฯ.}{ขมฺภกต\-วคฺโค ตติโย.}
\csromanmatikaanchor{mtk.25}
\csromantocmarkref{section}{4. ปิณฺฑปาตวคฺค}{mtk.25}
\bookmark[dest={mtk.25},level=1]{4. ปิณฺฑปาตวคฺค}
\csromanheader{1}{4. \textbf{ปิณฺฑปาตวคฺค}}
\proseitem{153}{อนาทริยํ ปฏิจฺจ อสกฺกจฺจํ ปิณฺฑปาตํ ภุญฺชนฺตสฺส ทุกฺกฏํ ฯเปฯ. เอกา ปญฺญตฺติ. เอเกน สมุฏฺฐาเนน สมุฏฺฐาติ, กายโต จ จิตฺตโต จ สมุฏฺฐาติ, น วาจโต ฯเปฯ.}

\prose{อนาทริยํ ปฏิจฺจ ตหํ ตหํ โอโลเกนฺเตน ปิณฺฑปาตํ ภุญฺชนฺตสฺส ทุกฺกฏํ ฯเปฯ. เอกา ปญฺญตฺติ. เอเกน สมุฏฺฐาเนน สมุฏฺฐาติ, กายโต จ จิตฺตโต จ สมุฏฺฐาติ, น วาจโต ฯเปฯ.}

\prose{อนาทริยํ ปฏิจฺจ ตหํ ตหํ โอมสิตฺวา ปิณฺฑปาตํ ภุญฺชนฺตสฺส ทุกฺกฏํ ฯเปฯ. เอกา ปญฺญตฺติ. เอเกน สมุฏฺฐาเนน สมุฏฺฐาติ, กายโต จ จิตฺตโต จ สมุฏฺฐาติ, น วาจโต ฯเปฯ.}

\prose{อนาทริยํ ปฏิจฺจ สูปญฺเญว พหุํ ภุญฺชนฺตสฺส ทุกฺกฏํ ฯเปฯ. เอกา ปญฺญตฺติ. เอเกน สมุฏฺฐาเนน สมุฏฺฐาติ, กายโต จ จิตฺตโต จ สมุฏฺฐาติ, น วาจโต ฯเปฯ.}

\prose{อนาทริยํ ปฏิจฺจ ถูปกโต โอมทฺทิตฺวา ปิณฺฑปาตํ ภุญฺชนฺตสฺส ทุกฺกฏํ ฯเปฯ. เอกา ปญฺญตฺติ. เอเกน สมุฏฺฐาเนน สมุฏฺฐาติ, กายโต จ จิตฺตโต จ สมุฏฺฐาติ, น วาจโต ฯเปฯ.}

\prose{อนาทริยํ ปฏิจฺจ สูปํ วา พฺยญฺชนํ วา โอทเนน ปฏิจฺฉาเทนฺตสฺส ทุกฺกฏํ ฯเปฯ. เอกา ปญฺญตฺติ. เอเกน สมุฏฺฐาเนน สมุฏฺฐาติ, กายโต จ จิตฺตโต จ สมุฏฺฐาติ, น วาจโต ฯเปฯ.}

\prose{อนาทริยํ ปฏิจฺจ สูปํ วา โอทนํ วา อคิลาโน อตฺตโน อตฺถาย วิญฺญาเปตฺวา ภุญฺชนฺตสฺส ทุกฺกฏํ กตฺถ ปญฺญตฺตนฺติ สาวตฺถิยํ ปญฺญตฺตํ. กํ อารพฺภาติ ฉพฺพคฺคิเย ภิกฺขู อารพฺภ. กิสฺมิํ วตฺถุสฺมินฺติ ฉพฺพคฺคิยา ภิกฺขู สูปมฺปิ โอทนมฺปิ อตฺตโน อตฺถาย วิญฺญาเปตฺวา ภุญฺชิํสุ, ตสฺมิํ วตฺถุสฺมิํ. \csromanfolio{49}เอกา ปญฺญตฺติ, เอกา อนุปญฺญตฺติ. ฉนฺนํ อาปตฺติ\-สมุฏฺฐานานํ ทฺวีหิ สมุฏฺฐาเนหิ สมุฏฺฐาติ, สิยา กายโต จ จิตฺตโต จ สมุฏฺฐาติ, น วาจโต, สิยา กายโต จ วาจโต จ จิตฺตโต จ สมุฏฺฐาติ ฯเปฯ.}

\prose{อนาทริยํ ปฏิจฺจ อุชฺฌานสญฺญินา ปเรสํ ปตฺตํ โอโลเกนฺตสฺส ทุกฺกฏํ ฯเปฯ. เอกา ปญฺญตฺติ. เอเกน สมุฏฺฐาเนน สมุฏฺฐาติ, กายโต จ จิตฺตโต จ สมุฏฺฐาติ, น วาจโต ฯเปฯ.}

\prose{อนาทริยํ ปฏิจฺจ มหนฺตํ กพฬํ กโรนฺตสฺส ทุกฺกฏํ ฯเปฯ. เอกา ปญฺญตฺติ. เอเกน สมุฏฺฐาเนน สมุฏฺฐาติ, กายโต จ จิตฺตโต จ สมุฏฺฐาติ, น วาจโต ฯเปฯ.}

\prosewithclosermidruled{อนาทริยํ ปฏิจฺจ ทีฆํ อาโลปํ กโรนฺตสฺส ทุกฺกฏํ ฯเปฯ. เอกา ปญฺญตฺติ. เอเกน สมุฏฺฐาเนน สมุฏฺฐาติ, กายโต จ จิตฺตโต จ สมุฏฺฐาติ, น วาจโต ฯเปฯ.}{ปิณฺฑปาตวคฺโค จตุตฺโถ.}
\csromanmatikaanchor{mtk.26}
\csromantocmarkref{section}{5. กพฬวคฺค}{mtk.26}
\bookmark[dest={mtk.26},level=1]{5. กพฬวคฺค}
\csromanheader{1}{5. \textbf{กพฬวคฺค}}
\proseitem{154}{อนาทริยํ ปฏิจฺจ อนาหเฏ กพเฬ มุขทฺวารํ วิวรนฺตสฺส ทุกฺกฏํ ฯเปฯ. เอกา ปญฺญตฺติ. เอเกน สมุฏฺฐาเนน สมุฏฺฐาติ, กายโต จ จิตฺตโต จ สมุฏฺฐาติ, น วาจโต ฯเปฯ.}

\prose{อนาทริยํ ปฏิจฺจ ภุญฺชมาเนน สพฺพํ หตฺถํ มุเข ปกฺขิปนฺตสฺส ทุกฺกฏํ ฯเปฯ. เอกา ปญฺญตฺติ. เอเกน สมุฏฺฐาเนน สมุฏฺฐาติ, กายโต จ จิตฺตโต จ สมุฏฺฐาติ, น วาจโต ฯเปฯ.}

\prose{อนาทริยํ ปฏิจฺจ สกพเฬน มุเขน พฺยาหรนฺตสฺส ทุกฺกฏํ กตฺถ ปญฺญตฺตนฺติ สาวตฺถิยํ ปญฺญตฺตํ. กํ อารพฺภาติ ฉพฺพคฺคิเย ภิกฺขู อารพฺภ. กิสฺมิํ วตฺถุสฺมินฺติ ฉพฺพคฺคิยา ภิกฺขู สกพเฬน มุเขน พฺยาหริํสุ, ตสฺมิํ วตฺถุสฺมิํ. เอกา ปญฺญตฺติ. ฉนฺนํ อาปตฺติ\-สมุฏฺฐานานํ เอเกน สมุฏฺฐาเนน สมุฏฺฐาติ, กายโต จ วาจโต จ}

\prose{\csromanfolio{50}อนาทริยํ ปฏิจฺจ ปิณฺฑุกฺเขปกํ ภุญฺชนฺตสฺส ทุกฺกฏํ ฯเปฯ. เอกา ปญฺญตฺติ. เอเกน สมุฏฺฐาเนน สมุฏฺฐาติ, กายโต จ จิตฺตโต จ สมุฏฺฐาติ, น วาจโต ฯเปฯ.}

\prose{อนาทริยํ ปฏิจฺจ กพฬาวจฺเฉทกํ ภุญฺชนฺตสฺส ทุกฺกฏํ ฯเปฯ. เอกา ปญฺญตฺติ. เอเกน สมุฏฺฐาเนน สมุฏฺฐาติ, กายโต จ จิตฺตโต จ สมุฏฺฐาติ, น วาจโต ฯเปฯ.}

\prose{อนาทริยํ ปฏิจฺจ อวคณฺฑการกํ ภุญฺชนฺตสฺส ทุกฺกฏํ ฯเปฯ. เอกา ปญฺญตฺติ. เอเกน สมุฏฺฐาเนน สมุฏฺฐาติ, กายโต จ จิตฺตโต จ สมุฏฺฐาติ, น วาจโต ฯเปฯ.}

\prose{อนาทริยํ ปฏิจฺจ หตฺถนิทฺธุนกํ ภุญฺชนฺตสฺส ทุกฺกฏํ ฯเปฯ. เอกา ปญฺญตฺติ. เอเกน สมุฏฺฐาเนน สมุฏฺฐาติ, กายโต จ จิตฺตโต จ สมุฏฺฐาติ, น วาจโต ฯเปฯ.}

\prose{อนาทริยํ ปฏิจฺจ สิตฺถาวการกํ ภุญฺชนฺตสฺส ทุกฺกฏํ ฯเปฯ. เอกา ปญฺญตฺติ. เอเกน สมุฏฺฐาเนน สมุฏฺฐาติ, กายโต จ จิตฺตโต จ สมุฏฺฐาติ, น วาจโต ฯเปฯ.}

\prose{อนาทริยํ ปฏิจฺจ ชิวฺหานิจฺฉารกํ ภุญฺชนฺตสฺส ทุกฺกฏํ ฯเปฯ. เอกา ปญฺญตฺติ. เอเกน สมุฏฺฐาเนน สมุฏฺฐาติ, กายโต จ จิตฺตโต จ สมุฏฺฐาติ, น วาจโต ฯเปฯ.}

\prosewithclosermidruled{อนาทริยํ ปฏิจฺจ จปุจปุการกํ ภุญฺชนฺตสฺส ทุกฺกฏํ ฯเปฯ. เอกา ปญฺญตฺติ. เอเกน สมุฏฺฐาเนน สมุฏฺฐาติ, กายโต จ จิตฺตโต จ สมุฏฺฐาติ, น วาจโต ฯเปฯ.}{กพฬวคฺโค ปญฺจโม.}
\csromanmatikaanchor{mtk.27}
\csromantocmarkref{section}{6. สุรุสุรุวคฺค}{mtk.27}
\bookmark[dest={mtk.27},level=1]{6. สุรุสุรุวคฺค}
\csromanheader{1}{6. \textbf{สุรุสุรุวคฺค}}
\proseitem{155}{อนาทริยํ ปฏิจฺจ สุรุสุรุการกํ ภุญฺชนฺตสฺส ทุกฺกฏํ กตฺถ ปญฺญตฺตนฺติ โกสมฺพิยํ ปญฺญตฺตํ. กํ อารพฺภาติ สมฺพหุเล ภิกฺขู อารพฺภ. กิสฺมิํ วตฺถุสฺมินฺติ สมฺพหุลา ภิกฺขู สุรุสุรุการกํ ขีรํ ปิวิํสุ, ตสฺมิํ \csromanfolio{51}วตฺถุสฺมิํ. เอกา ปญฺญตฺติ. ฉนฺนํ อาปตฺติ\-สมุฏฺฐานานํ เอเกน สมุฏฺฐาเนน สมุฏฺฐาติ, กายโต จ จิตฺตโต จ สมุฏฺฐาติ, น วาจโต ฯเปฯ.}

\prose{อนาทริยํ ปฏิจฺจ หตฺถนิลฺเลหกํ ภุญฺชนฺตสฺส ทุกฺกฏํ ฯเปฯ. เอกา ปญฺญตฺติ. เอเกน สมุฏฺฐาเนน สมุฏฺฐาติ, กายโต จ จิตฺตโต จ สมุฏฺฐาติ, น วาจโต ฯเปฯ.}

\prose{อนาทริยํ ปฏิจฺจ ปตฺตนิลฺเลหกํ ภุญฺชนฺตสฺส ทุกฺกฏํ ฯเปฯ. เอกา ปญฺญตฺติ. เอเกน สมุฏฺฐาเนน สมุฏฺฐาติ, กายโต จ จิตฺตโต จ สมุฏฺฐาติ, น วาจโต ฯเปฯ.}

\prose{อนาทริยํ ปฏิจฺจ โอฏฺฐนิลฺเลหกํ ภุญฺชนฺตสฺส ทุกฺกฏํ ฯเปฯ. เอกา ปญฺญตฺติ. เอเกน สมุฏฺฐาเนน สมุฏฺฐาติ, กายโต จ จิตฺตโต จ สมุฏฺฐาติ, น วาจโต ฯเปฯ.}

\prose{อนาทริยํ ปฏิจฺจ สามิเสน หตฺเถน ปานียถาลกํ ปฏิคฺคณฺหนฺตสฺส ทุกฺกฏํ กตฺถ ปญฺญตฺตนฺติ ภคฺเคสุ ปญฺญตฺตํ. กํ อารพฺภาติ สมฺพหุเล ภิกฺขู อารพฺภ. กิสฺมิํ วตฺถุสฺมินฺติ สมฺพหุลา ภิกฺขู สามิเสน หตฺเถน ปานียถาลกํ ปฏิคฺคเหสุํ, ตสฺมิํ วตฺถุสฺมิํ. เอกา ปญฺญตฺติ. ฉนฺนํ อาปตฺติ\-สมุฏฺฐานานํ เอเกน สมุฏฺฐาเนน สมุฏฺฐาติ, กายโต จ จิตฺตโต จ สมุฏฺฐาติ, น วาจโต ฯเปฯ.}

\prose{อนาทริยํ ปฏิจฺจ สสิตฺถกํ ปตฺตโธวนํ อนฺตรฆเร ฉฑฺเฑนฺตสฺส ทุกฺกฏํ กตฺถ ปญฺญตฺตนฺติ ภคฺเคสุ ปญฺญตฺตํ. กํ อารพฺภาติ สมฺพหุเล ภิกฺขู อารพฺภ. กิสฺมิํ วตฺถุสฺมินฺติ สมฺพหุลา ภิกฺขู สสิตฺถกํ ปตฺตโธวนํ อนฺตรฆเร ฉฑฺเฑสุํ, ตสฺมิํ วตฺถุสฺมิํ. เอกา ปญฺญตฺติ. ฉนฺนํ อาปตฺติ\-สมุฏฺฐานานํ เอเกน สมุฏฺฐาเนน สมุฏฺฐาติ, กายโต จ จิตฺตโต จ สมุฏฺฐาติ, น วาจโต ฯเปฯ.}

\prose{อนาทริยํ ปฏิจฺจ ฉตฺตปาณิสฺส ธมฺมํ เทเสนฺตสฺส ทุกฺกฏํ กตฺถ ปญฺญตฺตนฺติ สาวตฺถิยํ ปญฺญตฺตํ. กํ อารพฺภาติ ฉพฺพคฺคิเย ภิกฺขู อารพฺภ. กิสฺมิํ วตฺถุสฺมินฺติ ฉพฺพคฺคิยา ภิกฺขู ฉตฺตปาณิสฺส ธมฺมํ เทเสสุํ, ตสฺมิํ วตฺถุสฺมิํ. เอกา ปญฺญตฺติ, เอกา อนุปญฺญตฺติ. ฉนฺนํ อาปตฺติ\-สมุฏฺฐานานํ เอเกน สมุฏฺฐาเนน สมุฏฺฐาติ, วาจโต จ จิตฺตโต จ สมุฏฺฐาติ, น กายโต ฯเปฯ.}

\prose{\csromanfolio{52}อนาทริยํ ปฏิจฺจ ทณฺฑปาณิสฺส ธมฺมํ เทเสนฺตสฺส ทุกฺกฏํ ฯเปฯ. เอกา ปญฺญตฺติ, เอกา อนุปญฺญตฺติ. ฉนฺนํ อาปตฺติ\-สมุฏฺฐานานํ เอเกน สมุฏฺฐาเนน สมุฏฺฐาติ, วาจโต จ จิตฺตโต จ สมุฏฺฐาติ, น กายโต ฯเปฯ.}

\prose{อนาทริยํ ปฏิจฺจ สตฺถปาณิสฺส ธมฺมํ เทเสนฺตสฺส ทุกฺกฏํ ฯเปฯ. เอกา ปญฺญตฺติ, เอกา อนุปญฺญตฺติ. ฉนฺนํ อาปตฺติ\-สมุฏฺฐานานํ เอเกน สมุฏฺฐาเนน สมุฏฺฐาติ, วาจโต จ จิตฺตโต จ สมุฏฺฐาติ, น กายโต ฯเปฯ.}

\prosewithclosermidruled{อนาทริยํ ปฏิจฺจ อาวุธปาณิสฺส ธมฺมํ เทเสนฺตสฺส ทุกฺกฏํ ฯเปฯ. เอกา ปญฺญตฺติ, เอกา อนุปญฺญตฺติ. ฉนฺนํ อาปตฺติ\-สมุฏฺฐานานํ เอเกน สมุฏฺฐาเนน สมุฏฺฐาติ, วาจโต จ จิตฺตโต จ สมุฏฺฐาติ, น กายโต ฯเปฯ.}{สุรุสุรุวคฺโค ฉฏฺโฐ.}
\csromanmatikaanchor{mtk.28}
\csromantocmarkref{section}{7. ปาทุกวคฺค}{mtk.28}
\bookmark[dest={mtk.28},level=1]{7. ปาทุกวคฺค}
\csromanheader{1}{7. \textbf{ปาทุกวคฺค}}
\proseitem{156}{อนาทริยํ ปฏิจฺจ ปาทุการุฬฺหสฺส ธมฺมํ เทเสนฺตสฺส ทุกฺกฏํ ฯเปฯ. เอกา ปญฺญตฺติ, เอกา อนุปญฺญตฺติ. ฉนฺนํ อาปตฺติ\-สมุฏฺฐานานํ เอเกน สมุฏฺฐาเนน สมุฏฺฐาติ, วาจโต จ จิตฺตโต จ สมุฏฺฐาติ, น กายโต ฯเปฯ.}

\prose{อนาทริยํ ปฏิจฺจ อุปาหนา\-รุฬฺหสฺส ธมฺมํ เทเสนฺตสฺส ทุกฺกฏํ ฯเปฯ. เอกา ปญฺญตฺติ, เอกา อนุปญฺญตฺติ, ฉนฺนํ อาปตฺติ\-สมุฏฺฐานานํ เอเกน สมุฏฺฐาเนน สมุฏฺฐาติ, วาจโต จ จิตฺตโต จ สมุฏฺฐาติ, น กายโต ฯเปฯ.}

\prose{อนาทริยํ ปฏิจฺจ ยานคตสฺส ธมฺมํ เทเสนฺตสฺส ทุกฺกฏํ ฯเปฯ. เอกา ปญฺญตฺติ, เอกา อนุปญฺญตฺติ. ฉนฺนํ อาปตฺติ\-สมุฏฺฐานานํ เอเกน สมุฏฺฐาเนน สมุฏฺฐาติ, วาจโต จ จิตฺตโต จ สมุฏฺฐาติ, น กายโต ฯเปฯ.}

\prose{อนาทริยํ ปฏิจฺจ สยนคตสฺส ธมฺมํ เทเสนฺตสฺส ทุกฺกฏํ ฯเปฯ. เอกา ปญฺญตฺติ, เอกา อนุปญฺญตฺติ. ฉนฺนํ อาปตฺติ\-สมุฏฺฐานานํ เอเกน สมุฏฺฐาเนน สมุฏฺฐาติ, วาจโต จ จิตฺตโต จ สมุฏฺฐาติ, น กายโต ฯเปฯ.}

\prose{อนาทริยํ ปฏิจฺจ ปลฺลตฺถิกาย นิสินฺนสฺส ธมฺมํ เทเสนฺตสฺส ทุกฺกฏํ ฯเปฯ. เอกา ปญฺญตฺติ, เอกา อนุปญฺญตฺติ. ฉนฺนํ อาปตฺติ\-สมุฏฺฐานานํ เอเกน \csromanfolio{53}สมุฏฺฐาเนน สมุฏฺฐาติ, วาจโต จ จิตฺตโต จ สมุฏฺฐาติ, น กายโต ฯเปฯ.}

\prose{อนาทริยํ ปฏิจฺจ เวฐิตสีสสฺส ธมฺมํ เทเสนฺตสฺส ทุกฺกฏํ ฯเปฯ. เอกา ปญฺญตฺติ, เอกา อนุปญฺญตฺติ. ฉนฺนํ อาปตฺติ\-สมุฏฺฐานานํ เอเกน สมุฏฺฐาเนน สมุฏฺฐาติ, วาจโต จ จิตฺตโต จ สมุฏฺฐาติ, น กายโต ฯเปฯ.}

\prose{อนาทริยํ ปฏิจฺจ โอคุณฺฐิต\-สีสสฺส ธมฺมํ เทเสนฺตสฺส ทุกฺกฏํ ฯเปฯ. เอกา ปญฺญตฺติ, เอกา อนุปญฺญตฺติ. ฉนฺนํ อาปตฺติ\-สมุฏฺฐานานํ เอเกน สมุฏฺฐาเนน สมุฏฺฐาติ, วาจโต จ จิตฺตโต จ สมุฏฺฐาติ, น กายโต ฯเปฯ.}

\prose{อนาทริยํ ปฏิจฺจ ฉมายํ นิสีทิตฺวา อาสเน นิสินฺนสฺส ธมฺมํ เทเสนฺตสฺส ทุกฺกฏํ ฯเปฯ. เอกา ปญฺญตฺติ, เอกา อนุปญฺญตฺติ. ฉนฺนํ อาปตฺติ\-สมุฏฺฐานานํ เอเกน สมุฏฺฐาเนน สมุฏฺฐาติ, กายโต จ วาจโต จ}

\prose{อนาทริยํ ปฏิจฺจ นีเจ อาสเน นิสีทิตฺวา อุจฺเจ อาสเน นิสินฺนสฺส ธมฺมํ เทเสนฺตสฺส ทุกฺกฏํ ฯเปฯ. เอกา ปญฺญตฺติ, เอกา อนุปญฺญตฺติ. ฉนฺนํ อาปตฺติ\-สมุฏฺฐานานํ เอเกน สมุฏฺฐาเนน สมุฏฺฐาติ, กายโต จ วาจโต จ จิตฺตโต จ สมุฏฺฐาติ ฯเปฯ.}

\prose{อนาทริยํ ปฏิจฺจ ฐิเตน นิสินฺนสฺส ธมฺมํ เทเสนฺตสฺส ทุกฺกฏํ ฯเปฯ. เอกา ปญฺญตฺติ, เอกา อนุปญฺญตฺติ. ฉนฺนํ อาปตฺติ\-สมุฏฺฐานานํ เอเกน สมุฏฺฐาเนน สมุฏฺฐาติ, กายโต จ วาจโต จ จิตฺตโต จ สมุฏฺฐาติ ฯเปฯ.}

\prose{อนาทริยํ ปฏิจฺจ ปจฺฉโต คจฺฉนฺเตน ปุรโต คจฺฉนฺตสฺส ธมฺมํ เทเสนฺตสฺส ทุกฺกฏํ ฯเปฯ. เอกา ปญฺญตฺติ, เอกา อนุปญฺญตฺติ. ฉนฺนํ อาปตฺติ\-สมุฏฺฐานานํ เอเกน สมุฏฺฐาเนน สมุฏฺฐาติ, กายโต จ วาจโต จ จิตฺตโต จ สมุฏฺฐาติ ฯเปฯ.}

\prose{อนาทริยํ ปฏิจฺจ อุปฺปเถน คจฺฉนฺเตน ปเถน คจฺฉนฺตสฺส ธมฺมํ เทเสนฺตสฺส ทุกฺกฏํ ฯเปฯ. เอกา ปญฺญตฺติ, เอกา อนุปญฺญตฺติ. ฉนฺนํ อาปตฺติ\-สมุฏฺฐานานํ เอเกน สมุฏฺฐาเนน สมุฏฺฐาติ, กายโต จ วาจโต จ จิตฺตโต จ สมุฏฺฐาติ ฯเปฯ.}

\prose{อนาทริยํ ปฏิจฺจ ฐิเตน อุจฺจารํ วา ปสฺสาวํ วา กโรนฺตสฺส ทุกฺกฏํ ฯเปฯ. เอกา ปญฺญตฺติ, เอกา อนุปญฺญตฺติ. ฉนฺนํ อาปตฺติ\-สมุฏฺฐานานํ เอเกน \csromanfolio{54}สมุฏฺฐาเนน สมุฏฺฐาติ, กายโต จ จิตฺตโต จ สมุฏฺฐาติ, น วาจโต ฯเปฯ.}

\prose{อนาทริยํ ปฏิจฺจ หริเต อุจฺจารํ วา ปสฺสาวํ วา เขฬํ วา กโรนฺตสฺส ทุกฺกฏํ ฯเปฯ. เอกา ปญฺญตฺติ, เอกา อนุปญฺญตฺติ. ฉนฺนํ อาปตฺติ\-สมุฏฺฐานานํ เอเกน สมุฏฺฐาเนน สมุฏฺฐาติ, กายโต จ จิตฺตโต จ สมุฏฺฐาติ, น วาจโต ฯเปฯ.}

\prosewithclosermid{อนาทริยํ ปฏิจฺจ อุทเก อุจฺจารํ วา ปสฺสาวํ วา เขฬํ วา กโรนฺตสฺส ทุกฺกฏํ ฯเปฯ. กตฺถ ปญฺญตฺตนฺติ สาวตฺถิยํ ปญฺญตฺตํ. กํ อารพฺภาติ ฉพฺพคฺคิเย ภิกฺขู อารพฺภ. กิสฺมิํ วตฺถุสฺมินฺติ ฉพฺพคฺคิยา ภิกฺขู อุทเก อุจฺจารมฺปิ ปสฺสาวมฺปิ เขฬมฺปิ อกํสุ, ตสฺมิํ วตฺถุสฺมิํ. เอกา ปญฺญตฺติ, เอกา อนุปญฺญตฺติ. ฉนฺนํ อาปตฺติ\-สมุฏฺฐานานํ เอเกน สมุฏฺฐาเนน สมุฏฺฐาติ, กายโต จ จิตฺตโต จ สมุฏฺฐาติ, น วาจโต ฯเปฯ.}{ปาทุกวคฺโค สตฺตโม. ปญฺจสตฺตติ เสขิยา นิฏฺฐิตา.}
\csromancenter{ตสฺสุทฺทานํ.}
\setlength{\csromangathapagewidth}{0pt}
\setlength{\csromangathawakwidth}{0pt}
\csromangathameasuregroup{ปริมณฺฑลํ ปฏิจฺฉนฺนํ, \\ อุกฺขิตฺโต'ชฺชคฺฆิกา สทฺโท,}{\csromangathabat{ปริมณฺฑลํ ปฏิจฺฉนฺนํ,}{สุสํวุโต'กฺขิตฺตจกฺขุ.} \\ \csromangathabat{อุกฺขิตฺโต'ชฺชคฺฆิกา สทฺโท,}{ตโย เจว ปจาลนา.}}
\csromangathasetleft{ปริมณฺฑลํ ปฏิจฺฉนฺนํ, \\ อุกฺขิตฺโต'ชฺชคฺฆิกา สทฺโท,}
\csromangathagroup{\csromangathastanza{\csromangathabat{ปริมณฺฑลํ ปฏิจฺฉนฺนํ,}{สุสํวุโต'กฺขิตฺตจกฺขุ.} \\ \csromangathabat{อุกฺขิตฺโต'ชฺชคฺฆิกา สทฺโท,}{ตโย เจว ปจาลนา.}}}

\prose{ขมฺภํ โอคุณฺฐิโต เจวุ\-กฺกุฏิ\-ปลฺลตฺถิกาย จ.}

\setlength{\csromangathapagewidth}{0pt}
\setlength{\csromangathawakwidth}{0pt}
\csromangathameasuregroup{สกฺกจฺจํ ปตฺตสญฺญี จ, \\ สกฺกจฺจํ ปตฺตสญฺญี จ, \\ ถูปกโต ปฏิจฺฉนฺนํ, \\ น มหนฺตํ มณฺฑลํ ทฺวารํ, \\ อุกฺเขโป เฉทนา คณฺโฑ, \\ ชิวฺหานิจฺฉารกญฺเจว, \\ หตฺโถ ปตฺโต จ โอฏฺโฐ จ,}{\csromangathabat{สกฺกจฺจํ ปตฺตสญฺญี จ,}{สมสูปํ สมติตฺติกํ1.} \\ \csromangathabat{สกฺกจฺจํ ปตฺตสญฺญี จ,}{สปทานํ สมสูปกํ.} \\ \csromangathabat{ถูปกโต ปฏิจฺฉนฺนํ,}{วิญฺญตฺตุชฺฌานสญฺญินา.} \\ \csromangathabat{น มหนฺตํ มณฺฑลํ ทฺวารํ,}{สพฺพํ หตฺถํ น พฺยาหเร.} \\ \csromangathabat{อุกฺเขโป เฉทนา คณฺโฑ,}{ธุนํ สิตฺถาวการกํ.} \\ \csromangathabat{ชิวฺหานิจฺฉารกญฺเจว,}{จปุจปุ สุรุสุรุ.} \\ \csromangathabat{หตฺโถ ปตฺโต จ โอฏฺโฐ จ,}{สามิสํ สิตฺถเกน จ.}}
\csromangathasetleft{สกฺกจฺจํ ปตฺตสญฺญี จ, \\ สกฺกจฺจํ ปตฺตสญฺญี จ, \\ ถูปกโต ปฏิจฺฉนฺนํ, \\ น มหนฺตํ มณฺฑลํ ทฺวารํ, \\ อุกฺเขโป เฉทนา คณฺโฑ, \\ ชิวฺหานิจฺฉารกญฺเจว, \\ หตฺโถ ปตฺโต จ โอฏฺโฐ จ,}
\csromangathagroup{\csromangathastanza{\csromangathabat{สกฺกจฺจํ ปตฺตสญฺญี จ,}{สมสูปํ สมติตฺติกํ1.} \\ \csromangathabat{สกฺกจฺจํ ปตฺตสญฺญี จ,}{สปทานํ สมสูปกํ.}}\csromangathastanza{\csromangathabat{ถูปกโต ปฏิจฺฉนฺนํ,}{วิญฺญตฺตุชฺฌานสญฺญินา.} \\ \csromangathabat{น มหนฺตํ มณฺฑลํ ทฺวารํ,}{สพฺพํ หตฺถํ น พฺยาหเร.}}\csromangathastanza{\csromangathabat{อุกฺเขโป เฉทนา คณฺโฑ,}{ธุนํ สิตฺถาวการกํ.} \\ \csromangathabat{ชิวฺหานิจฺฉารกญฺเจว,}{จปุจปุ สุรุสุรุ.} \\ \csromangathabat{หตฺโถ ปตฺโต จ โอฏฺโฐ จ,}{สามิสํ สิตฺถเกน จ.}}}

\vspace{\tipitakaparskip}
\csromancenter{สมติตฺถิกํ (ก)}
\setlength{\csromangathapagewidth}{0pt}
\setlength{\csromangathawakwidth}{0pt}
\csromangathameasuregroup{ฉตฺตปาณิสฺส สทฺธมฺมํ, \\ เอวเมว ทณฺฑปาณิสฺส, \\ ปาทุกา อุปาหนา เจว, \\ ปลฺลตฺถิกา นิสินฺนสฺส, \\ ฉมา นีจาสเน ฐาเน, \\ ฐิตเกน น กาตพฺพํ,}{\csromangathabat{ฉตฺตปาณิสฺส สทฺธมฺมํ,}{น เทเสนฺติ ตถาคตา.} \\ \csromangathabat{เอวเมว ทณฺฑปาณิสฺส,}{สตฺถอาวุธปาณินํ.} \\ \csromangathabat{ปาทุกา อุปาหนา เจว,}{ยานเสยฺยาคตสฺส จ.} \\ \csromangathabat{ปลฺลตฺถิกา นิสินฺนสฺส,}{เวฐิโตคุณฺฐิตสฺส จ.} \\ \csromangathabat{ฉมา นีจาสเน ฐาเน,}{ปจฺฉโต อุปฺปเถน จ.} \\ \csromangathabat{ฐิตเกน น กาตพฺพํ,}{หริเต อุทกมฺหิ จาติ.}}
\csromangathasetleft{ฉตฺตปาณิสฺส สทฺธมฺมํ, \\ เอวเมว ทณฺฑปาณิสฺส, \\ ปาทุกา อุปาหนา เจว, \\ ปลฺลตฺถิกา นิสินฺนสฺส, \\ ฉมา นีจาสเน ฐาเน, \\ ฐิตเกน น กาตพฺพํ,}
\csromangathagroup{\csromangathastanza{\csromanfolio{55}\csromangathabat{ฉตฺตปาณิสฺส สทฺธมฺมํ,}{น เทเสนฺติ ตถาคตา.} \\ \csromangathabat{เอวเมว ทณฺฑปาณิสฺส,}{สตฺถอาวุธปาณินํ.}}\csromangathastanza{\csromangathabat{ปาทุกา อุปาหนา เจว,}{ยานเสยฺยาคตสฺส จ.} \\ \csromangathabat{ปลฺลตฺถิกา นิสินฺนสฺส,}{เวฐิโตคุณฺฐิตสฺส จ.}}\csromangathastanza{\csromangathabat{ฉมา นีจาสเน ฐาเน,}{ปจฺฉโต อุปฺปเถน จ.} \\ \csromangathabat{ฐิตเกน น กาตพฺพํ,}{หริเต อุทกมฺหิ จาติ.}}}

\vspace{\tipitakaparskip}
\csromancenter{เตสํ วคฺคานมุ\-ทฺทานํ.}
\setlength{\csromangathapagewidth}{0pt}
\setlength{\csromangathawakwidth}{0pt}
\csromangathameasuregroup{ปริมณฺฑลอุชฺชคฺฆิ, \\ กพฬา สุรุสุรุ จ,}{\csromangathabat{ปริมณฺฑลอุชฺชคฺฆิ,}{ขมฺภํ ปิณฺฑํ ตเถว จ.} \\ \csromangathabat{กพฬา สุรุสุรุ จ,}{ปาทุเกน จ สตฺตมาติ.}}
\csromangathasetleft{ปริมณฺฑลอุชฺชคฺฆิ, \\ กพฬา สุรุสุรุ จ,}
\csromangathagroupwithcloserruled{\csromangathastanza{\csromangathabat{ปริมณฺฑลอุชฺชคฺฆิ,}{ขมฺภํ ปิณฺฑํ ตเถว จ.} \\ \csromangathabat{กพฬา สุรุสุรุ จ,}{ปาทุเกน จ สตฺตมาติ.}}}{มหาวิภงฺเค กตฺถ\-ปญฺญตฺติวาโร นิฏฺฐิโต.}
\csromanmatikaanchor{mtk.29}
\csromanmatikaanchor{mtk.30}
\csromantocmarknopage{chapter}{2. กตาปตฺติวาร}{mtk.29}
\bookmark[dest={mtk.29},level=0]{2. กตาปตฺติวาร}
\csromantocmarkref{chapter}{1. ปาราชิกกณฺฑ}{mtk.30}
\bookmark[dest={mtk.30},level=0]{1. ปาราชิกกณฺฑ}
\chapterhead{2. \textbf{กตาปตฺติวาร 1. ปาราชิกกณฺฑ}}
\proseitem{157}{เมถุนํ ธมฺมํ ปฏิเสวนฺโต กติ อาปตฺติโย อาปชฺชติ. เมถุนํ ธมฺมํ ปฏิเสวนฺโต ติสฺโส อาปตฺติโย อาปชฺชติ. อกฺขายิเต สรีเร เมถุนํ ธมฺมํ ปฏิเสวติ, อาปตฺติ ปาราชิกสฺส. เยภุยฺเยน ขายิเต สรีเร เมถุนํ ธมฺมํ ปฏิเสวติ, อาปตฺติ ถุลฺลจฺจยสฺส. วฏฺฏกเต\csromansharedfootnote{f30706a5c7c4ec25}{วิวฏกเต (สฺยา)} มุเข อจฺฉุปนฺตํ องฺคชาตํ ปเวเสติ, อาปตฺติ ทุกฺกฏสฺส. เมถุนํ ธมฺมํ ปฏิเสวนฺโต อิมา ติสฺโส อาปตฺติโย อาปชฺชติ.}

\proseitem{158}{อทินฺนํ อาทิยนฺโต กติ อาปตฺติโย อาปชฺชติ. อทินฺนํ อาทิยนฺโต ติสฺโส อาปตฺติโย อาปชฺชติ. ปญฺจมาสกํ วา อติเรก\-ปญฺจ\-มาสกํ วา อคฺฆนกํ อทินฺนํ เถยฺยสงฺขาตํ อาทิยติ, อาปตฺติ ปาราชิกสฺส. อติเรกมาสกํ วา อูนปญฺจมาสกํ วา อคฺฆนกํ อทินฺนํ เถยฺยสงฺขาตํ อาทิยติ, อาปตฺติ ถุลฺลจฺจยสฺส. มาสกํ วา อูนมาสกํ วา อคฺฆนกํ อทินฺนํ เถยฺยสงฺขาตํ อาทิยติ, อาปตฺติ ทุกฺกฏสฺส. อทินฺนํ อาทิยนฺโต อิมา ติสฺโส อาปตฺติโย อาปชฺชติ.}

\proseitem{159}{\csromanfolio{56}สญฺจิจฺจ มนุสฺส\-วิคฺคหํ ชีวิตา โวโรเปนฺโต กติ อาปตฺติโย อาปชฺชติ. สญฺจิจฺจ มนุสฺส\-วิคฺคหํ ชีวิตา โวโรเปนฺโต ติสฺโส อาปตฺติโย อาปชฺชติ. มนุสฺสํ โอทิสฺส โอปาตํ ขณติ “ปปติตฺวา มริสฺสตี”ติ, อาปตฺติ ทุกฺกฏสฺส. ปปติเต ทุกฺขา เวทนา อุปฺปชฺชติ, อาปตฺติ ถุลฺลจฺจยสฺส. มรติ, อาปตฺติ ปาราชิกสฺส. สญฺจิจฺจ มนุสฺส\-วิคฺคหํ ชีวิตา โวโรเปนฺโต อิมา ติสฺโส อาปตฺติโย อาปชฺชติ.}

\proseitemwithcloserruled{160}{อสนฺตํ อภูตํ อุตฺตริ\-มนุสฺสธมฺมํ อุลฺลปนฺโต กติ อาปตฺติโย อาปชฺชติ. อสนฺตํ อภูตํ อุตฺตริ\-มนุสฺสธมฺมํ อุลฺลปนฺโต ติสฺโส อาปตฺติโย อาปชฺชติ. ปาปิจฺโฉ อิจฺฉาปกโต อสนฺตํ อภูตํ อุตฺตริ\-มนุสฺสธมฺมํ อุลฺลปติ, อาปตฺติ ปาราชิกสฺส. “โยเต วิหาเรวสติ, โส ภิกฺขุ อรหา”ติ ภณติ, ปฏิวิชานนฺตสฺส อาปตฺติ ถุลฺลจฺจยสฺส. น ปฏิวิชานนฺตสฺส อาปตฺติ ทุกฺกฏสฺส. อสนฺตํ อภูตํ อุตฺตริ\-มนุสฺสธมฺมํ อุลฺลปนฺโต อิมา ติสฺโส อาปตฺติโย อาปชฺชติ.}{จตฺตาโร ปาราชิกา นิฏฺฐิตา.}
\csromanmatikaanchor{mtk.31}
\csromantocmarkref{chapter}{2. สํฆาทิเสสกณฺฑ}{mtk.31}
\bookmark[dest={mtk.31},level=0]{2. สํฆาทิเสสกณฺฑ}
\chapterhead{2. สํฆาทิเสส\-กณฺฑ}
\proseitem{161}{อุปกฺกมิตฺวา อสุจิํ โมเจนฺโต ติสฺโส อาปตฺติโย อาปชฺชติ. เจเตติ อุปกฺกมติ มุจฺจติ, อาปตฺติ สํฆาทิเสสสฺส. เจเตติ อุปกฺกมติ น มุจฺจติ, อาปตฺติ ถุลฺลจฺจยสฺส. ปโยเค ทุกฺกฏํ.}

\prose{มาตุคาเมน สทฺธิํ กายสํสคฺคํ สมาปชฺชนฺโต ติสฺโส อาปตฺติโย อาปชฺชติ. กาเยน กายํ อามสติ, อาปตฺติ สํฆาทิเสสสฺส. กาเยน กาย\-ปฏิพทฺธํ อามสติ, อาปตฺติ ถุลฺลจฺจยสฺส. กาย\-ปฏิพทฺเธน กาย\-ปฏิพทฺธํ อามสติ, อาปตฺติ ทุกฺกฏสฺส.}

\prose{มาตุคามํ ทุฏฺฐุลฺลาหิ วาจาหิ โอภาเสนฺโต ติสฺโส อาปตฺติโย อาปชฺชติ. วจฺจมคฺคํ ปสฺสาวมคฺคํ อาทิสฺส วณฺณมฺปิ ภณติ, อวณฺณมฺปิ ภณติ, อาปตฺติ สํฆาทิเสสสฺส. วจฺจมคฺคํ ปสฺสาวมคฺคํ ฐเปตฺวา อธกฺขกํ อุพฺภ\-ชาณุมณฺฑลํ อาทิสฺส วณฺณมฺปิ ภณติ, อวณฺณมฺปิ ภณติ, อาปตฺติ ถุลฺลจฺจยสฺส. กาย\-ปฏิพทฺธํ อาทิสฺส วณฺณมฺปิ ภณติ, อวณฺณมฺปิ ภณติ, อาปตฺติ ทุกฺกฏสฺส.}

\prose{\csromanfolio{57}อตฺต\-กาม\-ปาริจริยาย วณฺณํ ภาสนฺโต ติสฺโส อาปตฺติโย อาปชฺชติ. มาตุคามสฺส สนฺติเก อตฺต\-กาม\-ปาริจริยาย วณฺณํ ภาสติ, อาปตฺติ สํฆาทิเสสสฺส. ปณฺฑกสฺส สนฺติเก อตฺต\-กาม\-ปาริจริยาย วณฺณํ ภาสติ, อาปตฺติ ถุลฺลจฺจยสฺส. ติรจฺฉาน\-คตสฺส สนฺติเก อตฺต\-กาม\-ปาริจริยาย วณฺณํ ภาสติ, อาปตฺติ ทุกฺกฏสฺส.}

\prose{สญฺจริตฺตํ สมาปชฺชนฺโต ติสฺโส อาปตฺติโย อาปชฺชติ. ปฏิคฺคณฺหาติ วีมํสติ ปจฺจาหรติ, อาปตฺติ สํฆาทิเสสสฺส. ปฏิคฺคณฺหาติ วีมํสติ น ปจฺจาหรติ, อาปตฺติ ถุลฺลจฺจยสฺส. ปฏิคฺคณฺหาติ น วีมํสติ น ปจฺจาหรติ, อาปตฺติ ทุกฺกฏสฺส.}

\prose{สญฺญาจิกาย กุฏิํ การาเปนฺโต ติสฺโส อาปตฺติโย อาปชฺชติ. การาเปติ, ปโยเค ทุกฺกฏํ. เอกํ ปิณฺฑํ อนาคเต อาปตฺติ ถุลฺลจฺจยสฺส. ตสฺมิํ ปิณฺเฑ อาคเต อาปตฺติ สํฆาทิเสสสฺส.}

\prose{มหลฺลกํ วิหารํ การาเปนฺโต ติสฺโส อาปตฺติโย อาปชฺชติ. การาเปติ, ปโยเค ทุกฺกฏํ. เอกํ ปิณฺฑํ อนาคเต อาปตฺติ ถุลฺลจฺจยสฺส. ตสฺมิํ ปิณฺเฑ อาคเต อาปตฺติ สํฆาทิเสสสฺส.}

\prose{ภิกฺขุํ อมูลเกน ปาราชิเกน ธมฺเมน อนุทฺธํเสนฺโต ติสฺโส อาปตฺติโย อาปชฺชติ. อโนกาสํ การาเปตฺวา จาวนาธิปฺปาโย วเทติ, อาปตฺติ สํฆาทิเสเสน ทุกฺกฏสฺส. โอกาสํ การาเปตฺวา อกฺโกสาธิปฺปาโย วเทติ, อาปตฺติ โอมสวาทสฺส.}

\prose{ภิกฺขุํ อญฺญภาคิยสฺส อธิกรณสฺส กิญฺจิ เทสํ เลสมตฺตํ อุปาทาย ปาราชิเกน ธมฺเมน อนุทฺธํเสนฺโต ติสฺโส อาปตฺติโย อาปชฺชติ. อโนกาสํ การาเปตฺวา จาวนาธิปฺปาโย วเทติ, อาปตฺติ สํฆาทิเสเสน ทุกฺกฏสฺส. โอกาสํ การาเปตฺวา อกฺโกสาธิปฺปาโย วเทติ, อาปตฺติ โอมสวาทสฺส.}

\prose{สํฆเภทโก ภิกฺขุ ยาวตติยํ สมนุภาสนาย น ปฏินิสฺสชฺชนฺโต ติสฺโส อาปตฺติโย อาปชฺชติ. ญตฺติยา ทุกฺกฏํ. ทฺวีหิ กมฺมวาจาหิ ถุลฺลจฺจยา. กมฺมวาจา\-ปริโยสาเน อาปตฺติ สํฆาทิเสเสสสฺส.}

\prose{เภทกา\-นุวตฺตกา ภิกฺขู ยาวตติยํ สมนุภาสนาย\csromansharedfootnote{b69b357f89e05230}{สมนุภาสิยมานา (สฺยา)} น ปฏินิสฺสชฺชนฺตา ติสฺโส อาปตฺติโย อาปชฺชนฺติ. ญตฺติยา ทุกฺกฏํ. ทฺวีหิ \csromanfolio{58}กมฺมวาจาหิ ถุลฺลจฺจยา. กมฺมวาจา\-ปริโยสาเน อาปตฺติ สํฆาทิเสสสฺส.}

\prose{ทุพฺพโจ ภิกฺขุ ยาวตติยํ สมนุภาสนาย\csromansharedfootnote{c18cbff5c753a546}{สมนุภาสิยมาโน (สฺยา)} น ปฏินิสฺสชฺชนฺโต ติสฺโส อาปตฺติโย อาปชฺชติ. ญตฺติยา ทุกฺกฏํ. ทฺวีหิ กมฺมวาจาหิ ถุลฺลจฺจยา. กมฺมวาจา\-ปริโยสาเน อาปตฺติ สํฆาทิเสสสฺส.}

\prosewithcloserruled{กุลทูสโก ภิกฺขุ ยาวตติยํ สมนุภาสนาย น ปฏินิสฺสชฺชนฺโต ติสฺโส อาปตฺติโย อาปชฺชติ. ญตฺติยา ทุกฺกฏํ. ทฺวีหิ กมฺมวาจาหิ ถุลฺลจฺจยา. กมฺมวาจา\-ปริโยสาเน อาปตฺติ สํฆาทิเสสสฺส.}{เตรส สํฆาทิเสสา นิฏฺฐิตา.}
\csromanmatikaanchor{mtk.32}
\csromanmatikaanchor{mtk.33}
\csromantocmarkref{chapter}{4. นิสฺสคฺคิยกณฺฑ}{mtk.32}
\bookmark[dest={mtk.32},level=0]{4. นิสฺสคฺคิยกณฺฑ}
\csromantocmarkref{section}{1. กถินวคฺค}{mtk.33}
\bookmark[dest={mtk.33},level=1]{1. กถินวคฺค}
\csromanmarkh{4. นิสฺสคฺคิย\-กณฺฑ}\csromanmarkhii{1. กถินวคฺค}\csromanheadernomark{1}{4. นิสฺสคฺคิย\-กณฺฑ 1. กถินวคฺค}
\proseitem{162}{อติเรกจีวรํ ทสาหํ อติกฺกาเมนฺโต เอกํ อาปตฺติํ อาปชฺชติ, นิสฺสคฺคิยํ ปาจิตฺติยํ.}

\prose{เอกรตฺตํ ติจีวเรน วิปฺปวสนฺโต เอกํ อาปตฺติํ อาปชฺชติ, นิสฺสคฺคิยํ ปาจิตฺติยํ.}

\prose{อกาลจีวรํ ปฏิคฺคเหตฺวา มาสํ อติกฺกาเมนฺโต เอกํ อาปตฺติํ อาปชฺชติ, นิสฺสคฺคิยํ ปาจิตฺติยํ.}

\prose{อญฺญาติกาย ภิกฺขุนิยา ปุราณจีวรํ โธวาเปนฺโต เทฺว อาปตฺติโย อาปชฺชติ. โธวาเปติ, ปโยเค ทุกฺกฏํ. โธวาปิเต นิสฺสคฺคิยํ ปาจิตฺติยํ.}

\prose{อญฺญาติกาย ภิกฺขุนิยา หตฺถโต จีวรํ ปฏิคฺคณฺหนฺโต เทฺว อาปตฺติโย อาปชฺชติ. คณฺหาติ, ปโยเค ทุกฺกฏํ. คหิเต นิสฺสคฺคิยํ ปาจิตฺติยํ.}

\prose{อญฺญาตกํ คหปติํ วา คหปตานิํ วา จีวรํ วิญฺญาเปนฺโต เทฺว อาปตฺติโย อาปชฺชติ. วิญฺญาเปติ, ปโยเค ทุกฺกฏํ. วิญฺญาปิเต นิสฺสคฺคิยํ ปาจิตฺติยํ.}

\prose{\csromanfolio{59}อญฺญาตกํ คหปติํ วา คหปตานิํ วา ตตุตฺตริ จิวรํ วิญฺญาเปนฺโต เทฺว อาปตฺติโย อาปชฺชติ. วิญฺญาเปติ, ปโยเค ทุกฺกฏํ. วิญฺญาปิเต นิสฺสคฺคิยํ ปาจิตฺติยํ.}

\prose{ปุพฺเพ อปฺปวาริโต อญฺญาตกํ คหปติกํ อุปสงฺกมิตฺวา จีวเร วิกปฺปํ อาปชฺชนฺโต เทฺว อาปตฺติโย อาปชฺชติ. วิกปฺปํ อาปชฺชติ, ปโยเค ทุกฺกฏํ. วิกปฺปํ อาปนฺเน นิสฺสคฺคิยํ ปาจิตฺติยํ.}

\prose{ปุพฺเพ อปฺปวาริโต อญฺญาตเก คหปติเก อุปสงฺกมิตฺวา จีวเร วิกปฺปํ อาปชฺชนฺโต เทฺว อาปตฺติโย อาปชฺชติ. วิกปฺปํ อาปชฺชติ, ปโยเค ทุกฺกฏํ. วิกปฺปํ อาปนฺเน นิสฺสคฺคิยํ ปาจิตฺติยํ.}

\prosewithclosermidruled{อติเรก\-ติ\-กฺขตฺตุํ โจทนาย อติเรก\-ฉ\-กฺขตฺตุํ ฐาเนน จีวรํ อภินิ\-ปฺผาเท\-นฺโต เทฺว อาปตฺติโย อาปชฺชติ. อภินิปฺผาเทติ, ปโยเค ทุกฺกฏํ. อภินิปฺผาทิเต นิสฺสคฺคิยํ ปาจิตฺติยํ.}{กถินวคฺโค ปฐโม.}
\csromanmatikaanchor{mtk.34}
\csromantocmarkref{section}{2. โกสิยวคฺค}{mtk.34}
\bookmark[dest={mtk.34},level=1]{2. โกสิยวคฺค}
\csromanheader{1}{2. \textbf{โกสิยวคฺค}}
\proseitem{163}{โกสิยมิสฺสกํ สนฺถตํ การาเปนฺโต เทฺว อาปตฺติโย อาปชฺชติ. การาเปติ, ปโยเค ทุกฺกฏํ. การาปิเต นิสฺสคฺคิยํ ปาจิตฺติยํ.}

\prose{สุทฺธ\-กาฬกานํ เอฬกโลมานํ สนฺถตํ การาเปนฺโต เทฺว อาปตฺติโย อาปชฺชติ. การาเปติ, ปโยเค ทุกฺกฏํ. การาปิเต นิสฺสคฺคิยํ ปาจิตฺติยํ.}

\prose{อนาทิยิตฺวา ตุลํ โอทาตานํ ตุลํ โคจริยานํ นวํ สนฺถตํ การาเปนฺโต เทฺว อาปตฺติโย อาปชฺชติ. การาเปติ, ปโยเค ทุกฺกฏํ. การาปิเต นิสฺสคฺคิยํ ปาจิตฺติยํ.}

\prose{อนุวสฺสํ สนฺถตํ การาเปนฺโต เทฺว อาปตฺติโย อาปชฺชติ. การาเปติ, ปโยเค ทุกฺกฏํ. การาปิเต นิสฺสคฺคิยํ ปาจิตฺติยํ.}

\prose{อนาทิยิตฺวา ปุราณ\-สนฺถตสฺส สามนฺตา สุคต\-วิทตฺถิํ นวํ นิสีทน\-สนฺถตํ การาเปนฺโต เทฺว อาปตฺติโย อาปชฺชติ. การาเปติ ปโยเค ทุกฺกฏํ. การาปิเต นิสฺสคฺคิยํ ปาจิตฺติยํ.}

\prose{\csromanfolio{60}เอฬกโลมานิ ปฏิคฺคเหตฺวา ติโยชนํ อติกฺกาเมนฺโต เทฺว อาปตฺติโย อาปชฺชติ. ปฐมํ ปาทํ ติโยชนํ อติกฺกาเมติ, อาปตฺติ ทุกฺกฏสฺส. ทุติยํ ปาทํ อติกฺกาเมติ, นิสฺสคฺคิยํ ปาจิตฺติยํ.}

\prose{อญฺญาติกาย ภิกฺขุนิยา เอฬกโลมานิ โธวาเปนฺโต เทฺว อาปตฺติโย อาปชฺชติ. โธวาเปติ, ปโยเค ทุกฺกฏํ. โธวาปิเต นิสฺสคฺคิยํ ปาจิตฺติยํ.}

\prose{รูปิยํ ปฏิคฺคณฺหนฺโต เทฺว อาปตฺติโย อาปชฺชติ. คณฺหาติ, ปโยเค ทุกฺกฏํ. คหิเต นิสฺสคฺคิยํ ปาจิตฺติยํ.}

\prose{นานปฺปการกํ รูปิย\-สํโวหารํ สมาปชฺชนฺโต เทฺว อาปตฺติโย อาปชฺชติ. สมาปชฺชติ, ปโยเค ทุกฺกฏํ. สมาปนฺเน นิสฺสคฺคิยํ ปาจิตฺติยํ.}

\prosewithclosermidruled{นานปฺปการกํ กยวิกฺกยํ สมาปชฺชนฺโต เทฺว อาปตฺติโย อาปชฺชติ. สมาปชฺชติ, ปโยเค ทุกฺกฏํ. สมาปนฺเน นิสฺสคฺคิยํ ปาจิตฺติยํ.}{โกสิยวคฺโค ทุติโย.}
\csromanmatikaanchor{mtk.35}
\csromantocmarkref{section}{3. ปตฺตวคฺค}{mtk.35}
\bookmark[dest={mtk.35},level=1]{3. ปตฺตวคฺค}
\csromanheader{1}{3. \textbf{ปตฺตวคฺค}}
\proseitem{164}{อติเรกปตฺตํ ทสาหํ อติกฺกาเมนฺโต เอกํ อาปตฺติํ อาปชฺชติ, นิสฺสคฺคิยํ ปาจิตฺติยํ.}

\prose{อูนปญฺจ\-พนฺธเนน ปตฺเตน อญฺญํ นวํ ปตฺตํ เจตาเปนฺโต เทฺว อาปตฺติโย อาปชฺชติ. เจตาเปติ, ปโยเค ทุกฺกฏํ. เจตาปิเต นิสฺสคฺคิยํ ปาจิตฺติยํ.}

\prose{เภสชฺชานิ ปฏิคฺคเหตฺวา สตฺตาหํ อติกฺกาเมนฺโต เอกํ อาปตฺติํ อาปชฺชติ, นิสฺสคฺคิยํ ปาจิตฺติยํ.}

\prose{อติเรกมาเส เสเส คิมฺหาเน วสฺสิก\-สาฏิก\-จีวรํ ปริเยสนฺโต เทฺว อาปตฺติโย อาปชฺชติ. ปริเยสติ, ปโยเค ทุกฺกฏํ. ปริยิฏฺเฐ นิสฺสคฺคิยํ ปาจิตฺติยํ.}

\prose{ภิกฺขุสฺส สามํ จีวรํ ทตฺวา กุปิโต อนตฺตมโน อจฺฉินฺทนฺโต เทฺว อาปตฺติโย อาปชฺชติ. อจฺฉินฺทติ, ปโยเค ทุกฺกฏํ. อจฺฉินฺเน นิสฺสคฺคิยํ ปาจิตฺติยํ.}

\prose{\csromanfolio{61}สามํ สุตฺตํ วิญฺญาเปตฺวา ตนฺตวาเยหิ จีวรํ วายาเปนฺโต เทฺว อาปตฺติโย อาปชฺชติ. วายาเปติ, ปโยเค ทุกฺกฏํ. วายาปิเต นิสฺสคฺคิยํ ปาจิตฺติยํ.}

\prose{ปุพฺเพ อปฺปวาริโต อญฺญาตกสฺส คหปติกสฺส ตนฺตวาเย อุปสงฺกมิตฺวา จีวเร วิกปฺปํ อาปชฺชนฺโต เทฺว อาปตฺติโย อาปชฺชติ. วิกปฺปํ อาปชฺชติ, ปโยเค ทุกฺกฏํ. วิกปฺปํ อาปนฺเน นิสฺสคฺคิยํ ปาจิตฺติยํ.}

\prose{อจฺเจกจีวรํ ปฏิคฺคเหตฺวา จีวร\-กาล\-สมยํ อติกฺกาเมนฺโต เอกํ อาปตฺติํ อาปชฺชติ, นิสฺสคฺคิยํ ปาจิตฺติยํ.}

\prose{ติณฺณํ จีวรานํ อญฺญตรํ จีวรํ อนฺตรฆเร นิกฺขิปิตฺวา อติเรกฉารตฺตํ วิปฺปวสนฺโต เอกํ อาปตฺติํ อาปชฺชติ, นิสฺสคฺคิยํ ปาจิตฺติยํ.}

\prosewithcloserruled{ชานํ สํฆิกํ ลาภํ ปริณตํ อตฺตโน ปริณาเมนฺโต เทฺว อาปตฺติโย อาปชฺชติ. ปริณาเมติ, ปโยเค ทุกฺกฏํ. ปริณามิเต นิสฺสคฺคิยํ ปาจิตฺติยํ.}{ปตฺตวคฺโค ตติโย. ติํส นิสฺสคฺคิยา ปาจิตฺติยา นิฏฺฐิตา.}
\csromanmatikaanchor{mtk.36}
\csromantocmarkref{chapter}{5. ปาจิตฺติยกณฺฑ}{mtk.36}
\bookmark[dest={mtk.36},level=0]{5. ปาจิตฺติยกณฺฑ}
\chapterhead{5. \textbf{ปาจิตฺติยกณฺฑ}}
\csromanmatikaanchor{mtk.37}
\csromantocmarkref{section}{1. มุสาวาทวคฺค}{mtk.37}
\bookmark[dest={mtk.37},level=1]{1. มุสาวาทวคฺค}
\csromanheader{1}{1. \textbf{มุสาวาทวคฺค}}
\proseitem{165}{สมฺปชาน\-มุสาวาทํ ภาสนฺโต กติ อาปตฺติโย อาปชฺชติ. สมฺปชาน\-มุสาวาทํ ภาสนฺโต ปญฺจ อาปตฺติโย อาปชฺชติ. ปาปิจฺโฉ อิจฺฉาปกโต อสนฺตํ อภูตํ อุตฺตริ\-มนุสฺสธมฺมํ อุลฺลปติ, อาปตฺติ ปาราชิกสฺส. ภิกฺขุํ อมูลเกน ปาราชิเกน ธมฺเมน อนุทฺธํเสติ, อาปตฺติ สํฆาทิเสสสฺส. “โย เต วิหาเร วสติ, โส ภิกฺขุ อรหา”ติ ภณติ, ปฏิวิชานนฺตสฺส อาปตฺติ ถุลฺลจฺจยสฺส. น ปฏิวิชานนฺตสฺส อาปตฺติ ทุกฺกฏสฺส. สมฺปชาน\-มุสาวาเท ปาจิตฺติยํ. สมฺปชาน\-มุสาวาทํ ภาสนฺโต อิมา ปญฺจ อาปตฺติโย อาปชฺชติ.}

\prose{\csromanfolio{62}โอมสนฺโต เทฺว อาปตฺติโย อาปชฺชติ. อุปสมฺปนฺนํ โอมสติ, อาปตฺติ ปาจิตฺติยสฺส. อนุปสมฺปนฺนํ โอมสติ, อาปตฺติ ทุกฺกฏสฺส.}

\prose{เปสุญฺญํ อุปสํหรนฺโต เทฺว อาปตฺติโย อาปชฺชติ. อุปสมฺปนฺนสฺส เปสุญฺญํ อุปสํหรติ, อาปตฺติ ปาจิตฺติยสฺส. อนุปสมฺปนฺนสฺส เปสุญฺญํ อุปสํหรติ, อาปตฺติ ทุกฺกฏสฺส.}

\prose{อนุปสมฺปนฺนํ ปทโส ธมฺมํ วาเจนฺโต เทฺว อาปตฺติโย อาปชฺชติ. วาเจติ, ปโยเค ทุกฺกฏํ. ปเท ปเท อาปตฺติ ปาจิตฺติยสฺส.}

\prose{อนุปสมฺปนฺเนน อุตฺตริทิรตฺตติรตฺตํ สหเสยฺยํ กปฺเปนฺโต เทฺว อาปตฺติโย อาปชฺชติ. นิปชฺชติ, ปโยเค ทุกฺกฏํ. นิปนฺเน อาปตฺติ ปาจิตฺติยสฺส.}

\prose{มาตุคาเมน สหเสยฺยํ กปฺเปนฺโต เทฺว อาปตฺติโย อาปชฺชติ. นิปชฺชติ, ปโยเค ทุกฺกฏํ. นิปนฺเน อาปตฺติ ปาจิตฺติยสฺส.}

\prose{มาตุคามสฺส อุตฺตริ\-ฉปฺปญฺจ\-วาจาหิ ธมฺมํ เทเสนฺโต เทฺว อาปตฺติโย อาปชฺชติ. เทเสติ, ปโยเค ทุกฺกฏํ. ปเท ปเท อาปตฺติ ปาจิตฺติยสฺส.}

\prose{อนุปสมฺปนฺนสฺส อุตฺตริ\-มนุสฺสธมฺมํ ภูตํ อาโรเจนฺโต เทฺว อาปตฺติโย อาปชฺชติ. อาโรเจติ, ปโยเค ทุกฺกฏํ. อาโรจิเต อาปตฺติ ปาจิตฺติยสฺส.}

\prose{ภิกฺขุสฺส ทุฏฺฐุลฺลํ อาปตฺติํ อนุปสมฺปนฺนสฺส อาโรเจนฺโต เทฺว อาปตฺติโย อาปชฺชติ. อาโรเจติ, ปโยเค ทุกฺกฏํ. อาโรจิเต อาปตฺติ ปาจิตฺติยสฺส.}

\prosewithclosermidruled{ปถวิํ ขณนฺโต เทฺว อาปตฺติโย อาปชฺชติ. ขณติ, ปโยเค ทุกฺกฏํ. ปหาเร ปหาเร อาปตฺติ ปาจิตฺติยสฺส.}{มุสาวาทวคฺโค ปฐโม.}
\csromanmatikaanchor{mtk.38}
\csromantocmarkref{section}{2. ภูตคามวคฺค}{mtk.38}
\bookmark[dest={mtk.38},level=1]{2. ภูตคามวคฺค}
\csromanheader{1}{2. \textbf{ภูตคามวคฺค}}
\proseitem{166}{ภูตคามํ ปาเตนฺโต เทฺว อาปตฺติโย อาปชฺชติ. ปาเตติ ปโยเค ทุกฺกฏํ. ปหาเร ปหาเร อาปตฺติ ปาจิตฺติยสฺส.}

\prose{\csromanfolio{63}อญฺเญนญฺญํ ปฏิจรนฺโต เทฺว อาปตฺติโย อาปชฺชติ. อนาโรปิเต อญฺญวาทเก อญฺเญนญฺญํ ปฏิจรติ, อาปตฺติ ทุกฺกฏสฺส. อาโรปิเต อญฺญวาทเก อญฺเญนญฺญํ ปฏิจรติ, อาปตฺติ ปาจิตฺติยสฺส.}

\prose{ภิกฺขุํ อุชฺฌาเปนฺโต เทฺว อาปตฺติโย อาปชฺชติ. อุชฺฌาเปติ, ปโยเค ทุกฺกฏํ. อุชฺฌาปิเต อาปตฺติ ปาจิตฺติยสฺส.}

\prose{สํฆิกํ มญฺจํ วา ปีฐํ วา ภิสิํ วา โกจฺฉํ วา อชฺโฌกาเส สนฺถริตฺวา อนุทฺธริตฺวา อนาปุจฺฉา ปกฺกมนฺโต เทฺว อาปตฺติโย อาปชฺชติ. ปฐมํ ปาทํ เลฑฺฑุปาตํ อติกฺกาเมติ, อาปตฺติ ทุกฺกฏสฺส. ทุติยํ ปาทํ อติกฺกาเมติ, อาปตฺติ ปาจิตฺติยสฺส.}

\prose{สํฆิเก วิหาเร เสยฺยํ สนฺถริตฺวา อนุทฺธริตฺวา อนาปุจฺฉา ปกฺกมนฺโต เทฺว อาปตฺติโย อาปชฺชติ. ปฐมํ ปาทํ ปริกฺเขปํ อติกฺกาเมติ, อาปตฺติ ทุกฺกฏสฺส. ทุติยํ ปาทํ อติกฺกาเมติ, อาปตฺติ ปาจิตฺติยสฺส.}

\prose{สํฆิเก วิหาเร ชานํ ปุพฺพุปคตํ ภิกฺขุํ อนุปขชฺช เสยฺยํ กปฺเปนฺโต เทฺว อาปตฺติโย อาปชฺชติ. นิปชฺชติ, ปโยเค ทุกฺกฏํ. นิปนฺเน อาปตฺติ ปาจิตฺติยสฺส.}

\prose{ภิกฺขุํ กุปิโต อนตฺตมโน สํฆิกา วิหารา นิกฺกฑฺเฒนฺโต เทฺว อาปตฺติโย อาปชฺชติ. นิกฺกฑฺฒติ, ปโยเค ทุกฺกฏํ. นิกฺกฑฺฒิเต อาปตฺติ ปาจิตฺติยสฺส.}

\prose{สํฆิเก วิหาเร อุปริ\-เวหาส\-กุฏิยา อาหจฺจปาทกํ มญฺจํ วา ปีฐํ วา อภินิสีทนฺโต เทฺว อาปตฺติโย อาปชฺชติ. อภินิสีทติ, ปโยเค ทุกฺกฏํ. อภินิสินฺเน อาปตฺติ ปาจิตฺติยสฺส.}

\prose{ทฺวตฺติปริยาเย อธิฏฺฐหิตฺวา ตตุตฺตริ อธิฏฺฐหนฺโต เทฺว อาปตฺติโย อาปชฺชติ. อธิฏฺเฐติ, ปโยเค ทุกฺกฏํ. อธิฏฺฐิเต อาปตฺติ ปาจิตฺติยสฺส.}

\prosewithclosermidruled{ชานํ สปฺปาณกํ อุทกํ ติณํ วา มตฺติกํ วา สิญฺจนฺโต เทฺว อาปตฺติโย อาปชฺชติ. สิญฺจติ, ปโยเค ทุกฺกฏํ. สิญฺจิเต อาปตฺติ ปาจิตฺติยสฺส.}{ภูตคามวคฺโค ทุติโย.}
\csromanmatikaanchor{mtk.39}
\csromantocmarkref{section}{3. โอวาทวคฺค}{mtk.39}
\bookmark[dest={mtk.39},level=1]{3. โอวาทวคฺค}
\csromanheader{1}{3. \textbf{โอวาทวคฺค}}
\proseitem{167}{\csromanfolio{64}อสมฺมโต ภิกฺขุนิโย โอวทนฺโต เทฺว อาปตฺติโย อาปชฺชติ. โอวทติ, ปโยเค ทุกฺกฏํ. โอวทิเต อาปตฺติ ปาจิตฺติยสฺส.}

\prose{อตฺถงฺคเต สูริเย ภิกฺขุนิโย โอวทนฺโต เทฺว อาปตฺติโย อาปชฺชติ. โอวทติ, ปโยเค ทุกฺกฏํ. โอวทิเต อาปตฺติ ปาจิตฺติยสฺส.}

\prose{ภิกฺขุนุ\-ปสฺสยํ อุปสงฺกมิตฺวา ภิกฺขุนิโย โอวทนฺโต เทฺว อาปตฺติโย อาปชฺชติ. โอวทติ, ปโยเค ทุกฺกฏํ. โอวทิเต อาปตฺติ ปาจิตฺติยสฺส.}

\prose{“อามิสเหตุ ภิกฺขู ภิกฺขุนิโย โอวทนฺตี”ติ ภณนฺโต เทฺว อาปตฺติโย อาปชฺชติ. ภณติ, ปโยเค ทุกฺกฏํ. ภณิเต อาปตฺติ ปาจิตฺติยสฺส.}

\prose{อญฺญาติกาย ภิกฺขุนิยา จีวรํ เทนฺโต เทฺว อาปตฺติโย อาปชฺชติ. เทติ, ปโยเค ทุกฺกฏํ. ทินฺเน อาปตฺติ ปาจิตฺติยสฺส.}

\prose{อญฺญาติกาย ภิกฺขุนิยา จีวรํ สิพฺเพนฺโต เทฺว อาปตฺติโย อาปชฺชติ. สิพฺเพติ, ปโยเค ทุกฺกฏํ. อาราปเถ อาราปเถ อาปตฺติ ปาจิตฺติยสฺส.}

\prose{ภิกฺขุนิยา สทฺธิํ สํวิธาย เอก\-ทฺธาน\-มคฺคํ ปฏิปชฺชนฺโต เทฺว อาปตฺติโย อาปชฺชติ. ปฏิปชฺชติ, ปโยเค ทุกฺกฏํ. ปฏิปนฺเน อาปตฺติ ปาจิตฺติยสฺส.}

\prose{ภิกฺขุนิยา สทฺธิํ สํวิธาย เอกํ นาวํ อภิรุหนฺโต เทฺว อาปตฺติโย อาปชฺชติ. อภิรุหติ, ปโยเค ทุกฺกฏํ. อภิรุเฬฺห อาปตฺติ ปาจิตฺติยสฺส.}

\prose{ชานํ ภิกฺขุนิ\-ปริปาจิตํ ปิณฺฑปาตํ ภุญฺชนฺโต เทฺว อาปตฺติโย อาปชฺชติ. “ภุญฺชิสฺสามี”ติ ปฏิคฺคณฺหาติ, อาปตฺติ ทุกฺกฏสฺส. อชฺโฌหาเร อชฺโฌหาเร อาปตฺติ ปาจิตฺติยสฺส.}

\prosewithclosermidruled{ภิกฺขุนิยา สทฺธิํ เอโก เอกาย รโห นิสชฺชํ กปฺเปนฺโต เทฺว อาปตฺติโย อาปชฺชติ. นิสีทติ, ปโยเค ทุกฺกฏํ. นิสินฺเน อาปตฺติ ปาจิตฺติยสฺส.}{โอวาทวคฺโค ตติโย.}
\csromanmatikaanchor{mtk.40}
\csromantocmarkref{section}{4. โภชนวคฺค}{mtk.40}
\bookmark[dest={mtk.40},level=1]{4. โภชนวคฺค}
\csromanheader{1}{4. \textbf{โภชนวคฺค}}
\proseitem{168}{\csromanfolio{65}ตตุตฺตริ อาวสถปิณฺฑํ ภุญฺชนฺโต เทฺว อาปตฺติโย อาปชฺชติ. “ภุญฺชิสฺสามี”ติ ปฏิคฺคณฺหาติ, อาปตฺติ ทุกฺกฏสฺส. อชฺโฌหาเร อชฺโฌหาเร อาปตฺติ ปาจิตฺติยสฺส.}

\prose{คณโภชนํ ภุญฺชนฺโต เทฺว อาปตฺติโย อาปชฺชติ. “ภุญฺชิสฺสามี”ติ ปฏิคฺคณฺหาติ, อาปตฺติ ทุกฺกฏสฺส. อชฺโฌหาเร อชฺโฌหาเร อาปตฺติ ปาจิตฺติยสฺส.}

\prose{ปรมฺปร\-โภชนํ ภุญฺชนฺโต เทฺว อาปตฺติโย อาปชฺชติ. “ภุญฺชิสฺสามี”ติ ปฏิคฺคณฺหาติ, อาปตฺติ ทุกฺกฏสฺส. อชฺโฌหาเร อชฺโฌหาเร อาปตฺติ ปาจิตฺติยสฺส.}

\prose{ทฺวตฺติ\-ปตฺตปูเร ปูเว ปฏิคฺคเหตฺวา ตตุตฺตริ ปฏิคฺคณฺหนฺโต เทฺว อาปตฺติโย อาปชฺชติ. คณฺหาติ, ปโยเค ทุกฺกฏํ. คหิเต อาปตฺติ ปาจิตฺติยสฺส.}

\prose{ภุตฺตาวี ปวาริโต อนติริตฺตํ ขาทนียํ วา โภชนียํ วา ภุญฺชนฺโต เทฺว อาปตฺติโย อาปชฺชติ. “ภุญฺชิสฺสามี”ติ ปฏิคฺคณฺหาติ, อาปตฺติ ทุกฺกฏสฺส. อชฺโฌหาเร อชฺโฌหาเร อาปตฺติ ปาจิตฺติยสฺส.}

\prose{ภิกฺขุํ ภุตฺตาวิํ ปวาริตํ อนติริตฺเตน ขาทนีเยน วา โภชนีเยน วา อภิหฏฺฐุํ ปวาเรนฺโต เทฺว อาปตฺติโย อาปชฺชติ. ตสฺส วจเนน “ขาทิสฺสามิ ภุญฺชิสฺสามี”ติ ปฏิคฺคณฺหาติ, อาปตฺติ ทุกฺกฏสฺส. โภชน\-ปริโยสาเน อาปตฺติ ปาจิตฺติยสฺส.}

\prose{วิกาเล ขาทนียํ วา โภชนียํ วา ภุญฺชนฺโต เทฺว อาปตฺติโย อาปชฺชติ. “ขาทิสฺสามิ ภุญฺชิสฺสามี”ติ ปฏิคฺคณฺหาติ, อาปตฺติ ทุกฺกฏสฺส. อชฺโฌหาเร อชฺโฌหาเร อาปตฺติ ปาจิตฺติยสฺส.}

\prose{สนฺนิธิการกํ ขาทนียํ วา โภชนียํ วา ภุญฺชนฺโต เทฺว อาปตฺติโย อาปชฺชติ. “ขาทิสฺสามิ ภุญฺชิสฺสามี”ติ ปฏิคฺคณฺหาติ, อาปตฺติ ทุกฺกฏสฺส. อชฺโฌหาเร อชฺโฌหาเร อาปตฺติ ปาจิตฺติยสฺส.}

\prose{ปณีต\-โภชนานิ อตฺตโน อตฺถาย วิญฺญาเปตฺวา ภุญฺชนฺโต เทฺว อาปตฺติโย อาปชฺชติ. “ภุญฺชิสฺสามี”ติ ปฏิคฺคณฺหาติ, อาปตฺติ ทุกฺกฏสฺส. อชฺโฌหาเร อชฺโฌหาเร อาปตฺติ ปาจิตฺติยสฺส.}

\prosewithclosermidruled{\csromanfolio{66}อทินฺนํ มุขทฺวารํ อาหารํ อาหรนฺโต เทฺว อาปตฺติโย อาปชฺชติ. “ภุญฺชิสฺสามี”ติ ปฏิคฺคณฺหาติ, อาปตฺติ ทุกฺกฏสฺส. อชฺโฌหาเร อชฺโฌหาเร อาปตฺติ ปาจิตฺติยสฺส.}{โภชนวคฺโค จตุตฺโถ.}
\csromanmatikaanchor{mtk.41}
\csromantocmarkref{section}{5. อเจลกวคฺค}{mtk.41}
\bookmark[dest={mtk.41},level=1]{5. อเจลกวคฺค}
\csromanheader{1}{5. \textbf{อเจลกวคฺค}}
\proseitem{169}{อเจลกสฺส วา ปริพฺพาชกสฺส วา ปริพฺพาชิกาย วา สหตฺถา ขาทนียํ วา โภชนียํ วา เทนฺโต เทฺว อาปตฺติโย อาปชฺชติ. เทติ, ปโยเค ทุกฺกฏํ. ทินฺเน อาปตฺติ ปาจิตฺติยสฺส.}

\prose{ภิกฺขุํ “เอหาวุโส, คามํ วา นิคมํ วา ปิณฺฑาย ปวิสิสฺสามา”ติ ตสฺส ทาเปตฺวา วา อทาเปตฺวา วา อุยฺโยเชนฺโต เทฺว อาปตฺติโย อาปชฺชติ. อุยฺโยเชติ, ปโยเค ทุกฺกฏํ. อุยฺโยชิเต อาปตฺติ ปาจิตฺติยสฺส.}

\prose{สโภชเน กุเล อนุปขชฺช นิสชฺชํ กปฺเปนฺโต เทฺว อาปตฺติโย อาปชฺชติ. นิสีทติ, ปโยเค ทุกฺกฏํ. นิสินฺเน อาปตฺติ ปาจิตฺติยสฺส.}

\prose{มาตุคาเมน สทฺธิํ รโห ปฏิจฺฉนฺเน อาสเน นิสชฺชํ กปฺเปนฺโต เทฺว อาปตฺติโย อาปชฺชติ. นิสีทติ, ปโยเค ทุกฺกฏํ. นิสินฺเน อาปตฺติ ปาจิตฺติยสฺส.}

\prose{มาตุคาเมน สทฺธิํ เอโก เอกาย รโห นิสชฺชํ กปฺเปนฺโต เทฺว อาปตฺติโย อาปชฺชติ. นิสีทติ, ปโยเค ทุกฺกฏํ. นิสินฺเน อาปตฺติ ปาจิตฺติยสฺส.}

\prose{นิมนฺติโต สภตฺโต สมาโน ปุเรภตฺตํ ปจฺฉาภตฺตํ กุเลสุ จาริตฺตํ อาปชฺชนฺโต เทฺว อาปตฺติโย อาปชฺชติ. ปฐมํ ปาทํ อุมฺมารํ อติกฺกาเมติ, อาปตฺติ ทุกฺกฏสฺส, ทุติยํ ปาทํ อติกฺกาเมติ, อาปตฺติ ปาจิตฺติยสฺส.}

\prose{ตตุตฺตริ เภสชฺชํ วิญฺญาเปนฺโต เทฺว อาปตฺติโย อาปชฺชติ. วิญฺญาเปติ, ปโยเค ทุกฺกฏํ. วิญฺญาปิเต อาปตฺติ ปาจิตฺติยสฺส.}

\prose{\csromanfolio{67}อุยฺยุตฺตํ เสนํ ทสฺสนาย คจฺฉนฺโต เทฺว อาปตฺติโย อาปชฺชติ. คจฺฉติ, อาปตฺติ ทุกฺกฏสฺส. ยตฺถ ฐิโต ปสฺสติ, อาปตฺติ ปาจิตฺติยสฺส.}

\prose{อติเรก\-ติรตฺตํ เสนาย วสนฺโต เทฺว อาปตฺติโย อาปชฺชติ. วสติ, ปโยเค ทุกฺกฏํ. วสิเต อาปตฺติ ปาจิตฺติยสฺส.}

\prosewithclosermidruled{อุยฺโยธิกํ คจฺฉนฺโต เทฺว อาปตฺติโย อาปชฺชติ. คจฺฉติ, อาปตฺติ ทุกฺกฏสฺส. ยตฺถ ฐิโต ปสฺสติ, อาปตฺติ ปาจิตฺติยสฺส.}{อเจลกวคฺโค ปญฺจโม.}
\csromanmatikaanchor{mtk.42}
\csromantocmarkref{section}{6. สุราเมรยวคฺค}{mtk.42}
\bookmark[dest={mtk.42},level=1]{6. สุราเมรยวคฺค}
\csromanheader{1}{6. สุราเมรย\-วคฺค}
\proseitem{170}{มชฺชํ ปิวนฺโต เทฺว อาปตฺติโย อาปชฺชติ. “ปิวิสฺสามี”ติ ปฏิคฺคณฺหาติ, อาปตฺติ ทุกฺกฏสฺส. อชฺโฌหาเร อชฺโฌหาเร อาปตฺติ ปาจิตฺติยสฺส.}

\prose{ภิกฺขุํ องฺคุลิ\-ปโตทเกน หาเสนฺโต เทฺว อาปตฺติโย อาปชฺชติ. หาเสติ, ปโยเค ทุกฺกฏํ. หสิเต อาปตฺติ ปาจิตฺติยสฺส.}

\prose{อุทเก กีฬนฺโต เทฺว อาปตฺติโย อาปชฺชติ. เหฏฺฐาโคปฺผเก อุทเก กีฬติ, อาปตฺติ ทุกฺกฏสฺส. อุปริโคปฺผเก กีฬติ, อาปตฺติ ปาจิตฺติยสฺส.}

\prose{อนาทริยํ กโรนฺโต เทฺว อาปตฺติโย อาปชฺชติ. กโรติ, ปโยเค ทุกฺกฏํ. กเต อาปตฺติ ปาจิตฺติยสฺส.}

\prose{ภิกฺขุํ ภิํสาเปนฺโต เทฺว อาปตฺติโย อาปชฺชติ. ภิํสาเปติ, ปโยเค ทุกฺกฏํ. ภิํสาปิเต อาปตฺติ ปาจิตฺติยสฺส.}

\prose{โชติํ สมาทหิตฺวา วิสิพฺเพนฺโต เทฺว อาปตฺติโย อาปชฺชติ. สมาทหติ, ปโยเค ทุกฺกฏํ. สมาทหิเต อาปตฺติ ปาจิตฺติยสฺส.}

\prose{โอเรนทฺธมาสํ นหายนฺโต เทฺว อาปตฺติโย อาปชฺชติ. นหายติ, ปโยเค ทุกฺกฏํ. นหาน\-ปริโยสาเน อาปตฺติ ปาจิตฺติยสฺส.}

\prose{อนาทิยิตฺวา ติณฺณํ ทุพฺพณฺณ\-กรณานํ อญฺญตรํ ทุพฺพณฺณ\-กรณํ นวํ จีวรํ ปริภุญฺชนฺโต เทฺว อาปตฺติโย อาปชฺชติ. ปริภุญฺชติ, ปโยเค ทุกฺกฏํ. ปริภุตฺเต อาปตฺติ ปาจิตฺติยสฺส.}

\prose{\csromanfolio{68}ภิกฺขุสฺส วา ภิกฺขุนิยา วา สิกฺขมานาย วา สามเณรสฺส วา สามเณริยา วา สามํ จีวรํ วิกปฺเปตฺวา อปฺปจฺจุทฺธารณํ ปริภุญฺชนฺโต เทฺว อาปตฺติโย อาปชฺชติ. ปริภุญฺชติ, ปโยเค ทุกฺกฏํ. ปริภุตฺเต อาปตฺติ ปาจิตฺติยสฺส.}

\prosewithclosermidruled{ภิกฺขุสฺส ปตฺตํ วา จีวรํ วา นิสีทนํ วา สูจิฆรํ วา กายพนฺธนํ วา อปนิเธนฺโต เทฺว อาปตฺติโย อาปชฺชติ. อปนิเธติ, ปโยเค ทุกฺกฏํ. อปนิธิเต อาปตฺติ ปาจิตฺติยสฺส.}{สุราเมรย\-วคฺโค ฉฏฺโฐ.}
\csromanmatikaanchor{mtk.43}
\csromantocmarkref{section}{7. สปฺปาณกวคฺค}{mtk.43}
\bookmark[dest={mtk.43},level=1]{7. สปฺปาณกวคฺค}
\csromanheader{1}{7. \textbf{สปฺปาณกวคฺค}}
\proseitem{171}{สญฺจิจฺจ ปาณํ ชีวิตา โวโรเปนฺโต กติ อาปตฺติโย อาปชฺชติ. สญฺจิจฺจ ปาณํ ชีวิตา โวโรเปนฺโต จตสฺโส อาปตฺติโย อาปชฺชติ. อโนทิสฺส โอปาตํ ขณติ “โย โกจิ ปปติตฺวา มริสฺสตี”ติ, อาปตฺติ ทุกฺกฏสฺส. มนุสฺโส ตสฺมิํ ปปติตฺวา มรติ, อาปตฺติ ปาราชิกสฺส. ยกฺโข วา เปโต วา ติรจฺฉานคต\-มนุสฺสวิคฺคโห วา ตสฺมิํ ปปติตฺวา มรติ, อาปตฺติ ถุลฺลจฺจยสฺส. ติรจฺฉานคโต ตสฺมิํ ปปติตฺวา มรติ, อาปตฺติ ปาจิตฺติยสฺส. สญฺจิจฺจ ปาณํ ชีวิตา โวโรเปนฺโต อิมา จตสฺโส อาปตฺติโย อาปชฺชติ.}

\prose{ชานํ สปฺปาณกํ อุทกํ ปริภุญฺชนฺโต เทฺว อาปตฺติโย อาปชฺชติ. ปริภุญฺชติ, ปโยเค ทุกฺกฏํ. ปริภุตฺเต อาปตฺติ ปาจิตฺติยสฺส.}

\prose{ชานํ ยถาธมฺมํ นิหตา\-ธิกรณํ ปุนกมฺมาย อุกฺโกเฏนฺโต เทฺว อาปตฺติโย อาปตฺติโย อาปชฺชติ. อุกฺโกเฏติ, ปโยเค ทุกฺกฏํ. อุกฺโกฏิเต อาปตฺติ ปาจิตฺติยสฺส.}

\prose{ภิกฺขุสฺส ชานํ ทุฏฺฐุลฺลํ อาปตฺติํ ปฏิจฺฉาเทนฺโต เอกํ อาปตฺติํ อาปชฺชติ, ปาจิตฺติยํ.}

\prose{ชานํ อูนวีสติวสฺสํ ปุคฺคลํ อุปสมฺปาเทนฺโต เทฺว อาปตฺติโย อาปชฺชติ. อุปสมฺปาเทติ, ปโยเค ทุกฺกฏํ. อุปสมฺปาทิเต อาปตฺติ ปาจิตฺติยสฺส.}

\prose{\csromanfolio{69}ชานํ เถยฺยสตฺเถน สทฺธิํ สํวิธาย เอก\-ทฺธาน\-มคฺคํ ปฏิปชฺชนฺโต เทฺว อาปตฺติโย อาปชฺชติ. ปฏิปชฺชติ, ปโยเค ทุกฺกฏํ. ปฏิปนฺเน อาปตฺติ ปาจิตฺติยสฺส.}

\prose{มาตุคาเมน สทฺธิํ สํวิธาย เอก\-ทฺธาน\-มคฺคํ ปฏิปชฺชนฺโต เทฺว อาปตฺติโย อาปชฺชติ. ปฏิปชฺชติ, ปโยเค ทุกฺกฏํ. ปฏิปนฺเน อาปตฺติ ปาจิตฺติยสฺส.}

\prose{ปาปิกาย ทิฏฺฐิยา ยาวตติยํ สมนุภาสนาย น ปฏินิสฺสชฺชนฺโต เทฺว อาปตฺติโย อาปชฺชติ. ญตฺติยา ทุกฺกฏํ. กมฺมวาจา\-ปริโยสาเน อาปตฺติ ปาจิตฺติยสฺส.}

\prose{ชานํ ตถาวาทินา ภิกฺขุนา อกฏา\-นุธมฺเมน ตํทิฏฺฐิํ อปฺปฏินิสฺสฏฺเฐน สทฺธิํ สมฺภุญฺชนฺโต เทฺว อาปตฺติโย อาปชฺชติ. สมฺภุญฺชติ, ปโยเค ทุกฺกฏํ. สมฺภุตฺเต อาปตฺติ ปาจิตฺติยสฺส.}

\prosewithclosermidruled{ชานํ ตถานาสิตํ สมณุทฺเทสํ อุปลาเปนฺโต เทฺว อาปตฺติโย อาปชฺชติ. อุปลาเปติ, ปโยเค ทุกฺกฏํ. อุปลาปิเต อาปตฺติ ปาจิตฺติยสฺส.}{สปฺปาณกวคฺโค สตฺตโม.}
\csromanmatikaanchor{mtk.44}
\csromantocmarkref{section}{8. สหธมฺมิกวคฺค}{mtk.44}
\bookmark[dest={mtk.44},level=1]{8. สหธมฺมิกวคฺค}
\csromanheader{1}{8. สหธมฺมิก\-วคฺค}
\proseitem{172}{ภิกฺขูหิ สหธมฺมิกํ วุจฺจมาโน “น ตาวาหํ อาวุโส เอตสฺมิํ สิกฺขาปเท สิกฺขิสฺสามิ, ยาว น อญฺญํ ภิกฺขุํ พฺยตฺตํ วินยธรํ ปริปุจฺฉิสฺสามี”ติ ภณนฺโต เทฺว อาปตฺติโย อาปชฺชติ. ภณติ, ปโยเค ทุกฺกฏํ. ภณิเต อาปตฺติ ปาจิตฺติยสฺส.}

\prose{วินยํ วิวณฺเณนฺโต เทฺว อาปตฺติโย อาปชฺชติ. วิวณฺเณติ, ปโยเค ทุกฺกฏํ. วิวณฺณิเต อาปตฺติ ปาจิตฺติยสฺส.}

\prose{โมเหนฺโต เทฺว อาปตฺติโย อาปชฺชติ. อนาโรปิเต โมเห โมเหติ, อาปตฺติ ทุกฺกฏสฺส. อาโรปิเต โมเห โมเหติ, อาปตฺติ ปาจิตฺติยสฺส.}

\prose{\csromanfolio{70}ภิกฺขุสฺส กุปิโต อนตฺตมโน ปหารํ เทนฺโต เทฺว อาปตฺติโย อาปชฺชติ. ปหรติ, ปโยเค ทุกฺกฏํ. ปหเต อาปตฺติ ปาจิตฺติยสฺส.}

\prose{ภิกฺขุสฺส กุปิโต อนตฺตมโน ตลสตฺติกํ อุคฺคิรนฺโต เทฺว อาปตฺติโย อาปชฺชติ. อุคฺคิรติ, ปโยเค ทุกฺกฏํ. อุคฺคิริเต อาปตฺติ ปาจิตฺติยสฺส.}

\prose{ภิกฺขุํ อมูลเกน สํฆาทิเสเสน อนุทฺธํเสนฺโต เทฺว อาปตฺติโย อาปชฺชติ. อนุทฺธํเสติ, ปโยเค ทุกฺกฏํ. อนุทฺธํสิเต อาปตฺติ ปาจิตฺติยสฺส.}

\prose{ภิกฺขุสฺส สญฺจิจฺจ กุกฺกุจฺจํ อุปทหนฺโต เทฺว อาปตฺติโย อาปชฺชติ. อุปทหติ, ปโยเค ทุกฺกฏํ. อุปทหิเต อาปตฺติ ปาจิตฺติยสฺส.}

\prose{ภิกฺขูนํ ภณฺฑน\-ชาตานํ กลหชาตานํ วิวาทาปนฺนานํ อุปสฺสุติํ ติฏฺฐนฺโต เทฺว อาปตฺติโย อาปชฺชติ. “โสสฺสามี”ติ คจฺฉติ, อาปตฺติ ทุกฺกฏสฺส. ยตฺถ ฐิโต สุณาติ, อาปตฺติ ปาจิตฺติยสฺส.}

\prose{ธมฺมิกานํ กมฺมานํ ฉนฺทํทตฺวา ปจฺฉา ขียนธมฺมํ อาปชฺชนฺโต เทฺว อาปตฺติโย อาปชฺชติ. ขิยฺยติ, ปโยเค ทุกฺกฏํ. ขิยฺยิเต อาปตฺติ ปาจิตฺติยสฺส.}

\prose{สํเฆ วินิจฺฉย\-กถาย วตฺตมานาย ฉนฺทํ อทตฺวา อุฏฺฐายาสนา ปกฺกมนฺโต เทฺว อาปตฺติโย อาปชฺชติ. ปริสาย หตฺถปาสํ วิชหนฺตสฺส อาปตฺติ ทุกฺกฏสฺส. วิชหิเต อาปตฺติ ปาจิตฺติยสฺส.}

\prose{สมคฺเคน สํเฆน จีวรํ ทตฺวา ปจฺฉา ขียนธมฺมํ อาปชฺชนฺโต เทฺว อาปตฺติโย อาปชฺชติ. ขิยฺยติ, ปโยเค ทุกฺกฏํ. ขิยฺยิเต อาปตฺติ ปาจิตฺติยสฺส.}

\prosewithclosermidruled{ชานํ สํฆิกํ ลาภํ ปริณตํ ปุคฺคลสฺส ปริณาเมนฺโต เทฺว อาปตฺติโย อาปชฺชติ. ปริณาเมติ, ปโยเค ทุกฺกฏํ. ปริณามิเต อาปตฺติ ปาจิตฺติยสฺส.}{สหธมฺมิก\-วคฺโค อฏฺฐโม.}
\csromanmatikaanchor{mtk.45}
\csromantocmarkref{section}{9. ราชวคฺค}{mtk.45}
\bookmark[dest={mtk.45},level=1]{9. ราชวคฺค}
\csromanheader{1}{9. \textbf{ราชวคฺค}}
\proseitem{173}{\csromanfolio{71}ปุพฺเพ อปฺปฏิสํวิทิโต รญฺโญ อนฺเตปุรํ ปวิสนฺโต เทฺว อาปตฺติโย อาปชฺชติ. ปฐมํ ปาทํ อุมฺมารํ อติกฺกาเมติ, อาปตฺติ ทุกฺกฏสฺส. ทุติยํ ปาทํ อติกฺกาเมติ, อาปตฺติ ปาจิตฺติยสฺส.}

\prose{รตนํ อุคฺคณฺหนฺโต เทฺว อาปตฺติโย อาปชฺชติ. คณฺหาติ, ปโยเค ทุกฺกฏํ. คหิเต อาปตฺติ ปาจิตฺติยสฺส.}

\prose{สนฺตํ ภิกฺขุํ อนาปุจฺฉา วิกาเล คามํ ปวิสนฺโต เทฺว อาปตฺติโย อาปชฺชติ. ปฐมํ ปาทํ ปริกฺเขปํ อติกฺกาเมติ, อาปตฺติ ทุกฺกฏสฺส. ทุติยํ ปาทํ อติกฺกาเมติ, อาปตฺติ ปาจิตฺติยสฺส.}

\prose{อฏฺฐิมยํ วา ทนฺตมยํ วา วิสาณมยํ วา สูจิฆรํ การาเปนฺโต เทฺว อาปตฺติโย อาปชฺชติ. การาเปติ, ปโยเค ทุกฺกฏํ. การาปิเต อาปตฺติ ปาจิตฺติยสฺส.}

\prose{ปมาณา\-ติกฺกนฺตํ มญฺจํ วา ปีฐํ วา การาเปนฺโต เทฺว อาปตฺติโย อาปชฺชติ. การาเปติ, ปโยเค ทุกฺกฏํ. การาปิเต อาปตฺติ ปาจิตฺติยสฺส.}

\prose{มญฺจํ วา ปีฐํ วา ตูโลนทฺธํ การาเปนฺโต เทฺว อาปตฺติโย อาปชฺชติ. การาเปติ, ปโยเค ทุกฺกฏํ. การาปิเต อาปตฺติ ปาจิตฺติยสฺส.}

\prose{ปมาณา\-ติกฺกนฺตํ นิสีทนํ การาเปนฺโต เทฺว อาปตฺติโย อาปชฺชติ. การาเปติ, ปโยเค ทุกฺกฏํ. การาปิเต อาปตฺติ ปาจิตฺติยสฺส.}

\prose{ปมาณา\-ติกฺกนฺตํ กณฺฑุ\-ปฺปฏิจฺฉาทิํ การาเปนฺโต เทฺว อาปตฺติโย อาปชฺชติ. การาเปติ, ปโยเค ทุกฺกฏํ. การาปิเต อาปตฺติ ปาจิตฺติยสฺส.}

\prose{ปมาณา\-ติกฺกนฺตํ วสฺสิกสาฏิกํ การาเปนฺโต เทฺว อาปตฺติโย อาปชฺชติ. การาเปติ, ปโยเค ทุกฺกฏํ. การาปิเต อาปตฺติ ปาจิตฺติยสฺส.}

\prosewithclosermidruled{สุคต\-จีวร\-ปฺปมาณํ จีวรํ การาเปนฺโต กติ อาปตฺติโย อาปชฺชติ. สุคต\-จีวร\-ปฺปมาณํ จีวรํ การาเปนฺโต เทฺว อาปตฺติโย อาปชฺชติ. \csromanfolio{72}การาเปติ, ปโยเค ทุกฺกฏํ. การาปิเต อาปตฺติ ปาจิตฺติยสฺส. สุคต จีวรปฺปมาณํ จีวรํ การาเปนฺโต อิมา เทฺว อาปตฺติโย อาปชฺชติ.}{ราชวคฺโค นวโม.}
\nitthitamruled{ขุทฺทกา นิฏฺฐิตา.}
\csromanmatikaanchor{mtk.46}
\csromantocmarkref{chapter}{6. ปาฏิเทสนียกณฺฑ}{mtk.46}
\bookmark[dest={mtk.46},level=0]{6. ปาฏิเทสนียกณฺฑ}
\chapterhead{6. ปาฏิเทสนีย\-กณฺฑ}
\proseitem{174}{อญฺญาติกาย ภิกฺขุนิยา อนฺตรฆรํ ปวิฏฺฐาย หตฺถโต ขาทนียํ วา โภชนียํ วา สหตฺถา ปฏิคฺคเหตฺวา ภุญฺชนฺโต กติ อาปตฺติโย อาปชฺชติ. อญฺญาติกาย ภิกฺขุนิยา อนฺตรฆรํ ปวิฏฺฐาย หตฺถโต ขาทนียํ วา โภชนียํ วา สหตฺถา ปฏิคฺคเหตฺวา ภุญฺชนฺโต เทฺว อาปตฺติโย อาปชฺชติ. “ภุญฺชิสฺสามี”ติ ปฏิคฺคณฺหาติ, อาปตฺติ ทุกฺกฏสฺส. อชฺโฌหาเร อชฺโฌหาเร อาปตฺติ ปาฏิเทสนียสฺส. อญฺญาติกาย ภิกฺขุนิยา อนฺตรฆรํ ปวิฏฺฐาย หตฺถโต ขาทนียํ วา โภชนียํ วา สหตฺถา ปฏิคฺคเหตฺวา ภุญฺชนฺโต อิมา เทฺว อาปตฺติโย อาปชฺชติ.}

\prose{ภิกฺขุนิยา โวสาสนฺติยา น นิวาเรตฺวา ภุญฺชนฺโต เทฺว อาปตฺติโย อาปชฺชติ. “ภุญฺชิสฺสามี”ติ ปฏิคฺคณฺหาติ, อาปตฺติ ทุกฺกฏสฺส. อชฺโฌหาเร อชฺโฌหาเร อาปตฺติ ปาฏิเทสนียสฺส.}

\prose{เสกฺข\-สมฺมเตสุ กุเลสุ ขาทนียํ วา โภชนียํ วา สหตฺถา ปฏิคฺคเหตฺวา ภุญฺชนฺโต เทฺว อาปตฺติโย อาปชฺชติ. “ภุญฺชิสฺสามี”ติ ปฏิคฺคณฺหาติ, อาปตฺติ ทุกฺกฏสฺส. อชฺโฌหาเร อชฺโฌหาเร อาปตฺติ ปาฏิเทสนียสฺส.}

\prosewithcloserruled{อารญฺญเกสุ เสนาสเนสุ ปุพฺเพ อปฺปฏิสํวิทิตํ ขาทนียํ วา โภชนียํ วา อชฺฌาราเม สหตฺถา ปฏิคฺคเหตฺวา ภุญฺชนฺโต กติ อาปตฺติโย อาปชฺชติ. อารญฺญเกสุ เสนาสเนสุ ปุพฺเพ อปฺปฏิสํวิทิตํ ขาทนียํ วา โภชนียํ วา อชฺฌาราเม สหตฺถา ปฏิคฺคเหตฺวา ภุญฺชนฺโต เทฺว อาปตฺติโย อาปชฺชติ. “ภุญฺชิสฺสามี”ติ ปฏิคฺคณฺหาติ, อาปตฺติ ทุกฺกฏสฺส. \csromanfolio{73}อชฺโฌหาเร อชฺโฌหาเร อาปตฺติ ปาฏิเทสนียสฺส. อารญฺญเกสุ เสนาสเนสุ ปุพฺเพ อปฺปฏิสํวิทิตํ ขาทนียํ วา โภชนียํ วา อชฺฌาราเม สหตฺถา ปฏิคฺคเหตฺวา ภุญฺชนฺโต อิมา เทฺว อาปตฺติโย อาปชฺชติ.}{จตฺตาโร ปาฏิเทสนียา นิฏฺฐิตา.}
\csromanmatikaanchor{mtk.47}
\csromanmatikaanchor{mtk.48}
\csromantocmarkref{chapter}{7. เสขิยกณฺฑ}{mtk.47}
\bookmark[dest={mtk.47},level=0]{7. เสขิยกณฺฑ}
\csromantocmarkref{section}{1. ปริมณฺฑลวคฺค}{mtk.48}
\bookmark[dest={mtk.48},level=1]{1. ปริมณฺฑลวคฺค}
\csromanmarkh{7. เสขิยกณฺฑ}\csromanmarkhii{1. ปริมณฺฑล\-วคฺค}\csromanheadernomark{1}{7. เสขิยกณฺฑ 1. ปริมณฺฑล\-วคฺค}
\proseitem{175}{อนาทริยํ ปฏิจฺจ ปุรโต วา ปจฺฉโต วา โอลมฺเพนฺโต นิวาเสนฺโต กติ อาปตฺติโย อาปชฺชติ. อนาทริยํ ปฏิจฺจ ปุรโต วา ปจฺฉโต วา โอลมฺเพนฺโต นิวาเสนฺโต เอกํ อาปตฺติํ อาปชฺชติ, ทุกฺกฏํ. อนาทริยํ ปฏิจฺจ ปุรโต วา ปจฺฉโต วา โอลมฺเพนฺโต นิวาเสนฺโต อิมํ เอกํ อาปตฺติํ อาปชฺชติ.}

\prose{อนาทริยํ ปฏิจฺจ ปุรโต วา ปจฺฉโต วา โอลมฺเพนฺโต ปารุปนฺโต เอกํ อาปตฺติํ อาปชฺชติ, ทุกฺกฏํ.}

\prose{อนาทริยํ ปฏิจฺจ กายํ วิวริตฺวา อนฺตรฆเร คจฺฉนฺโต เอกํ อาปตฺติํ อาปชฺชติ, ทุกฺกฏํ.}

\prose{อนาทริยํ ปฏิจฺจ กายํ วิวริตฺวา อนฺตรฆเร นิสีทนฺโต เอกํ อาปตฺติํ อาปชฺชติ, ทุกฺกฏํ.}

\prose{อนาทริยํ ปฏิจฺจ หตฺถํ วา ปาทํ วา กีฬาเปนฺโต อนฺตรฆเร คจฺฉนฺโต เอกํ อาปตฺติํ อาปชฺชติ, ทุกฺกฏํ.}

\prose{อนาทริยํ ปฏิจฺจ หตฺถํ วา ปาทํ วา กีฬาเปนฺโต อนฺตรฆเร นิสีทนฺโต เอกํ อาปตฺติํ อาปชฺชติ, ทุกฺกฏํ.}

\prose{อนาทริยํ ปฏิจฺจ ตหํ ตหํ โอโลเกนฺโต อนฺตรฆเร คจฺฉนฺโต เอกํ อาปตฺติํ อาปชฺชติ, ทุกฺกฏํ.}

\prose{อนาทริยํ ปฏิจฺจ ตหํ ตหํ โอโลเกนฺโต อนฺตรฆเร นิสีทนฺโต เอกํ อาปตฺติํ อาปชฺชติ, ทุกฺกฏํ.}

\prose{\csromanfolio{74}อนาทริยํ ปฏิจฺจ อุกฺขิตฺตกาย อนฺตรฆเร คจฺฉนฺโต เอกํ อาปตฺติํ อาปชฺชติ, ทุกฺกฏํ.}

\prosewithclosermidruled{อนาทริยํ ปฏิจฺจ อุกฺขิตฺตกาย อนฺตรฆเร นิสีทนฺโต เอกํ อาปตฺติํ อาปชฺชติ, ทุกฺกฏํ.}{ปริมณฺฑล\-วคฺโค ปฐโม.}
\csromanmatikaanchor{mtk.49}
\csromantocmarkref{section}{2. อุชฺชคฺฆิกวคฺค}{mtk.49}
\bookmark[dest={mtk.49},level=1]{2. อุชฺชคฺฆิกวคฺค}
\csromanheader{1}{2. อุชฺชคฺฆิก\-วคฺค}
\proseitem{176}{อนาทริยํ ปฏิจฺจ อุชฺชคฺฆิกาย อนฺตรฆเร คจฺฉนฺโต เอกํ อาปตฺติํ อาปชฺชติ, ทุกฺกฏํ.}

\prose{อนาทริยํ ปฏิจฺจ อุชฺชคฺฆิกาย อนฺตรฆเร นิสีทนฺโต เอกํ อาปตฺติํ อาปชฺชติ, ทุกฺกฏํ.}

\prose{อนาทริยํ ปฏิจฺจ อุจฺจาสทฺทํ มหาสทฺทํ กโรนฺโต อนฺตรฆเร คจฺฉนฺโต เอกํ อาปตฺติํ อาปชฺชติ, ทุกฺกฏํ.}

\prose{อนาทริยํ ปฏิจฺจ อุจฺจาสทฺทํ มหาสทฺทํ กโรนฺโต อนฺตรฆเร นิสีทนฺโต เอกํ อาปตฺติํ อาปชฺชติ, ทุกฺกฏํ.}

\prose{อนาทริยํ ปฏิจฺจ กายปฺปจาลกํ อนฺตรฆเร คจฺฉนฺโต เอกํ อาปตฺติํ อาปชฺชติ, ทุกฺกฏํ.}

\prose{อนาทริยํ ปฏิจฺจ กายปฺปจาลกํ อนฺตรฆเร นิสีทนฺโต เอกํ อาปตฺติํ อาปชฺชติ, ทุกฺกฏํ.}

\prose{อนาทริยํ ปฏิจฺจ พาหุปฺปจาลกํ อนฺตรฆเร คจฺฉนฺโต เอกํ อาปตฺติํ อาปชฺชติ, ทุกฺกฏํ.}

\prose{อนาทริยํ ปฏิจฺจ พาหุปฺปจาลกํ อนฺตรฆเร นิสีทนฺโต เอกํ อาปตฺติํ อาปชฺชติ, ทุกฺกฏํ.}

\prose{อนาทริยํ ปฏิจฺจ สีสปฺปจาลกํ อนฺตรฆเร คจฺฉนฺโต เอกํ อาปตฺติํ อาปชฺชติ, ทุกฺกฏํ.}

\prosewithclosermidruled{อนาทริยํ ปฏิจฺจ สีสปฺปจาลกํ อนฺตรฆเร นิสีทนฺโต เอกํ อาปตฺติํ อาปชฺชติ, ทุกฺกฏํ.}{อุชฺชคฺฆิก\-วคฺโค ทุติโย.}
\csromanmatikaanchor{mtk.50}
\csromantocmarkref{section}{3. ขมฺภกตวคฺค}{mtk.50}
\bookmark[dest={mtk.50},level=1]{3. ขมฺภกตวคฺค}
\csromanheader{1}{3. ขมฺภกต\-วคฺค}
\proseitem{177}{\csromanfolio{75}อนาทริยํ ปฏิจฺจ ขมฺภกโต อนฺตรฆเร คจฺฉนฺโต เอกํ อาปตฺติํ อาปชฺชติ, ทุกฺกฏํ.}

\prose{อนาทริยํ ปฏิจฺจ ขมฺภกโต อนฺตรฆเร นิสีทนฺโต เอกํ อาปตฺติํ อาปชฺชติ, ทุกฺกฏํ.}

\prose{อนาทริยํ ปฏิจฺจ โอคุณฺฐิโต อนฺตรฆเร คจฺฉนฺโต เอกํ อาปตฺติํ อาปชฺชติ, ทุกฺกฏํ.}

\prose{อนาทริยํ ปฏิจฺจ โอคุณฺฐิโต อนฺตรฆเร นิสีทนฺโต เอกํ อาปตฺติํ อาปชฺชติ, ทุกฺกฏํ.}

\prose{อนาทริยํ ปฏิจฺจ อุกฺกุฏิกาย อนฺตรฆเร คจฺฉนฺโต เอกํ อาปตฺติํ อาปชฺชติ, ทุกฺกฏํ.}

\prose{อนาทริยํ ปฏิจฺจ ปลฺลตฺถิกาย อนฺตรฆเร นิสีทนฺโต เอกํ อาปตฺติํ อาปชฺชติ, ทุกฺกฏํ.}

\prose{อนาทริยํ ปฏิจฺจ อสกฺกจฺจํ ปิณฺฑปาตํ ปฏิคฺคณฺหนฺโต เอกํ อาปตฺติํ อาปชฺชติ, ทุกฺกฏํ.}

\prose{อนาทริยํ ปฏิจฺจ ตหํ ตหํ โอโลเกนฺโต ปิณฺฑปาตํ ปฏิคฺคณฺหนฺโต เอกํ อาปตฺติํ อาปชฺชติ, ทุกฺกฏํ.}

\prose{อนาทริยํ ปฏิจฺจ สูปญฺเญว พหุํ ปฏิคฺคณฺหนฺโต เอกํ อาปตฺติํ อาปชฺชติ, ทุกฺกฏํ.}

\prosewithcloserruled{อนาทริยํ ปฏิจฺจ ถูปีกตํ ปิณฺฑปาตํ ปฏิคฺคณฺหนฺโต เอกํ อาปตฺติํ อาปชฺชติ, ทุกฺกฏํ.}{ขมฺภกตวคฺโคตติโย.}
\csromanmatikaanchor{mtk.51}
\csromantocmarkref{section}{4. ปิณฺฑปาตวคฺค}{mtk.51}
\bookmark[dest={mtk.51},level=1]{4. ปิณฺฑปาตวคฺค}
\csromanheader{1}{4. \textbf{ปิณฺฑปาตวคฺค}}
\proseitem{178}{อนาทริยํ ปฏิจฺจ อสกฺกจฺจํ ปิณฺฑปาตํ ภุญฺชนฺโต เอกํ อาปตฺติํ อาปชฺชติ, ทุกฺกฏํ.}

\prose{\csromanfolio{76}อนาทริยํ ปฏิจฺจ ตหํ ตหํ โอโลเกนฺโต ปิณฺฑปาตํ ภุญฺชนฺโต เอกํ อาปตฺติํ อาปชฺชติ, ทุกฺกฏํ.}

\prose{อนาทริยํ ปฏิจฺจ ตหํ ตหํ โอมสิตฺวา ปิณฺฑปาตํ ภุญฺชนฺโต เอกํ อาปตฺติํ อาปชฺชติ, ทุกฺกฏํ.}

\prose{อนาทริยํ ปฏิจฺจ สูปญฺเญว พหุํ ภุญฺชนฺโต เอกํ อาปตฺติํ อาปชฺชติ, ทุกฺกฏํ.}

\prose{อนาทริยํ ปฏิจฺจ ถูปกโต โอมทฺทิตฺวา ปิณฺฑปาตํ ภุญฺชนฺโต เอกํ อาปตฺติํ อาปชฺชติ, ทุกฺกฏํ.}

\prose{อนาทริยํ ปฏิจฺจ สูปํ วา พฺยญฺชนํ วา โอทเนน ปฏิจฺฉาเทนฺโต เอกํ อาปตฺติํ อาปชฺชติ, ทุกฺกฏํ.}

\prose{อนาทริยํ ปฏิจฺจ สูปํ วา โอทนํ วา อคิลาโน อตฺตโน อตฺถาย วิญฺญาเปตฺวา ภุญฺชนฺโต เอกํ อาปตฺติํ อาปชฺชติ, ทุกฺกฏํ.}

\prose{อนาทริยํ ปฏิจฺจ อุชฺฌานสญฺญี ปเรสํ ปตฺตํ โอโลเกนฺโต เอกํ อาปตฺติํ อาปชฺชติ, ทุกฺกฏํ.}

\prose{อนาทริยํ ปฏิจฺจ มหนฺตํ กพฬํ กโรนฺโต เอกํ อาปตฺติํ อาปชฺชติ, ทุกฺกฏํ.}

\prosewithclosermidruled{อนาทริยํ ปฏิจฺจ ทีฆํ อาโลปํ กโรนฺโต เอกํ อาปตฺติํ อาปชฺชติ, ทุกฺกฏํ.}{ปิณฺฑปาตวคฺโค จตุตฺโถ.}
\csromanmatikaanchor{mtk.52}
\csromantocmarkref{section}{5. กพลวคฺค}{mtk.52}
\bookmark[dest={mtk.52},level=1]{5. กพลวคฺค}
\csromanheader{1}{5. \textbf{กพฬวคฺค}}
\proseitem{179}{อนาทริยํ ปฏิจฺจ อนาหเฏ กพเฬ มุขทฺวารํ วิวรนฺโต เอกํ อาปตฺติํ อาปชฺชติ, ทุกฺกฏํ.}

\prose{อนาทริยํ ปฏิจฺจ ภุญฺชมาโน สพฺพํ หตฺถํ มุเข ปกฺขิปนฺโต เอกํ อาปตฺติํ อาปชฺชติ, ทุกฺกฏํ.}

\prose{อนาทริยํ ปฏิจฺจ สกพเฬน มุเขน พฺยาหรนฺโต เอกํ อาปตฺติํ อาปชฺชติ, ทุกฺกฏํ.}

\prose{\csromanfolio{77}อนาทริยํ ปฏิจฺจ ปิณฺฑุกฺเขปกํ ภุญฺชนฺโต เอกํ อาปตฺติํ อาปชฺชติ, ทุกฺกฏํ.}

\prose{อนาทริยํ ปฏิจฺจ กพฬาวจฺเฉทกํ ภุญฺชนฺโต เอกํ อาปตฺติํ อาปชฺชติ, ทุกฺกฏํ.}

\prose{อนาทริยํ ปฏิจฺจ อวคณฺฑการกํ ภุญฺชนฺโต เอกํ อาปตฺติํ อาปชฺชติ, ทุกฺกฏํ.}

\prose{อนาทริยํ ปฏิจฺจ หตฺถนิทฺธุนกํ ภุญฺชนฺโต เอกํ อาปตฺติํ อาปชฺชติ, ทุกฺกฏํ.}

\prose{อนาทริยํ ปฏิจฺจ สิตฺถาวการกํ ภุญฺชนฺโต เอกํ อาปตฺติํ อาปชฺชติ, ทุกฺกฏํ.}

\prose{อนาทริยํ ปฏิจฺจ ชิวฺหานิจฺฉารกํ ภุญฺชนฺโต เอกํ อาปตฺติํ อาปชฺชติ, ทุกฺกฏํ.}

\prosewithclosermidruled{อนาทริยํ ปฏิจฺจ จปุจปุการกํ ภุญฺชนฺโต เอกํ อาปตฺติํ อาปชฺชติ, ทุกฺกฏํ.}{กพฬวคฺโค ปญฺจโม.}
\csromanmatikaanchor{mtk.53}
\csromantocmarkref{section}{6. สุรุสุรุวคฺค}{mtk.53}
\bookmark[dest={mtk.53},level=1]{6. สุรุสุรุวคฺค}
\csromanheader{1}{6. \textbf{สุรุสุรุวคฺค}}
\proseitem{180}{อนาทริยํ ปฏิจฺจ สุรุสุรุการกํ ภุญฺชนฺโต เอกํ อาปตฺติํ อาปชฺชติ, ทุกฺกฏํ.}

\prose{อนาทริยํ ปฏิจฺจ หตฺถนิลฺเลหกํ ภุญฺชนฺโต เอกํ อาปตฺติํ อาปชฺชติ, ทุกฺกฏํ.}

\prose{อนาทริยํ ปฏิจฺจ ปตฺตนิลฺเลหกํ ภุญฺชนฺโต เอกํ อาปตฺติํ อาปชฺชติ, ทุกฺกฏํ.}

\prose{อนาทริยํ ปฏิจฺจ โอฏฺฐนิลฺเลหกํ ภุญฺชนฺโต เอกํ อาปตฺติํ อาปชฺชติ, ทุกฺกฏํ.}

\prose{อนาทริยํ ปฏิจฺจ สามิเสน หตฺเถน ปานียถาลกํ ปฏิคฺคณฺหนฺโต เอกํ อาปตฺติํ อาปชฺชติ, ทุกฺกฏํ.}

\prose{\csromanfolio{78}อนาทริยํ ปฏิจฺจ สสิตฺถกํ ปตฺตโธวนํ อนฺตรฆเร ฉฑฺเฑนฺโต เอกํ อาปตฺติํ อาปชฺชติ, ทุกฺกฏํ.}

\prose{อนาทริยํ ปฏิจฺจ ฉตฺตปาณิสฺส ธมฺมํ เทเสนฺโต เอกํ อาปตฺติํ อาปชฺชติ, ทุกฺกฏํ.}

\prose{อนาทริยํ ปฏิจฺจ ทณฺฑปาณิสฺส ธมฺมํ เทเสนฺโต เอกํ อาปตฺติํ อาปชฺชติ, ทุกฺกฏํ.}

\prose{อนาทริยํ ปฏิจฺจ สตฺถปาณิสฺส ธมฺมํ เทเสนฺโต เอกํ อาปตฺติํ อาปชฺชติ, ทุกฺกฏํ.}

\prosewithclosermidruled{อนาทริยํ ปฏิจฺจ อาวุธปาณิสฺส ธมฺมํ เทเสนฺโต เอกํ อาปตฺติํ อาปชฺชติ, ทุกฺกฏํ.}{สุรุสุรุวคฺโค ฉฏฺโฐ.}
\csromanmatikaanchor{mtk.54}
\csromantocmarkref{section}{7. ปาทุกวคฺค}{mtk.54}
\bookmark[dest={mtk.54},level=1]{7. ปาทุกวคฺค}
\csromanheader{1}{7. \textbf{ปาทุกวคฺค}}
\proseitem{181}{อนาทริยํ ปฏิจฺจ ปาทุการุฬฺหสฺส ธมฺมํ เทเสนฺโต เอกํ อาปตฺติํ อาปชฺชติ, ทุกฺกฏํ.}

\prose{อนาทริยํ ปฏิจฺจ อุปาหนา\-รุฬฺหสฺส ธมฺมํ เทเสนฺโต เอกํ อาปตฺติํ อาปชฺชติ, ทุกฺกฏํ.}

\prose{อนาทริยํ ปฏิจฺจ ยานคตสฺส ธมฺมํ เทเสนฺโต เอกํ อาปตฺติํ อาปชฺชติ, ทุกฺกฏํ.}

\prose{อนาทริยํ ปฏิจฺจ สยนคตสฺส ธมฺมํ เทเสนฺโต เอกํ อาปตฺติํ อาปชฺชติ, ทุกฺกฏํ.}

\prose{อนาทริยํ ปฏิจฺจ ปลฺลตฺถิกาย นิสินฺนสฺส ธมฺมํ เทเสนฺโต เอกํ อาปตฺติํ อาปชฺชติ, ทุกฺกฏํ.}

\prose{อนาทริยํ ปฏิจฺจ เวฐิตสีสสฺส ธมฺมํ เทเสนฺโต เอกํ อาปตฺติํ อาปชฺชติ, ทุกฺกฏํ.}

\prose{อนาทริยํ ปฏิจฺจ โอคุณฺฐิต\-สีสสฺส ธมฺมํ เทเสนฺโต เอกํ อาปตฺติํ อาปชฺชติ, ทุกฺกฏํ}

\prose{\csromanfolio{79}อนาทริยํ ปฏิจฺจ ฉมายํ นิสีทิตฺวา อาสเน นิสินฺนสฺส ธมฺมํ เทเสนฺโต เอกํ อาปตฺติํ อาปชฺชติ, ทุกฺกฏํ.}

\prose{อนาทริยํ ปฏิจฺจ นีเจ อาสเน นิสีทิตฺวา อุจฺเจ อาสเน นิสินฺนสฺส ธมฺมํ เทเสนฺโต เอกํ อาปตฺติํ อาปชฺชติ, ทุกฺกฏํ.}

\prose{อนาทริยํ ปฏิจฺจ ฐิโต นิสินฺนสฺส ธมฺมํ เทเสนฺโต เอกํ อาปตฺติํ อาปชฺชติ, ทุกฺกฏํ.}

\prose{อนาทริยํ ปฏิจฺจ ปจฺฉโต คจฺฉนฺโต ปุรโต คจฺฉนฺตสฺส ธมฺมํ เทเสนฺโต เอกํ อาปตฺติํ อาปชฺชติ, ทุกฺกฏํ.}

\prose{อนาทริยํ ปฏิจฺจ อุปฺปเถน คจฺฉนฺโต ปเถน คจฺฉนฺตสฺส ธมฺมํ เทเสนฺโต เอกํ อาปตฺติํ อาปชฺชติ, ทุกฺกฏํ.}

\prose{อนาทริยํ ปฏิจฺจ ฐิโต อุจฺจารํ วา ปสฺสาวํ วา กโรนฺโต เอกํ อาปตฺติํ อาปชฺชติ, ทุกฺกฏํ.}

\prosewithclosermidruled{อนาทริยํ ปฏิจฺจ หริเต อุจฺจารํ วา ปสฺสาวํ วา เขฬํ วา กโรนฺโต กติ อาปตฺติโย อาปชฺชติ. อนาทริยํ ปฏิจฺจ อุทเก อุจฺจารํ วา ปสฺสาวํ วา เขฬํ วา กโรนฺโต เอกํ อาปตฺติํ อาปชฺชติ, ทุกฺกฏํ. อนาทริยํ ปฏิจฺจ อุทเก อุจฺจารํ วา ปสฺสาวํ วา เขฬํ วา กโรนฺโต อิมํ เอกํ อาปตฺติํ อาปชฺชติ.}{ปาทุกวคฺโค สตฺตโม.}
\nitthitam{เสขิยา นิฏฺฐิตา.}
\nitthitamruled{กตาปตฺติวาโร นิฏฺฐิโต ทุติโย.}
\csromanmatikaanchor{mtk.55}
\csromantocmarkref{chapter}{3. วิปตฺติวาร}{mtk.55}
\bookmark[dest={mtk.55},level=0]{3. วิปตฺติวาร}
\chapterhead{3. \textbf{วิปตฺติวาร}}
\proseitem{182}{เมถุนํ ธมฺมํ ปฏิเสวนฺตสฺส อาปตฺติโย จตุนฺนํ วิปตฺตีนํ กติ วิปตฺติโย ภชนฺติ. เมถุนํ ธมฺมํ ปฏิเสวนฺตสฺส อาปตฺติโย จตุนฺนํ วิปตฺตีนํ เทฺว วิปตฺติโย ภชนฺติ, สิยา สีลวิปตฺติํ, สิยา อาจารวิปตฺติํ ฯเปฯ.}

\prosewithcloserruled{\csromanfolio{80}อนาทริยํ ปฏิจฺจ อุทเก อุจฺจารํ วา ปสฺสาวํ วา เขฬํ วา กโรนฺตสฺส อาปตฺติ จตุนฺนํ วิปตฺตีนํ กติ วิปตฺติโย ภชติ. อนาทริยํ ปฏิจฺจ อุทเก อุจฺจารํ วา ปสฺสาวํ วา เขฬํ วา กโรนฺตสฺส อาปตฺติ จตุนฺนํ วิปตฺตีนํ เอกํ วิปตฺติํ ภชติ, อาจารวิปตฺติํ.}{วิปตฺติวาโร นิฏฺฐิโต ตติโย.}
\csromanmatikaanchor{mtk.56}
\csromantocmarkref{chapter}{4. สงฺคหิตวาร}{mtk.56}
\bookmark[dest={mtk.56},level=0]{4. สงฺคหิตวาร}
\chapterhead{4. \textbf{สงฺคหิตวาร}}
\proseitem{183}{เมถุนํ ธมฺมํ ปฏิเสวนฺตสฺส อาปตฺติโย สตฺตนฺนํ อาปตฺติ\-กฺขนฺธานํ กติหิ อาปตฺติ\-กฺขนฺเธหิ สงฺคหิตา. เมถุนํ ธมฺมํ ปฏิเสวนฺตสฺส อาปตฺติโย สตฺตนฺนํ อาปตฺติ\-กฺขนฺธานํ ตีหิ อาปตฺติ\-กฺขนฺเธหิ สงฺคหิตา, สิยา ปาราชิกาปตฺติ\-กฺขนฺเธน, สิยา ถุลฺลจฺจยา\-ปตฺติกฺขนฺเธน, สิยา ทุกฺกฏาปตฺติ\-กฺขนฺเธน ฯเปฯ.}

\prosewithcloserruled{อนาทริยํ ปฏิจฺจ อุทเก อุจฺจารํ วา ปสฺสาวํ วา เขฬํ วา กโรนฺตสฺส อาปตฺติ สตฺตนฺนํ อาปตฺติ\-กฺขนฺธานํ กติหิ อาปตฺติ\-กฺขนฺเธหิ สงฺคหิตา. อนาทริยํ ปฏิจฺจ อุทเก อุจฺจารํ วา ปสฺสาวํ วา เขฬํ วา กโรนฺตสฺส อาปตฺติ สตฺตนฺนํ อาปตฺติ\-กฺขนฺธานํ เอเกน อาปตฺติ\-กฺขนฺเธน สงฺคหิตา, ทุกฺกฏาปตฺติ\-กฺขนฺเธน.}{สงฺคหิตวาโร นิฏฺฐิโต จตุตฺโถ.}
\csromanmatikaanchor{mtk.57}
\csromantocmarkref{chapter}{5. สมุฏฺฐานวาร}{mtk.57}
\bookmark[dest={mtk.57},level=0]{5. สมุฏฺฐานวาร}
\chapterhead{5. \textbf{สมุฏฺฐานวาร}}
\proseitem{184}{เมถุนํ ธมฺมํ ปฏิเสวนฺตสฺส อาปตฺติโย ฉนฺนํ อาปตฺติ\-สมุฏฺฐานานํ กติหิ สมุฏฺฐาเนหิ สมุฏฺฐนฺติ\csromansharedfootnote{ada65ed11f68b1b9}{สมุฏฺฐหนฺติ (สี, สฺยา)}. เมถุนํ ธมฺมํ ปฏิเสวนฺตสฺส อาปตฺติโย ฉนฺนํ อาปตฺติ\-สมุฏฺฐานานํ เอเกน สมุฏฺฐาเนน สมุฏฺฐนฺติ, กายโต จ จิตฺตโต จ สมุฏฺฐนฺติ, น วาจโต ฯเปฯ.}

\prosewithcloserruled{\csromanfolio{81}อนาทริยํ ปฏิจฺจ อุทเก อุจฺจารํ วา ปสฺสาวํ วาเขฬํ วา กโรนฺตสฺส อาปตฺติ ฉนฺนํ อาปตฺติ\-สมุฏฺฐานานํ กติหิ สมุฏฺฐาเนหิ สมุฏฺฐาติ. อนาทริยํ ปฏิจฺจ อุทเก อุจฺจารํ วา ปสฺสาวํ วา เขฬํ วา กโรนฺตสฺส อาปตฺติ ฉนฺนํ อาปตฺติ\-สมุฏฺฐานานํ เอเกน สมุฏฺฐาเนน สมุฏฺฐาติ, กายโต จ จิตฺตโต จ สมุฏฺฐาติ, น วาจโต.}{สมุฏฺฐานวาโร นิฏฺฐิโต ปญฺจโม.}
\csromanmatikaanchor{mtk.58}
\csromantocmarkref{chapter}{6. อธิกรณวาร}{mtk.58}
\bookmark[dest={mtk.58},level=0]{6. อธิกรณวาร}
\chapterhead{6. \textbf{อธิกรณวาร}}
\proseitem{185}{เมถุนํ ธมฺมํ ปฏิเสวนฺตสฺส อาปตฺติโย จตุนฺนํ อธิกรณานํ กตมํ อธิกรณํ. เมถุนํ ธมฺมํ ปฏิเสวนฺตสฺส อาปตฺติโย จตุนฺนํ อธิกรณานํ อาปตฺตา\-ธิกรณํ ฯเปฯ.}

\prosewithcloserruled{อนาทริยํ ปฏิจฺจ อุทเก อุจฺจารํ วา ปสฺสาวํ วา เขฬํ วา กโรนฺตสฺส อาปตฺติ จตุนฺนํ อธิกรณานํ กตมํ อธิกรณํ. อนาทริยํ ปฏิจฺจ อุทเก อุจฺจารํ วา ปสฺสาวํ วา เขฬํ วา กโรนฺตสฺส อาปตฺติ จตุนฺนํ อธิกรณานํ อาปตฺตา\-ธิกรณํ.}{อธิกรณวาโร นิฏฺฐิโต ฉฏฺโฐ.}
\csromanmatikaanchor{mtk.59}
\csromantocmarkref{chapter}{7. สมถวาร}{mtk.59}
\bookmark[dest={mtk.59},level=0]{7. สมถวาร}
\chapterhead{7. \textbf{สมถวาร}}
\proseitem{186}{เมถุนํ ธมฺมํ ปฏิเสวนฺตสฺส อาปตฺติโย สตฺตนฺนํ สมถานํ กติหิ สมเถหิ สมฺมนฺติ. เมถุนํ ธมฺมํ ปฏิเสวนฺตสฺส อาปตฺติโย สตฺตนฺนํ สมถานํ ตีหิ สมเถหิ สมฺมนฺติ, สิยา สมฺมุขา\-วินเยน จ ปฏิญฺญาต\-กรเณน จ, สิยา สมฺมุขา\-วินเยน จ ติณวตฺถารเกน จ ฯเปฯ.}

\prose{อนาทริยํ ปฏิจฺจ อุทเก อุจฺจารํ วา ปสฺสาวํ วา เขฬํ วา กโรนฺตสฺส อาปตฺติ สตฺตนฺนํ สมถานํ กติหิ สมเถหิ สมฺมติ. อนาทริยํ ปฏิจฺจ อุทเก อุจฺจารํ วา ปสฺสาวํ วา เขฬํ วา กโรนฺตสฺส อาปตฺติ สตฺตนฺนํ สมถานํ ตีหิ สมเถหิ สมฺมติ, สิยา สมฺมุขา\-วินเยน จ ปฏิญฺญาต\-กรเณน จ, สิยา สมฺมุขา\-วินเยน จ ติณวตฺถารเกน จ.}

\vspace{\tipitakaparskip}
\csromancenter{สมถวาโร นิฏฺฐิโต สตฺตโม.}
\csromanmatikaanchor{mtk.60}
\csromantocmarkref{chapter}{8. สมุจฺจยวาร}{mtk.60}
\bookmark[dest={mtk.60},level=0]{8. สมุจฺจยวาร}
\chapterheadpageread{8. \textbf{สมุจฺจยวาร}}
\proseitem{187}{\csromanfolio{82}เมถุนํ ธมฺมํ ปฏิเสวนฺโต กติ อาปตฺติโย อาปชฺชติ. เมถุนํ ธมฺมํ ปฏิเสวนฺโต ติสฺโส อาปตฺติโย อาปชฺชติ. อกฺขายิเต สรีเร เมถุนํ ธมฺมํ ปฏิเสวติ, อาปตฺติ ปาราชิกสฺส. เยภุยฺเยน ขายิเต สรีเร เมถุนํ ธมฺมํ ปฏิเสวติ, อาปตฺติ ถุลฺลจฺจยสฺส. วฏฺฏกเต มุเข อจฺฉุปนฺตํ องฺคชาตํ ปเวเสติ, อาปตฺติ ทุกฺกฏสฺส. เมถุนํ ธมฺมํ ปฏิเสวนฺโต อิมา ติสฺโส อาปตฺติโย อาปชฺชติ.}

\prose{ตา อาปตฺติโย จตุนฺนํ วิปตฺตีนํ กติ วิปตฺติโย ภชนฺติ, สตฺตนฺนํ อาปตฺติ\-กฺขนฺธานํ กติหิ อาปตฺติ\-กฺขนฺเธหิ สงฺคหิตา, ฉนฺนํ อาปตฺติ\-สมุฏฺฐานานํ กติหิ สมุฏฺฐาเนหิ สมุฏฺฐนฺติ, จตุนฺนํ อธิกรณานํ กตมํ อธิกรณํ, สตฺตนฺนํ สมถานํ กติหิ สมเถหิ สมฺมนฺติ. ตา อาปตฺติโย จตุนฺนํ วิปตฺตีนํ เทฺว วิปตฺติโย ภชนฺติ, สิยา สีลวิปตฺติํ, สิยา อาจารวิปตฺติํ. สตฺตนฺนํ อาปตฺติ\-กฺขนฺธานํ ตีหิ อาปตฺติ\-กฺขนฺเธหิ สงฺคหิตา, สิยา ปาราชิกาปตฺติ\-กฺขนฺเธน, สิยา ถุลฺลจฺจยา\-ปตฺติกฺขนฺเธน, สิยา ทุกฺกฏาปตฺติ\-กฺขนฺเธน. ฉนฺนํ อาปตฺติ\-สมุฏฺฐานานํ เอเกน สมุฏฺฐาเนน สมุฏฺฐนฺติ, กายโต จ จิตฺตโต จ สมุฏฺฐนฺติ, น วาจโต. จตุนฺนํ อธิกรณานํ อาปตฺตา\-ธิกรณํ. สตฺตนฺนํ สมถานํ ตีหิ สมเถหิ สมฺมนฺติ, สิยา สมฺมุขา\-วินเยน จ ปฏิญฺญาต\-กรเณน จ, สิยา สมฺมุขา\-วินเยน จ ติณวตฺถารเกน จ ฯเปฯ.}

\prose{อนาทริยํ ปฏิจฺจ อุทเก อุจฺจารํ วา ปสฺสาวํ วา เขฬํ วา กโรนฺโต กติ อาปตฺติโย อาปชฺชติ. อนาทริยํ ปฏิจฺจ อุทเก อุจฺจารํ วา ปสฺสาวํ วา เขฬํ วา กโรนฺโต เอกํ อาปตฺติํ อาปชฺชติ, ทุกฺกฏํ. อนาทริยํ ปฏิจฺจ อุทเก อุจฺจารํ วา ปสฺสาวํ วา เขฬํ วา กโรนฺโต อิมํ เอกํ อาปตฺติํ อาปชฺชติ.}

\prose{สา อาปตฺติ จตุนฺนํ วิปตฺตีนํ กติ วิปตฺติโย ภชติ, สตฺตนฺนํ อาปตฺติ\-กฺขนฺธานํ กติหิ อาปตฺติ\-กฺขนฺเธหิ สงฺคหิตา, ฉนฺนํ อาปตฺติ\-สมุฏฺฐานานํ กติหิ สมุฏฺฐาเนหิ สมุฏฺฐาติ, จตุนฺนํ อธิกรณานํ กตมํ อธิกรณํ, สตฺตนฺนํ สมถานํ กติหิ สมเถหิ สมฺมติ. สา อาปตฺติ จตุนฺนํ วิปตฺตีนํ เอกํ วิปตฺติํ ภชติ, อาจารวิปตฺติํ. สตฺตนฺนํ อาปตฺติ\-กฺขนฺธานํ เอเกน อาปตฺติ\-กฺขนฺเธน สงฺคหิตา, ทุกฺกฏาปตฺติ\-กฺขนฺเธน. \csromanfolio{83}ฉนฺนํ อาปตฺติ\-สมุฏฺฐานานํ เอเกน สมุฏฺฐาเนน สมุฏฺฐาติ, กายโต จ จิตฺตโต จ สมุฏฺฐาติ, น วาจโต. จตุนฺนํ อธิกรณานํ อาปตฺตา\-ธิกรณํ. สตฺตนฺนํ สมถานํ ตีหิ สมเถหิ สมฺมติ, สิยา สมฺมุขา\-วินเยน จ ปฏิญฺญาต\-กรเณน จ, สิยา สมฺมุขา\-วินเยน จ ติณวตฺถารเกน จ.}

\vspace{\tipitakaparskip}
\csromancenter{สมุจฺจยวาโร นิฏฺฐิโต อฏฺฐโม.}
\csromancenter{อิเม อฏฺฐ วารา สชฺฌาย\-มคฺเคน ลิขิตา.}
\csromancenter{ตสฺสุทฺทานํ.}
\setlength{\csromangathapagewidth}{0pt}
\setlength{\csromangathawakwidth}{0pt}
\csromangathameasuregroup{กตฺถปญฺญตฺติ กติ จ,}{\csromangathabat{กตฺถปญฺญตฺติ กติ จ,}{วิปตฺติสงฺคเหน จ.}}
\csromangathasetleft{กตฺถปญฺญตฺติ กติ จ,}
\csromangathagroup{\csromangathastanza{\csromangathabat{กตฺถปญฺญตฺติ กติ จ,}{วิปตฺติสงฺคเหน จ.}}}

\vspace{\tipitakaparskip}
\csromancenter{สมุฏฺฐานา\-ธิกรณา, สมโถ สมุจฺจเยน จาติ.}
\csromansectionrule
\csromanmatikaanchor{mtk.61}
\csromanmatikaanchor{mtk.62}
\csromantocmarknopage{section}{1. กตฺถปญฺญตฺติวาร}{mtk.61}
\bookmark[dest={mtk.61},level=1]{1. กตฺถปญฺญตฺติวาร}
\csromantocmarkref{chapter}{1. ปาราชิกกณฺฑ}{mtk.62}
\bookmark[dest={mtk.62},level=0]{1. ปาราชิกกณฺฑ}
\chapterhead{1. กตฺถ\-ปญฺญตฺติวาร 1. ปาราชิกกณฺฑ}
\proseitem{188}{ยํ เตน ภควตา ชานตา ปสฺสตา อรหตา สมฺมา\-สมฺพุทฺเธน เมถุนํ ธมฺมํ ปฏิเสวน\-ปจฺจยา ปาราชิกํ กตฺถ ปญฺญตฺตํ. กํ อารพฺภ. กิสฺมิํ วตฺถุสฺมิํ ฯเปฯ. เกนาภตนฺติ.}

\prose{ยํ เตน ภควตา ชานตา ปสฺสตา อรหตา สมฺมา\-สมฺพุทฺเธน เมถุนํ ธมฺมํ ปฏิเสวน\-ปจฺจยา ปาราชิกํ กตฺถ ปญฺญตฺตนฺติ เวสาลิยํ ปญฺญตฺตํ. กํ อารพฺภาติ สุทินฺนํ กลนฺทปุตฺตํ อารพฺภ. กิสฺมิํ วตฺถุสฺมินฺติ สุทินฺโน กลนฺทปุตฺโต ปุราณ\-ทุติยิกาย เมถุนํ ธมฺมํ ปฏิเสวิ, ตสฺมิํ วตฺถุสฺมิํ. อตฺถิ ตตฺถ ปญฺญตฺติ อนุปญฺญตฺติ อนุปฺปนฺนปญฺญตฺตีติ เอกา ปญฺญตฺติ, เทฺว อนุปญฺญตฺติโย, อนุปฺปนฺนปญฺญตฺติ ตสฺมิํ นตฺถิ. สพฺพตฺถ ปญฺญตฺติ ปเทสปญฺญตฺตีติ สพฺพตฺถ\-ปญฺญตฺติ. สาธารณ\-ปญฺญตฺติ อสาธารณปญฺญตฺตีติ สาธารณ\-ปญฺญตฺติ. เอกโตปญฺญตฺติ อุภโตปญฺญตฺตีติ อุภโตปญฺญตฺติ. ปญฺจนฺนํ ปาติโมกฺขุ\-ทฺเทสานํ กตฺโถคธํ กตฺถ ปริยาปนฺนนฺติ นิทาโนคธํ นิทาน\-ปริยาปนฺนํ. กตเมน อุทฺเทเสน อุทฺเทสํ อาคจฺฉตีติ ทุติเยน อุทฺเทเสน อุทฺเทสํ อาคจฺฉติ. จตุนฺนํ วิปตฺตีนํ กตมา วิปตฺตีติ สีลวิปตฺติ. สตฺตนฺนํ อาปตฺติ\-กฺขนฺธานํ กตโม อาปตฺติ\-กฺขนฺโธติ ปาราชิกาปตฺติ\-กฺขนฺโธ. ฉนฺนํ อาปตฺติ\-สมุฏฺฐานานํ กติหิ สมุฏฺฐาเนหิ สมุฏฺฐาตีติ เอเกน สมุฏฺฐาเนน สมุฏฺฐาติ, \csromanfolio{84}กายโต จ จิตฺตโต จ สมุฏฺฐาติ, น วาจโต ฯเปฯ, เกนาภตนฺติ ปรมฺปราภตํ.}

\setlength{\csromangathapagewidth}{0pt}
\setlength{\csromangathawakwidth}{0pt}
\csromangathameasuregroup{อุปาลิ ทาสโก เจว, \\ โมคฺคลิปุตฺเตน ปญฺจมา, \\ เอเต นาคา มหาปญฺญา, \\ วินยํ ทีเป ปกาเสสุํ,}{\csromangathabat{อุปาลิ ทาสโก เจว,}{โสณโก สิคฺคโว ตถา.} \\ \csromangathabat{โมคฺคลิปุตฺเตน ปญฺจมา,}{เอเต ชมฺพุสิริวฺหเย ฯเปฯ.} \\ \csromangathabat{เอเต นาคา มหาปญฺญา,}{วินยญฺญู มคฺคโกวิทา.} \\ \csromangathabat{วินยํ ทีเป ปกาเสสุํ,}{ปิฏกํ ตมฺพปณฺณิยาติ.}}
\csromangathasetleft{อุปาลิ ทาสโก เจว, \\ โมคฺคลิปุตฺเตน ปญฺจมา, \\ เอเต นาคา มหาปญฺญา, \\ วินยํ ทีเป ปกาเสสุํ,}
\csromangathagroup{\csromangathastanza{\csromangathabat{อุปาลิ ทาสโก เจว,}{โสณโก สิคฺคโว ตถา.} \\ \csromangathabat{โมคฺคลิปุตฺเตน ปญฺจมา,}{เอเต ชมฺพุสิริวฺหเย ฯเปฯ.}}\csromangathastanza{\csromangathabat{เอเต นาคา มหาปญฺญา,}{วินยญฺญู มคฺคโกวิทา.} \\ \csromangathabat{วินยํ ทีเป ปกาเสสุํ,}{ปิฏกํ ตมฺพปณฺณิยาติ.}}}

\proseitem{189}{ยํ เตน ภควตา ชานตา ปสฺสตา อรหตา สมฺมา\-สมฺพุทฺเธน อทินฺนํ อาทิยนปจฺจยา ปาราชิกํ กตฺถ ปญฺญตฺตนฺติ ราชคเห ปญฺญตฺตํ. กํ อารพฺภาติ ธนิยํ กุมฺภการ\-ปุตฺตํ อารพฺภ. กิสฺมิํ วตฺถุสฺมินฺติ ธนิโย กุมฺภการ\-ปุตฺโต รญฺโญ ทารูนิ อทินฺนํ อาทิยิ, ตสฺมิํ วตฺถุสฺมิํ. เอกา ปญฺญตฺติ, เอกา อนุปญฺญตฺติ. ฉนฺนํ อาปตฺติ\-สมุฏฺฐานานํ ตีหิ สมุฏฺฐาเนหิ สมุฏฺฐาติ, สิยา กายโต จ จิตฺตโต จ สมุฏฺฐาติ, น วาจโต, สิยา วาจโต จ จิตฺตโต จ สมุฏฺฐาติ, น กายโต, สิยา กายโต จ วาจโต จ จิตฺตโต จ}

\proseitem{190}{สญฺจิจฺจ มนุสฺส\-วิคฺคหํ ชีวิตา โวโรปน\-ปจฺจยา ปาราชิกํ กตฺถ ปญฺญตฺตนฺติ เวสาลิยํ ปญฺญตฺตํ. กํ อารพฺภาติ สมฺพหุเล ภิกฺขู อารพฺภ. กิสฺมิํ วตฺถุสฺมินฺติ สมฺพหุลา ภิกฺขู อญฺญมญฺญํ ชีวิตา โวโรเปสุํ, ตสฺมิํ วตฺถุสฺมิํ. เอกา ปญฺญตฺติ, เอกา อนุปญฺญตฺติ. ฉนฺนํ อาปตฺติ\-สมุฏฺฐานานํ ตีหิ สมุฏฺฐาเนหิ สมุฏฺฐาติ, สิยา กายโต จ จิตฺตโต จ สมุฏฺฐาติ, น วาจโต, สิยา วาจโต จ จิตฺตโต จ สมุฏฺฐาติ, น กายโต, สิยา กายโต จ วาจโต จ จิตฺตโต จ สมุฏฺฐาติ ฯเปฯ.}

\proseitemwithcloser{191}{อสนฺตํ อภูตํ อุตฺตริ\-มนุสฺสธมฺมํ อุลฺลปน\-ปจฺจยา ปาราชิกํ กตฺถ ปญฺญตฺตนฺติ เวสาลิยํ ปญฺญตฺตํ. กํ อารพฺภาติ วคฺคุมุทา\-ตีริเย ภิกฺขู อารพฺภ. กิสฺมิํ วตฺถุสฺมินฺติ วคฺคุมุทา\-ตีริยา ภิกฺขู คิหีนํ อญฺญมญฺญสฺส อุตฺตริ\-มนุสฺสธมฺมสฺส วณฺณํ ภาสิํสุ, ตสฺมิํ วตฺถุสฺมิํ. เอกา ปญฺญตฺติ, เอกา อนุปญฺญตฺติ. ฉนฺนํ อาปตฺติ\-สมุฏฺฐานานํ ตีหิ สมุฏฺฐาเนหิ สมุฏฺฐาติ, สิยา กายโต จ จิตฺตโต จ สมุฏฺฐาติ, น วาจโต, สิยา วาจโต จ จิตฺตโต จ สมุฏฺฐาติ, น กายโต, สิยา กายโต จ วาจโต จ จิตฺตโต จ สมุฏฺฐาติ ฯเปฯ.}{จตฺตาโร ปาราชิกา นิฏฺฐิตา.}
\csromanmatikaanchor{mtk.63}
\csromantocmarkref{chapter}{2. สํฆาทิเสสกณฺฑาทิ}{mtk.63}
\bookmark[dest={mtk.63},level=0]{2. สํฆาทิเสสกณฺฑาทิ}
\chapterheadpageread{2. \textbf{สํฆาทิเสสกณฺฑาทิ}}
\proseitem{192}{\csromanfolio{85}ยํ เตน ภควตา ชานตา ปสฺสตา อรหตา สมฺมา\-สมฺพุทฺเธน อุปกฺกมิตฺวา อสุจิํ โมจนปจฺจยา สํฆาทิเสโส กตฺถ ปญฺญตฺโต. กํ อารพฺภ. กิสฺมิํ วตฺถุสฺมิํ ฯเปฯ เกนาภตนฺติ.}

\prose{ยํ เตน ภควตา ชานตา ปสฺสตา อรหตา สมฺมา\-สมฺพุทฺเธน อุปกฺกมิตฺวา อสุจิํ โมจนปจฺจยา สํฆาทิเสโส กตฺถ ปญฺญตฺโตติ สาวตฺถิยํ ปญฺญตฺโต. กํ อารพฺภาติ อายสฺมนฺตํ เสยฺยสกํ อารพฺภ. กิสฺมิํ วตฺถุสฺมินฺติ อายสฺมา เสยฺยสโก อุปกฺกมิตฺวา อสุจิํ โมเจสิ, ตสฺมิํ วตฺถุสฺมิํ. อตฺถิ ตตฺถ ปญฺญตฺติ อนุปญฺญตฺติ อนุปฺปนฺนปญฺญตฺตีติ เอกา ปญฺญตฺติ, เอกา อนุปญฺญตฺติ, อนุปฺปนฺนปญฺญตฺติ ตสฺมิํ นตฺถิ. สพฺพตฺถ\-ปญฺญตฺติ ปเทสปญฺญตฺตีติ สพฺพตฺถ\-ปญฺญตฺติ. สาธารณ\-ปญฺญตฺติ อสาธารณปญฺญตฺตีติ อสาธารณ\-ปญฺญตฺติ. เอกโตปญฺญตฺติ อุภโตปญฺญตฺตีติ เอกโตปญฺญตฺติ. ปญฺจนฺนํ ปาติโมกฺขุ\-ทฺเทสานํ กตฺโถคธํ กตฺถ ปริยาปนฺนนฺติ นิทาโนคธํ นิทาน\-ปริยาปนฺนํ. กตเมน อุทฺเทเสน อุทฺเทสํ อาคจฺฉตีติ ตติเยน อุทฺเทเสน อุทฺเทสํ อาคจฺฉติ. จตุนฺนํ วิปตฺตีนํ กตมา วิปตฺตีติ สีลวิปตฺติ. สตฺตนฺนํ อาปตฺติ\-กฺขนฺธานํ กตโม อาปตฺติ\-กฺขนฺโธติ สํฆาทิเสสา\-ปตฺติกฺขนฺโธ. ฉนฺนํ อาปตฺติ\-สมุฏฺฐานานํ กติหิ สมุฏฺฐาเนหิ สมุฏฺฐาตีติ เอเกน สมุฏฺฐาเนน สมุฏฺฐาติ, กายโต จ จิตฺตโต จ สมุฏฺฐาติ น วาจโต ฯเปฯ เกนาภตนฺติ ปรมฺปราภตํ.}

\setlength{\csromangathapagewidth}{0pt}
\setlength{\csromangathawakwidth}{0pt}
\csromangathameasuregroup{อุปาลิ ทาสโก เจว, \\ โมคฺคลิปุตฺเตน ปญฺจมา, \\ เอเต นาคา มหาปญฺญา, \\ วินยํ ทีเป ปกาเสสุํ,}{\csromangathabat{อุปาลิ ทาสโก เจว,}{โสณโก สิคฺคโว ตถา.} \\ \csromangathabat{โมคฺคลิปุตฺเตน ปญฺจมา,}{เอเต ชมฺพุสิริวฺหเย ฯเปฯ.} \\ \csromangathabat{เอเต นาคา มหาปญฺญา,}{วินยญฺญู มคฺคโกวิทา.} \\ \csromangathabat{วินยํ ทีเป ปกาเสสุํ,}{ปิฏกํ ตมฺพปณฺณิยาติ.}}
\csromangathasetleft{อุปาลิ ทาสโก เจว, \\ โมคฺคลิปุตฺเตน ปญฺจมา, \\ เอเต นาคา มหาปญฺญา, \\ วินยํ ทีเป ปกาเสสุํ,}
\csromangathagroup{\csromangathastanza{\csromangathabat{อุปาลิ ทาสโก เจว,}{โสณโก สิคฺคโว ตถา.} \\ \csromangathabat{โมคฺคลิปุตฺเตน ปญฺจมา,}{เอเต ชมฺพุสิริวฺหเย ฯเปฯ.}}\csromangathastanza{\csromangathabat{เอเต นาคา มหาปญฺญา,}{วินยญฺญู มคฺคโกวิทา.} \\ \csromangathabat{วินยํ ทีเป ปกาเสสุํ,}{ปิฏกํ ตมฺพปณฺณิยาติ.}}}

\prose{มาตุคาเมน สทฺธิํ กายสํสคฺคํ สมาปชฺชน\-ปจฺจยา สํฆาทิเสโส กตฺถ ปญฺญตฺโตติ สาวตฺถิยํ ปญฺญตฺโต. กํ อารพฺภาติ อายสฺมนฺตํ อุทายิํ อารพฺภ. กิสฺมิํ วตฺถุสฺมินฺติ อายสฺมา อุทายี มาตุคาเมน สทฺธิํ กายสํสคฺคํ สมาปชฺชิ, ตสฺมิํ วตฺถุสฺมิํ. เอกา ปญฺญตฺติ. ฉนฺนํ อาปตฺติ\-สมุฏฺฐานานํ เอเกน สมุฏฺฐาเนน สมุฏฺฐาติ, กายโต จ จิตฺตโต จ สมุฏฺฐาติ, น วาจโต ฯเปฯ.}

\prose{\csromanfolio{86}มาตุคามํ ทุฏฺฐุลฺลาหิ วาจาหิ โอภาสน\-ปจฺจยา สํฆาทิเสโส กตฺถ ปญฺญตฺโตติ สาวตฺถิยํ ปญฺญตฺโต. กํ อารพฺภาติ อายสฺมนฺตํ อุทายิํ อารพฺภ. กิสฺมิํ วตฺถุสฺมินฺติ อายสฺมา อุทายี มาตุคามํ ทุฏฺฐาลฺลาหิ วาจาหิ โอภาสิ, ตสฺมิํ วตฺถุสฺมิํ. เอกา ปญฺญตฺติ. ฉนฺนํ อาปตฺติ\-สมุฏฺฐานานํ ตีหิ สมุฏฺฐาเนหิ สมุฏฺฐาติ, สิยา กายโต จ จิตฺตโต จ สมุฏฺฐาติ, น วาจโต, สิยา วาจโต จ จิตฺตโต จ สมุฏฺฐาติ, น กายโต, สิยา กายโต จ วาจโต จ}

\prose{มาตุคามสฺส สนฺติเก อตฺต\-กาม\-ปาริจริยาย วณฺณํ ภาสนปจฺจยา สํฆาทิเสโส กตฺถ ปญฺญตฺโตติ สาวตฺถิยํ ปญฺญตฺโต. กํ อารพฺภาติ อายสฺมนฺตํ อุทายิํ อารพฺภ. กิสฺมิํ วตฺถุสฺมินฺติ อายสฺมา อุทายี มาตุคามสฺส สนฺติเก อตฺต\-กาม\-ปาริจริยาย วณฺณํ ภาสิ, ตสฺมิํ วตฺถุสฺมิํ. เอกา ปญฺญตฺติ. ฉนฺนํ อาปตฺติ\-สมุฏฺฐานานํ ตีหิ สมุฏฺฐาเนหิ}

\prose{สญฺจริตฺตํ สมาปชฺชน\-ปจฺจยา สํฆาทิเสโส กตฺถ ปญฺญตฺโตติ สาวตฺถิยํ ปญฺญตฺโต. กํ อารพฺภาติ อายสฺมนฺตํ อุทายิํ อารพฺภ. กิสฺมิํ วตฺถุสฺมินฺติ อายสฺมา อุทายี สญฺจริตฺตํ สมาปชฺชิ, ตสฺมิํ วตฺถุสฺมิํ. เอกา ปญฺญตฺติ, เอกา อนุปญฺญตฺติ. ฉนฺนํ อาปตฺติ\-สมุฏฺฐานานํ ฉหิ สมุฏฺฐาเนหิ สมุฏฺฐาติ, สิยา กายโต สมุฏฺฐาติ, น วาจโต, น จิตฺติโต, สิยา วาจโต สมุฏฺฐาติ, น กายโต, น จิตฺตโต, สิยา กายโต จ วาจโต จ สมุฏฺฐาติ, น จิตฺตโต, สิยา กายโต จ จิตฺตโต จ สมุฏฺฐาติ, น วาจโต, สิยา วาจโต จ จิตฺตโต จ สมุฏฺฐาติ, น กายโต, สิยา กายโต จ วาจโต จ จิตฺตโต จ}

\prose{สญฺญาจิกาย กุฏิํ การาปน\-ปจฺจยา สํฆาทิเสโส กตฺถ ปญฺญตฺโตติ อาฬวิยํ ปญฺญตฺโต. กํ อารพฺภาติ อาฬวเก ภิกฺขู อารพฺภ. กิสฺมิํ วตฺถุสฺมินฺติ อาฬวกา ภิกฺขู สญฺญาจิกาย กุฏิโย การาเปสุํ, ตสฺมิํ วตฺถุสฺมิํ. เอกา ปญฺญตฺติ. ฉนฺนํ อาปตฺติ\-สมุฏฺฐานานํ ฉหิ สมุฏฺฐาเนหิ สมุฏฺฐาติ ฯเปฯ.}

\prose{\csromanfolio{87}มหลฺลกํ วิหารํ การาปน\-ปจฺจยา สํฆาทิเสโส กตฺถ ปญฺญตฺโตติ โกสมฺพิยํ ปญฺญตฺโต. กํ อารพฺภาติ อายสฺมนฺตํ ฉนฺนํ อารพฺภ. กิสฺมิํ วตฺถุสฺมินฺติ อายสฺมา ฉนฺโน วิหารวตฺถุํ โสเธนฺโต อญฺญตรํ เจติยรุกฺขํ เฉทาเปสิ, ตสฺมิํ วตฺถุสฺมิํ. เอกา ปญฺญตฺติ. ฉนฺนํ อาปตฺติ\-สมุฏฺฐานานํ ฉหิ สมุฏฺฐาเนหิ สมุฏฺฐาติ ฯเปฯ.}

\prose{ภิกฺขุํ อมูลเกน ปาราชิเกน ธมฺเมน อนุทฺธํสน\-ปจฺจยา สํฆาทิเสโส กตฺถ ปญฺญตฺโตติ ราชคเห ปญฺญตฺโต. กํ อารพฺภาติ เมตฺติย\-ภูมชเก ภิกฺขู อารพฺภ. กิสฺมิํ วตฺถุสฺมินฺติ เมตฺติย\-ภูมชกา ภิกฺขู อายสฺมนฺตํ ทพฺพํ มลฺลปุตฺตํ อมูลเกน ปาราชิเกน ธมฺเมน อนุทฺธํเสสุํ, ตสฺมิํ วตฺถุสฺมิํ. เอกา ปญฺญตฺติ. ฉนฺนํ อาปตฺติ\-สมุฏฺฐานานํ ตีหิ สมุฏฺฐาเนหิ สมุฏฺฐาติ ฯเปฯ.}

\prose{ภิกฺขุํ อญฺญภาคิยสฺส อธิกรณสฺส กิญฺจิ เทสํ เลสมตฺตํ อุปาทาย ปาราชิเกน ธมฺเมน อนุทฺธํสน\-ปจฺจยา สํฆาทิเสโส กตฺถ ปญฺญตฺโตติ ราชคเห ปญฺญตฺโต. กํ อารพฺภาติ เมตฺติย\-ภูมชเก ภิกฺขู อารพฺภ. กิสฺมิํ วตฺถุสฺมินฺติ เมตฺติย\-ภูมชกา ภิกฺขู อายสฺมนฺตํ ทพฺพํ มลฺลปุตฺตํ อญฺญภาคิยสฺส อธิกรณสฺส กิญฺจิ เทสํ เลสมตฺตํ อุปาทาย ปาราชิเกน ธมฺเมน อนุทฺธํเสสุํ, ตสฺมิํ วตฺถุสฺมิํ. เอกา ปญฺญตฺติ. ฉนฺนํ อาปตฺติ\-สมุฏฺฐานานํ ตีหิ สมุฏฺฐาเนหิ}

\prose{สํฆเภทกสฺส ภิกฺขุโน ยาวตติยํ สมนุภาสนาย น ปฏินิสฺสชฺชน\-ปจฺจยา สํฆาทิเสโส กตฺถ ปญฺญตฺโตติ ราชคเห ปญฺญตฺโต. กํ อารพฺภาติ เทวทตฺตํ อารพฺภ. กิสฺมิํ วตฺถุสฺมินฺติ เทวทตฺโต สมคฺคสฺส สํฆสฺส เภทาย ปรกฺกมิ, ตสฺมิํ วตฺถุสฺมิํ. เอกา ปญฺญตฺติ. ฉนฺนํ อาปตฺติ\-สมุฏฺฐานานํ เอเกน สมุฏฺฐาเนน สมุฏฺฐาติ, กายโต จ วาจโต จ จิตฺตโต จ สมุฏฺฐาติ ฯเปฯ.}

\prose{เภทกา\-นุวตฺตกานํ ภิกฺขูนํ ยาวตติยํ สมนุภาสนาย น ปฏินิสฺสชฺชน\-ปจฺจยา สํฆาทิเสโส กตฺถ ปญฺญตฺโตติ ราชคเห ปญฺญตฺโต. กํ อารพฺภาติ สมฺพหุเล ภิกฺขู อารพฺภ. กิสฺมิํ วตฺถุสฺมินฺติ สมฺพหุลา ภิกฺขู เทวทตฺตสฺส สํฆเภทาย ปรกฺกม\-นฺตสฺส อนุวตฺตกา อเหสุํ วคฺควาทกา, ตสฺมิํ วตฺถุสฺมิํ. เอกา ปญฺญตฺติ. ฉนฺนํ อาปตฺติ\-สมุฏฺฐานานํ \csromanfolio{88}เอเกน สมุฏฺฐาเนน สมุฏฺฐาติ, กายโต จ วาจโต จ จิตฺตโต จ สมุฏฺฐาติ ฯเปฯ.}

\prose{ทุพฺพจสฺส ภิกฺขุโน ยาวตติยํ สมนุภาสนาย น ปฏินิสฺสชฺชน\-ปจฺจยา สํฆาทิเสโส กตฺถ ปญฺญตฺโตติ โกสมฺพิยํ ปญฺญตฺโต. กํ อารพฺภาติ อายสฺมนฺตํ ฉนฺนํ อารพฺภ. กิสฺมิํ วตฺถุสฺมินฺติ อายสฺมา ฉนฺโน ภิกฺขูหิ สหธมฺมิกํ วุจฺจมาโน อตฺตานํ อวจนียํ อกาสิ, ตสฺมิํ วตฺถุสฺมิํ. เอกา ปญฺญตฺติ. ฉนฺนํ อาปตฺติ\-สมุฏฺฐานานํ เอเกน สมุฏฺฐาเนน สมุฏฺฐาติ, กายโต จ วาจโต จ}

\prose{กุลทูสกสฺส ภิกฺขุโน ยาวตติยํ สมนุภาสนาย น ปฏินิสฺสชฺชน\-ปจฺจยา สํฆาทิเสโส กตฺถ ปญฺญตฺโตติ สาวตฺถิยํ ปญฺญตฺโต. กํ อารพฺภาติ อสฺสชิปุนพฺพสุเก ภิกฺขู อารพฺภ. กิสฺมิํ วตฺถุสฺมินฺติ อสฺสชิ\-ปุนพฺพสุกา ภิกฺขู สํเฆน ปพฺพาชนีย\-กมฺมกตา ภิกฺขู ฉนฺทคามิตา โทสคามิตา โมหคามิตา ภยคามิตา ปาเปสุํ, ตสฺมิํ วตฺถุสฺมิํ. เอกา ปญฺญตฺติ. ฉนฺนํ อาปตฺติ\-สมุฏฺฐานานํ เอเกน สมุฏฺฐาเนน สมุฏฺฐาติ, กายโต จ วาจโต จ จิตฺตโต จ สมุฏฺฐาติ ฯเปฯ.}

\prosewithcloserruled{อนาทริยํ ปฏิจฺจ อุทเก อุจฺจารํ วา ปสฺสาวํ วา เขฬํ วา กรณปจฺจยา ทุกฺกฏํ กตฺถ ปญฺญตฺตนฺติ สาวตฺถิยํ ปญฺญตฺตํ. กํ อารพฺภาติ ฉพฺพคฺคิเย ภิกฺขู อารพฺภ. กิสฺมิํ วตฺถุสฺมินฺติ ฉพฺพคฺคิยา ภิกฺขู อุทเก อุจฺจารมฺปิ ปสฺสาวมฺปิ เขฬมฺปิ อกํสุ, ตสฺมิํ วตฺถุสฺมิํ. เอกา ปญฺญตฺติ, เอกา อนุปญฺญตฺติ. ฉนฺนํ อาปตฺติ\-สมุฏฺฐานานํ เอเกน สมุฏฺฐาเนน สมุฏฺฐาติ, กายโต จ จิตฺตโต จ สมุฏฺฐาติ, น วาจโต ฯเปฯ.}{กตฺถ\-ปญฺญตฺติวาโร นิฏฺฐิโต ปฐโม.}
\csromanmatikaanchor{mtk.64}
\csromanmatikaanchor{mtk.65}
\csromantocmarknopage{chapter}{2. กตาปตฺติวาร}{mtk.64}
\bookmark[dest={mtk.64},level=0]{2. กตาปตฺติวาร}
\csromantocmarkref{chapter}{1. ปาราชิกกณฺฑ}{mtk.65}
\bookmark[dest={mtk.65},level=0]{1. ปาราชิกกณฺฑ}
\chapterhead{2. \textbf{กตาปตฺติวาร 1. ปาราชิกกณฺฑ}}
\proseitem{193}{เมถุนํ ธมฺมํ ปฏิเสวน\-ปจฺจยา กติ อาปตฺติโย อาปชฺชติ. เมถุนํ ธมฺมํ ปฏิเสวน\-ปจฺจยา จตสฺโส อาปตฺติโย อาปชฺชติ. อกฺขายิเต สรีเร เมถุนํ ธมฺมํ ปฏิเสวติ, อาปตฺติ ปาราชิกสฺส. เยภุยฺเยน ขายิเต สรีเร เมถุนํ ธมฺมํ ปฏิเสวติ, อาปตฺติ \csromanfolio{89}ถุลฺลจฺจยสฺส. วฏฺฏกเต มุเข อจฺฉุปนฺตํ องฺคชาตํ ปเวเสติ, อาปตฺติ ทุกฺกฏสฺส. ชตุมฏฺฐเก ปาจิตฺติยํ. เมถุนํ ธมฺมํ ปฏิเสวน\-ปจฺจยา อิมา จตสฺโส อาปตฺติโย อาปชฺชติ.}

\prose{อทินฺนํ อาทิยนปจฺจยา กติ อาปตฺติโย อาปชฺชติ. อทินฺนํ อาทิยนปจฺจยา ติสฺโส อาปตฺติโย อาปชฺชติ. ปญฺจมาสกํ วา อติเรก\-ปญฺจ\-มาสกํ วา อคฺฆนกํ อทินฺนํ เถยฺยสงฺขาตํ อาทิยติ, อาปตฺติ ปาราชิกสฺส. อติเรกมาสกํ วา อูนปญฺจมาสกํ วา อคฺฆนกํ อทินฺนํ เถยฺยสงฺขาตํ อาทิยติ, อาปตฺติ ถุลฺลจฺจยสฺส. มาสกํ วา อูนมาสกํ วา อคฺฆนกํ อทินฺนํ เถยฺยสงฺขาตํ อาทิยติ, อาปตฺติ ทุกฺกฏสฺส. อทินฺนํ อาทิยนปจฺจยา อิมา ติสฺโส อาปตฺติโย อาปชฺชติ.}

\prose{สญฺจิจฺจ มนุสฺส\-วิคฺคหํ ชีวิตา โวโรปน\-ปจฺจยา กติ อาปตฺติโย อาปชฺชติ. สญฺจิจฺจ มนุสฺส\-วิคฺคหํ ชีวิตา โวโรปน\-ปจฺจยา ติสฺโส อาปตฺติโย อาปชฺชติ. มนุสฺสํ โอทิสฺส โอปาตํ ขณติ “ปปติตฺวา มริสฺสตี”ติ, อาปตฺติ ทุกฺกฏสฺส. ปปติเต ทุกฺขา เวทนา อุปฺปชฺชติ, อาปตฺติ ถุลฺลจฺจยสฺส. มรติ, อาปตฺติ ปาราชิกสฺส. สญฺจิจฺจ มนุสฺส\-วิคฺคหํ ชีวิตา โวโรปน\-ปจฺจยา อิมา ติสฺโส อาปตฺติโย อาปชฺชติ.}

\prosewithcloserruled{อสนฺตํ อภูตํ อุตฺตริ\-มนุสฺสธมฺมํ อุลฺลปน\-ปจฺจยา กติ อาปตฺติโย อาปชฺชติ. อสนฺตํ อภูตํ อุตฺตริ\-มนุสฺสธมฺมํ อุลฺลปน\-ปจฺจยา ติสฺโส อาปตฺติโย อาปชฺชติ. ปาปิจฺโฉ อิจฺฉาปกโต อสนฺตํ อภูตํ อุตฺตริ\-มนุสฺสธมฺมํ อุลฺลปติ, อาปตฺติ ปาราชิกสฺส. “โย เต วิหาเร วสติ, โส ภิกฺขุ อรหา”ติ ภณติ, ปฏิวิชานนฺตสฺส อาปตฺติ ถุลฺลจฺจยสฺส. น ปฏิวิชานนฺตสฺส อาปตฺติ ทุกฺกฏสฺส. อสนฺตํ อภูตํ อุตฺตริ\-มนุสฺสธมฺมํ อุลฺลปน\-ปจฺจยา อิมา ติสฺโส อาปตฺติโย อาปชฺชติ.}{จตฺตาโร ปาราชิกา นิฏฺฐิตา.}
\csromanmatikaanchor{mtk.66}
\csromantocmarkref{chapter}{2. สํฆาทิเสสกณฺฑาทิ}{mtk.66}
\bookmark[dest={mtk.66},level=0]{2. สํฆาทิเสสกณฺฑาทิ}
\chapterhead{2. \textbf{สํฆาทิเสสกณฺฑาทิ}}
\proseitem{194}{อุปกฺกมิตฺวา อสุจิโมจน\-ปจฺจยา กติ อาปตฺติโย อาปชฺชติ. อุปกฺกมิตฺวา อสุจิโมจน\-ปจฺจยา ติสฺโส อาปตฺติโย อาปชฺชติ. เจเตติ อุปกฺกมติ มุจฺจติ, อาปตฺติ สํฆาทิเสสสฺส. \csromanfolio{90}เจเตติ อุปกฺกมติ น มุจฺจติ, อาปตฺติ ถุลฺลจฺจยสฺส. ปโยเค ทุกฺกฏํ. อุปกฺกมิตฺวา อสุจิโมจน\-ปจฺจยา อิมา ติสฺโส อาปตฺติโย อาปชฺชติ.}

\prose{กายสํสคฺคํ สมาปชฺชน\-ปจฺจยา กติ อาปตฺติโย อาปชฺชติ. กายสํสคฺคํ สมาปชฺชน\-ปจฺจยา ปญฺจ อาปตฺติโย อาปชฺชติ. อวสฺสุตา ภิกฺขุนี อวสฺสุตสฺส ปุริส\-ปุคฺคลสฺส อธกฺขกํ อุพฺภ\-ชาณุมณฺฑลํ คหณํ สาทิยติ, อาปตฺติ ปาราชิกสฺส. ภิกฺขุ กาเยน กายํ อามสติ, อาปตฺติ สํฆาทิเสสสฺส. กาเยน กาย\-ปฏิพทฺธํ อามสติ, อาปตฺติ ถุลฺลจฺจยสฺส. กาย\-ปฏิพทฺเธน กาย\-ปฏิพทฺธํ อามสติ, อาปตฺติ ทุกฺกฏสฺส. องฺคุลิปโตทเก ปาจิตฺติยํ. กายสํสคฺคํ สมาปชฺชน\-ปจฺจยา อิมา ปญฺจ อาปตฺติโย อาปชฺชติ.}

\prose{มาตุคามํ ทุฏฺฐุลฺลาหิ วาจาหิ โอภาสน\-ปจฺจยา ติสฺโส อาปตฺติโย อาปชฺชติ. วจฺจมคฺคํ ปสฺสาวมคฺคํ อาทิสฺส วณฺณมฺปิ ภณติ อวณฺณมฺปิ ภณติ, อาปตฺติ สํฆาทิเสสสฺส. วจฺจมคฺคํ ปสฺสาวมคฺคํ ฐเปตฺวา อธกฺขกํ อุพฺภ\-ชาณุมณฺฑลํ อาทิสฺส วณฺณมฺปิ ภณติ อวณฺณมฺปิ ภณติ, อาปตฺติ ถุลฺลจฺจยสฺส. กาย\-ปฏิพทฺธํ อาทิสฺส วณฺณมฺปิ ภณติ อวณฺณมฺปิ ภณติ, อาปตฺติ ทุกฺกฏสฺส.}

\prose{อตฺต\-กาม\-ปาริจริยาย วณฺณํ ภาสนปจฺจยา ติสฺโส อาปตฺติโย อาปชฺชติ. มาตุคามสฺส สนฺติเก อตฺต\-กาม\-ปาริจริยาย วณฺณํ ภาสติ, อาปตฺติ สํฆาทิเสสสฺส. ปณฺฑกสฺส สนฺติเก อตฺต\-กาม\-ปาริจริยาย วณฺณํ ภาสติ, อาปตฺติ ถุลฺลจฺจยสฺส. ติรจฺฉาน\-คตสฺส สนฺติเก อตฺต\-กาม\-ปาริจริยาย วณฺณํ ภาสติ, อาปตฺติ ทุกฺกฏสฺส.}

\prose{สญฺจริตฺตํ สมาปชฺชน\-ปจฺจยา ติสฺโส อาปตฺติโย อาปชฺชติ. ปฏิคฺคณฺหาติ วีมํสติ ปจฺจาหรติ, อาปตฺติ สํฆาทิเสสสฺส. ปฏิคฺคณฺหาติ วีมํสติ น ปจฺจาหรติ, อาปตฺติ ถุลฺลจฺจยสฺส. ปฏิคฺคณฺหาติ น วีมํสติ น ปจฺจาหรติ, อาปตฺติ ทุกฺกฏสฺส.}

\prose{สญฺญาจิกาย กุฏิํ การาปน\-ปจฺจยา ติสฺโส อาปตฺติโย อาปชฺชติ. การาเปติ, ปโยเค ทุกฺกฏํ. เอกํ ปิณฺฑํ\csromansharedfootnote{e5b074e462d0d5c3}{เอกปิณฺเฑ (สฺยา)} อนาคเต อาปตฺติ ถุลฺลจฺจยสฺส. ตสฺมิํ ปิณฺเฑ อาคเต อาปตฺติ สํฆาทิเสสสฺส.}

\prose{\csromanfolio{91}มหลฺลกํ วิหารํ การาปน\-ปจฺจยา ติสฺโส อาปตฺติโย อาปชฺชติ. การาเปติ, ปโยเค ทุกฺกฏํ. เอกํ ปิณฺฑํ อนาคเต อาปตฺติ ถุลฺลจฺจยสฺส. ตสฺมิํ ปิณฺเฑ อาคเต อาปตฺติ สํฆาทิเสสสฺส.}

\prose{ภิกฺขุํ อมูลเกน ปาราชิเกน ธมฺเมน อนุทฺธํสน\-ปจฺจยา ติสฺโส อาปตฺติโย อาปชฺชติ. อโนกาสํ การาเปตฺวา จาวนาธิปฺปาโย วเทติ, อาปตฺติ สํฆาทิเสเสน ทุกฺกฏสฺส. โอกาสํ การาเปตฺวา อกฺโกสาธิปฺปาโย วเทติ, อาปตฺติ โอมสวาทสฺส.}

\prose{ภิกฺขุํ อญฺญภาคิยสฺส อธิกรณสฺส กิญฺจิ เทสํ เลสมตฺตํ อุปาทาย ปาราชิเกน ธมฺเมน อนุทฺธํสน\-ปจฺจยา ติสฺโส อาปตฺติโย อาปชฺชติ. อโนกาสํ การาเปตฺวา จาวนาธิปฺปาโย วเทติ, อาปตฺติ สํฆาทิเสเสน ทุกฺกฏสฺส. โอกาสํ การาเปตฺวา อกฺโกสาธิปฺปาโย วเทติ, อาปตฺติ โอมสวาทสฺส.}

\prose{สํฆเภทโก ภิกฺขุ ยาวตติยํ สมนุภาสนาย น ปฏินิสฺสชฺชน\-ปจฺจยา ติสฺโส อาปตฺติโย อาปชฺชติ. ญตฺติยา ทุกฺกฏํ. ทฺวีหิ กมฺมวาจาหิ ถุลฺลจฺจยา. กมฺมวาจา\-ปริโยสาเน อาปตฺติ สํฆาทิเสสสฺส.}

\prose{เภทกา\-นุวตฺตกา ภิกฺขู ยาวตติยํ สมนุภาสนาย น ปฏินิสฺสชฺชน\-ปจฺจยา ติสฺโส อาปตฺติโย อาปชฺชนฺติ. ญตฺติยา ทุกฺกฏํ. ทฺวีหิ กมฺมวาจาหิ ถุลฺลจฺจยา. กมฺมวาจา\-ปริโยสาเน อาปตฺติ สํฆาทิเสสสฺส.}

\prose{ทุพฺพโจ ภิกฺขุ ยาวตติยํ สมนุภาสนาย น ปฏินิสฺสชฺชน\-ปจฺจยา ติสฺโส อาปตฺติโย อาปชฺชติ. ญตฺติยา ทุกฺกฏํ. ทฺวีหิ กมฺมวาจาหิ ถุลฺลจฺจยา. กมฺมวาจา\-ปริโยสาเน อาปตฺติ สํฆาทิเสสสฺส.}

\prose{กุลทูสโก ภิกฺขุ ยาวตติยํ สมนุภาสนาย น ปฏินิสฺสชฺชน\-ปจฺจยา ติสฺโส อาปตฺติโย อาปชฺชติ. ญตฺติยา ทุกฺกฏํ. ทฺวีหิ กมฺมวาจาหิ ถุลฺลจฺจยา. กมฺมวาจา\-ปริโยสาเน อาปตฺติ สํฆาทิเสสสฺส ฯเปฯ.}

\prosewithcloserruled{อนาทริยํ ปฏิจฺจ อุทเก อุจฺจารํ วา ปสฺสาวํ วา เขฬํ วา กรณปจฺจยา กติ อาปตฺติโย อาปชฺชติ. อนาทริยํ ปฏิจฺจ อุทเก อุจฺจารํ วา \csromanfolio{92}ปสฺสาวํ วา เขฬํ วา กรณปจฺจยา เอกํ อาปตฺติํ อาปชฺชติ, ทุกฺกฏํ. อนาทริยํ ปฏิจฺจ อุทเก อุจฺจารํ วา ปสฺสาวํ วา เขฬํ วา กรณปจฺจยา อิมํ เอกํ อาปตฺติํ อาปชฺชติ.}{กตาปตฺติวาโร นิฏฺฐิโต ทุติโย.}
\csromanmatikaanchor{mtk.67}
\csromantocmarkref{chapter}{3. วิปตฺติวาร}{mtk.67}
\bookmark[dest={mtk.67},level=0]{3. วิปตฺติวาร}
\chapterhead{3. \textbf{วิปตฺติวาร}}
\proseitem{195}{เมถุนํ ธมฺมํ ปฏิเสวน\-ปจฺจยา อาปตฺติโย จตุนฺนํ วิปตฺตีนํ กติ วิปตฺติโย ภชนฺติ. เมถุนํ ธมฺมํ ปฏิเสวน\-ปจฺจยา อาปตฺติโย จตุนฺนํ วิปตฺตีนํ เทฺว วิปตฺติโย ภชนฺติ, สิยา สีลวิปตฺติํ, สิยา อาจารวิปตฺติํ ฯเปฯ.}

\prosewithcloserruled{อนาทริยํ ปฏิจฺจ อุทเก อุจฺจารํ วา ปสฺสาวํ วา เขฬํ วา กรณปจฺจยา อาปตฺติ จตุนฺนํ วิปตฺตีนํ กติ วิปตฺติโย ภชติ. อนาทริยํ ปฏิจฺจ อุทเก อุจฺจารํ วา ปสฺสาวํ ว เขฬํ วา กรณปจฺจยา อาปตฺติ จตุนฺนํ วิปตฺตีนํ เอกํ วิปตฺติํ ภชติ, อาจารวิปตฺติํ.}{วิปตฺติวาโร นิฏฺฐิโต ตติโย.}
\csromanmatikaanchor{mtk.68}
\csromantocmarkref{chapter}{4. สงฺคหิตวาร}{mtk.68}
\bookmark[dest={mtk.68},level=0]{4. สงฺคหิตวาร}
\chapterhead{4. \textbf{สงฺคหิตวาร}}
\proseitem{196}{เมถุนํ ธมฺมํ ปฏิเสวน\-ปจฺจยา อาปตฺติโย สตฺตนฺนํ อาปตฺติ\-กฺขนฺธานํ กติหิ อาปตฺติ\-กฺขนฺเธหิ สงฺคหิตา. เมถุนํ ธมฺมํ ปฏิเสวน\-ปจฺจยา อาปตฺติโย สตฺตนฺนํ อาปตฺติ\-กฺขนฺธานํ จตูหิ อาปตฺติ\-กฺขนฺเธหิ สงฺคหิตา, สิยา ปาราชิกาปตฺติ\-กฺขนฺเธน, สิยา ถุลฺลจฺจยา\-ปตฺติกฺขนฺเธน, สิยา ปาจิตฺติยาปตฺติ\-กฺขนฺเธน, สิยา ทุกฺกฏาปตฺติ\-กฺขนฺเธน ฯเปฯ.}

\prose{อนาทริยํ ปฏิจฺจ อุทเก อุจฺจารํ วา ปสฺสาวํ วา เขฬํ วา กรณปจฺจยา อาปตฺติ สตฺตนฺนํ อาปตฺติ\-กฺขนฺธานํ กติหิ อาปตฺติ\-กฺขนฺเธหิ สงฺคหิตา. อนาทริยํ ปฏิจฺจ อุทเก อุจฺจารํ วา ปสฺสาวํ วา เขฬํ วา กรณปจฺจยา อาปตฺติ สตฺตนฺนํ อาปตฺติ\-กฺขนฺธานํ เอเกน อาปตฺติ\-กฺขนฺเธน สงฺคหิตา, ทุกฺกฏาปตฺติ\-กฺขนฺเธน.}

\vspace{\tipitakaparskip}
\csromancenter{สงฺคหิตวาโร นิฏฺฐิโต จตุตฺโถ.}
\csromanmatikaanchor{mtk.69}
\csromantocmarkref{chapter}{5. สมุฏฺฐานวาร}{mtk.69}
\bookmark[dest={mtk.69},level=0]{5. สมุฏฺฐานวาร}
\chapterheadpageread{5. \textbf{สมุฏฺฐานวาร}}
\proseitem{197}{\csromanfolio{93}เมถุนํ ธมฺมํ ปฏิเสวน\-ปจฺจยา อาปตฺติโย ฉนฺนํ อาปตฺติ\-สมุฏฺฐานานํ กติหิ สมุฏฺฐาเนหิ สมุฏฺฐนฺติ. เมถุนํ ธมฺมํ ปฏิเสวน\-ปจฺจยา อาปตฺติโย ฉนฺนํ อาปตฺติ\-สมุฏฺฐานานํ เอเกน สมุฏฺฐาเนน สมุฏฺฐนฺติ, กายโต จ จิตฺตโต จ สมุฏฺฐนฺติ, น วาจโต ฯเปฯ.}

\prosewithcloserruled{อนาทริยํ ปฏิจฺจ อุทเก อุจฺจารํ วา ปสฺสาวํ วา เขฬํ วา กรณปจฺจยา อาปตฺติ ฉนฺนํ อาปตฺติ\-สมุฏฺฐานานํ กติหิ สมุฏฺฐาเนหิ สมุฏฺฐาติ. อนาทริยํ ปฏิจฺจ อุทเก อุจฺจรํ วา ปสฺสาวํ วา เขฬํ วา กรณปจฺจยา อาปตฺติ ฉนฺนํ อาปตฺติ\-สมุฏฺฐานานํ เอเกน สมุฏฺฐาเนน สมุฏฺฐาติ, กายโต จ จิตฺตโต จ สมุฏฺฐาติ, น วาจโต ฯเปฯ.}{สมุฏฺฐานวาโร นิฏฺฐิโต ปญฺจโม.}
\csromanmatikaanchor{mtk.70}
\csromantocmarkref{chapter}{6. อธิกรณวาร}{mtk.70}
\bookmark[dest={mtk.70},level=0]{6. อธิกรณวาร}
\chapterhead{6. \textbf{อธิกรณวาร}}
\proseitem{198}{เมถุนํ ธมฺมํ ปฏิเสวน\-ปจฺจยา อาปตฺติโย จตุนฺนํ อธิกรณานํ กตมํ อธิกรณํ. เมถุนํ ธมฺมํ ปฏิเสวน\-ปจฺจยา อาปตฺติโย จตุนฺนํ อธิกรณานํ อาปตฺตา\-ธิกรณํ ฯเปฯ.}

\prosewithcloserruled{อนาทริยํ ปฏิจฺจ อุทเก อุจฺจารํ วา ปสฺสาวํ วา เขฬํ วา กรณปจฺจยา อาปตฺติ จตุนฺนํ อธิกรณานํ กตมํ อธิกรณํ. อนาทริยํ ปฏิจฺจ อุทเก อุจฺจารํ วา ปสฺสาวํ วา เขฬํ วา กรณปจฺจยา อาปตฺติ จตุนฺนํ อธิกรณานํ อาปตฺตา\-ธิกรณํ.}{อธิกรณวาโร นิฏฺฐิโต ฉฏฺโฐ.}
\csromanmatikaanchor{mtk.71}
\csromantocmarkref{chapter}{7. สมถวาร}{mtk.71}
\bookmark[dest={mtk.71},level=0]{7. สมถวาร}
\chapterhead{7. \textbf{สมถวาร}}
\proseitem{199}{เมถุนํ ธมฺมํ ปฏิเสวน\-ปจฺจยา อาปตฺติโย สตฺตนฺนํ สมถานํ กติหิ สมเถหิ สมฺมนฺติ. เมถุนํ ธมฺมํ ปฏิเสวน\-ปจฺจยา อาปตฺติโย สตฺตนฺนํ สมถานํ ตีหิ สมเถหิ สมฺมนฺติ, สิยา สมฺมุขา\-วินเยน จ ปฏิญฺญาต\-กรเณน จ, สิยา สมฺมุขา\-วินเยน จ ติณวตฺถารเกน จ ฯเปฯ.}

\prosewithcloserruled{\csromanfolio{94}อนาทริยํ ปฏิจฺจ อุทเก อุจฺจารํ วา ปสฺสาวํ วา เขฬํ วา กรณปจฺจยา อาปตฺติ สตฺตนฺนํ สมถานํ กติหิ สมเถหิ สมฺมติ. อนาทริยํ ปฏิจฺจ อุทเก อุจฺจรํ วา ปสฺสาวํ วา เขฬํ วา กรณปจฺจยา อาปตฺติ สตฺตนฺนํ สมถานํ ตีหิ สมเถหิ สมฺมติ, สิยา สมฺมุขา\-วินเยน จ ปฏิญฺญาต\-กรเณน จ, สิยา สมฺมุขา\-วินเยน จ ติณวตฺถารเกน จาติ.}{สมถวาโร นิฏฺฐิโต สตฺตโม.}
\csromanmatikaanchor{mtk.72}
\csromantocmarkref{chapter}{8. สมุจฺจยวาร}{mtk.72}
\bookmark[dest={mtk.72},level=0]{8. สมุจฺจยวาร}
\chapterhead{8. \textbf{สมุจฺจยวาร}}
\proseitem{200}{เมถุนํ ธมฺมํ ปฏิเสวน\-ปจฺจยา กติ อาปตฺติโย อาปชฺชติ. เมถุนํ ธมฺมํ ปฏิเสวน\-ปจฺจยา จตสฺโส อาปตฺติโย อาปชฺชติ. อกฺขายิเต สรีเร เมถุนํ ธมฺมํ ปฏิเสวติ, อาปตฺติ ปาราชิกสฺส. เยภุยฺเยน ขายิเต สรีเร เมถุนํ ธมฺมํ ปฏิเสวติ, อาปตฺติ ถุลฺลจฺจยสฺส. วฏฺฏกเต มุเข อจฺฉุปนฺตํ องฺคชาตํ ปเวเสติ, อาปตฺติ ทุกฺกฏสฺส. ชตุมฏฺฐเก ปาจิตฺติยํ. เมถุนํ ธมฺมํ ปฏิเสวน\-ปจฺจยา อิมา จตสฺโส อาปตฺติโย อาปชฺชติ. ตา อาปตฺติโย จตุนฺนํ วิปตฺตีนํ กติ วิปตฺติโย ภชนฺติ, สตฺตนฺนํ อาปตฺติ\-กฺขนฺธานํ กติหิ อาปตฺติ\-กฺขนฺเธหิ สงฺคหิตา, ฉนฺนํ อาปตฺติ\-สมุฏฺฐานานํ กติหิ สมุฏฺฐาเนหิ สมุฏฺฐนฺติ, จตุนฺนํ อธิกรณานํ กตมํ อธิกรณํ, สตฺตนฺนํ สมถานํ กติหิ สมเถหิ สมฺมนฺติ. ตา อาปตฺติโย จตุนฺนํ วิปตฺตีนํ เทฺว วิปตฺติโย ภชนฺติ, สิยา สีลวิปตฺติํ, สิยา อาจารวิปตฺติํ. สตฺตนฺนํ อาปตฺติ\-กฺขนฺธานํ จตูหิ อาปตฺติ\-กฺขนฺเธหิ สงฺคหิตา, สิยา ปาราชิกาปตฺติ\-กฺขนฺเธน, สิยา ถุลฺลจฺจยา\-ปตฺติกฺขนฺเธน, สิยา ปาจิตฺติยาปตฺติ\-กฺขนฺเธน, สิยา ทุกฺกฏาปตฺติ\-กฺขนฺเธน. ฉนฺนํ อาปตฺติ\-สมุฏฺฐานานํ เอเกน สมุฏฺฐาเนน สมุฏฺฐนฺติ, กายโต จ จิตฺตโต จ สมุฏฺฐนฺติ, น วาจโต. จตุนฺนํ อธิกรณานํ อาปตฺตา\-ธิกรณํ. สตฺตนฺนํ สมถานํ ตีหิ สมเถหิ สมฺมนฺติ, สิยา สมฺมุขา\-วินเยน จ ปฏิญฺญาต\-กรเณน จ, สิยา สมฺมุขา\-วินเยน จ ติณวตฺถารเกน จ ฯเปฯ.}

\prose{อนาทริยํ ปฏิจฺจ อุทเก อุจฺจารํ วา ปสฺสาวํ วา เขฬํ วา กรณปจฺจยา กติ อาปตฺติโย อาปชฺชติ. อนาทริยํ ปฏิจฺจ อุทเก อุจฺจารํ วา ปสฺสาวํ วา เขฬํ วา กรณปจฺจยา เอกํ อาปตฺติํ อาปชฺชติ, ทุกฺกฏํ. อนาทริยํ \csromanfolio{95}ปฏิจฺจ อุทเก อุจฺจารํ วา ปสฺสาวํ วา เขฬํ วา กรณปจฺจยา อิมํ เอกํ อาปตฺติํ อาปชฺชติ. สา อาปตฺติ จตุนฺนํ วิปตฺตีนํ กติ วิปตฺติโย ภชติ, สตฺตนฺนํ อาปตฺติ\-กฺขนฺธานํ กติหิ อาปตฺติ\-กฺขนฺเธหิ สงฺคหิตา, ฉนฺนํ อาปตฺติ\-สมุฏฺฐานานํ กติหิ สมุฏฺฐาเนหิ สมุฏฺฐาติ, จตุนฺนํ อธิกรณานํ กตมํ อธิกรณํ, สตฺตนฺนํ สมถานํ กติหิ สมเถหิ สมฺมติ. สา อาปตฺติ จตุนฺนํ วิปตฺตีนํ เอกํ วิปตฺติํ ภชติ, อาจารวิปตฺติํ. สตฺตนฺนํ อาปตฺติ\-กฺขนฺธานํ เอเกน อาปตฺติ\-กฺขนฺเธน สงฺคหิตา, ทุกฺกฏาปตฺติ\-กฺขนฺเธน. ฉนฺนํ อาปตฺติ\-สมุฏฺฐานานํ เอเกน สมุฏฺฐาเนน สมุฏฺฐาติ, กายโต จ จิตฺตโต จ สมุฏฺฐาติ, น วาจโต. จตุนฺนํ อธิกรณานํ อาปตฺตา\-ธิกรณํ. สตฺตนฺนํ สมถานํ ตีหิ สมเถหิ สมฺมติ, สิยา สมฺมุขา\-วินเยน จ ปฏิญฺญาต\-กรเณน จ, สิยา สมฺมุขา\-วินเยน จ ติณวตฺถารเกน จาติ.}

\vspace{\tipitakaparskip}
\csromancenter{สมุจฺจยวาโร นิฏฺฐิโต อฏฺฐโม.}
\nitthitam{อฏฺฐ\-ปจฺจย\-วารา นิฏฺฐิตา.}
\nitthitam{มหาวิภงฺเค โสฬสมหาวารา นิฏฺฐิตา.}
\nitthitam{ภิกฺขุวิภงฺค\-มหาวาโร นิฏฺฐิโต.}
\csromanmatikaanchor{mtk.73}
\csromantocmarknopage{chapter}{ภิกฺขุนีวิภงฺค}{mtk.73}
\bookmark[dest={mtk.73},level=0]{ภิกฺขุนีวิภงฺค}
\nitthitamruled{ภิกฺขุนี\-วิภงฺค}
\csromanmatikaanchor{mtk.74}
\csromanmatikaanchor{mtk.75}
\csromantocmarknopage{chapter}{1. กตฺถปญฺญตฺติวาร}{mtk.74}
\bookmark[dest={mtk.74},level=0]{1. กตฺถปญฺญตฺติวาร}
\csromantocmarkref{chapter}{1. ปาราชิกกณฺฑ}{mtk.75}
\bookmark[dest={mtk.75},level=0]{1. ปาราชิกกณฺฑ}
\chapterhead{1. กตฺถ\-ปญฺญตฺติวาร 1. ปาราชิกกณฺฑ}
\proseitem{201}{\csromanfolio{96}ยํ เตน ภควตา ชานตา ปสฺสตา อรหตา สมฺมา\-สมฺพุทฺเธน ภิกฺขุนีนํ ปญฺจมํ ปาราชิกํ กตฺถ ปญฺญตฺตํ, กํ อารพฺภ, กิสฺมิํ วตฺถุสฺมิํ, อตฺถิ ตตฺถ ปญฺญตฺติ อนุปญฺญตฺติ อนุปฺปนฺนปญฺญตฺติ, สพฺพตฺถ\-ปญฺญตฺติ ปเทสปญฺญตฺติ, สาธารณ\-ปญฺญตฺติ อสาธารณ\-ปญฺญตฺติ, เอกโตปญฺญตฺติ อุภโตปญฺญตฺติ, จตุนฺนํ ปาติโมกฺขุ\-ทฺเทสานํ กตฺโถคธํ กตฺถ ปริยาปนฺนํ, กตเมน อุทฺเทเสน อุทฺเทสํ อาคจฺฉติ, จตุนฺนํ วิปตฺตีนํ กตมา วิปตฺติ, สตฺตนฺนํ อาปตฺติ\-กฺขนฺธานํ กตโม อาปตฺติกฺขนฺโธ, ฉนฺนํ อาปตฺติ\-สมุฏฺฐานานํ กติหิ สมุฏฺฐาเนหิ สมุฏฺฐาติ, จตุนฺนํ อธิกรณานํ กตมํ อธิกรณํ, สตฺตนฺนํ สมถานํ กติหิ สมเถหิ สมฺมติ, โก ตตฺถ วินโย, โก ตตฺถ อภิวินโย, กิํ ตตฺถ ปาติโมกฺขํ, กิํ ตตฺถ อธิปาติโมกฺขํ, กา วิปตฺติ, กา สมฺปตฺติ, กา ปฏิปตฺติ, กติ อตฺถวเส ปฏิจฺจ ภควตา ภิกฺขุนีนํ ปญฺจมํ ปาราชิกํ ปญฺญตฺตํ, กา สิกฺขนฺติ, กา สิกฺขิตสิกฺขา, กตฺถ ฐิตํ, กา ธาเรนฺติ, กสฺส วจนํ, เกนาภตนฺติ.}

\proseitem{202}{ยํ เตน ภควตา ชานตา ปสฺสตา อรหตา สมฺมา\-สมฺพุทฺเธน ภิกฺขุนีนํ ปญฺจมํ ปาราชิกํ กตฺถ ปญฺญตฺตนฺติ สาวตฺถิยํ ปญฺญตฺตํ. กํ อารพฺภาติ สุนฺทรีนนฺทํ ภิกฺขุนิํ อารพฺภ. กิสฺมิํ วตฺถุสฺมินฺติ สุนฺทรีนนฺทา ภิกฺขุนี อวสฺสุตา อวสฺสุตสฺส ปุริส\-ปุคฺคลสฺส กายสํสคฺคํ สาทิยิ, ตสฺมิํ วตฺถุสฺมิํ. อตฺถิ ตตฺถ ปญฺญตฺติ อนุปญฺญตฺติ อนุปฺปนฺนปญฺญตฺตีติ เอกา ปญฺญตฺติ, อนุปญฺญตฺติ อนุปฺปนฺนปญฺญตฺติ ตสฺมิํ นตฺถิ. สพฺพตฺถ\-ปญฺญตฺติ ปเทสปญฺญตฺตีติ สพฺพตฺถ\-ปญฺญตฺติ. สาธารณ\-ปญฺญตฺติ อสาธารณปญฺญตฺตีติ อสาธารณ\-ปญฺญตฺติ. เอกโตปญฺญตฺติ อุภโตปญฺญตฺตีติ เอกโตปญฺญตฺติ. จตุนฺนํ ปาติโมกฺขุ\-ทฺเทสานํ กตฺโถคธํ กตฺถ ปริยาปนฺนนฺติ นิทาโนคธํ นิทาน\-ปริยาปนฺนํ. กตเมน อุทฺเทเสน อุทฺเทสํ อาคจฺฉตีติ ทุติเยน อุทฺเทเสน อุทฺเทสํ อาคจฺฉติ. จตุนฺนํ วิปตฺตีนํ กตมา วิปตฺตีติ สีลวิปตฺติ. สตฺตนฺนํ อาปตฺติ\-กฺขนฺธานํ กตโม อาปตฺติ\-กฺขนฺโธติ ปาราชิกาปตฺติ\-กฺขนฺโธ. ฉนฺนํ อาปตฺติ\-สมุฏฺฐานานํ กติหิ สมุฏฺฐาเนหิ สมุฏฺฐาตีติ เอเกน \csromanfolio{97}สมุฏฺฐาเนน สมุฏฺฐาติ, กายโต จ จิตฺตโต จ สมุฏฺฐาติ, น วาจโต. จตุนฺนํ อธิกรณานํ กตมํ อธิกรณนฺติ อาปตฺตา\-ธิกรณํ. สตฺตนฺนํ สมถานํ กติหิ สมเถหิ สมฺมตีติ ทฺวีหิ สมเถหิ สมฺมติ สมฺมุขา\-วินเยน จ ปฏิญฺญาต\-กรเณน จ. โก ตตฺถ วินโย โก ตตฺถ อภิวินโยติ ปญฺญตฺติ วินโย, วิภตฺติ อภิวินโย. กิํ ตตฺถ ปาติโมกฺขํ กิํ ตตฺถ อธิปาติ\-โมกฺขนฺติ ปญฺญตฺติ ปาติโมกฺขํ, วิภตฺติ อธิปาติโมกฺขํ. กา วิปตฺตีติ อสํวโร วิปตฺติ. กา สมฺปตฺตีติ สํวโร สมฺปตฺติ. กา ปฏิปตฺตีติ “น เอวรูปํ กริสฺสามี”ติ ยาวชีวํ อาปาณโกฏิกํ สมาทาย สิกฺขติ สิกฺขาปเทสุ. * กติ อตฺถวเส ปฏิจฺจ ภควตา ภิกฺขุนีนํ ปญฺจมํ ปาราชิกํ ปญฺญตฺตนฺติ ทส อตฺถวเส ปฏิจฺจ ภควตา ภิกฺขุนีนํ ปญฺจมํ ปาราชิกํ ปญฺญตฺตํ สํฆสุฏฺฐุตาย สํฆผาสุตาย ทุมฺมงฺกูนํ ภิกฺขุนีนํ นิคฺคหาย เปสลานํ ภิกฺขุนีนํ ผาสุวิหาราย ทิฏฺฐ\-ธมฺมิกานํ อาสวานํ สํวราย สมฺปรายิกานํ อาสวานํ ปฏิฆาตาย อปฺปสนฺนานํ ปสาทาย ปสนฺนานํ ภิยฺโยภาวาย สทฺธมฺม\-ฏฺฐิติยา วินยา\-นุคฺคหาย. กา สิกฺขนฺตีติ เสกฺขา จ ปุถุชฺชน\-กลฺยาณิกา จ สิกฺขนฺติ. กา สิกฺขิต\-สิกฺขาติ อรหนฺติโย1 สิกฺขิตสิกฺขา. กตฺถ ฐิตนฺติ สิกฺขากามาสุ ฐิตํ. กา ธาเรนฺตีติ ยาสํ วตฺตติ, ตา ธาเรนฺติ. กสฺส วจนนฺติ ภควโต วจนํ อรหโต สมฺมา\-สมฺพุทฺธสฺส. เกนาภตนฺติ ปรมฺปราภตํ.}

\setlength{\csromangathapagewidth}{0pt}
\setlength{\csromangathawakwidth}{0pt}
\csromangathameasuregroup{อุปาลิ ทาสโก เจว, \\ โมคฺคลิปุตฺเตน ปญฺจมา, \\ ตโต มหินฺโท อิฏฺฏิโย, \\ วินยํ เต วาจยิํสุ, \\ นิกาเย ปญฺจ วาเจสุํ, \\ ตโต อริฏฺโฐ เมธาวี, \\ วิสารโท กาฬสุมโน,}{\csromangathabat{อุปาลิ ทาสโก เจว,}{โสณโก สิคฺคโว ตถา.} \\ \csromangathabat{โมคฺคลิปุตฺเตน ปญฺจมา,}{เอเต ชมฺพุสิริวฺหเย.} \\ \csromangathabat{ตโต มหินฺโท อิฏฺฏิโย,}{อุตฺติโย สมฺพโล ตถา.} \\ \csromangathabat{วินยํ เต วาจยิํสุ,}{ปิฏกํ ตมฺพปณฺณิยา.} \\ \csromangathabat{นิกาเย ปญฺจ วาเจสุํ,}{สตฺต เจว ปกรเณ.} \\ \csromangathabat{ตโต อริฏฺโฐ เมธาวี,}{ติสฺสทตฺโต จ ปณฺฑิโต.} \\ \csromangathabat{วิสารโท กาฬสุมโน,}{เถโร จ ทีฆนามโก.}}
\csromangathasetleft{อุปาลิ ทาสโก เจว, \\ โมคฺคลิปุตฺเตน ปญฺจมา, \\ ตโต มหินฺโท อิฏฺฏิโย, \\ วินยํ เต วาจยิํสุ, \\ นิกาเย ปญฺจ วาเจสุํ, \\ ตโต อริฏฺโฐ เมธาวี, \\ วิสารโท กาฬสุมโน,}
\csromangathagroup{\csromangathastanza{\csromangathabat{อุปาลิ ทาสโก เจว,}{โสณโก สิคฺคโว ตถา.} \\ \csromangathabat{โมคฺคลิปุตฺเตน ปญฺจมา,}{เอเต ชมฺพุสิริวฺหเย.}}\csromangathastanza{\csromangathabat{ตโต มหินฺโท อิฏฺฏิโย,}{อุตฺติโย สมฺพโล ตถา.} \\ \csromangathabat{วินยํ เต วาจยิํสุ,}{ปิฏกํ ตมฺพปณฺณิยา.}}\csromangathastanza{\csromangathabat{นิกาเย ปญฺจ วาเจสุํ,}{สตฺต เจว ปกรเณ.} \\ \csromangathabat{ตโต อริฏฺโฐ เมธาวี,}{ติสฺสทตฺโต จ ปณฺฑิโต.} \\ \csromangathabat{วิสารโท กาฬสุมโน,}{เถโร จ ทีฆนามโก.}}}

\prose{\csromansymbolmark{*}อํ 3. 311 ปิฏฺเฐปิ 1. อรหนฺตา (ก)}

\setlength{\csromangathapagewidth}{0pt}
\setlength{\csromangathawakwidth}{0pt}
\csromangathameasuregroup{ปุนเทว กาฬสุมโน, \\ ติสฺสตฺเถโร จ เมธาวี, \\ ปุนเทว สุมโน เมธาวี, \\ พหุสฺสุโต จูฬนาโค, \\ ธมฺมปาลิตนาโม จ, \\ ตสฺส สิสฺโส มหาปญฺโญ, \\ ทีเป ตารกราชาว, \\ อุปติสฺโส จ เมธาวี, \\ ปุนเทว สุมโน เมธาวี, \\ มหากถี มหาสิโว, \\ ปุนเทว อุปาลิ เมธาวี, \\ มหานาโค มหาปญฺโญ, \\ ปุนเทว อภโย เมธาวี, \\ ติสฺสตฺเถโร จ เมธาวี, \\ ตสฺส สิสฺโส มหาปญฺโญ, \\ สาสนํ อนุรกฺขนฺโต, \\ จูฬาภโย จ เมธาวี, \\ ติสฺสตฺเถโร จ เมธาวี, \\ จูฬเทโว จ เมธาวี, \\ สิวตฺเถโร จ เมธาวี, \\ เอเต นาคา มหาปญฺญา, \\ วินยํ ทีเป ปกาเสสุํ,}{\csromangathabat{ปุนเทว กาฬสุมโน,}{นาคตฺเถโร จ พุทฺธรกฺขิโต.} \\ \csromangathabat{ติสฺสตฺเถโร จ เมธาวี,}{เทวตฺเถโร จ ปณฺฑิโต.} \\ \csromangathabat{ปุนเทว สุมโน เมธาวี,}{วินเย จ วิสารโท.} \\ \csromangathabat{พหุสฺสุโต จูฬนาโค,}{คโชว ทุปฺปธํสิโย.} \\ \csromangathabat{ธมฺมปาลิตนาโม จ,}{โรหเณ สาธุปูชิโต.} \\ \csromangathabat{ตสฺส สิสฺโส มหาปญฺโญ,}{เขมนาโม ติเปฏโก.} \\ \csromangathabat{ทีเป ตารกราชาว,}{ปญฺญาย อติโรจถ.} \\ \csromangathabat{อุปติสฺโส จ เมธาวี,}{ผุสฺสเทโว มหากถี.} \\ \csromangathabat{ปุนเทว สุมโน เมธาวี,}{ปุปฺผนาโม พหุสฺสุโต.} \\ \csromangathabat{มหากถี มหาสิโว,}{ปิฏเก สพฺพตฺถ โกวิโท.} \\ \csromangathabat{ปุนเทว อุปาลิ เมธาวี,}{วินเย จ วิสารโท.} \\ \csromangathabat{มหานาโค มหาปญฺโญ,}{สทฺธมฺมวํสโกวิโท.} \\ \csromangathabat{ปุนเทว อภโย เมธาวี,}{ปิฏเก สพฺพตฺถ โกวิโท.} \\ \csromangathabat{ติสฺสตฺเถโร จ เมธาวี,}{วินเย จ วิสารโท.} \\ \csromangathabat{ตสฺส สิสฺโส มหาปญฺโญ,}{ปุปฺผนาโม พหุสฺสุโต.} \\ \csromangathabat{สาสนํ อนุรกฺขนฺโต,}{ชมฺพุทีเป ปติฏฺฐิโต.} \\ \csromangathabat{จูฬาภโย จ เมธาวี,}{วินเย จ วิสารโท.} \\ \csromangathabat{ติสฺสตฺเถโร จ เมธาวี,}{สทฺธมฺมวํสโกวิโท.} \\ \csromangathabat{จูฬเทโว จ เมธาวี,}{วินเย จ วิสารโท.} \\ \csromangathabat{สิวตฺเถโร จ เมธาวี,}{วินเย สพฺพตฺถ โกวิโท.} \\ \csromangathabat{เอเต นาคา มหาปญฺญา,}{วินยญฺญู มคฺคโกวิทา.} \\ \csromangathabat{วินยํ ทีเป ปกาเสสุํ,}{ปิฏกํ ตมฺพปณฺณิยาติ.}}
\csromangathasetleft{ปุนเทว กาฬสุมโน, \\ ติสฺสตฺเถโร จ เมธาวี, \\ ปุนเทว สุมโน เมธาวี, \\ พหุสฺสุโต จูฬนาโค, \\ ธมฺมปาลิตนาโม จ, \\ ตสฺส สิสฺโส มหาปญฺโญ, \\ ทีเป ตารกราชาว, \\ อุปติสฺโส จ เมธาวี, \\ ปุนเทว สุมโน เมธาวี, \\ มหากถี มหาสิโว, \\ ปุนเทว อุปาลิ เมธาวี, \\ มหานาโค มหาปญฺโญ, \\ ปุนเทว อภโย เมธาวี, \\ ติสฺสตฺเถโร จ เมธาวี, \\ ตสฺส สิสฺโส มหาปญฺโญ, \\ สาสนํ อนุรกฺขนฺโต, \\ จูฬาภโย จ เมธาวี, \\ ติสฺสตฺเถโร จ เมธาวี, \\ จูฬเทโว จ เมธาวี, \\ สิวตฺเถโร จ เมธาวี, \\ เอเต นาคา มหาปญฺญา, \\ วินยํ ทีเป ปกาเสสุํ,}
\csromangathagroup{\csromangathastanza{\csromanfolio{98}\csromangathabat{ปุนเทว กาฬสุมโน,}{นาคตฺเถโร จ พุทฺธรกฺขิโต.} \\ \csromangathabat{ติสฺสตฺเถโร จ เมธาวี,}{เทวตฺเถโร จ ปณฺฑิโต.}}\csromangathastanza{\csromangathabat{ปุนเทว สุมโน เมธาวี,}{วินเย จ วิสารโท.} \\ \csromangathabat{พหุสฺสุโต จูฬนาโค,}{คโชว ทุปฺปธํสิโย.}}\csromangathastanza{\csromangathabat{ธมฺมปาลิต\-นาโม จ,}{โรหเณ สาธุปูชิโต.} \\ \csromangathabat{ตสฺส สิสฺโส มหาปญฺโญ,}{เขมนาโม ติเปฏโก.}}\csromangathastanza{\csromangathabat{ทีเป ตารกราชาว,}{ปญฺญาย อติโรจถ.} \\ \csromangathabat{อุปติสฺโส จ เมธาวี,}{ผุสฺสเทโว มหากถี.}}\csromangathastanza{\csromangathabat{ปุนเทว สุมโน เมธาวี,}{ปุปฺผนาโม พหุสฺสุโต.} \\ \csromangathabat{มหากถี มหาสิโว,}{ปิฏเก สพฺพตฺถ โกวิโท.}}\csromangathastanza{\csromangathabat{ปุนเทว อุปาลิ เมธาวี,}{วินเย จ วิสารโท.} \\ \csromangathabat{มหานาโค มหาปญฺโญ,}{สทฺธมฺม\-วํส\-โกวิโท.}}\csromangathastanza{\csromangathabat{ปุนเทว อภโย เมธาวี,}{ปิฏเก สพฺพตฺถ โกวิโท.} \\ \csromangathabat{ติสฺสตฺเถโร จ เมธาวี,}{วินเย จ วิสารโท.}}\csromangathastanza{\csromangathabat{ตสฺส สิสฺโส มหาปญฺโญ,}{ปุปฺผนาโม พหุสฺสุโต.} \\ \csromangathabat{สาสนํ อนุรกฺขนฺโต,}{ชมฺพุทีเป ปติฏฺฐิโต.}}\csromangathastanza{\csromangathabat{จูฬาภโย จ เมธาวี,}{วินเย จ วิสารโท.} \\ \csromangathabat{ติสฺสตฺเถโร จ เมธาวี,}{สทฺธมฺม\-วํส\-โกวิโท.}}\csromangathastanza{\csromangathabat{จูฬเทโว จ เมธาวี,}{วินเย จ วิสารโท.} \\ \csromangathabat{สิวตฺเถโร จ เมธาวี,}{วินเย สพฺพตฺถ โกวิโท.}}\csromangathastanza{\csromangathabat{เอเต นาคา มหาปญฺญา,}{วินยญฺญู มคฺคโกวิทา.} \\ \csromangathabat{วินยํ ทีเป ปกาเสสุํ,}{ปิฏกํ ตมฺพปณฺณิยาติ.}}}

\proseitem{203}{ยํ เตน ภควตา ชานตา ปสฺสตา อรหตา สมฺมา\-สมฺพุทฺเธน ภิกฺขุนีนํ ฉฏฺฐํ ปาราชิกํ กตฺถ ปญฺญตฺตนฺติ สาวตฺถิยํ ปญฺญตฺตํ. กํ อารพฺภาติ ถุลฺลนนฺทํ ภิกฺขุนิํ อารพฺภ. กิสฺมิํ วตฺถุสฺมินฺติ ถุลฺลนนฺทา ภิกฺขุนี ชานํ ปาราชิกํ ธมฺมํ อชฺฌาปนฺนํ ภิกฺขุนิํ เนวตฺตนา ปฏิโจเทสิ, \csromanfolio{99}น คณสฺส อาโรเจสิ, ตสฺมิํ วตฺถุสฺมิํ. เอกา ปญฺญตฺติ. ฉนฺนํ อาปตฺติ\-สมุฏฺฐานานํ เอเกน สมุฏฺฐาเนน สมุฏฺฐาติ, กายโต จ วาจโต จ}

\proseitem{204}{ภิกฺขุนีนํ สตฺตมํ ปาราชิกํ กตฺถ ปญฺญตฺตนฺติ สาวตฺถิยํ ปญฺญตฺตํ. กํ อารพฺภาติ ถุลฺลนนฺทํ ภิกฺขุนิํ อารพฺภ. กิสฺมิํ วตฺถุสฺมินฺติ ถุลฺลนนฺทา ภิกฺขุนี สมคฺเคน สํเฆน อุกฺขิตฺตํ อริฏฺฐํ ภิกฺขุํ คทฺธพาธิ\-ปุพฺพํ อนุวตฺติ, ตสฺมิํ วตฺถุสฺมิํ. เอกา ปญฺญตฺติ. ฉนฺนํ อาปตฺติ\-สมุฏฺฐานานํ เอเกน สมุฏฺฐาเนน สมุฏฺฐาติ. ธุรนิกฺเขเป ฯเปฯ.}

\proseitemwithcloser{205}{ภิกฺขุนีนํ อฏฺฐมํ ปาราชิกํ กตฺถ ปญฺญตฺตนฺติ สาวตฺถิยํ ปญฺญตฺตํ. กํ อารพฺภาติ ฉพฺพคฺคิยา ภิกฺขุนิโย อารพฺภ. กิสฺมิํ วตฺถุสฺมินฺติ ฉพฺพคฺคิยา ภิกฺขุนิโย อฏฺฐมํ วตฺถุํ ปริปูเรสุํ, ตสฺมิํ วตฺถุสฺมิํ. เอกา ปญฺญตฺติ. ฉนฺนํ อาปตฺติ\-สมุฏฺฐานานํ เอเกน สมุฏฺฐาเนน สมุฏฺฐาติ. ธุรนิกฺเขเป ฯเปฯ.}{อฏฺฐ ปาราชิกา นิฏฺฐิตา.}
\csromancenter{ตสฺสุทฺทานํ.}
\setlength{\csromangathapagewidth}{0pt}
\setlength{\csromangathawakwidth}{0pt}
\csromangathameasuregroup{เมถุนาทินฺนาทานญฺจ, \\ กายสํสคฺคํ ฉาเทติ,}{\csromangathabat{เมถุนาทินฺนาทานญฺจ,}{มนุสฺสวิคฺคหุตฺตริ.} \\ \csromangathabat{กายสํสคฺคํ ฉาเทติ,}{อุกฺขิตฺตา อฏฺฐ วตฺถุกา.}}
\csromangathasetleft{เมถุนาทินฺนาทานญฺจ, \\ กายสํสคฺคํ ฉาเทติ,}
\csromangathagroup{\csromangathastanza{\csromangathabat{เมถุนาทินฺนาทานญฺจ,}{มนุสฺสวิคฺคหุตฺตริ.} \\ \csromangathabat{กายสํสคฺคํ ฉาเทติ,}{อุกฺขิตฺตา อฏฺฐ วตฺถุกา.}}}

\vspace{\tipitakaparskip}
\csromancenter{ปญฺญาเปสิ มหาวีโร, เฉชฺชวตฺถู อสํสยาติ.}
\csromansectionrule
\csromanmatikaanchor{mtk.76}
\csromantocmarkref{chapter}{2. สํฆาทิเสสกณฺฑ}{mtk.76}
\bookmark[dest={mtk.76},level=0]{2. สํฆาทิเสสกณฺฑ}
\chapterhead{2. สํฆาทิเสส\-กณฺฑ}
\proseitem{206}{ยํ เตน ภควตา ชานตา ปสฺสตา อรหตา สมฺมา\-สมฺพุทฺเธน อุสฺสยวาทิกาย ภิกฺขุนิยา อฑฺฑํ กโรนฺติยา สํฆาทิเสโส กตฺถ ปญฺญตฺโต, กํ อารพฺภ, กิสฺมิํ วตฺถุสฺมิํ ฯเปฯ เกนาภตนฺติ.}

\proseitem{207}{ยํ เตน ภควตา ชานตา ปสฺสตา อรหตา สมฺมา\-สมฺพุทฺเธน อุสฺสยวาทิกาย ภิกฺขุนิยา อฑฺฑํ กโรนฺติยา สํฆาทิเสโส กตฺถ ปญฺญตฺโตติ สาวตฺติยํ ปญฺญตฺโต. กํ อารพฺภาติ ถุลฺลนนฺทํ ภิกฺขุนิํ อารพฺภ. กิสฺมิํ วตฺถุสฺมินฺติ ถุลฺลนนฺทา ภิกฺขุนี อุสฺสยวาทิกา วิหริ, \csromanfolio{100}ตสฺมิํ วตฺถุสฺมิํ. อตฺถิ ตตฺถ ปญฺญตฺติ อนุปญฺญตฺติ อนุปฺปนฺนปญฺญตฺตีติ เอกา ปญฺญตฺติ, อนุปญฺญตฺติ อนุปฺปนฺนปญฺญตฺติ ตสฺมิํ นตฺถิ. สพฺพตฺถ\-ปญฺญตฺติ ปเทสปญฺญตฺตีติ สพฺพตฺถ\-ปญฺญตฺติ. สาธารณ\-ปญฺญตฺติ อสาธารณปญฺญตฺตีติ อสาธารณ\-ปญฺญตฺติ. เอกโตปญฺญตฺติ อุภโตปญฺญตฺตีติ เอกโตปญฺญตฺติ. จตุนฺนํ ปาติโมกฺขุ\-ทฺเทสานํ กตฺโถคธํ กตฺถ ปริยาปนฺนนฺติ นิทาโนคธํ นิทาน\-ปริยาปนฺนํ. กตเมน อุทฺเทเสน อุทฺเทสํ อาคจฺฉตีติ ตติเยน อุทฺเทเสน อุทฺเทสํ อาคจฺฉติ. จตุนฺนํ วิปตฺตีนํ กตมา วิปตฺตีติ สีลวิปตฺติ. สตฺตนฺนํ อาปตฺติ\-กฺขนฺธานํ กตโม อาปตฺติ\-กฺขนฺโธติ สํฆาทิเสสา\-ปตฺติกฺขนฺโธ. ฉนฺนํ อาปตฺติ\-สมุฏฺฐานานํ กติหิ สมุฏฺฐาเนหิ สมุฏฺฐาตีติ ทฺวีหิ สมุฏฺฐาเนหิ สมุฏฺฐาติ, สิยา กายโต จ วาจโต จ สมุฏฺฐาติ, น จิตฺตโต, สิยา กายโต จ วาจโต จ จิตฺตโต จ สมุฏฺฐาติ ฯเปฯ เกนาภตนฺติ ปรมฺปราภตํ.}

\setlength{\csromangathapagewidth}{0pt}
\setlength{\csromangathawakwidth}{0pt}
\csromangathameasuregroup{อุปาลิ ทาสโก เจว, \\ โมคฺคลิปุตฺเตน ปญฺจมา, \\ เอเต นาคา มหาปญฺญา, \\ วินยํ ทีเป ปกาเสสุํ,}{\csromangathabat{อุปาลิ ทาสโก เจว,}{โสณโก สิคฺคโว ตถา.} \\ \csromangathabat{โมคฺคลิปุตฺเตน ปญฺจมา,}{เอเต ชมฺพุสิริวฺหเย ฯเปฯ.} \\ \csromangathabat{เอเต นาคา มหาปญฺญา,}{วินยญฺญู มคฺคโกวิทา.} \\ \csromangathabat{วินยํ ทีเป ปกาเสสุํ,}{ปิฏกํ ตมฺพปณฺฑิยาติ.}}
\csromangathasetleft{อุปาลิ ทาสโก เจว, \\ โมคฺคลิปุตฺเตน ปญฺจมา, \\ เอเต นาคา มหาปญฺญา, \\ วินยํ ทีเป ปกาเสสุํ,}
\csromangathagroup{\csromangathastanza{\csromangathabat{อุปาลิ ทาสโก เจว,}{โสณโก สิคฺคโว ตถา.} \\ \csromangathabat{โมคฺคลิปุตฺเตน ปญฺจมา,}{เอเต ชมฺพุสิริวฺหเย ฯเปฯ.}}\csromangathastanza{\csromangathabat{เอเต นาคา มหาปญฺญา,}{วินยญฺญู มคฺคโกวิทา.} \\ \csromangathabat{วินยํ ทีเป ปกาเสสุํ,}{ปิฏกํ ตมฺพปณฺฑิยาติ.}}}

\proseitem{208}{โจริํ วุฏฺฐาเปนฺติยา สํฆาทิเสโส กตฺถ ปญฺญตฺโตติ สาวตฺติยํ ปญฺญตฺโต. กํ อารพฺภาติ ถุลฺลนนฺทํ ภิกฺขุนิํ อารพฺภ. กิสฺมิํ วตฺถุสฺมินฺติ ถุลฺลนนฺทา ภิกฺขุนี โจริํ วุฏฺฐาเปสิ, ตสฺมิํ วตฺถุสฺมิํ. เอกา ปญฺญตฺติ. ฉนฺนํ อาปตฺติ\-สมุฏฺฐานานํ ทฺวีหิ สมุฏฺฐาเนหิ สมุฏฺฐาติ, สิยา วาจโต จ จิตฺตโต จ สมุฏฺฐาติ, น กายโต, สิยา กายโต จ วาจโต จ}

\proseitem{209}{เอกาย คามนฺตรํ คจฺฉนฺติยา สํฆาทิเสโส กตฺถ ปญฺญตฺโตติ สาวตฺถิยํ ปญฺญตฺโต. กํ อารพฺภาติ อญฺญตรํ ภิกฺขุนิํ อารพฺภ. กิสฺมิํ วตฺถุสฺมินฺติ อญฺญตรา ภิกฺขุนี เอกา คามนฺตรํ คจฺฉิ, ตสฺมิํ วตฺถุสฺมิํ. เอกา ปญฺญตฺติ, ติสฺโส อนุปญฺญตฺติโย. ฉนฺนํ อาปตฺติ\-สมุฏฺฐานานํ เอเกน สมุฏฺฐาเนน สมุฏฺฐาติ, ปฐม\-ปาราชิเก ฯเปฯ.}

\proseitem{210}{สมคฺเคน สํเฆน อุกฺขิตฺตํ ภิกฺขุนิํ ธมฺเมน วินเยน สตฺถุสาสเนน อนปโลเกตฺวา การกสํฆํ อนญฺญาย คณสฺส \csromanfolio{101}ฉนฺทํ โอสาเรนฺติยา สํฆาทิเสโส กตฺถ ปญฺญตฺโตติ สาวตฺถิยํ ปญฺญตฺโต. กํ อารพฺภาติ ถุลฺลนนฺทํ ภิกฺขุนิํ อารพฺภ. กิสฺมิํ วตฺถุสฺมินฺติ ถุลฺลนนฺทา ภิกฺขุนี สมคฺเคน สํเฆน อุกฺขิตฺตํ ภิกฺขุนิํ ธมฺเมน วินเยน สตฺถุสาสเนน อนปโลเกตฺวา การกสํฆํ อนญฺญาย คณสฺส ฉนฺทํ โอสาเรสิ, ตสฺมิํ วตฺถุสฺมิํ. เอกา ปญฺญตฺติ. ฉนฺนํ อาปตฺติ\-สมุฏฺฐานานํ เอเกน สมุฏฺฐาเนน สมุฏฺฐาติ, ธุรนิกฺเขเป ฯเปฯ.}

\proseitem{211}{อวสฺสุตาย ภิกฺขุนิยา อวสฺสุตสฺส ปุริส\-ปุคฺคลสฺส หตฺถโต ขาทนียํ วา โภชนียํ วา สหตฺถา ปฏิคฺคเหตฺวา ภุญฺชนฺติยา สํฆาทิเสโส กตฺถ ปญฺญตฺโตติ สาวตฺถิยํ ปญฺญตฺโต. กํ อารพฺภาติ สุนฺทรีนนฺทํ ภิกฺขุนิํ อารพฺภ. กิสฺมิํ วตฺถุสฺมินฺติ สุนฺทรีนนฺทา ภิกฺขุนี อวสฺสุตา อวสฺสุตสฺส ปุริส\-ปุคฺคลสฺส หตฺถโต อามิสํ ปฏิคฺคเหสิ, ตสฺมิํ วตฺถุสฺมิํ. เอกา ปญฺญตฺติ. ฉนฺนํ อาปตฺติ\-สมุฏฺฐานานํ เอเกน สมุฏฺฐาเนน สมุฏฺฐาติ. ปฐม\-ปาราชิเก ฯเปฯ.}

\proseitem{212}{“กิํ เต อยฺเย เอโส ปุริสปุคฺคโล กริสฺสติ อวสฺสุโต วา อนวสฺสุโต วา, ยโต ตฺวํ อนวสฺสุตา, อิงฺฆ อยฺเย, ยํ เต เอโส ปุริสปุคฺคโล เทติ ขาทนียํ วา โภชนียํ วา, ตํ ตฺวํ สหตฺถา ปฏิคฺคเหตฺวา ขาท วา ภุญฺช วา”ติ อุยฺโยเชนฺติยา สํฆาทิเสโส กตฺถ ปญฺญตฺโตติ สาวตฺถิยํ ปญฺญตฺโต. กํ อารพฺภาติ อญฺญตรํ ภิกฺขุนิํ อารพฺภ. กิสฺมิํ วตฺถุสฺมินฺติ อญฺญตรา ภิกฺขุนี “กิํ เต อยฺเย เอโส ปุริสปุคฺคโล กริสฺสติ อวสฺสุโต วา อนวสฺสุโต วา, ยโต ตฺวํ อนวสฺสุตา, อิงฺฆ อยฺเย, ยํ เต เอโส ปุริสปุคฺคโล เทติ ขาทนียํ วา โภชนียํ วา, ตํ ตฺวํ สหตฺถา ปฏิคฺคเหตฺวา ขาท วา ภุญฺช วา”ติ อุยฺโยเชสิ, ตสฺมิํ วตฺถุสฺมิํ. เอกา ปญฺญตฺติ. ฉนฺนํ อาปตฺติ\-สมุฏฺฐานานํ ตีหิ สมุฏฺฐาเนหิ สมุฏฺฐาติ ฯเปฯ.}

\proseitem{213}{กุปิตาย อนตฺตมนาย ภิกฺขุนิยา ยาวตติยํ สมนุภาสนาย น ปฏินิสฺสชฺชนฺติยา สํฆาทิเสโส กตฺถ ปญฺญตฺโตติ สาวตฺถิยํ ปญฺญตฺโต. กํ อารพฺภาติ จณฺฑกาฬิํ ภิกฺขุนิํ อารพฺภ. กิสฺมิํ วตฺถุสฺมินฺติ จณฺฑกาฬี ภิกฺขุนี กุปิตา อนตฺตมนา เอวํ อวจ “พุทฺธํ ปจฺจาจิกฺขามิ, ธมฺมํ ปจฺจาจิกฺขามิ, สํฆํ ปจฺจาจิกฺขามิ, สิกฺขํ ปจฺจาจิกฺขามี”ติ, \csromanfolio{102}ตสฺมิํ วตฺถุสฺมิํ. เอกา ปญฺญตฺติ. ฉนฺนํ อาปตฺติ\-สมุฏฺฐานานํ เอเกน สมุฏฺฐาเนน สมุฏฺฐาติ, ธุรนิกฺเขเป ฯเปฯ.}

\proseitem{214}{กิสฺมิญฺจิเทว อธิกรเณ ปจฺจากตาย ภิกฺขุนิยา ยาวตติยํ สมนุภาสนาย น ปฏินิสฺสชฺชนฺติยา สํฆาทิเสโส กตฺถ ปญฺญตฺโตติ สาวตฺถิยํ ปญฺญตฺโต. กํ อารพฺภาติ จณฺฑกาฬิํ ภิกฺขุนิํ อารพฺภ. กิสฺมิํ วตฺถุสฺมินฺติ จณฺฑกาฬี ภิกฺขุนี กิสฺมิญฺจิเทว อธิกรเณ ปจฺจากตา กุปิตา อนตฺตมนา เอวํ อวจ “ฉนฺทคามินิโย จ ภิกฺขุนิโย, โทสคามินิโย จ ภิกฺขุนิโย, โมหคามินิโย จ ภิกฺขุนิโย, ภยคามินิโย จ ภิกฺขุนิโย”ติ, ตสฺมิํ วตฺถุสฺมิํ. เอกา ปญฺญตฺติ. ฉนฺนํ อาปตฺติ\-สมุฏฺฐานานํ เอเกน สมุฏฺฐาเนน สมุฏฺฐาติ, ธุรนิกฺเขเป ฯเปฯ.}

\proseitem{215}{สํสฏฺฐานํ ภิกฺขุนีนํ ยาวตติยํ สมนุภาสนาย น ปฏินิสฺสชฺชนฺตีนํ สํฆาทิเสโส กตฺถ ปญฺญตฺโตติ สาวตฺถิยํ ปญฺญตฺโต. กํ อารพฺภาติ สมฺพหุลา ภิกฺขุนิโย อารพฺภ. กิสฺมิํ วตฺถุสฺมินฺติ สมฺพหุลา ภิกฺขุนิโย สํสฏฺฐา วิหริํสุ, ตสฺมิํ วตฺถุสฺมิํ. เอกา ปญฺญตฺติ. ฉนฺนํ อาปตฺติ\-สมุฏฺฐานานํ เอเกน สมุฏฺฐาเนน สมุฏฺฐาติ, ธุรนิกฺเขเป ฯเปฯ.}

\proseitemwithcloser{216}{“สํสฏฺฐาว อยฺเย ตุเมฺห วิหรถ, มา ตุเมฺห นานา วิหริตฺถา”ติ อุยฺโยเชนฺติยา ยาวตติยํ สมนุภาสนาย น ปฏินิสฺสชฺชนฺติยา สํฆาทิเสโส กตฺถ ปญฺญตฺโตติ สาวตฺถิยํ ปญฺญตฺโต. กํ อารพฺภาติ ถุลฺลนนฺทํ ภิกฺขุนิํ อารพฺภ. กิสฺมิํ วตฺถุสฺมินฺติ ถุลฺลนนฺทา ภิกฺขุนี “สํสฏฺฐาว อยฺเย ตุเมฺห วิหรถ, มา ตุเมฺห นานา วิหริตฺถา”ติ อุยฺโยเชสิ, ตสฺมิํ วตฺถุสฺมิํ. เอกา ปญฺญตฺติ. ฉนฺนํ อาปตฺติ\-สมุฏฺฐานานํ เอเกน สมุฏฺฐาเนน สมุฏฺฐาติ, ธุรนิกฺเขเป ฯเปฯ.}{ทส สํฆาทิเสสา นิฏฺฐิตา.}
\csromancenter{ตสฺสุทฺทานํ.}
\setlength{\csromangathapagewidth}{0pt}
\setlength{\csromangathawakwidth}{0pt}
\csromangathameasuregroup{อุสฺสยโจริ คามนฺตํ, \\ กิํ เต กุปิตา กิสฺมิญฺจิ,}{\csromangathabat{อุสฺสยโจริ คามนฺตํ,}{อุกฺขิตฺตํ ขาทเนน จ.} \\ \csromangathabat{กิํ เต กุปิตา กิสฺมิญฺจิ,}{สํสฏฺฐา ญายเต ทสาติ.}}
\csromangathasetleft{อุสฺสยโจริ คามนฺตํ, \\ กิํ เต กุปิตา กิสฺมิญฺจิ,}
\csromangathagroup{\csromangathastanza{\csromangathabat{อุสฺสยโจริ คามนฺตํ,}{อุกฺขิตฺตํ ขาทเนน จ.} \\ \csromangathabat{กิํ เต กุปิตา กิสฺมิญฺจิ,}{สํสฏฺฐา ญายเต ทสาติ.}}}

\csromanmatikaanchor{mtk.77}
\csromantocmarkref{chapter}{3. นิสฺสคฺคิยกณฺฑ}{mtk.77}
\bookmark[dest={mtk.77},level=0]{3. นิสฺสคฺคิยกณฺฑ}
\chapterheadpageread{3. นิสฺสคฺคิย\-กณฺฑ}
\proseitem{217}{\csromanfolio{103}ยํ เตน ภควตา ชานตา ปสฺสตา อรหตา สมฺมา\-สมฺพุทฺเธน ปตฺต\-สนฺนิจยํ กโรนฺติยา นิสฺสคฺคิยํ ปาจิตฺติยํ กตฺถ ปญฺญตฺตนฺติ สาวตฺถิยํ ปญฺญตฺตํ. กํ อารพฺภาติ ฉพฺพคฺคิยา ภิกฺขุนิโย อารพฺภ. กิสฺมิํ วตฺถุสฺมินฺติ ฉพฺพคฺคิยา ภิกฺขุนิโย ปตฺต\-สนฺนิจยํ อกํสุ, ตสฺมิํ วตฺถุสฺมิํ. เอกา ปญฺญตฺติ. ฉนฺนํ อาปตฺติ\-สมุฏฺฐานานํ ทฺวีหิ สมุฏฺฐาเนหิ สมุฏฺฐาติ, กถินเก ฯเปฯ.}

\prose{อกาลจีวรํ “กาลจีวรนฺ”ติ อธิฏฺฐหิตฺวา ภาชาเปนฺติยา นิสฺสคฺคิยํ ปาจิตฺติยํ กตฺถ ปญฺญตฺตนฺติ สาวตฺถิยํ ปญฺญตฺตํ. กํ อารพฺภาติ ถุลฺลนนฺทํ ภิกฺขุนิํ อารพฺภ. กิสฺมิํ วตฺถุสฺมินฺติ ถุลฺลนนฺทา ภิกฺขุนี อกาลจีวรํ “กาลจีวรนฺ”ติ อธิฏฺฐหิตฺวา ภาชาเปสิ, ตสฺมิํ วตฺถุสฺมิํ. เอกา ปญฺญตฺติ. ฉนฺนํ อาปตฺติ\-สมุฏฺฐานานํ ตีหิ สมุฏฺฐาเนหิ}

\prose{ภิกฺขุนิยา สทฺธิํ จีวรํ ปริวตฺเตตฺวา อจฺฉินฺทนฺติยา นิสฺสคฺคิยํ ปาจิตฺติยํ กตฺถ ปญฺญตฺตนฺติ สาวตฺถิยํ ปญฺญตฺตํ. กํ อารพฺภาติ ถุลฺลนนฺทํ ภิกฺขุนิํ อารพฺภ. กิสฺมิํ วตฺถุสฺมินฺติ ถุลฺลนนฺทา ภิกฺขุนี ภิกฺขุนิยา สทฺธิํ จีวรํ ปริวตฺเตตฺวา อจฺฉินฺทิ, ตสฺมิํ วตฺถุสฺมิํ. เอกา ปญฺญตฺติ. ฉนฺนํ อาปตฺติ\-สมุฏฺฐานานํ ตีหิ สมุฏฺฐาเนหิ}

\prose{อญฺญํ วิญฺญาเปตฺวา อญฺญํ วิญฺญาเปนฺติยา นิสฺสคฺคิยํ ปาจิตฺติยํ กตฺถ ปญฺญตฺตนฺติ สาวตฺถิยํ ปญฺญตฺตํ. กํ อารพฺภาติ ถุลฺลนนฺทํ ภิกฺขุนิํ อารพฺภ. กิสฺมิํ วตฺถุสฺมินฺติ ถุลฺลนนฺทา ภิกฺขุนี อญฺญํ วิญฺญาเปตฺวา อญฺญํ วิญฺญาเปสิ, ตสฺมิํ วตฺถุสฺมิํ. เอกา ปญฺญตฺติ. ฉนฺนํ อาปตฺติ\-สมุฏฺฐานานํ ฉหิ สมุฏฺฐาเนหิ สมุฏฺฐาติ ฯเปฯ.}

\prose{อญฺญํ เจตาเปตฺวา อญฺญํ เจตาเปนฺติยา นิสฺสคฺคิยํ ปาจิตฺติยํ กตฺถ ปญฺญตฺตนฺติ สาวตฺถิยํ ปญฺญตฺตํ. กํ อารพฺภาติ ถุลฺลนนฺทํ ภิกฺขุนิํ อารพฺภ. กิสฺมิํ วตฺถุสฺมินฺติ ถุลฺลนนฺทา ภิกฺขุนี อญฺญํ เจตาเปตฺวา อญฺญํ เจตาเปสิ, ตสฺมิํ วตฺถุสฺมิํ. เอกา ปญฺญตฺติ. ฉนฺนํ อาปตฺติ\-สมุฏฺฐานานํ ฉหิ สมุฏฺฐาเนหิ สมุฏฺฐาติ ฯเปฯ.}

\prose{อญฺญาทตฺถิเกน ปริกฺขาเรน อญฺญุทฺทิสิเกน สํฆิเกน อญฺญํ เจตาเปนฺติยา นิสฺสคฺคิยํ ปาจิตฺติยํ กตฺถ ปญฺญตฺตนฺติ สาวตฺถิยํ ปญฺญตฺตํ. กํ \csromanfolio{104}อารพฺภาติ สมฺพหุลา ภิกฺขุนิโย อารพฺภ. กิสฺมิํ วตฺถุสฺมินฺติ สมฺพหุลา ภิกฺขุนิโย อญฺญทตฺถิเกน ปริกฺขาเรน อญฺญุทฺทิสิเกน สํฆิเกน อญฺญํ เจตาเปสุํ, ตสฺมิํ วตฺถุสฺมิํ. เอกา ปญฺญตฺติ. ฉนฺนํ อาปตฺติ\-สมุฏฺฐานานํ ฉหิ สมุฏฺฐาเนหิ สมุฏฺฐาติ ฯเปฯ.}

\prose{อญฺญทตฺถิเกน ปริกฺขาเรน อญฺญุทฺทิสิเกน สํฆิเกน สญฺญาจิเกน อญฺญํ เจตาเปนฺติยา นิสฺสคฺคิยํ ปาจิตฺติยํ กตฺถ ปญฺญตฺตนฺติ สาวตฺถิยํ ปญฺญตฺตํ. กํ อารพฺภาติ สมฺพหุลา ภิกฺขุนิโย อารพฺภ. กิสฺมิํ วตฺถุสฺมินฺติ สมฺพหุลา ภิกฺขุนิโย อญฺญทตฺถิเกน ปริกฺขาเรน อญฺญุทฺทิสิเกน สํฆิเกน สญฺญาจิเกน อญฺญํ เจตาเปสุํ, ตสฺมิํ วตฺถุสฺมิํ. เอกา ปญฺญตฺติ. ฉนฺนํ อาปตฺติ\-สมุฏฺฐานานํ ฉหิ สมุฏฺฐาเนหิ สมุฏฺฐาติ ฯเปฯ.}

\prose{อญฺญทตฺถิเกน ปริกฺขาเรน อญฺญุทฺทิสิเกน มหาชนิเกน อญฺญํ เจตาเปนฺติยา นิสฺสคฺคิยํ ปาจิตฺติยํ กตฺถ ปญฺญตฺตนฺติ สาวตฺถิยํ ปญฺญตฺตํ. กํ อารพฺภาติ สมฺพหุลา ภิกฺขุนิโย อารพฺภ. กิสฺมิํ วตฺถุสฺมินฺติ สมฺพหุลา ภิกฺขุนิโย อญฺญทตฺถิเกน ปริกฺขาเรน อญฺญุทฺทิสิเกน มหาชนิเกน อญฺญํ เจตาเปสุํ, ตสฺมิํ วตฺถุสฺมิํ. เอกา ปญฺญตฺติ. ฉนฺนํ อาปตฺติ\-สมุฏฺฐานานํ ฉหิ สมุฏฺฐาเนหิ สมุฏฺฐาติ ฯเปฯ.}

\prose{อญฺญทตฺถิเกน ปริกฺขาเรน อญฺญุทฺทิสิเกน มหาชนิเกน สญฺญาจิเกน อญฺญํ เจตาเปนฺติยา นิสฺสคฺคิยํ ปาจิตฺติยํ กตฺถ ปญฺญตฺตนฺติ สาวตฺถิยํ ปญฺญตฺตํ. กํ อารพฺภาติ สมฺพหุลา ภิกฺขุนิโย อารพฺภ. กิสฺมิํ วตฺถุสฺมินฺติ สมฺพหุลา ภิกฺขุนิโย อญฺญทตฺถิเกน ปริกฺขาเรน อญฺญุทฺทิสิเกน มหาชนิเกน สญฺญาจิเกน อญฺญํ เจตาเปสุํ, ตสฺมิํ วตฺถุสฺมิํ. เอกา ปญฺญตฺติ. ฉนฺนํ อาปตฺติ\-สมุฏฺฐานานํ ฉหิ สมุฏฺฐาเนหิ สมุฏฺฐาติ ฯเปฯ.}

\prose{อญฺญทตฺถิเกน ปริกฺขาเรน อญฺญุทฺทิสิเกน ปุคฺคลิเกน สญฺญาจิเกน อญฺญํ เจตาเปนฺติยา นิสฺสคฺคิยํ ปาจิตฺติยํ กตฺถ ปญฺญตฺตนฺติ สาวตฺถิยํ ปญฺญตฺตํ. กํ อารพฺภาติ ถุลฺลนนฺทํ ภิกฺขุนิํ อารพฺภ. กิสฺมิํ วตฺถุสฺมินฺติ ถุลฺลนนฺทา ภิกฺขุนี อญฺญทตฺถิเกน ปริกฺขาเรน อญฺญุทฺทิสิเกน ปุคฺคลิเกน สญฺญาจิเกน อญฺญํ เจตาเปสิ, ตสฺมิํ วตฺถุสฺมิํ. เอกา ปญฺญตฺติ. ฉนฺนํ อาปตฺติ\-สมุฏฺฐานานํ ฉหิ สมุฏฺฐาเนหิ สมุฏฺฐาติ ฯเปฯ.}

\prose{\csromanfolio{105}อติเรกจตุกฺกํส\-ปรมํ ครุปาวุรณํ เจตาเปนฺติยา นิสฺสคฺคิยํ ปาจิตฺติยํ กตฺถ ปญฺญตฺตนฺติ สาวตฺถิยํ ปญฺญตฺตํ. กํ อารพฺภาติ ถุลฺลนนฺทํ ภิกฺขุนิํ อารพฺภ. กิสฺมิํ วตฺถุสฺมินฺติ ถุลฺลนนฺทา ภิกฺขุนี ราชานํ กมฺพลํ วิญฺญาเปสิ, ตสฺมิํ วตฺถุสฺมิํ. เอกา ปญฺญตฺติ. ฉนฺนํ อาปตฺติ\-สมุฏฺฐานานํ ฉหิ สมุฏฺฐาเนหิ สมุฏฺฐาติ ฯเปฯ.}

\prosewithcloser{อติเรก\-อฑฺฒเตยฺย\-กํส\-ปรมํ ลหุปาวุรณํ เจตาเปนฺติยา นิสฺสคฺคิยํ ปาจิตฺติยํ กตฺถ ปญฺญตฺตนฺติ สาวตฺถิยํ ปญฺญตฺตํ. กํ อารพฺภาติ ถุลฺลนนฺทํ ภิกฺขุนิํ อารพฺภ. กิสฺมิํ วตฺถุสฺมินฺติ ถุลฺลนนฺทา ภิกฺขุนี ราชานํ โขมํ วิญฺญาเปสิ, ตสฺมิํ วตฺถุสฺมิํ. เอกา ปญฺญตฺติ. ฉนฺนํ อาปตฺติ\-สมุฏฺฐานานํ ฉหิ สมุฏฺฐาเนหิ สมุฏฺฐาติ ฯเปฯ.}{ทฺวาทส นิสฺสคฺคิยา ปาจิตฺติยา นิฏฺฐิตา.}
\csromancenter{ตสฺสุทฺทานํ.}
\setlength{\csromangathapagewidth}{0pt}
\setlength{\csromangathawakwidth}{0pt}
\csromangathameasuregroup{ปตฺตํ อกาลํ กาลญฺจ, \\ เจตาเปตฺวา อญฺญทตฺถิ,}{\csromangathabat{ปตฺตํ อกาลํ กาลญฺจ,}{ปริวตฺเต จ วิญฺญาเป.} \\ \csromangathabat{เจตาเปตฺวา อญฺญทตฺถิ,}{สํฆิกญฺจ มหาชนิกํ.}}
\csromangathasetleft{ปตฺตํ อกาลํ กาลญฺจ, \\ เจตาเปตฺวา อญฺญทตฺถิ,}
\csromangathagroup{\csromangathastanza{\csromangathabat{ปตฺตํ อกาลํ กาลญฺจ,}{ปริวตฺเต จ วิญฺญาเป.} \\ \csromangathabat{เจตาเปตฺวา อญฺญทตฺถิ,}{สํฆิกญฺจ มหาชนิกํ.}}}

\vspace{\tipitakaparskip}
\csromancenter{สญฺญาจิกา ปุคฺคลิกา, จตุกฺกํส\-ฑฺฒเตยฺย\-กาติ.}
\csromansectionrule
\csromanmatikaanchor{mtk.78}
\csromanmatikaanchor{mtk.79}
\csromantocmarkref{chapter}{4. ปาจิตฺติยกณฺฑ}{mtk.78}
\bookmark[dest={mtk.78},level=0]{4. ปาจิตฺติยกณฺฑ}
\csromantocmarkref{section}{1. ลสุณวคฺค}{mtk.79}
\bookmark[dest={mtk.79},level=1]{1. ลสุณวคฺค}
\csromanmarkh{4. ปาจิตฺติยกณฺฑ}\csromanmarkhii{1. ลสุณวคฺค}\csromanheadernomark{1}{4. \textbf{ปาจิตฺติยกณฺฑ} \textbf{1. ลสุณวคฺค}}
\proseitem{218}{ยํ เตน ภควตา ชานตา ปสฺสตา อรหตา สมฺมา\-สมฺพุทฺเธน ลสุณํ ขาทนฺติยา ปาจิตฺติยํ กตฺถ ปญฺญตฺตนฺติ สาวตฺถิยํ ปญฺญตฺตํ. กํ อารพฺภาติ ถุลฺลนนฺทํ ภิกฺขุนิํ อารพฺภ. กิสฺมิํ วตฺถุสฺมินฺติ ถุลฺลนนฺทา ภิกฺขุนี น มตฺตํ ชานิตฺวา ลสุณํ หราเปสิ, ตสฺมิํ วตฺถุสฺมิํ. เอกา ปญฺญตฺติ. ฉนฺนํ อาปตฺติ\-สมุฏฺฐานานํ ทฺวีหิ สมุฏฺฐาเนหิ สมุฏฺฐาติ, เอฬกโลมเก ฯเปฯ.}

\prose{สมฺพาเธ โลมํ สํหราเปนฺติยา ปาจิตฺติยํ กตฺถ ปญฺญตฺตนฺติ สาวตฺถิยํ ปญฺญตฺตํ. กํ อารพฺภาติ ฉพฺพคฺคิยา ภิกฺขุนิโย อารพฺภ. กิสฺมิํ วตฺถุสฺมินฺติ ฉพฺพคฺคิยา ภิกฺขุนิโย สมฺพาเธ โลมํ สํหราเปสุํ, ตสฺมิํ วตฺถุสฺมิํ. เอกา ปญฺญตฺติ. ฉนฺนํ อาปตฺติ\-สมุฏฺฐานานํ จตูหิ สมุฏฺฐาเนหิ สมุฏฺฐาติ ฯเปฯ.}

\prose{\csromanfolio{106}ตลฆาตเก ปาจิตฺติยํ กตฺถ ปญฺญตฺตนฺติ สาวตฺถิยํ ปญฺญตฺตํ. กํ อารพฺภาติ เทฺว ภิกฺขุนิโย อารพฺภ. กิสฺมิํ วตฺถุสฺมินฺติ เทฺว ภิกฺขุนิโย ตลฆาตกํ อกํสุ, ตสฺมิํ วตฺถุสฺมิํ. เอกา ปญฺญตฺติ. ฉนฺนํ อาปตฺติ\-สมุฏฺฐานานํ เอเกน สมุฏฺฐาเนน สมุฏฺฐาติ, ปฐม\-ปาราชิเก ฯเปฯ.}

\prose{ชตุมฏฺฐเก ปาจิตฺติยํ กตฺถ ปญฺญตฺตนฺติ สาวตฺถิยํ ปญฺญตฺตํ. กํ อารพฺภาติ อญฺญตรํ ภิกฺขุนิํ อารพฺภ. กิสฺมิํ วตฺถุสฺมินฺติ อญฺญตรา ภิกฺขุนี ชตุมฏฺฐกํ อาทิยิ, ตสฺมิํ วตฺถุสฺมิํ. เอกา ปญฺญตฺติ. ฉนฺนํ อาปตฺติ\-สมุฏฺฐานานํ เอเกน สมุฏฺฐาเนน สมุฏฺฐาติ, ปฐม\-ปาราชิเก ฯเปฯ.}

\prose{อติเรก\-ทฺวงฺคุลปพฺพ\-ปรมํ อุทกสุทฺธิกํ อาทิยนฺติยา ปาจิตฺติยํ กตฺถ ปญฺญตฺตนฺติ สกฺเกสุ ปญฺญตฺตํ. กํ อารพฺภาติ อญฺญตรํ ภิกฺขุนิํ อารพฺภ. กิสฺมิํ วตฺถุสฺมินฺติ อญฺญตรา ภิกฺขุนี อติคมฺภีรํ อุทกสุทฺธิกํ อาทิยิ, ตสฺมิํ วตฺถุสฺมิํ. เอกา ปญฺญตฺติ. ฉนฺนํ อาปตฺติ\-สมุฏฺฐานานํ เอเกน สมุฏฺฐาเนน สมุฏฺฐาติ, ปฐม\-ปาราชิเก ฯเปฯ.}

\prose{ภิกฺขุสฺส ภุญฺชนฺตสฺส ปานีเยน วา วิธูปเนน วา อุปติฏฺฐนฺติยา ปาจิตฺติยํ กตฺถ ปญฺญตฺตนฺติ สาวตฺถิยํ ปญฺญตฺตํ. กํ อารพฺภาติ อญฺญตรํ ภิกฺขุนิํ อารพฺภ. กิสฺมิํ วตฺถุสฺมินฺติ อญฺญตรา ภิกฺขุนี ภิกฺขุสฺส ภุญฺชนฺตสฺส ปานีเยน จ วิธูปเนน จ อุปติฏฺฐิ, ตสฺมิํ วตฺถุสฺมิํ. เอกา ปญฺญตฺติ. ฉนฺนํ อาปตฺติ\-สมุฏฺฐานานํ ทฺวีหิ สมุฏฺฐาเนหิ สมุฏฺฐาติ, เอฬกโลมเก ฯเปฯ.}

\prose{อามกธญฺญํ วิญฺญาเปตฺวา ภุญฺชนฺติยา ปาจิตฺติยํ กตฺถ ปญฺญตฺตนฺติ สาวตฺถิยํ ปญฺญตฺตํ. กํ อารพฺภาติ สมฺพหุลา ภิกฺขุนิโย อารพฺภ. กิสฺมิํ วตฺถุสฺมินฺติ สมฺพหุลา ภิกฺขุนิโย อามกธญฺญํ วิญฺญาเปตฺวา ภุญฺชิํสุ, ตสฺมิํ วตฺถุสฺมิํ. เอกา ปญฺญตฺติ. ฉนฺนํ อาปตฺติ\-สมุฏฺฐานานํ จตูหิ สมุฏฺฐาเนหิ สมุฏฺฐาติ ฯเปฯ.}

\prose{อุจฺจารํ วา ปสฺสาวํ วา สงฺการํ วา วิฆาสํ วา ติโรกุฏฺเฏ1 ฉฑฺเฑนฺติยา ปาจิตฺติยํ กตฺถ ปญฺญตฺตนฺติ สาวตฺถิยํ ปญฺญตฺตํ. กํ อารพฺภาติ อญฺญตรํ ภิกฺขุนิํ อารพฺภ. กิสฺมิํ วตฺถุสฺมินฺติ อญฺญตรา ภิกฺขุนี อุจฺจารํ2 ติโรกุฏฺเฏ ฉฑฺเฑสิ, ตสฺมิํ วตฺถุสฺมิํ. เอกา ปญฺญตฺติ. ฉนฺนํ อาปตฺติ\-สมุฏฺฐานานํ ฉหิ สมุฏฺฐาเนหิ สมุฏฺฐาติ ฯเปฯ.}

\proseitem{1}{ติโรกุฑฺเฑ (สี, สฺยา) 2. อุจฺจารมฺปิ ปสฺสาวมฺปิ สงฺการมฺปิ วิฆาสมฺปิ (ก)}

\prose{\csromanfolio{107}อุจฺจารํ วา ปสฺสาวํ วา สงฺการํ วา วิฆาสํ วา หริเต ฉฑฺเฑนฺติยา ปาจิตฺติยํ กตฺถ ปญฺญตฺตนฺติ สาวตฺถิยํ ปญฺญตฺตํ. กํ อารพฺภาติ สมฺพหุลา ภิกฺขุนิโย อารพฺภ. กิสฺมิํ วตฺถุสฺมินฺติ สมฺพหุลา ภิกฺขุนิโย อุจฺจารมฺปิ ปสฺสาวมฺปิ สงฺการมฺปิ วิฆาสมฺปิ หริเต ฉฑฺเฑสุํ, ตสฺมิํ วตฺถุสฺมิํ. เอกา ปญฺญตฺติ. ฉนฺนํ อาปตฺติ\-สมุฏฺฐานานํ ฉหิ สมุฏฺฐาเนหิ สมุฏฺฐาติ ฯเปฯ.}

\prosewithclosermidruled{นจฺจํ วา คีตํ วา วาทิตํ วา ทสฺสนาย คจฺฉนฺติยา ปาจิตฺติยํ กตฺถ ปญฺญตฺตนฺติ ราชคเห ปญฺญตฺตํ. กํ อารพฺภาติ ฉพฺพคฺคิยา ภิกฺขุนิโย อารพฺภ. กิสฺมิํ วตฺถุสฺมินฺติ ฉพฺพคฺคิยา ภิกฺขุนิโย นจฺจมฺปิ คีตมฺปิ วาทิตมฺปิ ทสฺสนาย อคมํสุ, ตสฺมิํ วตฺถุสฺมิํ. เอกา ปญฺญตฺติ. ฉนฺนํ อาปตฺติ\-สมุฏฺฐานานํ ทฺวีหิ สมุฏฺฐาเนหิ สมุฏฺฐาติ, เอฬกโลมเก ฯเปฯ.}{ลสุณวคฺโค ปฐโม.}
\csromanmatikaanchor{mtk.80}
\csromantocmarkref{section}{2. รตฺตนฺธการวคฺค}{mtk.80}
\bookmark[dest={mtk.80},level=1]{2. รตฺตนฺธการวคฺค}
\csromanheader{1}{2. รตฺตนฺธการ\-วคฺค}
\proseitem{219}{รตฺตนฺธกาเร อปฺปทีเป ปุริเสน สทฺธิํ เอเกเนกาย สนฺติฏฺฐนฺติยา ปาจิตฺติยํ กตฺถ ปญฺญตฺตนฺติ สาวตฺถิยํ ปญฺญตฺตํ. กํ อารพฺภาติ อญฺญตรํ ภิกฺขุนิํ อารพฺภ. กิสฺมิํ วตฺถุสฺมินฺติ อญฺญตรา ภิกฺขุนี รตฺตนฺธกาเร อปฺปทีเป ปุริเสน สทฺธิํ เอเกเนกา สนฺติฏฺฐิ, ตสฺมิํ วตฺถุสฺมิํ. เอกา ปญฺญตฺติ. ฉนฺนํ อาปตฺติ\-สมุฏฺฐานานํ ทฺวีหิ สมุฏฺฐาเนหิ สมุฏฺฐาติ, เถยฺยสตฺถเก ฯเปฯ.}

\prose{ปฏิจฺฉนฺเน โอกาเส ปุริเสน สทฺธิํ เอเกเนกาย สนฺติฏฺฐนฺติยา ปาจิตฺติยํ กตฺถ ปญฺญตฺตนฺติ สาวตฺถิยํ ปญฺญตฺตํ. กํ อารพฺภาติ อญฺญตรํ ภิกฺขุนิํ อารพฺภ. กิสฺมิํ วตฺถุสฺมินฺติ อญฺญตรา ภิกฺขุนี ปฏิจฺฉนฺเน โอกาเส ปุริเสน สทฺธิํ เอเกเนกา สนฺติฏฺฐิ, ตสฺมิํ วตฺถุสฺมิํ. เอกา ปญฺญตฺติ. ฉนฺนํ อาปตฺติ\-สมุฏฺฐานานํ ทฺวีหิ สมุฏฺฐาเนหิ สมุฏฺฐาติ, เถยฺยสตฺถเก ฯเปฯ.}

\prose{อชฺโฌกาเส ปุริเสน สทฺธิํ เอเกเนกาย สนฺติฏฺฐนฺติยา ปาจิตฺติยํ กตฺถ ปญฺญตฺตนฺติ สาวตฺถิยํ ปญฺญตฺตํ. กํ อารพฺภาติ อญฺญตรํ ภิกฺขุนิํ อารพฺภ. กิสฺมิํ วตฺถุสฺมินฺติ อญฺญตรา ภิกฺขุนี อชฺโฌกาเส ปุริเสน สทฺธิํ เอเกเนกา สนฺติฏฺฐิ, ตสฺมิํ วตฺถุสฺมิํ. เอกา ปญฺญตฺติ. ฉนฺนํ อาปตฺติ\-สมุฏฺฐานานํ ทฺวีหิ สมุฏฺฐาเนหิ สมุฏฺฐาติ, เถยฺยสตฺถเก ฯเปฯ.}

\prose{\csromanfolio{108}รถิกาย วา1 พฺยูเห วา สิงฺฆาฏเก วา ปุริเสน สทฺธิํ เอเกเนกาย สนฺติฏฺฐนฺติยา ปาจิตฺติยํ กตฺถ ปญฺญตฺตนฺติ สาวตฺถิยํ ปญฺญตฺตํ. กํ อารพฺภาติ ถุลฺลนนฺทํ ภิกฺขุนิํ อารพฺภ. กิสฺมิํ วตฺถุสฺมินฺติ ถุลฺลนนฺทา ภิกฺขุนี รถิกายปิ พฺยูเหปิ สิงฺฆาฏเกปิ ปุริเสน สทฺธิํ เอเกเนกา สนฺติฏฺฐิ, ตสฺมิํ วตฺถุสฺมิํ. เอกา ปญฺญตฺติ. ฉนฺนํ อาปตฺติ\-สมุฏฺฐานานํ ทฺวีหิ สมุฏฺฐาเนหิ สมุฏฺฐาติ, เถยฺยสตฺถเก. ปุเรภตฺตํ กุลานิ อุปสงฺกมิตฺวา อาสเน นิสีทิตฺวา สามิเก อนาปุจฺฉา ปกฺกมนฺติยา ปาจิตฺติยํ กตฺถ ปญฺญตฺตนฺติ สาวตฺถิยํ ปญฺญตฺตํ. กํ อารพฺภาติ อญฺญตรํ ภิกฺขุนิํ อารพฺภ. กิสฺมิํ วตฺถุสฺมินฺติ อญฺญตรา ภิกฺขุนี ปุเรภตฺตํ กุลานิ อุปสงฺกมิตฺวา อาสเน นิสีทิตฺวา สามิเก อนาปุจฺฉา ปกฺกามิ, ตสฺมิํ วตฺถุสฺมิํ. เอกา ปญฺญตฺติ. ฉนฺนํ อาปตฺติ\-สมุฏฺฐานานํ ทฺวีหิ สมุฏฺฐาเนหิ สมุฏฺฐาติ, กถินเก ฯเปฯ.}

\prose{ปจฺฉาภตฺตํ กุลานิ อุปสงฺกมิตฺวา สามิเก อนาปุจฺฉา อาสเน อภินิ\-สีท\-นฺติยา ปาจิตฺติยํ กตฺถ ปญฺญตฺตนฺติ สาวตฺถิยํ ปญฺญตฺตํ. กํ อารพฺภาติ ถุลฺลนนฺทํ ภิกฺขุนิํ อารพฺภ. กิสฺมิํ วตฺถุสฺมินฺติ ถุลฺลนนฺทา ภิกฺขุนี ปจฺฉาภตฺตํ กุลานิ อุปสงฺกมิตฺวา สามิเก อนาปุจฺฉา อาสเน อภินิสีทิ, ตสฺมิํ วตฺถุสฺมิํ. เอกา ปญฺญตฺติ. ฉนฺนํ อาปตฺติ\-สมุฏฺฐานานํ ทฺวีหิ สมุฏฺฐาเนหิ สมุฏฺฐาติ, กถินเก ฯเปฯ.}

\prose{วิกาเล กุลานิ อุปสงฺกมิตฺวา สามิเก อนาปุจฺฉา เสยฺยํ สนฺถริตฺวา วา สนฺถราเปตฺวา วา อภินิ\-สีท\-นฺติยา ปาจิตฺติยํ กตฺถ ปญฺญตฺตนฺติ สาวตฺถิยํ ปญฺญตฺตํ. กํ อารพฺภาติ สมฺพหุลา ภิกฺขุนิโย อารพฺภ. กิสฺมิํ วตฺถุสฺมินฺติ สมฺพหุลา ภิกฺขุนิโย วิกาเล กุลานิ อุปสงฺกมิตฺวา สามิเก อนาปุจฺฉา เสยฺยํ สนฺถริตฺวา อภินิสีทิํสุ, ตสฺมิํ วตฺถุสฺมิํ. เอกา ปญฺญตฺติ. ฉนฺนํ อาปตฺติ\-สมุฏฺฐานานํ ทฺวีหิ สมุฏฺฐาเนหิ สมุฏฺฐาติ, กถินเก ฯเปฯ.}

\prose{ทุคฺคหิเตน ทูปธาริเตน ปรํ อุชฺฌาเปนฺติยา ปาจิตฺติยํ กตฺถ ปญฺญตฺตนฺติ สาวตฺถิยํ ปญฺญตฺตํ. กํ อารพฺภาติ อญฺญตรํ ภิกฺขุนิํ อารพฺภ. กิสฺมิํ วตฺถุสฺมินฺติ อญฺญตรา ภิกฺขุนี ทุคฺคหิเตน ทูปธาริเตน ปรํ อุชฺฌาเปสิ, ตสฺมิํ วตฺถุสฺมิํ. เอกา ปญฺญตฺติ. ฉนฺนํ อาปตฺติ\-สมุฏฺฐานานํ ตีหิ สมุฏฺฐาเนหิ สมุฏฺฐาติ ฯเปฯ.}

\vspace{\tipitakaparskip}
\csromancenter{รถิยาย วา (ก)}
\prose{\csromanfolio{109}อตฺตานํ วา ปรํ วา นิรเยน วา พฺรหฺมจริเยน วา อภิสปนฺติยา ปาจิตฺติยํ กตฺถ ปญฺญตฺตนฺติ สาวตฺถิยํ ปญฺญตฺตํ. กํ อารพฺภาติ จณฺฑกาฬิํ ภิกฺขุนิํ อารพฺภ. กิสฺมิํ วตฺถุสฺมินฺติ จณฺฑกาฬี ภิกฺขุนี อตฺตานมฺปิ ปรมฺปิ นิรเยนปิ พฺรหฺมจริเยนปิ อภิสปิ, ตสฺมิํ วตฺถุสฺมิํ. เอกา ปญฺญตฺติ. ฉนฺนํ อาปตฺติ\-สมุฏฺฐานานํ ตีหิ สมุฏฺฐาเนหิ}

\prosewithclosermidruled{อตฺตานํ วธิตฺวา วธิตฺวา โรทนฺติยา ปาจิตฺติยํ กตฺถ ปญฺญตฺตนฺติ สาวตฺถิยํ ปญฺญตฺตํ. กํ อารพฺภาติ จณฺฑกาฬิํ ภิกฺขุนิํ อารพฺภ. กิสฺมิํ วตฺถุสฺมินฺติ จณฺฑกาฬี ภิกฺขุนี อตฺตานํ วธิตฺวา วธิตฺวา โรทิ, ตสฺมิํ วตฺถุสฺมิํ. เอกา ปญฺญตฺติ. ฉนฺนํ อาปตฺติ\-สมุฏฺฐานานํ เอเกน สมุฏฺฐาเนน สมุฏฺฐาติ, ธุรนิกฺเขเป ฯเปฯ.}{รตฺตนฺธการ\-วคฺโค ทุติโย.}
\csromanmatikaanchor{mtk.81}
\csromantocmarkref{section}{3. นหานวคฺค}{mtk.81}
\bookmark[dest={mtk.81},level=1]{3. นหานวคฺค}
\csromanheader{1}{3. \textbf{นหานวคฺค}}
\proseitem{220}{นคฺคาย นหายนฺติยา ปาจิตฺติยํ กตฺถ ปญฺญตฺตนฺติ สาวตฺถิยํ ปญฺญตฺตํ. กํ อารพฺภาติ สมฺพหุลา ภิกฺขุนิโย อารพฺภ. กิสฺมิํ วตฺถุสฺมินฺติ สมฺพหุลา ภิกฺขุนิโย นคฺคา นหายิํสุ, ตสฺมิํ วตฺถุสฺมิํ. เอกา ปญฺญตฺติ. ฉนฺนํ อาปตฺติ\-สมุฏฺฐานานํ ทฺวีหิ สมุฏฺฐาเนหิ สมุฏฺฐาติ, เอฬกโลมเก ฯเปฯ.}

\prose{ปมาณา\-ติกฺกนฺตํ อุทกสาฏิกํ การาเปนฺติยา ปาจิตฺติยํ กตฺถ ปญฺญตฺตนฺติ สาวตฺถิยํ ปญฺญตฺตํ. กํ อารพฺภาติ ฉพฺพคฺคิยา ภิกฺขุนิโย อารพฺภ. กิสฺมิํ วตฺถุสฺมินฺติ ฉพฺพคฺคิยา ภิกฺขุนิโย อปฺปมาณิกาโย อุทกสาฏิกาโย ธาเรสุํ, ตสฺมิํ วตฺถุสฺมิํ. เอกา ปญฺญตฺติ. ฉนฺนํ อาปตฺติ\-สมุฏฺฐานานํ ฉหิ สมุฏฺฐาเนหิ สมุฏฺฐาติ ฯเปฯ.}

\prose{ภิกฺขุนิยา จีวรํ วิสิพฺเพตฺวา วา วิสิพฺพาเปตฺวา วา เนว สิพฺเพนฺติยา น สิพฺพาปนาย อุสฺสุกฺกํ กโรนฺติยา ปาจิตฺติยํ กตฺถ ปญฺญตฺตนฺติ สาวตฺถิยํ ปญฺญตฺตํ. กํ อารพฺภาติ ถุลฺลนนฺทํ ภิกฺขุนิํ อารพฺภ. กิสฺมิํ วตฺถุสฺมินฺติ ถุลฺลนนฺทา ภิกฺขุนี ภิกฺขุนิยา จีวรํ วิสิพฺพาเปตฺวา เนว สิพฺเพสิ, น สิพฺพาปนาย อุสฺสุกฺกํ อกาสิ, ตสฺมิํ วตฺถุสฺมิํ. เอกา ปญฺญตฺติ. ฉนฺนํ อาปตฺติ\-สมุฏฺฐานานํ เอเกน สมุฏฺฐาเนน สมุฏฺฐาติ, ธุรนิกฺเขเป ฯเปฯ.}

\prose{\csromanfolio{110}ปญฺจาหิกํ สงฺฆาฏิจารํ อติกฺกาเมนฺติยา ปาจิตฺติยํ กตฺถ ปญฺญตฺตนฺติ สาวตฺถิยํ ปญฺญตฺตํ. กํ อารพฺภาติ สมฺพหุลา ภิกฺขุนิโย อารพฺภ. กิสฺมิํ วตฺถุสฺมินฺติ สมฺพหุลา ภิกฺขุนิโย ภิกฺขุนีนํ หตฺเถ จีวรํ นิกฺขิปิตฺวา สนฺตรุตฺตเรน ชนปท\-จาริกํ ปกฺกมิํสุ, ตสฺมิํ วตฺถุสฺมิํ. เอกา ปญฺญตฺติ. ฉนฺนํ อาปตฺติ\-สมุฏฺฐานานํ ทฺวีหิ สมุฏฺฐาเนหิ สมุฏฺฐาติ, กถินเก ฯเปฯ.}

\prose{จีวรสงฺกมนียํ ธาเรนฺติยา ปาจิตฺติยํ กตฺถ ปญฺญตฺตนฺติ สาวตฺถิยํ ปญฺญตฺตํ. กํ อารพฺภาติ อญฺญตรํ ภิกฺขุนิํ อารพฺภ. กิสฺมิํ วตฺถุสฺมินฺติ อญฺญตรา ภิกฺขุนี อญฺญตราย ภิกฺขุนิยา จีวรํ อนาปุจฺฉา ปารุปิ, ตสฺมิํ วตฺถุสฺมิํ. เอกา ปญฺญตฺติ. ฉนฺนํ อาปตฺติ\-สมุฏฺฐานานํ ทฺวีหิ สมุฏฺฐาเนหิ สมุฏฺฐาติ, กถินเก ฯเปฯ.}

\prose{คณสฺส จีวรลาภํ อนฺตรายํ กโรนฺติยา ปาจิตฺติยํ กตฺถ ปญฺญตฺตนฺติ สาวตฺถิยํ ปญฺญตฺตํ. กํ อารพฺภาติ ถุลฺลนนฺทํ ภิกฺขุนิํ อารพฺภ. กิสฺมิํ วตฺถุสฺมินฺติ ถุลฺลนนฺทา ภิกฺขุนี คณสฺส จีวรลาภํ อนฺตรายํ อกาสิ, ตสฺมิํ วตฺถุสฺมิํ. เอกา ปญฺญตฺติ. ฉนฺนํ อาปตฺติ\-สมุฏฺฐานานํ ตีหิ สมุฏฺฐาเนหิ สมุฏฺฐาติ ฯเปฯ.}

\prose{ธมฺมิกํ จีวรวิภงฺคํ ปฏิพาหนฺติยา ปาจิตฺติยํ กตฺถ ปญฺญตฺตนฺติ สาวตฺถิยํ ปญฺญตฺตํ. กํ อารพฺภาติ ถุลฺลนนฺทํ ภิกฺขุนิํ อารพฺภ. กิสฺมิํ วตฺถุสฺมินฺติ ถุลฺลนนฺทา ภิกฺขุนี ธมฺมิกํ จีวรวิภงฺคํ ปฏิพาหิ, ตสฺมิํ วตฺถุสฺมิํ. เอกา ปญฺญตฺติ. ฉนฺนํ อาปตฺติ\-สมุฏฺฐานานํ ตีหิ สมุฏฺฐาเนหิ สมุฏฺฐาติ ฯเปฯ.}

\prose{อคาริกสฺส วา ปริพฺพาชกสฺส วา ปริพฺพาชิกาย วา สมณจีวรํ เทนฺติยา ปาจิตฺติยํ กตฺถ ปญฺญตฺตนฺติ สาวตฺถิยํ ปญฺญตฺตํ. กํ อารพฺภาติ ถุลฺลนนฺทํ ภิกฺขุนิํ อารพฺภ. กิสฺมิํ วตฺถุสฺมินฺติ ถุลฺลนนฺทา ภิกฺขุนี อคาริกสฺส สมณจีวรํ อทาสิ, ตสฺมิํ วตฺถุสฺมิํ. เอกา ปญฺญตฺติ. ฉนฺนํ อาปตฺติ\-สมุฏฺฐานานํ ฉหิ สมุฏฺฐาเนหิ สมุฏฺฐาติ ฯเปฯ.}

\prose{ทุพฺพล\-จีวรปจฺจาสาย จีวร\-กาล\-สมยํ อติกฺกาเมนฺติยา ปาจิตฺติยํ กตฺถ ปญฺญตฺตนฺติ สาวตฺถิยํ ปญฺญตฺตํ. กํ อารพฺภาติ ถุลฺลนนฺทํ ภิกฺขุนิํ อารพฺภ. กิสฺมิํ วตฺถุสฺมินฺติ ถุลฺลนนฺทา ภิกฺขุนี ทุพฺพล\-จีวรปจฺจาสาย จีวร\-กาล\-สมยํ อติกฺกาเมสิ, ตสฺมิํ วตฺถุสฺมิํ. เอกา ปญฺญตฺติ. ฉนฺนํ อาปตฺติ\-สมุฏฺฐานานํ ตีหิ สมุฏฺฐาเนหิ}

\prosewithclosermidruled{\csromanfolio{111}ธมฺมิกํ กถินุทฺธารํ ปฏิพาหนฺติยา ปาจิตฺติยํ กตฺถ ปญฺญตฺตนฺติ สาวตฺถิยํ ปญฺญตฺตํ. กํ อารพฺภาติ ถุลฺลนนฺทํ ภิกฺขุนิํ อารพฺภ. กิสฺมิํ วตฺถุสฺมินฺติ ถุลฺลนนฺทา ภิกฺขุนี ธมฺมิกํ กถินุทฺธารํ ปฏิพาหิ, ตสฺมิํ วตฺถุสฺมิํ. เอกา ปญฺญตฺติ. ฉนฺนํ อาปตฺติ\-สมุฏฺฐานานํ ตีหิ สมุฏฺฐาเนหิ สมุฏฺฐาติ ฯเปฯ.}{นหานวคฺโค ตติโย.}
\csromanmatikaanchor{mtk.82}
\csromantocmarkref{section}{4. ตุวฏฺฏวคฺค}{mtk.82}
\bookmark[dest={mtk.82},level=1]{4. ตุวฏฺฏวคฺค}
\csromanheader{1}{4. \textbf{ตุวฏฺฏวคฺค}}
\proseitem{221}{ทฺวินฺนํ ภิกฺขุนีนํ เอกมญฺเจ ตุวฏฺเฏนฺตีนํ ปาจิตฺติยํ กตฺถ ปญฺญตฺตนฺติ สาวตฺถิยํ ปญฺญตฺตํ. กํ อารพฺภาติ สมฺพหุลา ภิกฺขุนิโย อารพฺภ. กิสฺมิํ วตฺถุสฺมินฺติ สมฺพหุลา ภิกฺขุนิโย เทฺว เอกมญฺเจ ตุวฏฺเฏสุํ, ตสฺมิํ วตฺถุสฺมิํ. เอกา ปญฺญตฺติ. ฉนฺนํ อาปตฺติ\-สมุฏฺฐานานํ ทฺวีหิ สมุฏฺฐาเนหิ สมุฏฺฐาติ, เอฬกโลมเก ฯเปฯ.}

\prose{ทฺวินฺนํ ภิกฺขุนีนํ เอก\-ตฺถรณ\-ปาวุรเณ ตุวฏฺเฏนฺตีนํ ปาจิตฺติยํ กตฺถ ปญฺญตฺตนฺติ สาวตฺถิยํ ปญฺญตฺตํ. กํ อารพฺภาติ สมฺพหุลา ภิกฺขุนิโย อารพฺภ. กิสฺมิํ วตฺถุสฺมินฺติ สมฺพหุลา ภิกฺขุนิโย เทฺว เอก\-ตฺถรณ\-ปาวุรณา ตุวฏฺเฏสุํ, ตสฺมิํ วตฺถุสฺมิํ. เอกา ปญฺญตฺติ. ฉนฺนํ อาปตฺติ\-สมุฏฺฐานานํ ทฺวีหิ สมุฏฺฐาเนหิ สมุฏฺฐาติ, เอฬกโลมเก ฯเปฯ.}

\prose{ภิกฺขุนิยา สญฺจิจฺจ อผาสุํ กโรนฺติยา ปาจิตฺติยํ กตฺถ ปญฺญตฺตนฺติ สาวตฺถิยํ ปญฺญตฺตํ. กํ อารพฺภาติ ถุลฺลนนฺทํ ภิกฺขุนิํ อารพฺภ. กิสฺมิํ วตฺถุสฺมินฺติ ถุลฺลนนฺทา ภิกฺขุนี ภิกฺขุนิยา สญฺจิจฺจ อผาสุํ อกาสิ, ตสฺมิํ วตฺถุสฺมิํ. เอกา ปญฺญตฺติ. ฉนฺนํ อาปตฺติ\-สมุฏฺฐานานํ ตีหิ สมุฏฺฐาเนหิ สมุฏฺฐาติ ฯเปฯ.}

\prose{ทุกฺขิตํ สหชีวินิํ เนว อุปฏฺเฐนฺติยา น อุปฏฺฐาปนาย อุสฺสุกฺกํ กโรนฺติยา ปาจิตฺติยํ กตฺถ ปญฺญตฺตนฺติ สาวตฺถิยํ ปญฺญตฺตํ. กํ อารพฺภาติ ถุลฺลนนฺทํ ภิกฺขุนิํ อารพฺภ. กิสฺมิํ วตฺถุสฺมินฺติ ถุลฺลนนฺทา ภิกฺขุนี ทุกฺขิตํ สหชีวินิํ เนว อุปฏฺเฐสิ, น อุปฏฺฐาปนาย อุสฺสุกฺกํ อกาสิ, ตสฺมิํ วตฺถุสฺมิํ. เอกา ปญฺญตฺติ. ฉนฺนํ อาปตฺติ\-สมุฏฺฐานานํ เอเกน สมุฏฺฐาเนน สมุฏฺฐาติ, ธุรนิกฺเขเป ฯเปฯ.}

\prose{\csromanfolio{112}ภิกฺขุนิยา อุปสฺสยํ ทตฺวา กุปิตาย อนตฺตมนาย นิกฺกฑฺฒนฺติยา ปาจิตฺติยํ กตฺถ ปญฺญตฺตนฺติ สาวตฺถิยํ ปญฺญตฺตํ. กํ อารพฺภาติ ถุลฺลนนฺทํ ภิกฺขุนิํ อารพฺภ. กิสฺมิํ วตฺถุสฺมินฺติ ถุลฺลนนฺทา ภิกฺขุนี ภิกฺขุนิยา อุปสฺสยํ ทตฺวา กุปิตา อนตฺตมนา นิกฺกฑฺฒิ, ตสฺมิํ วตฺถุสฺมิํ. เอกา ปญฺญตฺติ. ฉนฺนํ อาปตฺติ\-สมุฏฺฐานานํ ตีหิ สมุฏฺฐาเนหิ สมุฏฺฐาติ ฯเปฯ.}

\prose{สํสฏฺฐาย ภิกฺขุนิยา ยาวตติยํ สมนุภาสนาย น ปฏินิสฺสชฺชนฺติยา ปาจิตฺติยํ กตฺถ ปญฺญตฺตนฺติ สาวตฺถิยํ ปญฺญตฺตํ. กํ อารพฺภาติ จณฺฑกาฬิํ ภิกฺขุนิํ อารพฺภ. กิสฺมิํ วตฺถุสฺมินฺติ จณฺฑกาฬี ภิกฺขุนี สํสฏฺฐา วิหริ, ตสฺมิํ วตฺถุสฺมิํ. เอกา ปญฺญตฺติ. ฉนฺนํ อาปตฺติ\-สมุฏฺฐานานํ เอเกน สมุฏฺฐาเนน สมุฏฺฐาติ, ธุรนิกฺเขเป ฯเปฯ.}

\prose{อนฺโตรฏฺเฐ สาสงฺกสมฺมเต สปฺปฏิภเย อสตฺถิกาย จาริกํ จรนฺติยา ปาจิตฺติยํ กตฺถ ปญฺญตฺตนฺติ สาวตฺถิยํ ปญฺญตฺตํ. กํ อารพฺภาติ สมฺพหุลา ภิกฺขุนิโย อารพฺภ. กิสฺมิํ วตฺถุสฺมินฺติ สมฺพหุลา ภิกฺขุนิโย อนฺโตรฏฺเฐ สาสงฺกสมฺมเต สปฺปฏิภเย อสตฺถิกาโย จาริกํ จริํสุ, ตสฺมิํ วตฺถุสฺมิํ. เอกา ปญฺญตฺติ. ฉนฺนํ อาปตฺติ\-สมุฏฺฐานานํ ทฺวีหิ สมุฏฺฐาเนหิ สมุฏฺฐาติ, เอฬกโลมเก ฯเปฯ.}

\prose{ติโรรฏฺเฐ สาสงฺกสมฺมเต สปฺปฏิภเย อสตฺถิกาย จาริกํ จรนฺติยา ปาจิตฺติยํ กตฺถ ปญฺญตฺตนฺติ สาวตฺถิยํ ปญฺญตฺตํ. กํ อารพฺภาติ สมฺพหุลา ภิกฺขุนิโย อารพฺภ. กิสฺมิํ วตฺถุสฺมินฺติ สมฺพหุลา ภิกฺขุนิโย ติโรรฏฺเฐ สาสงฺกสมฺมเต สปฺปฏิภเย อสตฺถิกาโย จาริกํ จริํสุ, ตสฺมิํ วตฺถุสฺมิํ. เอกา ปญฺญตฺติ. ฉนฺนํ อาปตฺติ\-สมุฏฺฐานานํ ทฺวีหิ สมุฏฺฐาเนหิ สมุฏฺฐาติ, เอฬกโลมเก ฯเปฯ.}

\prose{อนฺโตวสฺสํ จาริกํ จรนฺติยา ปาจิตฺติยํ กตฺถ ปญฺญตฺตนฺติ ราชคเห ปญฺญตฺตํ. กํ อารพฺภาติ สมฺพหุลา ภิกฺขุนิโย อารพฺภ. กิสฺมิํ วตฺถุสฺมินฺติ สมฺพหุลา ภิกฺขุนิโย อนฺโตวสฺสํ จาริกํ จริํสุ, ตสฺมิํ วตฺถุสฺมิํ. เอกา ปญฺญตฺติ. ฉนฺนํ อาปตฺติ\-สมุฏฺฐานานํ ทฺวีหิ สมุฏฺฐาเนหิ สมุฏฺฐาติ, เอฬกโลมเก ฯเปฯ.}

\prosewithclosermidruled{วสฺสํวุฏฺฐาย\csromansharedfootnote{bd0ef7975ee1fbe5}{วสฺสํวุตฺถาย (สี, สฺยา)} ภิกฺขุนิยา จาริกํ น ปกฺกมนฺติยา ปาจิตฺติยํ กตฺถ ปญฺญตฺตนฺติ ราชคเห ปญฺญตฺตํ. กํ อารพฺภาติ สมฺพหุลา ภิกฺขุนิโย \csromanfolio{113}อารพฺภ. กิสฺมิํ วตฺถุสฺมินฺติ สมฺพหุลา ภิกฺขุนิโย วสฺสํวุฏฺฐา จาริกํ น ปกฺกมิํสุ, ตสฺมิํ วตฺถุสฺมิํ. เอกา ปญฺญตฺติ. ฉนฺนํ อาปตฺติ\-สมุฏฺฐานานํ เอเกน สมุฏฺฐาเนน สมุฏฺฐาติ, ปฐม\-ปาราชิเก ฯเปฯ.}{ตุวฏฺฏวคฺโค จตุตฺโถ.}
\csromanmatikaanchor{mtk.83}
\csromantocmarkref{section}{5. จิตฺตาคารวคฺค}{mtk.83}
\bookmark[dest={mtk.83},level=1]{5. จิตฺตาคารวคฺค}
\csromanheader{1}{5. \textbf{จิตฺตาคารวคฺค}}
\proseitem{222}{ราชาคารํ วา จิตฺตาคารํ วา อารามํ วา อุยฺยานํ วา โปกฺขรณิํ วา ทสฺสนาย คจฺฉนฺติยา ปาจิตฺติยํ กตฺถ ปญฺญตฺตนฺติ สาวตฺถิยํ ปญฺญตฺตํ. กํ อารพฺภาติ ฉพฺพคฺคิยา ภิกฺขุนิโย อารพฺภ. กิสฺมิํ วตฺถุสฺมินฺติ ฉพฺพคฺคิยา ภิกฺขุนิโย ราชาคารมฺปิ จิตฺตาคารมฺปิ ทสฺสนาย อคมํสุ, ตสฺมิํ วตฺถุสฺมิํ. เอกา ปญฺญตฺติ. ฉนฺนํ อาปตฺติ\-สมุฏฺฐานานํ ทฺวีหิ สมุฏฺฐาเนหิ สมุฏฺฐาติ, เอฬกโลมเก ฯเปฯ.}

\prose{อาสนฺทิํ วา ปลฺลงฺกํ วา ปริภุญฺช\-นฺติยา ปาจิตฺติยํ กตฺถ ปญฺญตฺตนฺติ สาวตฺถิยํ ปญฺญตฺตํ. กํ อารพฺภาติ สมฺพหุลา ภิกฺขุนิโย อารพฺภ. กิสฺมิํ วตฺถุสฺมินฺติ สมฺพหุลา ภิกฺขุนิโย อาสนฺทิมฺปิ ปลฺลงฺกมฺปิ ปริภุญฺชิํสุ, ตสฺมิํ วตฺถุสฺมิํ. เอกา ปญฺญตฺติ. ฉนฺนํ อาปตฺติ\-สมุฏฺฐานานํ ทฺวีหิ สมุฏฺฐาเนหิ สมุฏฺฐาติ, เอฬกโลมเก ฯเปฯ.}

\prose{สุตฺตํ กนฺตนฺติยา ปาจิตฺติยํ กตฺถ ปญฺญตฺตนฺติ สาวตฺถิยํ ปญฺญตฺตํ. กํ อารพฺภาติ ฉพฺพคฺคิยา ภิกฺขุนิโย อารพฺภ. กิสฺมิํ วตฺถุสฺมินฺติ ฉพฺพคฺคิยา ภิกฺขุนิโย สุตฺตํ กนฺติํสุ, ตสฺมิํ วตฺถุสฺมิํ. เอกา ปญฺญตฺติ. ฉนฺนํ อาปตฺติ\-สมุฏฺฐานานํ ทฺวีหิ สมุฏฺฐาเนหิ สมุฏฺฐาติ, เอฬกโลมเก ฯเปฯ.}

\prose{คิหิเวยฺยาวจฺจํ กโรนฺติยา ปาจิตฺติยํ กตฺถ ปญฺญตฺตนฺติ สาวตฺถิยํ ปญฺญตฺตํ. กํ อารพฺภาติ สมฺพหุลา ภิกฺขุนิโย อารพฺภ. กิสฺมิํ วตฺถุสฺมินฺติ สมฺพหุลา ภิกฺขุนิโย คิหิเวยฺยาวจฺจํ อกํสุ, ตสฺมิํ วตฺถุสฺมิํ. เอกา ปญฺญตฺติ. ฉนฺนํ อาปตฺติ\-สมุฏฺฐานานํ ทฺวีหิ สมุฏฺฐาเนหิ สมุฏฺฐาติ, เอฬกโลมเก ฯเปฯ.}

\prose{ภิกฺขุนิยา “เอหายฺเย, อิมํ อธิกรณํ วูปสเมหี”ติ วุจฺจมานาย “สาธู”ติ ปฏิสฺสุณิตฺวา เนว วูปสเมนฺติยา น วูปสมาย อุสฺสุกฺกํ \csromanfolio{114}กโรนฺติยา ปาจิตฺติยํ กตฺถ ปญฺญตฺตนฺติ สาวตฺถิยํ ปญฺญตฺตํ. กํ อารพฺภาติ ถุลฺลนนฺทํ ภิกฺขุนิํ อารพฺภ. กิสฺมิํ วตฺถุสฺมินฺติ ถุลฺลนนฺทา ภิกฺขุนี ภิกฺขุนิยา “เอหายฺเย, อิมํ อธิกรณํ วูปสเมหี”ติ วุจฺจมานา “สาธู”ติ ปฏิสฺสุณิตฺวา เนว วูปสเมสิ, น วูปสมาย อุสฺสุกฺกํ อกาสิ, ตสฺมิํ วตฺถุสฺมิํ. เอกา ปญฺญตฺติ. ฉนฺนํ อาปตฺติ\-สมุฏฺฐานานํ เอเกน สมุฏฺฐาเนน สมุฏฺฐาติ, ธุรนิกฺเขเป ฯเปฯ.}

\prose{อคาริกสฺส วา ปริพฺพาชกสฺส วา ปริพฺพาชิกาย วา สหตฺถา ขาทนียํ วา โภชนียํ วา เทนฺติยา ปาจิตฺติยํ กตฺถ ปญฺญตฺตนฺติ สาวตฺถิยํ ปญฺญตฺตํ. กํ อารพฺภาติ ถุลฺลนนฺทํ ภิกฺขุนิํ อารพฺภ. กิสฺมิํ วตฺถุสฺมินฺติ ถุลฺลนนฺทา ภิกฺขุนี อคาริกสฺส สหตฺถา ขาทนียมฺปิ โภชนียมฺปิ อทาสิ, ตสฺมิํ วตฺถุสฺมิํ. เอกา ปญฺญตฺติ. ฉนฺนํ อาปตฺติ\-สมุฏฺฐานานํ ทฺวีหิ สมุฏฺฐาเนหิ สมุฏฺฐาติ, เอฬกโลมเก ฯเปฯ.}

\prose{อาวสถจีวรํ อนิสฺสชฺชิตฺวา ปริภุญฺช\-นฺติยา ปาจิตฺติยํ กตฺถ ปญฺญตฺตนฺติ สาวตฺถิยํ ปญฺญตฺตํ. กํ อารพฺภาติ ถุลฺลนนฺทํ ภิกฺขุนิํ อารพฺภ. กิสฺมิํ วตฺถุสฺมินฺติ ถุลฺลนนฺทา ภิกฺขุนี อาวสถจีวรํ อนิสฺสชฺชิตฺวา ปริภุญฺชิ, ตสฺมิํ วตฺถุสฺมิํ. เอกา ปญฺญตฺติ. ฉนฺนํ อาปตฺติ\-สมุฏฺฐานานํ ทฺวีหิ สมุฏฺฐาเนหิ สมุฏฺฐาติ, กถินเก ฯเปฯ.}

\prose{อาวสถํ อนิสฺสชฺชิตฺวา จาริกํ ปกฺกมนฺติยา ปาจิตฺติยํ กตฺถ ปญฺญตฺตนฺติ สาวตฺถิยํ ปญฺญตฺตํ. กํ อารพฺภาติ ถุลฺลนนฺทํ ภิกฺขุนิํ อารพฺภ. กิสฺมิํ วตฺถุสฺมินฺติ ถุลฺลนนฺทา ภิกฺขุนี อาวสถํ อนิสฺสชฺชิตฺวา จาริกํ ปกฺกามิ, ตสฺมิํ วตฺถุสฺมิํ. เอกา ปญฺญตฺติ. ฉนฺนํ อาปตฺติ\-สมุฏฺฐานานํ ทฺวีหิ สมุฏฺฐาเนหิ สมุฏฺฐาติ, กถินเก ฯเปฯ.}

\prose{ติรจฺฉาน\-วิชฺชํ ปริยาปุณนฺติยา ปาจิตฺติยํ กตฺถ ปญฺญตฺตนฺติ สาวตฺถิยํ ปญฺญตฺตํ. กํ อารพฺภาติ ฉพฺพคฺคิยา ภิกฺขุนิโย อารพฺภ. กิสฺมิํ วตฺถุสฺมินฺติ ฉพฺพคฺคิยา ภิกฺขุนิโย ติรจฺฉาน\-วิชฺชํ ปริยาปุณิํสุ, ตสฺมิํ วตฺถุสฺมิํ. เอกา ปญฺญตฺติ. ฉนฺนํ อาปตฺติ\-สมุฏฺฐานานํ ทฺวีหิ สมุฏฺฐาเนหิ สมุฏฺฐาติ, ปทโสธมฺเม ฯเปฯ.}

\prosewithclosermidruled{ติรจฺฉาน\-วิชฺชํ วาเจนฺติยา ปาจิตฺติยํ กตฺถ ปญฺญตฺตนฺติ สาวตฺถิยํ ปญฺญตฺตํ. กํ อารพฺภาติ ฉพฺพคฺคิยา ภิกฺขุนิโย อารพฺภ. กิสฺมิํ วตฺถุสฺมินฺติ ฉพฺพคฺคิยา \csromanfolio{115}ภิกฺขุนิโย ติรจฺฉาน\-วิชฺชํ วาเจสุํ, ตสฺมิํ วตฺถุสฺมิํ. เอกา ปญฺญตฺติ. ฉนฺนํ อาปตฺติ\-สมุฏฺฐานานํ ทฺวีหิ สมุฏฺฐาเนหิ สมุฏฺฐาติ, ปทโสธมฺเม ฯเปฯ.}{จิตฺตาคารวคฺโค ปญฺจโม.}
\csromanmatikaanchor{mtk.84}
\csromantocmarkref{section}{6. อารามวคฺค}{mtk.84}
\bookmark[dest={mtk.84},level=1]{6. อารามวคฺค}
\csromanheader{1}{6. \textbf{อารามวคฺค}}
\proseitem{223}{ชานํ สภิกฺขุกํ อารามํ อนาปุจฺฉา ปวิสนฺติยา ปาจิตฺติยํ กตฺถ ปญฺญตฺตนฺติ สาวตฺถิยํ ปญฺญตฺตํ. กํ อารพฺภาติ สมฺพหุลา ภิกฺขุนิโย อารพฺภ. กิสฺมิํ วตฺถุสฺมินฺติ สมฺพหุลา ภิกฺขุนิโย อารามํ อนาปุจฺฉา ปวิสิํสุ, ตสฺมิํ วตฺถุสฺมิํ. เอกา ปญฺญตฺติ, เทฺว อนุปญฺญตฺติโย. ฉนฺนํ อาปตฺติ\-สมุฏฺฐานานํ เอเกน สมุฏฺฐาเนน สมุฏฺฐาติ, ธุรนิกฺเขเป ฯเปฯ.}

\prose{ภิกฺขุํ อกฺโกสนฺติยา ปริภาสนฺติยา ปาจิตฺติยํ กตฺถ ปญฺญตฺตนฺติ เวสาลิยํ ปญฺญตฺตํ. กํ อารพฺภาติ ฉพฺพคฺคิยา ภิกฺขุนิโย อารพฺภ. กิสฺมิํ วตฺถุสฺมินฺติ ฉพฺพคฺคิยา ภิกฺขุนิโย อายสฺมนฺตํ อุปาลิํ อกฺโกสิํสุ, ตสฺมิํ วตฺถุสฺมิํ. เอกา ปญฺญตฺติ. ฉนฺนํ อาปตฺติ\-สมุฏฺฐานานํ ตีหิ สมุฏฺฐาเนหิ สมุฏฺฐาติ ฯเปฯ.}

\prose{จณฺฑีกตาย คณํ ปริภาสนฺติยา ปาจิตฺติยํ กตฺถ ปญฺญตฺตนฺติ สาวตฺถิยํ ปญฺญตฺตํ. กํ อารพฺภาติ ถุลฺลนนฺทํ ภิกฺขุนิํ อารพฺภ. กิสฺมิํ วตฺถุสฺมินฺติ ถุลฺลนนฺทา ภิกฺขุนี จณฺฑีกตาย คณํ ปริภาสิ, ตสฺมิํ วตฺถุสฺมิํ. เอกา ปญฺญตฺติ. ฉนฺนํ อาปตฺติ\-สมุฏฺฐานานํ ตีหิ สมุฏฺฐาเนหิ สมุฏฺฐาติ ฯเปฯ.}

\prose{นิมนฺติตาย วา ปวาริตาย วา ขาทนียํ วา โภชนียํ วา อญฺญตฺร ภุญฺชนฺติยา ปาจิตฺติยํ กตฺถ ปญฺญตฺตนฺติ สาวตฺถิยํ ปญฺญตฺตํ. กํ อารพฺภาติ สมฺพหุลา ภิกฺขุนิโย อารพฺภ. กิสฺมิํ วตฺถุสฺมินฺติ สมฺพหุลา ภิกฺขุนิโย ภุตฺตาวินิโย ปวาริตา อญฺญตฺร ภุญฺชิํสุ, ตสฺมิํ วตฺถุสฺมิํ. เอกา ปญฺญตฺติ. ฉนฺนํ อาปตฺติ\-สมุฏฺฐานานํ จตูหิ สมุฏฺฐาเนหิ สมุฏฺฐาติ ฯเปฯ.}

\prose{กุลํ มจฺฉรา\-ย\-นฺติยา ปาจิตฺติยํ กตฺถ ปญฺญตฺตนฺติ สาวตฺถิยํ ปญฺญตฺตํ. กํ อารพฺภาติ อญฺญตรํ ภิกฺขุนิํ อารพฺภ. กิสฺมิํ วตฺถุสฺมินฺติ อญฺญตรา ภิกฺขุนี กุลํ มจฺฉรายิ, ตสฺมิํ วตฺถุสฺมิํ. เอกา ปญฺญตฺติ. ฉนฺนํ อาปตฺติ\-สมุฏฺฐานานํ ตีหิ สมุฏฺฐาเนหิ สมุฏฺฐาติ ฯเปฯ.}

\prose{\csromanfolio{116}อภิกฺขุเก อาวาเส วสฺสํ วสนฺติยา ปาจิตฺติยํ กตฺถ ปญฺญตฺตนฺติ สาวตฺถิยํ ปญฺญตฺตํ. กํ อารพฺภาติ สมฺพหุลา ภิกฺขุนิโย อารพฺภ. กิสฺมิํ วตฺถุสฺมินฺติ สมฺพหุลา ภิกฺขุนิโย อภิกฺขุเก อาวาเส วสฺสํ วสิํสุ, ตสฺมิํ วตฺถุสฺมิํ. เอกา ปญฺญตฺติ. ฉนฺนํ อาปตฺติ\-สมุฏฺฐานานํ ทฺวีหิ สมุฏฺฐาเนหิ สมุฏฺฐาติ, เอฬกโลมเก ฯเปฯ.}

\prose{วสฺสํวุฏฺฐาย ภิกฺขุนิยา อุภโตสํเฆ ตีหิ ฐาเนหิ น ปวาเรนฺติยา ปาจิตฺติยํ กตฺถ ปญฺญตฺตนฺติ สาวตฺถิยํ ปญฺญตฺตํ. กํ อารพฺภาติ สมฺพหุลา ภิกฺขุนิโย อารพฺภ. กิสฺมิํ วตฺถุสฺมินฺติ สมฺพหุลา ภิกฺขุนิโย วสฺสํวุฏฺฐา ภิกฺขุสํฆํ\csromansharedfootnote{87eff202b1aed955}{ภิกฺขุนิสํเฆ (ก)} น ปวาเรสุํ, ตสฺมิํ วตฺถุสฺมิํ. เอกา ปญฺญตฺติ. ฉนฺนํ อาปตฺติ\-สมุฏฺฐานานํ เอเกน สมุฏฺฐาเนน สมุฏฺฐาติ, ธุรนิกฺเขเป ฯเปฯ.}

\prose{โอวาทาย วา สํวาสาย วา น คจฺฉนฺติยา ปาจิตฺติยํ กตฺถ ปญฺญตฺตนฺติ สกฺเกสุ ปญฺญตฺตํ. กํ อารพฺภาติ ฉพฺพคฺคิยา ภิกฺขุนิโย อารพฺภ. กิสฺมิํ วตฺถุสฺมินฺติ ฉพฺพคฺคิยา ภิกฺขุนิโย โอวาทํ น คจฺฉิํสุ, ตสฺมิํ วตฺถุสฺมิํ. เอกา ปญฺญตฺติ. ฉนฺนํ อาปตฺติ\-สมุฏฺฐานานํ เอเกน สมุฏฺฐาเนน สมุฏฺฐาติ, ปฐม\-ปาราชิเก ฯเปฯ.}

\prose{อุโปสถมฺปิ น ปุจฺฉนฺติยา โอวาทมฺปิ น ยาจนฺติยา ปาจิตฺติยํ กตฺถ ปญฺญตฺตนฺติ สาวตฺถิยํ ปญฺญตฺตํ. กํ อารพฺภาติ สมฺพหุลา ภิกฺขุนิโย อารพฺภ. กิสฺมิํ วตฺถุสฺมินฺติ สมฺพหุลา ภิกฺขุนิโย อุโปสถมฺปิ น ปุจฺฉิํสุ, โอวาทมฺปิ น ยาจิํสุ, ตสฺมิํ วตฺถุสฺมิํ. เอกา ปญฺญตฺติ. ฉนฺนํ อาปตฺติ\-สมุฏฺฐานานํ เอเกน สมุฏฺฐาเนน สมุฏฺฐาติ, ธุรนิกฺเขเป ฯเปฯ.}

\prosewithclosermidruled{ปสาเข ชาตํ คณฺฑํ วา รุธิตํ\csromansharedfootnote{3b6122dcdcab6486}{รุหิตํ (สี, สฺยา)} วา อนปโลเกตฺวา สํฆํ วา คณํ วา ปุริเสน สทฺธิํ เอเกเนกาย เภทาเปนฺติยา ปาจิตฺติยํ กตฺถ ปญฺญตฺตนฺติ สาวตฺถิยํ ปญฺญตฺตํ. กํ อารพฺภาติ อญฺญตรํ ภิกฺขุนิํ อารพฺภ. กิสฺมิํ วตฺถุสฺมินฺติ อญฺญตรา ภิกฺขุนี ปสาเข ชาตํ คณฺฑํ ปุริเสน สทฺธิํ เอเกเนกา เภทาเปสิ, ตสฺมิํ วตฺถุสฺมิํ. เอกา ปญฺญตฺติ. ฉนฺนํ อาปตฺติ\-สมุฏฺฐานานํ ทฺวีหิ สมุฏฺฐาเนหิ สมุฏฺฐาติ, กถินเก ฯเปฯ.}{อารามวคฺโค ฉฏฺโฐ.}
\csromanmatikaanchor{mtk.85}
\csromantocmarkref{section}{7. คพฺภินีวคฺค}{mtk.85}
\bookmark[dest={mtk.85},level=1]{7. คพฺภินีวคฺค}
\csromanheader{1}{7. \textbf{คพฺภินีวคฺค}}
\proseitem{224}{\csromanfolio{117}คพฺภินิํ วุฏฺฐาเปนฺติยา ปาจิตฺติยํ กตฺถ ปญฺญตฺตนฺติ สาวตฺถิยํ ปญฺญตฺตํ. กํ อารพฺภาติ สมฺพหุลา ภิกฺขุนิโย อารพฺภ. กิสฺมิํ วตฺถุสฺมินฺติ สมฺพหุลา ภิกฺขุนิโย คพฺภินิํ วุฏฺฐาเปสุํ, ตสฺมิํ วตฺถุสฺมิํ. เอกา ปญฺญตฺติ. ฉนฺนํ อาปตฺติ\-สมุฏฺฐานานํ ตีหิ สมุฏฺฐาเนหิ สมุฏฺฐาติ ฯเปฯ.}

\prose{ปายนฺติํ วุฏฺฐาเปนฺติยา ปาจิตฺติยํ กตฺถ ปญฺญตฺตนฺติ สาวตฺถิยํ ปญฺญตฺตํ. กํ อารพฺภาติ สมฺพหุลา ภิกฺขุนิโย อารพฺภ. กิสฺมิํ วตฺถุสฺมินฺติ สมฺพหุลา ภิกฺขุนิโย ปายนฺติํ วุฏฺฐาเปสุํ, ตสฺมิํ วตฺถุสฺมิํ. เอกา ปญฺญตฺติ. ฉนฺนํ อาปตฺติ\-สมุฏฺฐานานํ ตีหิ สมุฏฺฐาเนหิ}

\prose{เทฺว วสฺสานิ ฉสุ ธมฺเมสุ อสิกฺขิตสิกฺขํ สิกฺขมานํ วุฏฺฐาเปนฺติยา ปาจิตฺติยํ กตฺถ ปญฺญตฺตนฺติ สาวตฺถิยํ ปญฺญตฺตํ. กํ อารพฺภาติ สมฺพหุลา ภิกฺขุนิโย อารพฺภ. กิสฺมิํ วตฺถุสฺมินฺติ สมฺพหุลา ภิกฺขุนิโย เทฺว วสฺสานิ ฉสุ ธมฺเมสุ อสิกฺขิตสิกฺขํ สิกฺขมานํ วุฏฺฐาเปสุํ, ตสฺมิํ วตฺถุสฺมิํ. เอกา ปญฺญตฺติ. ฉนฺนํ อาปตฺติ\-สมุฏฺฐานานํ ตีหิ สมุฏฺฐาเนหิ สมุฏฺฐาติ ฯเปฯ.}

\prose{เทฺว วสฺสานิ ฉสุ ธมฺเมสุ สิกฺขิตสิกฺขํ สิกฺขมานํ สํเฆน อสมฺมตํ วุฏฺฐาเปนฺติยา ปาจิตฺติยํ กตฺถ ปญฺญตฺตนฺติ สาวตฺถิยํ ปญฺญตฺตํ. กํ อารพฺภาติ สมฺพหุลา ภิกฺขุนิโย อารพฺภ. กิสฺมิํ วตฺถุสฺมินฺติ สมฺพหุลา ภิกฺขุนิโย เทฺว วสฺสานิ ฉสุ ธมฺเมสุ สิกฺขิตสิกฺขํ สิกฺขมานํ สํเฆน อสมฺมตํ วุฏฺฐาเปสุํ, ตสฺมิํ วตฺถุสฺมิํ. เอกา ปญฺญตฺติ. ฉนฺนํ อาปตฺติ\-สมุฏฺฐานานํ ตีหิ สมุฏฺฐาเนหิ}

\prose{อูนทฺวาทสวสฺสํ คิหิคตํ วุฏฺฐาเปนฺติยา ปาจิตฺติยํ กตฺถ ปญฺญตฺตนฺติ สาวตฺถิยํ ปญฺญตฺตํ. กํ อารพฺภาติ สมฺพหุลา ภิกฺขุนิโย อารพฺภ. กิสฺมิํ วตฺถุสฺมินฺติ สมฺพหุลา ภิกฺขุนิโย อูนทฺวาทสวสฺสํ คิหิคตํ วุฏฺฐาเปสุํ, ตสฺมิํ วตฺถุสฺมิํ. เอกา ปญฺญตฺติ. ฉนฺนํ อาปตฺติ\-สมุฏฺฐานานํ ตีหิ สมุฏฺฐาเนหิ สมุฏฺฐาติ ฯเปฯ.}

\prose{ปริปุณฺณ\-ทฺวาทสวสฺสํ คิหิคตํ เทฺว วสฺสานิ ฉสุ ธมฺเมสุ อสิกฺขิตสิกฺขํ วุฏฺฐาเปนฺติยา ปาจิตฺติยํ กตฺถ ปญฺญตฺตนฺติ สาวตฺถิยํ ปญฺญตฺตํ. กํ อารพฺภาติ สมฺพหุลา ภิกฺขุนิโย อารพฺภ. กิสฺมิํ วตฺถุสฺมินฺติ สมฺพหุลา ภิกฺขุนิโย ปริปุณฺณ\-ทฺวาทสวสฺสํ คิหิคตํ เทฺว วสฺสานิ ฉสุ ธมฺเมสุ อสิกฺขิตสิกฺขํ \csromanfolio{118}วุฏฺฐาเปสุํ, ตสฺมิํ วตฺถุสฺมิํ. เอกา ปญฺญตฺติ. ฉนฺนํ อาปตฺติ\-สมุฏฺฐานานํ ตีหิ สมุฏฺฐาเนหิ สมุฏฺฐาติ ฯเปฯ.}

\prose{ปริปุณฺณ\-ทฺวาทสวสฺสํ คิหิคตํ เทฺว วสฺสานิ ฉสุ ธมฺเมสุ สิกฺขิตสิกฺขํ สํเฆน อสมฺมตํ วุฏฺฐาเปนฺติยา ปาจิตฺติยํ กตฺถ ปญฺญตฺตนฺติ สาวตฺถิยํ ปญฺญตฺตํ. กํ อารพฺภาติ สมฺพหุลา ภิกฺขุนิโย อารพฺภ. กิสฺมิํ วตฺถุสฺมินฺติ สมฺพหุลา ภิกฺขุนิโย ปริปุณฺณ\-ทฺวาทสวสฺสํ คิหิคตํ เทฺว วสฺสานิ ฉสุ ธมฺเมสุ สิกฺขิตสิกฺขํ สํเฆน อสมฺมตํ วุฏฺฐาเปสุํ, ตสฺมิํ วตฺถุสฺมิํ. เอกา ปญฺญตฺติ. ฉนฺนํ อาปตฺติ\-สมุฏฺฐานานํ ตีหิ สมุฏฺฐาเนหิ สมุฏฺฐาติ - ป-.}

\prose{สหชีวินิํ วุฏฺฐาเปตฺวา เทฺว วสฺสานิ เนว อนุคฺคณฺหนฺติยา น อนุคฺคณฺหาเป\-นฺติยา ปาจิตฺติยํ กตฺถ ปญฺญตฺตนฺติ สาวตฺถิยํ ปญฺญตฺตํ. กํ อารพฺภาติ ถุลฺลนนฺทํ ภิกฺขุนิํ อารพฺภ. กิสฺมิํ วตฺถุสฺมินฺติ ถุลฺลนนฺทา ภิกฺขุนี สหชีวินิํ วุฏฺฐาเปตฺวา เทฺว วสฺสานิ เนว อนุคฺคเหสิ, น อนุคฺคณฺหาเปสิ, ตสฺมิํ วตฺถุสฺมิํ. เอกา ปญฺญตฺติ. ฉนฺนํ อาปตฺติ\-สมุฏฺฐานานํ เอเกน สมุฏฺฐาเนน สมุฏฺฐาติ, ธุรนิกฺเขเป ฯเปฯ.}

\prose{วุฏฺฐาปิตํ ปวตฺตินิํ เทฺว วสฺสานิ นานุ\-พนฺธ\-นฺติยา ปาจิตฺติยํ กตฺถ ปญฺญตฺตนฺติ สาวตฺถิยํ ปญฺญตฺตํ. กํ อารพฺภาติ สมฺพหุลา ภิกฺขุนิโย อารพฺภ. กิสฺมิํ วตฺถุสฺมินฺติ สมฺพหุลา ภิกฺขุนิโย วุฏฺฐาปิตํ ปวตฺตินิํ เทฺว วสฺสานิ นานุพนฺธิํสุ, ตสฺมิํ วตฺถุสฺมิํ. เอกา ปญฺญตฺติ. ฉนฺนํ อาปตฺติ\-สมุฏฺฐานานํ เอเกน สมุฏฺฐาเนน สมุฏฺฐาติ, ปฐม\-ปาราชิเก ฯเปฯ.}

\prosewithclosermidruled{สหชีวินิํ วุฏฺฐาเปตฺวา เนว วูปกาเสนฺติยา น วูปกาสาเป\-นฺติยา ปาจิตฺติยํ กตฺถ ปญฺญตฺตนฺติ สาวตฺถิยํ ปญฺญตฺตํ. กํ อารพฺภาติ ถุลฺลนนฺทํ ภิกฺขุนิํ อารพฺภ. กิสฺมิํ วตฺถุสฺมินฺติ ถุลฺลนนฺทา ภิกฺขุนี สหชีวินิํ วุฏฺฐาเปตฺวา เนว วูปกาเสสิ, น วูปกาสาเปสิ, ตสฺมิํ วตฺถุสฺมิํ. เอกา ปญฺญตฺติ. ฉนฺนํ อาปตฺติ\-สมุฏฺฐานานํ เอเกน สมุฏฺฐาเนน สมุฏฺฐาติ, ธุรนิกฺเขเป ฯเปฯ.}{คพฺภินิวคฺโค สตฺตโม.}
\csromanmatikaanchor{mtk.86}
\csromantocmarkref{section}{8. กุมารีภูตวคฺค}{mtk.86}
\bookmark[dest={mtk.86},level=1]{8. กุมารีภูตวคฺค}
\csromanheader{1}{8. กุมารีภูต\-วคฺค}
\proseitem{225}{อูนวีสติวสฺสํ กุมาริภูตํ วุฏฺฐาเปนฺติยา ปาจิตฺติยํ กตฺถ ปญฺญตฺตนฺติ สาวตฺถิยํ ปญฺญตฺตํ. กํ อารพฺภาติ สมฺพหุลา ภิกฺขุนิโย \csromanfolio{119}อารพฺภ. กิสฺมิํ วตฺถุสฺมินฺติ สมฺพหุลา ภิกฺขุนิโย อูนวีสติวสฺสํ กุมาริภูตํ วุฏฺฐาเปสุํ, ตสฺมิํ วตฺถุสฺมิํ. เอกา ปญฺญตฺติ. ฉนฺนํ อาปตฺติ\-สมุฏฺฐานานํ ตีหิ สมุฏฺฐาเนหิ สมุฏฺฐาติ ฯเปฯ.}

\prose{ปริปุณฺณ\-วีสติวสฺสํ กุมาริภูตํ เทฺว วสฺสานิ ฉสุ ธมฺเมสุ อสิกฺขิตสิกฺขํ วุฏฺฐาเปนฺติยา ปาจิตฺติยํ กตฺถ ปญฺญตฺตนฺติ สาวตฺถิยํ ปญฺญตฺตํ. กํ อารพฺภาติ สมฺพหุลา ภิกฺขุนิโย อารพฺภ. กิสฺมิํ วตฺถุสฺมินฺติ สมฺพหุลา ภิกฺขุนิโย ปริปุณฺณ\-วีสติวสฺสํ กุมาริภูตํ เทฺว วสฺสานิ ฉสุ ธมฺเมสุ อสิกฺขิตสิกฺขํ วุฏฺฐาเปสุํ, ตสฺมิํ วตฺถุสฺมิํ. เอกา ปญฺญตฺติ. ฉนฺนํ อาปตฺติ\-สมุฏฺฐานานํ ตีหิ สมุฏฺฐาเนหิ}

\prose{ปริปุณฺณ\-วีสติวสฺสํ กุมาริภูตํ เทฺว วสฺสานิ ฉสุ ธมฺเมสุ สิกฺขิตสิกฺขํ สํเฆน อสมฺมตํ วุฏฺฐาเปนฺติยา ปาจิตฺติยํ กตฺถ ปญฺญตฺตนฺติ สาวตฺถิยํ ปญฺญตฺตํ. กํ อารพฺภาติ สมฺพหุลา ภิกฺขุนิโย อารพฺภ. กิสฺมิํ วตฺถุสฺมินฺติ สมฺพหุลา ภิกฺขุนิโย ปริปุณฺณ\-วีสติวสฺสํ กุมาริภูตํ เทฺว วสฺสานิ ฉสุ ธมฺเมสุ สิกฺขิตสิกฺขํ สํเฆน อสมฺมตํ วุฏฺฐาเปสุํ, ตสฺมิํ วตฺถุสฺมิํ. เอกา ปญฺญตฺติ. ฉนฺนํ อาปตฺติ\-สมุฏฺฐานานํ ตีหิ สมุฏฺฐาเนหิ สมุฏฺฐาติ - ป-.}

\prose{อูนทฺวาทสวสฺสาย วุฏฺฐาเปนฺติยา ปาจิตฺติยํ กตฺถ ปญฺญตฺตนฺติ สาวตฺถิยํ ปญฺญตฺตํ. กํ อารพฺภาติ สมฺพหุลา ภิกฺขุนิโย อารพฺภ. กิสฺมิํ วตฺถุสฺมินฺติ สมฺพหุลา ภิกฺขุนิโย อูนทฺวาทสวสฺสา วุฏฺฐาเปสุํ, ตสฺมิํ วตฺถุสฺมิํ. เอกา ปญฺญตฺติ. ฉนฺนํ อาปตฺติ\-สมุฏฺฐานานํ ตีหิ สมุฏฺฐาเนหิ สมุฏฺฐาติ ฯเปฯ.}

\prose{ปริปุณฺณ\-ทฺวาทสวสฺสาย สํเฆน อสมฺมตาย วุฏฺฐาเปนฺติยา ปาจิตฺติยํ กตฺถ ปญฺญตฺตนฺติ สาวตฺถิยํ ปญฺญตฺตํ. กํ อารพฺภาติ สมฺพหุลา ภิกฺขุนิโย อารพฺภ. กิสฺมิํ วตฺถุสฺมินฺติ สมฺพหุลา ภิกฺขุนิโย ปริปุณฺณ\-ทฺวาทสวสฺสา สํเฆน อสมฺมตา วุฏฺฐาเปสุํ, ตสฺมิํ วตฺถุสฺมิํ. เอกา ปญฺญตฺติ. ฉนฺนํ อาปตฺติ\-สมุฏฺฐานานํ ตีหิ สมุฏฺฐาเนหิ สมุฏฺฐาติ, ทุติยปาราชิเก ฯเปฯ.}

\prose{“อลํ ตาว เต อยฺเย วุฏฺฐาปิเตนา”ติ วุจฺจมานาย “สาธู”ติ ปฏิสฺสุณิตฺวา ปจฺฉา ขียนธมฺมํ อาปชฺชนฺติยา ปาจิตฺติยํ กตฺถ ปญฺญตฺตนฺติ สาวตฺถิยํ ปญฺญตฺตํ. กํ อารพฺภาติ จณฺฑกาฬิํ ภิกฺขุนิํ อารพฺภ. กิสฺมิํ วตฺถุสฺมินฺติ จณฺฑกาฬี ภิกฺขุนี “อลํ ตาว เต อยฺเย วุฏฺฐาปิเตนา”ติ \csromanfolio{120}วุจฺจมานา “สาธู”ติ ปฏิสฺสุณิตฺวา ปจฺฉา ขียนธมฺมํ อาปชฺชิ, ตสฺมิํ วตฺถุสฺมิํ. เอกา ปญฺญตฺติ. ฉนฺนํ อาปตฺติ\-สมุฏฺฐานานํ ตีหิ สมุฏฺฐาเนหิ}

\prose{สิกฺขมานํ “สเจ เม ตฺวํ อยฺเย จีวรํ ทสฺสสิ, เอวาหํ ตํ วุฏฺฐาเปสฺสามี”ติ วตฺวา เนว วุฏฺฐาเปนฺติยา น วุฏฺฐาปนาย อุสฺสุกฺกํ กโรนฺติยา ปาจิตฺติยํ กตฺถ ปญฺญตฺตนฺติ สาวตฺถิยํ ปญฺญตฺตํ. กํ อารพฺภาติ ถุลฺลนนฺทํ ภิกฺขุนิํ อารพฺภ. กิสฺมิํ วตฺถุสฺมินฺติ ถุลฺลนนฺทา ภิกฺขุนี สิกฺขมานํ “สเจ เม ตฺวํ อยฺเย จีวรํ ทสฺสสิ, เอวาหํ ตํ วุฏฺฐาเปสฺสามี”ติ วตฺวา เนว วุฏฺฐาเปสิ, น วุฏฺฐาปนาย อุสฺสุกฺกํ อกาสิ, ตสฺมิํ วตฺถุสฺมิํ. เอกา ปญฺญตฺติ. ฉนฺนํ อาปตฺติ\-สมุฏฺฐานานํ เอเกน สมุฏฺฐาเนน สมุฏฺฐาติ, ธุรนิกฺเขเป ฯเปฯ.}

\prose{สิกฺขมานํ “สเจ มํ ตฺวํ อยฺเย เทฺว วสฺสานิ อนุพนฺธิสฺสสิ, เอวาหํ ตํ วุฏฺฐาเปสฺสามี”ติ วตฺวา เนว วุฏฺฐาเปนฺติยา น วุฏฺฐาปนาย อุสฺสุกฺกํ กโรนฺติยา ปาจิตฺติยํ กตฺถ ปญฺญตฺตนฺติ สาวตฺถิยํ ปญฺญตฺตํ. กํ อารพฺภาติ ถุลฺลนนฺทํ ภิกฺขุนิํ อารพฺภ. กิสฺมิํ วตฺถุสฺมินฺติ ถุลฺลนนฺทา ภิกฺขุนี สิกฺขมานํ “สเจ มํ ตฺวํ อยฺเย เทฺว วสฺสานิ อนุพนฺธิสฺสสิ, เอวาหํ ตํ วุฏฺฐาเปสฺสามี”ติ วตฺวา เนว วุฏฺฐาเปสิ, น วุฏฺฐาปนาย อุสฺสุกฺกํ อกาสิ, ตสฺมิํ วตฺถุสฺมิํ. เอกา ปญฺญตฺติ. ฉนฺนํ อาปตฺติ\-สมุฏฺฐานานํ เอเกน สมุฏฺฐาเนน สมุฏฺฐาติ. ธุรนิกฺเขเป ฯเปฯ.}

\prose{ปุริส\-สํสฏฺฐํ กุมารก\-สํสฏฺฐํ จณฺฑิํ โสกาวาสํ สิกฺขมานํ วุฏฺฐาเปนฺติยา ปาจิตฺติยํ กตฺถ ปญฺญตฺตนฺติ สาวตฺถิยํ ปญฺญตฺตํ. กํ อารพฺภาติ ถุลฺลนนฺทํ ภิกฺขุนิํ อารพฺภ. กิสฺมิํ วตฺถุสฺมินฺติ ถุลฺลนนฺทา ภิกฺขุนี ปุริส\-สํสฏฺฐํ กุมารก\-สํสฏฺฐํ จณฺฑิํ โสกาวาสํ สิกฺขมานํ วุฏฺฐาเปสิ, ตสฺมิํ วตฺถุสฺมิํ. เอกา ปญฺญตฺติ. ฉนฺนํ อาปตฺติ\-สมุฏฺฐานานํ ตีหิ สมุฏฺฐาเนหิ สมุฏฺฐาติ ฯเปฯ.}

\prose{มาตาปิตูหิ วา สามิเกน วา อนนุญฺญาตํ สิกฺขมานํ วุฏฺฐาเปนฺติยา ปาจิตฺติยํ กตฺถ ปญฺญตฺตนฺติ สาวตฺถิยํ ปญฺญตฺตํ. กํ อารพฺภาติ ถุลฺลนนฺทํ ภิกฺขุนิํ อารพฺภ. กิสฺมิํ วตฺถุสฺมินฺติ ถุลฺลนนฺทา ภิกฺขุนี มาตาปิตูหิปิ สามิเกนาปิ อนนุญฺญาตํ สิกฺขมานํ วุฏฺฐาเปสิ, ตสฺมิํ วตฺถุสฺมิํ. เอกา ปญฺญตฺติ. ฉนฺนํ อาปตฺติ\-สมุฏฺฐานานํ จตูหิ สมุฏฺฐาเนหิ สมุฏฺฐาติ. สิยา วาจโต สมุฏฺฐาติ, น กายโต, น จิตฺตโต, สิยา กายโต จ \csromanfolio{121}วาจโต จ สมุฏฺฐาติ, น จิตฺตโต, สิยา วาจโต จ จิตฺตโต จ สมุฏฺฐาติ, น กายโต, สิยา กายโต จ วาจโต จ จิตฺตโต จ สมุฏฺฐาติ ฯเปฯ.}

\prose{ปาริวาสิก\-ฉนฺททาเนน สิกฺขมานํ วุฏฺฐาเปนฺติยา ปาจิตฺติยํ กตฺถ ปญฺญตฺตนฺติ ราชคเห ปญฺญตฺตํ. กํ อารพฺภาติ ถุลฺลนนฺทํ ภิกฺขุนิํ อารพฺภ. กิสฺมิํ วตฺถุสฺมินฺติ ถุลฺลนนฺทา ภิกฺขุนี ปาริวาสิก\-ฉนฺททาเนน สิกฺขมานํ วุฏฺฐาเปสิ, ตสฺมิํ วตฺถุสฺมิํ. เอกา ปญฺญตฺติ. ฉนฺนํ อาปตฺติ\-สมุฏฺฐานานํ ตีหิ สมุฏฺฐาเนหิ สมุฏฺฐาติ ฯเปฯ.}

\prose{อนุวสฺสํ วุฏฺฐาเปนฺติยา ปาจิตฺติยํ กตฺถ ปญฺญตฺตนฺติ สาวตฺถิยํ ปญฺญตฺตํ. กํ อารพฺภาติ สมฺพหุลา ภิกฺขุนิโย อารพฺภ. กิสฺมิํ วตฺถุสฺมินฺติ สมฺพหุลา ภิกฺขุนิโย อนุวสฺสํ วุฏฺฐาเปสุํ, ตสฺมิํ วตฺถุสฺมิํ. เอกา ปญฺญตฺติ. ฉนฺนํ อาปตฺติ\-สมุฏฺฐานานํ ตีหิ สมุฏฺฐาเนหิ}

\prosewithclosermidruled{เอกํ วสฺสํ เทฺว วุฏฺฐาเปนฺติยา ปาจิตฺติยํ กตฺถ ปญฺญตฺตนฺติ สาวตฺถิยํ ปญฺญตฺตํ. กํ อารพฺภาติ สมฺพหุลา ภิกฺขุนิโย อารพฺภ. กิสฺมิํ วตฺถุสฺมินฺติ สมฺพหุลา ภิกฺขุนิโย เอกํ วสฺสํ เทฺว วุฏฺฐาเปสุํ, ตสฺมิํ วตฺถุสฺมิํ. เอกา ปญฺญตฺติ. ฉนฺนํ อาปตฺติ\-สมุฏฺฐานานํ ตีหิ สมุฏฺฐาเนหิ สมุฏฺฐาติ.}{กุมาริภูต\-วคฺโค อฏฺฐโม.}
\csromanmatikaanchor{mtk.87}
\csromantocmarkref{section}{9. ฉตฺตุปาหนวคฺค}{mtk.87}
\bookmark[dest={mtk.87},level=1]{9. ฉตฺตุปาหนวคฺค}
\csromanheader{1}{9. ฉตฺตุปาหน\-วคฺค}
\proseitem{226}{ฉตฺตุปาหนํ ธาเรนฺติยา ปาจิตฺติยํ กตฺถ ปญฺญตฺตนฺติ สาวตฺถิยํ ปญฺญตฺตํ. กํ อารพฺภาติ ฉพฺพคฺคิยา ภิกฺขุนิโย อารพฺภ. กิสฺมิํ วตฺถุสฺมินฺติ ฉพฺพคฺคิยา ภิกฺขุนิโย ฉตฺตุปาหนํ ธาเรสุํ, ตสฺมิํ วตฺถุสฺมิํ. เอกา ปญฺญตฺติ, เอกา อนุปญฺญตฺติ. ฉนฺนํ อาปตฺติ\-สมุฏฺฐานานํ ทฺวีหิ สมุฏฺฐาเนหิ สมุฏฺฐาติ, เอฬกโลมเก ฯเปฯ.}

\prose{ยาเนน ยายนฺติยา ปาจิตฺติยํ กตฺถ ปญฺญตฺตนฺติ สาวตฺถิยํ ปญฺญตฺตํ. กํ อารพฺภาติ ฉพฺพคฺคิยา ภิกฺขุนิโย อารพฺภ. กิสฺมิํ วตฺถุสฺมินฺติ ฉพฺพคฺคิยา ภิกฺขุนิโย ยาเนน ยายิํสุ, ตสฺมิํ วตฺถุสฺมิํ. เอกา ปญฺญตฺติ, เอกา อนุปญฺญตฺติ. ฉนฺนํ อาปตฺติ\-สมุฏฺฐานานํ ทฺวีหิ สมุฏฺฐาเนหิ สมุฏฺฐาติ, เอฬกโลมเก ฯเปฯ.}

\prose{\csromanfolio{122}สงฺฆาณิํ ธาเรนฺติยา ปาจิตฺติยํ กตฺถ ปญฺญตฺตนฺติ สาวตฺถิยํ ปญฺญตฺตํ. กํ อารพฺภาติ อญฺญตรํ ภิกฺขุนิํ อารพฺภ. กิสฺมิํ วตฺถุสฺมินฺติ อญฺญตรา ภิกฺขุนี สงฺฆาณิํ ธาเรสิ, ตสฺมิํ วตฺถุสฺมิํ. เอกา ปญฺญตฺติ. ฉนฺนํ อาปตฺติ\-สมุฏฺฐานานํ ทฺวีหิ สมุฏฺฐาเนหิ สมุฏฺฐาติ, เอฬกโลมเก ฯเปฯ.}

\prose{อิตฺถาลงฺการํ ธาเรนฺติยา ปาจิตฺติยํ กตฺถ ปญฺญตฺตนฺติ สาวตฺถิยํ ปญฺญตฺตํ. กํ อารพฺภาติ ฉพฺพคฺคิยา ภิกฺขุนิโย อารพฺภ. กิสฺมิํ วตฺถุสฺมินฺติ ฉพฺพคฺคิยา ภิกฺขุนิโย อิตฺถาลงฺการํ ธาเรสุํ, ตสฺมิํ วตฺถุสฺมิํ. เอกา ปญฺญตฺติ. ฉนฺนํ อาปตฺติ\-สมุฏฺฐานานํ ทฺวีหิ สมุฏฺฐาเนหิ สมุฏฺฐาติ, เอฬกโลมเก ฯเปฯ.}

\prose{คนฺธ\-วณฺณเกน นหายนฺติยา ปาจิตฺติยํ กตฺถ ปญฺญตฺตนฺติ สาวตฺถิยํ ปญฺญตฺตํ. กํ อารพฺภาติ ฉพฺพคฺคิยา ภิกฺขุนิโย อารพฺภ. กิสฺมิํ วตฺถุสฺมินฺติ ฉพฺพคฺคิยา ภิกฺขุนิโย คนฺธ\-วณฺณเกน นหายิํสุ, ตสฺมิํ วตฺถุสฺมิํ. เอกา ปญฺญตฺติ. ฉนฺนํ อาปตฺติ\-สมุฏฺฐานานํ ทฺวีหิ สมุฏฺฐาเนหิ สมุฏฺฐาติ, เอฬกโลมเก ฯเปฯ.}

\prose{วาสิตเกน ปิญฺญาเกน นหายนฺติยา ปาจิตฺติยํ กตฺถ ปญฺญตฺตนฺติ สาวตฺถิยํ ปญฺญตฺตํ. กํ อารพฺภาติ ฉพฺพคฺคิยา ภิกฺขุนิโย อารพฺภ. กิสฺมิํ วตฺถุสฺมินฺติ ฉพฺพคฺคิยา ภิกฺขุนิโย วาสิตเกน ปิญฺญาเกน นหายิํสุ, ตสฺมิํ วตฺถุสฺมิํ. เอกา ปญฺญตฺติ. ฉนฺนํ อาปตฺติ\-สมุฏฺฐานานํ ทฺวีหิ สมุฏฺฐาเนหิ สมุฏฺฐาติ, เอฬกโลมเก ฯเปฯ.}

\prose{ภิกฺขุนิยา อุมฺมทฺทาเปนฺติยา ปริมทฺทาเป\-นฺติยา ปาจิตฺติยํ กตฺถ ปญฺญตฺตนฺติ สาวตฺถิยํ ปญฺญตฺตํ. กํ อารพฺภาติ สมฺพหุลา ภิกฺขุนิโย อารพฺภ. กิสฺมิํ วตฺถุสฺมินฺติ สมฺพหุลา ภิกฺขุนิโย ภิกฺขุนิยา อุมฺมทฺทาเปสุํ ปริมทฺทาเปสุํ, ตสฺมิํ วตฺถุสฺมิํ. เอกา ปญฺญตฺติ. ฉนฺนํ อาปตฺติ\-สมุฏฺฐานานํ ทฺวีหิ สมุฏฺฐาเนหิ สมุฏฺฐาติ, เอฬกโลมเก ฯเปฯ.}

\prose{สิกฺขมานาย อุมฺมทฺทาเปนฺติยา ปริมทฺทาเป\-นฺติยา ปาจิตฺติยํ กตฺถ ปญฺญตฺตนฺติ สาวตฺถิยํ ปญฺญตฺตํ. กํ อารพฺภาติ สมฺพหุลา ภิกฺขุนิโย อารพฺภ. กิสฺมิํ วตฺถุสฺมินฺติ สมฺพหุลา ภิกฺขุนิโย สิกฺขมานาย อุมฺมทฺทาเปสุํ ปริมทฺทาเปสุํ, ตสฺมิํ วตฺถุสฺมิํ. เอกา ปญฺญตฺติ. ฉนฺนํ อาปตฺติ\-สมุฏฺฐานานํ ทฺวีหิ สมุฏฺฐาเนหิ สมุฏฺฐาติ, เอฬกโลมเก ฯเปฯ.}

\prose{สามเณริยา อุมฺมทฺทาเปนฺติยา ปริมทฺทาเป\-นฺติยา ปาจิตฺติยํ กตฺถ ปญฺญตฺตนฺติ สาวตฺถิยํ ปญฺญตฺตํ. กํ อารพฺภาติ สมฺพหุลา ภิกฺขุนิโย \csromanfolio{123}อารพฺภ. กิสฺมิํ วตฺถุสฺมินฺติ สมฺพหุลา ภิกฺขุนิโย สามเณริยา อุมฺมทฺทาเปสุํ ปริมทฺทาเปสุํ, ตสฺมิํ วตฺถุสฺมิํ. เอกา ปญฺญตฺติ. ฉนฺนํ อาปตฺติ\-สมุฏฺฐานานํ ทฺวีหิ สมุฏฺฐาเนหิ สมุฏฺฐาติ, เอฬกโลมเก ฯเปฯ.}

\prose{คิหินิยา อุมฺมทฺทาเปนฺติยา ปริมทฺทาเป\-นฺติยา ปาจิตฺติยํ กตฺถ ปญฺญตฺตนฺติ สาวตฺถิยํ ปญฺญตฺตํ. กํ อารพฺภาติ สมฺพหุลา ภิกฺขุนิโย อารพฺภ. กิสฺมิํ วตฺถุสฺมินฺติ สมฺพหุลา ภิกฺขุนิโย คิหินิยา อุมฺมทฺทาเปสุํ ปริมทฺทาเปสุํ, ตสฺมิํ วตฺถุสฺมิํ. เอกา ปญฺญตฺติ. ฉนฺนํ อาปตฺติ\-สมุฏฺฐานานํ ทฺวีหิ สมุฏฺฐาเนหิ สมุฏฺฐาติ, เอฬกโลมเก ฯเปฯ.}

\prose{ภิกฺขุสฺส ปุรโต อนาปุจฺฉา อาสเน นิสีทนฺติยา ปาจิตฺติยํ กตฺถ ปญฺญตฺตนฺติ สาวตฺถิยํ ปญฺญตฺตํ. กํ อารพฺภาติ สมฺพหุลา ภิกฺขุนิโย อารพฺภ. กิสฺมิํ วตฺถุสฺมินฺติ สมฺพหุลา ภิกฺขุนิโย ภิกฺขุสฺส ปุรโต อนาปุจฺฉา อาสเน นิสีทิํสุ, ตสฺมิํ วตฺถุสฺมิํ. เอกา ปญฺญตฺติ. ฉนฺนํ อาปตฺติ\-สมุฏฺฐานานํ ทฺวีหิ สมุฏฺฐาเนหิ สมุฏฺฐาติ, กถินเก ฯเปฯ.}

\prose{อโนกาสกตํ ภิกฺขุํ ปญฺหํ ปุจฺฉนฺติยา ปาจิตฺติยํ กตฺถ ปญฺญตฺตนฺติ สาวตฺถิยํ ปญฺญตฺตํ. กํ อารพฺภาติ สมฺพหุลา ภิกฺขุนิโย อารพฺภ. กิสฺมิํ วตฺถุสฺมินฺติ สมฺพหุลา ภิกฺขุนิโย อโนกาสกตํ ภิกฺขุํ ปญฺหํ ปุจฺฉิํสุ, ตสฺมิํ วตฺถุสฺมิํ. เอกา ปญฺญตฺติ. ฉนฺนํ อาปตฺติ\-สมุฏฺฐานานํ ทฺวีหิ สมุฏฺฐาเนหิ สมุฏฺฐาติ, ปทโสธมฺเม ฯเปฯ.}

\prosewithclosermidruled{อสํกจฺจิกาย\csromansharedfootnote{fd8e43a326d38e86}{อสงฺกจฺฉิกาย (สฺยา)} คามํ ปวิสนฺติยา ปาจิตฺติยํ กตฺถ ปญฺญตฺตนฺติ สาวตฺถิยํ ปญฺญตฺตํ. กํ อารพฺภาติ อญฺญตรํ ภิกฺขุนิํ อารพฺภ. กิสฺมิํ วตฺถุสฺมินฺติ อญฺญตรา ภิกฺขุนี อสํกจฺจิกา คามํ ปาวิสิ, ตสฺมิํ วตฺถุสฺมิํ. เอกา ปญฺญตฺติ. ฉนฺนํ อาปตฺติ\-สมุฏฺฐานานํ ทฺวีหิ สมุฏฺฐาเนหิ สมุฏฺฐาติ, สิยา กายโต สมุฏฺฐาติ, น วาจโต, น จิตฺตโต, สิยา กายโต จ จิตฺตโต จ สมุฏฺฐาติ, น วาจโต ฯเปฯ.}{ฉตฺตุปาหน\-วคฺโค นวโม.}
\nitthitam{นววคฺค\-ขุทฺทกา นิฏฺฐิตา.}
\csromancenter{ตสฺสุทฺทานํ.}
\setlength{\csromangathapagewidth}{0pt}
\setlength{\csromangathawakwidth}{0pt}
\csromangathameasuregroup{ลสุณํ สํหเร โลมํ, \\ ภุญฺชนฺตามกธญฺญานํ, \\ อนฺธกาเร ปฏิจฺฉนฺเน, \\ ปุเร ปจฺฉา วิกาเล จ, \\ นคฺโคทกา วิสิพฺเพตฺวา, \\ คณํ วิภงฺคสมณํ, \\ เอกมญฺจตฺถรเณน, \\ ทตฺวา สํสฏฺฐอนฺโต จ, \\ ราชา อาสนฺทิ สุตฺตญฺจ, \\ ทเท จีวราวสถํ, \\ อารามกฺโกสจณฺฑี จ, \\ วาเส ปวารโณวาทํ, \\ คพฺภี ปายนฺตี ฉ ธมฺเม, \\ ปริปุณฺณญฺจ สํเฆน, \\ กุมารี เทฺว จ สํเฆน, \\ อลํ สเจ จ เทฺววสฺสํ, \\ ปาริวาสิกานุวสฺสํ, \\ ฉตฺตยาเนน สงฺฆาณิ, \\ ปิญฺญากภิกฺขุนี เจว, \\ คิหิ ภิกฺขุสฺส ปุรโต,}{\csromangathabat{ลสุณํ สํหเร โลมํ,}{ตลมฏฺฐญฺจ สุทฺธิกํ.} \\ \csromangathabat{ภุญฺชนฺตามกธญฺญานํ,}{เทฺว วิฆาเสน ทสฺสนา.} \\ \csromangathabat{อนฺธกาเร ปฏิจฺฉนฺเน,}{อชฺโฌกาเส รถิกาย จ.} \\ \csromangathabat{ปุเร ปจฺฉา วิกาเล จ,}{ทุคฺคหิ นิรเย วธิ.} \\ \csromangathabat{นคฺโคทกา วิสิพฺเพตฺวา,}{ปญฺจาหิกํ สงฺกมนียํ.} \\ \csromangathabat{คณํ วิภงฺคสมณํ,}{ทุพฺพลํ กถิเนน จ.} \\ \csromangathabat{เอกมญฺจตฺถรเณน,}{สญฺจิจฺจ สหชีวินี.} \\ \csromangathabat{ทตฺวา สํสฏฺฐอนฺโต จ,}{ติโรวสฺสํ น ปกฺกเม.} \\ \csromangathabat{ราชา อาสนฺทิ สุตฺตญฺจ,}{คิหิ วูปสเมน จ.} \\ \csromangathabat{ทเท จีวราวสถํ,}{ปริยาปุณญฺจ วาจเย.} \\ \csromangathabat{อารามกฺโกสจณฺฑี จ,}{ภุญฺเชยฺย กุลมจฺฉรี.} \\ \csromangathabat{วาเส ปวารโณวาทํ,}{เทฺว ธมฺมา ปสาเขน จ.} \\ \csromangathabat{คพฺภี ปายนฺตี ฉ ธมฺเม,}{อสมฺมตูนทฺวาทส.} \\ \csromangathabat{ปริปุณฺณญฺจ สํเฆน,}{สห วุฏฺฐา ฉ ปญฺจ จ.} \\ \csromangathabat{กุมารี เทฺว จ สํเฆน,}{ทฺวาทส สมฺมเตน จ.} \\ \csromangathabat{อลํ สเจ จ เทฺววสฺสํ,}{สํสฏฺฐา สามิเกน จ.} \\ \csromangathabat{ปาริวาสิกานุวสฺสํ,}{ทุเว วุฏฺฐาปเนน จ.} \\ \csromangathabat{ฉตฺตยาเนน สงฺฆาณิ,}{อิตฺถาลงฺการวณฺณเก จ.} \\ \csromangathabat{ปิญฺญากภิกฺขุนี เจว,}{สิกฺขา จ สามเณริกา.} \\ \csromangathabat{คิหิ ภิกฺขุสฺส ปุรโต,}{อโนกาสํ สํกจฺจิกาติ.}}
\csromangathasetleft{ลสุณํ สํหเร โลมํ, \\ ภุญฺชนฺตามกธญฺญานํ, \\ อนฺธกาเร ปฏิจฺฉนฺเน, \\ ปุเร ปจฺฉา วิกาเล จ, \\ นคฺโคทกา วิสิพฺเพตฺวา, \\ คณํ วิภงฺคสมณํ, \\ เอกมญฺจตฺถรเณน, \\ ทตฺวา สํสฏฺฐอนฺโต จ, \\ ราชา อาสนฺทิ สุตฺตญฺจ, \\ ทเท จีวราวสถํ, \\ อารามกฺโกสจณฺฑี จ, \\ วาเส ปวารโณวาทํ, \\ คพฺภี ปายนฺตี ฉ ธมฺเม, \\ ปริปุณฺณญฺจ สํเฆน, \\ กุมารี เทฺว จ สํเฆน, \\ อลํ สเจ จ เทฺววสฺสํ, \\ ปาริวาสิกานุวสฺสํ, \\ ฉตฺตยาเนน สงฺฆาณิ, \\ ปิญฺญากภิกฺขุนี เจว, \\ คิหิ ภิกฺขุสฺส ปุรโต,}
\csromangathagroup{\csromangathastanza{\csromanfolio{124}\csromangathabat{ลสุณํ สํหเร โลมํ,}{ตลมฏฺฐญฺจ สุทฺธิกํ.} \\ \csromangathabat{ภุญฺชนฺตามกธญฺญานํ,}{เทฺว วิฆาเสน ทสฺสนา.}}\csromangathastanza{\csromangathabat{อนฺธกาเร ปฏิจฺฉนฺเน,}{อชฺโฌกาเส รถิกาย จ.} \\ \csromangathabat{ปุเร ปจฺฉา วิกาเล จ,}{ทุคฺคหิ นิรเย วธิ.}}\csromangathastanza{\csromangathabat{นคฺโคทกา วิสิพฺเพตฺวา,}{ปญฺจาหิกํ สงฺกมนียํ.} \\ \csromangathabat{คณํ วิภงฺคสมณํ,}{ทุพฺพลํ กถิเนน จ.}}\csromangathastanza{\csromangathabat{เอกมญฺจตฺถรเณน,}{สญฺจิจฺจ สหชีวินี.} \\ \csromangathabat{ทตฺวา สํสฏฺฐอนฺโต จ,}{ติโรวสฺสํ น ปกฺกเม.}}\csromangathastanza{\csromangathabat{ราชา อาสนฺทิ สุตฺตญฺจ,}{คิหิ วูปสเมน จ.} \\ \csromangathabat{ทเท จีวราวสถํ,}{ปริยาปุณญฺจ วาจเย.}}\csromangathastanza{\csromangathabat{อารามกฺโกสจณฺฑี จ,}{ภุญฺเชยฺย กุลมจฺฉรี.} \\ \csromangathabat{วาเส ปวารโณวาทํ,}{เทฺว ธมฺมา ปสาเขน จ.}}\csromangathastanza{\csromangathabat{คพฺภี ปายนฺตี ฉ ธมฺเม,}{อสมฺมตูนทฺวาทส.} \\ \csromangathabat{ปริปุณฺณญฺจ สํเฆน,}{สห วุฏฺฐา ฉ ปญฺจ จ.}}\csromangathastanza{\csromangathabat{กุมารี เทฺว จ สํเฆน,}{ทฺวาทส สมฺมเตน จ.} \\ \csromangathabat{อลํ สเจ จ เทฺววสฺสํ,}{สํสฏฺฐา สามิเกน จ.}}\csromangathastanza{\csromangathabat{ปาริวาสิกานุวสฺสํ,}{ทุเว วุฏฺฐาปเนน จ.} \\ \csromangathabat{ฉตฺตยาเนน สงฺฆาณิ,}{อิตฺถาลงฺการวณฺณเก จ.}}\csromangathastanza{\csromangathabat{ปิญฺญากภิกฺขุนี เจว,}{สิกฺขา จ สามเณริกา.} \\ \csromangathabat{คิหิ ภิกฺขุสฺส ปุรโต,}{อโนกาสํ สํกจฺจิกาติ.}}}

\vspace{\tipitakaparskip}
\csromancenter{เตสํ วคฺคานํ อุทฺทานํ.}
\setlength{\csromangathapagewidth}{0pt}
\setlength{\csromangathawakwidth}{0pt}
\csromangathameasuregroup{ลสุณนฺธการา นฺหานา, \\ อารามํ คพฺภินี เจว,}{\csromangathabat{ลสุณนฺธการา นฺหานา,}{ตุวฏฺฏา จิตฺตคารกา.} \\ \csromangathabat{อารามํ คพฺภินี เจว,}{กุมารี ฉตฺตุปาหนาติ.}}
\csromangathasetleft{ลสุณนฺธการา นฺหานา, \\ อารามํ คพฺภินี เจว,}
\csromangathagroup{\csromangathastanza{\csromangathabat{ลสุณนฺธการา นฺหานา,}{ตุวฏฺฏา จิตฺตคารกา.} \\ \csromangathabat{อารามํ คพฺภินี เจว,}{กุมารี ฉตฺตุปาหนาติ.}}}

\csromanmatikaanchor{mtk.88}
\csromantocmarkref{chapter}{5. ปาฏิเทสนียกณฺฑ}{mtk.88}
\bookmark[dest={mtk.88},level=0]{5. ปาฏิเทสนียกณฺฑ}
\chapterheadpageread{5. ปาฏิเทสนีย\-กณฺฑ}
\proseitem{227}{\csromanfolio{125}สปฺปิํ วิญฺญาเปตฺวา ภุญฺชนฺติยา ปาฏิเทสนียํ กตฺถ ปญฺญตฺตนฺติ สาวตฺถิยํ ปญฺญตฺตํ. กํ อารพฺภาติ ฉพฺพคฺคิยา ภิกฺขุนิโย อารพฺภ. กิสฺมิํ วตฺถุสฺมินฺติ ฉพฺพคฺคิยา ภิกฺขุนิโย สปฺปิํ วิญฺญาเปตฺวา ภุญฺชิํสุ, ตสฺมิํ วตฺถุสฺมิํ. เอกา ปญฺญตฺติ, เอกา อนุปญฺญตฺติ. ฉนฺนํ อาปตฺติ\-สมุฏฺฐานานํ จตูหิ สมุฏฺฐาเนหิ สมุฏฺฐาติ ฯเปฯ.}

\prose{เตลํ วิญฺญาเปตฺวา ภุญฺชนฺติยา ปาฏิเทสนียํ กตฺถ ปญฺญตฺตนฺติ สาวตฺถิยํ ปญฺญตฺตํ. กํ อารพฺภาติ ฉพฺพคฺคิยา ภิกฺขุนิโย อารพฺภ. กิสฺมิํ วตฺถุสฺมินฺติ ฉพฺพคฺคิยา ภิกฺขุนิโย เตลํ วิญฺญาเปตฺวา ภุญฺชิํสุ, ตสฺมิํ วตฺถุสฺมิํ. เอกา ปญฺญตฺติ, เอกา อนุปญฺญตฺติ. ฉนฺนํ อาปตฺติ\-สมุฏฺฐานานํ จตูหิ สมุฏฺฐาเนหิ สมุฏฺฐาติ ฯเปฯ.}

\prose{มธุํ วิญฺญาเปตฺวา ภุญฺชนฺติยา ปาฏิเทสนียํ กตฺถ ปญฺญตฺตนฺติ สาวตฺถิยํ ปญฺญตฺตํ. กํ อารพฺภาติ ฉพฺพคฺคิยา ภิกฺขุนิโย อารพฺภ. กิสฺมิํ วตฺถุสฺมินฺติ ฉพฺพคฺคิยา ภิกฺขุนิโย มธุํ วิญฺญาเปตฺวา ภุญฺชิํสุ, ตสฺมิํ วตฺถุสฺมิํ. เอกา ปญฺญตฺติ, เอกา อนุปญฺญตฺติ. ฉนฺนํ อาปตฺติ\-สมุฏฺฐานานํ จตูหิ สมุฏฺฐาเนหิ สมุฏฺฐาติ ฯเปฯ.}

\prose{ผาณิตํ วิญฺญาเปตฺวา ภุญฺชนฺติยา ปาฏิเทสนียํ กตฺถ ปญฺญตฺตนฺติ สาวตฺถิยํ ปญฺญตฺตํ. กํ อารพฺภาติ ฉพฺพคฺคิยา ภิกฺขุนิโย อารพฺภ. กิสฺมิํ วตฺถุสฺมินฺติ ฉพฺพคฺคิยา ภิกฺขุนิโย ผาณิตํ วิญฺญาเปตฺวา ภุญฺชิํสุ, ตสฺมิํ วตฺถุสฺมิํ. เอกา ปญฺญตฺติ, เอกา อนุปญฺญตฺติ. ฉนฺนํ อาปตฺติ\-สมุฏฺฐานานํ จตูหิ สมุฏฺฐาเนหิ สมุฏฺฐาติ ฯเปฯ.}

\prose{มจฺฉํ วิญฺญาเปตฺวา ภุญฺชนฺติยา ปาฏิเทสนียํ กตฺถ ปญฺญตฺตนฺติ สาวตฺถิยํ ปญฺญตฺตํ. กํ อารพฺภาติ ฉพฺพคฺคิยา ภิกฺขุนิโย อารพฺภ. กิสฺมิํ วตฺถุสฺมินฺติ ฉพฺพคฺคิยา ภิกฺขุนิโย มจฺฉํ วิญฺญาเปตฺวา ภุญฺชิํสุ, ตสฺมิํ วตฺถุสฺมิํ. เอกา ปญฺญตฺติ, เอกา อนุปญฺญตฺติ. ฉนฺนํ อาปตฺติ\-สมุฏฺฐานานํ จตูหิ สมุฏฺฐาเนหิ}

\prose{มํสํ วิญฺญาเปตฺวา ภุญฺชนฺติยา ปาฏิเทสนียํ กตฺถ ปญฺญตฺตนฺติ สาวตฺถิยํ ปญฺญตฺตํ. กํ อารพฺภาติ ฉพฺพคฺคิยา ภิกฺขุนิโย อารพฺภ. กิสฺมิํ วตฺถุสฺมินฺติ ฉพฺพคฺคิยา ภิกฺขุนิโย มํสํ วิญฺญาเปตฺวา ภุญฺชิํสุ, ตสฺมิํ วตฺถุสฺมิํ. เอกา ปญฺญตฺติ, เอกา อนุปญฺญตฺติ. ฉนฺนํ อาปตฺติ\-สมุฏฺฐานานํ จตูหิ สมุฏฺฐาเนหิ สมุฏฺฐาติ ฯเปฯ.}

\prose{\csromanfolio{126}ขีรํ วิญฺญาเปตฺวา ภุญฺชนฺติยา ปาฏิเทสนียํ กตฺถ ปญฺญตฺตนฺติ สาวตฺถิยํ ปญฺญตฺตํ. กํ อารพฺภาติ ฉพฺพคฺคิยา ภิกฺขุนิโย อารพฺภ. กิสฺมิํ วตฺถุสฺมินฺติ ฉพฺพคฺคิยา ภิกฺขุนิโย ขีรํ วิญฺญาเปตฺวา ภุญฺชิํสุ, ตสฺมิํ วตฺถุสฺมิํ. เอกา ปญฺญตฺติ, เอกา อนุปญฺญตฺติ. ฉนฺนํ อาปตฺติ\-สมุฏฺฐานานํ จตูหิ สมุฏฺฐาเนหิ สมุฏฺฐาติ ฯเปฯ.}

\prosewithcloser{ทธิํ วิญฺญาเปตฺวา ภุญฺชนฺติยา ปาฏิเทสนียํ กตฺถ ปญฺญตฺตนฺติ สาวตฺถิยํ ปญฺญตฺตํ. กํ อารพฺภาติ ฉพฺพคฺคิยา ภิกฺขุนิโย อารพฺภ. กิสฺมิํ วตฺถุสฺมินฺติ ฉพฺพคฺคิยา ภิกฺขุนิโย ทธิํ วิญฺญาเปตฺวา ภุญฺชิํสุ, ตสฺมิํ วตฺถุสฺมิํ. เอกา ปญฺญตฺติ, เอกา อนุปญฺญตฺติ. ฉนฺนํ อาปตฺติ\-สมุฏฺฐานานํ จตูหิ สมุฏฺฐาเนหิ สมุฏฺฐาติ, สิยา กายโต สมุฏฺฐาติ, น วาจโต, น จิตฺตโต, สิยา กายโต จ วาจโต จ สมุฏฺฐาติ, น จิตฺตโต, สิยา กายโต จ จิตฺตโต จ สมุฏฺฐาติ, น วาจโต, สิยา กายโต จ วาจโต จ จิตฺตโต จ}{อฏฺฐ ปาฏิเทสนียา นิฏฺฐิตา.}
\csromancenter{ตสฺสุทฺทานํ.}
\setlength{\csromangathapagewidth}{0pt}
\setlength{\csromangathawakwidth}{0pt}
\csromangathameasuregroup{สปฺปิํ เตลํ มธุญฺเจว, \\ มํสํ ขีรํ ทธิญฺจาปิ, \\ ปาฏิเทสนียา อฏฺฐ,}{\csromangathabat{สปฺปิํ เตลํ มธุญฺเจว,}{ผาณิตํ มจฺฉเมว จ.} \\ \csromangathabat{มํสํ ขีรํ ทธิญฺจาปิ,}{วิญฺญาเปตฺวาน ภิกฺขุนี.} \\ \csromangathabat{ปาฏิเทสนียา อฏฺฐ,}{สยํ พุทฺเธน เทสิตาติ.}}
\csromangathasetleft{สปฺปิํ เตลํ มธุญฺเจว, \\ มํสํ ขีรํ ทธิญฺจาปิ, \\ ปาฏิเทสนียา อฏฺฐ,}
\csromangathagroup{\csromangathastanza{\csromangathabat{สปฺปิํ เตลํ มธุญฺเจว,}{ผาณิตํ มจฺฉเมว จ.} \\ \csromangathabat{มํสํ ขีรํ ทธิญฺจาปิ,}{วิญฺญาเปตฺวาน ภิกฺขุนี.} \\ \csromangathabat{ปาฏิเทสนียา อฏฺฐ,}{สยํ พุทฺเธน เทสิตาติ.}}}

\prosewithcloserruled{เย สิกฺขาปทา ภิกฺขุ\-วิภงฺเค วิตฺถาริตา, เต สํขิตฺตา ภิกฺขุนิ\-วิภงฺเค.}{กตฺถ\-ปญฺญตฺติวาโร นิฏฺฐิโต ปฐโม.}
\csromanmatikaanchor{mtk.89}
\csromanmatikaanchor{mtk.90}
\csromantocmarknopage{chapter}{2. กตาปตฺติวาร}{mtk.89}
\bookmark[dest={mtk.89},level=0]{2. กตาปตฺติวาร}
\csromantocmarkref{chapter}{1. ปาราชิกกณฺฑ}{mtk.90}
\bookmark[dest={mtk.90},level=0]{1. ปาราชิกกณฺฑ}
\chapterhead{2. \textbf{กตาปตฺติวาร} \textbf{1. ปาราชิกกณฺฑ}}
\proseitem{228}{อวสฺสุตา ภิกฺขุนี อวสฺสุตสฺส ปุริส\-ปุคฺคลสฺส กายสํสคฺคํ สาทิยนฺตี กติ อาปตฺติโย อาปชฺชติ. อวสฺสุตา ภิกฺขุนี อวสฺสุตสฺส ปุริส\-ปุคฺคลสฺส กายสํสคฺคํ สาทิยนฺตี ติสฺโส อาปตฺติโย \csromanfolio{127}อาปชฺชติ. อธกฺขกํ อุพฺภ\-ชาณุมณฺฑลํ คหณํ สาทิยติ, อาปตฺติ ปาราชิกสฺส. อุพฺภกฺขกํ อโธ\-ชาณุมณฺฑลํ คหณํ สาทิยติ, อาปตฺติ ถุลฺลจฺจยสฺส. กาย\-ปฏิพทฺธํ คหณํ สาทิยติ, อาปตฺติ ทุกฺกฏสฺส. อวสฺสุตา ภิกฺขุนี อวสฺสุตสฺส ปุริส\-ปุคฺคลสฺส กายสํสคฺคํ สาทิยนฺตี อิมา ติสฺโส อาปตฺติโย อาปชฺชติ.}

\prose{วชฺชปฺปฏิจฺฉาทิกา ภิกฺขุนี วชฺชํ ปฏิจฺฉาเทนฺตี กติ อาปตฺติโย อาปชฺชติ. วชฺชปฺปฏิจฺฉาทิกา ภิกฺขุนี วชฺชํ ปฏิจฺฉาเทนฺตี ติสฺโส อาปตฺติโย อาปชฺชติ. ชานํ ปาราชิกํ ธมฺมํ ปฏิจฺฉาเทติ, อาปตฺติ ปาราชิกสฺส. เวมติกา ปฏิจฺฉาเทติ, อาปตฺติ ถุลฺลจฺจยสฺส. อาจารวิปตฺติํ ปฏิจฺฉาเทติ, อาปตฺติ ทุกฺกฏสฺส. วชฺชปฺปฏิจฺฉาทิกา ภิกฺขุนี วชฺชํ ปฏิจฺฉาเทนฺตี อิมา ติสฺโส อาปตฺติโย อาปชฺชติ.}

\prose{อุกฺขิตฺตานุวตฺติกา ภิกฺขุนี ยาวตติยํ สมนุภาสนาย น ปฏินิสฺสชฺชนฺตี กติ อาปตฺติโย อาปชฺชติ. อุกฺขิตฺตานุวตฺติกา ภิกฺขุนี ยาวตติยํ สมนุภาสนาย น ปฏินิสฺสชฺชนฺตี ติสฺโส อาปตฺติโย อาปชฺชติ. ญตฺติยา ทุกฺกฏํ. ทฺวีหิ กมฺมวาจาหิ ถุลฺลจฺจยา. กมฺมวาจา\-ปริโยสาเน อาปตฺติ ปาราชิกสฺส. อุกฺขิตฺตานุวตฺติกา ภิกฺขุนี ยาวตติยํ สมนุภาสนาย น ปฏินิสฺสชฺชนฺตี อิมา ติสฺโส อาปตฺติโย อาปชฺชติ.}

\prosewithcloserruled{อฏฺฐมํ วตฺถุํ ปริปูเรนฺตี กติ อาปตฺติโย อาปชฺชติ. อฏฺฐมํ วตฺถุํ ปริปูเรนฺตี ติสฺโส อาปตฺติโย อาปชฺชติ. ปุริเสน “อิตฺถนฺนามํ โอกาสํ1 อาคจฺฉา”ติ วุตฺตา คจฺฉติ, อาปตฺติ ทุกฺกฏสฺส. ปุริสสฺส หตฺถปาสํ โอกฺกนฺตมตฺเต อาปตฺติ ถุลฺลจฺจยสฺส. อฏฺฐมํ วตฺถุํ ปริปูเรติ, อาปตฺติ ปาราชิกสฺส. อฏฺฐมํ วตฺถุํ ปริปูเรนฺตี อิมา ติสฺโส อาปตฺติโย อาปชฺชติ.}{ปาราชิกานิฏฺฐิตา.}
\csromanmatikaanchor{mtk.91}
\csromantocmarkref{chapter}{2. สํฆาทิเสสกณฺฑ}{mtk.91}
\bookmark[dest={mtk.91},level=0]{2. สํฆาทิเสสกณฺฑ}
\chapterhead{2. สํฆาทิเสส\-กณฺฑ}
\proseitem{229}{อุสฺสยวาทิกา ภิกฺขุนี อฑฺฑํ กโรนฺตี ติสฺโส อาปตฺติโย อาปชฺชติ. เอกสฺส อาโรเจติ, อาปตฺติ ทุกฺกฏสฺส. ทุติยสฺส}

\vspace{\tipitakaparskip}
\csromancenter{อิตฺถนฺนามํ คหนํ (สี), อิตฺถนฺนามํ คพฺภํ (สฺยา)}
\prosecont{\csromanfolio{128}อาโรเจติ, อาปตฺติ ถุลฺลจฺจยสฺส. อฑฺฑปริโยสาเน อาปตฺติ สํฆาทิเสสสฺส.}

\prose{โจริํ วุฏฺฐาเปนฺตี ติสฺโส อาปตฺติโย อาปชฺชติ. ญตฺติยา ทุกฺกฏํ. ทฺวีหิ กมฺมวาจาหิ ถุลฺลจฺจยา. กมฺมวาจา\-ปริโยสาเน อาปตฺติ สํฆาทิเสสสฺส.}

\prose{เอกา คามนฺตรํ คจฺฉนฺตี ติสฺโส อาปตฺติโย อาปชฺชติ. คจฺฉติ อาปตฺติ ทุกฺกฏสฺส. ปฐมํ ปาทํ ปริกฺเขปํ อติกฺกาเมติ, อาปตฺติ ถุลฺลจฺจยสฺส. ทุติยํ ปาทํ อติกฺกาเมติ, อาปตฺติ สํฆาทิเสสสฺส.}

\prose{สมคฺเคน สํเฆน อุกฺขิตฺตํ ภิกฺขุนิํ ธมฺเมน วินเยน สตฺถุสาสเนน อนปโลเกตฺวา การกสํฆํ อนญฺญาย คณสฺส ฉนฺทํ โอสาเรนฺตี ติสฺโส อาปตฺติโย อาปชฺชติ. ญตฺติยา ทุกฺกฏํ. ทฺวีหิ กมฺมวาจาหิ ถุลฺลจฺจยา. กมฺมวาจา\-ปริโยสาเน อาปตฺติ สํฆาทิเสสสฺส.}

\prose{อวสฺสุตา ภิกฺขุนี อวสฺสุตสฺส ปุริส\-ปุคฺคลสฺส หตฺถโต ขาทนียํ วา โภชนียํ วา สหตฺถา ปฏิคฺคเหตฺวา ภุญฺชนฺตี ติสฺโส อาปตฺติโย อาปชฺชติ. “ขาทิสฺสามิ ภุญฺชิสฺสามี”ติ ปฏิคฺคณฺหาติ, อาปตฺติ ถุลฺลจฺจยสฺส. อชฺโฌหาเร อชฺโฌหาเร อาปตฺติ สํฆาทิเสสสฺส. อุทก\-ทนฺตโปนํ ปฏิคฺคณฺหาติ, อาปตฺติ ทุกฺกฏสฺส.}

\prose{“กิํ เต อยฺเย เอโส ปุริสปุคฺคโล กริสฺสติ อวสฺสุโต วา อนวสฺสุโต วา, ยโต ตฺวํ อนวสฺสุตา, อิงฺฆ อยฺเย, ยํ เต เอโส ปุริสปุคฺคโล เทติ ขาทนียํ วา โภชนียํ วา, ตํ ตฺวํ สหตฺถา ปฏิคฺคเหตฺวา ขาท วา ภุญฺช วา”ติ อุยฺโยเชนฺตี ติสฺโส อาปตฺติโย อาปชฺชติ. ตสฺสา วจเนน “ขาทิสฺสามิ ภุญฺชิสฺสามี”ติ ปฏิคฺคณฺหาติ, อาปตฺติ ทุกฺกฏสฺส. อชฺโฌหาเร อชฺโฌหาเร อาปตฺติ ถุลฺลจฺจยสฺส. โภชน\-ปริโยสาเน อาปตฺติ สํฆาทิเสสสฺส.}

\prose{กุปิตา ภิกฺขุนี ยาวตติยํ สมนุภาสนาย น ปฏินิสฺสชฺชนฺตี ติสฺโส อาปตฺติโย อาปชฺชติ. ญตฺติยา ทุกฺกฏํ. ทฺวีหิ กมฺมวาจาหิ ถุลฺลจฺจยา. กมฺมวาจา\-ปริโยสาเน อาปตฺติ สํฆาทิเสสสฺส.}

\prose{กิสฺมิญฺจิเทว อธิกรเณ ปจฺจากตา ภิกฺขุนี ยาวตติยํ สมนุภาสนาย น ปฏินิสฺสชฺชนฺตี ติสฺโส อาปตฺติโย อาปชฺชติ. ญตฺติยา \csromanfolio{129}ทุกฺกฏํ. ทฺวีหิ กมฺมวาจาหิ ถุลฺลจฺจยา. กมฺมวาจา\-ปริโยสาเน อาปตฺติ สํฆาทิเสสสฺส.}

\prose{สํสฏฺฐา ภิกฺขุนิโย ยาวตติยํ สมนุภาสนาย น ปฏินิสฺสชฺชนฺติโย ติสฺโส อาปตฺติโย อาปชฺชนฺติ. ญตฺติยา ทุกฺกฏํ. ทฺวีหิ กมฺมวาจาหิ ถุลฺลจฺจยา. กมฺมวาจา\-ปริโยสาเน อาปตฺติ สํฆาทิเสสสฺส.}

\prosewithcloserruled{“สํสฏฺฐาว อยฺเย ตุเมฺห วิหรถ, มา ตุเมฺห นานา วิหริตฺถา”ติ อุยฺโยเชนฺตี ยาวตติยํ สมนุภาสนาย น ปฏินิสฺสชฺชนฺตี ติสฺโส อาปตฺติโย อาปชฺชติ. ญตฺติยา ทุกฺกฏํ. ทฺวีหิ กมฺมวาจาหิ ถุลฺลจฺจยา. กมฺมวาจา\-ปริโยสาเน อาปตฺติ สํฆาทิเสสสฺส.}{สํฆาทิเสสา นิฏฺฐิตา.}
\csromanmatikaanchor{mtk.92}
\csromantocmarkref{chapter}{3. นิสฺสคฺคิยกณฺฑ}{mtk.92}
\bookmark[dest={mtk.92},level=0]{3. นิสฺสคฺคิยกณฺฑ}
\chapterhead{3. นิสฺสคฺคิย\-กณฺฑ}
\proseitem{230}{ปตฺต\-สนฺนิจยํ กโรนฺตี เอกํ อาปตฺติํ อาปชฺชติ นิสฺสคฺคิยํ ปาจิตฺติยํ.}

\prose{อกาลจีวรํ “กาลจีวรนฺ”ติ อธิฏฺฐหิตฺวา ภาชาเปนฺตี เทฺว อาปตฺติโย อาปชฺชติ. ภาชาเปติ, ปโยเค ทุกฺกฏํ. ภาชาปิเต นิสฺสคฺคิยํ ปาจิตฺติยํ.}

\prose{ภิกฺขุนิยา สทฺธิํ จีวรํ ปริวตฺเตตฺวา อจฺฉินฺทนฺตี เทฺว อาปตฺติโย อาปชฺชติ. อจฺฉินฺทติ, ปโยเค ทุกฺกฏํ. อจฺฉินฺเน นิสฺสคฺคิยํ ปาจิตฺติยํ.}

\prose{อญฺญํ วิญฺญาเปตฺวา อญฺญํ วิญฺญาเปนฺตี เทฺว อาปตฺติโย อาปชฺชติ. วิญฺญาเปติ, ปโยเค ทุกฺกฏํ. วิญฺญาปิเต นิสฺสคฺคิยํ ปาจิตฺติยํ.}

\prose{อญฺญํ เจตาเปตฺวา อญฺญํ เจตาเปนฺตี เทฺว อาปตฺติโย อาปชฺชติ. เจตาเปติ, ปโยเค ทุกฺกฏํ. เจตาปิเต นิสฺสคฺคิยํ ปาจิตฺติยํ.}

\prose{อญฺญทตฺถิเกน ปริกฺขาเรน อญฺญุทฺทิสิเกน สํฆิเกน อญฺญํ เจตาเปนฺตี เทฺว อาปตฺติโย อาปชฺชติ. เจตาเปติ, ปโยเค ทุกฺกฏํ. เจตาปิเต นิสฺสคฺคิยํ ปาจิตฺติยํ.}

\prose{\csromanfolio{130}อญฺญทตฺถิเกน ปริกฺขาเรน อญฺญุทฺทิสิเกน สํฆิเกน สํยาจิเกน อญฺญํ เจตาเปนฺตี เทฺว อาปตฺติโย อาปชฺชติ. เจตาเปติ, ปโยเค ทุกฺกฏํ. เจตาปิเต นิสฺสคฺคิยํ ปาจิตฺติยํ.}

\prose{อญฺญทตฺถิเกน ปริกฺขาเรน อญฺญุทฺทิสิเกน มหาชนิเกน อญฺญํ เจตาเปนฺตี เทฺว อาปตฺติโย อาปชฺชติ. เจตาเปติ, ปโยเค ทุกฺกฏํ. เจตาปิเต นิสฺสคฺคิยํ ปาจิตฺติยํ.}

\prose{อญฺญทตฺถิเกน ปริกฺขาเรน อญฺญุทฺทิสิเกน มหาชนิเกน สํยาจิเกน อญฺญํ เจตาเปนฺตี เทฺว อาปตฺติโย อาปชฺชติ. เจตาเปติ, ปโยเค ทุกฺกฏํ. เจตาปิเต นิสฺสคฺคิยํ ปาจิตฺติยํ.}

\prose{อญฺญทตฺถิเกน ปริกฺขาเรน อญฺญุทฺทิสิเกน ปุคฺคลิเกน สํยาจิเกน อญฺญํ เจตาเปนฺตี เทฺว อาปตฺติโย อาปชฺชติ. เจตาเปติ, ปโยเค ทุกฺกฏํ. เจตาปิเต นิสฺสคฺคิยํ ปาจิตฺติยํ.}

\prose{อติเรกจตุกฺกํส\-ปรมํ ครุปาวุรณํ เจตาเปนฺตี เทฺว อาปตฺติโย อาปชฺชติ. เจตาเปติ, ปโยเค ทุกฺกฏํ. เจตาปิเต นิสฺสคฺคิยํ ปาจิตฺติยํ.}

\prosewithcloserruled{อติเรกอฑฺฒเตยฺยกํ สปรมํ ลหุปาวุรณํ เจตาเปนฺตี เทฺว อาปตฺติโย อาปชฺชติ. เจตาเปติ, ปโยเค ทุกฺกฏํ. เจตาปิเต นิสฺสคฺคิยํ ปาจิตฺติยํ.}{นิสฺสคฺคิยา ปาจิตฺติยา นิฏฺฐิตา.}
\csromanmatikaanchor{mtk.93}
\csromanmatikaanchor{mtk.94}
\csromantocmarkref{chapter}{4. ปาจิตฺติยกณฺฑ}{mtk.93}
\bookmark[dest={mtk.93},level=0]{4. ปาจิตฺติยกณฺฑ}
\csromantocmarkref{section}{1. ลสุณวคฺค}{mtk.94}
\bookmark[dest={mtk.94},level=1]{1. ลสุณวคฺค}
\csromanmarkh{4. ปาจิตฺติยกณฺฑ}\csromanmarkhii{1. ลสุณวคฺค}\csromanheadernomark{1}{4. \textbf{ปาจิตฺติยกณฺฑ} \textbf{1. ลสุณวคฺค}}
\proseitem{231}{ลสุณํ ขาทนฺตี เทฺว อาปตฺติโย อาปชฺชติ. “ขาทิสฺสามี”ติ ปฏิคฺคณฺหาติ, อาปตฺติ ทุกฺกฏสฺส. อชฺโฌหาเร อชฺโฌหาเร อาปตฺติ ปาจิตฺติยสฺส.}

\prose{สมฺพาเธ โลมํ สํหราเปนฺตี เทฺว อาปตฺติโย อาปชฺชติ. สํหราเปติ, ปโยเค ทุกฺกฏํ. สํหราปิเต อาปตฺติ ปาจิตฺติยสฺส.}

\prose{\csromanfolio{131}ตลฆาตกํ กโรนฺตี เทฺว อาปตฺติโย อาปชฺชติ. กโรติ, ปโยเค ทุกฺกฏํ. กเต อาปตฺติ ปาจิตฺติยสฺส.}

\prose{ชตุมฏฺฐกํ อาทิยนฺตี เทฺว อาปตฺติโย อาปชฺชติ. อาทิยติ, ปโยเค ทุกฺกฏํ. อาทินฺเน อาปตฺติ ปาจิตฺติยสฺส.}

\prose{อติเรก\-ทฺวงฺคุลปพฺพ\-ปรมํ อุทกสุทฺธิกํ อาทิยนฺตี เทฺว อาปตฺติโย อาปชฺชติ. อาทิยติ, ปโยเค ทุกฺกฏํ. อาทินฺเน อาปตฺติ ปาจิตฺติยสฺส.}

\prose{ภิกฺขุสฺส ภุญฺชนฺตสฺส ปานีเยน วา วิธูปเนน วา อุปติฏฺฐนฺตี เทฺว อาปตฺติโย อาปชฺชติ. หตฺถปาเส ติฏฺฐติ, อาปตฺติ ปาจิตฺติยสฺส. หตฺถปาสํ วิชหิตฺวา ติฏฺฐติ, อาปตฺติ ทุกฺกฏสฺส.}

\prose{อามกธญฺญํ วิญฺญาเปตฺวา ภุญฺชนฺตี เทฺว อาปตฺติโย อาปชฺชติ. “ภุญฺชิสฺสามี”ติ ปฏิคฺคณฺหาติ, อาปตฺติ ทุกฺกฏสฺส. อชฺโฌหาเร อชฺโฌหาเร อาปตฺติ ปาจิตฺติยสฺส.}

\prose{อุจฺจารํ วา ปสฺสาวํ วา สงฺการํ วา วิฆาสํ วา ติโรกุฏฺเฏ วา ติโรปากาเร วา ฉฑฺเฑนฺตี เทฺว อาปตฺติโย อาปชฺชติ. ฉฑฺเฑติ, ปโยเค ทุกฺกฏํ. ฉฑฺฑิเต อาปตฺติ ปาจิตฺติยสฺส.}

\prose{อุจฺจารํ วา ปสฺสาวํ วา สงฺการํ วา วิฆาสํ วา หริเต ฉฑฺเฑนฺตี เทฺว อาปตฺติโย อาปชฺชติ. ฉฑฺเฑติ, ปโยเค ทุกฺกฏํ. ฉฑฺฑิเต อาปตฺติ ปาจิตฺติยสฺส.}

\prosewithclosermidruled{นจฺจํ วา คีตํ วา วาทิตํ วา ทสฺสนาย คจฺฉนฺตี เทฺว อาปตฺติโย อาปชฺชติ. คจฺฉติ, อาปตฺติ ทุกฺกฏสฺส. ยตฺถ ฐิตา ปสฺสติ วา สุณาติ วา, อาปตฺติ ปาจิตฺติยสฺส.}{ลสุณวคฺโค ปฐโม.}
\csromanmatikaanchor{mtk.95}
\csromantocmarkref{section}{2. รตฺตนฺธการวคฺค}{mtk.95}
\bookmark[dest={mtk.95},level=1]{2. รตฺตนฺธการวคฺค}
\csromanheader{1}{2. รตฺตนฺธการ\-วคฺค}
\proseitem{232}{รตฺตนฺธกาเร อปฺปทีเป ปุริเสน สทฺธิํ เอเกเนกา สนฺติฏฺฐนฺตี เทฺว อาปตฺติโย อาปชฺชติ. หตฺถปาเส ติฏฺฐติ, อาปตฺติ ปาจิตฺติยสฺส. หตฺถปาสํ วิชหิตฺวา ติฏฺฐติ, อาปตฺติ ทุกฺกฏสฺส.}

\prose{\csromanfolio{132}ปฏิจฺฉนฺเน โอกาเส ปุริเสน สทฺธิํ เอเกเนกา สนฺติฏฺฐนฺตี เทฺว อาปตฺติโย อาปชฺชติ. หตฺถปาเส ติฏฺฐติ, อาปตฺติ ปาจิตฺติยสฺส. หตฺถปาสํ วิชหิตฺวา ติฏฺฐติ, อาปตฺติ ทุกฺกฏสฺส.}

\prose{อชฺโฌกาเส ปุริเสน สทฺธิํ เอเกเนกา สนฺติฏฺฐนฺตี เทฺว อาปตฺติโย อาปชฺชติ. หตฺถปาเส ติฏฺฐติ, อาปตฺติ ปาจิตฺติยสฺส. หตฺถปาสํ วิชหิตฺวา ติฏฺฐติ, อาปตฺติ ทุกฺกฏสฺส.}

\prose{รถิกาย วาพฺยูเห วา สิงฺฆาฏเก วา ปุริเสน สทฺธิํ เอเกเนกา สนฺติฏฺฐนฺตี เทฺว อาปตฺติโย อาปชฺชติ. หตฺถปาเส ติฏฺฐติ, อาปตฺติ ปาจิตฺติยสฺส. หตฺถปาสํ วิชหิตฺวา ติฏฺฐติ, อาปตฺติ ทุกฺกฏสฺส.}

\prose{ปุเรภตฺตํ กุลานิ อุปสงฺกมิตฺวา อาสเน นิสีทิตฺวา สามิเก อนาปุจฺฉา ปกฺกมนฺตี เทฺว อาปตฺติโย อาปชฺชติ. ปฐมํ ปาทํ อโนวสฺสกํ อติกฺกาเมติ, อาปตฺติ ทุกฺกฏสฺส. ทุติยํ ปาทํ อติกฺกาเมติ, อาปตฺติ ปาจิตฺติยสฺส.}

\prose{ปจฺฉาภตฺตํ กุลานิ อุปสงฺกมิตฺวา สามิเก อนาปุจฺฉา อาสเน นิสีทนฺตี เทฺว อาปตฺติโย อาปชฺชติ. นิสีทติ, ปโยเค ทุกฺกฏํ. นิสินฺเน อาปตฺติ ปาจิตฺติยสฺส.}

\prose{วิกาเล กุลานิ อุปสงฺกมิตฺวา สามิเก อนาปุจฺฉา เสยฺยํ สนฺถริตฺวา วา สนฺถราเปตฺวา วา อภินิสีทนฺตี เทฺว อาปตฺติโย อาปชฺชติ. อภินิสีทติ, ปโยเค ทุกฺกฏํ. อภินิสินฺเน อาปตฺติ ปาจิตฺติยสฺส.}

\prose{ทุคฺคหิเตน ทูปธาริเตน ปรํ อุชฺฌาเปนฺตี เทฺว อาปตฺติโย อาปชฺชติ. อุชฺฌาเปติ, ปโยเค ทุกฺกฏํ. อุชฺฌาปิเต อาปตฺติ ปาจิตฺติยสฺส.}

\prose{อตฺตานํ วา ปรํ วา นิรเยน วา พฺรหฺมจริเยน วา อภิสปนฺตี เทฺว อาปตฺติโย อาปชฺชติ. อภิสปติ, ปโยเค ทุกฺกฏํ. อภิสปิเต อาปตฺติ ปาจิตฺติยสฺส.}

\prosewithclosermidruled{อตฺตานํ วธิตฺวา วธิตฺวา โรทนฺตี เทฺว อาปตฺติโย อาปชฺชติ. วธติ โรทติ, อาปตฺติ ปาจิตฺติยสฺส. วธติ, น โรทติ, อาปตฺติ ทุกฺกฏสฺส.}{รตฺตนฺธการ\-วคฺโค ทุติโย.}
\csromanmatikaanchor{mtk.96}
\csromantocmarkref{section}{3. นหานวคฺค}{mtk.96}
\bookmark[dest={mtk.96},level=1]{3. นหานวคฺค}
\csromanheader{1}{3. \textbf{นหานวคฺค}}
\proseitem{233}{\csromanfolio{133}นคฺคา นหายนฺตี เทฺว อาปตฺติโย อาปชฺชติ. นหายติ, ปโยเค ทุกฺกฏํ. นหาน\-ปริโยสาเน อาปตฺติ ปาจิตฺติยสฺส.}

\prose{ปมาณา\-ติกฺกนฺตํ อุทกสาฏิกํ การาเปนฺตี เทฺว อาปตฺติโย อาปชฺชติ. การาเปติ, ปโยเค ทุกฺกฏํ. การาปิเต อาปตฺติ ปาจิตฺติยสฺส.}

\prose{ภิกฺขุนิยา จีวรํ วิสิพฺเพตฺวา วา วิสิพฺพาเปตฺวา วา เนว สิพฺเพนฺตี น สิพฺพาปนาย อุสฺสุกฺกํ กโรนฺตี เอกํ อาปตฺติํ อาปชฺชติ ปาจิตฺติยํ.}

\prose{ปญฺจาหิกํ สงฺฆาฏิจารํ อติกฺกาเมนฺตี เอกํ อาปตฺติํ อาปชฺชติ ปาจิตฺติยํ.}

\prose{จีวรสงฺกมนียํ ธาเรนฺตี เทฺว อาปตฺติโย อาปชฺชติ. ธาเรติ, ปโยเค ทุกฺกฏํ. ธาริเต อาปตฺติ ปาจิตฺติยสฺส.}

\prose{คณสฺส จีวรลาภํ อนฺตรายํ กโรนฺตี เทฺว อาปตฺติโย อาปชฺชติ. กโรติ, ปโยเค ทุกฺกฏํ. กเต อาปตฺติ ปาจิตฺติยสฺส.}

\prose{ธมฺมิกํ จีวรวิภงฺคํ ปฏิพาหนฺตี เทฺว อาปตฺติโย อาปชฺชติ. ปฏิพาหติ, ปโยเค ทุกฺกฏํ. ปฏิพาหิเต อาปตฺติ ปาจิตฺติยสฺส.}

\prose{อคาริกสฺส วา ปริพฺพาชกสฺส วา ปริพฺพาชิกาย วา สมณจีวรํ เทนฺตี เทฺว อาปตฺติโย อาปชฺชติ. เทติ, ปโยเค ทุกฺกฏํ. ทินฺเน อาปตฺติ ปาจิตฺติยสฺส.}

\prose{ทุพฺพล\-จีวรปจฺจาสาย จีวร\-กาล\-สมยํ อติกฺกาเมนฺตี เทฺว อาปตฺติโย อาปชฺชติ. อติกฺกาเมติ, ปโยเค ทุกฺกฏํ. อติกฺกามิเต อาปตฺติ ปาจิตฺติยสฺส.}

\prosewithclosermidruled{ธมฺมิกํ กถินุทฺธารํ ปฏิพาหนฺตี เทฺว อาปตฺติโย อาปชฺชติ. ปฏิพาหติ, ปโยเค ทุกฺกฏํ. ปฏิพาหิเต อาปตฺติ ปาจิตฺติยสฺส.}{นหานวคฺโค ตติโย.}
\csromanmatikaanchor{mtk.97}
\csromantocmarkref{section}{4. ตุวฏฺฏวคฺค}{mtk.97}
\bookmark[dest={mtk.97},level=1]{4. ตุวฏฺฏวคฺค}
\csromanheader{1}{4. \textbf{ตุวฏฺฏวคฺค}}
\proseitem{234}{\csromanfolio{134}เทฺว ภิกฺขุนิโย เอกมญฺเจ ตุวฏฺเฏนฺติโย เทฺว อาปตฺติโย อาปชฺชนฺติ. นิปชฺชนฺติ, ปโยเค ทุกฺกฏํ. นิปนฺเน อาปตฺติ ปาจิตฺติยสฺส.}

\prose{เทฺว ภิกฺขุนิโย เอก\-ตฺถรณ\-ปาวุรณา ตุวฏฺเฏนฺติโย เทฺว อาปตฺติโย อาปชฺชนฺติ. นิปชฺชนฺติ, ปโยเค ทุกฺกฏํ. นิปนฺเน อาปตฺติ ปาจิตฺติยสฺส.}

\prose{ภิกฺขุนิยา สญฺจิจฺจ อผาสุํ กโรนฺตี เทฺว อาปตฺติโย อาปชฺชติ. กโรติ, ปโยเค ทุกฺกฏํ. กเต อาปตฺติ ปาจิตฺติยสฺส.}

\prose{ทุกฺขิตํ สหชีวินิํ เนว อุปฏฺเฐนฺตี น อุปฏฺฐาปนาย อุสฺสุกฺกํ กโรนฺตี เอกํ อาปตฺติํ อาปชฺชติ ปาจิตฺติยํ.}

\prose{ภิกฺขุนิยา อุปสฺสยํ ทตฺวา กุปิตา อนตฺตมนา นิกฺกฑฺฒนฺตี เทฺว อาปตฺติโย อาปชฺชติ. นิกฺกฑฺฒติ ปโยเค ทุกฺกฏํ. นิกฺกฑฺฒิเต อาปตฺติ ปาจิตฺติยสฺส.}

\prose{สํสฏฺฐา ภิกฺขุนี ยาวตติยํ สมนุภาสนาย น ปฏินิสฺสชฺชนฺตี เทฺว อาปตฺติโย อาปชฺชติ. ญตฺติยา ทุกฺกฏํ. กมฺมวาจา\-ปริโยสาเน อาปตฺติ ปาจิตฺติยสฺส.}

\prose{อนฺโตรฏฺเฐ สาสงฺกสมฺมเต สปฺปฏิภเย อสตฺถิกา จาริกํ จรนฺตี เทฺว อาปตฺติโย อาปชฺชติ. ปฏิปชฺชติ, ปโยเค ทุกฺกฏํ. ปฏิปนฺเน อาปตฺติ ปาจิตฺติยสฺส.}

\prose{ติโรรฏฺเฐ สาสงฺกสมฺมเต สปฺปฏิภเย อสตฺถิกา จาริกํ จรนฺตี เทฺว อาปตฺติโย อาปชฺชติ. ปฏิปชฺชติ, ปโยเค ทุกฺกฏํ. ปฏิปนฺเน อาปตฺติ ปาจิตฺติยสฺส.}

\prosewithclosermidruled{อนฺโตวสฺสํ จาริกํ จรนฺตี เทฺว อาปตฺติโย อาปชฺชติ. ปฏิปชฺชติ, ปโยเค ทุกฺกฏํ. ปฏิปนฺเน อาปตฺติ ปาจิตฺติยสฺส. วสฺสํวุฏฺฐา ภิกฺขุนี จาริกํ น ปกฺกมนฺตี เอกํ อาปตฺติํ อาปชฺชติ ปาจิตฺติยํ.}{ตุวฏฺฏวคฺโค จตุตฺโถ.}
\csromanmatikaanchor{mtk.98}
\csromantocmarkref{section}{5. จิตฺตาคารวคฺค}{mtk.98}
\bookmark[dest={mtk.98},level=1]{5. จิตฺตาคารวคฺค}
\csromanheader{1}{5. \textbf{จิตฺตาคารวคฺค}}
\proseitem{235}{\csromanfolio{135}ราชาคารํ วา จิตฺตาคารํ วา อารามํ วา อุยฺยานํ วา โปกฺขรณิํ วา ทสฺสนาย คจฺฉนฺตี เทฺว อาปตฺติโย อาปชฺชติ. คจฺฉติ, อาปตฺติ ทุกฺกฏสฺส. ยตฺถ ฐิตา ปสฺสติ, อาปตฺติ ปาจิตฺติยสฺส.}

\prose{อาสนฺทิํ วา ปลฺลงฺกํ วา ปริภุญฺชนฺตี เทฺว อาปตฺติโย อาปชฺชติ. ปริภุญฺชติ, ปโยเค ทุกฺกฏํ. ปริภุตฺเต อาปตฺติ ปาจิตฺติยสฺส.}

\prose{สุตฺตํ กนฺตนฺตี เทฺว อาปตฺติโย อาปชฺชติ. กนฺตติ, ปโยเค ทุกฺกฏํ. อุชฺชวุชฺชเว อาปตฺติ ปาจิตฺติยสฺส.}

\prose{คิหิเวยฺยาวจฺจํ กโรนฺตี เทฺว อาปตฺติโย อาปชฺชติ. กโรติ, ปโยเค ทุกฺกฏํ. กเต อาปตฺติ ปาจิตฺติยสฺส.}

\prose{ภิกฺขุนิยา “เอหายฺเย, อิมํ อธิกรณํ วูปสเมหี”ติ วุจฺจมานา “สาธู”ติ ปฏิสฺสุณิตฺวา เนว วูปสเมนฺตี น วูปสมาย อุสฺสุกฺกํ กโรนฺตี เอกํ อาปตฺติํ อาปชฺชติ ปาจิตฺติยํ.}

\prose{อคาริกสฺส วา ปริพฺพาชกสฺส วา ปริพฺพาชิกาย วา สหตฺถา ขาทนียํ วา โภชนียํ วา เทนฺตี เทฺว อาปตฺติโย อาปชฺชติ. เทติ, ปโยเค ทุกฺกฏํ. ทินฺเน อาปตฺติ ปาจิตฺติยสฺส.}

\prose{อาวสถจีวรํ อนิสฺสชฺชิตฺวา ปริภุญฺชนฺตี เทฺว อาปตฺติโย อาปชฺชติ. ปริภุญฺชติ, ปโยเค ทุกฺกฏํ. ปริภุตฺเต อาปตฺติ ปาจิตฺติยสฺส.}

\prose{อาวสถํ อนิสฺสชฺชิตฺวา จาริกํ ปกฺกมนฺตี เทฺว อาปตฺติโย อาปชฺชติ. ปฐมํ ปาทํ ปริกฺเขปํ อติกฺกาเมติ, อาปตฺติ ทุกฺกฏสฺส. ทุติยํ ปาทํ อติกฺกาเมติ, อาปตฺติ ปาจิตฺติยสฺส.}

\prose{ติรจฺฉาน\-วิชฺชํ ปริยาปุณนฺตี เทฺว อาปตฺติโย อาปชฺชติ. ปริยาปุณาติ, ปโยเค ทุกฺกฏํ. ปเท ปเท อาปตฺติ ปาจิตฺติยสฺส.}

\prosewithclosermidruled{ติรจฺฉาน\-วิชฺชํ วาเจนฺตี เทฺว อาปตฺติโย อาปชฺชติ. วาเจติ, ปโยเค ทุกฺกฏํ. ปเท ปเท อาปตฺติ ปาจิตฺติยสฺส.}{จิตฺตคารวคฺโค ปญฺจโม.}
\csromanmatikaanchor{mtk.99}
\csromantocmarkref{section}{6. อารามวคฺค}{mtk.99}
\bookmark[dest={mtk.99},level=1]{6. อารามวคฺค}
\csromanheader{1}{6. \textbf{อารามวคฺค}}
\proseitem{236}{\csromanfolio{136}ชานํ สภิกฺขุกํ อารามํ อนาปุจฺฉา ปวิสนฺตี เทฺว อาปตฺติโย อาปชฺชติ. ปฐมํ ปาทํ ปริกฺเขปํ อติกฺกาเมติ, อาปตฺติ ทุกฺกฏสฺส. ทุติยํ ปาทํ อติกฺกาเมติ, อาปตฺติ ปาจิตฺติยสฺส.}

\prose{ภิกฺขุํ อกฺโกสนฺตี ปริภาสนฺตี เทฺว อาปตฺติโย อาปชฺชติ. อกฺโกสติ, ปโยเค ทุกฺกฏํ. อกฺโกสิเต อาปตฺติ ปาจิตฺติยสฺส.}

\prose{จณฺฑีกตา คณํ ปริภาสนฺตี เทฺว อาปตฺติโย อาปชฺชติ. ปริภาสติ, ปโยเค ทุกฺกฏํ. ปริภาสิเต อาปตฺติ ปาจิตฺติยสฺส.}

\prose{นิมนฺติตา วา ปวาริตา วา ขาทนียํ วา โภชนียํ วา ภุญฺชนฺตี เทฺว อาปตฺติโย อาปชฺชติ. “ขาทิสฺสามิ ภุญฺชิสฺสามี”ติ ปฏิคฺคณฺหาติ, อาปตฺติ ทุกฺกฏสฺส. อชฺโฌหาเร อชฺโฌหาเร อาปตฺติ ปาจิตฺติยสฺส.}

\prose{กุลํ มจฺฉรายนฺตี เทฺว อาปตฺติโย อาปชฺชติ. มจฺฉรายติ, ปโยเค ทุกฺกฏํ. มจฺฉริเต อาปตฺติ ปาจิตฺติยสฺส.}

\prose{อภิกฺขุเก อาวาเส วสฺสํ วสนฺตี เทฺว อาปตฺติโย อาปชฺชติ. “วสฺสํ วสิสฺสามี”ติ เสนาสนํ ปญฺญเปติ, ปานียํ ปริโภชนียํ อุปฏฺฐเปติ, ปริเวณํ สมฺมชฺชติ, อาปตฺติ ทุกฺกฏสฺส. สห อรุณุคฺคมนา อาปตฺติ ปาจิตฺติยสฺส.}

\prose{วสฺสํวุฏฺฐา ภิกฺขุนี อุภโตสํเฆ ตีหิ ฐาเนหิ น ปวาเรนฺตี เอกํ อาปตฺติํ อาปชฺชติ, ปาจิตฺติยํ.}

\prose{โอวาทาย วา สํวาสาย วา น คจฺฉนฺตี เอกํ อาปตฺติํ อาปชฺชติ, ปาจิตฺติยํ.}

\prose{อุโปสถมฺปิ น ปุจฺฉนฺตี โอวาทมฺปิ น ยาจนฺตี เอกํ อาปตฺติํ อาปชฺชติ, ปาจิตฺติยํ.}

\prosewithclosermidruled{ปสาเข ชาตํ คณฺฑํ วา รุธิตํ วา อนปโลเกตฺวา สํฆํ วา คณํ วา ปุริเสน สทฺธิํ เอเกเนกา เภทาเปนฺตี เทฺว อาปตฺติโย อาปชฺชติ. เภทาเปติ, ปโยเค ทุกฺกฏํ. ภินฺเน อาปตฺติ ปาจิตฺติยสฺส.}{อารามวคฺโค ฉฏฺโฐ.}
\csromanmatikaanchor{mtk.100}
\csromantocmarkref{section}{7. คพฺภินีวคฺค}{mtk.100}
\bookmark[dest={mtk.100},level=1]{7. คพฺภินีวคฺค}
\csromanheader{1}{7. \textbf{คพฺภินีวคฺค}}
\proseitem{237}{\csromanfolio{137}คพฺภินิํ วุฏฺฐาเปนฺตี เทฺว อาปตฺติโย อาปชฺชติ. วุฏฺฐาเปติ, ปโยเค ทุกฺกฏํ. วุฏฺฐาปิเต อาปตฺติ ปาจิตฺติยสฺส.}

\prose{ปายนฺติํ วุฏฺฐาเปนฺตี เทฺว อาปตฺติโย อาปชฺชติ. วุฏฺฐาเปติ, ปโยเค ทุกฺกฏํ. วุฏฺฐาปิเต อาปตฺติ ปาจิตฺติยสฺส.}

\prose{เทฺว วสฺสานิ ฉสุ ธมฺเมสุ อสิกฺขิตสิกฺขํ สิกฺขมานํ วุฏฺฐาเปนฺตี เทฺว อาปตฺติโย อาปชฺชติ, วุฏฺฐาเปติ, ปโยเค ทุกฺกฏํ. วุฏฺฐาปิเต อาปตฺติ ปาจิตฺติยสฺส.}

\prose{เทฺว วสฺสานิ ฉสุ ธมฺเมสุ สิกฺขิตสิกฺขํ สิกฺขมานํ สํเฆน อสมฺมตํ วุฏฺฐาเปนฺตี เทฺว อาปตฺติโย อาปชฺชติ. วุฏฺฐาเปติ, ปโยเค ทุกฺกฏํ. วุฏฺฐาปิเต อาปตฺติ ปาจิตฺติยสฺส.}

\prose{อูนทฺวาทสวสฺสํ คิหิคตํ วุฏฺฐาเปนฺตี เทฺว อาปตฺติโย อาปชฺชติ. วุฏฺฐาเปติ, ปโยเค ทุกฺกฏํ. วุฏฺฐาปิเต อาปตฺติ ปาจิตฺติยสฺส.}

\prose{ปริปุณฺณ\-ทฺวาทสวสฺสํ คิหิคตํ เทฺว วสฺสานิ ฉสุ ธมฺเมสุ อสิกฺขิตสิกฺขํ วุฏฺฐาเปนฺตี เทฺว อาปตฺติโย อาปชฺชติ. วุฏฺฐาเปติ, ปโยเค ทุกฺกฏํ. วุฏฺฐาปิเต อาปตฺติ ปาจิตฺติยสฺส.}

\prose{ปริปุณฺณ\-ทฺวาทสวสฺสํ คิหิคตํ เทฺว วสฺสานิ ฉสุ ธมฺเมสุ สิกฺขิตสิกฺขํ สํเฆน อสมฺมตํ วุฏฺฐาเปนฺตี เทฺว อาปตฺติโย อาปชฺชติ. วุฏฺฐาเปติ, ปโยเค ทุกฺกฏํ. วุฏฺฐาปิเต อาปตฺติ ปาจิตฺติยสฺส.}

\prose{สหชีวินิํ วุฏฺฐาเปตฺวา เทฺว วสฺสานิ เนว อนุคฺคณฺหนฺตี นานุ\-คฺคณฺหาเปนฺตี เอกํ อาปตฺติํ อาปชฺชติ, ปาจิตฺติยํ.}

\prose{วุฏฺฐาปิตํ ปวตฺตินิํ เทฺว วสฺสานิ นานุพนฺธนฺตี เอกํ อาปตฺติํ อาปชฺชติ, ปาจิตฺติยํ.}

\prosewithclosermidruled{สหชีวินิํ วุฏฺฐาเปตฺวา เนว วูปกาเสนฺตี น วูปกาสาเปนฺตี เอกํ อาปตฺติํ อาปชฺชติ, ปาจิตฺติยํ.}{คพฺภินิวคฺโค สตฺตโม.}
\csromanmatikaanchor{mtk.101}
\csromantocmarkref{section}{8. กุมารีภูตวคฺค}{mtk.101}
\bookmark[dest={mtk.101},level=1]{8. กุมารีภูตวคฺค}
\csromanheader{1}{8. กุมารีภูต\-วคฺค}
\proseitem{238}{\csromanfolio{138}อูนวีสติวสฺสํ กุมาริภูตํ วุฏฺฐาเปนฺตี เทฺว อาปตฺติโย อาปชฺชติ. วุฏฺฐาเปติ, ปโยเค ทุกฺกฏํ. วุฏฺฐาปิเต อาปตฺติ ปาจิตฺติยสฺส.}

\prose{ปริปุณฺณ\-วีสติวสฺสํ กุมาริภูตํ เทฺว วสฺสานิ ฉสุ ธมฺเมสุ อสิกฺขิตสิกฺขํ วุฏฺฐาเปนฺตี เทฺว อาปตฺติโย อาปชฺชติ. วุฏฺฐาเปติ, ปโยเค ทุกฺกฏํ. วุฏฺฐาปิเต อาปตฺติ ปาจิตฺติยสฺส.}

\prose{ปริปุณฺณ\-วีสติวสฺสํ กุมาริภูตํ เทฺว วสฺสานิ ฉสุ ธมฺเมสุ สิกฺขิตสิกฺขํ สํเฆน อสมฺมตํ วุฏฺฐาเปนฺตี เทฺว อาปตฺติโย อาปชฺชติ. วุฏฺฐาเปติ, ปโยเค ทุกฺกฏํ. วุฏฺฐาปิเต อาปตฺติ ปาจิตฺติยสฺส.}

\prose{อูนทฺวาทสวสฺสา วุฏฺฐาเปนฺตี เทฺว อาปตฺติโย อาปชฺชติ. วุฏฺฐาเปติ, ปโยเค ทุกฺกฏํ. วุฏฺฐาปิเต อาปตฺติ ปาจิตฺติยสฺส.}

\prose{ปริปุณฺณ\-ทฺวาทสวสฺสา สํเฆน อสมฺมตา วุฏฺฐาเปนฺตี เทฺว อาปตฺติโย อาปชฺชติ. วุฏฺฐาเปติ, ปโยเค ทุกฺกฏํ. วุฏฺฐาปิเต อาปตฺติ ปาจิตฺติยสฺส.}

\prose{“อลํ ตาว เต อยฺเย วุฏฺฐาปิเตนา”ติ วุจฺจมานา “สาธู”ติ ปฏิสฺสุณิตฺวา ปจฺฉา ขียนธมฺมํ อาปชฺชนฺตี เทฺว อาปตฺติโย อาปชฺชติ. ขิยฺยติ, ปโยเค ทุกฺกฏํ. ขิยฺยิเต อาปตฺติ ปาจิตฺติยสฺส.}

\prose{สิกฺขมานํ “สเจ เม ตฺวํ อยฺเย จีวรํ ทสฺสสิ, เอวาหํ ตํ วุฏฺฐาเปสฺสามี”ติ วตฺวา เนว วุฏฺฐาเปนฺตี น วุฏฺฐาปนาย อุสฺสุกฺกํ กโรนฺตี เอกํ อาปตฺติํ อาปชฺชติ, ปาจิตฺติยํ.}

\prose{สิกฺขมานํ “สเจ เม ตฺวํ อยฺเย เทฺว วสฺสานิ อนุพนฺธิสฺสสิ, เอวาหํ ตํ วุฏฺฐาเปสฺสามี”ติ วตฺวา เนว วุฏฺฐาเปนฺตี น วุฏฺฐาปนาย อุสฺสุกฺกํ กโรนฺตี เอกํ อาปตฺติํ อาปชฺชติ, ปาจิตฺติยํ.}

\prose{ปุริส\-สํสฏฺฐํ กุมารก\-สํสฏฺฐํ จณฺฑิํ โสกาวาสํ สิกฺขมานํ วุฏฺฐาเปนฺตี เทฺว อาปตฺติโย อาปชฺชติ. วุฏฺฐาเปติ, ปโยเค ทุกฺกฏํ. วุฏฺฐาปิเต อาปตฺติ ปาจิตฺติยสฺส.}

\prose{มาตาปิตูหิ วา สามิเกน วา อนนุญฺญาตํ สิกฺขมานํ วุฏฺฐาเปนฺตี เทฺว อาปตฺติโย อาปชฺชติ. วุฏฺฐาเปติ, ปโยเค ทุกฺกฏํ. วุฏฺฐาปิเต อาปตฺติ ปาจิตฺติยสฺส.}

\prose{\csromanfolio{139}ปาริวาสิก\-ฉนฺททาเนน สิกฺขมานํ วุฏฺฐาเปนฺตี เทฺว อาปตฺติโย อาปชฺชติ. วุฏฺฐาเปติ, ปโยเค ทุกฺกฏํ. วุฏฺฐาปิเต อาปตฺติ ปาจิตฺติยสฺส.}

\prose{อนุวสฺสํ วุฏฺฐาเปนฺตี เทฺว อาปตฺติโย อาปชฺชติ. วุฏฺฐาเปติ, ปโยเค ทุกฺกฏํ. วุฏฺฐาปิเต อาปตฺติ ปาจิตฺติยสฺส.}

\prosewithclosermidruled{เอกํ วสฺสํ เทฺว วุฏฺฐาเปนฺตี เทฺว อาปตฺติโย อาปชฺชติ. วุฏฺฐาเปติ, ปโยเค ทุกฺกฏํ. วุฏฺฐาปิเต อาปตฺติ ปาจิตฺติยสฺส.}{กุมาริภูต\-วคฺโค อฏฺฐโม.}
\csromanmatikaanchor{mtk.102}
\csromantocmarkref{section}{9. ฉตฺตุปาหนวคฺค}{mtk.102}
\bookmark[dest={mtk.102},level=1]{9. ฉตฺตุปาหนวคฺค}
\csromanheader{1}{9. ฉตฺตุปาหน\-วคฺค}
\proseitem{239}{ฉตฺตุปาหนํ ธาเรนฺตี เทฺว อาปตฺติโย อาปชฺชติ. ธาเรติ, ปโยเค ทุกฺกฏํ. ธาริเต อาปตฺติ ปาจิตฺติยสฺส.}

\prose{ยาเนน ยายนฺตี เทฺว อาปตฺติโย อาปชฺชติ. ยายติ, ปโยเค ทุกฺกฏํ. ยายิเต อาปตฺติ ปาจิตฺติยสฺส.}

\prose{สงฺฆาณิํ ธาเรนฺตี เทฺว อาปตฺติโย อาปชฺชติ. ธาเรติ, ปโยเค ทุกฺกฏํ. ธาริเต อาปตฺติ ปาจิตฺติยสฺส.}

\prose{อิตฺถาลงฺการํ ธาเรนฺตี เทฺว อาปตฺติโย อาปชฺชติ. ธาเรติ, ปโยเค ทุกฺกฏํ. ธาริเต อาปตฺติ ปาจิตฺติยสฺส.}

\prose{คนฺธ\-วณฺณเกน นหายนฺตี เทฺว อาปตฺติโย อาปชฺชติ. นหายติ, ปโยเค ทุกฺกฏํ. นหาน\-ปริโยสาเน อาปตฺติ ปาจิตฺติยสฺส.}

\prose{วาสิตเกน ปิญฺญาเกน นหายนฺตี เทฺว อาปตฺติโย อาปชฺชติ. นหายติ, ปโยเค ทุกฺกฏํ. นหาน\-ปริโยสาเน อาปตฺติ ปาจิตฺติยสฺส.}

\prose{ภิกฺขุนิยา อุมฺมทฺทาเปนฺตี ปริมทฺทาเปนฺตี เทฺว อาปตฺติโย อาปชฺชติ. อุมฺมทฺทาเปติ, ปโยเค ทุกฺกฏํ. อุมฺมทฺทิเต อาปตฺติ ปาจิตฺติยสฺส.}

\prose{สิกฺขมานาย อุมฺมทฺทาเปนฺตี ปริมทฺทาเปนฺตี เทฺว อาปตฺติโย อาปชฺชติ. อุมฺมทฺทาเปติ, ปโยเค ทุกฺกฏํ. อุมฺมทฺทิเต อาปตฺติ ปาจิตฺติยสฺส.}

\prose{\csromanfolio{140}สามเณริยา อุมฺมทฺทาเปนฺตี ปริมทฺทาเปนฺตี เทฺว อาปตฺติโย อาปชฺชติ. อุมฺมทฺทาเปติ, ปโยเค ทุกฺกฏํ. อุมฺมทฺทิเต อาปตฺติ ปาจิตฺติยสฺส.}

\prose{คิหินิยา อุมฺมทฺทาเปนฺตี ปริมทฺทาเปนฺตี เทฺว อาปตฺติโย อาปชฺชติ. อุมฺมทฺทาเปติ, ปโยเค ทุกฺกฏํ. อุมฺมทฺทิเต อาปตฺติ ปาจิตฺติยสฺส.}

\prose{ภิกฺขุสฺส ปุรโต อนาปุจฺฉา อาสเน นิสีทนฺตี เทฺว อาปตฺติโย อาปชฺชติ. นิสีทติ, ปโยเค ทุกฺกฏํ. นิสินฺเน อาปตฺติ ปาจิตฺติยสฺส.}

\prose{อโนกาสกตํ ภิกฺขุํ ปญฺหํ ปุจฺฉนฺตี เทฺว อาปตฺติโย อาปชฺชติ. ปุจฺฉติ, ปโยเค ทุกฺกฏํ. ปุจฺฉิเต อาปตฺติ ปาจิตฺติยสฺส.}

\prosewithclosermidruled{อสํกจฺจิกา คามํ ปวิสนฺตี เทฺว อาปตฺติโย อาปชฺชติ. ปฐมํ ปาทํ ปริกฺเขปํ อติกฺกาเมติ, อาปตฺติ ทุกฺกฏสฺส. ทุติยํ ปาทํ อติกฺกาเมติ, อาปตฺติ ปาจิตฺติยสฺส.}{ฉตฺตุปาหน\-วคฺโค นวโม. ขุทฺทกํ นิฏฺฐิตํ.}
\csromanmatikaanchor{mtk.103}
\csromantocmarkref{chapter}{5. ปาฏิเทสนียกณฺฑ}{mtk.103}
\bookmark[dest={mtk.103},level=0]{5. ปาฏิเทสนียกณฺฑ}
\chapterhead{5. ปาฏิเทสนีย\-กณฺฑ}
\proseitem{240}{สปฺปิํ วิญฺญาเปตฺวา ภุญฺชนฺตี เทฺว อาปตฺติโย อาปชฺชติ. “ภุญฺชิสฺสามี”ติ ปฏิคฺคณฺหาติ, อาปตฺติ ทุกฺกฏสฺส. อชฺโฌหาเร อชฺโฌหาเร อาปตฺติ ปาฏิเทสนียสฺส.}

\prose{เตลํ วิญฺญาเปตฺวา ภุญฺชนฺตี เทฺว อาปตฺติโย อาปชฺชติ. “ภุญฺชิสฺสามี”ติ ปฏิคฺคณฺหาติ, อาปตฺติ ทุกฺกฏสฺส. อชฺโฌหาเร อชฺโฌหาเร อาปตฺติ ปาฏิเทสนียสฺส.}

\prose{มธุํ วิญฺญาเปตฺวา ภุญฺชนฺตี เทฺว อาปตฺติโย อาปชฺชติ. “ภุญฺชิสฺสามี”ติ ปฏิคฺคณฺหาติ, อาปตฺติ ทุกฺกฏสฺส. อชฺโฌหาเร อชฺโฌหาเร อาปตฺติ ปาฏิเทสนียสฺส.}

\prose{ผาณิตํ วิญฺญาเปตฺวา ภุญฺชนฺตี เทฺว อาปตฺติโย อาปชฺชติ “ภุญฺชิสฺสามี”ติ ปฏิคฺคณฺหาติ, อาปตฺติ ทุกฺกฏสฺส. อชฺโฌหาเร อชฺโฌหาเร อาปตฺติ ปาฏิเทสนียสฺส.}

\prose{\csromanfolio{141}มจฺฉํ วิญฺญาเปตฺวา ภุญฺชนฺตี เทฺว อาปตฺติโย อาปชฺชติ. “ภุญฺชิสฺสามี”ติ ปฏิคฺคณฺหาติ, อาปตฺติ ทุกฺกฏสฺส. อชฺโฌหาเร อชฺโฌหาเร อาปตฺติ ปาฏิเทสนียสฺส.}

\prose{มํสํ วิญฺญาเปตฺวา ภุญฺชนฺตี เทฺว อาปตฺติโย อาปชฺชติ. “ภุญฺชิสฺสามี”ติ ปฏิคฺคณฺหาติ, อาปตฺติ ทุกฺกฏสฺส. อชฺโฌหาเร อชฺโฌหาเร อาปตฺติ ปาฏิเทสนียสฺส.}

\prose{ขีรํ วิญฺญาเปตฺวา ภุญฺชนฺตี เทฺว อาปตฺติโย อาปชฺชติ. “ภุญฺชิสฺสามี”ติ ปฏิคฺคณฺหาติ, อาปตฺติ ทุกฺกฏสฺส. อชฺโฌหาเร อชฺโฌหาเร อาปตฺติ ปาฏิเทสนียสฺส.}

\prosewithcloser{ทธิํ วิญฺญาเปตฺวา ภุญฺชนฺตี เทฺว อาปตฺติโย อาปชฺชติ. “ภุญฺชิสฺสามี”ติ ปฏิคฺคณฺหาติ, อาปตฺติ ทุกฺกฏสฺส. อชฺโฌหาเร อชฺโฌหาเร อาปตฺติ ปาฏิเทสนียสฺส.}{อฏฺฐ ปาฏิเทสนียา นิฏฺฐิตา.}
\nitthitamruled{กตาปตฺติวาโร นิฏฺฐิโต ทุติโย.}
\csromanmatikaanchor{mtk.104}
\csromantocmarkref{chapter}{3. วิปตฺติวาร}{mtk.104}
\bookmark[dest={mtk.104},level=0]{3. วิปตฺติวาร}
\chapterhead{3. \textbf{วิปตฺติวาร}}
\proseitem{241}{อวสฺสุตาย ภิกฺขุนิยา อวสฺสุตสฺส ปุริส\-ปุคฺคลสฺส กายสํสคฺคํ สาทิยนฺติยา อาปตฺติโย จตุนฺนํ วิปตฺตีนํ กติ วิปตฺติโย ภชนฺติ. อวสฺสุตาย ภิกฺขุนิยา อวสฺสุตสฺส ปุริส\-ปุคฺคลสฺส กายสํสคฺคํ สาทิยนฺติยา อาปตฺติโย จตุนฺนํ วิปตฺตีนํ เทฺว วิปตฺติโย ภชนฺติ, สิยา สีลวิปตฺติํ, สิยา อาจารวิปตฺติํ ฯเปฯ.}

\prose{ทธิํ วิญฺญาเปตฺวา ภุญฺชนฺติยา อาปตฺติโย จตุนฺนํ วิปตฺตีนํ กติ วิปตฺติโย ภชนฺติ, ทธิํ วิญฺญาเปตฺวา ภุญฺชนฺติยา อาปตฺติโย จตุนฺนํ วิปตฺตีนํ เอกํวิปตฺติํ ภชนฺติ, อาจารวิปตฺติํ.}

\vspace{\tipitakaparskip}
\csromancenter{วิปตฺติวาโร นิฏฺฐิโต ตติโย.}
\csromanmatikaanchor{mtk.105}
\csromantocmarkref{chapter}{4. สงฺคหวาร}{mtk.105}
\bookmark[dest={mtk.105},level=0]{4. สงฺคหวาร}
\chapterheadpageread{4. \textbf{สงฺคหวาร}}
\proseitem{242}{\csromanfolio{142}อวสฺสุตาย ภิกฺขุนิยา อวสฺสุตสฺส ปุริส\-ปุคฺคลสฺส กายสํสคฺคํ สาทิยนฺติยา อาปตฺติโย สตฺตนฺนํ อาปตฺติ\-กฺขนฺธานํ กติหิ อาปตฺติ\-กฺขนฺเธหิ สงฺคหิตา. อวสฺสุตาย ภิกฺขุนิยา อวสฺสุตสฺส ปุริส\-ปุคฺคลสฺส กายสํสคฺคํ สาทิยนฺติยา อาปตฺติโย สตฺตนฺนํ อาปตฺติ\-กฺขนฺธานํ ตีหิ อาปตฺติ\-กฺขนฺเธหิ สงฺคหิตา, สิยา ปาราชิกาปตฺติ\-กฺขนฺเธน, สิยา ถุลฺลจฺจยา\-ปตฺติกฺขนฺเธน, สิยา ทุกฺกฏาปตฺติ\-กฺขนฺเธน ฯเปฯ.}

\prosewithcloserruled{ทธิํ วิญฺญาเปตฺวา ภุญฺชนฺติยา อาปตฺติโย สตฺตนฺนํ อาปตฺติ\-กฺขนฺธานํ กติหิ อาปตฺติ\-กฺขนฺเธหิ สงฺคหิตา. ทธิํ วิญฺญาเปตฺวา ภุญฺชนฺติยา อาปตฺติโย สตฺตนฺนํ อาปตฺติ\-กฺขนฺธานํ ทฺวีหิ อาปตฺติ\-กฺขนฺเธหิ สงฺคหิตา, สิยา ปาฏิเทสนียา\-ปตฺติกฺขนฺเธน, สิยา ทุกฺกฏาปตฺติ\-กฺขนฺเธน.}{สงฺคหวาโร นิฏฺฐิโต จตุตฺโถ.}
\csromanmatikaanchor{mtk.106}
\csromantocmarkref{chapter}{5. สมุฏฺฐานวาร}{mtk.106}
\bookmark[dest={mtk.106},level=0]{5. สมุฏฺฐานวาร}
\chapterhead{5. \textbf{สมุฏฺฐานวาร}}
\proseitem{243}{อวสฺสุตาย ภิกฺขุนิยา อวสฺสุตสฺส ปุริส\-ปุคฺคลสฺส กายสํสคฺคํ สาทิยนฺติยา อาปตฺติโย ฉนฺนํ อาปตฺติ\-สมุฏฺฐานานํ กติหิ สมุฏฺฐาเนหิ สมุฏฺฐนฺติ. อวสฺสุตาย ภิกฺขุนิยา อวสฺสุตสฺส ปุริส\-ปุคฺคลสฺส กายสํสคฺคํ สาทิยนฺติยา อาปตฺติโย ฉนฺนํ อาปตฺติ\-สมุฏฺฐานานํ เอเกน สมุฏฺฐาเนน สมุฏฺฐนฺติ, กายโต จ จิตฺตโต จ สมุฏฺฐนฺติ, น วาจโต ฯเปฯ.}

\prose{ทธิํ วิญฺญาเปตฺวา ภุญฺชนฺติยา อาปตฺติโย ฉนฺนํ อาปตฺติ\-สมุฏฺฐานานํ กติหิ สมุฏฺฐาเนหิ สมุฏฺฐนฺติ. ทธิํ วิญฺญาเปตฺวา ภุญฺชนฺติยา อาปตฺติโย ฉนฺนํ อาปตฺติ\-สมุฏฺฐานานํ จตูหิ สมุฏฺฐาเนหิ สมุฏฺฐนฺติ, สิยา กายโต สมุฏฺฐนฺติ, น วาจโต, น จิตฺตโต, สิยา กายโต จ วาจโต จ สมุฏฺฐนฺติ, น จิตฺตโต, สิยา กายโต จ จิตฺตโต จ สมุฏฺฐนฺติ, น วาจโต, สิยา กายโต จ วาจโต จ จิตฺตโต จ สมุฏฺฐนฺติ.}

\vspace{\tipitakaparskip}
\csromancenter{สมุฏฺฐานวาโร นิฏฺฐิโต ปญฺจโม.}
\csromanmatikaanchor{mtk.107}
\csromantocmarkref{chapter}{6. อธิกรณวาร}{mtk.107}
\bookmark[dest={mtk.107},level=0]{6. อธิกรณวาร}
\chapterheadpageread{6. \textbf{อธิกรณวาร}}
\proseitem{244}{\csromanfolio{143}อวสฺสุตาย ภิกฺขุนิยา อวสฺสุตสฺส ปุริส\-ปุคฺคลสฺส กายสํสคฺคํ สาทิยนฺติยา อาปตฺติโย จตุนฺนํ อธิกรณานํ กตมํ อธิกรณํ. อวสฺสุตาย ภิกฺขุนิยา อวสฺสุตสฺส ปุริส\-ปุคฺคลสฺส กายสํสคฺคํ สาทิยนฺติยา อาปตฺติโย จตุนฺนํ อธิกรณานํ อาปตฺตา\-ธิกรณํ ฯเปฯ.}

\prosewithcloserruled{ทธิํ วิญฺญาเปตฺวา ภุญฺชนฺติยา อาปตฺติโย จตุนฺนํ อธิกรณานํ กตมํ อธิกรณํ. ทธิํ วิญฺญาเปตฺวา ภุญฺชนฺติยา อาปตฺติโย จตุนฺนํ อธิกรณานํ อาปตฺตา\-ธิกรณํ.}{อธิกรณวาโร นิฏฺฐิโต ฉฏฺโฐ.}
\csromanmatikaanchor{mtk.108}
\csromantocmarkref{chapter}{7. สมถวาร}{mtk.108}
\bookmark[dest={mtk.108},level=0]{7. สมถวาร}
\chapterhead{7. \textbf{สมถวาร}}
\proseitem{245}{อวสฺสุตาย ภิกฺขุนิยา อวสฺสุตสฺส ปุริส\-ปุคฺคลสฺส กายสํสคฺคํ สาทิยนฺติยา อาปตฺติโย สตฺตนฺนํ สมถานํ กติหิ สมเถหิ สมฺมนฺติ. อวสฺสุตาย ภิกฺขุนิยา อวสฺสุตสฺส ปุริส\-ปุคฺคลสฺส กายสํสคฺคํ สาทิยนฺติยา อาปตฺติโย สตฺตนฺนํ สมถานํ ตีหิ สมเถหิ สมฺมนฺติ, สิยา สมฺมุขา\-วินเยน จ ปฏิญฺญาต\-กรเณน จ, สิยา สมฺมุขา\-วินเยน จ ติณวตฺถารเกน จ ฯเปฯ.}

\prosewithcloserruled{ทธิํ วิญฺญาเปตฺวา ภุญฺชนฺติยา อาปตฺติโย สตฺตนฺนํ สมถานํ กติหิ สมเถหิ สมฺมนฺติ. ทธิํ วิญฺญาเปตฺวา ภุญฺชนฺติยา อาปตฺติโย สตฺตนฺนํ สมถานํ ตีหิ สมเถหิ สมฺมนฺติ, สิยา สมฺมุขา\-วินเยน จ ปฏิญฺญาต\-กรเณน จ, สิยา สมฺมุขา\-วินเยน จ ติณวตฺถารเกน จ.}{สมถวาโร นิฏฺฐิโต สตฺตโม.}
\csromanmatikaanchor{mtk.109}
\csromantocmarkref{chapter}{8. สมุจฺจยวาร}{mtk.109}
\bookmark[dest={mtk.109},level=0]{8. สมุจฺจยวาร}
\chapterhead{8. \textbf{สมุจฺจยวาร}}
\proseitem{246}{อวสฺสุตา ภิกฺขุนี อวสฺสุตสฺส ปุริส\-ปุคฺคลสฺส กายสํสคฺคํ สาทิยนฺตี กติ อาปตฺติโย อาปชฺชติ. อวสฺสุตา ภิกฺขุนี อวสฺสุตสฺส \csromanfolio{144}ปุริส\-ปุคฺคลสฺส กายสํสคฺคํ สาทิยนฺตี ติสฺโส อาปตฺติโย อาปชฺชติ. อธกฺขกํ อุพฺภ\-ชาณุมณฺฑลํ คหณํ สาทิยติ, อาปตฺติ ปาราชิกสฺส. อุพฺภกฺขกํ อโธ\-ชาณุมณฺฑลํ คหณํ สาทิยติ, อาปตฺติ ถุลฺลจฺจยสฺส. กาย\-ปฏิพทฺธํ คหณํ สาทิยติ, อาปตฺติ ทุกฺกฏสฺส. อวสฺสุตา ภิกฺขุนี อวสฺสุตสฺส ปุริส\-ปุคฺคลสฺส กายสํสคฺคํ สาทิยนฺตี อิมา ติสฺโส อาปตฺติโย อาปชฺชติ.}

\prose{ตา อาปตฺติโย จตุนฺนํ วิปตฺตีนํ กติ วิปตฺติโย ภชนฺติ, สตฺตนฺนํ อาปตฺติ\-กฺขนฺธานํ กติหิ อาปตฺติ\-กฺขนฺเธหิ สงฺคหิตา, ฉนฺนํ อาปตฺติ\-สมุฏฺฐานานํ กติหิ สมุฏฺฐาเนหิ สมุฏฺฐนฺติ, จตุนฺนํ อธิกรณานํ กตมํ อธิกรณํ, สตฺตนฺนํ สมถานํ กติหิ สมเถหิ สมฺมนฺติ. ตา อาปตฺติโย จตุนฺนํ วิปตฺตีนํ เทฺว วิปตฺติโย ภชนฺติ, สิยา สีลวิปตฺติํ, สิยา อาจารวิปตฺติํ. สตฺตนฺนํ อาปตฺติ\-กฺขนฺธานํ ตีหิ อาปตฺติ\-กฺขนฺเธหิ สงฺคหิตา, สิยา ปาราชิกาปตฺติ\-กฺขนฺเธน, สิยา ถุลฺลจฺจยา\-ปตฺติกฺขนฺเธน, สิยา ทุกฺกฏาปตฺติ\-กฺขนฺเธน. ฉนฺนํ อาปตฺติ\-สมุฏฺฐานานํ เอเกน สมุฏฺฐาเนน สมุฏฺฐนฺติ, กายโต จ จิตฺตโต จ สมุฏฺฐนฺติ, น วาจโต. จตุนฺนํ อธิกรณานํ อาปตฺตา\-ธิกรณํ. สตฺตนฺนํ สมถานํ ตีหิ สมเถหิ สมฺมนฺติ, สิยา สมฺมุขา\-วินเยน จ ปฏิญฺญาต\-กรเณน จ, สิยา สมฺมุขา\-วินเยน จ ติณวตฺถารเกน จ ฯเปฯ.}

\prose{ทธิํ วิญฺญาเปตฺวา ภุญฺชนฺตี กติ อาปตฺติโย อาปชฺชติ. ทธิํ วิญฺญาเปตฺวา ภุญฺชนฺตี เทฺว อาปตฺติโย อาปชฺชติ. “ภุญฺชิสฺสามี”ติ ปฏิคฺคณฺหาติ, อาปตฺติ ทุกฺกฏสฺส. อชฺโฌหาเร อชฺโฌหาเร อาปตฺติ ปาฏิเทสนียสฺส. ทธิํ วิญฺญาเปตฺวา ภุญฺชนฺตี อิมา เทฺว อาปตฺติโย อาปชฺชติ.}

\prosewithcloserruled{ตา อาปตฺติโย จตุนฺนํ วิปตฺตีนํ กติ วิปตฺติโย ภชนฺติ, สตฺตนฺนํ อาปตฺติ\-กฺขนฺธานํ กติหิ อาปตฺติ\-กฺขนฺเธหิ สงฺคหิตา, ฉนฺนํ อาปตฺติ\-สมุฏฺฐานานํ กติหิ สมุฏฺฐาเนหิ สมุฏฺฐนฺติ, จตุนฺนํ อธิกรณานํ กตมํ อธิกรณํ, สตฺตนฺนํ สมถานํ กติหิ สมเถหิ สมฺมนฺติ. ตา อาปตฺติโย จตุนฺนํ วิปตฺตีนํ เอกํ วิปตฺติํ ภชนฺติ, อาจารวิปตฺติํ. สตฺตนฺนํ อาปตฺติ\-กฺขนฺธานํ ทฺวีหิ อาปตฺติ\-กฺขนฺเธหิ สงฺคหิตา, สิยา ปาฏิเทสนียา\-ปตฺติกฺขนฺเธน, สิยา ทุกฺกฏาปตฺติ\-กฺขนฺเธน. ฉนฺนํ อาปตฺติ\-สมุฏฺฐานานํ จตูหิ สมุฏฺฐาเนหิ สมุฏฺฐนฺติ, สิยา กายโต สมุฏฺฐนฺติ, น วาจโต, น จิตฺตโต, สิยา กายโต จ วาจโต จ สมุฏฺฐนฺติ, น จิตฺตโต, \csromanfolio{145}สิยา กายโต จ จิตฺตโต จ สมุฏฺฐนฺติ, น วาจโต, สิยา กายโต จ วาจโต จ จิตฺตโต จ สมุฏฺฐนฺติ. จตุนฺนํ อธิกรณานํ อาปตฺตา\-ธิกรณํ. สตฺตนฺนํ สมถานํ ตีหิ สมเถหิ สมฺมนฺติ, สิยา สมฺมุขา\-วินเยน จ ปฏิญฺญาต\-กรเณน จ, สิยา สมฺมุขา\-วินเยน จ ติณวตฺถารเกน จ.}{สมุจฺจยวาโร นิฏฺฐิโต อฏฺฐโม.}
\csromanmatikaanchor{mtk.110}
\csromanmatikaanchor{mtk.111}
\csromantocmarknopage{section}{1. กตฺถปญฺญตฺติวาร}{mtk.110}
\bookmark[dest={mtk.110},level=1]{1. กตฺถปญฺญตฺติวาร}
\csromantocmarkref{chapter}{1. ปาราชิกกณฺฑ}{mtk.111}
\bookmark[dest={mtk.111},level=0]{1. ปาราชิกกณฺฑ}
\chapterhead{1. กตฺถ\-ปญฺญตฺติวาร 1. ปาราชิกกณฺฑ}
\proseitem{247}{ยํ เตน ภควตา ชานตา ปสฺสตา อรหตา สมฺมา\-สมฺพุทฺเธน กายสํสคฺคํ สาทิยน\-ปจฺจยา ปาราชิกํ กตฺถ ปญฺญตฺตํ, กํ อารพฺภ, กิสฺมิํ วตฺถุสฺมิํ ฯเปฯ เกนาภตนฺติ.}

\prose{ยํ เตน ภควตา ชานตา ปสฺสตา อรหตา สมฺมา\-สมฺพุทฺเธน กายสํสคฺคํ สาทิยน\-ปจฺจยา ปาราชิกํ กตฺถ ปญฺญตฺตนฺติ สาวตฺถิยํ ปญฺญตฺตํ. กํ อารพฺภาติ สุนฺทรีนนฺทํ ภิกฺขุนิํ อารพฺภ. กิสฺมิํ วตฺถุสฺมินฺติ สุนฺทรีนนฺทา ภิกฺขุนี อวสฺสุตา อวสฺสุตสฺส ปุริส\-ปุคฺคลสฺส กายสํสคฺคํ สาทิยิ, ตสฺมิํ วตฺถุสฺมิํ. อตฺถิ ตตฺถ ปญฺญตฺติ อนุปญฺญตฺติ อนุปฺปนฺนปญฺญตฺตีติ เอกา ปญฺญตฺติ, อนุปญฺญตฺติ อนุปฺปนฺนปญฺญตฺติ ตสฺมิํ นตฺถิ. สพฺพตฺถ\-ปญฺญตฺติ ปเทสปญฺญตฺตีติ สพฺพตฺถ\-ปญฺญตฺติ. สาธารณ\-ปญฺญตฺติ อสาธารณปญฺญตฺตีติ อสาธารณ\-ปญฺญตฺติ. เอกโตปญฺญตฺติ อุภโตปญฺญตฺตีติ เอกโตปญฺญตฺติ. จตุนฺนํ ปาติโมกฺขุ\-ทฺเทสานํ กตฺโถคธํ กตฺถ ปริยาปนฺนนฺติ นิทาโนคธํ นิทาน\-ปริยาปนฺนํ. กตเมน อุทฺเทเสน อุทฺเทสํ อาคจฺฉตีติ ทุติยน อุทฺเทเสน อุทฺเทสํ อาคจฺฉติ. จตุนฺนํ วิปตฺตีนํ กตมา วิปตฺตีติ สีลวิปตฺติ. สตฺตนฺนํ อาปตฺติ\-กฺขนฺธานํ กตโม อาปตฺติ\-กฺขนฺโธติ ปาราชิกาปตฺติ\-กฺขนฺโธ. ฉนฺนํ อาปตฺติ\-สมุฏฺฐานานํ กติหิ สมุฏฺฐาเนหิ สมุฏฺฐาตีติ เอเกน สมุฏฺฐาเนน สมุฏฺฐาติ, กายโต จ จิตฺตโต จ สมุฏฺฐาติ, น วาจโต ฯเปฯ เกนาภตนฺติ ปรมฺปราภตํ.}

\setlength{\csromangathapagewidth}{0pt}
\setlength{\csromangathawakwidth}{0pt}
\csromangathameasuregroup{อุปาลิ ทาสโก เจว, \\ โมคฺคลิปุตฺเตน ปญฺจมา, \\ เอเต นาคา มหาปญฺญา, \\ วินยํ ทีเป ปกาเสสุํ,}{\csromangathabat{อุปาลิ ทาสโก เจว,}{โสณโก สิคฺคโว ตถา.} \\ \csromangathabat{โมคฺคลิปุตฺเตน ปญฺจมา,}{เอเต ชมฺพุสิริวฺหเย ฯเปฯ.} \\ \csromangathabat{เอเต นาคา มหาปญฺญา,}{วินยญฺญู มคฺคโกวิทา.} \\ \csromangathabat{วินยํ ทีเป ปกาเสสุํ,}{ปิฏกํ ตมฺพปณฺณิยาติ.}}
\csromangathasetleft{อุปาลิ ทาสโก เจว, \\ โมคฺคลิปุตฺเตน ปญฺจมา, \\ เอเต นาคา มหาปญฺญา, \\ วินยํ ทีเป ปกาเสสุํ,}
\csromangathagroup{\csromangathastanza{\csromanfolio{146}\csromangathabat{อุปาลิ ทาสโก เจว,}{โสณโก สิคฺคโว ตถา.} \\ \csromangathabat{โมคฺคลิปุตฺเตน ปญฺจมา,}{เอเต ชมฺพุสิริวฺหเย ฯเปฯ.}}\csromangathastanza{\csromangathabat{เอเต นาคา มหาปญฺญา,}{วินยญฺญู มคฺคโกวิทา.} \\ \csromangathabat{วินยํ ทีเป ปกาเสสุํ,}{ปิฏกํ ตมฺพปณฺณิยาติ.}}}

\prose{วชฺช\-ปฺปฏิจฺฉาทน\-ปจฺจยา ปาราชิกํ กตฺถ ปญฺญตฺตนฺติ สาวตฺถิยํ ปญฺญตฺตํ. กํ อารพฺภาติ ถุลฺลนนฺทํ ภิกฺขุนิํ อารพฺภ. กิสฺมิํ วตฺถุสฺมินฺติ ถุลฺลนนฺทา ภิกฺขุนี ชานํ ปาราชิกํ ธมฺมํ อชฺฌาปนฺนํ ภิกฺขุนิํ เนวตฺตนา ปฏิโจเทสิ, น คณสฺส อาโรเจสิ, ตสฺมิํ วตฺถุสฺมิํ. เอกา ปญฺญตฺติ. ฉนฺนํ อาปตฺติ\-สมุฏฺฐานานํ เอเกน สมุฏฺฐาเนน สมุฏฺฐาติ, ธุรนิกฺเขเป ฯเปฯ.}

\prose{ยาวตติยํ สมนุภาสนาย น ปฏินิสฺสชฺชน\-ปจฺจยา ปาราชิกํ กตฺถ ปญฺญตฺตนฺติ สาวตฺถิยํ ปญฺญตฺตํ. กํ อารพฺภาติ ถุลฺลนนฺทํ ภิกฺขุนิํ อารพฺภ. กิสฺมิํ วตฺถุสฺมินฺติ ถุลฺลนนฺทา ภิกฺขุนี สมคฺเคน สํเฆน อุกฺขิตฺตํ อริฏฺฐํ ภิกฺขุํ คทฺธพาธิ\-ปุพฺพํ อนุวตฺติ, ตสฺมิํ วตฺถุสฺมิํ. เอกา ปญฺญตฺติ. ฉนฺนํ อาปตฺติ\-สมุฏฺฐานานํ เอเกน สมุฏฺฐาเนน สมุฏฺฐาติ, ธุรนิกฺเขเป ฯเปฯ.}

\prosewithcloserruled{อฏฺฐมํ วตฺถุํ ปริปูรณ\-ปจฺจยา ปาราชิกํ กตฺถ ปญฺญตฺตนฺติ สาวตฺถิยํ ปญฺญตฺตํ. กํ อารพฺภาติ ฉพฺพคฺคิยา ภิกฺขุนิโย อารพฺภ. กิสฺมิํ วตฺถุสฺมินฺติ ฉพฺพคฺคิยา ภิกฺขุนิโย อฏฺฐมํ วตฺถุํ ปริปูเรสุํ, ตสฺมิํ วตฺถุสฺมิํ. เอกา ปญฺญตฺติ. ฉนฺนํ อาปตฺติ\-สมุฏฺฐานานํ เอเกน สมุฏฺฐาเนน สมุฏฺฐาติ, ธุรนิกฺเขเป ฯเปฯ.}{ปาราชิกา นิฏฺฐิตา.}
\csromanmatikaanchor{mtk.112}
\csromantocmarkref{chapter}{2. สํฆาทิเสสกณฺฑาทิ}{mtk.112}
\bookmark[dest={mtk.112},level=0]{2. สํฆาทิเสสกณฺฑาทิ}
\chapterhead{2. \textbf{สํฆาทิเสสกณฺฑาทิ}}
\proseitem{248}{ยํ เตน ภควตา ชานตา ปสฺสตา อรหตา สมฺมา\-สมฺพุทฺเธน อุสฺสยวาทิกาย ภิกฺขุนิยา อฑฺฑํ กรณปจฺจยา สํฆาทิเสโส กตฺถ ปญฺญตฺโต, กํ อารพฺภ, กิสฺมิํ วตฺถุสฺมิํ ฯเปฯ เกนาภตนฺติ.}

\prose{ยํ เตน ภควตา ชานตา ปสฺสตา อรหตา สมฺมา\-สมฺพุทฺเธน อุสฺสยวาทิกาย ภิกฺขุนิยา อฑฺฑํ กรณปจฺจยา สํฆาทิเสโส กตฺถ \csromanfolio{147}ปญฺญตฺโตติ สาวตฺถิยํ ปญฺญตฺโต. กํ อารพฺภาติ ถุลฺลนนฺทํ ภิกฺขุนิํ อารพฺภ. กิสฺมิํ วตฺถุสฺมินฺติ ถุลฺลนนฺทา ภิกฺขุนี อุสฺสยวาทิกา วิหริ, ตสฺมิํ วตฺถุสฺมิํ. อตฺถิ ตตฺถ ปญฺญตฺติ อนุปญฺญตฺติ อนุปฺปนฺนปญฺญตฺตีติ เอกา ปญฺญตฺติ, อนุปญฺญตฺติ อนุปฺปนฺนปญฺญตฺติ ตสฺมิํ นตฺถิ. สพฺพตฺถ\-ปญฺญตฺติ ปเทสปญฺญตฺตีติ สพฺพตฺถ\-ปญฺญตฺติ. สาธารณ\-ปญฺญตฺติ อสาธารณปญฺญตฺตีติ อสาธารณ\-ปญฺญตฺติ. เอกโตปญฺญตฺติ อุภโตปญฺญตฺตีติ เอกโตปญฺญตฺติ. จตุนฺนํ ปาติโมกฺขุ\-ทฺเทสานํ กตฺโถคธํ กตฺถ ปริยาปนฺนนฺติ นิทาโนคธํ นิทาน\-ปริยาปนฺนํ. กตเมน อุทฺเทเสน อุทฺเทสํ อาคจฺฉตีติ ตติเยน อุทฺเทเสน อุทฺเทสํ อาคจฺฉติ. จตุนฺนํ วิปตฺตีนํ กตมา วิปตฺตีติ สีลวิปตฺติ. สตฺตนฺนํ อาปตฺติ\-กฺขนฺธานํ กตโม อาปตฺติ\-กฺขนฺโธติ สํฆาทิเสสา\-ปตฺติกฺขนฺโธ. ฉนฺนํ อาปตฺติ\-สมุฏฺฐานานํ กติหิ สมุฏฺฐาเนหิ สมุฏฺฐาตีติ ทฺวีหิ สมุฏฺฐาเนหิ สมุฏฺฐาติ, สิยา กายโต จ วาจโต จ สมุฏฺฐาติ น จิตฺตโต, สิยา กายโต จ วาจโต จ จิตฺตโต จ สมุฏฺฐาติ ฯเปฯ เกนาภตนฺติ ปรมฺปราภตํ.}

\setlength{\csromangathapagewidth}{0pt}
\setlength{\csromangathawakwidth}{0pt}
\csromangathameasuregroup{อุปาลิ ทาสโก เจว, \\ โมคฺคลิปุตฺเตน ปญฺจมา, \\ เอเต นาคา มหาปญฺญา, \\ วินยํ ทีเป ปกาเสสุํ,}{\csromangathabat{อุปาลิ ทาสโก เจว,}{โสณโก สิคฺคโว ตถา.} \\ \csromangathabat{โมคฺคลิปุตฺเตน ปญฺจมา,}{เอเต ชมฺพุสิริวฺหเย ฯเปฯ.} \\ \csromangathabat{เอเต นาคา มหาปญฺญา,}{วินยญฺญู มคฺคโกวิทา.} \\ \csromangathabat{วินยํ ทีเป ปกาเสสุํ,}{ปิฏกํ ตมฺพปณฺณิยาติ.}}
\csromangathasetleft{อุปาลิ ทาสโก เจว, \\ โมคฺคลิปุตฺเตน ปญฺจมา, \\ เอเต นาคา มหาปญฺญา, \\ วินยํ ทีเป ปกาเสสุํ,}
\csromangathagroup{\csromangathastanza{\csromangathabat{อุปาลิ ทาสโก เจว,}{โสณโก สิคฺคโว ตถา.} \\ \csromangathabat{โมคฺคลิปุตฺเตน ปญฺจมา,}{เอเต ชมฺพุสิริวฺหเย ฯเปฯ.}}\csromangathastanza{\csromangathabat{เอเต นาคา มหาปญฺญา,}{วินยญฺญู มคฺคโกวิทา.} \\ \csromangathabat{วินยํ ทีเป ปกาเสสุํ,}{ปิฏกํ ตมฺพปณฺณิยาติ.}}}

\prose{โจริํ วุฏฺฐาปน\-ปจฺจยา สํฆาทิเสโส กตฺถ ปญฺญตฺโตติ สาวตฺถิยํ ปญฺญตฺโต. กํ อารพฺภาติ ถุลฺลนนฺทํ ภิกฺขุนิํ อารพฺภ. กิสฺมิํ วตฺถุสฺมินฺติ ถุลฺลนนฺทา ภิกฺขุนี โจริํ วุฏฺฐาเปสิ, ตสฺมิํ วตฺถุสฺมิํ. เอกา ปญฺญตฺติ. ฉนฺนํ อาปตฺติ\-สมุฏฺฐานานํ ทฺวีหิ สมุฏฺฐาเนหิ สมุฏฺฐาติ, สิยา วาจโต จ จิตฺตโต จ สมุฏฺฐาติ, น กายโต, สิยา กายโต จ วาจโต จ}

\prose{เอกา คามนฺตรํ คมนปจฺจยา สํฆาทิเสโส กตฺถ ปญฺญตฺโตติ สาวตฺถิยํ ปญฺญตฺโต. กํ อารพฺภาติ อญฺญตรํ ภิกฺขุนิํ อารพฺภ. กิสฺมิํ วตฺถุสฺมินฺติ อญฺญตรา ภิกฺขุนี เอกา คามนฺตรํ คจฺฉิ, ตสฺมิํ วตฺถุสฺมิํ. เอกา ปญฺญตฺติ, ติสฺโส อนุปญฺญตฺติโย. ฉนฺนํ อาปตฺติ\-สมุฏฺฐานานํ เอเกน สมุฏฺฐาเนน สมุฏฺฐาติ, ปฐม\-ปาราชิเก ฯเปฯ.}

\prose{\csromanfolio{148}สมคฺเคน สํเฆน อุกฺขิตฺตํ ภิกฺขุนิํ ธมฺเมน วินเยน สตฺถุสาสเนน อนปโลเกตฺวา การกสํฆํ อนญฺญาย คณํสฺส ฉนฺทํ โอสารณปจฺจยา สํฆาทิเสโส กตฺถ ปญฺญตฺโตติ สาวตฺถิยํ ปญฺญตฺโต. กํ อารพฺภาติ ถุลฺลนนฺทํ ภิกฺขุนิํ อารพฺภ. กิสฺมิํ วตฺถุสฺมินฺติ ถุลฺลนนฺทา ภิกฺขุนี สมคฺเคน สํเฆน อุกฺขิตฺตํ ภิกฺขุนิํ ธมฺเมน วินเยน สตฺถุสาสเนน อนปโลเกตฺวา การกสํฆํ อนญฺญาย คณสฺส ฉนฺทํ โอสาเรสิ, ตสฺมิํ วตฺถุสฺมิํ. เอกา ปญฺญตฺติ. ฉนฺนํ อาปตฺติ\-สมุฏฺฐานานํ เอเกน สมุฏฺฐาเนน สมุฏฺฐาติ, ธุรนิกฺเขเป ฯเปฯ.}

\prose{อวสฺสุตาย ภิกฺขุนิยา อวสฺสุตสฺส ปุริส\-ปุคฺคลสฺส หตฺถโต ขาทนียํ วา โภชนียํ วา สหตฺถา ปฏิคฺคเหตฺวา ภุญฺชน\-ปจฺจยา สํฆาทิเสโส กตฺถ ปญฺญตฺโตติ สาวตฺถิยํ ปญฺญตฺโต. กํ อารพฺภาติ สุนฺทรีนนฺทํ ภิกฺขุนิํ อารพฺภ. กิสฺมิํ วตฺถุสฺมินฺติ สุนฺทรีนนฺทา ภิกฺขุนี อวสฺสุตา อวสฺสุตสฺส ปุริส\-ปุคฺคลสฺส หตฺถโต อามิสํ ปฏิคฺคเหสิ, ตสฺมิํ วตฺถุสฺมิํ. เอกา ปญฺญตฺติ. ฉนฺนํ อาปตฺติ\-สมุฏฺฐานานํ เอเกน สมุฏฺฐาเนน สมุฏฺฐาติ, ปฐม\-ปาราชิเก ฯเปฯ.}

\prose{“กิํ เต อยฺเย เอโส ปุริสปุคฺคโล กริสฺสติ อวสฺสุโต วา อนวสฺสุโต วา, ยโต ตฺวํ อนวสฺสุตา, อิงฺฆ อยฺเย, ยํ เต เอโส ปุริสปุคฺคโล เทติ ขาทนียํ วา โภชนียํ วา, ตํ ตฺวํ สหตฺถา ปฏิคฺคเหตฺวา ขาท วา ภุญฺช วา”ติ อุยฺโยชน\-ปจฺจยา สํฆาทิเสโส กตฺถ ปญฺญตฺโตติ สาวตฺถิยํ ปญฺญตฺโต. กํ อารพฺภาติ อญฺญตรํ ภิกฺขุนิํ อารพฺภ. กิสฺมิํ วตฺถุสฺมินฺติ อญฺญตรา ภิกฺขุนี “กิํ เต อยฺเย เอโส ปุริสปุคฺคโล กริสฺสติ อวสฺสุโต วา อนวสฺสุโต วา, ยโต ตฺวํ อนวสฺสุตา, อิงฺฆ อยฺเย, ยํ เต เอโส ปุริสปุคฺคโล เทติ ขาทนียํ วา โภชนียํ วา, ตํ ตฺวํ สหตฺถา ปฏิคฺคเหตฺวา ขาท วา ภุญฺช วา”ติ อุยฺโยเชสิ, ตสฺมิํ วตฺถุสฺมิํ. เอกา ปญฺญตฺติ. ฉนฺนํ อาปตฺติ\-สมุฏฺฐานานํ ตีหิ สมุฏฺฐาเนหิ สมุฏฺฐาติ ฯเปฯ.}

\prose{กุปิตาย อนตฺตมนาย ภิกฺขุนิยา ยาวตติยํ สมนุภาสนาย น ปฏินิสฺสชฺชน\-ปจฺจยา สํฆาทิเสโส กตฺถ ปญฺญตฺโตติ สาวตฺถิยํ ปญฺญตฺโต. กํ อารพฺภาติ จณฺฑกาฬิํ ภิกฺขุนิํ อารพฺภ. กิสฺมิํ วตฺถุสฺมินฺติ จณฺฑกาฬี ภิกฺขุนี กุปิตา อนตฺตมนา เอวํ อวจ “พุทฺธํ ปจฺจาจิกฺขามิ, ธมฺมํ ปจฺจาจิกฺขามิ, สํฆํ ปจฺจาจิกฺขามิ, สิกฺขํ ปจฺจาจิกฺขามี”ติ, ตสฺมิํ \csromanfolio{149}วตฺถุสฺมิํ. เอกา ปญฺญตฺติ. ฉนฺนํ อาปตฺติ\-สมุฏฺฐานานํ เอเกน สมุฏฺฐาเนน สมุฏฺฐาติ, ธุรนิกฺเขเป ฯเปฯ.}

\prose{กิสฺมิญฺจิเทว อธิกรเณ ปจฺจากตาย ภิกฺขุนิยา ยาวตติยํ สมนุภาสนาย น ปฏินิสฺสชฺชน\-ปจฺจยา สํฆาทิเสโส กตฺถ ปญฺญตฺโตติ สาวตฺถิยํ ปญฺญตฺโต. กํ อารพฺภาติ จณฺฑกาฬิํ ภิกฺขุนิํ อารพฺภ. กิสฺมิํ วตฺถุสฺมินฺติ จณฺฑกาฬี ภิกฺขุนี กิสฺมิญฺจิเทว อธิกรเณ ปจฺจากตา กุปิตา อนตฺตมนา เอวํ อวจ “ฉนฺทคามินิโย จ ภิกฺขุนิโย, โทสคามินิโย จ ภิกฺขุนิโย, โมหคามินิโย จ ภิกฺขุนิโย, ภยคามินิโย จ ภิกฺขุนิโย”ติ, ตสฺมิํ วตฺถุสฺมิํ. เอกา ปญฺญตฺติ. ฉนฺนํ อาปตฺติ\-สมุฏฺฐานานํ เอเกน สมุฏฺฐาเนน สมุฏฺฐาติ, ธุรนิกฺเขเป ฯเปฯ.}

\prose{สํสฏฺฐานํ ภิกฺขุนีนํ ยาวตติยํ สมนุภาสนาย น ปฏินิสฺสชฺชน\-ปจฺจยา สํฆาทิเสโส กตฺถ ปญฺญตฺโตติ สาวตฺถิยํ ปญฺญตฺโต. กํ อาราพฺภาติ สมฺพหุลา ภิกฺขุนิโย อารพฺภ. กิสฺมิํ วตฺถุสฺมินฺติ สมฺพหุลา ภิกฺขุนิโย สํสฏฺฐา วิหริํสุ, ตสฺมิํ วตฺถุสฺมิํ. เอกา ปญฺญตฺติ. ฉนฺนํ อาปตฺติ\-สมุฏฺฐานานํ เอเกน สมุฏฺฐาเนน สมุฏฺฐาติ, ธุรนิกฺเขเป ฯเปฯ.}

\prose{“สํสฏฺฐาว อยฺเย ตุเมฺห วิหรถ, มา ตุเมฺห นานา วิหริตฺถา”ติ อุยฺโยเชนฺติยา ยาวตติยํ สมนุภาสนาย น ปฏินิสฺสชฺชน\-ปจฺจยา สํฆาทิเสโส กตฺถ ปญฺญตฺโตติ สาวตฺถิยํ ปญฺญตฺโต. กํ อารพฺภาติ ถุลฺลนนฺทํ ภิกฺขุนิํ อารพฺภ. กิสฺมิํ วตฺถุสฺมินฺติ ถุลฺลนนฺทา ภิกฺขุนี “สํสฏฺฐาว อยฺเย ตุเมฺห วิหรถ, มา ตุเมฺห นานา วิหริตฺถา”ติ อุยฺโยเชสิ, ตสฺมิํ วตฺถุสฺมิํ. เอกา ปญฺญตฺติ. ฉนฺนํ อาปตฺติ\-สมุฏฺฐานานํ เอเกน สมุฏฺฐาเนน สมุฏฺฐาติ, ธุรนิกฺเขเป ฯเปฯ.}

\prose{ทธิํ วิญฺญาเปตฺวา ภุญฺชน\-ปจฺจยา ปาฏิเทสนียํ กตฺถ ปญฺญตฺตนฺติ สาวตฺถิยํ ปญฺญตฺตํ. กํ อารพฺภาติ ฉพฺพคฺคิยา ภิกฺขุนิโย อารพฺภ. กิสฺมิํ วตฺถุสฺมินฺติ ฉพฺพคฺคิยา ภิกฺขุนิโย ทธิํ วิญฺญาเปตฺวา ภุญฺชิํสุ, ตสฺมิํ วตฺถุสฺมิํ. เอกา ปญฺญตฺติ, เอกา อนุปญฺญตฺติ. ฉนฺนํ อาปตฺติ\-สมุฏฺฐานานํ จตูหิ สมุฏฺฐาเนหิ สมุฏฺฐาติ ฯเปฯ.}

\vspace{\tipitakaparskip}
\csromancenter{กตฺถ\-ปญฺญตฺติวาโร นิฏฺฐิโต ปฐโม.}
\csromanmatikaanchor{mtk.113}
\csromanmatikaanchor{mtk.114}
\csromantocmarknopage{chapter}{2. กตาปตฺติวาร}{mtk.113}
\bookmark[dest={mtk.113},level=0]{2. กตาปตฺติวาร}
\csromantocmarkref{chapter}{1. ปาราชิกกณฺฑ}{mtk.114}
\bookmark[dest={mtk.114},level=0]{1. ปาราชิกกณฺฑ}
\chapterheadpageread{2. \textbf{กตาปตฺติวาร 1. ปาราชิกกณฺฑ}}
\proseitem{249}{\csromanfolio{150}กายสํสคฺคํ สาทิยน\-ปจฺจยา กติ อาปตฺติโย อาปชฺชติ. กายสํสคฺคํ สาทิยน\-ปจฺจยา ปญฺจ อาปตฺติโย อาปชฺชติ. อวสฺสุตา ภิกฺขุนี อวสฺสุตสฺส ปุริส\-ปุคฺคลสฺส อธกฺขกํ อุพฺภ\-ชาณุมณฺฑลํ คหณํ สาทิยติ, อาปตฺติ ปาราชิกสฺส. ภิกฺขุ กาเยน กายํ อามสติ, อาปตฺติ สํฆาทิเสสสฺส. กาเยน กาย\-ปฏิพทฺธํ อามสติ, อาปตฺติ ถุลฺลจฺจยสฺส. กาย\-ปฏิพทฺเธน กาย\-ปฏิพทฺธํ อามสติ, อาปตฺติ ทุกฺกฏสฺส. องฺคุลิปโตทเก ปาจิตฺติยํ. กายสํสคฺคํ สาทิยน\-ปจฺจยา อิมา ปญฺจ อาปตฺติโย อาปชฺชติ.}

\prose{วชฺช\-ปฺปฏิจฺฉาทน\-ปจฺจยา กติ อาปตฺติโย อาปชฺชติ. วชฺช\-ปฺปฏิจฺฉาทน\-ปจฺจยา จตสฺโส อาปตฺติโย อาปชฺชติ. ภิกฺขุนี ชานํ ปาราชิกํ ธมฺมํ\csromansharedfootnote{ee41948bf668df0d}{ธมฺมํ อชฺฌาปนฺนํ (สฺยา)} ปฏิจฺฉาเทติ, อาปตฺติ ปาราชิกสฺส. เวมติกา ปฏิจฺฉาเทติ, อาปตฺติ ถุลฺลจฺจยสฺส. ภิกฺขุ สํฆาทิเสสํ ปฏิจฺฉาเทติ, อาปตฺติ ปาจิตฺติยสฺส. อาจารวิปตฺติํ ปฏิจฺฉาเทติ, อาปตฺติ ทุกฺกฏสฺส. วชฺช\-ปฺปฏิจฺฉาทน\-ปจฺจยา อิมา จตสฺโส อาปตฺติโย อาปชฺชติ.}

\prose{ยาวตติยํ สมนุภาสนาย น ปฏินิสฺสชฺชน\-ปจฺจยา กติ อาปตฺติโย อาปชฺชติ. ยาวตติยํ สมนุภาสนาย น ปฏินิสฺสชฺชน\-ปจฺจยา ปญฺจ อาปตฺติโย อาปชฺชติ. อุกฺขิตฺตานุวตฺติกา ภิกฺขุนี ยาวตติยํ สมนุภาสนาย น ปฏินิสฺสชฺชติ, ญตฺติยา ทุกฺกฏํ. ทฺวีหิ กมฺมวาจาหิ ถุลฺลจฺจยา. กมฺมวาจา\-ปริโยสาเน อาปตฺติ ปาราชิกสฺส. เภทกานุวตฺติกา ภิกฺขุนี ยาวตติยํ สมนุภาสนาย น ปฏินิสฺสชฺชติ, อาปตฺติ สํฆาทิเสสสฺส. ปาปิกาย ทิฏฺฐิยา ยาวตติยํ สมนุภาสนาย น ปฏินิสฺสชฺชติ, อาปตฺติ ปาจิตฺติยสฺส. ยาวตติยํ สมนุภาสนาย น ปฏินิสฺสชฺชน\-ปจฺจยา อิมา ปญฺจ อาปตฺติโย อาปชฺชติ.}

\prosewithcloserruled{อฏฺฐมํ วตฺถุํ ปริปูรณ\-ปจฺจยา กติ อาปตฺติโย อาปชฺชติ. อฏฺฐมํ วตฺถุํ ปริปูรณ\-ปจฺจยา ติสฺโส อาปตฺติโย อาปชฺชติ. ปุริเสน “อิตฺถนฺนามํ โอกาสํ อาคจฺฉา”ติ วุตฺตา คจฺฉติ, อาปตฺติ ทุกฺกฏสฺส. \csromanfolio{151}ปุริสสฺส หตฺถปาสํ โอกฺกนฺตมตฺเต อาปตฺติ ถุลฺลจฺจยสฺส. อฏฺฐมํ วตฺถุํ ปริปูเรติ, อาปตฺติ ปาราชิกสฺส. อฏฺฐมํ วตฺถุํ ปริปูรณ\-ปจฺจยา อิมา ติสฺโส อาปตฺติโย อาปชฺชติ.}{ปาราชิกา นิฏฺฐิตา.}
\csromanmatikaanchor{mtk.115}
\csromantocmarkref{chapter}{2. สํฆาทิเสสกณฺฑาทิ}{mtk.115}
\bookmark[dest={mtk.115},level=0]{2. สํฆาทิเสสกณฺฑาทิ}
\chapterhead{2. \textbf{สํฆาทิเสสกณฺฑาทิ}}
\proseitem{250}{อุสฺสยวาทิกา ภิกฺขุนี อฑฺฑํ กรณปจฺจยา ติสฺโส อาปตฺติโย อาปชฺชติ. เอกสฺส อาโรเจติ, อาปตฺติ ทุกฺกฏสฺส. ทุติยสฺส อาโรเจติ, อาปตฺติ ถุลฺลจฺจยสฺส. อฑฺฑปริโยสาเน อาปตฺติ สํฆาทิเสสสฺส.}

\prose{โจริํ วุฏฺฐาปน\-ปจฺจยา ติสฺโส อาปตฺติโย อาปชฺชติ. ญตฺติยา ทุกฺกฏํ. ทฺวีหิ กมฺมวาจาหิ ถุลฺลจฺจยา. กมฺมวาจา\-ปริโยสาเน อาปตฺติ สํฆาทิเสสสฺส.}

\prose{เอกา คามนฺตรํ คมนปจฺจยา ติสฺโส อาปตฺติโย อาปชฺชติ. คจฺฉติ, อาปตฺติ ทุกฺกฏสฺส. ปฐมํ ปาทํ ปริกฺเขปํ อติกฺกาเมติ, อาปตฺติ ถุลฺลจฺจยสฺส. ทุติยํ ปาทํ อติกฺกาเมติ, อาปตฺติ สํฆาทิเสสสฺส.}

\prose{สมคฺเคน สํเฆน อุกฺขิตฺตํ ภิกฺขุนิํ ธมฺเมน วินเยน สตฺถุสาสเนน อนปโลเกตฺวา การกสํฆํ อนญฺญาย คณสฺส ฉนฺทํ โอสารณปจฺจยา ติสฺโส อาปตฺติโย อาปชฺชติ. ญตฺติยา ทุกฺกฏํ. ทฺวีหิ กมฺมวาจาหิ ถุลฺลจฺจยา. กมฺมวาจา\-ปริโยสาเน อาปตฺติ สํฆาทิเสสสฺส.}

\prose{อวสฺสุตา ภิกฺขุนี อวสฺสุตสฺส ปุริส\-ปุคฺคลสฺส หตฺถโต ขาทนียํ วา โภชนียํ วา สหตฺถา ปฏิคฺคเหตฺวา ภุญฺชน\-ปจฺจยา ติสฺโส อาปตฺติโย อาปชฺชติ. “ขาทิสฺสามิ ภุญฺชิสฺสามี”ติ ปฏิคฺคณฺหาติ, อาปตฺติ ถุลฺลจฺจยสฺส. อชฺโฌหาเร อชฺโฌหาเร อาปตฺติ สํฆาทิเสสสฺส. อุทก\-ทนฺตโปนํ ปฏิคฺคณฺหาติ, อาปตฺติ ทุกฺกฏสฺส.}

\prose{“กิํ เต อยฺเย เอโส ปุริสปุคฺคโล กริสฺสติ อวสฺสุโต วา อนวสฺสุโต วา, ยโต ตฺวํ อนวสฺสุตา, อิงฺฆ อยฺเย, ยํ เต เอโส \csromanfolio{152}ปุริสปุคฺคโล เทติ ขาทนียํ วา โภชนียํ วา, ตํ ตฺวํ สหตฺถา ปฏิคฺคเหตฺวา ขาท วา ภุญฺช วา”ติ อุยฺโยชน\-ปจฺจยา ติสฺโส อาปตฺติโย อาปชฺชติ. ตสฺสา วจเนน “ขาทิสฺสามิ ภุญฺชิสฺสามี”ติ ปฏิคฺคณฺหาติ, อาปตฺติ ทุกฺกฏสฺส. อชฺโฌหาเร อชฺโฌหาเร อาปตฺติ ถุลฺลจฺจยสฺส. โภชน\-ปริโยสาเน อาปตฺติ สํฆาทิเสสสฺส.}

\prose{กุปิตา ภิกฺขุนี ยาวตติยํ สมนุภาสนาย น ปฏินิสฺสชฺชน\-ปจฺจยา ติสฺโส อาปตฺติโย อาปชฺชติ. ญตฺติยา ทุกฺกฏํ. ทฺวีหิ กมฺมวาจาหิ ถุลฺลจฺจยา. กมฺมวาจา\-ปริโยสาเน อาปตฺติ สํฆาทิเสสสฺส.}

\prose{กิสฺมิญฺจิเทว อธิกรเณ ปจฺจากตา ภิกฺขุนี ยาวตติยํ สมนุภาสนาย น ปฏินิสฺสชฺชน\-ปจฺจยา ติสฺโส อาปตฺติโย อาปชฺชติ. ญตฺติยา ทุกฺกฏํ. ทฺวีหิ กมฺมวาจาหิ ถุลฺลจฺจยา. กมฺมวาจา\-ปริโยสาเน อาปตฺติ สํฆาทิเสสสฺส.}

\prose{สํสฏฺฐา ภิกฺขุนี ยาวติยํ สมนุภาสนาย น ปฏินิสฺสชฺชน\-ปจฺจยา ติสฺโส อาปตฺติโย อาปชฺชติ. ญตฺติยา ทุกฺกฏํ. ทฺวีหิ กมฺมวาจาหิ ถุลฺลจฺจยา. กมฺมวาจา\-ปริโยสาเน อาปตฺติ สํฆาทิเสสสฺส.}

\prose{“สํสฏฺฐาว อยฺเย ตุเมฺห วิหรถ, มา ตุเมฺห นานา วิหริตฺถา”ติ อุยฺโยเชนฺตี ยาวตติยํ สมนุภาสนาย น ปฏินิสฺสชฺชน\-ปจฺจยา ติสฺโส อาปตฺติโย อาปชฺชติ. ญตฺติยา ทุกฺกฏํ. ทฺวีหิ กมฺมวาจาหิ ถุลฺลจฺจยา. กมฺมวาจา\-ปริโยสาเน อาปตฺติ สํฆาทิเสสสฺส.}

\vspace{\tipitakaparskip}
\csromancenter{ทส สํฆาทิเสโส นิฏฺฐิตา ฯเปฯ.}
\csromancenter{(ยถา เหฏฺฐา ตถา วิตฺถาเรตพฺพา, ปจฺจยเมว นานากรณํ)}
\prose{ทธิํ วิญฺญาเปตฺวา ภุญฺชน\-ปจฺจยา กติ อาปตฺติโย อาปชฺชติ. ทธิํ วิญฺญาเปตฺวา ภุญฺชน\-ปจฺจยา เทฺว อาปตฺติโย อาปชฺชติ. “ภุญฺชิสฺสามี”ติ ปฏิคฺคณฺหาติ, อาปตฺติ ทุกฺกฏสฺส. อชฺโฌหาเร อชฺโฌหาเร อาปตฺติ ปาฏิเทสนียสฺส. ทธิํ วิญฺญาเปตฺวา ภุญฺชน\-ปจฺจยา อิมา เทฺว อาปตฺติโย อาปชฺชติ.}

\vspace{\tipitakaparskip}
\csromancenter{กตาปตฺติวาโร นิฏฺฐิโต ทุติโย.}
\csromanmatikaanchor{mtk.116}
\csromantocmarkref{chapter}{3. วิปตฺติวาร}{mtk.116}
\bookmark[dest={mtk.116},level=0]{3. วิปตฺติวาร}
\chapterheadpageread{3. \textbf{วิปตฺติวาร}}
\proseitemwithcloserruled{251}{\csromanfolio{153}กายสํสคฺคํ สาทิยน\-ปจฺจยา อาปตฺติโย จตุนฺนํ วิปตฺตีนํ กติ วิปตฺติโย ภชนฺติ. กายสํสคฺคํ สาทิยน\-ปจฺจยา อาปตฺติโย จตุนฺนํ วิปตฺตีนํ เทฺว วิปตฺติโย ภชนฺติ, สิยา สีลวิปตฺติํ, สิยา อาจารวิปตฺติํ ฯเปฯ. ทธิํ วิญฺญาเปตฺวา ภุญฺชน\-ปจฺจยา อาปตฺติโย จตุนฺนํ วิปตฺตีนํ กติ วิปตฺติโย ภชนฺติ. ทธิํ วิญฺญาเปตฺวา ภุญฺชน\-ปจฺจยา อาปตฺติโย จตุนฺนํ วิปตฺตีนํ เอกํ วิปตฺติํ ภชนฺติ, อาจารวิปตฺติํ.}{วิปตฺติวาโร นิฏฺฐิโต ตติโย.}
\csromanmatikaanchor{mtk.117}
\csromantocmarkref{chapter}{4. สงฺคหวาร}{mtk.117}
\bookmark[dest={mtk.117},level=0]{4. สงฺคหวาร}
\chapterhead{4. \textbf{สงฺคหวาร}}
\proseitem{252}{กายสํสคฺคํ สาทิยน\-ปจฺจยา อาปตฺติโย สตฺตนฺนํ อาปตฺติ\-กฺขนฺธานํ กติหิ อาปตฺติ\-กฺขนฺเธหิ สงฺคหิตา. กายสํสคฺคํ สาทิยน\-ปจฺจยา อาปตฺติโย สตฺตนฺนํ อาปตฺติ\-กฺขนฺธานํ ปญฺจหิ อาปตฺติ\-กฺขนฺเธหิ สงฺคหิตา, สิยา ปาราชิกาปตฺติ\-กฺขนฺเธน, สิยา สํฆาทิเสสา\-ปตฺติกฺขนฺเธน, สิยา ถุลฺลจฺจยา\-ปตฺติกฺขนฺเธน, สิยา ปาจิตฺติยาปตฺติ\-กฺขนฺเธน, สิยา ทุกฺกฏาปตฺติ\-กฺขนฺเธน ฯเปฯ.}

\prosewithcloserruled{ทธิํ วิญฺญาเปตฺวา ภุญฺชน\-ปจฺจยา อาปตฺติโย สตฺตนฺนํ อาปตฺติ\-กฺขนฺธานํ กติหิ อาปตฺติ\-กฺขนฺเธหิ สงฺคหิตา. ทธิํ วิญฺญาเปตฺวา ภุญฺชน\-ปจฺจยา อาปตฺติโย สตฺตนฺนํ อาปตฺติ\-กฺขนฺธานํ ทฺวีหิ อาปตฺติ\-กฺขนฺเธหิ สงฺคหิตา, สิยา ปาฏิเทสนียา\-ปตฺติกฺขนฺเธน, สิยา ทุกฺกฏาปตฺติ\-กฺขนฺเธน.}{สงฺคหวาโร นิฏฺฐิโต จตุตฺโถ.}
\csromanmatikaanchor{mtk.118}
\csromantocmarkref{chapter}{5. สมุฏฺฐานวาร}{mtk.118}
\bookmark[dest={mtk.118},level=0]{5. สมุฏฺฐานวาร}
\chapterhead{5. \textbf{สมุฏฺฐานวาร}}
\proseitem{253}{กายสํสคฺคํ สาทิยน\-ปจฺจยา อาปตฺติโย ฉนฺนํ อาปตฺติ\-สมุฏฺฐานานํ กติหิ สมุฏฺฐาเนหิ สมุฏฺฐนฺติ. กายสํสคฺคํ สาทิยน\-ปจฺจยา อาปตฺติโย ฉนฺนํ อาปตฺติ\-สมุฏฺฐานานํ เอเกน สมุฏฺฐาเนน สมุฏฺฐนฺติ, กายโต จ จิตฺตโต จ สมุฏฺฐนฺติ, น วาจโต ฯเปฯ.}

\prosewithcloserruled{\csromanfolio{154}ทธิํ วิญฺญาเปตฺวา ภุญฺชน\-ปจฺจยา อาปตฺติโย ฉนฺนํ อาปตฺติ\-สมุฏฺฐานานํ กติหิ สมุฏฺฐาเนหิ สมุฏฺฐนฺติ. ทธิํ วิญฺญาเปตฺวา ภุญฺชน\-ปจฺจยา อาปตฺติโย ฉนฺนํ อาปตฺติ\-สมุฏฺฐานานํ จตูหิ สมุฏฺฐาเนหิ สมุฏฺฐนฺติ, สิยา กายโต สมุฏฺฐนฺติ, น วาจโต, น จิตฺตโต, สิยา กายโต จ วาจโต จ สมุฏฺฐนฺติ, น จิตฺตโต, สิยา กายโต จ จิตฺตโต จ สมุฏฺฐนฺติ, น วาจโต, สิยา กายโต จ วาจโต จ จิตฺตโต จ สมุฏฺฐนฺติ.}{สมุฏฺฐานวาโร นิฏฺฐิโต ปญฺจโม.}
\csromanmatikaanchor{mtk.119}
\csromantocmarkref{chapter}{6. อธิกรณวาร}{mtk.119}
\bookmark[dest={mtk.119},level=0]{6. อธิกรณวาร}
\chapterhead{6. \textbf{อธิกรณวาร}}
\proseitem{254}{กายสํสคฺคํ สาทิยน\-ปจฺจยา อาปตฺติโย จตุนฺนํ อธิกรณานํ กตมํ อธิกรณํ. กายสํสคฺคํ สาทิยน\-ปจฺจยา อาปตฺติโย จตุนฺนํ อธิกรณานํ อาปตฺตา\-ธิกรณํ ฯเปฯ.}

\prosewithcloserruled{ทธิํ วิญฺญาเปตฺวา ภุญฺชน\-ปจฺจยา อาปตฺติโย จตุนฺนํ อธิกรณานํ กตมํ อธิกรณํ. ทธิํ วิญฺญาเปตฺวา ภุญฺชน\-ปจฺจยา อาปตฺติโย จตุนฺนํ อธิกรณานํ อาปตฺตา\-ธิกรณํ.}{อธิกรณวาโร นิฏฺฐิโต ฉฏฺโฐ.}
\csromanmatikaanchor{mtk.120}
\csromantocmarkref{chapter}{7. สมถวาร}{mtk.120}
\bookmark[dest={mtk.120},level=0]{7. สมถวาร}
\chapterhead{7. \textbf{สมถวาร}}
\proseitem{255}{กายสํสคฺคํ สาทิยน\-ปจฺจยา อาปตฺติโย สตฺตนฺนํ สมถานํ กติหิ สมเถหิ สมฺมนฺติ. กายสํสคฺคํ สาทิยน\-ปจฺจยา อาปตฺติโย สตฺตนฺนํ สมถานํ ตีหิ สมเถหิ สมฺมนฺติ, สิยา สมฺมุขา\-วินเยน จ ปฏิญฺญาต\-กรเณน จ, สิยา สมฺมุขา\-วินเยน จ ติณวตฺถารเกน จ ฯเปฯ.}

\prose{ทธิํ วิญฺญาเปตฺวา ภุญฺชน\-ปจฺจยา อาปตฺติโย สตฺตนฺนํ สมถานํ กติหิ สมเถหิ สมฺมนฺติ. ทธิํ วิญฺญาเปตฺวา ภุญฺชน\-ปจฺจยา อาปตฺติโย สตฺตนฺนํ สมถานํ ตีหิ สมเถหิ สมฺมนฺติ, สิยา สมฺมุขา\-วินเยน จ ปฏิญฺญาต\-กรเณน จ, สิยา สมฺมุขา\-วินเยน จ ติณวตฺถารเกน จ.}

\vspace{\tipitakaparskip}
\csromancenter{สมถวาโร นิฏฺฐิโต สตฺตโม.}
\csromanmatikaanchor{mtk.121}
\csromantocmarkref{chapter}{8. สมุจฺจยวาร}{mtk.121}
\bookmark[dest={mtk.121},level=0]{8. สมุจฺจยวาร}
\chapterheadpageread{8. \textbf{สมุจฺจยวาร}}
\proseitem{256}{\csromanfolio{155}กายสํสคฺคํ สาทิยน\-ปจฺจยา กติ อาปตฺติโย อาปชฺชติ. กายสํสคฺคํ สาทิยน\-ปจฺจยา ปญฺจ อาปตฺติโย อาปชฺชติ. อวสฺสุตา ภิกฺขุนี อวสฺสุตสฺส ปุริส\-ปุคฺคลสฺส อธกฺขกํ อุพฺภ\-ชาณุมณฺฑลํ คหณํ สาทิยติ, อาปตฺติ ปาราชิกสฺส. ภิกฺขุ กาเยน กายํ อามสติ, อาปตฺติ สํฆาทิเสสสฺส. กาเยน กาย\-ปฏิพทฺธํ อามสติ, อาปตฺติ ถุลฺลจฺจยสฺส. กาย\-ปฏิพทฺเธน กาย\-ปฏิพทฺธํ อามสติ, อาปตฺติ ทุกฺกฏสฺส. องฺคุลิปโตทเก ปาจิตฺติยํ. กายสํสคฺคํ สาทิยน\-ปจฺจยา อิมา ปญฺจ อาปตฺติโย อาปชฺชติ.}

\prose{ตา อาปตฺติโย จตุนฺนํ วิปตฺตีนํ กติ วิปตฺติโย ภชนฺติ, สตฺตนฺนํ อาปตฺติ\-กฺขนฺธานํ กติหิ อาปตฺติ\-กฺขนฺเธหิ สงฺคหิตา, ฉนฺนํ อาปตฺติ\-สมุฏฺฐานานํ กติหิ สมุฏฺฐาเนหิ สมุฏฺฐนฺติ, จตุนฺนํ อธิกรณานํ กตมํ อธิกรณํ, สตฺตนฺนํ สมถานํ กติหิ สมเถหิ สมฺมนฺติ. ตา อาปตฺติโย จตุนฺนํ วิปตฺตีนํ เทฺว วิปตฺติโย ภชนฺติ, สิยา สีลวิปตฺติํ, สิยา อาจารวิปตฺติํ. สตฺตนฺนํ อาปตฺติ\-กฺขนฺธานํ ปญฺจหิ อาปตฺติ\-กฺขนฺเธหิ สงฺคหิตา, สิยา ปาราชิกาปตฺติ\-กฺขนฺเธน, สิยา สํฆาทิเสสา\-ปตฺติกฺขนฺเธน, สิยา ถุลฺลจฺจยา\-ปตฺติกฺขนฺเธน, สิยา ปาจิตฺติยาปตฺติ\-กฺขนฺเธน, สิยา ทุกฺกฏาปตฺติ\-กฺขนฺเธน. ฉนฺนํ อาปตฺติ\-สมุฏฺฐานานํ เอเกน สมุฏฺฐาเนน สมุฏฺฐนฺติ, กายโต จ จิตฺตโต จ สมุฏฺฐนฺติ, น วาจโต. จตุนฺนํ อธิกรณานํ อาปตฺตา\-ธิกรณํ. สตฺตนฺนํ สมถานํ ตีหิ สมเถหิ สมฺมนฺติ, สิยา สมฺมุขา\-วินเยน จ ปฏิญฺญาต\-กรเณน จ, สิยา สมฺมุขา\-วินเยน จ ติณวตฺถารเกน จ ฯเปฯ.}

\prose{ทธิํ วิญฺญาเปตฺวา ภุญฺชน\-ปจฺจยา กติ อาปตฺติโย อาปชฺชติ. ทธิํ วิญฺญาเปตฺวา ภุญฺชน\-ปจฺจยา เทฺว อาปตฺติโย อาปชฺชติ. ภุญฺชิสฺสามีติ ปฏิคฺคณฺหาติ, อาปตฺติ ทุกฺกฏสฺส. อชฺโฌหาเร อชฺโฌหาเร อาปตฺติ ปาฏิเทสนียสฺส. ทธิํ วิญฺญาเปตฺวา ภุญฺชน\-ปจฺจยา อิมา เทฺว อาปตฺติโย อาปชฺชติ.}

\prose{ตา อาปตฺติโย จตุนฺนํ วิปตฺตีนํ กติ วิปตฺติโย ภชนฺติ, สตฺตนฺนํ อาปตฺติ\-กฺขนฺธานํ กติหิ อาปตฺติ\-กฺขนฺเธหิ สงฺคหิตา, ฉนฺนํ อาปตฺติ\-สมุฏฺฐานานํ กติหิ สมุฏฺฐาเนหิ สมุฏฺฐนฺติ, จตุนฺนํ อธิกรณานํ กตมํ \csromanfolio{156}อธิกรณํ, สตฺตนฺนํ สมถานํ กติหิ สมเถหิ สมฺมนฺติ. ตา อาปตฺติโย จตุนฺนํ วิปตฺตีนํ เอกํ วิปตฺติํ ภชนฺติ, อาจารวิปตฺติํ. สตฺตนฺนํ อาปตฺติ\-กฺขนฺธานํ ทฺวีหิ อาปตฺติ\-กฺขนฺเธหิ สงฺคหิตา, สิยา ปาฏิเทสนียา\-ปตฺติกฺขนฺเธน, สิยา ทุกฺกฏาปตฺติ\-กฺขนฺเธน. ฉนฺนํ อาปตฺติ\-สมุฏฺฐานานํ จตูหิ สมุฏฺฐาเนหิ สมุฏฺฐนฺติ, สิยา กายโต สมุฏฺฐนฺติ, น วาจโต, น จิตฺตโต, สิยา กายโต จ วาจโต จ สมุฏฺฐนฺติ น จิตฺตโต, สิยา กายโต จ จิตฺตโต จ สมุฏฺฐนฺติ, น วาจโต, สิยา กายโต จ วาจโต จ จิตฺตโต จ สมุฏฺฐนฺติ. จตุนฺนํ อธิกรณานํ อาปตฺตา\-ธิกรณํ. สตฺตนฺนํ สมถานํ ตีหิ สมเถหิ สมฺมนฺติ, สิยา สมฺมุขา\-วินเยน จ ปฏิญฺญาต\-กรเณน จ, สิยา สมฺมุขา\-วินเยน จ ติณวตฺถารเกน จาติ.}

\vspace{\tipitakaparskip}
\csromancenter{สมุจฺจยวาโร นิฏฺฐิโต อฏฺฐโม\csromansharedfootnote{3acc069bba1d730d}{อิมสฺสานนฺตรํ โปราณสีหฬโปตฺถเก เอวมฺปิ ทิสฺสติ\csromandash{} “เตสมุทฺทานํ\csromandash{} กติวาโร จ วิปตฺติวาโร จ, สงฺคหวาโร จ สมุฏฺฐานวาโร. อธิกรณวาโร จ สมถวาโร จ, สมุจฺจยวาโร จ อิเม สตฺต วารา. อาทิโต ปญฺญตฺติวาเรน สห อฏฺฐวาราติ” อิติ.}.}
\nitthitam{อฏฺฐ ปจฺจยวารา นิฏฺฐิตา.}
\nitthitam{ภิกฺขุนิ\-วิภงฺเค โสฬส มหาวารา นิฏฺฐิตา.}
\csromanmatikaanchor{mtk.122}
\csromantocmarknopage{section}{สมุฏฺฐานสีสสงฺเขป}{mtk.122}
\bookmark[dest={mtk.122},level=1]{สมุฏฺฐานสีสสงฺเขป}
\csromanheader{1}{สมุฏฺฐาน\-สีส\-สงฺเขป}
\csromanheader{2}{สมุฏฺฐานสฺสุ\-ทฺทานํ}
\csromanmatikaanchor{mtk.123}
\csromantocmarkref{section}{อุทฺทานคาถา}{mtk.123}
\bookmark[dest={mtk.123},level=1]{อุทฺทานคาถา}
\proseitem{257}{\csromanfolio{157}อนิจฺจา สพฺเพ สงฺขารา, ทุกฺขานตฺตา จ สงฺขตา.}

\setlength{\csromangathapagewidth}{0pt}
\setlength{\csromangathawakwidth}{0pt}
\csromangathameasuregroup{นิพฺพานญฺเจว ปญฺญตฺติ, \\ พุทฺธจนฺเท อนุปฺปนฺเน, \\ เตสํ สภาคธมฺมานํ, \\ ทุกฺกรํ วิวิธํ กตฺวา, \\ อุปฺปชฺชนฺติ มหาวีรา, \\ เต เทสยนฺติ สทฺธมฺมํ, \\ องฺคีรโส สกฺยมุนิ, \\ สพฺพสตฺตุตฺตโม สีโห, \\ สุตฺตนฺตมภิธมฺมญฺจ, \\ เอวํ นียติ สทฺธมฺโม,}{\csromangathabat{นิพฺพานญฺเจว ปญฺญตฺติ,}{อนตฺตา อิติ นิจฺฉยา.} \\ \csromangathabat{พุทฺธจนฺเท อนุปฺปนฺเน,}{พุทฺธาทิจฺเจ อนุคฺคเต.} \\ \csromangathabat{เตสํ สภาคธมฺมานํ,}{นามมตฺตํ น นายติ.} \\ \csromangathabat{ทุกฺกรํ วิวิธํ กตฺวา,}{ปูรยิตฺวาน ปารมี.} \\ \csromangathabat{อุปฺปชฺชนฺติ มหาวีรา,}{จกฺขุภูตา สพฺรหฺมเก.} \\ \csromangathabat{เต เทสยนฺติ สทฺธมฺมํ,}{ทุกฺขหานิํ สุขาวหํ.} \\ \csromangathabat{องฺคีรโส สกฺยมุนิ,}{สพฺพภูตานุกมฺปโก.} \\ \csromangathabat{สพฺพสตฺตุตฺตโม สีโห,}{ปิฏเก ตีณิ เทสยิ.} \\ \csromangathabat{สุตฺตนฺตมภิธมฺมญฺจ,}{วินยญฺจ มหาคุณํ.} \\ \csromangathabat{เอวํ นียติ สทฺธมฺโม,}{วินโย ยทิ ติฏฺฐติ.}}
\csromangathasetleft{นิพฺพานญฺเจว ปญฺญตฺติ, \\ พุทฺธจนฺเท อนุปฺปนฺเน, \\ เตสํ สภาคธมฺมานํ, \\ ทุกฺกรํ วิวิธํ กตฺวา, \\ อุปฺปชฺชนฺติ มหาวีรา, \\ เต เทสยนฺติ สทฺธมฺมํ, \\ องฺคีรโส สกฺยมุนิ, \\ สพฺพสตฺตุตฺตโม สีโห, \\ สุตฺตนฺตมภิธมฺมญฺจ, \\ เอวํ นียติ สทฺธมฺโม,}
\csromangathagroup{\csromangathastanza{\csromangathabat{นิพฺพานญฺเจว ปญฺญตฺติ,}{อนตฺตา อิติ นิจฺฉยา.} \\ \csromangathabat{พุทฺธจนฺเท อนุปฺปนฺเน,}{พุทฺธาทิจฺเจ อนุคฺคเต.}}\csromangathastanza{\csromangathabat{เตสํ สภาคธมฺมานํ,}{นามมตฺตํ น นายติ.} \\ \csromangathabat{ทุกฺกรํ วิวิธํ กตฺวา,}{ปูรยิตฺวาน ปารมี.}}\csromangathastanza{\csromangathabat{อุปฺปชฺชนฺติ มหาวีรา,}{จกฺขุภูตา สพฺรหฺมเก.} \\ \csromangathabat{เต เทสยนฺติ สทฺธมฺมํ,}{ทุกฺขหานิํ สุขาวหํ.}}\csromangathastanza{\csromangathabat{องฺคีรโส สกฺยมุนิ,}{สพฺพภูตา\-นุกมฺปโก.} \\ \csromangathabat{สพฺพสตฺตุ\-ตฺตโม สีโห,}{ปิฏเก ตีณิ เทสยิ.}}\csromangathastanza{\csromangathabat{สุตฺตนฺตมภิธมฺมญฺจ,}{วินยญฺจ มหาคุณํ.} \\ \csromangathabat{เอวํ นียติ สทฺธมฺโม,}{วินโย ยทิ ติฏฺฐติ.}}}

\csromanheader{2}{อุภโต จ วิภงฺคานิ, ขนฺธกา ยา จ มาติกา.}
\setlength{\csromangathapagewidth}{0pt}
\setlength{\csromangathawakwidth}{0pt}
\csromangathameasuregroup{มาลา สุตฺตคุเณเนว, \\ ตสฺเสว ปริวารสฺส, \\ สมฺเภทํ นิทานญฺจญฺญํ, \\ ตสฺมา สิกฺเข ปริวารํ,}{\csromangathabat{มาลา สุตฺตคุเณเนว,}{ปริวาเรน คนฺถิตา.} \\ \csromangathabat{ตสฺเสว ปริวารสฺส,}{สมุฏฺฐานํ นิยโต กตํ.} \\ \csromangathabat{สมฺเภทํ นิทานญฺจญฺญํ,}{สุตฺเต ทิสฺสนฺติ อุปริ.} \\ \csromangathabat{ตสฺมา สิกฺเข ปริวารํ,}{ธมฺมกาโม สุเปสโลติ.}}
\csromangathasetleft{มาลา สุตฺตคุเณเนว, \\ ตสฺเสว ปริวารสฺส, \\ สมฺเภทํ นิทานญฺจญฺญํ, \\ ตสฺมา สิกฺเข ปริวารํ,}
\csromangathagroup{\csromangathastanza{\csromangathabat{มาลา สุตฺตคุเณเนว,}{ปริวาเรน คนฺถิตา.} \\ \csromangathabat{ตสฺเสว ปริวารสฺส,}{สมุฏฺฐานํ นิยโต กตํ.}}\csromangathastanza{\csromangathabat{สมฺเภทํ นิทานญฺจญฺญํ,}{สุตฺเต ทิสฺสนฺติ อุปริ.} \\ \csromangathabat{ตสฺมา สิกฺเข ปริวารํ,}{ธมฺมกาโม สุเปสโลติ.}}}

\csromanmatikaanchor{mtk.124}
\csromantocmarkref{section}{เตรสสมุฏฺฐาน}{mtk.124}
\bookmark[dest={mtk.124},level=1]{เตรสสมุฏฺฐาน}
\csromanheader{1}{เตรส\-สมุฏฺฐาน}
\setlength{\csromangathapagewidth}{0pt}
\setlength{\csromangathawakwidth}{0pt}
\csromangathameasuregroup{วิภงฺเค ทฺวีสุ ปญฺญตฺตํ, \\ ปวกฺขามิ สมุฏฺฐานํ, \\ ปาราชิกํ ยํ ปฐมํ, \\ สญฺจริตฺตานุภาสนญฺจ, \\ โลมานิ ปทโสธมฺโม, \\ เถยฺยเทสนโจรี จ, \\ เตรเสเต สมุฏฺฐาน, \\ เอเกกสฺมิํ สมุฏฺฐาเน,}{\csromangathabat{วิภงฺเค ทฺวีสุ ปญฺญตฺตํ,}{อุทฺทิสนฺติ อุโปสเถ.} \\ \csromangathabat{ปวกฺขามิ สมุฏฺฐานํ,}{ยถาญายํ สุณาถ เม.} \\ \csromangathabat{ปาราชิกํ ยํ ปฐมํ,}{ทุติยญฺจ ตโต ปรํ.} \\ \csromangathabat{สญฺจริตฺตานุภาสนญฺจ,}{อติเรกญฺจ จีวรํ.} \\ \csromangathabat{โลมานิ ปทโสธมฺโม,}{ภูตํ สํวิธาเนน จ.} \\ \csromangathabat{เถยฺยเทสนโจรี จ,}{อนนุญฺญาตาย เตรส.} \\ \csromangathabat{เตรเสเต สมุฏฺฐาน,}{นยา วิญฺญูหิ จินฺติตา.} \\ \csromangathabat{เอเกกสฺมิํ สมุฏฺฐาเน,}{สทิสา อิธ ทิสฺสเร.}}
\csromangathasetleft{วิภงฺเค ทฺวีสุ ปญฺญตฺตํ, \\ ปวกฺขามิ สมุฏฺฐานํ, \\ ปาราชิกํ ยํ ปฐมํ, \\ สญฺจริตฺตานุภาสนญฺจ, \\ โลมานิ ปทโสธมฺโม, \\ เถยฺยเทสนโจรี จ, \\ เตรเสเต สมุฏฺฐาน, \\ เอเกกสฺมิํ สมุฏฺฐาเน,}
\csromangathagroup{\csromangathastanza{\csromangathabat{วิภงฺเค ทฺวีสุ ปญฺญตฺตํ,}{อุทฺทิสนฺติ อุโปสเถ.} \\ \csromangathabat{ปวกฺขามิ สมุฏฺฐานํ,}{ยถาญายํ สุณาถ เม.}}\csromangathastanza{\csromangathabat{ปาราชิกํ ยํ ปฐมํ,}{ทุติยญฺจ ตโต ปรํ.} \\ \csromangathabat{สญฺจริตฺตานุภาสนญฺจ,}{อติเรกญฺจ จีวรํ.}}\csromangathastanza{\csromangathabat{โลมานิ ปทโสธมฺโม,}{ภูตํ สํวิธาเนน จ.} \\ \csromangathabat{เถยฺยเทสนโจรี จ,}{อนนุญฺญาตาย เตรส.}}\csromangathastanza{\csromangathabat{เตรเสเต สมุฏฺฐาน,}{นยา วิญฺญูหิ จินฺติตา.} \\ \csromangathabat{เอเกกสฺมิํ สมุฏฺฐาเน,}{สทิสา อิธ ทิสฺสเร.}}}

\csromanmatikaanchor{mtk.125}
\csromantocmarkref{section}{1. ปฐมปาราชิกสมุฏฺฐาน}{mtk.125}
\bookmark[dest={mtk.125},level=1]{1. ปฐมปาราชิกสมุฏฺฐาน}
\csromanheader{1}{1. ปฐมปาราชิก\-สมุฏฺฐาน}
\setlength{\csromangathapagewidth}{0pt}
\setlength{\csromangathawakwidth}{0pt}
\csromangathameasuregroup{เมถุนํ สุกฺกสํสคฺโค, \\ ปุพฺพูปปริปาจิตา, \\ สโภชเน รโห เทฺว จ, \\ ปหาเร อุคฺคิเร เจว, \\ อธกฺขคามาวสฺสุตา, \\ วสฺสํวุฏฺฐา จ โอวาทํ, \\ ฉสตฺตติ อิเม สิกฺขา, \\ สพฺเพ เอกสมุฏฺฐานา,}{\csromangathabat{เมถุนํ สุกฺกสํสคฺโค,}{อนิยตา ปฐมิกา.} \\ \csromangathabat{ปุพฺพูปปริปาจิตา,}{รโห ภิกฺขุนิยา สห.} \\ \csromangathabat{สโภชเน รโห เทฺว จ,}{องฺคุลิ อุทเก หสํ.} \\ \csromangathabat{ปหาเร อุคฺคิเร เจว,}{เตปญฺญาสา จ เสขิยา.} \\ \csromangathabat{อธกฺขคามาวสฺสุตา,}{ตลมฏฺฐญฺจ สุทฺธิกา.} \\ \csromangathabat{วสฺสํวุฏฺฐา จ โอวาทํ,}{นานุพนฺเธ ปวตฺตินิํ.} \\ \csromangathabat{ฉสตฺตติ อิเม สิกฺขา,}{กายมานสิกา กตา.} \\ \csromangathabat{สพฺเพ เอกสมุฏฺฐานา,}{ปฐมํ ปาราชิกํ ยถา.}}
\csromangathasetleft{เมถุนํ สุกฺกสํสคฺโค, \\ ปุพฺพูปปริปาจิตา, \\ สโภชเน รโห เทฺว จ, \\ ปหาเร อุคฺคิเร เจว, \\ อธกฺขคามาวสฺสุตา, \\ วสฺสํวุฏฺฐา จ โอวาทํ, \\ ฉสตฺตติ อิเม สิกฺขา, \\ สพฺเพ เอกสมุฏฺฐานา,}
\csromangathagroupwithcloserruled{\csromangathastanza{\csromanfolio{158}\csromangathaitembat{258}{เมถุนํ สุกฺกสํสคฺโค,}{อนิยตา ปฐมิกา.} \\ \csromangathacontbat{ปุพฺพูปปริปาจิตา,}{รโห ภิกฺขุนิยา สห.} \\ \csromangathacontbat{สโภชเน รโห เทฺว จ,}{องฺคุลิ อุทเก หสํ.} \\ \csromangathacontbat{ปหาเร อุคฺคิเร เจว,}{เตปญฺญาสา จ เสขิยา.} \\ \csromangathacontbat{อธกฺขคามาวสฺสุตา,}{ตลมฏฺฐญฺจ สุทฺธิกา.} \\ \csromangathacontbat{วสฺสํวุฏฺฐา จ โอวาทํ,}{นานุพนฺเธ ปวตฺตินิํ.} \\ \csromangathacontbat{ฉสตฺตติ อิเม สิกฺขา,}{กายมานสิกา กตา.} \\ \csromangathacontbat{สพฺเพ เอกสมุฏฺฐานา,}{ปฐมํ ปาราชิกํ ยถา.}}}{ปฐมปาราชิก\-สมุฏฺฐานํ นิฏฺฐิตํ.}
\csromanmatikaanchor{mtk.126}
\csromantocmarkref{section}{2. ทุติยปาราชิกสมุฏฺฐาน}{mtk.126}
\bookmark[dest={mtk.126},level=1]{2. ทุติยปาราชิกสมุฏฺฐาน}
\csromanheader{1}{2. ทุติยปาราชิก\-สมุฏฺฐาน}
\setlength{\csromangathapagewidth}{0pt}
\setlength{\csromangathawakwidth}{0pt}
\csromangathameasuregroup{อทินฺนํ วิคฺคหุตฺตริ, \\ อมูลา อญฺญภาคิยา, \\ อจฺฉินฺเท ปริณามเน, \\ ทุฏฺฐุลฺลา ปถวีขเณ, \\ นิกฺกฑฺฒนํ สิญฺจนญฺจ, \\ เอหิ อนาทริ ภิํสา, \\ ชานํ สปฺปาณกํ กมฺมํ, \\ สหธมฺมิกวิเลขา, \\ กุกฺกุจฺจํ ธมฺมิกํ จีวรํ ทตฺวา\textsuperscript{1}, \\ กิํ เต อกาลํ อจฺฉินฺเท, \\ คณํ วิภงฺคํ ทุพฺพลํ, \\ อกฺโกสจณฺฑี มจฺฉรี, \\ เทฺววสฺสํ สิกฺขา สํเฆน, \\ กุมาริภูตา ติสฺโส จ, \\ อลํ ตาว โสกาวาสํ, \\ สิกฺขาปทา สตฺตติเม, \\ กายจิตฺเตน น วาจา, \\ ตีหิ ทฺวาเรหิ ชายนฺติ,}{\csromangathabat{อทินฺนํ วิคฺคหุตฺตริ,}{ทุฏฺฐุลฺลา อตฺตกามินํ.} \\ \csromangathabat{อมูลา อญฺญภาคิยา,}{อนิยตา ทุติยิกา.} \\ \csromangathabat{อจฺฉินฺเท ปริณามเน,}{มุสา โอมสเปสุณา.} \\ \csromangathabat{ทุฏฺฐุลฺลา ปถวีขเณ,}{ภูตํ อญฺญาย อุชฺฌาเป.} \\ \csromangathabat{นิกฺกฑฺฒนํ สิญฺจนญฺจ,}{อามิสเหตุ ภุตฺตาวี.} \\ \csromangathabat{เอหิ อนาทริ ภิํสา,}{อปนิเธ จ ชีวิตํ.} \\ \csromangathabat{ชานํ สปฺปาณกํ กมฺมํ,}{อูนสํวาสนาสนา.} \\ \csromangathabat{สหธมฺมิกวิเลขา,}{โมโห อมูลเกน จ.} \\ \csromangathabat{กุกฺกุจฺจํ ธมฺมิกํ จีวรํ ทตฺวา\textsuperscript{1},}{ปริณาเมยฺย ปุคฺคเล.} \\ \csromangathabat{กิํ เต อกาลํ อจฺฉินฺเท,}{ทุคฺคหี นิรเยน จ.} \\ \csromangathabat{คณํ วิภงฺคํ ทุพฺพลํ,}{กถินาผาสุปสฺสยํ.} \\ \csromangathabat{อกฺโกสจณฺฑี มจฺฉรี,}{คพฺภินี จ ปายนฺติยา.} \\ \csromangathabat{เทฺววสฺสํ สิกฺขา สํเฆน,}{ตโย เจว คิหีคตา.} \\ \csromangathabat{กุมาริภูตา ติสฺโส จ,}{อูนทฺวาทสสมฺมตา.} \\ \csromangathabat{อลํ ตาว โสกาวาสํ,}{ฉนฺทา อนุวสฺสา จ เทฺว.} \\ \csromangathabat{สิกฺขาปทา สตฺตติเม,}{สมุฏฺฐานา ติกา กตา.} \\ \csromangathabat{กายจิตฺเตน น วาจา,}{วาจาจิตฺตํ น กายิกํ.} \\ \csromangathabat{ตีหิ ทฺวาเรหิ ชายนฺติ,}{ปาราชิกํ ทุติยํ ยถา.}}
\csromangathasetleft{อทินฺนํ วิคฺคหุตฺตริ, \\ อมูลา อญฺญภาคิยา, \\ อจฺฉินฺเท ปริณามเน, \\ ทุฏฺฐุลฺลา ปถวีขเณ, \\ นิกฺกฑฺฒนํ สิญฺจนญฺจ, \\ เอหิ อนาทริ ภิํสา, \\ ชานํ สปฺปาณกํ กมฺมํ, \\ สหธมฺมิกวิเลขา, \\ กุกฺกุจฺจํ ธมฺมิกํ จีวรํ ทตฺวา\textsuperscript{1}, \\ กิํ เต อกาลํ อจฺฉินฺเท, \\ คณํ วิภงฺคํ ทุพฺพลํ, \\ อกฺโกสจณฺฑี มจฺฉรี, \\ เทฺววสฺสํ สิกฺขา สํเฆน, \\ กุมาริภูตา ติสฺโส จ, \\ อลํ ตาว โสกาวาสํ, \\ สิกฺขาปทา สตฺตติเม, \\ กายจิตฺเตน น วาจา, \\ ตีหิ ทฺวาเรหิ ชายนฺติ,}
\csromangathagroupwithcloserruled{\csromangathastanza{\csromangathaitembat{259}{อทินฺนํ วิคฺคหุตฺตริ,}{ทุฏฺฐุลฺลา อตฺตกามินํ.} \\ \csromangathacontbat{อมูลา อญฺญภาคิยา,}{อนิยตา ทุติยิกา.} \\ \csromangathacontbat{อจฺฉินฺเท ปริณามเน,}{มุสา โอมสเปสุณา.} \\ \csromangathacontbat{ทุฏฺฐุลฺลา ปถวีขเณ,}{ภูตํ อญฺญาย อุชฺฌาเป.} \\ \csromangathacontbat{นิกฺกฑฺฒนํ สิญฺจนญฺจ,}{อามิสเหตุ ภุตฺตาวี.} \\ \csromangathacontbat{เอหิ อนาทริ ภิํสา,}{อปนิเธ จ ชีวิตํ.} \\ \csromangathacontbat{ชานํ สปฺปาณกํ กมฺมํ,}{อูนสํวาสนาสนา.} \\ \csromangathacontbat{สหธมฺมิกวิเลขา,}{โมโห อมูลเกน จ.} \\ \csromangathacontbat{กุกฺกุจฺจํ ธมฺมิกํ จีวรํ ทตฺวา\csromansharedfootnote{43673e94ac800aa1}{กุกฺกุจฺจํ ธมฺมิกํ จีวรํ (สี), กุกฺกุจฺจํ ธมฺมิกํ ทตฺวา (สฺยา)},}{ปริณาเมยฺย ปุคฺคเล.} \\ \csromangathacontbat{กิํ เต อกาลํ อจฺฉินฺเท,}{ทุคฺคหี นิรเยน จ.} \\ \csromangathacontbat{คณํ วิภงฺคํ ทุพฺพลํ,}{กถินาผาสุปสฺสยํ.} \\ \csromangathacontbat{อกฺโกสจณฺฑี มจฺฉรี,}{คพฺภินี จ ปายนฺติยา.}}\csromangathastanza{\csromanfolio{159}\csromangathabat{เทฺววสฺสํ สิกฺขา สํเฆน,}{ตโย เจว คิหีคตา.} \\ \csromangathabat{กุมาริภูตา ติสฺโส จ,}{อูนทฺวาทสสมฺมตา.}}\csromangathastanza{\csromangathabat{อลํ ตาว โสกาวาสํ,}{ฉนฺทา อนุวสฺสา จ เทฺว.} \\ \csromangathabat{สิกฺขาปทา สตฺตติเม,}{สมุฏฺฐานา ติกา กตา.}}\csromangathastanza{\csromangathabat{กายจิตฺเตน น วาจา,}{วาจาจิตฺตํ น กายิกํ.} \\ \csromangathabat{ตีหิ ทฺวาเรหิ ชายนฺติ,}{ปาราชิกํ ทุติยํ ยถา.}}}{ทุติยปาราชิก\-สมุฏฺฐานํ นิฏฺฐิตํ.}
\csromanmatikaanchor{mtk.127}
\csromantocmarkref{section}{3. สญฺจริตฺตสมุฏฺฐาน}{mtk.127}
\bookmark[dest={mtk.127},level=1]{3. สญฺจริตฺตสมุฏฺฐาน}
\csromanheader{1}{3. สญฺจริตฺต\-สมุฏฺฐาน}
\setlength{\csromangathapagewidth}{0pt}
\setlength{\csromangathawakwidth}{0pt}
\csromangathameasuregroup{สญฺจรี กุฏิ วิหาโร, \\ วิญฺญตฺตุตฺตริ อภิหฏฺฐุํ, \\ โกสิยา สุทฺธเทฺวภาคา, \\ ริญฺจนฺติ รูปิกา เจว, \\ อูนพนฺธนวสฺสิกา, \\ ทฺวารทานสิพฺพานิ\textsuperscript{1} จ, \\ รตนํ สูจิ มญฺโจ จ, \\ วสฺสิกา จ สุคเตน, \\ เทฺว สํฆิกา มหาชนิกา เทฺว, \\ เทฺว วิฆาสา สาฏิกา จ, \\ สมปญฺญาสิเม ธมฺมา, \\ กายโต น วาจาจิตฺตา, \\ กายวาจา น จ จิตฺตา\textsuperscript{1}, \\ วาจาจิตฺตา น กาเยน, \\ ฉสมุฏฺฐานิกา เอเต,}{\csromangathabat{สญฺจรี กุฏิ วิหาโร,}{โธวนญฺจ ปฏิคฺคโห.} \\ \csromangathabat{วิญฺญตฺตุตฺตริ อภิหฏฺฐุํ,}{อุภินฺนํ ทูตเกน จ.} \\ \csromangathabat{โกสิยา สุทฺธเทฺวภาคา,}{ฉพฺพสฺสานิ นิสีทนํ.} \\ \csromangathabat{ริญฺจนฺติ รูปิกา เจว,}{อุโภ นานปฺปการกา.} \\ \csromangathabat{อูนพนฺธนวสฺสิกา,}{สุตฺตํ วิกปฺปเนน จ.} \\ \csromangathabat{ทฺวารทานสิพฺพานิ\textsuperscript{1} จ,}{ปูวปจฺจยโชติ จ.} \\ \csromangathabat{รตนํ สูจิ มญฺโจ จ,}{ตูลํ นิสีทนกณฺฑุ จ.} \\ \csromangathabat{วสฺสิกา จ สุคเตน,}{วิญฺญตฺติ อญฺญํ เจตาปนา.} \\ \csromangathabat{เทฺว สํฆิกา มหาชนิกา เทฺว,}{ปุคฺคลลหุกา ครุ.} \\ \csromangathabat{เทฺว วิฆาสา สาฏิกา จ,}{สมณจีวเรน จ.} \\ \csromangathabat{สมปญฺญาสิเม ธมฺมา,}{ฉหิ ฐาเนหิ ชายเร.} \\ \csromangathabat{กายโต น วาจาจิตฺตา,}{วาจโต น กายมนา.} \\ \csromangathabat{กายวาจา น จ จิตฺตา\textsuperscript{1},}{กายจิตฺตา น วาจิกา\textsuperscript{1}.} \\ \csromangathabat{วาจาจิตฺตา น กาเยน,}{ตีหิ ทฺวาเรหิ\textsuperscript{1} ชายเร.} \\ \csromangathabat{ฉสมุฏฺฐานิกา เอเต,}{สญฺจริตฺเตน สาทิสา.}}
\csromangathasetleft{สญฺจรี กุฏิ วิหาโร, \\ วิญฺญตฺตุตฺตริ อภิหฏฺฐุํ, \\ โกสิยา สุทฺธเทฺวภาคา, \\ ริญฺจนฺติ รูปิกา เจว, \\ อูนพนฺธนวสฺสิกา, \\ ทฺวารทานสิพฺพานิ\textsuperscript{1} จ, \\ รตนํ สูจิ มญฺโจ จ, \\ วสฺสิกา จ สุคเตน, \\ เทฺว สํฆิกา มหาชนิกา เทฺว, \\ เทฺว วิฆาสา สาฏิกา จ, \\ สมปญฺญาสิเม ธมฺมา, \\ กายโต น วาจาจิตฺตา, \\ กายวาจา น จ จิตฺตา\textsuperscript{1}, \\ วาจาจิตฺตา น กาเยน, \\ ฉสมุฏฺฐานิกา เอเต,}
\csromangathagroupwithcloser{\csromangathastanza{\csromangathaitembat{260}{สญฺจรี กุฏิ วิหาโร,}{โธวนญฺจ ปฏิคฺคโห.} \\ \csromangathacontbat{วิญฺญตฺตุตฺตริ อภิหฏฺฐุํ,}{อุภินฺนํ ทูตเกน จ.} \\ \csromangathacontbat{โกสิยา สุทฺธเทฺวภาคา,}{ฉพฺพสฺสานิ นิสีทนํ.} \\ \csromangathacontbat{ริญฺจนฺติ รูปิกา เจว,}{อุโภ นานปฺปการกา.} \\ \csromangathacontbat{อูนพนฺธนวสฺสิกา,}{สุตฺตํ วิกปฺปเนน จ.} \\ \csromangathacontbat{ทฺวารทานสิพฺพานิ\csromansharedfootnote{5cd277c6e3366d54}{ทฺวารทานสิพฺพินี (สี, สฺยา)} จ,}{ปูวปจฺจยโชติ จ.} \\ \csromangathacontbat{รตนํ สูจิ มญฺโจ จ,}{ตูลํ นิสีทนกณฺฑุ จ.} \\ \csromangathacontbat{วสฺสิกา จ สุคเตน,}{วิญฺญตฺติ อญฺญํ เจตาปนา.} \\ \csromangathacontbat{เทฺว สํฆิกา มหาชนิกา เทฺว,}{ปุคฺคลลหุกา ครุ.} \\ \csromangathacontbat{เทฺว วิฆาสา สาฏิกา จ,}{สมณจีวเรน จ.} \\ \csromangathacontbat{สมปญฺญาสิเม ธมฺมา,}{ฉหิ ฐาเนหิ ชายเร.} \\ \csromangathacontbat{กายโต น วาจาจิตฺตา,}{วาจโต น กายมนา.} \\ \csromangathacontbat{กายวาจา น จ จิตฺตา\csromansharedfootnote{5153d8713b8a99f9}{น จิตฺตโต (สี, สฺยา)},}{กายจิตฺตา น วาจิกา\csromansharedfootnote{842875173ee623be}{น วาจโต (สี, สฺยา)}.} \\ \csromangathacontbat{วาจาจิตฺตา น กาเยน,}{ตีหิ ทฺวาเรหิ\csromansharedfootnote{951be36857920615}{ตีหิ ฐาเนหิ (สฺยา)} ชายเร.} \\ \csromangathacontbat{ฉสมุฏฺฐานิกา เอเต,}{สญฺจริตฺเตน สาทิสา.}}}{สญฺจริตฺต\-สมุฏฺฐานํ นิฏฺฐิตํ.}
\csromanmatikaanchor{mtk.128}
\csromantocmarkref{section}{4. สมนุภาสนาสมุฏฺฐาน}{mtk.128}
\bookmark[dest={mtk.128},level=1]{4. สมนุภาสนาสมุฏฺฐาน}
\csromanheader{1}{4. สมนุภาสนา\-สมุฏฺฐาน}
\setlength{\csromangathapagewidth}{0pt}
\setlength{\csromangathawakwidth}{0pt}
\csromangathameasuregroup{เภทานุวตฺตทุพฺพจ, \\ ฉนฺทํ อุชฺชคฺฆิกา เทฺว จ, \\ ฉมา นีจาสเน ฐานํ, \\ วชฺชานุวตฺติคหณา, \\ กิสฺมิํ สํสฏฺฐา เทฺว วธิ, \\ ปุน สํสฏฺฐา น วูปสเม, \\ อนฺวทฺธํ\textsuperscript{1} สห ชีวินิํ เทฺว, \\ สตฺตติํส อิเม ธมฺมา, \\ สพฺเพ เอกสมุฏฺฐานา,}{\csromangathabat{เภทานุวตฺตทุพฺพจ,}{ทูสทุฏฺฐุลฺลทิฏฺฐิ จ.} \\ \csromangathabat{ฉนฺทํ อุชฺชคฺฆิกา เทฺว จ,}{เทฺว จ สทฺทา น พฺยาหเร.} \\ \csromangathabat{ฉมา นีจาสเน ฐานํ,}{ปจฺฉโต อุปฺปเถน จ.} \\ \csromangathabat{วชฺชานุวตฺติคหณา,}{โอสาเร ปจฺจาจิกฺขนา.} \\ \csromangathabat{กิสฺมิํ สํสฏฺฐา เทฺว วธิ,}{วิสิพฺเพ ทุกฺขิตาย จ.} \\ \csromangathabat{ปุน สํสฏฺฐา น วูปสเม,}{อารามญฺจ ปวารณา.} \\ \csromangathabat{อนฺวทฺธํ\textsuperscript{1} สห ชีวินิํ เทฺว,}{จีวรํ อนุพนฺธนา.} \\ \csromangathabat{สตฺตติํส อิเม ธมฺมา,}{กายวาจาย จิตฺตโต.} \\ \csromangathabat{สพฺเพ เอกสมุฏฺฐานา,}{สมนุภาสนา ยถา.}}
\csromangathasetleft{เภทานุวตฺตทุพฺพจ, \\ ฉนฺทํ อุชฺชคฺฆิกา เทฺว จ, \\ ฉมา นีจาสเน ฐานํ, \\ วชฺชานุวตฺติคหณา, \\ กิสฺมิํ สํสฏฺฐา เทฺว วธิ, \\ ปุน สํสฏฺฐา น วูปสเม, \\ อนฺวทฺธํ\textsuperscript{1} สห ชีวินิํ เทฺว, \\ สตฺตติํส อิเม ธมฺมา, \\ สพฺเพ เอกสมุฏฺฐานา,}
\csromangathagroupwithcloserruled{\csromangathastanza{\csromanfolio{160}\csromangathaitembat{261}{เภทา\-นุวตฺต\-ทุพฺพจ,}{ทูสทุฏฺฐุลฺล\-ทิฏฺฐิ จ.} \\ \csromangathacontbat{ฉนฺทํ อุชฺชคฺฆิกา เทฺว จ,}{เทฺว จ สทฺทา น พฺยาหเร.} \\ \csromangathacontbat{ฉมา นีจาสเน ฐานํ,}{ปจฺฉโต อุปฺปเถน จ.} \\ \csromangathacontbat{วชฺชานุวตฺติคหณา,}{โอสาเร ปจฺจาจิกฺขนา.} \\ \csromangathacontbat{กิสฺมิํ สํสฏฺฐา เทฺว วธิ,}{วิสิพฺเพ ทุกฺขิตาย จ.} \\ \csromangathacontbat{ปุน สํสฏฺฐา น วูปสเม,}{อารามญฺจ ปวารณา.} \\ \csromangathacontbat{อนฺวทฺธํ\csromansharedfootnote{cd98119baaef1e6b}{อนฺวทฺธมาสํ (สี, สฺยา)} สห ชีวินิํ เทฺว,}{จีวรํ อนุพนฺธนา.} \\ \csromangathacontbat{สตฺตติํส อิเม ธมฺมา,}{กายวาจาย จิตฺตโต.}}\csromangathastanza{\csromangathabat{สพฺเพ เอกสมุฏฺฐานา,}{สมนุภาสนา ยถา.}}}{สนุภาสนาสมุฏฺฐานํ นิฏฺฐิตํ.}
\csromanmatikaanchor{mtk.129}
\csromantocmarkref{section}{5. กถินสมุฏฺฐาน}{mtk.129}
\bookmark[dest={mtk.129},level=1]{5. กถินสมุฏฺฐาน}
\csromanheader{1}{5. กถิน\-สมุฏฺฐาน}
\setlength{\csromangathapagewidth}{0pt}
\setlength{\csromangathawakwidth}{0pt}
\csromangathameasuregroup{อุพฺภตํ กถินํ ตีณิ, \\ อจฺเจกํ จาปิ สาสงฺกํ, \\ อุปสฺสยํ ปรมฺปรา, \\ วิกปฺปํ รญฺโญ วิกาเล, \\ อุสฺสยาสนฺนิจยญฺจ, \\ ปญฺจาหิกา สงฺกมนี, \\ ปสาเข อาสเน เจว, \\ กายวาจา น จ จิตฺตา\textsuperscript{1}, \\ ทฺวิสมุฏฺฐานิกา สพฺเพ,}{\csromangathabat{อุพฺภตํ กถินํ ตีณิ,}{ปฐมํ ปตฺตเภสชฺชํ.} \\ \csromangathabat{อจฺเจกํ จาปิ สาสงฺกํ,}{ปกฺกมนฺเตน วา ทุเว.} \\ \csromangathabat{อุปสฺสยํ ปรมฺปรา,}{อนติริตฺตํ นิมนฺตนา.} \\ \csromangathabat{วิกปฺปํ รญฺโญ วิกาเล,}{โวสาสารญฺญเกน จ.} \\ \csromangathabat{อุสฺสยาสนฺนิจยญฺจ,}{ปุเร ปจฺฉา วิกาเล จ.} \\ \csromangathabat{ปญฺจาหิกา สงฺกมนี,}{เทฺวปิ อาวสเถน จ.} \\ \csromangathabat{ปสาเข อาสเน เจว,}{ติํส เอกูนกา อิเม.} \\ \csromangathabat{กายวาจา น จ จิตฺตา\textsuperscript{1},}{ตีหิ ทฺวาเรหิ ชายเร.} \\ \csromangathabat{ทฺวิสมุฏฺฐานิกา สพฺเพ,}{กถิเนน สหาสมา.}}
\csromangathasetleft{อุพฺภตํ กถินํ ตีณิ, \\ อจฺเจกํ จาปิ สาสงฺกํ, \\ อุปสฺสยํ ปรมฺปรา, \\ วิกปฺปํ รญฺโญ วิกาเล, \\ อุสฺสยาสนฺนิจยญฺจ, \\ ปญฺจาหิกา สงฺกมนี, \\ ปสาเข อาสเน เจว, \\ กายวาจา น จ จิตฺตา\textsuperscript{1}, \\ ทฺวิสมุฏฺฐานิกา สพฺเพ,}
\csromangathagroupwithcloser{\csromangathastanza{\csromangathaitembat{262}{อุพฺภตํ กถินํ ตีณิ,}{ปฐมํ ปตฺตเภสชฺชํ.} \\ \csromangathacontbat{อจฺเจกํ จาปิ สาสงฺกํ,}{ปกฺกมนฺเตน วา ทุเว.} \\ \csromangathacontbat{อุปสฺสยํ ปรมฺปรา,}{อนติริตฺตํ นิมนฺตนา.} \\ \csromangathacontbat{วิกปฺปํ รญฺโญ วิกาเล,}{โวสาสารญฺญเกน จ.} \\ \csromangathacontbat{อุสฺสยาสนฺนิจยญฺจ,}{ปุเร ปจฺฉา วิกาเล จ.} \\ \csromangathacontbat{ปญฺจาหิกา สงฺกมนี,}{เทฺวปิ อาวสเถน จ.} \\ \csromangathacontbat{ปสาเข อาสเน เจว,}{ติํส เอกูนกา อิเม.} \\ \csromangathacontbat{กายวาจา น จ จิตฺตา\csromansharedfootnote{b3ea7c47299d28bc}{น จิตฺตโต (สฺยา)},}{ตีหิ ทฺวาเรหิ ชายเร.} \\ \csromangathacontbat{ทฺวิสมุฏฺฐานิกา สพฺเพ,}{กถิเนน สหาสมา.}}}{กถิน\-สมุฏฺฐานํ นิฏฺฐิตํ.}
\csromanheader{2}{6. \textbf{เอฬกโลม สมุฏฺฐาน}}
\csromanmatikaanchor{mtk.130}
\csromantocmarkref{section}{6. เอฬกโลมสมุฏฺฐาน}{mtk.130}
\bookmark[dest={mtk.130},level=1]{6. เอฬกโลมสมุฏฺฐาน}
\setlength{\csromangathapagewidth}{0pt}
\setlength{\csromangathawakwidth}{0pt}
\csromangathameasuregroup{เอฬกโลมา เทฺว เสยฺยา, \\ คณวิกาลสนฺนิธิ, \\ อุยฺยุตฺตํ เสนํ\textsuperscript{1} อุยฺโยธิ, \\ ทุพฺพณฺเณ เทฺว เทสนิกา, \\ นฺหานมตฺถรณํ เสยฺยา, \\ อนฺโตวสฺสํ จิตฺตาคารํ, \\ เวยฺยาวจฺจํ สหตฺถา จ, \\ ฉตฺตํ ยานญฺจ สงฺฆาณิํ, \\ ภิกฺขุนี สิกฺขมานา จ, \\ อสํกจฺจิกา อาปตฺติ, \\ กาเยน น วาจาจิตฺเตน, \\ ทฺวีสมุฏฺฐานิกา สพฺเพ,}{\csromangathabat{เอฬกโลมา เทฺว เสยฺยา,}{อาหจฺจ ปิณฺฑโภชนํ.} \\ \csromangathabat{คณวิกาลสนฺนิธิ,}{ทนฺตโปเนน เจลกา.} \\ \csromangathabat{อุยฺยุตฺตํ เสนํ\textsuperscript{1} อุยฺโยธิ,}{สุรา โอเรน นฺหายนา.} \\ \csromangathabat{ทุพฺพณฺเณ เทฺว เทสนิกา,}{ลสุณุปติฏฺเฐ นจฺจนา.} \\ \csromangathabat{นฺหานมตฺถรณํ เสยฺยา,}{อนฺโตรฏฺเฐ ตถา พหิ.} \\ \csromangathabat{อนฺโตวสฺสํ จิตฺตาคารํ,}{อาสนฺทิ สุตฺตกนฺตนา.} \\ \csromangathabat{เวยฺยาวจฺจํ สหตฺถา จ,}{อภิกฺขุกาวาเสน จ.} \\ \csromangathabat{ฉตฺตํ ยานญฺจ สงฺฆาณิํ,}{อลงฺการํ คนฺธวาสิตํ.} \\ \csromangathabat{ภิกฺขุนี สิกฺขมานา จ,}{สามเณรี คิหินิยา.} \\ \csromangathabat{อสํกจฺจิกา อาปตฺติ,}{จตฺตารีสา จตุตฺตริ.} \\ \csromangathabat{กาเยน น วาจาจิตฺเตน,}{กายจิตฺเตน น วาจโต.} \\ \csromangathabat{ทฺวีสมุฏฺฐานิกา สพฺเพ,}{สมา เอฬกโลมิกาติ.}}
\csromangathasetleft{เอฬกโลมา เทฺว เสยฺยา, \\ คณวิกาลสนฺนิธิ, \\ อุยฺยุตฺตํ เสนํ\textsuperscript{1} อุยฺโยธิ, \\ ทุพฺพณฺเณ เทฺว เทสนิกา, \\ นฺหานมตฺถรณํ เสยฺยา, \\ อนฺโตวสฺสํ จิตฺตาคารํ, \\ เวยฺยาวจฺจํ สหตฺถา จ, \\ ฉตฺตํ ยานญฺจ สงฺฆาณิํ, \\ ภิกฺขุนี สิกฺขมานา จ, \\ อสํกจฺจิกา อาปตฺติ, \\ กาเยน น วาจาจิตฺเตน, \\ ทฺวีสมุฏฺฐานิกา สพฺเพ,}
\csromangathagroupwithcloserruled{\csromangathastanza{\csromanfolio{161}\csromangathaitembat{263}{เอฬกโลมา เทฺว เสยฺยา,}{อาหจฺจ ปิณฺฑโภชนํ.} \\ \csromangathacontbat{คณวิกาลสนฺนิธิ,}{ทนฺตโปเนน เจลกา.} \\ \csromangathacontbat{อุยฺยุตฺตํ เสนํ\csromansharedfootnote{5669f4d2c455fd1a}{อุยฺยุตฺตํ วเส (สฺยา)} อุยฺโยธิ,}{สุรา โอเรน นฺหายนา.} \\ \csromangathacontbat{ทุพฺพณฺเณ เทฺว เทสนิกา,}{ลสุณุปติฏฺเฐ นจฺจนา.} \\ \csromangathacontbat{นฺหานมตฺถรณํ เสยฺยา,}{อนฺโตรฏฺเฐ ตถา พหิ.} \\ \csromangathacontbat{อนฺโตวสฺสํ จิตฺตาคารํ,}{อาสนฺทิ สุตฺตกนฺตนา.} \\ \csromangathacontbat{เวยฺยาวจฺจํ สหตฺถา จ,}{อภิกฺขุกาวาเสน จ.} \\ \csromangathacontbat{ฉตฺตํ ยานญฺจ สงฺฆาณิํ,}{อลงฺการํ คนฺธวาสิตํ.} \\ \csromangathacontbat{ภิกฺขุนี สิกฺขมานา จ,}{สามเณรี คิหินิยา.} \\ \csromangathacontbat{อสํกจฺจิกา อาปตฺติ,}{จตฺตารีสา จตุตฺตริ.} \\ \csromangathacontbat{กาเยน น วาจาจิตฺเตน,}{กายจิตฺเตน น วาจโต.} \\ \csromangathacontbat{ทฺวีสมุฏฺฐานิกา สพฺเพ,}{สมา เอฬกโลมิกาติ.}}}{เอฬกโลม\-สมุฏฺฐานํ นิฏฺฐิตํ.}
\csromanmatikaanchor{mtk.131}
\csromantocmarkref{section}{7. ปทโสธมฺมสมุฏฺฐาน}{mtk.131}
\bookmark[dest={mtk.131},level=1]{7. ปทโสธมฺมสมุฏฺฐาน}
\csromanheader{1}{7. ปทโสธมฺม\-สมุฏฺฐาน}
\proseitem{264}{ปทญฺญตฺร อสมฺมตา, ตถา อตฺถงฺคเตน จ.}

\setlength{\csromangathapagewidth}{0pt}
\setlength{\csromangathawakwidth}{0pt}
\csromangathameasuregroup{ติรจฺฉานวิชฺชา เทฺว วุตฺตา, \\ สตฺต สิกฺขาปทา เอเต, \\ วาจาจิตฺเตน ชายนฺติ,}{\csromangathabat{ติรจฺฉานวิชฺชา เทฺว วุตฺตา,}{อโนกาโส\textsuperscript{1} จ ปุจฺฉนา.} \\ \csromangathabat{สตฺต สิกฺขาปทา เอเต,}{วาจา น กายจิตฺตโต\textsuperscript{1}.} \\ \csromangathabat{วาจาจิตฺเตน ชายนฺติ,}{น ตุ กาเยน ชายเร.}}
\csromangathasetleft{ติรจฺฉานวิชฺชา เทฺว วุตฺตา, \\ สตฺต สิกฺขาปทา เอเต, \\ วาจาจิตฺเตน ชายนฺติ,}
\csromangathagroupwithcloser{\csromangathastanza{\csromangathabat{ติรจฺฉาน\-วิชฺชา เทฺว วุตฺตา,}{อโนกาโส\csromansharedfootnote{8c1d74f3441f204e}{อโนกาเส (สี, สฺยา)} จ ปุจฺฉนา.} \\ \csromangathabat{สตฺต สิกฺขาปทา เอเต,}{วาจา น กายจิตฺตโต\csromansharedfootnote{8c1d74f3441f204e}{อโนกาเส (สี, สฺยา)}.} \\ \csromangathabat{วาจาจิตฺเตน ชายนฺติ,}{น ตุ กาเยน ชายเร.}}}{ทฺวิสมุฏฺฐานิกา สพฺเพ, ปทโสธมฺม\-สทิสา. ปทโสธมฺม\-สมุฏฺฐานํ นิฏฺฐิตํ.}
\csromanmatikaanchor{mtk.132}
\csromantocmarkref{section}{8. อทฺธานสมุฏฺฐาน}{mtk.132}
\bookmark[dest={mtk.132},level=1]{8. อทฺธานสมุฏฺฐาน}
\csromanheader{1}{8. อทฺธาน\-สมุฏฺฐาน}
\setlength{\csromangathapagewidth}{0pt}
\setlength{\csromangathawakwidth}{0pt}
\csromangathameasuregroup{อทฺธานนาวํ ปณีตํ, \\ ธญฺญํ นิมนฺติตา เจว, \\ สิกฺขา ปนฺนรส เอเต, \\ กายวาจาหิ ชายนฺติ, \\ กายจิตฺเตน ชายนฺติ, \\ กายวาจาหิ จิตฺเตน, \\ ปญฺญาตฺตา พุทฺธญาเณน,}{\csromangathabat{อทฺธานนาวํ ปณีตํ,}{มาตุคาเมน สํหเร.} \\ \csromangathabat{ธญฺญํ นิมนฺติตา เจว,}{อฏฺฐ จ ปาฏิเทสนี.} \\ \csromangathabat{สิกฺขา ปนฺนรส เอเต,}{กายา น วาจา น มนา.} \\ \csromangathabat{กายวาจาหิ ชายนฺติ,}{น เต จิตฺเตน ชายเร.} \\ \csromangathabat{กายจิตฺเตน ชายนฺติ,}{น เต ชายนฺติ วาจโต.} \\ \csromangathabat{กายวาจาหิ จิตฺเตน,}{สมุฏฺฐานา จตุพฺพิธา.} \\ \csromangathabat{ปญฺญาตฺตา พุทฺธญาเณน,}{อทฺธาเนน สหา สมา\textsuperscript{1}.}}
\csromangathasetleft{อทฺธานนาวํ ปณีตํ, \\ ธญฺญํ นิมนฺติตา เจว, \\ สิกฺขา ปนฺนรส เอเต, \\ กายวาจาหิ ชายนฺติ, \\ กายจิตฺเตน ชายนฺติ, \\ กายวาจาหิ จิตฺเตน, \\ ปญฺญาตฺตา พุทฺธญาเณน,}
\csromangathagroupwithcloserruled{\csromangathastanza{\csromanfolio{162}\csromangathaitembat{265}{อทฺธานนาวํ ปณีตํ,}{มาตุคาเมน สํหเร.} \\ \csromangathacontbat{ธญฺญํ นิมนฺติตา เจว,}{อฏฺฐ จ ปาฏิเทสนี.} \\ \csromangathacontbat{สิกฺขา ปนฺนรส เอเต,}{กายา น วาจา น มนา.} \\ \csromangathacontbat{กายวาจาหิ ชายนฺติ,}{น เต จิตฺเตน ชายเร.} \\ \csromangathacontbat{กายจิตฺเตน ชายนฺติ,}{น เต ชายนฺติ วาจโต.} \\ \csromangathacontbat{กายวาจาหิ จิตฺเตน,}{สมุฏฺฐานา จตุพฺพิธา.} \\ \csromangathacontbat{ปญฺญาตฺตา พุทฺธญาเณน,}{อทฺธาเนน สหา สมา\csromansharedfootnote{dcad73ce49d12f05}{สมานยา (สฺยา)}.}}}{อทฺธาน\-สมุฏฺฐานํ นิฏฺฐิตํ.}
\csromanmatikaanchor{mtk.133}
\csromantocmarkref{section}{9. เถยฺยสตฺถสมุฏฺฐาน}{mtk.133}
\bookmark[dest={mtk.133},level=1]{9. เถยฺยสตฺถสมุฏฺฐาน}
\csromanheader{1}{9. เถยฺยสตฺถ\-สมุฏฺฐาน}
\setlength{\csromangathapagewidth}{0pt}
\setlength{\csromangathawakwidth}{0pt}
\csromangathameasuregroup{เถยฺยสตฺถํ อุปสฺสุติ, \\ รตฺติฉนฺนญฺจ โอกาสํ, \\ กายจิตฺเตน ชายนฺติ, \\ ตีหิ ทฺวาเรหิ ชายนฺติ, \\ เถยฺยสตฺถสมุฏฺฐานา,}{\csromangathabat{เถยฺยสตฺถํ อุปสฺสุติ,}{สูปวิญฺญาปเนน จ.} \\ \csromangathabat{รตฺติฉนฺนญฺจ โอกาสํ,}{เอเต พฺยูเหน สตฺตมา.} \\ \csromangathabat{กายจิตฺเตน ชายนฺติ,}{น เต ชายนฺติ วาจโต.} \\ \csromangathabat{ตีหิ ทฺวาเรหิ ชายนฺติ,}{ทฺวิสมุฏฺฐานิกา อิเม.} \\ \csromangathabat{เถยฺยสตฺถสมุฏฺฐานา,}{เทสิตาทิจฺจพนฺธุนา.}}
\csromangathasetleft{เถยฺยสตฺถํ อุปสฺสุติ, \\ รตฺติฉนฺนญฺจ โอกาสํ, \\ กายจิตฺเตน ชายนฺติ, \\ ตีหิ ทฺวาเรหิ ชายนฺติ, \\ เถยฺยสตฺถสมุฏฺฐานา,}
\csromangathagroupwithcloserruled{\csromangathastanza{\csromangathaitembat{266}{เถยฺยสตฺถํ อุปสฺสุติ,}{สูปวิญฺญาปเนน จ.} \\ \csromangathacontbat{รตฺติฉนฺนญฺจ โอกาสํ,}{เอเต พฺยูเหน สตฺตมา.} \\ \csromangathacontbat{กายจิตฺเตน ชายนฺติ,}{น เต ชายนฺติ วาจโต.} \\ \csromangathacontbat{ตีหิ ทฺวาเรหิ ชายนฺติ,}{ทฺวิสมุฏฺฐานิกา อิเม.} \\ \csromangathacontbat{เถยฺยสตฺถสมุฏฺฐานา,}{เทสิตาทิจฺจพนฺธุนา.}}}{เถยฺยสตฺถ\-สมุฏฺฐานํ นิฏฺฐิตํ.}
\csromanmatikaanchor{mtk.134}
\csromantocmarkref{section}{10. ธมฺมเทสนาสมุฏฺฐาน}{mtk.134}
\bookmark[dest={mtk.134},level=1]{10. ธมฺมเทสนาสมุฏฺฐาน}
\csromanheader{1}{10. ธมฺมเทสนา\-สมุฏฺฐาน}
\setlength{\csromangathapagewidth}{0pt}
\setlength{\csromangathawakwidth}{0pt}
\csromangathameasuregroup{ฉตฺตปาณิสฺส สทฺธมฺมํ, \\ เอวเมว\textsuperscript{1} ทณฺฑปาณิสฺส, \\ ปาทุกุปาหนา ยานํ, \\ เวฐิโตคุณฺฐิโต เจว, \\ วาจาจิตฺเตน ชายนฺติ, \\ สพฺเพ เอกสมุฏฺฐานา,}{\csromangathabat{ฉตฺตปาณิสฺส สทฺธมฺมํ,}{น เทเสนฺติ ตถาคตา.} \\ \csromangathabat{เอวเมว\textsuperscript{1} ทณฺฑปาณิสฺส,}{สตฺถอาวุธปาณินํ.} \\ \csromangathabat{ปาทุกุปาหนา ยานํ,}{เสยฺยปลฺลตฺถิกาย จ.} \\ \csromangathabat{เวฐิโตคุณฺฐิโต เจว,}{เอกาทสมนูนกา.} \\ \csromangathabat{วาจาจิตฺเตน ชายนฺติ,}{น เต ชายนฺติ กายโต.} \\ \csromangathabat{สพฺเพ เอกสมุฏฺฐานา,}{สมกา ธมฺมเทสเน.}}
\csromangathasetleft{ฉตฺตปาณิสฺส สทฺธมฺมํ, \\ เอวเมว\textsuperscript{1} ทณฺฑปาณิสฺส, \\ ปาทุกุปาหนา ยานํ, \\ เวฐิโตคุณฺฐิโต เจว, \\ วาจาจิตฺเตน ชายนฺติ, \\ สพฺเพ เอกสมุฏฺฐานา,}
\csromangathagroupwithcloser{\csromangathastanza{\csromangathaitembat{267}{ฉตฺตปาณิสฺส สทฺธมฺมํ,}{น เทเสนฺติ ตถาคตา.} \\ \csromangathacontbat{เอวเมว\csromansharedfootnote{79fc19c77904a235}{ตเถว (สี, สฺยา)} ทณฺฑปาณิสฺส,}{สตฺถอาวุธปาณินํ.} \\ \csromangathacontbat{ปาทุกุปาหนา ยานํ,}{เสยฺยปลฺลตฺถิกาย จ.} \\ \csromangathacontbat{เวฐิโตคุณฺฐิโต เจว,}{เอกาทสมนูนกา.} \\ \csromangathacontbat{วาจาจิตฺเตน ชายนฺติ,}{น เต ชายนฺติ กายโต.} \\ \csromangathacontbat{สพฺเพ เอกสมุฏฺฐานา,}{สมกา ธมฺมเทสเน.}}}{ธมฺมเทสนา\-สมุฏฺฐานํ นิฏฺฐิตํ.}
\csromanmatikaanchor{mtk.135}
\csromantocmarkref{section}{11. ภูตาโรจนสมุฏฺฐาน}{mtk.135}
\bookmark[dest={mtk.135},level=1]{11. ภูตาโรจนสมุฏฺฐาน}
\csromanheader{1}{11. ภูตาโรจน\-สมุฏฺฐาน}
\setlength{\csromangathapagewidth}{0pt}
\setlength{\csromangathawakwidth}{0pt}
\csromangathameasuregroup{ภูตํ กาเยน ชายติ, \\ วาจโต จ สมุฏฺฐาติ, \\ กายวาจาย ชายติ, \\ ภูตาโรจนกา นาม,}{\csromangathabat{ภูตํ กาเยน ชายติ,}{น วาจา น จ จิตฺตโต.} \\ \csromangathabat{วาจโต จ สมุฏฺฐาติ,}{น กายา น จ จิตฺตโต.} \\ \csromangathabat{กายวาจาย ชายติ,}{น ตุ ชายติ จิตฺตโต.} \\ \csromangathabat{ภูตาโรจนกา นาม,}{ตีหิ ฐาเนหิ ชายติ.}}
\csromangathasetleft{ภูตํ กาเยน ชายติ, \\ วาจโต จ สมุฏฺฐาติ, \\ กายวาจาย ชายติ, \\ ภูตาโรจนกา นาม,}
\csromangathagroupwithcloserruled{\csromangathastanza{\csromanfolio{163}\csromangathaitembat{268}{ภูตํ กาเยน ชายติ,}{น วาจา น จ จิตฺตโต.} \\ \csromangathacontbat{วาจโต จ สมุฏฺฐาติ,}{น กายา น จ จิตฺตโต.} \\ \csromangathacontbat{กายวาจาย ชายติ,}{น ตุ ชายติ จิตฺตโต.} \\ \csromangathacontbat{ภูตาโรจนกา นาม,}{ตีหิ ฐาเนหิ ชายติ.}}}{ภูตาโรจน\-สมุฏฺฐานํ นิฏฺฐิตํ.}
\csromanmatikaanchor{mtk.136}
\csromantocmarkref{section}{12. โจริวุฏฺฐาปนสมุฏฺฐาน}{mtk.136}
\bookmark[dest={mtk.136},level=1]{12. โจริวุฏฺฐาปนสมุฏฺฐาน}
\csromanheader{1}{12. โจริ\-วุฏฺฐาปน\-สมุฏฺฐาน}
\setlength{\csromangathapagewidth}{0pt}
\setlength{\csromangathawakwidth}{0pt}
\csromangathameasuregroup{โจรี วาจาย จิตฺเตน, \\ ชายติ ตีหิ ทฺวาเรหิ, \\ อกตํ ทฺวิสมุฏฺฐานํ,}{\csromangathabat{โจรี วาจาย จิตฺเตน,}{น ตํ ชายติ กายโต.} \\ \csromangathabat{ชายติ ตีหิ ทฺวาเรหิ,}{โจริวุฏฺฐาปนํ อิทํ.} \\ \csromangathabat{อกตํ ทฺวิสมุฏฺฐานํ,}{ธมฺมราเชน ภาสิตํ\textsuperscript{1}.}}
\csromangathasetleft{โจรี วาจาย จิตฺเตน, \\ ชายติ ตีหิ ทฺวาเรหิ, \\ อกตํ ทฺวิสมุฏฺฐานํ,}
\csromangathagroupwithcloserruled{\csromangathastanza{\csromangathaitembat{269}{โจรี วาจาย จิตฺเตน,}{น ตํ ชายติ กายโต.} \\ \csromangathacontbat{ชายติ ตีหิ ทฺวาเรหิ,}{โจริวุฏฺฐาปนํ อิทํ.} \\ \csromangathacontbat{อกตํ ทฺวิสมุฏฺฐานํ,}{ธมฺมราเชน ภาสิตํ\csromansharedfootnote{844421e919f1e5b4}{ฐปิตํ (สฺยา)}.}}}{โจริ\-วุฏฺฐาปน\-สมุฏฺฐานํ นิฏฺฐิตํ.}
\csromanmatikaanchor{mtk.137}
\csromantocmarkref{section}{13. อนนุญฺญาตสมุฏฺฐาน}{mtk.137}
\bookmark[dest={mtk.137},level=1]{13. อนนุญฺญาตสมุฏฺฐาน}
\csromanheader{1}{13. อนนุญฺญาต\-สมุฏฺฐาน}
\setlength{\csromangathapagewidth}{0pt}
\setlength{\csromangathawakwidth}{0pt}
\csromangathameasuregroup{อนนุญฺญาตํ วาจาย, \\ ชายติ กายวาจาย, \\ ชายติ วาจาจิตฺเตน, \\ ชายติ ตีหิ ทฺวาเรหิ,}{\csromangathabat{อนนุญฺญาตํ วาจาย,}{น กายา น จ จิตฺตโต.} \\ \csromangathabat{ชายติ กายวาจาย,}{น ตํ ชายติ จิตฺตโต.} \\ \csromangathabat{ชายติ วาจาจิตฺเตน,}{น ตํ ชายติ กายโต.} \\ \csromangathabat{ชายติ ตีหิ ทฺวาเรหิ,}{อกตํ จตุฐานิกํ.}}
\csromangathasetleft{อนนุญฺญาตํ วาจาย, \\ ชายติ กายวาจาย, \\ ชายติ วาจาจิตฺเตน, \\ ชายติ ตีหิ ทฺวาเรหิ,}
\csromangathagroupwithcloser{\csromangathastanza{\csromangathaitembat{270}{อนนุญฺญาตํ วาจาย,}{น กายา น จ จิตฺตโต.} \\ \csromangathacontbat{ชายติ กายวาจาย,}{น ตํ ชายติ จิตฺตโต.} \\ \csromangathacontbat{ชายติ วาจาจิตฺเตน,}{น ตํ ชายติ กายโต.} \\ \csromangathacontbat{ชายติ ตีหิ ทฺวาเรหิ,}{อกตํ จตุฐานิกํ.}}}{อนนุญฺญาต\-สมุฏฺฐานํ นิฏฺฐิตํ.}
\setlength{\csromangathapagewidth}{0pt}
\setlength{\csromangathawakwidth}{0pt}
\csromangathameasuregroup{สมุฏฺฐานํ หิ สงฺเขปํ, \\ อสมฺโมหกรํ ฐานํ, \\ ธารยนฺโต อิมํ วิญฺญู,}{\csromangathabat{สมุฏฺฐานํ หิ สงฺเขปํ,}{ทส ตีณี สุเทสิตํ.} \\ \csromangathabat{อสมฺโมหกรํ ฐานํ,}{เนตฺติธมฺมานุโลมิกํ.} \\ \csromangathabat{ธารยนฺโต อิมํ วิญฺญู,}{สมุฏฺฐาเน น มุยฺหตีติ.}}
\csromangathasetleft{สมุฏฺฐานํ หิ สงฺเขปํ, \\ อสมฺโมหกรํ ฐานํ, \\ ธารยนฺโต อิมํ วิญฺญู,}
\csromangathagroupwithcloser{\csromangathastanza{\csromangathabat{สมุฏฺฐานํ หิ สงฺเขปํ,}{ทส ตีณี สุเทสิตํ.} \\ \csromangathabat{อสมฺโมหกรํ ฐานํ,}{เนตฺติธมฺมานุโลมิกํ.} \\ \csromangathabat{ธารยนฺโต อิมํ วิญฺญู,}{สมุฏฺฐาเน น มุยฺหตีติ.}}}{สมุฏฺฐาน\-สีส\-สงฺเขโป นิฏฺฐิโต.}
\csromanmatikaanchor{mtk.138}
\csromantocmarknopage{section}{อนฺตรเปยฺยาล}{mtk.138}
\bookmark[dest={mtk.138},level=1]{อนฺตรเปยฺยาล}
\csromanheader{1}{\textbf{อนฺตรเปยฺยาล}}
\csromanmatikaanchor{mtk.139}
\csromantocmarkref{section}{กติปุจฺฉาวาร}{mtk.139}
\bookmark[dest={mtk.139},level=1]{กติปุจฺฉาวาร}
\csromanheader{1}{\textbf{กติปุจฺฉาวาร}}
\proseitem{271}{\csromanfolio{164}กติ อาปตฺติโย, กติ อาปตฺติกฺขนฺธา, กติ วินีตวตฺถูนิ, กติ อคารวา, กติ คารวา, กติ วินีตวตฺถูนิ, กติ วิปตฺติโย, กติ อาปตฺติ\-สมุฏฺฐานา, กติ วิวาทมูลานิ, กติ อนุวาทมูลานิ, กติ สารณียา ธมฺมา, กติ เภทกร\-วตฺถูนิ, กติ อธิกรณานิ, กติ สมถา.}

\prose{ปญฺจ อาปตฺติโย, ปญฺจ อาปตฺติกฺขนฺธา, ปญฺจ วินีตวตฺถูนิ, สตฺต อาปตฺติโย, สตฺต อาปตฺติกฺขนฺธา, สตฺต วินีตวตฺถูนิ, ฉ อคารวา, ฉ คารวา, ฉ วินีตวตฺถูนิ, จตสฺโส วิปตฺติโย, ฉ อาปตฺติ\-สมุฏฺฐานา, ฉ วิวาทมูลานิ, ฉ อนุวาทมูลานิ, ฉ สารณียา ธมฺมา, อฏฺฐารส เภทกร\-วตฺถูนิ, จตฺตาริ อธิกรณานิ, สตฺต สมถา.}

\prose{ตตฺถ กตมา ปญฺจ อาปตฺติโย. ปาราชิกาปตฺติ, สํฆาทิเสสาปตฺติ, ปาจิตฺติยาปตฺติ, ปาฏิเทสนียา\-ปตฺติ, ทุกฺกฏาปตฺติ. อิมา ปญฺจ อาปตฺติโย.}

\prose{ตตฺถ กตเม ปญฺจ อาปตฺติกฺขนฺธา. ปาราชิกาปตฺติ\-กฺขนฺโธ, สํฆาทิเสสา\-ปตฺติกฺขนฺโธ, ปาจิตฺติยาปตฺติ\-กฺขนฺโธ, ปาฏิเทสนียา\-ปตฺติกฺขนฺโธ, ทุกฺกฏาปตฺติ\-กฺขนฺโธ. อิเม ปญฺจ อาปตฺติกฺขนฺธา.}

\prose{ตตฺถ กตมานิ ปญฺจ วินีตวตฺถูนิ. ปญฺจหิ อาปตฺติ\-กฺขนฺเธหิ อารติ วิรติ ปฏิวิรติ เวรมณี อกิริยา อกรณํ อนชฺฌาปตฺติ เวลาอนติกฺกโม เสตุฆาโต. อิมานิ ปญฺจ วินีตวตฺถูนิ.}

\prose{ตตฺถ กตมา สตฺต อาปตฺติโย. ปาราชิกาปตฺติ, สํฆาทิเสสาปตฺติ, ถุลฺลจฺจยา\-ปตฺติ, ปาจิตฺติยาปตฺติ, ปาฏิเทสนียา\-ปตฺติ, ทุกฺกฏาปตฺติ, ทุพฺภาสิตา\-ปตฺติ. อิมา สตฺต อาปตฺติโย.}

\prose{ตตฺถ กตเม สตฺต อาปตฺติกฺขนฺธา. ปาราชิกาปตฺติ\-กฺขนฺโธ, สฆาทิเสสาปตฺติกฺขนฺโธ, ถุลฺลจฺจยา\-ปตฺติกฺขนฺโธ, ปาจิตฺติยาปตฺติ\-กฺขนฺโธ, ปาฏิเทสนียา\-ปตฺติกฺขนฺโธ, ทุกฺกฏาปตฺติ\-กฺขนฺโธ, ทุพฺภาสิตา\-ปตฺติกฺขนฺโธ. อิเม สตฺต อาปตฺติกฺขนฺธา.}

\prose{\csromanfolio{165}ตตฺถ กตมานิ สตฺต วินีตวตฺถูนิ. สตฺตหิ อาปตฺติ\-กฺขนฺเธหิ อารติ วิรติ ปฏิวิรติ เวรมณี อกิริยา อกรณํ อนชฺฌาปตฺติ เวลาอนติกฺกโม เสตุฆาโต. อิมานิ สตฺต วินีตวตฺถูนิ.}

\prose{ตตฺถ กตเม ฉ อคารวา. พุทฺเธ อคารโว, ธมฺเม อคารโว, สํเฆ อคารโว, สิกฺขาย อคารโว, อปฺปมาเท อคารโว, ปฏิสนฺธาเร อคารโว. อิเม ฉ อคารวา.}

\prose{ตตฺถ กตเม ฉ คารวา. พุทฺเธ คารโว, ธมฺเม คารโว, สํเฆ คารโว, สิกฺขาย คารโว, อปฺปมาเท คารโว, ปฏิสนฺธาเร คารโว. อิเม ฉ คารวา.}

\prose{ตตฺถ กตมานิ ฉ วินีตวตฺถูนิ. ฉหิ อคารเวหิ อารติ วิรติ ปฏิวิรติ เวรมณี อกิริยา อกรณํ อนชฺฌาปตฺติ เวลาอนติกฺกโม เสตุฆาโต. อิมานิ ฉ วินีตวตฺถูนิ.}

\prose{ตตฺถ กตมา จตสฺโส วิปตฺติโย. สีลวิปตฺติ, อาจารวิปตฺติ, ทิฏฺฐิวิปตฺติ, อาชีววิปตฺติ. อิมา จตสฺโส วิปตฺติโย.}

\prose{ตตฺถ กตเม ฉ อาปตฺติ\-สมุฏฺฐานา. อตฺถาปตฺติ กายโต สมุฏฺฐาติ, น วาจโต, น จิตฺตโต, อตฺถาปตฺติ วาจโต สมุฏฺฐาติ, น กายโต, น จิตฺตโต, อตฺถาปตฺติ กายโต จ วาจโต จ สมุฏฺฐาติ, น จิตฺตโต, อตฺถาปตฺติ กายโต จ จิตฺตโต จ สมุฏฺฐาติ, น วาจโต, อาตฺถาปตฺติ วาจโต จ จิตฺตโต จ สมุฏฺฐาติ, น กายโต, อตฺถาปตฺติ กายโต จ วาจโต จ จิตฺตโต จ สมุฏฺฐาติ. อิเม ฉ อาปตฺติ\-สมุฏฺฐานา.}

\proseitem{272}{ตตฺถ\csromansymbolfootnote{*}{วิ 4. 212; อํ 2. 294; ม 3. 33; ที 3. 204 ปิฏฺเฐสุปิ.}กตมานิ ฉ วิวาทมูลานิ. อิธ ภิกฺขุ โกธโน โหติ อุปนาหี, โย โส ภิกฺขุ โกธโน โหติ อุปนาหี, โส สตฺถริปิ อคารโว วิหรติ อปฺปติสฺโส, ธมฺเมปิ อคารโว วิหรติ อปฺปติสฺโส, สํเฆปิ อคารโว วิหรติ อปฺปติสฺโส, สิกฺขายปิ น ปริปูรการี โหติ, โย โส ภิกฺขุ สตฺถริ อคารโว วิหรติ อปฺปติสฺโส, ธมฺเม ฯเปฯ สํเฆ ฯเปฯ สิกฺขาย น ปริปูรการี, โส สํเฆ วิวาทํ ชเนติ, โย โหติ วิวาโท พหุชนาหิตาย พหุชนา\-สุขาย พหุโน ชนสฺส อนตฺถาย อหิตาย ทุกฺขาย เทวมนุสฺสานํ, เอวรูปญฺเจ ตุเมฺห วิวาทมูลํ อชฺฌตฺตํ วา พหิทฺธา วา สมนุปสฺเสยฺยาถ, ตตฺร \csromanfolio{166}ตุเมฺห ตสฺเสว ปาปกสฺส วิวาทมูลสฺส ปหานาย วายเมยฺยาถ, เอวรูปญฺเจ ตุเมฺห วิวาทมูลํ อชฺฌตฺตํ วา พหิทฺธา วา น สมนุปสฺเสยฺยาถ, ตตฺร ตุเมฺห ตสฺเสว ปาปกสฺส วิวาทมูลสฺส อายติํ อนวสฺสวาย ปฏิปชฺเชยฺยาถ, เอวเมตสฺส ปาปกสฺส วิวาทมูลสฺส ปหานํ โหติ, เอวเมตสฺส ปาปกสฺส วิวาทมูลสฺส อายติํ อนวสฺสโว โหติ.}

\prose{\csromansymbolfootnote{*}{วิ 4. 213 ปิฏฺเฐปิ.}ปุน จปรํ ภิกฺขุ มกฺขี โหติ ปฬาสี ฯเปฯ. อิสฺสุกี โหติ มจฺฉรี. สโฐ โหติ มายาวี. ปาปิจฺโฉ โหติ มิจฺฉาทิฏฺฐิ. สนฺทิฏฺฐิ\-ปรามาสี โหติ อาธานคฺคาหี ทุปฺปฏินิสฺสคฺคี, โย โส ภิกฺขุ สนฺทิฏฺฐิ\-ปรามาสี โหติ อาธานคฺคาหี ทุปฺปฏินิสฺสคฺคี, โส สตฺถริปิ อคารโว วิหรติ อปฺปติสฺโส, ธมฺเมปิ อคารโว วิหรติ อปฺปติสฺโส, สํเฆปิ อคารโว วิหรติ อปฺปติสฺโส, สิกฺขายปิ น ปริปูรการี โหติ, โย โส ภิกฺขุ สตฺถริ อคารโว วิหรติ อปฺปติสฺโส. ธมฺเม ฯเปฯ สํเฆ ฯเปฯ สิกฺขาย น ปริปูรการี, โส สํเฆ วิวาทํ ชเนติ, โย โหติ วิวาโท พหุชนาหิตาย พหุชนา\-สุขาย พหุโน ชนสฺส อนตฺถาย อหิตาย ทุกฺขาย เทวมนุสฺสานํ, เอวรูปญฺเจ ตุเมฺห วิวาทมูลํ อชฺฌตฺตํ วา พหิทฺธา วา สมนุปสฺเสยฺยาถ, ตตฺร ตุเมฺห ตสฺเสว ปาปกสฺส วิวาทมูลสฺส ปหานาย วายเมยฺยาถ, เอวรูปญฺเจ ตุเมฺห วิวาทมูลํ อชฺฌตฺตํ วา พหิทฺธา วา น สมนุปสฺเสยฺยาถ, ตตฺร ตุเมฺห ตสฺเสว ปาปกสฺส วิวาทมูลสฺส อายติํ อนวสฺสวาย ปฏิปชฺเชยฺยาถ, เอวเมตสฺส ปาปกสฺส วิวาทมูลสฺส ปหานํ โหติ, เอวเมตสฺส ปาปกสฺส วิวาทมูลสฺส อายติํ อนวสฺสโว โหติ. อิมานิ ฉ วิวาทมูลานิ.}

\proseitem{273}{\csromansymbolfootnote{+}{วิ 4. 214 ปิฏฺเฐปิ.}ตตฺถ กตมานิ ฉ อนุวาทมูลานิ. อิธ ภิกฺขุ โกธโน โหติ อุปนาหี, โย โส ภิกฺขุ โกธโน โหติ อุปนาหี, โส สตฺถริปิ อคารโว วิหรติ อปฺปติสฺโส, ธมฺเมปิ อคารโว วิหรติ อปฺปติสฺโส, สํเฆปิ อคารโว วิหรติ อปฺปติสฺโส, สิกฺขายปิ น ปริปูรการี โหติ, โย โส ภิกฺขุ สตฺถริ อคารโว วิหรติ อปฺปติสฺโส, ธมฺเม ฯเปฯ สํเฆ ฯเปฯ สิกฺขาย น ปริปูรการี, โส สํเฆ อนุวาทํ ชเนติ, โย โหติ อนุวาโท พหุชนาหิตาย พหุชนา\-สุขาย พหุโน ชนสฺส อนตฺถาย อหิตาย ทุกฺขาย เทวมนุสฺสานํ, \csromanfolio{167}เอวรูปญฺเจ ตุเมฺห อนุวาทมูลํ อชฺฌตฺตํ วา พหิทฺธา วา สมนุปสฺเสยฺยาถ, ตตฺร ตุเมฺห ตสฺเสว ปาปกสฺส อนุวาทมูลสฺส ปหานาย วายเมยฺยาถ, เอวรูปญฺเจ ตุเมฺห อนุวาทมูลํ อชฺฌตฺตํ วา พหิทฺธา วา น สมนุปสฺเสยฺยาถ, ตตฺร ตุเมฺห ตสฺเสว ปาปกสฺส อนุวาทมูลสฺส อายติํ อนวสฺสวาย ปฏิปชฺเชยฺยาถ, เอวเมตสฺส ปาปกสฺส อนุวาทมูลสฺส ปหานํ โหติ, เอวเมตสฺส ปาปกสฺส อนุวาทมูลสฺส อายติํ อนวสฺสโว โหติ.}

\prose{ปุน จปรํ ภิกฺขุ มกฺขี โหติ ปลาสี ฯเปฯ. อิสฺสุกี โหติ มจฺฉรี. สโฐ โหติ มายาวี. ปาปิจฺโฉ โหติ มิจฺฉาทิฏฺฐิ. สนฺทิฏฺฐิ\-ปรามาสี โหติ อาธานคฺคาหี ทุปฺปฏินิสฺสคฺคี, โย โส ภิกฺขุ สนฺทิฏฺฐิ\-ปรามาสี โหติ อาธานคฺคาหี ทุปฺปฏินิสฺสคฺคี, โส สตฺถริปิ อคารโว วิหรติ อปฺปติสฺโส, ธมฺเมปิ อคารโว วิหรติ อปฺปติสฺโส, สํเฆปิ อคารโว วิหรติ อปฺปติสฺโส, สิกฺขายปิ น ปริปูรการี โหติ, โย โส ภิกฺขุ สตฺถริ อคารโว วิหรติ อปฺปติสฺโส, ธมฺเม ฯเปฯ สํเฆ ฯเปฯ สิกฺขาย น ปริปูรการี, โส สํเฆ อนุวาทํ ชเนติ, โย โหติ อนุวาโท พหุชนาหิตาย พหุชนา\-สุขาย พหุโน ชนสฺส อนตฺถาย อหิตาย ทุกฺขาย เทวมนุสฺสานํ, เอวรูปญฺเจ ตุเมฺห อนุวาทมูลํ อชฺฌตฺตํ วา พหิทฺธา วา สมนุปสฺเสยฺยาถ, ตตฺร ตุเมฺห ตสฺเสว ปาปกสฺส อนุวาทมูลสฺส ปหานาย วายเมยฺยาถ, เอวรูปญฺเจ ตุเมฺห อนุวาทมูลํ อชฺฌตฺตํ วา พหิทฺธา วา น สมนุปสฺเสยฺยาถ, ตตฺร ตุเมฺห ตสฺเสว ปาปกสฺส อนุวาทมูลสฺส อายติํ อนวสฺสวาย ปฏิปชฺเชยฺยาถ, เอวเมตสฺส ปาปกสฺส อนุวาทมูลสฺส ปหานํ โหติ, เอวเมตสฺส ปาปกสฺส อนุวาทมูลสฺส อายติํ อนวสฺสโว โหติ. อิมานิ ฉ อนุวาทมูลานิ.}

\proseitem{274}{\csromansymbolfootnote{*}{ที 3. 203; อํ 2. 255 ปิฏฺเฐสุปิ.}ตตฺถ กตเม ฉ สารณียา ธมฺมา. อิธ ภิกฺขุโน เมตฺตํ กายกมฺมํ ปจฺจุปฏฺฐิตํ โหติ สพฺรหฺมจารีสุ อาวิ เจว รโห จ, อยมฺปิ ธมฺโม สารณีโย ปิยกรโณ ครุกรโณ สงฺคหาย อวิวาทาย สามคฺคิยา เอกีภาวาย สํวตฺตติ.}

\prose{ปุน จปรํ ภิกฺขุโน เมตฺตํ วจีกมฺมํ ปจฺจุปฏฺฐิตํ โหติ สพฺรหฺมจารีสุ อาวิ เจว รโห จ, อยมฺปิ ธมฺโม สารณีโย ปิยกรโณ ครุกรโณ สงฺคหาย อวิวาทาย สามคฺคิยา เอกี ภาวาย สํวตฺตติ.}

\prose{\csromanfolio{168}ปุน จปรํ ภิกฺขุโน เมตฺตํ มโนกมฺมํ ปจฺจุปฏฺฐิตํ โหติ สพฺรหฺมจารีสุ อาวิ เจว รโห จ, อยมฺปิ ธมฺโม สารณีโย ปิยกรโณ ครุกรโณ สงฺคหาย อวิวาทาย สามคฺคิยา เอกีภาวาย สํวตฺตติ.}

\prose{ปุน จปรํ ภิกฺขุ เย เต ลาภา ธมฺมิกา ธมฺมลทฺธา อนฺตมโส ปตฺต\-ปริยาปนฺน\-มตฺตมฺปิ, ตถารูเปหิ ลาเภหิ อปฺปฏิวิภตฺต\-โภคี โหติ สีลวนฺเตหิ สพฺรหฺมจารีหิ สาธารณโภคี, อยมฺปิ ธมฺโม สารณีโย ปิยกรโณ ครุกรโณ สงฺคหาย อวิวาทาย สามคฺคิยา เอกีภาวาย สํวตฺตติ.}

\prose{ปุน จปรํ ภิกฺขุ ยานิ ตานิ สีลานิ อขณฺฑานิ อจฺฉิทฺทานิ อสพลานิ อกมฺมาสานิ ภุชิสฺสานิ วิญฺญุ\-ปสตฺถานิ อปรามฏฺฐานิ สมาธิ\-สํวตฺตนิกานิ, ตถารูเปสุ สีเลสุ สีลสามญฺญคโต วิหรติ สพฺรหฺมจารีหิ อาวิ เจว รโห จ, อยมฺปิ ธมฺโม สารณีโย ปิยกรโณ ครุกรโณ สงฺคหาย อวิวาทาย สามคฺคิยา เอกีภาวาย สํวตฺตติ.}

\prose{ปุน จปรํ ภิกฺขุ ยายํ ทิฏฺฐิ อริยา นิยฺยานิกา นิยฺยาติ ตกฺกรสฺส สมฺมา ทุกฺขกฺขยาย, ตถารูปาย ทิฏฺฐิยา ทิฏฺฐิ\-สามญฺญคโต วิหรติ สพฺรหฺมจารีหิ อาวิ เจว รโห จ, อยมฺปิ ธมฺโม สารณีโย ปิยกรโณ ครุกรโณ สงฺคหาย อวิวาทาย สามคฺคิยา เอกีภาวาย สํวตฺตติ. อิเม ฉ สารณียา ธมฺมา.}

\csromanmatikaanchor{mtk.140}
\csromantocmarkref{subsection}{1. ฉอาปตฺติสมุฏฺฐานวาร}{mtk.140}
\bookmark[dest={mtk.140},level=2]{1. ฉอาปตฺติสมุฏฺฐานวาร}
\proseitem{275}{ตตฺถ กตมานิ อฏฺฐารส เภทกร\-วตฺถูนิ. อิธ ภิกฺขุ อธมฺมํ “ธมฺโม”ติ ทีเปติ, ธมฺมํ “อธมฺโม”ติ ทีเปติ, อวินยํ “วินโย”ติ ทีเปติ, วินยํ “อวินโย”ติ ทีเปติ, อภาสิตํ อลปิตํ ตถาคเตน “ภาสิตํ ลปิตํ ตถาคเตนา”ติ ทีเปติ, ภาสิตํ ลปิตํ ตถาคเตน “อภาสิตํ อลปิตํ ตถาคเตนา”ติ ทีเปติ, อนาจิณฺณํ ตถาคเตน “อาจิณฺณํ ตถาคเตนา”ติ ทีเปติ, อาจิณฺณํ ตถาคเตน “อนาจิณฺณํ ตถาคเตนา”ติ ทีเปติ, อปญฺญตฺตํ ตถาคเตน “ปญฺญตฺตํ ตถาคเตนา”ติ ทีเปติ, ปญฺญตฺตํ ตถาคเตน “อปญฺญตฺตํ ตถาคเตนา”ติ ทีเปติ, อาปตฺติํ “อนาปตฺตี”ติ ทีเปติ, อนาปตฺติํ “อาปตฺตี”ติ ทีเปติ, ลหุกํ อาปตฺติํ “ครุกา อาปตฺตี”ติ ทีเปติ, ครุกํ อาปตฺติํ “ลหุกา อาปตฺตี”ติ ทีเปติ, สาวเสสํ อาปตฺติํ “อนวเสสา อาปตฺตี”ติ ทีเปติ, อนวเสสํ อาปตฺติํ “สาวเสสา \csromanfolio{169}อาปตฺตี”ติ ทีเปติ, ทุฏฺฐุลฺลํ อาปตฺติํ “อทุฏฺฐุลฺลา อาปตฺตี”ติ ทีเปติ, อทุฏฺฐุลฺลํ อาปตฺติํ “ทุฏฺฐุลฺลา อาปตฺตี”ติ ทีเปติ. อิมานิ อฏฺฐารส เภทกร\-วตฺถูนิ.}

\prose{\csromansymbolfootnote{*}{วิ 4. 211; วิ 5. 264 ปิฏฺเฐสุปิ.}ตตฺถ กตมานิ จตฺตาริ อธิกรณานิ. วิวาทา\-ธิกรณํ อนุวาทา\-ธิกรณํ อาปตฺตา\-ธิกรณํ กิจฺจา\-ธิกรณํ. อิมานิ จตฺตาริ อธิกรณานิ.}

\prosewithcloser{ตตฺถ กตเม สตฺต สมถา. สมฺมุขาวินโย, สติวินโย, อมูฬฺหวินโย, ปฏิญฺญาต\-กรณํ, เยภุยฺยสิกา, ตสฺสปาปิยสิกา, ติณวตฺถารโก. อิเม สตฺต สมถา.}{กติปุจฺฉาวาโร นิฏฺฐิโต.}
\csromancenter{ตสฺสุทฺทานํ.}
\prose{อาปตฺติ อาปตฺติกฺขนฺธา, วินีตา สตฺตธา ปุน.}

\prose{วินีตาคารวา เจว, คารวา มูลเมว จ.}

\prose{ปุน วินีตา วิปตฺติ, สมุฏฺฐานา วิวาทนา.}

\prose{อนุวาทา สารณียํ, เภทา\-ธิกรเณน จ.}

\prose{สตฺเตว สมถา วุตฺตา, ปทา สตฺตรสา อิเมติ.}
\csromansectionrule

\prose{1. ฉอาปตฺติสมุฏฺฐาน\-วาร}

\proseitem{276}{ปฐเมน อาปตฺติ\-สมุฏฺฐาเนน ปาราชิกํ อาปชฺเชยฺยาติ น หีติ วตฺตพฺพํ. สํฆาทิเสสํ อาปชฺเชยฺยาติ สิยาติ วตฺตพฺพํ. ถุลฺลจฺจยํ อาปชฺเชยฺยาติ สิยาติ วตฺตพฺพํ. ปาจิตฺติยํ อาปชฺเชยฺยาติ สิยาติ วตฺตพฺพํ. ปาฏิเทสนียํ อาปชฺเชยฺยาติ สิยาติ วตฺตพฺพํ. ทุกฺกฏํ อาปชฺเชยฺยาติ สิยาติ วตฺตพฺพํ. ทุพฺภาสิตํ อาปชฺเชยฺยาติ น หีติ วตฺตพฺพํ.}

\prose{ทุติเยน อาปตฺติ\-สมุฏฺฐาเนน ปาราชิกํ อาปชฺเชยฺยาติ น หีติ วตฺตพฺพํ. สํฆาทิเสสํ อาปชฺเชยฺยาติ สิยาติ วตฺตพฺพํ. ถุลฺลจฺจยํ อาปชฺเชยฺยาติ สิยาติ วตฺตพฺพํ. ปาจิตฺติยํ อาปชฺเชยฺยาติ สิยาติ วตฺตพฺพํ. ปาฏิเทสนียํ อาปชฺเชยฺยาติ น หีติ วตฺตพฺพํ. ทุกฺกฏํ อาปชฺเชยฺยาติ สิยาติ วตฺตพฺพํ. ทุพฺภาสิตํ อาปชฺเชยฺยาติ น หีติ วตฺตพฺพํ.}

\prose{\csromanfolio{170}ตติเยน อาปตฺติ\-สมุฏฺฐาเนน ปาราชิกํ อาปชฺเชยฺยาติ น หีติ วตฺตพฺพํ. สํฆาทิเสสํ อาปชฺเชยฺยาติ สิยาติ วตฺตพฺพํ. ถุลฺลจฺจยํ อาปชฺเชยฺยาติ สิยาติ วตฺตพฺพํ. ปาจิตฺติยํ อาปชฺเชยฺยาติ สิยาติ วตฺตพฺพํ. ปาฏิเทสนียํ อาปชฺเชยฺยาติ สิยาติ วตฺตพฺพํ. ทุกฺกฏํ อาปชฺเชยฺยาติ สิยาติ วตฺตพฺพํ. ทุพฺภาสิตํ อาปชฺเชยฺยาติ น หีติ วตฺตพฺพํ.}

\prose{จตุตฺเถน อาปตฺติ\-สมุฏฺฐาเนน ปาราชิกํ อาปชฺเชยฺยาติ สิยาติ วตฺตพฺพํ. สํฆาทิเสสํ อาปชฺเชยฺยาติ สิยาติ วตฺตพฺพํ. ถุลฺลจฺจยํ อาปชฺเชยฺยาติ สิยาติ วตฺตพฺพํ. ปาจิตฺติยํ อาปชฺเชยฺยาติ สิยาติ วตฺตพฺพํ. ปาฏิเทสนียํ อาปชฺเชยฺยาติ สิยาติ วตฺตพฺพํ. ทุกฺกฏํ อาปชฺเชยฺยาติ สิยาติ วตฺตพฺพํ. ทุพฺภาสิตํ อาปชฺเชยฺยาติ น หีติ วตฺตพฺพํ.}

\prose{ปญฺจเมน อาปตฺติ\-สมุฏฺฐาเนน ปาราชิกํ อาปชฺเชยฺยาติ สิยาติ วตฺตพฺพํ. สํฆาทิเสสํ อาปชฺเชยฺยาติ สิยาติ วตฺตพฺพํ. ถุลฺลจฺจยํ อาปชฺเชยฺยาติ สิยาติ วตฺตพฺพํ. ปาจิตฺติยํ อาปชฺเชยฺยาติ สิยาติ วตฺตพฺพํ. ปาฏิเทสนียํ อาปชฺเชยฺยาติ น หีติ วตฺตพฺพํ. ทุกฺกฏํ อาปชฺเชยฺยาติ สิยาติ วตฺตพฺพํ. ทุพฺภาสิตํ อาปชฺเชยฺยาติ สิยาติ วตฺตพฺพํ.}

\prosewithcloserruled{ฉฏฺเฐน อาปตฺติ\-สมุฏฺฐาเนน ปาราชิกํ อาปชฺเชยฺยาติ สิยาติ วตฺตพฺพํ. สํฆาทิเสสํ อาปชฺเชยฺยาติ สิยาติ วตฺตพฺพํ. ถุลฺลจฺจยํ อาปชฺเชยฺยาติ สิยาติ วตฺตพฺพํ. ปาจิตฺติยํ อาปชฺเชยฺยาติ สิยาติ วตฺตพฺพํ. ปาฏิเทสนียํ อาปชฺเชยฺยาติ สิยาติ วตฺตพฺพํ. ทุกฺกฏํ อาปชฺเชยฺยาติ สิยาติ วตฺตพฺพํ. ทุพฺภาสิตํ อาปชฺเชยฺยาติ น หีติ วตฺตพฺพํ.}{ฉ อาปตฺติ\-สมุฏฺฐานวาโร นิฏฺฐิโต ปฐโม.}
\csromanmatikaanchor{mtk.141}
\csromantocmarkref{subsection}{2. กตาปตฺติวาร}{mtk.141}
\bookmark[dest={mtk.141},level=2]{2. กตาปตฺติวาร}
\csromanheader{2}{2. \textbf{กตาปตฺติวาร}}
\proseitem{277}{ปฐเมน อาปตฺติ\-สมุฏฺฐาเนน กติ อาปตฺติโย อาปชฺชติ. ปฐเมน อาปตฺติ\-สมุฏฺฐาเนน ปญฺจ อาปตฺติโย อาปชฺชติ. ภิกฺขุ กปฺปิยสญฺญี สญฺญาจิกาย กุฏิํ กโรติ อเทสิต\-วตฺถุกํ ปมาณา\-ติกฺกนฺตํ สารมฺภํ อปริกฺกมนํ, ปโยเค ทุกฺกฏํ. เอกํ ปิณฺฑํ อนาคเต \csromanfolio{171}อาปตฺติ ถุลฺลจฺจยสฺส. ตสฺมิํ ปิณฺเฑ อาคเต อาปตฺติ สํฆาทิเสสสฺส. ภิกฺขุ กปฺปิยสญฺญี วิกาเล โภชนํ ภุญฺชติ, อาปตฺติ ปาจิตฺติยสฺส. ภิกฺขุ กปฺปิยสญฺญี อญฺญาติกาย ภิกฺขุนิยา อนฺตรฆรํ ปวิฏฺฐาย หตฺถโต ขาทนียํ วา โภชนียํ วา สหตฺถา ปฏิคฺคเหตฺวา ภุญฺชติ, อาปตฺติ ปาฏิเทสนียสฺส. ปฐเมน อาปตฺติ\-สมุฏฺฐาเนน อิมา ปญฺจ อาปตฺติโย อาปชฺชติ.}

\prose{ตา อาปตฺติโย จตุนฺนํ วิปตฺตีนํ กติ วิปตฺติโย ภชนฺติ, สตฺตนฺนํ อาปตฺติ\-กฺขนฺธานํ กติหิ อาปตฺติ\-กฺขนฺเธหิ สงฺคหิตา, ฉนฺนํ อาปตฺติ\-สมุฏฺฐานานํ กติหิ สมุฏฺฐาเนหิ สมุฏฺฐนฺติ, จตุนฺนํ อธิกรณานํ กตมํ อธิกรณํ, สตฺตนฺนํ สมถานํ กติหิ สมเถหิ สมฺมนฺติ. ตา อาปตฺติโย จตุนฺนํ วิปตฺตีนํ เทฺว วิปตฺติโย ภชนฺติ, สิยา สีลวิปตฺติํ, สิยา อาจารวิปตฺติํ. สตฺตนฺนํ อาปตฺติ\-กฺขนฺธานํ ปญฺจหิ อาปตฺติ\-กฺขนฺเธหิ สงฺคหิตา, สิยา สํฆาทิเสสา\-ปตฺติกฺขนฺเธน, สิยา ถุลฺลจฺจยา\-ปตฺติกฺขนฺเธน, สิยา ปาจิตฺติยาปตฺติ\-กฺขนฺเธน, สิยา ปาฏิเทสนียา\-ปตฺติกฺขนฺเธน, สิยา ทุกฺกฏาปตฺติ\-กฺขนฺเธน. ฉนฺนํ อาปตฺติ\-สมุฏฺฐานานํ เอเกน สมุฏฺฐาเนน สมุฏฺฐนฺติ, กายโต สมุฏฺฐนฺติ, น วาจโต, น จิตฺตโต. จตุนฺนํ อธิกรณานํ อาปตฺตา\-ธิกรณํ. สตฺตนฺนํ สมถานํ ตีหิ สมเถหิ สมฺมนฺติ, สิยา สมฺมุขา\-วินเยน จ ปฏิญฺญาต\-กรเณน จ, สิยา สมฺมุขา\-วินเยน จ ติณวตฺถารเกน จ.}

\proseitem{278}{ทุติเยน อาปตฺติ\-สมุฏฺฐาเนน กติ อาปตฺติโย อาปชฺชติ. ทุติเยน อาปตฺติ\-สมุฏฺฐาเนน จตสฺโส อาปตฺติโย อาปชฺชติ. ภิกฺขุ กปฺปิยสญฺญี สมาทิสติ “กุฏิํ เม กโรถา”ติ, ตสฺส กุฏิํ กโรนฺติ อเทสิต\-วตฺถุกํ ปมาณา\-ติกฺกนฺตํ สารมฺภํ อปริกฺกมนํ, ปโยเค ทุกฺกฏํ. เอกํ ปิณฺฑํ อนาคเต อาปตฺติ ถุลฺลจฺจยสฺส. ตสฺมิํ ปิณฺเฑ อาคเต อาปตฺติ สํฆาทิเสสสฺส. ภิกฺขุ กปฺปิยสญฺญี อนุปสมฺปนฺนํ ปทโส ธมฺมํ วาเจติ, อาปตฺติ ปาจิตฺติยสฺส. ทุติเยน อาปตฺติ\-สมุฏฺฐาเนน อิมา จตสฺโส อาปตฺติโย อาปชฺชติ.}

\prose{ตา อาปตฺติโย จตุนฺนํ วิปตฺตีนํ กติ วิปตฺติโย ภชนฺติ ฯเปฯ. สตฺตนฺนํ สมถานํ กติหิ สมเถหิ สมฺมนฺติ. ตา อาปตฺติโย จตุนฺนํ วิปตฺตีนํ เทฺว วิปตฺติโย ภชนฺติ, สิยา สีลวิปตฺติํ, สิยา อาจารวิปตฺติํ. สตฺตนฺนํ \csromanfolio{172}อาปตฺติ\-กฺขนฺธานํ จตูหิ อาปตฺติ\-กฺขนฺเธหิ สงฺคหิตา, สิยา สํฆาทิเสสา\-ปตฺติกฺขนฺเธน, สิยา ถุลฺลจฺจยา\-ปตฺติกฺขนฺเธน, สิยา ปาจิตฺติยาปตฺติ\-กฺขนฺเธน, สิยา ทุกฺกฏาปตฺติ\-กฺขนฺเธน. ฉนฺนํ อาปตฺติ\-สมุฏฺฐานานํ เอเกน สมุฏฺฐาเนน สมุฏฺฐนฺติ, วาจโต สมุฏฺฐนฺติ, น กายโต, น จิตฺตโต. จตุนฺนํ อธิกรณานํ อาปตฺตา\-ธิกรณํ. สตฺตนฺนํ สมถานํ ตีหิ สมเถหิ สมฺมนฺติ, สิยา สมฺมุขา\-วินเยน จ ปฏิญฺญาต\-กรเณน จ, สิยา สมฺมุขา\-วินเยน จ ติณวตฺถารเกน จ.}

\proseitem{279}{ตติเยน อาปตฺติ\-สมุฏฺฐาเนน กติ อาปตฺติโย อาปชฺชติ. ตติเยน อาปตฺติ\-สมุฏฺฐาเนน ปญฺจ อาปตฺติโย อาปชฺชติ. ภิกฺขุ กปฺปิยสญฺญี สํวิทหิตฺวา กุฏิํ กโรติ อเทสิต\-วตฺถุกํ ปมาณา\-ติกฺกนฺตํ สารมฺภํ อปริกฺกมนํ, ปโยเค ทุกฺกฏํ. เอกํ ปิณฺฑํ อนาคเต อาปตฺติ ถุลฺลจฺจยสฺส. ตสฺมิํ ปิณฺเฑ อาคเต อาปตฺติ สํฆาทิเสสสฺส. ภิกฺขุ กปฺปิยสญฺญี ปณีต\-โภชนานิ อตฺตโน อตฺถาย วิญฺญาเปตฺวา ภุญฺชติ, อาปตฺติ ปาจิตฺติยสฺส. ภิกฺขุ กปฺปิยสญฺญี ภิกฺขุนิยา โวสาสนฺติยา น นิวาเรตฺวา ภุญฺชติ, อาปตฺติ ปาฏิเทสนียสฺส. ตติเยน อาปตฺติ\-สมุฏฺฐาเนน อิมา ปญฺจ อาปตฺติโย อาปชฺชติ.}

\prose{ตา อาปตฺติโย จตุนฺนํ วิปตฺตีนํ กติ วิปตฺติโย ภชนฺติ ฯเปฯ. สตฺตนฺนํ สมถานํ กติหิ สมเถหิ สมฺมนฺติ. ตา อาปตฺติโย จตุนฺนํ วิปตฺตีนํ เทฺว วิปตฺติโย ภชนฺติ, สิยา สีลวิปตฺติํ, สิยา อาจารวิปตฺติํ. สตฺตนฺนํ อาปตฺติ\-กฺขนฺธานํ ปญฺจหิ อาปตฺติ\-กฺขนฺเธหิ สงฺคหิตา, สิยา สํฆาทิเสสา\-ปตฺติกฺขนฺเธน, สิยา ถุลฺลจฺจยา\-ปตฺติกฺขนฺเธน, สิยา ปาจิตฺติยาปตฺติ\-กฺขนฺเธน, สิยา ปาฏิเทสนียา\-ปตฺติกฺขนฺเธน, สิยา ทุกฺกฏาปตฺติ\-กฺขนฺเธน. ฉนฺนํ อาปตฺติ\-สมุฏฺฐานานํ เอเกน สมุฏฺฐาเนน สมุฏฺฐนฺติ, กายโต จ วาจโต จ สมุฏฺฐนฺติ, น จิตฺตโต. จตุนฺนํ อธิกรณานํ อาปตฺตา\-ธิกรณํ. สตฺตนฺนํ สมถานํ ตีหิ สมเถหิ สมฺมนฺติ, สิยา สมฺมุขา\-วินเยน จ ปฏิญฺญาต\-กรเณน จ, สิยา สมฺมุขา\-วินเยน จ ติณวตฺถารเกน จ.}

\proseitem{280}{จตุตฺเถน อาปตฺติ\-สมุฏฺฐาเนน กติ อาปตฺติโย อาปชฺชติ. จตุตฺเถน อาปตฺติ\-สมุฏฺฐาเนน ฉ อาปตฺติโย อาปชฺชติ. ภิกฺขุ เมถุนํ ธมฺมํ ปฏิเสวติ, อาปตฺติ ปาราชิกสฺส. ภิกฺขุ อกปฺปิยสญฺญี สญฺญาจิกาย กุฏิํ กโรติ อเทสิต\-วตฺถุกํ ปมาณา\-ติกฺกนฺตํ สารมฺภํ อปริกฺกมนํ, \csromanfolio{173}ปโยเค ทุกฺกฏํ. เอกํ ปิณฺฑํ อนาคเต อาปตฺติ ถุลฺลจฺจยสฺส. ตสฺมิํ ปิณฺเฑ อาคเต อาปตฺติ สํฆาทิเสสสฺส. ภิกฺขุ อกปฺปิยสญฺญี วิกาเล โภชนํ ภุญฺชติ, อาปตฺติ ปาจิตฺติยสฺส. ภิกฺขุ อกปฺปิยสญฺญี อญฺญาติกาย ภิกฺขุนิยา อนฺตรฆรํ ปวิฏฺฐาย หตฺถโต ขาทนียํ วา โภชนียํ วา สหตฺถา ปฏิคฺคเหตฺวา ภุญฺชติ, อาปตฺติ ปาฏิเทสนียสฺส. จตุตฺเถน อาปตฺติ\-สมุฏฺฐาเนน อิมา ฉ อาปตฺติโย อาปชฺชติ.}

\prose{ตา อาปตฺติโย จตุนฺนํ วิปตฺตีนํ กติ วิปตฺติโย ภชนฺติ ฯเปฯ. สตฺตนฺนํ สมถานํ กติหิ สมเถหิ สมฺมนฺติ. ตา อาปตฺติโย จตุนฺนํ วิปตฺตีนํ เทฺว วิปตฺติโย ภชนฺติ, สิยา สีลวิปตฺติํ, สิยา อาจารวิปตฺติํ. สตฺตนฺนํ อาปตฺติ\-กฺขนฺธานํ ฉหิ อาปตฺติ\-กฺขนฺเธหิ สงฺคหิตา, สิยา ปาราชิกาปตฺติ\-กฺขนฺเธน, สิยา สํฆาทิเสสา\-ปตฺติกฺขนฺเธน, สิยา ถุลฺลจฺจยา\-ปตฺติกฺขนฺเธน, สิยา ปาจิตฺติยาปตฺติ\-กฺขนฺเธน, สิยา ปาฏิเทสนียา\-ปตฺติกฺขนฺเธน, สิยา ทุกฺกฏาปตฺติ\-กฺขนฺเธน. ฉนฺนํ อาปตฺติ\-สมุฏฺฐานานํ เอเกน สมุฏฺฐาเนน สมุฏฺฐนฺติ, กายโต จ จิตฺตโต จ สมุฏฺฐนฺติ, น วาจโต. จตุนฺนํ อธิกรณานํ อาปตฺตา\-ธิกรณํ. สตฺตนฺนํ สมถานํ ตีหิ สมเถหิ สมฺมนฺติ, สิยา สมฺมุขา\-วินเยน จ ปฏิญฺญาต\-กรเณน จ, สิยา สมฺมุขา\-วินเยน จ ติณวตฺถารเกน จ.}

\proseitem{281}{ปญฺจเมน อาปตฺติ\-สมุฏฺฐาเนน กติ อาปตฺติโย อาปชฺชติ. ปญฺจเมน อาปตฺติ\-สมุฏฺฐาเนน ฉ อาปตฺติโย อาปชฺชติ. ภิกฺขุ ปาปิจฺโฉ อิจฺฉาปกโต อสนฺตํ อภูตํ อุตฺตริ\-มนุสฺสธมฺมํ อุลฺลปติ, อาปตฺติ ปาราชิกสฺส. ภิกฺขุ อกปฺปิยสญฺญี สมาทิสติ “กุฏิํ เม กโรถา”ติ, ตสฺส กุฏิํ กโรนฺติ อเทสิต\-วตฺถุกํ ปมาณา\-ติกฺกนฺตํ สารมฺภํ อปริกฺกมนํ, ปโยเค ทุกฺกฏํ. เอกํ ปิณฺฑํ อนาคเต อาปตฺติ ถุลฺลจฺจยสฺส. ตสฺมิํ ปิณฺเฑ อาคเต อาปตฺติ สํฆาทิเสสสฺส. ภิกฺขุ อกปฺปิยสญฺญี อนุปสมฺปนฺนํ ปทโส ธมฺมํ วาเจติ, อาปตฺติ ปาจิตฺติยสฺส. น ขุํเสตุกาโม น วมฺเภตุกาโม น มงฺกุกตฺตุกาโม ทวกมฺยตา หีเนน หีนํ วเทติ, อาปตฺติ ทุพฺภาสิตสฺส. ปญฺจเมน อาปตฺติ\-สมุฏฺฐาเนน อิมา ฉ อาปตฺติโย อาปชฺชติ.}

\prose{ตา อาปตฺติโย จตุนฺนํ วิปตฺตีนํ กติ วิปตฺติโย ภชนฺติ ฯเปฯ. สตฺตนฺนํ สมถานํ กติหิ สมเถหิ สมฺมนฺติ. ตา อาปตฺติโย จตุนฺนํ \csromanfolio{174}วิปตฺตีนํ เทฺว วิปตฺติโย ภชนฺติ, สิยา สีลวิปตฺติํ, สิยา อาจารวิปตฺติํ. สตฺตนฺนํ อาปตฺติ\-กฺขนฺธานํ ฉหิ อาปตฺติ\-กฺขนฺเธหิ สงฺคหิตา, สิยา ปาราชิกาปตฺติ\-กฺขนฺเธน, สิยา สํฆาทิเสสา\-ปตฺติกฺขนฺเธน, สิยา ถุลฺลจฺจยา\-ปตฺติกฺขนฺเธน, สิยา ปาจิตฺติยาปตฺติ\-กฺขนฺเธน, สิยา ทุกฺกฏาปตฺติ\-กฺขนฺเธน, สิยา ทุพฺภาสิตา\-ปตฺติกฺขนฺเธน. ฉนฺนํ อาปตฺติ\-สมุฏฺฐานานํ เอเกน สมุฏฺฐาเนน สมุฏฺฐนฺติ, วาจโต จ จิตฺตโต จ สมุฏฺฐนฺติ, น กายโต. จตุนฺนํ อธิกรณานํ อาปตฺตา\-ธิกรณํ. สตฺตนฺนํ สมถานํ ตีหิ สมเถหิ สมฺมนฺติ, สิยา สมฺมุขา\-วินเยน จ ปฏิญฺญาต\-กรเณน จ, สิยา สมฺมุขา\-วินเยน จ ติณวตฺถารเกน จ.}

\proseitem{282}{ฉฏฺเฐน อาปตฺติ\-สมุฏฺฐาเนน กติ อาปตฺติโย อาปชฺชติ. ฉฏฺเฐน อาปตฺติ\-สมุฏฺฐาเนน ฉ อาปตฺติโย อาปชฺชติ. ภิกฺขุ สํวิทหิตฺวา ภณฺฑํ อวหรติ, อาปตฺติ ปาราชิกสฺส. ภิกฺขุ อกปฺปิยสญฺญี สํวิทหิตฺวา กุฏิํ กโรติ อเทสิต\-วตฺถุกํ ปมาณา\-ติกฺกนฺตํ สารมฺภํ อปริกฺกมนํ, ปโยเค ทุกฺกฏํ. เอกํ ปิณฺฑํ อนาคเต อาปตฺติ ถุลฺลจฺจยสฺส. ตสฺมิํ ปิณฺเฑ อาคเต อาปตฺติ สํฆาทิเสสสฺส. ภิกฺขุ อกปฺปิยสญฺญี ปณีต\-โภชนานิ อตฺตโน อตฺถาย วิญฺญาเปตฺวา ภุญฺชติ, อาปตฺติ ปาจิตฺติยสฺส. ภิกฺขุ อกปฺปิยสญฺญี ภิกฺขุนิยา โวสาสนฺติยา น นิวาเรตฺวา ภุญฺชติ, อาปตฺติ ปาฏิเทสนียสฺส. ฉฏฺเฐน อาปตฺติ\-สมุฏฺฐาเนน อิมา ฉ อาปตฺติโย อาปชฺชติ.}

\prose{ตา อาปตฺติโย จตุนฺนํ วิปตฺตีนํ กติ วิปตฺติโย ภชนฺติ, สตฺตนฺนํ อาปตฺติ\-กฺขนฺธานํ กติหิ อาปตฺติ\-กฺขนฺเธหิ สงฺคหิตา, ฉนฺนํ อาปตฺติ\-สมุฏฺฐานานํ กติหิ อาปตฺติ\-สมุฏฺฐาเนหิ สมุฏฺฐนฺติ, จตุนฺนํ อธิกรณานํ กตมํ อธิกรณํ, สตฺตนฺนํ สมถานํ กติหิ สมเถหิ สมฺมนฺติ. ตา อาปตฺติโย จตุนฺนํ วิปตฺตีนํ เทฺว วิปตฺติโย ภชนฺติ, สิยา สีลวิปตฺติํ, สิยา อาจารวิปตฺติํ. สตฺตนฺนํ อาปตฺติ\-กฺขนฺธานํ ฉหิ อาปตฺติ\-กฺขนฺเธหิ สงฺคหิตา, สิยา ปาราชิกาปตฺติ\-กฺขนฺเธน, สิยา สํฆาทิเสสา\-ปตฺติกฺขนฺเธน, สิยา ถุลฺลจฺจยา\-ปตฺติกฺขนฺเธน, สิยา ปาจิตฺติยาปตฺติ\-กฺขนฺเธน, สิยา ปาฏิเทสนียา\-ปตฺติกฺขนฺเธน, สิยา ทุกฺกฏาปตฺติ\-กฺขนฺเธน. ฉนฺนํ อาปตฺติ\-สมุฏฺฐานานํ เอเกน สมุฏฺฐาเนน สมุฏฺฐนฺติ, กายโต จ วาจโต จ จิตฺตโต จ สมุฏฺฐนฺติ. จตุนฺนํ อธิกรณานํ อาปตฺตา\-ธิกรณํ. \csromanfolio{175}สตฺตนฺนํ สมถานํ ตีหิ สมเถหิ สมฺมนฺติ, สิยา สมฺมุขา\-วินเยน จ ปฏิญฺญาต\-กรเณน จ, สิยา สมฺมุขา\-วินเยน จ ติณวตฺถารเกน จาติ.}

\csromanheader{2}{ฉนฺนํ อาปตฺติ\-สมุฏฺฐานานํ}
\nitthitamruled{กตาปตฺติวาโรนิฏฺฐิโต ทุติโย.}
\csromanmatikaanchor{mtk.142}
\csromantocmarkref{subsection}{3. อาปตฺติสมุฏฺฐานคาถา}{mtk.142}
\bookmark[dest={mtk.142},level=2]{3. อาปตฺติสมุฏฺฐานคาถา}
\csromanheader{2}{3. อาปตฺติสมุฏฺฐาน\-คาถา}
\setlength{\csromangathapagewidth}{0pt}
\setlength{\csromangathawakwidth}{0pt}
\csromangathameasuregroup{}{สมุฏฺฐานา กายิกา อนนฺตทสฺสินา, \\ อกฺขาตา โลกหิเตน วิเวกทสฺสินา. \\ อาปตฺติโย เตน สมุฏฺฐิตา กติ, \\ ปุจฺฉามิ ตํ พฺรูหิ วิภงฺคโกวิท. \\ สมุฏฺฐานา กายิกา อนนฺตทสฺสินา, \\ อกฺขาตา โลกหิเตน วิเวกทสฺสินา. \\ อาปตฺติโย เตน สมุฏฺฐิตา ปญฺจ, \\ เอตํ เต อกฺขามิ วิภงฺคโกวิท. \\ สมุฏฺฐานา วาจสิกา อนนฺตทสฺสินา, \\ อกฺขาตา โลกหิเตน วิเวกทสฺสินา. \\ อาปตฺติโย เตน สมุฏฺฐิตา กติ, \\ ปุจฺฉามิ ตํ พฺรูหิ วิภงฺคโกวิท. \\ สมุฏฺฐานา วาจสิกา อนนฺตทสฺสินา, \\ อกฺขาตาโลกหิเตน วิเวกทสฺสินา. \\ อาปตฺติโย เตน สมุฏฺฐิตา จตสฺโส, \\ เอตํ เต อกฺขามิ วิภงฺคโกวิท. \\ สมุฏฺฐานา กายิกา วาจสิกา อนนฺตทสฺสินา, \\ อกฺขาตา โลกหิเตน วิเวกทสฺสินา. \\ อาปตฺติโย เตน สมุฏฺฐิตา กติ, \\ ปุจฺฉามิ ตํ พฺรูหิ วิภงฺคโกวิท. \\ สมุฏฺฐานา กายิกา วาจสิกา อนนฺตทสฺสินา, \\ อกฺขาตา โลกหิเตน วิเวกทสฺสินา. \\ อาปตฺติโย เตน สมุฏฺฐิตา ปญฺจ, \\ เอตํ เต อกฺขามิ วิภงฺคโกวิท. \\ สมุฏฺฐานา กายิกา มานสิกา อนนฺตทสฺสินา, \\ อกฺขาตา โลกหิเตน วิเวกทสฺสินา. \\ อาปตฺติโย เตน สมุฏฺฐิตา กติ, \\ ปุจฺฉามิ ตํ พฺรูหิ วิภงฺคโกวิท. \\ สมุฏฺฐานา กายิกา มานสิกา อนนฺตทสฺสินา, \\ อกฺขาตา โลกหิเตน วิเวกทสฺสินา. \\ อาปตฺติโย เตน สมุฏฺฐิตา ฉ, \\ เอตํ เต อกฺขามิ วิภงฺคโกวิท. \\ สมุฏฺฐานา วาจสิกา มานสิกา อนนฺตทสฺสินา, \\ อกฺขาตา โลกหิเตน วิเวกทสฺสินา. \\ อาปตฺติโย เตน สมุฏฺฐิตา กติ, \\ ปุจฺฉามิ ตํ พฺรูหิ วิภงฺคโกวิท. \\ สมุฏฺฐานา วาจสิกา มานสิกา อนนฺตทสฺสินา, \\ อกฺขาตา โลกหิเตน วิเวกทสฺสินา. \\ อาปตฺติโย เตน สมุฏฺฐิตา ฉ, \\ เอตํ เต อกฺขามิ วิภงฺคโกวิท. \\ สมุฏฺฐานา กายิกา วาจสิกา มานสิกา อนนฺตทสฺสินา, \\ อกฺขาตา โลกหิเตน วิเวกทสฺสินา. \\ อาปตฺติโย เตน สมุฏฺฐิตา กติ, \\ ปุจฺฉามิ ตํ พฺรูหิ วิภงฺคโกวิท. \\ สมุฏฺฐานา กายิกา วาจสิกา มานสิกา อนนฺตทสฺสินา, \\ อกฺขาตา โลกหิเตน วิเวกทสฺสินา. \\ อาปตฺติโย เตน สมุฏฺฐิตา ฉ, \\ เอตํ เต อกฺขามิ วิภงฺคโกวิทาติ.}
\csromangathameasurewak{สมุฏฺฐานา กายิกา อนนฺตทสฺสินา, \\ อกฺขาตา โลกหิเตน วิเวกทสฺสินา. \\ อาปตฺติโย เตน สมุฏฺฐิตา กติ, \\ ปุจฺฉามิ ตํ พฺรูหิ วิภงฺคโกวิท. \\ สมุฏฺฐานา กายิกา อนนฺตทสฺสินา, \\ อกฺขาตา โลกหิเตน วิเวกทสฺสินา. \\ อาปตฺติโย เตน สมุฏฺฐิตา ปญฺจ, \\ เอตํ เต อกฺขามิ วิภงฺคโกวิท. \\ สมุฏฺฐานา วาจสิกา อนนฺตทสฺสินา, \\ อกฺขาตา โลกหิเตน วิเวกทสฺสินา. \\ อาปตฺติโย เตน สมุฏฺฐิตา กติ, \\ ปุจฺฉามิ ตํ พฺรูหิ วิภงฺคโกวิท. \\ สมุฏฺฐานา วาจสิกา อนนฺตทสฺสินา, \\ อกฺขาตาโลกหิเตน วิเวกทสฺสินา. \\ อาปตฺติโย เตน สมุฏฺฐิตา จตสฺโส, \\ เอตํ เต อกฺขามิ วิภงฺคโกวิท. \\ สมุฏฺฐานา กายิกา วาจสิกา อนนฺตทสฺสินา, \\ อกฺขาตา โลกหิเตน วิเวกทสฺสินา. \\ อาปตฺติโย เตน สมุฏฺฐิตา กติ, \\ ปุจฺฉามิ ตํ พฺรูหิ วิภงฺคโกวิท. \\ สมุฏฺฐานา กายิกา วาจสิกา อนนฺตทสฺสินา, \\ อกฺขาตา โลกหิเตน วิเวกทสฺสินา. \\ อาปตฺติโย เตน สมุฏฺฐิตา ปญฺจ, \\ เอตํ เต อกฺขามิ วิภงฺคโกวิท. \\ สมุฏฺฐานา กายิกา มานสิกา อนนฺตทสฺสินา, \\ อกฺขาตา โลกหิเตน วิเวกทสฺสินา. \\ อาปตฺติโย เตน สมุฏฺฐิตา กติ, \\ ปุจฺฉามิ ตํ พฺรูหิ วิภงฺคโกวิท. \\ สมุฏฺฐานา กายิกา มานสิกา อนนฺตทสฺสินา, \\ อกฺขาตา โลกหิเตน วิเวกทสฺสินา. \\ อาปตฺติโย เตน สมุฏฺฐิตา ฉ, \\ เอตํ เต อกฺขามิ วิภงฺคโกวิท. \\ สมุฏฺฐานา วาจสิกา มานสิกา อนนฺตทสฺสินา, \\ อกฺขาตา โลกหิเตน วิเวกทสฺสินา. \\ อาปตฺติโย เตน สมุฏฺฐิตา กติ, \\ ปุจฺฉามิ ตํ พฺรูหิ วิภงฺคโกวิท. \\ สมุฏฺฐานา วาจสิกา มานสิกา อนนฺตทสฺสินา, \\ อกฺขาตา โลกหิเตน วิเวกทสฺสินา. \\ อาปตฺติโย เตน สมุฏฺฐิตา ฉ, \\ เอตํ เต อกฺขามิ วิภงฺคโกวิท. \\ สมุฏฺฐานา กายิกา วาจสิกา มานสิกา อนนฺตทสฺสินา, \\ อกฺขาตา โลกหิเตน วิเวกทสฺสินา. \\ อาปตฺติโย เตน สมุฏฺฐิตา กติ, \\ ปุจฺฉามิ ตํ พฺรูหิ วิภงฺคโกวิท. \\ สมุฏฺฐานา กายิกา วาจสิกา มานสิกา อนนฺตทสฺสินา, \\ อกฺขาตา โลกหิเตน วิเวกทสฺสินา. \\ อาปตฺติโย เตน สมุฏฺฐิตา ฉ, \\ เอตํ เต อกฺขามิ วิภงฺคโกวิทาติ.}
\setlength{\csromangathaleftcol}{0pt}
\csromangathagroup{\csromangathastanza{\csromangathawak{\csromangathaitemline{283}{สมุฏฺฐานา กายิกา อนนฺตทสฺสินา,} \\ \csromangathacontline{อกฺขาตา โลกหิเตน วิเวกทสฺสินา.} \\ \csromangathacontline{อาปตฺติโย เตน สมุฏฺฐิตา กติ,} \\ \csromangathacontline{ปุจฺฉามิ ตํ พฺรูหิ วิภงฺคโกวิท.} \\ \csromangathacontline{สมุฏฺฐานา กายิกา อนนฺตทสฺสินา,} \\ \csromangathacontline{อกฺขาตา โลกหิเตน วิเวกทสฺสินา.} \\ \csromangathacontline{อาปตฺติโย เตน สมุฏฺฐิตา ปญฺจ,} \\ \csromangathacontline{เอตํ เต อกฺขามิ วิภงฺคโกวิท.} \\ \csromangathacontline{สมุฏฺฐานา วาจสิกา อนนฺตทสฺสินา,} \\ \csromangathacontline{อกฺขาตา โลกหิเตน วิเวกทสฺสินา.} \\ \csromangathacontline{อาปตฺติโย เตน สมุฏฺฐิตา กติ,} \\ \csromangathacontline{ปุจฺฉามิ ตํ พฺรูหิ วิภงฺคโกวิท.} \\ \csromangathacontline{สมุฏฺฐานา วาจสิกา อนนฺตทสฺสินา,} \\ \csromangathacontline{อกฺขาตาโลกหิเตน วิเวกทสฺสินา.} \\ \csromangathacontline{อาปตฺติโย เตน สมุฏฺฐิตา จตสฺโส,} \\ \csromangathacontline{เอตํ เต อกฺขามิ วิภงฺคโกวิท.} \\ \csromangathacontline{สมุฏฺฐานา กายิกา วาจสิกา อนนฺตทสฺสินา,} \\ \csromangathacontline{อกฺขาตา โลกหิเตน วิเวกทสฺสินา.} \\ \csromangathacontline{อาปตฺติโย เตน สมุฏฺฐิตา กติ,} \\ \csromangathacontline{ปุจฺฉามิ ตํ พฺรูหิ วิภงฺคโกวิท.}}}\csromangathastanza{\csromangathawak{\csromanfolio{176}สมุฏฺฐานา กายิกา วาจสิกา อนนฺตทสฺสินา, \\ อกฺขาตา โลกหิเตน วิเวกทสฺสินา. \\ อาปตฺติโย เตน สมุฏฺฐิตา ปญฺจ, \\ เอตํ เต อกฺขามิ วิภงฺคโกวิท.}}\csromangathastanza{\csromangathawak{สมุฏฺฐานา กายิกา มานสิกา อนนฺตทสฺสินา, \\ อกฺขาตา โลกหิเตน วิเวกทสฺสินา. \\ อาปตฺติโย เตน สมุฏฺฐิตา กติ, \\ ปุจฺฉามิ ตํ พฺรูหิ วิภงฺคโกวิท.}}\csromangathastanza{\csromangathawak{สมุฏฺฐานา กายิกา มานสิกา อนนฺตทสฺสินา, \\ อกฺขาตา โลกหิเตน วิเวกทสฺสินา. \\ อาปตฺติโย เตน สมุฏฺฐิตา ฉ, \\ เอตํ เต อกฺขามิ วิภงฺคโกวิท.}}\csromangathastanza{\csromangathawak{สมุฏฺฐานา วาจสิกา มานสิกา อนนฺตทสฺสินา, \\ อกฺขาตา โลกหิเตน วิเวกทสฺสินา. \\ อาปตฺติโย เตน สมุฏฺฐิตา กติ, \\ ปุจฺฉามิ ตํ พฺรูหิ วิภงฺคโกวิท.}}\csromangathastanza{\csromangathawak{สมุฏฺฐานา วาจสิกา มานสิกา อนนฺตทสฺสินา, \\ อกฺขาตา โลกหิเตน วิเวกทสฺสินา. \\ อาปตฺติโย เตน สมุฏฺฐิตา ฉ, \\ เอตํ เต อกฺขามิ วิภงฺคโกวิท.}}\csromangathastanza{\csromangathawak{สมุฏฺฐานา กายิกา วาจสิกา มานสิกา อนนฺตทสฺสินา, \\ อกฺขาตา โลกหิเตน วิเวกทสฺสินา. \\ อาปตฺติโย เตน สมุฏฺฐิตา กติ, \\ ปุจฺฉามิ ตํ พฺรูหิ วิภงฺคโกวิท.}}\csromangathastanza{\csromangathawak{สมุฏฺฐานา กายิกา วาจสิกา มานสิกา อนนฺตทสฺสินา, \\ อกฺขาตา โลกหิเตน วิเวกทสฺสินา. \\ อาปตฺติโย เตน สมุฏฺฐิตา ฉ, \\ เอตํ เต อกฺขามิ วิภงฺคโกวิทาติ.}}}

\vspace{\tipitakaparskip}
\csromancenter{อาปตฺติสมุฏฺฐาน\-คาถา นิฏฺฐิตา ตติยา.}
\csromanmatikaanchor{mtk.143}
\csromantocmarkref{subsection}{4. วิปตฺติปจฺจยวาร}{mtk.143}
\bookmark[dest={mtk.143},level=2]{4. วิปตฺติปจฺจยวาร}
\csromanheader{2}{4. วิปตฺติ\-ปจฺจยวาร}
\proseitem{284}{\csromanfolio{177}สีลวิปตฺติ\-ปจฺจยา กติ อาปตฺติโย อาปชฺชติ. สีลวิปตฺติ\-ปจฺจยา จตสฺโส อาปตฺติโย อาปชฺชติ. ภิกฺขุนี ชานํ ปาราชิกํ ธมฺมํ ปฏิจฺฉาเทติ, อาปตฺติ ปาราชิกสฺส. เวมติกา ปฏิจฺฉาเทติ, อาปตฺติ ถุลฺลจฺจยสฺส. ภิกฺขุ สํฆาทิเสสํ ปฏิจฺฉาเทติ, อาปตฺติ ปาจิตฺติยสฺส. อตฺตโน ทุฏฺฐุลฺลํ อาปตฺติํ ปฏิจฺฉาเทติ, อาปตฺติ ทุกฺกฏสฺส. สีลวิปตฺติ\-ปจฺจยา อิมา จตสฺโส อาปตฺติโย อาปชฺชติ.}

\prose{ตา อาปตฺติโย จตุนฺนํ วิปตฺตีนํ กติ วิปตฺติโย ภชนฺติ ฯเปฯ. สตฺตนฺนํ สมถานํ กติหิ สมเถหิ สมฺมนฺติ. ตา อาปตฺติโย จตุนฺนํ วิปตฺตีนํ เทฺว วิปตฺติโย ภชนฺติ, สิยา สีลวิปตฺติํ, สิยา อาจารวิปตฺติํ. สตฺตนฺนํ อาปตฺติ\-กฺขนฺธานํ จตูหิ อาปตฺติ\-กฺขนฺเธหิ สงฺคหิตา, สิยา ปาราชิกาปตฺติ\-กฺขนฺเธน, สิยา ถุลฺลจฺจยา\-ปตฺติกฺขนฺเธน, สิยา ปาจิตฺติยาปตฺติ\-กฺขนฺเธน, สิยา ทุกฺกฏาปตฺติ\-กฺขนฺเธน. ฉนฺนํ อาปตฺติ\-สมุฏฺฐานานํ เอเกน สมุฏฺฐาเนน สมุฏฺฐนฺติ, กายโต จ วาจโต จ จิตฺตโต จ สมุฏฺฐนฺติ. จตุนฺนํ อธิกรณานํ อาปตฺตา\-ธิกรณํ. สตฺตนฺนํ สมถานํ ตีหิ สมเถหิ สมฺมนฺติ, สิยา สมฺมุขา\-วินเยน จ ปฏิญฺญาต\-กรเณน จ, สิยา สมฺมุขา\-วินเยน จ ติณวตฺถารเกน จ.}

\proseitem{285}{อาจารวิปตฺติ\-ปจฺจยา กติ อาปตฺติโย อาปชฺชติ. อาจารวิปตฺติ\-ปจฺจยา เอกํ อาปตฺติํ อาปชฺชติ. อาจารวิปตฺติํ ปฏิจฺฉาเทติ, อาปตฺติ ทุกฺกฏสฺส. อาจารวิปตฺติ\-ปจฺจยา อิมํ เอกํ อาปตฺติํ อาปชฺชติ.}

\prose{สา อาปตฺติ จตุนฺนํ วิปตฺตีนํ กติ วิปตฺติโย ภชติ ฯเปฯ. สตฺตนฺนํ สมถานํ กติหิ สมเถหิ สมฺมติ. สา อาปตฺติ จตุนฺนํ วิปตฺตีนํ เอกํ วิปตฺติํ ภชติ, อาจารวิปตฺติํ. สตฺตนฺนํ อาปตฺติ\-กฺขนฺธานํ เอเกน อาปตฺติ\-กฺขนฺเธน สงฺคหิตา, ทุกฺกฏาปตฺติ\-กฺขนฺเธน. ฉนฺนํ อาปตฺติ\-สมุฏฺฐานานํ เอเกน สมุฏฺฐาเนน สมุฏฺฐาติ, กายโต จ วาจโต จ จิตฺตโต จ สมุฏฺฐาติ. จตุนฺนํ อธิกรณานํ อาปตฺตา\-ธิกรณํ. สตฺตนฺนํ สมถานํ ตีหิ สมเถหิ สมฺมติ, สิยา สมฺมุขา\-วินเยน จ ปฏิญฺญาต\-กรเณน จ, สิยา สมฺมุขา\-วินเยน จ ติณวตฺถารเกน จ.}

\proseitem{286}{\csromanfolio{178}ทิฏฺฐิวิปตฺติ\-ปจฺจยา กติ อาปตฺติโย อาปชฺชติ. ทิฏฺฐิวิปตฺติ\-ปจฺจยา เทฺว อาปตฺติโย อาปชฺชติ. ปาปิกาย ทิฏฺฐิยา ยาวตติยํ สมนุภาสนาย น ปฏินิสฺสชฺชติ, ญตฺติยา ทุกฺกฏํ\csromansharedfootnote{a4c3e164029c2cbd}{ญตฺติยา ทุกฺกฏํ, ทฺวีหิ กมฺมวาจาหิ ทุกฺกฏา (สฺยา, ก)}. กมฺมวาจา\-ปริโยสาเน อาปตฺติ ปาจิตฺติยสฺส. ทิฏฺฐิวิปตฺติ\-ปจฺจยา อิมา เทฺว อาปตฺติโย อาปชฺชติ.}

\prose{ตา อาปตฺติโย จตุนฺนํ วิปตฺตีนํ กติ วิปตฺติโย ภชนฺติ ฯเปฯ. สตฺตนฺนํ สมถานํ กติหิ สมเถหิ สมฺมนฺติ. ตา อาปตฺติโย จตุนฺนํ วิปตฺตีนํ เอกํ วิปตฺติํ ภชนฺติ, อาจารวิปตฺติํ. สตฺตนฺนํ อาปตฺติ\-กฺขนฺธานํ ทฺวีหิ อาปตฺติ\-กฺขนฺเธหิ สงฺคหิตา, สิยา ปาจิตฺติยาปตฺติ\-กฺขนฺเธน, สิยา ทุกฺกฏาปตฺติ\-กฺขนฺเธน. ฉนฺนํ อาปตฺติ\-สมุฏฺฐานานํ เอเกน สมุฏฺฐาเนน สมุฏฺฐนฺติ, กายโต จ วาจโต จ จิตฺตโต จ สมุฏฺฐนฺติ. จตุนฺนํ อธิกรณานํ อาปตฺตา\-ธิกรณํ. สตฺตนฺนํ สมถานํ ตีหิ สมเถหิ สมฺมนฺติ, สิยา สมฺมุขา\-วินเยน จ ปฏิญฺญาต\-กรเณน จ, สิยา สมฺมุขา\-วินเยน จ ติณวตฺถารเกน จ.}

\proseitem{287}{อาชีววิปตฺติ\-ปจฺจยา กติ อาปตฺติโย อาปชฺชติ. อาชีววิปตฺติ\-ปจฺจยา ฉ อาปตฺติโย อาปชฺชติ. อาชีวเหตุ อาชีวการณา ปาปิจฺโฉ อิจฺฉาปกโต อสนฺตํ อภูตํ อุตฺตริ\-มนุสฺสธมฺมํ อุลฺลปติ, อาปตฺติ ปาราชิกสฺส. อาชีวเหตุ อาชีวการณา สญฺจริตฺตํ สมาปชฺชติ, อาปตฺติ สํฆาทิเสสสฺส. อาชีวเหตุ อาชีวการณา “โย เต วิหาเร วสติ, โส ภิกฺขุ อรหา”ติ ภณติ, ปฏิวิชานนฺตสฺส อาปตฺติ ถุลฺลจฺจยสฺส. อาชีวเหตุ อาชีวการณา ภิกฺขุ ปณีต\-โภชนานิ อตฺตโน อตฺถาย วิญฺญาเปตฺวา ภุญฺชติ, อาปตฺติ ปาจิตฺติยสฺส. อาชีวเหตุ อาชีวการณา ภิกฺขุนี ปณีต\-โภชนานิ อตฺตโน อตฺถาย วิญฺญาเปตฺวา ภุญฺชติ, อาปตฺติ ปาฏิเทสนียสฺส. อาชีวเหตุ อาชีวการณา สูปํ วา โอทนํ วา อคิลาโน อตฺตโน อตฺถาย วิญฺญาเปตฺวา ภุญฺชติ, อาปตฺติ ทุกฺกฏสฺส. อาชีววิปตฺติ\-ปจฺจยา อิมา ฉ อาปตฺติโย อาปชฺชติ.}

\prosewithcloserruled{ตา อาปตฺติโย จตุนฺนํ วิปตฺตีนํ กติ วิปตฺติโย ภชนฺติ ฯเปฯ. สตฺตนฺนํ สมถานํ กติหิ สมเถหิ สมฺมนฺติ. ตา อาปตฺติโย จตุนฺนํ วิปตฺตีนํ \csromanfolio{179}เทฺว วิปตฺติโย ภชนฺติ, สิยา สีลวิปตฺติํ, สิยา อาจารวิปตฺติํ. สตฺตนฺนํ อาปตฺติ\-กฺขนฺธานํ ฉหิ อาปตฺติ\-กฺขนฺเธหิ สงฺคหิตา, สิยา ปาราชิกาปตฺติ\-กฺขนฺเธน, สิยา สํฆาทิเสสาปตฺติกฺขานฺเธน, สิยา ถุลฺลจฺจยา\-ปตฺติกฺขนฺเธน, สิยา ปาจิตฺติยาปตฺติ\-กฺขนฺเธน, สิยา ปาฏิเทสนียา\-ปตฺติกฺขนฺเธน, สิยา ทุกฺกฏาปตฺติ\-กฺขนฺเธน. ฉนฺนํ อาปตฺติ\-สมุฏฺฐานานํ ฉหิ สมุฏฺฐาเนหิ สมุฏฺฐนฺติ, สิยา กายโต สมุฏฺฐนฺติ, น วาจโต, น จิตฺตโต, สิยา วาจโต สมุฏฺฐนฺติ, น กายโต, น จิตฺตโต, สิยา กายโต จ วาจโต จ สมุฏฺฐนฺติ, น จิตฺตโต, สิยา กายโต จ จิตฺตโต จ สมุฏฺฐนฺติ, น วาจโต, สิยา วาจโต จ จิตฺตโต จ สมุฏฺฐนฺติ, น กายโต, สิยา กายโต จ วาจโต จ จิตฺตโต จ สมุฏฺฐนฺติ. จตุนฺนํ อธิกรณานํ อาปตฺตา\-ธิกรณํ. สตฺตนฺนํ สมถานํ ตีหิ สมเถหิ สมฺมนฺติ, สิยา สมฺมุขา\-วินเยน จ ปฏิญฺญาต\-กรเณน จ, สิยา สมฺมุขา\-วินเยน จ ติณวตฺถารเกน จ.}{วิปตฺติ\-ปจฺจยวาโร นิฏฺฐิโต จตุตฺโถ.}
\csromanmatikaanchor{mtk.144}
\csromantocmarkref{subsection}{5. อธิกรณปจฺจยวาร}{mtk.144}
\bookmark[dest={mtk.144},level=2]{5. อธิกรณปจฺจยวาร}
\csromanheader{2}{5. อธิกรณปจฺจย\-วาร}
\proseitem{288}{วิวาทาธิกรณ\-ปจฺจยา กติ อาปตฺติโย อาปชฺชติ. วิวาทาธิกรณ\-ปจฺจยา เทฺว อาปตฺติโย อาปชฺชติ. อุปสมฺปนฺนํ โอมสติ, อาปตฺติ ปาจิตฺติยสฺส. อนุปสมฺปนฺนํ โอมสติ, อาปตฺติ ทุกฺกฏสฺส. วิวาทาธิกรณ\-ปจฺจยา อิมา เทฺว อาปตฺติโย อาปชฺชติ.}

\prose{ตา อาปตฺติโย จตุนฺนํ วิปตฺตีนํ กติ วิปตฺติโย ภชนฺติ ฯเปฯ. สตฺตนฺนํ สมถานํ กติหิ สมเถหิ สมฺมนฺติ. ตา อาปตฺติโย จตุนฺนํ วิปตฺตีนํ เอกํ วิปตฺติํ ภชนฺติ, อาจารวิปตฺติํ. สตฺตนฺนํ อาปตฺติ\-กฺขนฺธานํ ทฺวีหิ อาปตฺติ\-กฺขนฺเธหิ สงฺคหิตา, สิยา ปาจิตฺติยาปตฺติ\-กฺขนฺเธน, สิยา ทุกฺกฏาปตฺติ\-กฺขนฺเธน. ฉนฺนํ อาปตฺติ\-สมุฏฺฐานานํ ตีหิ สมุฏฺฐาเนหิ สมุฏฺฐนฺติ, สิยา กายโต จ จิตฺตโต จ สมุฏฺฐนฺติ, น วาจโต, สิยา วาจโต จ จิตฺตโต จ สมุฏฺฐนฺติ, น กายโต, สิยา กายโต จ วาจโต จ \csromanfolio{180}จิตฺตโต จ สมุฏฺฐนฺติ. จตุนฺนํ อธิกรณานํ อาปตฺตา\-ธิกรณํ. สตฺตนฺนํ สมถานํ ตีหิ สมเถหิ สมฺมนฺติ, สิยา สมฺมุขา\-วินเยน จ ปฏิญฺญาต\-กรเณน จ, สิยา สมฺมุขา\-วินเยน จ ติณวตฺถารเกน จ.}

\proseitem{289}{อนุวาทาธิกรณ\-ปจฺจยา กติ อาปตฺติโย อาปชฺชติ. อนุวาทาธิกรณ\-ปจฺจยา ติสฺโส อาปตฺติโย อาปชฺชติ. ภิกฺขุํ อมูลเกน ปาราชิเกน ธมฺเมน อนุทฺธํเสติ, อาปตฺติ สํฆาทิเสสสฺส. อมูลเกน สํฆาทิเสเสน อนุทฺธํเสติ, อาปตฺติ ปาจิตฺติยสฺส. อมูลิกาย อาจารวิปตฺติยา อนุทฺธํเสติ, อาปตฺติ ทุกฺกฏสฺส. อนุวาทาธิกรณ\-ปจฺจยา อิมา ติสฺโส อาปตฺติโย อาปชฺชติ.}

\prose{ตา อาปตฺติโย จตุนฺนํ วิปตฺตีนํ กติ วิปตฺติโย ภชนฺติ ฯเปฯ. สตฺตนฺนํ สมถานํ กติหิ สมเถหิ สมฺมนฺติ. ตา อาปตฺติโย จตุนฺนํ วิปตฺตีนํ เทฺว วิปตฺติโย ภชนฺติ, สิยา สีลวิปตฺติํ, สิยา อาจารวิปตฺติํ. สตฺตนฺนํ อาปตฺติ\-กฺขนฺธานํ ตีหิ อาปตฺติ\-กฺขนฺเธหิ สงฺคหิตา, สิยา สํฆาทิเสสา\-ปตฺติกฺขนฺเธน, สิยา ปาจิตฺติยาปตฺติ\-กฺขนฺเธน, สิยา ทุกฺกฏาปตฺติ\-กฺขนฺเธน. ฉนฺนํ อาปตฺติ\-สมุฏฺฐานานํ ตีหิ สมุฏฺฐาเนหิ สมุฏฺฐนฺติ, สิยา กายโต จ จิตฺตโต จ สมุฏฺฐนฺติ, น วาจโต, สิยา วาจโต จ จิตฺตโต จ สมุฏฺฐนฺติ, น กายโต, สิยา กายโต จ วาจโต จ จิตฺตโต จ สมุฏฺฐนฺติ. จตุนฺนํ อธิกรณานํ อาปตฺตา\-ธิกรณํ. สตฺตนฺนํ สมถานํ ตีหิ สมเถหิ สมฺมนฺติ, สิยา สมฺมุขา\-วินเยน จ ปฏิญฺญาต\-กรเณน จ, สิยา สมฺมุขา\-วินเยน จ ติณวตฺถารเกน จ.}

\proseitem{290}{อาปตฺตาธิกรณ\-ปจฺจยา กติ อาปตฺติโย อาปชฺชติ. อาปตฺตาธิกรณ\-ปจฺจยา จตสฺโส อาปตฺติโย อาปชฺชติ. ภิกฺขุนี ชานํ ปาราชิกํ ธมฺมํ ปฏิจฺฉาเทติ, อาปตฺติ ปาราชิกสฺส. เวมติกา ปฏิจฺฉาเทติ, อาปตฺติ ถุลฺลจฺจยสฺส. ภิกฺขุ สํฆาทิเสสํ ปฏิจฺฉาเทติ, อาปตฺติ ปาจิตฺติยสฺส. อาจารวิปตฺติํ ปฏิจฺฉาเทติ, อาปตฺติ ทุกฺกฏสฺส. อาปตฺตาธิกรณ\-ปจฺจยา อิมา จตสฺโส อาปตฺติโย อาปชฺชติ.}

\prose{ตา อาปตฺติโย จตุนฺนํ วิปตฺตีนํ กติ วิปตฺติโย ภชนฺติ ฯเปฯ. สตฺตนฺนํ สมถานํ กติหิ สมเถหิ สมฺมนฺติ. ตา อาปตฺติโย จตุนฺนํ วิปตฺตีนํ เทฺว วิปตฺติโย ภชนฺติ, สิยา สีลวิปตฺติํ, สิยา อาจารวิปตฺติํ. \csromanfolio{181}สตฺตนฺนํ อาปตฺติ\-กฺขนฺธานํ จตูหิ อาปตฺติ\-กฺขนฺเธหิ สงฺคหิตา, สิยา ปาราชิกาปตฺติ\-กฺขนฺเธน, สิยา ถุลฺลจฺจยา\-ปตฺติกฺขนฺเธน, สิยา ปาจิตฺติยา ปตฺติกฺขนฺเธน, สิยา ทุกฺกฏาปตฺติ\-กฺขนฺเธน. ฉนฺนํ อาปตฺติ\-สมุฏฺฐานานํ เอเกน สมุฏฺฐาเนน สมุฏฺฐนฺติ, กายโต จ วาจโต จ จิตฺตโต จ สมุฏฺฐนฺติ. จตุนฺนํ อธิกรณานํ อาปตฺตา\-ธิกรณํ. สตฺตนฺนํ สมถานํ ตีหิ สมเถหิ สมฺมนฺติ, สิยา สมฺมุขา\-วินเยน จ ปฏิญฺญาต\-กรเณน จ, สิยา สมฺมุขา\-วินเยน จ ติณวตฺถารเกน จ.}

\proseitem{291}{กิจฺจาธิกรณ\-ปจฺจยา กติ อาปตฺติโย อาปชฺชติ. กิจฺจาธิกรณ\-ปจฺจยา ปญฺจ อาปตฺติโย อาปชฺชติ. อุกฺขิตฺตานุวตฺติกา ภิกฺขุนี ยาวตติยํ สมนุภาสนาย น ปฏินิสฺสชฺชติ, ญตฺติยา ทุกฺกฏํ. ทฺวีหิ กมฺมวาจาหิ ถุลฺลจฺจยา. กมฺมวาจา\-ปริโยสาเน อาปตฺติ ปาราชิกสฺส. เภทกา\-นุวตฺตกา ภิกฺขู ยาวตติยํ สมนุภาสนาย น ปฏินิสฺสชฺชนฺติ, อาปตฺติ สํฆาทิเสสสฺส. ปาปิกาย ทิฏฺฐิยา ยาวตติยํ สมนุภาสนาย น ปฏินิสฺสชฺชติ, อาปตฺติ ปาจิตฺติยสฺส. กิจฺจาธิกรณ\-ปจฺจยา อิมา ปญฺจ อาปตฺติโย อาปชฺชติ.}

\prose{ตา อาปตฺติโย จตุนฺนํ วิปตฺตีนํ กติ วิปตฺติโย ภชนฺติ ฯเปฯ. สตฺตนฺนํ สมถานํ กติหิ สมเถหิ สมฺมนฺติ. ตา อาปตฺติโย จตุนฺนํ วิปตฺตีนํ เทฺว วิปตฺติโย ภชนฺติ, สิยา สีลวิปตฺติํ, สิยา อาจารวิปตฺติํ. สตฺตนฺนํ อาปตฺติ\-กฺขนฺธานํ ปญฺจหิ อาปตฺติ\-กฺขนฺเธหิ สงฺคหิตา, สิยา ปาราชิกาปตฺติ\-กฺขนฺเธน, สิยา สํฆาทิเสสา\-ปตฺติกฺขนฺเธน, สิยา ถุลฺลจฺจยา\-ปตฺติกฺขนฺเธน, สิยา ปาจิตฺติยาปตฺติ\-กฺขนฺเธน, สิยา ทุกฺกฏาปตฺติ\-กฺขนฺเธน. ฉนฺนํ อาปตฺติ\-สมุฏฺฐานานํ เอเกน สมุฏฺฐาเนน สมุฏฺฐนฺติ, กายโต จ วาจโต จ จิตฺตโต จ สมุฏฺฐนฺติ, จตุนฺนํ อธิกรณานํ อาปตฺตา\-ธิกรณํ. สตฺตนฺนํ สมถานํ ตีหิ สมเถหิ สมฺมนฺติ, สิยา สมฺมุขา\-วินเยน จ ปฏิญฺญาต\-กรเณน จ, สิยา สมฺมุขา\-วินเยน จ ติณวตฺถารเกน จ.}

\prose{ฐเปตฺวา สตฺต อาปตฺติโย สตฺต อาปตฺติกฺขนฺเธ อวเสสา อาปตฺติโย จตุนฺนํ วิปตฺตีนํ กติ วิปตฺติโย ภชนฺติ, สตฺตนฺนํ อาปตฺติ\-กฺขนฺธานํ กติหิ อาปตฺติ\-กฺขนฺเธหิ สงฺคหิตา, ฉนฺนํ อาปตฺติ\-สมุฏฺฐานานํ กติหิ สมุฏฺฐาเนหิ สมุฏฺฐนฺติ, จตุนฺนํ อธิกรณานํ กตมํ อธิกรณํ, สตฺตนฺนํ สมถานํ กติหิ สมเถหิ สมฺมนฺติ. ฐเปตฺวา สตฺต อาปตฺติโย \csromanfolio{182}สตฺต อาปตฺติกฺขนฺเธ อวเสสา อาปตฺติโย จตุนฺนํ วิปตฺตีนํ น กตมํ วิปตฺติํ ภชนฺติ, สตฺตนฺนํ อาปตฺติ\-กฺขนฺธานํ น กตเมน อาปตฺติ\-กฺขนฺเธน สงฺคหิตา, ฉนฺนํ อาปตฺติ\-สมุฏฺฐานานํ น กตเมน อาปตฺติ\-สมุฏฺฐาเนน สมุฏฺฐนฺติ, จตุนฺนํ อธิกรณานํ น กตมํ อธิกรณํ, สตฺตนฺนํ สมถานํ น กตเมน สมเถน สมฺมนฺติ. ตํ กิสฺส เหตุ, ฐเปตฺวา สตฺต อาปตฺติโย สตฺต อาปตฺติกฺขนฺเธ นตฺถญฺญา อาปตฺติโยติ.}

\vspace{\tipitakaparskip}
\csromancenter{อธิกรณปจฺจย\-วาโร นิฏฺฐิโต ปญฺจโม.}
\nitthitam{อนฺตรเปยฺยาลํ\csromansharedfootnote{f28fe478ea008160}{อนนฺตรเปยฺยาลํ (สี, สฺยา)} นิฏฺฐิตํ.}
\csromancenter{ตสฺสุทฺทานํ.}
\setlength{\csromangathapagewidth}{0pt}
\setlength{\csromangathawakwidth}{0pt}
\csromangathameasuregroup{กติปุจฺฉา สมุฏฺฐานา, \\ สมุฏฺฐานา วิปตฺติ จ,}{\csromangathabat{กติปุจฺฉา สมุฏฺฐานา,}{กตาปตฺติ ตเถว จ.} \\ \csromangathabat{สมุฏฺฐานา วิปตฺติ จ,}{ตถาธิกรเณน จาติ.}}
\csromangathasetleft{กติปุจฺฉา สมุฏฺฐานา, \\ สมุฏฺฐานา วิปตฺติ จ,}
\csromangathagroup{\csromangathastanza{\csromangathabat{กติปุจฺฉา สมุฏฺฐานา,}{กตาปตฺติ ตเถว จ.} \\ \csromangathabat{สมุฏฺฐานา วิปตฺติ จ,}{ตถาธิกรเณน จาติ.}}}

\csromanmatikaanchor{mtk.145}
\csromantocmarknopage{section}{สมถเภท}{mtk.145}
\bookmark[dest={mtk.145},level=1]{สมถเภท}
\csromanheader{1}{\textbf{สมถเภท}}
\csromanmatikaanchor{mtk.146}
\csromantocmarkref{subsection}{6. อธิกรนปริยายวาร}{mtk.146}
\bookmark[dest={mtk.146},level=2]{6. อธิกรนปริยายวาร}
\csromanheader{2}{6. อธิกรณ\-ปริยาย\-วาร}
\proseitem{292}{\csromanfolio{183}วิวาทา\-ธิกรณสฺส กิํ ปุพฺพงฺคมํ, กติ ฐานานิ, กติ วตฺถูนิ, กติ ภูมิโย, กติ เหตู, กติ มูลานิ, กติหากาเรหิ วิวทติ, วิวาทา\-ธิกรณํ กติหิ สมเถหิ สมฺมติ.}

\prose{อนุวาทา\-ธิกรณสฺส กิํ ปุพฺพงฺคมํ, กติ ฐานานิ, กติ วตฺถูนิ, กติ ภูมิโย, กติ เหตู, กติ มูลานิ, กติหากาเรหิ อนุวทติ, อนุวาทา\-ธิกรณํ กติหิ สมเถหิ สมฺมติ.}

\prose{อาปตฺตา\-ธิกรณสฺส กิํ ปุพฺพงฺคมํ, กติ ฐานานิ, กติ วตฺถูนิ, กติ ภูมิโย, กติ เหตู, กติ มูลานิ, กติหากาเรหิ อาปตฺติํ อาปชฺชติ, อาปตฺตา\-ธิกรณํ กติหิ สมเถหิ สมฺมติ.}

\prose{กิจฺจา\-ธิกรณสฺส กิํ ปุพฺพงฺคมํ, กติ ฐานานิ, กติ วตฺถูนิ, กติ ภูมิโย, กติ เหตู, กติ มูลานิ, ตติหากาเรหิ กิจฺจํ ชายติ, กิจฺจา\-ธิกรณํ กติหิ สมเถหิ สมฺมติ.}

\proseitem{293}{วิวาทา\-ธิกรณสฺส กิํ ปุพฺพงฺคมนฺติ โลโภ ปุพฺพงฺคโม, โทโส ปุพฺพงฺคโม, โมโห ปุพฺพงฺคโม, อโลโภ ปุพฺพงฺคโม, อโทโส ปุพฺพงฺคโม, อโมโห ปุพฺพงฺคโม. กติ ฐานานีติ อฏฺฐารส เภทกร\-วตฺถูนิ ฐานานิ. กติ วตฺถูนีติ อฏฺฐารส เภทกร\-วตฺถูนิ. กติ ภูมิโยติ อฏฺฐารส เภทกร\-วตฺถูนิ ภูมิโย. กติ เหตูติ นว เหตู, ตโย กุสลเหตู, ตโย อกุสลเหตู, ตโย อพฺยากตเหตู. กติ มูลานีติ ทฺวาทส มูลานิ. กติหากาเรหิ วิวทตีติ ทฺวีหากาเรหิ วิวทติ ธมฺมทิฏฺฐิ วา อธมฺมทิฏฺฐิ วา.\csromansymbolfootnote{*}{วิ 4. 220 ปิฏฺเฐปิ.}วิวาทา\-ธิกรณํ กติหิ สมเถหิ สมฺมตีติ วิวาทา\-ธิกรณํ ทฺวีหิ สมเถหิ สมฺมติ สมฺมุขา\-วินเยน จ เยภุยฺยสิกาย จ.}

\proseitem{294}{อนุวาทา\-ธิกรณสฺส กิํ ปุพฺพงฺคมนฺติ โลโภ ปุพฺพงฺคโม, โทโส ปุพฺพงฺคโม, โมโห ปุพฺพงฺคโม, อโลโภ ปุพฺพงฺคโม, อโทโส ปุพฺพงฺคโม, อโมโห ปุพฺพงฺคโม. กติ ฐานานีติ จตสฺโส \csromanfolio{184}วิปตฺติโย ฐานานิ. กติ วตฺตูนีติ จตสฺโส วิปตฺติโย วตฺถูนิ. กติ ภูมิโยติ จตสฺโส วิปตฺติโย ภูมิโย. กติเหตูติ นว เหตู, ตโย กุสลเหตู, ตโย อกุสลเหตู, ตโย อพฺยากตเหตู. กติ มูลานีติ จุทฺทส มูลานิ. กติหากาเรหิ อนุวทตีติ ทฺวีหากาเรหิ อนุวทติ วตฺถุโต วา อาปตฺติโต วา.\csromansymbolfootnote{*}{วิ 4. 229 ปิฏฺเฐปิ.}อนุวาทา\-ธิกรณํ กติหิ สมเถหิ สมฺมตีติ อนุวาทา\-ธิกรณํ จตูหิ สมเถหิ สมฺมติ สมฺมุขา\-วินเยน จ สติวินเยน จ อมูฬฺหวินเยน จ ตสฺสปาปิยสิกาย จ.}

\proseitem{295}{อาปตฺตา\-ธิกรณสฺส กิํ ปุพฺพงฺคมนฺติ โลโภ ปุพฺพงฺคโม, โทโส ปุพฺพงฺคโม, โมโห ปุพฺพงฺคโม, อโลโภ ปุพฺพงฺคโม, อโทโส ปุพฺพงฺคโม, อโมโห ปุพฺพงฺคโม. กติ ฐานานีติ สตฺต อาปตฺติกฺขนฺธา ฐานานิ. กติ วตฺถูนีติ สตฺต อาปตฺติกฺขนฺธา วตฺถูนิ. กติ ภูมิโยติ สตฺต อาปตฺติกฺขนฺธา ภูมิโย. กติ เหตูติ ฉ เหตู\csromansharedfootnote{3aba616e17e08d83}{กติ เหตูติ นว เหตู, ตโย กุสลเหตู (สี, สฺยา)}, ตโย อกุสลเหตู, ตโย อพฺยากตเหตู. กติ มูลานีติ ฉ อาปตฺติ\-สมุฏฺฐานานิ มูลานิ. กติหากาเรหิ อาปตฺติํ อาปชฺชตีติ ฉหากาเรหิ อาปตฺติํ อาปชฺชติ, อลชฺชิตา, อญฺญาณตา, กุกฺกุจฺจ\-ปกตตา, อกปฺปิเย กปฺปิยสญฺญิตา, กปฺปิเย อกปฺปิย\-สญฺญิตา, สติสมฺโมสา.\csromansymbolfootnote{+}{วิ 4. 233 ปิฏฺเฐปิ.}อาปตฺตา\-ธิกรณํ กติหิ สมเถหิ สมฺมตีติ อาปตฺตา\-ธิกรณํ ตีหิ สมเถหิ สมฺมติ สมฺมุขา\-วินเยน จ ปฏิญฺญาต กรเณน จ สมฺมุขา\-วินเยน จ ติณวตฺถารเกน จ.}

\proseitem{296}{กิจฺจา\-ธิกรณสฺส กิํ ปุพฺพงฺคมนฺติ โลโภ ปุพฺพงฺคโม, โทโส ปุพฺพงฺคโม, โมโห ปุพฺพงฺคโม, อโลโภ ปุพฺพงฺคโม, อโทโส ปุพฺพงฺคโม, อโมโห ปุพฺพงฺคโม. กติ ฐานานีติ จตฺตาริ กมฺมานิ ฐานานิ. กติ วตฺถูนีติ จตฺตาริ กมฺมานิ วตฺถูนิ. กติ ภูมิโยติ จตฺตาริ กมฺมานิ ภูมิโย. กติ เหตูติ นว เหตู, ตโย กุสลเหตู, ตโย อกุสลเหตู, ตโย อพฺยากตเหตู. กติ มูลานีติ เอกํ มูลํ สํโฆ. กติหากาเรหิ กิจฺจํ ชายตีติ ทฺวีหากาเรหิ กิจฺจํ ชายติ ญตฺติโต วา อปโลกนโต วา.}

\prose{\csromanfolio{185}\csromansymbolfootnote{*}{วิ 4. 238 ปิฏฺเฐปิ.}กิจฺจา\-ธิกรณํ กติหิ สมเถหิ สมฺมตีติ กิจฺจา\-ธิกรณํ เอเกน สมเถน สมฺมติ สมฺมุขา\-วินเยน.}

\prose{กติ สมถา. สตฺต สมถา สมฺมุขาวินโย สติวินโย อมูฬฺหวินโย ปฏิญฺญาต\-กรณํ เยภุยฺยสิกา ตสฺสปาปิยสิกา ติณวตฺถารโก, อิเม สตฺต สมถา.}

\prose{สิยา อิเม สตฺต สมถา ทส สมถา โหนฺติ, ทส สมถา สตฺต สมถา โหนฺติ วตฺถุวเสน ปริยาเยน? สิยาติ.}

\prosewithcloserruled{กถญฺจ สิยา? วิวาทา\-ธิกรณสฺส เทฺว สมถา, อนุวาทา\-ธิกรณสฺส จตฺตาโร สมถา, อาปตฺตา\-ธิกรณสฺส ตโย สมถา, กิจฺจา\-ธิกรณสฺส เอโก สมโถ. เอวํ อิเม สตฺต สมถา ทส สมถา โหนฺติ, ทส สมถา สตฺต สมถา โหนฺติ วตฺถุวเสน ปริยาเยน.}{ปริยายวาโร นิฏฺฐิโต ฉฏฺโฐ.}
\csromanmatikaanchor{mtk.147}
\csromantocmarkref{subsection}{7. สาธารณวาร}{mtk.147}
\bookmark[dest={mtk.147},level=2]{7. สาธารณวาร}
\csromanheader{2}{7. \textbf{สาธารณวาร}}
\proseitem{297}{กติ สมถา วิวาทา\-ธิกรณสฺส สาธารณา, กติ สมถา วิวาทา\-ธิกรณสฺส อสาธารณา, กติ สมถา อนุวาทา\-ธิกรณสฺส สาธารณา, กติ สมถา อนุวาทา\-ธิกรณสฺส อสาธารณา, กติ สมถา อาปตฺตา\-ธิกรณสฺส สาธารณา, กติ สมถา อาปตฺตา\-ธิกรณสฺส อสาธารณา, กติ สมถา กิจฺจา\-ธิกรณสฺส สาธารณา, กติ สมถา กิจฺจา\-ธิกรณสฺส อสาธารณา.}

\prose{เทฺว สมถา วิวาทา\-ธิกรณสฺส สาธารณา สมฺมุขาวินโย เยภุยฺยสิกา. ปญฺจ สมถา วิวาทา\-ธิกรณสฺส อสาธารณา สติวินโย อมูฬฺหวินโย ปฏิญฺญาต\-กรณํ ตสฺสปาปิยสิกา ติณวตฺถารโก.}

\prose{จตฺตาโร สมถา อนุวาทา\-ธิกรณสฺส สาธารณา สมฺมุขาวินโย สติวินโย อมูฬฺหวินโย ตสฺสปาปิยสิกา. ตโย สมถา อนุวาทา\-ธิกรณสฺส อสาธารณา เยภุยฺยสิกา ปฏิญฺญาต\-กรณํ ติณวตฺถารโก.}

\prose{\csromanfolio{186}ตโย สมถา อาปตฺตา\-ธิกรณสฺส สาธารณา สมฺมุขาวินโย ปฏิญฺญาต\-กรณํ ติณวตฺถารโก. จตฺตาโร สมถา อาปตฺตา\-ธิกรณสฺส อสาธารณา เยภุยฺยสิกา สติวินโย อมูฬฺหวินโย ตสฺสปาปิยสิกา.}

\prosewithcloserruled{เอโก สมโถ กิจฺจา\-ธิกรณสฺส สาธารโณ สมฺมุขาวินโย. ฉ สมถา กิจฺจา\-ธิกรณสฺส อสาธารณา เยภุยฺยสิกา สติวินโย อมูฬฺหวินโย ปฏิญฺญาต\-กรณํ ตสฺสปาปิยสิกา ติณวตฺถารโก.}{สาธารณวาโร นิฏฺฐิโต สตฺตโม.}
\csromanmatikaanchor{mtk.148}
\csromantocmarkref{subsection}{8. ตพฺภาคิยวาร}{mtk.148}
\bookmark[dest={mtk.148},level=2]{8. ตพฺภาคิยวาร}
\csromanheader{2}{8. \textbf{ตพฺภาคิยวาร}}
\proseitem{298}{กติ สมถา วิวาทา\-ธิกรณสฺส ตพฺภาคิยา, กติ สมถา วิวาทา\-ธิกรณสฺส อญฺญภาคิยา, กติ สมถา อนุวาทา\-ธิกรณสฺส ตพฺภาคิยา, กติ สมถา อนุวาทา\-ธิกรณสฺส อญฺญภาคิยา, กติ สมถา อาปตฺตา\-ธิกรณสฺส ตพฺภาคิยา, กติ สมถา อาปตฺตา\-ธิกรณสฺส อญฺญภาคิยา, กติ สมถา กิจฺจา\-ธิกรณสฺส ตพฺภาคิยา, กติ สมถา กิจฺจา\-ธิกรณสฺส อญฺญภาคิยา.}

\prose{เทฺว สมถา วิวาทา\-ธิกรณสฺส ตพฺภาคิยา สมฺมุขาวินโย เยภุยฺยสิกา. ปญฺจ สมถา วิวาทา\-ธิกรณสฺส อญฺญภาคิยา สติวินโย อมูฬฺหวินโย ปฏิญฺญาต\-กรณํ ตสฺสปาปิยสิกา ติณวตฺถารโก.}

\prose{จตฺตาโร สมถา อนุวาทา\-ธิกรณสฺส ตพฺภาคิยา สมฺมุขาวินโย สติวินโย อมูฬฺหวินโย ตสฺสปาปิยสิกา. ตโย สมถา อนุวาทา\-ธิกรณสฺส อญฺญภาคิยา เยภุยฺยสิกา ปฏิญฺญาต\-กรณํ ติณวตฺถารโก.}

\prose{ตโย สมถา อาปตฺตา\-ธิกรณสฺส ตพฺภาคิยา สมฺมุขาวินโย ปฏิญฺญาต\-กรณํ ติณวตฺถารโก. จตฺตาโร สมถา อาปตฺตา\-ธิกรณสฺส อญฺญภาคิยา เยภุยฺยสิกา สติวินโย อมูฬฺหวินโย ตสฺสปาปิยสิกา.}

\prosewithcloserruled{\csromanfolio{187}เอโก สมโถ กิจฺจา\-ธิกรณสฺส ตพฺภาคิโย สมฺมุขาวินโย. ฉ สมถา กิจฺจา\-ธิกรณสฺส อญฺญภาคิยา เยภุยฺยสิกา สติวินโย อมูฬฺหวินโย ปฏิญฺญาต\-กรณํ ตสฺสปาปิยสิกา ติณวตฺถารโก.}{ตพฺภาคิยวาโร นิฏฺฐิโต อฏฺฐโม.}
\csromanmatikaanchor{mtk.149}
\csromantocmarkref{subsection}{9. สมถา สมถสฺส สาธารณวาร}{mtk.149}
\bookmark[dest={mtk.149},level=2]{9. สมถา สมถสฺส สาธารณวาร}
\csromanheader{2}{9. \textbf{สมถา สมถสฺส สาธารณวาร}}
\proseitem{299}{สมถา สมถสฺส สาธารณา, สมถา สมถสฺส อสาธารณา. สิยา สมถา สมถสฺส สาธารณา, สิยา สมถา สมถสฺส อสาธารณา.}

\prose{กถํ สิยา สมถา สมถสฺส สาธารณา, กถํ สิยา สมถา สมถสฺส อสาธารณา, เยภุยฺยสิกา สมฺมุขา\-วินยสฺส สาธารณา, สติวินยสฺส อมูฬฺห\-วินยสฺส ปฏิญฺญาต\-กรณสฺส ตสฺสปาปิยสิกาย ติณวตฺถารกสฺส อสาธารณา.}

\prose{สติวินโย สมฺมุขา\-วินยสฺส สาธารโณ, อมูฬฺห\-วินยสฺส ปฏิญฺญาต\-กรณสฺส ตสฺสปาปิยสิกาย ติณวตฺถารกสฺส เยภุยฺย สิกาย อสาธารโณ.}

\prose{อมูฬฺหวินโย สมฺมุขา\-วินยสฺส สาธารโณ, ปฏิญฺญาต\-กรณสฺส ตสฺสปาปิยสิกาย ติณวตฺถารกสฺส เยภุยฺยสิกาย สติวินยสฺส อสาธารโณ.}

\prose{ปฏิญฺญาต\-กรณํ สมฺมุขา\-วินยสฺส สาธารณํ, ตสฺสปาปิยสิกาย ติณวตฺถารกสฺสเยภุยฺยสิกาย สติวินยสฺส อมูฬฺห\-วินยสฺส อสาธารณํ.}

\prose{ตสฺสปาปิยสิกา สมฺมุขา\-วินยสฺส สาธารณา, ติณวตฺถารกสฺส เยภุยฺยสิกาย สติวินยสฺส อมูฬฺห\-วินยสฺส ปฏิญฺญาต\-กรณสฺส อสาธารณา.}

\prosewithcloserruled{ติณวตฺถารโก สมฺมุขา\-วินยสฺส สาธารโณ, เยภุยฺยสิกาย สติวินยสฺส อมูฬฺห\-วินยสฺส ปฏิญฺญาต\-กรณสฺส ตสฺสปาปิยสิกาย \csromanfolio{188}อสาธารโณ. เอวํ สิยา สมถา สมถสฺส สาธารณา, เอวํ สิยา สมถา สมถสฺส อสาธารณา.}{สมถา สมถสฺส สาธารณวาโร นิฏฺฐิโต นวโม.}
\csromanmatikaanchor{mtk.150}
\csromantocmarkref{subsection}{10. สมถา สมถสฺส ตพฺภาคิยวาร}{mtk.150}
\bookmark[dest={mtk.150},level=2]{10. สมถา สมถสฺส ตพฺภาคิยวาร}
\csromanheader{2}{10. \textbf{สมถา สมถสฺส ตพฺภาคิยวาร}}
\proseitem{300}{\textbf{สมถา สมถสฺส ตพฺพฺ}หาคิยา, สมถา สมถสฺส อญฺญภาคิยา. สิยา สมถา สมถสฺส ตพฺภาคิยา, สิยา สมถา สมถสฺส อญฺญภาคิยา.}

\prose{กถํ สิยา สมถา สมถสฺส ตพฺภาคิยา, กถํ สิยา สมถา สมถสฺส อญฺญภาคิยา. เยภุยฺยสิกา สมฺมุขา\-วินยสฺส ตพฺภาคิยา, สติวินยสฺส อมูฬฺห\-วินยสฺส ปฏิญฺญาต\-กรณสฺส ตสฺสปาปิยสิกาย ติณวตฺถารกสฺส อญฺญภาคิยา.}

\prose{สติวินโย สมฺมุขา\-วินยสฺส ตพฺภาคิโย, อมูฬฺห\-วินยสฺส ปฏิญฺญาต\-กรณสฺส ตสฺสปาปิยสิกาย ติณวตฺถารกสฺส เยภุยฺยสิกาย อญฺญภาคิโย.}

\prose{อมูฬฺหวินโย สมฺมุขา\-วินยสฺส ตพฺภาคิโย, ปฏิญฺญาต\-กรณสฺส ตสฺสปาปิยสิกาย ติณวตฺถารกสฺส เยภุยฺยสิกาย สติวินยสฺส อญฺญภาคิโย.}

\prose{ปฏิญฺญาต\-กรณํ สมฺมุขา\-วินยสฺส ตพฺภาคิยํ, ตสฺสปาปิยสิกาย ติณวตฺถารกสฺส เยภุยฺยสิกาย สติวินยสฺส อมูฬฺห\-วินยสฺส อญฺญภาคิยํ.}

\prose{ตสฺสปาปิยสิกา สมฺมุขา\-วินยสฺส ตพฺภาคิยา, ติณวตฺถารกสฺส เยภุยฺยสิกาย สติวินยสฺส อมูฬฺห\-วินยสฺส ปฏิญฺญาต\-กรณสฺส อญฺญภาคิยา.}

\prose{ติณวตฺถารโก สมฺมุขา\-วินยสฺส ตพฺภาคิโย, เยภุยฺยสิกาย สติวินยสฺส อมูฬฺห\-วินยสฺส ปฏิญฺญาต\-กรณสฺส ตสฺสปาปิยสิกาย อญฺญภาคิโย. เอวํ สิยา สมถา สมถสฺส ตพฺภาคิยา, เอวํ สิยา สมถา สมถสฺส อญฺญภาคิยา.}

\vspace{\tipitakaparskip}
\csromancenter{\textbf{สมถา สมถสฺส ตพฺพฺ}หาคิยวาโร นิฏฺฐิโต ทสโม.}
\csromanmatikaanchor{mtk.151}
\csromantocmarkref{subsection}{11. สมถสมฺมุขาวินยวาร}{mtk.151}
\bookmark[dest={mtk.151},level=2]{11. สมถสมฺมุขาวินยวาร}
\csromanheader{2}{11. สมถ\-สมฺมุขาวินย\-วาร}
\proseitem{301}{\csromanfolio{189}สมโถ สมฺมุขาวินโย, สมฺมุขาวินโย สมโถ. สมโถ เยภุยฺยสิกา, เยภุยฺยสิกา สมโถ. สมโถ สติวินโย, สติวินโย สมโถ. สมโถ อมูฬฺหวินโย, อมูฬฺหวินโย สมโถ. สมโถ ปฏิญฺญาต\-กรณํ, ปฏิญฺญาต\-กรณํ สมโถ. สมโถ ตสฺสปาปิยสิกา, ตสฺสปาปิยสิกา สมโถ. สมโถ ติณวตฺถารโก, ติณวตฺถารโก สมโถ.}

\prose{เยภุยฺยสิกา สติวินโย อมูฬฺหวินโย ปฏิญฺญาต\-กรณํ ตสฺสปาปิยสิกา ติณวตฺถารโก, อิเม สมถา สมถา, โน สมฺมุขาวินโย, สมฺมุขาวินโย สมโถ เจว สมฺมุขาวินโย จ.}

\prose{สติวินโย อมูฬฺหวินโย ปฏิญฺญาต\-กรณํ ตสฺสปาปิยสิกา ติณวตฺถารโก สมฺมุขาวินโย, อิเม สมถา สมถา, โน เยภุยฺยสิกา, เยภุยฺยสิกา สมโถ เจว เยภุยฺยสิกา จ.}

\prose{อมูฬฺหวินโย ปฏิญฺญาต\-กรณํ ตสฺสปาปิยสิกา ติณวตฺถารโก สมฺมุขาวินโย เยภุยฺยสิกา, อิเม สมถา สมถา, โน สติวินโย, สติวินโย สมโถ เจว สติวินโย จ.}

\prose{ปฏิญฺญาต\-กรณํ ตสฺสปาปิยสิกา ติณวตฺถารโก สมฺมุขาวินโย เยภุยฺยสิกา สติวินโย, อิเม สมถา สมถา, โน อมูฬฺหวินโย, อมูฬฺหวินโย สมโถ เจว อมูฬฺหวินโย จ.}

\prose{ตสฺสปาปิยสิกา ติณวตฺถารโก สมฺมุขาวินโย เยภุยฺยสิกา สติวินโย อมูฬฺหวินโย, อิเม สมถา สมถา, โน ปฏิญฺญาต\-กรณํ, ปฏิญฺญาต\-กรณํ สมโถ เจว ปฏิญฺญาต\-กรณญฺจ.}

\prose{ติณวตฺถารโก สมฺมุขาวินโย เยภุยฺยสิกา สติวินโย อมูฬฺหวินโย ปฏิญฺญาต\-กรณํ, อิเม สมถา สมถา, โน ตสฺสปาปิยสิกา, ตสฺสปาปิยสิกา สมโถ เจว ตสฺสปาปิยสิกา จ.}

\prose{สมฺมุขาวินโย เยภุยฺยสิกา สติวินโย อมูฬฺหวินโย ปฏิญฺญาต\-กรณํ ตสฺสปาปิยสิกา, อิเม สมถา สมถา, โน ติณวตฺถารโก, ติณวตฺถารโก สมโถ เจว ติณวตฺถารโก จ.}

\vspace{\tipitakaparskip}
\csromancenter{สมถ\-สมฺมุขาวินย\-วาโร นิฏฺฐิโต เอกาทสโม.}
\csromanmatikaanchor{mtk.152}
\csromantocmarkref{subsection}{12. วินยวาร}{mtk.152}
\bookmark[dest={mtk.152},level=2]{12. วินยวาร}
\csromanheader{2}{12. \textbf{วินยวาร}}
\proseitem{302}{\csromanfolio{190}วินโย สมฺมุขาวินโย, สมฺมุขาวินโย วินโย. วินโย เยภุยฺยสิกา, เยภุยฺยาสิกา วินโย. วินโย สติวินโย, สติวินโย วินโย. วินโย อมูฬฺหวินโย, อมูฬฺหวินโย วินโย. วินโย ปฏิญฺญาต\-กรณํ, ปฏิญฺญาต\-กรณํ วินโย. วินโย ตสฺสปาปิยสิกา, ตสฺสปาปิยสิกา วินโย. วินโย ติณวตฺถารโก, ติณวตฺถารโก วินโย.}

\prose{วินโย สิยา สมฺมุขาวินโย, สิยา น สมฺมุขาวินโย, สมฺมุขาวินโย วินโย เจว สมฺมุขาวินโย จ.}

\prose{วินโย สิยา เยภุยฺยสิกา, สิยา น เยภุยฺยสิกา, เยภุยฺยสิกา วินโย เจว เยภุยฺยสิกา จ.}

\prose{วินโย สิยา สติวินโย, สิยา น สติวินโย, สติวินโย วินโย เจว สติวินโย จ.}

\prose{วินโย สิยา อมูฬฺหวินโย, สิยา น อมูฬฺหวินโย, อมูฬฺหวินโย วินโย เจว อมูฬฺหวินโย จ.}

\prose{วินโย สิยา ปฏิญฺญาต\-กรณํ, สิยา น ปฏิญฺญาต\-กรณํ, ปฏิญฺญาต\-กรณํ วินโย เจว ปฏิญฺญาต\-กรณญฺจ.}

\prose{วินโย สิยา ตสฺสปาปิยสิกา, สิยา น ตสฺสปาปิยสิกา, ตสฺสปาปิยสิกา วินโย เจว ตสฺสปาปิยสิกา จ.}

\prosewithcloserruled{วินโย สิยา ติณวตฺถารโก, สิยา น ติณวตฺถารโก, ติณวตฺถารโก วินโย เจว ติณวตฺถารโก จ.}{วินยวาโร นิฏฺฐิโต ทฺวาทสโม.}
\csromanmatikaanchor{mtk.153}
\csromantocmarkref{subsection}{13. กุสลวาร}{mtk.153}
\bookmark[dest={mtk.153},level=2]{13. กุสลวาร}
\csromanheader{2}{13. \textbf{กุสลวาร}}
\proseitem{303}{สมฺมุขาวินโย กุสโล อกุสโล อพฺยากโต, เยภุยฺยสิกา กุสลา อกุสลา อพฺยากตา, สติวินโย กุสโล อกุสโล อพฺยากโต, อมูฬฺหวินโย กุสโล \csromanfolio{191}อกุสโล อพฺยากโต, ปฏิญฺญาต\-กรณํ กุสลํ อกุสลํ อพฺยากตํ, ตสฺสปาปิยสิกา กุสลา อกุสลา อพฺยากตา, ติณวตฺถารโก กุสโล อกุสโล อพฺยากโต.}

\prose{สมฺมุขาวินโย สิยา กุสโล, สิยา อพฺยากโต, นตฺถิ สมฺมุขาวินโย อกุสโล.}

\prose{เยภุยฺยสิกา สิยา กุสลา, สิยา อกุสลา, สิยา อพฺยากตา.}

\prose{สติวินโย สิยา กุสโล, สิยา อกุสโล, สิยา อพฺยากโต.}

\prose{อมูฬฺหวินโย สิยา กุสโล, สิยา อกุสโล, สิยา อพฺยากโต.}

\prose{ปฏิญฺญาต\-กรณํ สิยา กุสลํ, สิยา อกุสลํ, สิยา อพฺยากตํ.}

\prose{ตสฺสปาปิยสิกา สิยา กุสลา, สิยา อกุสลา, สิยา อพฺยากตา.}

\prose{ติณวตฺถารโก สิยา กุสโล, สิยา อกุสโล, สิยา อพฺยากโต.}

\prose{วิวาทา\-ธิกรณํ กุสลํ อกุสลํ อพฺยากตํ, อนุวาทา\-ธิกรณํ กุสลํ อกุสลํ อพฺยากตํ, อาปตฺตา\-ธิกรณํ กุสลํ อกุสลํ อพฺยากตํ, กิจฺจา\-ธิกรณํ กุสลํ อกุสลํ อพฺยากตํ.}

\prose{วิวาทา\-ธิกรณํ สิยา กุสลํ, สิยา อกุสลํ, สิยา อพฺยากตํ.}

\prose{อนุวาทา\-ธิกรณํ สิยา กุสลํ, สิยา อกุสลํ, สิยา อพฺยากตํ.}

\prose{อาปตฺตา\-ธิกรณํ สิยา อกุสลํ, สิยา อพฺยากตํ, นตฺถิ อาปตฺตา\-ธิกรณํ กุสลํ.}

\prose{กิจฺจา\-ธิกรณํ สิยา กุสลํ, สิยา อกุสลํ, สิยา อพฺยากตํ.}

\vspace{\tipitakaparskip}
\csromancenter{กุสลวาโร นิฏฺฐิโต เตรสโม.}
\csromancenter{\textbf{ยตฺถวาร, ปุจฺฉาวาร}}
\csromanmatikaanchor{mtk.154}
\csromantocmarkref{subsection}{14. ยตฺถวาร, ปุจฺฉาวาร}{mtk.154}
\bookmark[dest={mtk.154},level=2]{14. ยตฺถวาร, ปุจฺฉาวาร}
\proseitem{304}{\csromanfolio{192}ยตฺถ เยภุยฺยสิกา ลพฺภติ, ตตฺถ สมฺมุขาวินโย ลพฺภติ, ยตฺถ สมฺมุขาวินโย ลพฺภติ, ตตฺถ เยภุยฺยสิกา ลพฺภติ, น ตตฺถ สติวินโย ลพฺภติ, น ตตฺถ อมูฬฺหวินโย ลพฺภติ, น ตตฺถ ปฏิญฺญาต\-กรณํ ลพฺภติ, น ตตฺถ ตสฺสปาปิยสิกา ลพฺภติ, น ตตฺถ ติณวตฺถารโก ลพฺภติ.}

\prose{ยตฺถ สติวินโย ลพฺภติ, ตตฺถ สมฺมุขาวินโย ลพฺภติ, ยตฺถ สมฺมุขาวินโย ลพฺภติ, ตตฺถ สติวินโย ลพฺภติ, น ตตฺถ อมูฬฺห วินโย ลพฺภติ, น ตตฺถ ปฏิญฺญาต\-กรณํ ลพฺภติ, น ตตฺถ ตสฺสปาปิยสิกา ลพฺภติ, น ตตฺถ ติณวตฺถารโก ลพฺภติ, น ตตฺถ เยภุยฺยสิกา ลพฺภติ.}

\prose{ยตฺถ อมูฬฺหวินโย ลพฺภติ, ตตฺถ สมฺมุขาวินโย ลพฺภติ, ยตฺถ สมฺมุขาวินโย ลพฺภติ, ตตฺถ อมูฬฺหวินโย ลพฺภติ, น ตตฺถ ปฏิญฺญาต\-กรณํ ลพฺภติ, น ตตฺถ ตสฺสปาปิยสิกา ลพฺภติ, น ตตฺถ ติณวตฺถารโก ลพฺภติ, น ตตฺถ เยภุยฺยสิกา ลพฺภติ, น ตตฺถ สติวินโย ลพฺภติ.}

\prose{ยตฺถ ปฏิญฺญาต\-กรณํ ลพฺภติ, ตตฺถ สมฺมุขาวินโย ลพฺภติ, ยตฺถ สมฺมุขาวินโย ลพฺภติ, ตตฺถ ปฏิญฺญาต\-กรณํ ลพฺภติ, น ตตฺถ ตสฺสปาปิยสิกา ลพฺภติ, น ตตฺถ ติณวตฺถารโก ลพฺภติ, น ตตฺถ เยภุยฺยสิกา ลพฺภติ, น ตตฺถ สติวินโย ลพฺภติ, น ตตฺถ อมูฬฺหวินโย ลพฺภติ.}

\prose{ยตฺถ ตสฺสปาปิยสิกา ลพฺภติ, ตตฺถ สมฺมุขาวินโย ลพฺภติ, ยตฺถ สมฺมุขาวินโย ลพฺภติ, ตตฺถ ตสฺสปาปิยสิกา ลพฺภติ, น ตตฺถ ติณวตฺถารโก ลพฺภติ, น ตตฺถ เยภุยฺยสิกา ลพฺภติ, น ตตฺถ สติวินโย ลพฺภติ, น ตตฺถ อมูฬฺหวินโย ลพฺภติ, น ตตฺถ ปฏิญฺญาต\-กรณํ ลพฺภติ.}

\csromanmatikaanchor{mtk.155}
\csromantocmarkref{subsection}{15. สมถวาร, วิสชฺชนาวาร}{mtk.155}
\bookmark[dest={mtk.155},level=2]{15. สมถวาร, วิสชฺชนาวาร}
\prose{ยตฺถ ติณวตฺถารโก ลพฺภติ, ตตฺถ สมฺมุขาวินโย ลพฺภติ, ยตฺถ สมฺมุขาวินโย ลพฺภติ, ตตฺถ ติณวตฺถารโก ลพฺภติ, น ตตฺถ เยภุยฺยสิกา ลพฺภติ, น ตตฺถ สติวินโย ลพฺภติ, น ตตฺถ \csromanfolio{193}อมูฬฺหวินโย ลพฺภติ, น ตตฺถ ปฏิญฺญาต\-กรณํ ลพฺภติ, น ตตฺถ ตสฺสปาปิยสิกา ลพฺภติ.}

\prose{ยตฺถ เยภุยฺยสิกา, ตตฺถ สมฺมุขาวินโย, ยตฺถ สมฺมุขาวินโย, ตตฺถ เยภุยฺยสิกา, น ตตฺถ สติวินโย, น ตตฺถ อมูฬฺหวินโย, น ตตฺถ ปฏิญฺญาต\-กรณํ, น ตตฺถ ตสฺสปาปิยสิกา, น ตตฺถ ติณวตฺถารโก.}

\prose{ยตฺถ สติวินโย, ตตฺถ สมฺมุขาวินโย, ยตฺถ สมฺมุขาวินโย, ตตฺถ สติวินโย, น ตตฺถ อมูฬฺหวินโย, น ตตฺถ ปฏิญฺญาต\-กรณํ, น ตตฺถ ตสฺสปาปิยสิกา, น ตตฺถ ติณวตฺถารโก, น ตตฺถ เยภุยฺยสิกา. (สมฺมุขา\-วินยํ กาตุน มูลํ ฯเปฯ.)}

\prose{ยตฺถ ติณวตฺถารโก, ตตฺถ สมฺมุขาวินโย, ยตฺถ สมฺมุขาวินโย, ตตฺถ ติณวตฺถารโก, น ตตฺถ เยภุยฺยสิกา, น ตตฺถ สติวินโย, น ตตฺถ อมูฬฺหวินโย, น ตตฺถ ปฏิญฺญาต\-กรณํ, น ตตฺถ ตสฺสปาปิยสิกา.}

\vspace{\tipitakaparskip}
\csromancenter{จกฺกเปยฺยาลํ.}
\nitthitamruled{ยตฺถวาโร นิฏฺฐิโต จุทฺทสโม.}
\csromancenter{\textbf{สมถวาร, วิสฺสชฺชนาวาร}}
\proseitem{305}{ยสฺมิํ สมเย สมฺมุขา\-วินเยน จ เยภุยฺยสิกาย จ อธิกรณํ วูปสมฺมติ, ยตฺถ เยภุยฺยสิกา ลพฺภติ, ตตฺถ สมฺมุขาวินโย ลพฺภติ, ยตฺถ สมฺมุขาวินโย ลพฺภติ, ตตฺถ เยภุยฺยสิกา ลพฺภติ, น ตตฺถ สติวินโย ลพฺภติ, น ตตฺถ อมูฬฺหวินโย ลพฺภติ, น ตตฺถ ปฏิญฺญาต\-กรณํ ลพฺภติ, น ตตฺถ ตสฺสปาปิยสิกา ลพฺภติ, น ตตฺถ ติณวตฺถารโก ลพฺภติ.}

\prose{ยสฺมิํ สมเย สมฺมุขา\-วินเยน จ สติวินเยน จ อธิกรณํ วูปสมฺมติ, ยตฺถ สติวินโย ลพฺภติ, ตตฺถ สมฺมุขาวินโย ลพฺภติ, ยตฺถ สมฺมุขาวินโย ลพฺภติ, ตตฺถ สติวินโย ลพฺภติ, น ตตฺถ อมูฬฺหวินโย ลพฺภติ, น ตตฺถ ปฏิญฺญาต\-กรณํ ลพฺภติ, น ตตฺถ \csromanfolio{194}ตสฺสปาปิยสิกา ลพฺภติ, น ตตฺถ ติณวตฺถารโก ลพฺภติ, น ตตฺถ เยภุยฺยสิกา ลพฺภติ.}

\prose{ยสฺมิํ สมเย สมฺมุขา\-วินเยน จ อมูฬฺหวินเยน จ อธิกรณํ วูปสมฺมติ, ยตฺถ อมูฬฺหวินโย ลพฺภติ, ตตฺถ สมฺมุขาวินโย ลพฺภติ, ยตฺถ สมฺมุขาวินโย ลพฺภติ, ตตฺถ อมูฬฺหวินโย ลพฺภติ, น ตตฺถ ปฏิญฺญาต\-กรณํ ลพฺภติ, น ตตฺถ ตสฺสปาปิยสิกา ลพฺภติ, น ตตฺถ ติณวตฺถารโก ลพฺภติ, น ตตฺถ เยภุยฺยสิกา ลพฺภติ, น ตตฺถ สติวินโย ลพฺภติ.}

\prose{ยสฺมิํ สมเย สมฺมุขา\-วินเยน จ ปฏิญฺญาต\-กรเณน จ อธิกรณํ วูปสมฺมติ, ยตฺถ ปฏิญฺญาต\-กรณํ ลพฺภติ, ตตฺถ สมฺมุขาวินโย ลพฺภติ, ยตฺถ สมฺมุขาวินโย ลพฺภติ, ตตฺถ ปฏิญฺญาต\-กรณํ ลพฺภติ, น ตตฺถ ตสฺสปาปิยสิกา ลพฺภติ, น ตตฺถ ติณวตฺถารโก ลพฺภติ, น ตตฺถ เยภุยฺยสิกา ลพฺภติ, น ตตฺถ สติวินโย ลพฺภติ, น ตตฺถ อมูฬฺหวินโย ลพฺภติ.}

\prose{ยสฺมิํ สมเย สมฺมุขา\-วินเยน จ ตสฺสปาปิยสิกาย จ อธิกรณํ วูปสมฺมติ, ยตฺถ ตสฺสปาปิยสิกา ลพฺภติ, ตตฺถ สมฺมุขาวินโย ลพฺภติ, ยตฺถ สมฺมุขาวินโย ลพฺภติ, ตตฺถ ตสฺสปาปิยสิกา ลพฺภติ, น ตตฺถ ติณวตฺถารโก ลพฺภติ, น ตตฺถ เยภุยฺยสิกา ลพฺภติ, น ตตฺถ สติวินโย ลพฺภติ, น ตตฺถ อมูฬฺหวินโย ลพฺภติ, น ตตฺถ ปฏิญฺญาต\-กรณํ ลพฺภติ.}

\prose{ยสฺมิํ สมเย สมฺมุขา\-วินเยน จ ติณวตฺถารเกน จ อธิกรณํ วูปสมฺมติ, ยตฺถ ติณวตฺถารโก ลพฺภติ, ตตฺถ สมฺมุขาวินโย ลพฺภติ, ยตฺถ สมฺมุขาวินโย ลพฺภติ, ตตฺถ ติณวตฺถารโก ลพฺภติ, น ตตฺถ เยภุยฺยสิกา ลพฺภติ, น ตตฺถ สติวินโย ลพฺภติ, น ตตฺถ อมูฬฺหวินโย ลพฺภติ, น ตฺถ ปฏิญฺญาต\-กรณํ ลพฺภติ, น ตตฺถ ตสฺสปาปิยสิกา ลพฺภติ.}

\vspace{\tipitakaparskip}
\csromancenter{สมถวาโร นิฏฺฐิโต ปนฺนรสโม.}
\csromanmatikaanchor{mtk.156}
\csromantocmarkref{subsection}{16. สํสฏฺฐวาร}{mtk.156}
\bookmark[dest={mtk.156},level=2]{16. สํสฏฺฐวาร}
\csromanheader{2}{16. \textbf{สํสฏฺฐวาร}}
\proseitem{306}{\csromanfolio{195}อธิกรณนฺติ วา สมถาติ วา อิเม ธมฺมา สํสฏฺฐา, อุทาหุ วิสํสฏฺฐา, ลพฺภา จ ปนิเมสํ ธมฺมานํ วินิพฺภุชิตฺวา วินิพฺภุชิตฺวา\csromansharedfootnote{10cfad6389cff880}{วินิพฺภุชฺชิตฺวา วินิพฺภุชฺชิตฺวา (ก), ฏีกายํ เอกปทเมว ทิสฺสติ.} นานากรณํ ปญฺญาเปตุนฺติ.}

\prosewithcloserruled{อธิกรณนฺติ วา สมถาติ วา อิเม ธมฺมา วิสํสฏฺฐา, โน สํสฏฺฐา, ลพฺภา จ ปนิเมสํ ธมฺมานํ วินิพฺภุชิตฺวา วินิพฺภุชิตฺวา นานากรณํ ปญฺญาเปตุนฺติ, โส “มา เหวนฺ”ติสฺส วจนีโย, อธิกรณนฺติ วา สมถาติ วา อิเม ธมฺมา สํสฏฺฐา, โน วิสํสฏฺฐา, โน จ ลพฺภา\csromansharedfootnote{7229db784204bce1}{น จ ลพฺภา(ก)} อิเมสํ ธมฺมานํ วินิพฺภุชิตฺวา วินิพฺภุชิตฺวา นานากรณํ ปญฺญาเปตุํ. ตํ กิสฺส เหตุ? นนุ วุตฺตํ ภควตา จตฺตาริมานิ ภิกฺขเว อธิกรณานิ สตฺต สมถา, อธิกรณา สมเถหิ สมฺมนฺติ, สมถา อธิกรเณหิ สมฺมนฺติ. เอวํ อิเม ธมฺมา สํสฏฺฐา, โน วิสํสฏฺฐา, โน จ ลพฺภา อิเมสํ ธมฺมานํ วินิพฺภุชิตฺวา วินิพฺภุชิตฺวา นานากรณํ ปญฺญาเปตุนฺติ.}{สํสฏฺฐวาโร นิฏฺฐิโต โสฬสโม.}
\csromanmatikaanchor{mtk.157}
\csromantocmarkref{subsection}{17. สมฺมติวาร}{mtk.157}
\bookmark[dest={mtk.157},level=2]{17. สมฺมติวาร}
\csromanheader{2}{17. \textbf{สมฺมติวาร}}
\proseitem{307}{วิวาทา\-ธิกรณํ กติหิ สมเถหิ สมฺมติ, อนุวาทา\-ธิกรณํ กติหิ สมเถหิ สมฺมติ, อาปตฺตา\-ธิกรณํ กติหิ สมเถหิ สมฺมติ, กิจฺจา\-ธิกรณํ กติหิ สมเถหิ สมฺมติ.}

\prose{\csromansymbolfootnote{*}{วิ 4. 220, 229 ปิฏฺเฐสุปิ.}วิวาทา\-ธิกรณํ ทฺวีหิ สมเถหิ สมฺมติ, สมฺมุขา\-วินเยน จ เยภุยฺยสิกาย จ.}

\prose{\csromansymbolmark{*}อนุวาทา\-ธิกรณํ จตูหิ สมเถหิ สมฺมติ, สมฺมุขา\-วินเยน จ สติวินเยน จ อมูฬฺหวินเยน จ ตสฺสปาปิยสิกาย จ.}

\prose{\csromansymbolfootnote{+}{วิ 4. 233-238 ปิฏฺเฐสุปิ.}อาปตฺตา\-ธิกรณํ ตีหิ สมเถหิ สมฺมติ, สมฺมุขา\-วินเยน จ ปฏิญฺญาต\-กรเณน จ ติณวตฺถารเกน จ.}

\prose{\csromansymbolmark{+}กิจฺจา\-ธิกรณํ เอเกน สมเถน สมฺมติ, สมฺมุขา\-วินเยน.}

\prose{\csromanfolio{196}วิวาทาธิกรณญฺจ อนุวาทาธิกรณญฺจ กติหิ สมเถหิ สมฺมนฺติ. วิวาทาธิกรณญฺจ อนุวาทาธิกรณญฺจ ปญฺจหิ สมเถหิ สมฺมนฺติ สมฺมุขา\-วินเยน จ เยภุยฺยสิกาย จ สติวินเยน จ อมูฬฺหวินเยน จ ตสฺสปาปิยสิกาย จ.}

\prose{วิวาทาธิกรณญฺจ อาปตฺตาธิกรณญฺจ กติหิ สมเถหิ สมฺมนฺติ. วิวาทาธิกรณญฺจ อาปตฺตาธิกรณญฺจ จตูหิ สมเถหิ สมฺมนฺติ สมฺมุขา\-วินเยน จ เยภุยฺยสิกาย จ ปฏิญฺญาต\-กรเณน จ ติณวตฺถารเกน จ.}

\prose{วิวาทาธิกรณญฺจ กิจฺจาธิกรณญฺจ กติหิ สมเถหิ สมฺมนฺติ. วิวาทาธิกรณญฺจ กิจฺจาธิกรณญฺจ ทฺวีหิ สมเถหิ สมฺมนฺติ สมฺมุขา\-วินเยน จ เยภุยฺยสิกาย จ.}

\prose{อนุวาทาธิกรณญฺจ อาปตฺตาธิกรณญฺจ กติหิ สมเถหิ สมฺมนฺติ. อนุวาทาธิกรณญฺจ อาปตฺตาธิกรณญฺจ ฉหิ สมเถหิ สมฺมนฺติ สมฺมุขา\-วินเยน จ สติวินเยน จ อมูฬฺหวินเยน จ ปฏิญฺญาต\-กรเณน จ ตสฺสปาปิยสิกาย จ ติณวตฺถารเกน จ.}

\prose{อนุวาทาธิกรณญฺจ กิจฺจาธิกรณญฺจ กติหิ สมเถหิ สมฺมนฺติ. อนุวาทาธิกรณญฺจ กิจฺจาธิกรณญฺจ จตูหิ สมเถหิ สมฺมนฺติ สมฺมุขา\-วินเยน จ สติวินเยน จ อมูฬฺหวินเยน จ ตสฺสปาปิยสิกาย จ.}

\prose{อาปตฺตาธิกรณญฺจ กิจฺจาธิกรณญฺจ กติหิ สมเถหิ สมฺมนฺติ. อาปตฺตาธิกรณญฺจ กิจฺจาธิกรณญฺจ ตีหิ สมเถหิ สมฺมนฺติ สมฺมุขา\-วินเยน จ ปฏิญฺญาต\-กรเณน จ ติณวตฺถารเกน จ.}

\prose{วิวาทาธิกรณญฺจ อนุวาทาธิกรณญฺจ อาปตฺตาธิกรณญฺจ กติหิ สมเถหิ สมฺมนฺติ. วิวาทาธิกรณญฺจ อนุวาทาธิกรณญฺจ อาปตฺตาธิกรณญฺจ สตฺตหิ สมเถหิ สมฺมนฺติ สมฺมุขา\-วินเยน จ เยภุยฺยสิกาย จ สติวินเยน จ อมูฬฺหวินเยน จ ปฏิญฺญาต\-กรเณน จ ตสฺสปาปิยสิกาย จ ติณวตฺถารเกน จ.}

\prose{วิวาทาธิกรณญฺจ อนุวาทาธิกรณญฺจ กิจฺจาธิกรณญฺจ กติหิ สมเถหิ สมฺมนฺติ. วิวาทาธิกรณญฺจ อนุวาทาธิกรณญฺจ กิจฺจาธิกรณญฺจ ปญฺจหิ \csromanfolio{197}สมเถหิ สมฺมนฺติ สมฺมุขา\-วินเยน จ เยภุยฺยสิกาย จ สติวินเยน จ อมูฬฺหวินเยน จ ตสฺสปาปิยสิกาย จ.}

\prose{อนุวาทาธิกรณญฺจ อาปตฺตาธิกรณญฺจ กิจฺจาธิกรณญฺจ กติหิ สมเถหิ สมฺมนฺติ. อนุวาทาธิกรณญฺจ อาปตฺตาธิกรณญฺจ กิจฺจาธิกรณญฺจ ฉหิ สมเถหิ สมฺมนฺติ สมฺมุขา\-วินเยน จ สติวินเยน จ อมูฬฺหวินเยน จ ปฏิญฺญาต\-กรเณน จ ตสฺสปาปิยสิกาย จ ติณวตฺถารเกน จ.}

\prosewithcloserruled{วิวาทาธิกรณญฺจ อนุวาทาธิกรณญฺจ อาปตฺตาธิกรณญฺจ กิจฺจาธิกรณญฺจ กติหิ สมเถหิ สมฺมนฺติ. วิวาทาธิกรณญฺจ อนุวาทาธิกรณญฺจ อาปตฺตาธิกรณญฺจ กิจฺจาธิกรณญฺจ สตฺตหิ สมเถหิ สมฺมนฺติ สมฺมุขา\-วินเยน จ เยภุยฺยสิกาย จ สติวินเยน จ อมูฬฺหวินเยน จ ปฏิญฺญาต\-กรเณน จ ตสฺสปาปิยสิกาย จ ติณวตฺถารเกน จ.}{สมฺมติวาโร นิฏฺฐิโต สตฺตรสโม.}
\csromanmatikaanchor{mtk.158}
\csromantocmarkref{subsection}{18. สมฺมนฺติ น สมฺมนฺติวาร}{mtk.158}
\bookmark[dest={mtk.158},level=2]{18. สมฺมนฺติ น สมฺมนฺติวาร}
\csromanheader{2}{18. \textbf{สมฺมนฺติ น สมฺมนฺติวาร}}
\proseitem{308}{วิวาทา\-ธิกรณํ กติหิ สมเถหิ สมฺมติ, กติหิ สมเถหิ น สมฺมติ. อนุวาทา\-ธิกรณํ กติหิ สมเถหิ สมฺมติ, กติหิ สมเถหิ น สมฺมติ. อาปตฺตา\-ธิกรณํ กติหิ สมเถหิ สมฺมติ, กติหิ สมเถหิ น สมฺมติ. กิจฺจา\-ธิกรณํ กติหิ สมเถหิ สมฺมติ, กติหิ สมเถหิ น สมฺมติ.}

\prose{วิวาทา\-ธิกรณํ ทฺวีหิ สมเถหิ สมฺมติ, สมฺมุขา\-วินเยน จ เยภุยฺยสิกาย จ. ปญฺจหิ สมเถหิ น สมฺมติ, สติวินเยน จ อมูฬฺหวินเยน จ ปฏิญฺญาต\-กรเณน จ ตสฺสปาปิยสิกาย จ ติณวตฺถารเกน จ.}

\prose{อนุวาทา\-ธิกรณํ จตูหิ สมเถหิ สมฺมติ, สมฺมุขา\-วินเยน จ สติวินเยน จ อมูฬฺหวินเยน จ ตสฺสปาปิยสิกาย จ. ตีหิ สมเถหิ น สมฺมติ, เยภุยฺยสิกาย จ ปฏิญฺญาต\-กรเณน จ ติณวตฺถารเกน จ.}

\prose{\csromanfolio{198}\csromansymbolfootnote{*}{วิ 4. 233 ปิฏฺเฐปิ.}อาปตฺตา\-ธิกรณํ ตีหิ สมเถหิ สมฺมติ สมฺมุขา\-วินเยน จ ปฏิญฺญาต\-กรเณน จ ติณวตฺถารเกน จ. จตูหิ สมเถหิ น สมฺมติ เยภุยฺยสิกาย จ สติวินเยน จ อมูฬฺหวินเยน จ ตสฺสปาปิยสิกาย จ.}

\prose{กิจฺจา\-ธิกรณํ เอเกน สมเถน สมฺมติ สมฺมุขา\-วินเยน. ฉหิ สมเถหิ น สมฺมติ เยภุยฺยสิกาย จ สติวินเยน จ อมูฬฺหวินเยน จ ปฏิญฺญาต\-กรเณน จ ตสฺสปาปิยสิกาย จ ติณวตฺถารเกน จ.}

\prose{วิวาทาธิกรณญฺจ อนุวาทาธิกรณญฺจ กติหิ สมเถหิ สมฺมนฺติ, กติหิ สมเถหิ น สมฺมนฺติ. วิวาทาธิกรณญฺจ อนุวาทาธิกรณญฺจ ปญฺจหิ สมเถหิ สมฺมนฺติ สมฺมุขา\-วินเยน จ เยภุยฺยสิกาย จ สติวินเยน จ อมูฬฺหวินเยน จ ตสฺสปาปิยสิกาย จ. ทฺวีหิ สมเถหิ น สมฺมนฺติ ปฏิญฺญาต\-กรเณน จ ติณวตฺถารเกน จ.}

\prose{วิวาทาธิกรณญฺจ อาปตฺตาธิกรณญฺจ กติหิ สมเถหิ สมฺมนฺติ, กติหิ สมเถหิ น สมฺมนฺติ. วิวาทาธิกรณญฺจ อาปตฺตาธิกรณญฺจ จตูหิ สมเถหิ สมฺมนฺติ สมฺมุขา\-วินเยน จ เยภุยฺยสิกาย จ ปฏิญฺญาต\-กรเณน จ ติณวตฺถารเกน จ. ตีหิ สมเถหิ น สมฺมนฺติ สติวินเยน จ อมูฬฺหวินเยน จ ตสฺสปาปิยสิกาย จ.}

\prose{วิวาทาธิกรณญฺจ กิจฺจาธิกรณญฺจ กติหิ สมเถหิ สมฺมนฺติ, กติหิ สมเถหิ น สมฺมนฺติ. วิวาทาธิกรณญฺจ กิจฺจาธิกรณญฺจ ทฺวีหิ สมเถหิ สมฺมนฺติ สมฺมุขา\-วินเยน จ เยภุยฺยสิกาย จ. ปญฺจหิ สมเถหิ น สมฺมนฺติ สติวินเยน จ อมูฬฺหวินเยน จ ปฏิญฺญาต\-กรเณน จ ตสฺสปาปิยสิกาย จ ติณวตฺถารเกน จ.}

\prose{อนุวาทาธิกรณญฺจ อาปตฺตาธิกรณญฺจ กติหิ สมเถหิ สมฺมนฺติ, กติหิ สมเถหิ น สมฺมนฺติ. อนุวาทาธิกรณญฺจ อาปตฺตาธิกรณญฺจ ฉหิ สมเถหิ สมฺมนฺติ สมฺมุขา\-วินเยน จ สติวินเยน จ อมูฬฺหวินเยน จ ปฏิญฺญาต\-กรเณน จ ตสฺสปาปิยสิกาย จ ติณวตฺถารเกน จ. เอเกน สมเถน น สมฺมนฺติ เยภุยฺยสิกาย.}

\prose{อนุวาทาธิกรณญฺจ กิจฺจาธิกรณญฺจ กติหิ สมเถหิ สมฺมนฺติ, กติหิ สมเถหิ น สมฺมนฺติ. อนุวาทาธิกรณญฺจ กิจฺจาธิกรณญฺจ จตูหิ สมเถหิ สมฺมนฺติ สมฺมุขา\-วินเยน จ สติวินเยน จ อมูฬฺหวินเยน จ \csromanfolio{199}ตสฺสปาปิยสิกาย จ. ตีหิ สมเถหิ น สมฺมนฺติ เยภุยฺยสิกาย จ ปฏิญฺญาต\-กรเณน จ ติณวตฺถารเกน จ.}

\prose{อาปตฺตาธิกรณญฺจ กิจฺจาธิกรณญฺจ กติหิ สมเถหิ สมฺมนฺติ, กติหิ สมเถหิ น สมฺมนฺติ. อาปตฺตาธิกรณญฺจ กิจฺจาธิกรณญฺจ ตีหิ สมเถหิ สมฺมนฺติ สมฺมุขา\-วินเยน จ ปฏิญฺญาต\-กรเณน จ ติณวตฺถารเกน จ. จตูหิ สมเถหิ น สมฺมนฺติ เยภุยฺยสิกาย จ สติวินเยน จ อมูฬฺหวินเยน จ ตสฺสปาปิยสิกาย จ.}

\prose{วิวาทาธิกรณญฺจ อนุวาทาธิกรณญฺจ อาปตฺตาธิกรณญฺจ กติหิ สมเถหิ สมฺมนฺติ, กติหิ สมเถหิ น สมฺมนฺติ. วิวาทาธิกรณญฺจ อนุวาทาธิกรณญฺจ อาปตฺตาธิกรณญฺจ สตฺตหิ สมเถหิ สมฺมนฺติ สมฺมุขา\-วินเยน จ เยภุยฺยสิกาย จ สติวินเยน จ อมูฬฺหวินเยน จ ปฏิญฺญาต\-กรเณน จ ตสฺสปาปิยสิกาย จ ติณวตฺถารเกน จ.}

\prose{วิวาทาธิกรณญฺจ อนุวาทาธิกรณญฺจ กิจฺจาธิกรณญฺจ กติหิ สมเถหิ สมฺมนฺติ, กติหิ สมเถหิ น สมฺมนฺติ. วิวาทาธิกรณญฺจ อนุวาทาธิกรณญฺจ กิจฺจาธิกรณญฺจ ปญฺจหิ สมเถหิ สมฺมนฺติ สมฺมุขา\-วินเยน จ เยภุยฺยสิกาย จ สติวินเยน จ อมูฬฺหวินเยน จ ตสฺสปาปิยสิกาย จ. ทฺวีหิ สมเถหิ น สมฺมนฺติ ปฏิญฺญาต\-กรเณน จ ติณวตฺถารเกน จ.}

\prose{อนุวาทาธิกรณญฺจ อาปตฺตาธิกรณญฺจ กิจฺจาธิกรณญฺจ กติหิ สมเถหิ สมฺมนฺติ, กติหิ สมเถหิ น สมฺมนฺติ. อนุวาทาธิกรณญฺจ อาปตฺตาธิกรณญฺจ กิจฺจาธิกรณญฺจ ฉหิ สมเถหิ สมฺมนฺติ สมฺมุขา\-วินเยน จ สติวินเยน จ อมูฬฺหวินเยน จ ปฏิญฺญาต\-กรเณน จ ตสฺสปาปิยสิกาย จ ติณวตฺถารเกน จ. เอเกน สมเถน น สมฺมนฺติ เยภุยฺยสิกาย.}

\prose{วิวาทาธิกรณญฺจ อนุวาทาธิกรณญฺจ อาปตฺตาธิกรณญฺจ กิจฺจาธิกรณญฺจ กติหิ สมเถหิ สมฺมนฺติ, กติหิ สมเถหิ น สมฺมนฺติ. วิวาทาธิกรณญฺจ อนุวาทาธิกรณญฺจ อาปตฺตาธิกรณญฺจ กิจฺจาธิกรณญฺจ สตฺตหิ สมเถหิ สมฺมนฺติ สมฺมุขา\-วินเยน จ เยภุยฺยสิกาย จ สติวินเยน จ อมูฬฺหวินเยน จ ปฏิญฺญาต\-กรเณน จ ตสฺสปาปิยสิกาย จ ติณวตฺถารเกน จ.}

\vspace{\tipitakaparskip}
\csromancenter{สมฺมนฺติ น สมฺมนฺติวาโร นิฏฺฐิโต อฏฺฐารสโม.}
\csromanmatikaanchor{mtk.159}
\csromantocmarkref{subsection}{19. สมถาธิกรณวาร}{mtk.159}
\bookmark[dest={mtk.159},level=2]{19. สมถาธิกรณวาร}
\csromanheader{2}{19. สมถา\-ธิกรณ\-วาร}
\proseitem{309}{\csromanfolio{200}สมถา สมเถหิ สมฺมนฺติ, สมถา อธิกรเณหิ สมฺมนฺติ, อธิกรณา สมเถหิ สมฺมนฺติ, อธิกรณา อธิกรเณหิ สมฺมนฺติ.}

\prose{สิยา สมถา สมเถหิ สมฺมนฺติ, สิยา สมถา สมเถหิ น สมฺมนฺติ, สิยา สมถา อธิกรเณหิ สมฺมนฺติ, สิยา สมถา อธิกรเณหิ น สมฺมนฺติ, สิยา อธิกรณา สมเถหิ สมฺมนฺติ, สิยา อธิกรณา สมเถหิ น สมฺมนฺติ, สิยา อธิกรณา อธิกรเณหิ สมฺมนฺติ, สิยา อธิกรณา อธิกรเณหิ น สมฺมนฺติ.}

\proseitem{310}{กถํ สิยา สมถา สมเถหิ สมฺมนฺติ, กถํ สิยา สมถา สมเถหิ น สมฺมนฺติ, เยภุยฺยสิกา สมฺมุขา\-วินเยน สมฺมติ, สติวินเยน น สมฺมติ, อมูฬฺหวินเยน น สมฺมติ, ปฏิญฺญาต\-กรเณน น สมฺมติ, ตสฺสปาปิยสิกาย น สมฺมติ, ติณวตฺถารเกน น สมฺมติ.}

\prose{สติวินโย สมฺมุขา\-วินเยน สมฺมติ, อมูฬฺหวินเยน น สมฺมติ, ปฏิญฺญาต\-กรเณน น สมฺมติ, ตสฺสปาปิยสิกาย น สมฺมติ, ติณวตฺถารเกน น สมฺมติ, เยภุยฺยสิกาย น สมฺมติ.}

\prose{อมูฬฺหวินโย สมฺมุขา\-วินเยน สมฺมติ, ปฏิญฺญาต\-กรเณน น สมฺมติ, ตสฺสปาปิยสิกาย น สมฺมติ, ติณวตฺถารเกน น สมฺมติ, เยภุยฺยสิกาย น สมฺมติ, สติวินเยน น สมฺมติ.}

\prose{ปฏิญฺญาต\-กรณํ สมฺมุขา\-วินเยน สมฺมติ, ตสฺสปาปิยสิกาย น สมฺมติ, ติณวตฺถารเกน น สมฺมติ, เยภุยฺยสิกาย น สมฺมติ, สติวินเยน น สมฺมติ, อมูฬฺหวินเยน น สมฺมติ.}

\prose{ตสฺสปาปิยสิกา สมฺมุขา\-วินเยน สมฺมติ, ติณวตฺถารเกน น สมฺมติ, เยภุยฺยสิกาย น สมฺมติ, สติวินเยน น สมฺมติ, อมูฬฺหวินเยน น สมฺมติ, ปฏิญฺญาต\-กรเณน น สมฺมติ.}

\prose{ติณวตฺถารโก สมฺมุขา\-วินเยน สมฺมติ, เยภุยฺยสิกาย น สมฺมติ, สติวินเยน น สมฺมติ, อมูฬฺหวินเยน น สมฺมติ, ปฏิญฺญาต\-กรเณน น สมฺมติ, ตสฺสปาปิยสิกาย น สมฺมติ, เอวํ สิยา สมถา สมเถหิ สมฺมนฺติ, เอวํ สิยา สมถา สมเถหิ น สมฺมนฺติ.}

\proseitem{311}{\csromanfolio{201}กถํ สิยา สมถา อธิกรเณหิ สมฺมนฺติ, กถํ สิยา สมถา อธิกรเณหิ น สมฺมนฺติ. สมฺมุขาวินโย วิวาทา\-ธิกรเณน น สมฺมติ, อนุวาทา\-ธิกรเณน น สมฺมติ, อาปตฺตา\-ธิกรเณน น สมฺมติ, กิจฺจา\-ธิกรเณน สมฺมติ.}

\prose{เยภุยฺยสิกา วิวาทา\-ธิกรเณน น สมฺมติ, อนุวาทา\-ธิกรเณน น สมฺมติ, อาปตฺตา\-ธิกรเณน น สมฺมติ, กิจฺจา\-ธิกรเณน สมฺมติ.}

\prose{สติวินโย วิวาทา\-ธิกรเณน น สมฺมติ, อนุวาทา\-ธิกรเณน น สมฺมติ. อาปตฺตา\-ธิกรเณน น สมฺมติ, กิจฺจา\-ธิกรเณน สมฺมติ.}

\prose{อมูฬฺหวินโย วิวาทา\-ธิกรเณน น สมฺมติ, อนุวาทา\-ธิกรเณน น สมฺมติ, อาปตฺตา\-ธิกรเณน น สมฺมติ, กิจฺจา\-ธิกรเณน สมฺมติ.}

\prose{ปฏิญฺญาต\-กรณํ วิวาทา\-ธิกรเณน น สมฺมติ, อนุวาทา\-ธิกรเณน น สมฺมติ, อาปตฺตา\-ธิกรเณน น สมฺมติ, กิจฺจา\-ธิกรเณน สมฺมติ.}

\prose{ตสฺสปาปิยสิกา วิวาทา\-ธิกรเณน น สมฺมติ, อนุวาทา\-ธิกรเณน น สมฺมติ, อาปตฺตา\-ธิกรเณน น สมฺมติ, กิจฺจา\-ธิกรเณน สมฺมติ.}

\prose{ติณวตฺถารโก วิวาทา\-ธิกรเณน น สมฺมติ, อนุวาทา\-ธิกรเณน น สมฺมติ, อาปตฺตา\-ธิกรเณน น สมฺมติ, กิจฺจา\-ธิกรเณน สมฺมติ. เอวํ สิยา สมถา อธิกรเณหิ สมฺมนฺติ. เอวํ สิยา สมถา อธิกรเณหิ น สมฺมนฺติ.}

\proseitem{312}{กถํ สิยา อธิกรณา สมเถหิ สมฺมนฺติ, กถํ สิยา อธิกรณา สมเถหิ น สมฺมนฺติ. วิวาทา\-ธิกรณํ สมฺมุขา\-วินเยน จ เยภุยฺยสิกาย จ สมฺมติ, สติวินเยน จ อมูฬฺหวินเยน จ ปฏิญฺญาต\-กรเณน จ ตสฺสปาปิยสิกาย จ ติณวตฺถารเกน จ น สมฺมติ.}

\prose{อนุวาทา\-ธิกรณํ สมฺมุขา\-วินเยน จ สติวินเยน จ อมูฬฺหวินเยน จ ตสฺสปาปิยสิกาย จ สมฺมติ, เยภุยฺยสิกาย จ ปฏิญฺญาต\-กรเณน จ ติณวตฺถารเกน จ น สมฺมติ.}

\prose{อาปตฺตา\-ธิกรณํ สมฺมุขา\-วินเยน จ ปฏิญฺญาต\-กรเณน จ ติณวตฺถารเกน จ สมฺมติ, เยภุยฺยสิกาย จ สติวินเยน จ อมูฬฺหวินเยน จ ตสฺสปาปิยสิกาย จ น สมฺมติ.}

\prose{\csromanfolio{202}กิจฺจา\-ธิกรณํ สมฺมุขา\-วินเยน สมฺมติ, เยภุยฺยสิกาย จ สติวินเยน จ อมูฬฺหวินเยน จ ปฏิญฺญาต\-กรเณน จ ตสฺสปาปิยสิกาย จ ติณวตฺถารเกน จ น สมฺมติ. เอวํ สิยา อธิกรณา สมเถหิ สมฺมนฺติ, เอวํ สิยา อธิกรณา สมเถหิ น สมฺมนฺติ.}

\proseitem{313}{กถํ สิยา อธิกรณา อธิกรเณหิ สมฺมนฺติ, กถํ สิยา อธิกรณา อธิกรเณหิ น สมฺมนฺติ. วิวาทา\-ธิกรณํ วิวาทา\-ธิกรเณน น สมฺมติ, อนุวาทา\-ธิกรเณน น สมฺมติ, อาปตฺตา\-ธิกรเณน น สมฺมติ, กิจฺจา\-ธิกรเณน สมฺมติ.}

\prose{อนุวาทา\-ธิกรณํ วิวาทา\-ธิกรเณน น สมฺมติ, อนุวาทา\-ธิกรเณน น สมฺมติ, อาปตฺตา\-ธิกรเณน น สมฺมติ, กิจฺจา\-ธิกรเณน สมฺมติ.}

\prose{อาปตฺตา\-ธิกรณํ วิวาทา\-ธิกรเณน น สมฺมติ, อนุวาทา\-ธิกรเณน น สมฺมติ, อาปตฺตา\-ธิกรเณน น สมฺมติ, กิจฺจา\-ธิกรเณน สมฺมติ.}

\prose{กิจฺจา\-ธิกรณํ วิวาทา\-ธิกรเณน น สมฺมติ, อนุวาทา\-ธิกรเณน น สมฺมติ, อาปตฺตา\-ธิกรเณน น สมฺมติ, กิจฺจา\-ธิกรเณน สมฺมติ. เอวํ สิยา อธิกรณา อธิกรเณหิ สมฺมนฺติ, เอวํ สิยา อธิกรณา อธิกรเณหิ น สมฺมนฺติ.}

\prosewithcloserruled{ฉาปิ สมถา จตฺตาโรปิ อธิกรณา สมฺมุขา\-วินเยน สมฺมนฺติ, สมฺมุขาวินโย น เกนจิ สมฺมติ.}{สมถา\-ธิกรณ\-วาโร นิฏฺฐิโต เอกูนวีสติโม.}
\csromanmatikaanchor{mtk.160}
\csromantocmarkref{subsection}{20. สมุฏฺฐาปนวาร}{mtk.160}
\bookmark[dest={mtk.160},level=2]{20. สมุฏฺฐาปนวาร}
\csromanheader{2}{20. สมุฏฺฐาปน\-วาร}
\proseitem{314}{วิวาทา\-ธิกรณํ จตุนฺนํ อธิกรณานํ กตมํ อธิกรณํ สมุฏฺฐาเปติ. วิวาทา\-ธิกรณํ จตุนฺนํ อธิกรณานํ น กตมํ อธิกรณํ สมุฏฺฐาเปติ, อปิจ วิวาทาธิกรณ\-ปจฺจยา จตฺตาโร อธิกรณา ชายนฺติ. ยถา กถํ วิย,\csromansymbolfootnote{*}{วิ 4. 211; วิ 5. 270 ปิฏฺเฐสุปิ.}อิธ ภิกฺขู วิวทนฺติ ธมฺโมติ วา อธมฺโมติ วา, วินโยติ วา อวินโยติ วา, อภาสิตํ อลปิตํ ตถาคเตนาติ วา ภาสิตํ ลปิตํ ตถาคเตนาติ วา, อนาจิณฺณํ ตถาคเตนาติ วา อาจิณฺณํ \csromanfolio{203}ตถาคเตนาติ วา, อปญฺญตฺตํ ตถาคเตนาติ วา ปญฺญตฺตํ ตถาคเตนาติ วา, อาปตฺตีติ วา อนาปตฺตีติ วา, ลหุกา อาปตฺตีติ วา ครุกา อาปตฺตีติ วา, สาวเสสา อาปตฺตีติ วา อนวเสสา อาปตฺตีติ วา, ทุฏฺฐุลฺลา อาปตฺตีติ วา อทุฏฺฐุลฺลา อาปตฺตีติ วา. ยํ ตตฺถ ภณฺฑนํ กลโห วิคฺคโห วิวาโท นานาวาโท อญฺญถาวาโท วิปจฺจตาย โวหาโร เมธกํ\csromansharedfootnote{71574eab26838910}{เมธคํ (ก)}, อิทํ วุจฺจติ วิวาทา\-ธิกรณํ. วิวาทา\-ธิกรเณ สํโฆ วิวทติ วิวาทา\-ธิกรณํ, วิวทมาโน อนุวทติ อนุวาทา\-ธิกรณํ, อนุวทมาโน อาปตฺติํ อาปชฺชติ อาปตฺตา\-ธิกรณํ, ตาย อาปตฺติยา สํโฆ กมฺมํ กโรติ กิจฺจา\-ธิกรณํ. เอวํ วิวาทาธิกรณ\-ปจฺจยา จตฺตาโร อธิกรณา ชายนฺติ.}

\proseitem{315}{อนุวาทา\-ธิกรณํ จตุนฺนํ อธิกรณานํ กตมํ อธิกรณํ สมุฏฺฐาเปติ. อนุวาทา\-ธิกรณํ จตุนฺนํ อธิกรณานํ น กตมํ อธิกรณํ สมุฏฺฐาเปติ, อปิจ อนุวาทาธิกรณ\-ปจฺจยา จตฺตาโร อธิกรณา ชายนฺติ. ยถา กถํ วิย,\csromansymbolfootnote{*}{วิ 4. 212; 5. 271 ปิฏฺเฐสุปิ.}อิธ ภิกฺขู ภิกฺขุํ อนุวทนฺติ สีลวิปตฺติยา วา อาจารวิปตฺติยา วา ทิฏฺฐิ\-วิปตฺติยา วา อาชีววิปตฺติยา วา, โย ตตฺถ อนุวาโท อนุวทนา อนุลฺลปนา อนุภณนา อนุสมฺปวงฺกตา อพฺภุสฺสหนตา อนุพล\-ปฺปทานํ, อิทํ วุจฺจติ อนุวาทา\-ธิกรณํ. อนุวาทา\-ธิกรเณ สํโฆ วิวทติ วิวาทา\-ธิกรณํ. วิวทมาโน อนุวทติ อนุวาทา\-ธิกรณํ. อนุวทมาโน อาปตฺติํ อาปชฺชติ อาปตฺตา\-ธิกรณํ, ตาย อาปตฺติยา สํโฆ กมฺมํ กโรติ กิจฺจา\-ธิกรณํ. เอวํ อนุวาทาธิกรณ\-ปจฺจยา จตฺตาโร อธิกรณา ชายนฺติ.}

\proseitem{316}{อาปตฺตา\-ธิกรณํ จตุนฺนํ อธิกรณานํ กตมํ อธิกรณํ สมุฏฺฐาเปติ. อาปตฺตา\-ธิกรณํ จตุนฺนํ อธิกรณานํ น กตมํ อธิกรณํ สมุฏฺฐาเปติ, อปิจ อาปตฺตาธิกรณ\-ปจฺจยา จตฺตาโร อธิกรณา ชายนฺติ. ยถา กถํ วิย, ปญฺจปิ อาปตฺติกฺขนฺธา อาปตฺตา\-ธิกรณํ, สตฺตปิ อาปตฺติกฺขนฺธา อาปตฺตา\-ธิกรณํ, อิทํ วุจฺจติ อาปตฺตา\-ธิกรณํ. อาปตฺตา\-ธิกรเณ สํโฆ วิวทติ วิวาทา\-ธิกรณํ. วิวทมาโน อนุวทติ อนุวาทา\-ธิกรณํ. อนุวทมาโน อาปตฺติํ อาปชฺชติ อาปตฺตา\-ธิกรณํ \csromanfolio{204}ตาย อาปตฺติยา สํโฆ กมฺมํ กโรติ กิจฺจา\-ธิกรณํ. เอวํ อาปตฺตาธิกรณ\-ปจฺจยา จตฺตาโร อธิกรณา ชายนฺติ.}

\proseitemwithcloserruled{317}{กิจฺจา\-ธิกรณํ จตุนฺนํ อธิกรณานํ กตมํ อธิกรณํ สมุฏฺฐาเปติ. กิจฺจา\-ธิกรณํ จตุนฺนํ อธิกรณานํ น กตมํ อธิกรณํ สมุฏฺฐาเปติ, อปิจ กิจฺจาธิกรณ\-ปจฺจยา จตฺตาโร อธิกรณา ชายนฺติ. ยถา กถํ วิย, ยา สํฆสฺส กิจฺจยตา กรณียตา อปโลกน\-กมฺมํ ญตฺติกมฺมํ ญตฺติ\-ทุติย\-กมฺมํ ญตฺติ\-จตุตฺถ\-กมฺมํ, อิทํ วุจฺจติ กิจฺจา\-ธิกรณํ. กิจฺจาธิกรเณ สํโฆ วิวทติ วิวาทา\-ธิกรณํ. วิวทมาโน อนุวทติ อนุวาทา\-ธิกรณํ. อนุวทมาโน อาปตฺติํ อาปชฺชติ อาปตฺตา\-ธิกรณํ, ตาย อาปตฺติยา สํโฆ กมฺมํ กโรติ กิจฺจา\-ธิกรณํ. เอวํ กิจฺจาธิกรณ\-ปจฺจยา จตฺตาโร อธิกรณา ชายนฺติ.}{สมุฏฺฐาปน\-วาโร นิฏฺฐิโต วีสติโม.}
\csromanmatikaanchor{mtk.161}
\csromantocmarkref{subsection}{21. ภชติวาร}{mtk.161}
\bookmark[dest={mtk.161},level=2]{21. ภชติวาร}
\csromanheader{2}{21. \textbf{ภชติวาร}}
\proseitem{318}{วิวาทา\-ธิกรณํ จตุนฺนํ อธิกรณานํ กตมํ อธิกรณํ ภชติ, กตมํ อธิกรณํ อุปนิสฺสิตํ, กตมํ อธิกรณํ ปริยาปนฺนํ, กตเมน อธิกรเณน สงฺคหิตํ.}

\prose{อนุวาทา\-ธิกรณํ จตุนฺนํ อธิกรณานํ กตมํ อธิกรณํ ภชติ, กตมํ อธิกรณํ อุปนิสฺสิตํ, กตมํ อธิกรณํ ปริยาปนฺนํ, กตเมน อธิกรเณน สงฺคหิตํ.}

\prose{อาปตฺตา\-ธิกรณํ จตุนฺนํ อธิกรณานํ กตมํ อธิกรณํ ภชติ, กตมํ อธิกรณํ อุปนิสฺสิตํ, กตมํ อธิกรณํ ปริยาปนฺนํ, กตเมน อธิกรเณน สงฺคหิตํ.}

\prose{กิจฺจา\-ธิกรณํ จตุนฺนํ อธิกรณานํ กตมํ อธิกรณํ ภชติ, กตมํ อธิกรณํ อุปนิสฺสิตํ, กตมํ อธิกรณํ ปริยาปนฺนํ, กตเมน อธิกรเณน สงฺคหิตํ.}

\prose{วิวาทา\-ธิกรณํ จตุนฺนํ อธิกรณานํ วิวาทา\-ธิกรณํ ภชติ, วิวาทา\-ธิกรณํ อุปนิสฺสิตํ, วิวาทา\-ธิกรณํ ปริยาปนฺนํ, วิวาทา\-ธิกรเณน สงฺคหิตํ.}

\prose{\csromanfolio{205}อนุวาทา\-ธิกรณํ จตุนฺนํ อธิกรณานํ อนุวาทา\-ธิกรณํ ภชติ, อนุวาทา\-ธิกรณํ อุปนิสฺสิตํ, อนุวาทา\-ธิกรณํ ปริยาปนฺนํ, อนุวาทา\-ธิกรเณน สงฺคหิตํ.}

\prose{อาปตฺตา\-ธิกรณํ จตุนฺนํ อธิกรณานํ อาปตฺตา\-ธิกรณํ ภชติ, อาปตฺตา\-ธิกรณํ อุปนิสฺสิตํ, อาปตฺตา\-ธิกรณํ ปริยาปนฺนํ, อาปตฺตา\-ธิกรเณน สงฺคหิตํ.}

\prose{กิจฺจา\-ธิกรณํ จตุนฺนํ อธิกรณานํ กิจฺจา\-ธิกรณํ ภชติ, กิจฺจา\-ธิกรณํ อุปนิสฺสิตํ, กิจฺจา\-ธิกรณํ ปริยาปนฺนํ, กิจฺจา\-ธิกรเณน สงฺคหิตํ.}

\proseitem{319}{วิวาทา\-ธิกรณํ สตฺตนฺนํ สมถานํ กติ สมเถ ภชติ, กติ สมเถ อุปนิสฺสิยํ, กติ สมเถ ปริยาปนฺนํ, กติหิ สมเถหิ สงฺคหิตํ, กติหิ สมเถหิ สมฺมติ.}

\prose{อนุวาทา\-ธิกรณํ สตฺตนฺนํ สมถานํ กติ สมเถ ภชติ, กติ สมเถ อุปนิสฺสิตํ, กติ สมเถ ปริยาปนฺนํ, กติหิ สมเถหิ สงฺคหิตํ, กติหิ สมเถหิ สมฺมติ.}

\prose{อาปตฺตา\-ธิกรณํ สตฺตนฺนํ สมถานํ กติ สมเถ ภชติ, กติ สมเถ อุปนิสฺสิตํ, กติ สมเถ ปริยาปนฺนํ, กติหิ สมเถหิ สงฺคหิตํ, กติหิ สมเถหิ สมฺมติ.}

\prose{กิจฺจา\-ธิกรณํ สตฺตนฺนํ สมถานํ กติ สมเถ ภชติ, กติ สมเถ อุปนิสฺสิยํ, กติ สมเถ ปริยาปนฺนํ, กติหิ สมเถหิ สงฺคหิตํ, กติหิ สมเถหิ สมฺมติ.}

\prose{วิวาทา\-ธิกรณํ สตฺตนฺนํ สมถานํ เทฺว สมเถ ภชติ, เทฺว สมเถ อุปนิสฺสิตํ, เทฺว สมเถ ปริยาปนฺนํ, ทฺวีหิ สมเถหิ สงฺคหิตํ, ทฺวีหิ สมเถหิ สมฺมติ สมฺมุขา\-วินเยน จ เยภุยฺยสิกาย จ.}

\prose{อนุวาทา\-ธิกรณํ สตฺตนฺนํ สมถานํ จตฺตาโร สมเถ ภชติ, จตฺตาโร สมเถ อุปนิสฺสิตํ, จตฺตาโร สมเถ ปริยาปนฺนํ, จตูหิ สมเถหิ สงฺคหิตํ, จตูหิ สมเถหิ สมฺมติ สมฺมุขา\-วินเยน จ สติวินเยน จ อมูฬฺหวินเยน จ ตสฺสปาปิยสิกาย จ.}

\prose{\csromanfolio{206}อาปตฺตา\-ธิกรณํ สตฺตนฺนํ สมถานํ ตโย สมเถ ภชติ, ตโย สมเถ อุปนิสฺสิตํ, ตโย สมเถ ปริยาปนฺนํ, ตีหิ สมเถหิ สงฺคหิตํ, ตีหิ สมเถหิ สมฺมติ สมฺมุขา\-วินเยน จ ปฏิญฺญาต\-กรเณน จ ติณวตฺถารเกน จ.}

\prose{กิจฺจา\-ธิกรณํ สตฺตนฺนํ สมถานํ เอกํ สมถํ ภชติ, เอกํ สมถํ อุปนิสฺสิตํ, เอกํ สมถํ ปริยาปนฺนํ, เอเกน สมเถน สงฺคหิตํ, เอเกน สมเถน สมฺมติ สมฺมุขาวินเยนาติ.}

\vspace{\tipitakaparskip}
\csromancenter{ภชติวาโร นิฏฺฐิโต เอกวีสติโม.}
\nitthitam{สมถเภโท นิฏฺฐิโต.}
\csromancenter{ตสฺสุทฺทานํ.}
\setlength{\csromangathapagewidth}{0pt}
\setlength{\csromangathawakwidth}{0pt}
\csromangathameasuregroup{อธิกรณํ ปริยายํ, \\ สมถา สาธารณิกา, \\ สมถา สมฺมุขา เจว, \\ ยตฺถ สมถสํสฏฺฐา, \\ สมถาธิกรณญฺเจว,}{\csromangathabat{อธิกรณํ ปริยายํ,}{สาธารณา จ ภาคิยา.} \\ \csromangathabat{สมถา สาธารณิกา,}{สมถสฺส ตพฺภาคิยา.} \\ \csromangathabat{สมถา สมฺมุขา เจว,}{วินเยน กุสเลน จ.} \\ \csromangathabat{ยตฺถ สมถสํสฏฺฐา,}{สมฺมนฺติ น สมฺมนฺติ จ.} \\ \csromangathabat{สมถาธิกรณญฺเจว,}{สมุฏฺฐานํ ภชนฺติ จาติ.}}
\csromangathasetleft{อธิกรณํ ปริยายํ, \\ สมถา สาธารณิกา, \\ สมถา สมฺมุขา เจว, \\ ยตฺถ สมถสํสฏฺฐา, \\ สมถาธิกรณญฺเจว,}
\csromangathagroup{\csromangathastanza{\csromangathabat{อธิกรณํ ปริยายํ,}{สาธารณา จ ภาคิยา.} \\ \csromangathabat{สมถา สาธารณิกา,}{สมถสฺส ตพฺภาคิยา.}}\csromangathastanza{\csromangathabat{สมถา สมฺมุขา เจว,}{วินเยน กุสเลน จ.} \\ \csromangathabat{ยตฺถ สมถสํสฏฺฐา,}{สมฺมนฺติ น สมฺมนฺติ จ.} \\ \csromangathabat{สมถาธิกรณญฺเจว,}{สมุฏฺฐานํ ภชนฺติ จาติ.}}}

\csromanmatikaanchor{mtk.162}
\csromantocmarkref{section}{ขนฺธกปุจฺฉาวาร}{mtk.162}
\bookmark[dest={mtk.162},level=1]{ขนฺธกปุจฺฉาวาร}
\csromanheader{1}{ขนฺธก\-ปุจฺฉา\-วาร}
\setlength{\csromangathapagewidth}{0pt}
\setlength{\csromangathawakwidth}{0pt}
\csromangathameasuregroup{}{อุปสมฺปทํ ปุจฺฉิสฺสํ สนิทานํ สนิทฺเทสํ, \\ สมุกฺกฏฺฐปทานํ กติ อาปตฺติโย. \\ อุปสมฺปทํ วิสฺสชฺชิสฺสํ สนิทานํ สนิทฺเทสํ, \\ สมุกฺกฏฺฐปทานํ เทฺว อาปตฺติโย. \\ อุโปสถํ ปุจฺฉิสฺสํ สนิทานํ สนิทฺเทสํ, \\ สมุกฺกฏฺฐปทานํ กติ อาปตฺติโย. \\ อุโปสถํ วิสฺสชฺชิสฺสํ สนิทานํ สนิทฺเทสํ, \\ สมุกฺกฏฺฐปทานํ ติสฺโส อาปตฺติโย. \\ วสฺสูปนายิกํ ปุจฺฉิสฺสํ สนิทานํ สนิทฺเทสํ, \\ สมุกฺกฏฺฐปทานํ กติ อาปตฺติโย. \\ วสฺสูปนายิกํ วิสฺสชฺชิสฺสํ สนิทานํ สนิทฺเทสํ, \\ สมุกฺกฏฺฐปทานํ เอกา อาปตฺติ. \\ ปวารณํ ปุจฺฉิสฺสํ สนิทานํ สนิทฺเทสํ, \\ สมุกฺกฏฺฐปทานํ กติ อาปตฺติโย. \\ ปวารณํ วิสฺสชฺชิสฺสํ สนิทานํ สนิทฺเทสํ, \\ สมุกฺกฏฺฐปทานํ ติสฺโส อาปตฺติโย. \\ จมฺมสญฺญุตฺตํ ปุจฺฉิสฺสํ สนิทานํ สนิทฺเทสํ, \\ สมุกฺกฏฺฐปทานํ กติ อาปตฺติโย. \\ จมฺมสญฺญุตฺตํ วิสฺสชฺชิสฺสํ สนิทานํ สนิทฺเทสํ, \\ สมุกฺกฏฺฐปทานํ ติสฺโส อาปตฺติโย. \\ เภสชฺชํ ปุจฺฉิสฺสํ สนิทานํ สนิทฺเทสํ, \\ สมุกฺกฏฺฐปทานํ กติ อาปตฺติโย. \\ เภสชฺชํ วิสฺสชฺชิสฺสํ สนิทานํ สนิทฺเทสํ, \\ สมุกฺกฏฺฐปทานํ ติสฺโส อาปตฺติโย. \\ กถินกํ ปุจฺฉิสฺสํ สนิทานํ สนิทฺเทสํ, \\ สมุกฺกฏฺฐปทานํ กติ อาปตฺติโย. \\ กถินกํ วิสฺสชฺชิสฺสํ สนิทานํ สนิทฺเทสํ, \\ สมุกฺกฏฺฐปทานํ นตฺถิ ตตฺถ อาปตฺติ\textsuperscript{1}. \\ จีวรสญฺญุตฺตํ ปุจฺฉิสฺสํ สนิทานํ สนิทฺเทสํ, \\ สมุกฺกฏฺฐปทานํ กติ อาปตฺติโย. \\ จีวรสญฺญุตฺตํ วิสฺสชฺชิสฺสํ สนิทานํ สนิทฺเทสํ, \\ สมุกฺกฏฺฐปทานํ ติสฺโส อาปตฺติโย. \\ จมฺเปยฺยกํ ปุจฺฉิสฺสํ สนิทานํ สนิทฺเทสํ, \\ สมุกฺขฏฺฐปทานํ กติ อาปตฺติโย. \\ จมฺเปยฺยกํ วิสฺสชฺชิสฺสํ สนิทานํ สนิทฺเทสํ, \\ สมุกฺกฏฺฐปทานํ เอกา อาปตฺติ. \\ โกสมฺพกํ ปุจฺฉิสฺสํ สนิทานํ สนิทฺเทสํ, \\ สมุกฺกฏฺฐปทานํ กติ อาปตฺติโย. \\ โกสมฺพกํ วิสฺสชฺชิสฺสํ สนิทานํ สนิทฺเทสํ, \\ สมุกฺกฏฺฐปทานํ เอกา อาปตฺติ. \\ กมฺมกฺขนฺธกํ ปุจฺฉิสฺสํ สนิทานํ สนิทฺเทสํ, \\ สมุกฺกฏฺฐปทานํ กติ อาปตฺติโย. \\ กมฺมกฺขนฺธกํ วิสฺสชฺชิสฺสํ สนิทานํ สนิทฺเทสํ, \\ สมุกฺกฏฺฐปทานํ เอกา อาปตฺติ. \\ ปาริวาสิกํ ปุจฺฉิสฺสํ สนิทานํ สนิทฺเทสํ, \\ สมุกฺกฏฺฐปทานํ กติ อาปตฺติโย. \\ ปาริวาสิกํ วิสฺสชฺชิสฺสํ สนิทานํ สนิทฺเทสํ, \\ สมุกฺกฏฺฐปทานํ เอกา อาปตฺติ. \\ สมุจฺจยํ ปุจฺฉิสฺสํ สนิทานํ สนิทฺเทสํ, \\ สมุกฺกฏฺฐปทานํ กติ อาปตฺติโย. \\ สมุจฺจยํ วิสฺสชฺชิสฺสํ สนิทานํ สนิทฺเทสํ, \\ สมุกฺกฏฺฐปทานํ เอกา อาปตฺติ. \\ สมถํ ปุจฺฉิสฺสํ สนิทานํ สนิทฺเทสํ, \\ สมุกฺกฏฺฐปทานํ กติ อาปตฺติโย. \\ สมถํ วิสฺสชฺชิสฺสํ สนิทานํ สนิทฺเทสํ, \\ สมุกฺกฏฺฐปทานํ เทฺว อาปตฺติโย. \\ ขุทฺทกวตฺถุกํ ปุจฺฉิสฺสํ สนิทานํ สนิทฺเทสํ, \\ สมุกฺกฏฺฐปทานํ กติ อาปตฺติโย. \\ ขุทฺทกวตฺถุกํ วิสฺสชฺชิสฺสํ สนิทานํ สนิทฺเทสํ, \\ สมุกฺกฏฺฐปทานํ ติสฺโส อาปตฺติโย. \\ เสนาสนํ ปุจฺฉิสฺสํ สนิทานํ สนิทฺเทสํ, \\ สมุกฺกฏฺฐปทานํ กติ อาปตฺติโย. \\ เสนาสนํ วิสฺสชฺชิสฺสํ สนิทานํ สนิทฺเทสํ, \\ สมุกฺกฏฺฐปทานํ ติสฺโส อาปตฺติโย. \\ สํฆเภทํ ปุจฺฉิสฺสํ สนิทานํ สนิทฺเทสํ, \\ สมุกฺกฏฺฐปทานํ กติ อาปตฺติโย. \\ สํฆเภทํ วิสฺสชฺชิสฺสํ สนิทานํ สนิทฺเทสํ, \\ สมุกฺกฏฺฐปทานํ เทฺว อาปตฺติโย. \\ สมาจารํ ปุจฺฉิสฺสํ สนิทานํ สนิทฺเทสํ, \\ สมุกฺกฏฺฐปทานํ กติ อาปตฺติโย. \\ สมาจารํ วิสฺสชฺชิสฺสํ สนิทานํ สนิทฺเทสํ, \\ สมุกฺกฏฺฐปทานํ เอกา อาปตฺติ. \\ ฐปนํ ปุจฺฉิสฺสํ สนิทานํ สนิทฺเทสํ, \\ สมุกฺกฏฺฐปทานํ กติ อาปตฺติโย. \\ ฐปนํ วิสฺสชฺชิสฺสํ สนิทานํ สนิทฺเทสํ, \\ สมุกฺกฏฺฐปทานํ เอกา อาปตฺติ. \\ ภิกฺขุนิกฺขนฺธกํ ปุจฺฉิสฺสํ สนิทานํ สนิทฺเทสํ, \\ สมุกฺกฏฺฐปทานํ กติ อาปตฺติโย. \\ ภิกฺขุนิกฺขนฺธกํ วิสฺสชฺชิสฺสํ สนิทานํ สนิทฺเทสํ, \\ สมุกฺกฏฺฐปทานํ เทฺว อาปตฺติโย. \\ ปญฺจสติกํ ปุจฺฉิสฺสํ สนิทานํ สนิทฺเทสํ, \\ สมุกฺกฏฺฐปทานํ กติ อาปตฺติโย. \\ ปญฺจสติกํ วิสฺสชฺชิสฺสํ สนิทานํ สนิทฺเทสํ, \\ สมุกฺกฏฺฐปทานํ นตฺถิ ตตฺถ อาปตฺติ. \\ สตฺตสติกํ ปุจฺฉิสฺสํ สนิทานํ สนิทฺเทสํ, \\ สมุกฺกฏฺฐปทานํ กติ อาปตฺติโย. \\ สตฺตสติกํ วิสฺสชฺชิสฺสํ สนิทานํ สนิทฺเทสํ, \\ สมุกฺกฏฺฐปทานํ นตฺถิ ตตฺถ อาปตฺตีติ. \\ ขนฺธกปุจฺฉาวาโร นิฏฺฐิโต ปฐโม.}
\csromangathameasurewak{อุปสมฺปทํ ปุจฺฉิสฺสํ สนิทานํ สนิทฺเทสํ, \\ สมุกฺกฏฺฐปทานํ กติ อาปตฺติโย. \\ อุปสมฺปทํ วิสฺสชฺชิสฺสํ สนิทานํ สนิทฺเทสํ, \\ สมุกฺกฏฺฐปทานํ เทฺว อาปตฺติโย. \\ อุโปสถํ ปุจฺฉิสฺสํ สนิทานํ สนิทฺเทสํ, \\ สมุกฺกฏฺฐปทานํ กติ อาปตฺติโย. \\ อุโปสถํ วิสฺสชฺชิสฺสํ สนิทานํ สนิทฺเทสํ, \\ สมุกฺกฏฺฐปทานํ ติสฺโส อาปตฺติโย. \\ วสฺสูปนายิกํ ปุจฺฉิสฺสํ สนิทานํ สนิทฺเทสํ, \\ สมุกฺกฏฺฐปทานํ กติ อาปตฺติโย. \\ วสฺสูปนายิกํ วิสฺสชฺชิสฺสํ สนิทานํ สนิทฺเทสํ, \\ สมุกฺกฏฺฐปทานํ เอกา อาปตฺติ. \\ ปวารณํ ปุจฺฉิสฺสํ สนิทานํ สนิทฺเทสํ, \\ สมุกฺกฏฺฐปทานํ กติ อาปตฺติโย. \\ ปวารณํ วิสฺสชฺชิสฺสํ สนิทานํ สนิทฺเทสํ, \\ สมุกฺกฏฺฐปทานํ ติสฺโส อาปตฺติโย. \\ จมฺมสญฺญุตฺตํ ปุจฺฉิสฺสํ สนิทานํ สนิทฺเทสํ, \\ สมุกฺกฏฺฐปทานํ กติ อาปตฺติโย. \\ จมฺมสญฺญุตฺตํ วิสฺสชฺชิสฺสํ สนิทานํ สนิทฺเทสํ, \\ สมุกฺกฏฺฐปทานํ ติสฺโส อาปตฺติโย. \\ เภสชฺชํ ปุจฺฉิสฺสํ สนิทานํ สนิทฺเทสํ, \\ สมุกฺกฏฺฐปทานํ กติ อาปตฺติโย. \\ เภสชฺชํ วิสฺสชฺชิสฺสํ สนิทานํ สนิทฺเทสํ, \\ สมุกฺกฏฺฐปทานํ ติสฺโส อาปตฺติโย. \\ กถินกํ ปุจฺฉิสฺสํ สนิทานํ สนิทฺเทสํ, \\ สมุกฺกฏฺฐปทานํ กติ อาปตฺติโย. \\ กถินกํ วิสฺสชฺชิสฺสํ สนิทานํ สนิทฺเทสํ, \\ สมุกฺกฏฺฐปทานํ นตฺถิ ตตฺถ อาปตฺติ\textsuperscript{1}. \\ จีวรสญฺญุตฺตํ ปุจฺฉิสฺสํ สนิทานํ สนิทฺเทสํ, \\ สมุกฺกฏฺฐปทานํ กติ อาปตฺติโย. \\ จีวรสญฺญุตฺตํ วิสฺสชฺชิสฺสํ สนิทานํ สนิทฺเทสํ, \\ สมุกฺกฏฺฐปทานํ ติสฺโส อาปตฺติโย. \\ จมฺเปยฺยกํ ปุจฺฉิสฺสํ สนิทานํ สนิทฺเทสํ, \\ สมุกฺขฏฺฐปทานํ กติ อาปตฺติโย. \\ จมฺเปยฺยกํ วิสฺสชฺชิสฺสํ สนิทานํ สนิทฺเทสํ, \\ สมุกฺกฏฺฐปทานํ เอกา อาปตฺติ. \\ โกสมฺพกํ ปุจฺฉิสฺสํ สนิทานํ สนิทฺเทสํ, \\ สมุกฺกฏฺฐปทานํ กติ อาปตฺติโย. \\ โกสมฺพกํ วิสฺสชฺชิสฺสํ สนิทานํ สนิทฺเทสํ, \\ สมุกฺกฏฺฐปทานํ เอกา อาปตฺติ. \\ กมฺมกฺขนฺธกํ ปุจฺฉิสฺสํ สนิทานํ สนิทฺเทสํ, \\ สมุกฺกฏฺฐปทานํ กติ อาปตฺติโย. \\ กมฺมกฺขนฺธกํ วิสฺสชฺชิสฺสํ สนิทานํ สนิทฺเทสํ, \\ สมุกฺกฏฺฐปทานํ เอกา อาปตฺติ. \\ ปาริวาสิกํ ปุจฺฉิสฺสํ สนิทานํ สนิทฺเทสํ, \\ สมุกฺกฏฺฐปทานํ กติ อาปตฺติโย. \\ ปาริวาสิกํ วิสฺสชฺชิสฺสํ สนิทานํ สนิทฺเทสํ, \\ สมุกฺกฏฺฐปทานํ เอกา อาปตฺติ. \\ สมุจฺจยํ ปุจฺฉิสฺสํ สนิทานํ สนิทฺเทสํ, \\ สมุกฺกฏฺฐปทานํ กติ อาปตฺติโย. \\ สมุจฺจยํ วิสฺสชฺชิสฺสํ สนิทานํ สนิทฺเทสํ, \\ สมุกฺกฏฺฐปทานํ เอกา อาปตฺติ. \\ สมถํ ปุจฺฉิสฺสํ สนิทานํ สนิทฺเทสํ, \\ สมุกฺกฏฺฐปทานํ กติ อาปตฺติโย. \\ สมถํ วิสฺสชฺชิสฺสํ สนิทานํ สนิทฺเทสํ, \\ สมุกฺกฏฺฐปทานํ เทฺว อาปตฺติโย. \\ ขุทฺทกวตฺถุกํ ปุจฺฉิสฺสํ สนิทานํ สนิทฺเทสํ, \\ สมุกฺกฏฺฐปทานํ กติ อาปตฺติโย. \\ ขุทฺทกวตฺถุกํ วิสฺสชฺชิสฺสํ สนิทานํ สนิทฺเทสํ, \\ สมุกฺกฏฺฐปทานํ ติสฺโส อาปตฺติโย. \\ เสนาสนํ ปุจฺฉิสฺสํ สนิทานํ สนิทฺเทสํ, \\ สมุกฺกฏฺฐปทานํ กติ อาปตฺติโย. \\ เสนาสนํ วิสฺสชฺชิสฺสํ สนิทานํ สนิทฺเทสํ, \\ สมุกฺกฏฺฐปทานํ ติสฺโส อาปตฺติโย. \\ สํฆเภทํ ปุจฺฉิสฺสํ สนิทานํ สนิทฺเทสํ, \\ สมุกฺกฏฺฐปทานํ กติ อาปตฺติโย. \\ สํฆเภทํ วิสฺสชฺชิสฺสํ สนิทานํ สนิทฺเทสํ, \\ สมุกฺกฏฺฐปทานํ เทฺว อาปตฺติโย. \\ สมาจารํ ปุจฺฉิสฺสํ สนิทานํ สนิทฺเทสํ, \\ สมุกฺกฏฺฐปทานํ กติ อาปตฺติโย. \\ สมาจารํ วิสฺสชฺชิสฺสํ สนิทานํ สนิทฺเทสํ, \\ สมุกฺกฏฺฐปทานํ เอกา อาปตฺติ. \\ ฐปนํ ปุจฺฉิสฺสํ สนิทานํ สนิทฺเทสํ, \\ สมุกฺกฏฺฐปทานํ กติ อาปตฺติโย. \\ ฐปนํ วิสฺสชฺชิสฺสํ สนิทานํ สนิทฺเทสํ, \\ สมุกฺกฏฺฐปทานํ เอกา อาปตฺติ. \\ ภิกฺขุนิกฺขนฺธกํ ปุจฺฉิสฺสํ สนิทานํ สนิทฺเทสํ, \\ สมุกฺกฏฺฐปทานํ กติ อาปตฺติโย. \\ ภิกฺขุนิกฺขนฺธกํ วิสฺสชฺชิสฺสํ สนิทานํ สนิทฺเทสํ, \\ สมุกฺกฏฺฐปทานํ เทฺว อาปตฺติโย. \\ ปญฺจสติกํ ปุจฺฉิสฺสํ สนิทานํ สนิทฺเทสํ, \\ สมุกฺกฏฺฐปทานํ กติ อาปตฺติโย. \\ ปญฺจสติกํ วิสฺสชฺชิสฺสํ สนิทานํ สนิทฺเทสํ, \\ สมุกฺกฏฺฐปทานํ นตฺถิ ตตฺถ อาปตฺติ. \\ สตฺตสติกํ ปุจฺฉิสฺสํ สนิทานํ สนิทฺเทสํ, \\ สมุกฺกฏฺฐปทานํ กติ อาปตฺติโย. \\ สตฺตสติกํ วิสฺสชฺชิสฺสํ สนิทานํ สนิทฺเทสํ, \\ สมุกฺกฏฺฐปทานํ นตฺถิ ตตฺถ อาปตฺตีติ. \\ ขนฺธกปุจฺฉาวาโร นิฏฺฐิโต ปฐโม.}
\setlength{\csromangathaleftcol}{0pt}
\csromangathagroup{\csromangathastanza{\csromangathawak{\csromanfolio{207}\csromangathaitemline{320}{อุปสมฺปทํ ปุจฺฉิสฺสํ สนิทานํ สนิทฺเทสํ,} \\ \csromangathacontline{สมุกฺกฏฺฐปทานํ กติ อาปตฺติโย.} \\ \csromangathacontline{อุปสมฺปทํ วิสฺสชฺชิสฺสํ สนิทานํ สนิทฺเทสํ,} \\ \csromangathacontline{สมุกฺกฏฺฐปทานํ เทฺว อาปตฺติโย.} \\ \csromangathacontline{อุโปสถํ ปุจฺฉิสฺสํ สนิทานํ สนิทฺเทสํ,} \\ \csromangathacontline{สมุกฺกฏฺฐปทานํ กติ อาปตฺติโย.} \\ \csromangathacontline{อุโปสถํ วิสฺสชฺชิสฺสํ สนิทานํ สนิทฺเทสํ,} \\ \csromangathacontline{สมุกฺกฏฺฐปทานํ ติสฺโส อาปตฺติโย.} \\ \csromangathacontline{วสฺสูปนายิกํ ปุจฺฉิสฺสํ สนิทานํ สนิทฺเทสํ,} \\ \csromangathacontline{สมุกฺกฏฺฐปทานํ กติ อาปตฺติโย.} \\ \csromangathacontline{วสฺสูปนายิกํ วิสฺสชฺชิสฺสํ สนิทานํ สนิทฺเทสํ,} \\ \csromangathacontline{สมุกฺกฏฺฐปทานํ เอกา อาปตฺติ.} \\ \csromangathacontline{ปวารณํ ปุจฺฉิสฺสํ สนิทานํ สนิทฺเทสํ,} \\ \csromangathacontline{สมุกฺกฏฺฐปทานํ กติ อาปตฺติโย.} \\ \csromangathacontline{ปวารณํ วิสฺสชฺชิสฺสํ สนิทานํ สนิทฺเทสํ,} \\ \csromangathacontline{สมุกฺกฏฺฐปทานํ ติสฺโส อาปตฺติโย.} \\ \csromangathacontline{จมฺมสญฺญุตฺตํ ปุจฺฉิสฺสํ สนิทานํ สนิทฺเทสํ,} \\ \csromangathacontline{สมุกฺกฏฺฐปทานํ กติ อาปตฺติโย.} \\ \csromangathacontline{จมฺมสญฺญุตฺตํ วิสฺสชฺชิสฺสํ สนิทานํ สนิทฺเทสํ,} \\ \csromangathacontline{สมุกฺกฏฺฐปทานํ ติสฺโส อาปตฺติโย.} \\ \csromangathacontline{เภสชฺชํ ปุจฺฉิสฺสํ สนิทานํ สนิทฺเทสํ,} \\ \csromangathacontline{สมุกฺกฏฺฐปทานํ กติ อาปตฺติโย.} \\ \csromangathacontline{เภสชฺชํ วิสฺสชฺชิสฺสํ สนิทานํ สนิทฺเทสํ,} \\ \csromangathacontline{สมุกฺกฏฺฐปทานํ ติสฺโส อาปตฺติโย.}}}\csromangathastanza{\csromangathawak{\csromanfolio{208}กถินกํ ปุจฺฉิสฺสํ สนิทานํ สนิทฺเทสํ, \\ สมุกฺกฏฺฐปทานํ กติ อาปตฺติโย. \\ กถินกํ วิสฺสชฺชิสฺสํ สนิทานํ สนิทฺเทสํ, \\ สมุกฺกฏฺฐปทานํ นตฺถิ ตตฺถ อาปตฺติ\csromansharedfootnote{c1d0c9ebbd9ea164}{น กตมา อาปตฺติโย (ก)}.}}\csromangathastanza{\csromangathawak{จีวรสญฺญุตฺตํ ปุจฺฉิสฺสํ สนิทานํ สนิทฺเทสํ, \\ สมุกฺกฏฺฐปทานํ กติ อาปตฺติโย. \\ จีวรสญฺญุตฺตํ วิสฺสชฺชิสฺสํ สนิทานํ สนิทฺเทสํ, \\ สมุกฺกฏฺฐปทานํ ติสฺโส อาปตฺติโย.}}\csromangathastanza{\csromangathawak{จมฺเปยฺยกํ ปุจฺฉิสฺสํ สนิทานํ สนิทฺเทสํ, \\ สมุกฺขฏฺฐปทานํ กติ อาปตฺติโย. \\ จมฺเปยฺยกํ วิสฺสชฺชิสฺสํ สนิทานํ สนิทฺเทสํ, \\ สมุกฺกฏฺฐปทานํ เอกา อาปตฺติ.}}\csromangathastanza{\csromangathawak{โกสมฺพกํ ปุจฺฉิสฺสํ สนิทานํ สนิทฺเทสํ, \\ สมุกฺกฏฺฐปทานํ กติ อาปตฺติโย. \\ โกสมฺพกํ วิสฺสชฺชิสฺสํ สนิทานํ สนิทฺเทสํ, \\ สมุกฺกฏฺฐปทานํ เอกา อาปตฺติ.}}\csromangathastanza{\csromangathawak{กมฺมกฺขนฺธกํ ปุจฺฉิสฺสํ สนิทานํ สนิทฺเทสํ, \\ สมุกฺกฏฺฐปทานํ กติ อาปตฺติโย. \\ กมฺมกฺขนฺธกํ วิสฺสชฺชิสฺสํ สนิทานํ สนิทฺเทสํ, \\ สมุกฺกฏฺฐปทานํ เอกา อาปตฺติ.}}\csromangathastanza{\csromangathawak{ปาริวาสิกํ ปุจฺฉิสฺสํ สนิทานํ สนิทฺเทสํ, \\ สมุกฺกฏฺฐปทานํ กติ อาปตฺติโย. \\ ปาริวาสิกํ วิสฺสชฺชิสฺสํ สนิทานํ สนิทฺเทสํ, \\ สมุกฺกฏฺฐปทานํ เอกา อาปตฺติ.}}\csromangathastanza{\csromangathawak{สมุจฺจยํ ปุจฺฉิสฺสํ สนิทานํ สนิทฺเทสํ, \\ สมุกฺกฏฺฐปทานํ กติ อาปตฺติโย. \\ สมุจฺจยํ วิสฺสชฺชิสฺสํ สนิทานํ สนิทฺเทสํ, \\ สมุกฺกฏฺฐปทานํ เอกา อาปตฺติ.}}\csromangathastanza{\csromangathawak{\csromanfolio{209}สมถํ ปุจฺฉิสฺสํ สนิทานํ สนิทฺเทสํ, \\ สมุกฺกฏฺฐปทานํ กติ อาปตฺติโย. \\ สมถํ วิสฺสชฺชิสฺสํ สนิทานํ สนิทฺเทสํ, \\ สมุกฺกฏฺฐปทานํ เทฺว อาปตฺติโย.}}\csromangathastanza{\csromangathawak{ขุทฺทกวตฺถุกํ ปุจฺฉิสฺสํ สนิทานํ สนิทฺเทสํ, \\ สมุกฺกฏฺฐปทานํ กติ อาปตฺติโย. \\ ขุทฺทกวตฺถุกํ วิสฺสชฺชิสฺสํ สนิทานํ สนิทฺเทสํ, \\ สมุกฺกฏฺฐปทานํ ติสฺโส อาปตฺติโย.}}\csromangathastanza{\csromangathawak{เสนาสนํ ปุจฺฉิสฺสํ สนิทานํ สนิทฺเทสํ, \\ สมุกฺกฏฺฐปทานํ กติ อาปตฺติโย. \\ เสนาสนํ วิสฺสชฺชิสฺสํ สนิทานํ สนิทฺเทสํ, \\ สมุกฺกฏฺฐปทานํ ติสฺโส อาปตฺติโย.}}\csromangathastanza{\csromangathawak{สํฆเภทํ ปุจฺฉิสฺสํ สนิทานํ สนิทฺเทสํ, \\ สมุกฺกฏฺฐปทานํ กติ อาปตฺติโย. \\ สํฆเภทํ วิสฺสชฺชิสฺสํ สนิทานํ สนิทฺเทสํ, \\ สมุกฺกฏฺฐปทานํ เทฺว อาปตฺติโย.}}\csromangathastanza{\csromangathawak{สมาจารํ ปุจฺฉิสฺสํ สนิทานํ สนิทฺเทสํ, \\ สมุกฺกฏฺฐปทานํ กติ อาปตฺติโย. \\ สมาจารํ วิสฺสชฺชิสฺสํ สนิทานํ สนิทฺเทสํ, \\ สมุกฺกฏฺฐปทานํ เอกา อาปตฺติ.}}\csromangathastanza{\csromangathawak{ฐปนํ ปุจฺฉิสฺสํ สนิทานํ สนิทฺเทสํ, \\ สมุกฺกฏฺฐปทานํ กติ อาปตฺติโย. \\ ฐปนํ วิสฺสชฺชิสฺสํ สนิทานํ สนิทฺเทสํ, \\ สมุกฺกฏฺฐปทานํ เอกา อาปตฺติ.}}\csromangathastanza{\csromangathawak{ภิกฺขุนิกฺขนฺธกํ ปุจฺฉิสฺสํ สนิทานํ สนิทฺเทสํ, \\ สมุกฺกฏฺฐปทานํ กติ อาปตฺติโย. \\ ภิกฺขุนิกฺขนฺธกํ วิสฺสชฺชิสฺสํ สนิทานํ สนิทฺเทสํ, \\ สมุกฺกฏฺฐปทานํ เทฺว อาปตฺติโย.}}\csromangathastanza{\csromangathawak{\csromanfolio{210}ปญฺจสติกํ ปุจฺฉิสฺสํ สนิทานํ สนิทฺเทสํ, \\ สมุกฺกฏฺฐปทานํ กติ อาปตฺติโย. \\ ปญฺจสติกํ วิสฺสชฺชิสฺสํ สนิทานํ สนิทฺเทสํ, \\ สมุกฺกฏฺฐปทานํ นตฺถิ ตตฺถ อาปตฺติ.}}\csromangathastanza{\csromangathawak{สตฺตสติกํ ปุจฺฉิสฺสํ สนิทานํ สนิทฺเทสํ, \\ สมุกฺกฏฺฐปทานํ กติ อาปตฺติโย. \\ สตฺตสติกํ วิสฺสชฺชิสฺสํ สนิทานํ สนิทฺเทสํ, \\ สมุกฺกฏฺฐปทานํ นตฺถิ ตตฺถ อาปตฺตีติ.}}\csromangathastanza{\csromangathawak{ขนฺธกปุจฺฉาวาโร นิฏฺฐิโต ปฐโม.}}}

\vspace{\tipitakaparskip}
\csromancenter{ตสฺสุทฺทานํ.}
\setlength{\csromangathapagewidth}{0pt}
\setlength{\csromangathawakwidth}{0pt}
\csromangathameasuregroup{อุปสมฺปทูโปสโถ, \\ จมฺมเภสชฺชกถินา, \\ โกสมฺพกฺขนฺธกํ กมฺมํ, \\ สมถขุทฺทกา เสนา, \\ ฐปนํ ภิกฺขุนิกฺขนฺธํ,}{\csromangathabat{อุปสมฺปทูโปสโถ,}{วสฺสูปนายิกปวารณา.} \\ \csromangathabat{จมฺมเภสชฺชกถินา,}{จีวรํ จมฺเปยฺยเกน จ.} \\ \csromangathabat{โกสมฺพกฺขนฺธกํ กมฺมํ,}{ปาริวาสิสมุจฺจยา.} \\ \csromangathabat{สมถขุทฺทกา เสนา,}{สํฆเภทํ สมาจาโร.} \\ \csromangathabat{ฐปนํ ภิกฺขุนิกฺขนฺธํ,}{ปญฺจสตฺตสเตน จาติ.}}
\csromangathasetleft{อุปสมฺปทูโปสโถ, \\ จมฺมเภสชฺชกถินา, \\ โกสมฺพกฺขนฺธกํ กมฺมํ, \\ สมถขุทฺทกา เสนา, \\ ฐปนํ ภิกฺขุนิกฺขนฺธํ,}
\csromangathagroup{\csromangathastanza{\csromangathabat{อุปสมฺปทูโปสโถ,}{วสฺสูปนายิกปวารณา.} \\ \csromangathabat{จมฺมเภสชฺชกถินา,}{จีวรํ จมฺเปยฺยเกน จ.}}\csromangathastanza{\csromangathabat{โกสมฺพกฺขนฺธกํ กมฺมํ,}{ปาริวาสิสมุจฺจยา.} \\ \csromangathabat{สมถขุทฺทกา เสนา,}{สํฆเภทํ สมาจาโร.} \\ \csromangathabat{ฐปนํ ภิกฺขุนิกฺขนฺธํ,}{ปญฺจสตฺตสเตน จาติ.}}}

\csromanmatikaanchor{mtk.163}
\csromantocmarknopage{section}{เอกุตฺตริกนย}{mtk.163}
\bookmark[dest={mtk.163},level=1]{เอกุตฺตริกนย}
\csromanheader{1}{\textbf{เอกุตฺตริกนย}}
\csromanmatikaanchor{mtk.164}
\csromantocmarkref{paragraph}{1. เอกกวาร}{mtk.164}
\bookmark[dest={mtk.164},level=4]{1. เอกกวาร}
\csromanheader{4}{1. \textbf{เอกกวาร}}
\proseitemwithcloser{321}{\csromanfolio{211}อาปตฺติกรา ธมฺมา ชานิตพฺพา, อนาปตฺติกรา ธมฺมา ชานิตพฺพา, อาปตฺติ ชานิตพฺพา, อนาปตฺติ ชานิตพฺพา, ลหุกา อาปตฺติ ชานิตพฺพา, ครุกา อาปตฺติ ชานิตพฺพา, สาวเสสา อาปตฺติ ชานิตพฺพา, อนวเสสา อาปตฺติ ชานิตพฺพา, ทุฏฺฐุลฺลา อาปตฺติ ชานิตพฺพา, อทุฏฺฐุลฺลา อาปตฺติ ชานิตพฺพา, สปฺปฏิกมฺมา อาปตฺติ ชานิตพฺพา, อปฺปฏิกมฺมา อาปตฺติ ชานิตพฺพา, เทสนาคามินี อาปตฺติ ชานิตพฺพา, อเทสนาคามินี อาปตฺติ ชานิตพฺพา, อนฺตรายิกา อาปตฺติ ชานิตพฺพา, อนนฺตรายิกา อาปตฺติ ชานิตพฺพา, สาวชฺช\-ปญฺญตฺติ อาปตฺติ ชานิตพฺพา, อนวชฺช\-ปญฺญตฺติ อาปตฺติ ชานิตพฺพา, กิริยโต สมุฏฺฐิตา อาปตฺติ ชานิตพฺพา, อกิริยโต สมุฏฺฐิตา อาปตฺติ ชานิตพฺพา, กิริยากิริยโต สมุฏฺฐิตา อาปตฺติ ชานิตพฺพา, ปุพฺพาปตฺติ ชานิตพฺพา, อปราปตฺติ ชานิตพฺพา, ปุพฺพาปตฺตีนํ อนฺตราปตฺติ ชานิตพฺพา, อปราปตฺตีนํ อนฺตราปตฺตี ชานิตพฺพา, เทสิตา คณนูปคา อาปตฺติ ชานิตพฺพา, เทสิตา น คณนูปคา อาปตฺติ ชานิตพฺพา, ปญฺญตฺติ ชานิตพฺพา, อนุปญฺญตฺติ ชานิตพฺพา, อนุปฺปนฺนปญฺญตฺติ ชานิตพฺพา, สพฺพตฺถ\-ปญฺญตฺติ ชานิตพฺพา, ปเทสปญฺญตฺติ ชานิตพฺพา, สาธารณ\-ปญฺญตฺติ ชานิตพฺพา, อสาธารณ\-ปญฺญตฺติ ชานิตพฺพา, เอกโตปญฺญตฺติ ชานิตพฺพา, อุภโตปญฺญตฺติ ชานิตพฺพา, ถุลฺลวชฺชา อาปตฺติ ชานิตพฺพา, อถุลฺลวชฺชา อาปตฺติ ชานิตพฺพา, คิหิปฏิสํยุตฺตา อาปตฺติ ชานิตพฺพา, น คิหิปฏิสํยุตฺตา อาปตฺติ ชานิตพฺพา, นิยตา อาปตฺติ ชานิตพฺพา, อนิยตา อาปตฺติ ชานิตพฺพา, อาทิกโร ปุคฺคโล ชานิตพฺโพ, อนาทิกโร ปุคฺคโล ชานิตพฺโพ, อธิจฺจาปตฺติโก ปุคฺคโล ชานิตพฺโพ, อภิณฺหาปตฺติโก ปุคฺคโล ชานิตพฺโพ, โจทโก ปุคฺคโล ชานิตพฺโพ, จุทิตโก ปุคฺคโล ชานิตพฺโพ, อธมฺมโจทโก ปุคฺคโล ชานิตพฺโพ, อธมฺม\-จุทิตโก ปุคฺคโล ชานิตพฺโพ, ธมฺมโจทโก ปุคฺคโล ชานิตพฺโพ, ธมฺมจุทิตโก ปุคฺคโล ชานิตพฺโพ, นิยโต ปุคฺคโล ชานิตพฺโพ, อนิยโต ปุคฺคโล ชานิตพฺโพ, อภพฺพาปตฺติโก ปุคฺคโล ชานิตพฺโพ, ภพฺพาปตฺติโก ปุคฺคโล ชานิตพฺโพ, อุกฺขิตฺตโก ปุคฺคโล ชานิตพฺโพ, อนุกฺขิตฺตโก ปุคฺคโล ชานิตพฺโพ, นาสิตโก \csromanfolio{212}ปุคฺคโล ชานิตพฺโพ, อนาสิตโก ปุคฺคโล ชานิตพฺโพ, สมาน\-สํวาสโก ปุคฺคโล ชานิตพฺโพ, อสมาน\-สํวาสโก ปุคฺคโล ชานิตพฺโพ, ฐปนํ ชานิตพฺพนฺติ.}{เอกกํ นิฏฺฐิตํ.}
\csromancenter{ตสฺสุทฺทานํ.}
\setlength{\csromangathapagewidth}{0pt}
\setlength{\csromangathawakwidth}{0pt}
\csromangathameasuregroup{กรา อาปตฺติ ลหุกา, \\ ปฏิกมฺมเทสนา จ, \\ กิริยากิริยํ ปุพฺพา, \\ ปญฺญตฺติ อนานุปฺปนฺน, \\ ถุลฺลคิหินิยตา จ, \\ อธมฺมธมฺมนิยโต,}{\csromangathabat{กรา อาปตฺติ ลหุกา,}{สาวเสสา จ ทุฏฺฐุลฺลา.} \\ \csromangathabat{ปฏิกมฺมเทสนา จ,}{อนฺตรา วชฺชกิริยํ.} \\ \csromangathabat{กิริยากิริยํ ปุพฺพา,}{อนฺตรา คณนูปคา.} \\ \csromangathabat{ปญฺญตฺติ อนานุปฺปนฺน,}{สพฺพสาธารณา จ เอกโต\textsuperscript{1}.} \\ \csromangathabat{ถุลฺลคิหินิยตา จ,}{อาทิ อธิจฺจโจทโก.} \\ \csromangathabat{อธมฺมธมฺมนิยโต,}{อภพฺโพกฺขิตฺตนาสิตา.}}
\csromangathasetleft{กรา อาปตฺติ ลหุกา, \\ ปฏิกมฺมเทสนา จ, \\ กิริยากิริยํ ปุพฺพา, \\ ปญฺญตฺติ อนานุปฺปนฺน, \\ ถุลฺลคิหินิยตา จ, \\ อธมฺมธมฺมนิยโต,}
\csromangathagroup{\csromangathastanza{\csromangathabat{กรา อาปตฺติ ลหุกา,}{สาวเสสา จ ทุฏฺฐุลฺลา.} \\ \csromangathabat{ปฏิกมฺมเทสนา จ,}{อนฺตรา วชฺชกิริยํ.}}\csromangathastanza{\csromangathabat{กิริยากิริยํ ปุพฺพา,}{อนฺตรา คณนูปคา.} \\ \csromangathabat{ปญฺญตฺติ อนานุปฺปนฺน,}{สพฺพสาธารณา จ เอกโต\csromansharedfootnote{7caf631444c50a23}{ปญฺญตฺตานุปฺปนฺนา สพฺพา, สาธารณา จ เอกโต (สฺยา)}.}}\csromangathastanza{\csromangathabat{ถุลฺลคิหินิยตา จ,}{อาทิ อธิจฺจโจทโก.} \\ \csromangathabat{อธมฺมธมฺมนิยโต,}{อภพฺโพกฺขิตฺตนาสิตา.}}}

\vspace{\tipitakaparskip}
\csromancenter{สมานํ ฐปนํ เจว, อุทฺทานํ เอกเก อิทนฺติ.}
\csromansectionrule
\csromanmatikaanchor{mtk.165}
\csromantocmarkref{paragraph}{2. ทุกวาร}{mtk.165}
\bookmark[dest={mtk.165},level=4]{2. ทุกวาร}
\csromanheader{4}{2. \textbf{ทุกวาร}}
\proseitem{322}{อตฺถาปตฺติ สญฺญา วิโมกฺขา, อตฺถาปตฺติ โน สญฺญาวิโมกฺขา. อตฺถาปตฺติ ลทฺธ\-สมาปตฺติกสฺส, อตฺถาปตฺติ น ลทฺธ\-สมาปตฺติกสฺส. อตฺถาปตฺติ สทฺธมฺม\-ปฏิสญฺญุตฺตา, อตฺถาปตฺติ อสทฺธมฺม\-ปฏิสญฺญุตฺตา. อตฺถาปตฺติ สปริกฺขาร\-ปฏิสญฺญุตฺตา, อตฺถาปตฺติ ปรปริกฺขาร\-ปฏิสญฺญุตฺตา. อตฺถาปตฺติ สปุคฺคล\-ปฏิสญฺญุตฺตา, อตฺถาปตฺติ ปรปุคฺคล\-ปฏิสญฺญุตฺตา. อตฺถิ สจฺจํ ภณนฺโต ครุกํ อาปตฺติํ อาปชฺชติ มุสา ภณนฺโต ลหุกํ, อตฺถิ มุสา ภณนฺโต ครุกํ อาปตฺติํ อาปชฺชติ สจฺจํ ภณนฺโต ลหุกํ. อตฺถาปตฺติ ภูมิคโต อาปชฺชติ โน เวหาสคโต, อตฺถาปตฺติ เวหาสคโต อาปชฺชติ โน ภูมิคโต. อตฺถาปตฺติ นิกฺขมนฺโต อาปชฺชติ โน ปวิสนฺโต, อตฺถาปตฺติ ปวิสนฺโต อาปชฺชติ โน นิกฺขมนฺโต. อตฺถาปตฺติ อาทิยนฺโต อาปชฺชติ, อตฺถาปตฺติ อนาทิยนฺโต อาปชฺชติ. \csromanfolio{213}อตฺถาปตฺติ สมาทิยนฺโต อาปชฺชติ, อตฺถาปตฺติ น สมาทิยนฺโต อาปชฺชติ. อตฺถาปตฺติ กโรนฺโต อาปชฺชติ, อตฺถาปตฺติ น กโรนฺโต อาปชฺชติ. อตฺถาปตฺติ เทนฺโต อาปชฺชติ, อตฺถาปตฺติ น เทนฺโต อาปชฺชติ. (อตฺถาปตฺติ เทเสนฺโต อาปชฺชติ, อตฺถาปตฺติ น เทเสนฺโต อาปชฺชติ.)\csromansharedfootnote{7b8eb3a0940593fb}{( ) (นตฺถิ กตฺถจิ)} อตฺถาปตฺติ ปฏิคฺคณฺหนฺโต อาปชฺชติ, อตฺถาปตฺติ น ปฏิคฺคณฺหนฺโต อาปชฺชติ. อตฺถาปตฺติ ปริโภเคน อาปชฺชติ, อตฺถาปตฺติ น ปริโภเคน อาปชฺชติ. อตฺถาปตฺติ รตฺติํ อาปชฺชติ, โน ทิวา, อตฺถาปตฺติ ทิวา อาปชฺชติ โน รตฺติํ. อตฺถาปตฺติ อรุณุคฺเค อาปชฺชติ, อตฺถาปตฺติ น อรุณุคฺเค อาปชฺชติ. อตฺถาปตฺติ ฉินฺทนฺโต อาปชฺชติ, อตฺถาปตฺติ น ฉินฺทนฺโต อาปชฺชติ. อตฺถาปตฺติ ฉาเทนฺโต อาปชฺชติ, อตฺถาปตฺติ น ฉาเทนฺโต อาปชฺชติ. อตฺถาปตฺติ ธาเรนฺโต อาปชฺชติ, อตฺถาปตฺติ น ธาเรนฺโต อาปชฺชติ.}

\prose{เทฺว อุโปสถา จาตุทฺทสิโก จ ปนฺนรสิโก จ. เทฺว ปวารณา จาตุทฺทสิกา จ ปนฺนรสิกา จ. เทฺว กมฺมานิ อปโลกน\-กมฺมํ ญตฺติกมฺมํ. อปรานิปิ เทฺว กมฺมานิ ญตฺติ\-ทุติย\-กมฺมํ ญตฺติ\-จตุตฺถ\-กมฺมํ. เทฺว กมฺมวตฺถูนิ อปโลกน\-กมฺมสฺส วตฺถุ ญตฺติกมฺมสฺส วตฺถุ. อปรานิปิ เทฺว กมฺมวตฺถูนิ ญตฺติ\-ทุติย\-กมฺมสฺส วตฺถุ ญตฺติ\-จตุตฺถ\-กมฺมสฺส วตฺถุ. เทฺว กมฺมโทสา อปโลกน\-กมฺมสฺส โทโส ญตฺติกมฺมสฺส โทโส. อปเรปิ เทฺว กมฺมโทสา ญตฺติ\-ทุติย\-กมฺมสฺส โทโส ญตฺติ\-จตุตฺถ\-กมฺมสฺส โทโส. เทฺว กมฺม\-สมฺปตฺติโย อปโลกน\-กมฺมสฺส สมฺปตฺติ ญตฺติกมฺมสฺส สมฺปตฺติ. อปราปิ เทฺว กมฺม\-สมฺปตฺติโย ญตฺติ\-ทุติย\-กมฺมสฺส สมฺปตฺติ ญตฺติ\-จตุตฺถ\-กมฺมสฺส สมฺปตฺติ. เทฺว นานาสํวาสก\-ภูมิโย อตฺตนา วา อตฺตานํ นานาสํวาสกํ กโรติ, สมคฺโค วา นํ สํโฆ อุกฺขิปติ อทสฺสเน วา อปฺปฏิกมฺเม วา อปฺปฏินิสฺสคฺเค วา. เทฺว สมานสํวาสก\-ภูมิโย อตฺตนา วา อตฺตานํ สมาน\-สํวาสกํ กโรติ, สมคฺโค วา นํ สํโฆ อุกฺขิตฺตํ โอสาเรติ ทสฺสเน วา ปฏิกมฺเม วา ปฏินิสฺสคฺเค วา. เทฺว ปาราชิกา ภิกฺขูนญฺจ ภิกฺขุนีนญฺจ. เทฺว สํฆาทิเสสา. เทฺว ถุลฺลจฺจยา. เทฺว ปาจิตฺติยา. เทฺว ปาฏิเทสนียา. เทฺว ทุกฺกฏา. เทฺว ทุพฺภาสิตา ภิกฺขูนญฺจ ภิกฺขุนีนญฺจ. สตฺต อาปตฺติโย, สตฺต อาปตฺติกฺขนฺธา. ทฺวีหากาเรหิ สํโฆ ภิชฺชติ กมฺเมน วา สลากคฺคาเหน วา.}

\prose{\csromanfolio{214}เทฺว ปุคฺคลา น อุปสมฺ\-ปาเท\-ตพฺพา อทฺธานหีโน องฺคหีโน. อปเรปิ เทฺว ปุคฺคลา น อุปสมฺ\-ปาเท\-ตพฺพา วตฺถุวิปนฺโน กรณ\-ทุกฺกฏโก. อปเรปิ เทฺว ปุคฺคลา น อุปสมฺ\-ปาเท\-ตพฺพา อปริปูโร ปริปูโร โน จ ยาจติ. ทฺวินฺนํ ปุคฺคลานํ นิสฺสาย น วตฺถพฺพํ อลชฺชิสฺส จ พาลสฺส จ. ทฺวินฺนํ ปุคฺคลานํ นิสฺสโย น ทาตพฺโพ อลชฺชิสฺส จ ลชฺชิโน จ น ยาจติ. ทฺวินฺนํ ปุคฺคลานํ นิสฺสโย ทาตพฺโพ พาลสฺส จ ลชฺชิสฺส จ ยาจติ. เทฺว ปุคฺคลา อภพฺพา อาปตฺติํ อาปชฺชิตุํ พุทฺธา จ ปจฺเจกพุทฺธา จ. เทฺว ปุคฺคลา ภพฺพา อาปตฺติํ อาปชฺชิตุํ ภิกฺขู จ ภิกฺขุนิโย จ. เทฺว ปุคฺคลา อภพฺพา สญฺจิจฺจ อาปตฺติํ อาปชฺชิตุํ ภิกฺขู จ ภิกฺขุนิโย จ อริยปุคฺคลา. เทฺว ปุคฺคลา ภพฺพา สญฺจิจฺจ อาปตฺติํ อาปชฺชิตุํ ภิกฺขู จ ภิกฺขุนิโย จ ปุถุชฺชนา. เทฺว ปุคฺคลา อภพฺพา สญฺจิจฺจ สาติสารํ วตฺถุํ อชฺฌาจริตุํ ภิกฺขู จ ภิกฺขุนิโย จ อริยปุคฺคลา. เทฺว ปุคฺคลา ภพฺพา สญฺจิจฺจ สาติสารํ วตฺถุํ อชฺฌาจริตุํ ภิกฺขู จ ภิกฺขุนิโย จ ปุถุชฺชนา.}

\prose{เทฺว ปฏิกฺโกสา กาเยน วา ปฏิกฺโกสติ, วาจาย วา ปฏิกฺโกสติ. เทฺว นิสฺสารณา อตฺถิ ปุคฺคโล อปฺปตฺโต นิสฺสารณํ, ตํ เจ สํโฆ นิสฺสาเรติ, เอกจฺโจ สุนิสฺสาริโต เอกจฺโจ ทุนฺนิสฺสาริโต. เทฺว โอสารณา อตฺถิ ปุคฺคโล อปฺปตฺโต โอสารณํ, ตํ เจ สํโฆ โอสาเรติ, เอกจฺโจ โสสาริโต เอกจฺโจ โทสาริโต. เทฺว ปฏิญฺญา กาเยน วา ปฏิชานาติ, วาจาย วา ปฏิชานาติ. เทฺว ปฏิคฺคหา กาเยน วา ปฏิคฺคณฺหาติ, กาย\-ปฏิพทฺเธน วา ปฏิคฺคณฺหาติ. เทฺว ปฏิกฺเขปา กาเยน วา ปฏิกฺขิปติ, วาจาย วา ปฏิกฺขิปติ. เทฺว อุปฆาติกา สิกฺขู\-ปฆาติกา จ โภคูปฆาติกา จ. เทฺว โจทนา กาเยน วา โจเทติ, วาจาย วา โจเทติ. เทฺว กถินสฺส ปลิโพธา อาวาสปลิโพโธ จ จีวร\-ปลิโพโธ จ. เทฺว กถินสฺส อปลิโพธา อาวาส\-อปลิโพโธ จ จีวร\-อปลิโพโธ จ. เทฺว จีวรานิ คหปติกญฺจ ปํสุกูลญฺจ. เทฺว ปตฺตา อโยปตฺโต มตฺติกาปตฺโต. เทฺว มณฺฑลานิ ติปุมยํ สีสมยํ. เทฺว ปตฺตสฺส อธิฏฺฐานา กาเยน วา อธิฏฺเฐติ, วาจาย วา อธิฏฺเฐติ. เทฺว จีวรสฺส อธิฏฺฐานา กาเยน วา อธิฏฺเฐติ, วาจาย วา อธิฏฺเฐติ. เทฺว วิกปฺปนา สมฺมุขา\-วิกปฺปนา จ ปรมฺมุขา\-วิกปฺปนา จ. เทฺว วินยา ภิกฺขูนญฺจ ภิกฺขุนีนญฺจ. \csromanfolio{215}เทฺว เวนยิกา ปญฺญตฺตญฺจ ปญฺญตฺตานุโลมญฺจ. เทฺว วินยสฺส สลฺเลขา อกปฺปิเย เสตุฆาโต กปฺปิเย มตฺตการิตา. ทฺวีหากาเรหิ อาปตฺติํ อาปชฺชติ กาเยน วา อาปชฺชติ, วาจาย วา อาปชฺชติ. ทฺวีหากาเรหิ อาปตฺติยา วุฏฺฐาติ กาเยน วา วุฏฺฐาติ, วาจาย วา วุฏฺฐาติ. เทฺว ปริวาสา ปฏิจฺฉนฺน\-ปริวาโส อปฺปฏิจฺฉนฺน\-ปริวาโส. อปเรปิ เทฺว ปริวาสา สุทฺธนฺต\-ปริวาโส สโมธาน\-ปริวาโส. เทฺว มานตฺตา ปฏิจฺฉนฺน\-มานตฺตํ อปฺปฏิจฺฉนฺนมานตํ. อปเรปิ เทฺว มานตฺตา ปกฺขมานตฺตํ สโมธาน\-มานตฺตํ. ทฺวินฺนํ ปุคฺคลานํ รตฺติจฺเฉโท ปาริวาสิกสฺส จ มานตฺต\-จาริกสฺส จ. เทฺว อนาทริยานิ ปุคฺคลานาทริยญฺจ ธมฺมานาทริยญฺจ. เทฺว โลณานิ ชาติมญฺจ การิมญฺจ\csromansharedfootnote{bdec5b03373a3fa7}{ชาติมยญฺจ ขาริมยญฺจ (สฺยา)}. อปรานิปิ เทฺว โลณานิ สามุทฺทํ กาฬโลณํ. อปรานิปิ เทฺว โลณานิ สินฺธวํ อุพฺภิทํ\csromansharedfootnote{a3738bb796a4883e}{อุพฺภิรํ (อิติปิ)}. อปรานิปิ เทฺว โลณานิ โรมกํ ปกฺกาลกํ. เทฺว ปริโภคา อพฺภนฺตร\-ปริโภโค จ พาหิร\-ปริโภโค จ. เทฺว อกฺโกสา หีโน จ อกฺโกโส อุกฺกฏฺโฐ จ อกฺโกโส. ทฺวีหากาเรหิ เปสุญฺญํ โหติ ปิยกมฺยสฺส วา เภทา\-ธิปฺปายสฺส วา. ทฺวีหากาเรหิ คณโภชนํ ปสวติ นิมนฺตนโต วา วิญฺญตฺติโต วา. เทฺว วสฺสูปนายิกา ปุริมิกา ปจฺฉิมิกา. เทฺว อธมฺมิกานิ ปาติโมกฺข\-ฏฺฐปนานิ. เทฺว ธมฺมิกานิ ปาติโมกฺข\-ฏฺฐปนานิ.}

\prose{\csromansymbolfootnote{*}{อํ 1. 82 ปิฏฺฐาทีสุปิ.}เทฺว ปุคฺคลา พาลา โย จ อนาคตํ ภารํ วหติ, โย จ อาคตํ ภารํ น วหติ. เทฺว ปุคฺคลา ปณฺฑิตา โย จ อนาคตํ ภารํ น วหติ, โย จ อาคตํ ภารํ วหติ. อปเรปิ เทฺว ปุคฺคลา พาลา โย จ อกปฺปิเย กปฺปิยสญฺญี, โย จ กปฺปิเย อกปฺปิยสญฺญี. เทฺว ปุคฺคลา ปณฺฑิตา โย จ อกปฺปิเย อกปฺปิยสญฺญี, โย จ กปฺปิเย กปฺปิยสญฺญี. อปเรปิ เทฺว ปุคฺคลา พาลา โย จ อนาปตฺติยา อาปตฺติสญฺญี, โย จ อาปตฺติยา อนาปตฺติสญฺญี. เทฺว ปุคฺคลา ปณฺฑิตา โย จ อาปตฺติยา อาปตฺติสญฺญี, โย จ อนาปตฺติยา อนาปตฺติสญฺญี. อปเรปิ เทฺว ปุคฺคลา พาลา โย จ อธมฺเม ธมฺมสญฺญี, โย จ ธมฺเม อธมฺมสญฺญี. เทฺว ปุคฺคลา ปณฺฑิตา โย จ อธมฺเม อธมฺมสญฺญี, โย จ ธมฺเม ธมฺมสญฺญี. อปเรปิ เทฺว ปุคฺคลา พาลา โย จ อวินเย วินยสญฺญี, \csromanfolio{216}โย จ วินเย อวินยสญฺญี. เทฺว ปุคฺคลา ปณฺฑิตา โย จ อวินเย อวินยสญฺญี, โย จ วินเย วินยสญฺญี.}

\prosewithcloser{\csromansymbolfootnote{*}{อํ 1. 84 ปิฏฺเฐปิ.}ทฺวินฺนํ ปุคฺคลานํ อาสวา วฑฺฒนฺติ โย จ น กุกฺกุจฺจายิ\-ตพฺพํ กุกฺกุจฺจายติ, โย จ กุกฺกุจฺจายิ\-ตพฺพํ น กุกฺกุจฺจายติ. ทฺวินฺนํ ปุคฺคลานํ อาสวา น วฑฺฒนฺติ โย จ น กุกฺกุจฺจายิ\-ตพฺพํ น กุกฺกุจฺจายติ, โย จ กุกฺกุจฺจายิ\-ตพฺพํ กุกฺกุจฺจายติ. อปเรสมฺปิ ทฺวินฺนํ ปุคฺคลานํ อาสวา วฑฺฒนฺติ โย จ อกปฺปิเย กปฺปิยสญฺญี, โย จ กปฺปิเย อกปฺปิยสญฺญี. ทฺวินฺนํ ปุคฺคลานํ อาสวา น วฑฺฒนฺติ โย จ อกปฺปิเย อกปฺปิยสญฺญี, โย จ กปฺปิเย กปฺปิยสญฺญี. อปเรสมฺปิ ทฺวินฺนํ ปุคฺคลานํ อาสวา วฑฺฒนฺติ โย จ อนาปตฺติยา อาปตฺติสญฺญี, โย จ อาปตฺติยา อนาปตฺติสญฺญี. ทฺวินฺนํ ปุคฺคลานํ อาสวา น วฑฺฒนฺติ โย จ อนาปตฺติยา อนาปตฺติสญฺญี, โย จ อาปตฺติยา อาปตฺติสญฺญี, อปเรสมฺปิ ทฺวินฺนํ ปุคฺคลานํ อาสวา วฑฺฒนฺติ. โย จ อธมฺเม ธมฺมสญฺญี, โย จ ธมฺเม อธมฺมสญฺญี. ทฺวินฺนํ ปุคฺคลานํ อาสวา น วฑฺฒนฺติ โย จ อธมฺเม อธมฺมสญฺญี, โย จ ธมฺเม ธมฺมสญฺญี. อปเรสมฺปิ ทฺวินฺนํ ปุคฺคลานํ อาสวา วฑฺฒนฺติ โย จ อวินเย วินยสญฺญี, โย จ วินเย อวินยสญฺญี. ทฺวินฺนํ ปุคฺคลานํ อาสวา น วฑฺฒนฺติ โย จ อวินเย อวินยสญฺญี, โย จ วินเย วินยสญฺญี.}{ทุกา นิฏฺฐิตา.}
\csromancenter{ตสฺสุทฺทานํ.}
\setlength{\csromangathapagewidth}{0pt}
\setlength{\csromangathawakwidth}{0pt}
\csromangathameasuregroup{สญฺญา ลทฺธา จ สทฺธมฺมา, \\ สจฺจํ ภูมิ นิกฺขมนฺโต, \\ กโรนฺโต เทนฺโต คณฺหนฺโต, \\ อรุณาฉินฺทํ ฉาเทนฺโต, \\ ปวารณา กมฺมาปรา, \\ อปรา เทฺว จ สมฺปตฺติ, \\ ปาราชิสํฆถุลฺลจฺจย, \\ ทุกฺกฏา ทุพฺภาสิตา เจว, \\ ภิชฺชติ อุปสมฺปทา, \\ น วตฺถพฺพํ น ทาตพฺพํ, \\ สญฺจิจฺจ สาติสารา จ, \\ โอสารณา ปฏิญฺญา จ, \\ อุปฆาติ โจทนา จ, \\ จีวรา ปตฺตมณฺฑลา, \\ วิกปฺปนา จ วินยา, \\ อาปชฺชติ จ วุฏฺฐาติ, \\ เทฺว มานตฺตา อปเร จ, \\ เทฺว โลณา ตโย อปเร, \\ เปสุญฺโญ จ คณาวสฺส,}{\csromangathabat{สญฺญา ลทฺธา จ สทฺธมฺมา,}{ปริกฺขารา จ ปุคฺคลา.} \\ \csromangathabat{สจฺจํ ภูมิ นิกฺขมนฺโต,}{อาทิยนฺโต สมาทิยํ.} \\ \csromangathabat{กโรนฺโต เทนฺโต คณฺหนฺโต,}{ปริโภเคน รตฺติ จ.} \\ \csromangathabat{อรุณาฉินฺทํ ฉาเทนฺโต,}{ธาเรนฺโต จ อุโปสถา.} \\ \csromangathabat{ปวารณา กมฺมาปรา,}{วตฺถุ อปรา โทสา จ.} \\ \csromangathabat{อปรา เทฺว จ สมฺปตฺติ,}{นานา สมานเมว จ.} \\ \csromangathabat{ปาราชิสํฆถุลฺลจฺจย,}{ปาจิตฺติ ปาฏิเทสนา.} \\ \csromangathabat{ทุกฺกฏา ทุพฺภาสิตา เจว,}{สตฺต อาปตฺติกฺขนฺธา จ.} \\ \csromangathabat{ภิชฺชติ อุปสมฺปทา,}{ตเถว อปเร ทุเว.} \\ \csromangathabat{น วตฺถพฺพํ น ทาตพฺพํ,}{อภพฺพาภพฺพเมว จ.} \\ \csromangathabat{สญฺจิจฺจ สาติสารา จ,}{ปฏิกฺโกสา นิสฺสารณา.} \\ \csromangathabat{โอสารณา ปฏิญฺญา จ,}{ปฏิคฺคหา ปฏิกฺขิปา.} \\ \csromangathabat{อุปฆาติ โจทนา จ,}{กถินา จ ทุเว ตถา.} \\ \csromangathabat{จีวรา ปตฺตมณฺฑลา,}{อธิฏฺฐานา ตเถว เทฺว.} \\ \csromangathabat{วิกปฺปนา จ วินยา,}{เวนยิกา จ สลฺเลขา.} \\ \csromangathabat{อาปชฺชติ จ วุฏฺฐาติ,}{ปริวาสาปเร ทุเว.} \\ \csromangathabat{เทฺว มานตฺตา อปเร จ,}{รตฺติจฺเฉโท อนาทริ.} \\ \csromangathabat{เทฺว โลณา ตโย อปเร,}{ปริโภคกฺโกเสน จ.} \\ \csromangathabat{เปสุญฺโญ จ คณาวสฺส,}{ฐปนา ภารกปฺปิยํ.}}
\csromangathasetleft{สญฺญา ลทฺธา จ สทฺธมฺมา, \\ สจฺจํ ภูมิ นิกฺขมนฺโต, \\ กโรนฺโต เทนฺโต คณฺหนฺโต, \\ อรุณาฉินฺทํ ฉาเทนฺโต, \\ ปวารณา กมฺมาปรา, \\ อปรา เทฺว จ สมฺปตฺติ, \\ ปาราชิสํฆถุลฺลจฺจย, \\ ทุกฺกฏา ทุพฺภาสิตา เจว, \\ ภิชฺชติ อุปสมฺปทา, \\ น วตฺถพฺพํ น ทาตพฺพํ, \\ สญฺจิจฺจ สาติสารา จ, \\ โอสารณา ปฏิญฺญา จ, \\ อุปฆาติ โจทนา จ, \\ จีวรา ปตฺตมณฺฑลา, \\ วิกปฺปนา จ วินยา, \\ อาปชฺชติ จ วุฏฺฐาติ, \\ เทฺว มานตฺตา อปเร จ, \\ เทฺว โลณา ตโย อปเร, \\ เปสุญฺโญ จ คณาวสฺส,}
\csromangathagroup{\csromangathastanza{\csromangathabat{สญฺญา ลทฺธา จ สทฺธมฺมา,}{ปริกฺขารา จ ปุคฺคลา.} \\ \csromangathabat{สจฺจํ ภูมิ นิกฺขมนฺโต,}{อาทิยนฺโต สมาทิยํ.}}\csromangathastanza{\csromangathabat{กโรนฺโต เทนฺโต คณฺหนฺโต,}{ปริโภเคน รตฺติ จ.} \\ \csromangathabat{อรุณาฉินฺทํ ฉาเทนฺโต,}{ธาเรนฺโต จ อุโปสถา.}}\csromangathastanza{\csromangathabat{ปวารณา กมฺมาปรา,}{วตฺถุ อปรา โทสา จ.} \\ \csromangathabat{อปรา เทฺว จ สมฺปตฺติ,}{นานา สมานเมว จ.}}\csromangathastanza{\csromangathabat{ปาราชิสํฆถุลฺลจฺจย,}{ปาจิตฺติ ปาฏิเทสนา.} \\ \csromangathabat{ทุกฺกฏา ทุพฺภาสิตา เจว,}{สตฺต อาปตฺติกฺขนฺธา จ.}}\csromangathastanza{\csromanfolio{217}\csromangathabat{ภิชฺชติ อุปสมฺปทา,}{ตเถว อปเร ทุเว.} \\ \csromangathabat{น วตฺถพฺพํ น ทาตพฺพํ,}{อภพฺพาภพฺพเมว จ.}}\csromangathastanza{\csromangathabat{สญฺจิจฺจ สาติสารา จ,}{ปฏิกฺโกสา นิสฺสารณา.} \\ \csromangathabat{โอสารณา ปฏิญฺญา จ,}{ปฏิคฺคหา ปฏิกฺขิปา.}}\csromangathastanza{\csromangathabat{อุปฆาติ โจทนา จ,}{กถินา จ ทุเว ตถา.} \\ \csromangathabat{จีวรา ปตฺตมณฺฑลา,}{อธิฏฺฐานา ตเถว เทฺว.}}\csromangathastanza{\csromangathabat{วิกปฺปนา จ วินยา,}{เวนยิกา จ สลฺเลขา.} \\ \csromangathabat{อาปชฺชติ จ วุฏฺฐาติ,}{ปริวาสาปเร ทุเว.}}\csromangathastanza{\csromangathabat{เทฺว มานตฺตา อปเร จ,}{รตฺติจฺเฉโท อนาทริ.} \\ \csromangathabat{เทฺว โลณา ตโย อปเร,}{ปริโภคกฺโกเสน จ.} \\ \csromangathabat{เปสุญฺโญ จ คณาวสฺส,}{ฐปนา ภารกปฺปิยํ.}}}

\vspace{\tipitakaparskip}
\csromancenter{อนาปตฺติ อธมฺมธมฺมา, วินเย อาสเว ตถาติ.}
\csromansectionrule
\csromanmatikaanchor{mtk.166}
\csromantocmarkref{paragraph}{3. ติกวาร}{mtk.166}
\bookmark[dest={mtk.166},level=4]{3. ติกวาร}
\csromanheader{4}{3. \textbf{ติกวาร}}
\proseitem{323}{อตฺถาปตฺติ ติฏฺฐนฺเต ภควติ อาปชฺชติ, โน ปรินิพฺพุเต, อตฺถาปตฺติ ปรินิพฺพุเต ภควติ อาปชฺชติ, โน ติฏฺฐนฺเต, อตฺถาปตฺติ ติฏฺฐนฺเตปิ ภควติ อาปชฺชติ ปรินิพฺพุเตปิ. อตฺถาปตฺติ กาเล อาปชฺชติ, โน วิกาเล, อตฺถาปตฺติ วิกาเล อาปชฺชติ, โน กาเล, อตฺถาปตฺติ กาเล เจว อาปชฺชติ วิกาเล จ. อตฺถาปตฺติ รตฺติํ อาปชฺชติ, โน ทิวา, อตฺถาปตฺติ ทิวา อาปชฺชติ, โน รตฺติํ, อตฺถาปตฺติ รตฺติํ เจว อาปชฺชติ ทิวา จ. อตฺถาปตฺติ ทสวสฺโส อาปชฺชติ, โน อูนทสวสฺโส, อตฺถาปตฺติ อูนทสวสฺโส อาปชฺชติ, โน ทสวสฺโส, อตฺถาปตฺติ ทสวสฺโส เจว อาปชฺชติ อูนทสวสฺโส จ. อตฺถาปตฺติ ปญฺจวสฺโส อาปชฺชติ, โน อูนปญฺจวสฺโส, อตฺถาปตฺติ อูนปญฺจวสฺโส อาปชฺชติ, โน ปญฺจวสฺโส, อตฺถาปตฺติ ปญฺจวสฺโส เจว อาปชฺชติ อูนปญฺจวสฺโส จ. อตฺถาปตฺติ กุสลจิตฺโต อาปชฺชติ, อตฺถาปตฺติ อกุสลจิตฺโต อาปชฺชติ, อตฺถาปตฺติ อพฺยากตจิตฺโต อาปชฺชติ. อตฺถาปตฺติ สุขเวทนา\-สมงฺคี อาปชฺชติ, อตฺถาปตฺติ ทุกฺขเวทนา\-สมงฺคี อาปชฺชติ, อตฺถาปตฺติ อทุกฺขมสุข\-เวทนาสมงฺคี อาปชฺชติ. ตีณิ \csromanfolio{218}โจทนาวตฺถูนิ ทิฏฺเฐน สุเตน ปริสงฺกาย. ตโย สลากคฺคาหา คุฬฺหโก วิวฏโก\csromansharedfootnote{9bcc894a4b1230e3}{วิวฏฺฏโก (ก)} สกณฺณชปฺปโก. ตโย ปฏิกฺเขปา มหิจฺฉตา อสนฺตุฏฺฐิตา\csromansharedfootnote{26a85f5958fac9e7}{อสนฺตุฏฺฐตา (สฺยา)} อสลฺเลขตา. ตโย อนุญฺญาตา อปฺปิจฺฉตา สนฺตุฏฺฐิตา สลฺเลขตา. อปเรปิ ตโย ปฏิกฺเขปา มหิจฺฉตา อสนฺตุฏฺฐิตา อมตฺตญฺญุตา. ตโย อนุญฺญาตา อปฺปิจฺฉตา สนฺตุฏฺฐิตา มตฺตญฺญุตา. ติสฺโส ปญฺญตฺติโย ปญฺญตฺติ อนุปญฺญตฺติ อนุปฺปนฺนปญฺญตฺติ. อปราปิ ติสฺโส ปญฺญตฺติโย สพฺพตฺถ\-ปญฺญตฺติ ปเทสปญฺญตฺติ สาธารณ\-ปญฺญตฺติ. อปราปิ ติสฺโส ปญฺญตฺติโย อสาธารณ\-ปญฺญตฺติ เอกโตปญฺญตฺติ อุภโตปญฺญตฺติ.}

\prose{อตฺถาปตฺติ พาโล อาปชฺชติ, โน ปณฺฑิโต, อตฺถาปตฺติ ปณฺฑิโต อาปชฺชติ, โน พาโล, อตฺถาปตฺติ พาโล เจว อาปชฺชติ ปณฺฑิโต จ. อตฺถาปตฺติ กาเฬ อาปชฺชติ, โน ชุเณฺห, อตฺถาปตฺติ ชุเณฺห อาปชฺชติ, โน กาเฬ, อตฺถาปตฺติ กาเฬ เจว อาปชฺชติ ชุเณฺห จ. อตฺถิ กาเฬ กปฺปติ, โน ชุเณฺห, อตฺถิ ชุเณฺห กปฺปติ, โน กาเฬ, อตฺถิ กาเฬ เจว กปฺปติ ชุเณฺห จ. อตฺถาปตฺติ เหมนฺเต อาปชฺชติ, โน คิเมฺห, โน วสฺเส, อตฺถาปตฺติ คิเมฺห อาปชฺชติ, โน เหมนฺเต, โน วสฺเส, อตฺถาปตฺติ วสฺเส อาปชฺชติ, โน เหมนฺเต, โน คิเมฺห. อตฺถาปตฺติ สํโฆ อาปชฺชติ, น คโณ, น ปุคฺคโล, อตฺถาปตฺติ คโณ, อาปชฺชติ, น สํโฆ, น ปุคฺคโล, อตฺถาปตฺติ ปุคฺคโล อาปชฺชติ, น สํโฆ, น คโณ. อตฺถิ สํฆสฺส กปฺปติ, น คณสฺส, น ปุคฺคลสฺส, อตฺถิ คณสฺส กปฺปติ, น สํฆสฺส, น ปุคฺคลสฺส, อตฺถิ ปุคฺคลสฺส กปฺปติ, น สํฆสฺส, น คณสฺส. ติสฺโส ฉาทนา วตฺถุํ ฉาเทติ, โน อาปตฺติํ, อาปตฺติํ ฉาเทติ, โน วตฺถุํ, วตฺถุํ เจว ฉาเทติ อาปตฺติญฺจ. ติสฺโส ปฏิจฺฉาทิโย ชนฺตาฆร\-ปฏิจฺฉาทิ อุทก\-ปฏิจฺฉาทิ วตฺถ\-ปฏิจฺฉาทิ.\csromansymbolfootnote{*}{อํ 1. 286 ปิฏฺเฐปิ.}ตีณิ ปฏิจฺฉนฺนานิ วหนฺติ, โน วิวฏานิ, มาตุคาโม ปฏิจฺฉนฺโน วหติ, โน วิวโฏ, พฺราหฺมณานํ มนฺตา ปฏิจฺฉนฺนา วหนฺติ, โน วิวฏา, มิจฺฉาทิฏฺฐิ ปฏิจฺฉนฺนา วหติ, โน วิวฏา. ตีณิ วิวฏานิ วิโรจนฺติ, โน ปฏิจฺฉนฺนานิ, จนฺทมณฺฑลํ วิวฏํ วิโรจติ, โน ปฏิจฺฉนฺนํ, สูริยมณฺฑลํ วิวฏํ วิโรจติ, โน ปฏิจฺฉนฺนํ, ตถาคต\-ปฺปเวทิโต ธมฺมวินโย วิวโฏ วิโรจติ, โน ปฏิจฺฉนฺโน. ตโย เสนาสนคฺคาหา \csromanfolio{219}ปุริมโก ปจฺฉิมโก อนฺตรามุตฺตโก. อตฺถาปตฺติ คิลาโน อาปชฺชติ, โน อคิลาโน, อตฺถาปตฺติ อคิลาโน อาปชฺชติ, โน คิลาโน, อตฺถาปตฺติ คิลาโน เจว อาปชฺชติ อคิลาโน จ.}

\prose{ตีณิ อธมฺมิกานิ ปาติโมกฺข\-ฏฺฐปนานิ. ตีณิ ธมฺมิกานิ ปาติโมกฺข\-ฏฺฐปนานิ. ตโย ปริวาสา ปฏิจฺฉนฺน\-ปริวาโส อปฺปฏิจฺฉนฺน\-ปริวาโส สุทฺธนฺต\-ปริวาโส. ตโย มานตฺตา ปฏิจฺฉนฺน\-มานตฺตํ อปฺปฏิจฺฉนฺน\-มานตฺตํ ปกฺขมานตฺตํ. ตโย ปาริวาสิกสฺส ภิกฺขุโน รตฺติจฺเฉทา สหวาโส วิปฺปวาโส อนาโรจนา. อตฺถาปตฺติ อนฺโต อาปชฺชติ, โน พหิ, อตฺถาปตฺติ พหิ อาปชฺชติ, โน อนฺโต, อตฺถาปตฺติ อนฺโต เจว อาปชฺชติ พหิ จ. อตฺถาปตฺติ อนฺโตสีมาย อาปชฺชติ, โน พหิสีมาย, อตฺถาปตฺติ พหิสีมาย อาปชฺชติ, โน อนฺโตสีมาย, อตฺถาปตฺติ อนฺโตสีมาย เจว อาปชฺชติ พหิสีมาย จ. ตีหากาเรหิ อาปตฺติํ อาปชฺชติ, กาเยน อาปชฺชติ, วาจาย อาปชฺชติ, กาเยน วาจาย อาปชฺชติ. อปเรหิปิ ตีหากาเรหิ อาปตฺติํ อาปชฺชติ สํฆมชฺเฌ คณมชฺเฌ ปุคฺคลสฺส สนฺติเก. ตีหากาเรหิ อาปตฺติยา วุฏฺฐาติ, กาเยน วุฏฺฐาติ, วาจาย วุฏฺฐาติ, กาเยน วาจาย วุฏฺฐาติ. อปเรหิปิ ตีหากาเรหิ อาปตฺติยา วุฏฺฐาติ สํฆมชฺเฌ คณมชฺเฌ ปุคฺคลสฺส สนฺติเก. ตีณิ อธมฺมิกานิ อมูฬฺห\-วินยสฺส ทานานิ. ตีณิ ธมฺมิกานิ อมูฬฺห\-วินยสฺส ทานานิ.}

\prose{\csromansymbolfootnote{*}{วิ 4. 9 ปิฏฺเฐปิ.}ตีหงฺเคหิ สมนฺนาคตสฺส ภิกฺขุโน อากงฺขมาโน สํโฆ ตชฺชนีย\-กมฺมํ กเรยฺย, ภณฺฑนการโก โหติ กลหการโก วิวาทการโก ภสฺสการโก สํเฆ อธิกรณ\-การโก, พาโลโหติ อพฺยตฺโต อาปตฺติพหุโล อนปทาโน, คิหิสํสฏฺโฐ วิหรติ อนนุโลมิเกหิ คิหิสํสคฺเคหิ.\csromansymbolfootnote{+}{วิ 4. 18 ปิฏฺเฐปิ. ++ วิ 4. 31 ปิฏฺเฐปิ.}ตีหงฺเคหิ สมนฺนาคตสฺส ภิกฺขุโน อากงฺขมาโน สํโฆ นิยสฺสกมฺมํ กเรยฺย, ภณฺฑนการโก โหติ ฯเปฯ สํเฆ อธิกรณ\-การโก, พาโล โหติ อพฺยตฺโต อาปตฺติพหุโล อนปทาโน, คิหิสํสฏฺโฐ วิหรติ อนนุโลมิเกหิ คิหิสํสคฺเคหิ. ++ ตีหงฺเคหิ สมนฺนาคตสฺส ภิกฺขุโน อากงฺขมาโน สํโฆ ปพฺพาชนีย\-กมฺมํ กเรยฺย, ภณฺฑนการโก โหติ ฯเปฯ สํเฆ อธิกรณ\-การโก, พาโล โหติ อพฺยตฺโต อาปตฺติพหุโล อนปทาโน, กุลทูกโก โหติ ปาปสมาจาโร ปาปสมาจารา \csromanfolio{220}ทิสฺสนฺติ เจว สุยฺยนฺติ จ. ตีหงฺเคหิ สมนฺนาคตสฺส ภิกฺขุโน อากงฺขมาโน สํโฆ ปฏิสารณีย\-กมฺมํ กเรยฺย, ภณฺฑนการโก โหติ - ป- สํเฆ อธิกรณ\-การโก, พาโล โหติ อพฺยตฺโต อาปตฺติพหุโล อนปทาโน, คิหี อกฺโกสติ ปริภาสติ.\csromansymbolfootnote{*}{วิ 4. 51 ปิฏฺเฐปิ.}ตีหงฺเคหิ สมนฺนาคตสฺส ภิกฺขุโน อากงฺขมาโน สํโฆ อาปตฺติยา อทสฺสเน อุกฺเขปนีย\-กมฺมํ กเรยฺย, ภณฺฑนการโก โหติ ฯเปฯ สํเฆ อธิกรณ\-การโก, พาโล โหติ อพฺยตฺโต อาปตฺติพหุโล อนปทาโน, อาปตฺติํ อาปชฺชิตฺวา น อิจฺฉติ อาปตฺติํ ปสฺสิตุํ.\csromansymbolfootnote{+}{วิ 4. 62 ปิฏฺเฐปิ. ++ วิ 4. 73 ปิฏฺเฐปิ.}ตีหงฺเคหิ สมนฺนาคตสฺส ภิกฺขุโน อากงฺกมาโน สํโฆ อาปตฺติยา อปฺปฏิกมฺเม อุกฺเขปนีย\-กมฺมํ กเรยฺย, ภณฺฑนการโก โหติ ฯเปฯ สํเฆ อธิกรณ\-การโก, พาโล โหติ อพฺยตฺโต อาปตฺติพหุโล อนปทาโน, อาปตฺติํ อาปชฺชิตฺวา น อิจฺฉติ อาปตฺติํ ปฏิกาตุํ. ++ ตีหงฺเคหิ สมนฺนาคตสฺส ภิกฺขุโน อากงฺขมาโน สํโฆ ปาปิกาย ทิฏฺฐิยา อปฺปฏินิสฺสคฺเค อุกฺเขปนีย\-กมฺมํ กเรยฺย, ภณฺฑนการโก โหติ ฯเปฯ สํเฆ อธิกรณ\-การโก, พาโล โหติ อพฺยตฺโต อาปตฺติพหุโล อนปทาโน, น อิจฺฉติ ปาปิกํ ทิฏฺฐิํ ปฏินิสฺสชฺชิตุํ.}

\prose{ตีหงฺเคหิ สมนฺนาคตสฺส ภิกฺขุโน อากงฺขมาโน สํโฆ อาคาฬฺหาย เจเตยฺย, ภณฺฑนการโก โหติ กลหการโก วิวาทการโก ภสฺสการโก สํเฆ อธิกรณ\-การโก, พาโล โหติ อพฺยตฺโต อาปตฺติพหุโล อนปทาโน, คิหิสํสฏฺโฐ วิหรติ อนนุโลมิเกหิ คิหิสํสคฺเคหิ. ตีหงฺเคหิ สมนฺนาคตสฺส ภิกฺขุโน กมฺมํ กาตพฺพํ, อลชฺชี จ โหติ พาโล จ อปกตตฺโต จ. อปเรหิปิ ตีหงฺเคหิ สมนฺนาคตสฺส ภิกฺขุโน กมฺมํ กาตพฺพํ, อธิสีเล สีลวิปนฺโน โหติ, อชฺฌาจาเร อาจารวิปนฺโน โหติ, อติทิฏฺฐิยา ทิฏฺฐิวิปนฺโน โหติ. อปเรหิปิ ตีหงฺเคหิ สมนฺนาคตสฺส ภิกฺขุโน กมฺมํ กาตพฺพํ, กายิเกน ทเวน สมนฺนาคโต โหติ, วาจสิเกน ทเวน สมนฺนาคโต โหติ, กายิก\-วาจสิเกน ทเวน สมนฺนาคโต โหติ. อปเรหิปิ ตีหงฺเคหิ สมนฺนาคตสฺส ภิกฺขุโน กมฺมํ กาตพฺพํ, กายิเกน อนาจาเรน สมนฺนาคโต โหติ, วาจสิเกน อนาจาเรน สมนฺนาคโต โหติ, กายิก\-วาจสิเกน อนาจาเรน สมนฺนาคโต โหติ. อปเรหิปิ ตีหงฺเคหิ สมนฺนาคตสฺส \csromanfolio{221}ภิกฺขุโน กมฺมํ กาตพฺพํ, กายิเกน อุปฆาติเกน สมนฺนาคโต โหติ, วาจสิเกน อุปฆาติเกน สมนฺนาคโต โหติ, กายิก\-วาจสิเกน อุปฆาติเกน สมนฺนาคโต โหติ. อปเรหิปิ ตีหงฺเคหิ สมนฺนาคตสฺส ภิกฺขุโน กมฺมํ กาตพฺพํ, กายิเกน มิจฺฉาชีเวน สมนฺนาคโต โหติ, วาจสิเกน มิจฺฉาชีเวน สมนฺนาคโต โหติ, กายิก\-วาจสิเกน มิจฺฉาชีเวน สมนฺนาคโต โหติ. อปเรหิปิ ตีหงฺเคหิ สมนฺนาคตสฺส ภิกฺขุโน กมฺมํ กาตพฺพํ, อาปตฺติํ อาปนฺโน กมฺมกโต อุปสมฺปาเทติ, นิสฺสยํ เทติ, สามเณรํ อุปฏฺฐาเปติ. อปเรหิปิ ตีหงฺเคหิ สมนฺนาคตสฺส ภิกฺขุโน กมฺมํ กาตพฺพํ, ยาย อาปตฺติยา สํเฆน กมฺมํ กตํ โหติ, ตํ อาปตฺติํ อาปชฺชติ, อญฺญํ วา ตาทิสิกํ ตโต วา ปาปิฏฺฐตรํ. อปเรหิปิ ตีหงฺเคหิ สมนฺนาคตสฺส ภิกฺขุโน กมฺมํ กาตพฺพํ, พุทฺธสฺส อวณฺณํ ภาสติ, ธมฺมสฺส อวณฺณํ ภาสติ, สํฆสฺส อวณฺณํ ภาสติ.}

\prose{ตีหงฺเคหิ สมนฺนาคตสฺส ภิกฺขุโน สํฆมชฺเฌ อุโปสถํ ฐเปนฺตสฺส “อลํ ภิกฺขุ, มา ภณฺฑนํ มา กลหํ มา วิคฺคหํ มา วิวาทนฺ”ติ โอมทฺทิตฺวา สํเฆน อุโปสโถ กาตพฺโพ, อลชฺชี จ โหติ พาโล จ อปกตตฺโต จ. ตีหงฺเคหิ สมนฺนาคตสฺส ภิกฺขุโน สํฆมชฺเฌ ปวารณํ ฐเปนฺตสฺส “อลํ ภิกฺขุ, มา ภณฺฑนํ มา กลหํ มา วิคฺคหํ มา วิวาทนฺ”ติ โอมทฺทิตฺวา สํเฆน ปวาเรตพฺพํ, อลชฺชี จ โหติ พาโล จ อปกตตฺโต จ. ตีหงฺเคหิ สมนฺนาคตสฺส ภิกฺขุโน น กาจิ สํฆสมฺมุติ ทาตพฺพา, อลชฺชี จ โหติ พาโล จ อปกตตฺโต จ. ตีหงฺเคหิ สมนฺนาคเตน ภิกฺขุนา สํเฆน โวหริตพฺพํ\csromansharedfootnote{977bc1480d91aff5}{สํโฆ น โวหริตพฺโพ (สฺยา)}, อลชฺชี จ โหติ พาโล จ อปกตตฺโต จ. ตีหงฺเคหิ สมนฺนาคโต ภิกฺขุ น กิสฺมิญฺจิ ปจฺเจกฏฺฐาเน ฐเปตพฺโพ, อลชฺชี จ โหติ พาโล จ อปกตตฺโต จ. ตีหงฺเคหิ สมนฺนาคตสฺส ภิกฺขุโน นิสฺสาย น วตฺถพฺพํ, อลชฺชี จ โหติ พาโล จ อปกตตฺโต จ. ตีหงฺเคหิ สมนฺนาคตสฺส ภิกฺขุโน นิสฺสโย น ทาตพฺโพ, อลชฺชี จ โหติ พาโล จ อปกตตฺโต จ. ตีหงฺเคหิ สมนฺนาคตสฺส ภิกฺขุโน โอกาสกมฺมํ การาเปนฺตสฺส นาลํ โอกาสกมฺมํ กาตุํ, อลชฺชี จ โหติ พาโล จ อปกตตฺโต จ. ตีหงฺเคหิ สมนฺนาคตสฺส ภิกฺขุโน สวจนียํ นาทาตพฺพํ, อลชฺชี จ โหติ \csromanfolio{222}พาโล จ อปกตตฺโต จ. ตีหงฺเคหิ สมนฺนาคตสฺส ภิกฺขุโน วินโย น ปุจฺฉิตพฺโพ, อลชฺชี จ โหติ พาโล จ อปกตตฺโต จ. ตีหงฺเคหิ สมนฺนาคเตน ภิกฺขุนา วินโย น ปุจฺฉิตพฺโพ, อลชฺชี จ โหติ พาโล จ อปกตตฺโต จ. ตีหงฺเคหิ สมนฺนาคตสฺส ภิกฺขุโน วินโย น วิสฺสชฺเชตพฺโพ, อลชฺชี จ โหติ พาโล จ อปกตตฺโต จ. ตีหงฺเคหิ สมนฺนาคเตน ภิกฺขุนา วินโย น วิสฺสชฺเชตพฺโพ, อลชฺชี จ โหติ พาโล จ อปกตตฺโต จ. ตีหงฺเคหิ สมนฺนาคตสฺส ภิกฺขุโน อนุโยโค น ทาตพฺโพ, อลชฺชี จ โหติ พาโล จ อปกตตฺโต จ. ตีหงฺเคหิ สมนฺนาคเตน ภิกฺขุนา สทฺธิํ วินโย น สากจฺฉิตพฺโพ\csromansharedfootnote{6ab36080b2a58bf4}{สากจฺฉาตพฺโพ (ก)}, อลชฺชี จ โหติ พาโล จ อปกตตฺโต จ. ตีหงฺเคหิ สมนฺนาคเตน ภิกฺขุนา น อุปสมฺ\-ปาเท\-ตพฺพํ, น นิสฺสโย ทาตพฺโพ, น สามเณโร อุปฏฺฐาเปตพฺโพ, อลชฺชี จ โหติ พาโล จ อปกตตฺโต จ.}

\prose{ตโย อุโปสถา จาตุทฺทสิโก ปนฺนรสิโก สามคฺคิ\-อุโปสโถ. อปเรปิ ตโย อุโปสถา สํเฆอุโปสโถ คเณอุโปสโถ ปุคฺคเล\-อุโปสโถ. อปเรปิ ตโย อุโปสถา สุตฺตุทฺเทโส อุโปสโถ ปาริสุทฺธิ\-อุโปสโถ อธิฏฺฐานุ\-โปสโถ. ติสฺโส ปวารณา จาตุทฺทสิกา ปนฺนรสิกา สามคฺคิ\-ปวารณา. อปราปิ ติสฺโส ปวารณา สํเฆปวารณา คเณปวารณา ปุคฺคเล\-ปวารณา. อปราปิ ติสฺโส ปวารณา เตวาจิกา\-ปวารณา เทฺว วาจิกาปวารณา สมานวสฺสิกา\-ปวารณา.}

\prosewithcloser{ตโย อาปายิกา เนรยิกา อิทมปฺปหาย, โย จ อพฺรหฺมจารี พฺรหฺมจาริ\-ปฏิญฺโญ, โย จ สุทฺธํ พฺรหฺมจาริํ ปริสุทฺธ\-พฺรหฺมจริยํ จรนฺตํ อมูลเกน อพฺรหฺมจริเยน อนุทฺธํเสติ, โย จายํ เอวํวาที เอวํทิฏฺฐิ\csromansharedfootnote{a66455bc93388679}{เอวํทิฏฺฐี (สี)} “นตฺถิ กาเมสุ โทโส”ติ, โส กาเมสุ ปาตพฺยตํ อาปชฺชติ.\csromansymbolfootnote{*}{อํ 1. 202 ปิฏฺเฐปิ.}ตีณิ อกุสลมูลานิ โลโภ อกุสลมูลํ โทโส อกุสลมูลํ โมโห อกุสลมูลํ.\csromansymbolfootnote{+}{อํ 1. 203 ปิฏฺเฐปิ.}ตีณิ กุสลมูลานิ อโลโภ กุสลมูลํ อโทโส กุสลมูลํ อโมโห กุสลมูลํ. ตีณิ ทุจฺจริตานิ กายทุจฺจริตํ วจีทุจฺจริตํ มโนทุจฺจริตํ. ตีณิ สุจริตานิ กายสุจริตํ วจีสุจริตํ มโนสุจริตํ. ตโย อตฺถวเส ปฏิจฺจ ภควตา กุเลสุ ติกโภชนํ ปญฺญตฺตํ, ทุมฺมงฺกูนํ ปุคฺคลานํ นิคฺคหาย, เปสลานํ \csromanfolio{223}ภิกฺขูนํ ผาสุวิหาราย “มา ปาปิจฺฉา ปกฺขํ นิสฺสาย สํฆํ ภินฺเทยฺยูนฺ”ติ, กุลา\-นุทฺทยตาย จ.\csromansymbolfootnote{*}{วิ 4. 366 ปิฏฺเฐปิ.}ตีหิ อสทฺธมฺเมหิ อภิภูโต ปริยาทินฺน\-จิตฺโต เทวทตฺโต อาปายิโก เนรยิโก กปฺปฏฺโฐ อเตกิจฺโฉ, ปาปิจฺฉตา ปาปมิตฺตตา โอรมตฺตเกน วิเสสา\-ธิคเมน อนฺตรา โวสานํ อาปาทิ. ติสฺโส สมฺมุติโย ทณฺฑสมฺมุติ สิกฺกาสมฺมุติ ทณฺฑ\-สิกฺกา\-สมฺมุติ. ติสฺโส ปาทุกา ธุวฏฺฐานิกา อสงฺกมนียา, วจฺจปาทุกา ปสฺสาวปาทุกา อาจมนปาทุกา. ติสฺโส ปาทฆํสนิโย, สกฺขรํ กถลา สมุทฺท\-เผณโกติ.}{ติกํ นิฏฺฐิตํ.}
\csromancenter{ตสฺสุทฺทานํ.}
\setlength{\csromangathapagewidth}{0pt}
\setlength{\csromangathawakwidth}{0pt}
\csromangathameasuregroup{ติฏฺฐนฺเต กาเล รตฺติํ จ, \\ เวทนา โจทนา วตฺถุ, \\ ปญฺญตฺติ อปเร เทฺว จ, \\ เหมนฺเต สํโฆ สํฆสฺส, \\ ปฏิจฺฉนฺนา วิวฏา จ, \\ ปาติโมกฺขํ ปริวาสํ, \\ อนฺโต อนฺโต จ สีมาย, \\ วุฏฺฐาติ อปเร เจว, \\ ตชฺชนียา นิยสฺสา จ, \\ อทสฺสนา ปฏิกมฺเม, \\ อาคาฬฺหกมฺมาธิสีเล, \\ อาชีวาปนฺนา ตาทิสิกา, \\ ปวารณา สมฺมุติ จ, \\ น วตฺถพฺพํ น ทาตพฺพํ, \\ น กเร สวจนียํ, \\ น วิสฺสชฺเช ทุเว เจว, \\ สากจฺฉา อุปสมฺปทา, \\ อุโปสถติกา ตีณิ, \\ อาปายิกา อกุสลา, \\ ติกโภชนสทฺธมฺเม,}{\csromangathabat{ติฏฺฐนฺเต กาเล รตฺติํ จ,}{ทส ปญฺจ กุสเลน.} \\ \csromangathabat{เวทนา โจทนา วตฺถุ,}{สลากา เทฺว ปฏิกฺขิปา.} \\ \csromangathabat{ปญฺญตฺติ อปเร เทฺว จ,}{พาโล กาเฬ จ กปฺปติ.} \\ \csromangathabat{เหมนฺเต สํโฆ สํฆสฺส,}{ฉาทนา จ ปฏิจฺฉาทิ.} \\ \csromangathabat{ปฏิจฺฉนฺนา วิวฏา จ,}{เสนาสนคิลายนา.} \\ \csromangathabat{ปาติโมกฺขํ ปริวาสํ,}{มานตฺตา ปาริวาสิกา.} \\ \csromangathabat{อนฺโต อนฺโต จ สีมาย,}{อาปชฺชติ ปุนาปเร.} \\ \csromangathabat{วุฏฺฐาติ อปเร เจว,}{อมูฬฺหวินยา ทุเว.} \\ \csromangathabat{ตชฺชนียา นิยสฺสา จ,}{ปพฺพาชปฏิสารณี.} \\ \csromangathabat{อทสฺสนา ปฏิกมฺเม,}{อนิสฺสคฺเค จ ทิฏฺฐิยา.} \\ \csromangathabat{อาคาฬฺหกมฺมาธิสีเล,}{ทวานาจารฆาติกา.} \\ \csromangathabat{อาชีวาปนฺนา ตาทิสิกา,}{อวณฺณุโปสเถน จ.} \\ \csromangathabat{ปวารณา สมฺมุติ จ,}{โวหารปจฺเจเกน จ.} \\ \csromangathabat{น วตฺถพฺพํ น ทาตพฺพํ,}{โอกาสํ น กเร ตถา.} \\ \csromangathabat{น กเร สวจนียํ,}{น ปุจฺฉิตพฺพกา ทุเว.} \\ \csromangathabat{น วิสฺสชฺเช ทุเว เจว,}{อนุโยคมฺปิ โน ทเท.} \\ \csromangathabat{สากจฺฉา อุปสมฺปทา,}{นิสฺสยสามเณรา จ.} \\ \csromangathabat{อุโปสถติกา ตีณิ,}{ปวารณติกา ตโย.} \\ \csromangathabat{อาปายิกา อกุสลา,}{กุสลา จริตา ทุเว.} \\ \csromangathabat{ติกโภชนสทฺธมฺเม,}{สมฺมุติ ปาทุเกน จ.}}
\csromangathasetleft{ติฏฺฐนฺเต กาเล รตฺติํ จ, \\ เวทนา โจทนา วตฺถุ, \\ ปญฺญตฺติ อปเร เทฺว จ, \\ เหมนฺเต สํโฆ สํฆสฺส, \\ ปฏิจฺฉนฺนา วิวฏา จ, \\ ปาติโมกฺขํ ปริวาสํ, \\ อนฺโต อนฺโต จ สีมาย, \\ วุฏฺฐาติ อปเร เจว, \\ ตชฺชนียา นิยสฺสา จ, \\ อทสฺสนา ปฏิกมฺเม, \\ อาคาฬฺหกมฺมาธิสีเล, \\ อาชีวาปนฺนา ตาทิสิกา, \\ ปวารณา สมฺมุติ จ, \\ น วตฺถพฺพํ น ทาตพฺพํ, \\ น กเร สวจนียํ, \\ น วิสฺสชฺเช ทุเว เจว, \\ สากจฺฉา อุปสมฺปทา, \\ อุโปสถติกา ตีณิ, \\ อาปายิกา อกุสลา, \\ ติกโภชนสทฺธมฺเม,}
\csromangathagroup{\csromangathastanza{\csromangathabat{ติฏฺฐนฺเต กาเล รตฺติํ จ,}{ทส ปญฺจ กุสเลน.} \\ \csromangathabat{เวทนา โจทนา วตฺถุ,}{สลากา เทฺว ปฏิกฺขิปา.}}\csromangathastanza{\csromangathabat{ปญฺญตฺติ อปเร เทฺว จ,}{พาโล กาเฬ จ กปฺปติ.} \\ \csromangathabat{เหมนฺเต สํโฆ สํฆสฺส,}{ฉาทนา จ ปฏิจฺฉาทิ.}}\csromangathastanza{\csromangathabat{ปฏิจฺฉนฺนา วิวฏา จ,}{เสนาสนคิลายนา.} \\ \csromangathabat{ปาติโมกฺขํ ปริวาสํ,}{มานตฺตา ปาริวาสิกา.}}\csromangathastanza{\csromangathabat{อนฺโต อนฺโต จ สีมาย,}{อาปชฺชติ ปุนาปเร.} \\ \csromangathabat{วุฏฺฐาติ อปเร เจว,}{อมูฬฺหวินยา ทุเว.}}\csromangathastanza{\csromangathabat{ตชฺชนียา นิยสฺสา จ,}{ปพฺพาชปฏิสารณี.} \\ \csromangathabat{อทสฺสนา ปฏิกมฺเม,}{อนิสฺสคฺเค จ ทิฏฺฐิยา.}}\csromangathastanza{\csromangathabat{อาคาฬฺหกมฺมาธิสีเล,}{ทวานาจารฆาติกา.} \\ \csromangathabat{อาชีวาปนฺนา ตาทิสิกา,}{อวณฺณุโปสเถน จ.}}\csromangathastanza{\csromangathabat{ปวารณา สมฺมุติ จ,}{โวหารปจฺเจเกน จ.} \\ \csromangathabat{น วตฺถพฺพํ น ทาตพฺพํ,}{โอกาสํ น กเร ตถา.}}\csromangathastanza{\csromangathabat{น กเร สวจนียํ,}{น ปุจฺฉิตพฺพกา ทุเว.} \\ \csromangathabat{น วิสฺสชฺเช ทุเว เจว,}{อนุโยคมฺปิ โน ทเท.}}\csromangathastanza{\csromanfolio{224}\csromangathabat{สากจฺฉา อุปสมฺปทา,}{นิสฺสยสามเณรา จ.} \\ \csromangathabat{อุโปสถติกา ตีณิ,}{ปวารณติกา ตโย.}}\csromangathastanza{\csromangathabat{อาปายิกา อกุสลา,}{กุสลา จริตา ทุเว.} \\ \csromangathabat{ติกโภชนสทฺธมฺเม,}{สมฺมุติ ปาทุเกน จ.}}}

\vspace{\tipitakaparskip}
\csromancenter{ปาทฆํสนิกา เจว, อุทฺทานํ ติกเก อิทนฺติ.}
\csromansectionrule
\csromanmatikaanchor{mtk.167}
\csromantocmarkref{paragraph}{4. จตุกฺกวาร}{mtk.167}
\bookmark[dest={mtk.167},level=4]{4. จตุกฺกวาร}
\csromanheader{4}{4. \textbf{จตุกฺกวาร}}
\proseitem{324}{อตฺถาปตฺติ สกวาจาย อาปชฺชติ, ปรวาจาย วุฏฺฐาติ, อตฺถาปตฺติ ปรวาจาย อาปชฺชติ, สกวาจาย วุฏฺฐาติ, อตฺถาปตฺติ สกวาจาย อาปชฺชติ, สกวาจาย วุฏฺฐาติ, อตฺถาปตฺติ ปรวาจาย อาปชฺชติ, ปรวาจาย วุฏฺฐาติ. อตฺถาปตฺติ กาเยน อาปชฺชติ, วาจาย วุฏฺฐาติ, อตฺถาปตฺติ วาจาย อาปชฺชติ, กาเยน วุฏฺฐาติ, อตฺถาปตฺติ กาเยน อาปชฺชติ, กาเยน วุฏฺฐาติ, อตฺถาปตฺติ วาจาย อาปชฺชติ, วาจาย วุฏฺฐาติ. อตฺถาปตฺติ ปสุตฺโต อาปชฺชติ, ปฏิพุทฺโธ วุฏฺฐาติ, อตฺถาปตฺติ ปฏิพุทฺโธ อาปชฺชติ, ปสุตฺโต วุฏฺฐาติ, อตฺถาปตฺติ ปสุตฺโต อาปชฺชติ, ปสุตฺโต วุฏฺฐาติ, อตฺถาปตฺติ ปฏิพุทฺโธ อาปชฺชติ, ปฏิพุทฺโธ วุฏฺฐาติ, อตฺถาปตฺติ อจิตฺตโก อาปชฺชติ, สจิตฺตโก วุฏฺฐาติ, อตฺถาปตฺติ สจิตฺตโก อาปชฺชติ, อจิตฺตโก วุฏฺฐาติ, อตฺถาปตฺติ อจิตฺตโก อาปชฺชติ, อจิตฺตโก วุฏฺฐาติ, อตฺถาปตฺติ สจิตฺตโก อาปชฺชติ, สจิตฺตโก วุฏฺฐาติ. อตฺถาปตฺติ อาปชฺชนฺโต เทเสติ, เทเสนฺโต อาปชฺชติ. อตฺถาปตฺติ อาปชฺชนฺโต วุฏฺฐาติ, วุฏฺฐหนฺโต อาปชฺชติ. อตฺถาปตฺติ กมฺเมน อาปชฺชติ, อกมฺเมน วุฏฺฐาติ, อตฺถาปตฺติ อกมฺเมน อาปชฺชติ, กมฺเมน วุฏฺฐาติ, อตฺถาปตฺติ กมฺเมน อาปชฺชติ, กมฺเมน วุฏฺฐาติ, อตฺถาปตฺติ อกมฺเมน อาปชฺชติ, อกมฺเมน วุฏฺฐาติ.}

\prose{\csromansymbolfootnote{*}{อํ 1. 568; ที 3. 193; อภิ 2. 390 ปิฏฺเฐสุปิ.}จตฺตาโร อนริยโวหารา, อทิฏฺเฐ ทิฏฺฐวาทิตา, อสฺสุเต สุตวาทิตา, อมุเต มุตวาทิตา, อวิญฺญาเต วิญฺญาตวาทิตา. จตฺตาโร อริยโวหารา, อทิฏฺเฐ อทิฏฺฐวาทิตา, อสฺสุเต อสฺสุตวาทิตา, อมุเต อมุตวาทิตา, อวิญฺญาเต อวิญฺญาตวาทิตา. อปเรปิ จตฺตาโร อนริยโวหารา, ทิฏฺเฐ อทิฏฺฐวาทิตา, สุเต อสฺสุตวาทิตา, มุเต อมุตวาทิตา, วิญฺญาเต อวิญฺญาตวาทิตา. \csromanfolio{225}จตฺตาโร อริยโวหารา, ทิฏฺเฐ ทิฏฺฐวาทิตา, สุเต สุตวาทิตา, มุเต มุตวาทิตา, วิญฺญาเต วิญฺญาตวาทิตา.}

\prose{จตฺตาโร ปาราชิกา ภิกฺขูนํ ภิกฺขุนีหิ สาธารณา. จตฺตาโร ปาราชิกา ภิกฺขุนีนํ ภิกฺขูหิ อสาธารณา. จตฺตาโร ปริกฺขารา อตฺถิ ปริกฺขาโร รกฺขิตพฺโพ โคเปตพฺโพ มมายิตพฺโพ ปริภุญฺชิ\-ตพฺโพ, อตฺถิ ปริกฺขาโร รกฺขิตพฺโพ โคเปตพฺโพ มมายิตพฺโพ น ปริภุญฺชิ\-ตพฺโพ, อตฺถิ ปริกฺขาโร รกฺขิตพฺโพ โคเปตพฺโพ น มมายิตพฺโพ น ปริภุญฺชิ\-ตพฺโพ, อตฺถิ ปริกฺขาโร น รกฺขิตพฺโพ น โคเปตพฺโพ น มมายิตพฺโพ น ปริภุญฺชิ\-ตพฺโพ. อตฺถาปตฺติ สมฺมุขา อาปชฺชติ, ปรมฺมุขา วุฏฺฐาติ, อตฺถาปตฺติ ปรมฺมุขา อาปชฺชติ, สมฺมุขา วุฏฺฐาติ, อตฺถาปตฺติ สมฺมุขา อาปชฺชติ, สมฺมุขา วุฏฺฐาติ, อตฺถาปตฺติ ปรมฺมุขา อาปชฺชติ, ปรมฺมุขา วุฏฺฐาติ. อตฺถาปตฺติ อชาน นฺโต อาปชฺชติ, ชานนฺโต วุฏฺฐาติ, อตฺถาปตฺติ ชานนฺโต อาปชฺชติ, อชานนฺโต วุฏฺฐาติ, อตฺถาปตฺติ อชานนฺโต อาปชฺชติ, อชานนฺโต วุฏฺฐาติ, อตฺถาปตฺติ ชานนฺโต อาปชฺชติ, ชานนฺโต วุฏฺฐาติ.}

\prose{จตูหากาเรหิ อาปตฺติํ อาปชฺชติ, กาเยน อาปชฺชติ, วาจาย อาปชฺชติ, กาเยน วาจาย อาปชฺชติ, กมฺมวาจาย อาปชฺชติ. อปเรหิปิ จตูหากาเรหิ อาปตฺติํ อาปชฺชติ สํฆมชฺเฌ คณมชฺเฌ ปุคฺคลสฺส สนฺติเก ลิงฺค\-ปาตุภาเวน. จตูหากาเรหิ อาปตฺติยา วุฏฺฐาติ, กาเยน วุฏฺฐาติ, วาจาย วุฏฺฐาติ, กาเยน วาจาย วุฏฺฐาติ, กมฺมวาจาย วุฏฺฐาติ. อปเรหิปิ จตูหากาเรหิ อาปตฺติยา วุฏฺฐาติ สํฆมชฺเฌ คณมชฺเฌ ปุคฺคลสฺส สนฺติเก ลิงฺค\-ปาตุภาเวน. สห ปฏิลาเภน ปุริมํ ชหติ, ปจฺฉิเม ปติฏฺฐาติ, วิญฺญตฺติโย ปฏิปฺปสฺสมฺภนฺติ, ปณฺณตฺติโย นิรุชฺฌนฺติ. สห ปฏิลาเภน ปจฺฉิมํ ชหติ, ปุริเม ปติฏฺฐาติ, วิญฺญตฺติโย ปฏิปฺปสฺสมฺภนฺติ, ปณฺณตฺติโย นิรุชฺฌนฺติ. จตสฺโส โจทนา, สีลวิปตฺติยา โจเทติ, อาจารวิปตฺติยา โจเทติ, ทิฏฺฐิ\-วิปตฺติยา โจเทติ, อาชีววิปตฺติยา โจเทติ. จตฺตาโร ปริวาสา ปฏิจฺฉนฺน\-ปริวาโส อปฺปฏิจฺฉนฺน\-ปริวาโส สุทฺธนฺต\-ปริวาโส สโมธาน\-ปริวาโส. จตฺตาโร มานตฺตา ปฏิจฺฉนฺน\-มานตฺตํ อปฺปฏิจฺฉนฺน\-มานตฺตํ ปกฺขมานตฺตํ สโมธาน\-มานตฺตํ. จตฺตาโร มานตฺต\-จาริกสฺส ภิกฺขุโน รตฺติจฺเฉทา สหวาโส วิปฺปวาโส อนาโรจนา อูเน คเณ จรติ. จตฺตาโร \csromanfolio{226}สามุกฺกํสา. จตฺตาโร ปฏิคฺคหิต\-ปริโภคา ยาวกาลิกํ ยามกาลิกํ สตฺตาหกาลิกํ ยาวชีวิกํ. จตฺตาริ มหาวิกฏานิ คูโถ มุตฺตํ ฉาริกา มตฺติกา. จตฺตาริ กมฺมานิ อปโลกน\-กมฺมํ ญตฺติกมฺมํ ญตฺติ\-ทุติย\-กมฺมํ ญตฺติ\-จตุตฺถ\-กมฺมํ. อปรานิปิ จตฺตาริ กมฺมานิ, อธมฺเมน วคฺคกมฺมํ, อธมฺเมน สมคฺคกมฺมํ, ธมฺเมน วคฺคกมฺมํ, ธมฺเมน สมคฺคกมฺมํ. จตสฺโส วิปตฺติโย สีลวิปตฺติ อาจารวิปตฺติ ทิฏฺฐิวิปตฺติ อาชีววิปตฺติ. จตฺตาริ อธิกรณานิ วิวาทา\-ธิกรณํ อนุวาทา\-ธิกรณํ อาปตฺตา\-ธิกรณํ กิจฺจา\-ธิกรณํ.\csromansymbolfootnote{*}{อํ 1. 547-8 ปิฏฺเฐปิ.}จตฺตาโร ปริสทูสนา, ภิกฺขุ ทุสฺสีโล ปาปธมฺโม ปริสทูสโน, ภิกฺกุนี ทุสฺสีลา ปาปธมฺมา ปริสทูสนา, อุปาสโก ทุสฺสีโล ปาปธมฺโม ปริสทูสโน, อุปาสิกา ทุสฺสีลา ปาปธมฺมา ปริสทูสนา. จตฺตาโร ปริสโสภนา\csromansharedfootnote{5e2e3627d6024020}{ปริสโสภณา (สฺยา, ก)}, ภิกฺขุ สีลวา กลฺยาณธมฺมา ปริสโสภโน, ภิกฺขุนี สีลวตี กลฺยาณธมฺมา ปริสโสภนา, อุปาสโก สีลวา กลฺยาณธมฺโม ปริสโสภโน, อุปาสิกา สีลวตี กลฺยาณธมฺมา ปริสโสภนา.}

\prose{อตฺถาปตฺติ อาคนฺตุโก อาปชฺชติ, โน อาวาสิโก, อตฺถาปตฺติ อาวาสิโก อาปชฺชติ, โน อาคนฺตุโก, อตฺถาปตฺติ อาคนฺตุโก เจว อาปชฺชติ อาวาสิโก จ, อตฺถาปตฺติ เนว อาคนฺตุโก อาปชฺชติ, โน อาวาสิโก. อตฺถาปตฺติ คมิโก อาปชฺชติ, โน อาวาสิโก, อตฺถาปตฺติ อาวาสิโก อาปชฺชติ, โน คมิโก, อตฺถาปตฺติ คมิโก เจว อาปชฺชติ อาวาสิโก จ, อตฺถาปตฺติ เนว คมิโก อาปชฺชติ, โน อาวาสิโก. อตฺถิ วตฺถุ\-นานตฺตตา โน อาปตฺติ\-นานตฺตตา, อตฺถิ อาปตฺติ\-นานตฺตตา โน วตฺถุ\-นานตฺตตา, อตฺถิ วตฺถุ\-นานตฺตตา เจว อาปตฺติ\-นานตฺตตา จ, อตฺถิ เนว วตฺถุ\-นานตฺตตา โน อาปตฺติ\-นานตฺตตา. อตฺถิ วตฺถุ\-สภาคตา โน อาปตฺติ\-สภาคตา, อตฺถิ อาปตฺติ\-สภาคตา โน วตฺถุ\-สภาคตา, อตฺถิ วตฺถุ\-สภาคตา เจว อาปตฺติ\-สภาคตา จ, อตฺถิ เนว วตฺถุ\-สภาคตา โน อาปตฺติ\-สภาคตา. อตฺถาปตฺติ อุปชฺฌาโย อาปชฺชติ, โน สทฺธิวิหาริโก, อตฺถาปตฺติ สทฺธิวิหาริโก อาปชฺชติ, โน อุปชฺฌาโย, อตฺถาปตฺติ อุปชฺฌาโย เจว อาปชฺชติ สทฺธิวิหาริโก จ, อตฺถาปตฺติ เนว อุปชฺฌาโย อาปชฺชติ, โน สทฺธิวิหาริโก. อตฺถาปตฺติ อาจริโย \csromanfolio{227}อาปชฺชติ โน อนฺเตวาสิโก, อตฺถาปตฺติ อนฺเตวาสิโก อาปชฺชติ, โน อาจริโย, อตฺถาปตฺติ อาจริโย เจว อาปชฺชติ อนฺเตวาสิโก จ, อตฺถาปตฺติ เนว อาจริโย อาปชฺชติ, โน อนฺเตวาสิโก. จตฺตาโร ปจฺจยา อนาปตฺติ วสฺสจฺเฉทสฺส, สํโฆ วา ภินฺโน โหติ, สํฆํ วา ภินฺทิตุกามา โหนฺติ, ชีวิตนฺตราโย วา โหติ, พฺรหฺมจริย\-นฺตราโย วา โหติ. จตฺตาริ วจีทุจฺจริตานิ มุสาวาโท ปิสุณา วาจา ผรุสา วาจา สมฺผปฺปลาโป. จตฺตาริ วจีสุจริตานิ สจฺจวาจา อปิสุณา วาจา สณฺหา วาจา มนฺตาภาสา. อตฺถิ อาทิยนฺโต ครุกํ อาปตฺติํ อาปชฺชติ, ปโยเชนฺโต ลหุกํ, อตฺถิ อาทิยนฺโต ลหุกํ อาปตฺติํ อาปชฺชติ, ปโยเชนฺโต ครุกํ, อตฺถิ อาทิยนฺโตปิ ปโยเชนฺโตปิ ครุกํ อาปตฺติํ อาปชฺชติ, อตฺถิ อาทิยนฺโตปิ ปโยเชนฺโตปิ ลหุกํ อาปตฺติํ อาปชฺชติ.}

\prose{อตฺถิ ปุคฺคโล อภิวาทนารโห โน ปจฺจุฏฺฐานา\-รโห, อตฺถิ ปุคฺคโล ปจฺจุฏฺฐานา\-รโห โน อภิวาทนารโห, อตฺถิ ปุคฺคโล อภิวาทนารโห เจว ปจฺจุฏฺฐานา\-รโห จ, อตฺถิ ปุคฺคโล เนว อภิวาทนารโห โน ปจฺจุฏฺฐานา\-รโห. อตฺถิ ปุคฺคโล อาสนารโห โน อภิวาทนารโห, อตฺถิ ปุคฺคโล อภิวาทนารโห โน อาสนารโห, อตฺถิ ปุคฺคโล อาสนารโห เจว อภิวาทนารโห จ, อตฺถิ ปุคฺคโล เนว อาสนารโห โน อภิวาทนารโห. อตฺถาปตฺติ กาเล อาปชฺชติ, โน วิกาเล, อตฺถาปตฺติ วิกาเล อาปชฺชติ โน กาเล, อตฺถาปตฺติ กาเล เจว อาปชฺชติ วิกาเล จ, อตฺถาปตฺติ เนว กาเล, อาปชฺชติ, โน วิกาเล. อตฺถิ ปฏิคฺคหิตํ กาเล กปฺปติ, โน วิกาเล, อตฺถิ ปฏิคฺคหิตํ วิกาเล กปฺปติ, โน กาเล, อตฺถิ ปฏิคฺคหิตํ กาเล เจว กปฺปติ วิกาเล จ, อตฺถิ ปฏิคฺคหิตํ เนว กาเล กปฺปติ, โน วิกาเล. อตฺถาปตฺติ ปจฺจนฺติเมสุ ชนปเทสุ อาปชฺชติ, โน มชฺฌิเมสุ, อตฺถาปตฺติ มชฺฌิเมสุ ชนปเทสุ อาปชฺชติ, โน ปจฺจนฺติเมสุ, อตฺถาปติ ปจฺจนฺติเมสุ เจว ชนปเทสุ อาปชฺชติ มชฺฌิเมสุ จ, อตฺถาปตฺติ เนว ปจฺจนฺติเมสุ ชนปเทสุ อาปชฺชติ, โน มชฺฌิเมสุ. อตฺถิ ปจฺจนฺติเมสุ ชนปเทสุ กปฺปติ, โน มชฺฌิเมสุ, อตฺถิ มชฺฌิเมสุ ชนปเทสุ กปฺปติ, โน ปจฺจนฺติเมสุ, อตฺถิ ปจฺจนฺติเมสุ เจว ชนปเทสุ กปฺปติ มชฺฌิเมสุ จ, อตฺถิ เนว ปจฺจนฺติเมสุ ชนปเทสุ กปฺปติ, \csromanfolio{228}โน มชฺฌิเมสุ. อตฺถาปตฺติ อนฺโต อาปชฺชติ โน พหิ, อตฺถาปตฺติ พหิ อาปชฺชติ, โน อนฺโต, อตฺถาปตฺติ อนฺโต เจว อาปชฺชติ พหิ จ, อตฺถาปตฺติ เนว อนฺโต อาปชฺชติ, โน พหิ. อตฺถาปตฺติ อนฺโตสีมาย อาปชฺชติ, โน พหิสีมาย, อตฺถาปตฺติ พหิสีมาย อาปชฺชติ, โน อนฺโตสีมาย, อตฺถาปตฺติ อนฺโตสีมาย เจว อาปชฺชติ, พหิสีมาย จ, อตฺถาปตฺติ เนว อนฺโตสีมาย อาปชฺชติ, โน พหิสีมาย. อตฺถาปตฺติ คาเม อาปชฺชติ, โน อรญฺเญ, อตฺถาปตฺติ อรญฺเญ อาปชฺชติ, โน คาเม, อตฺถาปตฺติ คาเม เจว อาปชฺชติ อรญฺเญ จ, อตฺถาปตฺติ เนว คาเม อาปชฺชติ, โน อรญฺเญ.}

\prosewithcloser{จตสฺโส โจทนา วตฺถุ\-สนฺทสฺสนา อาปตฺติสนฺทสฺสานา สํวาส\-ปฏิกฺเขโป สามีจิ\-ปฏิกฺเขโป. จตฺตาโร ปุพฺพกิจฺจา. จตฺตาโร ปตฺตกลฺลา. จตฺตาริ อนญฺญ\-ปาจิตฺติยานิ. จตสฺโส ภิกฺขุ\-สมฺมุติโย.\csromansymbolfootnote{*}{อํ 1. 325, 6; อภิ 2. 390; ที 3. 191 ปิฏฺเฐสุปิ.}จตฺตาริ อคติคมนานิ, ฉนฺทาคติํ คจฺฉติ, โทสาคติํ คจฺฉติ, โมหคติํ คจฺฉติ, ภยาคติํ คจฺฉติ. จตฺตาริ นาคติคมนานิ, น ฉนฺทาคติํ คจฺฉติ, น โทสาคติํ คจฺฉติ, น โมหาคติํ คจฺฉติ, น ภยาคติํ คจฺฉติ. จตูหงฺเคหิ สมนฺนาคโต อลชฺชี ภิกฺขุ สํฆํ ภินฺทติ ฉนฺทาคติํ คจฺฉนฺโต โทสาคติํ คจฺฉนฺโต โมหาคติํ คจฺฉนฺโต ภยาคติํ คจฺฉนฺโต. จตูหงฺเคหิ สมนฺนาคโต เปสโล ภิกฺขุ ภินฺนํ สํฆํ สมคฺคํ กโรติ น ฉนฺทาคติํ คจฺฉนฺโต น โทสาคติํ คจฺฉนฺโต น โมหาคติํ คจฺฉนฺโต น ภยาคติํ คจฺฉนฺโต. จตูหงฺเคหิ สมนฺนาคตสฺส ภิกฺขุโน วินโย น ปุจฺฉิตพฺโพ, ฉนฺทาคติํ คจฺฉติ, โทสาคติํ คจฺฉติ, โมหาคติํ คจฺฉติ, ภยาคติํ คจฺฉติ. จตูหงฺเคหิ สมนฺนาคเตน ภิกฺขุนา วินโย น ปุจฺฉิตพฺโพ, ฉนฺทาคติํ คจฺฉติ, โทสาคติํ คจฺฉติ, โมหาคติํ คจฺฉติ, ภยาคติํ คจฺฉติ. จตูหงฺเคหิ สมนฺนาคตสฺส ภิกฺขุโน วินโย น วิสฺสชฺเชตพฺโพ, ฉนฺทาคติํ คจฺฉติ, โทสาคติํ คจฺฉติ, โมหาคติํ คจฺฉติ, ภยาคติํ คจฺฉติ. จตูหงฺเคหิ สมนฺนาคเตน ภิกฺขุนา วินโย น วิสฺสชฺเชตพฺโพ, ฉนฺทาคติํ คจฺฉติ, โทสาคติํ คจฺฉติ, โมหาคติํ คจฺฉติ, ภยาคติํ คจฺฉติ. จตูหงฺเคหิ สมนฺนาคตสฺส ภิกฺขุโน อนุโยโค น ทาตพฺโพ, ฉนฺทาคติํ คจฺฉติ, โทสาคติํ คจฺฉติ, โมหาคติํ คจฺฉติ, ภยาคติํ คจฺฉติ. \csromanfolio{229}จตูหงฺเคหิ สมนฺนาคเตน ภิกฺขุนา สทฺธิํ วินโย น สากจฺฉิตพฺโพ, ฉนฺทาคติํ คจฺฉติ, โทสาคติํ คจฺฉติ, โมหาคติํ คจฺฉติ, ภยาคติํ คจฺฉติ. อตฺถาปตฺติ คิลาโน อาปชฺชติ, โน อคิลาโน, อตฺถาปตฺติ อคิลาโน อาปชฺชติ, โน คิลาโน, อตฺถาปตฺติ คิลาโน เจว อาปชฺชติ อคิลาโน จ, อตฺถาปตฺติ เนว คิลาโน อาปชฺชติ, โน อคิลาโน. จตฺตาริ อธมฺมิกานิ ปาติโมกฺข\-ฏฺฐปนานิ. จตฺตาริ ธมฺมิกานิ ปาติโมกฺข\-ฏฺฐปนานิ.}{จตุกฺกํ นิฏฺฐิตํ.}
\csromancenter{ตสฺสุทฺทานํ.}
\setlength{\csromangathapagewidth}{0pt}
\setlength{\csromangathawakwidth}{0pt}
\csromangathameasuregroup{สกวาจาย กาเยน, \\ อาปชฺชนฺโต จ กมฺเมน, \\ ภิกฺขูนํ ภิกฺขุนีนญฺจ, \\ อชานกาเย มชฺเฌ จ, \\ ปฏิลาเภน โจทนา, \\ มานตฺตจาริกา จาปิ, \\ มหาวิกฏกมฺมานิ, \\ อธิกรณา ทุสฺสีลา จ, \\ คมิโก วตฺถุนานตฺตา, \\ อาจริโย ปจฺจยา วา, \\ อาทิยนฺโต ปุคฺคโล จ, \\ กาเล จ กปฺปติ เจว, \\ อนฺโต อนฺโต จ สีมาย, \\ ปุพฺพกิจฺจํ ปตฺตกลฺลํ, \\ อคติ นาคติ เจว, \\ ปุจฺฉิตพฺพา ทุเว เจว, \\ อนุโยโค จ สากจฺฉา,}{\csromangathabat{สกวาจาย กาเยน,}{ปสุตฺโต จ อจิตฺตโก.} \\ \csromangathabat{อาปชฺชนฺโต จ กมฺเมน,}{โวหารา จตุโร ตถา.} \\ \csromangathabat{ภิกฺขูนํ ภิกฺขุนีนญฺจ,}{ปริกฺขาโร จ สมฺมุขา.} \\ \csromangathabat{อชานกาเย มชฺเฌ จ,}{วุฏฺฐาติ ทุวิธา ตถา.} \\ \csromangathabat{ปฏิลาเภน โจทนา,}{ปริวาสา จ วุจฺจติ.} \\ \csromangathabat{มานตฺตจาริกา จาปิ,}{สามุกฺกํสา ปฏิคฺคหิ.} \\ \csromangathabat{มหาวิกฏกมฺมานิ,}{ปุน กมฺเม วิปตฺติโย.} \\ \csromangathabat{อธิกรณา ทุสฺสีลา จ,}{โสภนาคนฺตุเกน จ.} \\ \csromangathabat{คมิโก วตฺถุนานตฺตา,}{สภาคุปชฺฌาเยน จ.} \\ \csromangathabat{อาจริโย ปจฺจยา วา,}{ทุจฺจริตํ สุจริตํ.} \\ \csromangathabat{อาทิยนฺโต ปุคฺคโล จ,}{อรโห อาสเนน จ.} \\ \csromangathabat{กาเล จ กปฺปติ เจว,}{ปจฺจนฺติเมสุ กปฺปติ.} \\ \csromangathabat{อนฺโต อนฺโต จ สีมาย,}{คาเม จ โจทนาย จ.} \\ \csromangathabat{ปุพฺพกิจฺจํ ปตฺตกลฺลํ,}{อนญฺญา สมฺมุติโย จ.} \\ \csromangathabat{อคติ นาคติ เจว,}{อลชฺชี เปสเลน จ.} \\ \csromangathabat{ปุจฺฉิตพฺพา ทุเว เจว,}{วิสฺสชฺเชยฺยา ตถา ทุเว.} \\ \csromangathabat{อนุโยโค จ สากจฺฉา,}{คิลาโน ฐปเนน จาติ.}}
\csromangathasetleft{สกวาจาย กาเยน, \\ อาปชฺชนฺโต จ กมฺเมน, \\ ภิกฺขูนํ ภิกฺขุนีนญฺจ, \\ อชานกาเย มชฺเฌ จ, \\ ปฏิลาเภน โจทนา, \\ มานตฺตจาริกา จาปิ, \\ มหาวิกฏกมฺมานิ, \\ อธิกรณา ทุสฺสีลา จ, \\ คมิโก วตฺถุนานตฺตา, \\ อาจริโย ปจฺจยา วา, \\ อาทิยนฺโต ปุคฺคโล จ, \\ กาเล จ กปฺปติ เจว, \\ อนฺโต อนฺโต จ สีมาย, \\ ปุพฺพกิจฺจํ ปตฺตกลฺลํ, \\ อคติ นาคติ เจว, \\ ปุจฺฉิตพฺพา ทุเว เจว, \\ อนุโยโค จ สากจฺฉา,}
\csromangathagroup{\csromangathastanza{\csromangathabat{สกวาจาย กาเยน,}{ปสุตฺโต จ อจิตฺตโก.} \\ \csromangathabat{อาปชฺชนฺโต จ กมฺเมน,}{โวหารา จตุโร ตถา.}}\csromangathastanza{\csromangathabat{ภิกฺขูนํ ภิกฺขุนีนญฺจ,}{ปริกฺขาโร จ สมฺมุขา.} \\ \csromangathabat{อชานกาเย มชฺเฌ จ,}{วุฏฺฐาติ ทุวิธา ตถา.}}\csromangathastanza{\csromangathabat{ปฏิลาเภน โจทนา,}{ปริวาสา จ วุจฺจติ.} \\ \csromangathabat{มานตฺตจาริกา จาปิ,}{สามุกฺกํสา ปฏิคฺคหิ.}}\csromangathastanza{\csromangathabat{มหาวิกฏกมฺมานิ,}{ปุน กมฺเม วิปตฺติโย.} \\ \csromangathabat{อธิกรณา ทุสฺสีลา จ,}{โสภนาคนฺตุเกน จ.}}\csromangathastanza{\csromangathabat{คมิโก วตฺถุนานตฺตา,}{สภาคุปชฺฌาเยน จ.} \\ \csromangathabat{อาจริโย ปจฺจยา วา,}{ทุจฺจริตํ สุจริตํ.}}\csromangathastanza{\csromangathabat{อาทิยนฺโต ปุคฺคโล จ,}{อรโห อาสเนน จ.} \\ \csromangathabat{กาเล จ กปฺปติ เจว,}{ปจฺจนฺติเมสุ กปฺปติ.}}\csromangathastanza{\csromangathabat{อนฺโต อนฺโต จ สีมาย,}{คาเม จ โจทนาย จ.} \\ \csromangathabat{ปุพฺพกิจฺจํ ปตฺตกลฺลํ,}{อนญฺญา สมฺมุติโย จ.}}\csromangathastanza{\csromangathabat{อคติ นาคติ เจว,}{อลชฺชี เปสเลน จ.} \\ \csromangathabat{ปุจฺฉิตพฺพา ทุเว เจว,}{วิสฺสชฺเชยฺยา ตถา ทุเว.} \\ \csromangathabat{อนุโยโค จ สากจฺฉา,}{คิลาโน ฐปเนน จาติ.}}}

\csromanmatikaanchor{mtk.168}
\csromantocmarkref{paragraph}{5. ปญฺจกวาร}{mtk.168}
\bookmark[dest={mtk.168},level=4]{5. ปญฺจกวาร}
\csromanheader{4}{5. \textbf{ปญฺจกวาร}}
\proseitem{325}{\csromanfolio{230}ปญฺจ อาปตฺติโย. ปญฺจ อาปตฺติกฺขนฺธา. ปญฺจ วินีตวตฺถูนิ. ปญฺจ กมฺมานิ อานนฺตริกานิ. ปญฺจ ปุคฺคลา นิยตา. ปญฺจ เฉทนกา อาปตฺติโย. ปญฺจหากาเรหิ อาปตฺติํ อาปชฺชติ. ปญฺจ อาปตฺติโย มุสาวาท\-ปจฺจยา. ปญฺจหากาเรหิ กมฺมํ น อุเปติ, สยํ วา กมฺมํ น กโรติ, ปรํ วา น อชฺเฌสติ, ฉนฺทํ วา ปาริสุทฺธิํ วา น เทติ, กยิรมาเน กมฺเม ปฏิกฺโกสติ, กเต วา ปน กมฺเม อธมฺมทิฏฺฐิ โหติ. ปญฺจหากาเรหิ กมฺมํ อุเปติ, สยํ วา กมฺมํ กโรติ, ปรํ วา อชฺเฌสติ, ฉนฺทํ วา ปาริสุทฺธิํ วา เทติ, กยิรมาเน กมฺเม นปฺปฏิกฺโกสติ, กเต วา ปน กมฺเม ธมฺมทิฏฺฐิ โหติ. ปญฺจ ปิณฺฑปาติกสฺส ภิกฺขุโน กปฺปนฺติ อนามนฺตจาโร คณโภชนํ ปรมฺปร\-โภชนํ อนธิฏฺฐานํ อวิกปฺปนา.\csromansymbolfootnote{*}{อํ 2. 112 ปิฏฺเฐ.}ปญฺจหงฺเคหิ สมนฺนาคโต ภิกฺขุ อุสฺสงฺกิตปริสงฺกิโต โหติ ปาปภิกฺขุปิ อกุปฺปธมฺโมปิ, เวสิยาโคจโร วา โหติ, วิธวาโคจโร วา โหติ, ถุลฺลกุมาริ\-โคจโร วา โหติ, ปณฺฑกโคจโร วา โหติ, ภิกฺขุนิ\-โคจโร วา โหติ. ปญฺจ เตลานิ ติลเตลํ สาสปเตลํ มธุกเตลํ เอรณฺฑกเตลํ วสาเตลํ. ปญฺจ วสานิ อจฺฉวสํ มจฺฉวสํ สุสุกาวสํ สูกรวสํ คทฺรภวสํ.\csromansymbolfootnote{+}{อํ 2. 129, 131; ที 3. 196 ปิฏฺเฐสุปิ.}ปญฺจ พฺยสนานิ ญาติพฺยสนํ โภคพฺยสนํ โรคพฺยสนํ สีลพฺยสนํ ทิฏฺฐิพฺยสนํ. + ปญฺจ สมฺปทา ญาติ สมฺปทา โภคสมฺปทา อโรคสมฺปทา สีลสมฺปทา ทิฏฺฐิสมฺปทา. ปญฺจ นิสฺสย\-ปฏิปฺปสฺสทฺธิโย อุปชฺฌายมฺหา, อุปชฺฌาโย ปกฺกนฺโต วา โหติ, วิพฺภนฺโต วา, กาลงฺกโต วา, ปกฺขสงฺกนฺโต วา, อาณตฺติเยว ปญฺจมี. ปญฺจ ปุคฺคลา น อุปสมฺ\-ปาเท\-ตพฺพา อทฺธานหีโน องฺคหีโน วตฺถุวิปนฺโน กรณ\-ทุกฺกฏโก อปริปูโร. ปญฺจ ปํสุกูลานิ โสสานิกํ ปาปณิกํ อุนฺทูรกฺขายิกํ อุปจิก\-กฺขายิกํ อคฺคิทฑฺฒํ. อปรานิปิ ปญฺจ ปํสุกูลานิ โคขายิกํ อชกฺขายิกํ ถูปจีวรํ อาภิเสกิกํ คตปฏิยาคตํ\csromansharedfootnote{29255efcd498e46a}{ภตปฏิยาภตํ (ก)}. ปญฺจ อวหารา เถยฺยาวหาโร ปสยฺหาวหาโร ปริกปฺปา\-วหาโร ปฏิจฺฉนฺนา\-วหาโร กุสาวหาโร. ปญฺจ มหาโจรา สนฺโต สํวิชฺชมานา โลกสฺมิํ. ปญฺจ อวิสฺสชฺชนิยานิ. ปญฺจ อเวภงฺคิยานิ. ปญฺจาปตฺติโย กายโต สมุฏฺฐนฺติ, น วาจโต, นจิตฺตโต. ปญฺจาปตฺติโย กายโต จ วาจโต จ สมุฏฺฐนฺติ, \csromanfolio{231}น จิตฺตโต. ปญฺจาปตฺติโย เทสนาคามินิโย. ปญฺจ สํฆา. ปญฺจ ปาติโมกฺขุ\-ทฺเทสา. สพฺพ\-ปจฺจนฺติเมสุ ชนปเทสุ วินยธร\-ปญฺจเมน คเณน อุปสมฺ\-ปาเท\-ตพฺพํ. ปญฺจานิสํสา กถินตฺถาเร. ปญฺจ กมฺมานิ. ยาวตติยเก ปญฺจ อาปตฺติโย. ปญฺจหากาเรหิ อทินฺนํ อาทิยนฺตสฺส อาปตฺติ ปาราชิกสฺส. ปญฺจหากาเรหิ อทินฺนํ อาทิยนฺตสฺส อาปตฺติ ถุลฺลจฺจยสฺส. ปญฺจหากาเรหิ อทินฺนํ อาทิยนฺตสฺส อาปตฺติ ทุกฺกฏสฺส. ปญฺจ อกปฺปิยานิ น ปริภุญฺชิ\-ตพฺพานิ, อทินฺนญฺจ โหติ, อวิทิตญฺจ โหติ, อกปฺปิยญฺจ โหติ, อปฺปฏิคฺคหิตญฺจ โหติ, อกตาติริตฺตญฺจ โหติ. ปญฺจ กปฺปิยานิ ปริภุญฺชิ\-ตพฺพานิ, ทินฺนญฺจ โหติ, วิทิตญฺจ โหติ, กปฺปิยญฺจ โหติ, ปฏิคฺคหิตญฺจ โหติ, กตาติริตฺตญฺจ โหติ. ปญฺจ ทานานิ อปุญฺญานิ ปุญฺญสมฺมตานิ โลกสฺมิํ\csromansharedfootnote{7ab6f9f1c22228a1}{โลกสฺส (สฺยา)} มชฺชทานํ สมชฺชทานํ อิตฺถิทานํ อุสภทานํ จิตฺตกมฺม\-ทานํ. ปญฺจ อุปฺปนฺนา ทุปฺปฏิวิโนทยา\csromansharedfootnote{083382fba26937fb}{วิโนทิยา (สฺยา)} อุปฺปนฺโน ราโค ทุปฺปฏิวิโนทโย, อุปฺปนฺโน โทโส ทุปฺปฏิวิโนทโย, อุปฺปนฺโน โมโห ทุปฺปฏิวิโนทโย, อุปฺปนฺนํ ปฏิภานํ ทุปฺปฏิวิโนทยํ, อุปฺปนฺนํ คมิยจิตฺตํ ทุปฺปฏิวิโนทยํ. ปญฺจานิสํสา สมฺมชฺชนิยา, สกจิตฺตํ ปสีทติ, ปรจิตฺตํ ปสีทติ, เทวตา อตฺตมนา โหนฺติ, ปาสาทิกสํวตฺตนิก\-กมฺมํ อุปจินติ, กายสฺส เภทา ปรํ มรณา สุคติํ สคฺคํ โลกํ อุปปชฺชติ. อปเรปิ ปญฺจานิสํสา สมฺมชฺชนิยา, สกจิตฺตํ ปสีทติ, ปรจิตฺตํ ปสีทติ, เทวตา อตฺตมนา โหนฺติ, สตฺถุสาสนํ กตํ โหติ, ปจฺฉิมา ชนตา ทิฏฺฐานุคติํ อาปชฺชติ.}

\prose{ปญฺจหงฺเคหิ สมนฺนาคโต วินยธโร “พาโล”เตฺวว สงฺขํ คจฺฉติ, อตฺตโน ภาส\-ปริยนฺตํ น อุคฺคณฺหาติ, ปรสฺส ภาส\-ปริยนฺตํ น อุคฺคณฺหาติ, อตฺตโน ภาส\-ปริยนฺตํ น อุคฺคเหตฺวา ปรสฺส ภาส\-ปริยนฺตํ น อุคฺคเหตฺวา อธมฺเมน กาเรติ อปฺปฏิญฺญาย. ปญฺจหงฺเคหิ สมนฺนาคโต วินยธโร “ปณฺฑิโต”เตฺวว สงฺขํ คจฺฉติ, อตฺตโน ภาส\-ปริยนฺตํ อุคฺคณฺหาติ, ปรสฺส ภาส\-ปริยนฺตํ อุคฺคณฺหาติ, อตฺตโน ภาส\-ปริยนฺตํ อุคฺคเหตฺวา ปรสฺส ภาส\-ปริยนฺตํ อุคฺคเหตฺวา ธมฺเมน กาเรติ ปฏิญฺญาย. อปเรหิปิ ปญฺจหงฺเคหิ สมนฺนาคโต วินยธโร “พาโล”เตฺวว สงฺขํ คจฺฉติ, อาปตฺติํ น ชานาติ, อาปตฺติยา มูลํ น ชานาติ, \csromanfolio{232}อาปตฺติ\-สมุทยํ น ชานาติ, อาปตฺตินิโรธํ น ชานาติ, อาปตฺติ\-นิโรธคามินิํ ปฏิปทํ น ชานาติ. ปญฺจหงฺเคหิ สมนฺนาคโต วินยธโร “ปณฺฑิโต”เตฺวว สงฺขํ คจฺฉติ, อาปตฺติํ ชานาติ, อาปตฺติยา มูลํ ชานาติ, อาปตฺติ\-สมุทยํ ชานาติ, อาปตฺตินิโรธํ ชานาติ, อาปตฺติ\-นิโรธคามินิํ ปฏิปทํ ชานาติ. อปเรหิปิ ปญฺจหงฺเคหิ สมนฺนาคโต วินยธโร “พาโล”เตฺวว สงฺขํ คจฺฉติ, อธิกรณํ น ชานาติ, อธิกรณสฺส มูลํ น ชานาติ, อธิกรณ\-สมุทยํ น ชานาติ, อธิกรณ\-นิโรธํ น ชานาติ, อธิกรณนิโรธ\-คามินิํ ปฏิปทํ น ชานาติ. ปญฺจหงฺเคหิ สมนฺนาคโต วินยธโร “ปณฺฑิโต”เตฺวว สงฺขํ คจฺฉติ, อธิกรณํ ชานาติ, อธิกรณสฺส มูลํ ชานาติ, อธิกรณ\-สมุทยํ ชานาติ, อธิกรณ\-นิโรธํ ชานาติ, อธิกรณนิโรธ\-คามินิํ ปฏิปทํ ชานาติ. อปเรหิปิ ปญฺจหงฺเคหิ สมนฺนาคโต วินยธโร “พาโล”เตฺวว สงฺขํ คจฺฉติ, วตฺถุํ น ชานาติ, นิทานํ น ชานาติ, ปญฺญตฺติํ น ชานาติ, อนุปญฺญตฺติํ น ชานาติ, อนุสนฺธิ\-วจนปถํ น ชานาติ. ปญฺจหงฺเคหิ สมนฺนาคโต วินยธโร “ปณฺฑิโต”เตฺวว สงฺขํ คจฺฉติ, วตฺถุํ ชานาติ, นิทานํ ชานาติ, ปญฺญตฺติํ ชานาติ, อนุปญฺญตฺติํ ชานาติ, อนุสนฺธิ\-วจนปถํ ชานาติ. อปเรหิปิ ปญฺจหงฺเคหิ สมนฺนาคโต วินยธโร “พาโล”เตฺวว สงฺขํ คจฺฉติ, ญตฺติํ น ชานาติ, ญตฺติยา กรณํ น ชานาติ, น ปุพฺพกุสโล โหติ, น อปรกุสโล โหติ, อกาลญฺญู จ โหติ. ปญฺจหงฺเคหิ สมนฺนาคโต วินยธโร “ปณฺฑิโต”เตฺวว สงฺขํ คจฺฉติ, ญตฺติํ ชานาติ, ญตฺติยา กรณํ ชานาติ, ปุพฺพกุสโล โหติ, อปรกุสโล โหติ, กาลญฺญู จ โหติ. อปเรหิปิ ปญฺจหงฺเคหิ สมนฺนาคโต วินยธโร “พาโล”เตฺวว สงฺขํ คจฺฉติ, อาปตฺตานาปตฺติํ น ชานาติ, ลหุกครุกํ อาปตฺติํ น ชานาติ, สาวเสสา\-นวเสสํ อาปตฺติํ น ชานาติ, ทุฏฺฐุลฺลา\-ทุฏฺฐุลฺลํ อาปตฺติํ น ชานาติ, อาจริยปรมฺปรา โข ปนสฺส น สุคฺคหิตา โหติ น สุมนสิกตา น สูปธาริตา. ปญฺจหงฺเคหิ สมนฺนาคโต วินยธโร “ปณฺฑิโต”เตฺวว สงฺขํ คจฺฉติ, อาปตฺตานาปตฺติํ ชานาติ, ลหุกครุกํ อาปตฺติํ ชานาติ, สาวเสสา\-นวเสสํ อาปตฺติํ ชานาติ, ทุฏฺฐุลฺลา\-ทุฏฺฐุลฺลํ อาปตฺติํ ชานาติ, อาจริยปรมฺปรา โข ปนสฺส สุคฺคหิตา โหติ สุมนสิกตา สูปธาริตา. อปเรหิปิ ปญฺจหงฺเคหิ สมนฺนาคโต วินยธโร “พาโล”เตฺวว สงฺขํ คจฺฉติ, อาปตฺตานาปตฺติํ น ชานาติ, ลหุกครุกํ อาปตฺติํ น ชานาติ, \csromanfolio{233}สาวเสสา\-นวเสสํ อาปตฺติํ น ชานาติ, ทุฏฺฐุลฺลา\-ทุฏฺฐุลฺลํ อาปตฺติํ น ชานาติ, อุภยานิ โข ปนสฺส ปาติโมกฺขานิ น วิตฺถาเรน สฺวาคตานิ โหนฺติ น สุวิภตฺตานิ น สุปฺปวตฺตีนิ น สุวินิจฺฉิตานิ สุตฺตโส อนุพฺยญฺชนโส. ปญฺจหงฺเคหิ สมนฺนาคโต วินยธโร “ปณฺฑิโต”เตฺวว สงฺขํ คจฺฉติ, อาปตฺตานาปตฺติํ ชานาติ, ลหุกครุกํ อาปตฺติํ ชานาติ, สาวเสสา\-นวเสสํ อาปตฺติํ ชานาติ, ทุฏฺฐุลฺลา\-ทุฏฺฐุลฺลํ อาปตฺติํ ชานาติ, อุภยานิ โข ปนสฺส ปาติโมกฺขานิ วิตฺถาเรน สฺวาคตานิ โหนฺติ สุวิภตฺตานิ สุปฺปวตฺตีนิ สุวินิจฺฉิตานิ สุตฺตโส อนุพฺยญฺชนโส. อปเรหิปิ ปญฺจหงฺเคหิ สมนฺนาคโต วินยธโร “พาโล”เตฺวว สงฺขํ คจฺฉติ, อาปตฺตานาปตฺติํ น ชานาติ, ลหุกครุกํ อาปตฺติํ น ชานาติ, สาวเสสา\-นวเสสํ อาปตฺติํ น ชานาติ, ทุฏฺฐุลฺลา\-ทุฏฺฐุลฺลํ อาปตฺติํ น ชานาติ, อธิกรเณ จ น วินิจฺฉย\-กุสโล โหติ. ปญฺจหงฺเคหิ สมนฺนาคโต วินยธโร “ปณฺฑิโต”เตฺวว สงฺขํ คจฺฉติ, อาปตฺตานาปตฺติํ ชานาติ, ลหุกครุกํ อาปตฺติํ ชานาติ, สาวเสสา\-นวเสสํ อาปตฺติํ ชานาติ, ทุฏฺฐุลฺลา\-ทุฏฺฐุลฺลํ อาปตฺติํ ชานาติ, อธิกรเณ จ วินิจฺฉย\-กุสโล โหติ.}

\prose{\csromansymbolfootnote{*}{อํ 2. 193; วิ 5. 335 ปิฏฺเฐสุปิ.}ปญฺจ อารญฺญิกา, มนฺทตฺตา โมมูหตฺตา อารญฺญิโก โหติ, ปาปิจฺโฉ อิจฺฉาปกโต อารญฺญิโก โหติ, อุมฺมาทา จิตฺตกฺเขปา อารญฺญิโก โหติ, “วณฺณิตํ พุทฺเธหิ พุทฺธสาวเกหี”ติ อารญฺญิโก โหติ, อปิ จ อปฺปิจฺฉญฺเญว นิสฺสาย สนฺตุฏฺฐิญฺเญว นิสฺสาย สลฺเลขญฺเญว นิสฺสาย ปวิเวกญฺเญว นิสฺสาย อิทมตฺถิตญฺเญว\csromansharedfootnote{49d011ca9d933377}{อิทมฏฺฐิตญฺเญว (สี, สฺยา)} นิสฺสาย อารญฺญิโก โหติ. ปญฺจ ปิณฺฑปาติกา ฯเปฯ. ปญฺจ ปํสุกูลิกา ฯเปฯ. ปญฺจ รุกฺขมูลิกา ฯเปฯ. ปญฺจ โสสานิกา ฯเปฯ. ปญฺจ อพฺโภกาสิกา ฯเปฯ. ปญฺจ เตจีวริกา ฯเปฯ. ปญฺจ สปทานจาริกา ฯเปฯ. ปญฺจ เนสชฺชิกา ฯเปฯ. ปญฺจ ยถา\-สนฺถติกา ฯเปฯ. ปญฺจ เอกาสนิกา ฯเปฯ. ปญฺจ ขลุปจฺฉาภตฺติกา ฯเปฯ. ปญฺจ ปตฺตปิณฺฑิกา, มนฺทตฺตา โมมูหตฺตา ปตฺตปิณฺฑิโก โหติ, ปาปิจฺโฉ อิจฺฉาปกโต ปตฺตปิณฺฑิโก โหติ, อุมฺมาทา จิตฺตกฺเขปา ปตฺตปิณฺฑิโก โหติ, “วณฺณิตํ พุทฺเธหิ พุทฺธสาวเกหี”ติ ปตฺตปิณฺฑิโก โหติ, อปิ จ อปฺปิจฺฉญฺเญว นิสฺสาย สนฺตุฏฺฐิญฺเญว นิสฺสาย สลฺเลขญฺเญว นิสฺสาย ปวิเวกญฺเญว นิสฺสาย อิทมตฺถิตญฺเญว นิสฺสาย ปตฺตปิณฺฑิโก โหติ.}

\prose{\csromanfolio{234}\csromansymbolfootnote{*}{วิ 5. 131 ปิฏฺเฐปิ.}ปญฺจหงฺเคหิ สมนฺนาคเตน ภิกฺขุนา นานิสฺสิเตน วตฺถพฺพํ, อุโปสถํ น ชานาติ, อุโปสถกมฺมํ น ชานาติ, ปาติโมกฺขํ น ชานาติ, ปาติโมกฺขุ\-ทฺเทสํ น ชานาติ, อูนปญฺจวสฺโส โหติ. ปญฺจหงฺเคหิ สมนฺนาคเตน ภิกฺขุนา อนิสฺสิเตน วตฺถพฺพํ, อุโปสถํ ชานาติ, อุโปสถกมฺมํ ชานาติ, ปาติโมกฺขํ ชานาติ, ปาติโมกฺขุ\-ทฺเทสํ ชานาติ, ปญฺจวสฺโส วา โหติ อติเรก\-ปญฺจวสฺโส วา. อปเรหิปิ ปญฺจหงฺเคหิ สมนฺนาคเตน ภิกฺขุนา นานิสฺสิเตน วตฺถพฺพํ, ปวารณํ น ชานาติ, ปวารณากมฺมํ น ชานาติ, ปาติโมกฺขํ น ชานาติ, ปาติโมกฺขุ\-ทฺเทสํ น ชานาติ, อูนปญฺจวสฺโส โหติ. ปญฺจหงฺเคหิ สมนฺนาคเตน ภิกฺขุนา อนิสฺสิเตน วตฺถพฺพํ, ปวารณํ ชานาติ, ปวารณากมฺมํ ชานาติ, ปาติโมกฺขํ ชานาติ, ปาติโมกฺขุ\-ทฺเทสํ ชานาติ, ปญฺจวสฺโส วา โหติ อติเรก\-ปญฺจวสฺโส วา. อปเรหิปิ ปญฺจหงฺเคหิ สมนฺนาคเตน ภิกฺขุนา นานิสฺสิเตน วตฺถพฺพํ, อาปตฺตานาปตฺติํ น ชานาติ, ลหุกครุกํ อาปตฺติํ น ชานาติ, สาวเสสา\-นวเสสํ อาปตฺติํ น ชานาติ, ทุฏฺฐุลฺลา\-ทุฏฺฐุลฺลํ อาปตฺติํ น ชานาติ, อูนปญฺจวสฺโส โหติ. ปญฺจหงฺเคหิ สมนฺนาคเตน ภิกฺขุนา อนิสฺสิเตน วตฺถพฺพํ, อาปตฺตานาปตฺติํ ชานาติ, ลหุกครุกํ อาปตฺติํ ชานาติ, สาวเสสา\-นวเสสํ อาปตฺติํ ชานาติ, ทุฏฺฐุลฺลา\-ทุฏฺฐุลฺลํ อาปตฺติํ ชานาติ, ปญฺจวสฺโส วา โหติ อติเรก\-ปญฺจวสฺโส วา. ปญฺจหงฺเคหิ สมนฺนาคตาย ภิกฺขุนิยา นานิสฺสิตาย วตฺถพฺพํ, อุโปสถํ น ชานาติ, อุโปสถกมฺมํ น ชานาติ, ปาติโมกฺขํ น ชานาติ, ปาติโมกฺขุ\-ทฺเทสํ น ชานาติ, อูนปญฺจวสฺสา โหติ. ปญฺจหงฺเคหิ สมนฺนาคตาย ภิกฺขุนิยา อนิสฺสิตาย วตฺถพฺพํ, อุโปสถํ ชานาติ, อุโปสถกมฺมํ ชานาติ, ปาติโมกฺขํ ชานาติ, ปาติโมกฺขุ\-ทฺเทสํ ชานาติ, ปญฺจวสฺสา วา โหติ อติเรก\-ปญฺจวสฺสา วา. อปเรหิปิ ปญฺจหงฺเคหิ สนฺนาคตาย ภิกฺขุนิยา นานิสฺสิตาย วตฺตพฺพํ, ปวารณํ น ชานาติ, ปวารณากมฺมํ น ชานาติ, ปาติโมกฺขํ น ชานาติ, ปาติโมกฺขุ\-ทฺเทสํ น ชานาติ, อูนปญฺจวสฺสา โหติ. ปญฺจหงฺเคหิ สมนฺนาคตาย ภิกฺขุนิยา อนิสฺสิตาย วตฺถพฺพํ, ปวารณํ ชานาติ, ปวารณากมฺมํ ชานาติ, ปาติโมกฺขํ ชานาติ, ปาติโมกฺขุ\-ทฺเทสํ ชานาติ, ปญฺจวสฺสา วา โหติ อติเรก\-ปญฺจวสฺสา วา. อปเรหิปิ ปญฺจหงฺเคหิ สมนฺนาคตาย ภิกฺขุนิยา นานิสฺสิตาย \csromanfolio{235}วตฺถพฺพํ, อาปตฺตานาปตฺติํ น ชานาติ, ลหุกครุกํ อาปตฺติํ น ชานาติ, สาวเสสา\-นวเสสํ อาปตฺติํ น ชานาติ, ทุฏฺฐุลฺลา\-ทุฏฺฐุลฺลํ อาปตฺติํ น ชานาติ, อูนปญฺจวสฺสา โหติ. ปญฺจหงฺเคหิ สมนฺนาคตาย ภิกฺขุนิยา อนิสฺสิตาย วตฺถพฺพํ, อาปตฺตานาปตฺติํ ชานาติ, ลหุกครุกํ อาปตฺติํ ชานาติ, สาวเสสา\-นวเสสํ อาปตฺติํ ชานาติ, ทุฏฺฐุลฺลา\-ทุฏฺฐุลฺลํ อาปตฺติํ ชานาติ, ปญฺจวสฺสา วา โหติ อติเรก\-ปญฺจวสฺสา วา.}

\prose{ปญฺจ อาทีนวา อปาสาทิเก, อตฺตาปิ อตฺตานํ อุปวทติ, อนุวิจฺจปิ วิญฺญู ครหนฺติ, ปาปโก กิตฺติสทฺโท อพฺภุคฺคจฺฉติ, สมฺมูโฬฺห กาลํ กโรติ, กายสฺส เภทา ปรํ มรณา อปายํ ทุคฺคติํ วินิปาตํ นิรยํ อุปปชฺชติ. ปญฺจานิสํสา ปาสาทิเก, อตฺตาปิ อตฺตานํ น อุปวทติ, อนุวิจฺจปิ วิญฺญู ปสํสนฺติ, กลฺยาโณ กิตฺติสทฺโท อพฺภุคฺคจฺฉติ, อสมฺมูโฬฺห กาลํ กโรติ, กายสฺส เภทา ปรํ มรณา สุคติํ สคฺคํ โลกํ อุปปชฺชติ. อปเรปิ ปญฺจ อาทีนวา อปาสาทิเก, อปฺปสนฺนา น ปสีทนฺติ, ปสนฺนานํ เอกจฺจานํ อญฺญถตฺตํ โหติ, สตฺถุสาสนํ อกตํ โหติ, ปจฺฉิมา ชนตา ทิฏฺฐานุคติํ นาปชฺชติ, จิตฺตมสฺส น ปสีทติ. ปญฺจานิสํสา ปาสาทิเก, อปฺปสนฺนา ปสีทนฺติ, ปสนฺนานํ ภิยฺโยภาวาย โหติ, สตฺถุสาสนํ กตํ โหติ, ปจฺฉิมา ชนตา ทิฏฺฐานุคติํ อาปชฺชติ, จิตฺตมสฺส ปสีทติ. ปญฺจ อาทีนวา กุลูปเก, อนามนฺตจาเร อาปชฺชติ, รโห นิสชฺชาย อาปชฺชติ, ปฏิจฺฉนฺเน อาสเน อาปชฺชติ, มาตุคามสฺส อุตฺตริ\-ฉปฺปญฺจ\-วาจาหิ ธมฺมํ เทเสนฺโต อาปชฺชติ, กามสงฺกปฺป\-พหุโล จ วิหรติ. ปญฺจ อาทีนวา กุลูปกสฺส ภิกฺขุโน อติเวลํ กุเลสุ สํสฏฺฐสฺส วิหรโต, มาตุคามสฺส อภิณฺหทสฺสนํ, ทสฺสาเน สติ สํสคฺโค, สํสคฺเค สติ วิสฺสาโส, วิสฺสาเส สติ โอตาโร, โอติณฺณจิตฺตสฺเสตํ ภิกฺขุโน ปาฏิกงฺขํ อนภิรโต วา พฺรหฺมจริยํ จริสฺสติ, อญฺญตรํ วา สํกิลิฏฺฐํ อาปตฺติํ อาปชฺชิสฺสติ, สิกฺขํ วา ปจฺจกฺขาย หีนายาวตฺติสฺสติ.}

\prosewithcloser{\csromansymbolfootnote{*}{วิ 2. 52 ปิฏฺเฐ.}ปญฺจ พีชชาตานิ มูลพีชํ ขนฺธพีชํ ผฬุพีชํ อคฺคพีชํ พีชพีชญฺเญว\csromansharedfootnote{b82f19efd045c64a}{พีชพีชเมว (ก)} ปญฺจมํ. ปญฺจหิ สมณกปฺเปหิ ผลํ ปริภุญฺชิ\-ตพฺพํ อคฺคิปริชิตํ สตฺถ\-ปริชิตํ นขปริชิตํ อพีชํ นิพฺพตฺต\-พีชญฺเญว\csromansharedfootnote{178386dc3513dcdd}{นิพฺพฏฺฏพีชญฺเญว (สี), นิพฺพฏพีชญฺเญว (สฺยา), นิปฺปฏฺฏพีชญฺเญว (ก)} ปญฺจมํ. ปญฺจ วิสุทฺธิโย, นิทานํ อุทฺทิสิตฺวา \csromanfolio{236}อวเสสํ สุเตน สาเวตพฺพํ, อยํ ปฐมา วิสุทฺธิ, นิทานํ อุทฺทิสิตฺวา จตฺตาริ ปาราชิกานิ อุทฺทิสิตฺวา อวเสสํ สุเตน สาเวตพฺพํ, อยํ ทุติยา วิสุทฺธิ, นิทานํ อุทฺทิสิตฺวา จตฺตาริ ปาราชิกานิ อุทฺทิสิตฺวา เตรส สํฆาทิเสเส อุทฺทิสิตฺวา อวเสสํ สุเตน สาเวตพฺพํ, อยํ ตติยา วิสุทฺธิ, นิทานํ อุทฺทิสิตฺวา จตฺตาริ ปาราชิกานิ อุทฺทิสิตฺวา เตรส สํฆาทิเสเส อุทฺทิสิตฺวา เทฺว อนิยเต อุทฺทิสิตฺวา อวเสสํ สุเตน สาเวตพฺพํ, อยํ จตุตฺถา วิสุทฺธิ, วิตฺถาเรเนว ปญฺจมี. อปราปิ ปญฺจ วิสุทฺธิโย สุตฺตุทฺเทโส ปาริสุทฺธิ\-อุโปสโถ อธิฏฺฐานุ\-โปสโถ ปวารณา สามคฺคี\-อุโปสโถเยว ปญฺจโม. ปญฺจานิสํสา วินยธเร, อตฺตโน สีลกฺขนฺโธ สุคุตฺโต โหติ สุรกฺขิโต, กุกฺกุจฺจ\-ปกตานํ ปฏิสรณํ โหติ, วิสารโท สํฆมชฺเฌ โวหรติ, ปจฺจตฺถิเก สหธมฺเมน สุนิคฺคหิตํ นิคฺคณฺหาติ, สทฺธมฺม\-ฏฺฐิติยา ปฏิปนฺโน โหติ. ปญฺจ อธมฺมิกานิ ปาติโมกฺข\-ฏฺฐปนานิ. ปญฺจ ธมฺมิกานิ ปาติโมกฺขฏฺฐปนานีติ.}{ปญฺจกํ นิฏฺฐิตํ.}
\csromancenter{ตสฺสุทฺทานํ.}
\setlength{\csromangathapagewidth}{0pt}
\setlength{\csromangathawakwidth}{0pt}
\csromangathameasuregroup{อาปตฺติ อาปตฺติกฺขนฺธา, \\ ปุคฺคลา เฉทนา เจว, \\ น อุเปติ อุเปติ จ, \\ วสํ พฺยสนา สมฺปทา, \\ โสสานิกํ ขายิตญฺจ, \\ อวิสฺสชฺชิ อเวภงฺคิ, \\ เทสนา สํฆํ อุทฺเทสํ, \\ กมฺมานิ ยาวตติยํ, \\ อกปฺปิยํ กปฺปิยญฺจ, \\ สมฺมชฺชนี อปเร จ, \\ อธิกรณํ วตฺถุ ญตฺติ, \\ ลหุกฏฺฐมกา เอเต, \\ อรญฺญํ ปิณฺฑปาตญฺจ, \\ อพฺโภกาโส จีวรญฺจ, \\ สนฺถติ ขลุ ปจฺฉาปิ, \\ อุโปสถํ ปวารณํ, \\ กณฺหสุกฺกปทา เอเต, \\ อปาสาทิกปาสาทิ, \\ กุลูปเก อติเวลํ, \\ วิสุทฺธิ อปเร เจว,}{\csromangathabat{อาปตฺติ อาปตฺติกฺขนฺธา,}{วินีตานนฺตเรน จ.} \\ \csromangathabat{ปุคฺคลา เฉทนา เจว,}{อาปชฺชติ จ ปจฺจยา.} \\ \csromangathabat{น อุเปติ อุเปติ จ,}{กปฺปนฺตุสงฺกิเตลญฺจ.} \\ \csromangathabat{วสํ พฺยสนา สมฺปทา,}{ปสฺสทฺธิ ปุคฺคเลน จ.} \\ \csromangathabat{โสสานิกํ ขายิตญฺจ,}{เถยฺยํ โจโร จ วุจฺจติ.} \\ \csromangathabat{อวิสฺสชฺชิ อเวภงฺคิ,}{กายโต กายวาจโต.} \\ \csromangathabat{เทสนา สํฆํ อุทฺเทสํ,}{ปจฺจนฺติกถิเนน จ.} \\ \csromangathabat{กมฺมานิ ยาวตติยํ,}{ปาราชิถุลฺลทุกฺกฏํ.} \\ \csromangathabat{อกปฺปิยํ กปฺปิยญฺจ,}{อปุญฺญทุวิโนทยา.} \\ \csromangathabat{สมฺมชฺชนี อปเร จ,}{ภาสํ อาปตฺติเมว จ.} \\ \csromangathabat{อธิกรณํ วตฺถุ ญตฺติ,}{อาปตฺตา อุภยานิ จ.} \\ \csromangathabat{ลหุกฏฺฐมกา เอเต,}{กณฺหสุกฺกา วิชานถ.} \\ \csromangathabat{อรญฺญํ ปิณฺฑปาตญฺจ,}{ปํสุรุกฺขสุสานิกา.} \\ \csromangathabat{อพฺโภกาโส จีวรญฺจ,}{สปทาโน นิสชฺชิโก.} \\ \csromangathabat{สนฺถติ ขลุ ปจฺฉาปิ,}{ปตฺตปิณฺฑิกเมว จ.} \\ \csromangathabat{อุโปสถํ ปวารณํ,}{อาปตฺตานาปตฺติปิ จ.} \\ \csromangathabat{กณฺหสุกฺกปทา เอเต,}{ภิกฺขุนีนมฺปิ เต ตถา.} \\ \csromangathabat{อปาสาทิกปาสาทิ,}{ตเถว อปเร ทุเว.} \\ \csromangathabat{กุลูปเก อติเวลํ,}{พีชํ สมณกปฺปิ จ.} \\ \csromangathabat{วิสุทฺธิ อปเร เจว,}{วินยาธมฺมิเกน จ.}}
\csromangathasetleft{อาปตฺติ อาปตฺติกฺขนฺธา, \\ ปุคฺคลา เฉทนา เจว, \\ น อุเปติ อุเปติ จ, \\ วสํ พฺยสนา สมฺปทา, \\ โสสานิกํ ขายิตญฺจ, \\ อวิสฺสชฺชิ อเวภงฺคิ, \\ เทสนา สํฆํ อุทฺเทสํ, \\ กมฺมานิ ยาวตติยํ, \\ อกปฺปิยํ กปฺปิยญฺจ, \\ สมฺมชฺชนี อปเร จ, \\ อธิกรณํ วตฺถุ ญตฺติ, \\ ลหุกฏฺฐมกา เอเต, \\ อรญฺญํ ปิณฺฑปาตญฺจ, \\ อพฺโภกาโส จีวรญฺจ, \\ สนฺถติ ขลุ ปจฺฉาปิ, \\ อุโปสถํ ปวารณํ, \\ กณฺหสุกฺกปทา เอเต, \\ อปาสาทิกปาสาทิ, \\ กุลูปเก อติเวลํ, \\ วิสุทฺธิ อปเร เจว,}
\csromangathagroup{\csromangathastanza{\csromangathabat{อาปตฺติ อาปตฺติกฺขนฺธา,}{วินีตานนฺตเรน จ.} \\ \csromangathabat{ปุคฺคลา เฉทนา เจว,}{อาปชฺชติ จ ปจฺจยา.}}\csromangathastanza{\csromangathabat{น อุเปติ อุเปติ จ,}{กปฺปนฺตุสงฺกิเตลญฺจ.} \\ \csromangathabat{วสํ พฺยสนา สมฺปทา,}{ปสฺสทฺธิ ปุคฺคเลน จ.}}\csromangathastanza{\csromangathabat{โสสานิกํ ขายิตญฺจ,}{เถยฺยํ โจโร จ วุจฺจติ.} \\ \csromangathabat{อวิสฺสชฺชิ อเวภงฺคิ,}{กายโต กายวาจโต.}}\csromangathastanza{\csromangathabat{เทสนา สํฆํ อุทฺเทสํ,}{ปจฺจนฺติกถิเนน จ.} \\ \csromangathabat{กมฺมานิ ยาวตติยํ,}{ปาราชิถุลฺลทุกฺกฏํ.}}\csromangathastanza{\csromangathabat{อกปฺปิยํ กปฺปิยญฺจ,}{อปุญฺญทุวิโนทยา.} \\ \csromangathabat{สมฺมชฺชนี อปเร จ,}{ภาสํ อาปตฺติเมว จ.}}\csromangathastanza{\csromangathabat{อธิกรณํ วตฺถุ ญตฺติ,}{อาปตฺตา อุภยานิ จ.} \\ \csromangathabat{ลหุกฏฺฐมกา เอเต,}{กณฺหสุกฺกา วิชานถ.}}\csromangathastanza{\csromangathabat{อรญฺญํ ปิณฺฑปาตญฺจ,}{ปํสุรุกฺขสุสานิกา.} \\ \csromangathabat{อพฺโภกาโส จีวรญฺจ,}{สปทาโน นิสชฺชิโก.}}\csromangathastanza{\csromanfolio{237}\csromangathabat{สนฺถติ ขลุ ปจฺฉาปิ,}{ปตฺตปิณฺฑิกเมว จ.} \\ \csromangathabat{อุโปสถํ ปวารณํ,}{อาปตฺตานาปตฺติปิ จ.}}\csromangathastanza{\csromangathabat{กณฺหสุกฺกปทา เอเต,}{ภิกฺขุนีนมฺปิ เต ตถา.} \\ \csromangathabat{อปาสาทิกปาสาทิ,}{ตเถว อปเร ทุเว.}}\csromangathastanza{\csromangathabat{กุลูปเก อติเวลํ,}{พีชํ สมณกปฺปิ จ.} \\ \csromangathabat{วิสุทฺธิ อปเร เจว,}{วินยาธมฺมิเกน จ.}}}

\vspace{\tipitakaparskip}
\csromancenter{ธมฺมิกา จ ตถา วุตฺตา, นิฏฺฐิตา สุทฺธิ\-ปญฺจกาติ.}
\csromansectionrule
\csromanmatikaanchor{mtk.169}
\csromantocmarkref{paragraph}{6. ฉกฺกวาร}{mtk.169}
\bookmark[dest={mtk.169},level=4]{6. ฉกฺกวาร}
\csromanheader{4}{6. \textbf{ฉกฺกวาร}}
\proseitem{326}{\csromansymbolfootnote{*}{ที 3. 202 ปิฏฺเฐปิ.}ฉ อคารวา. ฉ คารวา. ฉ วินีตวตฺถูนิ. ฉ สามีจิโย. ฉ อาปตฺติ\-สมุฏฺฐานา. ฉจฺเฉทนกา อาปตฺติโย. ฉหากาเรหิ อาปตฺติํ อาปชฺชติ. ฉานิสํสา วินยธเร. ฉ ปรมานิ. ฉารตฺตํ ติจีวเรน วิปฺปวสิตพฺพํ. ฉ จีวรานิ. ฉ รชนานิ. ฉ อาปตฺติโย กายโต จ จิตฺตโต จ สมุฏฺฐนฺติ, น วาจโต. ฉ อาปตฺติโย วาจโต จ จิตฺตโต จ สมุฏฺฐนฺติ, น กายโต. ฉ อาปตฺติโย กายโต จ วาจโต จ จิตฺตโต จ สมุฏฺฐนฺติ. ฉ กมฺมานิ. ฉ วิวาทมูลานิ. ฉ อนุวาทมูลานิ. ฉ สารณียา ธมฺมา. ทีฆโส ฉ วิทตฺถิโย สุคต\-วิทตฺถิยา. ติริยํ ฉ วิทตฺถิโย. ฉ นิสฺสย\-ปฏิปฺปสฺสทฺธิโย อาจริยมฺหา. ฉ นหาเน อนุปญฺญตฺติโย. วิปฺปกต\-จีวรํ อาทาย ปกฺกมติ. วิปฺปกต\-จีวรํ สมาทาย ปกฺกมติ.}

\prose{\csromansymbolfootnote{+}{วิ 3. 94 ปิฏฺฐาทีสุปิ.}ฉหงฺเคหิ สมนฺนาคเตน ภิกฺขุนา อุปสมฺ\-ปาเท\-ตพฺพํ, นิสฺสโย ทาตพฺโพ, สามเณโร อุปฏฺฐาเปตพฺโพ. อเสกฺเขน สีลกฺขนฺเธน สมนฺนาคโต โหติ, อเสกฺเขน สมาธิ\-กฺขนฺเธน สมนฺนาคโต โหติ, อเสกฺเขน ปญฺญ\-กฺขนฺเธน สมนฺนาคโต โหติ, อเสกฺเขน วิมุตฺติ\-กฺขนฺเธน สมนฺนาคโต โหติ, อเสกฺเขน วิมุตฺติ\-ญาณ\-ทสฺสนกฺขนฺเธน สมนฺนาคโต โหติ, ทสวสฺโส วา โหติ อติเรก\-ทสวสฺโส วา.}

\prose{อปเรหิปิ ฉหงฺเคหิ สมนฺนาคเตน ภิกฺขุนา อุปสมฺ\-ปาเท\-ตพฺพํ, นิสฺสโย ทาตพฺโพ, สามเณโร อุปฏฺฐาเปตพฺโพ. อตฺตนา อเสกฺเขน \csromanfolio{238}สีลกฺขนฺเธน สมนฺนาคโต โหติ ปรํ อเสกฺเข สีลกฺขนฺเธ สมาทเปตา, อตฺตนา อเสกฺเขน สมาธิ\-กฺขนฺเธน สมนฺนาคโต โหติ ปรํ อเสกฺเข สมาธิ\-กฺขนฺเธ สมาทเปตา, อตฺตนา อเสกฺเขน ปญฺญ\-กฺขนฺเธน สมนฺนาคโต โหติ ปรํ อเสกฺเข ปญฺญกฺขนฺเธ สมาทเปตา, อตฺตนา อเสกฺเขน วิมุตฺติ\-กฺขนฺเธน สมนฺนาคโต โหติ ปรํ อเสกฺเข วิมุตฺติ\-กฺขนฺเธ สมาทเปตา, อตฺตนา อเสกฺเขน วิมุตฺติ\-ญาณ\-ทสฺสนกฺขนฺเธน สมนฺนาคโต โหติ ปรํ อเสกฺเข วิมุตฺติ\-ญาณ\-ทสฺสนกฺขนฺเธ สมาทเปตา, ทสวสฺโส วา โหติ อติเรก\-ทสวสฺโส วา.}

\prose{อปเรหิปิ ฉหงฺเคหิ สมนฺนาคเตน ภิกฺขุนา อุปสมฺ\-ปาเท\-ตพฺพํ, นิสฺสโย ทาตพฺโพ, สามเณโร อุปฏฺฐาเปตพฺโพ. สทฺโธ โหติ, หิริมา โหติ, โอตฺตปฺปี โหติ, อารทฺธวีริโย โหติ, อุปฏฺฐิตสฺสติ โหติ, ทสวสฺโส วา โหติ อติเรก\-ทสวสฺโส วา.}

\prose{อปเรหิปิ ฉหงฺเคหิ สมนฺนาคเตน ภิกฺขุนา อุปสมฺ\-ปาเท\-ตพฺพํ, นิสฺสโย ทาตพฺโพ, สามเณโร อุปฏฺฐาเปตพฺโพ. น อธิสีเล สีลวิปนฺโน โหติ, น อชฺฌาจาเร อาจารวิปนฺโน โหติ, น อติทิฏฺฐิยา ทิฏฺฐิวิปนฺโน โหติ, พหุสฺสุโต โหติ, ปญฺญวา โหติ, ทสวสฺโส วา โหติ อติเรก\-ทสวสฺโส วา.}

\prose{อปเรหิปิ ฉหงฺเคหิ สมนฺนาคเตน ภิกฺขุนา อุปสมฺ\-ปาเท\-ตพฺพํ, นิสฺสโย ทาตพฺโพ, สามเณโร อุปฏฺฐาเปตพฺโพ. ปฏิพโล โหติ อนฺเตวาสิํ วา สทฺธิวิหาริํ วา คิลานํ อุปฏฺฐากุํ วา อุปฏฺฐาเปตุํ วา, อนภิรตํ วูปกาเสตุํ วา วูปกาสาเปตุํ วา\csromansharedfootnote{5fcbadbc5fec1926}{อนภิรติํ (ก)}, อุปฺปนฺนํ กุกฺกุจฺจํ ธมฺมโต วิโนเทตุํ\csromansharedfootnote{b1b1318b6b96f20d}{วิโนเทตุํ วา วิโนทาเปตุํ วา (สฺยา) วิ 3. 96 ปิฏฺเฐปิ.}, อาปตฺติํ ชานาติ, อาปตฺติ\-วุฏฺฐานํ ชานาติ, ทสวสฺโส วา โหติ อติเรก\-ทสวสฺโส วา.}

\prose{อปเรหิปิ ฉหงฺเคหิ สมนฺนาคเตน ภิกฺขุนา อุปสมฺ\-ปาเท\-ตพฺพํ, นิสฺสโย ทาตพฺโพ, สามเณโร อุปฏฺฐาเปตพฺโพ. ปฏิพโล โหติ อนฺเตวาสิํ วา สทฺธิวิหาริํ วา อาภิสมาจาริกาย สิกฺขาย สิกฺขาเปตุํ, อาทิพฺรหฺมจาริยกาย สิกฺขาย วิเนตุํ, อภิธมฺเม วิเนตุํ, อภิวินเย วิเนตุํ, อุปฺปนฺนํ ทิฏฺฐิคตํ ธมฺมโต วิเวเจตุํ, ทสวสฺโส วา โหติ อติเรก\-ทสวสฺโส วา.}

\prose{\csromanfolio{239}อปเรหิปิ ฉหงฺเคหิ สมนฺนาคเตน ภิกฺขุนา อุปสมฺ\-ปาเท\-ตพฺพํ, นิสฺสโย ทาตพฺโพ, สามเณโร อุปฏฺฐาเปตพฺโพ. อปตฺติํ ชานาติ, อนาปตฺติํ ชานาติ, ลหุกํ อาปตฺติํ ชานาติ, ครุกํ อาปตฺติํ ชานาติ, อุภยานิ โข ปนสฺส ปาติโมกฺขานิ วิตฺถาเรน สฺวาคตานิ โหนฺติ สุวิภตฺตานิ สุปฺปวตฺตีนิ สุวินิจฺฉิตานิ สุตฺตโส อนุพฺยญฺชนโส, ทสวสฺโส วา โหติ อติเรก\-ทสวสฺโส วา.}

\prosewithcloser{ฉ อธมฺมิกานิ ปาติโมกฺข\-ฏฺฐปนานิ. ฉ ธมฺมิกานิ ปาติโมกฺขฏฺฐปนานีติ.}{ฉกฺกํ นิฏฺฐิตํ.}
\csromancenter{ตสฺสุทฺทานํ.}
\setlength{\csromangathapagewidth}{0pt}
\setlength{\csromangathawakwidth}{0pt}
\csromangathameasuregroup{อคารวา คารวา จ, \\ สมุฏฺฐานา เฉทนา เจว, \\ ปรมานิ จ ฉารตฺตํ, \\ กายโต จิตฺตโต จาปิ, \\ กายวาจาจิตฺตโต จ, \\ อนุวาทา ทีฆโส จ, \\ อนุปญฺญตฺติ อาทาย, \\ อเสกฺเข สมาทเปตา,}{\csromangathabat{อคารวา คารวา จ,}{วินีตา สามีจิปิ จ.} \\ \csromangathabat{สมุฏฺฐานา เฉทนา เจว,}{อาการานิสํเสน จ.} \\ \csromangathabat{ปรมานิ จ ฉารตฺตํ,}{จีวรํ รชนานิ จ.} \\ \csromangathabat{กายโต จิตฺตโต จาปิ,}{วาจโต จิตฺตโตปิ จ.} \\ \csromangathabat{กายวาจาจิตฺตโต จ,}{กมฺมวิวาทเมว จ.} \\ \csromangathabat{อนุวาทา ทีฆโส จ,}{ติริยํ นิสฺสเยน จ.} \\ \csromangathabat{อนุปญฺญตฺติ อาทาย,}{สมาทาย ตเถว จ.} \\ \csromangathabat{อเสกฺเข สมาทเปตา,}{สทฺโธ อธิสีเลน จ.}}
\csromangathasetleft{อคารวา คารวา จ, \\ สมุฏฺฐานา เฉทนา เจว, \\ ปรมานิ จ ฉารตฺตํ, \\ กายโต จิตฺตโต จาปิ, \\ กายวาจาจิตฺตโต จ, \\ อนุวาทา ทีฆโส จ, \\ อนุปญฺญตฺติ อาทาย, \\ อเสกฺเข สมาทเปตา,}
\csromangathagroup{\csromangathastanza{\csromangathabat{อคารวา คารวา จ,}{วินีตา สามีจิปิ จ.} \\ \csromangathabat{สมุฏฺฐานา เฉทนา เจว,}{อาการานิสํเสน จ.}}\csromangathastanza{\csromangathabat{ปรมานิ จ ฉารตฺตํ,}{จีวรํ รชนานิ จ.} \\ \csromangathabat{กายโต จิตฺตโต จาปิ,}{วาจโต จิตฺตโตปิ จ.}}\csromangathastanza{\csromangathabat{กายวาจาจิตฺตโต จ,}{กมฺมวิวาทเมว จ.} \\ \csromangathabat{อนุวาทา ทีฆโส จ,}{ติริยํ นิสฺสเยน จ.}}\csromangathastanza{\csromangathabat{อนุปญฺญตฺติ อาทาย,}{สมาทาย ตเถว จ.} \\ \csromangathabat{อเสกฺเข สมาทเปตา,}{สทฺโธ อธิสีเลน จ.}}}

\vspace{\tipitakaparskip}
\csromancenter{คิลานาภิ\-สมา\-จารี, อาปตฺตา\-ธมฺม\-ธมฺมิกาติ.}
\csromansectionrule
\csromanmatikaanchor{mtk.170}
\csromantocmarkref{paragraph}{7. สตฺตกวาร}{mtk.170}
\bookmark[dest={mtk.170},level=4]{7. สตฺตกวาร}
\csromanheader{4}{7. \textbf{สตฺตกวาร}}
\proseitem{327}{สตฺตาปตฺติโย. สตฺตา\-ปตฺติกฺขนฺธา. สตฺต วินีตวตฺถูนิ. สตฺต สามีจิโย. สตฺต อธมฺมิกา ปฏิญฺญาต\-กรณา. สตฺต ธมฺมิกา ปฏิญฺญาต\-กรณา. สตฺตนฺนํ อนาปตฺติ สตฺตาห\-กรณีเยน คนฺตุํ. สตฺตานิสํสา วินยธเร. สตฺต ปรมานิ. สตฺตเม อรุณุคฺคมเน นิสฺสคฺคิยํ โหติ. สตฺต สมถา. สตฺต กมฺมานิ. สตฺต อามกธญฺญานิ. ติริยํ สตฺตนฺตรา. \csromanfolio{240}คณโภชเน สตฺต อนุปญฺญตฺติโย. เภสชฺชานิ ปฏิคฺคเหตฺวา สตฺตาหปรมํ สนฺนิธิการกํ ปริภุญฺชิ\-ตพฺพานิ. กตจีวรํ อาทาย ปกฺกมติ. กตจีวรํ สมาทาย ปกฺกมติ. ภิกฺขุสฺส น โหติ อาปตฺติ ทฏฺฐพฺพา. ภิกฺขุสฺส โหติ อาปตฺติ ทฏฺฐพฺพา. ภิกฺขุสฺส โหติ อาปตฺติ ทฏฺฐพฺพา\csromansharedfootnote{c86806b0803b5d80}{ปฏิกาตพฺพา (สพฺพตฺต) อฏฺฐกถา จ จมฺเปยฺยกฺขนฺธเก อธมฺมกมฺมาทิกถา จ โอโลเกตพฺพา.}. สตฺต อธมฺมิกานิ ปาติโมกฺข\-ฏฺฐปนานิ. สตฺต ธมฺมิกานิ ปาติโมกฺข\-ฏฺฐปนานิ.}

\prose{\csromansymbolfootnote{*}{อํ 2. 504 ปิฏฺฐาทีสุ จ.}สตฺตหงฺเคหิ สมนฺนาคโต ภิกฺขุ วินยธโร โหติ. อาปตฺติํ ชานาติ, อนาปตฺติํ ชานาติ, ลหุกํ อาปตฺติํ ชานาติ, ครุกํ อาปตฺติํ ชานาติ, สีลวา โหติ ปาติโมกฺข\-สํวร\-สํวุโต วิหรติ อาจารโคจร\-สมฺปนฺโน อณุมตฺเตสุ วชฺเชสุ ภยทสฺสาวี สมาทาย สิกฺขติ สิกฺขาปเทสุ, จตุนฺนํ ฌานานํ อาภิเจตสิกานํ ทิฏฺฐ\-ธมฺม\-สุข\-วิหารานํ นิกามลาภี โหติ อกิจฺฉลาภี อกสิรลาภี, อาสวานญฺจ ขยา อนาสวํ เจโตวิมุตฺติํ ปญฺญาวิมุตฺติํ ทิฏฺเฐว ธมฺเม สยํ อภิญฺญา สจฺฉิกตฺวา อุปสมฺปชฺช วิหรติ.}

\prose{อปเรหิปิ สตฺตหงฺเคหิ สมนฺนาคโต ภิกฺขุ วินยธโร โหติ. อาปตฺติํ ชานาติ, อนาปตฺติํ ชานาติ, ลหุกํ อาปตฺติํ ชานาติ, ครุกํ อาปตฺติํ ชานาติ, พหุสฺสุโต โหติ สุตธโร สุตสนฺนิจโย, เย เต ธมฺมา อาทิกลฺยาณา มชฺเฌกลฺยาณา ปริโยสาน\-กลฺยาณา สาตฺถํ สพฺยญฺชนํ เกวล\-ปริปุณฺณํ ปริสุทฺธํ พฺรหฺมจริยํ อภิวทนฺติ, ตถารูปสฺส ธมฺมา พหุสฺสุตา โหนฺติ ธาตา\csromansharedfootnote{1c25e3e705596ff0}{ธตา (สี, สฺยา)} วจสา ปริจิตา มนสา\-นุเปกฺขิตา ทิฏฺฐิยา สุปฺปฏิวิทฺธา, จตุนฺนํ ฌานานํ อาภิเจตสิกานํ ทิฏฺฐ\-ธมฺม\-สุข\-วิหารานํ นิกามลาภี โหติ อกิจฺฉลาภี อกสิรลาภี, อาสวานญฺจ ขยา อนาสวํ เจโตวิมุตฺติํ ปญฺญาวิมุตฺติํ ทิฏฺเฐว ธมฺเม สยํ อภิญฺญา สจฺฉิกตฺวา อุปสมฺปชฺช วิหรติ.}

\prose{อปเรหิปิ สตหงฺเคหิ สมนฺนาคโต ภิกฺขุ วินยธโร โหติ. อาปตฺติํ ชานาติ, อนาปตฺติํ ชานาติ, ลหุกํ อาปตฺติํ ชานาติ, ครุกํ อาปตฺติํ ชานาติ, อุภยานิ โข ปนสฺส ปาติโมกฺขานิ วิตฺถาเรน สฺวาคตานิ โหนฺติ สุวิภตฺตานิ สุปฺปวตฺตีนิ สุวินิจฺฉิตานิ สุตฺตโส อนุพฺยญฺชนโส, จตุนฺนํ ฌานานํ อาภิเจตสิกานํ ทิฏฺฐิธมฺมสุขวิหารานํ นิกามลาภี โหติ อกิจฺฉลาภี อกสิรลาภี, \csromanfolio{241}อาสวานญฺจ ขยา อนาสวํ เจโตวิมุตฺติํ ปญฺญาวิมุตฺติํ ทิฏฺเฐว ธมฺเม สยํ อภิญฺญา สจฺฉิกตฺวา อุปสมฺปชฺช วิหรติ.}

\prose{อปเรหิปิ สตฺตหงฺเคหิ สมนฺนาคโต ภิกฺขุ วินยธโร โหติ. อาปตฺติํ ชานาติ, อนาปตฺติํ ชานาติ, ลหุกํ อาปตฺติํ ชานาติ, ครุกํ อาปตฺติํ ชานาติ, อเนกวิหิตํ ปุพฺเพนิวาสํ อนุสฺสรติ, เสยฺยถิทํ, เอกมฺปิ ชาติํ เทฺวปิ ชาติโย ติสฺโสปิ ชาติโย จตสฺโสปิ ชาติโย ปญฺจปิ ชาติโย ทสปิ ชาติโย วีสมฺปิ ชาติโย ติํสมฺปิ ชาติโย จตฺตาลีสมฺปิ ชาติโย ปญฺญาสมฺปิ ชาติโย ชาติสตมฺปิ ชาติสหสฺสมฺปิ ชาติ\-สต\-สหสฺสมฺปิ อเนเกปิ สํวฏฺฏกปฺเป อเนเกปิ วิวฏฺฏกปฺเป อเนเกปิ สํวฏฺฏ\-วิวฏฺฏ\-กปฺเป “อมุตฺราสิํ เอวํนาโม เอวํโคตฺโต เอวํวณฺโณ เอวมาหาโร เอวํ\-สุขทุกฺข\-ปฺปฏิสํเวที เอวมา\-ยุปริยนฺโต, โส ตโต จุโต อมุตฺร อุทปาทิํ, ตตฺราปาสิํ เอวํนาโม เอวํโคตฺโต เอวํวณฺโณ เอวมาหาโร เอวํ\-สุขทุกฺข\-ปฺปฏิสํเวที เอวมา\-ยุปริยนฺโต, โส ตโต จุโต อิธูปปนฺโน”ติ, อิติ สาการํ สอุทฺเทสํ อเนกวิหิตํ ปุพฺเพนิวาสํ อนุสฺสรติ, ทิพฺเพน จกฺขุนา วิสุทฺเธน อติกฺกนฺต\-มานุสเกน\csromansharedfootnote{9698df013f0ee02e}{อติกฺกนฺตมานุสฺสเกน (ก)} สตฺเต ปสฺสติ จวมาเน อุปปชฺชมาเน หีเน ปณีเต สุวณฺเณ ทุพฺพณฺเณ สุคเต ทุคฺคเต, ยถากมฺมูปเค สตฺเต ปชานาติ, “อิเม วต โภนฺโต สตฺตา กาย\-ทุจฺจริเตน สมนฺนาคตา, วจีทุจฺจริเตน สมนฺนาคตา, มโน\-ทุจฺจริเตน สมนฺนาคตา, อริยานํ อุปวาทกา, มิจฺฉาทิฏฺฐิกา, มิจฺฉาทิฏฺฐิ\-กมฺม\-สมาทานา, เต กายสฺส เภทา ปรํ มรณา อปายํ ทุคฺคติํ วินิปาตํ นิรยํ อุปปนฺนา. อิเม วา ปน โภนฺโต สตฺตา กายสุจริเตน สมนฺนาคตา, วจีสุจริเตน สมนฺนาคตา, มโนสุจริเตน สมนฺนาคตา, อริยานํ อนุปวาทกา, สมฺมาทิฏฺฐิกา, สมฺมา\-ทิฏฺฐิ\-กมฺม\-สมาทานา, เต กายสฺส เภทา ปรํ มรณา สุคติํ สคฺคํ โลกํ อุปปนฺนา”ติ. อิติ ทิพฺเพน จกฺขุนา วิสุทฺเธน อติกฺกนฺต\-มานุสเกน สตฺเต ปสฺสติ จวมาเน อุปปชฺชมาเน หีเน ปณีเต สุวณฺเณ ทุพฺพณฺเณ สุคเต ทุคฺคเต, ยถากมฺมูปเค สตฺเต ปชานาติ, อาสวานญฺจ ขยา อนาสวํ เจโตวิมุตฺติํ ปญฺญาวิมุตฺติํ ทิฏฺเฐว ธมฺเม สยํ อภิญฺญา สจฺฉิกตฺวา อุปสมฺปชฺช วิหรติ.}

\prose{\csromanfolio{242}สตฺตหงฺเคหิ สมนฺนาคโต วินยธโร โสภติ. อาปตฺติํ ชานาติ, อนาปตฺติํ ชานาติ, ลหุกํ อาปตฺติํ ชานาติ, ครุกํ อาปตฺติํ ชานาติ, สีลวา โหติ ฯเปฯ สมาทาย สิกฺขติ สิกฺขาปเทสุ, จตุนฺนํ ฌานานํ อาภิเจตสิกานํ ทิฏฺฐ\-ธมฺม\-สุข\-วิหารานํ นิกามลาภี โหติ อกิจฺฉลาภี อกสิรลาภี, อาสวานญฺจ ขยา อนาสวํ เจโตวิมุตฺติํ ปญฺญาวิมุตฺติํ ทิฏฺเฐว ธมฺเม สยํ อภิญฺญา สจฺฉิกตฺวา อุปสมฺปชฺช วิหรติ.}

\prose{อปเรหิปิ สตฺตหงฺเคหิ สมนฺนาคโต วินยธโร โสภติ. อาปตฺติํ ชานาติ, อนาปตฺติํ ชานาติ, ลหุกํ อาปตฺติํ ชานาติ, ครุกํ อาปตฺติํ ชานาติ, พหุสฺสุโต โหติ ฯเปฯ ทิฏฺฐิยา สุปฺปฏิวิทฺโธ, จตุนฺนํ ฌานานํ อาภิเจตสิกานํ ทิฏฺฐ\-ธมฺม\-สุข\-วิหารานํ นิกามลาภี โหติ อกิจฺฉลาภี อกสิรลาภี, อาสวานญฺจ ขยา อนาสวํ เจโตวิมุตฺติํ ปญฺญาวิมุตฺติํ ทิฏฺเฐว ธมฺเม สยํ อภิญฺญา สจฺฉิกตฺวา อุปสมฺปชฺช วิหรติ.}

\prose{อปเรหิปิ สตฺตหงฺเคหิ สมนฺนาคโต วินยธโร โสภติ. อาปตฺติํ ชานาติ, อนาปตฺติํ ชานาติ, ลหุกํ อาปตฺติํ ชานาติ, ครุกํ อาปตฺติํ ชานาติ, อุภยานิ โข ปนสฺส ปาติโมกฺขานิ วิตฺถาเรน สฺวาคตานิ โหนฺติ สุวิภตฺตานิ สุปฺปวตฺตีนิ สุวินิจฺฉิตานิ สุตฺตโส อนุพฺยญฺชนโส, จตุนฺนํ ฌานานํ อาภิเจตสิกานํ ทิฏฺฐ\-ธมฺม\-สุข\-วิหารานํ นิกามลาภี โหติ อกิจฺฉลาภี อกสิรลาภี, อาสวานญฺจ ขยา อนาสวํ เจโตวิมุตฺติํ ปญฺญาวิมุตฺติํ ทิฏฺเฐว ธมฺเม สยํ อภิญฺญา สจฺฉิกตฺวา อุปสมฺปชฺช วิหรติ.}

\prose{อปเรหิปิ สตฺตหงฺเคหิ สมนฺนาคโต วินยธโร โสภติ. อาปตฺติํ ชานาติ, อนาปตฺติํ ชานาติ, ลหุกํ อาปตฺติํ ชานาติ, ครุกํ อาปตฺติํ ชานาติ, อเนกวิหิตํ ปุพฺเพนิวาสํ อนุสฺสรติ, เสยฺยถิทํ, เอกมฺปิ ชาติํ เทฺวปิ ชาติโย ฯเปฯ อิติ สาการํ สอุทฺเทสํ อเนกวิหิตํ ปุพฺเพนิวาสํ อนุสฺสรติ, ทิพฺเพน จกฺขุนา วิสุทฺเธน อติกฺกนฺต\-มานุสเกน สตฺเต ปสฺสติ จวมาเน อุปปชฺชมาเน หีเน ปณีเต สุวณฺเณ ทุพฺพณฺเณ สุคเต ทุคฺคเต, ยถากมฺมูปเค สตฺเต ปชานาติ ฯเปฯ อิติ ทิพฺเพน จกฺขุนา วิสุทฺเธน อติกฺกนฺต\-มานุสเกน สตฺเต ปสฺสติ จวมาเน อุปปชฺชมาเน หีเน ปณีเต สุวณฺเณ ทุพฺพณฺเณ สุคเต ทุคฺคเต, \csromanfolio{243}ยถากมฺมูปเค สตฺเต ปชานาติ, อาสวานญฺจ ขยา อนาสวํ เจโตวิมุตฺติํ ปญฺญาวิมุตฺติํ ทิฏฺเฐว ธมฺเม สยํ อภิญฺญา สจฺฉิกตฺวา อุปสมฺปชฺช วิหรติ.}

\prose{\csromansymbolfootnote{*}{ที 3. 208 ปิฏฺเฐปิ.}สตฺต อสทฺธมฺมา. อสฺสทฺโธ โหติ, อหิริโก โหติ, อโนตฺตปฺปี โหติ, อปฺปสฺสุโต โหติ, กุสีโต โหติ, มุฏฺฐสฺสติ โหติ, ทุปฺปญฺโญ โหติ.}

\prosewithcloser{\csromansymbolmark{*}สตฺต สทฺธมฺมา. สทฺโธ โหติ, หิริมา โหติ, โอตฺตปฺปี โหติ, พหุสฺสุโต โหติ, อารทฺธวีริโย โหติ, อุปฏฺฐิตสฺสติ โหติ, ปญฺญวา โหตีติ.}{สตฺตกํ นิฏฺฐิตํ.}
\csromancenter{ตสฺสุทฺทานํ.}
\setlength{\csromangathapagewidth}{0pt}
\setlength{\csromangathawakwidth}{0pt}
\csromangathameasuregroup{อาปตฺติ อาปตฺติกฺขนฺธา, \\ อธมฺมิกา ธมฺมิกา จ, \\ อานิสํสา ปรมานิ, \\ กมฺมา อามกธญฺญา จ, \\ สตฺตาหปรมํ อาทาย, \\ น โหติ โหติ โหติ จ, \\ จตุโร วินยธรา,}{\csromangathabat{อาปตฺติ อาปตฺติกฺขนฺธา,}{วินีตา สามีจิปิ จ.} \\ \csromangathabat{อธมฺมิกา ธมฺมิกา จ,}{อนาปตฺติ จ สตฺตาหํ.} \\ \csromangathabat{อานิสํสา ปรมานิ,}{อรุณสมเถน จ.} \\ \csromangathabat{กมฺมา อามกธญฺญา จ,}{ติริยํ คณโภชเน.} \\ \csromangathabat{สตฺตาหปรมํ อาทาย,}{สมาทาย ตเถว จ.} \\ \csromangathabat{น โหติ โหติ โหติ จ,}{อธมฺมา ธมฺมิกานิ จ.} \\ \csromangathabat{จตุโร วินยธรา,}{จตุภิกฺขู จ โสภเน.}}
\csromangathasetleft{อาปตฺติ อาปตฺติกฺขนฺธา, \\ อธมฺมิกา ธมฺมิกา จ, \\ อานิสํสา ปรมานิ, \\ กมฺมา อามกธญฺญา จ, \\ สตฺตาหปรมํ อาทาย, \\ น โหติ โหติ โหติ จ, \\ จตุโร วินยธรา,}
\csromangathagroup{\csromangathastanza{\csromangathabat{อาปตฺติ อาปตฺติกฺขนฺธา,}{วินีตา สามีจิปิ จ.} \\ \csromangathabat{อธมฺมิกา ธมฺมิกา จ,}{อนาปตฺติ จ สตฺตาหํ.}}\csromangathastanza{\csromangathabat{อานิสํสา ปรมานิ,}{อรุณสมเถน จ.} \\ \csromangathabat{กมฺมา อามกธญฺญา จ,}{ติริยํ คณโภชเน.}}\csromangathastanza{\csromangathabat{สตฺตาหปรมํ อาทาย,}{สมาทาย ตเถว จ.} \\ \csromangathabat{น โหติ โหติ โหติ จ,}{อธมฺมา ธมฺมิกานิ จ.} \\ \csromangathabat{จตุโร วินยธรา,}{จตุภิกฺขู จ โสภเน.}}}

\vspace{\tipitakaparskip}
\csromancenter{สตฺต เจว อสทฺธมฺมา, สตฺต สทฺธมฺมา เทสิตาติ.}
\csromansectionrule
\csromanmatikaanchor{mtk.171}
\csromantocmarkref{paragraph}{8. อฏฺฐกวาร}{mtk.171}
\bookmark[dest={mtk.171},level=4]{8. อฏฺฐกวาร}
\csromanheader{4}{8. \textbf{อฏฺฐกวาร}}
\proseitemwithcloser{328}{อฏฺฐานิสํเส สมฺปสฺสมาเนน น โส ภิกฺขุ อาปตฺติยา อทสฺสเน อุกฺขิปิตพฺโพ. อฏฺฐานิสํเส สมฺปสฺสมาเนน ปเรสมฺปิ สทฺธาย สา อาปตฺติ เทเสตพฺพา. อฏฺฐ ยาวตติยกา. อฏฺฐหากาเรหิ กุลานิ ทูเสติ. อฏฺฐ มาติกา จีวรสฺส อุปฺปาทาย. อฏฺฐ มาติกา กถินสฺส อุพฺภาราย. อฏฺฐ ปานานิ.\csromansymbolfootnote{+}{วิ 4. 365; อํ 3. 10 ปิฏฺเฐสุปิ.}อฏฺฐหิ อสทฺธมฺเมหิ อภิภูโต ปริยาทินฺน\-จิตฺโต เทวทตฺโต อาปายิโก เนรยิโก กปฺปฏฺโฐ \csromanfolio{244}อเตกิจฺโฉ. อฏฺฐ โลกธมฺมา. อฏฺฐ ครุธมฺมา. อฏฺฐ ปาฏิเทสนียา. อฏฺฐงฺคิโก มุสาวาโท. อฏฺฐ อุโปสถงฺคานิ. อฏฺฐ ทูเตยฺยงฺคานิ. อฏฺฐ ติตฺถิย\-วตฺตานิ. อฏฺฐ อจฺฉริยา อพฺภุตธมฺมา มหาสมุทฺเท. อฏฺฐ อจฺฉริยา อพฺภุตธมฺมา อิมสฺมิํ ธมฺมวินเย. อฏฺฐ อนติริตฺตา. อฏฺฐ อติริตฺตา. อฏฺฐเม อรุณุคฺคมเน นิสฺสคฺคิยํ โหติ. อฏฺฐ ปาราชิกา. อฏฺฐมํ วตฺถุํ ปริปูเรนฺตี นาเสตพฺพา. อฏฺฐมํ วตฺถุํ ปริปูเรนฺติยา เทสิตาปิ อเทสิตา โหติ. อฏฺฐวาจิกา อุปสมฺปทา. อฏฺฐนฺนํ ปจฺจุฏฺฐาตพฺพํ. อฏฺฐนฺนํ อาสนํ ทาตพฺพํ. อุปาสิกา อฏฺฐ วรานิ ยาจติ. อฏฺฐหงฺเคหิ สมนฺนาคโต ภิกฺขุ ภิกฺขุโน\-วาทโก สมฺมนฺนิตพฺโพ\csromansharedfootnote{7f044bd59346eed7}{สมฺมนิตพฺโพ (ก)}. อฏฺฐานิสํสา วินยธเร. อฏฺฐ ปรมานิ. ตสฺส ปาปิยสิกกมฺมกเตน ภิกฺขุนา อฏฺฐสุ ธมฺเมสุ สมฺมา วตฺติตพฺพํ. อฏฺฐ อธมฺมิกานิ ปาติโมกฺข\-ฏฺฐปนานิ. อฏฺฐ ธมฺมิกานิ ปาติโมกฺขฏฺฐปนานีติ.}{อฏฺฐกํ นิฏฺฐิตํ.}
\csromancenter{ตสฺสุทฺทานํ.}
\setlength{\csromangathapagewidth}{0pt}
\setlength{\csromangathawakwidth}{0pt}
\csromangathameasuregroup{น โส ภิกฺขุ ปเรสํปิ,}{\csromangathabat{น โส ภิกฺขุ ปเรสํปิ,}{ยาวตติยทูสนา.}}
\csromangathasetleft{น โส ภิกฺขุ ปเรสํปิ,}
\csromangathagroup{\csromangathastanza{\csromangathabat{น โส ภิกฺขุ ปเรสํปิ,}{ยาวตติยทูสนา.}}}

\csromanheader{2}{มาติกา กถินุพฺภารา, ปานา อภิภูเตน จ.}
\setlength{\csromangathapagewidth}{0pt}
\setlength{\csromangathawakwidth}{0pt}
\csromangathameasuregroup{โลกธมฺมา ครุธมฺมา, \\ อุโปสถา จ ทูตงฺคา, \\ อพฺภุตา อนติริตฺตํ, \\ ปาราชิกฏฺฐมํ วตฺถุ, \\ ปจฺจุฏฺฐานาสนญฺเจว, \\ อานิสํสา ปรมานิ,}{\csromangathabat{โลกธมฺมา ครุธมฺมา,}{ปาฏิเทสนียา มุสา.} \\ \csromangathabat{อุโปสถา จ ทูตงฺคา,}{ติตฺถิยา สมุทฺเทปิ จ.} \\ \csromangathabat{อพฺภุตา อนติริตฺตํ,}{อติริตฺตํ นิสฺสคฺคิยํ.} \\ \csromangathabat{ปาราชิกฏฺฐมํ วตฺถุ,}{อเทสิตูปสมฺปทา.} \\ \csromangathabat{ปจฺจุฏฺฐานาสนญฺเจว,}{วรํ โอวาทเกน จ.} \\ \csromangathabat{อานิสํสา ปรมานิ,}{อฏฺฐธมฺเมสุ วตฺตนา.}}
\csromangathasetleft{โลกธมฺมา ครุธมฺมา, \\ อุโปสถา จ ทูตงฺคา, \\ อพฺภุตา อนติริตฺตํ, \\ ปาราชิกฏฺฐมํ วตฺถุ, \\ ปจฺจุฏฺฐานาสนญฺเจว, \\ อานิสํสา ปรมานิ,}
\csromangathagroup{\csromangathastanza{\csromangathabat{โลกธมฺมา ครุธมฺมา,}{ปาฏิเทสนียา มุสา.} \\ \csromangathabat{อุโปสถา จ ทูตงฺคา,}{ติตฺถิยา สมุทฺเทปิ จ.}}\csromangathastanza{\csromangathabat{อพฺภุตา อนติริตฺตํ,}{อติริตฺตํ นิสฺสคฺคิยํ.} \\ \csromangathabat{ปาราชิกฏฺฐมํ วตฺถุ,}{อเทสิตูปสมฺปทา.}}\csromangathastanza{\csromangathabat{ปจฺจุฏฺฐานาสนญฺเจว,}{วรํ โอวาทเกน จ.} \\ \csromangathabat{อานิสํสา ปรมานิ,}{อฏฺฐธมฺเมสุ วตฺตนา.}}}

\vspace{\tipitakaparskip}
\csromancenter{อธมฺมิกา ธมฺมิกา จ, อฏฺฐกา สุปฺปกาสิตาติ.}
\csromansectionrule
\csromanmatikaanchor{mtk.172}
\csromantocmarkref{paragraph}{9. นวกวาร}{mtk.172}
\bookmark[dest={mtk.172},level=4]{9. นวกวาร}
\csromanheader{4}{9. \textbf{นวกวาร}}
\proseitemwithcloser{329}{นว อาฆาตวตฺถูนิ. นว อาฆาต\-ปฏิวินยา. นว วินีตวตฺถูนิ. นว ปฐมาปตฺติกา. นวหิ สํโฆ ภิชฺชติ. นว ปณีต\-โภชนานิ. \csromanfolio{245}นวมํเสหิ ทุกฺกฏํ. นว ปาติโมกฺขุ\-ทฺเทสา. นว ปรมานิ. นว ตณฺหามูลกา ธมฺมา. นว วิธมานา. นว จีวรานิ อธิฏฺฐาตพฺพานิ. นว จีวรานิ น วิกปฺเปตพฺพานิ. ทีฆโส นว วิทตฺถิโย สุคต\-วิทตฺถิยา. นว อธมฺมิกานิ ทานานิ. นว อธมฺมิกา ปฏิคฺคหา. นว อธมฺมิกา ปริโภคา. ตีณิ ธมฺมิกานิ ทานานิ, ตโย ธมฺมิกา ปฏิคฺคหา, ตโย ธมฺมิกา ปริโภคา. นว อธมฺมิกา สญฺญตฺติโย\csromansharedfootnote{e1034f407605cd6b}{ปญฺญตฺติโย (สพฺพตฺถ) อฏฺฐกถา จ สมถกฺขนฺธเก กณฺหปกฺขนวกญฺจ โอโลเกตพฺพํ.}. นว ธมฺมิกา สญฺญตฺติโย. อธมฺมกมฺเม เทฺว นวกานิ. ธมฺมกมฺเม เทฺว นวกานิ. นว อธมฺมิกานิ ปาติโมกฺข\-ฏฺฐปนานิ. นว ธมฺมิกานิ ปาติโมกฺขฏฺฐปนานีติ.}{นวกํ นิฏฺฐิตํ.}
\csromancenter{ตสฺสุทฺทานํ.}
\setlength{\csromangathapagewidth}{0pt}
\setlength{\csromangathawakwidth}{0pt}
\csromangathameasuregroup{อาฆาตวตฺถุวินยา, \\ ภิชฺชติ จ ปณีตญฺจ, \\ ตณฺหา มานา อธิฏฺฐานา, \\ ทานา ปฏิคฺคหา โภคา, \\ อธมฺมธมฺมสญฺญตฺติ,}{\csromangathabat{อาฆาตวตฺถุวินยา,}{วินีตา ปฐเมน จ.} \\ \csromangathabat{ภิชฺชติ จ ปณีตญฺจ,}{มํสุทฺเทสปรมานิ จ.} \\ \csromangathabat{ตณฺหา มานา อธิฏฺฐานา,}{วิกปฺเป จ วิทตฺถิโย.} \\ \csromangathabat{ทานา ปฏิคฺคหา โภคา,}{ติวิธา ปุน ธมฺมิกา.} \\ \csromangathabat{อธมฺมธมฺมสญฺญตฺติ,}{ทุเว เทฺว นวกานิ จ.}}
\csromangathasetleft{อาฆาตวตฺถุวินยา, \\ ภิชฺชติ จ ปณีตญฺจ, \\ ตณฺหา มานา อธิฏฺฐานา, \\ ทานา ปฏิคฺคหา โภคา, \\ อธมฺมธมฺมสญฺญตฺติ,}
\csromangathagroup{\csromangathastanza{\csromangathabat{อาฆาตวตฺถุวินยา,}{วินีตา ปฐเมน จ.} \\ \csromangathabat{ภิชฺชติ จ ปณีตญฺจ,}{มํสุทฺเทสปรมานิ จ.}}\csromangathastanza{\csromangathabat{ตณฺหา มานา อธิฏฺฐานา,}{วิกปฺเป จ วิทตฺถิโย.} \\ \csromangathabat{ทานา ปฏิคฺคหา โภคา,}{ติวิธา ปุน ธมฺมิกา.} \\ \csromangathabat{อธมฺมธมฺมสญฺญตฺติ,}{ทุเว เทฺว นวกานิ จ.}}}

\vspace{\tipitakaparskip}
\csromancenter{ปาติโมกฺข\-ฏฺฐปนานิ, อธมฺม\-ธมฺมิกานิ จาติ.}
\csromansectionrule
\csromanmatikaanchor{mtk.173}
\csromantocmarkref{paragraph}{10. ทสกวาร}{mtk.173}
\bookmark[dest={mtk.173},level=4]{10. ทสกวาร}
\csromanheader{4}{10. \textbf{ทสกวาร}}
\proseitem{330}{ทส อาฆาตวตฺถูนิ. ทส อาฆาต\-ปฏิวินยา. ทส วินีตวตฺถูนิ. ทสวตฺถุกา มิจฺฉาทิฏฺฐิ. ทสวตฺถุกา สมฺมาทิฏฺฐิ. ทส อนฺตคฺคาหิกา ทิฏฺฐิ. ทส มิจฺฉตฺตา. ทส สมฺมตฺตา. ทส อกุสล\-กมฺมปถา. ทส กุสล\-กมฺม\-ปถา. ทส อธมฺมิกา สลากคฺคาหา. ทส ธมฺมิกา สลากคฺคาหา. สามเณรานํ ทส สิกฺขาปทานิ. ทสหงฺเคหิ สมนฺนาคโต สามเณโร นาเสตพฺโพ.}

\prose{ทสหงฺเคหิ สมนฺนาคโต วินยธโร “พาโล”เตฺวว สงฺขํ คจฺฉติ. อตฺตโน ภาส\-ปริยนฺตํ น อุคฺคณฺหาติ, ปรสฺส ภาส\-ปริยนฺตํ น \csromanfolio{246}อุคฺคณฺหาติ, อตฺตโน ภาส\-ปริยนฺตํ อนุคฺคเหตฺวา, ปรสฺส ภาส\-ปริยนฺตํ อนุคฺคเหตฺวา, อธมฺเมน กาเรติ อปฺปฏิญฺญาย, อาปตฺติํ น ชานาติ, อาปตฺติยา มูลํ น ชานาติ, อาปตฺติ\-สมุทยํ น ชานาติ, อาปตฺตินิโรธํ น ชานาติ, อาปตฺติ\-นิโรธคามินิํ ปฏิปทํ น ชานาติ.}

\prose{ทสหงฺเคหิ สมนฺนาคโต วินยธโร “ปณฺฑิโต”เตฺวว สงฺขํ คจฺฉติ. อตฺตโน ภาส\-ปริยนฺตํ อุคฺคณฺหาติ, ปรสฺส ภาส\-ปริยนฺตํ อุคฺคณฺหาติ, อตฺตโน ภาส\-ปริยนฺตํ อุคฺคเหตฺวา, ปรสฺส ภาส\-ปริยนฺตํ อุคฺคเหตฺวา, ธมฺเมน กาเรติ ปฏิญฺญาย, อาปตฺติํ ชานาติ, อาปตฺติยา มูลํ ชานาติ, อาปตฺติ\-สมุทยํ ชานาติ, อาปตฺตินิโรธํ ชานาติ, อาปตฺติ\-นิโรธคามินิํ ปฏิปทํ ชานาติ.}

\prose{อปเรหิปิ ทสหงฺเคหิ สมนฺนาคโต วินยธโร “พาโล”เตฺวว สงฺขํ คจฺฉติ. อธิกรณํ น ชานาติ, อธิกรณสฺส มูลํ น ชานาติ, อธิกรณ\-สมุทยํ น ชานาติ, อธิกรณ\-นิโรธํ น ชานาติ, อธิกรณนิโรธ\-คามินิํ ปฏิปทํ น ชานาติ, วตฺถุํ น ชานาติ, นิทานํ\csromansharedfootnote{43d83cc9219b44e4}{อุทฺทานํ (ก)} น ชนาติ, ปญฺญตฺติํ น ชานาติ, อนุปญฺญตฺติํ น ชานาติ, อนุสนฺธิ\-วจนปถํ น ชานาติ.}

\prose{ทสหงฺเคหิ สมนฺนาคโต วินยธโร “ปณฺฑิโต”เตฺวว สงฺขํ คจฺฉติ. อธิกรณํ ชานาติ, อธิกรณสฺส มูลํ ชานาติ, อธิกรณ\-สมุทยํ ชานาติ, อธิกรณ\-นิโรธํ ชานาติ, อธิกรณนิโรธ\-คามินิํ ปฏิปทํ ชานาติ, วตฺถุํ ชานาติ, นิทานํ ชานาติ, ปญฺญตฺติํ ชานาติ, อนุปญฺญตฺติํ ชานาติ, อนุสนฺธิ\-วจนปถํ ชานาติ.}

\prose{อปเรหิปิ ทสหงฺเคหิ สมนฺนาคโต วินยธโร “พาโล”เตฺวว สงฺขํ คจฺฉติ. ญตฺติํ น ชานาติ, ญตฺติยา กรณํ น ชานาติ, น ปุพฺพกุสโล โหติ, น อปรกุสโล โหติ, อกาลญฺญู จ โหติ, อาปตฺตานาปตฺติํ น ชานาติ, ลหุกครุกํ อาปตฺติํ น ชานาติ, สาวเสสา\-นวเสสํ อาปตฺติํ น ชานาติ, ทุฏฺฐุลฺลา\-ทุฏฺฐุลฺลํ อาปตฺติํ น ชานาติ, อาจริยปรมฺปรา โข ปนสฺส น สุคฺคหิตา โหติ น สุมนสิกตา น สูปธาริตา.}

\prose{\csromanfolio{247}ทสหงฺเคหิ สมนฺนาคโต วินยธโร “ปณฺฑิโต”เตฺวว สงฺขํ คจฺฉติ. ญตฺติํ ชานาติ, ญตฺติยา กรณํ ชานาติ, ปุพฺพกุสโล โหติ, อปรกุสโล โหติ, กาลญฺญู จ โหติ, อาปตฺตานาปตฺติํ ชานาติ, ลหุกครุกํ อาปตฺติํ ชานาติ, สาวเสสา\-นวเสสํ อาปตฺติํ ชานาติ, ทุฏฺฐุลฺลา\-ทุฏฺฐุลฺลํ อาปตฺติํ ชานาติ, อาจริยปรมฺปรา โข ปนสฺส สุคฺคหิตา โหติ สุมนสิกตา สูปธาริตา.}

\prose{อปเรหิปิ ทสหงฺเคหิ สมนฺนาคโต วินยธโร “พาโล”เตฺวว สงฺขํ คจฺฉติ. อาปตฺตานาปตฺติํ น ชานาติ, ลหุกครุกํ อาปตฺติํ น ชานาติ, สาวเสสา\-นวเสสํ อาปตฺติํ น ชานาติ, ทุฏฺฐุลฺลา\-ทุฏฺฐุลฺลํ อาปตฺติํ น ชานาติ, อุภยานิ โข ปนสฺส ปาติโมกฺขานิ วิตฺถาเรน น สฺวาคตานิ โหนฺติ น สุวิภตฺตานิ น สุปฺปวตฺตีนิ น สุวินิจฺฉิตานิ สุตฺตโส อนุพฺยญฺชนโส, อาปตฺตานาปตฺติํ น ชานาติ, ลหุกครุกํ อาปตฺติํ น ชานาติ, สาวเสสา\-นวเสสํ อาปตฺติํ น ชานาติ, ทุฏฺฐุลฺลา\-ทุฏฺฐุลฺลํ อาปตฺติํ น ชานาติ, อธิกรเณ จ น วินิจฺฉย\-กุสโล โหติ.}

\prose{ทสหงฺเคหิ สมนฺนาคโต วินยธโร “ปณฺฑิโต”เตฺวว สงฺขํ คจฺฉติ, อาปตฺตานาปตฺติํ ชานาติ, ลหุกครุกํ อาปตฺติํ ชานาติ, สาวเสสา\-นวเสสํ อาปตฺติํ ชานาติ, ทุฏฺฐุลฺลา\-ทุฏฺฐุลฺลํ อาปตฺติํ ชานาติ, อุภยานิ โข ปนสฺส ปาติโมกฺขานิ วิตฺถาเรน สฺวาคตานิ โหนฺติ สุวิภตฺตานิ สุปฺปวตฺตีนิ สุวินิจฺฉิตานิ สุตฺตโส อนุพฺยญฺชนโส, อาปตฺตานาปตฺติํ ชานาติ, ลหุกครุกํ อาปตฺติํ ชานาติ, สาวเสสา\-นวเสสํ อาปตฺติํ ชานาติ, ทุฏฺฐุลฺลา\-ทุฏฺฐุลฺลํ อาปตฺติํ ชานาติ, อธิกรเณ จ วินิจฺฉย\-กุสโล โหติ.}

\prosewithcloser{\csromansymbolfootnote{*}{อํ 3. 312 ปิฏฺเฐปิ.}ทสหงฺเคหิ สมนฺนาคโต ภิกฺขุ อุพฺพาหิกาย สมฺมนฺนิตพฺโพ. ทส อตฺถวเส ปฏิจฺจ ตถาคเตน สาวกานํ สิกฺขาปทํ ปญฺญตฺตํ. ทส อาทีนวา ราชนฺเตปุร\-ปฺปเวสเน. ทส ทานวตฺถูนิ. ทส รตนานิ. ทสวคฺโค ภิกฺขุสํโฆ. ทสวคฺเคน คเณน อุปสมฺ\-ปาเท\-ตพฺพํ. ทส ปํสุกูลานิ. ทส จีวรธารณา\csromansharedfootnote{e6c79bd930dbfcf3}{จีวรธารณานิ (สฺยา)}. ทสาหปรมํ อติเรกจีวรํ \csromanfolio{248}ธาเรตพฺพํ. ทส สุกฺกานิ. ทส อิตฺถิโย. ทส ภริยาโย. เวสาลิยา ทส วตฺถูนิ ทีเปนฺติ. ทส ปุคฺคลา อวนฺทิยา. ทส อกฺโกสวตฺถูนิ. ทสหากาเรหิ เปสุญฺญํ อุปสํหรติ. ทส เสนาสนานิ. ทส วรานิ ยาจิํสุ. ทส อธมฺมิกานิ ปาติโมกฺข\-ฏฺฐปนานิ. ทส ธมฺมิกานิ ปาติโมกฺข\-ฏฺฐปนานิ. ทสานิสํสา ยาคุยา. ทส มํสา อกปฺปิยา. ทส ปรมานิ. ทสวสฺเสน ภิกฺขุนา พฺยตฺเตน ปฏิพเลน ปพฺพาเชตพฺพํ, อุปสมฺ\-ปาเท\-ตพฺพํ, นิสฺสโย ทาตพฺโพ, สามเณโร อุปฏฺฐาเปตพฺโพ. ทสวสฺสาย ภิกฺขุนิยา พฺยตฺตาย ปฏิพลาย ปพฺพาเชตพฺพํ, อุปสมฺ\-ปาเท\-ตพฺพํ, นิสฺสโย ทาตพฺโพ, สามเณรี อุปฏฺฐาเปตพฺพา. ทสวสฺสาย ภิกฺขุนิยา พฺยตฺตาย ปฏิพลาย วุฏฺฐาปน\-สมฺมุติ สาทิตพฺพา. ทสวสฺสาย คิหิคตาย สิกฺขา ทาตพฺพาติ.}{ทสกํ นิฏฺฐิตํ.}
\csromancenter{ตสฺสุทฺทานํ.}
\setlength{\csromangathapagewidth}{0pt}
\setlength{\csromangathawakwidth}{0pt}
\csromangathameasuregroup{อาฆาตํ วินยํ วตฺถุ, \\ มิจฺฉตฺตา เจว สมฺมตฺตา, \\ สลากา อธมฺมา ธมฺมา, \\ ภาสาธิกรณญฺเจว, \\ ลหุกา ครุกา เอเต, \\ อุพฺพาหิกา จ สิกฺขา จ, \\ รตนํ ทสวคฺโค จ, \\ ปํสุกูลธารณา จ, \\ ภริยา ทส วตฺถูนิ, \\ เปสุญฺญญฺเจว เสนานิ, \\ ธมฺมิกา ยาคุมํสา จ, \\ วุฏฺฐาปนา คิหิคตา,}{\csromangathabat{อาฆาตํ วินยํ วตฺถุ,}{มิจฺฉา สมฺมา จ อนฺตคา.} \\ \csromangathabat{มิจฺฉตฺตา เจว สมฺมตฺตา,}{อกุสลา กุสลาปิ จ.} \\ \csromangathabat{สลากา อธมฺมา ธมฺมา,}{สามเณรา จ นาสนา.} \\ \csromangathabat{ภาสาธิกรณญฺเจว,}{ญตฺติลหุกเมว จ.} \\ \csromangathabat{ลหุกา ครุกา เอเต,}{กณฺหสุกฺกา วิชานถ.} \\ \csromangathabat{อุพฺพาหิกา จ สิกฺขา จ,}{อนฺเตปุเร จ วตฺถูนิ.} \\ \csromangathabat{รตนํ ทสวคฺโค จ,}{ตเถว อุปสมฺปทา.} \\ \csromangathabat{ปํสุกูลธารณา จ,}{ทสาหสุกฺกอิตฺถิโย.} \\ \csromangathabat{ภริยา ทส วตฺถูนิ,}{อวนฺทิยกฺโกเสน จ.} \\ \csromangathabat{เปสุญฺญญฺเจว เสนานิ,}{วรานิ จ อธมฺมิกา.} \\ \csromangathabat{ธมฺมิกา ยาคุมํสา จ,}{ปรมา ภิกฺขุ ภิกฺขุนี.} \\ \csromangathabat{วุฏฺฐาปนา คิหิคตา,}{ทสกา สุปฺปกาสิตาติ.}}
\csromangathasetleft{อาฆาตํ วินยํ วตฺถุ, \\ มิจฺฉตฺตา เจว สมฺมตฺตา, \\ สลากา อธมฺมา ธมฺมา, \\ ภาสาธิกรณญฺเจว, \\ ลหุกา ครุกา เอเต, \\ อุพฺพาหิกา จ สิกฺขา จ, \\ รตนํ ทสวคฺโค จ, \\ ปํสุกูลธารณา จ, \\ ภริยา ทส วตฺถูนิ, \\ เปสุญฺญญฺเจว เสนานิ, \\ ธมฺมิกา ยาคุมํสา จ, \\ วุฏฺฐาปนา คิหิคตา,}
\csromangathagroup{\csromangathastanza{\csromangathabat{อาฆาตํ วินยํ วตฺถุ,}{มิจฺฉา สมฺมา จ อนฺตคา.} \\ \csromangathabat{มิจฺฉตฺตา เจว สมฺมตฺตา,}{อกุสลา กุสลาปิ จ.}}\csromangathastanza{\csromangathabat{สลากา อธมฺมา ธมฺมา,}{สามเณรา จ นาสนา.} \\ \csromangathabat{ภาสาธิกรณญฺเจว,}{ญตฺติลหุกเมว จ.}}\csromangathastanza{\csromangathabat{ลหุกา ครุกา เอเต,}{กณฺหสุกฺกา วิชานถ.} \\ \csromangathabat{อุพฺพาหิกา จ สิกฺขา จ,}{อนฺเตปุเร จ วตฺถูนิ.}}\csromangathastanza{\csromangathabat{รตนํ ทสวคฺโค จ,}{ตเถว อุปสมฺปทา.} \\ \csromangathabat{ปํสุกูลธารณา จ,}{ทสาหสุกฺกอิตฺถิโย.}}\csromangathastanza{\csromangathabat{ภริยา ทส วตฺถูนิ,}{อวนฺทิยกฺโกเสน จ.} \\ \csromangathabat{เปสุญฺญญฺเจว เสนานิ,}{วรานิ จ อธมฺมิกา.}}\csromangathastanza{\csromangathabat{ธมฺมิกา ยาคุมํสา จ,}{ปรมา ภิกฺขุ ภิกฺขุนี.} \\ \csromangathabat{วุฏฺฐาปนา คิหิคตา,}{ทสกา สุปฺปกาสิตาติ.}}}

\csromanmatikaanchor{mtk.174}
\csromantocmarkref{paragraph}{11. เอกาทสกวาร}{mtk.174}
\bookmark[dest={mtk.174},level=4]{11. เอกาทสกวาร}
\csromanheader{4}{11. \textbf{เอกาทสกวาร}}
\proseitemwithcloser{331}{\csromanfolio{249}เอกาทส ปุคฺคลา อนุปสมฺปนฺนา น อุปสมฺ\-ปาเท\-ตพฺพา, อุปสมฺปนฺนา นาเสตพฺพา. เอกาทส ปาทุกา อกปฺปิยา. เอกาทส ปตฺตา อกปฺปิยา. เอกาทส จีวรานิ อกปฺปิยานิ. เอกาทส ยาวตติยกา. ภิกฺขุนีนํ เอกาทส อนฺตรายิกา ธมฺมา ปุจฺฉิตพฺพา. เอกาทส จีวรานิ อธิฏฺฐาตพฺพานิ. เอกาทส จีวรานิ น วิกปฺเปตพฺพานิ. เอกาทเส อรุณุคฺคมเน นิสฺสคฺคิยํ โหติ. เอกา ทส คณฺฐิกา กปฺปิยา. เอกาทส วิธา\csromansharedfootnote{0e445aaa01da027c}{วีถา (สฺยา)} กปฺปิยา. เอกาทส ปถวี อกปฺปิยา. เอกาทส ปถวี กปฺปิยา. เอกาทส นิสฺสย\-ปฏิปฺปสฺสทฺธิโย. เอกาทส ปุคฺคลา อวนฺทิยา. เอกาทส ปรมานิ. เอกาทส วรานิ ยาจิํสุ. เอกาทส สีมาโทสา. อกฺโกสก\-ปริภาสเก ปุคฺคเล เอกาทสาทีนวา ปาฏิกงฺขา.\csromansymbolfootnote{*}{อํ 3. 542; ขุ 11. 195 ปิฏฺเฐสุปิ.}เมตฺตาย เจโตวิมุตฺติยา อาเสวิตาย ภาวิตาย พหุลีกตาย ยานีกตาย วตฺถุกตาย อนุฏฺฐิตาย ปริจิตาย สุสมารทฺธาย เอกาทสานิสํสา ปาฏิกงฺขา. สุขํ สุปติ, สุขํ ปฏิพุชฺฌติ, น ปาปกํ สุปินํ ปสฺสติ, มนุสฺสานํ ปิโย โหติ, อมนุสฺสานํ ปิโย โหติ, เทวตา รกฺขนฺติ, นาสฺส อคฺคิ วา วิสํ วา สตฺถํ วา กมติ, ตุวฏฺฏํ จิตฺตํ สมาธิยติ, มุขวณฺโณ วิปฺปสีทติ, อสมฺมูโฬฺห กาลํ กโรติ, อุตฺตริ อปฺปฏิวิชฺฌนฺโต พฺรหฺมโลกูปโค โหติ. เมตฺตาย เจโตวิมุตฺติยา อาเสวิตาย ภาวิตาย พหุลีกตาย ยานีกตาย วตฺถุกตาย อนุฏฺฐิตาย ปริจิตาย สุสมารทฺธาย อิเม เอกาทสานิสํสา ปาฏิกงฺขาติ.}{เอกาทสกํ นิฏฺฐิตํ.}
\csromancenter{ตสฺสุทฺทานํ.}
\setlength{\csromangathapagewidth}{0pt}
\setlength{\csromangathawakwidth}{0pt}
\csromangathameasuregroup{นาเสตพฺพา ปาทุกา จ, \\ ตติยา ปุจฺฉิตพฺพา จ, \\ อรุณา คณฺฐิกา วิธา, \\ นิสฺสยาวนฺทิยา เจว, \\ สีมาโทสา จ อกฺโกสา,}{\csromangathabat{นาเสตพฺพา ปาทุกา จ,}{ปตฺตา จ จีวรานิ จ.} \\ \csromangathabat{ตติยา ปุจฺฉิตพฺพา จ,}{อธิฏฺฐานวิกปฺปนา.} \\ \csromangathabat{อรุณา คณฺฐิกา วิธา,}{อกปฺปิยา จ กปฺปิยํ.} \\ \csromangathabat{นิสฺสยาวนฺทิยา เจว,}{ปรมานิ วรานิ จ.} \\ \csromangathabat{สีมาโทสา จ อกฺโกสา,}{เมตฺตาเยกาทสา กตาติ.}}
\csromangathasetleft{นาเสตพฺพา ปาทุกา จ, \\ ตติยา ปุจฺฉิตพฺพา จ, \\ อรุณา คณฺฐิกา วิธา, \\ นิสฺสยาวนฺทิยา เจว, \\ สีมาโทสา จ อกฺโกสา,}
\csromangathagroup{\csromangathastanza{\csromangathabat{นาเสตพฺพา ปาทุกา จ,}{ปตฺตา จ จีวรานิ จ.} \\ \csromangathabat{ตติยา ปุจฺฉิตพฺพา จ,}{อธิฏฺฐานวิกปฺปนา.}}\csromangathastanza{\csromangathabat{อรุณา คณฺฐิกา วิธา,}{อกปฺปิยา จ กปฺปิยํ.} \\ \csromangathabat{นิสฺสยาวนฺทิยา เจว,}{ปรมานิ วรานิ จ.} \\ \csromangathabat{สีมาโทสา จ อกฺโกสา,}{เมตฺตาเยกาทสา กตาติ.}}}

\vspace{\tipitakaparskip}
\csromancenter{เอกุตฺตริกํ.}
\csromancenter{ตสฺสุทฺทานํ.}
\setlength{\csromangathapagewidth}{0pt}
\setlength{\csromangathawakwidth}{0pt}
\csromangathameasuregroup{เอกกา จ ทุกา เจว, \\ ฉสตฺตฏฺฐนวกา จ, \\ หิตาย สพฺพสตฺตานํ, \\ เอกุตฺตริกา วิมลา,}{\csromangathabat{เอกกา จ ทุกา เจว,}{ติกา จ จตุปญฺจกา.} \\ \csromangathabat{ฉสตฺตฏฺฐนวกา จ,}{ทส เอกาทสานิ จ.} \\ \csromangathabat{หิตาย สพฺพสตฺตานํ,}{ญาตธมฺเมน ตาทินา.} \\ \csromangathabat{เอกุตฺตริกา วิมลา,}{มหาวีเรน เทสิตาติ.}}
\csromangathasetleft{เอกกา จ ทุกา เจว, \\ ฉสตฺตฏฺฐนวกา จ, \\ หิตาย สพฺพสตฺตานํ, \\ เอกุตฺตริกา วิมลา,}
\csromangathagroupwithcloser{\csromangathastanza{\csromanfolio{250}\csromangathabat{เอกกา จ ทุกา เจว,}{ติกา จ จตุปญฺจกา.} \\ \csromangathabat{ฉสตฺตฏฺฐนวกา จ,}{ทส เอกาทสานิ จ.}}\csromangathastanza{\csromangathabat{หิตาย สพฺพสตฺตานํ,}{ญาตธมฺเมน ตาทินา.} \\ \csromangathabat{เอกุตฺตริกา วิมลา,}{มหาวีเรน เทสิตาติ.}}}{\textbf{เอกุตฺตริกํ นิฏฺฐิตํ}}
\csromanheader{2}{อุโปสถา\-ทิ\-ปุจฺฉา\-วิสฺสชฺชนา}
\csromanmatikaanchor{mtk.175}
\csromanmatikaanchor{mtk.176}
\csromantocmarknopage{section}{อุโปสถาทิปุจฺฉาวิสชฺชนา}{mtk.175}
\bookmark[dest={mtk.175},level=1]{อุโปสถาทิปุจฺฉาวิสชฺชนา}
\csromantocmarkref{section}{อาทิมชฺฌนฺตปุจฺฉนํ}{mtk.176}
\bookmark[dest={mtk.176},level=1]{อาทิมชฺฌนฺตปุจฺฉนํ}
\csromanheader{1}{อาทิ\-มชฺฌ\-นฺต\-ปุจฺฉนํ}
\csromanmatikaanchor{mtk.177}
\csromantocmarkref{section}{อาทิมชฺฌนฺตวิสชฺชนา}{mtk.177}
\bookmark[dest={mtk.177},level=1]{อาทิมชฺฌนฺตวิสชฺชนา}
\proseitem{332}{\csromanfolio{251}อุโปสถ\-กมฺมสฺส โก อาทิ, กิํ มชฺเฌ, กิํ ปริโยสานํ. ปวารณา\-กมฺมสฺส โก อาทิ, กิํ มชฺเฌ, กิํ ปริโยสานํ. ตชฺชนีย\-กมฺมสฺส โก อาทิ, กิํ มชฺเฌ, กิํ ปริโยสานํ. นิยสฺส\-กมฺมสฺส ฯเปฯ. ปพฺพาชนีย\-กมฺมสฺส ฯเปฯ. ปฏิสารณีย\-กมฺมสฺส ฯเปฯ. อุกฺเขปนีย\-กมฺมสฺส ฯเปฯ. ปริวาส\-ทานสฺส ฯเปฯ. มูลาย\-ปฏิกสฺสนาย ฯเปฯ. มานตฺตทานสฺส ฯเปฯ. อพฺภานสฺส ฯเปฯ. อุปสมฺปทา\-กมฺมสฺส โก อาทิ, กิํ มชฺเฌ, กิํ ปริโยสานํ. ตชฺชนีย\-กมฺมสฺส ปฏิปฺปสฺสทฺธิยา โก อาทิ, กิํ มชฺเฌ, กิํ ปริโยสานํ. นิยสฺส\-กมฺมสฺส ปฏิปฺปสฺสทฺธิยา โก อาทิ, กิํ มชฺเฌ, กิํ ปริโยสานํ. ปพฺพาชนีย\-กมฺมสฺส ปฏิปฺปสฺสทฺธิยา โก อาทิ, กิํ มชฺเฌ, กิํ ปริโยสานํ. ปฏิสารณีย\-กมฺมสฺส ปฏิปฺปสฺสทฺธิยาโก อาทิ, กิํ มชฺเฌ, กิํ ปริโยสานํ. อุกฺเขปนีย\-กมฺมสฺส ปฏิปฺปสฺสทฺธิยา โก อาทิ, กิํ มชฺเฌ, กิํ ปริโยสานํ. สติวินยสฺส โก อาทิ, กิํ มชฺเฌ, กิํ ปริโยสานํ. อมูฬฺห\-วินยสฺส โก อาทิ, กิํ มชฺเฌ, กิํ ปริโยสานํ. ตสฺสปาปิยสิกาย โก อาทิ, กิํ มชฺเฌ, กิํ ปริโยสานํ. ติณวตฺถารกสฺส โก อาทิ, กิํ มชฺเฌ, กิํ ปริโยสานํ. ภิกฺขุโน วาทกสมฺมุติยา โก อาทิ, กิํ มชฺเฌ, กิํ ปริโยสานํ. ติจีวเรน อวิปฺปวาส\-สมฺมุติยา โก อาทิ, กิํ มชฺเฌ, กิํ ปริโยสนํ. สนฺถต\-สมฺมุติยา โก อาทิ, กิํ มชฺเฌ, กิํ ปริโยสานํ. รูปิยฉฑฺฑก\-สมฺมุติยา โก อาทิ, กิํ มชฺเฌ, กิํ ปริโยสานํ. สาฏิยคฺคาหาปก\-สมฺมุติยา โก อาทิ, กิํ มชฺเฌ, กิํ ปริโยสานํ. ปตฺต\-คฺคาหาปก\-สมฺมุติยา โก อาทิ, กิํ มชฺเฌ, กิํ ปริโยสานํ. ทณฺฑสมฺมุติยา โก อาทิ, กิํ มชฺเฌ, กิํ ปริโยสานํ. สิกฺกาสมฺมุติยา โก อาทิ, กิํ มชฺเฌ, กิํ ปริโยสานํ. ทณฺฑ\-สิกฺกา\-สมฺมุติยา โก อาทิ, กิํ มชฺเฌ, กิํ ปริโยสานํ.}

\csromanheader{2}{อาทิ\-มชฺฌ\-นฺต\-วิสฺสชฺชนา}
\proseitem{333}{อุโปสถ\-กมฺมสฺส โก อาทิ, กิํ มชฺเฌ, กิํ ปริโยสานนฺติ อุโปสถ\-กมฺมสฺส สามคฺคี อาทิ, กิริยา มชฺเฌ, นิฏฺฐานํ ปริโยสานํ.}

\prose{\csromanfolio{252}ปวารณา\-กมฺมสฺส โก อาทิ, กิํ มชฺเฌ, กิํ ปริโยสานนฺติ ปวารณา\-กมฺมสฺส สามคฺคี อาทิ, กิริยา มชฺเฌ, นิฏฺฐานํ ปริโยสานํ.}

\prose{ตชฺชนีย\-กมฺมสฺส โก อาทิ, กิํ มชฺเฌ, กิํ ปริโยสานนฺติ ตชฺชนีย\-กมฺมสฺส วตฺถุ จ ปุคฺคโล จ อาทิ, ญตฺติ มชฺเฌ, กมฺมวาจา ปริโยสานํ.}

\prose{นิยสฺส\-กมฺมสฺส ฯเปฯ. ปพฺพาชนีย\-กมฺมสฺส ฯเปฯ. ปฏิสารณีย\-กมฺมสฺส ฯเปฯ. อุกฺเขปนีย\-กมฺมสฺส ฯเปฯ. ปริวาส\-ทานสฺส ฯเปฯ. มูลาย ปฏิกสฺสนาย ฯเปฯ. มานตฺตทานสฺส ฯเปฯ. อพฺภานสฺส โก อาทิ, กิํ มชฺเฌ, กิํ ปริโยสานนฺติ อพฺภานสฺส วตฺถุ จ ปุคฺคโล จ อาทิ, ญตฺติ มชฺเฌ, กมฺมวาจา ปริโยสานํ.}

\prose{อุปสมฺปทา\-กมฺมสฺส โก อาทิ, กิํ มชฺเฌ, กิํ ปริโยสานนฺติ อุปสมฺปทา\-กมฺมสฺส ปุคฺคโล อาทิ, ญตฺติ มชฺเฌ, กมฺมวาจา ปริโยสานํ.}

\prose{ตชฺชนีย\-กมฺมสฺส ปฏิปฺปสฺสทฺธิยา โก อาทิ, กิํ มชฺเฌ, กิํ ปริโยสานนฺติ ตชฺชนีย\-กมฺมสฺส ปฏิปฺปสฺสทฺธิยา สมฺมาวตฺตนา อาทิ, ญตฺติ มชฺเฌ, กมฺมวาจา ปริโยสานํ.}

\prose{นิยสฺส\-กมฺมสฺส ฯเปฯ. ปพฺพาชนีย\-กมฺมสฺส ฯเปฯ. ปฏิสารณีย\-กมฺมสฺส ฯเปฯ. อุกฺเขปนีย\-กมฺมสฺส ปฏิปฺปสฺสทฺธิยา โก อาทิ, กิํ มชฺเฌ, กิํ ปริโยสานนฺติ อุกฺเขปนีย\-กมฺมสฺส ปฏิปฺปสฺสทฺธิยา สมฺมาวตฺตนา อาทิ, ญตฺติ มชฺเฌ, กมฺมวาจา ปริโยสานํ.}

\prose{สติวินยสฺส โก อาทิ, กิํ มชฺเฌ, กิํ ปริโยสนนฺติ สติวินยสฺส วตฺถุ จ ปุคฺคโล จ อาทิ, ญตฺติ มชฺเฌ, กมฺมวาจา ปริโยสานํ.}

\prosewithcloser{อมูฬฺห\-วินยสฺส ฯเปฯ. ตสฺสปาปิยสิกาย ฯเปฯ. ติณวตฺถารกสฺส ฯเปฯ. ภิกฺขุโนวาทก\-สมฺมุติยา ฯเปฯ. ติจีวเรน อวิปฺปวาส\-สมฺมุติยา ฯเปฯ. สนฺถต\-สมฺมุติยา ฯเปฯ. รูปิยฉฑฺฑก\-สมฺมุติยา ฯเปฯ. สาฏิยคฺคาหาปก\-สมฺมุติยา ฯเปฯ. ปตฺต\-คฺคาหาปก\-สมฺมุติยา ฯเปฯ. ทณฺฑสมฺมุติยา ฯเปฯ. สิกฺกาสมฺมุติยา ฯเปฯ. ทณฺฑ\-สิกฺกา\-สมฺมุติยา โก อาทิ, กิํ มชฺเฌ, กิํ ปริโยสานนฺติ ทณฺฑ\-สิกฺกา\-สมฺมุติยา วตฺถุ จ ปุคฺคโล จ อาทิ, ญตฺติ มชฺเฌ, กมฺมวาจา ปริโยสานํ.}{อุโปสถา\-ทิ\-ปุจฺฉา\-วิสฺสชฺชนา นิฏฺฐิตา.}
\csromanmatikaanchor{mtk.178}
\csromantocmarkref{section}{อตฺถวสปกรณ}{mtk.178}
\bookmark[dest={mtk.178},level=1]{อตฺถวสปกรณ}
\csromanheader{1}{อตฺถวส\-ปกรณ}
\proseitem{334}{\csromanfolio{253}\csromansymbolfootnote{*}{อํ 3. 311; วิ 5. 12 ปิฏฺฐาทีสุปิ.}ทส อตฺถวเส ปฏิจฺจ ตถาคเตน สาวกานํ สิกฺขาปทํ ปญฺญตฺตํ. สํฆสุฏฺฐุตาย, สํฆผาสุตาย, ทุมฺมงฺกูนํ ปุคฺคลานํ นิคฺคหาย, เปสลานํ ภิกฺขูนํ ผาสุวิหาราย, ทิฏฺฐ\-ธมฺมิกานํ อาสวานํ สํวราย, สมฺปรายิกานํ อาสวานํ ปฏิฆาตาย, อปฺปสนฺนานํ ปสาทาย, ปสนฺนานํ ภิยฺโยภาวาย, สทฺธมฺม\-ฏฺฐิติยา, วินยา\-นุคฺคหาย.}

\prose{ยํ สํฆสุฏฺฐุ, ตํ สํฆผาสุ. ยํ สํฆผาสุ, ตํ ทุมฺมงฺกูนํ ปุคฺคลานํ นิคฺคหาย. ยํ ทุมฺมงฺกูนํ ปุคฺคลานํ นิคฺคหาย, ตํ เปสลานํ ภิกฺขุนํ ผาสุวิหาราย. ยํ เปสลานํ ภิกฺขูนํ ผาสุวิหาราย, ตํ ทิฏฺฐ\-ธมฺมิกานํ อาสวานํ สํวราย. ยํ ทิฏฺฐ\-ธมฺมิกานํ อาสวานํ สํวราย, ตํ สมฺปรายิกานํ อาสวานํ ปฏิฆาตาย. ยํ สมฺปรายิกานํ อาสวานํ ปฏิฆาตาย, ตํ อปฺปสนฺนานํ ปสาทาย. ยํ อปฺปสนฺนานํ ปสาทาย, ตํ ปสนฺนานํ ภิยฺโยภาวาย. ยํ ปสนฺนานํ ภิยฺโยภาวาย, ตํ สทฺธมฺม\-ฏฺฐิติยา. ยํ สทฺธมฺม\-ฏฺฐิติยา, ตํ วินยา\-นุคฺคหาย.}

\prose{ยํ สํฆสุฏฺฐุ, ตํ สํฆผาสุ. ยํ สํฆสุฏฺฐุ, ตํ ทุมฺมงฺกูนํ ปุคฺคลานํ นิคฺคหาย. ยํ สํฆสุฏฺฐุ, ตํ เปสลานํ ภิกฺขูนํ ผาสุวิหาราย. ยํ สํฆสุฏฺฐุ, ตํ ทิฏฺฐ\-ธมฺมิกานํ อาสวานํ สํวราย. ยํ สํฆสุฏฺฐุ, ตํ สมฺปรายิกานํ อาสวานํ ปฏิฆาตาย. ยํ สํฆสุฏฺฐุ, ตํ อปฺปสนฺนานํ ปสาทาย. ยํ สํฆสุฏฺฐุ, ตํ ปสนฺนานํ ภิยฺโยภาวาย. ยํ สํฆสุฏฺฐุ, ตํ สทฺธมฺม\-ฏฺฐิติยา. ยํ สํฆสุฏฺฐุ, ตํ วินยา\-นุคฺคหาย.}

\prose{ยํ สํฆผาสุ, ตํ ทุมฺมงฺกูนํ ปุคฺคลานํ นิคฺคหาย. ยํ สํฆผาสุ, ตํ เปสลานํ ภิกฺขูนํ ผาสุวิหาราย. ยํ สํฆผาสุ, ตํ ทิฏฺฐ\-ธมฺมิกานํ อาสวานํ สํวราย. ยํ สํฆผาสุ, ตํ สมฺปรายิกานํ อาสวานํ ปฏิฆาตาย. ยํ สํฆผาสุ, ตํ อปฺปสนฺนานํ ปสาทาย. ยํ สํฆผาสุ, ตํ ปสนฺนานํ ภิยฺโยภาวาย. ยํ สํฆผาสุ, ตํ สทฺธมฺม\-ฏฺฐิติยา. ยํ สํฆผาสุ, ตํ วินยา\-นุคฺคหาย. ยํ สํฆผาสุ, ตํ สํฆสุฏฺฐุ.}

\prose{ยํ ทุมฺมงฺกูนํ ปุคฺคลานํ นิคฺคหาย ฯเปฯ. ยํ เปสลานํ ภิกฺขูนํ ผาสุวิหาราย. ยํ ทิฏฺฐ\-ธมฺมิกานํ อาสวานํ สํวราย. ยํ สมฺปรายิกานํ \csromanfolio{254}อาสวานํ ปฏิฆาตาย. ยํ อปฺปสนฺนานํ ปสาทาย. ยํ ปสนฺนานํ ภิยฺโยภาวาย. ยํ สทฺธมฺม\-ฏฺฐิติยา. ยํ วินยา\-นุคฺคหาย, ตํ สํฆสุฏฺฐุ. ยํ วินยา\-นุคฺคหาย, ตํ สํฆผาสุ. ยํ วินยา\-นุคฺคหาย, ตํ ทุมฺมงฺกูนํ ปุคฺคลานํ นิคฺคหาย. ยํ วินยา\-นุคฺคหาย, ตํ เปสลานํ ภิกฺขูนํ ผาสุวิหาราย. ยํ วินยา\-นุคฺคหาย, ตํ ทิฏฺฐ\-ธมฺมิกานํ อาสวานํ สํวราย. ยํ วินยา\-นุคฺคหาย, ตํ สมฺปรายิกานํ อาสวานํ ปฏิฆาตาย. ยํ วินยา\-นุคฺคหาย, ตํ อปฺปสนฺนานํ ปสาทาย. ยํ วินยา\-นุคฺคหาย, ตํ ปสนฺนานํ ภิยฺโยภาวาย. ยํ วินยา\-นุคฺคหาย, ตํ สทฺธมฺม\-ฏฺฐิติยาติ.}

\setlength{\csromangathapagewidth}{0pt}
\setlength{\csromangathawakwidth}{0pt}
\csromangathameasuregroup{อตฺถสตํ ธมฺมสตํ, \\ จตฺตาริ ญาณสตานิ,}{\csromangathabat{อตฺถสตํ ธมฺมสตํ,}{เทฺว จ นิรุตฺติสตานิ.} \\ \csromangathabat{จตฺตาริ ญาณสตานิ,}{อตฺถวเส ปกรเณติ.}}
\csromangathasetleft{อตฺถสตํ ธมฺมสตํ, \\ จตฺตาริ ญาณสตานิ,}
\csromangathagroupwithclosermid{\csromangathastanza{\csromangathabat{อตฺถสตํ ธมฺมสตํ,}{เทฺว จ นิรุตฺติสตานิ.} \\ \csromangathabat{จตฺตาริ ญาณสตานิ,}{อตฺถวเส ปกรเณติ.}}}{อตฺถวส\-ปกรณํ นิฏฺฐิตํ. มหาวคฺโค นิฏฺฐิโต.}
\csromancenter{ตสฺสุทฺทานํ.}
\setlength{\csromangathapagewidth}{0pt}
\setlength{\csromangathawakwidth}{0pt}
\csromangathameasuregroup{ปฐมํ อฏฺฐปุจฺฉายํ, \\ ภิกฺขูนํ โสฬส เอเต, \\ เปยฺยาลอนฺตรา เภทา, \\ ปวารณตฺถวสิกา,}{\csromangathabat{ปฐมํ อฏฺฐปุจฺฉายํ,}{ปจฺจเยสุ ปุนฏฺฐ จ.} \\ \csromangathabat{ภิกฺขูนํ โสฬส เอเต,}{ภิกฺขุนีนญฺจ โสฬส.} \\ \csromangathabat{เปยฺยาลอนฺตรา เภทา,}{เอกุตฺตริกเมว จ.} \\ \csromangathabat{ปวารณตฺถวสิกา,}{มหาวคฺคสฺส สงฺคโหติ.}}
\csromangathasetleft{ปฐมํ อฏฺฐปุจฺฉายํ, \\ ภิกฺขูนํ โสฬส เอเต, \\ เปยฺยาลอนฺตรา เภทา, \\ ปวารณตฺถวสิกา,}
\csromangathagroupwithcloser{\csromangathastanza{\csromangathabat{ปฐมํ อฏฺฐปุจฺฉายํ,}{ปจฺจเยสุ ปุนฏฺฐ จ.} \\ \csromangathabat{ภิกฺขูนํ โสฬส เอเต,}{ภิกฺขุนีนญฺจ โสฬส.}}\csromangathastanza{\csromangathabat{เปยฺยาลอนฺตรา เภทา,}{เอกุตฺตริกเมว จ.} \\ \csromangathabat{ปวารณตฺถวสิกา,}{มหาวคฺคสฺส สงฺคโหติ.}}}{อตฺถวส\-ปกรณํ นิฏฺฐิตํ.}
\csromanmatikaanchor{mtk.179}
\csromantocmarknopage{section}{คาถาสงฺคณิก}{mtk.179}
\bookmark[dest={mtk.179},level=1]{คาถาสงฺคณิก}
\csromanheader{1}{\textbf{คาถาสงฺคณิก}}
\csromanmatikaanchor{mtk.180}
\csromantocmarkref{subparagraph}{1. สตฺตนคเรสุ ปญฺญตฺตสิกฺขาปท}{mtk.180}
\bookmark[dest={mtk.180},level=5]{1. สตฺตนคเรสุ ปญฺญตฺตสิกฺขาปท}
\csromanheader{5}{1. สตฺตนคเรสุ ปญฺญตฺต\-สิกฺขาปท}
\setlength{\csromangathapagewidth}{0pt}
\setlength{\csromangathawakwidth}{0pt}
\csromangathameasuregroup{ทฺวีสุ วินเยสุ เย ปญฺญตฺตา, \\ กติ เต สิกฺขาปทา โหนฺติ, \\ ทฺวีสุ วินเยสุ เย ปญฺญตฺตา, \\ อฑฺฒุฑฺฒสตานิ เต โหนฺติ, \\ กตเมสุ สตฺตสุ นคเรสุ ปญฺญตฺตา, \\ ตํ วจนปถํ\textsuperscript{1} นิสามยิตฺวา, \\ ทส เวสาลิยํ ปญฺญตฺตา, \\ ฉอูน ตีณิ สตานิ, \\ ฉ อาฬวิยํ ปญฺญตฺตา, \\ อฏฺฐ สกฺเกสุ วุจฺจนฺติ, \\ เย เวสาลิยํ ปญฺญตฺตา, \\ เมถุนวิคฺคหุตฺตริ, \\ ภูตํ ปรมฺปรภตฺตํ, \\ ภิกฺขุนีสุ จ อกฺโกโส, \\ เย ราชคเห ปญฺญตฺตา, \\ อนฺตรวาสกํ รูปิยํ สุตฺตํ, \\ คณโภชนํ วิกาเล จ, \\ จีวรํ ทตฺวา โวสาสนฺติ, \\ คิรคฺคจริยา ตตฺเถว, \\ เย สาวตฺถิยํ ปญฺญตฺตา, \\ ปาราชิกานิ จตฺตาริ, \\ อนิยตา จ เทฺว โหนฺติ, \\ ฉปญฺญาสสตญฺเจว, \\ ทสเยว จ คารยฺหา, \\ ฉอูน ตีณิสตานิ, \\ เย อาฬวิยํ ปญฺญตฺตา, \\ กุฏิโกสิยเสยฺยา จ, \\ สปฺปาณกญฺจ สิญฺจนฺติ, \\ เย โกสมฺพิยํ ปญฺญตฺตา, \\ มหาวิหาโร โทวจสฺสํ, \\ อนาทริยํ สหธมฺโม, \\ เย สกฺเกสุ ปญฺญตฺตา, \\ เอฬกโลมานิ ปตฺโต จ, \\ สูจิ อารญฺญิโก เจว, \\ อุทกสุทฺธิยา โอวาโท, \\ เย ภคฺเคสุ ปญฺญตฺตา, \\ สมาทหิตฺวา วิสิพฺเพนฺติ, \\ ปาราชิกานิ จตฺตาริ, \\ สตฺต จ นิสฺสคฺคิยานิ, \\ เทฺว คารยฺหา ตโย เสกฺขา, \\ ฉสุ นคเรสุ ปญฺญตฺตา, \\ ฉอูน ตีณิ สตานิ, \\ การุณิเกน พุทฺเธน,}{เอกํสํ จีวรํ กตฺวา, ปคฺคณฺหิตฺวาน อญฺชลิํ. \\ อาสีสมานรูโปว\textsuperscript{1}, กิสฺส ตฺวํ อิธ มาคโต. \\ \csromangathabat{ทฺวีสุ วินเยสุ เย ปญฺญตฺตา,}{อุทฺเทสํ อาคจฺฉนฺติ อุโปสเถสุ.} \\ \csromangathabat{กติ เต สิกฺขาปทา โหนฺติ,}{กติสุ นคเรสุ ปญฺญตฺตา.} \\ ภทฺทโก เต อุมฺมงฺโค, โยนิโส ปริปุจฺฉสิ. \\ ตคฺฆ เต อหมกฺขิสฺสํ, ยถาสิ กุสโล ตถา. \\ \csromangathabat{ทฺวีสุ วินเยสุ เย ปญฺญตฺตา,}{อุทฺเทสํ อาคจฺฉนฺติ อุโปสเถสุ.} \\ \csromangathabat{อฑฺฒุฑฺฒสตานิ เต โหนฺติ,}{สตฺตสุ นคเรสุ ปญฺญตฺตา.} \\ \csromangathabat{กตเมสุ สตฺตสุ นคเรสุ ปญฺญตฺตา,}{อิงฺฆ เม ตฺวํ พฺยากร นํ\textsuperscript{1}.} \\ \csromangathabat{ตํ วจนปถํ\textsuperscript{1} นิสามยิตฺวา,}{ปฏิปชฺเชม หิตาย โน สิยา.} \\ เวสาลิยํ ราชคเห, สาวตฺถิยํ จ อาฬวิยํ. \\ โกสมฺพิยํ จ สกฺเกสุ, ภคฺเคสุ เจว ปญฺญตฺตา. \\ กติ เวสาลิยํ ปญฺญตฺตา, กติ ราชคเห กตา. \\ สาวตฺถิยํ กติ โหนฺติ, กติ อาฬวิยํ กตา. \\ กติ โกสมฺพิยํ ปญฺญตฺตา, กติ สกฺเกสุ วุจฺจนฺติ. \\ กติ ภคฺเคสุ ปญฺญตฺตา, ตํ เม อกฺขาหิ ปุจฺฉิโต. \\ \csromangathabat{ทส เวสาลิยํ ปญฺญตฺตา,}{เอกวีส ราชคเห กตา.} \\ \csromangathabat{ฉอูน ตีณิ สตานิ,}{สพฺเพ สาวตฺถิยํ กตา.} \\ \csromangathabat{ฉ อาฬวิยํ ปญฺญตฺตา,}{อฏฺฐ โกสมฺพิยํ กตา.} \\ \csromangathabat{อฏฺฐ สกฺเกสุ วุจฺจนฺติ,}{ตโย ภคฺเคสุ ปญฺญตฺตา.} \\ \csromangathabat{เย เวสาลิยํ ปญฺญตฺตา,}{เต สุโณหิ ยถาตถํ\textsuperscript{1}.} \\ \csromangathabat{เมถุนวิคฺคหุตฺตริ,}{อติเรกญฺจ กาฬกํ.} \\ \csromangathabat{ภูตํ ปรมฺปรภตฺตํ,}{ทนฺตโปเนน\textsuperscript{1} อเจลโก.} \\ \csromangathabat{ภิกฺขุนีสุ จ อกฺโกโส,}{ทเสเต เวสาลิยํ กตา.} \\ \csromangathabat{เย ราชคเห ปญฺญตฺตา,}{เต สุโณหิ ยถาตถํ.} \\ อทินฺนาทานํ ราชคเห, \\ เทฺว อนุทฺธํสนา เทฺวปิ จ เภทา. \\ \csromangathabat{อนฺตรวาสกํ รูปิยํ สุตฺตํ,}{อุชฺฌาปเนน จ ปาจิตปิณฺฑํ.} \\ \csromangathabat{คณโภชนํ วิกาเล จ,}{จาริตฺตํ นหานํ อูนวีสติ.} \\ \csromangathabat{จีวรํ ทตฺวา โวสาสนฺติ,}{เอเต ราชคเห กตา.} \\ \csromangathabat{คิรคฺคจริยา ตตฺเถว,}{ฉนฺททาเนน เอกวีสติ.} \\ \csromangathabat{เย สาวตฺถิยํ ปญฺญตฺตา,}{เต สุโณหิ ยถาตถํ.} \\ \csromangathabat{ปาราชิกานิ จตฺตาริ,}{สํฆาทิเสสา ภวนฺติ โสฬส.} \\ \csromangathabat{อนิยตา จ เทฺว โหนฺติ,}{นิสฺสคฺคิยา จตุวีสติ.} \\ \csromangathabat{ฉปญฺญาสสตญฺเจว,}{ขุทฺทกานิ ปวุจฺจนฺติ.} \\ \csromangathabat{ทสเยว จ คารยฺหา,}{เทฺวสตฺตติ จ เสขิยา.} \\ \csromangathabat{ฉอูน ตีณิสตานิ,}{สพฺเพ สาวตฺถิยํ กตา.} \\ \csromangathabat{เย อาฬวิยํ ปญฺญตฺตา,}{เต สุโณหิ ยถาตถํ.} \\ \csromangathabat{กุฏิโกสิยเสยฺยา จ,}{ขณเน คจฺฉ เทวเต.} \\ \csromangathabat{สปฺปาณกญฺจ สิญฺจนฺติ,}{ฉ เอเต อาฬวิยํ กตา.} \\ \csromangathabat{เย โกสมฺพิยํ ปญฺญตฺตา,}{เต สุโณหิ ยถาตถํ.} \\ \csromangathabat{มหาวิหาโร โทวจสฺสํ,}{อญฺญํ ทฺวารํ สุราย จ.} \\ \csromangathabat{อนาทริยํ สหธมฺโม,}{ปโยปาเนน อฏฺฐมํ.} \\ \csromangathabat{เย สกฺเกสุ ปญฺญตฺตา,}{เต สุโณหิ ยถาตถํ.} \\ \csromangathabat{เอฬกโลมานิ ปตฺโต จ,}{โอวาโท เจว เภสชฺชํ.} \\ \csromangathabat{สูจิ อารญฺญิโก เจว,}{อฏฺเฐเต\textsuperscript{1} กาปิลวตฺถเว.} \\ \csromangathabat{อุทกสุทฺธิยา โอวาโท,}{ภิกฺขุนีสุ ปวุจฺจนฺติ.} \\ \csromangathabat{เย ภคฺเคสุ ปญฺญตฺตา,}{เต สุโณหิ ยถาตถํ.} \\ \csromangathabat{สมาทหิตฺวา วิสิพฺเพนฺติ,}{สามิเสน สสิตฺถกํ.} \\ \csromangathabat{ปาราชิกานิ จตฺตาริ,}{สํฆาทิเสสานิ ภวนฺติ.} \\ \csromangathabat{สตฺต จ นิสฺสคฺคิยานิ,}{อฏฺฐ ทฺวตฺติํส ขุทฺทกา.} \\ \csromangathabat{เทฺว คารยฺหา ตโย เสกฺขา,}{ฉปฺปญฺญาส สิกฺขาปทา.} \\ \csromangathabat{ฉสุ นคเรสุ ปญฺญตฺตา,}{พุทฺเธนาทิจฺจพนฺธุนา.} \\ \csromangathabat{ฉอูน ตีณิ สตานิ,}{สพฺเพ สาวตฺถิยํ กตา.} \\ \csromangathabat{การุณิเกน พุทฺเธน,}{โคตเมน ยสสฺสินา.}}
\csromangathasetleft{ทฺวีสุ วินเยสุ เย ปญฺญตฺตา, \\ กติ เต สิกฺขาปทา โหนฺติ, \\ ทฺวีสุ วินเยสุ เย ปญฺญตฺตา, \\ อฑฺฒุฑฺฒสตานิ เต โหนฺติ, \\ กตเมสุ สตฺตสุ นคเรสุ ปญฺญตฺตา, \\ ตํ วจนปถํ\textsuperscript{1} นิสามยิตฺวา, \\ ทส เวสาลิยํ ปญฺญตฺตา, \\ ฉอูน ตีณิ สตานิ, \\ ฉ อาฬวิยํ ปญฺญตฺตา, \\ อฏฺฐ สกฺเกสุ วุจฺจนฺติ, \\ เย เวสาลิยํ ปญฺญตฺตา, \\ เมถุนวิคฺคหุตฺตริ, \\ ภูตํ ปรมฺปรภตฺตํ, \\ ภิกฺขุนีสุ จ อกฺโกโส, \\ เย ราชคเห ปญฺญตฺตา, \\ อนฺตรวาสกํ รูปิยํ สุตฺตํ, \\ คณโภชนํ วิกาเล จ, \\ จีวรํ ทตฺวา โวสาสนฺติ, \\ คิรคฺคจริยา ตตฺเถว, \\ เย สาวตฺถิยํ ปญฺญตฺตา, \\ ปาราชิกานิ จตฺตาริ, \\ อนิยตา จ เทฺว โหนฺติ, \\ ฉปญฺญาสสตญฺเจว, \\ ทสเยว จ คารยฺหา, \\ ฉอูน ตีณิสตานิ, \\ เย อาฬวิยํ ปญฺญตฺตา, \\ กุฏิโกสิยเสยฺยา จ, \\ สปฺปาณกญฺจ สิญฺจนฺติ, \\ เย โกสมฺพิยํ ปญฺญตฺตา, \\ มหาวิหาโร โทวจสฺสํ, \\ อนาทริยํ สหธมฺโม, \\ เย สกฺเกสุ ปญฺญตฺตา, \\ เอฬกโลมานิ ปตฺโต จ, \\ สูจิ อารญฺญิโก เจว, \\ อุทกสุทฺธิยา โอวาโท, \\ เย ภคฺเคสุ ปญฺญตฺตา, \\ สมาทหิตฺวา วิสิพฺเพนฺติ, \\ ปาราชิกานิ จตฺตาริ, \\ สตฺต จ นิสฺสคฺคิยานิ, \\ เทฺว คารยฺหา ตโย เสกฺขา, \\ ฉสุ นคเรสุ ปญฺญตฺตา, \\ ฉอูน ตีณิ สตานิ, \\ การุณิเกน พุทฺเธน,}
\csromangathagroup{\csromangathastanza{\csromanfolio{255}\csromangathaitemline{335}{เอกํสํ จีวรํ กตฺวา, ปคฺคณฺหิตฺวาน อญฺชลิํ.} \\ \csromangathacontline{อาสีสมานรูโปว\csromansharedfootnote{9869bf47d1c716d9}{อาสิํสมานรูโปว (สี, สฺยา)}, กิสฺส ตฺวํ อิธ มาคโต.} \\ \csromangathacontbat{ทฺวีสุ วินเยสุ เย ปญฺญตฺตา,}{อุทฺเทสํ อาคจฺฉนฺติ อุโปสเถสุ.} \\ \csromangathacontbat{กติ เต สิกฺขาปทา โหนฺติ,}{กติสุ นคเรสุ ปญฺญตฺตา.} \\ \csromangathacontline{ภทฺทโก เต อุมฺมงฺโค, โยนิโส ปริปุจฺฉสิ.} \\ \csromangathacontline{ตคฺฆ เต อหมกฺขิสฺสํ, ยถาสิ กุสโล ตถา.} \\ \csromangathacontbat{ทฺวีสุ วินเยสุ เย ปญฺญตฺตา,}{อุทฺเทสํ อาคจฺฉนฺติ อุโปสเถสุ.} \\ \csromangathacontbat{อฑฺฒุฑฺฒสตานิ เต โหนฺติ,}{สตฺตสุ นคเรสุ ปญฺญตฺตา.} \\ \csromangathacontbat{กตเมสุ สตฺตสุ นคเรสุ ปญฺญตฺตา,}{อิงฺฆ เม ตฺวํ พฺยากร นํ\csromansharedfootnote{e0b7e5f0231f7c3d}{อิงฺฆ เม ตํ พฺยากร (ก)}.} \\ \csromangathacontbat{ตํ วจนปถํ\csromansharedfootnote{6fb8195c2136ae53}{ตว วจนปถํ (สฺยา)} นิสามยิตฺวา,}{ปฏิปชฺเชม หิตาย โน สิยา.} \\ \csromangathacontline{เวสาลิยํ ราชคเห, สาวตฺถิยํ จ อาฬวิยํ.} \\ \csromangathacontline{โกสมฺพิยํ จ สกฺเกสุ, ภคฺเคสุ เจว ปญฺญตฺตา.} \\ \csromangathacontline{กติ เวสาลิยํ ปญฺญตฺตา, กติ ราชคเห กตา.} \\ \csromangathacontline{สาวตฺถิยํ กติ โหนฺติ, กติ อาฬวิยํ กตา.} \\ \csromangathacontline{กติ โกสมฺพิยํ ปญฺญตฺตา, กติ สกฺเกสุ วุจฺจนฺติ.} \\ \csromangathacontline{กติ ภคฺเคสุ ปญฺญตฺตา, ตํ เม อกฺขาหิ ปุจฺฉิโต.}}\csromangathastanza{\csromanfolio{256}\csromangathabat{ทส เวสาลิยํ ปญฺญตฺตา,}{เอกวีส ราชคเห กตา.} \\ \csromangathabat{ฉอูน ตีณิ สตานิ,}{สพฺเพ สาวตฺถิยํ กตา.}}\csromangathastanza{\csromangathabat{ฉ อาฬวิยํ ปญฺญตฺตา,}{อฏฺฐ โกสมฺพิยํ กตา.} \\ \csromangathabat{อฏฺฐ สกฺเกสุ วุจฺจนฺติ,}{ตโย ภคฺเคสุ ปญฺญตฺตา.}}\csromangathastanza{\csromangathabat{เย เวสาลิยํ ปญฺญตฺตา,}{เต สุโณหิ ยถาตถํ\csromansharedfootnote{0b522285ac617f96}{ยถากถํ (สี, สฺยา, เอวมุปริปิ)}.} \\ \csromangathabat{เมถุนวิคฺคหุตฺตริ,}{อติเรกญฺจ กาฬกํ.}}\csromangathastanza{\csromangathabat{ภูตํ ปรมฺปรภตฺตํ,}{ทนฺตโปเนน\csromansharedfootnote{676fccc5251f2f5a}{ทนฺตโปเณน (ก)} อเจลโก.} \\ \csromangathabat{ภิกฺขุนีสุ จ อกฺโกโส,}{ทเสเต เวสาลิยํ กตา.}}\csromangathastanza{\csromangathabat{เย ราชคเห ปญฺญตฺตา,}{เต สุโณหิ ยถาตถํ.} \\ อทินฺนาทานํ ราชคเห, \\ เทฺว อนุทฺธํสนา เทฺวปิ จ เภทา.}\csromangathastanza{\csromangathabat{อนฺตรวาสกํ รูปิยํ สุตฺตํ,}{อุชฺฌาปเนน จ ปาจิตปิณฺฑํ.} \\ \csromangathabat{คณโภชนํ วิกาเล จ,}{จาริตฺตํ นหานํ อูนวีสติ.}}\csromangathastanza{\csromangathabat{จีวรํ ทตฺวา โวสาสนฺติ,}{เอเต ราชคเห กตา.} \\ \csromangathabat{คิรคฺคจริยา ตตฺเถว,}{ฉนฺททาเนน เอกวีสติ.}}\csromangathastanza{\csromangathabat{เย สาวตฺถิยํ ปญฺญตฺตา,}{เต สุโณหิ ยถาตถํ.} \\ \csromangathabat{ปาราชิกานิ จตฺตาริ,}{สํฆาทิเสสา ภวนฺติ โสฬส.}}\csromangathastanza{\csromangathabat{อนิยตา จ เทฺว โหนฺติ,}{นิสฺสคฺคิยา จตุวีสติ.} \\ \csromangathabat{ฉปญฺญาสสตญฺเจว,}{ขุทฺทกานิ ปวุจฺจนฺติ.}}\csromangathastanza{\csromangathabat{ทสเยว จ คารยฺหา,}{เทฺวสตฺตติ จ เสขิยา.} \\ \csromangathabat{ฉอูน ตีณิสตานิ,}{สพฺเพ สาวตฺถิยํ กตา.}}\csromangathastanza{\csromangathabat{เย อาฬวิยํ ปญฺญตฺตา,}{เต สุโณหิ ยถาตถํ.} \\ \csromangathabat{กุฏิโกสิยเสยฺยา จ,}{ขณเน คจฺฉ เทวเต.}}\csromangathastanza{\csromangathabat{สปฺปาณกญฺจ สิญฺจนฺติ,}{ฉ เอเต อาฬวิยํ กตา.} \\ \csromangathabat{เย โกสมฺพิยํ ปญฺญตฺตา,}{เต สุโณหิ ยถาตถํ.}}\csromangathastanza{\csromangathabat{มหาวิหาโร โทวจสฺสํ,}{อญฺญํ ทฺวารํ สุราย จ.} \\ \csromangathabat{อนาทริยํ สหธมฺโม,}{ปโยปาเนน อฏฺฐมํ.}}\csromangathastanza{\csromanfolio{257}\csromangathabat{เย สกฺเกสุ ปญฺญตฺตา,}{เต สุโณหิ ยถาตถํ.} \\ \csromangathabat{เอฬกโลมานิ ปตฺโต จ,}{โอวาโท เจว เภสชฺชํ.}}\csromangathastanza{\csromangathabat{สูจิ อารญฺญิโก เจว,}{อฏฺเฐเต\csromansharedfootnote{88262a876b82bea0}{ฉ เอเต (สพฺพตฺถ)} กาปิลวตฺถเว.} \\ \csromangathabat{อุทกสุทฺธิยา โอวาโท,}{ภิกฺขุนีสุ ปวุจฺจนฺติ.}}\csromangathastanza{\csromangathabat{เย ภคฺเคสุ ปญฺญตฺตา,}{เต สุโณหิ ยถาตถํ.} \\ \csromangathabat{สมาทหิตฺวา วิสิพฺเพนฺติ,}{สามิเสน สสิตฺถกํ.}}\csromangathastanza{\csromangathabat{ปาราชิกานิ จตฺตาริ,}{สํฆาทิเสสานิ ภวนฺติ.} \\ \csromangathabat{สตฺต จ นิสฺสคฺคิยานิ,}{อฏฺฐ ทฺวตฺติํส ขุทฺทกา.}}\csromangathastanza{\csromangathabat{เทฺว คารยฺหา ตโย เสกฺขา,}{ฉปฺปญฺญาส สิกฺขาปทา.} \\ \csromangathabat{ฉสุ นคเรสุ ปญฺญตฺตา,}{พุทฺเธนา\-ทิจฺจพนฺธุนา.}}\csromangathastanza{\csromangathabat{ฉอูน ตีณิ สตานิ,}{สพฺเพ สาวตฺถิยํ กตา.} \\ \csromangathabat{การุณิเกน พุทฺเธน,}{โคตเมน ยสสฺสินา.}}}

\csromanmatikaanchor{mtk.181}
\csromantocmarkref{subparagraph}{2. จตุวิปตฺติ}{mtk.181}
\bookmark[dest={mtk.181},level=5]{2. จตุวิปตฺติ}
\csromanheader{5}{2. \textbf{จตุวิปตฺติ}}
\setlength{\csromangathapagewidth}{0pt}
\setlength{\csromangathawakwidth}{0pt}
\csromangathameasuregroup{ยํ ตํ\textsuperscript{1} ปุจฺฉิมฺห อกิตฺตยิ โน, \\ อญฺญํ ตํ ปุจฺฉามิ, ตทิงฺฆ พฺรูหิ, \\ อนวเสสํ ทุฏฺฐุลฺลญฺจ อทุฏฺฐุลฺลํ, \\ สพฺพานิเปตานิ วิยากโรหิ, \\ ถุลฺลจฺจยํ ปาจิตฺติยา, \\ ทุพฺภาสิตํ โย จายํ,}{\csromangathabat{ยํ ตํ\textsuperscript{1} ปุจฺฉิมฺห อกิตฺตยิ โน,}{ตํ ตํ พฺยากตํ อนญฺญถา.} \\ \csromangathabat{อญฺญํ ตํ ปุจฺฉามิ, ตทิงฺฆ พฺรูหิ,}{ครุก ลหุกญฺจาปิ สาวเสสํ.} \\ \csromangathabat{อนวเสสํ ทุฏฺฐุลฺลญฺจ อทุฏฺฐุลฺลํ,}{เย จ ยาวตติยกา.} \\ สาธารณํ อสาธารณํ, \\ วิภตฺติโย จ\textsuperscript{1} เยหิ สมเถหิ สมฺมนฺติ. \\ \csromangathabat{สพฺพานิเปตานิ วิยากโรหิ,}{หนฺท วากฺยํ สุโณม เต.} \\ เอกติํสา เย ครุกา, อฏฺเฐตฺถ อนวเสสา. \\ เย ครุกา เต ทุฏฺฐุลฺลา, เย ทุฏฺฐุลฺลา สา สีลวิปตฺติ. \\ ปาราชิกํ สํฆาทิเสโส, “สีลวิปตฺตี”ติ วุจฺจติ. \\ \csromangathabat{ถุลฺลจฺจยํ ปาจิตฺติยา,}{ปาฏิเทสนียํ ทุกฺกฏํ.} \\ \csromangathabat{ทุพฺภาสิตํ โย จายํ,}{อกฺโกสติ หสาธิปฺปาโย.}}
\csromangathasetleft{ยํ ตํ\textsuperscript{1} ปุจฺฉิมฺห อกิตฺตยิ โน, \\ อญฺญํ ตํ ปุจฺฉามิ, ตทิงฺฆ พฺรูหิ, \\ อนวเสสํ ทุฏฺฐุลฺลญฺจ อทุฏฺฐุลฺลํ, \\ สพฺพานิเปตานิ วิยากโรหิ, \\ ถุลฺลจฺจยํ ปาจิตฺติยา, \\ ทุพฺภาสิตํ โย จายํ,}
\csromangathagroup{\csromangathastanza{\csromangathaitembat{336}{ยํ ตํ\csromansharedfootnote{fe3336ec2c3d2573}{ยํ ยํ (ก)} ปุจฺฉิมฺห อกิตฺตยิ โน,}{ตํ ตํ พฺยากตํ อนญฺญถา.} \\ \csromangathacontbat{อญฺญํ ตํ ปุจฺฉามิ, ตทิงฺฆ พฺรูหิ,}{ครุก ลหุกญฺจาปิ สาวเสสํ.} \\ \csromangathacontbat{อนวเสสํ ทุฏฺฐุลฺลญฺจ อทุฏฺฐุลฺลํ,}{เย จ ยาวตติยกา.} \\ \csromangathacontline{สาธารณํ อสาธารณํ,} \\ \csromangathacontline{วิภตฺติโย จ\csromansharedfootnote{fed2c465ca0be845}{วิปตฺติโย จ (สี, สฺยา)} เยหิ สมเถหิ สมฺมนฺติ.} \\ \csromangathacontbat{สพฺพานิเปตานิ วิยากโรหิ,}{หนฺท วากฺยํ สุโณม เต.} \\ \csromangathacontline{เอกติํสา เย ครุกา, อฏฺเฐตฺถ อนวเสสา.} \\ \csromangathacontline{เย ครุกา เต ทุฏฺฐุลฺลา, เย ทุฏฺฐุลฺลา สา สีลวิปตฺติ.} \\ \csromangathacontline{ปาราชิกํ สํฆาทิเสโส, “สีลวิปตฺตี”ติ วุจฺจติ.}}\csromangathastanza{\csromanfolio{258}\csromangathabat{ถุลฺลจฺจยํ ปาจิตฺติยา,}{ปาฏิเทสนียํ ทุกฺกฏํ.} \\ \csromangathabat{ทุพฺภาสิตํ โย จายํ,}{อกฺโกสติ หสาธิปฺปาโย.}}}

\prose{อยํ สา อาจารวิปตฺติ\-สมฺมตา.}

\setlength{\csromangathapagewidth}{0pt}
\setlength{\csromangathawakwidth}{0pt}
\csromangathameasuregroup{วิปรีตทิฏฺฐิํ คณฺหนฺติ, \\ อพฺภาจิกฺขนฺติ สมฺพุทฺธํ,}{\csromangathabat{วิปรีตทิฏฺฐิํ คณฺหนฺติ,}{อสทฺธมฺเมหิ ปุรกฺขตา.} \\ \csromangathabat{อพฺภาจิกฺขนฺติ สมฺพุทฺธํ,}{ทุปฺปญฺญา โมหปารุตา.}}
\csromangathasetleft{วิปรีตทิฏฺฐิํ คณฺหนฺติ, \\ อพฺภาจิกฺขนฺติ สมฺพุทฺธํ,}
\csromangathagroup{\csromangathastanza{\csromangathabat{วิปรีตทิฏฺฐิํ คณฺหนฺติ,}{อสทฺธมฺเมหิ ปุรกฺขตา.} \\ \csromangathabat{อพฺภาจิกฺขนฺติ สมฺพุทฺธํ,}{ทุปฺปญฺญา โมหปารุตา.}}}

\prose{อยํ สา ทิฏฺฐิวิปตฺติ\-สมฺมตา.}

\prose{อาชีวเหตุ อาชีวการณา ปาปิจฺโฉ อิจฺฉาปกโต อสนฺตํ อภูตํ อุตฺตริ\-มนุสฺสธมฺมํ อุลฺลปติ. อาชีวเหตุ อาชีวการณา สญฺจริตฺตํ สมาปชฺชติ. อาชีวเหตุ อาชีวการณา “โย เต วิหาเร วสติ, โส ภิกฺขุ อรหา”ติ ภณติ. อาชีวเหตุ อาชีวการณา ภิกฺขุ ปณีต\-โภชนานิ อตฺตโน อตฺถาย วิญฺญาเปตฺวา ภุญฺชติ. อาชิวเหตุ อาชีวการณา ภิกฺขุนี ปณีต\-โภชนานิ อตฺตโน อตฺถาย วิญฺญาเปตฺวา ภุญฺชติ. อาชีวเหตุ อาชีวการณา สูปํ วา โอทนํ วา อคิลาโน อตฺตโน อตฺถาย วิญฺญาเปตฺวา ภุญฺชติ. อยํ สา อาชีววิปตฺติ สมฺมตา.}

\setlength{\csromangathapagewidth}{0pt}
\setlength{\csromangathawakwidth}{0pt}
\csromangathameasuregroup{เอกาทส ยาวตติยกา, \\ อุกฺขิตฺตานุวตฺติกา, \\ อริฏฺโฐ จณฺฑกาฬี จ,}{\csromangathabat{เอกาทส ยาวตติยกา,}{เต สุโณหิ ยถาตถํ.} \\ \csromangathabat{อุกฺขิตฺตานุวตฺติกา,}{อฏฺฐ ยาวตติยกา.} \\ \csromangathabat{อริฏฺโฐ จณฺฑกาฬี จ,}{อิเม เต ยาวตติยกา.}}
\csromangathasetleft{เอกาทส ยาวตติยกา, \\ อุกฺขิตฺตานุวตฺติกา, \\ อริฏฺโฐ จณฺฑกาฬี จ,}
\csromangathagroup{\csromangathastanza{\csromangathabat{เอกาทส ยาวตติยกา,}{เต สุโณหิ ยถาตถํ.} \\ \csromangathabat{อุกฺขิตฺตานุวตฺติกา,}{อฏฺฐ ยาวตติยกา.} \\ \csromangathabat{อริฏฺโฐ จณฺฑกาฬี จ,}{อิเม เต ยาวตติยกา.}}}

\csromanmatikaanchor{mtk.182}
\csromantocmarkref{subparagraph}{3. เฉทนกาทิ}{mtk.182}
\bookmark[dest={mtk.182},level=5]{3. เฉทนกาทิ}
\csromanheader{5}{3. \textbf{เฉทนกาทิ}}
\proseitem{337}{กติ เฉทนกานิ. กติ เภทนกานิ. กติ อุทฺทาลนกานิ. กติ อนญฺญ\-ปาจิตฺติยานิ. กติ ภิกฺขุ\-สมฺมุติโย. กติ สามีจิโย. กติ ปรมานิ.}

\setlength{\csromangathapagewidth}{0pt}
\setlength{\csromangathawakwidth}{0pt}
\csromangathameasuregroup{กติ ชานนฺติ ปญฺญตฺตา,}{\csromangathabat{กติ ชานนฺติ ปญฺญตฺตา,}{พุทฺเธนาทิจฺจพนฺธุนา.}}
\csromangathasetleft{กติ ชานนฺติ ปญฺญตฺตา,}
\csromangathagroup{\csromangathastanza{\csromangathabat{กติ ชานนฺติ ปญฺญตฺตา,}{พุทฺเธนา\-ทิจฺจพนฺธุนา.}}}

\prose{ฉ เฉทนกานิ. เอกํ เภทนกํ. เอกํ อุทฺทาลนกํ. จตฺตาริ อนญฺญ\-ปาจิตฺติยานิ. จตสฺโส ภิกฺขุ\-สมฺมุติโย. สตฺต สามีจิโย. จุทฺทส ปรมานิ.}

\setlength{\csromangathapagewidth}{0pt}
\setlength{\csromangathawakwidth}{0pt}
\csromangathameasuregroup{โสทส\textsuperscript{1} ชานนฺติ ปญฺญตฺตา,}{\csromangathabat{โสทส\textsuperscript{1} ชานนฺติ ปญฺญตฺตา,}{พุทฺเธนาทิจฺจพนฺธุนา.}}
\csromangathasetleft{โสทส\textsuperscript{1} ชานนฺติ ปญฺญตฺตา,}
\csromangathagroup{\csromangathastanza{\csromangathabat{โสทส\csromansharedfootnote{bed619f3af68800d}{โสฬส (สี, สฺยา) อฏฺฐกถา โอโลเกตพฺพา.} ชานนฺติ ปญฺญตฺตา,}{พุทฺเธนา\-ทิจฺจพนฺธุนา.}}}

\csromanmatikaanchor{mtk.183}
\csromantocmarkref{subparagraph}{4. อสาธารณาทิ}{mtk.183}
\bookmark[dest={mtk.183},level=5]{4. อสาธารณาทิ}
\csromanheader{5}{4. \textbf{อสาธารณาทิ}}
\setlength{\csromangathapagewidth}{0pt}
\setlength{\csromangathawakwidth}{0pt}
\csromangathameasuregroup{ฉ สํฆาทิเสสา, \\ นิสฺสคฺคิยานิ ทฺวาทส, \\ เทฺววีสติ ขุทฺทกา, \\ สตํ ติํสา จ ภิกฺขุนีนํ,}{วีสํ เทฺว สตานิ ภิกฺขูนํ สิกฺขาปทานิ, \\ อุทฺเทสํ อาคจฺฉนฺติ อุโปสเถสุ. \\ ตีณิ สตานิ จตฺตาริ ภิกฺขุนีนํ สิกฺขาปทานิ, \\ อุทฺเทสํ อาคจฺฉนฺติ อุโปสเถสุ. \\ ฉจตฺตารีสา ภิกฺขูนํ, ภิกฺขุนีหิ อสาธารณา. \\ สตํ ติํสา จ ภิกฺขุนีนํ, ภิกฺขูหิ อสาธารณา. \\ สตํ สตฺตติ ฉจฺเจว, อุภินฺนํ อสาธารณา. \\ สตํ สตฺตติ จตฺตาริ, อุภินฺนํ สมสิกฺขตา. \\ วีสํ เทฺวสตานิ ภิกฺขูนํ สิกฺขาปทานิ, \\ อุทฺเทสํ อาคจฺฉนฺติ อุโปสเถสุ. \\ เต สุโณหิ ยถาตถํ. \\ ปาราชิกานิ จตฺตาริ, สํฆาทิเสสานิ ภวนฺติ เตรส. \\ อนิยตา เทฺว โหนฺติ. \\ นิสฺสคฺคิยานิ ติํเสว, เทฺวนวุติ จ ขุทฺทกา. \\ จตฺตาโร ปาฏิเทสนียา, ปญฺจสตฺตติ เสขิยา. \\ วีสํ เทฺว สตานิ จิเม โหนฺติ ภิกฺขูนํ สิกฺขาปทานิ, \\ อุทฺเทสํ อาคจฺฉนฺติ อุโปสเถสุ. \\ ตีณิ สตานิ จตฺตาริ, ภิกฺขุนีนํ สิกฺขาปทานิ. \\ อุทฺเทสํ อาคจฺฉนฺติ อุโปสเถสุ, เต สุโณหิ ยถาตถํ. \\ ปาราชิกานิ อฏฺฐ, สํฆาทิเสสานิ ภวนฺติ สตฺตรส. \\ นิสฺสคฺคิยานิ ติํเสว, \\ สตํ สฏฺฐิ ฉ เจว ขุทฺทกานิ  ปวุจฺจนฺติ. \\ อฏฺฐ ปาฏิเทสนียา, ปญฺจสตฺตติ เสขิยา. \\ ตีณิ สตานิ จตฺตาริ จิเม โหนฺติ ภิกฺขุนีนํ สิกฺขาปทานิ, \\ อุทฺเทสํ อาคจฺฉนฺติ อุโปสเถสุ. \\ ฉจตฺตารีสา ภิกฺขูนํ, ภิกฺขุนีหิ อสาธารณา. \\ เต สุโณหิ ยถาตถํ. \\ \csromangathabat{ฉ สํฆาทิเสสา,}{เทฺว อนิยเตหิ อฏฺฐ.} \\ \csromangathabat{นิสฺสคฺคิยานิ ทฺวาทส,}{เตหิ เต โหนฺติ วีสติ.} \\ \csromangathabat{เทฺววีสติ ขุทฺทกา,}{จตุโร ปาฏิเทสนียา.} \\ ฉจตฺตารีสา จิเม โหนฺติ, \\ ภิกฺขูนํ ภิกฺขุนีหิ อสาธารณา. \\ \csromangathabat{สตํ ติํสา จ ภิกฺขุนีนํ,}{ภิกฺขูหิ อสาธารณา.}}
\csromangathasetleft{ฉ สํฆาทิเสสา, \\ นิสฺสคฺคิยานิ ทฺวาทส, \\ เทฺววีสติ ขุทฺทกา, \\ สตํ ติํสา จ ภิกฺขุนีนํ,}
\csromangathagroup{\csromangathastanza{\csromanfolio{259}\csromangathaitemline{338}{วีสํ เทฺว สตานิ ภิกฺขูนํ สิกฺขาปทานิ,} \\ \csromangathacontline{อุทฺเทสํ อาคจฺฉนฺติ อุโปสเถสุ.} \\ \csromangathacontline{ตีณิ สตานิ จตฺตาริ ภิกฺขุนีนํ สิกฺขาปทานิ,} \\ \csromangathacontline{อุทฺเทสํ อาคจฺฉนฺติ อุโปสเถสุ.} \\ \csromangathacontline{ฉจตฺตารีสา ภิกฺขูนํ, ภิกฺขุนีหิ อสาธารณา.} \\ \csromangathacontline{สตํ ติํสา จ ภิกฺขุนีนํ, ภิกฺขูหิ อสาธารณา.} \\ \csromangathacontline{สตํ สตฺตติ ฉจฺเจว, อุภินฺนํ อสาธารณา.} \\ \csromangathacontline{สตํ สตฺตติ จตฺตาริ, อุภินฺนํ สมสิกฺขตา.} \\ \csromangathacontline{วีสํ เทฺวสตานิ ภิกฺขูนํ สิกฺขาปทานิ,} \\ \csromangathacontline{อุทฺเทสํ อาคจฺฉนฺติ อุโปสเถสุ.} \\ \csromangathacontline{เต สุโณหิ ยถาตถํ.} \\ \csromangathacontline{ปาราชิกานิ จตฺตาริ, สํฆาทิเสสานิ ภวนฺติ เตรส.} \\ \csromangathacontline{อนิยตา เทฺว โหนฺติ.} \\ \csromangathacontline{นิสฺสคฺคิยานิ ติํเสว, เทฺวนวุติ จ ขุทฺทกา.} \\ \csromangathacontline{จตฺตาโร ปาฏิเทสนียา, ปญฺจสตฺตติ เสขิยา.} \\ \csromangathacontline{วีสํ เทฺว สตานิ จิเม โหนฺติ ภิกฺขูนํ สิกฺขาปทานิ,} \\ \csromangathacontline{อุทฺเทสํ อาคจฺฉนฺติ อุโปสเถสุ.} \\ \csromangathacontline{ตีณิ สตานิ จตฺตาริ, ภิกฺขุนีนํ สิกฺขาปทานิ.} \\ \csromangathacontline{อุทฺเทสํ อาคจฺฉนฺติ อุโปสเถสุ, เต สุโณหิ ยถาตถํ.} \\ \csromangathacontline{ปาราชิกานิ อฏฺฐ, สํฆาทิเสสานิ ภวนฺติ สตฺตรส.} \\ \csromangathacontline{นิสฺสคฺคิยานิ ติํเสว,} \\ \csromangathacontline{สตํ สฏฺฐิ ฉ เจว ขุทฺทกานิ  ปวุจฺจนฺติ.} \\ \csromangathacontline{อฏฺฐ ปาฏิเทสนียา, ปญฺจสตฺตติ เสขิยา.} \\ \csromangathacontline{ตีณิ สตานิ จตฺตาริ จิเม โหนฺติ ภิกฺขุนีนํ สิกฺขาปทานิ,} \\ \csromangathacontline{อุทฺเทสํ อาคจฺฉนฺติ อุโปสเถสุ.} \\ \csromangathacontline{ฉจตฺตารีสา ภิกฺขูนํ, ภิกฺขุนีหิ อสาธารณา.} \\ \csromangathacontline{เต สุโณหิ ยถาตถํ.}}\csromangathastanza{\csromanfolio{260}\csromangathabat{ฉ สํฆาทิเสสา,}{เทฺว อนิยเตหิ อฏฺฐ.} \\ \csromangathabat{นิสฺสคฺคิยานิ ทฺวาทส,}{เตหิ เต โหนฺติ วีสติ.}}\csromangathastanza{\csromangathabat{เทฺววีสติ ขุทฺทกา,}{จตุโร ปาฏิเทสนียา.} \\ ฉจตฺตารีสา จิเม โหนฺติ, \\ ภิกฺขูนํ ภิกฺขุนีหิ อสาธารณา. \\ \csromangathabat{สตํ ติํสา จ ภิกฺขุนีนํ,}{ภิกฺขูหิ อสาธารณา.}}}

\prose{เต สุโณหิ ยถาตถํ.}

\setlength{\csromangathapagewidth}{0pt}
\setlength{\csromangathawakwidth}{0pt}
\csromangathameasuregroup{ปาราชิกานิ จตฺตาริ, \\ นิสฺสคฺคิยานิ ทฺวาทส,}{\csromangathabat{ปาราชิกานิ จตฺตาริ,}{สํฆมฺหา ทส นิสฺสเร.} \\ \csromangathabat{นิสฺสคฺคิยานิ ทฺวาทส,}{ฉนฺนวุติ จ ขุทฺทกา.}}
\csromangathasetleft{ปาราชิกานิ จตฺตาริ, \\ นิสฺสคฺคิยานิ ทฺวาทส,}
\csromangathagroup{\csromangathastanza{\csromangathabat{ปาราชิกานิ จตฺตาริ,}{สํฆมฺหา ทส นิสฺสเร.} \\ \csromangathabat{นิสฺสคฺคิยานิ ทฺวาทส,}{ฉนฺนวุติ จ ขุทฺทกา.}}}

\prose{อฏฺฐ ปาฏิเทสนียา.}

\setlength{\csromangathapagewidth}{0pt}
\setlength{\csromangathawakwidth}{0pt}
\csromangathameasuregroup{สตํ ติํสา จิเม โหนฺติ, \\ สตํ สตฺตติ ฉจฺเจว,}{\csromangathabat{สตํ ติํสา จิเม โหนฺติ,}{ภิกฺขุนีนํ ภิกฺขูหิ อสาธารณา.} \\ \csromangathabat{สตํ สตฺตติ ฉจฺเจว,}{อุภินฺนํ อสาธารณา.}}
\csromangathasetleft{สตํ ติํสา จิเม โหนฺติ, \\ สตํ สตฺตติ ฉจฺเจว,}
\csromangathagroup{\csromangathastanza{\csromangathabat{สตํ ติํสา จิเม โหนฺติ,}{ภิกฺขุนีนํ ภิกฺขูหิ อสาธารณา.} \\ \csromangathabat{สตํ สตฺตติ ฉจฺเจว,}{อุภินฺนํ อสาธารณา.}}}

\prose{เต สุโณหิ ยถาตถํ.}

\setlength{\csromangathapagewidth}{0pt}
\setlength{\csromangathawakwidth}{0pt}
\csromangathameasuregroup{ปาราชิกานิ จตฺตาริ, \\ อนิยตา เทฺว โหนฺติ, \\ สตํ อฏฺฐารสา เจว,}{\csromangathabat{ปาราชิกานิ จตฺตาริ,}{สํฆาทิเสสานิ ภวนฺติ โสฬส.} \\ \csromangathabat{อนิยตา เทฺว โหนฺติ,}{นิสฺสคฺคิยานิ จตุวีสติ.} \\ \csromangathabat{สตํ อฏฺฐารสา เจว,}{ขุทฺทกานิ ปวุจฺจนฺติ.}}
\csromangathasetleft{ปาราชิกานิ จตฺตาริ, \\ อนิยตา เทฺว โหนฺติ, \\ สตํ อฏฺฐารสา เจว,}
\csromangathagroup{\csromangathastanza{\csromangathabat{ปาราชิกานิ จตฺตาริ,}{สํฆาทิเสสานิ ภวนฺติ โสฬส.} \\ \csromangathabat{อนิยตา เทฺว โหนฺติ,}{นิสฺสคฺคิยานิ จตุวีสติ.} \\ \csromangathabat{สตํ อฏฺฐารสา เจว,}{ขุทฺทกานิ ปวุจฺจนฺติ.}}}

\prose{ทฺวาทส ปาฏิเทสนียา.}

\setlength{\csromangathapagewidth}{0pt}
\setlength{\csromangathawakwidth}{0pt}
\csromangathameasuregroup{สตํ สตฺตติ ฉจฺเจวิเม โหนฺติ, \\ สตํ สตฺตติ จตฺตาริ,}{\csromangathabat{สตํ สตฺตติ ฉจฺเจวิเม โหนฺติ,}{อุภินฺนํ อสาธารณา.} \\ \csromangathabat{สตํ สตฺตติ จตฺตาริ,}{อุภินฺนํ สมสิกฺขตา.}}
\csromangathasetleft{สตํ สตฺตติ ฉจฺเจวิเม โหนฺติ, \\ สตํ สตฺตติ จตฺตาริ,}
\csromangathagroup{\csromangathastanza{\csromangathabat{สตํ สตฺตติ ฉจฺเจวิเม โหนฺติ,}{อุภินฺนํ อสาธารณา.} \\ \csromangathabat{สตํ สตฺตติ จตฺตาริ,}{อุภินฺนํ สมสิกฺขตา.}}}

\prose{เต สุโณหิ ยถาตถํ.}

\setlength{\csromangathapagewidth}{0pt}
\setlength{\csromangathawakwidth}{0pt}
\csromangathameasuregroup{ปาราชิกานิ จตฺตาริ, \\ นิสฺสคฺคิยานิ อฏฺฐรส,}{\csromangathabat{ปาราชิกานิ จตฺตาริ,}{สํฆาทิเสสานิ ภวนฺติ สตฺต.} \\ \csromangathabat{นิสฺสคฺคิยานิ อฏฺฐรส,}{สมสตฺตติ ขุทฺทกา.}}
\csromangathasetleft{ปาราชิกานิ จตฺตาริ, \\ นิสฺสคฺคิยานิ อฏฺฐรส,}
\csromangathagroup{\csromangathastanza{\csromangathabat{ปาราชิกานิ จตฺตาริ,}{สํฆาทิเสสานิ ภวนฺติ สตฺต.} \\ \csromangathabat{นิสฺสคฺคิยานิ อฏฺฐรส,}{สมสตฺตติ ขุทฺทกา.}}}

\prose{ปญฺจสตฺตติ เสขิยานิ.}

\setlength{\csromangathapagewidth}{0pt}
\setlength{\csromangathawakwidth}{0pt}
\csromangathameasuregroup{สตํ สตฺตติ จตฺตาริ จิเม โหนฺติ, \\ อฏฺเฐว ปาราชิกา เย ทุราสทา, \\ ปณฺฑุปลาโส ปุถุสิลา, \\ ตาโลว มตฺถกจฺฉินฺโน, \\ เตวีสติ สํฆาทิเสสา,}{\csromangathabat{สตํ สตฺตติ จตฺตาริ จิเม โหนฺติ,}{อุภินฺนํ สมสิกฺขตา.} \\ \csromangathabat{อฏฺเฐว ปาราชิกา เย ทุราสทา,}{ตาลวตฺถุสมูปมา.} \\ \csromangathabat{ปณฺฑุปลาโส ปุถุสิลา,}{สีสจฺฉินฺโนว โส นโร.} \\ \csromangathabat{ตาโลว มตฺถกจฺฉินฺโน,}{อวิรุฬฺหี ภวนฺติ เต.} \\ \csromangathabat{เตวีสติ สํฆาทิเสสา,}{เทฺว อนิยตา.}}
\csromangathasetleft{สตํ สตฺตติ จตฺตาริ จิเม โหนฺติ, \\ อฏฺเฐว ปาราชิกา เย ทุราสทา, \\ ปณฺฑุปลาโส ปุถุสิลา, \\ ตาโลว มตฺถกจฺฉินฺโน, \\ เตวีสติ สํฆาทิเสสา,}
\csromangathagroup{\csromangathastanza{\csromangathabat{สตํ สตฺตติ จตฺตาริ จิเม โหนฺติ,}{อุภินฺนํ สมสิกฺขตา.} \\ \csromangathabat{อฏฺเฐว ปาราชิกา เย ทุราสทา,}{ตาลวตฺถุสมูปมา.}}\csromangathastanza{\csromangathabat{ปณฺฑุปลาโส ปุถุสิลา,}{สีสจฺฉินฺโนว โส นโร.} \\ \csromangathabat{ตาโลว มตฺถก\-จฺฉินฺโน,}{อวิรุฬฺหี ภวนฺติ เต.} \\ \csromangathabat{เตวีสติ สํฆาทิเสสา,}{เทฺว อนิยตา.}}}

\prose{\csromanfolio{261}เทฺว จตฺตารีส นิสฺสคฺคิยา.}

\setlength{\csromangathapagewidth}{0pt}
\setlength{\csromangathawakwidth}{0pt}
\csromangathameasuregroup{อฏฺฐาสีติสตํ ปาจิตฺติยา, \\ ปญฺจสตฺตติ เสขิยา, \\ สมฺมุขา จ ปฏิญฺญาย, \\ เทฺว อุโปสถา เทฺว ปวารณา,}{\csromangathabat{อฏฺฐาสีติสตํ ปาจิตฺติยา,}{ทฺวาทส ปาฏิเทสนียา.} \\ \csromangathabat{ปญฺจสตฺตติ เสขิยา,}{ตีหิ สมเถหิ สมฺมนฺติ.} \\ \csromangathabat{สมฺมุขา จ ปฏิญฺญาย,}{ติณวตฺถารเกน จ.} \\ \csromangathabat{เทฺว อุโปสถา เทฺว ปวารณา,}{จตฺตาริ กมฺมานิ ชิเนน เทสิตา.} \\ ปญฺเจว อุทฺเทสา จตุโร ภวนฺติ, \\ อนญฺญถา อาปตฺติกฺขนฺธา จ ภวนฺติ สตฺต. \\ อธิกรณานิ จตฺตาริ สตฺตหิ สมเถหิ สมฺมนฺติ. \\ ทฺวีหิ จตูหิ ตีหิ กิจฺจํ เอเกน สมฺมติ.}
\csromangathasetleft{อฏฺฐาสีติสตํ ปาจิตฺติยา, \\ ปญฺจสตฺตติ เสขิยา, \\ สมฺมุขา จ ปฏิญฺญาย, \\ เทฺว อุโปสถา เทฺว ปวารณา,}
\csromangathagroup{\csromangathastanza{\csromangathabat{อฏฺฐาสีติสตํ ปาจิตฺติยา,}{ทฺวาทส ปาฏิเทสนียา.} \\ \csromangathabat{ปญฺจสตฺตติ เสขิยา,}{ตีหิ สมเถหิ สมฺมนฺติ.}}\csromangathastanza{\csromangathabat{สมฺมุขา จ ปฏิญฺญาย,}{ติณวตฺถารเกน จ.} \\ \csromangathabat{เทฺว อุโปสถา เทฺว ปวารณา,}{จตฺตาริ กมฺมานิ ชิเนน เทสิตา.}}\csromangathastanza{ปญฺเจว อุทฺเทสา จตุโร ภวนฺติ, \\ อนญฺญถา อาปตฺติกฺขนฺธา จ ภวนฺติ สตฺต. \\ อธิกรณานิ จตฺตาริ สตฺตหิ สมเถหิ สมฺมนฺติ. \\ ทฺวีหิ จตูหิ ตีหิ กิจฺจํ เอเกน สมฺมติ.}}

\csromanmatikaanchor{mtk.184}
\csromantocmarkref{subparagraph}{5. ปาราชิกาทิอาปตฺติ}{mtk.184}
\bookmark[dest={mtk.184},level=5]{5. ปาราชิกาทิอาปตฺติ}
\csromanheader{5}{5. ปาราชิกา\-ทิ\-อาปตฺติ}
\setlength{\csromangathapagewidth}{0pt}
\setlength{\csromangathawakwidth}{0pt}
\csromangathameasuregroup{‘ปาราชิกนฺ’ติ ยํ วุตฺตํ, \\ จุโต ปรทฺโธ ภฏฺโฐ จ, \\ สํวาโสปิ ตหิํ นตฺถิ, \\ ‘สํฆาทิเสโส’ติ ยํ วุตฺตํ, \\ สํโฆว เทติ ปริวาสํ, \\ มานตฺตํ เทติ อพฺเภติ, \\ ‘อนิยโต’ติ ยํ วุตฺตํ, \\ อนิยโต น นิยโต, \\ ติณฺณมญฺญตรํ ฐานํ, \\ ‘ถุลฺลจฺจยนฺ’ติ ยํ วุตฺตํ, \\ เอกสฺส มูเล โย เทเสติ, \\ อจฺจโย เตน สโม นตฺถิ, \\ ‘นิสฺสคฺคิยนฺ’ติ ยํ วุตฺตํ, \\ สํฆมชฺเฌ คณมชฺเฌ, \\ นิสฺสชฺชิตฺวาน เทเสติ, \\ ‘ปาจิตฺติยนฺ’ติ ยํ วุตฺตํ, \\ ปาเตติ กุสลํ ธมฺมํ, \\ จิตฺตสํโมหนฏฺฐานํ, \\ ‘ปาฏิเทสนียนฺ’ติ ยํ วุตฺตํ, \\ ภิกฺขุ อญฺญาตโก สนฺโต, \\ สามํ คเหตฺวา ภุญฺเชยฺย, \\ นิมนฺตนาสุ ภุญฺชนฺตา ฉนฺทาย, \\ อนิวาเรตฺวา ตหิํ ภุญฺเช, \\ สทฺธาจิตฺตํ กุลํ คนฺตฺวา, \\ อคิลาโน ตหิํ ภุญฺเช, \\ โย เจ อรญฺเญ วิหรนฺโต, \\ อวิทิตํ ตหิํ ภุญฺเช, \\ ภิกฺขุนี อญฺญาติกา สนฺตา, \\ สปฺปิ เตลํ มธุํ ผาณิตํ, \\ ทธิํ สยํ วิญฺญาเปยฺย ภิกฺขุนี, \\ ‘ทุกฺกฏนฺ’ติ ยํ วุตฺตํ, \\ อปรทฺธํ วิรทฺธญฺจ, \\ ยํ มนุสฺโส กเร ปาปํ, \\ ‘ทุกฺกฏนฺ’ติ ปเวเทนฺติ, \\ ‘ทุพฺภาสิตนฺ’ติ ยํ วุตฺตํ, \\ ทุพฺภาสิตํ ทุราภฏฺฐํ, \\ ยญฺจ วิญฺญู ครหนฺติ, \\ ‘เสขิยนฺ’ติ ยํ วุตฺตํ, \\ เสกฺขสฺส สิกฺขมานสฺส, \\ อาทิ เจตํ จรณญฺจ, \\ สิกฺขา เอตาทิสี นตฺถิ, \\ \textsuperscript{*}ฉนฺนมติวสฺสติ, \\ ตสฺมา ฉนฺนํ วิวเรถ, \\ คติ มิคานํ ปวนํ,}{\csromangathabat{‘ปาราชิกนฺ’ติ ยํ วุตฺตํ,}{ตํ สุโณหิ ยถาตถํ.} \\ \csromangathabat{จุโต ปรทฺโธ ภฏฺโฐ จ,}{สทฺธมฺมา หิ นิรงฺกโต.} \\ \csromangathabat{สํวาโสปิ ตหิํ นตฺถิ,}{เตเนตํ อิติ วุจฺจติ.} \\ \csromangathabat{‘สํฆาทิเสโส’ติ ยํ วุตฺตํ,}{ตํ สุโณหิ ยถาตถํ.} \\ \csromangathabat{สํโฆว เทติ ปริวาสํ,}{มูลาย ปฏิกสฺสติ.} \\ \csromangathabat{มานตฺตํ เทติ อพฺเภติ,}{เตเนตํ อิติ วุจฺจติ.} \\ \csromangathabat{‘อนิยโต’ติ ยํ วุตฺตํ,}{ตํ สุโณหิ ยถาตถํ.} \\ \csromangathabat{อนิยโต น นิยโต,}{อเนกํสิกตํ ปทํ.} \\ \csromangathabat{ติณฺณมญฺญตรํ ฐานํ,}{‘อนิยโต’ติ ปวุจฺจติ.} \\ \csromangathabat{‘ถุลฺลจฺจยนฺ’ติ ยํ วุตฺตํ,}{ตํ สุโณหิ ยถาตถํ.} \\ \csromangathabat{เอกสฺส มูเล โย เทเสติ,}{โย จ ตํ ปฏิคณฺหติ.} \\ \csromangathabat{อจฺจโย เตน สโม นตฺถิ,}{เตเนตํ อิติ วุจฺจติ.} \\ \csromangathabat{‘นิสฺสคฺคิยนฺ’ติ ยํ วุตฺตํ,}{ตํ สุโณหิ ยถาตถํ.} \\ \csromangathabat{สํฆมชฺเฌ คณมชฺเฌ,}{เอกสฺเสว จ เอกโต.} \\ \csromangathabat{นิสฺสชฺชิตฺวาน เทเสติ,}{เตเนตํ อิติ วุจฺจติ.} \\ \csromangathabat{‘ปาจิตฺติยนฺ’ติ ยํ วุตฺตํ,}{ตํ สุโณหิ ยถาตถํ.} \\ \csromangathabat{ปาเตติ กุสลํ ธมฺมํ,}{อริยมคฺคํ อปรชฺฌติ.} \\ \csromangathabat{จิตฺตสํโมหนฏฺฐานํ,}{เตเนตํ อิติ วุจฺจติ.} \\ \csromangathabat{‘ปาฏิเทสนียนฺ’ติ ยํ วุตฺตํ,}{ตํ สุโณหิ ยถาตถํ.} \\ \csromangathabat{ภิกฺขุ อญฺญาตโก สนฺโต,}{กิจฺฉา ลทฺธาย โภชนํ.} \\ \csromangathabat{สามํ คเหตฺวา ภุญฺเชยฺย,}{‘คารยฺหนฺ’ติ ปวุจฺจติ.} \\ \csromangathabat{นิมนฺตนาสุ ภุญฺชนฺตา ฉนฺทาย,}{โวสาสติ ตตฺถ ภิกฺขุนิํ.} \\ \csromangathabat{อนิวาเรตฺวา ตหิํ ภุญฺเช,}{‘คารยฺหนฺ’ติ ปวุจฺจติ.} \\ \csromangathabat{สทฺธาจิตฺตํ กุลํ คนฺตฺวา,}{อปฺปโภคํ อนาฬิยํ\textsuperscript{1}.} \\ \csromangathabat{อคิลาโน ตหิํ ภุญฺเช,}{‘คารยฺหนฺ’ติ ปวุจฺจติ.} \\ \csromangathabat{โย เจ อรญฺเญ วิหรนฺโต,}{สาสงฺเก สภยานเก.} \\ \csromangathabat{อวิทิตํ ตหิํ ภุญฺเช,}{‘คารยฺหนฺ’ติ ปวุจฺจติ.} \\ \csromangathabat{ภิกฺขุนี อญฺญาติกา สนฺตา,}{ยํ ปเรสํ มมายิตํ.} \\ \csromangathabat{สปฺปิ เตลํ มธุํ ผาณิตํ,}{มจฺฉมํสํ อโถ ขีรํ.} \\ \csromangathabat{ทธิํ สยํ วิญฺญาเปยฺย ภิกฺขุนี,}{คารยฺหปตฺตา สุคตสฺส สาสเน.} \\ \csromangathabat{‘ทุกฺกฏนฺ’ติ ยํ วุตฺตํ,}{ตํ สุโณหิ ยถาตถํ.} \\ \csromangathabat{อปรทฺธํ วิรทฺธญฺจ,}{ขลิตํ ยญฺจ ทุกฺกฏํ.} \\ \csromangathabat{ยํ มนุสฺโส กเร ปาปํ,}{อาวิ วา ยทิ วา รโห.} \\ \csromangathabat{‘ทุกฺกฏนฺ’ติ ปเวเทนฺติ,}{เตเนตํ อิติ วุจฺจติ.} \\ \csromangathabat{‘ทุพฺภาสิตนฺ’ติ ยํ วุตฺตํ,}{ตํ สุโณหิ ยถาตถํ.} \\ \csromangathabat{ทุพฺภาสิตํ ทุราภฏฺฐํ,}{สํกิลิฏฺฐญฺจ ยํ ปทํ.} \\ \csromangathabat{ยญฺจ วิญฺญู ครหนฺติ,}{เตเนตํ อิติ วุจฺจติ.} \\ \csromangathabat{‘เสขิยนฺ’ติ ยํ วุตฺตํ,}{ตํ สุโณหิ ยถาตถํ.} \\ \csromangathabat{เสกฺขสฺส สิกฺขมานสฺส,}{อุชุมคฺคานุสาริโน.} \\ \csromangathabat{อาทิ เจตํ จรณญฺจ,}{มุขํ สญฺญมสํวโร.} \\ \csromangathabat{สิกฺขา เอตาทิสี นตฺถิ,}{เตเนตํ อิติ วุจฺจติ.} \\ \csromangathabat{\textsuperscript{*}ฉนฺนมติวสฺสติ,}{วิวฏํ นาติวสฺสติ.} \\ \csromangathabat{ตสฺมา ฉนฺนํ วิวเรถ,}{เอวํ ตํ นาติวสฺสติ.} \\ \csromangathabat{คติ มิคานํ ปวนํ,}{อากาโส ปกฺขินํ คติ.}}
\csromangathasetleft{‘ปาราชิกนฺ’ติ ยํ วุตฺตํ, \\ จุโต ปรทฺโธ ภฏฺโฐ จ, \\ สํวาโสปิ ตหิํ นตฺถิ, \\ ‘สํฆาทิเสโส’ติ ยํ วุตฺตํ, \\ สํโฆว เทติ ปริวาสํ, \\ มานตฺตํ เทติ อพฺเภติ, \\ ‘อนิยโต’ติ ยํ วุตฺตํ, \\ อนิยโต น นิยโต, \\ ติณฺณมญฺญตรํ ฐานํ, \\ ‘ถุลฺลจฺจยนฺ’ติ ยํ วุตฺตํ, \\ เอกสฺส มูเล โย เทเสติ, \\ อจฺจโย เตน สโม นตฺถิ, \\ ‘นิสฺสคฺคิยนฺ’ติ ยํ วุตฺตํ, \\ สํฆมชฺเฌ คณมชฺเฌ, \\ นิสฺสชฺชิตฺวาน เทเสติ, \\ ‘ปาจิตฺติยนฺ’ติ ยํ วุตฺตํ, \\ ปาเตติ กุสลํ ธมฺมํ, \\ จิตฺตสํโมหนฏฺฐานํ, \\ ‘ปาฏิเทสนียนฺ’ติ ยํ วุตฺตํ, \\ ภิกฺขุ อญฺญาตโก สนฺโต, \\ สามํ คเหตฺวา ภุญฺเชยฺย, \\ นิมนฺตนาสุ ภุญฺชนฺตา ฉนฺทาย, \\ อนิวาเรตฺวา ตหิํ ภุญฺเช, \\ สทฺธาจิตฺตํ กุลํ คนฺตฺวา, \\ อคิลาโน ตหิํ ภุญฺเช, \\ โย เจ อรญฺเญ วิหรนฺโต, \\ อวิทิตํ ตหิํ ภุญฺเช, \\ ภิกฺขุนี อญฺญาติกา สนฺตา, \\ สปฺปิ เตลํ มธุํ ผาณิตํ, \\ ทธิํ สยํ วิญฺญาเปยฺย ภิกฺขุนี, \\ ‘ทุกฺกฏนฺ’ติ ยํ วุตฺตํ, \\ อปรทฺธํ วิรทฺธญฺจ, \\ ยํ มนุสฺโส กเร ปาปํ, \\ ‘ทุกฺกฏนฺ’ติ ปเวเทนฺติ, \\ ‘ทุพฺภาสิตนฺ’ติ ยํ วุตฺตํ, \\ ทุพฺภาสิตํ ทุราภฏฺฐํ, \\ ยญฺจ วิญฺญู ครหนฺติ, \\ ‘เสขิยนฺ’ติ ยํ วุตฺตํ, \\ เสกฺขสฺส สิกฺขมานสฺส, \\ อาทิ เจตํ จรณญฺจ, \\ สิกฺขา เอตาทิสี นตฺถิ, \\ \textsuperscript{*}ฉนฺนมติวสฺสติ, \\ ตสฺมา ฉนฺนํ วิวเรถ, \\ คติ มิคานํ ปวนํ,}
\csromangathagroupwithcloser{\csromangathastanza{\csromangathaitembat{339}{‘ปาราชิกนฺ’ติ ยํ วุตฺตํ,}{ตํ สุโณหิ ยถาตถํ.} \\ \csromangathacontbat{จุโต ปรทฺโธ ภฏฺโฐ จ,}{สทฺธมฺมา หิ นิรงฺกโต.} \\ \csromangathacontbat{สํวาโสปิ ตหิํ นตฺถิ,}{เตเนตํ อิติ วุจฺจติ.} \\ \csromangathacontbat{‘สํฆาทิเสโส’ติ ยํ วุตฺตํ,}{ตํ สุโณหิ ยถาตถํ.} \\ \csromangathacontbat{สํโฆว เทติ ปริวาสํ,}{มูลาย ปฏิกสฺสติ.} \\ \csromangathacontbat{มานตฺตํ เทติ อพฺเภติ,}{เตเนตํ อิติ วุจฺจติ.} \\ \csromangathacontbat{‘อนิยโต’ติ ยํ วุตฺตํ,}{ตํ สุโณหิ ยถาตถํ.} \\ \csromangathacontbat{อนิยโต น นิยโต,}{อเนกํสิกตํ ปทํ.} \\ \csromangathacontbat{ติณฺณมญฺญตรํ ฐานํ,}{‘อนิยโต’ติ ปวุจฺจติ.} \\ \csromangathacontbat{‘ถุลฺลจฺจยนฺ’ติ ยํ วุตฺตํ,}{ตํ สุโณหิ ยถาตถํ.} \\ \csromangathacontbat{เอกสฺส มูเล โย เทเสติ,}{โย จ ตํ ปฏิคณฺหติ.} \\ \csromangathacontbat{อจฺจโย เตน สโม นตฺถิ,}{เตเนตํ อิติ วุจฺจติ.} \\ \csromangathacontbat{‘นิสฺสคฺคิยนฺ’ติ ยํ วุตฺตํ,}{ตํ สุโณหิ ยถาตถํ.} \\ \csromangathacontbat{สํฆมชฺเฌ คณมชฺเฌ,}{เอกสฺเสว จ เอกโต.} \\ \csromangathacontbat{นิสฺสชฺชิตฺวาน เทเสติ,}{เตเนตํ อิติ วุจฺจติ.}}\csromangathastanza{\csromanfolio{262}\csromangathabat{‘ปาจิตฺติยนฺ’ติ ยํ วุตฺตํ,}{ตํ สุโณหิ ยถาตถํ.} \\ \csromangathabat{ปาเตติ กุสลํ ธมฺมํ,}{อริยมคฺคํ อปรชฺฌติ.}}\csromangathastanza{\csromangathabat{จิตฺตสํโมหนฏฺฐานํ,}{เตเนตํ อิติ วุจฺจติ.} \\ \csromangathabat{‘ปาฏิเทสนียนฺ’ติ ยํ วุตฺตํ,}{ตํ สุโณหิ ยถาตถํ.}}\csromangathastanza{\csromangathabat{ภิกฺขุ อญฺญาตโก สนฺโต,}{กิจฺฉา ลทฺธาย โภชนํ.} \\ \csromangathabat{สามํ คเหตฺวา ภุญฺเชยฺย,}{‘คารยฺหนฺ’ติ ปวุจฺจติ.}}\csromangathastanza{\csromangathabat{นิมนฺตนาสุ ภุญฺชนฺตา ฉนฺทาย,}{โวสาสติ ตตฺถ ภิกฺขุนิํ.} \\ \csromangathabat{อนิวาเรตฺวา ตหิํ ภุญฺเช,}{‘คารยฺหนฺ’ติ ปวุจฺจติ.}}\csromangathastanza{\csromangathabat{สทฺธาจิตฺตํ กุลํ คนฺตฺวา,}{อปฺปโภคํ อนาฬิยํ\csromansharedfootnote{e5d5f0395c4b3144}{อนาฬิยํ (สี, สฺยา)}.} \\ \csromangathabat{อคิลาโน ตหิํ ภุญฺเช,}{‘คารยฺหนฺ’ติ ปวุจฺจติ.}}\csromangathastanza{\csromangathabat{โย เจ อรญฺเญ วิหรนฺโต,}{สาสงฺเก สภยานเก.} \\ \csromangathabat{อวิทิตํ ตหิํ ภุญฺเช,}{‘คารยฺหนฺ’ติ ปวุจฺจติ.}}\csromangathastanza{\csromangathabat{ภิกฺขุนี อญฺญาติกา สนฺตา,}{ยํ ปเรสํ มมายิตํ.} \\ \csromangathabat{สปฺปิ เตลํ มธุํ ผาณิตํ,}{มจฺฉมํสํ อโถ ขีรํ.}}\csromangathastanza{\csromangathabat{ทธิํ สยํ วิญฺญาเปยฺย ภิกฺขุนี,}{คารยฺหปตฺตา สุคตสฺส สาสเน.} \\ \csromangathabat{‘ทุกฺกฏนฺ’ติ ยํ วุตฺตํ,}{ตํ สุโณหิ ยถาตถํ.}}\csromangathastanza{\csromangathabat{อปรทฺธํ วิรทฺธญฺจ,}{ขลิตํ ยญฺจ ทุกฺกฏํ.} \\ \csromangathabat{ยํ มนุสฺโส กเร ปาปํ,}{อาวิ วา ยทิ วา รโห.}}\csromangathastanza{\csromangathabat{‘ทุกฺกฏนฺ’ติ ปเวเทนฺติ,}{เตเนตํ อิติ วุจฺจติ.} \\ \csromangathabat{‘ทุพฺภาสิตนฺ’ติ ยํ วุตฺตํ,}{ตํ สุโณหิ ยถาตถํ.}}\csromangathastanza{\csromangathabat{ทุพฺภาสิตํ ทุราภฏฺฐํ,}{สํกิลิฏฺฐญฺจ ยํ ปทํ.} \\ \csromangathabat{ยญฺจ วิญฺญู ครหนฺติ,}{เตเนตํ อิติ วุจฺจติ.}}\csromangathastanza{\csromangathabat{‘เสขิยนฺ’ติ ยํ วุตฺตํ,}{ตํ สุโณหิ ยถาตถํ.} \\ \csromangathabat{เสกฺขสฺส สิกฺขมานสฺส,}{อุชุมคฺคานุสาริโน.}}\csromangathastanza{\csromangathabat{อาทิ เจตํ จรณญฺจ,}{มุขํ สญฺญมสํวโร.} \\ \csromangathabat{สิกฺขา เอตาทิสี นตฺถิ,}{เตเนตํ อิติ วุจฺจติ.}}\csromangathastanza{\csromanfolio{263}\csromangathabat{\csromansymbolfootnote{*}{ขุ 1. 144 ปิฏฺเฐ อุทาเนปิ.}ฉนฺนม\-ติวสฺสติ,}{วิวฏํ นาติวสฺสติ.} \\ \csromangathabat{ตสฺมา ฉนฺนํ วิวเรถ,}{เอวํ ตํ นาติวสฺสติ.} \\ \csromangathabat{คติ มิคานํ ปวนํ,}{อากาโส ปกฺขินํ คติ.}}}{วิภโว คติ ธมฺมานํ, นิพฺพานํ อรหโต คตีติ. คาถา\-สงฺคณิกํ นิฏฺฐิตํ.}
\csromancenter{ตสฺสุทฺทานํ.}
\setlength{\csromangathapagewidth}{0pt}
\setlength{\csromangathawakwidth}{0pt}
\csromangathameasuregroup{สตฺตนคเรสุ ปญฺญตฺตา, \\ ภิกฺขูนํ ภิกฺขุนีนญฺจ, \\ สาสนํ อนุคฺคหาย,}{\csromangathabat{สตฺตนคเรสุ ปญฺญตฺตา,}{วิปตฺติ จตุโรปิ จ.} \\ \csromangathabat{ภิกฺขูนํ ภิกฺขุนีนญฺจ,}{สาธารณา อสาธารณา.} \\ \csromangathabat{สาสนํ อนุคฺคหาย,}{คาถาสงฺคณิกํ อิทนฺติ.}}
\csromangathasetleft{สตฺตนคเรสุ ปญฺญตฺตา, \\ ภิกฺขูนํ ภิกฺขุนีนญฺจ, \\ สาสนํ อนุคฺคหาย,}
\csromangathagroupwithcloser{\csromangathastanza{\csromangathabat{สตฺตนคเรสุ ปญฺญตฺตา,}{วิปตฺติ จตุโรปิ จ.} \\ \csromangathabat{ภิกฺขูนํ ภิกฺขุนีนญฺจ,}{สาธารณา อสาธารณา.} \\ \csromangathabat{สาสนํ อนุคฺคหาย,}{คาถา\-สงฺคณิกํ อิทนฺติ.}}}{คาถา\-สงฺคณิกํ นิฏฺฐิตํ.}
\csromanmatikaanchor{mtk.185}
\csromantocmarknopage{section}{อธิกรณเภท}{mtk.185}
\bookmark[dest={mtk.185},level=1]{อธิกรณเภท}
\csromanheader{1}{อธิกรณ\-เภท}
\csromanmatikaanchor{mtk.186}
\csromantocmarkref{subparagraph}{1. อุกฺโกฏนเภทาทิ}{mtk.186}
\bookmark[dest={mtk.186},level=5]{1. อุกฺโกฏนเภทาทิ}
\csromanheader{6}{1. อุกฺโกฏน\-เภทาทิ}
\proseitem{340}{\csromanfolio{264}\csromansymbolfootnote{*}{วิ 4. 211; วิ 5. 169 ปิฏฺเฐสุปิ.}จตฺตาริ อธิกรณานิ. วิวาทา\-ธิกรณํ, อนุวาทา\-ธิกรณํ, อาปตฺตา\-ธิกรณํ, กิจฺจา\-ธิกรณํ. อิมานิ จตฺตาริ อธิกรณานิ.}

\prose{อิเมสํ จตุนฺนํ อธิกรณานํ กติ อุกฺโกฏา. อิเมสํ จตุนฺนํ อธิกรณานํ ทส อุกฺโกฏา. วิวาทา\-ธิกรณสฺส เทฺว อุกฺโกฏา, อนุวาทา\-ธิกรณสฺส จตฺตาโร อุกฺโกฏา, อาปตฺตา\-ธิกรณสฺส ตโย อุกฺโกฏา, กิจฺจา\-ธิกรณสฺส เอโก อุกฺโกโฏ. อิเมสํ จตุนฺนํ อธิกรณานํ อิเม ทส อุกฺโกฏา.}

\prose{วิวาทา\-ธิกรณํ อุกฺโกเฏนฺโต กติ สมเถ อุกฺโกเฏติ. อนุวาทา\-ธิกรณํ อุกฺโกเฏนฺโต กติ สมเถ อุกฺโกเฏติ. อาปตฺตา\-ธิกรณํ อุกฺโกเฏนฺโต กติ สมเถ อุกฺโกเฏติ. กิจฺจา\-ธิกรณํ อุกฺโกเฏนฺโต กติ สมเถ อุกฺโกเฏติ.}

\prose{วิวาทา\-ธิกรณํ อุกฺโกเฏนฺโต เทฺว สมเถ อุกฺโกเฏติ. อนุวาทา\-ธิกรณํ อุกฺโกเฏนฺโต จตฺตาโร สมเถ อุกฺโกเฏติ. อาปตฺตา\-ธิกรณํ อุกฺโกเฏนฺโต ตโย สมเถ อุกฺโกเฏติ. กิจฺจา\-ธิกรณํ อุกฺโกเฏนฺโต เอกํ สมถํ อุกฺโกเฏติ.}

\proseitem{341}{กติ อุกฺโกฏา. กติหากาเรหิ อุกฺโกฏนํ ปสวติ. กติหงฺเคหิ สมนฺนาคโต ปุคฺคโล อธิกรณํ อุกฺโกเฏติ. กติ ปุคฺคลา อธิกรณํ อุกฺโกเฏนฺตา อาปตฺติํ อาปชฺชนฺติ.}

\prose{ทฺวาทส อุกฺโกฏา. ทสหากาเรหิ อุกฺโกฏนํ ปสวติ. จตูหงฺเคหิ สมนฺนาคโต ปุคฺคโล อธิกรณํ อุกฺโกเฏติ. จตฺตาโร ปุคฺคลา อธิกรณํ อุกฺโกเฏนา อาปตฺติํ อาปชฺชนฺติ.}

\prose{กตเม ทฺวาทส อุกฺโกฏา. อกตํ กมฺมํ, ทุกฺกฏํ กมฺมํ, ปุน กาตพฺพํ กมฺมํ, อนิหตํ, ทุนฺนิหตํ, ปุน นิหนิตพฺพํ, อวินิจฺฉิตํ, ทุวินิจฺฉิตํ, ปุน วินิจฺฉิตพฺพํ, อวูปสนฺตํ, ทุวูปสนฺตํ, ปุน วูปสเม\-ตพฺพนฺติ อิเม ทฺวาทส อุกฺโกฏา.}

\prose{\csromanfolio{265}กตเมหิ ทสหากาเรหิ อุกฺโกฏนํ ปสวติ. ตตฺถ ชาตกํ อธิกรณํ อุกฺโกเฏติ, ตตฺถ ชาตกํ วูปสนฺตํ อธิกรณํ อุกฺโกเฏติ, อนฺตรามคฺเค อธิกรณํ อุกฺโกเฏติ, อนฺตรามคฺเค วูปสนฺตํ อธิกรณํ อุกฺโกเฏติ, ตตฺถ คตํ อธิกรณํ อุกฺโกเฏติ, ตตฺถ คตํ วูปสนฺตํ อธิกรณํ อุกฺโกเฏติ, ตตฺถ คตํ วูปสนฺตํ อธิกรณํ อุกฺโกเฏติ, สติวินยํ อุกฺโกเฏติ, อมูฬฺหวินยํ อุกฺโกเฏติ, ตสฺสปาปิยสิกํ อุกฺโกเฏติ, ติณวตฺถารกํ อุกฺโกเฏติ. อิเมหิ ทสหากาเรหิ อุกฺโกฏนํ ปสวติ.}

\prose{กตเมหิ จตูหงฺเคหิ สมนฺนาคโต ปุคฺคโล อธิกรณํ อุกฺโกเฏติ. ฉนฺทาคติํ คจฺฉนฺโต อธิกรณํ อุกฺโกเฏติ, โทสาคติํ คจฺฉนฺโต อธิกรณํ อุกฺโกเฏติ, โมหาคติํ คจฺฉนฺโต อธิกรณํ อุกฺโกเฏติ, ภยาคติํ คจฺฉนฺโต อธิกรณํ อุกฺโกเฏติ. อิเมหิ จตูหงฺเคหิ สมนฺนาคโต ปุคฺคโล อธิกรณํ อุกฺโกเฏติ.}

\prose{กตเม จตฺตาโร ปุคฺคลา อธิกรณํ อุกฺโกเฏนฺตา อาปตฺติํ อาปชฺชนฺติ. ตทหุ\-ปสมฺปนฺโน อุกฺโกเฏติ, อุกฺโกฏนกํ ปาจิตฺติยํ. อาคนฺตุโก อุกฺโกเฏติ, อุกฺโกฏนกํ ปาจิตฺติยํ. การโก อุกฺโกเฏติ, อุกฺโกฏนกํ ปาจิตฺติยํ. ฉนฺททายโก อุกฺโกเฏติ, อุกฺโกฏนกํ ปาจิตฺติยํ. อิเม จตฺตาโร ปุคฺคลา อธิกรณํ อุกฺโกเฏนฺตา อาปตฺติํ อาปชฺชนฺติ.}

\csromanmatikaanchor{mtk.187}
\csromantocmarkref{subparagraph}{2. อธิกรณนิทานาทิ}{mtk.187}
\bookmark[dest={mtk.187},level=5]{2. อธิกรณนิทานาทิ}
\csromanheader{6}{2. อธิกรณ\-นิทานาทิ}
\proseitem{342}{วิวาทา\-ธิกรณํ กิํนิทานํ กิํสมุทยํ กิํ ชาติกํ กิํปภวํ กิํสมฺภารํ กิํสมุฏฺฐานํ. อนุวาทา\-ธิกรณํ กิํนิทานํ กิํสมุทยํ กิํชาติกํ กิํปภวํ กิํสมฺภารํ กิํสมุฏฺฐานํ. อาปตฺตา\-ธิกรณํ กิํนิทานํ กิํสมุทยํ กิํชาติกํ กิํปภวํ กิํสมฺภารํ กิํสมุฏฺฐานํ. กิจฺจา\-ธิกรณํ กิํนิทานํ กิํสมุทยํ กิํชาติกํ กิํปภวํ กิํสมฺภารํ กิํสมุฏฺฐานํ.}

\prose{วิวาทา\-ธิกรณํ วิวาทนิทานํ วิวาท\-สมุทยํ วิวาทชาติกํ วิวาทปภวํ วิวาท\-สมฺภารํ วิวาท\-สมุฏฺฐานํ. อนุวาทา\-ธิกรณํ อนุวาทนิทานํ อนุวาท\-สมุทยํ อนุวาทชาติกํ อนุวาท\-ปภวํ อนุวาท\-สมฺภารํ อนุวาท\-สมุฏฺฐานํ. อาปตฺตา\-ธิกรณํ อาปตฺตินิทานํ อาปตฺติ\-สมุทยํ อาปตฺติชาติกํ อาปตฺติปภวํ อาปตฺติ\-สมฺภารํ อาปตฺติ\-สมุฏฺฐานํ. กิจฺจา\-ธิกรณํ กิจฺจยนิทานํ กิจฺจยสมุทยํ กิจฺจยชาติกํ กิจฺจย\-ปภวํ กิจฺจย\-สมฺภารํ กิจฺจย\-สมุฏฺฐานํ.}

\prose{\csromanfolio{266}วิวาทา\-ธิกรณํ กิํนิทานํ กิํสมุทยํ กิํชาติกํ กิํปภวํ กิํสมฺภารํ กิํสมุฏฺฐานํ. อนุวาทา\-ธิกรณํ ฯเปฯ. อาปตฺตา\-ธิกรณํ ฯเปฯ. กิจฺจา\-ธิกรณํ กิํนิทานํ กิํสมุทยํ กิํชาติกํ กิํปภวํ กิํสมฺภารํ กิํสมุฏฺฐานํ.}

\prose{วิวาทา\-ธิกรณํ เหตุนิทานํ เหตุสมุทยํ เหตุชาติกํ เหตุปภวํ เหตุสมฺภารํ เหตุ\-สมุฏฺฐานํ. อนุวาทา\-ธิกรณํ ฯเปฯ. อาปตฺตา\-ธิกรณํ ฯเปฯ. กิจฺจา\-ธิกรณํ เหตุนิทานํ เหตุสมุทยํ เหตุชาติกํ เหตุปภวํ เหตุสมฺภารํ เหตุ\-สมุฏฺฐานํ.}

\prose{วิวาทา\-ธิกรณํ กิํนิทานํ กิํ สมุทยํ กิํ ชาติกํ กิํปภวํ กิํ สมฺภารํ กิํสมุฏฺฐานํ. อนุวาทา\-ธิกรณํ ฯเปฯ. อาปตฺตา\-ธิกรณํ ฯเปฯ. กิจฺจา\-ธิกรณํ กิํนิทานํ กิํสมุทยํ กิํชาติกํ กิํปภวํ กิํสมฺภารํ กิํสมุฏฺฐานํ.}

\prose{วิวาทา\-ธิกรณํ ปจฺจยนิทานํ ปจฺจย\-สมุทยํ ปจฺจยชาติกํ ปจฺจย\-ปภวํ ปจฺจย\-สมฺภารํ ปจฺจย\-สมุฏฺฐานํ. อนุวาทา\-ธิกรณํ ฯเปฯ. อาปตฺตา\-ธิกรณํ ฯเปฯ. กิจฺจา\-ธิกรณํ ปจฺจยนิทานํ ปจฺจย\-สมุทยํ ปจฺจยชาติกํ ปจฺจย\-ปภวํ ปจฺจย\-สมฺภารํ ปจฺจย\-สมุฏฺฐานํ.}

\csromanmatikaanchor{mtk.188}
\csromantocmarkref{subparagraph}{3. อธิกรณมูลาทิ}{mtk.188}
\bookmark[dest={mtk.188},level=5]{3. อธิกรณมูลาทิ}
\csromanheader{6}{3. อธิกรณ\-มูลาทิ}
\proseitem{343}{จตุนฺนํ อธิกรณานํ กติ มูลานิ. กติ สมุฏฺฐานา. จตุนฺนํ อธิกรณานํ เตตฺติํส มูลานิ. เตตฺติํส สมุฏฺฐานา.}

\prose{จตุนฺนํ อธิกรณานํ กตมานิ เตตฺติํส มูลานิ. วิวาทา\-ธิกรณสฺส ทฺวาทส มูลานิ, อนุวาทา\-ธิกรณสฺส จุทฺทส มูลานิ, อาปตฺตา\-ธิกรณสฺส ฉ มูลานิ, กิจฺจา\-ธิกรณสฺส เอกํ มูลํ สํโฆ. จตุนฺนํ อธิกรณานํ อิมานิ เตตฺติํส มูลานิ.}

\prose{จตุนฺนํ อธิกรณานํ กตเม เตตฺติํส สมุฏฺฐานา. วิวาทา\-ธิกรณสฺส อฏฺฐารส\-เภทกร\-วตฺถูนิ สมุฏฺฐานา, อนุวาทา\-ธิกรณสฺส จตสฺโส วิปตฺติโย สมุฏฺฐานา, อาปตฺตา\-ธิกรณสฺส สตฺตา\-ปตฺติกฺขนฺธา สมุฏฺฐานา, กิจฺจธิกรณสฺส จตฺตาริ กมฺมานิ สมุฏฺฐานา. จตุนฺนํ อธิกรณานํ อิเม เตตฺติํส สมุฏฺฐานา.}

\csromanmatikaanchor{mtk.189}
\csromantocmarkref{subparagraph}{4. อธิกรณปจฺจยาปตฺติ}{mtk.189}
\bookmark[dest={mtk.189},level=5]{4. อธิกรณปจฺจยาปตฺติ}
\csromanheader{6}{4. อธิกรณ\-ปจฺจยา\-ปตฺติ}
\proseitem{344}{\csromanfolio{267}วิวาทา\-ธิกรณํ อาปตฺตานาปตฺตีติ, วิวาทา\-ธิกรณํ น อาปตฺติ. กิํ ปน วิวาทาธิกรณ\-ปจฺจยา อาปตฺติํ อาปชฺเชยฺยาติ, อาม วิวาทาธิกรณ\-ปจฺจยา อาปตฺติํ อาปชฺเชยฺย. วิวาทาธิกรณ\-ปจฺจยา กติ อาปตฺติโย อาปชฺชติ, วิวาทาธิกรณ\-ปจฺจยา เทฺว อาปตฺติโย อาปชฺชติ. อุปสมฺปนฺนํ โอมสติ, อาปตฺติ ปาจิตฺติยสฺส. อนุปสมฺปนฺนํ โอมสติ, อาปตฺติ ทุกฺกฏสฺส. วิวาทาธิกรณ\-ปจฺจยา อิมา เทฺว อาปตฺติโย อาปชฺชติ.}

\prose{ตา อาปตฺติโย จตุนฺนํ วิปตฺตีนํ กติ วิปตฺติโย ภชนฺติ, จตุนฺนํ อธิกรณานํ กตมํ อธิกรณํ, สตฺตนฺนํ อาปตฺติ\-กฺขนฺธานํ กติหิ อาปตฺติ\-กฺขนฺเธหิ สงฺคหิตา, ฉนฺนํ อาปตฺติ\-สมุฏฺฐานานํ กติหิ สมุฏฺฐาเนหิ สมุฏฺฐนฺติ, กติหิ อธิกรเณหิ, กติสุ ฐาเนสุ, กติหิ สมเถหิ สมฺมนฺติ.}

\prose{ตา อาปตฺติโย จตุนฺนํ วิปตฺตีนํ เอกํ วิปตฺติํ ภชนฺติ อาจารวิปตฺติํ. จตุนฺนํ อธิกรณานํ อาปตฺตา\-ธิกรณํ. สตฺตนฺนํ อาปตฺติ\-กฺขนฺธานํ ทฺวีหิ อาปตฺติ\-กฺขนฺเธหิ สงฺคหิตา, สิยา ปาจิตฺติยาปตฺติ\-กฺขนฺเธน, สิยา ทุกฺกฏาปตฺติ\-กฺขนฺเธน. ฉนฺนํ อาปตฺติ\-สมุฏฺฐานานํ ตีหิ สมุฏฺฐาเนหิ สมุฏฺฐนฺติ. เอเกน อธิกรเณน กิจฺจา\-ธิกรเณน. ตีสุ ฐาเนสุ สํฆมชฺเฌ คณมชฺเฌ ปุคฺคลสฺส สนฺติเก. ตีหิ สมเถหิ สมฺมนฺติ สิยา สมฺมุขา\-วินเยน จ ปฏิญฺญาต\-กรเณน จ, สิยา สมฺมุขา\-วินเยน จ ติณวตฺถารเกน จ.}

\proseitem{345}{อนุวาทา\-ธิกรณํ อาปตฺตานาปตฺตีติ, อนุวาทา\-ธิกรณํ น อาปตฺติ. กิํ ปน อนุวาทาธิกรณ\-ปจฺจยา อาปตฺติํ อาปชฺเชยฺยาติ, อาม อนุวาทาธิกรณ\-ปจฺจยา อาปตฺติํ อาปชฺเชยฺย. อนุวาทาธิกรณ\-ปจฺจยา กติ อาปตฺติโย อาปชฺชติ, อนุวาทาธิกรณ\-ปจฺจยา ติสฺโส อาปตฺติโย อาปชฺชติ. ภิกฺขุํ อมูลเกน ปาราชิเกน ธมฺเมน อนุทฺธํเสติ, อาปตฺติ สํฆาทิเสสสฺส. อมูลเกน สํฆาทิเสเสน อนุทฺธํเสติ, อาปตฺติ ปาจิตฺติยสฺส. อมูลิกาย อาจารวิปตฺติยา อนุทฺธํเสติ, อาปตฺติ ทุกฺกฏสฺส. อนุวาทาธิกรณ\-ปจฺจยา อิมา ติสฺโส อาปตฺติโย อาปชฺชติ.}

\prose{\csromanfolio{268}ตา อาปตฺติโย จตุนฺนํ วิปตฺตีนํ กติ วิปตฺติโย ภชนฺติ, จตุนฺนํ อธิกรณานํ กตมํ อธิกรณํ, สตฺตนฺนํ อาปตฺติ\-กฺขนฺธานํ กติหิ อาปตฺติ\-กฺขนฺเธหิ สงฺคหิตา, ฉนฺนํ อาปตฺติ\-สมุฏฺฐานานํ กติหิ สมุฏฺฐาเนหิ สมุฏฺฐนฺติ, กติหิ อธิกรเณหิ, กติสุ ฐาเนสุ, กติหิ สมเถหิ สมฺมนฺติ.}

\prose{ตา อาปตฺติโย จตุนฺนํ วิปตฺตีนํ เทฺว วิปตฺติโย ภชนฺติ สิยา สีลวิปตฺติํ, สิยา อาจารวิปตฺติํ. จตุนฺนํ อธิกรณานํ อาปตฺตา\-ธิกรณํ. สตฺตนฺนํ อาปตฺติ\-กฺขนฺธานํ ตีหิ อาปตฺติ\-กฺขนฺเธหิ สงฺคหิตา, สิยา สํฆาทิเสสา\-ปตฺติกฺขนฺเธน,สิยา ปาจิตฺติยาปตฺติ\-กฺขนฺเธน, สิยา ทุกฺกฏาปตฺติ\-กฺขนฺเธน. ฉนฺนํ อาปตฺติ\-สมุฏฺฐานานํ ตีหิ สมุฏฺฐาเนหิ สมุฏฺฐนฺติ. ยา ตา อาปตฺติโย ครุกา, ตา อาปตฺติโย เอเกน อธิกรเณน, กิจฺจา\-ธิกรเณน. เอกมฺหิ ฐาเน สํฆมชฺเฌ. ทฺวีหิ สมเถหิ สมฺมนฺติ สมฺมุขา\-วินเยน จ ปฏิญฺญาต\-กรเณน จ. ยา ตา อาปตฺติโย ลหุกา, ตา อาปตฺติโย เอเกน อธิกรเณน กิจฺจา\-ธิกรเณน. ตีสุ ฐาเนสุ สํฆมชฺเฌ คณมชฺเฌ ปุคฺคลสฺส สนฺติเก. ตีหิ สมเถหิ สมฺมนฺติ สิยา สมฺมุขา\-วินเยน จ ปฏิญฺญาต\-กรเณน จ, สิยา สมฺมุขา\-วินเยน จ ติณวตฺถารเกน จ.}

\proseitem{346}{อาปตฺตา\-ธิกรณํ อาปตฺตานาปตฺตีติ, อาปตฺตา\-ธิกรณํ อาปตฺติ. กิํ ปน อาปตฺตาธิกรณ\-ปจฺจยา อาปตฺติํ อาปชฺเชยฺยาติ, อาม อาปตฺตาธิกรณ\-ปจฺจยา อาปตฺติํ อาปชฺเชยฺย. อาปตฺตาธิกรณ\-ปจฺจยา กติ อาปติโย อาปชฺชติ, อาปตฺตาธิกรณ\-ปจฺจยา จตสฺโส อาปตฺติโย อาปชฺชติ. ภิกฺขุนี ชานํ ปาราชิกํ ธมฺมํ\csromansharedfootnote{355cfbe14bc98f66}{ปาราชิกํ ธมฺมํ อชฺฌาปนฺนํ (สฺยา)} ปฏิจฺฉาเทติ, อาปตฺติ ปาราชิกสฺส. เวมติกา ปฏิจฺฉาเทติ, อาปตฺติ ถุลฺลจฺจยสฺส. ภิกฺขุ สํฆาทิเสสํ ปฏิจฺฉาเทติ, อาปตฺติ ปาจิตฺติยสฺส. อาจารวิปตฺติํ ปฏิจฺฉาเทติ, อาปตฺติ ทุกฺกฏสฺส. อาปตฺตาธิกรณ\-ปจฺจยา อิมา จตสฺโส อาปตฺติโย อาปชฺชติ.}

\prose{ตา อาปตฺติโย จตุนฺนํ วิปตฺตีนํ กติ วิปตฺติโย ภชนฺติ, จตุนฺนํ อธิกรณานํ กตมํ อธิกรณํ, สตฺตนฺนํ อาปตฺติ\-กฺขนฺธานํ กติหิ อาปตฺติ\-กฺขนฺเธหิ สงฺคหิตา, ฉนฺนํ อาปตฺติ\-สมุฏฺฐานานํ กติหิ สมุฏฺฐาเนหิ \csromanfolio{269}สมุฏฺฐนฺติ, กติหิ อธิกรเณหิ, กติสุ ฐาเนสุ, กติหิ สมเถหิ สมฺมนฺติ.}

\prose{ตา อาปตฺติโย จตุนฺนํ วิปตฺตีนํ เทฺว วิปตฺติโย ภชนฺติ สิยา สีลวิปตฺติํ, สิยา อาจารวิปตฺติํ. จตุนฺนํ อธิกรณานํ อาปตฺตา\-ธิกรณํ. สตฺตนฺนํ อาปตฺติ\-กฺขนฺธานํ จตูหิ อาปตฺติ\-กฺขนฺเธหิ สงฺคหิตา สิยา ปาราชิกาปตฺติ\-กฺขนฺเธน, สิยา ถุลฺลจฺจยา\-ปตฺติกฺขนฺเธน, สิยา ปาจิตฺติยาปตฺติ\-กฺขนฺเธน, สิยา ทุกฺกฏาปตฺติ\-กฺขนฺเธน. ฉนฺนํ อาปตฺติ\-สมุฏฺฐานานํ เอเกน สมุฏฺฐาเนน สมุฏฺฐนฺติ, กายโต จ วาจโต จ จิตฺตโต จ สมุฏฺฐนฺติ. ยา สา อาปตฺติ อนวเสสา, สา อาปตฺติ น กตเมน อธิกรเณน. น กตมมฺหิ ฐาเน. น กตเมน สมเถน สมฺมติ. ยา ตา อาปตฺติโย ลหุกา, ตา อาปตฺติโย เอเกน อธิกรเณน, กิจฺจา\-ธิกรเณน. ตีสุ ฐาเนสุ สํฆมชฺเฌ คณมชฺเฌ ปุคฺคลสฺส สนฺติเก. ตีหิ สมเถหิ สมฺมนฺติ สยา สมฺมุขา\-วินเยน จ ปฏิญฺญาต\-กรเณน จ, สิยา สมฺมุขา\-วินเยน จ ติณวตฺถารเกน จ.}

\proseitem{347}{กิจฺจา\-ธิกรณํ อาปตฺตานาปตฺตีติ, กิจฺจา\-ธิกรณํ น อาปตฺติ. กิํ ปน กิจฺจาธิกรณ\-ปจฺจยา อาปตฺติํ อาปชฺเชยฺยาติ, อาม กิจฺจาธิกรณ\-ปจฺจยา อาปตฺติํ อาปชฺเชยฺย. กิจฺจาธิกรณ\-ปจฺจยา กติ อาปตฺติโย อาปชฺชติ, กิจฺจาธิกรณ\-ปจฺจยา ปญฺจ อาปตฺติโย อาปชฺชติ. อุกฺขิตฺตานุวตฺติกา ภิกฺขุนี ยาวตติยํ สมนุภาสนาย น ปฏินิสฺสชฺชติ. ญตฺติยา ทุกฺกฏํ. ทฺวีหิ กมฺมวาจาหิ ถุลฺลจฺจยา. กมฺมวาจา\-ปริโยสาเน อาปตฺติ ปาราชิกสฺส. เภทกา\-นุวตฺตกา ภิกฺขู ยาวตติยํ สมนุภาสนาย น ปฏินิสฺสชฺชนฺติ, อาปตฺติ สํฆาทิเสสสฺส. ปาปิกาย ทิฏฺฐิยา ยาวตติยํ สมนุภาสนาย น ปฏินิสฺสชฺชนฺติ, อาปตฺติ ปาจิตฺติยสฺส. กิจฺจาธิกรณ\-ปจฺจยา อิมา ปญฺจ อาปตฺติโย อาปชฺชติ.}

\prose{ตา อาปตฺติโย จตุนฺนํ วิปตฺตีนํ กติ วิปตฺติโย ภชนฺติ, จตุนฺนํ อธิกรณานํ กตมํ อธิกรณํ, สตฺตนฺนํ อาปตฺติ\-กฺขนฺธานํ กติหิ อาปตฺติ\-กฺขนฺเธหิ สงฺคหิตา, ฉนฺนํ อาปตฺติ\-สมุฏฺฐานานํ กติหิ สมุฏฺฐาเนหิ สมุฏฺฐนฺติ, กติหิ อธิกรเณหิ, กติสุ ฐาเนสุ, กติหิ สมเถหิ สมฺมนฺติ.}

\prose{\csromanfolio{270}ตา อาปตฺติโย จตุนฺนํ วิปตฺตีนํ เทฺว วิปตฺติโย ภชนฺติ สิยา สีลวิปตฺติํ, สิยา อาจารวิปตฺติํ. จตุนฺนํ อธิกรณานํ อาปตฺตา\-ธิกรณํ. สตฺตนฺนํ อาปตฺติ\-กฺขนฺธานํ ปญฺจหิ อาปตฺติ\-กฺขนฺเธหิ สงฺคหิตา สิยา ปาราชิกาปตฺติ\-กฺขนฺเธน, สิยา สํฆาทิเสสา\-ปตฺติกฺขนฺเธน, สิยา ถุลฺลจฺจยา\-ปตฺติกฺขนฺเธน, สิยา ปาจิตฺติยาปตฺติ\-กฺขนฺเธน, สิยา ทุกฺกฏาปตฺติ\-กฺขนฺเธน. ฉนฺนํ อาปตฺติ\-สมุฏฺฐานานํ เอเกน สมุฏฺฐาเนน สมุฏฺฐนฺติ, กายโต จ วาจโต จ จิตฺตโต จ สมุฏฺฐนฺติ. ยา สา อาปตฺติ อนวเสสา, สา อาปตฺติ น กตเมน อธิกรเณน. น กตมมฺหิ ฐาเน. น กตเมน สมเถน สมฺมติ. ยา สา อาปตฺติ ครุกา, สา อาปตฺติ เอเกน อธิกรเณน กิจฺจา\-ธิกรเณน. เอกมฺหิ ฐาเน สํฆมชฺเฌ. ทฺวีหิ สมเถหิ สมฺมติ สมฺมุขา\-วินเยน จ ปฏิญฺญาต\-กรเณน จ. ยา ตา อาปตฺติโย ลหุกา, ตา อาปตฺติโย เอเกน อธิกรเณน กิจฺจา\-ธิกรเณน. ตีสุ ฐาเนสุ สํฆมชฺเฌ คณมชฺเฌ ปุคฺคลสฺส สนฺติเก. ตีหิ สมเถหิ สมฺมนฺติ สิยา สมฺมุขา\-วินเยน จ ปฏิญฺญาต\-กรเณน จ, สิยา สมฺมุขา\-วินเยน จ ติณวตฺถารเกน จ.}

\csromanmatikaanchor{mtk.190}
\csromantocmarkref{subparagraph}{5. อธิกรณาธิปฺปาย}{mtk.190}
\bookmark[dest={mtk.190},level=5]{5. อธิกรณาธิปฺปาย}
\csromanheader{6}{5. อธิกรณา\-ธิปฺปาย}
\proseitem{348}{วิวาทา\-ธิกรณํ โหติ อนุวาทา\-ธิกรณํ, โหติ อาปตฺตา\-ธิกรณํ, โหติ กิจฺจา\-ธิกรณํ. วิวาทา\-ธิกรณํ น โหติ อนุวาทา\-ธิกรณํ, น โหติ อาปตฺตา\-ธิกรณํ, น โหติ กิจฺจา\-ธิกรณํ, อปิ จ วิวาทาธิกรณ\-ปจฺจยา โหติ อนุวาทา\-ธิกรณํ, โหติ อาปตฺตา\-ธิกรณํ, โหติ กิจฺจา\-ธิกรณํ. ยถา กถํ วิย,\csromansymbolfootnote{*}{วิ 4. 211; วิ 5. 202 ปิฏฺเฐสุปิ.}อิธ ภิกฺขู วิวทนฺติ ธมฺโมติ วา อธมฺโมติ วา ฯเปฯ ทุฏฺฐุลฺลาปตฺตีติ วา อทุฏฺฐุลฺลาปตฺตีติ วา, ยํ ตตฺถ ภณฺฑนํ กลโห วิคฺคโห วิวาโท นานาวาโท อญฺญถาวาโท วิปจฺจตาย โวหาโร เมธกํ, อิทํ วุจฺจติ วิวาทา\-ธิกรณํ. วิวาทา\-ธิกรเณ สาโฆ วิวทติ วิวาทา\-ธิกรณํ, วิวทมาโน อนุวทติ อนุวาทา\-ธิกรณํ. อนุวทมาโน อาปตฺติํ อาปชฺชติ อาปตฺตา\-ธิกรณํ. ตาย อาปตฺติยา สํโฆ กมฺมํ กโรติ กิจฺจา\-ธิกรณํ. เอวํ วิวาทาธิกรณ\-ปจฺจยา โหติ อนุวาทา\-ธิกรณํ, โหติ อาปตฺตา\-ธิกรณํ, โหติ กิจฺจา\-ธิกรณํ.}

\prose{\csromanfolio{271}อนุวาทา\-ธิกรณํ โหติ อาปตฺตา\-ธิกรณํ, โหติ กิจฺจา\-ธิกรณํ, โหติ วิวาทา\-ธิกรณํ. อนุวาทา\-ธิกรณํ น โหติ อาปตฺตา\-ธิกรณํ, น โหติ กิจฺจา\-ธิกรณํ, น โหติ วิวาทา\-ธิกรณํ, อปิ จ อนุวาทาธิกรณ\-ปจฺจยา โหติ อาปตฺตา\-ธิกรณํ, โหติ กิจฺจา\-ธิกรณํ, โหติ วิวาทา\-ธิกรณํ. ยถา กถํ วิย,\csromansymbolfootnote{*}{วิ 4. 212; วิ 5. 203 ปิฏฺเฐสุปิ.}อิธ ภิกฺขู ภิกฺขุํ อนุวทนฺติ สีลวิปตฺติยา วา อาจารวิปตฺติยา วา ทิฏฺฐิ\-วิปตฺติยา วา อาชีววิปตฺติยา วา, โย ตตฺถ อนุวาโท อนุวทนา อนุลฺลปนา อนุภณนา อนุสมฺปวงฺกตา อพฺภุสฺสหนตา อนุพล\-ปฺปทานํ, อิทํ วุจฺจติ อนุวาทา\-ธิกรณํ. อนุวาทา\-ธิกรเณ สํโฆ วิวทติ วิวาทา\-ธิกรณํ. วิวทมาโน อนุวทติ อนุวาทา\-ธิกรณํ. อนุวทมาโน อาปตฺติํ อาปชฺชติ อาปตฺตา\-ธิกรณา. ตาย อาปตฺติยา สํโฆ กมฺมํ กโรติ กิจฺจา\-ธิกรณํ, เอวํ อนุวาทาธิกรณ\-ปจฺจยา โหติ อาปตฺตา\-ธิกรณํ, โหติ กิจฺจา\-ธิกรณํ, โหติ วิวาทา\-ธิกรณํ.}

\prose{อาปตฺตา\-ธิกรณํ โหติ กิจฺจา\-ธิกรณํ, โหติ วิวาทา\-ธิกรณํ, โหติ อนุวาทา\-ธิกรณํ. อาปตฺตา\-ธิกรณํ น โหติ กิจฺจา\-ธิกรณํ, น โหติ วิวาทา\-ธิกรณํ, น โหติ อนุวาทา\-ธิกรณํ, อปิ จ อาปตฺตาธิกรณ\-ปจฺจยา โหติ กิจฺจา\-ธิกรณํ, โหติ วิวาทา\-ธิกรณํ, โหติ อนุวาทา\-ธิกรณํ. ยถา กถํ วิย,\csromansymbolmark{*}ปญฺจปิ อาปตฺติกฺขนฺธา อาปตฺตา\-ธิกรณํ, สตฺตปิ อาปตฺติกฺขนฺธา อาปตฺตา\-ธิกรณํ, อิทํ วุจฺจติ อาปตฺตา\-ธิกรณํ. อาปตฺตา\-ธิกรเณ สํโฆ วิวทติ วิวาทา\-ธิกรณํ. วิวทมาโน อนุวทติ อนุวาทา\-ธิกรณํ. อนุวทมาโน อาปตฺติํ อาปชฺชติ อาปตฺตา\-ธิกรณํ. ตาย อาปตฺติยา สํโฆ กมฺมํ กโรติ กิจฺจา\-ธิกรณํ. เอวํ อาปตฺตาธิกรณ\-ปจฺจยา โหติ กิจฺจา\-ธิกรณํ, โหติ วิวาทา\-ธิกรณํ, โหติ อนุวาทา\-ธิกรณํ.}

\csromanmatikaanchor{mtk.191}
\csromanmatikaanchor{mtk.192}
\csromantocmarkref{subparagraph}{6. ปุจฺฉาวาร}{mtk.191}
\bookmark[dest={mtk.191},level=5]{6. ปุจฺฉาวาร}
\csromantocmarkref{subparagraph}{7. วิสชฺชนาวาร}{mtk.192}
\bookmark[dest={mtk.192},level=5]{7. วิสชฺชนาวาร}
\prose{กิจฺจา\-ธิกรณํ โหติ วิวาทา\-ธิกรณํ, โหติ อนุวาทา\-ธิกรณํ, โหติ อาปตฺตา\-ธิกรณํ. กิจฺจา\-ธิกรณํ น โหติ วิวาทา\-ธิกรณํ, น โหติ อนุวาทา\-ธิกรณํ, น โหติ อาปตฺตา\-ธิกรณํ, อปิ จ กิจฺจาธิกรณ\-ปจฺจยา โหติ วิวาทา\-ธิกรณํ, โหติ อนุวาทา\-ธิกรณํ, โหติ อาปตฺตา\-ธิกรณํ. ยถา กถํ วิย, ยา สํฆสฺส กิจฺจยตา กรณียตา อปโลกน\-กมฺมํ ญตฺติกมฺมํ ญตฺติ\-ทุติย\-กมฺมํ ญตฺติ\-จตุตฺถ\-กมฺมํ, \csromanfolio{272}อิทํ วุจฺจติ กิจฺจา\-ธิกรณํ. กิจฺจาธิกรเณ สํโฆ วิวทติ วิวาทา\-ธิกรณํ. วิวทมาโน อนุวทติ อนุวาทา\-ธิกรณํ. อนุวทมาโน อาปตฺติํ อาปชฺชติ อาปตฺตา\-ธิกรณํ. ตาย อาปตฺติยา สํโฆ กมฺมํ กโรติ กิจฺจา\-ธิกรณํ. เอวํ กิจฺจาธิกรณ\-ปจฺจยา โหติ วิวาทา\-ธิกรณํ, โหติ อนุวาทา\-ธิกรณํ, โหติ อาปตฺตา\-ธิกรณํ.}

\csromanheader{6}{6. \textbf{ปุจฺฉาวาร}}
\proseitem{349}{ยตฺถ สติวินโย ตตฺถ สมฺมุขาวินโย, ยตฺถ สมฺมุขาวินโย ตตฺถ สติวินโย. ยตฺถ อมูฬฺหวินโย ตตฺถ สมฺมุขาวินโย, ยตฺถ สมฺมุขาวินโย ตตฺถ อมูฬฺหวินโย. ยตฺถ ปฏิญฺญาต\-กรณํ ตตฺถ สมฺมุขาวินโย, ยตฺถ สมฺมุขาวินโย ตตฺถ ปฏิญฺญาต\-กรณํ. ยตฺถ เยภุยฺยสิกา ตตฺถ สมฺมุขาวินโย, ยตฺถ สมฺมุขาวินโย ตตฺถ เยภุยฺยสิกา. ยตฺถ ตสฺสปาปิยสิกา ตตฺถ สมฺมุขาวินโย, ยตฺถ สมฺมุขาวินโย ตตฺถ ตสฺสปาปิยสิกา. ยตฺถ ติณวตฺถารโก ตตฺถ สมฺมุขาวินโย, ยตฺถ สมฺมุขาวินโย ตตฺถ ติณวตฺถารโก.}

\csromanheader{2}{7. \textbf{วิสฺสชฺชนาวาร}}
\proseitem{350}{ยสฺมิํ สมเย สมฺมุขา\-วินเยน จ สติวินเยน จ อธิกรณํ วูปสมฺมติ, ยตฺถ สติวินโย ตตฺถ สมฺมุขาวินโย, ยตฺถ สมฺมุขาวินโย ตตฺถ สติวินโย, น ตตฺถ อมูฬฺหวินโย, น ตตฺถ ปฏิญฺญาต\-กรณํ, น ตตฺถ เยภุยฺยสิกา, น ตตฺถ ตสฺสปาปิยสิกา, น ตตฺถ ติณวตฺถารโก. ยสฺมิํ สมเย สมฺมุขา\-วินเยน จ อมูฬฺหวินเยน จ ฯเปฯ สมฺมุขา\-วินเยน จ ปฏิญฺญาต\-กรเณน จ ฯเปฯ สมฺมุขา วินเยน จ เยภุยฺยสิกาย จ ฯเปฯ สมฺมุขา\-วินเยน จ ตสฺสปาปิยสิกาย จ ฯเปฯ สมฺมุขา\-วินเยน จ ติณวตฺถารเกน จ อธิกรณํ วูปสมฺมติ, ยตฺถ ติณวตฺถารโก ตตฺถ สมฺมุขาวินโย, ยตฺถ สมฺมุขาวินโย ตตฺถ ติณวตฺถารโก, น ตตฺถ สติวินโย, น ตตฺถ อมูฬฺหวินโย, น ตตฺถ ปฏิญฺญาต\-กรณํ, น ตตฺถ เยภุยฺยสิกา, น ตตฺถ ตสฺสปาปิยสิกา.}

\csromanmatikaanchor{mtk.193}
\csromantocmarkref{subparagraph}{8. สํสฏฺฐวาร}{mtk.193}
\bookmark[dest={mtk.193},level=5]{8. สํสฏฺฐวาร}
\csromanheader{6}{8. \textbf{สํสฏฺฐวาร}}
\proseitem{351}{สมฺมุขา\-วินโยติ วา สติวินโยติ วา อิเม ธมฺมาสํสฏฺฐา, อุทาหุ วิสํสฏฺฐา, ลพฺภา จ ปนิเมสํ ธมฺมานํ วินิพฺภุชิตฺวา วินิพฺภุชิตฺวา \csromanfolio{273}นานากรณํ ปญฺญาเปตุํ. สมฺมุขา\-วินโยติ วา อมูฬฺหวินโยติ วา ฯเปฯ. สมฺมุขา\-วินโยติ วา ปฏิญฺญาต\-กรณนฺติ วา. สมฺมุขา\-วินโยติ วา เยภุยฺยสิกาติ วา. สมฺมุขา\-วินโยติ วา ตสฺสปาปิยสิกาติ วา. สมฺมุขา\-วินโยติ วา ติณวตฺถารโกติ วา อิเม ธมฺมา สํสฏฺฐา, อุทาหุ วิสํสฏฺฐา, ลพฺภา จ ปนิเมสํ ธมฺมานํ วินิพฺภุชิตฺวา วินิพฺภุชิตฺวา นานากรณํ ปญฺญาเปตุํ.}

\prose{สมฺมุขา\-วินโยติ วา สติวินโยติ วา อิเม ธมฺมา สํสฏฺฐา, โน วิสํสฏฺฐา, น จ ลพฺภา อิเมสํ ธมฺมานํ วินิพฺภุชิตฺวา วินิพฺภุชิตฺวา นานากรณํ ปญฺญาเปตุํ. สมฺมุขา\-วินโยติ วา อมูฬฺหวินโยติ วา ฯเปฯ. สมฺมุขา\-วินโยติ วา ปฏิญฺญาต\-กรณนฺติ วา. สมฺมุขา\-วินโยติ วา เยภุยฺยสิกาติ วา. สมฺมุขา\-วินโยติ วา ตสฺสปาปิยสิกาติ วา. สมฺมุขา\-วินโยติ วา ติณวตฺถารโกติ วา อิเม ธมฺมา สํสฏฺฐา, โน วิสํสฏฺฐา, น จ ลพฺภา อิเมสํ ธมฺมานํ วินิพฺภุชิตฺวา วินิพฺภุชิตฺวา นานากรณํ ปญฺญาเปตุํ.}

\csromanmatikaanchor{mtk.194}
\csromantocmarkref{subparagraph}{9. สตฺตสมถนิทาน}{mtk.194}
\bookmark[dest={mtk.194},level=5]{9. สตฺตสมถนิทาน}
\csromanheader{6}{9. สตฺต\-สมถ\-นิทาน}
\proseitem{352}{สมฺมุขาวินโย กิํนิทาโน กิํสมุทโย กิํชาติโก กิํปภโว กิํสมฺภาโร กิํสมุฏฺฐาโน. สติวินโย กิํนิทาโน กิํสมุทโย กิํชาติโก กิํปภโว กิํสมฺภาโร กิํสมุฏฺฐาโน. อมูฬฺหวินโย กิํนิทาโน กิํสมุทโย กิํชาติโก กิํปภโว กิํสมฺภาโร กิํสมุฏฺฐาโน. ปฏิญฺญาต\-กรณํ กิํนิทานํ กิํสมุทยํ กิํชาติกํ กิํปภวํ กิํสมฺภารํ กิํสมุฏฺฐานํ. เยภุยฺยสิกา กิํนิทานา กิํสมุทยา กิํชาติกา กิํปภวา กิํสมฺภารา กิํสมุฏฺฐานา. ตสฺสปาปิยสิกา กิํนิทานา กิํสมุทยา กิํชาติกา กิํปภวา กิํสมฺภารา กิํสมุฏฺฐานา. ติณวตฺถารโก กิํนิทาโน กิํสมุทโย กิํชาติโก กิํปภโว กิํสมฺภาโร กิํสมุฏฺฐาโน.}

\prose{สมฺมุขาวินโย นิทานนิทาโน นิทานสมุทโย นิทานชาติโก นิทานปภโว นิทานสมฺภาโร นิทาน\-สมุฏฺฐาโน. สติวินโย ฯเปฯ. อมูฬฺหวินโย ฯเปฯ. ปฏิญฺญาต\-กรณํ นิทานนิทานํ นิทาน\-สมุทยํ นิทานชาติกํ นิทานปภวํ นิทาน\-สมฺภารํ นิทาน\-สมุฏฺฐานํ. เยภุยฺยสิกา ฯเปฯ. ตสฺสปาปิยสิกา นิทานนิทานา นิทาน สมุทยา นิทานชาติกา นิทานปภวา \csromanfolio{274}นิทานสมฺภารา นิทาน\-สมุฏฺฐานา. ติณวตฺถารโก นิทานนิทาโน นิทานสมุทโย นิทานชาติโก นิทานปภโว นิทานสมฺภาโร นิทาน\-สมุฏฺฐาโน.}

\prose{สมฺมุขาวินโย กิํนิทาโน กิํสมุทโย กิํชาติโก กิํปภโว กิํสมฺภาโร กิํสมุฏฺฐาโน. สติวินโย ฯเปฯ. อมูฬฺหวินโย ฯเปฯ. ปฏิญฺญาต\-กรณํ ฯเปฯ. เยภุยฺยสิกา ฯเปฯ. ตสฺสปาปิยสิกา ฯเปฯ. ติณวตฺถารโก กิํนิทาโน กิํสมุทโย กิํชาติโก กิํปภโว กิํสมฺภาโร กิํ สมุฏฺฐาโน.}

\prose{สมฺมุขาวินโย เหตุนิทาโน เหตุสมุทโย เหตุชาติโก เหตุปภโว เหตุสมฺภาโร เหตุสมุฏฺฐาโน. สติวินโย ฯเปฯ. อมูฬฺหวินโย ฯเปฯ. ปฏิญฺญาต\-กรณํ เหตุนิทานํ เหตุสมุทยํ เหตุชาติกํ เหตุปภวํ เหตุสมฺภารํ เหตุ\-สมุฏฺฐานํ. เยภุยฺยสิกา ฯเปฯ. ตสฺสปาปิยสิกา เหตุนิทานา เหตุสมุทยา เหตุชาติกา เหตุปภวา เหตุสมฺภารา เหตุสมุฏฺฐานา. ติณวตฺถารโก เหตุนิทาโน เหตุสมุทโย เหตุชาติโก เหตุปภโว เหตุสมฺภาโร เหตุสมุฏฺฐาโน.}

\prose{สมฺมุขาวินโย กิํนิทาโน กิํสมุทโย กิํชาติโก กิํปภโว กิํสมฺภาโร กิํสมุฏฺฐาโน. สติวินโย ฯเปฯ. อมูฬฺหวินโย ฯเปฯ. ปฏิญฺญาต\-กรณํ ฯเปฯ. เยภุยฺยสิกา ฯเปฯ. ตสฺสปาปิยสิกา ฯเปฯ. ติณวตฺถารโก กิํนิทาโน กิํสมุทโย กิํชาติโก กิํปภโว กิํสมฺภาโร กิํสมุฏฺฐาโน. สมฺมุขาวินโย ปจฺจยนิทาโน ปจฺจย\-สมุทโย ปจฺจยชาติโก ปจฺจยปภโว ปจฺจย\-สมฺภาโร ปจฺจย\-สมุฏฺฐาโน. สติวินโย ฯเปฯ. อมูฬฺหวินโย ฯเปฯ. ปฏิญฺญาต\-กรณํ ปจฺจยนิทานํ ปจฺจย\-สมุทยํ ปจฺจยชาติกํ ปจฺจย\-ปภวํ ปจฺจย\-สมฺภารํ ปจฺจย\-สมุฏฺฐานํ. เยภุยฺยสิกา ฯเปฯ. ตสฺสปาปิยสิกา ปจฺจยนิทานา ปจฺจย\-สมุทยา ปจฺจยชาติกา ปจฺจยปภวา ปจฺจย\-สมฺภารา ปจฺจย\-สมุฏฺฐานา. ติณวตฺถารโก ปจฺจยนิทาโน ปจฺจย\-สมุทโย ปจฺจยชาติโก ปจฺจยปภโว ปจฺจย\-สมฺภาโร ปจฺจย\-สมุฏฺฐาโน.}

\csromanmatikaanchor{mtk.195}
\csromantocmarkref{subparagraph}{10. สตฺตสมถนานตฺถาทิ}{mtk.195}
\bookmark[dest={mtk.195},level=5]{10. สตฺตสมถนานตฺถาทิ}
\proseitem{353}{สตฺตนฺนํ สมถานํ กติ มูลานิ, กติ สมุฏฺฐานา. สตฺตนฺนํ สมถานํ ฉพฺพีส มูลานิ, ฉตฺติํส สมุฏฺฐานา. สตฺตนฺนํ สมถานํ กตมานิ \csromanfolio{275}ฉพฺพีส มูลานิ. สมฺมุขา\-วินยสฺส จตฺตาริ มูลานิ สํฆสมฺมุขตา ธมฺม\-สมฺมุขตา วินย\-สมฺมุขตา ปุคฺคล\-สมฺมุขตา. สติวินยสฺส จตฺตาริ มูลานิ. อมูฬฺห\-วินยสฺส จตฺตาริ มูลานิ. ปฏิญฺญาต\-กรณสฺส เทฺว มูลานิ โย จ เทเสติ, ยสฺส จ เทเสติ. เยภุยฺยสิกาย จตฺตาริ มูลานิ. ตสฺสปาปิยสิกาย จตฺตาริ มูลานิ. ติณวตฺถารกสฺส จตฺตาริ มูลานิ สํฆสมฺมุขตา ธมฺม\-สมฺมุขตา วินย\-สมฺมุขตา ปุคฺคล\-สมฺมุขตา. สตฺตนฺนํ สมถานํ อิมานิ ฉพฺพีส มูลานิ.}

\prose{สตฺตนฺนํ สมถานํ กตเม ฉตฺติํส สมุฏฺฐานา. สติวินยสฺส กมฺมสฺส กิริยา, กรณํ, อุปคมนํ, อชฺฌุปคมนํ, อธิวาสนา, อปฺปฏิกฺโกสนา. อมูฬฺห\-วินยสฺส กมฺมสฺส ฯเปฯ. ปฏิญฺญาต\-กรณสฺส กมฺมสฺส. เยภุยฺยสิกาย กมฺมสฺส. ตสฺสปาปิยสิกาย กมฺมสฺส. ติณวตฺถารกสฺส กมฺมสฺส กิริยา, กรณํ, อุปคมนํ, อชฺฌุปคมนํ, อธิวาสนา, อปฺปฏิกฺโกสนา. สตฺตนฺนํ สมถานํ อิเม ฉตฺติํส สมุฏฺฐานา.}

\csromanheader{2}{10. \textbf{สตฺตสมถานานตฺถาทิ}}
\proseitem{354}{สมฺมุขา\-วินโยติ วา สติวินโยติ วา อิเม ธมฺมา นานตฺถา นานาพฺยญฺชนา, อุทาหุ เอกตฺถา พฺยญฺชนเมว นานํ. สมฺมุขา\-วินโยติ วา อมูฬฺหวินโยติ วา ฯเปฯ. สมฺมุขา\-วินโยติ วา ปฏิญฺญาต\-กรณนฺติ วา. สมฺมุขา\-วินโยติ วา เยภุยฺยสิกาติ วา. สมฺมุขา\-วินโยติ วา ตสฺสปาปิยสิกาติ วา. สมฺมุขา\-วินโยติ วา ติณวตฺถารโกติ วา อิเม ธมฺมา นานตฺถา นานาพฺยญฺชนา, อุทาหุ เอกตฺถา พฺยญฺชนเมว นานํ. สมฺมุขา\-วินโยติ วา สติวินโยติ วา อิเม ธมฺมา นานตฺถา เจว นานา พฺยญฺชนา จ. สมฺมุขา\-วินโยติ วา อมูฬฺหวินโยติ วา ฯเปฯ. สมฺมุขา\-วินโยติ วา ปฏิญฺญาต\-กรณนฺติ วา. สมฺมุขา\-วินโยติ วา เยภุยฺยสิกาติ วา. สมฺมุขา\-วินโยติ วา ตสฺสปาปิยสิกาติ วา. สมฺมุขา\-วินโยติ วา ติณวตฺถารโกติ วา อิเม ธมฺมา นานตฺถา เจว นานา พฺยญฺชนา จ.}

\proseitem{355}{\csromansymbolfootnote{*}{วิ 4. 218 ปิฏฺเฐปิ.}วิวาโท วิวาทา\-ธิกรณํ, วิวาโท โน อธิกรณํ, อธิกรณํ โน วิวาโท, อธิกรณ\-ญฺเจว วิวาโท จ. สิยา วิวาโท วิวาทา\-ธิกรณํ, สิยา วิวาโท โน อธิกรณํ, สิยา อธิกรณํ โน วิวาโท, สิยา อธิกรณ\-ญฺเจว วิวาโท จ.}

\prose{\csromanfolio{276}\csromansymbolfootnote{*}{วิ 4. 218 ปิฏฺฐาทีสุปิ.}ตตฺถ กตโม วิวาโท วิวาทา\-ธิกรณํ. อิธ ภิกฺขู วิวทนฺติ ธมฺโมติ วา อธมฺโมติ วา ฯเปฯ ทุฏฺฐุลฺลา อาปตฺตีติ วา อทุฏฺฐุลฺลา อาปตฺตีติ วา, ยํ ตตฺถ ภณฺฑนํ กลโห วิคฺคโห วิวาโท นานาวาโท อญฺญถาวาโท วิปจฺจตาย โวหาโร เมธกํ, อยํ วิวาโท วิวาทา\-ธิกรณํ.}

\prose{ตตฺถ กตโม วิวาโท โน อธิกรณํ. มาตาปิ ปุตฺเตน วิวทติ, ปุตฺโตปิ มาตรา วิวทติ, ปิตาปิ ปุตฺเตน วิวทติ, ปุตฺโตปิ ปิตรา วิวทติ, ภาตาปิ ภาตรา วิวทติ, ภาตาปิ ภคินิยา วิวทติ, ภคินีปิ ภาตรา วิวทติ, สหาโยปิ สหาเยน วิวทติ, อยํ วิวาโท โน อธิกรณํ.}

\prose{ตตฺถ กตมํ อธิกรณํ โน วิวาโท. อนุวาทา\-ธิกรณํ อาปตฺตา\-ธิกรณํ กิจฺจา\-ธิกรณํ, อิทํ อธิกรณํ โน วิวาโท.}

\prose{ตตฺถ กตมํ อธิกรณ\-ญฺเจว วิวาโท จ. วิวาทา\-ธิกรณํ อธิกรณญฺจว วิวาโท จ.}

\proseitem{356}{\csromansymbolmark{*}อนุวาโท อนุวาทา\-ธิกรณํ, อนุวาโท โน อธิกรณํ, อธิกรณํ โน อนุวาโท, อธิกรณ\-ญฺเจว อนุวาโท จ. สิยา อนุวาโท อนุวาทา\-ธิกรณํ, สิยา อนุวาโท โน อธิกรณํ, สิยา อธิกรณํ โน อนุวาโท, สิยา อธิกรณ\-ญฺเจว อนุวาโท จ.}

\prose{ตตฺถ กตโม อนุวาโท อนุวาทา\-ธิกรณํ. อิธ ภิกฺขู ภิกฺขุํ อนุวทนฺติ สีลวิปตฺติยา วา อาจารวิปตฺติยา วา ทิฏฺฐิ\-วิปตฺติยา วา อาชีววิปตฺติยา วา, โย ตตฺถ อนุวาโท อนุวทนา อนุลฺลปนา อนุภณนา อนุสมฺปวงฺกตา อพฺภุสฺสหนตา อนุพล\-ปฺปทานํ, อยํ อนุวาโท อนุวาทา\-ธิกรณํ.}

\prose{ตตฺถ กตโม อนุวาโท โน อธิกรณํ. มาตาปิ ปุตฺตํ อนุวทติ, ปุตฺโตปิ มาตรํ อนุวทติ, ปิตาปิ ปุตฺตํ อนุวทติ, ปุตฺโตปิ ปิตรํ อนุวทติ, ภาตาปิ ภาตรํ อนุวทติ, ภาตาปิ ภคินิํ อนุวทติ, ภคินีปิ ภาตรํ อนุวทติ, สหาโยปิ สหายํ อนุวทติ, อยํ อนุวาโท โน อธิกรณํ.}

\prose{ตตฺถ กตมํ อธิกรณํ โน อนุวาโท. อาปตฺตา\-ธิกรณํ กิจฺจา\-ธิกรณํ วิวาทา\-ธิกรณํ, อิทํ อธิกรณํ โน อนุวาโท.}

\prose{\csromanfolio{277}ตตฺถ กตมํ อธิกรณ\-ญฺเจว อนุวาโท จ. อนุวาทา\-ธิกรณํ อธิกรณ\-ญฺเจว อนุวาโท จ.}

\proseitem{357}{อาปตฺติ อาปตฺตา\-ธิกรณํ, อาปตฺติ โน อธิกรณํ, อธิกรณํ โน อาปตฺติ, อธิกรณ\-ญฺเจว อาปตฺติ จ. สิยา อาปตฺติ อาปตฺตา\-ธิกรณํ, สิยา อาปตฺติ โน อธิกรณํ, สิยา อธิกรณํ โน อาปตฺติ, สิยา อธิกรณ\-ญฺเจว อาปตฺติ จ.}

\prose{ตตฺถ กตมา อาปตฺติ อาปตฺตา\-ธิกรณํ. ปญฺจปิ อาปตฺติกฺขนฺธา อาปตฺตา\-ธิกรณํ, สตฺตปิ อาปตฺติกฺขนฺธา อาปตฺตา\-ธิกรณํ. อยํ อาปตฺติ อาปตฺตา\-ธิกรณํ.}

\prose{ตตฺถ กตมา อาปตฺติ โน อธิกรณํ. โสตาปตฺติ สมาปตฺติ, อยํ อาปตฺติ โน อธิกรณํ.}

\prose{ตตฺถ กตมํ อธิกรณํ โน อาปตฺติ. กิจฺจา\-ธิกรณํ วิวาทา\-ธิกรณํ อนุวาทา\-ธิกรณํ, อิทํ อธิกรณํ โน อาปตฺติ.}

\prose{ตตฺถ กตมํ อธิกรณ\-ญฺเจว อาปตฺติ จ. อาปตฺตา\-ธิกรณํ อธิกรณ\-ญฺเจว อาปตฺติ จ.}

\proseitem{385}{\csromansymbolfootnote{*}{วิ 4. 219 ปิตฺเถปิ.}กิจฺจํ กิจฺจา\-ธิกรณํ, กิจฺจํ โน อธิกรณํ, อธิกรณํ โน กิจฺจํ, อธิกรณ\-ญฺเจว กิจฺจญฺจ. สิยา กิจฺจํ กิจฺจา\-ธิกรณํ, สิยา กิจฺจํ โน อธิกรณํ, สิยา อธิกรณํ โน กิจฺจํ, สิยา อธิกรณ\-ญฺเจว กิจฺจญฺจ.}

\prose{ตตฺถ กตมํ กิจฺจํ กิจฺจา\-ธิกรณํ. ยา สํฆสฺส กิจฺจยตา กรณียตา อปโลกน\-กมฺมํ ญตฺติกมฺมํ ญตฺติ\-ทุติย\-กมฺมํ ญตฺติ\-จตุตฺถ\-กมฺมํ, อิทํ กิจฺจํ กิจฺจา\-ธิกรณํ.}

\prose{ตตฺถ กตมํ กิจฺจํ โน อธิกรณํ. อาจริยกิจฺจํ อุปชฺฌาย\-กิจฺจํ\csromansharedfootnote{411b53fece2d731d}{อุปชฺฌายกิจฺจํ สกิจฺจํ (ก)} สมานุ\-ปชฺฌาย\-กิจฺจํ สมานา\-จริย\-กิจฺจํ, อิทํ กิจฺจํ โน อธิกรณํ.}

\prose{ตตฺถ กตมํ อธิกรณํ โน กิจฺจํ. วิวาทา\-ธิกรณํ อนุวาทา\-ธิกรณํ อาปตฺตา\-ธิกรณํ, อิทํ อธิกรณํ โน กิจฺจํ.}

\prosewithcloser{\csromanfolio{278}ตตฺถ กตมํ อธิกรณ\-ญฺเจว กิจฺจญฺจ. กิจฺจา\-ธิกรณํ อธิกรณ\-ญฺเจว กิจฺจญฺจาติ.}{อธิกรณ\-เภทํ นิฏฺฐิตํ.}
\csromancenter{ตสฺสุทฺทานํ.}
\setlength{\csromangathapagewidth}{0pt}
\setlength{\csromangathawakwidth}{0pt}
\csromangathameasuregroup{อธิกรณํ อุกฺโกฏา, \\ นิทานเหตุปจฺจยา, \\ อาปตฺติ โหติ ยตฺถ จ, \\ เหตุปจฺจยมูลานิ, \\ วิวาโท อธิกรณนฺติ,}{\csromangathabat{อธิกรณํ อุกฺโกฏา,}{อาการา ปุคฺคเลน จ.} \\ \csromangathabat{นิทานเหตุปจฺจยา,}{มูลํ สมุฏฺฐาเนน จ.} \\ \csromangathabat{อาปตฺติ โหติ ยตฺถ จ,}{สํสฏฺฐา นิทาเนน จ\textsuperscript{1}.} \\ \csromangathabat{เหตุปจฺจยมูลานิ,}{สมุฏฺฐาเนน พฺยญฺชนา.} \\ \csromangathabat{วิวาโท อธิกรณนฺติ,}{เภทาธิกรเณ อิทนฺติ.}}
\csromangathasetleft{อธิกรณํ อุกฺโกฏา, \\ นิทานเหตุปจฺจยา, \\ อาปตฺติ โหติ ยตฺถ จ, \\ เหตุปจฺจยมูลานิ, \\ วิวาโท อธิกรณนฺติ,}
\csromangathagroup{\csromangathastanza{\csromangathabat{อธิกรณํ อุกฺโกฏา,}{อาการา ปุคฺคเลน จ.} \\ \csromangathabat{นิทานเหตุปจฺจยา,}{มูลํ สมุฏฺฐาเนน จ.}}\csromangathastanza{\csromangathabat{อาปตฺติ โหติ ยตฺถ จ,}{สํสฏฺฐา นิทาเนน จ\csromansharedfootnote{a98a7d63ce39883c}{สํสฏฺฐา นิทานปภวา (สี)}.} \\ \csromangathabat{เหตุปจฺจยมูลานิ,}{สมุฏฺฐาเนน พฺยญฺชนา.} \\ \csromangathabat{วิวาโท อธิกรณนฺติ,}{เภทาธิกรเณ อิทนฺติ.}}}

\csromanmatikaanchor{mtk.196}
\csromantocmarknopage{section}{อปรคาถาสงฺคณิก}{mtk.196}
\bookmark[dest={mtk.196},level=1]{อปรคาถาสงฺคณิก}
\csromanheader{1}{อปร\-คาถาสงฺคณิก}
\csromanheader{2}{1. โจทนา\-ทิ\-ปุจฺฉา\-วิสฺสชฺชนา}
\csromanmatikaanchor{mtk.197}
\csromantocmarkref{subparagraph}{1. โจทนาทิปุจฺฉาวิสชฺชนา}{mtk.197}
\bookmark[dest={mtk.197},level=5]{1. โจทนาทิปุจฺฉาวิสชฺชนา}
\setlength{\csromangathapagewidth}{0pt}
\setlength{\csromangathawakwidth}{0pt}
\csromangathameasuregroup{สจฺจํ อหมฺปิ ชานามิ, \\ อญฺญญฺจ ตาหํ ปุจฺฉามิ,}{โจทนา กิมตฺถาย, สารณา กิสฺส การณา. \\ สํโฆ กิมตฺถาย, มติกมฺมํ ปน กิสฺส การณา. \\ โจทนา สารณตฺถาย, นิคฺคหตฺถาย สารณา. \\ สํโฆ ปริคฺคหตฺถาย, มติกมฺมํ ปน ปาฏิเยกฺกํ. \\ มา โข ตุริโต อภณิ, มา โข จณฺฑิกโต ภณิ. \\ มา โข ปฏิฆํ ชนยิ, สเจ อนุวิชฺชโก ตุวํ. \\ มา โข สหสา อภณิ, กถํ วิคฺคาหิกํ อนตฺถสํหิตํ. \\ สุตฺเต วินเย อนุโลเม, ปญฺญตฺเต อนุโลมิเก. \\ อนุโยควตฺตํ นิสามย, กุสเลน พุทฺธิมตา กตํ. \\ สุวุตฺตํ สิกฺขาปทานุโลมิกํ, คติํ น นาเสนฺโต สมฺปรายิกํ. \\ หิเตสี อนุยุญฺชสฺสุ, กาเลนตฺถูปสํหิตํ. \\ จุทิตสฺส จ โจทกสฺส จ, \\ สหสา โวหารํ มา ปธาเรสิ. \\ โจทโก อาห อาปนฺโนติ, \\ จุทิตโก อาห อนาปนฺโนติ. \\ อุโภ อนุกฺขิปนฺโต, ปฏิญฺญานุสนฺธิเตน การเย. \\ ปฏิญฺญา ลชฺชีสุ กตา, อลชฺชีสุ เอวํ น วิชฺชติ. \\ พหุมฺปิ อลชฺชี ภาเสยฺย, วตฺตานุสนฺธิเตน\textsuperscript{1} การเย. \\ อลชฺชี กีทิโส โหติ, ปฏิญฺญา ยสฺส น รูหติ. \\ เอตญฺจ\textsuperscript{1} ตาหํ ปุจฺฉามิ, กีทิโส วุจฺจติ อลชฺชีปุคฺคโล. \\ สญฺจิจฺจ อาปตฺติํ อาปชฺชติ, อาปตฺติํ ปริคูหติ. \\ อคติคมนญฺจ คจฺฉติ, เอทิโส วุจฺจติ อลชฺชีปุคฺคโล. \\ สจฺจํ อหมฺปิ ชานามิ, เอทิโส วุจฺจติ อลชฺชีปุคฺคโล. \\ อญฺญญฺจ ตาหํ ปุจฺฉามิ, กีทิโส วุจฺจติ ลชฺชีปุคฺคโล. \\ สญฺจิจฺจ อาปตฺติํ นาปชฺชติ, อาปตฺติํ น ปริคูหติ. \\ อคติคมนํ น คจฺฉติ, เอทิโส วุจฺจติ ลชฺชีปุคฺคโล. \\ \csromangathabat{สจฺจํ อหมฺปิ ชานามิ,}{เอทิโส วุจฺจติ วุจฺจติ ลชฺชีปุคฺคโล.} \\ \csromangathabat{อญฺญญฺจ ตาหํ ปุจฺฉามิ,}{กีทิโส วุจฺจติ อธมฺมโจทโก.}}
\csromangathameasurewak{โจทนา กิมตฺถาย, สารณา กิสฺส การณา. \\ สํโฆ กิมตฺถาย, มติกมฺมํ ปน กิสฺส การณา. \\ โจทนา สารณตฺถาย, นิคฺคหตฺถาย สารณา. \\ สํโฆ ปริคฺคหตฺถาย, มติกมฺมํ ปน ปาฏิเยกฺกํ. \\ มา โข ตุริโต อภณิ, มา โข จณฺฑิกโต ภณิ. \\ มา โข ปฏิฆํ ชนยิ, สเจ อนุวิชฺชโก ตุวํ. \\ มา โข สหสา อภณิ, กถํ วิคฺคาหิกํ อนตฺถสํหิตํ. \\ สุตฺเต วินเย อนุโลเม, ปญฺญตฺเต อนุโลมิเก. \\ อนุโยควตฺตํ นิสามย, กุสเลน พุทฺธิมตา กตํ. \\ สุวุตฺตํ สิกฺขาปทานุโลมิกํ, คติํ น นาเสนฺโต สมฺปรายิกํ. \\ หิเตสี อนุยุญฺชสฺสุ, กาเลนตฺถูปสํหิตํ. \\ จุทิตสฺส จ โจทกสฺส จ, \\ สหสา โวหารํ มา ปธาเรสิ. \\ โจทโก อาห อาปนฺโนติ, \\ จุทิตโก อาห อนาปนฺโนติ. \\ อุโภ อนุกฺขิปนฺโต, ปฏิญฺญานุสนฺธิเตน การเย. \\ ปฏิญฺญา ลชฺชีสุ กตา, อลชฺชีสุ เอวํ น วิชฺชติ. \\ พหุมฺปิ อลชฺชี ภาเสยฺย, วตฺตานุสนฺธิเตน\textsuperscript{1} การเย. \\ อลชฺชี กีทิโส โหติ, ปฏิญฺญา ยสฺส น รูหติ. \\ เอตญฺจ\textsuperscript{1} ตาหํ ปุจฺฉามิ, กีทิโส วุจฺจติ อลชฺชีปุคฺคโล. \\ สญฺจิจฺจ อาปตฺติํ อาปชฺชติ, อาปตฺติํ ปริคูหติ. \\ อคติคมนญฺจ คจฺฉติ, เอทิโส วุจฺจติ อลชฺชีปุคฺคโล. \\ สจฺจํ อหมฺปิ ชานามิ, เอทิโส วุจฺจติ อลชฺชีปุคฺคโล. \\ อญฺญญฺจ ตาหํ ปุจฺฉามิ, กีทิโส วุจฺจติ ลชฺชีปุคฺคโล. \\ สญฺจิจฺจ อาปตฺติํ นาปชฺชติ, อาปตฺติํ น ปริคูหติ. \\ อคติคมนํ น คจฺฉติ, เอทิโส วุจฺจติ ลชฺชีปุคฺคโล.}
\csromangathasetleft{สจฺจํ อหมฺปิ ชานามิ, \\ อญฺญญฺจ ตาหํ ปุจฺฉามิ,}
\csromangathagroup{\csromangathastanza{\csromangathawak{\csromanfolio{279}\csromangathaitemline{359}{โจทนา กิมตฺถาย, สารณา กิสฺส การณา.} \\ \csromangathacontline{สํโฆ กิมตฺถาย, มติกมฺมํ ปน กิสฺส การณา.} \\ \csromangathacontline{โจทนา สารณตฺถาย, นิคฺคหตฺถาย สารณา.} \\ \csromangathacontline{สํโฆ ปริคฺคหตฺถาย, มติกมฺมํ ปน ปาฏิเยกฺกํ.} \\ \csromangathacontline{มา โข ตุริโต อภณิ, มา โข จณฺฑิกโต ภณิ.} \\ \csromangathacontline{มา โข ปฏิฆํ ชนยิ, สเจ อนุวิชฺชโก ตุวํ.} \\ \csromangathacontline{มา โข สหสา อภณิ, กถํ วิคฺคาหิกํ อนตฺถ\-สํหิตํ.} \\ \csromangathacontline{สุตฺเต วินเย อนุโลเม, ปญฺญตฺเต อนุโลมิเก.} \\ \csromangathacontline{อนุโยควตฺตํ นิสามย, กุสเลน พุทฺธิมตา กตํ.} \\ \csromangathacontline{สุวุตฺตํ สิกฺขาปทานุโลมิกํ, คติํ น นาเสนฺโต สมฺปรายิกํ.} \\ \csromangathacontline{หิเตสี อนุยุญฺชสฺสุ, กาเลนตฺถูปสํหิตํ.} \\ \csromangathacontline{จุทิตสฺส จ โจทกสฺส จ,} \\ \csromangathacontline{สหสา โวหารํ มา ปธาเรสิ.} \\ \csromangathacontline{โจทโก อาห อาปนฺโนติ,} \\ \csromangathacontline{จุทิตโก อาห อนาปนฺโนติ.} \\ \csromangathacontline{อุโภ อนุกฺขิปนฺโต, ปฏิญฺญานุสนฺธิเตน การเย.} \\ \csromangathacontline{ปฏิญฺญา ลชฺชีสุ กตา, อลชฺชีสุ เอวํ น วิชฺชติ.} \\ \csromangathacontline{พหุมฺปิ อลชฺชี ภาเสยฺย, วตฺตานุสนฺธิเตน\csromansharedfootnote{b66ff660eb574b37}{วุตฺตานุสนฺธิเตน (สี, สฺยา, ก)} การเย.} \\ \csromangathacontline{อลชฺชี กีทิโส โหติ, ปฏิญฺญา ยสฺส น รูหติ.} \\ \csromangathacontline{เอตญฺจ\csromansharedfootnote{3b1ec6abb7659122}{เอวญฺจ (ก)} ตาหํ ปุจฺฉามิ, กีทิโส วุจฺจติ อลชฺชีปุคฺคโล.} \\ \csromangathacontline{สญฺจิจฺจ อาปตฺติํ อาปชฺชติ, อาปตฺติํ ปริคูหติ.} \\ \csromangathacontline{อคติคมนญฺจ คจฺฉติ, เอทิโส วุจฺจติ อลชฺชีปุคฺคโล.} \\ \csromangathacontline{สจฺจํ อหมฺปิ ชานามิ, เอทิโส วุจฺจติ อลชฺชีปุคฺคโล.} \\ \csromangathacontline{อญฺญญฺจ ตาหํ ปุจฺฉามิ, กีทิโส วุจฺจติ ลชฺชีปุคฺคโล.} \\ \csromangathacontline{สญฺจิจฺจ อาปตฺติํ นาปชฺชติ, อาปตฺติํ น ปริคูหติ.} \\ \csromangathacontline{อคติคมนํ น คจฺฉติ, เอทิโส วุจฺจติ ลชฺชีปุคฺคโล.}}}\csromangathastanza{\csromanfolio{280}\csromangathabat{สจฺจํ อหมฺปิ ชานามิ,}{เอทิโส วุจฺจติ วุจฺจติ ลชฺชีปุคฺคโล.} \\ \csromangathabat{อญฺญญฺจ ตาหํ ปุจฺฉามิ,}{กีทิโส วุจฺจติ อธมฺมโจทโก.}}}

\prose{อกาเลน โจเทติ อภูเตน.}

\prose{ผรุเสน อนตฺถ\-สํหิเตน.}

\prose{โทสนฺตโร โจเทติ โน เมตฺตาจิตฺโต.}

\prose{เอทิโส วุจฺจติ อธมฺมโจทโก.}

\setlength{\csromangathapagewidth}{0pt}
\setlength{\csromangathawakwidth}{0pt}
\csromangathameasuregroup{สจฺจํ อหมฺปิ ชานามิ, \\ อญฺญญฺจ ตาหํ ปุจฺฉามิ, \\ กาเลน โจเทติ ภูเตน, \\ เมตฺตาจิตฺโต โจเทติ โน โทสนฺตโร, \\ สจฺจํ อหมฺปิ ชานามิ, \\ อญฺญญฺจ ตาหํ ปุจฺฉามิ, \\ ปุพฺพาปรํ น ชานาติ, \\ อนุสนฺธิวจนปถํ น ชานาติ, \\ สจฺจํ อหมฺปิ ชานามิ, \\ อญฺญญฺจ ตาหํ ปุจฺฉามิ, \\ สีลวิปตฺติยา โจเทติ, \\ อาชีเวนปิ โจเทติ,}{\csromangathabat{สจฺจํ อหมฺปิ ชานามิ,}{เอทิโส วุจฺจติ อธมฺมโจทโก.} \\ \csromangathabat{อญฺญญฺจ ตาหํ ปุจฺฉามิ,}{กีทิโส วุจฺจติ ธมฺมโจทโก.} \\ \csromangathabat{กาเลน โจเทติ ภูเตน,}{สเณฺหน อตฺถสํหิเตน.} \\ \csromangathabat{เมตฺตาจิตฺโต โจเทติ โน โทสนฺตโร,}{เอทิโส วุจฺจติ ธมฺมโจทโก.} \\ \csromangathabat{สจฺจํ อหมฺปิ ชานามิ,}{เอทิโส วุจฺจติ ธมฺมโจทโก.} \\ \csromangathabat{อญฺญญฺจ ตาหํ ปุจฺฉามิ,}{กีทิโส วุจฺจติ พาลโจทโก.} \\ \csromangathabat{ปุพฺพาปรํ น ชานาติ,}{ปุพฺพาปรสฺส อโกวิโท.} \\ \csromangathabat{อนุสนฺธิวจนปถํ น ชานาติ,}{อนุสนฺธิวจนปถสฺส อโกวิโท.} \\ เอทิโส วุจฺจติ พาลโจทโก. \\ สจฺจํ อหมฺปิ ชานามิ, \\ เอทิโส วุจฺจติ พาลโจทโก. \\ อญฺญญฺจ ตาหํ ปุจฺฉามิ, \\ กีทิโส วุจฺจติ ปณฺฑิตโจทโก. \\ ปุพฺพาปรมฺปิ ชานาติ, \\ ปุพฺพาปรสฺส โกวิโท. \\ อนุสนฺธิวจนปถํ ชานาติ, \\ อนุสนฺธิวจนปถสฺส โกวิโท. \\ เอทิโส วุจฺจติ ปณฺฑิตโจทโก. \\ \csromangathabat{สจฺจํ อหมฺปิ ชานามิ,}{เอทิโส วุจฺจติ ปณฺฑิตโจทโก.} \\ \csromangathabat{อญฺญญฺจ ตาหํ ปุจฺฉามิ,}{โจทนา กินฺติ วุจฺจติ.} \\ \csromangathabat{สีลวิปตฺติยา โจเทติ,}{อโถ อาจารทิฏฺฐิยา.} \\ \csromangathabat{อาชีเวนปิ โจเทติ,}{โจทนา เตน วุจฺจตีติ.}}
\csromangathameasurewak{เอทิโส วุจฺจติ พาลโจทโก. \\ สจฺจํ อหมฺปิ ชานามิ, \\ เอทิโส วุจฺจติ พาลโจทโก. \\ อญฺญญฺจ ตาหํ ปุจฺฉามิ, \\ กีทิโส วุจฺจติ ปณฺฑิตโจทโก. \\ ปุพฺพาปรมฺปิ ชานาติ, \\ ปุพฺพาปรสฺส โกวิโท. \\ อนุสนฺธิวจนปถํ ชานาติ,}
\csromangathasetleft{สจฺจํ อหมฺปิ ชานามิ, \\ อญฺญญฺจ ตาหํ ปุจฺฉามิ, \\ กาเลน โจเทติ ภูเตน, \\ เมตฺตาจิตฺโต โจเทติ โน โทสนฺตโร, \\ สจฺจํ อหมฺปิ ชานามิ, \\ อญฺญญฺจ ตาหํ ปุจฺฉามิ, \\ ปุพฺพาปรํ น ชานาติ, \\ อนุสนฺธิวจนปถํ น ชานาติ, \\ สจฺจํ อหมฺปิ ชานามิ, \\ อญฺญญฺจ ตาหํ ปุจฺฉามิ, \\ สีลวิปตฺติยา โจเทติ, \\ อาชีเวนปิ โจเทติ,}
\csromangathagroupwithcloser{\csromangathastanza{\csromangathabat{สจฺจํ อหมฺปิ ชานามิ,}{เอทิโส วุจฺจติ อธมฺมโจทโก.} \\ \csromangathabat{อญฺญญฺจ ตาหํ ปุจฺฉามิ,}{กีทิโส วุจฺจติ ธมฺมโจทโก.}}\csromangathastanza{\csromangathabat{กาเลน โจเทติ ภูเตน,}{สเณฺหน อตฺถสํหิเตน.} \\ \csromangathabat{เมตฺตาจิตฺโต โจเทติ โน โทสนฺตโร,}{เอทิโส วุจฺจติ ธมฺมโจทโก.}}\csromangathastanza{\csromangathabat{สจฺจํ อหมฺปิ ชานามิ,}{เอทิโส วุจฺจติ ธมฺมโจทโก.} \\ \csromangathabat{อญฺญญฺจ ตาหํ ปุจฺฉามิ,}{กีทิโส วุจฺจติ พาลโจทโก.}}\csromangathastanza{\csromangathabat{ปุพฺพาปรํ น ชานาติ,}{ปุพฺพาปรสฺส อโกวิโท.} \\ \csromangathabat{อนุสนฺธิ\-วจนปถํ น ชานาติ,}{อนุสนฺธิวจนปถสฺส อโกวิโท.}}\csromangathastanza{\csromangathawak{เอทิโส วุจฺจติ พาลโจทโก. \\ สจฺจํ อหมฺปิ ชานามิ, \\ เอทิโส วุจฺจติ พาลโจทโก. \\ อญฺญญฺจ ตาหํ ปุจฺฉามิ,}}\csromangathastanza{\csromangathawak{กีทิโส วุจฺจติ ปณฺฑิตโจทโก. \\ ปุพฺพาปรมฺปิ ชานาติ, \\ ปุพฺพาปรสฺส โกวิโท. \\ อนุสนฺธิ\-วจนปถํ ชานาติ,}}\csromangathastanza{อนุสนฺธิวจนปถสฺส โกวิโท. \\ เอทิโส วุจฺจติ ปณฺฑิตโจทโก. \\ \csromangathabat{สจฺจํ อหมฺปิ ชานามิ,}{เอทิโส วุจฺจติ ปณฺฑิตโจทโก.}}\csromangathastanza{\csromangathabat{อญฺญญฺจ ตาหํ ปุจฺฉามิ,}{โจทนา กินฺติ วุจฺจติ.} \\ \csromangathabat{สีลวิปตฺติยา โจเทติ,}{อโถ อาจารทิฏฺฐิยา.} \\ \csromangathabat{อาชีเวนปิ โจเทติ,}{โจทนา เตน วุจฺจตีติ.}}}{อปรํ คาถา\-สงฺคณิกํ นิฏฺฐิตํ.}
\csromanmatikaanchor{mtk.198}
\csromantocmarknopage{chapter}{โจทนากณฺฑ}{mtk.198}
\bookmark[dest={mtk.198},level=0]{โจทนากณฺฑ}
\chapterheadpageread{\textbf{โจทนากณฺฑ}}
\csromanmatikaanchor{mtk.199}
\csromantocmarkref{chapter}{1. อนุวิชฺชกอนุโยค}{mtk.199}
\bookmark[dest={mtk.199},level=0]{1. อนุวิชฺชกอนุโยค}
\chapterhead{1. \textbf{อนุวิชฺชก อนุโยค}}
\proseitem{360}{\csromanfolio{281}อนุวิชฺชเกน โจทโก ปุจฺฉิตพฺโพ “ยํ โข ตฺวํ อาวุโส อิมํ ภิกฺขุํ โจเทสิ, กิมฺหิ นํ โจเทสิ, สีลวิปตฺติยา วา โจเทสิ, อาจารวิปตฺติยา วา โจเทสิ, ทิฏฺฐิ\-วิปตฺติยา วา โจเทสี”ติ. โส เจ เอวํ วเทยฺย “สีลวิปตฺติยา วา โจเทมิ, อาจารวิปตฺติยา วา โจเทมิ, ทิฏฺฐิ\-วิปตฺติยา วา โจเทมี”ติ. โส เอวมสฺส วจนีโย “ชานาสิ ปนายสฺมา สีลวิปตฺติํ, ชานาสิ อาจารวิปตฺติํ, ชานาสิ ทิฏฺฐิวิปตฺตินฺ”ติ. โส เจ เอวํ วเทยฺย “ชานามิ โข อหํ อาวุโส สีลวิปตฺติํ, ชานามิ อาจารวิปตฺติํ, ชานามิ ทิฏฺฐิวิปตฺตินฺ”ติ. โส เอวมสฺส วจนีโย “กตมา ปนาวุโส สีลวิปตฺติ, กตมา อาจารวิปตฺติ, กตมา ทิฏฺฐิวิปตฺตี”ติ. โส เจ เอวํ วเทยฺย “จตฺตาริ จ ปาราชิกานิ เตรส จ สํฆาทิเสสา, อยํ สีลวิปตฺติ. ถุลฺลจฺจยํ ปาจิตฺติยํ ปาฏิเทสนียํ ทุกฺกฏํ ทุพฺภาสิตํ, อยํ อาจารวิปตฺติ. มิจฺฉาทิฏฺฐิ อนฺตคฺคาหิกา ทิฏฺฐิ, อยํ ทิฏฺฐิวิปตฺตี”ติ. โส เอวมสฺส วจนีโย “ยํ โข ตฺวํ อาวุโส อิมํ ภิกฺขุํ โจเทสิ, ทิฏฺเฐน วา โจเทสิ, สุเตน วา โจเทสิ, ปริสงฺกาย วา โจเทสี”ติ. โส เจ เอวํ วเทยฺย “ทิฏฺเฐน วา โจเทมิ, สุเตน วา โจเทมิ, ปริสงฺกาย วา โจเทมี”ติ. โส เอวมสฺส วจนีโย “ยํ โข ตฺวํ อาวุโส อิมํ ภิกฺขุํ ทิฏฺเฐน โจเทสิ, กิํ เต ทิฏฺฐํ, กินฺติ เต ทิฏฺฐํ, กทา เต ทิฏฺฐํ, กตฺถ เต ทิฏฺฐํ, ปาราชิกํ อชฺฌาปชฺชนฺโต ทิฏฺโฐ, สํฆาทิเสสํ อชฺฌาปชฺชนฺโต ทิฏฺโฐ, ถุลฺลจฺจยํ. ปาจิตฺติยํ. ปาฏิเทสนียํ. ทุกฺกฏํ. ทุพฺภาสิตํ อชฺฌาปชฺชนฺโต ทิฏฺโฐ, กตฺถ จ ตฺวํ อโหสิ, กตฺถ จายํ ภิกฺขุ อโหสิ, กิญฺจ ตฺวํ กโรสิ, กิญฺจายํ ภิกฺขุ กโรตี”ติ. โส เจ เอวํ วเทยฺย “น โข อหํ อาวุโส อิมํ ภิกฺขุํ ทิฏฺเฐน โจเทมิ, อปิ จ สุเตน โจเทมี”ติ. โส เอวมสฺส วจนีโย “ยํ โข ตฺวํ อาวุโส อิมํ ภิกฺขุํ สุเตน โจเทสิ, กิํ เต สุตํ, กินฺติ เต สุตํ, กทา เต สุตํ, กตฺถ เต สุตํ, ‘ปาราชิกํ อชฺฌาปนฺโน’ติ สุตํ, สํฆาทิเสสํ. ถุลฺลจฺจยํ. ปาจิตฺติยํ. ปาฏิเทสนียํ. ทุกฺกฏํ. ‘ทุพฺภาสิตํ อชฺฌาปนฺโน’ติ สุตํ, ภิกฺขุสฺส สุตํ, ภิกฺขุนิยา สุตํ, สิกฺขมานาย \csromanfolio{282}สุตํ, สามเณรสฺส สุตํ, สามเณริยา สุตํ, อุปาสกสฺส สุตํ, อุปาสิกาย สุตํ, ราชูนํ สุตํ, ราช\-มหามตฺตานํ สุตํ, ติตฺถิยานํ สุตํ, ติตฺถิย\-สาวกานํ สุตนฺ”ติ. โส เจ เอวํ วเทยฺย “น โข อหํ อาวุโส อิมํ ภิกฺขุํ สุเตน โจเทมิ, อปิ จ ปริสงฺกาย โจเทมี”ติ. โส เอวมสฺส วจนีโย “ยํ โข ตฺวํ อาวุโส อิมํ ภิกฺขุํ ปริสงฺกาย โจเทสิ, กิํ ปริสงฺกสิ, กินฺติ ปริสงฺกสิ, กทา ปริสงฺกสิ, กตฺถ ปริสงฺกสิ, ‘ปาราชิกํ อชฺฌาปนฺโน’ติ ปริสงฺกสิ, ‘สํฆาทิเสสํ อชฺฌาปนฺโน’ติ ปริสงฺกสิ, ถลฺลจฺจยํ. ปาจิตฺติยํ. ปาฏิเทสนียํ. ทุกฺกฏํ. ‘ทุพฺภาสิตํ อชฺฌาปนฺโน’ติ ปริสงฺกสิ, ภิกฺขุสฺส สุตฺวา ปริสงฺกสิ, ภิกฺขุนิยา สุตฺวา ปริสงฺกสิ, สิกฺขมานาย สุตฺวา ปริสงฺกสิ, สามเณรสฺส สุตฺวา ปริสงฺกสิ, สามเณริยา สุตฺวา ปริสงฺกสิ, อุปาสกสฺส สุตฺวา ปริสงฺกสิ, อุปาสิกาย สุตฺวา ปริสงฺกสิ, ราชูนํ สุตฺวา ปริสงฺกสิ, ราช\-มหามตฺตานํ สุตฺวา ปริสงฺกสิ, ติตฺถิยานํ สุตฺวา ปริสงฺกสิ, ติตฺถิย\-สาวกานํ สุตฺวา ปริสงฺกสี”ติ.}

\setlength{\csromangathapagewidth}{0pt}
\setlength{\csromangathawakwidth}{0pt}
\csromangathameasuregroup{ทิฏฺฐํ ทิฏฺเฐน สเมติ, \\ ทิฏฺฐํ ปฏิจฺจ น อุเปติ, \\ โส ปุคฺคโล ปฏิญฺญาย, \\ สุตํ สุเตน สเมติ, \\ สุตํ ปฏิจฺจ น อุเปติ, \\ โส ปุคฺคโล ปฏิญฺญาย, \\ มุตํ มุเตน สเมติ, \\ มุตํ ปฏิจฺจ น อุเปติ, \\ โส ปุคฺคโล ปฏิญฺญาย,}{\csromangathabat{ทิฏฺฐํ ทิฏฺเฐน สเมติ,}{ทิฏฺเฐน สํสนฺทเต ทิฏฺฐํ.} \\ \csromangathabat{ทิฏฺฐํ ปฏิจฺจ น อุเปติ,}{อสุทฺธปริสงฺกิโต.} \\ \csromangathabat{โส ปุคฺคโล ปฏิญฺญาย,}{กาตพฺโพ เตนุโปสโถ.} \\ \csromangathabat{สุตํ สุเตน สเมติ,}{สุเตน สํสนฺทเต สุตํ.} \\ \csromangathabat{สุตํ ปฏิจฺจ น อุเปติ,}{อสุทฺธปริสงฺกิโต.} \\ \csromangathabat{โส ปุคฺคโล ปฏิญฺญาย,}{กาตพฺโพ เตนุโปสโถ.} \\ \csromangathabat{มุตํ มุเตน สเมติ,}{มุเตน สํสนฺทเต มุตํ.} \\ \csromangathabat{มุตํ ปฏิจฺจ น อุเปติ,}{อสุทฺธปริสงฺกิโต.} \\ \csromangathabat{โส ปุคฺคโล ปฏิญฺญาย,}{กาตพฺโพ เตนุโปสโถ.}}
\csromangathasetleft{ทิฏฺฐํ ทิฏฺเฐน สเมติ, \\ ทิฏฺฐํ ปฏิจฺจ น อุเปติ, \\ โส ปุคฺคโล ปฏิญฺญาย, \\ สุตํ สุเตน สเมติ, \\ สุตํ ปฏิจฺจ น อุเปติ, \\ โส ปุคฺคโล ปฏิญฺญาย, \\ มุตํ มุเตน สเมติ, \\ มุตํ ปฏิจฺจ น อุเปติ, \\ โส ปุคฺคโล ปฏิญฺญาย,}
\csromangathagroup{\csromangathastanza{\csromangathaitembat{361}{ทิฏฺฐํ ทิฏฺเฐน สเมติ,}{ทิฏฺเฐน สํสนฺทเต ทิฏฺฐํ.} \\ \csromangathacontbat{ทิฏฺฐํ ปฏิจฺจ น อุเปติ,}{อสุทฺธปริสงฺกิโต.} \\ \csromangathacontbat{โส ปุคฺคโล ปฏิญฺญาย,}{กาตพฺโพ เตนุโปสโถ.} \\ \csromangathacontbat{สุตํ สุเตน สเมติ,}{สุเตน สํสนฺทเต สุตํ.} \\ \csromangathacontbat{สุตํ ปฏิจฺจ น อุเปติ,}{อสุทฺธปริสงฺกิโต.} \\ \csromangathacontbat{โส ปุคฺคโล ปฏิญฺญาย,}{กาตพฺโพ เตนุโปสโถ.} \\ \csromangathacontbat{มุตํ มุเตน สเมติ,}{มุเตน สํสนฺทเต มุตํ.} \\ \csromangathacontbat{มุตํ ปฏิจฺจ น อุเปติ,}{อสุทฺธปริสงฺกิโต.} \\ \csromangathacontbat{โส ปุคฺคโล ปฏิญฺญาย,}{กาตพฺโพ เตนุโปสโถ.}}}

\proseitem{362}{โจทนาย โก อาทิ, กิํ มชฺเฌ, กิํ ปริโยสานํ. โจทนาย โอกาสกมฺมํ อาทิ, กิริยา มชฺเฌ, สมโถ ปริโยสานํ. โจทนาย กติ มูลานิ, กติ วตฺถูนิ, กติ ภูมิโย, กติหากาเรหิ โจเทติ. โจทนาย เทฺว มูลานิ, ตีณิ วตฺถูนิ, ปญฺจ ภูมิโย, ทฺวีหากาเรหิ โจเทติ. โจทนาย กตมานิ เทฺว มูลานิ. สมูลิกา วา อมูลิกา วา. โจทนาย อิมานิ เทฺว มูลานิ. โจทนาย กตมานิ ตีณิ วตฺถูนิ. \csromanfolio{283}ทิฏฺเฐน สุเตน ปริสงฺกาย. โจทนาย อิมานิ ตีณิ วตฺถูนิ. โจทนาย กตมา ปญฺจ ภูมิโย. กาเลน วกฺขามิ, โน อกาเลน, ภูเตน วกฺขามิ, โน อภูเตน, สเณฺหน วกฺขามิ, โน ผรุเสน, อตฺถสํหิเตน วกฺขามิ, โน อนตฺถ\-สํหิเตน, เมตฺตาจิตฺโต วกฺขามิ, โน โทสนฺตโรติ. โจทนาย อิมา ปญฺจ ภูมิโย.}

\prose{กตเมหิ ทฺวีหากาเรหิ โจเทติ. กาเยน วา โจเทติ, วาจาย วา โจเทติ. อิเมหิ ทฺวีหากาเรหิ โจเทติ.}

\csromanmatikaanchor{mtk.200}
\csromantocmarkref{chapter}{2. โจทกาทิปฏิปตฺติ}{mtk.200}
\bookmark[dest={mtk.200},level=0]{2. โจทกาทิปฏิปตฺติ}
\chapterhead{2. โจทกา\-ทิ\-ปฏิปตฺติ}
\proseitem{363}{โจทเกน กถํ ปฏิปชฺชิตพฺพํ, จุทิตเกน กถํ ปฏิปชฺชิตพฺพํ, สํเฆน กถํ ปฏิปชฺชิตพฺพํ, อนุวิชฺชเกน กถํ ปฏิปชฺชิตพฺพํ. โจทเกน กถํ ปฏิปชฺชิตพฺพนฺติ. โจทเกน ปญฺจสุ ธมฺเมสุ ปติฏฺฐาย ปโร โจเทตพฺโพ. กาเลน วกฺขามิ, โน อกาเลน, ภูเตน วกฺขามิ, โน อภูเตน, สเณฺหน วกฺขามิ, โน ผรุเสน, อตฺถสํหิเตน วกฺขามิ, โน อนตฺถ\-สํหิเตน, เมตฺตาจิตฺโต วกฺขามิ, โน โทสนฺตโรติ. โจทเกน เอวํ ปฏิปชฺชิตพฺพํ. จุทิตเกน กถํ ปฏิปชฺชิตพฺพนฺติ. จุทิตเกน ทฺวีสุ ธมฺเมสุ ปฏิปชฺชิตพฺพํ สจฺเจ จ อกุปฺเป จ. จุทิตเกน เอวํ ปฏิปชฺชิตพฺพํ. สํเฆน กถํ ปฏิปชฺชิตพฺพนฺติ. สํเฆน โอติณฺณาโนติณฺณํ ชานิตพฺพํ. สํเฆน เอวํ ปฏิปชฺชิตพฺพํ. อนุวิชฺชเกน กถํ ปฏิปชฺชิตพฺพนฺติ. อนุวิชฺชเกน เยน ธมฺเมน เยน วินเยน เยน สตฺถุสาสเนน ตํ อธิกรณํ วูปสมฺมติ, ตถา ตํ อธิกรณํ วูปสเมตพฺพํ. อนุวิชฺชเกน เอวํ ปฏิปชฺชิตพฺพํ.}

\setlength{\csromangathapagewidth}{0pt}
\setlength{\csromangathawakwidth}{0pt}
\csromangathameasuregroup{อุโปสโถ กิมตฺถาย, \\ ปริวาโส กิมตฺถาย, \\ มานตฺตํ กิมตฺถาย, \\ อุโปสโถ สามคฺคตฺถาย, \\ ปริวาโส มานตฺตตฺถาย, \\ มานตฺตํ อพฺภานตฺถาย, \\ ฉนฺทา โทสา ภยา โมหา, \\ กายสฺส เภทา ทุปฺปญฺโญ, \\ นิรยํ คจฺฉติ ทุมฺเมโธ, \\ น จ อามิสํ นิสฺสาย,}{\csromangathabat{อุโปสโถ กิมตฺถาย,}{ปวารณา กิสฺส การณา.} \\ \csromangathabat{ปริวาโส กิมตฺถาย,}{มูลายปฏิกสฺสนา กิสฺส การณา.} \\ \csromangathabat{มานตฺตํ กิมตฺถาย,}{อพฺภานํ กิสฺส การณา.} \\ \csromangathabat{อุโปสโถ สามคฺคตฺถาย,}{วิสุทฺธตฺถาย ปวารณา.} \\ \csromangathabat{ปริวาโส มานตฺตตฺถาย,}{มูลายปฏิกสฺสนา นิคฺคหตฺถาย.} \\ \csromangathabat{มานตฺตํ อพฺภานตฺถาย,}{วิสุทฺธตฺถาย อพฺภานํ.} \\ \csromangathabat{ฉนฺทา โทสา ภยา โมหา,}{เถเร จ ปริภาสติ.} \\ \csromangathabat{กายสฺส เภทา ทุปฺปญฺโญ,}{ขโต อุปหตินฺทฺริโย.} \\ \csromangathabat{นิรยํ คจฺฉติ ทุมฺเมโธ,}{น จ สิกฺขาย คารโว.} \\ \csromangathabat{น จ อามิสํ นิสฺสาย,}{น จ นิสฺสาย ปุคฺคลํ.} \\ อุโภ เอเต วิวชฺเชตฺวา, \\ ยถาธมฺโม ตถา กเร.}
\csromangathameasurewak{อุโภ เอเต วิวชฺเชตฺวา, \\ ยถาธมฺโม ตถา กเร.}
\csromangathasetleft{อุโปสโถ กิมตฺถาย, \\ ปริวาโส กิมตฺถาย, \\ มานตฺตํ กิมตฺถาย, \\ อุโปสโถ สามคฺคตฺถาย, \\ ปริวาโส มานตฺตตฺถาย, \\ มานตฺตํ อพฺภานตฺถาย, \\ ฉนฺทา โทสา ภยา โมหา, \\ กายสฺส เภทา ทุปฺปญฺโญ, \\ นิรยํ คจฺฉติ ทุมฺเมโธ, \\ น จ อามิสํ นิสฺสาย,}
\csromangathagroup{\csromangathastanza{\csromangathaitembat{364}{อุโปสโถ กิมตฺถาย,}{ปวารณา กิสฺส การณา.} \\ \csromangathacontbat{ปริวาโส กิมตฺถาย,}{มูลาย\-ปฏิกสฺสนา กิสฺส การณา.} \\ \csromangathacontbat{มานตฺตํ กิมตฺถาย,}{อพฺภานํ กิสฺส การณา.} \\ \csromangathacontbat{อุโปสโถ สามคฺคตฺถาย,}{วิสุทฺธตฺถาย ปวารณา.} \\ \csromangathacontbat{ปริวาโส มานตฺตตฺถาย,}{มูลาย\-ปฏิกสฺสนา นิคฺคหตฺถาย.} \\ \csromangathacontbat{มานตฺตํ อพฺภานตฺถาย,}{วิสุทฺธตฺถาย อพฺภานํ.}}\csromangathastanza{\csromanfolio{284}\csromangathabat{ฉนฺทา โทสา ภยา โมหา,}{เถเร จ ปริภาสติ.} \\ \csromangathabat{กายสฺส เภทา ทุปฺปญฺโญ,}{ขโต อุปหตินฺทฺริโย.}}\csromangathastanza{\csromangathabat{นิรยํ คจฺฉติ ทุมฺเมโธ,}{น จ สิกฺขาย คารโว.} \\ \csromangathabat{น จ อามิสํ นิสฺสาย,}{น จ นิสฺสาย ปุคฺคลํ.}}\csromangathastanza{\csromangathawak{อุโภ เอเต วิวชฺเชตฺวา, \\ ยถาธมฺโม ตถา กเร.}}}

\csromanmatikaanchor{mtk.201}
\csromantocmarkref{chapter}{3. โจทกสฺสอตฺตฌาปน}{mtk.201}
\bookmark[dest={mtk.201},level=0]{3. โจทกสฺสอตฺตฌาปน}
\chapterhead{3. โจทกสฺส\-อตฺต\-ฌาปน}
\setlength{\csromangathapagewidth}{0pt}
\setlength{\csromangathawakwidth}{0pt}
\csromangathameasuregroup{อกาเลน โจเทติ อภูเตน,}{โกธโน อุปนาหี จ, \\ จณฺโฑ จ ปริภาสโก. \\ อนาปตฺติยา ‘อาปตฺตี’ติ โรเปติ, \\ ตาทิโส โจทโก ฌาเปติ อตฺตานํ. \\ อุปกณฺณกํ ชปฺปติ ชิมฺหํ เปกฺขติ, \\ วีติหรติ กุมฺมคฺคํ ปฏิเสวติ. \\ อนาปตฺติยา ‘อาปตฺตี’ติ โรเปติ, \\ ตาทิโส โจทโก ฌาเปติ อตฺตานํ. \\ \csromangathabat{อกาเลน โจเทติ อภูเตน,}{ผรุเสน อนตฺถสํหิเตน.} \\ โทสนฺตโร โจเทติ โน เมตฺตาจิตฺโต. \\ อนาปตฺติยา ‘อาปตฺตี’ติ โรเปติ, \\ ตาทิโส โจทโก ฌาเปติ อตฺตานํ. \\ ธมฺมาธมฺมํ น ชานาติ, \\ ธมฺมาธมฺมสฺส อโกวิโท. \\ อนาปตฺติยา ‘อาปตฺตี’ติ โรเปติ, \\ ตาทิโส โจทโก ฌาเปติ อตฺตานํ. \\ วินยาวินยํ น ชานาติ, \\ วินยาวินยสฺส อโกวิโท. \\ อนาปตฺติยา ‘อาปตฺตี’ติ โรเปติ, \\ ตาทิโส โจทโก ฌาเปติ อตฺตานํ. \\ ภาสิตาภาสิตํ น ชานาติ, \\ ภาสิตาภาสิตสฺส อโกวิโท. \\ อนาปตฺติยา ‘อาปตฺตี’ติ โรเปติ, \\ ตาทิโส โจทโก ฌาเปติ อตฺตานํ. \\ อาจิณฺณานาจิณฺณํ น ชานาติ, \\ อาจิณฺณานาจิณฺณสฺส อโกวิโท. \\ อนาปตฺติยา ‘อาปตฺตี’ติ โรเปติ, \\ ตาทิโส โจทโก ฌาเปติ อตฺตานํ. \\ ปญฺญตฺตาปญฺญตฺตํ น ชานาติ, \\ ปญฺญตฺตาปญฺญตฺตสฺส อโกวิโท. \\ อนาปตฺติยา ‘อาปตฺตี’ติ โรเปติ, \\ ตาทิโส โจทโก ฌาเปติ อตฺตานํ. \\ อาปตฺตานาปตฺติํ น ชานาติ, \\ อาปตฺตานาปตฺติยา อโกวิโท. \\ อนาปตฺติยา ‘อาปตฺตี’ติ โรเปติ, \\ ตาทิโส เจทโก ฌาเปติ อตฺตานํ. \\ ลหุกครุกํ น ชานาติ, \\ ลหุกครุกสฺส อโกวิโท. \\ อนาปตฺติยา ‘อาปตฺตี’ติ โรเปติ, \\ ตาทิโส โจทโก ฌาเปติ อตฺตานํ. \\ สาวเสสานวเสสํ น ชานาติ, \\ สาวเสสานวเสสสฺส อโกวิโท. \\ อนาปตฺติยา ‘อาปตฺตี’ติ โรเปติ, \\ ตาทิโส โจทโก ฌาเปติ อตฺตานํ. \\ ทุฏฺฐุลฺลาทุฏฺฐุลฺลํ น ชานาติ, \\ ทุฏฺฐุลฺลาทุฏฺฐุลฺลสฺส อโกวิโท. \\ อนาปตฺติยา ‘อาปตฺตี’ติ โรเปติ, \\ ตาทิโส โจทโก ฌาเปติ อตฺตานํ. \\ ปุพฺพาปรํ น ชานาติ, \\ ปุพฺพาปรสฺส อโกวิโท. \\ อนาปตฺติยา ‘อาปตฺตี’ติ โรเปติ, \\ ตาทิโส โจทโก ฌาเปติ อตฺตานํ. \\ อนุสนฺธิวจนปถํ น ชานาติ, \\ อนุสนฺธิวจนปถสฺส อโกวิโท. \\ อนาปตฺติยา ‘อาปตฺตี’ติ โรเปติ, \\ ตาทิโส โจทโก ฌาเปติ อตฺตานนฺติ.}
\csromangathameasurewak{โกธโน อุปนาหี จ, \\ จณฺโฑ จ ปริภาสโก. \\ อนาปตฺติยา ‘อาปตฺตี’ติ โรเปติ, \\ ตาทิโส โจทโก ฌาเปติ อตฺตานํ. \\ อุปกณฺณกํ ชปฺปติ ชิมฺหํ เปกฺขติ, \\ วีติหรติ กุมฺมคฺคํ ปฏิเสวติ. \\ อนาปตฺติยา ‘อาปตฺตี’ติ โรเปติ, \\ ตาทิโส โจทโก ฌาเปติ อตฺตานํ. \\ ตาทิโส โจทโก ฌาเปติ อตฺตานํ. \\ ธมฺมาธมฺมํ น ชานาติ, \\ ธมฺมาธมฺมสฺส อโกวิโท. \\ อนาปตฺติยา ‘อาปตฺตี’ติ โรเปติ, \\ ตาทิโส โจทโก ฌาเปติ อตฺตานํ. \\ วินยาวินยํ น ชานาติ, \\ วินยาวินยสฺส อโกวิโท. \\ อนาปตฺติยา ‘อาปตฺตี’ติ โรเปติ, \\ ตาทิโส โจทโก ฌาเปติ อตฺตานํ. \\ ภาสิตาภาสิตํ น ชานาติ, \\ ภาสิตาภาสิตสฺส อโกวิโท. \\ อนาปตฺติยา ‘อาปตฺตี’ติ โรเปติ, \\ ตาทิโส โจทโก ฌาเปติ อตฺตานํ. \\ อาจิณฺณานาจิณฺณํ น ชานาติ, \\ อาจิณฺณานาจิณฺณสฺส อโกวิโท. \\ อนาปตฺติยา ‘อาปตฺตี’ติ โรเปติ, \\ ตาทิโส โจทโก ฌาเปติ อตฺตานํ. \\ ปญฺญตฺตาปญฺญตฺตํ น ชานาติ, \\ ปญฺญตฺตาปญฺญตฺตสฺส อโกวิโท. \\ อนาปตฺติยา ‘อาปตฺตี’ติ โรเปติ, \\ ตาทิโส โจทโก ฌาเปติ อตฺตานํ. \\ อาปตฺตานาปตฺติํ น ชานาติ, \\ อาปตฺตานาปตฺติยา อโกวิโท. \\ อนาปตฺติยา ‘อาปตฺตี’ติ โรเปติ, \\ ตาทิโส เจทโก ฌาเปติ อตฺตานํ. \\ ลหุกครุกํ น ชานาติ, \\ ลหุกครุกสฺส อโกวิโท. \\ อนาปตฺติยา ‘อาปตฺตี’ติ โรเปติ, \\ ตาทิโส โจทโก ฌาเปติ อตฺตานํ. \\ สาวเสสานวเสสํ น ชานาติ, \\ สาวเสสานวเสสสฺส อโกวิโท. \\ อนาปตฺติยา ‘อาปตฺตี’ติ โรเปติ, \\ ตาทิโส โจทโก ฌาเปติ อตฺตานํ. \\ ทุฏฺฐุลฺลาทุฏฺฐุลฺลํ น ชานาติ, \\ ทุฏฺฐุลฺลาทุฏฺฐุลฺลสฺส อโกวิโท. \\ อนาปตฺติยา ‘อาปตฺตี’ติ โรเปติ, \\ ตาทิโส โจทโก ฌาเปติ อตฺตานํ. \\ ปุพฺพาปรํ น ชานาติ, \\ ปุพฺพาปรสฺส อโกวิโท. \\ อนาปตฺติยา ‘อาปตฺตี’ติ โรเปติ, \\ ตาทิโส โจทโก ฌาเปติ อตฺตานํ. \\ อนุสนฺธิวจนปถํ น ชานาติ, \\ อนุสนฺธิวจนปถสฺส อโกวิโท. \\ อนาปตฺติยา ‘อาปตฺตี’ติ โรเปติ, \\ ตาทิโส โจทโก ฌาเปติ อตฺตานนฺติ.}
\csromangathasetleft{อกาเลน โจเทติ อภูเตน,}
\csromangathagroupwithclosermajor{\csromangathastanza{\csromangathawak{โกธโน อุปนาหี จ, \\ จณฺโฑ จ ปริภาสโก. \\ อนาปตฺติยา ‘อาปตฺตี’ติ โรเปติ, \\ ตาทิโส โจทโก ฌาเปติ อตฺตานํ.}}\csromangathastanza{\csromangathawak{อุปกณฺณกํ ชปฺปติ ชิมฺหํ เปกฺขติ, \\ วีติหรติ กุมฺมคฺคํ ปฏิเสวติ. \\ อนาปตฺติยา ‘อาปตฺตี’ติ โรเปติ, \\ ตาทิโส โจทโก ฌาเปติ อตฺตานํ.}}\csromangathastanza{\csromangathabat{อกาเลน โจเทติ อภูเตน,}{ผรุเสน อนตฺถ\-สํหิเตน.} \\ โทสนฺตโร โจเทติ โน เมตฺตาจิตฺโต. \\ อนาปตฺติยา ‘อาปตฺตี’ติ โรเปติ,}\csromangathastanza{\csromangathawak{ตาทิโส โจทโก ฌาเปติ อตฺตานํ. \\ ธมฺมาธมฺมํ น ชานาติ, \\ ธมฺมาธมฺมสฺส อโกวิโท. \\ อนาปตฺติยา ‘อาปตฺตี’ติ โรเปติ,}}\csromangathastanza{\csromangathawak{ตาทิโส โจทโก ฌาเปติ อตฺตานํ. \\ วินยาวินยํ น ชานาติ, \\ วินยาวินยสฺส อโกวิโท. \\ อนาปตฺติยา ‘อาปตฺตี’ติ โรเปติ,}}\csromangathastanza{\csromangathawak{ตาทิโส โจทโก ฌาเปติ อตฺตานํ. \\ ภาสิตาภาสิตํ น ชานาติ, \\ ภาสิตาภาสิตสฺส อโกวิโท. \\ อนาปตฺติยา ‘อาปตฺตี’ติ โรเปติ,}}\csromangathastanza{\csromangathawak{\csromanfolio{285}ตาทิโส โจทโก ฌาเปติ อตฺตานํ. \\ อาจิณฺณานาจิณฺณํ น ชานาติ, \\ อาจิณฺณานาจิณฺณสฺส อโกวิโท. \\ อนาปตฺติยา ‘อาปตฺตี’ติ โรเปติ,}}\csromangathastanza{\csromangathawak{ตาทิโส โจทโก ฌาเปติ อตฺตานํ. \\ ปญฺญตฺตา\-ปญฺญตฺตํ น ชานาติ, \\ ปญฺญตฺตาปญฺญตฺตสฺส อโกวิโท. \\ อนาปตฺติยา ‘อาปตฺตี’ติ โรเปติ,}}\csromangathastanza{\csromangathawak{ตาทิโส โจทโก ฌาเปติ อตฺตานํ. \\ อาปตฺตานาปตฺติํ น ชานาติ, \\ อาปตฺตานาปตฺติยา อโกวิโท. \\ อนาปตฺติยา ‘อาปตฺตี’ติ โรเปติ,}}\csromangathastanza{\csromangathawak{ตาทิโส เจทโก ฌาเปติ อตฺตานํ. \\ ลหุกครุกํ น ชานาติ, \\ ลหุกครุกสฺส อโกวิโท. \\ อนาปตฺติยา ‘อาปตฺตี’ติ โรเปติ,}}\csromangathastanza{\csromangathawak{ตาทิโส โจทโก ฌาเปติ อตฺตานํ. \\ สาวเสสา\-นวเสสํ น ชานาติ, \\ สาวเสสานวเสสสฺส อโกวิโท. \\ อนาปตฺติยา ‘อาปตฺตี’ติ โรเปติ,}}\csromangathastanza{\csromangathawak{ตาทิโส โจทโก ฌาเปติ อตฺตานํ. \\ ทุฏฺฐุลฺลา\-ทุฏฺฐุลฺลํ น ชานาติ, \\ ทุฏฺฐุลฺลาทุฏฺฐุลฺลสฺส อโกวิโท. \\ อนาปตฺติยา ‘อาปตฺตี’ติ โรเปติ,}}\csromangathastanza{\csromangathawak{ตาทิโส โจทโก ฌาเปติ อตฺตานํ. \\ ปุพฺพาปรํ น ชานาติ, \\ ปุพฺพาปรสฺส อโกวิโท. \\ อนาปตฺติยา ‘อาปตฺตี’ติ โรเปติ,}}\csromangathastanza{\csromangathawak{\csromanfolio{286}ตาทิโส โจทโก ฌาเปติ อตฺตานํ. \\ อนุสนฺธิ\-วจนปถํ น ชานาติ, \\ อนุสนฺธิวจนปถสฺส อโกวิโท. \\ อนาปตฺติยา ‘อาปตฺตี’ติ โรเปติ,}}\csromangathastanza{\csromangathawak{ตาทิโส โจทโก ฌาเปติ อตฺตานนฺติ.}}}{โจทนากณฺฑํ นิฏฺฐิตํ.}
\csromancenter{ตสฺสุทฺทานํ.}
\setlength{\csromangathapagewidth}{0pt}
\setlength{\csromangathawakwidth}{0pt}
\csromangathameasuregroup{โจทนา อนุวิชฺชา จ, \\ คติ โจทนกณฺฑมฺหิ,}{\csromangathabat{โจทนา อนุวิชฺชา จ,}{อาทิ มูเลนุโปสโถ.} \\ \csromangathabat{คติ โจทนกณฺฑมฺหิ,}{สาสนํ ปติฏฺฐาปยนฺติ.}}
\csromangathasetleft{โจทนา อนุวิชฺชา จ, \\ คติ โจทนกณฺฑมฺหิ,}
\csromangathagroup{\csromangathastanza{\csromangathabat{โจทนา อนุวิชฺชา จ,}{อาทิ มูเลนุโปสโถ.} \\ \csromangathabat{คติ โจทนกณฺฑมฺหิ,}{สาสนํ ปติฏฺฐาปยนฺติ.}}}

\csromanmatikaanchor{mtk.202}
\csromantocmarknopage{section}{จูฬสงฺคาม}{mtk.202}
\bookmark[dest={mtk.202},level=1]{จูฬสงฺคาม}
\csromanheader{1}{\textbf{จูฬสงฺคาม}}
\csromanmatikaanchor{mtk.203}
\csromantocmarkref{section}{1. อนุวิชฺชกสฺสปฏิปตฺติ}{mtk.203}
\bookmark[dest={mtk.203},level=1]{1. อนุวิชฺชกสฺสปฏิปตฺติ}
\csromanheader{1}{1. อนุวิชฺชกสฺส\-ปฏิปตฺติ}
\proseitem{365}{\csromanfolio{287}\csromansymbolfootnote{*}{วิ 5. 318 ปิฏฺเฐปิ.}สงฺคามา\-วจเรน ภิกฺขุนา สํฆํ อุปส\-งฺกม\-นฺเตน นีจจิตฺเตน สํโฆ อุปสงฺกมิตพฺโพ รโชหรณ\-สเมน จิตฺเตน, อาสนกุสเลน ภวิตพฺพํ นิสชฺช\-กุสเลน, เถเร ภิกฺขู อนุปขชฺชนฺเตน, นเว ภิกฺขู อาสเนน อปฺปฏิพาหนฺเตน ยถาปติรูเป อาสเน นิสีทิตพฺพํ, อนานากถิเกน ภวิตพฺพํ อติรจฺฉานกถิเกน, สามํ วา ธมฺโม ภาสิตพฺโพ, ปโร วา อชฺเฌสิตพฺโพ, อริโย วา ตุณฺหีภาโว นาติมญฺญิตพฺโพ.}

\prose{สํเฆน อนุมเตน ปุคฺคเลน อนุวิชฺชเกน อนุวิชฺชิตุกาเมน น อุปชฺฌาโย ปุจฺฉิตพฺโพ, น อาจริโย ปุจฺฉิตพฺโพ, น สทฺธิวิหาริโก ปุจฺฉิตพฺโพ, น อนฺเตวาสิโก ปุจฺฉิตพฺโพ, น สมานุ\-ปชฺฌายโก ปุจฺฉิตพฺโพ, น สมานาจริยโก ปุจฺฉิตพฺโพ, น ชาติ ปุจฺฉิตพฺพา, น นามํ ปุจฺฉิตพฺพํ, น โคตฺตํ ปุจฺฉิตพฺพํ, น อาคโม ปุจฺฉิตพฺโพ, น กุลปเทโส ปุจฺฉิตพฺโพ, น ชาติภูมิ ปุจฺฉิตพฺพา, ตํ กิํ การณา, อตฺรสฺส เปมํ วา โทโส วา, เปเม วา สติ โทเส วา ฉนฺทาปิ คจฺเฉยฺย, โทสาปิ คจฺเฉยฺย, โมหาปิ คจฺเฉยฺย, ภยาปิ คจฺเฉยฺย.}

\prose{สํเฆน อนุมเตน ปุคฺคเลน อนุวิชฺชเกน อนุวิชฺชิตุกาเมน สํฆครุเกน ภวิตพฺพํ, โน ปุคฺคล\-ครุเกน, สทฺธมฺมครุเกน ภวิตพฺพํ, โน อามิสครุเกน, อตฺถวสิเกน ภวิตพฺพํ, โน ปริส\-กปฺปิเกน, กาเลน อนุวิชฺชิตพฺพํ, โน อกาเลน, ภูเตน อนุวิชฺชิตพฺพํ, โน อภูเตน, สเณฺหน อนุวิชฺชิตพฺพํ, โน ผรุเสน, อตฺถสํหิเตน อนุวิชฺชิตพฺพํ, โน อนตฺถ\-สํหิเตน, เมตฺตาจิตฺเตน อนุวิชฺชิตพฺพํ, โน โทสนฺตเรน, น อุปกณฺณก\-ชปฺปินา ภวิตพฺพํ, น ชิมฺหํ เปกฺขิตพฺพํ, น อกฺขิ นิขณิตพฺพํ, น ภมุกํ อุกฺขิปิตพฺพํ, น สีสํ อุกฺขิปิตพฺพํ, น หตฺถวิกาโร กาตพฺโพ, น หตฺถมุทฺทา ทสฺเสตพฺพา.}

\prose{อาสนกุสเลน ภวิตพฺพํ นิสชฺช\-กุสเลน, ยุคมตฺตํ เปกฺขนฺเตน อตฺถํ อนุวิธิ\-ยนฺเตน สเก อาสเน นิสีทิตพฺพํ, น จ อาสนา \csromanfolio{288}วุฏฺฐาตพฺพํ, น วีติหาตพฺพํ, น กุมฺมคฺโค เสวิตพฺโพ, น พาหา\-วิกฺเขปกํ\csromansharedfootnote{cc2705a3ca395c95}{น วาจา วิกฺเขปกํ (สฺยา)} ภณิตพฺพํ, อตุริเตน ภวิตพฺพํ อสาหสิเกน, อจณฺฑิกเตน ภวิตพฺพํ วจนกฺขเมน, เมตฺตาจิตฺเตน ภวิตพฺพํ หิตานุกมฺปินา, การุณิเกน ภวิตพฺพํ หิตปริสกฺกินา, อสมฺผปฺปลาปินา ภวิตพฺพํ ปริยนฺต\-ภาณินา, อเวรวสิเกน ภวิตพฺพํ อนสุรุตฺเตน, อตฺตา ปริคฺคเหตพฺโพ, ปโร ปริคฺคเหตพฺโพ, โจทโก ปริคฺคเหตพฺโพ, จุทิตโก ปริคฺคเหตพฺโพ, อธมฺมโจทโก ปริคฺคเหตพฺโพ, อธมฺม\-จุทิตโก ปริคฺคเหตพฺโพ, ธมฺมโจทโก ปริคฺคเหตพฺโพ, ธมฺมจุทิตโก ปริคฺคเหตพฺโพ, วุตฺตํ อหาเปนฺเตน อวุตฺตํ อปกาเสนฺเตน โอติณฺณานิ ปทพฺยญฺชนานิ สาธุกํ ปริคฺคเหตฺวา ปโร ปฏิปุจฺฉิตฺวา ยถา ปฏิญฺญาย กาเรตพฺโพ, มนฺโท หาเสตพฺโพ\csromansharedfootnote{e0522a150f31e524}{เวโป ปหาเสตพฺโพ (สฺยา)}, ภีรู อสฺสาเสตพฺโพ, จณฺโฑ นิเสเธตพฺโพ, อสุจิ วิภาเวตพฺโพ, อุชุมทฺทเวน, น ฉนฺทาคติํ คนฺตพฺพํ, น โทสาคติํ คนฺตพฺพํ, น โมหาคติํ คนฺตพฺพํ, น ภยาคติํ คนฺตพฺพํ, มชฺฌตฺเตน ภวิตพฺพํ ธมฺเมสุ จ ปุคฺคเลสุ จ, เอวญฺจ ปน อนุวิชฺชโก อนุวิชฺชมาโน สตฺถุ เจว สาสนกโร โหติ, วิญฺญูนญฺจ สพฺรหฺมจารีนํ ปิโย จ โหติ มนาโป จ ครุ จ ภาวนีโย จ.}

\proseitem{366}{สุตฺตํ สํสนฺทน\-ตฺถาย, โอปมฺมํ นิทสฺสนตฺถาย, อตฺโถ วิญฺญาปน\-ตฺถาย, ปฏิปุจฺฉา ฐปนตฺถาย, โอกาสกมฺมํ โจทนตฺถาย, โจทนา สารณตฺถาย, สารณา สวจนียตฺถาย, สวจนียํ ปลิโพธ\-ตฺถาย, ปลิโพโธ วินิจฺฉย\-ตฺถาย, วินิจฺฉโย สนฺตีรณ\-ตฺถาย, สนฺตีรณํ ฐานาฐาน\-คมนตฺถาย, ฐานาฐาน\-คมนํ ทุมฺมงฺกูนํ ปุคฺคลานํ นิคฺคหตฺถาย, เปสลานํ ภิกฺขูนํ สมฺปคฺคหตฺถาย, สํโฆ สมฺปริคฺคหสมฺปฏิจฺฉนตฺถาย, สํเฆน อนุมตา ปุคฺคลา ปจฺเจก\-ฏฺฐายิโน อวิสํวาทก\-ฏฺฐายิโน.}

\prose{วินโย สํวรตฺถาย, สํวโร อวิปฺปฏิสาร\-ตฺถาย, อวิปฺปฏิสาโร ปามุชฺชตฺถาย, ปามุชฺชํ ปีตตฺถาย, ปีติ ปสฺสทฺธ\-ตฺถาย, ปสฺสทฺธิ สุขตฺถาย, สุขํ สมาธตฺถาย, สมาธิ ยถา\-ภูต\-ญาณ\-ทสฺสน\-ตฺถาย, ยถา\-ภูต\-ญาณ\-ทสฺสนํ นิพฺพิทตฺถาย, นิพฺพิทา วิราคตฺถาย, วิราโค วิมุตฺตตฺถาย, วิมุตฺติ วิมุตฺติ\-ญาณ\-ทสฺสน\-ตฺถาย, วิมุตฺติ\-ญาณ\-ทสฺสนํ อนุปาทา\-ปรินิพฺพาน\-ตฺถาย, \csromanfolio{289}เอตทตฺถา กถา, เอตทตฺถา มนฺตนา, เอตทตฺถา อุปนิสา, เอตทตฺถํ โสตาวธานํ, ยทิทํ อนุปาทา\-จิตฺตสฺส วิโมกฺโขติ.}

\setlength{\csromangathapagewidth}{0pt}
\setlength{\csromangathawakwidth}{0pt}
\csromangathameasuregroup{อนุโยควตฺตํ นิสามย, \\ สุวุตฺตํ สิกฺขาปทานุโลมิกํ, \\ วตฺถุํ วิปตฺติํ อาปตฺติํ, \\ ปุพฺพาปรํ น ชานาติ, \\ กมฺมญฺจ อธิกรณญฺจ, \\ รตฺโต ทุฏฺโฐ จ มูโฬฺห จ, \\ น จ สญฺญตฺติกุสโล, \\ ลทฺธปกฺโข อหิริโก, \\ ส เว\textsuperscript{1} ตาทิสโก ภิกฺขุ, \\ วตฺถุํ วิปตฺติํ อาปตฺติํ, \\ ปุพฺพาปรญฺจ ชานาติ, \\ กมฺมญฺจ อธิกรณญฺจ, \\ อรตฺโต อทุฏฺโฐ อมูโฬฺห, \\ สญฺญตฺติยา จ กุสโล, \\ ลทฺธปกฺโข หิริมโน, \\ ส เว ตาทิสโก ภิกฺขุ,}{\csromangathabat{อนุโยควตฺตํ นิสามย,}{กุสเลน พุทฺธิมตา กตํ.} \\ \csromangathabat{สุวุตฺตํ สิกฺขาปทานุโลมิกํ,}{คติํ น นาเสนฺโต สมฺปรายิกํ.} \\ \csromangathabat{วตฺถุํ วิปตฺติํ อาปตฺติํ,}{นิทานํ อาการอโกวิโท.} \\ \csromangathabat{ปุพฺพาปรํ น ชานาติ,}{กตากตํ สเมน จ.} \\ \csromangathabat{กมฺมญฺจ อธิกรณญฺจ,}{สมเถ จาปิ อโกวิโท.} \\ \csromangathabat{รตฺโต ทุฏฺโฐ จ มูโฬฺห จ,}{ภยา โมหา จ คจฺฉติ.} \\ \csromangathabat{น จ สญฺญตฺติกุสโล,}{นิชฺฌตฺติยา จ อโกวิโท.} \\ \csromangathabat{ลทฺธปกฺโข อหิริโก,}{กณฺหกมฺโม อนาทโร.} \\ \csromangathabat{ส เว\textsuperscript{1} ตาทิสโก ภิกฺขุ,}{‘อปฺปฏิกฺโข’ติ วุจฺจติ.} \\ \csromangathabat{วตฺถุํ วิปตฺติํ อาปตฺติํ,}{นิทานํ อาการโกวิโท.} \\ \csromangathabat{ปุพฺพาปรญฺจ ชานาติ,}{กตากตํ สเมน จ.} \\ \csromangathabat{กมฺมญฺจ อธิกรณญฺจ,}{สมเถ จาปิ โกวิโท.} \\ \csromangathabat{อรตฺโต อทุฏฺโฐ อมูโฬฺห,}{ภยา โมหา น คจฺฉติ.} \\ \csromangathabat{สญฺญตฺติยา จ กุสโล,}{นิชฺฌตฺติยา จ โกวิโท.} \\ \csromangathabat{ลทฺธปกฺโข หิริมโน,}{สุกฺกกมฺโม สคารโว.} \\ \csromangathabat{ส เว ตาทิสโก ภิกฺขุ,}{‘สปฺปฏิกฺโข’ติ วุจฺจตีติ.}}
\csromangathasetleft{อนุโยควตฺตํ นิสามย, \\ สุวุตฺตํ สิกฺขาปทานุโลมิกํ, \\ วตฺถุํ วิปตฺติํ อาปตฺติํ, \\ ปุพฺพาปรํ น ชานาติ, \\ กมฺมญฺจ อธิกรณญฺจ, \\ รตฺโต ทุฏฺโฐ จ มูโฬฺห จ, \\ น จ สญฺญตฺติกุสโล, \\ ลทฺธปกฺโข อหิริโก, \\ ส เว\textsuperscript{1} ตาทิสโก ภิกฺขุ, \\ วตฺถุํ วิปตฺติํ อาปตฺติํ, \\ ปุพฺพาปรญฺจ ชานาติ, \\ กมฺมญฺจ อธิกรณญฺจ, \\ อรตฺโต อทุฏฺโฐ อมูโฬฺห, \\ สญฺญตฺติยา จ กุสโล, \\ ลทฺธปกฺโข หิริมโน, \\ ส เว ตาทิสโก ภิกฺขุ,}
\csromangathagroupwithcloser{\csromangathastanza{\csromangathaitembat{367}{อนุโยควตฺตํ นิสามย,}{กุสเลน พุทฺธิมตา กตํ.} \\ \csromangathacontbat{สุวุตฺตํ สิกฺขาปทานุโลมิกํ,}{คติํ น นาเสนฺโต สมฺปรายิกํ.} \\ \csromangathacontbat{วตฺถุํ วิปตฺติํ อาปตฺติํ,}{นิทานํ อาการอโกวิโท.} \\ \csromangathacontbat{ปุพฺพาปรํ น ชานาติ,}{กตากตํ สเมน จ.} \\ \csromangathacontbat{กมฺมญฺจ อธิกรณญฺจ,}{สมเถ จาปิ อโกวิโท.} \\ \csromangathacontbat{รตฺโต ทุฏฺโฐ จ มูโฬฺห จ,}{ภยา โมหา จ คจฺฉติ.} \\ \csromangathacontbat{น จ สญฺญตฺติกุสโล,}{นิชฺฌตฺติยา จ อโกวิโท.} \\ \csromangathacontbat{ลทฺธปกฺโข อหิริโก,}{กณฺหกมฺโม อนาทโร.} \\ \csromangathacontbat{ส เว\csromansharedfootnote{3b180d75344e4b4b}{สเจ (ก)} ตาทิสโก ภิกฺขุ,}{‘อปฺปฏิกฺโข’ติ วุจฺจติ.} \\ \csromangathacontbat{วตฺถุํ วิปตฺติํ อาปตฺติํ,}{นิทานํ อาการโกวิโท.} \\ \csromangathacontbat{ปุพฺพาปรญฺจ ชานาติ,}{กตากตํ สเมน จ.} \\ \csromangathacontbat{กมฺมญฺจ อธิกรณญฺจ,}{สมเถ จาปิ โกวิโท.} \\ \csromangathacontbat{อรตฺโต อทุฏฺโฐ อมูโฬฺห,}{ภยา โมหา น คจฺฉติ.} \\ \csromangathacontbat{สญฺญตฺติยา จ กุสโล,}{นิชฺฌตฺติยา จ โกวิโท.} \\ \csromangathacontbat{ลทฺธปกฺโข หิริมโน,}{สุกฺกกมฺโม สคารโว.} \\ \csromangathacontbat{ส เว ตาทิสโก ภิกฺขุ,}{‘สปฺปฏิกฺโข’ติ วุจฺจตีติ.}}}{จูฬสงฺคามํ นิฏฺฐิตํ.}
\csromancenter{ตสฺสุทฺทานํ.}
\setlength{\csromangathapagewidth}{0pt}
\setlength{\csromangathawakwidth}{0pt}
\csromangathameasuregroup{นีจจิตฺเตน ปุจฺเฉยฺย, \\ สุตฺตํ สํสนฺทนตฺถาย, \\ อุทฺทานํ จูฬสงฺคาเม,}{\csromangathabat{นีจจิตฺเตน ปุจฺเฉยฺย,}{ครุ สํเฆ น ปุคฺคเล.} \\ \csromangathabat{สุตฺตํ สํสนฺทนตฺถาย,}{วินยานุคฺคเหน จ.} \\ \csromangathabat{อุทฺทานํ จูฬสงฺคาเม,}{เอกุทฺเทโส\textsuperscript{1} อิทํ กตนฺติ.}}
\csromangathasetleft{นีจจิตฺเตน ปุจฺเฉยฺย, \\ สุตฺตํ สํสนฺทนตฺถาย, \\ อุทฺทานํ จูฬสงฺคาเม,}
\csromangathagroup{\csromangathastanza{\csromangathabat{นีจจิตฺเตน ปุจฺเฉยฺย,}{ครุ สํเฆ น ปุคฺคเล.} \\ \csromangathabat{สุตฺตํ สํสนฺทน\-ตฺถาย,}{วินยานุคฺคเหน จ.} \\ \csromangathabat{อุทฺทานํ จูฬสงฺคาเม,}{เอกุทฺเทโส\csromansharedfootnote{dbed1ef3f92fec7d}{เอกุทฺเทสํ (สี, สฺยา)} อิทํ กตนฺติ.}}}

\csromanmatikaanchor{mtk.204}
\csromantocmarknopage{section}{มหาสงฺคาม}{mtk.204}
\bookmark[dest={mtk.204},level=1]{มหาสงฺคาม}
\csromanheader{1}{\textbf{มหาสงฺคาม}}
\csromanheader{2}{1. \textbf{โวหรนฺเตน ชานิตพฺพาทิ}}
\csromanmatikaanchor{mtk.205}
\csromantocmarkref{subsection}{1. โวหรนฺเตนชานิตพฺพาทิ}{mtk.205}
\bookmark[dest={mtk.205},level=2]{1. โวหรนฺเตนชานิตพฺพาทิ}
\proseitem{368}{\csromanfolio{290}สงฺคามา\-วจเรน ภิกฺขุนา สํเฆ โวหรนฺเตน วตฺถุ ชานิตพฺพํ. วิปตฺติ ชานิตพฺพา, อาปตฺติ ชานิตพฺพา, นิทานํ ชานิตพฺพํ, อากาโร ชานิตพฺโพ, ปุพฺพาปรํ ชานิตพฺพํ, กตากตํ ชานิตพฺพํ, กมฺมํ ชานิตพฺพํ, อธิกรณํ ชานิตพฺพํ, สมโถ ชานิตพฺโพ, น ฉนฺทาคติ คนฺตพฺพา, น โทสาคติ คนฺตพฺพา, น โมหาคติ คนฺตพฺพา, น ภยาคติ คนฺตพฺพา, สญฺญาปนีเย ฐาเน สญฺญาเปตพฺพํ, นิชฺฌาปนีเย ฐาเน นิชฺฌาเปตพฺพํ, เปกฺขนีเย ฐาเน เปกฺขิตพฺพํ, ปสาทนีเย ฐาเน ปสาเทตพฺพํ, “ลทฺธปกฺโขมฺหี”ติ ปรปกฺโข นาวชานิตพฺโพ, “พหุสฺสุโตมฺหี”ติ อปฺปสฺสุโต นาวชานิตพฺโพ, “เถรตโรมฺหี”ติ นวกตโร นาวชานิตพฺโพ, อสมฺปตฺตํ นพฺยาหาตพฺพํ, สมฺปตฺตํ ธมฺมโต วินยโต น ปริหาเปตพฺพํ, เยน ธมฺเมน เยน วินเยน เยน สตฺถุสาสเนน ตํ อธิกรณํ วูปสมฺมติ, ตถา ตํ อธิกรณํ วูปสเมตพฺพํ.}

\proseitem{369}{วตฺถุ ชานิตพฺพนฺติ อฏฺฐ\-ปาราชิกานํ\csromansharedfootnote{39b55fab4efcdac2}{อฏฺฐนฺนํ ปาราชิกานํ (สี, ก)} วตฺถุ ชานิตพฺพํ, เตวีสสํฆาทิเสสานํ วตฺถุ ชานิตพฺพํ, เทฺวอนิยตานํ วตฺถุ ชานิตพฺพํ, เทฺว\-จตฺตารีส\-นิสฺสคฺคิยานํ วตฺถุ ชานิตพฺพํ, อฏฺฐาสีติสต\-ปาจิตฺติยานํ วตฺถุ ชานิตพฺพํ, ทฺวาทส\-ปาฏิเทสนียานํ วตฺถุ ชานิตพฺพํ, ทุกฺกฏานํ วตฺถุ ชานิตพฺพํ, ทุพฺภาสิตานํ วตฺถุ ชานิตพฺพํ.}

\proseitem{370}{\textbf{วิปตฺติ ชานิตพฺพา}ติ สีลวิปตฺติ ชานิตพฺพา, อาจารวิปตฺติ ชานิตพฺพา, ทิฏฺฐิวิปตฺติ ชานิตพฺพา, อาชีววิปตฺติ ชานิตพฺพา.}

\proseitem{371}{อาปตฺติ ชานิตพฺพาติ ปาราชิกาปตฺติ ชานิตพฺพา, สํฆาทิเสสาปตฺติ ชานิตพฺพา, ถุลฺลจฺจยา\-ปตฺติ ชานิตพฺพา, ปาจิตฺติยาปตฺติ ชานิตพฺพา, ปาฏิเทสนียา\-ปตฺติ ชานิตพฺพา, ทุกฺกฏาปตฺติ ชานิตพฺพา, ทุพฺภาสิตา\-ปตฺติ ชานิตพฺพา.}

\proseitem{372}{\csromanfolio{291}นิทานํ ชานิตพฺพนฺติ อฏฺฐ\-ปาราชิกานํ นิทานํ ชานิตพฺพํ, เตวีสสํฆาทิเสสานํ นิทานํ ชานิตพฺพํ, เทฺวอนิยตานํ นิทานํ ชานิตพฺพํ, เทฺว\-จตฺตารีส\-นิสฺสคฺคิยานํ นิทานํ ชานิตพฺพํ, อฏฺฐาสีติสต\-ปาจิตฺติยานํ นิทานํ ชานิตพฺพํ, ทฺวาทส\-ปาฏิเทสนียานํ นิทานํ ชานิตพฺพํ, ทุกฺกฏานํ นิทานํ ชานิตพฺพํ, ทุพฺภาสิตานํ นิทานํ ชานิตพฺพํ.}

\proseitem{373}{\textbf{อากาโร ชานิตพฺโพ}ติ สํโฆ อาการโต ชานิตพฺโพ, คโณ อาการโต ชานิตพฺโพ, ปุคฺคโล อาการโต ชานิตพฺโพ, โจทโก อาการโต ชานิตพฺโพ, จุทิตโก อาการโต ชานิตพฺโพ. \textbf{สํโฆ อาการโต ชานิตพฺโพ}ติ “ปฏิพโล นุ โข อยํ สํโฆ อิมํ อธิกรณํ วูปสเมตุํ ธมฺเมน วินเยน สตฺถุสาสเนน, อุทาหุ โน”ติเอวํ สํโฆ อาการโต ชานิตพฺโพ. คโณ อาการโต ชานิตพฺโพติ “ปฏิพโล นุ โข อยํ คโณ อิมํ อธิกรณํ วูปสเมตุํ ธมฺเมน วินเยน สตฺถุสาสเนน, อุทาหุ โน”ติเอวํ คโณ อาการโต ชานิตพฺโพ. ปุคฺคโล อาการโต ชานิตพฺโพติ “ปฏิพโล นุ โข อยํ ปุคฺคโล อิมํ อธิกรณํ วูปเสเมตุํ ธมฺเมน วินเยน สตฺถุสาสเนน, อุทาหุ โน”ติเอวํ ปุคฺคโล อาการโต ชานิตพฺโพ. โจทโก อาการโต ชานิตพฺโพติ “กจฺจิ นุ โข อยมายสฺมา ปญฺจสุ ธมฺเมสุ ปติฏฺฐาย ปรํ โจเทติ, อุกาหุ โน”ติเอวํ โจทโก อาการโต ชานิตพฺโพ. จุทิตโก อาการโต ชานิตพฺโพติ “กจฺจิ นุ โข อยมายสฺมา ทฺวีสุ ธมฺเมสุ ปติฏฺฐิโต สจฺเจ จ อกุปฺเป จ, อุทาหุ โน”ติเอวํ จุทิตโก อาการโต ชานิตพฺโพ.}

\proseitem{374}{\textbf{ปุพฺพาปรํ ชานิตพฺพนฺ}ติ “กจฺจิ นุ โข อยมายสฺมา วตฺถุโต วา วตฺถุํ สงฺกมติ, วิปตฺติโต วา วิปตฺติํ สงฺกมติ, อาปตฺติโต วา อาปตฺติํ สงฺกมติ, อวชานิตฺวา วา ปฏิชานาติ, ปฏิชานิตฺวา วา อวชานาติ, อญฺเญน วา อญฺญํ ปฏิจรติ, อุทาหุ โน”ติเอวํ ปุพฺพาปรํ ชานิตพฺพํ.}

\csromanmatikaanchor{mtk.206}
\csromantocmarkref{subsection}{2. อคติอคนฺตพฺพ}{mtk.206}
\bookmark[dest={mtk.206},level=2]{2. อคติอคนฺตพฺพ}
\proseitem{375}{กตากตํ ชานิตพฺพนฺติ เมถุนธมฺโม ชานิตพฺโพ, เมถุน\-ธมฺมสฺส อนุโลมํ ชานิตพฺพํ, เมถุน\-ธมฺมสฺส ปุพฺพภาโค ชานิตพฺโพ. เมถุนธมฺโม ชานิตพฺโพติ ทฺวยํ ทฺวยสมาปตฺติ ชานิตพฺพา. เมถุน\-ธมฺมสฺส \csromanfolio{292}อนุโลมํ ชานิตพฺพนฺติ ภิกฺขุ อตฺตโน มุเขน ปรสฺส องฺคชาตํ คณฺหาติ. เมถุน\-ธมฺมสฺส ปุพฺพภาโค ชานิตพฺโพติ วณฺณาวณฺโณ\csromansharedfootnote{1680b08d4e64c989}{วณฺโณ อวณฺโณ (สฺยา)}. กายสํสคฺโค, ทุฏฺฐุลฺลวาจา, อตฺต\-กาม\-ปาริจริยา, วจนม\-นุปฺปทานํ\csromansharedfootnote{4a0dd830cc3012bd}{ธนมนุปฺปทานํ (ก)}.}

\proseitem{376}{กมฺมํ ชานิตพฺพนฺติ โสฬสกมฺมานิ ชานิตพฺพานิ, จตฺตาริ อปโลกน\-กมฺมานิ ชานิตพฺพานิ, จตฺตาริ ญตฺติกมฺมานิ ชานิตพฺพานิ, จตฺตาริ ญตฺติ\-ทุติย\-กมฺมานิ ชานิตพฺพานิ, จตฺตาริ ญตฺติ\-จตุตฺถ\-กมฺมานิ ชานิตพฺพานิ.}

\proseitem{377}{อธิกรณํ ชานิตพฺพนฺติ จตฺตาริ อธิกรณานิ ชานิตพฺพานิ, วิวาทา\-ธิกรณํ ชานิตพฺพํ, อนุวาทา\-ธิกรณํ ชานิตพฺพํ, อาปตฺตา\-ธิกรณํ ชานิตพฺพํ, กิจฺจา\-ธิกรณํ ชานิตพฺพํ.}

\proseitem{378}{สมโถ ชานิตพฺโพติ สตฺต สมถา ชานิตพฺพา, สมฺมุขาวินโย ชานิตพฺโพ, สติวินโย ชานิตพฺโพ, อมูฬฺหวินโย ชานิตพฺโพ, ปฏิญฺญาต\-กรณํ ชานิตพฺพํ, เยภุยฺยสิกา ชานิตพฺพา, ตสฺสปาปิยสิกา ชานิตพฺพา, ติณวตฺถารโก ชานิตพฺโพ.}

\prose{2. อคติอคนฺตพฺพ}

\proseitem{379}{น ฉนฺทาคติ คนฺตพฺพาติ ฉนฺทาคติํ คจฺฉนฺโต กถํ ฉนฺทาคติํ คจฺฉติ, อิเธกจฺโจ “อยํ เม อุปชฺฌาโย วา อาจริโย วา สทฺธิวิหาริโก วา อนฺเตวาสิโก วา สมานุ\-ปชฺฌายโก วา สมานาจริยโก วา สนฺทิฏฺโฐ วา สมฺภตฺโต วา ญาติสาโลหิโต วา”ติ ตสฺสานุกมฺปาย ตสฺสา\-นุรกฺขาย อธมฺมํ “ธมฺโม”ติ ทีเปติ, ธมฺมํ “อธมฺโม”ติ ทีเปติ, อวินยํ “วินโย”ติ ทีเปติ, วินยํ “อวินโย”ติ ทีเปติ, อภาสิตํ อลปิตํ ตถาคเตน “ภาสิตํ ลปิตํ ตถาคเตนา”ติ ทีเปติ, ภาสิตํ ลปิตํ ตถาคเตน “อภาสิตํ อลปิตํ ตถาคเตนา”ติ ทีเปติ, อนาจิณฺณํ ตถาคเตน “อาจิณฺณํ ตถาคเตนา”ติ ทีเปติ, อาจิณฺณํ ตถาคเตน “อนาจิณฺณํ ตถาคเตนา”ติ ทีเปติ, อปญฺญตฺตํ ตถาคเตน “ปญฺญตฺตํ ตถาคเตนา”ติ ทีเปติ, ปญฺญตฺตํ ตถาคเตน “อปญฺญตฺตํ ตถาคเตนา”ติ ทีเปติ, \csromanfolio{293}อนาปตฺติํ “อาปตี”ติ ทีเปติ, อาปตฺติํ “อนาปตฺตี”ติ ทีเปติ, ลหุกํ อาปตฺติํ “ครุกา อาปตฺตี”ติ ทีเปติ, ครุกํ อาปตฺติํ “ลหุกา อาปตฺตี”ติ ทีเปติ, สาวเสสํ อาปตฺติํ “อนวเสสา อาปตฺตี”ติ ทีเปติ, อนวเสสํ อาปตฺติํ “สาวเสสา อาปตฺตี”ติ ทีเปติ, ทุฏฺฐุลฺลํ อาปตฺติํ “อทุฏฺฐุลฺลา อาปตฺตี”ติ ทีเปติ, อทุฏฺฐุลฺลํ อาปตฺติํ “ทุฏฺฐุลฺลา อาปตฺตี”ติ ทีเปติ. อิเมหิ อฏฺฐารสหิ วตฺถูหิ ฉนฺทาคติํ คจฺฉนฺโต พหุชนาหิตาย ปฏิปนฺโน โหติ พหุชนา\-สุขาย พหุโน ชนสฺส อนตฺถาย อหิตาย ทุกฺขาย เทวมนุสฺสานํ. อิเมหิ อฏฺฐารสหิ วตฺถูหิ ฉนฺทาคติํ คจฺฉนฺโต ขตํ อุปหตํ อตฺตานํ ปริหรติ, สาวชฺโช จ โหติ, สานุวชฺโช จ วิญฺญูนํ, พหุญฺจ อปุญฺญํ ปสวติ. ฉนฺทาคติํ คจฺฉนฺโต เอวํ ฉนฺทาคติํ คจฺฉติ.}

\proseitem{380}{น โทสาคติ คนฺตพฺพาติ โทสาคติํ คจฺฉนฺโต กถํ โทสาคติํ คจฺฉติ, อิเธกจฺโจ “อนตฺถํ เม อจรี”ติ อาฆาตํ พนฺธติ, “อนตฺถํ เม จรตี”ติ อาฆาตํ พนฺธติ, “อนตฺถํ เม จริสฺสตี”ติ อาฆาตํ พนฺธติ, “ปิยสฺส เม มนาปสฺส อนตฺถํ อจริ, อนตฺถํ จรติ, อนตฺถํ จริสฺสตี”ติ อาฆาตํ พนฺธติ, “อปฺปิยสฺส เม อมนาปสฺส อตฺถํ อจริ, อตฺถํ จรติ, อตฺถํ จริสฺสตี”ติ อาฆาตํ พนฺธติ. อิเมหิ นวหิ อาฆาตวตฺถูหิ อาฆาโต ปฏิฆาโต กุทฺโธ โกธาภิภูโต อธมฺมํ “ธมฺโม”ติ ทีเปติ, ธมฺมํ “อธมฺโม”ติ ทีเปติ ฯเปฯ ทุฏฺฐุลฺลํ อาปตฺติํ “อทุฏฺฐุลฺลา อาปตฺตี”ติ ทิเปติ, อทุฏฺฐุลฺลํ อาปตฺติํ “ทุฏฺฐุลฺลา อาปตฺตี”ติ ทีเปติ. อิเมหิ อฏฺฐารสหิ วตฺถูหิ โทสาคติํ คจฺฉนฺโต พหุชนา หิตาย ปฏิปนฺโน โหติ พหุชนา\-สุขาย พหุโน ชนสฺส อนตฺถาย อหิตาย ทุกฺขาย เทวมนุสฺสานํ. อิเมหิ อฏฺฐารสหิ วตฺถูหิ โทสาคติํ คจฺฉนฺโต ขตํ อุปหตํ อตฺตานํ ปริหรติ, สาวชฺโช จ โหติ, สานุวชฺโช จ วิญฺญูนํ, พหุญฺจ อปุญฺญํ ปสวติ. โทสาคติํ คจฺฉนฺโต เอวํ โทสาคติํ คจฺฉติ.}

\proseitem{381}{\textbf{น โมหาคติ คนฺตพฺพา}ติ โมหาคติํ คจฺฉนฺโต กถํ โมหาคติํ คจฺฉติ, รตฺโต ราควเสน คจฺฉติ, ทุฏฺโฐ โทสวเสน คจฺฉติ, มูโฬฺห โมหวเสน คจฺฉติ, ปรามฏฺโฐ ทิฏฺฐิวเสน คจฺฉติ, มูโฬฺห สํมูโฬฺห โมหาภิภูโต อธมฺมํ “ธมฺโม”ติ ทีเปติ, ธมฺมํ \csromanfolio{294}“อธมฺโม”ติ ทีเปติ ฯเปฯ ทุฏฺฐุลฺลํ อาปตฺติํ “อทุฏฺฐุลฺลา อาปตฺตี”ติ ทีเปติ, อทุฏฺฐุลฺลํ อาปตฺติํ “ทุฏฺฐุลฺลา อาปตฺตี”ติ ทีเปติ. อิเมหิ อฏฺฐารสหิ วตฺถูหิ โมหาคติํ คจฺฉนฺโต พหุชนาหิตาย ปฏิปนฺโน โหติ พหุชนา\-สุขาย พหุโน ชนสฺส อนตฺถาย อหิตาย ทุกฺขาย เทวมนุสฺสานํ. อิเมหิ อฏฺฐารสหิ วตฺถูหิ โมหาคติํ คจฺฉนฺโต ขตํ อุปหตํ อตฺตานํ ปริหรติ, สาวชฺโช จ โหติ, สานุวชฺโช จ วิญฺญูนํ, พหุญฺจ อปุญฺญํ ปสวติ. โมหาคติํ คจฺฉนฺโต เอวํ โมหาคติํ คจฺฉติ.}

\proseitem{382}{น ภยาคติ คนฺตพฺพาติ ภยาคติํ คจฺฉนฺโต กถํ ภยาคติํ คจฺฉติ, อิเธกจฺโจ “อยํ วิสมนิสฺสิโต วา คหนนิสฺสิโต วา พลวนิสฺสิโต วา กกฺขโฬ ผรุโส ชีวิตนฺตรายํ วา พฺรหฺมจริย\-นฺตรายํ วา กริสฺสตี”ติ ตสฺส ภยา ภีโต อธมฺมํ “ธมฺโม”ติ ทีเปติ, ธมฺมํ “อธมฺโม”ติ ทีเปติ, อวินยํ “วินโย”ติ ทีเปติ, วินยํ “อวินโย”ติ ทีเปติ, อภาสิตํ อลปิตํ ตถาคเตน “ภาสิตํ ลปิตํ ตถาคเตนา”ติ ทีเปติ, ภาสิตํ ลปิตํ ตถาคเตน “อภาสิตํ อลปิตํ ตถาคเตนา”ติ ทีเปติ, อนาจิณฺณํ ตถาคเตน “อาจิณฺณํ ตถาคเตนา”ติ ทีเปติ, อาจิณฺณํ ตถาคเตน “อนาจิณฺณํ ตถาคเตนา”ติ ทีเปติ, อปญฺญตฺตํ ตถาคเตน “ปญฺญตฺตํ ตถาคเตนา”ติ ทีเปติ, ปญฺญตฺตํ ตถาคเตน “อปญฺญตฺตํ ตถาคเตนา”ติ ทีเปติ, อนาปตฺติํ “อาปตฺตี”ติ ทีเปติ, อาปตฺติํ “อนาปตฺตี”ติ ทีเปติ, ลหุกํ อาปตฺติํ “ครุกา อาปตฺตี”ติ ทีเปติ, ครุกํ อาปตฺติํ “ลหุกา อาปตฺตี”ติ ทีเปติ, สาวเสสํ อาปตฺติํ “อนวเสสา อาปตฺตี”ติ ทีเปติ, อนวเสสํ อาปตฺติํ “สาวเสสา อาปตฺตี”ติ ทีเปติ, ทุฏฺฐุลฺลํ อาปตฺติํ “อทุฏฺฐุลฺลา อาปตฺตี”ติ ทีเปติ, อทุฏฺฐุลฺลํ อาปตฺติํ “ทุฏฺฐุลฺลา อาปตฺตี”ติ ทีเปติ. อิเมหิ อฏฺฐารสหิ วตฺถูหิ ภยาคติํ คจฺฉนฺโต พหุชนาหิตาย ปฏิปนฺโน โหติ พหุชนา\-สุขาย พหุโน ชนสฺส อนตฺถาย อหิตาย ทุกฺขาย เทวมนุสฺสานํ. อิเมหิ อฏฺฐารสหิ วตฺถูหิ ภยาคติํ คจฺฉนฺโต ขตํ อุปหตํ อตฺตานํ ปริหรติ, สาวชฺโช จ โหติ, สานุวชฺโช จ วิญฺญูนํ, พหุญฺจ อปุญฺญํ ปสวติ. ภยาคติํ คจฺฉนฺโต เอวํ ภยาคติํ คจฺฉติ.}

\setlength{\csromangathapagewidth}{0pt}
\setlength{\csromangathawakwidth}{0pt}
\csromangathameasuregroup{ฉนฺทา โทสา ภยา โมหา, \\ นิหียติ ตสฺส ยโส,}{\csromangathabat{ฉนฺทา โทสา ภยา โมหา,}{โย ธมฺมํ อติวตฺตติ.} \\ \csromangathabat{นิหียติ ตสฺส ยโส,}{กาฬปกฺเขว จนฺทิมาติ.}}
\csromangathasetleft{ฉนฺทา โทสา ภยา โมหา, \\ นิหียติ ตสฺส ยโส,}
\csromangathagroup{\csromangathastanza{\csromangathabat{ฉนฺทา โทสา ภยา โมหา,}{โย ธมฺมํ อติวตฺตติ.} \\ \csromangathabat{นิหียติ ตสฺส ยโส,}{กาฬปกฺเขว จนฺทิมาติ.}}}

\csromanmatikaanchor{mtk.207}
\csromantocmarkref{subsection}{3. อคติอคมน}{mtk.207}
\bookmark[dest={mtk.207},level=2]{3. อคติอคมน}
\prose{\csromanfolio{295}3. อคติอคมน}

\proseitem{383}{\textbf{กถํ น ฉนฺทาคติํ คจฺฉติ,} อธมฺมํ “อธมฺโม”ติ ทีเปนฺโต น ฉนฺทาคติํ คจฺฉติ, ธมฺมํ “ธมฺโม”ติ ทีเปนฺโต น ฉนฺทาคติํ คจฺฉติ, อวินยํ “อวินโย”ติ ทีเปนฺโต น ฉนฺทาคติํ คจฺฉติ, วินยํ “วินโย”ติ ทีเปนฺโต น ฉนฺทาคติํ คจฺฉติ, อภาสิตํ อลปิตํ ตถาคเตน “อภาสิตํ อลปิตํ ตถาคเตนา”ติ ทีเปนฺโต น ฉนฺทาคติํ คจฺฉติ, ภาสิตํ ลปิตํ ตถาคเตน “ภาสิตํ ลปิตํ ตถาคเตนา”ติ ทีเปนฺโต น ฉนฺทาคติํ คจฺฉติ, อนาจิณฺณํ ตถาคเตน “อนาจิณฺณํ ตถาคเตนา”ติ ทีเปนฺโต น ฉนฺทาคติํ คจฺฉติ, อาจิณฺณํ ตถาคเตน “อาจิณฺณํ ตถาคเตนา”ติ ทีเปนฺโต น ฉนฺทาคติํ คจฺฉติ, อปญฺญตฺตํ ตถาคเตน “อปญฺญตฺตํ ตถาคเตนา”ติ ทีเปนฺโต น ฉนฺทาคติํ คจฺฉติ, ปญฺญตฺตํ ตถาคเตน “ปญฺญตฺตํ ตถาคเตนา”ติ ทีเปนฺโต น ฉนฺทาคติํ คจฺฉติ, อนาปตฺติํ “อนาปตฺตี”ติ ทีเปนฺโต น ฉนฺทาคติํ คจฺฉติ, อาปตฺติํ “อาปตฺตี”ติ ทีเปนฺโต น ฉนฺทาคติํ คจฺฉติ, ลหุกํ อาปตฺติํ “ลหุกา อาปตฺตี”ติ ทีเปนฺโต น ฉนฺทาคติํ คจฺฉติ, ครุกํ อาปติํ “ครุกา อาปตฺตี”ติ ทีเปนฺโต น ฉนฺทาคติํ คจฺฉติ, สาวเสสํ อาปตฺติํ “สาวเสสา อาปตฺตี”ติ ทีเปนฺโต น ฉนฺทาคติํ คจฺฉติ, อนวเสสํ อาปตฺติํ “อนวเสสา อาปตฺตี”ติ ทีเปนฺโต น ฉนฺทาคติํ คจฺฉติ, ทุฏฺฐุลฺลํ อาปตฺติํ “ทุฏฺฐุลฺลา อาปตฺตี”ติ ทีเปนฺโต น ฉนฺทาคติํ คจฺฉติ, อทุฏฺฐุลฺลํ อาปตฺติํ “อทุฏฺฐุลฺลา อาปตฺตี”ติ ทีเปนฺโต น ฉนฺทาคติํ คจฺฉติ. เอวํ น ฉนฺทาคติํ คจฺฉติ.}

\proseitem{384}{\textbf{กถํ น โทสาคติํ คจฺฉติ,} อธมฺมํ “อธมฺโม”ติ ทีเปนฺโต น โทสาคติํ คจฺฉติ, ธมฺมํ “ธมฺโม”ติ ทีเปนฺโต น โทสาคติํ คจฺฉติ ฯเปฯ ทุฏฺฐุลฺลํ อาปตฺติํ “ทุฏฺฐุลฺลา อาปตฺตี”ติ ทีเปนฺโต น โทสาคติํ คจฺฉติ, อทุฏฺฐุลฺลํ อาปตฺติํ “อทุฏฺฐุลฺลา อาปตฺตี”ติ ทีเปนฺโต น โทสาคติํ คจฺฉติ. เอวํ น โทสาคติํ คจฺฉติ.}

\proseitem{385}{\textbf{กถํ น โมหาคติํ คจฺฉติ,} อธมฺมํ “อธมฺโม”ติ ทีเปนฺโต น โมหาคติํ คจฺฉติ, ธมฺมํ “ธมฺโม”ติ ทีเปนฺโต น โมหาคติํ คจฺฉติ ฯเปฯ. ทุฏฺฐุลฺลํ อาปตฺติํ “ทุฏฺฐุลฺลา อาปตฺตี”ติ ทีเปนฺโต น \csromanfolio{296}โมหาคติํ คจฺฉติ, อทุฏฺฐุลฺลํ อาปตฺติํ “อทุฏฺฐุลฺลา อาปตฺตี”ติ ทีเปนฺโต น โมหาคติํ คจฺฉติ. เอวํ น โมหาคติํ คจฺฉติ.}

\proseitem{386}{\textbf{กถํ น ภยาคติํ คจฺฉติ,} อธมฺมํ “อธมฺโม”ติ ทีเปนฺโต น ภยาคติํ คจฺฉติ, ธมฺมํ “ธมฺโม”ติ ทีเปนฺโต น ภยาคติํ คจฺฉติ, อวินยํ “อวินโย”ติ ทีเปนฺโต น ภยาคติํ คจฺฉติ, วินยํ “วินโย”ติ ทีเปนฺโต น ภยาคติํ คจฺฉติ, อภาสิตํ อลปิตํ ตถาคเตน “อภาสิตํ อลปิตํ ตถาคเตนา”ติ ทีเปนฺโต น ภยาคติํ คจฺฉติ, ภาสิตํ ลปิตํ ตถาคเตน “ภาสิตํ ลปิตํ ตถาคเตนา”ติ ทีเปนฺโต น ภยาคติํ คจฺฉติ, อนาจิณฺณํ ตถาคเตน “อนาจิณฺณํ ตถาคเตนา”ติ ทีเปนฺโต น ภยาคติํ คจฺฉติ, อาจิณฺณํ ตถาคเตน “อาจิณฺณํ ตถาคเตนา”ติ ทีเปนฺโต น ภยาคติํ คจฺฉติ, อปญฺญตฺตํ ตถาคเตน “อปญฺญตฺตํ ตถาคเตนา”ติ ทีเปนฺโตน ภยาคติํ คจฺฉติ, ปญฺญตฺตํ ตถาคเตน “ปญฺญตฺตํ ตถาคเตนา”ติ ทีเปนฺโต น ภยาคติํ คจฺฉติ, อนาปตฺติํ “อนาปตฺตี”ติ ทีเปนฺโต น ภยาคติํ คจฺฉติ, อาปตฺติํ “อาปตฺตี”ติ ทีเปนฺโต น ภยาคติํ คจฺฉติ, ลหุกํ อาปตฺติํ “ลหุกา อาปตฺตี”ติ ทีเปนฺโต น ภยาคติํ คจฺฉติ, ครุกํ อาปตฺติํ “ครุกา อาปตฺตี”ติ ทีเปนฺโต น ภยาคติํ คจฺฉติ, สาวเสสํ อาปตฺติํ “สาวเสสา อาปตฺตี”ติ ทีเปนฺโต น ภยาคติํ คจฺฉติ, อนวเสสํ อาปตฺติํ “อนวเสสา อาปตฺตี”ติ ทีเปนฺโต น ภยาคติํ คจฺฉติ, ทุฏฺฐุลฺลํ อาปตฺติํ “ทุฏฺฐุลฺลา อาปตฺตี”ติ ทีเปนฺโต น ภายาคติํ คจฺฉติ, อทุฏฺฐุลฺลํ อาปตฺติํ “อทุฏฺฐุลฺลา อาปตฺตี”ติ ทีเปนฺโต น ภยาคติํ คจฺฉติ. เอวํ น ภยาคติํ คจฺฉติ.}

\setlength{\csromangathapagewidth}{0pt}
\setlength{\csromangathawakwidth}{0pt}
\csromangathameasuregroup{ฉนฺทา โทสา ภยา โมหา, \\ อาปูรติ ตสฺส ยโส,}{\csromangathabat{ฉนฺทา โทสา ภยา โมหา,}{โย ธมฺมํ นาติวตฺตติ.} \\ \csromangathabat{อาปูรติ ตสฺส ยโส,}{สุกฺกปกฺเขว จนฺทิมาติ.}}
\csromangathasetleft{ฉนฺทา โทสา ภยา โมหา, \\ อาปูรติ ตสฺส ยโส,}
\csromangathagroup{\csromangathastanza{\csromangathabat{ฉนฺทา โทสา ภยา โมหา,}{โย ธมฺมํ นาติวตฺตติ.} \\ \csromangathabat{อาปูรติ ตสฺส ยโส,}{สุกฺกปกฺเขว จนฺทิมาติ.}}}

\csromanmatikaanchor{mtk.208}
\csromantocmarkref{subsection}{4. สญฺญาปนียาทิ}{mtk.208}
\bookmark[dest={mtk.208},level=2]{4. สญฺญาปนียาทิ}
\csromanheader{2}{4. \textbf{สญฺญาปนียาทิ}}
\proseitem{387}{กถํ สญฺญาปนีเย ฐาเน สญฺญาเปติ, อธมฺมํ “อธมฺโม”ติ ทีเปนฺโต สญฺญาปนีเย ฐาเน สญฺญาเปติ, ธมฺมํ “ธมฺโม”ติ ทีเปนฺโต สญฺญาปนีเย ฐาเน สญฺญาเปติ ฯเปฯ ทุฏฺฐุลฺลํ อาปตฺติํ “ทุฏฺฐุลฺลา อาปตฺตี”ติ ทีเปนฺโต สญฺญาปนีเย ฐาเน สญฺญาเปติ, อทุฏฺฐุลฺลํ อาปตฺติํ \csromanfolio{297}“อทุฏฺฐุลฺลา อาปตฺตี”ติ ทีเปนฺโต สญฺญาปนีเย ฐาเน สญฺญาเปติ. เอวํ สญฺญาปนีเย ฐาเน สญฺญาเปติ.}

\proseitem{388}{\textbf{กถํ นิชฺฌาปนีเย ฐาเน นิชฺฌาเปติ,} อธมฺมํ “อธมฺโม”ติ ทีเปนฺโต นิชฺฌาปนีเย ฐาเน นิชฺฌาเปติ, ธมฺมํ “ธมฺโม”ติ ทีเปนฺโต นิชฺฌาปนีเย ฐาเน นิชฺฌาเปติ ฯเปฯ ทุฏฺฐุลฺลํ อาปตฺติํ “ทุฏฺฐุลฺลา อาปตฺตี”ติ ทีเปนฺโต นิชฺฌาปนีเย ฐาเน นิชฺฌาเปติ, อทุฏฺฐุลฺลํ อาปตฺติํ “อทุฏฺฐุลฺลา อาปตฺตี”ติ ทีเปนฺโต นิชฺฌาปนีเย ฐาเน นิชฺฌาเปติ. เอวํ นิชฺฌาปนีเย ฐาเน นิชฺฌาเปติ.}

\proseitem{389}{\textbf{กถํ เปกฺขนีเย ฐาเน เปกฺขติ,} อธมฺมํ “อธมฺโม”ติ ทีเปนฺโต เปกฺขนีเย ฐาเน เปกฺขติ, ธมฺมํ “ธมฺโม”ติ ทีเปนฺโต เปกฺขนีเย ฐาเน เปกฺขติ ฯเปฯ ทุฏฺฐุลฺลํ อาปตฺติํ “ทุฏฺฐุลฺลา อาปตฺตี”ติ ทีเปนฺโต เปกฺขนีเย ฐาเน เปกฺขติ, อทุฏฺฐุลฺลํ อาปตฺติํ “อทุฏฺฐุลฺลา อาปตฺตี”ติ ทีเปนฺโต เปกฺขนีเย ฐาเน เปกฺขติ. เอวํ เปกฺขนีเย ฐาเน เปกฺขติ.}

\proseitem{390}{\textbf{กถํ ปสาทนีเย ฐาเน ปสาเทติ,} อธมฺมํ “อธมฺโม”ติ ทีเปนฺโต ปสาทนีเย ฐาเน ปสาเทติ, ธมฺมํ “ธมฺโม”ติ ทีเปนฺโต ปสาทนีเย ฐาเน ปสาเทติ ฯเปฯ ทุฏฺฐุลฺลํ อาปตฺติํ “ทุฏฺฐุลฺลา อาปตฺตี”ติ ทีเปนฺโต ปสาทนีเย ฐาเน ปสาเทติ, อทุฏฺฐุลฺลํ อาปตฺติํ “อทุฏฺฐุลฺลา อาปตฺตี”ติ ทีเปนฺโต ปสาทนีเย ฐาเน ปสาเทติ. เอวํ ปสาทนีเย ฐาเน ปสาเทติ.}

\csromanmatikaanchor{mtk.209}
\csromantocmarkref{subsection}{5. ปรปกฺขาทิอวชานน}{mtk.209}
\bookmark[dest={mtk.209},level=2]{5. ปรปกฺขาทิอวชานน}
\csromanheader{2}{5. ปรปกฺขา\-ทิ\-อวชานน}
\proseitem{391}{กถํ “ลทฺธปกฺโขมฺหี”ติ ปรปกฺขํ อวชานาติ, อิเธกจฺโจ ลทฺธปกฺโข โหติ ลทฺธปริวาโร ปกฺขวา ญาติมา “อยํ อลทฺธปกฺโข อลทฺธ\-ปริวาโร น ปกฺขวา น ญาติมา”ติ ตสฺส อวชานนฺโต อธมฺมํ “ธมฺโม”ติ ทีเปติ, ธมฺมํ “อธมฺโม”ติ ทีเปติ ฯเปฯ ทุฏฺฐุลฺลํ อาปตฺติํ “อทุฏฺฐุลฺลา อาปตฺตี”ติ ทีเปติ, อทุฏฺฐุลฺลํ อาปตฺติํ “ทุฏฺฐุลฺลา อาปตฺตี”ติ ทีเปติ. เอวํ “ลทฺธปกฺโขมฺหี”ติ ปรปกฺขํ อวชานาติ.}

\proseitem{392}{\textbf{กถํ “พหุสฺสุโตมฺหี”ติ อปฺปสฺสุตํ อวชานาติ,} อิเธกจฺโจ พหุสฺสุโต โหติ สุตธโร สุตสนฺนิจโย, “อยํ อปฺปสฺสุโต \csromanfolio{298}อปฺปาคโม อปฺปธโร”ติ ตสฺส อวชานนฺโต อธมฺมํ “ธมฺโม”ติ ทีเปติ, ธมฺมํ “อธมฺโม”ติ ทีเปติ ฯเปฯ ทุฏฺฐุลฺลํ อาปตฺติํ “อทุฏฺฐุลฺลา อาปตฺตี”ติ ทีเปติ, อทุฏฺฐุลฺลํ อาปตฺติํ “ทุฏฺฐุลฺลา อาปตฺตี”ติ ทีเปติ. เอวํ “พหุสฺสุโตมฺหี”ติ อปฺปสฺสุตํ อวชานาติ.}

\proseitem{393}{\textbf{กถํ “เถรตโรมฺหี”ติ นวกตรํ อวชานาติ,} อิเธกจฺโจ เถโร โหติ รตฺตญฺญู จิรปพฺพชิโต “อยํ นวโก อปฺปญฺญาโต อปฺปกตญฺญู, อิมสฺส วจนํ อกตํ ภวิสฺสตี”ติ ตสฺส อวชานนฺโต อธมฺมํ “ธมฺโม”ติ ทีเปติ, ธมฺมํ “อธมฺโม”ติ ทีเปติ ฯเปฯ ทุฏฺฐุลฺลํ อาปตฺติํ “อทุฏฺฐุลฺลา อาปตฺตี”ติ ทีเปติ, อทุฏฺฐุลฺลํ อาปตฺติํ “ทุฏฺฐุลฺลา อาปตฺตี”ติ ทีเปติ. เอวํ “เถรตโรมฺหี”ติ นวกตรํ อวชานาติ.}

\proseitem{394}{อสมฺปตฺตํ น พฺยาหริตพฺพนฺติ อโนติณฺณํ ภารํ\csromansharedfootnote{cf2bc75d1904cc60}{ภาสํ (สฺยา)} น โอตาเรตพฺพํ. สมฺปตฺตํ ธมฺมโต วินยโต น ปริหาเป\-ตพฺพนฺติ ยํอตฺถาย สํโฆ สนฺนิปติโต โหติ, ตํ อตฺถํ ธมฺมโต วินยโต น ปริหาเปตพฺพํ.}

\proseitem{395}{เยน ธมฺเมนาติ ภูเตน วตฺถุนา. เยน วินเยนาติ โจเทตฺวา สาเรตฺวา. เยน สตฺถุ\-สาสเน\-นาติ ญตฺติสมฺปทาย อนุสฺสาวน\-สมฺปทาย, เยน ธมฺเมน เยน วินเยน เยน สตฺถุสาสเนน ตํ อธิกรณํ วูปสมฺมติ, ตถา ตํ อธิกรณํ วูปสเม\-ตพฺพนฺติ.}

\csromanmatikaanchor{mtk.210}
\csromantocmarkref{subsection}{6. อนุวิชฺชกสฺสอนุโยค}{mtk.210}
\bookmark[dest={mtk.210},level=2]{6. อนุวิชฺชกสฺสอนุโยค}
\csromanheader{2}{6. \textbf{อนุวิชฺชกสฺส อนุโยค}}
\proseitem{396}{อนุวิชฺชเกน โจทโก ปุจฺฉิตพฺโพ “ยํ โข ตฺวํ อาวุโส อิมสฺส ภิกฺขุโน ปวารณํ ฐเปสิ, กิมฺหิ นํ ฐเปสิ, สีลวิปตฺติยา วา ฐเปสิ, อาจารวิปตฺติยา วา ฐเปสิ, ทิฏฺฐิ\-วิปตฺติยา วา ฐเปสี”ติ, โส เจ เอวํ วเทยฺย “สีลวิปตฺติยา วา ฐเปมิ, อาจารวิปตฺติยา วา ฐเปมิ, ทิฏฺฐิ\-วิปตฺติยา วา ฐเปมี”ติ. โส เอวมสฺส วจนีโย “ชานาติ ปนายสฺมา สีลวิปตฺติํ, ชานาติ อาจารวิปตฺติํ, ชานาติ ทิฏฺฐิวิปตฺตินฺ”ติ. โส เจ เอวํ วเทยฺย “ชานามิ โข อหํ อาวุโส สีลวิปตฺติํ, ชานามิ อาจารวิปตฺติํ, ชานามิ ทิฏฺฐิวิปตฺตินฺ”ติ. โส เอวมสฺส วจนีโย \csromanfolio{299}“กตมา ปนาวุโส สีลวิปตฺติ, กตมา อาจารวิปตฺติ, กตมา ทิฏฺฐิวิปตฺตี”ติ. โส เจ เอวํ วเทยฺย “จตฺตาริ ปาราชิกานิ เตรส สํฆาทิเสสา, อยํ สีลวิปตฺติ. ถุลฺลจฺจยํ ปาจิตฺติยํ ปฏิเทสนียํ ทุกฺกฏํ ทุพฺภาสิตํ, อยํ อาจารวิปตฺติ. มิจฺฉาทิฏฺฐิ อนฺตคฺคาหิกา\-ทิฏฺฐิ, อยํ ทิฏฺฐิวิปตฺตี”ติ. โส เอวมสฺส วจนีโย “ยํ โข ตฺวํ อาวุโส อิมสฺส ภิกฺขุโน ปวารณํ ฐเปสิ, ทิฏฺเฐน วา ฐเปสิ, สุเตน วา ฐเปสิ, ปริสงฺกาย วา ฐเปสี”ติ. โส เจ เอวํ วเทยฺย “ทิฏฺเฐน วา ฐเปมิ, สุเตน วา ฐเปมิ, ปริสงฺกาย วา ฐเปมี”ติ. โส เอวมสฺส วจนีโย “ยํ โข ตฺวํ อาวุโส อิมสฺส ภิกฺขุโน ทิฏฺเฐน ปวารณํ ฐเปสิ, กิํ เต ทิฏฺฐํ, กินฺติ เต ทิฏฺฐํ, กทา เต ทิฏฺฐํ, กตฺถ เต ทิฏฺฐํ, ปาราชิกํ อชฺฌาปชฺชนฺโต ทิฏฺโฐ, สํฆาทิเสสํ อชฺฌาปชฺชนฺโต ทิฏฺโฐ, ถุลฺลจฺจยํ. ปาจิตฺติยํ. ปาฏิเทสนียํ. ทุกฺกฏํ. ทุพฺภาสิตํ อชฺฌาปชฺชนฺโต ทิฏฺโฐ, กตฺถ จ ตฺวํ อโหสิ, กตฺถ จายํ ภิกฺขุ อโหสิ, กิญฺจ ตฺวํ กโรสิ, กิญฺจายํ ภิกฺขุ กโรตี”ติ. โส เจ เอวํ วเทยฺย “น โข อหํ อาวุโส อิมสฺส ภิกฺขุโน ทิฏฺเฐน ปวารณํ ฐเปมิ, อปิ จ สุเตน ปวารณํ ฐเปมี”ติ. โส เอวมสฺส วจนีโย “ยํ โข ตฺวํ อาวุโส อิมสฺส ภิกฺขุโน สุเตน ปวารณํ ฐเปสิ, กิํ เต สุตํ, กินฺติ เต สุตํ, กทา เต สุตํ, กตฺถ เต สุตํ, ‘ปาราชิกํ อชฺฌาปนฺโน’ติ สุตํ, ‘สํฆาทิเสสํ อชฺฌาปนฺโน’ติ สุตํ, ถุลฺลจฺจยํ. ปาจิตฺติยํ. ปาฏิเทสนียํ. ทุกฺกฏํ. ‘ทุพฺภาสิตํ อชฺฌาปนฺโน’ติ สุตํ, ภิกฺขุสฺส สุตํ, ภิกฺขุนิยา สุตํ, สิกฺขมานาย สุตํ, สามเณรสฺส สุตํ, สามเณริยา สุตํ, อุปาสกสฺส สุตํ, อุปาสิกาย สุตํ, ราชูนํ สุตํ, ราชามหามตฺตานํ สุตํ, ติตฺถิยานํ สุตํ, ติตฺถิย\-สาวกานํ สุตนฺ”ติ. โส เจ เอวํ วเทยฺย “น โข อหํ อาวุโส อิมสฺส ภิกฺขุโน สุเตน ปวารณํ ฐเปมิ, อปิ จ ปริสงฺกาย ปวารณํ ฐเปมี”ติ, โส เอวมสฺส วจนีโย “ยํ โข ตฺวํ อาวุโส อิมสฺส ภิกฺขุโน ปริสงฺกาย ปวารณํ ฐเปสิ, กิํ ปริสงฺกสิ, กินฺติ ปริสงฺกสิ, กทา ปริสงฺกสิ, กตฺถ ปริสงฺกสิ, ‘ปาราชิกํ อชฺฌาปนฺโน’ติ ปริสงฺกสิ, ‘สํฆาทิเสสํ อชฺฌาปนฺโน’ติ ปริสงฺกสิ, ถุลฺลจฺจยํ. ปาจิตฺติยํ. ปาฏิเทสนียํ. ทุกฺกฏํ. ‘ทุพฺภาสิตํ อชฺฌาปนฺโน’ติ ปริสงฺกสิ, ภิกฺขุสฺส สุตฺวา ปริสงฺกสิ, ภิกฺขุนิยา สุตฺวา ปริสงฺกสิ, สิกฺขมานาย สุตฺวา ปริสงฺกสิ, สามเณรสฺส สุตฺวา ปริสงฺกสิ, สามเณริยา สุตฺวา ปริสงฺกสิ, อุปาสกสฺส สุตฺวา ปริสงฺกสิ, อุปาสิกาย \csromanfolio{300}สุตฺวา ปริสงฺกสิ, ราชูนํ สุตฺวา ปริสงฺกสิ, ราช\-มหามตฺตานํ สุตฺวา ปริสงฺกสิ, ติตฺถิยานํ สุตฺวา ปริสงฺกสิ, ติตฺถิย\-สาวกานํ สุตฺวา ปริสงฺกสี”ติ.}

\setlength{\csromangathapagewidth}{0pt}
\setlength{\csromangathawakwidth}{0pt}
\csromangathameasuregroup{ทิฏฺฐํ ทิฏฺเฐน สเมติ, \\ ทิฏฺฐํ ปฏิจฺจ น อุเปติ, \\ โส ปุคฺคโล ปฏิญฺญาย, \\ สุตํ สุเตน สเมติ, \\ สุตํ ปฏิจฺจ น อุเปติ, \\ โส ปุคฺคโล ปฏิญฺญาย, \\ มุตํ มุเตน สเมติ, \\ มุตํ ปฏิจฺจ น อุเปติ, \\ โส ปุคฺคโล ปฏิญฺญาย,}{\csromangathabat{ทิฏฺฐํ ทิฏฺเฐน สเมติ,}{ทิฏฺเฐน สํสนฺทเต ทิฏฺฐํ.} \\ \csromangathabat{ทิฏฺฐํ ปฏิจฺจ น อุเปติ,}{อสุทฺธปริสงฺกิโต.} \\ \csromangathabat{โส ปุคฺคโล ปฏิญฺญาย,}{กาตพฺพา เตน ปวารณา.} \\ \csromangathabat{สุตํ สุเตน สเมติ,}{สุเตน สํสนฺทเต สุตํ.} \\ \csromangathabat{สุตํ ปฏิจฺจ น อุเปติ,}{อสุทฺธปริสงฺกิโต.} \\ \csromangathabat{โส ปุคฺคโล ปฏิญฺญาย,}{กาตพฺพา เตน ปวารณา.} \\ \csromangathabat{มุตํ มุเตน สเมติ,}{มุเตน สํสนฺทเต มุตํ.} \\ \csromangathabat{มุตํ ปฏิจฺจ น อุเปติ,}{อสุทฺธปริสงฺกิโต.} \\ \csromangathabat{โส ปุคฺคโล ปฏิญฺญาย,}{กาตพฺพา เตน ปวารณาติ.}}
\csromangathasetleft{ทิฏฺฐํ ทิฏฺเฐน สเมติ, \\ ทิฏฺฐํ ปฏิจฺจ น อุเปติ, \\ โส ปุคฺคโล ปฏิญฺญาย, \\ สุตํ สุเตน สเมติ, \\ สุตํ ปฏิจฺจ น อุเปติ, \\ โส ปุคฺคโล ปฏิญฺญาย, \\ มุตํ มุเตน สเมติ, \\ มุตํ ปฏิจฺจ น อุเปติ, \\ โส ปุคฺคโล ปฏิญฺญาย,}
\csromangathagroup{\csromangathastanza{\csromangathaitembat{397}{ทิฏฺฐํ ทิฏฺเฐน สเมติ,}{ทิฏฺเฐน สํสนฺทเต ทิฏฺฐํ.} \\ \csromangathacontbat{ทิฏฺฐํ ปฏิจฺจ น อุเปติ,}{อสุทฺธปริสงฺกิโต.} \\ \csromangathacontbat{โส ปุคฺคโล ปฏิญฺญาย,}{กาตพฺพา เตน ปวารณา.} \\ \csromangathacontbat{สุตํ สุเตน สเมติ,}{สุเตน สํสนฺทเต สุตํ.} \\ \csromangathacontbat{สุตํ ปฏิจฺจ น อุเปติ,}{อสุทฺธปริสงฺกิโต.} \\ \csromangathacontbat{โส ปุคฺคโล ปฏิญฺญาย,}{กาตพฺพา เตน ปวารณา.} \\ \csromangathacontbat{มุตํ มุเตน สเมติ,}{มุเตน สํสนฺทเต มุตํ.} \\ \csromangathacontbat{มุตํ ปฏิจฺจ น อุเปติ,}{อสุทฺธปริสงฺกิโต.} \\ \csromangathacontbat{โส ปุคฺคโล ปฏิญฺญาย,}{กาตพฺพา เตน ปวารณาติ.}}}

\csromanmatikaanchor{mtk.211}
\csromantocmarkref{subsection}{7. ปุจฺฉาวิภาค}{mtk.211}
\bookmark[dest={mtk.211},level=2]{7. ปุจฺฉาวิภาค}
\csromanheader{2}{7. \textbf{ปุจฺฉาวิภาค}}
\proseitem{398}{กิํ เต ทิฏฺฐนฺติ กตมา ปุจฺฉา, กินฺติ เต ทิฏฺฐนฺติ กตมา ปุจฺฉา, กทา เต ทิฏฺฐนฺติ กตมา ปุจฺฉา, กตฺถ เต ทิฏฺฐนฺติ กตมา ปุจฺฉา.}

\proseitem{399}{กิํ เต ทิฏฺฐนฺติ วตฺถุปุจฺฉา, วิปตฺติปุจฺฉา, อาปตฺติปุจฺฉา, อชฺฌาจาร\-ปุจฺฉา. วตฺถุปุจฺฉาติ อฏฺฐ\-ปาราชิกานํ วตฺถุปุจฺฉา, เตวีสสํฆาทิเสสานํ วตฺถุปุจฺฉา, เทฺวอนิยตานํ วตฺถุปุจฺฉา, เทฺว\-จตฺตาริ\-สนิสฺสคฺคิยานํ วตฺถุปุจฺฉา, อฏฺฐาสีติสต\-ปาจิตฺติยานํ วตฺถุปุจฺฉา, ทฺวาทส\-ปาฏิเทสนียานํ วตฺถุปุจฺฉา, ทุกฺกฏานํ วตฺถุปุจฺฉา, ทุพฺภาสิตานํ วตฺถุปุจฺฉา. วิปตฺติ\-ปุจฺฉาติ สีลวิปตฺติ\-ปุจฺฉา, อาจารวิปตฺติ\-ปุจฺฉา, ทิฏฺฐิวิปตฺติ\-ปุจฺฉา, อาชีววิปตฺติ\-ปุจฺฉา. อาปตฺติปุจฺฉาติ ปาราชิกาปตฺติ\-ปุจฺฉา, สํฆาทิเสสาปตฺติปุจฺฉา, ถุลฺลจฺจยา\-ปตฺติปุจฺฉา, ปาจิตฺติยาปตฺติ\-ปุจฺฉา, ปาฏิเทสนียา\-ปตฺติปุจฺฉา, ทุกฺกฏาปตฺติ\-ปุจฺฉา, ทุพฺภาสิตา\-ปตฺติปุจฺฉา. อชฺฌาจาร\-ปุจฺฉาติ ทฺวยํ\-ทฺวย\-สมาปตฺติปุจฺฉา.}

\proseitem{400}{กินฺติ เต ทิฏฺฐนฺติ ลิงฺคปุจฺฉา, อิริยาปถ\-ปุจฺฉา, อาการปุจฺฉา, วิปฺปการ\-ปุจฺฉา. ลิงฺคปุจฺฉาติ ทีฆํ วา รสฺสํ วา กณฺหํ วา โอทาตํ วา. อิริยาปถ\-ปุจฺฉาติ คจฺฉนฺตํ วา ฐิตํ วา นิสินฺนํ วา นิปนฺนํ วา. อาการปุจฺฉาติ \csromanfolio{301}คิหิลิงฺเค วา ติตฺถิยลิงฺเค วา ปพฺพชิตลิงฺเค วา. วิปฺปการ\-ปุจฺฉาติ คจฺฉนฺตํ วา ฐิตํ วา นิสินฺนํ วา นิปนฺนํ วา.}

\proseitem{401}{กทา เต ทิฏฺฐนฺติ กาลปุจฺฉา, สมยปุจฺฉา, ทิวสปุจฺฉา, อุตุปุจฺฉา. กาลปุจฺฉาติ ปุพฺพณฺหกาเล วา มชฺฌนฺหิก\-กาเล วา สายนฺหกาเล วา. สมยปุจฺฉาติ ปุพฺพณฺหสมเย วา มชฺฌนฺหิก\-สมเย วา สายนฺหสมเย วา. ทิวสปุจฺฉาติ ปุเรภตฺตํ วา ปจฺฉาภตฺตํ วา รตฺติํ วา ทิวา วา กาเฬ วา ชุเณฺห วา. อุตุปุจฺฉาติ เหมนฺเต วา คิเมฺห วา วสฺเส วา\csromansharedfootnote{ab3ba278850c0abd}{วสฺสาเน วา (สี)}.}

\proseitemwithcloser{402}{\textbf{กตฺถ เต ทิฏฺฐนฺ}ติ ฐานปุจฺฉา, ภูมิปุจฺฉา, โอกาสปุจฺฉา, ปเทสปุจฺฉา. ฐานปุจฺฉาติ ภูมิยา วา ปถวิยา วา ธรณิยา วา ชคติยา วา. ภูมิปุจฺฉาติ ภูมิยา วา ปถวิยา วา ปพฺพเต วา ปาสาเณ วา ปาสาเท วา. โอกาสปุจฺฉาติ ปุรตฺถิเม วา โอกาเส ปจฺฉิเม วา โอกาเส อุตฺตเร วา โอกาเส ทกฺขิเณ วา โอกาเส. ปเทสปุจฺฉาติ ปุรตฺถิเม วา ปเทเส ปจฺฉิเม วา ปเทเส อุตฺตเร วา ปเทเส ทกฺขิเณ วา ปเทเสติ.}{มหาสงฺคามํ นิฏฺฐิตํ.}
\csromancenter{ตสฺสุทฺทานํ.}
\setlength{\csromangathapagewidth}{0pt}
\setlength{\csromangathawakwidth}{0pt}
\csromangathameasuregroup{วตฺถุ นิทานํ อากาโร, \\ กมฺมาธิกรณญฺเจว, \\ โทสา โมหา ภยา เจว, \\ เปกฺขา ปสาเท ปกฺโขมฺหิ, \\ อสมฺปตฺตญฺจ สมฺปตฺตํ, \\ สตฺถุสฺส สาสเนนาปิ,}{\csromangathabat{วตฺถุ นิทานํ อากาโร,}{ปุพฺพาปรํ กตากตํ.} \\ \csromangathabat{กมฺมาธิกรณญฺเจว,}{สมโถ ฉนฺทคามิ จ.} \\ \csromangathabat{โทสา โมหา ภยา เจว,}{สญฺญา นิชฺฌาปเนน จ.} \\ \csromangathabat{เปกฺขา ปสาเท ปกฺโขมฺหิ,}{สุตเถรตเรน จ.} \\ \csromangathabat{อสมฺปตฺตญฺจ สมฺปตฺตํ,}{ธมฺเมน วินเยน จ.} \\ \csromangathabat{สตฺถุสฺส สาสเนนาปิ,}{มหาสงฺคามญาปนาติ.}}
\csromangathasetleft{วตฺถุ นิทานํ อากาโร, \\ กมฺมาธิกรณญฺเจว, \\ โทสา โมหา ภยา เจว, \\ เปกฺขา ปสาเท ปกฺโขมฺหิ, \\ อสมฺปตฺตญฺจ สมฺปตฺตํ, \\ สตฺถุสฺส สาสเนนาปิ,}
\csromangathagroup{\csromangathastanza{\csromangathabat{วตฺถุ นิทานํ อากาโร,}{ปุพฺพาปรํ กตากตํ.} \\ \csromangathabat{กมฺมาธิกรณญฺเจว,}{สมโถ ฉนฺทคามิ จ.}}\csromangathastanza{\csromangathabat{โทสา โมหา ภยา เจว,}{สญฺญา นิชฺฌาปเนน จ.} \\ \csromangathabat{เปกฺขา ปสาเท ปกฺโขมฺหิ,}{สุตเถรตเรน จ.}}\csromangathastanza{\csromangathabat{อสมฺปตฺตญฺจ สมฺปตฺตํ,}{ธมฺเมน วินเยน จ.} \\ \csromangathabat{สตฺถุสฺส สาสเนนาปิ,}{มหาสงฺคามญาปนาติ.}}}

\csromanmatikaanchor{mtk.212}
\csromantocmarknopage{section}{กถินเภท}{mtk.212}
\bookmark[dest={mtk.212},level=1]{กถินเภท}
\csromanheader{1}{\textbf{กถินเภท}}
\csromanmatikaanchor{mtk.213}
\csromantocmarkref{paragraph}{1. กถินอตฺถตาทิ}{mtk.213}
\bookmark[dest={mtk.213},level=4]{1. กถินอตฺถตาทิ}
\csromanheader{4}{1. กถิน\-อตฺถ\-ตาทิ}
\proseitem{403}{\csromanfolio{302}กสฺส กถินํ\csromansharedfootnote{735cdf96b0305619}{กฐินํ(สี, สฺยา)} อนตฺถตํ, กสฺส กถินํ อตฺถตํ, กินฺติ กถินํ อนตฺถตํ, กินฺติ กถินํ อตฺถตํ.}

\prose{\textbf{กสฺส กถินํ อนตฺถตนฺ}ติ ทฺวินฺนํ ปุคฺคลานํ อนตฺถตํ โหติ กถินํ อนตฺถารกสฺส จ อนนุโมทกสฺส จ. อิเมสํ ทฺวินฺนํ ปุคฺคลานํ อนตฺถตํ โหติ กถินํ.}

\prose{\textbf{กสฺส กถินํ อตฺถตนฺ}ติ ทฺวินฺนํ ปุคฺคลานํ อตฺถตํ โหติ กถินํ อตฺถารกสฺส จ อนุโมทกสฺส จ. อิเมสํ ทฺวินฺนํ ปุคฺคลานํ อตฺถตํ โหติ กถินํ.}

\prose{กินฺติ กถินํ อนตฺถตนฺติ จตุวีสติยา อากาเรหิ อนตฺถตํ โหติ กถินํ, น อุลฺลิขิต\-มตฺเตน อตฺถตํ โหติ กถินํ, น โธวนมตฺเตน อตฺถตํ โหติ กถินํ, น จีวรวิจารณ\-มตฺเตน อตฺถตํ โหติ กถินํ, น เฉทนมตฺเตน อตฺถตํ โหติ กถินํ, น พนฺธน\-มตฺเตน อตฺถตํ โหติ กถินํ, น โอวฏฺฏิย\-กรณ\-มตฺเตน อตฺถตํ โหติ กถินํ, น กณฺฑุส\-กรณ\-มตฺเตน\csromansharedfootnote{45513aeaaa2087d1}{น คณฺฑุสกรณมตฺเตน (ก)} อตฺถตํ โหติ กถินํ, น ทฬฺหีกมฺม\-กรณ\-มตฺเตน อตฺถตํ โหติ กถินํ, น อนุวาต\-กรณ\-มตฺเตน อตฺถตํ โหติ กถินํ, น ปริภณฺฑ\-กรณ\-มตฺเตน อตฺถตํ โหติ กถินํ, น โอวทฺเธยฺย\-กรณ\-มตฺเตน\csromansharedfootnote{ae8546808d609fe6}{น โอวฏฺเฏยฺยกรณมตฺเตน (สี, สฺยา), น โอวเทยฺยกรณมตฺเตน (ก)} อตฺถตํ โหติ กถินํ, น กมฺพล\-มทฺทน\-มตฺเตน อตฺถตํ โหติ กถินํ, น นิมิตฺตกเตน อตฺถตํ โหติ กถินํ, น ปริกถา\-กเตน อตฺถตํ โหติ กถินํ, น กุกฺกุกเตน อตฺถตํ โหติ กถินํ, น สนฺนิธิกเตน อตฺถตํ โหติ กถินํ, น นิสฺสคฺคิเยน อตฺถตํ โหติ กถินํ, น อกปฺปกเตน อตฺถตํ โหติ กถินํ, น อญฺญตฺร สงฺฆาฏิยา อตฺถตํ โหติ กถินํ, น อญฺญตฺร อุตฺตราสงฺเคน อตฺถตํ โหติ กถินํ, น อญฺญตฺร อนฺตรวาสเกน อตฺถตํ โหติ กถินํ, น อญฺญตฺร ปญฺจเกน วา อติเรก\-ปญฺจเกน วา ตทเหว สญฺฉินฺเนน สมณฺฑลีกเตน อตฺถตํ โหติ กถินํ, น อญฺญตฺร ปุคฺคลสฺส อตฺถารา \csromanfolio{303}อตฺถตํ โหติ กถินํ, สมฺมา เจ อตฺถตํ โหติ กถินํ, ตญฺเจ นิสฺสีมฏฺโฐ อนุโมทติ, เอวมฺปิ อนตฺถตํ โหติ กถินํ.}

\prose{นิมิตฺตกมฺมํ นาม นิมิตฺตํ กโรติ “อิมินา ทุสฺเสน กถินํ อตฺถริสฺสามี”ติ. ปริกถา นาม ปริกถํ กโรติ “อิมาย ปริกถาย กถินทุสฺสํ นิพฺพตฺเตสฺสามี”ติ. กุกฺกุกตํ นาม อนาทิยทานํ วุจฺจติ. สนฺนิธิ นาม เทฺว สนฺนิธิโย กรณสนฺนิธิ วา นิจยสนฺนิธิ วา. นิสฺสาคฺคิยํ นาม กยิรมาเน อรุณํ อุฏฺฐหติ\csromansharedfootnote{ba85837a553c53a6}{อุทฺริยติ (สี, สฺยา)}. อิเมหิ จตุวีสติยา อากาเรหิ อนตฺถตํ โหติ กถินํ.}

\prose{กินฺติ กถินํ อตฺถตนฺติ สตฺตรสหิ อากาเรหิ อตฺถตํ โหติ กถินํ, อหเตน อตฺถตํ โหติ กถินํ, อหตกปฺเปน อตฺถตํ โหติ กถินํ, ปิโลติกาย อตฺถตํ โหติ กถินํ, ปํสุกูเลน อตฺถตํ โหติ กถินํ, ปาปณิเกน อตฺถตํ โหติ กถินํ, อนิมิตฺตกเตน อตฺถตํ โหติ กถินํ, อปริกถากเตน อตฺถตํ โหติ กถินํ, อกุกฺกุกเตน อตฺถตํ โหติ กถินํ, อสนฺนิธิกเตน อตฺถตํ โหติ กถินํ, อนิสฺสคฺคิเยน อตฺถตํ โหติ กถินํ, กปฺปกเตน อตฺถตํ โหติ กถินํ, สํฆาฏิยา อตฺถตํ โหติ กถินํ, อุตฺตราสงฺเคน อตฺถตํ โหติ กถินํ, อนฺตรวาสเกน อตฺถตํ โหติ กถินํ, ปญฺจเกน วา อติเรก\-ปญฺจเกน วา ตทเหว สญฺฉิ นฺเนน สมณฺฑลีกเตน อตฺถตํ โหติ กถินํ, ปุคฺคลสฺส อตฺถารา อตฺถตํ โหติ กถินํ, สมฺมา เจ อตฺถตํ โหติ กถินํ, ตญฺเจ สีมฏฺโฐ อนุโมทติ, เอวมฺปิ อตฺถตํ โหติ กถินํ. อิเมหิ สตฺตรสหิ อากาเรหิ อตฺถตํ โหติ กถินํ.}

\prose{สห กถินสฺส อตฺถารา กติ ธมฺมา ชายนฺติ. สห กถินสฺส อตฺถารา ปนฺนรส ธมฺมา ชายนฺติ, อฏฺฐ มาติกา เทฺว ปลิโพธา ปญฺจานิสํสา. สห กถินสฺส อตฺถารา อิเม ปนฺนรส ธมฺมา ชายนฺติ.}

\csromanmatikaanchor{mtk.214}
\csromantocmarkref{paragraph}{2. กถินอนนฺตรปจฺจยาทิ}{mtk.214}
\bookmark[dest={mtk.214},level=4]{2. กถินอนนฺตรปจฺจยาทิ}
\csromanheader{4}{2. กถิน\-อนนฺตรปจฺจยาทิ}
\proseitem{404}{ปโยคสฺส กตเม ธมฺมา อนนฺตร\-ปจฺจเยน ปจฺจโย, สมนนฺตร\-ปจฺจเยน ปจฺจโย, นิสฺสย\-ปจฺจเยน ปจฺจโย, อุปนิสฺสย\-ปจฺจเยน \csromanfolio{304}ปจฺจโย, ปุเรชาต\-ปจฺจเยน ปจฺจโย, ปจฺฉาชาต\-ปจฺจเยน ปจฺจโย, สหชาต\-ปจฺจเยน ปจฺจโย. ปุพฺพกรณสฺส กตเม ธมฺมา อนนฺตร\-ปจฺจเยน ปจฺจโย ฯเปฯ. ปจฺจุทฺธารสฺส กตเม ธมฺมา. อธิฏฺฐานสฺส กตเม ธมฺมา. อตฺถารสฺส กตเม ธมฺมา. มาติกานญฺจ ปลิโพธานญฺจ กตเม ธมฺมา. วตฺถุสฺส กเตเม ธมฺมา อนนฺตร\-ปจฺจเยน ปจฺจโย, สมนนฺตร\-ปจฺจเยน ปจฺจโย, นิสฺสย\-ปจฺจเยน ปจฺจโย, อุปนิสฺสย\-ปจฺจเยน ปจฺจโย, ปุเรชาต\-ปจฺจเยน ปจฺจโย, ปจฺฉาชาต\-ปจฺจเยน ปจฺจโย, สหชาต\-ปจฺจเยน ปจฺจโย.}

\prose{ปุพฺพกรณํ ปโยคสฺส อนนฺตร\-ปจฺจเยน ปจฺจโย, สมนนฺตร\-ปจฺจเยน ปจฺจโย, นิสฺสย\-ปจฺจเยน ปจฺจโย, อุปนิสฺสย\-ปจฺจเยน ปจฺจโย. ปโยโค ปุพฺพกรณสฺส ปุเรชาต\-ปจฺจเยน ปจฺจโย. ปุพฺพกรณํ ปโยคสฺส ปจฺฉาชาต\-ปจฺจเยน ปจฺจโย. ปนฺนรส ธมฺมา สหชาต\-ปจฺจเยน ปจฺจโย. ปจฺจุทฺธาโร ปุพฺพกรณสฺส อนนฺตร\-ปจฺจเยน ปจฺจโย, สมนนฺตร\-ปจฺจเยน ปจฺจโย, นิสฺสย\-ปจฺจเยน ปจฺจโย, อุปนิสฺสย\-ปจฺจเยน ปจฺจโย. ปุพฺพกรณํ ปจฺจุทฺธารสฺส ปุเรชาต\-ปจฺจเยน ปจฺจโย. ปจฺจุทฺธาโร ปุพฺพกรณสฺส ปจฺฉาชาต\-ปจฺจเยน ปจฺจโย. ปนฺนรส ธมฺมา สหชาต\-ปจฺจเยน ปจฺจโย. อธิฏฺฐานํ ปจฺจุทฺธารสฺส อนนฺตร\-ปจฺจเยน ปจฺจโย, สมนนฺตร\-ปจฺจเยน ปจฺจโย, นิสฺสย\-ปจฺจเยน ปจฺจโย, อุปนิสฺสย\-ปจฺจเยน ปจฺจโย. ปจฺจุทฺธาโร อธิฏฺฐานสฺส ปุเรชาต\-ปจฺจเยน ปจฺจโย. อธิฏฺฐานํ ปจฺจุทฺธารสฺส ปจฺฉาชาต\-ปจฺจเยน ปจฺจโย. ปนฺนรส ธมฺมา สหชาต\-ปจฺจเยน ปจฺจโย. อตฺถาโร อธิฏฺฐานสฺส อนนฺตร\-ปจฺจเยน ปจฺจโย, สมนนฺตร\-ปจฺจเยน ปจฺจโย, นิสฺสย\-ปจฺจเยน ปจฺจโย, อุปนิสฺสยจฺจเยน ปจฺจโย. อธิฏฺฐานํ อตฺถารสฺส ปุเรชาต\-ปจฺจเยน ปจฺจโย. อตฺถาโร อธิฏฺฐานสฺส ปจฺฉาชาต\-ปจฺจเยน ปจฺจโย. ปนฺนรส ธมฺมา สหชาต\-ปจฺจเยน ปจฺจโย. มาติกา จ ปลิโพธา จ อตฺถารสฺส อนนฺตร\-ปจฺจเยน ปจฺจโย, สมนนฺตร\-ปจฺจเยน ปจฺจโย, นิสฺสย\-ปจฺจเยน ปจฺจโย, อุปนิสฺสย\-ปจฺจเยน ปจฺจโย. อตฺถาโร มาติกานญฺจ ปลิโพธานญฺจ ปุเรชาต\-ปจฺจเยน ปจฺจโย. มาติกา จ ปลิโพธา จ อตฺถารสฺส ปจฺฉาชาต\-ปจฺจเยน ปจฺจโย. ปนฺนรส ธมฺมา สหชาต\-ปจฺจเยน ปจฺจโย. อาสา จ อนาสา จ วตฺถุสฺส อนนฺตร\-ปจฺจเยน ปจฺจโย, สมนนฺตร\-ปจฺจเยน ปจฺจโย, นิสฺสย\-ปจฺจเยน \csromanfolio{305}ปจฺจโย, อุปนิสฺสย\-ปจฺจเยน ปจฺจโย. วตฺถุ อาสานญฺจ อนาสานญฺจ ปุเรชาต\-ปจฺจเยน ปจฺจโย. อาสา จ อนาสา จ วตฺถุสฺส ปจฺฉาชาต\-ปจฺจเยน ปจฺจโย. ปนฺนรส ธมฺมา สหชาต\-ปจฺจเยน ปจฺจโย.}

\csromanmatikaanchor{mtk.215}
\csromantocmarkref{paragraph}{3. ปุพฺพกรณนิทานาทิวิภาค}{mtk.215}
\bookmark[dest={mtk.215},level=4]{3. ปุพฺพกรณนิทานาทิวิภาค}
\csromanheader{4}{3. ปุพฺพกรณ\-นิทานา\-ทิวิ\-ภาค}
\proseitem{405}{ปุพฺพกรณํ กิํนิทานํ กิํ สมุทยํ กิํชาติกํ กิํปภวํ กิํ สมฺภารํ กิํ สมุฏฺฐานํ. ปจฺจุทฺธาโร กิํนิทาโน กิํสมุทโย กิํชาติโก กิํปภโว กิํสมฺภาโร กิํสมุฏฺฐาโน. อธิฏฺฐานํ กิํนิทานํ กิํสมุทยํ กิํชาติกํ กิํปภวํ กิํสมฺภารํ กิํสมุฏฺฐานํ. อตฺถาโร กิํนิทาโน กิํสมุทโย กิํชาติโก กิํปภโว กิํสมฺภาโร กิํสมุฏฺฐาโน. มาติกา จ ปลิโพธา จ กิํนิทานา กิํสมุทยา กิํชาติกา กิํปภวา กิํสมฺภารา กิํ สมุฏฺฐานา. อาสา จ อนาสา จ กิํนิทานา กิํสมุทยา กิํชาติกา กิํปภวา กิํสมฺภารา กิํสมุฏฺฐานา.}

\prose{ปุพฺพกรณํ ปโยคนิทานํ ปโยค\-สมุทยํ ปโยคชาติกํ ปโยคปภวํ ปโยค\-สมฺภารํ ปโยค\-สมุฏฺฐานํ. ปจฺจุทฺธาโร ปุพฺพกรณ\-นิทาโน ปุพฺพกรณ\-สมุทโย ปุพฺพกรณ\-ชาติโก ปุพฺพกรณ\-ปภโว ปุพฺพกรณ\-สมฺภาโร ปุพฺพกรณ\-สมุฏฺฐาโน. อธิฏฺฐานํ ปจฺจุทฺธาร\-นิทานํ ปจฺจุทฺธาร\-สมุทยํ ปจฺจุทฺธาร\-ชาติกํ ปจฺจุทฺธาร\-ปภวํ ปจฺจุทฺธาร\-สมฺภารํ ปจฺจุทฺธาร\-สมุฏฺฐานํ. อตฺถาโร อธิฏฺฐาน\-นิทาโน อธิฏฺฐาน\-สมุทโย, อธิฏฺฐาน\-ชาติโก, อธิฏฺฐาน\-ปภโว อธิฏฺฐาน\-สมฺภาโร อธิฏฺฐาน\-สมุฏฺฐาโน. มาติกา จ ปลิโพธา จ อตฺถารนิทานา อตฺถาร\-สมุทยา อตฺถารชาติกา อตฺถารปภวา อถารนิทานา อตฺถาร\-สมุทยา อตฺถารชาติกา อตฺถารปภวา อตฺถาร\-สมฺภารา อตฺถาร\-สมุฏฺฐานา. อาสา จ อนาสา จ วตฺถุนิทานา วตฺถุสมุทยา วตฺถุชาติกา วตฺถุปภวา วตฺถุสมฺภารา วตฺถุ\-สมุฏฺฐานา.}

\proseitem{406}{ปโยโค กิํนิทาโน กิํสมุทโย กิํชาติโก กิํปภโว กิํสมฺภาโร กิํสมุฏฺฐาโน. ปุพฺพกรณํ ฯเปฯ. ปจฺจุทฺธาโร. อธิฏฺฐานํ. อตฺถาโร. มาติกา จ ปลิโพธา จ. วตฺถุ. อาสา จ อนาสา จ กิํนิทานา กิํสมุทยา กิํชาติกา กิํปภวา กิํสมฺภารา กิํสมุฏฺฐานา.}

\prose{\csromanfolio{306}ปโยโค เหตุนิทาโน เหตุสมุทโย เหตุชาติโก เหตุปภโว เหตุสมฺภาโร เหตุสมุฏฺฐาโน. ปุพฺพกรณํ ฯเปฯ. ปจฺจุทฺธาโร. อธิฏฺฐานํ. อตฺถาโร. มาติกา จ ปลิโพธา จ. วตฺถุ. อาสา จ อนาสา จ เหตุนิทานา เหตุสมุทยา เหตุชาติกา เหตุปภวา เหตุสมฺภารา เหตุสมุฏฺฐานา.}

\proseitem{407}{ปโยโค กิํนิทาโน กิํสมุทโย กิํชาติโก กิํปภโว กิํสมฺภาโร กิํสมุฏฺฐาโน. ปุพฺพกรณํ ฯเปฯ. ปจฺจุทฺธาโร. อธิฏฺฐานํ. อตฺถาโร. มาติกา จ ปลิโพธา จ. วตฺถุ. อาสา จ อนาสา จ กิํนิทานา กิํสมุทยา กิํชาติกา กิํปภวา กิํสมฺภารา กิํสมุฏฺฐานา.}

\prose{ปโยโค ปจฺจยนิทาโน ปจฺจย\-สมุทโย ปจฺจยชาติโก ปจฺจยปภโว ปจฺจย\-สมฺภาโร ปจฺจย\-สมุฏฺฐาโน. ปุพฺพกรณํ ฯเปฯ. ปจฺจุทฺธาโร. อธิฏฺฐานํ. อตฺถาโร. มาติกา จ ปลิโพธา จ. วตฺถุ. อาสา จ อนาสา จ ปจฺจยนิทานา ปจฺจย\-สมุทยา ปจฺจยชาติกา ปจฺจยปภวา ปจฺจย\-สมฺภารา ปจฺจย\-สมุฏฺฐานา.}

\proseitem{408}{ปุพฺพกรณํ กติหิ ธมฺเมหิ สงฺคหิตํ. ปุพฺพกรณํ สตฺตหิ ธมฺเมหิ สงฺคหิตํ โธวเนน วิจารเณน เฉทเนน พนฺธเนน สิพฺพเนน รชเนน กปฺปกรเณน, ปุพฺพกรณํ อิเมหิ สตฺตหิ ธมฺเมหิ สงฺคหิตํ.}

\prose{ปจฺจุทฺธาโร กติหิ ธมฺเมหิ สงฺคหิโต. ปจฺจุทฺธาโร ตีหิ ธมฺเมหิ สงฺคหิโต สํฆาฏิยา อุตฺตราสงฺเคน อนฺตรวาสเกน.}

\prose{อธิฏฺฐานํ กติหิ ธมฺเมหิ สงฺคหิตํ. อธิฏฺฐานํ ตีหิ ธมฺเมหิ สงฺคหิตํ สํฆาฏิยา อุตฺตราสงฺเคน อนฺตรวาสเกน.}

\prose{อตฺถาโร กติหิ ธมฺเมหิ สงฺคหิโต. อตฺถาโร เอเกน ธมฺเมน สงฺคหิโต วจีเภเทน.}

\prose{กถินสฺส กติ มูลานิ, กติ วตฺถูนิ, กติ ภูมิโย. กถินสฺส เอกํ มูลํ สํโฆ. ตีณิ วตฺถูนิ สํฆาฏิ อุตฺตราสงฺโค อนฺตรวาสโก. ฉ ภูมิโย โขมํ กปฺปาสิกํ โกเสยฺยํ กมฺพลํ สาณํ ภงฺคํ.}

\prose{\csromanfolio{307}กถินสฺส โก อาทิ, กิํ มชฺเฌ, กิํ ปริโยสานํ. กถินสฺส ปุพฺพกรณํ อาทิ, กฺริยา มชฺเฌ, อตฺถาโร ปริโยสานํ.}

\proseitem{409}{กถิหงฺเคหิ สมนฺนาคโต ปุคฺคโล อภพฺโพ กถินํ อตฺถริตุํ. กติหงฺเคหิ สมนฺนาคโต ปุคฺคโล ภพฺโพ กถินํ อตฺถริตุํ. อฏฺฐหงฺเคหิ สมนฺนาคโต ปุคฺคโล อภพฺโพ กถินํ อตฺถริตุํ. อฏฺฐหงฺเคหิ สมนฺนาคโต ปุคฺคโล ภพฺโพ กถินํ อตฺถริตุํ. กตเมหิ อฏฺฐหงฺเคหิ สมนฺนาคโต ปุคฺคโล อภพฺโพ กถินํ อตฺถริตุํ. ปุพฺพกรณํ น ชานาติ, ปจฺจุทฺธารํ น ชานาติ, อธิฏฺฐานํ น ชานาติ, อตฺถารํ น ชานาติ, มาติกํ น ชานาติ, ปลิโพธํ น ชานาติ, อุทฺธารํ น ชานาติ, อานิสํสํ น ชานาติ. อิเมหิ อฏฺฐหงฺเคหิ สมนฺนาคโต ปุคฺคโล อภพฺโพ กถินํ อตฺถริตุํ. กตเมหิ อฏฺฐหงฺเคหิ สมนฺนาคโต ปุคฺคโล ภพฺโพ กถินํ อตฺถริตุํ. ปุพฺพกรณํ ชานาติ, ปจฺจุทฺธารํ ชานาติ, อธิฏฺฐานํ ชานาติ, อตฺถารํ ชานาติ, มาติกํ ชานาติ, ปลิโพธํ ชานาติ, อุทฺธารํ ชานาติ, อานิสํสํ ชานาติ. อิเมหิ อฏฺฐหงฺเคหิ สมนฺนาคโต ปุคฺคโล ภพฺโพ กถินํ อตฺถริตุํ.}

\proseitem{410}{กตินํ ปุคฺคลานํ กถินตฺถารา น รุหนฺติ, กตินํ ปุคฺคลานํ กถินตฺถารา รุหนฺติ. ติณฺณํ ปุคฺคลานํ กถินตฺถารา น รุหนฺติ, ติณฺณํ ปุคฺคลานํ กถินตฺถารา รุหนฺติ. กตเมสํ ติณฺณํ ปุคฺคลานํ กถินตฺถารา น รุหนฺติ. นิสฺสีมฏฺโฐ อนุโมทติ, อนุโมเทนฺโต น วาจํ ภินฺทติ, วาจํ ภินฺทนฺโต น ปรํ วิญฺญาเปติ, อิเมสํ ติณฺณํ ปุคฺคลานํ กถินตฺถารา น รุหนฺติ. กตเมสํ ติณฺณํ ปุคฺคลานํ กถินตฺถารา รุหนฺติ. สีมฏฺโฐ อนุโมทติ, อนุโมเทนฺโต วาจํ ภินฺทติ, วาจํ ภินฺทนฺโต ปรํ วิญฺญาเปติ, อิเมสํ ติณฺณํ ปุคฺคลานํ กถินตฺถารา รุหนฺติ.}

\proseitem{411}{กติ กถินตฺถารา น รุหนฺติ, กติ กถินตฺถารา รุหนฺติ. ตโย กถินตฺถารา น รุหนฺติ, ตโย กถินตฺถารา รุหนฺติ. กตเม ตโย กถินตฺถารา น รุหนฺติ, วตฺถุวิปนฺน\-ญฺเจว โหติ กาลวิปนฺนญฺจ กรณ\-วิปนฺนญฺจ, อิเม ตโย กถินตฺถารา น รุหนฺติ. กตเม ตโย กถินตฺถารา รุหนฺติ, วตฺถุ\-สมฺปนฺน\-ญฺเจว โหติ กาล\-สมฺปนฺนญฺจ กรณ\-สมฺปนฺนญฺจ, อิเม ตโย กถินตฺถารา รุหนฺติ.}

\csromanmatikaanchor{mtk.216}
\csromantocmarkref{paragraph}{4. กถินาทิชานิตพฺพวิภาค}{mtk.216}
\bookmark[dest={mtk.216},level=4]{4. กถินาทิชานิตพฺพวิภาค}
\csromanheader{4}{4. กถินาทิ\-ชานิตพฺพ\-วิภาค}
\proseitem{412}{\csromanfolio{308}กถินํ ชานิตพฺพํ, กถินตฺถาโร ชานิตพฺโพ, กถินสฺส อตฺถารมาโส ชานิตพฺโพ, กถินสฺส อตฺถารวิปตฺติ ชานิตพฺพา, กถินสฺส อตฺถาร\-สมฺปตฺติ ชานิตพฺพา, นิมิตฺตกมฺมํ ชานิตพฺพํ, ปริกถา ชานิตพฺพา, กุกฺกุกตํ ชานิตพฺพํ, สนฺนิธิ ชานิตพฺพา, นิสฺสคฺคิยํ ชานิตพฺพํ.}

\prose{\textbf{กถินํ ชานิตพฺพนฺ}ติ เตสญฺเญว ธมฺมานํ สงฺคโห สมวาโย นามํ นามกมฺมํ นามเธยฺยํ นิรุตฺติ พฺยญฺชนํ อภิลาโป ยทิทํ กถินนฺติ.}

\prose{\textbf{กถินสฺส อตฺถารมาโส ชานิตพฺโพ}ติ วสฺสานสฺส ปจฺฉิโม มาโส ชานิตพฺโพ.}

\prose{\textbf{กถินสฺส อตฺถารวิปตฺติ ชานิตพฺพา}ติ จตุวีสติยา อากาเรหิ กถินสฺส อตฺถารวิปตฺติ ชานิตพฺพา.}

\prose{กถินสฺส อตฺถาร\-สมฺปตฺติ ชานิตพฺพาติ สตฺตรสหิ อากาเรหิ กถินสฺส อตฺถาร\-สมฺปตฺติ ชานิตพฺพา.}

\prose{\textbf{นิมิตฺตกมฺมํ ชานิตพฺพนฺ}ติ นิมิตฺตํ กโรติ “อิมินา ทุสฺเสน กถินํ อตฺถริสฺสามี”ติ.}

\prose{\textbf{ปริกถา ชานิตพฺพา}ติ ปริกถํ กโรติ “อิมาย ปริกถาย กถินทุสฺสํ นิพฺพตฺเตสฺสามี”ติ.}

\prose{\textbf{กุกฺกุกตํ ชานิตพฺพนฺ}ติ อนาทิยทานํ ชานิตพฺพํ.}

\prose{\textbf{สนฺนิธิ ชานิตพฺพา}ติ เทฺว สนฺนิธิโย ชานิตพฺพา, กรณสนฺนิธิ วา นิจยสนฺนิธิ วา.}

\prose{\textbf{นิสฺสคฺคิยํ ชานิตพฺพนฺ}ติ กริยมาเน อรุณํ อุฏฺฐหติ.}

\prose{\textbf{กถินตฺถาโร ชานิตพฺโพ}ติ สเจ สํฆสฺส กถินทุสฺสํ อุปฺปนฺนํ โหติ, สํเฆน กถํ ปฏิปชฺชิตพฺพํ, อตฺถารเกน กถํ ปฏิปชฺชิตพฺพํ, อนุโมทเกน กถํ ปฏิปชฺชิตพฺพํ.}

\proseitem{413}{สํเฆน ญตฺติทุติเยน กมฺเมน กถินตฺถาร\-กสฺส ภิกฺขุโน ทาตพฺพํ, เตน กถินตฺถาร\-เกน ภิกฺขุนา ตทเหว โธวิตฺวา วิมชฺชิตฺวา วิจาเรตฺวา ฉินฺทิตฺวา สิพฺเพตฺวา รชิตฺวา กปฺปํ กตฺวา กถินํ อตฺถริตพฺพํ. \csromanfolio{309}สเจ สํฆาฏิยา กถินํ อตฺถริตุกาโม โหติ, โปราณิกา สํฆาฏิ ปจฺจุทฺธริตพฺพา, น วา สํฆาฏิ อธิฏฺฐาตพฺพา. “อิมาย สํฆาฏิยา กถินํ อตฺถรามี”ติ วาจา ภินฺทิตพฺพา. สเจ อุตฺตราสงฺเคน กถินํ อตฺถริตุกาโม โหติ, โปราณโก อุตฺตราสงฺโค ปจฺจุทฺธริตพฺโพ, นโว อุตฺตราสงฺโค อธิฏฺฐาตพฺโพ. “อิมินา อุตฺตราสงฺเคน กถินํ อตฺถรามี”ติ วาจา ภินฺทิตพฺพา. สเจ อนฺตรวาสเกน กถินํ อตฺถริตุกาโม โหติ, โปราณโก อนฺตรวาสโก ปจฺจุทฺธริตพฺโพ, นโว อนฺตรวาสโก อธิฏฺฐาตพฺโพ. “อิมินา อนฺตรวาสเกน กถินํ อตฺถรามี”ติ วาจา ภินฺทิตพฺพา. เตน กถินตฺถาร\-เกน ภิกฺขุนา สํฆํ อุปสงฺกมิตฺวา เอกํสํ อุตฺตราสงฺคํ กริตฺวา อญฺชลิํ ปคฺคเหตฺวา เอวมสฺส วจนีโย “อตฺถตํ ภนฺเต สํฆสฺส กถินํ, ธมฺมิโก กถินตฺถาโร, อนุโมทถา”ติ. เตหิ อนุโมทเกหิ ภิกฺขูหิ เอกํสํ อุตฺตราสงฺคํ กริตฺวา อญฺชลิํ ปคฺคเหตฺวา เอวมสฺส วจนีโย “อตฺถตํ อาวุโส สํฆสฺส กถินํ, ธมฺมิโก กถินตฺถาโร, อนุโมทามา”ติ. เตน กถินตฺถาร\-เกน ภิกฺขุนา สมฺพหุเล ภิกฺขู อุปสงฺกมิตฺวา เอกํสํ อุตฺตราสงฺคํ กริตฺวา อญฺชลิํ ปคฺคเหตฺวา เอวมสฺสุ วจนียา “อตฺถตํ ภนฺเต สํฆสฺส กถินํ, ธมฺมิโก กถินตฺถาโร, อนุโมทถา”ติ. เตหิ อนุโมทเกหิ ภิกฺขูหิ เอกํสํ อุตฺตราสงฺคํ กริตฺวา อญฺชลิํ ปคฺคเหตฺวา เอวมสฺส วจนีโย “อตฺถตํ อาวุโส สํฆสฺส กถินํ, ธมฺมิโก กถินตฺถาโร, อนุโมทามา”ติ. เตน กถินตฺถาร\-เกน ภิกฺขุนา เอกํ ภิกฺขุํ อุปสงฺกมิตฺวา เอกํสํ อุตฺตราสงฺคํ กริตฺวา อญฺชลิํ ปคฺคเหตฺวา เอวมสฺส วจนีโย “อตฺถตํ อาวุโส สํฆสฺส กถินํ, ธมฺมิโก กถินตฺถาโร, อนุโมทาหี”ติ. เตน อนุโมทเกน ภิกฺขุนา เอกํสํ อุตฺตราสงฺคํ กริตฺวา อญฺชลิํ ปคฺคเหตฺวา เอวมสฺส วจนีโย “อตฺถตํ อาวุโส สํฆสฺส กถินํ, ธมฺมิโก กถินตฺถาโร, อนุโมทามี”ติ.}

\csromanmatikaanchor{mtk.217}
\csromantocmarkref{paragraph}{5. ปุคฺคลสฺเสวกถินตฺถาร}{mtk.217}
\bookmark[dest={mtk.217},level=4]{5. ปุคฺคลสฺเสวกถินตฺถาร}
\csromanheader{4}{5. ปุคฺคลสฺเส\-ว\-กถินตฺถาร}
\proseitem{414}{สํโฆ กถินํ อตฺถรติ, คโณ กถินํ อตฺถรติ, ปุคฺคโล กถินํ อตฺถรตีติ? น สํโฆ กถินํ อตฺถรติ, น คโณ กถินํ อตฺถรติ, ปุคฺคโล กถินํ อตฺถรตีติ. หญฺจิ น สํโฆ กถินํ อตฺถรติ, น คโณ กถินํ อตฺถรติ, ปุคฺคโล กถินํ อตฺถรติ. สํฆสฺส \csromanfolio{310}อนตฺถตํ โหติ กถินํ, คณสฺส อนตฺถตํ โหติ กถินํ, ปุคฺคลสฺส อตฺถตํ โหติ กถินํ. สํโฆ ปาติโมกฺขํ อุทฺทิสติ, คโณ ปาติโมกฺขํ อุทฺทิสติ, ปุคฺคโล ปาติโมกฺขํ อุทฺทิสตีติ? น สํโฆ ปาติโมกฺขํ อุทฺทิสติ, น คโณ ปาติโมกฺขํ อุทฺทิสติ, ปุคฺคโล ปาติโมกฺขํ อุทฺทิสตีติ. หญฺจิ น สํโฆ ปาติโมกฺขํ อุทฺทิสติ, น คโณ ปาติโมกฺขํ อุทฺทิสติ, ปุคฺคโล ปาติโมกฺขํ อุทฺทิสติ. สํฆสฺส อนุทฺทิฏฺฐํ โหติ ปาติโมกฺขํ, คณสฺส อนุทฺทิฏฺฐํ โหติ ปาติโมกฺขํ, ปุคฺคลสฺส อุทฺทิฏฺฐํ โหติ ปาติโมกฺขํ. สํฆสฺส สามคฺคิยา คณสฺส สามคฺคิยา ปุคฺคลสฺส อุทฺเทสา สํฆสฺส อุทฺทิฏฺฐํ โหติ ปาติโมกฺขํ, คณสฺส อุทฺทิฏฺฐํ โหติ ปาติโมกฺขํ, ปุคฺคลสฺส อุทฺทิฏฺฐํ โหติ ปาติโมกฺขํ. เอวเมว น สํโฆ กถินํ อตฺถรติ, น คโณ กถินํ อตฺถรติ, ปุคฺคโล กถินํ อตฺถรติ. สํฆสฺส อนุโมทนาย คณสฺส อนุโมทนาย ปุคฺคลสฺส อตฺถารา สํฆสฺส อตฺถตํ โหติ กถินํ, คณสฺส อตฺถตํ โหติ กถินํ, ปุคฺคลสฺส อตฺถตํ โหติ กถินนฺติ.}

\csromanmatikaanchor{mtk.218}
\csromantocmarkref{paragraph}{6. ปลิโพธปญฺหาพฺยากรณ}{mtk.218}
\bookmark[dest={mtk.218},level=4]{6. ปลิโพธปญฺหาพฺยากรณ}
\csromanheader{4}{6. ปลิโพธ\-ปญฺหาพฺยากรณ}
\setlength{\csromangathapagewidth}{0pt}
\setlength{\csromangathawakwidth}{0pt}
\csromangathameasuregroup{ปกฺกมนนฺติโก กถินุทฺธาโร, \\ เอตญฺจ ตาหํ ปุจฺฉามิ, \\ ปกฺกมนนฺติโก กถินุทฺธาโร, \\ เอตญฺจ ตาหํ วิสฺสชฺชิสฺสํ,}{\csromangathabat{ปกฺกมนนฺติโก กถินุทฺธาโร,}{วุตฺโต อาทิจฺจพนฺธุนา.} \\ \csromangathabat{เอตญฺจ ตาหํ ปุจฺฉามิ,}{กตโม ปลิโพโธ ปฐมํ ฉิชฺชติ.} \\ \csromangathabat{ปกฺกมนนฺติโก กถินุทฺธาโร,}{วุตฺโต อาทิจฺจพนฺธุนา.} \\ \csromangathabat{เอตญฺจ ตาหํ วิสฺสชฺชิสฺสํ,}{จีวรปลิโพโธ ปฐมํ ฉิชฺชติ.}}
\csromangathasetleft{ปกฺกมนนฺติโก กถินุทฺธาโร, \\ เอตญฺจ ตาหํ ปุจฺฉามิ, \\ ปกฺกมนนฺติโก กถินุทฺธาโร, \\ เอตญฺจ ตาหํ วิสฺสชฺชิสฺสํ,}
\csromangathagroup{\csromangathastanza{\csromangathaitembat{415}{ปกฺกมนนฺติโก กถินุทฺธาโร,}{วุตฺโต อาทิจฺจพนฺธุนา.} \\ \csromangathacontbat{เอตญฺจ ตาหํ ปุจฺฉามิ,}{กตโม ปลิโพโธ ปฐมํ ฉิชฺชติ.}}\csromangathastanza{\csromangathaitembat{415}{ปกฺกมนนฺติโก กถินุทฺธาโร,}{วุตฺโต อาทิจฺจพนฺธุนา.} \\ \csromangathacontbat{เอตญฺจ ตาหํ วิสฺสชฺชิสฺสํ,}{จีวร\-ปลิโพโธ ปฐมํ ฉิชฺชติ.}}}

\prose{ตสฺส สห พหิสีมคมนา, อาวาสปลิโพโธ ฉิชฺชติ. (1)}

\setlength{\csromangathapagewidth}{0pt}
\setlength{\csromangathawakwidth}{0pt}
\csromangathameasuregroup{นิฏฺฐานนฺติโก กถินุทฺธาโร, \\ เอตญฺจ ตาหํ ปุจฺฉามิ, \\ นิฏฺฐานนฺติโก กถินุทฺธาโร, \\ เอตญฺจ ตาหํ วิสฺสชฺชิสฺสํ,}{\csromangathabat{นิฏฺฐานนฺติโก กถินุทฺธาโร,}{วุตฺโต อาทิจฺจพนฺธุนา.} \\ \csromangathabat{เอตญฺจ ตาหํ ปุจฺฉามิ,}{กตโม ปลิโพโธ ปฐมํ ฉิชฺชติ.} \\ \csromangathabat{นิฏฺฐานนฺติโก กถินุทฺธาโร,}{วุตฺโต อาทิจฺจพนฺธุนา.} \\ \csromangathabat{เอตญฺจ ตาหํ วิสฺสชฺชิสฺสํ,}{อาวาสปลิโพโธ ปฐมํ ฉิชฺชติ.}}
\csromangathasetleft{นิฏฺฐานนฺติโก กถินุทฺธาโร, \\ เอตญฺจ ตาหํ ปุจฺฉามิ, \\ นิฏฺฐานนฺติโก กถินุทฺธาโร, \\ เอตญฺจ ตาหํ วิสฺสชฺชิสฺสํ,}
\csromangathagroup{\csromangathastanza{\csromangathabat{นิฏฺฐานนฺติโก กถินุทฺธาโร,}{วุตฺโต อาทิจฺจพนฺธุนา.} \\ \csromangathabat{เอตญฺจ ตาหํ ปุจฺฉามิ,}{กตโม ปลิโพโธ ปฐมํ ฉิชฺชติ.}}\csromangathastanza{\csromangathabat{นิฏฺฐานนฺติโก กถินุทฺธาโร,}{วุตฺโต อาทิจฺจพนฺธุนา.} \\ \csromangathabat{เอตญฺจ ตาหํ วิสฺสชฺชิสฺสํ,}{อาวาสปลิโพโธ ปฐมํ ฉิชฺชติ.}}}

\prose{จีวเร นิฏฺฐิเต จีวร\-ปลิโพโธ ฉิชฺชติ. (2)}

\setlength{\csromangathapagewidth}{0pt}
\setlength{\csromangathawakwidth}{0pt}
\csromangathameasuregroup{สนฺนิฏฺฐานนฺติโก กถินุทฺธาโร, \\ เอตญฺจ ตาหํ ปุจฺฉามิ, \\ สนฺนิฏฺฐานนฺติโก กถินุทฺธาโร,}{\csromangathabat{สนฺนิฏฺฐานนฺติโก กถินุทฺธาโร,}{วุตฺโต อาทิจฺจพนฺธุนา.} \\ \csromangathabat{เอตญฺจ ตาหํ ปุจฺฉามิ,}{กตโม ปลิโพโธ ปฐมํ ฉิชฺชติ.} \\ \csromangathabat{สนฺนิฏฺฐานนฺติโก กถินุทฺธาโร,}{วุตฺโต อาทิจฺจพนฺธุนา.}}
\csromangathasetleft{สนฺนิฏฺฐานนฺติโก กถินุทฺธาโร, \\ เอตญฺจ ตาหํ ปุจฺฉามิ, \\ สนฺนิฏฺฐานนฺติโก กถินุทฺธาโร,}
\csromangathagroup{\csromangathastanza{\csromangathabat{สนฺนิฏฺฐานนฺติโก กถินุทฺธาโร,}{วุตฺโต อาทิจฺจพนฺธุนา.} \\ \csromangathabat{เอตญฺจ ตาหํ ปุจฺฉามิ,}{กตโม ปลิโพโธ ปฐมํ ฉิชฺชติ.} \\ \csromangathabat{สนฺนิฏฺฐานนฺติโก กถินุทฺธาโร,}{วุตฺโต อาทิจฺจพนฺธุนา.}}}

\prose{เอตญฺจ ตาหํ วิสฺสชฺชิสฺสํ,}

\prose{เทฺว ปลิโพธา อปุพฺพํ อจริมํ ฉิชฺชนฺติ. (3)}

\setlength{\csromangathapagewidth}{0pt}
\setlength{\csromangathawakwidth}{0pt}
\csromangathameasuregroup{นาสนนฺติโก กถินุทฺธาโร, \\ เอตญฺจ ตาหํ ปุจฺฉามิ, \\ นาสนนฺติโก กถินุทฺธาโร, \\ เอตญฺจ ตาหํ วิสฺสชฺชิสฺสํ,}{\csromangathabat{นาสนนฺติโก กถินุทฺธาโร,}{วุตฺโต อาทิจฺจพนฺธุนา.} \\ \csromangathabat{เอตญฺจ ตาหํ ปุจฺฉามิ,}{กตโม ปลิโพโธ ปฐมํ ฉิชฺชติ.} \\ \csromangathabat{นาสนนฺติโก กถินุทฺธาโร,}{วุตฺโต อาทิจฺจพนฺธุนา.} \\ \csromangathabat{เอตญฺจ ตาหํ วิสฺสชฺชิสฺสํ,}{อาวาสปลิโพโธ ปฐมํ ฉิชฺชติ.}}
\csromangathasetleft{นาสนนฺติโก กถินุทฺธาโร, \\ เอตญฺจ ตาหํ ปุจฺฉามิ, \\ นาสนนฺติโก กถินุทฺธาโร, \\ เอตญฺจ ตาหํ วิสฺสชฺชิสฺสํ,}
\csromangathagroup{\csromangathastanza{\csromanfolio{311}\csromangathabat{นาสนนฺติโก กถินุทฺธาโร,}{วุตฺโต อาทิจฺจพนฺธุนา.} \\ \csromangathabat{เอตญฺจ ตาหํ ปุจฺฉามิ,}{กตโม ปลิโพโธ ปฐมํ ฉิชฺชติ.}}\csromangathastanza{\csromangathabat{นาสนนฺติโก กถินุทฺธาโร,}{วุตฺโต อาทิจฺจพนฺธุนา.} \\ \csromangathabat{เอตญฺจ ตาหํ วิสฺสชฺชิสฺสํ,}{อาวาสปลิโพโธ ปฐมํ ฉิชฺชติ.}}}

\prose{จีวเร นฏฺเฐ จีวร\-ปลิโพโธ ฉิชฺชติ. (4)}

\setlength{\csromangathapagewidth}{0pt}
\setlength{\csromangathawakwidth}{0pt}
\csromangathameasuregroup{สวนนฺติโก กถินุทฺธาโร, \\ เอตญฺจ ตาหํ ปุจฺฉามิ, \\ สวนนฺติโก กถินุทฺธาโร, \\ เอตญฺจ ตาหํ วิสฺสชฺชิสฺสํ,}{\csromangathabat{สวนนฺติโก กถินุทฺธาโร,}{วุตฺโต อาทิจฺจพนฺธุนา.} \\ \csromangathabat{เอตญฺจ ตาหํ ปุจฺฉามิ,}{กตโม ปลิโพโธ ปฐมํ ฉิชฺชติ.} \\ \csromangathabat{สวนนฺติโก กถินุทฺธาโร,}{วุตฺโต อาทิจฺจพนฺธุนา.} \\ \csromangathabat{เอตญฺจ ตาหํ วิสฺสชฺชิสฺสํ,}{จีวรปลิโพโธ ปฐมํ ฉิชฺชติ.}}
\csromangathasetleft{สวนนฺติโก กถินุทฺธาโร, \\ เอตญฺจ ตาหํ ปุจฺฉามิ, \\ สวนนฺติโก กถินุทฺธาโร, \\ เอตญฺจ ตาหํ วิสฺสชฺชิสฺสํ,}
\csromangathagroup{\csromangathastanza{\csromangathabat{สวนนฺติโก กถินุทฺธาโร,}{วุตฺโต อาทิจฺจพนฺธุนา.} \\ \csromangathabat{เอตญฺจ ตาหํ ปุจฺฉามิ,}{กตโม ปลิโพโธ ปฐมํ ฉิชฺชติ.}}\csromangathastanza{\csromangathabat{สวนนฺติโก กถินุทฺธาโร,}{วุตฺโต อาทิจฺจพนฺธุนา.} \\ \csromangathabat{เอตญฺจ ตาหํ วิสฺสชฺชิสฺสํ,}{จีวร\-ปลิโพโธ ปฐมํ ฉิชฺชติ.}}}

\prose{ตสฺส สห สวเนน อาวาสปลิโพโธ ฉิชฺชติ. (5)}

\setlength{\csromangathapagewidth}{0pt}
\setlength{\csromangathawakwidth}{0pt}
\csromangathameasuregroup{อาสาวจฺเฉทิโก กถินุทฺธาโร, \\ เอตญฺจ ตาหํ ปุจฺฉามิ, \\ อาสาวจฺเฉทิโก กถินุทฺธาโร, \\ เอตญฺจ ตาหํ วิสฺสชฺชิสฺสํ,}{\csromangathabat{อาสาวจฺเฉทิโก กถินุทฺธาโร,}{วุตฺโต อาทิจฺจพนฺธุนา.} \\ \csromangathabat{เอตญฺจ ตาหํ ปุจฺฉามิ,}{กตโม ปลิโพโธ ปฐมํ ฉิชฺชติ.} \\ \csromangathabat{อาสาวจฺเฉทิโก กถินุทฺธาโร,}{วุตฺโต อาทิจฺจพนฺธุนา.} \\ \csromangathabat{เอตญฺจ ตาหํ วิสฺสชฺชิสฺสํ,}{อาวาสปลิโพโธ ปฐมํ ฉิชฺชติ.}}
\csromangathasetleft{อาสาวจฺเฉทิโก กถินุทฺธาโร, \\ เอตญฺจ ตาหํ ปุจฺฉามิ, \\ อาสาวจฺเฉทิโก กถินุทฺธาโร, \\ เอตญฺจ ตาหํ วิสฺสชฺชิสฺสํ,}
\csromangathagroup{\csromangathastanza{\csromangathabat{อาสาวจฺเฉทิโก กถินุทฺธาโร,}{วุตฺโต อาทิจฺจพนฺธุนา.} \\ \csromangathabat{เอตญฺจ ตาหํ ปุจฺฉามิ,}{กตโม ปลิโพโธ ปฐมํ ฉิชฺชติ.}}\csromangathastanza{\csromangathabat{อาสาวจฺเฉทิโก กถินุทฺธาโร,}{วุตฺโต อาทิจฺจพนฺธุนา.} \\ \csromangathabat{เอตญฺจ ตาหํ วิสฺสชฺชิสฺสํ,}{อาวาสปลิโพโธ ปฐมํ ฉิชฺชติ.}}}

\prose{จีวราสาย อุปจฺฉินฺนาย จีวร\-ปลิโพโธ ฉิชฺชติ. (6)}

\setlength{\csromangathapagewidth}{0pt}
\setlength{\csromangathawakwidth}{0pt}
\csromangathameasuregroup{สีมาติกฺกมนนฺติโก กถินุทฺธาโร, \\ เอตญฺจ ตาหํ ปุจฺฉามิ, \\ สีมาติกฺกมนนฺติโก กถินุทฺธาโร, \\ เอตญฺจ ตาหํ วิสฺสชฺชิสฺสํ,}{\csromangathabat{สีมาติกฺกมนนฺติโก กถินุทฺธาโร,}{วุตฺโต อาทิจฺจพนฺธุนา.} \\ \csromangathabat{เอตญฺจ ตาหํ ปุจฺฉามิ,}{กตโม ปลิโพโธ ปฐมํ ฉิชฺชติ.} \\ \csromangathabat{สีมาติกฺกมนนฺติโก กถินุทฺธาโร,}{วุตฺโต อาทิจฺจพนฺธุนา.} \\ \csromangathabat{เอตญฺจ ตาหํ วิสฺสชฺชิสฺสํ,}{จีวรปลิโพโธ ปฐมํ ฉิชฺชติ.}}
\csromangathasetleft{สีมาติกฺกมนนฺติโก กถินุทฺธาโร, \\ เอตญฺจ ตาหํ ปุจฺฉามิ, \\ สีมาติกฺกมนนฺติโก กถินุทฺธาโร, \\ เอตญฺจ ตาหํ วิสฺสชฺชิสฺสํ,}
\csromangathagroup{\csromangathastanza{\csromangathabat{สีมา\-ติกฺกมน\-นฺติโก กถินุทฺธาโร,}{วุตฺโต อาทิจฺจพนฺธุนา.} \\ \csromangathabat{เอตญฺจ ตาหํ ปุจฺฉามิ,}{กตโม ปลิโพโธ ปฐมํ ฉิชฺชติ.}}\csromangathastanza{\csromangathabat{สีมา\-ติกฺกมน\-นฺติโก กถินุทฺธาโร,}{วุตฺโต อาทิจฺจพนฺธุนา.} \\ \csromangathabat{เอตญฺจ ตาหํ วิสฺสชฺชิสฺสํ,}{จีวร\-ปลิโพโธ ปฐมํ ฉิชฺชติ.}}}

\prose{ตสฺส พหิสีเม\csromansharedfootnote{ee0fec7096726e0f}{พหิสีมคตสฺส (สี, สฺยา)} อาวาสปลิโพโธ ฉิชฺชติ. (7)}

\setlength{\csromangathapagewidth}{0pt}
\setlength{\csromangathawakwidth}{0pt}
\csromangathameasuregroup{สหุพฺภาโร\textsuperscript{1} กถินุทฺธาโร, \\ เอตญฺจ ตาหํ ปุจฺฉามิ, \\ สหุพฺภาโร กถินุทฺธาโร,}{\csromangathabat{สหุพฺภาโร\textsuperscript{1} กถินุทฺธาโร,}{วุตฺโต อาทิจฺจพนฺธุนา.} \\ \csromangathabat{เอตญฺจ ตาหํ ปุจฺฉามิ,}{กตโม ปลิโพโธ ปฐมํ ฉิชฺชติ.} \\ \csromangathabat{สหุพฺภาโร กถินุทฺธาโร,}{วุตฺโต อาทิจฺจพนฺธุนา.}}
\csromangathasetleft{สหุพฺภาโร\textsuperscript{1} กถินุทฺธาโร, \\ เอตญฺจ ตาหํ ปุจฺฉามิ, \\ สหุพฺภาโร กถินุทฺธาโร,}
\csromangathagroup{\csromangathastanza{\csromangathabat{สหุพฺภาโร\csromansharedfootnote{8a0bf473028d6afc}{สอุพฺภาโร (ก)} กถินุทฺธาโร,}{วุตฺโต อาทิจฺจพนฺธุนา.} \\ \csromangathabat{เอตญฺจ ตาหํ ปุจฺฉามิ,}{กตโม ปลิโพโธ ปฐมํ ฉิชฺชติ.} \\ \csromangathabat{สหุพฺภาโร กถินุทฺธาโร,}{วุตฺโต อาทิจฺจพนฺธุนา.}}}

\prose{เอตญฺจ ตาหํ วิสฺสชฺชิสฺสํ,}

\prose{เทฺว ปลิโพธา อปุพฺพํ อจริมํ ฉิชฺชนฺตีติ. (8)}

\proseitem{416}{กติ กถินุทฺธารา สํฆาธีนา. กติ กถินุทฺธารา ปุคฺคลาธีนา. กติ กถินุทฺธารา เนว สํฆาธีนา น ปุคฺคลาธีนา. เอโก กถินุทฺธาโร \csromanfolio{312}สํฆาธีโน, อนฺตรุพฺภาโร. จตฺตาโร กถินุทฺธารา ปุคฺคลาธีนา ปกฺกมนนฺติโก นิฏฺฐานนฺติโก สนฺนิฏฺฐานนฺติโก สีมา\-ติกฺกมน\-นฺติโก. จตฺตาโร กถินุทฺธารา เนว สํฆาธีนา น ปุคฺคลาธีนา, นาสนนฺติโก สวนนฺติโก อาสาวจฺเฉทิโก สหุพฺภาโร. กติ กถินุทฺธารา อนฺโตสีมาย อุทฺธริยฺยนฺติ. กติ กถินุทฺธารา พหิสีมาย อุทฺธริยฺยนฺติ. กติ กถินุทฺธารา สิยา อนฺโตสีมาย อุทฺธริยฺยนฺติ, สิยา พหิสีมาย อุทฺธริยฺยนฺติ. เทฺว กถินุทฺธารา อนฺโตสีมาย อุทฺธริยฺยนฺติ อนฺตรุพฺภาโร สหุพฺภาโร. ตโย กถินุทฺธารา พหิสีมาย อุทฺธริยฺยนฺติ ปกฺกมนนฺติโก สวนนฺติโก สีมา\-ติกฺกมน\-นฺติโก. จตฺตาโร กถินุทฺธารา สิยา อนฺโตสีมาย อุทฺธริยฺยนฺติ สิยา พหิสีมาย อุทฺธริยฺยนฺติ, นิฏฺฐานนฺติโก สนฺนิฏฺฐานนฺติโก นาสนนฺติโก อาสาวจฺเฉทิโก.}

\prosewithcloser{กติ กถินุทฺธารา เอกุปฺปาทา เอกนิโรธา. กติ กถินุทฺธารา เอกุปฺปาทา นานานิโรธา. เทฺว กถินุทฺธารา เอกุปฺปาทา เอกนิโรธา อนฺตรุพฺภาโร สหุพฺภาโร. อวเสสา กถินุทฺธารา เอกุปฺปาทา นานานิโรธาติ.}{กถินเภโท นิฏฺฐิโต.}
\csromancenter{ตสฺสุทฺทานํ.}
\setlength{\csromangathapagewidth}{0pt}
\setlength{\csromangathawakwidth}{0pt}
\csromangathameasuregroup{กสฺส กินฺติ ปนฺนรส, \\ ปจฺจยสงฺคหมูลา, \\ ติณฺณํ\textsuperscript{1} ตโย ชานิตพฺพํ, \\ ปลิโพธาธินา สีมาย,}{\csromangathabat{กสฺส กินฺติ ปนฺนรส,}{ธมฺมา นิทานเหตุ จ.} \\ \csromangathabat{ปจฺจยสงฺคหมูลา,}{อาทิ จ อตฺถารปุคฺคลา\textsuperscript{1}.} \\ \csromangathabat{ติณฺณํ\textsuperscript{1} ตโย ชานิตพฺพํ,}{อตฺถารํ อุทฺเทเสน จ.} \\ \csromangathabat{ปลิโพธาธินา สีมาย,}{อุปฺปาทนิโรเธน จาติ\textsuperscript{1}.}}
\csromangathasetleft{กสฺส กินฺติ ปนฺนรส, \\ ปจฺจยสงฺคหมูลา, \\ ติณฺณํ\textsuperscript{1} ตโย ชานิตพฺพํ, \\ ปลิโพธาธินา สีมาย,}
\csromangathagroup{\csromangathastanza{\csromangathabat{กสฺส กินฺติ ปนฺนรส,}{ธมฺมา นิทานเหตุ จ.} \\ \csromangathabat{ปจฺจยสงฺคหมูลา,}{อาทิ จ อตฺถารปุคฺคลา\csromansharedfootnote{71444bd7ef86d2dc}{อฏฺฐปุคฺคลา (สี)}.}}\csromangathastanza{\csromangathabat{ติณฺณํ\csromansharedfootnote{44a169ce5883ee2d}{เภทโต ติณฺณํ (ก)} ตโย ชานิตพฺพํ,}{อตฺถารํ อุทฺเทเสน จ.} \\ \csromangathabat{ปลิโพธาธินา สีมาย,}{อุปฺปาทนิโรเธน จาติ\csromansharedfootnote{44a169ce5883ee2d}{เภทโต ติณฺณํ (ก)}.}}}

\csromanmatikaanchor{mtk.219}
\csromantocmarknopage{section}{อุปาลิปญฺจก}{mtk.219}
\bookmark[dest={mtk.219},level=1]{อุปาลิปญฺจก}
\csromanheader{1}{\textbf{อุปาลิปญฺจก}}
\csromanmatikaanchor{mtk.220}
\csromantocmarkref{subparagraph}{1. อนิสฺสิตวคฺค}{mtk.220}
\bookmark[dest={mtk.220},level=5]{1. อนิสฺสิตวคฺค}
\csromanheader{5}{1. \textbf{อนิสฺสิตวคฺค}}
\proseitem{417}{\csromanfolio{313}เตน สมเยน พุทฺโธ ภควา สาวตฺถิยํ วิหรติ เชตวเน อนาถ\-ปิณฺฑิกสฺส อาราเม. อถ โข อายสฺมา อุปาลิ เยน ภควา เตนุปสงฺกมิ, อุปสงฺกมิตฺวา ภควนฺตํ อภิวาเทตฺวา เอกมนฺตํ นิสีทิ, เอกมนฺตํ นิสินฺโน โข อายสฺมา อุปาลิ ภควนฺตํ เอตทโวจ “กติหิ นุ โข ภนฺเต องฺเคหิ สมนฺนาคเตน ภิกฺขุนา ยาวชีวํ นานิสฺสิเตน วตฺถพฺพนฺ”ติ.}

\prose{ปญฺจหุปาลิ องฺเคหิ สมนฺนาคเตน ภิกฺขุนา ยาวชีวํ นานิสฺสิเตน วตฺถพฺพํ.\csromansymbolfootnote{*}{วิ 5. 234 ปิฏฺเฐปิ.}กตเมหิ ปญฺจหิ, อุโปสถํ น ชานาติ, อุโปสถกมฺมํ น ชานาติ, ปาติโมกฺขํ น ชานาติ, ปาติโมกฺขุ\-ทฺเทสํ น ชานาติ, อูนปญฺจวสฺโส โหติ. อิเมหิ โข อุปาลิ ปญฺจหงฺเคหิ สมนฺนาคเตน ภิกฺขุนา ยาวชีวํ นานิสฺสิเตน วตฺถพฺพํ. ปญฺจหุปาลิ องฺเคหิ สมนฺนาคเตน ภิกฺขุนา ยาวชีวํ อนิสฺสิเตน วตฺถพฺพํ. กตเมหิ ปญฺจหิ, อุโปสถํ ชานาติ, อุโปสถกมฺมํ ชานาติ, ปาติโมกฺขํ ชานาติ, ปาติโมกฺขุ\-ทฺเทสํ ชานาติ, ปญฺจวสฺโส วา โหติ อติเรก\-ปญฺจวสฺโส วา. อิเมหิ โข อุปาลิ ปญฺจหงฺเคหิ สมนฺนาคเตน ภิกฺขุนา ยาวชีวํ อนิสฺสิเตน วตฺถพฺพํ.}

\prose{อปเรหิปิ อุปาลิ ปญฺจหงฺเคหิ สมนฺนาคเตน ภิกฺขุนา ยาวชีวํ นานิสฺสิเตน วตฺถพฺพํ. กตเมหิ ปญฺจหิ, ปวารณํ น ชานาติ, ปวารณากมฺมํ น ชานาติ, ปาติโมกฺขํ น ชานาติ, ปาติโมกฺขุ\-ทฺเทสํ น ชานาติ, อูนปญฺจวสฺโส โหติ. อิเมหิ โข อุปาลิ ปญฺจหงฺเคหิ สมนฺนาคเตน ภิกฺขุนา ยาวชีวํ นานิสฺสิเตน วตฺถพฺพํ. ปญฺจหุปาลิ องฺเคหิ สมนฺนาคเตน ภิกฺขุนา ยาวชีวํ อนิสฺสิเตน วตฺถพฺพํ. กตเมหิ ปญฺจหิ, ปวารณํ ชานาติ, ปวารณากมฺมํ ชานาติ, ปาติโมกฺขํ ชานาติ, ปาติโมกฺขุ\-ทฺเทสํ ชานาติ, ปญฺจวสฺโส วา โหติ อติเรก\-ปญฺจวสฺโส วา. อิเมหิ โข อุปาลิ ปญฺจหงฺเคหิ สมนฺนาคเตน ภิกฺขุนา ยาวชีวํ อนิสฺสิเตน วตฺถพฺพํ.}

\prose{อปเรหิปิ อุปาลิ ปญฺจหงฺเคหิ สมนฺนาคเตน ภิกฺขุนา ยาวชีวํ นานิสฺสิเตน วตฺถพฺพํ. กตเมหิ ปญฺจหิ, อาปตฺตานาปตฺติํ น ชานาติ, \csromanfolio{314}ลหุกครุกํ อาปตฺติํ น ชานาติ, สาวเสสา\-นวเสสํ อาปตฺติํ น ชานาติ, ทุฏฺฐุลฺลา\-ทุฏฺฐุลฺลํ อาปตฺติํ น ชานาติ, อูนปญฺจวสฺโส โหติ. อิเมหิ โข อุปาลิ ปญฺจหงฺเคหิ สมนฺนาคเตน ภิกฺขุนา ยาวชีวํ นานิสฺสิเตน วตฺถพฺพํ. ปญฺจหุปาลิ องฺเคหิ สมนฺนาคเตน ภิกฺขุนา ยาวชีวํ อนิสฺสิเตน วตฺถพฺพํ. กตเมหิ ปญฺจหิ, อาปตฺตานาปตฺติํ ชานาติ, ลหุกครุกํ อาปตฺติํ ชานาติ, สาวเสสา\-นวเสสํ อาปตฺติํ ชานาติ, ทุฏฺฐุลฺลา\-ทุฏฺฐุลฺลํ อาปตฺติํ ชานาติ, ปญฺจวสฺโส วา โหติ อติเรก\-ปญฺจวสฺโส วา. อิเมหิ โข อุปาลิ ปญฺจหงฺเคหิ สมนฺนาคเตน ภิกฺขุนา ยาวชีวํ อนิสฺสิเตน วตฺถพฺพํ.}

\proseitem{418}{กติหิ นุ โข ภนฺเต องฺเคหิ สมนฺนาคเตน ภิกฺขุนา น อุปสมฺ\-ปาเท\-ตพฺพํ, น นิสฺสโย ทาตพฺโพ, น สามเณโร อุปฏฺฐาเป\-ตพฺโพติ.}

\prose{\csromansymbolmark{*}ปญฺจหุปาลิ องฺเคหิ สมนฺนาคเตน ภิกฺขุนา น อุปสมฺ\-ปาเท\-ตพฺพํ, น นิสฺสโย ทาตพฺโพ, น สามเณโร อุปฏฺฐาเปตพฺโพ. กตเมหิ ปญฺจหิ, น ปฏิพโล โหติ อนฺเตวาสิํ วา สทฺธิวิหาริํ วา คิลานํ อุปฏฺฐาตุํ วา อุปฏฺฐาเปตุํ วา, อนภิรตํ วูปกาเสตุํ วา วูปกาสาเปตุํ วา, อุปฺปนฺนํ กุกฺกุจฺจํ ธมฺมโต วิโนเทตุํ1, อภิธมฺเม วิเนตุํ, อภิวินเย วิเนตุํ. อิเมหิ โข อุปาลิ ปญฺจหงฺเคหิ สมนฺนาคเตน ภิกฺขุนา น อุปสมฺ\-ปาเท\-ตพฺพํ, น นิสฺสโย ทาตพฺโพ, น สามเณโร อุปฏฺฐาเปตพฺโพ. ปญฺจหุปาลิ องฺเคหิ สมนฺนาคเตน ภิกฺขุนา อุปสมฺ\-ปาเท\-ตพฺพํ, นิสฺสโย ทาตพฺโพ, สามเณโร อุปฏฺฐาเปตพฺโพ. กตเมหิ ปญฺจหิ, ปฏิพโล โหติ อนฺเตวาสิํ วา สทฺธิวิหาริํ วา คิลานํ อุปฏฺฐาตุํ วา อุปฏฺฐาเปตุํ วา, อนภิรตํ วูปกาเสตุํ วา วูปกาสาเปตุํ วา, อุปฺปนฺนํ กุกฺกุจฺจํ ธมฺมโต วิโนเทตุํ, อภิธมฺเม วิเนตุํ, อภิวินเย วิเนตุํ. อิเมหิ โข อุปาลิ ปญฺจหงฺเคหิ สมนฺนาคเตน ภิกฺขุนา อุปสมฺ\-ปาเท\-ตพฺพํ, นิสฺสโย ทาตพฺโพ, สามเณโร อุปฏฺฐาเปตพฺโพ.}

\prose{อปเรหิปิ อุปาลิ ปญฺจหงฺเคหิ สมนฺนาคเตน ภิกฺขุนา น อุปสมฺ\-ปาเท\-ตพฺพํ, น นิสฺสโย ทาตพฺโพ, น สามเณโร อุปฏฺฐาเปตพฺโพ. กตเมหิ}

\prose{\csromansymbolmark{*}วิ 3. 92 ปิฏฺเฐปิ.}

\proseitem{1}{อิมสฺมิํ ฐาเน สพฺพตฺถ “วิโนเทตุํ วา วิโนทาเปตุํ วา”ติ ปาโฐ ทิสฺสติ,}

\prose{วิมติ\-วิโนทนี\-ฏีกาย มหาวคฺค\-วณฺณนา โอโลเกตพฺพา. \csromanfolio{315}ปญฺจหิ, น ปฏิพโล โหติ อนฺเตวาสิํ วา สทฺธิวิหาริํ วา อภิสมาจาริกาย สิกฺขาย สิกฺขาเปตุํ, อาทิพฺรหฺมจาริยกาย สิกฺขาย วิเนตุํ, อธิสีเล วิเนตุํ, อธิจิตฺเต วิเนตุํ, อธิปญฺญาย วิเนตุํ. อิเมหิ โข อุปาลิ ปญฺจหงฺเคหิ สมนฺนาคเตน ภิกฺขุนา น อุปสมฺ\-ปาเท\-ตพฺพํ, น นิสฺสโย ทาตพฺโพ, น สามเณโร อุปฏฺฐาเปตพฺโพ. ปญฺจหุปาลิ องฺเคหิ สมนฺนาคเตน ภิกฺขุนา อุปสมฺ\-ปาเท\-ตพฺพํ, นิสฺสโย ทาตพฺโพ, สามเณโร อุปฏฺฐาเปตพฺโพ. กตเมหิ ปญฺจหิ, ปฏิพโล โหติ อนฺเตวาสิํ วา สทฺธิวิหาริํ วา อภิสมาจาริกาย สิกฺขาย สิกฺขาเปตุํ, อาทิพฺรหฺมจาริยกาย สิกฺขาย วิเนตุํ, อธิสีเล วิเนตุํ, อธิจิตฺเต วิเนตุํ, อธิปญฺญาย วิเนตุํ. อิเมหิ โข อุปาลิ ปญฺจหงฺเคหิ สมนฺนาคเตน ภิกฺขุนา อุปสมฺ\-ปาเท\-ตพฺพํ, นิสฺสโย ทาตพฺโพ, สามเณโร อุปฏฺฐาเป\-ตพฺโพติ.}

\proseitem{419}{กติหิ นุ โข ภนฺเต องฺเคหิ สมนฺนาคตสฺส ภิกฺขุโน กมฺมํ กาตพฺพนฺติ.}

\prose{ปญฺจหุปาลิ องฺเคหิ สมนฺนาคตสฺส ภิกฺขุโน กมฺมํ กาตพฺพํ. กตเมหิ ปญฺจหิ, อลชฺชี จ โหติ พาโล จ อปกตตฺโต จ มิจฺฉาทิฏฺฐิโก จ โหติ อาชีววิปนฺโน จ. อิเมหิ โข อุปาลิ ปญฺจหงฺเคหิ สมนฺนาคตสฺส ภิกฺขุโน กมฺมํ กาตพฺพํ.}

\prose{อปเรหิปิ อุปาลิ ปญฺจหงฺเคหิ สมนฺนาคตสฺส ภิกฺขุโน กมฺมํ กาตพฺพํ. กตเมหิ ปญฺจหิ, อธิสีเล สีลวิปนฺโน โหติ, อชฺฌาจาเร อาจารวิปนฺโน โหติ, อติทิฏฺฐิยา ทิฏฺฐิวิปนฺโน โหติ, มิจฺฉาทิฏฺฐิโก จ โหติ อาชีววิปนฺโน จ. อิเมหิ โข อุปาลิ ปญฺจหงฺเคหิ สมนฺนาคตสฺส ภิกฺขุโน กมฺมํ กาตพฺพํ.}

\prose{อปเรหิปิ อุปาลิ ปญฺจหงฺเคหิ สมนฺนาคตสฺส ภิกฺขุโน กมฺมํ กาตพฺพํ. กตเมหิ ปญฺจหิ, กายิเกน ทเวน สมนฺนาคโต โหติ, วาจสิเกน ทเวน สมนฺนาคโต โหติ, กายิก\-วาจสิเกน ทเวน สมนฺนาคโต โหติ, มิจฺฉาทิฏฺฐิโก จ โหติ, อาชีววิปนฺโน จ. อิเมหิ โข อุปาลิ ปญฺจหงฺเคหิ สมนฺนาคตสฺส ภิกฺขุโน กมฺมํ กาตพฺพํ.}

\prose{\csromanfolio{316}อปเรหิปิ อุปาลิ ปญฺจหงฺเคหิ สมนฺนาคตสฺส ภิกฺขุโน กมฺมํ กาตพฺพํ. กตเมหิ ปญฺจหิ, กายิเกน อนาจาเรน สมนฺนาคโต โหติ, วาจสิเกน อนาจาเรน สมนฺนาคโต โหติ, กายิก\-วาจสิเกน อนาจาเรน สมนฺนาคโต โหติ, มิจฺฉาทิฏฺฐิโก จ โหติ, อาชีววิปนฺโน จ. อิเมหิ โข อุปาลิ ปญฺจหงฺเคหิ สมนฺนาคตสฺส ภิกฺขุโน กมฺมํ กาตพฺพํ.}

\prose{อปเรหิปิ อุปาลิ ปญฺจหงฺเคหิ สมนฺนาคตสฺส ภิกฺขุโน กมฺมํ กาตพฺพํ. กตเมหิ ปญฺจหิ, กายิเกน อุปฆาติเกน สมนฺนาคโต โหติ, วาจสิเกน อุปฆาติเกน สมนฺนาคโต โหติ, กายิก\-วาจสิเกน อุปฆาติเกน สมนฺนาคโต โหติ, มิจฺฉาทิฏฺฐิโก จ โหติ, อาชีววิปนฺโน จ. อิเมหิ โข อุปาลิ ปญฺจหงฺเคหิ สมนฺนาคตสฺส ภิกฺขุโน กมฺมํ กาตพฺพํ.}

\prose{อปเรหิปิ อุปาลิ ปญฺจหงฺเคหิ สมนฺนาคตสฺส ภิกฺขุโน กมฺมํ กาตพฺพํ. กตเมหิ ปญฺจหิ, กายิเกน มิจฺฉาชีเวน สมนฺนาคโต โหติ, วาจสิเกน มิจฺฉาชีเวน สมนฺนาคโต โหติ, กายิก\-วาจสิเกน มิจฺฉาชีเวน สมนฺนาคโต โหติ, มิจฺฉาทิฏฺฐิโก จ โหติ, อาชีววิปนฺโน จ. อิเมหิ โข อุปาลิ ปญฺจหงฺเคหิ สมนฺนาคตสฺส ภิกฺขุโน กมฺมํ กาตพฺพํ.}

\prose{อปเรหิปิ อุปาลิ ปญฺจหงฺเคหิ สมนฺนาคตสฺส ภิกฺขุโน กมฺมํ กาตพฺพํ. กตเมหิ ปญฺจหิ, อาปตฺติํ อาปนฺโน กมฺมกโต อุปสมฺปาเทติ, นิสฺสยํ เทติ, สามเณรํ อุปฏฺฐาเปติ, ภิกฺขุโนวาทก\-สมฺมุติํ สาทิยติ, สมฺมโตปิ ภิกฺขุนิโย โอวทติ. อิเมหิ โข อุปาลิ ปญฺจหงฺเคหิ สมนฺนาคตสฺส ภิกฺขุโน กมฺมํ กาตพฺพํ.}

\prose{อปเรหิปิ อุปาลิ ปญฺจหงฺเคหิ สมนฺนาคตสฺส ภิกฺขุโน กมฺมํ กาตพฺพํ. กตเมหิ ปญฺจหิ, ยาย อาปตฺติยา สํเฆน กมฺมํ กตํ โหติ ตํ อาปตฺติํ อาปชฺชติ, อญฺญํ วา ตาทิสิกํ, ตโต วา ปาปิฏฺฐตรํ, กมฺมํ ครหติ, กมฺมิเก ครหติ. อิเมหิ โข อุปาลิ ปญฺจหงฺเคหิ สมนฺนาคตสฺส ภิกฺขุโน กมฺมํ กาตพฺพํ.}

\prose{\csromanfolio{317}อปเรหิปิ อุปาลิ ปญฺจหงฺเคหิ สมนฺนาคตสฺส ภิกฺขุโน กมฺมํ กาตพฺพํ. กตเมหิ ปญฺจหิ, พุทฺธสฺส อวณฺณํ ภาสติ, ธมฺมสฺส อวณฺณํ ภาสติ, สํฆสฺส อวณฺณํ ภาสติ, มิจฺฉาทิฏฺฐิโก จ โหติ, อาชีววิปนฺโน จ. อิเมหิ โข อุปาลิ ปญฺจหงฺเคหิ สมนฺนาคตสฺส ภิกฺขุโน กมฺมํ กาตพฺพนฺติ.}

\vspace{\tipitakaparskip}
\csromancenter{อนิสฺสิตวคฺโค นิฏฺฐิโต ปฐโม.}
\csromancenter{ตสฺสุทฺทานํ.}
\setlength{\csromangathapagewidth}{0pt}
\setlength{\csromangathawakwidth}{0pt}
\csromangathameasuregroup{อุโปสถํ ปวารณํ, \\ อภิสมาจารลชฺชี จ, \\ อนาจารํ อุปฆาติ,}{\csromangathabat{อุโปสถํ ปวารณํ,}{อาปตฺติ จ คิลานกํ.} \\ \csromangathabat{อภิสมาจารลชฺชี จ,}{อธิสีเล ทเวน จ.} \\ \csromangathabat{อนาจารํ อุปฆาติ,}{มิจฺฉา อาปตฺติเมว จ.}}
\csromangathasetleft{อุโปสถํ ปวารณํ, \\ อภิสมาจารลชฺชี จ, \\ อนาจารํ อุปฆาติ,}
\csromangathagroup{\csromangathastanza{\csromangathabat{อุโปสถํ ปวารณํ,}{อาปตฺติ จ คิลานกํ.} \\ \csromangathabat{อภิสมาจารลชฺชี จ,}{อธิสีเล ทเวน จ.} \\ \csromangathabat{อนาจารํ อุปฆาติ,}{มิจฺฉา อาปตฺติเมว จ.}}}

\vspace{\tipitakaparskip}
\csromancenter{ยายาปตฺติยา พุทฺธสฺส, ปฐโม วคฺคสงฺคโหติ.}
\csromansectionrule
\csromanmatikaanchor{mtk.221}
\csromantocmarkref{subparagraph}{2. นปฺปฏิปฺปสฺสมฺภนวคฺค}{mtk.221}
\bookmark[dest={mtk.221},level=5]{2. นปฺปฏิปฺปสฺสมฺภนวคฺค}
\csromanheader{5}{2. นปฺปฏิปฺปสฺสมฺภน\-วคฺค}
\proseitem{420}{กติหิ นุ โข ภนฺเต องฺเคหิ สมนฺนาคตสฺส ภิกฺขุโน กมฺมํ นปฺปฏิปฺปสฺสมฺเภตพฺพนฺติ.}

\prose{ปญฺจหุปาลิ องฺเคหิ สมนฺนาคตสฺส ภิกฺขุโน กมฺมํ นปฺปฏิปฺปสฺสมฺเภตพฺพํ. กตเมหิ ปญฺจหิ, อาปตฺติํ อาปนฺโน กมฺมกโต อุปสมฺปาเทติ, นิสฺสยํ เทติ, สามเณรํ อุปฏฺฐาเปติ, ภิกฺขุโนวาทก\-สมฺมุติํ สาทิยติ, สมฺมโตปิ ภิกฺขุนิโย โอวทติ. อิเมหิ โข อุปาลิ ปญฺจหงฺเคหิ สมนฺนาคตสฺส ภิกฺขุโน กมฺมํ นปฺปฏิปฺปสฺสมฺเภตพฺพํ.}

\prose{\csromansymbolmark{*}อปเรหิปิ อุปาลิ ปญฺจหงฺเคหิ สมนฺนาคตสฺส ภิกฺขุโน กมฺมํ นปฺปฏิปฺปสฺสมฺเภตพฺพํ. กตเมหิ ปญฺจหิ, ยาย อาปตฺติยา สํเฆน กมฺมํ กตํ โหติ ตํ อาปตฺติํ อาปชฺชติ, อญฺญํ วา ตาทิสิกํ, ตโต วา ปาปิฏฺฐตรํ, กมฺมํ ครหติ, กมฺมิเก ครหติ. อิเมหิ โข อุปาลิ ปญฺจหงฺเคหิ สมนฺนาคตสฺส ภิกฺขุโน กมฺมํ นปฺปฏิปฺปสฺสมฺเภตพฺพํ.}

\prose{อปเรหิปิ อุปาลิ ปญฺจหงฺเคหิ สมนฺนาคตสฺส ภิกฺขุโน กมฺมํ นปฺปฏิปฺปสฺสมฺเภตพฺพํ. กตเมหิ ปญฺจหิ, พุทฺธสฺส อวณฺณํ ภาสติ, ธมฺมสฺส อวณฺณํ ภาสติ, สํฆสฺส อวณฺณํ ภาสติ, มิจฺฉาทิฏฺฐิโก จ โหติ,}

\vspace{\tipitakaparskip}
\csromancenter{\csromansymbolmark{*}วิ 4. 11, 20, 35, 45, 54, 65, 76 ปิฏฺเฐสุปิ.}
\prosecont{\csromanfolio{318}อาชีววิปนฺโน จ. อิเมหิ โข อุปาลิ ปญฺจหงฺเคหิ สมนฺนาคตสฺส ภิกฺขุโน กมฺมํ นปฺปฏิปฺปสฺสมฺเภตพฺพํ.}

\prose{อปเรหิปิ อุปาลิ ปญฺจหงฺเคหิ สมนฺนาคตสฺส ภิกฺขุโน กมฺมํ นปฺปฏิปฺปสฺสมฺเภตพฺพํ. กตเมหิ ปญฺจหิ, อลชฺชี จ โหติ, พาโล จ, อปกตตฺโต จ, โอมทฺทการโก จ โหติ, วตฺเตสุ สิกฺขาย จ น ปริปูรการี. อิเมหิ โข อุปาลิ ปญฺจหงฺเคหิ สมนฺนาคตสฺส ภิกฺขุโน กมฺมํ นปฺปฏิปฺปสฺสมฺเภตพฺพนฺติ.}

\proseitem{421}{สงฺคามา\-วจเรน ภนฺเต ภิกฺขุนา สํฆํ อุปส\-งฺกม\-นฺเตน กติ ธมฺเม อชฺฌตฺตํ อุปฏฺฐาเปตฺวา สํโฆ อุปสงฺกมิตพฺโพติ.}

\prose{\csromansymbolmark{*}สงฺคามา\-วจเรน อุปาลิ ภิกฺขุนา สํฆํ อุปส\-งฺกม\-นฺเตน ปญฺจ ธมฺเม อชฺฌตฺตํ อุปฏฺฐาเปตฺวา สํโฆ อุปสงฺกมิตพฺโพ. กตเม ปญฺจ, สงฺคามา\-วจเรน อุปาลิ ภิกฺขุนา สํฆํ อุปส\-งฺกม\-นฺเตน นีจจิตฺเตน สํโฆ อุปสงฺกมิตพฺโพ รโชหรณ\-สเมน จิตฺเตน, อาสนกุสเลน ภวิตพฺพํ นิสชฺช\-กุสเลน, เถเร ภิกฺขู อนุปขชฺชนฺเตน นเว ภิกฺกู อาสเนน อปฺปฏิพาหนฺเตน ยถาปติรูเป อาสเน นิสีทิตพฺพํ, อนานากถิเกน ภวิตพฺพํ อติรจฺฉานกถิเกน, สามํ วา ธมฺโม ภาสิตพฺโพ, ปโร วา อชฺเฌสิตพฺโพ, อริโย วา ตุณฺหิภาโว นาติมญฺญิตพฺโพ, สเจ อุปาลิ สํโฆ สมคฺค\-กรณียานิ กมฺมานิ กโรติ, ตตฺร เจ อุปาลิ ภิกฺขุโน นกฺขมติ, อปิ ทิฏฺฐาวิกมฺมํ กตฺวา ญาเปตพฺพา สามคฺคี, ตํ กิสฺสเหตุ, มาหํ สํเฆน นานตฺโต อสฺสนฺติ. สงฺคามา\-วจเรนุ\-ปาลิ ภิกฺขุนา สํฆํ อุปส\-งฺกม\-นฺเตน อิเม ปญฺจ ธมฺเม อชฺฌตฺตํ อุปฏฺฐาเปตฺวา สํโฆ อุปสงฺกมิตพฺโพติ.}

\proseitem{422}{กติหิ นุ โข ภนฺเต องฺเคหิ สมนฺนาคโต ภิกฺขุ สํเฆ โวหรนฺโต พหุชนอกนฺโต จ โหติ พหุชน\-อมนาโป จ พหุชน\-อรุจิโต จาติ.}

\prose{ปญฺจหุปาลิ องฺเคหิ สมนฺนาคโต ภิกฺขุ สํเฆ โวหรนฺโต พหุชนอกนฺโต จ โหติ พหุชน\-อมนาโป จ พหุชน\-อรุจิโต จ. กตเมหิ ปญฺจหิ, อุสฺสิตมนฺตี จ โหติ, นิสฺสิตชปฺปี จ, น จ}

\vspace{\tipitakaparskip}
\csromancenter{\csromansymbolmark{*}วิ 5. 287 ปิฏฺเฐปิ.}
\prosecont{\csromanfolio{319}ภาสา\-นุสนฺธิ\-กุสโล โหติ, น ยถาธมฺเม ยถาวินเย ยถา ปตฺติยา โจเทตา โหติ, น ยถาธมฺเม ยถาวินเย ยถาปตฺติยา กาเรตา โหติ, อิเมหิ โข อุปาลิ ปญฺจหงฺเคหิ สมนฺนาคโต ภิกฺขุ สํเฆ โวหรนฺโต พหุชนอกนฺโต จ โหติ พหุชน\-อมนาโป จ พหุชน\-อรุจิโต จ. ปญฺจหุปาลิ องฺเคหิ สมนฺนาคโต ภิกฺขุ สํเฆ โวหรนฺโต พหุชนกนฺโต จ โหติ พหุชนมนาโป จ พหุชนรุจิโต จ. กตเมหิ ปญฺจหิ, น อุสฺสิตมนฺตี จ โหติ, น นิสฺสิตชปฺปี จ, ภาสา\-นุสนฺธิ\-กุสโล จ โหติ, ยถาธมฺเม ยถาวินเย ยถา ปตฺติยา โจเทตา โหติ, ยถาธมฺเม ยถาวินเย ยถาปตฺติยา กาเรตา โหติ. อิเมหิ โข อุปาลิ ปญฺจหงฺเคหิ สมนฺนาคโต ภิกฺขุ สํเฆ โวหรนฺโต พหุชนกนฺโต จ โหติ พหุชนมนาโป จ พหุชนรุจิโต จ.}

\prose{อปเรหิปิ อุปาลิ ปญฺจหงฺเคหิ สมนฺนาคโต ภิกฺขุ สํเฆ โวหรนฺโต พหุชนอกนฺโต จ โหติ พหุชน\-อมนาโป จ พหุชน\-อรุจิโต จ. กตเมหิ ปญฺจหิ, อุสฺสาเทตา จ โหติ, อปสาเทตา จ, อธมฺมํ คณฺหาติ, ธมฺมํ ปฏิพาหติ, สมฺผญฺจ พหุํ ภาสติ. อิเมหิ โข อุปาลิ ปญฺจหงฺเคหิ สมนฺนาคโต ภิกฺขุ สํเฆ โวหรนฺโต พหุชนอกนฺโต จ โหติ พหุชน\-อมนาโป จ พหุชน\-อรุจิโต จ. ปญฺจหุปาลิ องฺเคหิ สมนฺนาคโต ภิกฺขุ สํเฆ โวหรนฺโต พหุชนกนฺโต จ โหติ พหุชนมนาโป จ พหุชนรุจิโต จ. กตเมหิ ปญฺจหิ, น อุสฺสาเทตา จ โหติ, น อปสาเทตา จ, ธมฺมํ คณฺหาติ, อธมฺมํ ปฏิพาหติ, สมฺผญฺจ น พหุํ ภาสติ. อิเมหิ โข อุปาลิ ปญฺจหงฺเคหิ สมนฺนาคโต ภิกฺขุ สํเฆ โวหรนฺโต พหุชนกนฺโต จ โหติ พหุชนมนาโป จ พหุชนรุจิโต จ.}

\csromanmatikaanchor{mtk.222}
\csromantocmarkref{subparagraph}{3. โวหารวคฺค}{mtk.222}
\bookmark[dest={mtk.222},level=5]{3. โวหารวคฺค}
\prose{อปเรหิปิ อุปาลิ ปญฺจหงฺเคหิ สมนฺนาคโต ภิกฺขุ สํเฆ โวหรนฺโต พหุชนอกนฺโต จ โหติ พหุชน\-อมนาโป จ พหุชน\-อรุจิโต จ. กตเมหิ ปญฺจหิ, ปสยฺหปวตฺตา โหติ, อโนกาสกมฺมํ กาเรตฺวา ปวตฺตา โหติ, น ยถาธมฺเม ยถาวินเย ยถาปตฺติยา โจเทตา โหติ, น ยถาธมฺเม ยถาวินเย ยถาปตฺติยา กาเรตา โหติ, น ยถาทิฏฺฐิยา พฺยากตา โหติ. อิเมหิ โข อุปาลิ ปญฺจหงฺเคหิ สมนฺนาคโต ภิกฺขุสํเฆ โวหรนฺโต พหุชนอกนฺโต จ โหติ \csromanfolio{320}พหุชน\-อมนาโป จ พหุชน\-อรุจิโต จ. ปญฺจหุปาลิ องฺเคหิ สมนฺนาคโต ภิกฺขุ สํเฆ โวหรนฺโต พหุชนกนฺโต จ โหติ พหุชนมนาโป จ พหุชนรุจิโต จ. กตเมหิ ปญฺจหิ, น ปสยฺหปวตฺตา โหติ, โอกาสกมฺมํ กาเรตฺวา ปวตฺตา โหติ, ยถาธมฺเม ยถาวินเย ยถาปตฺติยา โจเทตา โหติ, ยถาธมฺเม ยถาวินเย ยถาปตฺติยา กาเรตา โหติ, ยถาทิฏฺฐิยา พฺยากตา โหติ. อิเมหิ โข อุปาลิ ปญฺจหงฺเคหิ สมนฺนาคโต ภิกฺขุ สํเฆ โวหรนฺโต พหุชนกนฺโต จ โหติ พหุชนมนาโป จ พหุชนรุจิโต จาติ.}

\proseitem{423}{กติ นุ โข ภนฺเต อานิสํสา วินย\-ปริยตฺติยาติ.}

\prose{ปญฺจิเม อุปาลิ อานิสํสา วินย\-ปริยตฺติยา. กตเม ปญฺจ, อตฺตโน สีลกฺขนฺโธ สุคุตฺโต โหติ สุรกฺขิโต, กุกฺกุจฺจ\-ปกตานํ ปฏิสรณํ โหติ, วิสารโท สํฆมชฺเฌ โวหรติ, ปจฺจตฺถิเก สหธมฺเมน สุนิคฺคหิตํ นิคฺคณฺหาติ, สทฺธมฺม\-ฏฺฐิติยา ปฏิปนฺโน โหติ. อิเมหิ โข อุปาลิ ปญฺจานิสํสา วินย\-ปริยตฺติยา.}

\vspace{\tipitakaparskip}
\csromancenter{นปฺปฏิปฺปสฺสมฺภน\-วคฺโค นิฏฺฐิโต ทุติโย.}
\csromancenter{ตสฺสุทฺทานํ.}
\setlength{\csromangathapagewidth}{0pt}
\setlength{\csromangathawakwidth}{0pt}
\csromangathameasuregroup{อาปนฺโน ยายวณฺณํ จ, \\ อุสฺสิตา อุสฺสาเทตา จ,}{\csromangathabat{อาปนฺโน ยายวณฺณํ จ,}{อลชฺชี สงฺคาเมน จ.} \\ \csromangathabat{อุสฺสิตา อุสฺสาเทตา จ,}{ปสยฺห ปริยตฺติยาติ.}}
\csromangathasetleft{อาปนฺโน ยายวณฺณํ จ, \\ อุสฺสิตา อุสฺสาเทตา จ,}
\csromangathagroupwithcloserruled{\csromangathastanza{\csromangathabat{อาปนฺโน ยายวณฺณํ จ,}{อลชฺชี สงฺคาเมน จ.} \\ \csromangathabat{อุสฺสิตา อุสฺสาเทตา จ,}{ปสยฺห ปริยตฺติยาติ.}}}{ปฐม\-ยมก\-ปญฺญตฺติ.}
\csromancenter{\textbf{โวหารวคฺค} 424. กติหิ นุ โข ภนฺเต องฺเคหิ สมนฺนาคเตน ภิกฺขุนา สํเฆ น โวหริตพฺพนฺติ.}
\prose{ปญฺจหุปาลิ องฺเคหิ สมนฺนาคเตน ภิกฺขุนา สํเฆ น โวหริตพฺพํ. กตเมหิ ปญฺจหิ, อาปตฺติํ น ชานาติ, อาปตฺติ\-สมุฏฺฐานํ น ชานาติ, อาปตฺติยา ปโยคํ น ชานาติ, อาปตฺติยา วูปสมํ น ชานาติ, อาปตฺติยา น วินิจฺฉย\-กุสโล โหติ. อิเมหิ โข อุปาลิ \csromanfolio{321}ปญฺจหงฺเคหิ สมนฺนาคเตน ภิกฺขุนา สํเฆ น โวหริตพฺพํ. ปญฺจหุปาลิ องฺเคหิ สมนฺนาคเตน ภิกฺขุนา สํเฆ โวหริตพฺพํ. กตเมหิ ปญฺจหิ, อาปตฺติํ ชานาติ, อาปตฺติ\-สมุฏฺฐานํ ชานาติ, อาปตฺติยา ปโยคํ ชานาติ, อาปตฺติยา วูปสมํ ชานาติ, อาปตฺติยา วินิจฺฉย\-กุสโล โหติ. อิเมหิ โข อุปาลิ ปญฺจหงฺเคหิ สมนฺนาคเตน ภิกฺขุนา สํเฆ โวหริตพฺพํ.}

\prose{อปเรหิปิ อุปาลิ ปญฺจหงฺเคหิ สมนฺนาคเตน ภิกฺขุนา สํเฆ น โวหริตพฺพํ. กตเมหิ ปญฺจหิ, อธิกรณํ น ชานาติ, อธิกรณ\-สมุฏฺฐานํ น ชานาติ, อธิกรณสฺส ปโยคํ น ชานาติ, อธิกรณสฺส วูปสมํ น ชานาติ, อธิกรณสฺส น วินิจฺฉย\-กุสโล โหติ. อิเมหิ โข อุปาลิ ปญฺจหงฺเคหิ สมนฺนาคเตน ภิกฺขุนา สํเฆ น โวหริตพฺพํ. ปญฺจหุปาลิ องฺเคหิ สมนฺนาคเตน ภิกฺขุนา สํเฆ โวหริตพฺพํ. กตเมหิ ปญฺจหิ, อธิกรณํ ชานาติ, อธิกรณ\-สมุฏฺฐานํ ชานาติ, อธิกรณสฺส ปโยคํ ชานาติ, อธิกรณ\-สมุฏฺฐานํ ชานาติ, อธิกรณสฺส ปโยคํ ชานาติ, อธิกรณสฺส วูปสมํ ชานาติ, อธิกรณสฺส วินิจฺฉย\-กุสโล โหติ. อิเมหิ โข อุปาลิ ปญฺจหงฺเคหิ สมนฺนาคเตน ภิกฺขุนา สํเฆ โวหริตพฺพํ.}

\prose{อปเรหิปิ อุปาลิ ปญฺจหงฺเคหิ สมนฺนาคเตน ภิกฺขุนา สํเฆ น โวหริตพฺพํ. กตเมหิ ปญฺจหิ, ปสยฺหปวตฺตา โหติ, อโนกาสกมฺมํ กาเรตฺวา ปวตฺตา โหติ, น ยถาธมฺเม ยถาวินเย ยถาปตฺติยา โจเทตา โหติ, น ยถาธมฺเม ยถาวินเย ยถาปตฺติยา กาเรตา โหติ, น ยถาทิฏฺฐิยา พฺยากตา โหติ. อิเมหิ โข อุปาลิ ปญฺจหงฺเคหิ สมนฺนาคเตน ภิกฺขุนา สํเฆ น โวหริตพฺพํ. ปญฺจหุปาลิ องฺเคหิ สมนฺนาคเตน ภิกฺขุนา สํเฆ โวหริตพฺพํ. กตเมหิ ปญฺจหิ, น ปสยฺหปวตฺตา โหติ, โอกาสกมฺมํ กาเรตฺวา ปวตฺตา โหติ, ยถาธมฺเม ยถาวินเย ยถาปตฺติยา โจเทตา โหติ, ยถาธมฺเม ยถาวินเย ยถาปตฺติยา กาเรตา โหติ, ยถาทิฏฺฐิยา พฺยากตา โหติ. อิเมหิ โข อุปาลิ ปญฺจหงฺเคหิ สมนฺนาคเตน ภิกฺขุนา สํเฆ โวหริตพฺพํ.}

\prose{อปเรหิปิ อุปาลิ ปญฺจหงฺเคหิ สมนฺนาคเตน ภิกฺขุนา สํเฆ น โวหริตพฺพํ. กตเมหิ ปญฺจหิ, อาปตฺตานาปตฺติํ น ชานาติ, ลหุกครุกํ อาปตฺติํ น ชานาติ, สาวเสสา\-นวเสสํ อาปตฺติํ น ชานาติ, \csromanfolio{322}ทุฏฺฐุลฺลา\-ทุฏฺฐุลฺลํ อาปตฺติํ น ชานาติ, สปฺปฏิกมฺมํ อปฺปฏิกมฺมํ อาปตฺติํ น ชานาติ. อิเมหิ โข อุปาลิ ปญฺจหงฺเคหิ สมนฺนาคเตน ภิกฺขุนา สํเฆ น โวหริตพฺพํ. ปญฺจหุปาลิ องฺเคหิ สมนฺนาคเตน ภิกฺขุนา สํเฆ โวหริตพฺพํ. กตเมหิ ปญฺจหิ, อาปตฺตานาปตฺติํ ชานาติ, ลหุกครุกํ อาปตฺติํ ชานาติ, สาวเสสา\-นวเสสํ อาปตฺติํ ชานาติ, ทุฏฺฐุลฺลทุฏฺฐุลฺลํ อาปตฺติํ ชานาติ, สปฺปฏิกมฺมํ อปฺปฏิกมฺมํ อาปตฺติํ ชานาติ. อิเมหิ โข อุปาลิ ปญฺจหงฺเคหิ สมนฺนาคเตน ภิกฺขุนา สํเฆ โวหริตพฺพํ.}

\prose{อปเรหิปิ อุปาลิ ปญฺจหงฺเคหิ สมนฺนาคเตน ภิกฺขุนา สํเฆ น โวหริตพฺพํ. กตเมหิ ปญฺจหิ, กมฺมํ น ชานาติ, กมฺมสฺส กรณํ น ชานาติ, กมฺมสฺส วตฺถุํ น ชานาติ, กมฺมสฺส วตฺตํ น ชานาติ, กมฺมสฺส วุปสมํ น ชานาติ. อิเมหิ โข อุปาลิ ปญฺจหงฺเคหิ สมนฺนาคเตน ภิกฺขุนา สํเฆ น โวหริตพฺพํ. ปญฺจหุปาลิ องฺเคหิ สมนฺนาคเตน ภิกฺขุนา สํเฆ โวหริตพฺพํ. กตเมหิ ปญฺจหิ, กมฺมํ ชานาติ, กมฺมสฺส กรณํ ชานาติ, กมฺมสฺส วตฺถุํ ชานาติ, กมฺมสฺส วตฺตํ ชานาติ, กมฺมสฺส วูปสมํ ชานาติ. อิเมหิ โข อุปาลิ ปญฺจหงฺเคหิ สมนฺนาคเตน ภิกฺขุนา สํเฆ โวหริตพฺพํ.}

\prose{อปเรหิปิ อุปาลิ ปญฺจหงฺเคหิ สมนฺนาคเตน ภิกฺขุนา สํเฆ น โวหริตพฺพํ. กตเมหิ ปญฺจหิ, วตฺถุํ น ชานาติ, นิทานํ น ชานาติ, ปญฺญตฺติํ น ชานาติ, ปทปจฺจา\-ภฏฺฐํ น ชานาติ, อนุสนฺธิ\-วจนปถํ น ชานาติ. อิเมหิ โข อุปาลิ ปญฺจหงฺเคหิ สมนฺนาคเตน ภิกฺขุนา สํเฆ น โวหริตพฺพํ. ปญฺจหุปาลิ องฺเคหิ สมนฺนาคเตน ภิกฺขุนา สํเฆ โวหริตพฺพํ. กตเมหิ ปญฺจหิ, วตฺถุํ ชานาติ, นิทานํ ชานาติ, ปญฺญตฺติํ ชานาติ, ปทปจฺจา\-ภฏฺฐํ ชานาติ, อนุสนฺธิ\-วจนปถํ ชานาติ. อิเมหิ โข อุปาลิ ปญฺจหงฺเคหิ สมนฺนาคเตน ภิกฺขุนา สํเฆ โวหริตพฺพํ.}

\prose{อปเรหิปิ อุปาลิ ปญฺจหงฺเคหิ สมนฺนาคเตน ภิกฺขุนา สํเฆ น โวหริตพฺพํ. กตเมหิ ปญฺจหิ, ฉนฺทาคติํ คจฺฉติ, โทสาคติํ คจฺฉติ, โมหาคติํ คจฺฉติ, ภยาคติํ คจฺฉติ, อลชฺชี จ โหติ. อิเมหิ โข อุปาลิ ปญฺจหงฺเคหิ สมนฺนาคเตน ภิกฺขุนา สํเฆ น โวหริตพฺพํ. ปญฺจหุปาลิ องฺเคหิ สมนฺนาคเตน ภิกฺขุนา สํเฆ โวหริตพฺพํ. กตเมหิ ปญฺจหิ, น ฉนฺทาคติํ คจฺฉติ, น โทสาคติํ คจฺฉติ, น \csromanfolio{323}โมหาคติํ คจฺฉติ, น ภยาคติํ คจฺฉติ, ลชฺชี จ โหติ. อิเมหิ โข อุปาลิ ปญฺจหงฺเคหิ สมนฺนาคเตน ภิกฺขุนา สํเฆ โวหริตพฺพํ.}

\prose{อปเรหิปิ อุปาลิ ปญฺจหงฺเคหิ สมนฺนาคเตน ภิกฺขุนา สํเฆ น โวหริตพฺพํ. กตเมหิ ปญฺจหิ, ฉนฺทาคติํ คจฺฉติ, โทสาคติํ คจฺฉติ, โมหาคติํ คจฺฉติ, ภยาคติํ คจฺฉติ, อกุสโล จ โหติ วินเย. อิเมหิ โข อุปาลิ ปญฺจหงฺเคหิ สมนฺนาคเตน ภิกฺขุนา สํเฆ น โวหริตพฺพํ. ปญฺจหุปาลิ องฺเคหิ สมนฺนาคเตน ภิกฺขุนา สํเฆ โวหริตพฺพํ. กตเมหิ ปญฺจหิ, น ฉนฺทาคติํ คจฺฉติ, น โทสาคติํ คจฺฉติ, น โมหาคติํ คจฺฉติ, น ภยาคติํ คจฺฉติ, กุสโล จ โหติ วินเย. อิเมหิ โข อุปาลิ ปญฺจหงฺเคหิ สมนฺนาคเตน ภิกฺขุนา สํเฆ โวหริตพฺพํ.}

\prose{อปเรหิปิ อุปาลิ ปญฺจหงฺเคหิ สมนฺนาคเตน ภิกฺขุนา สํเฆ น โวหริตพฺพํ. กตเมหิ ปญฺจหิ, ญตฺติํ น ชานาติ, ญตฺติยา กรณํ น ชานาติ, ญตฺติยา อนุสฺสาวนํ น ชานาติ, ญตฺติยา สมถํ น ชานาติ, ญตฺติยา วูปสมํ น ชานาติ. อิเมหิ โข อุปาลิ ปญฺจหงฺเคหิ สมนฺนาคเตน ภิกฺขุนา สํเฆ น โวหริตพฺพํ. ปญฺจหุปาลิ องฺเคหิ สมนฺนาคเตน ภิกฺขุนา สํเฆ โวหริตพฺพํ. กตเมหิ ปญฺจหิ, ญตฺติํ ชานาติ, ญตฺติยา กรณํ ชานาติ, ญตฺติยา อนุสฺสาวนํ ชานาติ, ญตฺติยา สมถํ ชานาติ, ญตฺติยา วูปสมํ ชานาติ. อิเมหิ โข อุปาลิ ปญฺจหงฺเคหิ สมนฺนาคเตน ภิกฺขุนา สํเฆ โวหริตพฺพํ.}

\prose{อปเรหิปิ อุปาลิ ปญฺจหงฺเคหิ สมนฺนาคเตน ภิกฺขุนา สํเฆ น โวหริตพฺพํ. กตเมหิ ปญฺจหิ, สุตฺตํ น ชานาติ, สุตฺตานุโลมํ น ชานาติ, วินยํ น ชานาติ, วินยานุโลมํ น ชานาติ, น จ ฐานาฐาน\-กุสโล โหติ. อิเมหิ โข อุปาลิ ปญฺจหงฺเคหิ สมนฺนาคเตน ภิกฺขุนา สํเฆ น โวหริตพฺพํ. ปญฺจหุปาลิ องฺเคหิ สมนฺนาคเตน ภิกฺขุนา สํเฆ โวหริตพฺพํ. กตเมหิ ปญฺจหิ, สุตฺตํ ชานาติ, สุตฺตานุโลมํ ชานาติ, วินยํ ชานาติ, วินยานุโลมํ ชานาติ, ฐานาฐาน\-กุสโล จ โหติ. อิเมหิ โข อุปาลิ ปญฺจหงฺเคหิ สมนฺนาคเตน ภิกฺขุนา สํเฆ โวหริตพฺพํ.}

\prose{\csromanfolio{324}อปเรหิปิ อุปาลิ ปญฺจหงฺเคหิ สมนฺนาคเตน ภิกฺขุนา สํเฆ น โวหริตพฺพํ. กตเมหิ ปญฺจหิ, ธมฺมํ น ชานาติ, ธมฺมานุโลมํ น ชานาติ, วินยํ น ชานาติ, วินยานุโลมํ น ชานาติ, น จ ปุพฺพาปร\-กุสโล โหติ. อิเมหิ โข อุปาลิ ปญฺจหงฺเคหิ สมนฺนาคเตน ภิกฺขุนา สํเฆ น โวหริตพฺพํ. ปญฺจหุปาลิ องฺเคหิ สมนฺนาคเตน ภิกฺขุนา สํเฆ โวหริตพฺพํ. กตเมหิ ปญฺจหิ, ธมฺมํ ชานาติ, ธมฺมานุโลมํ ชานาติ, วินยํ ชานาติ, วินยานุโลมํ ชานาติ, ปุพฺพาปร\-กุสโล จ โหติ. อิเมหิ โข อุปาลิ ปญฺจหงฺเคหิ สมนฺนาคเตน ภิกฺขุนา สํเฆ โวหริตพฺพนฺติ.}

\vspace{\tipitakaparskip}
\csromancenter{โวหารวคฺโค นิฏฺฐิโต ตติโย.}
\csromancenter{ตสฺสุทฺทานํ.}
\setlength{\csromangathapagewidth}{0pt}
\setlength{\csromangathawakwidth}{0pt}
\csromangathameasuregroup{อาปตฺติ อธิกรณํ, \\ กมฺมํ วตฺถุํ อลชฺชี จ,}{\csromangathabat{อาปตฺติ อธิกรณํ,}{ปสยฺหาปตฺติ ชานนา.} \\ \csromangathabat{กมฺมํ วตฺถุํ อลชฺชี จ,}{อกุสโล จ ญตฺติยา.}}
\csromangathasetleft{อาปตฺติ อธิกรณํ, \\ กมฺมํ วตฺถุํ อลชฺชี จ,}
\csromangathagroup{\csromangathastanza{\csromangathabat{อาปตฺติ อธิกรณํ,}{ปสยฺหาปตฺติ ชานนา.} \\ \csromangathabat{กมฺมํ วตฺถุํ อลชฺชี จ,}{อกุสโล จ ญตฺติยา.}}}

\vspace{\tipitakaparskip}
\csromancenter{สุตฺตํ น ชานาติ ธมฺมํ, ตติโย วคฺคสงฺคโหติ.}
\csromansectionrule
\csromanmatikaanchor{mtk.223}
\csromantocmarkref{subparagraph}{4. ทิฏฺฐาวิกมฺมวคฺค}{mtk.223}
\bookmark[dest={mtk.223},level=5]{4. ทิฏฺฐาวิกมฺมวคฺค}
\csromanheader{5}{4. ทิฏฺฐาวิกมฺม\-วคฺค}
\proseitem{425}{กติ นุ โข ภนฺเต อธมฺมิกา ทิฏฺฐา\-วิ\-กมฺมาติ. ปญฺจิเม อุปาลิ อธมฺมิกา ทิฏฺฐาวิกมฺมา. กตเม ปญฺจ, อนาปตฺติยา ทิฏฺฐิํ อาวิ กโรติ, อเทสนาคามินิยา อาปตฺติยา ทิฏฺฐิํ อาวิ กโรติ, เทสิตาย อาปตฺติยา ทิฏฺฐิํ อาวิ กโรติ, จตูหิ ปญฺจหิ ทิฏฺฐิํ อาวิ กโรติ, มโน มานเสน ทิฏฺฐิํ อาวิ กโรติ. อิเม โข อุปาลิ ปญฺจ อธมฺมิกา ทิฏฺฐาวิกมฺมา.}

\prose{ปญฺจิเม อุปาลิ ธมฺมิกา ทิฏฺฐาวิกมฺมา. กตเม ปญฺจ, อาปตฺติยา ทิฏฺฐิํ อาวิ กโรติ, เทสนาคามินิยา อาปตฺติยา ทิฏฺฐิํ อาวิ กโรติ, อเทสิตาย อาปตฺติยา ทิฏฺฐิํ อาวิ กโรติ, น จตูหิ ปญฺจหิ ทิฏฺฐิํ อาวิ กโรติ, น มโน มานเสน ทิฏฺฐิํ อาวิ กโรติ. อิเม โข อุปาลิ ปญฺจ ธมฺมิกา ทิฏฺฐาวิกมฺมา.}

\prose{\csromanfolio{325}อปเรปิ อุปาลิ ปญฺจ อธมฺมิกา ทิฏฺฐาวิกมฺมา. กตเม ปญฺจ, นานา\-สํวาสกสฺส สนฺติเก ทิฏฺฐิํ อาวิ กโรติ, นานาสีมาย ฐิตสฺส สนฺติเก ทิฏฺฐิํ อาวิ กโรติ, อปกตตฺตสฺส สนฺติเก ทิฏฺฐิํ อาวิ กโรติ, จตูหิ ปญฺจหิ ทิฏฺฐิํ อาวิ กโรติ, มโน มานเสน ทิฏฺฐิํ อาวิ กโรติ. อิเม โข อุปาลิ ปญฺจ อธมฺมิกา ทิฏฺฐาวิกมฺมา.}

\prose{ปญฺจิเม อุปาลิ ธมฺมิกา ทิฏฺฐาวิกมฺมา. กตเม ปญฺจ, สมาน\-สํวาส\-กสฺส สนฺติเก ทิฏฺฐิํ อาวิ กโรติ, สมานสีมาย ฐิตสฺส สนฺติเก ทิฏฺฐิํ อาวิ กโรติ, ปกตตฺตสฺส สนฺติเก ทิฏฺฐิํ อาวิ กโรติ, น จตูหิ ปญฺจหิ ทิฏฺฐิํ อาวิ กโรติ, น มโน มานเสน ทิฏฺฐิํ อาวิ กโรติ. อิเม โข อุปาลิ ปญฺจ ธมฺมิกา ทิฏฺฐา\-วิ\-กมฺมาติ.}

\proseitem{426}{กติ นุ โข ภนฺเต อธมฺมิกา ปฏิคฺคหาติ. ปญฺจิเม อุปาลิ อธมฺมิกา ปฏิคฺคหา. กตเม ปญฺจ, กาเยน ทิยฺยมานํ กาเยน อปฺปฏิคฺคหิตํ, กาเยน ทิยฺยมานํ กาย\-ปฺปฏิพทฺเธน อปฺปฏิคฺคหิตํ, กาย\-ปฺปฏิพทฺเธน ทิยฺยมานํ กาเยน อปฺปฏิคฺคหิตํ, กาย\-ปฺปฏิพทฺเธน ทิยฺยมานํ กาย\-ปฺปฏิพทฺเธน อปฺปฏิคฺคหิตํ, นิสฺสคฺคิเยน ทิยฺยมานํ กาเยน วา กาย\-ปฺปฏิพทฺเธน วา อปฺปฏิคฺคหิตํ. อิเม โข อุปาลิ ปญฺจ อธมฺมิกา ปฏิคฺคหา.}

\prose{ปญฺจิเม อุปาลิ ธมฺมิกา ปฏิคฺคหา. กตเม ปญฺจ, กาเยน ทิยฺยมานํ กาเยน ปฏิคฺคหิตํ, กาเยน ทิยฺยมานํ กาย\-ปฺปฏิพทฺเธน ปฏิคฺคหิตํ, กาย\-ปฺปฏิพทฺเธน ทิยฺยมานํ กาเยน ปฏิคฺคหิตํ, กาย\-ปฺปฏิพทฺเธน ทิยฺยมานํ กาย\-ปฺปฏิพทฺเธน ปฏิคฺคหิตํ, นิสฺสคฺคิเยน ทิยฺยมานํ กาเยน วา กาย\-ปฺปฏิพทฺเธน วา ปฏิคฺคหิตํ. อิเม โข อุปาลิ ปญฺจ ธมฺมิกา ปฏิคฺคหาติ.}

\proseitem{427}{กติ นุ โข ภนฺเต อนติริตฺตาติ. ปญฺจิเม อุปาลิ อนติริตฺตา. กตเม ปญฺจ,\csromansymbolfootnote{*}{วิ 2. 111 ปิฏฺเฐปิ.}อกปฺปิยกตํ โหติ, อปฺปฏิคฺคหิตกตํ โหติ, อนุจฺจาริตกตํ โหติ, อหตฺถปาเส กตํ โหติ, “อลเมตํ สพฺพนฺ”ติ อวุตฺตํ โหติ. อิเม โข อุปาลิ ปญฺจ อนติริตฺตา.}

\prose{ปญฺจิเม อุปาลิ อติริตฺตา. กตเม ปญฺจ,\csromansymbolmark{*}กปฺปิยกตํ โหติ, ปฏิคฺคหิตกตํ โหติ, อุจฺจาริตกตํ โหติ, หตฺถปาเส กตํ โหติ, “อลเมตํ สพฺพนฺ”ติ วุตฺตํ โหติ. อิเม โข อุปาลิ ปญฺจ อติริตฺตาติ.}

\proseitem{428}{\csromanfolio{326}กติหิ นุ โข ภนฺเต อากาเรหิ ปวารณา ปญฺญายตีติ. ปญฺจหุปาลิ อากาเรหิ ปวารณา ปญฺญายติ. กตเมหิ ปญฺจหิ,\csromansymbolfootnote{*}{วิ 2. 111 ปิฏฺเฐปิ.}อสนํ ปญฺญายติ, โภชนํ ปญฺญายติ, หตฺถปาเส ฐิโต อภิหรติ, ปฏิกฺเขโป ปญฺญายติ. อิเมหิ โข อุปาลิ ปญฺจหากาเรหิ ปวารณา ปญฺญายตีติ.}

\proseitem{429}{กติ นุ โข ภนฺเต อธมฺมิกา ปฏิญฺญาต\-กรณาติ. ปญฺจิเม อุปาลิ อธมฺมิกา ปฏิญฺญาต\-กรณา. กตเม ปญฺจ, ภิกฺขุ ปาราชิกํ อชฺฌาปนฺโน โหติ, ปาราชิเกน โจทิยมาโน สํฆาทิเสสํ อชฺฌาปนฺโน ปฏิชานาติ, ตํ สํโฆ สํฆาทิเสเสน กาเรติ, อธมฺมิกํ ปฏิญฺญาต\-กรณํ. ภิกฺขุ ปาราชิกํ อชฺฌาปนฺโน โหติ, ปาราชิเกน โจทิยมาโน ปาจิตฺติยํ ฯเปฯ. ปาฏิเทสนียํ. ทุกฺกฏํ อชฺฌาปนฺโน ปฏิชานาติ, ตํ สํโฆ ทุกฺกเฏน กาเรติ, อธมฺมิกํ ปฏิญฺญาต\-กรณํ. ภิกฺขุ สํฆาทิเสสํ ฯเปฯ. ปาจิตฺติยํ. ปาฏิเทสนียํ. ทุกฺกฏํ อชฺฌาปนฺโน โหติ. ทุกฺกเฏน โจทิยมาโน ปาราชิกํ อชฺฌาปนฺโน ปฏิชานาติ, ตํ สํโฆ ปาราชิเกน กาเรติ, อธมฺมิกํ ปฏิญฺญาต\-กรณํ. ภิกฺขุ ทุกฺกฏํ อชฺฌาปนฺโน โหติ, ทุกฺกเฏน โจทิยมาโน สํฆาทิเสสํ ฯเปฯ ปาจิตฺติยํ, ปาฏิเทสนียํ อชฺฌาปนฺโน ปฏิชานาติ, ตํ สํโฆ ปาฏิเทสนีเยน กาเรติ, อธมฺมิกํ ปฏิญฺญาต\-กรณํ. อิเม โข อุปาลิ ปญฺจ อธมฺมิกา ปฏิญฺญาต\-กรณา.}

\prose{ปญฺจิเม อุปาลิ ธมฺมิกา ปฏิญฺญาต\-กรณา. กตเม ปญฺจ, ภิกฺขุ ปาราชิกํ อชฺฌาปนฺโน โหติ, ปาราชิเกน โจทิยมาโน ปาราชิกํ อชฺฌาปนฺโน ปฏิชานาติ, ตํ สํโฆ ปาราชิเกน กาเรติ, ธมฺมิกํ ปฏิญฺญาต\-กรณํ. ภิกฺขุ สํฆาทิเสสํ ฯเปฯ ปาจิตฺติยํ. ปาฏิเทสนียํ. ทุกฺกฏํ อชฺฌาปนฺโน โหติ, ทุกฺกเฏน โจทิยมาโน ทุกฺกฏํ อชฺฌาปนฺโน ปฏิชานาติ, ตํ สํโฆ ทุกฺกเฏน กาเรติ, ธมฺมิกํ ปฏิญฺญาต\-กรณํ. อิเม โข อุปาลิ ปญฺจ ธมฺมิกา ปฏิญฺญาต\-กรณาติ.}

\proseitem{430}{กติหิ นุ โข ภนฺเต องฺเคหิ สมนฺนาคตสฺส ภิกฺขุโน โอกาสกมฺมํ การาเปนฺตสฺส นาลํ โอกาสกมฺมํ กาตุนฺติ. ปญฺจหุปาลิ องฺเคหิ สมนฺนาคตสฺส ภิกฺขุโน โอกาสกมฺมํ การาเปนฺตสฺส นาลํ \csromanfolio{327}โอกาสกมฺมํ กาตุํ. กตเมหิ ปญฺจหิ, อลชฺชี จ โหติ, พาโล จ, อปกตตฺโต จ, จาวนาธิปฺปาโย วตฺตา โหติ, โน วุฏฺฐานา\-ธิปฺปาโย. อิเมหิ โข อุปาลิ ปญฺจหงฺเคหิ สมนฺนาคตสฺส ภิกฺขุโน โอกาสกมฺมํ การาเปนฺตสฺส นาลํ โอกาสกมฺมํ กาตุํ.}

\prose{ปญฺจหุปาลิ องฺเคหิ สมนฺนาคตสฺส ภิกฺขุโน โอกาสกมฺมํ การาเปนฺตสฺส อลํ โอกาสกมฺมํ กาตุํ. กตเมหิ ปญฺจหิ, ลชฺชี จ โหติ, ปณฺฑิโต จ, ปกตตฺโต จ, วุฏฺฐานา\-ธิปฺปาโย วตฺตา โหติ, โน จาวนาธิปฺปาโย. อิเมหิ โข อุปาลิ ปญฺจหงฺเคหิ สมนฺนาคตสฺส ภิกฺขุโน โอกาสกมฺมํ การาเปนฺตสฺส อหํ โอกาสกมฺมํ กาตุนฺติ.}

\proseitem{431}{กติหิ นุ โข ภนฺเต องฺเคหิ สมนฺนาคเตน ภิกฺขุนา สทฺธิํ วินโย น สากจฺฉิ ตพฺโพติ. ปญฺจหุปาลิ องฺเคหิ สมนฺนาคเตน ภิกฺขุนา สทฺธิํ วินโย น สากจฺฉิตพฺโพ. กตเมหิ ปญฺจหิ, วตฺถุํ น ชานาติ, นิทานํ น ชานาติ, ปญฺญตฺติํ น ชานาติ, ปทปจฺจา\-ภฏฺฐํ น ชานาติ, อนุสนฺธิ\-วจนปถํ น ชานาติ. อิเมหิ โข อุปาลิ ปญฺจหงฺเคหิ สมนฺนาคเตน ภิกฺขุนา สทฺธิํ วินโย น สากจฺฉิตพฺโพ.}

\prose{ปญฺจหุปาลิ องฺเคหิ สมนฺนาคเตน ภิกฺขุนา สทฺธิํ วินโย สากจฺฉิตพฺโพ. กตเมหิ ปญฺจหิ, วตฺถุํ ชานาติ, นิทานํ ชานาติ, ปญฺญตฺติํ ชานาติ, ปทปจฺจา\-ภฏฺฐํ ชานาติ, อนุสนฺธิ\-วจนปถํ ชานาติ. อิเมหิ โข อุปาลิ ปญฺจหงฺเคหิ สมนฺนาคเตน ภิกฺขุนา สทฺธิํ วินโย สากจฺฉิตพฺโพติ.}

\proseitem{432}{กติ นุ โข ภนฺเต ปญฺหาปุจฺฉาติ. ปญฺจิมา อุปาลิ ปญฺหาปุจฺฉา. กตมา ปญฺจ, มนฺทตฺตา โมมูหตฺตา ปญฺหํ ปุจฺฉติ, ปาปิจฺโฉ อิจฺฉาปกโต ปญฺหํ ปุจฺฉติ, ปริภวา ปญฺหํ ปุจฺฉติ, อญฺญาตุกาโม ปญฺหํ ปุจฺฉติ, “สเจ เม ปญฺหํ ปุฏฺโฐ สมฺมเทว พฺยากริสฺสติ, อิจฺเจตํ กุสลํ, โน เจ ปญฺหํ ปุฏฺโฐ สมฺมเทว พฺยากริสฺสติ, อหมสฺส สมฺมเทว พฺยากริสฺสามี”ติ ปญฺหํ ปุจฺฉติ. อิมา โข อุปาลิ ปญฺจ ปญฺหาปุจฺฉาติ.}

\proseitem{433}{กติ นุ โข ภนฺเต อญฺญ\-พฺยากรณาติ. ปญฺจิเม อุปาลิ อญฺญพฺยากรณา. กตเม ปญฺจ, มนฺทตฺตา โมมูหตฺตา อญฺญํ พฺยากโรติ, ปาปิจฺโฉ อิจฺฉาปกโต อญฺญํ พฺยากโรติ, อุมฺมาทา จิตฺตกฺเขปา อญฺญํ \csromanfolio{328}พฺยากโรติ, อธิมาเนน อญฺญํ พฺยากโรติ, ภูตํ อญฺญํ พฺยากโรติ. อิเม โข อุปาลิ ปญฺจ อญฺญ\-พฺยากรณาติ.}

\proseitem{434}{กติ นุ โข ภนฺเต วิสุทฺธิโยติ. ปญฺจิมา อุปาลิ วิสุทฺธิโย. กตมา ปญฺจ, นิทานํ อุทฺทิสิตฺวา อวเสสํ สุเตน สาเวตพฺพํ, อยํ ปฐมา วิสุทฺธิ, นิทานํ อุทฺทิสิตฺวา จตฺตาริ ปาราชิกานิ อุทฺทิสิตฺวา อวเสสํ สุเตน สาเวตพฺพํ, อยํ ทุติยา วิสุทฺธิ, นิทานํ อุทฺทิสิตฺวา จตฺตาริ ปาราชิกานิ อุทฺทิสิตฺวา เตรส สํฆาทิเสเส อุทฺทิสิตฺวา อวเสสํ สุเตน สาเวตพฺพํ, อยํ ตติยา วิสุทฺธิ, นิทานํ อุทฺทิสิตฺวา จตฺตาริ ปาราชิกานิ อุทฺทิสิตฺวา เตรส สํฆาทิเสเส อุทฺทิสิตฺวา เทฺว อนิยเต อุทฺทิสิตฺวา อวเสสํ สุเตน สาเวตพฺพํ, อยํ จตุตฺถา วิสุทฺธิ, วิตฺถาเรเนว ปญฺจมี. อิมา โข อุปาลิ ปญฺจ วิสุทฺธิโยติ.}

\proseitem{435}{กติ นุโข ภนฺเต โภชนาติ. ปญฺจิเม อุปาลิ โภชนา. กตเม ปญฺจ, โอทโน กุมฺมาโส สตฺตุ มจฺโฉ มํสํ. อิเม โข อุปาลิ ปญฺจ โภชนาติ.}

\vspace{\tipitakaparskip}
\csromancenter{ทิฏฺฐาวิกมฺม\-วคฺโค นิฏฺฐิโต จตุตฺโถ.}
\csromancenter{ตสฺสุทฺทานํ.}
\setlength{\csromangathapagewidth}{0pt}
\setlength{\csromangathawakwidth}{0pt}
\csromangathameasuregroup{ทิฏฺฐาวิกมฺมา อปเร, \\ ปวารณา ปฏิญฺญาตํ,}{\csromangathabat{ทิฏฺฐาวิกมฺมา อปเร,}{ปฏิคฺคหานติริตฺตา.} \\ \csromangathabat{ปวารณา ปฏิญฺญาตํ,}{โอกาสํ สากจฺเฉน จ.}}
\csromangathasetleft{ทิฏฺฐาวิกมฺมา อปเร, \\ ปวารณา ปฏิญฺญาตํ,}
\csromangathagroup{\csromangathastanza{\csromangathabat{ทิฏฺฐาวิกมฺมา อปเร,}{ปฏิคฺคหานติริตฺตา.} \\ \csromangathabat{ปวารณา ปฏิญฺญาตํ,}{โอกาสํ สากจฺเฉน จ.}}}

\vspace{\tipitakaparskip}
\csromancenter{ปญฺหํ อญฺญพฺยากรณา, วิสุทฺธิ จาปิ โภชนาติ.}
\csromansectionrule
\csromanmatikaanchor{mtk.224}
\csromantocmarkref{subparagraph}{5. อตฺตาทานวคฺค}{mtk.224}
\bookmark[dest={mtk.224},level=5]{5. อตฺตาทานวคฺค}
\csromanheader{5}{5. \textbf{อตฺตาทานวคฺค}}
\proseitem{436}{\csromansymbolfootnote{*}{วิ 4. 436; อํ 3. 318 ปิฏฺเฐสุปิ.}โจทเกน ภนฺเต ภิกฺขุนา ปรํ โจเทตุกาเมน กติ ธมฺเม อชฺฌตฺตํ ปจฺจเวกฺขิตฺวา ปโร โจเทตพฺโพติ. โจทเกนุปาลิ ภิกฺขุนา ปรํ โจเทตุกาเมน ปญฺจ ธมฺเม อชฺฌตฺตํ ปจฺจเวกฺขิตฺวา ปโร โจเทตพฺโพ. กตเม ปญฺจ, โจทเกนุปาลิ ภิกฺขุนา ปรํ โจเทตุกาเมน เอวํ ปจฺจเวกฺขิตพฺพํ “ปริสุทฺธ\-กาย\-สมาจาโร นุ โขมฺหิ, ปริสุทฺเธนมฺหิ กายสมาจาเรน สมนฺนาคโต อจฺฉิทฺเทน อปฺปฏิมํเสน, สํวิชฺชติ นุ โข \csromanfolio{329}เม เอโส ธมฺโม, อุทาหุ โน”ติ. โน เจ อุปาลิ ภิกฺขุ ปริสุทฺธ\-กาย\-สมาจาโร โหติ ปริสุทฺเธน กายสมาจาเรน สมนฺนาคโต อจฺฉิทฺเทน อปฺปฏิมํเสน, ตสฺส ภวนฺติ วตฺตาโร “อิงฺฆ ตาว อายสฺมา กายิกํ สิกฺขสฺสู”ติ, อิติสฺส ภวนฺติ วตฺตาโร.}

\prose{ปุน จปรํ อุปาลิ โจทเกน ภิกฺขุนา ปรํ โจเทตุกาเมน เอวํ ปจฺจเวกฺขิตพฺพํ “ปริสุทฺธ\-วจี\-สมาจาโร นุ โขมฺหิ, ปริสุทฺเธนมฺหิ วจีสมาจาเรน สมนฺนาคโต อจฺฉิทฺเทน อปฺปฏิมํเสน, สํวิชฺชติ นุ โข เม เอโส ธมฺโม, อุทาหุ โน”ติ. โน เจ อุปาลิ ภิกฺขุ ปริสุทฺธ\-วจี\-สมาจาโร โหติ ปริสุทฺเธน วจีสมาจาเรน สมนฺนาคโต อจฺฉิทฺเทน อปฺปฏิมํเสน, ตสฺส ภวนฺติ วตฺตาโร “อิงฺฆ ตาว อายสฺมา วาจสิกํ สิกฺขสฺสู”ติ, อิติสฺส ภวนฺติ วตฺตาโร.}

\prose{ปุน จปรํ อุปาลิ โจทเกน ภิกฺขุนา ปรํ โจเทตุกาเมน เอวํ ปจฺจเวกฺขิตพฺพํ “เมตฺตํ นุ โข เม จิตฺตํ ปจฺจุปฏฺฐิตํ สพฺรหฺมจารีสุ อนาฆาตํ, สํวิชฺชติ นุ โข เม เอโส ธมฺโม, อุทาหุ โน”ติ. โน เจ อุปาลิ ภิกฺขุโน เมตฺตํ จิตฺตํ ปจฺจุปฏฺฐิตํ โหติ สพฺรหฺมจารีสุ อนาฆาตํ. ตสฺส ภวนฺติ วตฺตาโร “อิงฺฆ ตาว อายสฺมา สพฺรหฺมจารีสุ เมตฺตํ จิตฺตํ อุปฏฺฐาเปหี”ติ, อิติสฺส ภวนฺติ วตฺตาโร.}

\prose{ปุน จปรํ อุปาลิ โจทเกน ภิกฺขุนา ปรํ โจเทตุกาเมน เอวํ ปจฺจเวกฺขิตพฺพํ “พหุสฺสุโต นุ โขมฺหิ สุตธโร สุตสนฺนิจโย, เย เต ธมฺมา อาทิกลฺยาณา มชฺเฌกลฺยาณา ปริโยสาน\-กลฺยาณา สาตฺถํ สพฺยญฺชนํ เกวล\-ปริปุณฺณํ ปริสุทฺธํ พฺรหฺมจริยํ อภิวทนฺติ, ตถารูปา เม ธมฺมา พหุสฺสุตา โหนฺติ ธาตา วจสา ปริจิตา มนสา\-นุเปกฺขิตา ทิฏฺฐิยา สุปฺปฏิวิทฺธา, สํวิชฺชติ นุ โข เม เอโส ธมฺโม, อุทาหุ โน”ติ. โน เจ อุปาลิ ภิกฺขุ พหุสฺสุโต โหติ สุตธโร สุตสนฺนิจโย, เย เต ธมฺมา อาทิกลฺยาณา มชฺเฌกลฺยาณา ปริโยสาน\-กลฺยาณา สาตฺถํ สพฺยญฺชนํ เกวล\-ปริปุณฺณํ ปริสุทฺธํ พฺรหฺมจริยํ อภิวทนฺติ, ตถารูปสฺส ธมฺมา น พหุสฺสุตา โหนฺติ ธาตา วจสา ปริจิตา มนสา\-นุเปกฺขิตา ทิฏฺฐิยา สุปฺปฏิวิทฺธา, ตสฺส ภวนฺติ วตฺตาโร “อิงฺฆ ตาว อายสฺมา อาคมํ ปริยาปุณสฺสู”ติ, อิติสฺส ภวนฺติ วตฺตาโร.}

\prose{\csromanfolio{330}ปุน จปรํ อุปาลิ โจทเกน ภิกฺขุนา ปรํ โจเทตุกาเมน เอวํ ปจฺจเวกฺขิตพฺพํ “อุภยานิ โข เม ปาติโมกฺขานิ วิตฺถาเรน สฺวาคตานิ โหนฺติ สุวิภตฺตานิ สุปฺปวตฺตีนิ สุวินิจฺฉิตานิ สุตฺตโส อนุพฺยญฺชนโส, สํวิชฺชติ นุ โข เม เอโส ธมฺโม, อุทาหุ โน”ติ. โน เจ อุปาลิ ภิกฺขุโน อุภยานิ ปาติโมกฺขานิ วิตฺถาเรน สฺวาคตานิ โหนฺติ สุวิภตฺตานิ สุปฺปวตฺตีนิ สุวินิจฺฉิตานิ สุตฺตโส อนุพฺยญฺชนโส, อิทํ ปนาวุโส กตฺถ วุตฺตํ ภควตาติ อิติ ปุฏฺโฐ น สมฺปายติ\csromansharedfootnote{a8960b3f402cf3b1}{น สมฺปาทยติ (ก), น สมฺปาเทติ (สฺยา)}, ตสฺส ภวนฺติ วตฺตาโร “อิงฺฆ ตาว อายสฺมา วินยํ ปริยาปุณสฺสู”ติ, อิติสฺส ภวนฺติ วตฺตาโร. โจทเกนุปาลิ ภิกฺขุนา ปรํ โจเทตุกาเมน อิเม ปญฺจ ธมฺเม อชฺฌตฺตํ ปจฺจเวกฺขิตฺวา ปโร โจเทตพฺโพติ.}

\proseitem{437}{โจทเกน ภนฺเต ภิกฺขุนา ปรํ โจเทตุกาเมน กติ ธมฺเม อชฺฌตฺตํ อุปฏฺฐาเปตฺวา ปโร โจเทตพฺโพติ. โจทเกนุปาลิ ภิกฺขุนา ปรํ โจเทตุกาเมน ปญฺจ ธมฺเม อชฺฌตฺตํ อุปฏฺฐาเปตฺวา ปโร โจเทตพฺโพ. กตเม ปญฺจ, กาเลน วกฺขามิ โน อกาเลน, ภูเตน วกฺขามิ โน อภูเตน, สเณฺหน วกฺขามิ โน ผรุเสน, อตฺถสํหิเตน วกฺขามิ โน อนตฺถ\-สํหิเตน, เมตฺตาจิตฺโต วกฺขามิ โน โทสนฺตโรติ. โจทเกนุปาลิ ภิกฺขุนา ปรํ โจเทตุกาเมน อิเม ปญฺจ ธมฺเม อชฺฌตฺตํ อุปฏฺฐาเปตฺวา ปโร โจเทตพฺโพติ.}

\proseitem{438}{\csromansymbolfootnote{*}{วิ 4. 437 ปิฏฺเฐปิ.}โจทเกน ภนฺเต ภิกฺขุนา ปรํ โจเทตุกาเมน กติ ธมฺเม อชฺฌตฺตํ มนสิ กริตฺวา ปโร โจเทตพฺโพติ. โจทเกนุปาลิ ภิกฺขุนา ปรํ โจเทตุกาเมน ปญฺจ ธมฺเม อชฺฌตฺตํ มนสิ กริตฺวา ปโร โจเทตพฺโพ. กตเม ปญฺจ, การุญฺญตา, หิเตสิตา, อนุกมฺปตา, อาปตฺติวุฏฺฐานตา, วินยปุเรกฺขารตา. โจทเกนุปาลิ ภิกฺขุนา ปรํ โจเทตุกาเมน อิเม ปญฺจ ธมฺเม อชฺฌตฺตํ มนสิ กริตฺวา ปโร โจเทตพฺโพติ.}

\proseitem{439}{กติหิ นุ โข ภนฺเต องฺเคหิ สมนฺนาคตสฺส ภิกฺขุโน โอกาสกมฺมํ การาเปนฺตสฺส นาลํ โอกาสกมฺมํ กาตุนฺติ. ปญฺจหุปาลิ องฺเคหิ สมนฺนาคตสฺส ภิกฺขุโน โอกาสกมฺมํ การาเปนฺตสฺส นาลํ โอกาสกมฺมํ กาตุํ. กตเมหิ ปญฺจหิ, อปริสุทฺธ\-กายสมาจาโร \csromanfolio{331}โหติ, อปริสุทฺธ\-วจีสมาจาโร โหติ, อปริสุทฺธา\-ชีโว โหติ, พาโล โหติ อพฺยตฺโต, น ปฏิพโล อนุยุญฺชิยมาโน อนุโยคํ ทาตุํ. อิเมหิ โข อุปาลิ ปญฺจหงฺเคหิ สมนฺนาคตสฺส ภิกฺขุโน โอกาสกมฺมํ การาเปนฺตสฺส นาลํ โอกาสกมฺมํ กาตุํ.}

\prose{ปญฺจหุปาลิ องฺเคหิ สมนฺนาคตสฺส ภิกฺขุโน โอกาสกมฺมํ การาเปนฺตสฺส อลํ โอกาสกมฺมํ กาตุํ. กตเมหิ ปญฺจหิ, ปริสุทฺธ\-กาย\-สมาจาโร โหติ, ปริสุทฺธ\-วจี\-สมาจาโร โหติ, ปริสุทฺธาชีโว โหติ, ปณฺฑิโต โหติ พฺยตฺโต, ปฏิพโล อนุยุญฺชิยมาโน อนุโยคํ ทาตุํ. อิเมหิ โข อุปาลิ ปญฺจหงฺเคหิ สมนฺนาคตสฺส ภิกฺขุโน โอกาสกมฺมํ การาเปนฺตสฺส อลํ โอกาสกมฺมํ กาตุนฺติ.}

\proseitem{440}{\csromansymbolfootnote{*}{วิ 4. 434 ปิฏฺเฐปิ.}อตฺตาทานํ อาทาตุกาเมน ภนฺเต ภิกฺขุนา กติหงฺเคหิ สมนฺนาคตํ อตฺตาทานํ อาทาตพฺพนฺติ. อตฺตาทานํ อาทาตุกาเมนุ\-ปาลิ ภิกฺขุนา ปญฺจ\-งฺค\-สมนฺนาคตํ\csromansharedfootnote{e3acf7008800d806}{ปญฺจหงฺเคหิ สมนฺนาคตํ (สี, สฺยา)} อตฺตาทานํ อาทาตพฺพํ. กตเม ปญฺจ\csromansharedfootnote{fe56b050d1375859}{กตเมหิ ปญฺจหิ (สฺยา)}, อตฺถาทานํ อาทาตุกาเมน อุปาลิ ภิกฺขุนา เอวํ ปจฺจเวกฺขิตพฺพํ “ยํ โข อหํ อิมํ อตฺตาทานํ อาทาตุกาโม, กาโล นุ โข อิมํ อตฺตาทานํ อาทาตุํ, อุทาหุ โน”ติ. สเจ อุปาลิ ภิกฺขุ ปจฺจเวกฺขมาโน เอวํ ชานาติ “อกาโล อิมํ อตฺตาทานํ อาทาตุํ, โน กาโล”ติ. น ตํ อุปาลิ อตฺตาทานํ อาทาตพฺพํ.}

\prose{สเจ ปนุปาลิ ภิกฺขุ ปจฺจเวกฺขมาโน เอวํ ชานาติ “กาโล อิมํ อตฺตาทานํ อาทาตุํ, โน อกาโล”ติ. เตนุปาลิ ภิกฺขุนา อุตฺตริ ปจฺจเวกฺขิตพฺพํ “ยํ โข อหํ อิมํ อตฺตาทานํ อาทาตุกาโม, ภูตํ นุ โข อิทํ อตฺตาทานํ, อุทาหุ โน”ติ. สเจ อุปาลิ ภิกฺขุ ปจฺจเวกฺขมาโน เอวํ ชานาติ “อภูตํ อิทํ อตฺตาทานํ, โน ภูตนฺ”ติ. น ตํ อุปาลิ อตฺตาทานํ อาทาตพฺพํ.}

\prose{สเจ ปนุปาลิ ภิกฺขุ ปจฺจเวกฺขมาโน เอวํ ชานาติ “ภูตํ อิทํ อตฺตาทานํ, โน อภูตนฺ”ติ. เตนุปาลิ ภิกฺขุนา อุตฺตริ ปจฺจเวกฺขิตพฺพํ “ยํ โข อหํ อิมํ อตฺตาทานํ อาทาตุกาโม, อตฺถสํหิตํ นุ โข อิทํ อตฺตาทานํ, อุทาหุ โน”ติ. สเจ อุปาลิ ภิกฺขุ ปจฺจเวกฺขมาโน เอวํ \csromanfolio{332}ชานาติ “อนตฺถ\-สํหิตํ อิทํ อตฺตาทานํ, โน อตฺถสํหิตนฺ”ติ. น ตํ อุปาลิ อตฺตาทานํ อาทาตพฺพํ.}

\prose{สเจ ปนุปาลิ ภิกฺขุ ปจฺจเวกฺขมาโน เอวํ ชานาติ “อตฺถสํหิตํ อิทํ อตฺตาทานํ, โน อนตฺถสํหิตนฺ”ติ. เตนุปาลิ ภิกฺขุนา อุตฺตริ ปจฺจเวกฺขิตพฺพํ “อิมํ โข อหํ อตฺตาทานํ อาทิยมาโน ลภิสฺสามิ สนฺทิฏฺเฐ สมฺภตฺเต ภิกฺขู ธมฺมโต วินยโต ปกฺเข, อุทาหุ โน”ติ. สเจ อุปาลิ ภิกฺขุ ปจฺจเวกฺขมาโน เอวํ ชานาติ “อิมํ โข อหํ อตฺตาทานํ อาทิยมาโน น ลภิสฺสามิ สนฺทิฏฺเฐ สมฺภตฺเต ภิกฺขู ธมฺมโต วินยโต ปกฺเข”ติ. น ตํ อุปาลิ อตฺตาทานํ อาทาตพฺพํ.}

\prose{สเจ ปนุปาลิ ภิกฺขุ ปจฺจเวกฺขมาโน เอวํ ชานาติ “อิทํ โข อหํ อตฺตาทานํ อาทิยมาโน ลภิสฺสามิ สนฺทิฏฺเฐ สมฺภตฺเต ภิกฺขู ธมฺมโต วินยโต ปกฺเข”ติ. เตนุปาลิ ภิกฺขุนา อุตฺตริ ปจฺจเวกฺขิตพฺพํ “อิมํ โข เม อตฺตาทานํ อาทิยโต ภวิสฺสติ สํฆสฺส ตโตนิทานํ ภณฺฑนํ กลโห วิคฺคโห วิวาโท สํฆเภโท สํฆราชิ สํฆววตฺถานํ สํฆนานากรณํ, อุทาหุ โน”ติ. สเจ อุปาลิ ภิกฺขุ ปจฺจเวกฺขมาโน เอวํ ชานาติ “อิมํ โข เม อตฺตาทานํ อาทิยโต ภวิสฺสติ สํฆสฺส ตโตนิทานํ ภณฺฑนํ กลโห วิคฺคโห วิวาโท สํฆเภโท สํฆราชิ สํฆววตฺถานํ สํฆนานากรณนฺ”ติ. น ตํ อุปาลิ อตฺตาทานํ อาทาตพฺพํ.}

\prose{สเจ ปนุปาลิ ภิกฺขุ ปจฺจเวกฺขมาโน เอวํ ชานาติ “อิมํ โข เม อตฺตาทานํ อาทิยโต น ภวิสฺสติ สํฆสฺส ตโตนิทานํ ภณฺฑนํ กลโห วิคฺคโห วิวาโท สํฆเภโท สํฆราชิ สํฆววตฺถานํ สํฆนานากรณนฺ”ติ. ตํ อาทาตพฺพํ อุปาลิ อตฺตาทานํ. เอวํ ปญฺจ\-งฺค\-สมนฺนาคตํ โข อุปาลิ อตฺตาทานํ อาทินฺนํ ปจฺฉาปิ อวิปฺปฏิสารกรํ ภวิสฺสตีติ.}

\proseitem{441}{กติหิ นุ โข ภนฺเต องฺเคหิ สมนฺนาคโต ภิกฺขุ อธิกรณ\-ชาตานํ ภิกฺขุนํ พหูปกาโร โหตีติ. ปญฺจหุปาลิ องฺเคหิ สมนฺนาคโต ภิกฺขุ อธิกรณ\-ชาตานํ ภิกฺขูนํ พหูปกาโร โหติ. กตเมหิ ปญฺจหิ, สีลวา โหติ ปาติโมกฺข\-สํวร\-สํวุโต วิหรติ อาจารโคจร\-สมฺปนฺโน อณุมตฺเตสุ วชฺเชสุ ภยทสฺสาวี สมาทาย \csromanfolio{333}สิกฺขติ สิกฺขาปเทสุ, พหุสฺสุโต โหติ สุตธโร สุตสนฺนิจโย, เย เต ธมฺมา อาทิกลฺยาณา มชฺเฌกลฺยาณา ปริโยสาน\-กลฺยาณา สาตฺถํ สพฺยญฺชนํ เกวล\-ปริปุณฺณํ ปริสุทฺธํ พฺรหฺมจริยํ อภิวทนฺติ, ตถารูปสฺส ธมฺมา พหุสฺสุตา โหนฺติ ธาตา วจสา ปริจิตา มนสา\-นุเปกฺขิตา ทิฏฺฐิยา สุปฺปฏิวิทฺธา, อุภยานิ โข ปนสฺส ปาติโมกฺขานิ วิตฺถาเรน สฺวาคตานิ โหนฺติ สุวิภตฺตานิ สุปฺปวตฺตีนิ สุวินิจฺฉิตานิ สุตฺตโส อนุพฺยญฺชนโส, วินเย โข ปน ฐิโต โหติ อสํหีโร. ปฏิพโล โหติ อุโภ อตฺถ\-ปจฺจตฺถิเก อสฺสาเสตุํ สญฺญาเปตุํ นิชฺฌาเปตุํ เปกฺเขตุํ ปสาเทตุํ. อิเมหิ โข อุปาลิ ปญฺจหงฺเคหิ สมนฺนาคโต ภิกฺขุ อธิกรณ\-ชาตานํ ภิกฺขูนํ พหูปกาโร โหติ.}

\prose{อปเรหิปิ อุปาลิ ปญฺจหงฺเคหิ สมนฺนาคโต ภิกฺขุ อธิกรณ\-ชาตานํ ภิกฺขูนํ พหูปกาโร โหติ. กตเมหิ ปญฺจหิ, ปริสุทฺธ\-กาย\-สมาจาโร โหติ, ปริสุทฺธ\-วจี\-สมาจาโร โหติ, ปริสุทฺธชีโว โหติ, ปณฺฑิโต โหติ พฺยตฺโต, ปฏิพโล อนุยุญฺชิยมาโน อนุโยคํ ทาตุํ. อิเมหิ โข อุปาลิ ปญฺจหงฺเคหิ สมนฺนาคโต ภิกฺขุ อธิกรณ\-ชาตานํ ภิกฺขูนํ พหูปกาโร โหติ.}

\prose{อปเรหิปิ อุปาลิ ปญฺจหงฺเคหิ สมนฺนาคโต ภิกฺขุ อธิกรณ\-ชาตานํ ภิกฺขูนํ พหูปกาโร โหติ. กตเมหิ ปญฺจหิ, วตฺถุํ ชานาติ, นิทานํ ชานาติ, ปญฺญตฺติํ ชานาติ, ปทปจฺจา\-ภฏฺฐํ ชานาติ, อนุสนฺธิ\-วจนปถํ ชานาติ. อิเมหิ โข อุปาลิ ปญฺจหงฺเคหิ สมนฺนาคโต ภิกฺขุ อธิกรณ\-ชาตานํ ภิกฺขูนํ พหูปกาโร โหตีติ.}

\proseitem{442}{กติหิ นุ โข ภนฺเต องฺเคหิ สมนฺนาคเตน ภิกฺขุนา นานุ\-ยุญฺชิ\-ตพฺพนฺติ. ปญฺจหุปาลิ องฺเคหิ สมนฺนาคเตน ภิกฺขุนา นานุ\-ยุญฺชิ\-ตพฺพํ. กตเมหิ ปญฺจหิ, สุตฺตํ น ชานาติ, สุตฺตานุโลมํ น ชานาติ, วินยํ น ชานาติ, วินยานุโลมํ น ชานาติ, น จ ฐานาฐาน\-กุสโล โหติ. อิเมหิ โข อุปาลิ ปญฺจหงฺเคหิ สมนฺนาคเตน ภิกฺขุนา นานุ\-ยุญฺชิ\-ตพฺพํ.}

\prose{ปญฺจหุปาลิ องฺเคหิ สมนฺนาคเตน ภิกฺขุนา อนุยุญฺชิตพฺพํ. กตเมหิ ปญฺจหิ, สุตฺตํ ชานาติ, สุตฺตานุโลมํ ชานาติ, วินยํ ชานาติ, \csromanfolio{334}วินยานุโลมํ ชานาติ, ฐานาฐาน\-กุสโล จ โหติ. อิเมหิ โข อุปาลิ ปญฺจหงฺเคหิ สมนฺนาคเตน ภิกฺขุนา อนุยุญฺชิตพฺพํ.}

\prose{อปเรหิปิ อุปาลิ ปญฺจหงฺเคหิ สมนฺนาคเตน ภิกฺขุนา นานุ\-ยุญฺชิ\-ตพฺพํ. กตเมหิ ปญฺจหิ, ธมฺมํ น ชานาติ, ธมฺมานุโลมํ น ชานาติ, วินยํ น ชานาติ, วินยานุโลมํ น ชานาติ. น จ ปุพฺพาปร\-กุสโล โหติ. อิเมหิ โข อุปาลิ ปญฺจหงฺเคหิ สมนฺนาคเตน ภิกฺขุนา นานุ\-ยุญฺชิ\-ตพฺพํ.}

\prose{ปญฺจหุปาลิ องฺเคหิ สมนฺนาคเตน ภิกฺขุนา อนุยุญฺชิตพฺพํ. กตเมหิ ปญฺจหิ, ธมฺมํ ชานาติ, ธมฺมานุโลมํ ชานาติ, วินยํ ชานาติ, วินยานุโลมํ ชานาติ, ปุพฺพาปร\-กุสโล จ โหติ. อิเมหิ โข อุปาลิ ปญฺจหงฺเคหิ สมนฺนาคเตน ภิกฺขุนา อนุยุญฺชิตพฺพํ.}

\prose{อปเรหิปิ อุปาลิ ปญฺจหงฺเคหิ สมนฺนาคเตน ภิกฺขุนา นานุ\-ยุญฺชิ\-ตพฺพํ. กตเมหิ ปญฺจหิ, วตฺถุํ น ชานาติ, นิทานํ น ชานาติ, ปญฺญตฺติํ น ชานาติ, ปทปจฺจา\-ภฏฺฐํ น ชานาติ, อนุสนฺธิ\-วจนปถํ น ชานาติ. อิเมหิ โข อุปาลิ ปญฺจหงฺเคหิ สมนฺนาคเตน ภิกฺขุนา นานุ\-ยุญฺชิ\-ตพฺพํ.}

\prose{ปญฺจหุปาลิ องฺเคหิ สมนฺนาคเตน ภิกฺขุนา อนุยุญฺชิตพฺพํ. กตเมหิ ปญฺจหิ, วตฺถุํ ชานาติ, นิทานํ ชานาติ, ปญฺญตฺติํ ชานาติ. ปทปจฺจภฏฺฐํ ชานาติ, อนุสนฺธิ\-วจนปถํ ชานาติ. อิเมหิ โข อุปาลิ ปญฺจหงฺเคหิ สมนฺนาคเตน ภิกฺขุนา อนุยุญฺชิตพฺพํ.}

\prose{อปเรหิปิ อุปาลิ ปญฺจหงฺเคหิ สมนฺนาคเตน ภิกฺขุนา นานุ\-ยุญฺชิ\-ตพฺพํ. กตเมหิ ปญฺจหิ, อาปตฺติํ น ชานาติ, อาปตฺติ\-สมุฏฺฐานํ น ชานาติ, อาปตฺติยา ปโยคํ น ชานาติ, อาปตฺติยา วูปสมํ น ชานาติ, อาปตฺติยา น วินิจฺฉย\-กุสโล โหติ. อิเมหิ โข อุปาลิ ปญฺจหงฺเคหิ สมนฺนาคเตน ภิกฺขุนา นานุ\-ยุญฺชิ\-ตพฺพํ.}

\prose{ปญฺจหุปาลิ องฺเคหิ สมนฺนาคเตน ภิกฺขุนา อนุยุญฺชิตพฺพํ. กตเมหิ ปญฺจหิ, อาปตฺติํ ชานาติ, อาปตฺติ\-สมุฏฺฐานํ ชานาติ, อาปตฺติยา ปโยคํ ชานาติ, อาปตฺติยา วูปสมํ ชานาติ, อาปตฺติยา วินิจฺฉย\-กุสโล โหติ. อิเมหิ โข อุปาลิ ปญฺจหงฺเคหิ สมนฺนาคเตน ภิกฺขุนา อนุยุญฺชิตพฺพํ.}

\prose{\csromanfolio{335}อปเรหิปิ อุปาลิ ปญฺจหงฺเคหิ สมนฺนาคเตน ภิกฺขุนา นานุ\-ยุญฺชิ\-ตพฺพํ. กตเมหิ ปญฺจหิ, อธิกรณํ น ชานาติ, อธิกรณ\-สมุฏฺฐานํ น ชานาติ, อธิกรณสฺส ปโยคํ น ชานาติ, อธิกรณสฺส วูปสมํ น ชานาติ, อธิกรณสฺส น วินิจฺฉย\-กุสโล โหติ. อิเมหิ โข อุปาลิ ปญฺจหงฺเคหิ สมนฺนาคเตน ภิกฺขุนา นานุ\-ยุญฺชิ\-ตพฺพํ.}

\prose{ปญฺจหุปาลิ องฺเคหิ สมนฺนาคเตน ภิกฺขุนา อนุยุญฺชิตพฺพํ. กตเมหิ ปญฺจหิ, อธิกรณํ ชานาติ, อธิกรณ\-สมุฏฺฐานํ ชานาติ, อธิกรณสฺส ปโยคํ ชานาติ, อธิกรณสฺส วูปสมํ ชานาติ, อธิกรณสฺส วินิจฺฉย\-กุสโล โหติ. อิเมหิ โข อุปาลิ ปญฺจหงฺเคหิ สมนฺนาคเตน ภิกฺขุนา อนุยุญฺชิ\-ตพฺพนฺติ.}

\vspace{\tipitakaparskip}
\csromancenter{อตฺตาทานวคฺโค นิฏฺฐิโต ปญฺจโม.}
\csromancenter{ตสฺสุทฺทานํ.}
\setlength{\csromangathapagewidth}{0pt}
\setlength{\csromangathawakwidth}{0pt}
\csromangathameasuregroup{ปริสุทฺธญฺจ กาเลน, \\ อตฺตาทานํ อธิกรณํ,}{\csromangathabat{ปริสุทฺธญฺจ กาเลน,}{การุญฺญํ โอกาเสน จ.} \\ \csromangathabat{อตฺตาทานํ อธิกรณํ,}{อปเรหิปิ วตฺถุญฺจ.}}
\csromangathasetleft{ปริสุทฺธญฺจ กาเลน, \\ อตฺตาทานํ อธิกรณํ,}
\csromangathagroup{\csromangathastanza{\csromangathabat{ปริสุทฺธญฺจ กาเลน,}{การุญฺญํ โอกาเสน จ.} \\ \csromangathabat{อตฺตาทานํ อธิกรณํ,}{อปเรหิปิ วตฺถุญฺจ.}}}

\prose{สุตฺตํ ธมฺมํ ปุน วตฺถุญฺจ, อาปตฺติ อธิกรเณน จาติ.}
\csromansectionrule

\csromanmatikaanchor{mtk.225}
\csromantocmarkref{subparagraph}{6. ธุตงฺควคฺค}{mtk.225}
\bookmark[dest={mtk.225},level=5]{6. ธุตงฺควคฺค}
\csromanheader{5}{6. \textbf{ธุตงฺควคฺค}}
\proseitem{443}{กติ นุ โข ภนฺเต อารญฺญิกาติ. ปญฺจิเม อุปาลิ อารญฺญิกา.\csromansymbolfootnote{*}{อํ 2. 193; วิ 5. 233 ปิฏฺเฐสุปิ.}กตเม ปญฺจ, มนฺทตฺตา โมมูหตฺตา อารญฺญิโก โหติ, ปาปิจฺโฉ อิจฺฉาปกโต อารญฺญิโก โหติ, อุมฺมาทา จิตฺตกฺเขปา อารญฺญิโก โหติ, “วณฺณิตํ พุทฺเธหิ พุทฺธสาวเกหี”ติ อารญฺญิโก โหติ, อปิ จ อปฺปิจฺฉญฺเญว นิสฺสาย สนฺตุฏฺฐิญฺเญว นิสฺสาย สลฺเลขญฺเญว นิสฺสาย ปวิเวกญฺเญว นิสฺสาย อิทมตฺถิตญฺเญว นิสฺสาย อารญฺญิโก โหติ. อิเม โข อุปาลิ ปญฺจ อารญฺญิกาติ.}

\prose{กติ นุ โข ภนฺเต ปิณฺฑปาติกาติ ฯเปฯ. กติ นุ โข ภนฺเต ปํสุกูลิกาติ ฯเปฯ. กติ นุ โข ภนฺเต รุกฺขมูลิกาติ ฯเปฯ. กติ นุ โข ภนฺเต โสสานิกาติ ฯเปฯ. กติ นุ โข ภนฺเต อพฺโภกาสิกาติ ฯเปฯ. \csromanfolio{336}กติ นุ โข ภนฺติ เตจีวริกาติ ฯเปฯ. กติ นุ โข ภนฺเต สปทานจาริกาติ ฯเปฯ. กติ นุ โข ภนฺเต เนสชฺชิกาติ ฯเปฯ. กติ นุ โข ภนฺเต ยถา\-สนฺถติ\-กาติ ฯเปฯ. กติ นุ โข ภนฺเต เอกาสนิกาติ ฯเปฯ. กติ นุ โข ภนฺเต ขลุปจฺฉาภตฺติ\-กาติ ฯเปฯ. กติ นุ โข ภนฺเต ปตฺตปิณฺฑิกาติ. ปญฺจิเม อุปาลิ ปตฺตปิณฺฑิกา. กตเม ปญฺจ, มนฺทตฺตา โมมูหตฺตา ปตฺตปิณฺฑิโก โหติ, ปาปิจฺโฉ อิจฺฉาปกโต ปตฺตปิณฺฑิโก โหติ, อุมฺมาทา จิตฺตกฺเขปา ปตฺตปิณฺฑิโก โหติ, “วณฺณิตํ พุทฺเธหิ พุทฺธสาวเกหี”ติ ปตฺตปิณฺฑิโก โหติ, อปิ จ อปฺปิจฺฉญฺเญว นิสฺสาย สนฺตุฏฺฐิญฺเญว นิสฺสาย สลฺเลขญฺเญว นิสฺสาย ปวิเวกญฺเญว นิสฺสาย อิทมตฺถิตญฺเญว นิสฺสาย ปตฺตปิณฺฑิโก โหติ. อิเม โข อุปาลิ ปญฺจ ปตฺตปิณฺฑิกาติ.}

\vspace{\tipitakaparskip}
\csromancenter{ธุตงฺเควคฺโค นิฏฺฐิโต ฉฏฺโฐ.}
\csromancenter{ตสฺสุทฺทานํ.}
\setlength{\csromangathapagewidth}{0pt}
\setlength{\csromangathawakwidth}{0pt}
\csromangathameasuregroup{อารญฺญิโก ปิณฺฑิปํสุ, \\ อพฺโภ เตจีวรญฺเจว,}{\csromangathabat{อารญฺญิโก ปิณฺฑิปํสุ,}{รุกฺขสุสานปญฺจมํ.} \\ \csromangathabat{อพฺโภ เตจีวรญฺเจว,}{สปทานเนสชฺชิกา.}}
\csromangathasetleft{อารญฺญิโก ปิณฺฑิปํสุ, \\ อพฺโภ เตจีวรญฺเจว,}
\csromangathagroup{\csromangathastanza{\csromangathabat{อารญฺญิโก ปิณฺฑิปํสุ,}{รุกฺขสุสานปญฺจมํ.} \\ \csromangathabat{อพฺโภ เตจีวรญฺเจว,}{สปทานเนสชฺชิกา.}}}

\vspace{\tipitakaparskip}
\csromancenter{สนฺถเต\-กาสน\-ญฺเจว, ขลุปจฺฉา ปตฺตปิณฺฑิกาติ.}
\csromansectionrule
\csromanmatikaanchor{mtk.226}
\csromantocmarkref{subparagraph}{7. มุสาวาทวคฺค}{mtk.226}
\bookmark[dest={mtk.226},level=5]{7. มุสาวาทวคฺค}
\csromanheader{5}{7. \textbf{มุสาวาทวคฺค}}
\proseitem{444}{กติ นุ โข ภนฺเต มุสาวาทาติ. ปญฺจิเม อุปาลิ มุสาวาทา. กตเม ปญฺจ, อตฺถิ มุสาวาโท ปาราชิกคามี, อตฺถิ มุสาวาโท สํฆาทิเสสคามี, อตฺถิ มุสาวาโท ถุลฺลจฺจย\-คามี, อตฺถิ มุสาวาโท ปาจิตฺติยคามี, อตฺถิ มุสาวาโท ทุกฺกฏคามี. อิเม โข อุปาลิ ปญฺจ มุสาวาทาติ.}

\proseitem{445}{กติหิ นุ โข ภนฺเต องฺเคหิ สมนฺนาคตสฺส ภิกฺขุโน สํฆมชฺเฌ อุโปสถํ วา ปวารณํ วา ฐเปนฺตสฺส “อลํ ภิกฺขุ, มา ภณฺฑนํ มา กลหํ มา วิคฺคหํ มา วิวาทนฺ”ติ โอมทฺทิตฺวา สํเฆน อุโปสโถ วา ปวารณา วา กาตพฺพาติ. ปญฺจหุปาลิ องฺเคหิ สมนฺนาคตสฺส ภิกฺขุโน \csromanfolio{337}สํฆมชฺเฌ อุโปสถํ วา ปวารณํ วา ฐเปนฺตสฺส “อลํ ภิกฺขุ, มา ภณฺฑนํ มา กลหํ มา วิคฺคหํ มา วิวาทนฺ”ติ โอมทฺทิตฺวา สํเฆน อุโปสโถ วา ปวารณา วา กาตพฺพา. กตเมหิ ปญฺจหิ, อลชฺชี จ โหติ, พาโล จ, อปกตตฺโต จ, จาวนาธิปฺปาโย วตฺตา โหติ, โน วุฏฺฐานา\-ธิปฺปาโย. อิเมหิ โข อุปาลิ ปญฺจหงฺเคหิ สมนฺนาคตสฺส ภิกฺขุโน สํฆมชฺเฌ อุโปสถํ วา ปวารณํ วา ฐเปนฺตสฺส “อลํ ภิกฺขุ, มา ภณฺฑนํ มา กลหํ มา วิคฺคหํ มา วิวาทนฺ”ติ โอมทฺทิตฺวา สํเฆน อุโปสโถ วา ปวารณา วา กาตพฺพา.}

\prose{อปเรหิปิ อุปาลิ ปญฺจหงฺเคหิ สมนฺนาคตสฺส ภิกฺขุโน สํฆมชฺเฌ อุโปสถํ วา ปวารณํ วา ฐเปนฺตสฺส “อลํ ภิกฺขุ, มา ภณฺฑนํ มา กลหํ มา วิคฺคหํ มา วิวาทนฺ”ติ โอมทฺทิตฺวา สํเฆน อุโปสโถ วา ปวารณา วา กาตพฺพา. กตเมหิ ปญฺจหิ, อปริสุทฺธ\-กายสมาจาโร โหติ, อปริสุทฺธ\-วจีสมาจาโร โหติ, อปริสุทฺธา\-ชีโว โหติ, พาโล โหติ อพฺยตฺโต, ภณฺฑนการโก โหติ กลหการโก. อิเมหิ โข อุปาลิ ปญฺจหงฺเคหิ สมนฺนาคตสฺส ภิกฺขุโน สํฆมชฺเฌ อุโปสถํ วา ปวารณํ วา ฐเปนฺตสฺส “อลํ ภิกฺขุ, มา ภณฺฑนํ มา กลหํ มา วิคฺคหํ มา วิวาทนฺ”ติ โอมทฺทิตฺวา สํเฆน อุโปสโถ วา ปวารณา วา กาตพฺพาติ.}

\proseitem{446}{กติหิ นุ โข ภนฺเต องฺเคหิ สมนฺนาคตสฺส ภิกฺขุโน อนุโยโค น ทาตพฺโพติ. ปญฺจหุปาลิ องฺเคหิ สมนฺนาคตสฺส ภิกฺขุโน อนุโยโค น ทาตพฺโพ. กตเมหิ ปญฺจหิ, อาปตฺตานาปตฺติํ น ชานติ, ลหุกครุกํ อาปตฺติํ น ชานาติ, สาวเสสา\-นวเสสํ อาปตฺติํ น ชานาติ, ทุฏฺฐุลฺลา\-ทุฏฺฐุลฺลํ อาปตฺติํ น ชานาติ, สปฺปฏิกมฺมา\-ปฏิกมฺมํ อาปตฺติํ น ชานาติ. อิเมหิ โข อุปาลิ ปญฺจหงฺเคหิ สมนฺนาคตสฺส ภิกฺขุโน อนุโยโค น ทาตพฺโพ.}

\prose{ปญฺจหุปาลิ องฺเคหิ สมนฺนาคตสฺส ภิกฺขุโน อนุโยโค ทาตพฺโพ. กตเมหิ ปญฺจหิ, อาปตฺตานาปตฺติํ ชานาติ, ลหุกครุกํ อาปตฺติํ ชานาติ, สาวเสสา\-นวเสสํ อาปตฺติํ ชานาติ, ทุฏฺฐุลฺลา\-ทุฏฺฐุลฺลํ อาปตฺติํ ชานาติ, สปฺปฏิกมฺมา\-ปฏิกมฺมํ อาปตฺติํ ชานาติ. อิเมหิ โข อุปาลิ ปญฺจหงฺเคหิ สมนฺนาคตสฺส ภิกฺขุโน อนุโยโค ทาตพฺโพติ.}

\proseitem{447}{\csromanfolio{338}กติหิ นุ โข ภนฺเต อากาเรหิ ภิกฺขุ อาปตฺติํ อาปชฺชตีติ. ปญฺจหุปาลิ อากาเรหิ ภิกฺขุ อาปตฺติํ อาปชฺชติ. กตเมหิ ปญฺจหิ, อลชฺชิตา, อญฺญาณตา, กุกฺกุจฺจ\-ปกตตา, อกปฺปิเย กปฺปิยสญฺญิตา, กปฺปิเย อกปฺปิย\-สญฺญิตา. อิเมหิ โข อุปาลิ ปญฺจหากาเรหิ ภิกฺขุ อาปตฺติํ อาปชฺชติ.}

\prose{อปเรหิปิ อุปาลิ ปญฺจหากาเรหิ ภิกฺขุ อาปตฺติํ อาปชฺชติ. กตเมหิ ปญฺจหิ, อทสฺสเนน, อสฺสวเนน, ปสุตฺตกตา, ตถาสญฺญี, สติสมฺโมสา. อิเมหิ โข อุปาลิ ปญฺจหากาเรหิ ภิกฺขุ อาปตฺติํ อาปชฺชตีติ.}

\proseitem{448}{กติ นุ โข ภนฺเต เวราติ. ปญฺจิเม อุปาลิ เวรา. กตเม ปญฺจ, ปาณาติปาโต, อทินฺนาทานํ, กาเมสุ\-มิจฺฉาจาโร, มุสาวาโท, สุรา\-เมรย\-มชฺช\-ปฺปมาทฏฺฐานํ. อิเม โข อุปาลิ ปญฺจ เวราติ.}

\prose{กติ นุ โข ภนฺเต เวรมณิโยติ. ปญฺจิมา อุปาลิ เวรมณิโย. กตมา ปญฺจ, ปาณาติปาตา เวรมณี\csromansharedfootnote{4e794017f3930854}{เวรมณิ(ก)}, อทินฺนาทานา เวรมณี, กาเมสุ\-มิจฺฉาจารา เวรมณี, มุสาวาทา เวรมณี, สุรา\-เมรย\-มชฺช\-ปฺปมาทฏฺฐานา เวรมณี. อิมา โข อุปาลิ ปญฺจ เวรมณิโยติ.}

\proseitem{449}{กติ นุ โข ภนฺเต พฺยสนานีติ. ปญฺจิมานิ อุปาลิ พฺยสนานิ.\csromansymbolfootnote{*}{วิ 5. 230; อํ 2. 129 ปิฏฺฐาทีสุปิ.}กตมานิ ปญฺจ, ญาติพฺยาสนํ, โภคพฺยสนํ, โรคพฺยสนํ, สีลพฺยสนํ, ทิฏฺฐิพฺยสนํ. อิมานิ โข อุปาลิ ปญฺจ พฺยสนานีติ.}

\prose{กติ นุ โข ภนฺเต สมฺปทาติ. ปญฺจิมา อุปาลิ สมฺปทา.\csromansymbolmark{*}กตมา ปญฺจ, ญาติสมฺปทา, โภคสมฺปทา, อาโรคฺยสมฺปทา, สีลสมฺปทา, ทิฏฺฐิสมฺปทา. อิมา โข อุปาลิ ปญฺจ สมฺปทาติ.}

\vspace{\tipitakaparskip}
\csromancenter{มุสาวาทวคฺโค นิฏฺฐิโต สตฺตโม.}
\csromancenter{ตสฺสุทฺทานํ.}
\setlength{\csromangathapagewidth}{0pt}
\setlength{\csromangathawakwidth}{0pt}
\csromangathameasuregroup{มุสาวาโท จ โอมทฺทิ, \\ อาปตฺติญฺจ อปเรหิ, \\ พฺยสนํ สมฺปทา เจว,}{\csromangathabat{มุสาวาโท จ โอมทฺทิ,}{อปเรหิ อนุโยโค.} \\ \csromangathabat{อาปตฺติญฺจ อปเรหิ,}{เวรา เวรมณีปิ จ.} \\ \csromangathabat{พฺยสนํ สมฺปทา เจว,}{สตฺตโม วคฺคสงฺคโหติ.}}
\csromangathasetleft{มุสาวาโท จ โอมทฺทิ, \\ อาปตฺติญฺจ อปเรหิ, \\ พฺยสนํ สมฺปทา เจว,}
\csromangathagroup{\csromangathastanza{\csromangathabat{มุสาวาโท จ โอมทฺทิ,}{อปเรหิ อนุโยโค.} \\ \csromangathabat{อาปตฺติญฺจ อปเรหิ,}{เวรา เวรมณีปิ จ.} \\ \csromangathabat{พฺยสนํ สมฺปทา เจว,}{สตฺตโม วคฺคสงฺคโหติ.}}}

\csromanmatikaanchor{mtk.227}
\csromantocmarkref{subparagraph}{8. ภิกฺขุโนวาทวคฺค}{mtk.227}
\bookmark[dest={mtk.227},level=5]{8. ภิกฺขุโนวาทวคฺค}
\csromanheader{5}{8. ภิกฺขุโน\-วาทวคฺค}
\proseitem{450}{\csromanfolio{339}กติหิ นุ โข ภนฺเต องฺเคหิ สมนฺนาคตสฺส ภิกฺขุโน ภิกฺขุนิสํเฆเนว กมฺมํ กาตพฺพนฺติ. ปญฺจหุปาลิ องฺเคหิ สมนฺนาคตสฺส ภิกฺขุโน ภิกฺขุนิสํเฆเนว กมฺมํ กาตพฺพํ, อวนฺทิโย โส ภิกฺขุ ภิกฺขุนิ\-สํเฆน. กตเมหิ ปญฺจหิ, วิวริตฺวา กายํ ภิกฺขุนีนํ ทสฺเสติ, อูรุํ ทสฺเสติ, องฺคชาตํ ทสฺเสติ, อุโภ อํสกูเฏ ทสฺเสติ, โอภาสติ คิหี สมฺปโยเชติ. อิเมหิ โข อุปาลิ ปญฺจหงฺเคหิ สมนฺนาคตสฺส ภิกฺขุโน ภิกฺขุนิสํเฆเนว กมฺมํ กาตพฺพํ, อวนฺทิโย โส ภิกฺขุ ภิกฺขุนิ\-สํเฆน.}

\prose{อปเรหิปิ อุปาลิ ปญฺจหงฺเคหิ สมนฺนาคตสฺส ภิกฺขุโน ภิกฺขุนิสํเฆเนว กมฺมํ กาตพฺพํ, อวนฺทิโย โส ภิกฺขุ ภิกฺขุนิ\-สํเฆน. กตเมหิ ปญฺจหิ, ภิกฺขุนีนํ อลาภาย ปริสกฺกติ, ภิกฺขุนีนํ อนตฺถาย ปริสกฺกติ, ภิกฺขุนีนํ อวาสาย ปริสกฺกติ, ภิกฺขุนิโย อกฺโกสติ ปริภาสติ, ภิกฺขู ภิกฺขุนีหิ เภเทติ. อิเมหิ โข อุปาลิ ปญฺจหงฺเคหิ สมนฺนาคตสฺส ภิกฺขุโน ภิกฺขุนิสํเฆเนว กมฺมํ กาตพฺพํ, อวนฺทิโย โส ภิกฺขุ ภิกฺขุนิ\-สํเฆน.}

\prose{อปเรหิปิ อุปาลิ ปญฺจหงฺเคหิ สมนฺนาคตสฺส ภิกฺขุโน ภิกฺขุนิสํเฆเนว กมฺมํ กาตพฺพํ, อวนฺทิโย โส ภิกฺขุ ภิกฺขุนิ\-สํเฆน. กตเมหิ ปญฺจหิ, ภิกฺขุนีนํ อลาภาย ปริสกฺกติ, ภิกฺขุนีนํ อนตฺถาย ปริสกฺกติ, ภิกฺขุนีนํ อวาสาย ปริสกฺกติ, ภิกฺขุนิโย อกฺโกสติ ปริภาสติ, ภิกฺขู ภิกฺขุนีหิ สมฺปโยเชติ. อิเมหิ โข อุปาลิ ปญฺจหงฺเคหิ สมนฺนาคตสฺส ภิกฺขุโน ภิกฺขุนิสํเฆเนว กมฺมํ กาตพฺพํ, อวนฺทิโย โส ภิกฺขุ ภิกฺขุนิ\-สํเฆนาติ.}

\proseitem{451}{กติหิ นุ โข ภนฺเต องฺเคหิ สมนฺนาคตาย ภิกฺขุนิยา กมฺมํ กาตพฺพนฺติ. ปญฺจหุปาลิ องฺเคหิ สมนฺนาคตาย ภิกฺขุนิยา กมฺมํ กาตพฺพํ. กตเมหิ ปญฺจหิ, วิวริตฺวา กายํ ภิกฺขูนํ ทสฺเสติ, อูรุํ ทสฺเสติ, องฺคชาตํ ทสฺเสติ, อุโภ อํสกูเฏ ทสฺเสติ, โอภาสติ คิหี สมฺปโยเชติ. อิเมหิ โข อุปาลิ ปญฺจหงฺเคหิ สมนฺนาคตาย ภิกฺขุนิยา กมฺมํ กาตพฺพํ.}

\prose{\csromanfolio{340}อปเรหิปิ อุปาลิ ปญฺจหงฺเคหิ สมนฺนาคตาย ภิกฺขุนิยา กมฺมํ กาตพฺพํ. กตเมหิ ปญฺจหิ, ภิกฺขูนํ อลาภาย ปริสกฺกติ, ภิกฺขูนํ อนตฺถาย ปริสกฺกติ, ภิกฺขูนํ อวาสาย ปริสกฺกติ, ภิกฺขู อกฺโกสติ ปริภาสติ, ภิกฺขุนิโย ภิกฺขูหิ เภเทติ. อิเมหิ โข อุปาลิ ปญฺจหงฺเคหิ สมนฺนาคตาย ภิกฺขุนิยา กมฺมํ กาตพฺพํ.}

\prose{อปเรหิปิ อุปาลิ ปญฺจหงฺเคหิ สมนฺนาคตาย ภิกฺขุนิยา กมฺมํ กาตพฺพํ. กตเมหิ ปญฺจหิ, ภิกฺขูนํ อลาภาย ปริสกฺกติ, ภิกฺขูนํ อนตฺถาย ปริสกฺกติ, ภิกฺขูนํ อวาสาย ปริสกฺกติ, ภิกฺขู อกฺโกสติ ปริภาสติ, ภิกฺขุนิโย ภิกฺขูหิ สมฺปโยเชติ. อิเมหิ โข อุปาลิ ปญฺจหงฺเคหิ สมนฺนาคตาย ภิกฺขุนิยา กมฺมํ กาตพฺพนฺติ.}

\proseitem{452}{กติหิ นุ โข ภนฺเต องฺเคหิ สมนฺนาคเตน ภิกฺขุนา ภิกฺขุนีนํ โอวาโท น ฐเปตพฺโพติ. ปญฺจหุปาลิ องฺเคหิ สมนฺนาคเตน ภิกฺขุนา ภิกฺขุนีนํ โอวาโท น ฐเปตพฺโพ. กตเมหิ ปญฺจหิ, อลชฺชี จ โหติ, พาโล จ, อปกตตฺโต จ, จาวนาธิปฺปาโย วตฺตา โหติ, โน วุฏฺฐานา\-ธิปฺปาโย. อิเมหิ โข อุปาลิ ปญฺจหงฺเคหิ สมนฺนาคเตน ภิกฺขุนา ภิกฺขุนีนํ โอวาโท น ฐเปตพฺโพ.}

\prose{อปเรหิปิ อุปาลิ ปญฺจหงฺเคหิ สมนฺนาคเตน ภิกฺขุนา ภิกฺขุนีนํ โอวาโท น ฐเปตพฺโพ. กตเมหิ ปญฺจหิ, อปริสุทฺธ\-กายสมาจาโร โหติ, อปริสุทฺธ\-วจีสมาจาโร โหติ, อปริสุทฺธา\-ชีโว โหติ, พาโล โหติ อพฺยตฺโต, น ปฏิพโล อนุยุญฺชิยมาโน อนุโยคํ ทาตุํ. อิเมหิ โข อุปาลิ ปญฺจหงฺเคหิ สมนฺนาคเตน ภิกฺขุนา ภิกฺขุนีนํ โอวาโท น ฐเปตพฺโพ.}

\prose{อปเรหิปิ อุปาลิ ปญฺจหงฺเคหิ สมนฺนาคเตน ภิกฺขุนา ภิกฺขุนีนํ โอวาโท น ฐเปตพฺโพ. กตเมหิ ปญฺจหิ, กายิเกน อนาจาเรน สมนฺนาคโต โหติ, วาจสิเกน อนาจาเรน สมนฺนาคโต โหติ, กายิก\-วาจสิเกน อนาจาเรน สมนฺนาคโต โหติ, ภิกฺขุนีนํ อกฺโกสก\-ปริภาสโก โหติ, ภิกฺขุนีหิ สทฺธิํ สํสฏฺโฐ วิหรติ อนนุโลมิเกน สํสคฺเคน. อิเมหิ โข อุปาลิ ปญฺจหงฺเคหิ สมนฺนาคเตน ภิกฺขุนา ภิกฺขุนีนํ โอวาโท น ฐเปตพฺโพ.}

\prose{\csromanfolio{341}อปเรหิปิ อุปาลิ ปญฺจหงฺเคหิ สมนฺนาคเตน ภิกฺขุนา ภิกฺขุนีนํ โอวาโท น ฐเปตพฺโพ. กตเมหิ ปญฺจหิ, อลชฺชี จ โหติ, พาโล จ, อปกตตฺโต จ, ภณฺฑนการโก จ โหติ กลหการโก, สิกฺขาย จ น ปริปูริการี. อิเมหิ โข อุปาลิ ปญฺจหงฺเคหิ สมนฺนาคเตน ภิกฺขุนา ภิกฺขุนีนํ โอวาโท น ฐเปตพฺโพติ.}

\proseitem{453}{กติหิ นุ โข ภนฺเต องฺเคหิ สมนฺนาคเตน ภิกฺขุนา ภิกฺขุนีนํ โอวาโท น คเตตพฺโพติ. ปญฺจหุปาลิ องฺเคหิ สมนฺนาคเตน ภิกฺขุนา ภิกฺขุนีนํ โอวาโท น คเหตพฺโพ. กตเมหิ ปญฺจหิ, กายิเกน อนาจาเรน สมนฺนาคโต โหติ,วาจสิเกน อนาจาเรน สมนฺนาคโต โหติ, กายิก\-วาจสิเกน อนาจาเรน สมนฺนาคโต โหติ, ภิกฺขุนีนํ อกฺโกสก\-ปริภาสโก โหติ, ภิกฺขุนีหิ สทฺธิํ สํสฏฺโฐ วิหรติ อนนุโลมิเกน สํสคฺเคน. อิเมหิ โข อุปาลิ ปญฺจหงฺเคหิ สมนฺนาคเตน ภิกฺขุนา ภิกฺขุนีนํ โอวาโท น คเหตพฺโพ.}

\prose{อปเรหิปิ อุปาลิ ปญฺจหงฺเคหิ สมนฺนาคเตน ภิกฺขุนา ภิกฺขุนีนํ โอวาโท น คเหตพฺโพ. กตเมหิ ปญฺจหิ, อลชฺชี จ โหติ, พาโล จ, อปกตตฺโต จ, คมิโก วา โหติ, คิลาโน วา. อิเมหิ โข อุปาลิ ปญฺจหงฺเคหิ สมนฺนาคเตน ภิกฺขุนา ภิกฺขุนีนํ โอวาโท น คเหตพฺโพติ.}

\proseitem{454}{กติหิ นุ โข ภนฺเต องฺเคหิ สมนฺนาคเตน ภิกฺขุนา สทฺธิํ น สากจฺฉิตพฺโพติ. ปญฺจหุปาลิ องฺเคหิ สมนฺนาคเตน ภิกฺขุนา สทฺธิํ น สากจฺฉิตพฺโพ. กตเมหิ ปญฺจหิ, น อเสกฺเขน สีลกฺขนฺเธน สมนฺนาคโต โหติ, น อเสกฺเขน สมาธิ\-กฺขนฺเธน สมนฺนาคโต โหติ, น อเสกฺเขน ปญฺญา\-กฺขนฺเธน สมนฺนาคโต โหติ, น อเสกฺเขน วิมุตฺติ\-กฺขนฺเธน สมนฺนาคโต โหติ, น อเสกฺเขน วิมุตฺติ\-ญาณ\-ทสฺสนกฺขนฺเธน สมนฺนาคโต โหติ. อิเมหิ โข อุปาลิ ปญฺจหงฺเคหิ สมนฺนาคเตน ภิกฺขุนา สทฺธิํ น สากจฺฉิตพฺโพ. ปญฺจหุปาลิ องฺเคหิ สมนฺนาคเตน ภิกฺขุนา สทฺธิํ สากจฺฉาตพฺโพ. กตเมหิ ปญฺจหิ, อเสกฺเขน สีลกฺขนฺเธน สมนฺนาคโต โหติ, อเสกฺเขน สมาธิ\-กฺขนฺเธน สมนฺนาคโต โหติ, อเสกฺเขน ปญฺญา\-กฺขนฺเธน สมนฺนาคโต \csromanfolio{342}โหติ, อเสกฺเขน วิมุตฺติ\-กฺขนฺเธน สมนฺนาคโต โหติ, อเสกฺเขน วิมุตฺติ\-ญาณ\-ทสฺสนกฺขนฺเธน สมนฺนาคโต โหติ. อิเมหิ โข อุปาลิ ปญฺจหงฺเคหิ สมนฺนาคเตน ภิกฺขุนา สทฺธิํ สากจฺฉิตพฺโพ.}

\prose{อปเรหิปิ อุปาลิ ปญฺจหงฺเคหิ สมนฺนาคเตน ภิกฺขุนา สทฺธิํ น สากจฺฉิตพฺโพ. กตเมหิ ปญฺจหิ, น อตฺถปฏิสมฺภิทา\-ปตฺโต โหติ, น ธมฺมปฏิสมฺภิทา\-ปตฺโต โหติ, น นิรุตฺติปฏิสมฺภิทา\-ปตฺโต โหติ, น ปฏิภาน\-ปฏิสมฺภิทาปตฺโต โหติ, ยถาวิมุตฺตํ จิตฺตํ น ปจฺจเวกฺขิตา. อิเมหิ โข อุปาลิ ปญฺจหงฺเคหิ สมนฺนาคเตน ภิกฺขุนา สทฺธิํ น สากจฺฉิตพฺโพ. ปญฺจหุปาลิ องฺเคหิ สมนฺนาคเตน ภิกฺขุนา สทฺธิํ สากจฺฉิตพฺโพ. กตเมหิ ปญฺจหิ, อตฺถปฏิสมฺภิทา\-ปตฺโต โหติ, ธมฺมปฏิสมฺภิ ทาปตฺโต โหติ, นิรุตฺติปฏิสมฺภิทา\-ปตฺโต โหติ, ปฏิภาน\-ปฏิสมฺภิทาปตฺโต โหติ, ยถาวิมุตฺตํ จิตฺตํ ปจฺจเวกฺขิตา. อิเมหิ โข อุปาลิ ปญฺจหงฺเคหิ สมนฺนาคเตน ภิกฺขุนา สทฺธิํ สากจฺฉิตพฺโพติ.}

\vspace{\tipitakaparskip}
\csromancenter{ภิกฺขุโน\-วาทวคฺโค นิฏฺฐิโต อฏฺฐโม.}
\csromancenter{ตสฺสุทฺทานํ.}
\setlength{\csromangathapagewidth}{0pt}
\setlength{\csromangathawakwidth}{0pt}
\csromangathameasuregroup{ภิกฺขุนีเหว กาตพฺพํ, \\ ภิกฺขุนีนํ ตโย กมฺมา,}{\csromangathabat{ภิกฺขุนีเหว กาตพฺพํ,}{อปเรหิ ตถา ทุเว.} \\ \csromangathabat{ภิกฺขุนีนํ ตโย กมฺมา,}{น ฐเปตพฺพา เทฺว ทุกา.}}
\csromangathasetleft{ภิกฺขุนีเหว กาตพฺพํ, \\ ภิกฺขุนีนํ ตโย กมฺมา,}
\csromangathagroup{\csromangathastanza{\csromangathabat{ภิกฺขุนีเหว กาตพฺพํ,}{อปเรหิ ตถา ทุเว.} \\ \csromangathabat{ภิกฺขุนีนํ ตโย กมฺมา,}{น ฐเปตพฺพา เทฺว ทุกา.}}}

\vspace{\tipitakaparskip}
\csromancenter{น คเหตพฺพา เทฺว วุตฺตา, สากจฺฉาสุ จ เทฺว ทุกาติ.}
\csromansectionrule
\csromanmatikaanchor{mtk.228}
\csromantocmarkref{subparagraph}{9. อุพฺพาหิกวคฺค}{mtk.228}
\bookmark[dest={mtk.228},level=5]{9. อุพฺพาหิกวคฺค}
\csromanheader{5}{9. \textbf{อุพฺพาหิกวคฺค}}
\proseitem{455}{กติหิ นุ โข ภนฺเต องฺเคหิ สมนฺนาคโต ภิกฺขุ อุพฺพาหิกาย น สมฺมนฺนิ\-ตพฺโพติ. ปญฺจหุปาลิ องฺเคหิ สมนฺนาคโต ภิกฺขุ อุพฺพาหิกาย น สมฺมนฺนิตพฺโพ. กตเมหิ ปญฺจหิ, น อตฺถกุสโล โหติ, น ธมฺมกุสโล โหติ, น นิรุตฺติกุสโล โหติ, น พฺยญฺชนกุสโล โหติ, น ปุพฺพาปร\-กุสโล โหติ. อิเมหิ โข อุปาลิ ปญฺจหงฺเคหิ สมนฺนาคโต ภิกฺขุ อุพฺพาหิกาย น สมฺมนฺนิตพฺโพ. ปญฺจหุปาลิ องฺเคหิ สมนฺนาคโต ภิกฺขุ อุพฺพาหิกาย สมฺมนฺนิตพฺโพ. กตเมหิ ปญฺจหิ, อตฺถกุสโล โหติ, ธมฺมกุสโล โหติ, \csromanfolio{343}นิรุตฺติกุสโล โหติ, พฺยญฺชนกุสโล โหติ, ปุพฺพาปร\-กุสโล โหติ. อิเมหิ โข อุปาลิ ปญฺจหงฺเคหิ สมนฺนาคโต ภิกฺขุ อุพฺพาหิกาย สมฺมนฺนิตพฺโพ.}

\prose{อปเรหิปิ อุปาลิ ปญฺจหงฺเคหิ สมนฺนาคโต ภิกฺขุ อุพฺพาหิกาย น สมฺมนฺนิตพฺโพ. กตเมหิ ปญฺจหิ, โกธโน โหติ โกธาภิภูโต, มกฺขี โหติ มกฺขาภิภูโต, ปฬาสี โหติ ปฬาสาภิภูโต, อิสฺสุกี โหติ อิสฺสาภิภูโต, สนฺทิฏฺฐิ\-ปรามาสี โหติ อาธานคฺคาหี ทุปฺปฏินิสฺสคฺคี. อิเมหิ โข อุปาลิ ปญฺจหงฺเคหิ สมนฺนาคโต ภิกฺขุ อุพฺพาหิกาย น สมฺมนฺนิตพฺโพ. ปญฺจหุปาลิ องฺเคหิ สมนฺนาคโต ภิกฺขุ อุพฺพาหิกาย สมฺมนฺนิตพฺโพ. กตเมหิ ปญฺจหิ, น โกธโน โหติ น โกธาภิภูโต, น มกฺขี โหติ น มกฺขาภิภูโต, น ปฬาสี โหติ น ปฬาสาภิภูโต, น อิสฺสุกี โหติ น อิสฺสาภิภูโต, อสนฺทิฏฺฐิ\-ปรามาสี โหติ อนาธานคฺคาหี สุปฺปฏินิสฺสคฺคี. อิเมหิ โข อุปาลิ ปญฺจหงฺเคหิ สมนฺนาคโต ภิกฺขุ อุพฺพาหิกาย สมฺมนฺนิตพฺโพ.}

\prose{อปเรหิปิ อุปาลิ ปญฺจหงฺเคหิ สมนฺนาคโต ภิกฺขุ อุพฺพาหิกาย น สมฺมนฺนิตพฺโพ. กตเมหิ ปญฺจหิ, กุปฺปติ, พฺยาปชฺชติ, ปติฏฺฐิยติ, โกปํ ชเนติ, อขโม โหติ อปทกฺขิณ\-คฺคาหี อนุสาสนิํ. อิเมหิ โข อุปาลิ ปญฺจหงฺเคหิ สมนฺนาคโต ภิกฺขุ อุพฺพาหิกาย น สมฺมนฺนิตพฺโพ. ปญฺจหุปาลิ องฺเคหิ สมนฺนาคโต ภิกฺขุ อุพฺภาหิกาย สมฺมนฺนิตพฺโพ. กตเมหิ ปญฺจหิ, น กุปฺปติ, น พฺยาปชฺชติ, น ปติฏฺฐิยติ, น โกปํ ชเนติ, ขโม โหติ ปทกฺขิณ\-คฺคาหี อนุสาสนิํ. อิเมหิ โข อุปาลิ ปญฺจหงฺเคหิ สมนฺนาคโต ภิกฺขุ อุพฺพาหิกาย สมฺมนฺนิตพฺโพ.}

\prose{อปเรหิปิ อุปาลิ ปญฺจหงฺเคหิ สมนฺนาคโต ภิกฺขุ อุพฺพาหิกาย น สมฺมนฺนิตพฺโพ. กตเมหิ ปญฺจหิ, ปสาเรตา โหติ, โน สาเรตา, อโนกาสกมฺมํ กาเรตฺวา ปวตฺตา โหติ, น ยถาธมฺเม ยถาวินเย ยถาปตฺติยา โจเทตา โหติ, น ยถาธมฺเม ยถาวินเย ยถาปตฺติยา กาเรตา โหติ, น ยถาทิฏฺฐิยา พฺยากตา โหติ. อิเมหิ โข อุปาลิ ปญฺจหงฺเคหิ สมนฺนาคโต ภิกฺขุ อุพฺพาหิกาย น สมฺมนฺนิตพฺโพ. ปญฺจหุปาลิ องฺเคหิ สมนฺนาคโต ภิกฺขุ อุพฺพหิกาย สมฺมนฺนิตพฺโพ. กตเมหิ ปญฺจหิ, สาเรตา โหติ, โน ปสาเรตา, \csromanfolio{344}โอกาสกมฺมํ กาเรตฺวา ปวตฺตา โหติ, ยถาธมฺเม ยถาวินเย ยถาปตฺติยา โจเทตา โหติ, ยถาธมฺเม ยถาวินเย ยถาปตฺติยา กาเรตา โหติ, ยถาทิฏฺฐิยา พฺยากตา โหติ. อิเมหิ โข อุปาลิ ปญฺจหงฺเคหิ สมนฺนาคโต ภิกฺขุ อุพฺพาหิกาย สมฺมนฺนิตพฺโพ.}

\prose{อปเรหิปิ อุปาลิ ปญฺจหงฺเคหิ สมนฺนาคโต ภิกฺขุ อุพฺพาหิกาย น สมฺมนฺนิตพฺโพ. กตเมหิ ปญฺจหิ, ฉนฺทาคติํ คจฺฉติ, โทสาคติํ คจฺฉติ, โมหาคติํ คจฺฉติ, ภยาคติํ คจฺฉติ, อลชฺชี จ โหติ. อิเมหิ โข อุปาลิ ปญฺจหงฺเคหิ สมนฺนาคโต ภิกฺขุ อุพฺพาหิกาย น สมฺมนฺนิตพฺโพ. ปญฺจหุปาลิ องฺเคหิ สมนฺนาคโต ภิกฺขุ อุพฺพาหิกาย สมฺมนฺนิตพฺโพ. กตเมหิ ปญฺจหิ, น ฉนฺทาคติํ คจฺฉติ, น โทสาคติํ คจฺฉติ, น โมหาคติํ คจฺฉติ, น ภยาคติํ คจฺฉติ, ลชฺชี จ โหติ. อิเมหิ โข อุปาลิ ปญฺจหงฺเคหิ สมนฺนาคโต ภิกฺขุ อุพฺพาหิกาย สมฺมนฺนิตพฺโพ.}

\prose{อปเรหิปิ อุปาลิ ปญฺจหงฺเคหิ สมนฺนาคโต ภิกฺขุ อุพฺพาหิกาย น สมฺมนฺนิตพฺโพ. กตเมหิ ปญฺจหิ, ฉนฺทาคติํ คจฺฉติ, โทสาคติํ คจฺฉติ, โมหาคติํ คจฺฉติ, ภยาคติํ คจฺฉติ, อกุสโล จ โหติ วินเย. อิเมหิ โข อุปาลิ ปญฺจหงฺเคหิ สมนฺนาคโต ภิกฺขุ อุพฺพาหิกาย น สมฺมนฺนิตพฺโพ. ปญฺจหุปาลิ องฺเคหิ สมนฺนาคโต ภิกฺขุ อุพฺพาหิกาย สมฺมนฺนิตพฺโพ. กตเมหิ ปญฺจหิ, น ฉนฺทาคติํ คจฺฉติ, น โทสาคติํ คจฺฉติ, น โมหาคติํ คจฺฉติ, น ภยาคติํ คจฺฉติ, กุสโล จ โหติ วินเย. อิเมหิ โข อุปาลิ ปญฺจหงฺเคหิ สมนฺนาคโต ภิกฺขุ อุพฺพาหิกาย สมฺมนฺนิ\-ตพฺโพติ.}

\proseitem{456}{กติหิ นุ โข ภนฺเต องฺเคหิ สมนฺนาคโต ภิกฺขุ “พาโล”เตฺวว สงฺขํ คจฺฉตีติ. ปญฺจหุปาลิ องฺเคหิ สมนฺนาคโต ภิกฺขุ พาโลเตฺวว สงฺขํ คจฺฉติ. กตเมหิ ปญฺจหิ, สุตฺตํ น ชานาติ, สุตฺตานุโลมํ น ชานาติ, วินยํ น ชานาติ, วินยานุโลมํ น ชานาติ, น จ ฐานาฐาน\-กุสโล โหติ. อิเมหิ โข อุปาลิ ปญฺจหงฺเคหิ สมนฺนาคโต ภิกฺขุ “พาโล”เตฺวว สงฺขํ คจฺฉติ. ปญฺจหุปาลิ องฺเคหิ สมนฺนาคโต ภิกฺขุ “ปณฺฑิโต”เตฺวว สงฺขํ คจฺฉติ. กตเมหิ ปญฺจหิ, สุตฺตํ ชานาติ, สุตฺตานุโลมํ ชานาติ, วินยํ ชานาติ, \csromanfolio{345}วินยานุโลมํ ชานาติ, ฐานาฐาน\-กุสโล จ โหติ. อิเมหิ โข อุปาลิ ปญฺจหงฺเคหิ สมนฺนาคโต ภิกฺขุ “ปณฺฑิโต”เตฺวว สงฺขํ คจฺฉติ.}

\prose{อปเรหิปิ อุปาลิ ปญฺจหงฺเคหิ สมนฺนาคโต ภิกฺขุ “พาโล”เตฺวว สงฺขํ คจฺฉติ. กตเมหิ ปญฺจหิ, ธมฺมํ น ชานาติ, ธมฺมานุโลมํ น ชานาติ. วินยํ น ชานาติ, วินยานุโลมํ น ชานาติ, น จ ปุพฺพาปร\-กุสโล โหติ. อิเมหิ โข อุปาลิ ปญฺจหงฺเคหิ สมนฺนาคโต ภิกฺขุ “พาโล”เตฺวว สงฺขํ คจฺฉติ. ปญฺจหุปาลิ องฺเคหิ สมนฺนาคโต ภิกฺขุ “ปณฺฑิโต”เตฺวว สงฺขํ คจฺฉติ. กตเมหิ ปญฺจหิ, ธมฺมํ ชานาติ, ธมฺมานุโลมํ ชานาติ, วินยํ ชานาติ, วินยานุโลมํ ชานาติ, ปุพฺพาปร\-กุสโล จ โหติ. อิเมหิ โข อุปาลิ ปญฺจหงฺเคหิ สมนฺนาคโต ภิกฺขุ “ปณฺฑิโต”เตฺวว สงฺขํ คจฺฉติ.}

\prose{อปเรหิปิ อุปาลิ ปญฺจหงฺเคหิ สมนฺนาคโต ภิกฺขุ “พาโล”เตฺวว สงฺขํ คจฺฉติ. กตเมหิ ปญฺจหิ, วตฺถุํ น ชานาติ, นิทานํ น ชานาติ, ปญฺญตฺติํ น ชานาติ, ปทปจฺจา\-ภฏฺฐํ น ชานาติ, อนุสนฺธิ\-วจนปถํ น ชานาติ. อิเมหิ โข อุปาลิ ปญฺจหงฺเคหิ สมนฺนาคโต ภิกฺขุ “พาโล”เตฺวว สงฺขํ คจฺฉติ. ปญฺจหุปาลิ องฺเคหิ สมนฺนาคโต ภิกฺขุ “ปณฺฑิโต”เตฺวว สงฺขํ คจฺฉติ. กตเมหิ ปญฺจหิ, วตฺถุํ ชานาติ, นิทานํ ชานาติ, ปญฺญตฺติํ ชานาติ, ปทปจฺจา\-ภฏฺฐํ ชานาติ, อนุสนฺธิ\-วจนปถํ ชานาติ. อิเมหิ โข อุปาลิ ปญฺจหงฺเคหิ สมนฺนาคโต ภิกฺขุ “ปณฺฑิโต”เตฺวว สงฺขํ คจฺฉติ.}

\prose{อปเรหิปิ อุปาลิ ปญฺจหงฺเคหิ สมนฺนาคโต ภิกฺขุ “พาโล”เตฺวว สงฺขํ คจฺฉติ. กตเมหิ ปญฺจหิ, อาปตฺติํ น ชานาติ, อาปตฺติ\-สมุฏฺฐานํ น ชานาติ, อาปตฺติยา ปโยคํ น ชานาติ, อาปตฺติยา วูปสมํ น ชานาติ, อาปตฺติยา น วินิจฺฉย\-กุสโล โหติ. อิเมหิ โข อุปาลิ ปญฺจหงฺเคหิ สมนฺนาคโต ภิกฺขุ “พาโล”เตฺวว สงฺขํ คจฺฉติ. ปญฺจหุปาลิ องฺเคหิ สมนฺนาคโต ภิกฺขุ “ปณฺฑิโต”เตฺวว สงฺขํ คจฺฉติ. กตเมหิ ปญฺจหิ, อาปตฺติํ ชานาติ, อาปตฺติ\-สมุฏฺฐานํ ชานาติ, อาปตฺติยา ปโยคํ ชานาติ, อาปตฺติยา วูปสมํ ชานาติ, อาปตฺติยา วินิจฺฉย\-กุสโล โหติ. อิเมหิ โข อุปาลิ ปญฺจหงฺเคหิ สมนฺนาคโต ภิกฺขุ “ปณฺฑิโต”เตฺวว สงฺขํ คจฺฉติ.}

\prose{\csromanfolio{346}อปเรหิปิ อุปาลิ ปญฺจหงฺเคหิ สมนฺนาคโต ภิกฺขุ “พาโล”เตฺวว สงฺขํ คจฺฉติ. กตเมหิ ปญฺจหิ, อธิกรณํ น ชานาติ, อธิกรณ\-สมุฏฺฐานํ น ชานาติ, อธิกรณสฺส ปโยคํ น ชานาติ, อธิกรณสฺส วูปสมํ น ชานาติ, อธิกรณสฺส น วินิจฺฉย\-กุสโล โหติ. อิเมหิ โข อุปาลิ ปญฺจหงฺเคหิ สมนฺนาคโต ภิกฺขุ “พาโล”เตฺวว สงฺขํ คจฺฉติ. ปญฺจหุปาลิ องฺเคหิ สมนฺนาคโต ภิกฺขุ “ปณฺฑิโต”เตฺวว สงฺขํ คจฺฉติ. กตเมหิ ปญฺจหิ, อธิกรณํ ชานาติ, อธิกรณ\-สมุฏฺฐานํ ชานาติ, อธิกรณสฺส ปโยคํ ชานาติ, อธิกรณสฺส วูปสมํ ชานาติ, อธิกรณสฺส วินิจฺฉย\-กุสโล โหติ. อิเมหิ โข อุปาลิ ปญฺจหงฺเคหิ สมนฺนาคโต ภิกฺขุ “ปณฺฑิโต”เตฺวว สงฺขํ คจฺฉตีติ.}

\vspace{\tipitakaparskip}
\csromancenter{อุพฺพาหิกวคฺโค นิฏฺฐิโต นวโม.}
\csromancenter{ตสฺสุทฺทานํ.}
\setlength{\csromangathapagewidth}{0pt}
\setlength{\csromangathawakwidth}{0pt}
\csromangathameasuregroup{น อตฺถกุสโล เจว, \\ ปสาเรตา ฉนฺทาคติํ, \\ สุตฺตํ ธมฺมญฺจ วตฺถุญฺจ,}{\csromangathabat{น อตฺถกุสโล เจว,}{โกธโน กุปฺปตี จ โย.} \\ \csromangathabat{ปสาเรตา ฉนฺทาคติํ,}{อกุสโล ตเถว จ.} \\ \csromangathabat{สุตฺตํ ธมฺมญฺจ วตฺถุญฺจ,}{อาปตฺติํ อธิกรณํ.}}
\csromangathasetleft{น อตฺถกุสโล เจว, \\ ปสาเรตา ฉนฺทาคติํ, \\ สุตฺตํ ธมฺมญฺจ วตฺถุญฺจ,}
\csromangathagroup{\csromangathastanza{\csromangathabat{น อตฺถกุสโล เจว,}{โกธโน กุปฺปตี จ โย.} \\ \csromangathabat{ปสาเรตา ฉนฺทาคติํ,}{อกุสโล ตเถว จ.} \\ \csromangathabat{สุตฺตํ ธมฺมญฺจ วตฺถุญฺจ,}{อาปตฺติํ อธิกรณํ.}}}

\vspace{\tipitakaparskip}
\csromancenter{เทฺว เทฺว ปกาสิตา สพฺเพ, กณฺหสุกฺกํ วิชานถาติ.}
\csromansectionrule
\csromanmatikaanchor{mtk.229}
\csromantocmarkref{subparagraph}{10. อธิกรณวูปสมวคฺค}{mtk.229}
\bookmark[dest={mtk.229},level=5]{10. อธิกรณวูปสมวคฺค}
\csromanheader{5}{10. อธิกรณ\-วูปสม\-วคฺค}
\proseitem{457}{กติหิ นุ โข ภนฺเต องฺเคหิ สมนฺนาคโต ภิกฺขุ นาลํ อธิกรณํ วูปสเมตุนฺติ. ปญฺจหุปาลิ องฺเคหิ สมนฺนาคโต ภิกฺขุ นาลํ อธิกรณํ วูปสเมตุํ. กตเมหิ ปญฺจหิ, อาปตฺติํ น ชานาติ, อาปตฺติ\-สมุฏฺฐานํ น ชานาติ, อาปตฺติยา ปโยคํ น ชานาติ, อาปตฺติยา วูปสมํ น ชานาติ. อาปตฺติยา น วินิจฺฉย\-กุสโล โหติ. อิเมหิ โข อุปาลิ ปญฺจหงฺเคหิ สมนฺนาคโต ภิกฺขุ นาลํ อธิกรณํ วูปสเมตุํ. ปญฺจหุปาลิ องฺเคหิ สมนฺนาคโต ภิกฺขุ อลํ อธิกรณํ วูปสเมตุํ. กตเมหิ ปญฺจหิ, อาปตฺติํ ชานาติ, อาปตฺติ\-สมุฏฺฐานํ ชานาติ, อาปตฺติยา ปโยคํ ชานาติ, อาปตฺติยา วูปสมํ ชานาติ, \csromanfolio{347}อาปตฺติยา วินิจฺฉย\-กุสโล โหติ. อิเมหิ โข อุปาลิ ปญฺจหงฺเคหิ สมนฺนาคโต ภิกฺขุ อลํ อธิกรณํ วูปสเมตุํ.}

\prose{อปเรหิปิ อุปาลิ ปญฺจหงฺเคหิ สมนฺนาคโต ภิกฺขุ นาลํ อธิกรณํ วูปสเมตุํ. กตเมหิ ปญฺจหิ, อธิกรณํ น ชานาติ, อธิกรณ\-สมุฏฺฐานํ น ชานาติ, อธิกรณสฺส ปโยคํ น ชานาติ, อธิกรณสฺส วูปสมํ น ชานาติ, อธิกรณสฺส น วินิจฺฉย\-กุสโล โหติ. อิเมหิ โข อุปาลิ ปญฺจหงฺเคหิ สมนฺนาคโต ภิกฺขุ นาลํ อธิกรณํ วูปสเมตุํ.}

\prose{ปญฺจหุปาลิ องฺเคหิ สมนฺนาคโต ภิกฺขุ อลํ อธิกรณํ วูปสเมตุํ. กตเมหิ ปญฺจหิ, อธิกรณํ ชานาติ, อธิกรณ\-สมุฏฺฐานํ ชานาติ, อธิกรณสฺส ปโยคํ ชานาติ, อธิกรณสฺส วูปสมํ ชานาติ, อธิกรณสฺส วินิจฺฉย\-กุสโล โหติ. อิเมหิ โข อุปาลิ ปญฺจหงฺเคหิ สมนฺนาคโต ภิกฺขุ อลํ อธิกรณํ วูปสเมตุํ.}

\prose{อปเรหิปิ อุปาลิ ปญฺจหงฺเคหิ สมนฺนาคโต ภิกฺขุ นาลํ อธิกรณํ วูปสเมตุํ. กตเมหิ ปญฺจหิ, ฉนฺทาคติํ คจฺฉติ, โทสาคติํ คจฺฉติ, โมหาคติํ คจฺฉติ, ภยาคติํ คจฺฉติ, อลชฺชี จ โหติ. อิเมหิ โข อุปาลิ ปญฺจหงฺเคหิ สมนฺนาคโต ภิกฺขุ นาลํ อธิกรณํ วูปสเมตุํ. ปญฺจหุปาลิ องฺเคหิ สมนฺนาคโต ภิกฺขุ อลํ อธิกรณํ วูปสเมตุํ. กตเมหิ ปญฺจหิ, น ฉนฺทาคติํ คจฺฉติ, น โทสาคติํ คจฺฉติ, น โมหาคติํ คจฺฉติ, น ภยาคติํ คจฺฉติ, ลชฺชี จ โหติ. อิเมหิ โข อุปาลิ ปญฺจหงฺเคหิ สมนฺนาคโต ภิกฺขุ อลํ อธิกรณํ วูปสเมตุํ.}

\prose{อปเรหิปิ อุปาลิ ปญฺจหงฺเคหิ สมนฺนาคโต ภิกฺขุ นาลํ อธิกรณํ วูปสเมตุํ. กตเมหิ ปญฺจหิ, ฉนฺทาคติํ คจฺฉติ, โทสาคติํ คจฺฉติ, โมหาคติํ คจฺฉติ, ภยาคติํ คจฺฉติ, อปฺปสฺสุโต จ โหติ. อิเมหิ โข อุปาลิ ปญฺจหงฺเคหิ สมนฺนาคโต ภิกฺขุ นาลํ อธิกรณํ วูปสเมตุํ. ปญฺจหุปาลิ องฺเคหิ สมนฺนาคโต ภิกฺขุ อลํ อธิกรณํ วูปสเมตุํ. กตเมหิ ปญฺจ, น ฉนฺทาคติํ คจฺฉติ, น โทสาคติํ คจฺฉติ, น โมหาคติํ คจฺฉติ, น ภยาคติํ คจฺฉติ, พหุสฺสุโต จ โหติ. อิเมหิ โข อุปาลิ ปญฺจหงฺเคหิ สมนฺนาคโต ภิกฺขุ อลํ อธิกรณํ วูปสเมตุํ.}

\prose{\csromanfolio{348}อปเรหิปิ อุปาลิ ปญฺจหงฺเคหิ สมนฺนาคโต ภิกฺขุ นาลํ อธิกรณํ วูปสเมตุํ. กตเมหิ ปญฺจหิ, วตฺถุํ น ชานาติ, นิทานํ น ชานาติ, ปญฺญตฺติํ น ชานาติ, ปทปจฺจา\-ภฏฺฐํ น ชานาติ, อนุสนฺธิ\-วจนปถํ น ชานาติ. อิเมหิ โข อุปาลิ ปญฺจหงฺเคหิ สมนฺนาคโต ภิกฺขุ นาลํ อธิกรณํ วูปสเมตุํ. ปญฺจหุปาลิ องฺเคหิ สมนฺนาคโต ภิกฺขุ อลํ อธิกรณํ วูปสเมตุํ. กตเมหิ ปญฺจหิ, วตฺถุํ ชานาติ, นิทานํ ชานาติ, ปญฺญตฺติํ ชานาติ, ปทปจฺจา\-ภฏฺฐํ ชานาติ, อนุสนฺธิ\-วจนปถํ ชานาติ. อิเมหิ โข อุปาลิ ปญฺจหงฺเคหิ สมนฺนาคโต ภิกฺขุ อลํ อธิกรณํ วูปสเมตุํ.}

\prose{อปเรหิปิ อุปาลิ ปญฺจหงฺเคหิ สมนฺนาคโต ภิกฺขุ นาลํ อธิกรณํ วูปสเมตุํ. กตเมหิ ปญฺจหิ, ฉนฺทาคติํ คจฺฉติ, โทสาคติํ คจฺฉติ, โมหาคติํ คจฺฉติ, ภยาคติํ คจฺฉติ, อกุสโล จ โหติ วินเย. อิเมหิ โข อุปาลิ ปญฺจหงฺเคหิ สมนฺนาคโต ภิกฺขุ นาลํ อธิกรณํ วูปสเมตุํ. ปญฺจหุปาลิ องฺเคหิ สมนฺนาคโต ภิกฺขุ อลํ อธิกรณํ วูปสเมตุํ. กตเมหิ ปญฺจหิ, น ฉนฺทาคติํ คจฺฉติ, น โทสาคติํ คจฺฉติ, น โมหาคติํ คจฺฉติ, น ภยาคติํ คจฺฉติ, กุสโล จ โหติ วินเย. อิเมหิ โข อุปาลิ ปญฺจหงฺเคหิ สมนฺนาคโต ภิกฺขุ อลํ อธิกรณํ วูปสเมตุํ.}

\prose{อปเรหิปิ อุปาลิ ปญฺจหงฺเคหิ สมนฺนาคโต ภิกฺขุ นาลํ อธิกรณํ วูปสเมตุํ. กตเมหิ ปญฺจหิ, ฉนฺทาคติํ คจฺฉติ, โทสาคติํ คจฺฉติ, โมหาคติํ คจฺฉติ, ภยาคติํ คจฺฉติ, ปุคฺคลครุ โหติ, โน สํฆครุ. อิเมหิ โข อุปาลิ ปญฺจหงฺเคหิ สมนฺนาคโต ภิกฺขุ นาลํ อธิกรณํ วูปสเมตุํ. ปญฺจหุปาลิ องฺเคหิ สมนฺนาคโต ภิกฺขุ อลํ อธิกรณํ วูปสเมตุํ. กตเมหิ ปญฺจหิ, น ฉนฺทาคติํ คจฺฉติ, น โทสาคติํ คจฺฉติ, น โมหาคติํ คจฺฉติ, น ภยาคติํ คจฺฉติ, สํฆครุ โหติ, โน ปุคฺคลครุ. อิเมหิ โข อุปาลิ ปญฺจหงฺเคหิ สมนฺนาคโต ภิกฺขุ อลํ อธิกรณํ วูปสเมตุํ.}

\prose{อปเรหิปิ อุปาลิ ปญฺจหงฺเคหิ สมนฺนาคโต ภิกฺขุ นาลํ อธิกรณํ วูปสเมตุํ. กตเมหิ ปญฺจหิ, ฉนฺทาคติํ คจฺฉติ, โทสาคติํ คจฺฉติ, โมหาคติํ คจฺฉติ, ภยาคติํ คจฺฉติ, อามิสครุ โหติ, โน สทฺธมฺมครุ. อิเมหิ โข อุปาลิ ปญฺจหงฺเคหิ สมนฺนาคโต ภิกฺขุ นาลํ \csromanfolio{349}อธิกรณํ วูปสเมตุํ. ปญฺจหุปาลิ องฺเคหิ สมนฺนาคโต ภิกฺขุ อลํ อธิกรณํ วูปสเมตุํ. กตเมหิ ปญฺจหิ, น ฉนฺทาคติํ คจฺฉติ, น โทสาคติํ คจฺฉติ, น โมหาคติํ คจฺฉติ, น ภยาคติํ คจฺฉติ, สทฺธมฺมครุ โหติ, โน อามิสครุ. อิเมหิ โข อุปาลิ ปญฺจหงฺเคหิ สมนฺนาคโต ภิกฺขุ อลํ อธิกรณํ วูปสเมตุนฺติ.}

\proseitem{458}{กติหิ นุ โข ภนฺเต อากาเรหิ สํโฆ ภิชฺชตีติ. ปญฺจหุปาลิ อากาเรหิ สํโฆ ภิชฺชติ. กตเมหิ ปญฺจหิ, กมฺเมน, อุทฺเทเสน, โวหรนฺโต, อนุสฺสาวเนน, สลากคฺคาเหน, อิเมหิ โข อุปาลิ ปญฺจหากาเรหิ สํโฆ ภิชฺชตีติ.}

\prose{“สํฆราชิ สํฆราชี”ติ ภนฺเต วุจฺจติ, กิตฺตาวตา นุ โข ภนฺเต สํฆราชิ โหติ โน จ สํฆเภโท, กิตฺตาวตา จ ปน สํฆราชิ เจว โหติ สํฆเภโท จาติ. ปญฺญตฺเตตํ อุปาลิ มยา อาคนฺตุกานํ ภิกฺขูนํ อาคนฺตุกวตฺตํ, เอวํ สุปญฺญตฺเต โข อุปาลิ มยา สิกฺขาปเท อาคนฺตุกา ภิกฺขู อาคนฺตุกวตฺเต น วตฺตนฺติ, เอวมฺปิ โข อุปาลิ สํฆราชิ โหติ โน จ สํฆเภโท. ปญฺญตฺเตตํ อุปาลิ มยา อาวาสิกานํ ภิกฺขูนํ อาวาสิกวตฺตํ, เอวํ สุปญฺญตฺเต โข อุปาลิ มยา สิกฺขาปเท อาวาสิกา ภิกฺขู อาวาสิกวตฺเต น วตฺตนฺติ, เอวมฺปิ โข อุปาลิ สํฆราชิ โหติ โน จ สํฆเภโท. ปญฺญตฺเตตํ อุปาลิ มยา ภิกฺขูนํ ภตฺตคฺเค ภตฺตคฺค\-วตฺตํ ยถาวุฑฺฒํ ยถารตฺตํ ยถาปติรูปํ อคฺคาสนํ อคฺโคทกํ อคฺคปิณฺฑํ, เอวํ สุปญฺญตฺเต โข อุปาลิ มยา สิกฺขาปเท นวา ภิกฺขู ภตฺตคฺเค เถรานํ ภิกฺขูนํ อาสนํ ปฏิพาหนฺติ, เอวมฺปิ โข อุปาลิ สํฆราชิ โหติ โน จ สํฆเภโท. ปญฺญตฺเตตํ อุปาลิ มยา ภิกฺขูนํ เสนาสเน เสนาสนวตฺตํ ยถาวุฑฺฒํ ยถารตฺตํ ยถาปติรูปํ, เอวํ สุปญฺญตฺเต โข อุปาลิ มยา สิกฺขาปเท นวา ภิกฺขู เถรานํ ภิกฺขูนํ เสนาสนํ ปฏิพาหนฺติ, เอวมฺปิ โข อุปาลิ สํฆราชิ โหติ โน จ สํเฆเภโท. ปญฺญตฺเตตํ อุปาลิ มยา ภิกฺขูนํ อนฺโตสีมาย เอกํ อุโปสถํ เอกํ ปวารณํ เอกํ สํฆกมฺมํ เอกํ กมฺมากมฺมํ, เอวํ สุปญฺญตฺเต โข อุปาลิ มยา สิกฺขาปเท ตตฺเถว อนฺโตสีมาย อาเวนิภาวํ\csromansharedfootnote{49a9357f5b8e7492}{อาเวณิภาวํ (สี, สฺยา)} กริตฺวา คณํ \csromanfolio{350}พนฺธิตฺวา อาเวนิํ\csromansharedfootnote{c79a44111c78d734}{อาเวณิ (สี, สฺยา)} อุโปสถํ กโรนฺติ, อาเวนิํ ปวารณํ กโรนฺติ, อาเวนิํ สํฆกมฺมํ กโรนฺติ, อาเวนิํ กมฺมากมฺมานิ กโรนฺติ, เอวํ โข อุปาลิ สํฆราชิ เจว โหติ สํฆเภโท จาติ.}

\vspace{\tipitakaparskip}
\csromancenter{อธิกรณ\-วูปสม\-วคฺโค นิฏฺฐิโต ทสโม.}
\csromancenter{ตสฺสุทฺทานํ.}
\setlength{\csromangathapagewidth}{0pt}
\setlength{\csromangathawakwidth}{0pt}
\csromangathameasuregroup{อาปตฺติํ อธิกรณํ, \\ วตฺถุํ จ อกุสโล จ,}{\csromangathabat{อาปตฺติํ อธิกรณํ,}{ฉนฺทา อปฺปสฺสุเตน จ.} \\ \csromangathabat{วตฺถุํ จ อกุสโล จ,}{ปุคฺคโล อามิเสน จ.}}
\csromangathasetleft{อาปตฺติํ อธิกรณํ, \\ วตฺถุํ จ อกุสโล จ,}
\csromangathagroup{\csromangathastanza{\csromangathabat{อาปตฺติํ อธิกรณํ,}{ฉนฺทา อปฺปสฺสุเตน จ.} \\ \csromangathabat{วตฺถุํ จ อกุสโล จ,}{ปุคฺคโล อามิเสน จ.}}}

\vspace{\tipitakaparskip}
\csromancenter{ภิชฺชติ สํฆราชิ จ, สํฆเภโท ตเถว จาติ.}
\csromansectionrule
\csromanmatikaanchor{mtk.230}
\csromantocmarkref{subparagraph}{11. สํฆเภทกวคฺค}{mtk.230}
\bookmark[dest={mtk.230},level=5]{11. สํฆเภทกวคฺค}
\csromanheader{5}{11. \textbf{สํฆเภทกวคฺค}}
\proseitem{459}{\csromansymbolfootnote{*}{วิ 4. 369, 370; อํ 3. 315 ปิฏฺฐาทีสุปิ.}กติหิ นุ โข ภนฺเต องฺเคหิ สมนฺนาคโต สํฆเภทโก อาปายิโก เนรยิโก กปฺปฏฺโฐ อเตกิจฺโฉติ. ปญฺจหุปาลิ องฺเคหิ สมนฺนาคโต สํฆเภทโก อาปายิโก เนรยิโก กปฺปฏฺโฐ อเตกิจฺโฉ. กตเมหิ ปญฺจหิ, อิธุปาลิ ภิกฺขุ อธมฺมํ “ธมฺโม”ติ ทีเปติ, ธมฺมํ “อธมฺโม”ติ ทีเปติ, อวินยํ “วินโย”ติ ทีเปติ, วินยํ “อวินโย”ติ ทีเปติ, วินิธาย ทิฏฺฐิํ กมฺเมน. อิเมหิ โข อุปาลิ ปญฺจหงฺเคหิ สมนฺนาคโต สํฆเภทโก อาปายิโก เนรยิโก กปฺปฏฺโฐ อเตกิจฺโฉ.}

\prose{อปเรหิปิ อุปาลิ ปญฺจหงฺเคหิ สมนฺนาคโต สํฆเภทโก อาปายิโก เนรยิโก กปฺปฏฺโฐ อเตกิจฺโฉ. กตเมหิ ปญฺจหิ, อิธุปาลิ ภิกฺขุ อธมฺมํ “ธมฺโม”ติ ทีเปติ, ธมฺมํ “อธมฺโม”ติ ทีเปติ, อวินยํ “วินโย”ติ ทีเปติ, วินยํ “อวินโย”ติ ทีเปติ, วินิธาย ทิฏฺฐิํ อุทฺเทเสน. อิเมหิ โข อุปาลิ ปญฺจหงฺเคหิ สมนฺนาคโต สํฆเภทโก อาปายิโก เนรยิโก กปฺปฏฺโฐ อเตกิจฺโฉ.}

\prose{อปเรหิปิ อุปาลิ ปญฺจหงฺเคหิ สมนฺนาคโต สํฆเภทโก อาปายิโก เนรยิโก กปฺปฏฺโฐ อเตกิจฺโฉ. กตเมหิ ปญฺจหิ, อิธุปาลิ ภิกฺขุ อธมฺมํ “ธมฺโม”ติ ทีเปติ, ธมฺมํ “อธมฺโม”ติ ทีเปติ, \csromanfolio{351}อวินยํ “วินโย”ติ ทีเปติ, วินยํ “อวินโย”ติ ทีเปติ, วินิธาย ทิฏฺฐิํ โวหรนฺโต. อิเมหิ โข อุปาลิ ปญฺจหงฺเคหิ สมนฺนาคโต สํฆเภทโก อาปายิโก เนรยิโก กปฺปฏฺโฐ อเตกิจฺโฉ.}

\prose{อปเรหิปิ อุปาลิ ปญฺจหงฺเคหิ สมนฺนาคโต สํฆเภทโก อาปายิโก เนรยิโก กปฺปฏฺโฐ อเตกิจฺโฉ. กตเมหิ ปญฺจหิ, อิธุปาลิ ภิกฺขุ อธมฺมํ “ธมฺโม”ติ ทีเปติ, ธมฺมํ “อธมฺโม”ติ ทีเปติ, อวินยํ “วินโย”ติ ทีเปติ, วินยํ “อวินโย”ติ ทีเปติ, วินิธาย ทิฏฺฐิํ อนุสฺสาวเนน. อิเมหิ โข อุปาลิ ปญฺจหงฺเคหิ สมนฺนาคโต สํฆเภทโก อาปายิโก เนรยิโก กปฺปฏฺโฐ อเตกิจฺโฉ.}

\prose{อปเรหิปิ อุปาลิ ปญฺจหงฺเคหิ สมนฺนาคโต สํฆเภทโก อาปายิโก เนรยิโก กปฺปฏฺโฐ อเตกิจฺโฉ. กตเมหิ ปญฺจหิ, อิธุปาลิ ภิกฺขุ อธมฺมํ “ธมฺโม”ติ ทีเปติ, ธมฺมํ “อธมฺโม”ติ ทีเปติ, อวินยํ “วินโย”ติ ทีเปติ, วินยํ “อวินโย”ติ ทีเปติ, วินิธาย ทิฏฺฐิํ สลากคฺคาเหน. อิเมหิ โข อุปาลิ ปญฺจหงฺเคหิ สมนฺนาคโต สํฆเภทโก อาปายิโก เนรยิโก กปฺปฏฺโฐ อเตกิจฺโฉ.}

\prose{อปเรหิปิ อุปาลิ ปญฺจหงฺเคหิ สมนฺนาคโต สํฆเภทโก อาปายิโก เนรยิโก กปฺปฏฺโฐ อเตกิจฺโฉ. กตเมหิ ปญฺจหิ, อิธุปาลิ ภิกฺขุ อธมฺมํ “ธมฺโม”ติ ทีเปติ, ธมฺมํ “อธมฺโม”ติ ทีเปติ, อวินยํ “วินโย”ติ ทีเปติ, วินยํ “อวินโย”ติ ทีเปติ, วินิธาย ขนฺติํ กมฺเมน ฯเปฯ วินิธาย ขนฺติํ อุทฺเทเสน ฯเปฯ วินิธาย ขนฺติํ โวหรนฺโต ฯเปฯ วินิธาย ขนฺติํ อนุสฺสาวเนน ฯเปฯ วินิธาย ขนฺติํ สลากคฺคาเหน. อิเมหิ โข อุปาลิ ปญฺจหงฺเคหิ สมนฺนาคโต สํฆเภทโก อาปายิโก เนรยิโก กปฺปฏฺโฐ อเตกิจฺโฉ.}

\prose{อปเรหิปิ อุปาลิ ปญฺจหงฺเคหิ สมนฺนาคโต สํฆเภทโก อาปายิโก เนรยิโก กปฺปฏฺโฐ อเตกิจฺโฉ. กตเมหิ ปญฺจหิ, อิธุปาลิ ภิกฺขุ อธมฺมํ “ธมฺโม”ติ ทีเปติ, ธมฺมํ “อธมฺโม”ติ ทีเปติ, อวินยํ “วินโย”ติ ทีเปติ, วินยํ “อวินโย”ติ ทีเปติ, วินิธาย รุจิํ กมฺเมน ฯเปฯ วินิธาย รุจิํ อุทฺเทเสน ฯเปฯ วินิธาย รุจิํ โวหรนฺโต ฯเปฯ วินิธาย รุจิํ อนุสฺสาวเนน ฯเปฯ วินิธาย รุจิํ สลากคฺคาเหน. อิเมหิ โข อุปาลิ \csromanfolio{352}ปญฺจหงฺเคหิ สมนฺนาคโต สํฆเภทโก อาปายิโก เนรยิโก กปฺปฏฺโฐ อเตกิจฺโฉ.}

\prose{อปเรหิปิ อุปาลิ ปญฺจหงฺเคหิ สมนฺนาคโต สํฆเภทโก อาปายิโก เนรยิโก กปฺปฏฺโฐ อเตกิจฺโฉ. กตเมหิ ปญฺจหิ, อิธุปาลิ ภิกฺขุ อธมฺมํ “ธมฺโม”ติ ทีเปติ, ธมฺมํ “อธมฺโม”ติ ทีเปติ, อวินยํ “วินโย”ติ ทีเปติ, วินยํ “อวินโย”ติ ทีเปติ, วินิธาย สญฺญํ กมฺเมน ฯเปฯ วินิธาย สญฺญํ อุทฺเทเสน ฯเปฯ วินิธาย สญฺญํ โวหรนฺโต ฯเปฯ วินิธาย สญฺญํ อนุสฺสาวเนน ฯเปฯ วินิธาย สญฺญํ สลากคฺคาเหน. อิเมหิ โข อุปาลิ ปญฺจหงฺเคหิ สมนฺนาคโต สํฆเภทโก อาปายิโก เนรยิโก กปฺปฏฺโฐ อเตกิจฺโฉติ.}

\vspace{\tipitakaparskip}
\csromancenter{สํฆเภทกวคฺโค นิฏฺฐิโต เอกาทสโม.}
\csromancenter{ตสฺสุทฺทานํ.}
\setlength{\csromangathapagewidth}{0pt}
\setlength{\csromangathawakwidth}{0pt}
\csromangathameasuregroup{วินิธาย ทิฏฺฐิํ กมฺเมน, \\ อนุสฺสาวเน สลาเกน,}{\csromangathabat{วินิธาย ทิฏฺฐิํ กมฺเมน,}{อุทฺเทเส โวหเรน จ.} \\ \csromangathabat{อนุสฺสาวเน สลาเกน,}{ปญฺเจเต ทิฏฺฐินิสฺสิตา.}}
\csromangathasetleft{วินิธาย ทิฏฺฐิํ กมฺเมน, \\ อนุสฺสาวเน สลาเกน,}
\csromangathagroup{\csromangathastanza{\csromangathabat{วินิธาย ทิฏฺฐิํ กมฺเมน,}{อุทฺเทเส โวหเรน จ.} \\ \csromangathabat{อนุสฺสาวเน สลาเกน,}{ปญฺเจเต ทิฏฺฐินิสฺสิตา.}}}

\vspace{\tipitakaparskip}
\csromancenter{ขนฺติํ รุจิํ จ สญฺญญฺจ, ตโย เต ปญฺจธา นยาติ.}
\csromansectionrule
\csromanmatikaanchor{mtk.231}
\csromantocmarkref{subparagraph}{12. ทุติยสํฆเภทกวคฺค}{mtk.231}
\bookmark[dest={mtk.231},level=5]{12. ทุติยสํฆเภทกวคฺค}
\csromanheader{5}{12. \textbf{ทุติยสํฆเภทกวคฺค}}
\proseitem{460}{กติหิ นุ โข ภนฺเต องฺเคหิ สมนฺนาคโต สํฆเภทโก น อาปายิโก น เนรยิโก น กปฺปฏฺโฐ น อเตกิจฺโฉติ. ปญฺจหุปาลิ องฺเคหิ สมนฺนาคโต สํฆเภทโก น อาปายิโก น เนรยิโก น กปฺปฏฺโฐ น อเตกิจฺโฉ. กตเมหิ ปญฺจหิ, อิธุปาลิ ภิกฺขุ อธมฺมํ “ธมฺโม”ติ ทีเปติ, ธมฺมํ “อธมฺโม”ติ ทีเปติ, อวินยํ “วินโย”ติ ทีเปติ, วินยํ “อวินโย”ติ ทีเปติ, อวินิธาย ทิฏฺฐิํ กมฺเมน. อิเมหิ โข อุปาลิ ปญฺจหงฺเคหิ สมนฺนาคโต สํฆเภทโก น อาปายิโก น เนรยิโก น กปฺปฏฺโฐ น อเตกิจฺโฉ.}

\prose{อปเรหิปิ อุปาลิ ปญฺจหงฺเคหิ สมนฺนาคโต สํฆเภทโก น อาปายิโก น เนรยิโก น กปฺปฏฺโฐ น อเตกิจฺโฉ. กตเมหิ ปญฺจหิ, อิธุปาลิ ภิกฺขุ อธมฺมํ “ธมฺโม”ติ ทีเปติ, ธมฺมํ “อธมฺโม”ติ ทีเปติ, \csromanfolio{353}อวินยํ “วินโย”ติ ทีเปติ, วินยํ “อวินโย”ติ ทีเปติ, อวินิธาย ทิฏฺฐิํ อุทฺเทเสน. อิเมหิ โข อุปาลิ ปญฺจหงฺเคหิ สมนฺนาคโต สํฆเภทโก น อาปายิโก น เนรยิโก น กปฺปฏฺโฐ น อเตกิจฺโฉ.}

\prose{อปเรหิปิ อุปาลิ ปญฺจหงฺเคหิ สมนฺนาคโต สํฆเภทโก น อาปายิโก น เนรยิโก น กปฺปฏฺโฐ น อเตกิจฺโฉ. กตเมหิ ปญฺจหิ, อิธุปาลิ ภิกฺขุ อธมฺมํ “ธมฺโม”ติ ทีเปติ, ธมฺมํ “อธมฺโม”ติ ทีเปติ, อวินยํ “วินโย”ติ ทีเปติ, วินยํ “อวินโย”ติ ทีเปติ, อวินิธาย ทิฏฺฐิํ โวหรนฺโต. อิเมหิ โข อุปาลิ ปญฺจหงฺเคหิ สมนฺนาคโต สํฆเภทโก น อาปายิโก น เนรยิโก น กปฺปฏฺโฐ น อเตกิจฺโฉ.}

\prose{อปเรหิปิ อุปาลิ ปญฺจหงฺเคหิ สมนฺนาคโต สํฆเภทโก น อาปายิโก น เนรยิโก น กปฺปฏฺโฐ น อเตกิจฺโฉ. กตเมหิ ปญฺจหิ, อิธุปาลิ ภิกฺขุ อธมฺมํ “ธมฺโม”ติ ทีเปติ, ธมฺมํ “อธมฺโม”ติ ทีเปติ, อวินยํ “วินโย”ติ ทีเปติ, วินยํ “อวินโย”ติ ทีเปติ, อวินิธาย ทิฏฺฐิํ อนุสฺสาวเนน. อิเมหิ โข อุปาลิ ปญฺจหงฺเคหิ สมนฺนาคโต สํฆเภทโก น อาปายิโก น เนรยิโก น กปฺปฏฺโฐ น อเตกิจฺโฉ.}

\prose{อปเรหิปิ อุปาลิ ปญฺจหงฺเคหิ สมนฺนาคโต สํฆเภทโก น อาปายิโก น เนรยิโก น กปฺปฏฺโฐ น อเตกิจฺโฉ. กตเมหิ ปญฺจหิ, อิธุปาลิ ภิกฺขุ อธมฺมํ “ธมฺโม”ติ ทีเปติ, ธมฺมํ “อธมฺโม”ติ ทีเปติ, อวินยํ “วินโย”ติ ทีเปติ, วินยํ “อวินโย”ติ ทีเปติ, อวินิธาย ทิฏฺฐิํ สลากคฺคาเหน. อิเมหิ โข อุปาลิ ปญฺจหงฺเคหิ สมนฺนาคโต สํฆเภทโก น อาปายิโก น เนรยิโก น กปฺปฏฺโฐ น อเตกิจฺโฉ.}

\prose{อปเรหิปิ อุปาลิ ปญฺจหงฺเคหิ สมนฺนาคโต สํฆเภทโก น อาปายิโก น เนรยิโก น กปฺปฏฺโฐ น อเตกิจฺโฉ. กตเมหิ ปญฺจหิ, อิธุปาลิ ภิกฺขุ อธมฺมํ “ธมฺโม”ติ ทีเปติ, ธมฺมํ “อธมฺโม”ติ ทีเปติ, อวินยํ “วินโย”ติ ทีเปติ, วินยํ “อวินโย”ติ ทีเปติ, อวินิธาย ขนฺติํ กมฺเมน ฯเปฯ อวินิธาย ขนฺติํ อุทฺเทเสน ฯเปฯ อวินิธาย ขนฺติํ โวหรนฺโต ฯเปฯ อวินิธาย ขนฺติํ อนุสฺสาวเนน ฯเปฯ อวินิธาย ขนฺติํ สลากคฺคาเหน. อิเมหิ โข อุปาลิ ปญฺจหงฺเคหิ สมนฺนาคโต สํฆเภทโก น อาปายิโก น เนรยิโก น กปฺปฏฺโฐ น อเตกิจฺโฉ.}

\prose{\csromanfolio{354}อปเรหิปิ อุปาลิ ปญฺจหงฺเคหิ สมนฺนาคโต สํฆเภทโก น อาปายิโก น เนรยิโก น กปฺปฏฺโฐ น อเตกิจฺโฉ. กตเมหิ ปญฺจหิ, อิธุปาลิ ภิกฺขุ อธมฺมํ “ธมฺโม”ติ ทีเปติ, ธมฺมํ “อธมฺโม”ติ ทีเปติ, อวินยํ “วินโย”ติ ทีเปติ, วินยํ “อวินโย”ติ ทีเปติ, อวินิธาย รุจิํ กมฺเมน ฯเปฯ อวินิธาย รุจิํ อุทฺเทเสน ฯเปฯ อวินิธาย รุจิํ โวหรนฺโต ฯเปฯ อวินิธาย รุจิํ อนุสฺสาวเนน ฯเปฯ อวินิธาย รุจิํ สลากคฺคาเหน. อิเมหิ โข อุปาลิ ปญฺจหงฺเคหิ สมนฺนาคโต สํฆเภทโก น อาปายิโก น เนรยิโก น กปฺปฏฺโฐ น อเตกิจฺโฉ.}

\prose{อปเรหิปิ อุปาลิ ปญฺจหงฺเคหิ สมนฺนาคโต สํฆเภทโก น อาปายิโก น เนรยิโก น กปฺปฏฺโฐ น อเตกิจฺโฉ. กตเมหิ ปญฺจหิ, อิธุปาลิ ภิกฺขุ อธมฺมํ “ธมฺโม”ติ ทีเปติ, ธมฺมํ “อธมฺโม”ติ ทีเปติ, อวินยํ “วินโย”ติ ทีเปติ, วินยํ “อวินโย”ติ ทีเปติ, อวินิธาย สญฺญํ กมฺเมน ฯเปฯ อวินิธาย สญฺญํ อุทฺเทเสน ฯเปฯ อวินิธาย สญฺญํ โวหรนฺโต ฯเปฯ อวินิธาย สญฺญํ อนุสฺสาวเนน ฯเปฯ อวินิธาย สญฺญํ สลากคฺคาเหน. อิเมหิ โข อุปาลิ ปญฺจหงฺเคหิ สมนฺนาคโต สํฆเภทโก น อาปายิโก น เนรยิโก น กปฺปฏฺโฐ น อเตกิจฺโฉติ.}

\vspace{\tipitakaparskip}
\csromancenter{ทุติยสํฆเภทกวคฺโค นิฏฺฐิโต ทฺวาทสโม.}
\csromancenter{ตสฺสุทฺทานํ.}
\setlength{\csromangathapagewidth}{0pt}
\setlength{\csromangathawakwidth}{0pt}
\csromangathameasuregroup{อวินิธาย ทิฏฺฐิํ กมฺเมน, \\ อนุสฺสาวเน สลาเกน, \\ ขนฺติํ รุจิํ จ สญฺญญฺจ, \\ เหฏฺฐิเม กณฺหปกฺขมฺหิ,}{\csromangathabat{อวินิธาย ทิฏฺฐิํ กมฺเมน,}{อุทฺเทเส โวหเรน จ.} \\ \csromangathabat{อนุสฺสาวเน สลาเกน,}{ปญฺเจเต ทิฏฺฐินิสฺสิตา.} \\ \csromangathabat{ขนฺติํ รุจิํ จ สญฺญญฺจ,}{ตโย เต ปญฺจธา นยา.} \\ \csromangathabat{เหฏฺฐิเม กณฺหปกฺขมฺหิ,}{สมวีสติ วิธี ยถา.}}
\csromangathasetleft{อวินิธาย ทิฏฺฐิํ กมฺเมน, \\ อนุสฺสาวเน สลาเกน, \\ ขนฺติํ รุจิํ จ สญฺญญฺจ, \\ เหฏฺฐิเม กณฺหปกฺขมฺหิ,}
\csromangathagroup{\csromangathastanza{\csromangathabat{อวินิธาย ทิฏฺฐิํ กมฺเมน,}{อุทฺเทเส โวหเรน จ.} \\ \csromangathabat{อนุสฺสาวเน สลาเกน,}{ปญฺเจเต ทิฏฺฐินิสฺสิตา.}}\csromangathastanza{\csromangathabat{ขนฺติํ รุจิํ จ สญฺญญฺจ,}{ตโย เต ปญฺจธา นยา.} \\ \csromangathabat{เหฏฺฐิเม กณฺหปกฺขมฺหิ,}{สมวีสติ วิธี ยถา.}}}

\vspace{\tipitakaparskip}
\csromancenter{ตเถว สุกฺกปกฺขมฺหิ, สมวีสติ ชานถาติ.}
\csromansectionrule
\csromanmatikaanchor{mtk.232}
\csromantocmarkref{subparagraph}{13. อาวาสิกวคฺค}{mtk.232}
\bookmark[dest={mtk.232},level=5]{13. อาวาสิกวคฺค}
\csromanheader{5}{13. \textbf{อาวาสิกวคฺค}}
\proseitem{461}{กติหิ นุ โข ภนฺเต องฺเคหิ สมนฺนาคโต อาวาสิโก ภิกฺขุ ยถาภตํ นิกฺขิตฺโต, เอวํ นิรเยติ. ปญฺจหุปาลิ องฺเคหิ สมนฺนาคโต อาวาสิโก ภิกฺขุ ยถาภตํ นิกฺขิตฺโต, เอวํ \csromanfolio{355}นิรเย. กตเมหิ ปญฺจหิ, ฉนฺทาคติํ คจฺฉติ, โทสาคติํ คจฺฉติ, โมหาคติํ คจฺฉติ, ภยาคติํ คจฺฉติ, สํฆิกํ ปุคฺคลิก\-ปริโภเคน ปริภุญฺชติ. อิเมหิ โข อุปาลิ ปญฺจหงฺเคหิ สมนฺนาคโต อาวาสิโก ภิกฺขุ ยถาภตํ นิกฺขิตฺโต, เอวํ นิรเย.}

\prose{ปญฺจหุปาลิ องฺเคหิ สมนฺนาคโต อาวาสิโก ภิกฺขุ ยถาภตํ นิกฺขิตฺโต, เอวํ สคฺเค. กตเมหิ ปญฺจหิ, น ฉนฺทาคติํ คจฺฉติ, น โทสาคติํ คจฺฉติ, น โมหาคติํ คจฺฉติ, น ภยาคติํ คจฺฉติ, สํฆิกํ น ปุคฺคลิก\-ปริโภเคน ปริภุญฺชติ, อิเมหิ โข อุปาลิ ปญฺจหงฺเคหิ สมนฺนาคโต อาวาสิโก ภิกฺขุ ยถาภตํ นิกฺขิตฺโต, เอวํ สคฺเคติ.}

\proseitem{462}{กติ นุ โข ภนฺเต อธมฺมิกา วินย\-พฺยากรณาติ, ปญฺจิเม อุปาลิ อธมฺมิกา วินย\-พฺยากรณา. กตเม ปญฺจ, อิธุปาลิ ภิกฺขุ อธมฺมํ “ธมฺโม”ติ ปริณาเมติ, ธมฺมํ “อธมฺโม”ติ ปริณาเมติ, อวินยํ “วินโย”ติ ปริณาเมติ, วินยํ “อวินโย”ติ ปริณาเมติ, อปญฺญตฺตํ ปญฺญเปติ, ปญฺญตฺตํ สมุจฺฉินฺทติ. อิเม โข อุปาลิ ปญฺจ อธมฺมิกา วินย\-พฺยากรณา. ปญฺจิเม อุปาลิ ธมฺมิกา วินย\-พฺยากรณา. กตเม ปญฺจ, อิธุปาลิ ภิกฺขุ อธมฺมํ “อธมฺโม”ติ ปริณาเมติ, ธมฺมํ “ธมฺโม”ติ ปริณาเมติ, อวินยํ “อวินโย”ติ ปริณาเมติ, วินยํ “วินโย”ติ ปริณาเมติ, อปญฺญตฺตํ น ปญฺญเปติ, ปญฺญตฺตํ น สมุจฺฉินฺทติ. อิเม โข อุปาลิ ปญฺจ ธมฺมิกา วินย\-พฺยากรณาติ.}

\proseitem{463}{\csromansymbolfootnote{*}{อํ 2. 241 ปิฏฺเฐปิ.}กติหิ นุ โข ภนฺเต องฺเคหิ สมนฺนาคโต ภตฺตุทฺเทสโก ยถาภตํ นิกฺขิตฺโต, เอวํ นิรเยติ. ปญฺจหุปาลิ องฺเคหิ สมนฺนาคโต ภตฺตุทฺเทสโก ยถาภตํ นิกฺขิตฺโต, เอวํ นิรเย. กตเมหิ ปญฺจหิ, ฉนฺทาคติํ คจฺฉติ, โทสาคติํ คจฺฉติ, โมหาคติํ คจฺฉติ, ภยาคติํ คจฺฉติ, อุทฺทิฏฺฐา\-นุทฺทิฏฺฐํ น ชานาติ. อิเมหิ โข อุปาลิ ปญฺจหงฺเคหิ สมนฺนาคโต ภตฺตุทฺเทสโก ยถาภตํ นิกฺขิตฺโต, เอวํ นิรเย.}

\prose{\csromansymbolmark{*}ปญฺจหุปาลิ องฺเคหิ สมนฺนาคโต ภตฺตุทฺเทสโก ยถาภตํ นิกฺขิตฺโต, เอวํ สคฺเค. กตเมหิ ปญฺจหิ, น ฉนฺทาคติํ คจฺฉติ, น โทสาคติํ คจฺฉติ, น โมหาคติํ คจฺฉติ, น ภยาคติํ คจฺฉติ, \csromanfolio{356}อุทฺทิฏฺฐา\-นุทฺทิฏฺฐํ ชานาติ. อิเมหิ โข อุปาลิ ปญฺจหงฺเคหิ สมนฺนาคโต ภตฺตุทฺเทสโก ยถาภตํ นิกฺขิตฺโต, เอวํ สคฺเคติ.}

\proseitem{464}{กติหิ นุ โข ภนฺเต องฺเคหิ สมนฺนาคโต เสนาสน\-ปญฺญาปโก ฯเปฯ. ภณฺฑาคาริโก ฯเปฯ. จีวรปฏิคฺคาหโก ฯเปฯ. จีวรภาชโก ฯเปฯ. ยาคุภาชโก ฯเปฯ. ผลภาชโก ฯเปฯ. ขชฺชภาชโก ฯเปฯ. อปฺปมตฺตก\-วิสฺสชฺชโก ฯเปฯ. สาฏิย\-คฺคาหาปโก ฯเปฯ. ปตฺตคฺคาหาปโก ฯเปฯ. อารามิกเปสโก ฯเปฯ. สามเณรเปสโก ยถาภตํ นิกฺขิตฺโต, เอวํ นิรเยติ. ปญฺจหุปาลิ องฺเคหิ สมนฺนาคโต สามเณรเปสโก ยถาภตํ นิกฺขิตฺโต, เอวํ นิรเย. กตเมหิ ปญฺจหิ, ฉนฺทาคติํ คจฺฉติ, โทสาคติํ คจฺฉติ, โมหาคติํ คจฺฉติ, ภยาคติํ คจฺฉติ, เปสิตาเปสิตํ น ชานาติ. อิเมหิ โข อุปาลิ ปญฺจหงฺเคหิ สมนฺนาคโต สามเณรเปสโก ยถาภตํ นิกฺขิตฺโต, เอวํ นิรเย. ปญฺจหุปาลิ องฺเคหิ สมนฺนาคโต สามเณรเปสโก ยถาภตํ นิกฺขิตฺโต, เอวํ สคฺเค. กตเมหิ ปญฺจหิ, น ฉนฺทาคติํ คจฺฉติ, น โทสาคติํ คจฺฉติ, น โมหาคติํ คจฺฉติ, น ภยาคติํ คจฺฉติ, เปสิตาเปสิตํ ชานาติ. อิเมหิ โข อุปาลิ ปญฺจหงฺเคหิ สมนฺนาคโต สามเณรเปสโก ยถาภตํ นิกฺขิตฺโต, เอวํ สคฺเคติ.}

\vspace{\tipitakaparskip}
\csromancenter{อาวาสิกวคฺโค นิฏฺฐิโต เตรสโม.}
\csromancenter{ตสฺสุทฺทานํ.}
\setlength{\csromangathapagewidth}{0pt}
\setlength{\csromangathawakwidth}{0pt}
\csromangathameasuregroup{อาวาสิกพฺยากรณา, \\ ภณฺฑจีวรคฺคาโห จ, \\ ยาคุ ผลํ ขชฺชกญฺจ,}{\csromangathabat{อาวาสิกพฺยากรณา,}{ภตฺตุเสนาสนานิ จ.} \\ \csromangathabat{ภณฺฑจีวรคฺคาโห จ,}{จีวรสฺส จ ภาชโก.} \\ \csromangathabat{ยาคุ ผลํ ขชฺชกญฺจ,}{อปฺปสาฏิยคาหโก.}}
\csromangathasetleft{อาวาสิกพฺยากรณา, \\ ภณฺฑจีวรคฺคาโห จ, \\ ยาคุ ผลํ ขชฺชกญฺจ,}
\csromangathagroup{\csromangathastanza{\csromangathabat{อาวาสิกพฺยากรณา,}{ภตฺตุเสนาสนานิ จ.} \\ \csromangathabat{ภณฺฑจีวรคฺคาโห จ,}{จีวรสฺส จ ภาชโก.} \\ \csromangathabat{ยาคุ ผลํ ขชฺชกญฺจ,}{อปฺปสาฏิยคาหโก.}}}

\prose{ปตฺโต อารามิโก เจว, สามเณเรน เปสโกติ.}
\csromansectionrule

\csromanmatikaanchor{mtk.233}
\csromantocmarkref{subparagraph}{14. กถินตฺถารวคฺค}{mtk.233}
\bookmark[dest={mtk.233},level=5]{14. กถินตฺถารวคฺค}
\csromanheader{5}{14. กถินตฺถาร\-วคฺค}
\proseitem{465}{กติ นุ โข ภนฺเต อานิสํสา กถิน\-ตฺถาเรติ. ปญฺจิเม อุปาลิ อานิสํสา กถินตฺถาเร. กตเม ปญฺจ, อนามนฺตจาโร, อสมาทานจาโร, คณโภชนํ, ยาวทตฺถ\-จีวรํ, โย ตตฺถ \csromanfolio{357}จีวรุปฺปาโท โส เนสํ ภวิสฺสติ. อิเม โข อุปาลิ ปญฺจ อานิสํสา กถิน\-ตฺถาเรติ.}

\proseitem{466}{กติ นุ โข ภนฺเต อาทีนวา มุฏฺฐสฺสติสฺส อสมฺปชานสฺส นิทฺทํ โอกฺกมโตติ. ปญฺจิเม อุปาลิ อาทีนวา มุฏฺฐสฺสติสฺส อสมฺปชานสฺส นิทฺทํ โอกฺกมโต. กตเม ปญฺจ, ทุกฺขํ สุปติ, ทุกฺขํ ปฏิพุชฺฌติ, ปาปกํ สุปินํ ปสฺสติ, เทวตา น รกฺขนฺติ, อสุจิ มุจฺจติ. อิเม โข อุปาลิ ปญฺจ อาทีนวา มุฏฺฐสฺสติสฺส อสมฺปชานสฺส นิทฺทํ โอกฺกมโต. ปญฺจิเม อุปาลิ อานิสํสา อุปฏฺฐิต\-สฺสติสฺส สมฺปชานสฺส นิทฺทํ โอกฺกมโต. กตเม ปญฺจ, สุขํ สุปติ, สุขํ ปฏิพุชฺฌติ, น ปาปกํ สุปินํ ปสฺสติ, เทวตา รกฺขนฺติ, อสุจิ นมุจฺจติ. อิเม โข อุปาลิ ปญฺจ อานิสํสา อุปฏฺฐิต\-สฺสติสฺส สมฺปชานสฺส นิทฺทํ โอกฺกมโตติ.}

\proseitem{467}{กติ นุ โข ภนฺเต อวนฺทิยาติ. ปญฺจิเม อุปาลิ อวนฺทิยา. กตเม ปญฺจ, อนฺตรฆรํ ปวิฏฺโฐ อวนฺทิโย, รจฺฉคโต อวนฺทิโย, โอตมสิโก อวนฺทิโย, อสมนฺนาหรนฺโต อวนฺทิโย, สุตฺโต อวนฺทิโย. อิเม โข อุปาลิ ปญฺจ อวนฺทิยา.}

\prose{อปเรปิ อุปาลิ ปญฺจ อวนฺทิยา. กตเม ปญฺจ, ยาคุปาเน อวนฺทิโย, ภตฺตคฺเค อวนฺทิโย, เอกาวตฺโต อวนฺทิโย, อญฺญวิหิโต อวนฺทิโย, นคฺโค อวนฺทิโย. อิเม โข อุปาลิ ปญฺจ อวนฺทิยา.}

\prose{อปเรปิ อุปาลิ ปญฺจ อวนฺทิยา. กตเม ปญฺจ, ขาทนฺโต อวนฺทิโย, ภุญฺชนฺโต อวนฺทิโย, อุจฺจารํ กโรนฺโต อวนฺทิโย, ปสฺสาวํ กโรนฺโต อวนฺทิโย, อุกฺขิตฺตโก อวนฺทิโย. อิเม โข อุปาลิ ปญฺจ อวนฺทิยา.}

\prose{อปเรปิ อุปาลิ ปญฺจ อวนฺทิยา. กตเม ปญฺจ, ปุเร อุปสมฺปนฺเนน ปจฺฉา อุปสมฺปนฺโน อวนฺทิโย, อนุปสมฺปนฺโน อวนฺทิโย, นานาสํวาสโก วุฑฺฒตโร อธมฺมวาที อวนฺทิโย, มาตุคาโม อวนฺทิโย, ปณฺฑโก อวนฺทิโย. อิเม โข อุปาลิ ปญฺจ อวนฺทิยา.}

\prose{อปเรปิ อุปาลิ ปญฺจ อวนฺทิยา. กตเม ปญฺจ, ปาริวาสิโก อวนฺทิโย, มูลาย\-ปฏิกสฺสนา\-รโห อวนฺทิโย, มานตฺตารโห อวนฺทิโย, มานตฺตจาริโก อวนฺทิโย, อพฺภานารโห อวนฺทิโย. อิเม โข อุปาลิ ปญฺจ อวนฺทิยาติ.}

\proseitem{468}{\csromanfolio{358}กติ นุ โข ภนฺเต วนฺทิยาติ. ปญฺจิเม อุปาลิ วนฺทิยา. กตเม ปญฺจ, ปจฺฉา อุปสมฺปนฺเนน ปุเร อุปสมฺปนฺโน วนฺทิโย, นานาสํวาสโก วุฑฺฒตโร ธมฺมวาที วนฺทิโย, อาจริโย วนฺทิโย, อุปชฺฌาโย วนฺทิโย, สเทวเก โลเก สมารเก สพฺรหฺมเก สสฺสมณพฺรหฺมณิยา ปชาย สเทว\-มนุสฺสาย ตถาคโต อรหํ สมฺมาสมฺพุทฺโธ วนฺทิโย. อิเม โข อุปาลิ ปญฺจ วนฺทิยาติ.}

\proseitem{469}{นวกตเรน ภนฺเต ภิกฺขุนา วุฑฺฒตรสฺส ภิกฺขุโน ปาเท วนฺทนฺเตน กติ ธมฺเม อชฺฌตฺตํ อุปฏฺฐาเปตฺวา ปาทา วนฺทิตพฺพาติ. นวกตเรนุ\-ปาลิ ภิกฺขุนา วุฑฺฒตรสฺส ภิกฺขุโน ปาเท วนฺทนฺเตน ปญฺจ ธมฺเม อชฺฌตฺตํ อุปฏฺฐาเปตฺวา ปาทา วนฺทิตพฺพา. กตเม ปญฺจ, นวกตเรนุ\-ปาลิ ภิกฺขุนา วุฑฺฒตรสฺส ภิกฺขุโน ปาเท วนฺทนฺเตน เอกํสํ อุตฺตราสงฺคํ กริตฺวา อญฺชลิํ ปคฺคเหตฺวา อุโภหิ ปาณิตเลหิ ปาทานิ ปริสมฺ\-พาหนฺเตน เปมญฺจ คารวญฺจ อุปฏฺฐาเปตฺวา ปาทา วนฺทิตพฺพา. นวกตเรนุ\-ปาลิ ภิกฺขุนา วุฑฺฒตรสฺส ภิกฺขุโน ปาเท วนฺทนฺเตน อิเม ปญฺจ ธมฺเม อชฺฌตฺตํ อุปฏฺฐาเปตฺวา ปาทา วนฺทิตพฺพาติ.}

\vspace{\tipitakaparskip}
\csromancenter{กถินตฺถาร\-วคฺโค นิฏฺฐิโต จุทฺทสโม.}
\csromancenter{ตสฺสุทฺทานํ.}
\setlength{\csromangathapagewidth}{0pt}
\setlength{\csromangathawakwidth}{0pt}
\csromangathameasuregroup{กถินตฺถารนิทฺทา จ, \\ ปุเร จ ปาริวาสิ จ,}{\csromangathabat{กถินตฺถารนิทฺทา จ,}{อนฺตรา ยาคุขาทเน.} \\ \csromangathabat{ปุเร จ ปาริวาสิ จ,}{วนฺทิโย วนฺทิตพฺพกนฺติ.}}
\csromangathasetleft{กถินตฺถารนิทฺทา จ, \\ ปุเร จ ปาริวาสิ จ,}
\csromangathagroupwithcloser{\csromangathastanza{\csromangathabat{กถินตฺถารนิทฺทา จ,}{อนฺตรา ยาคุขาทเน.} \\ \csromangathabat{ปุเร จ ปาริวาสิ จ,}{วนฺทิโย วนฺทิตพฺพกนฺติ.}}}{อุปาลิปญฺจกํ นิฏฺฐิตํ.}
\csromancenter{เตสํ วคฺคานํ อุทฺทานํ.}
\setlength{\csromangathapagewidth}{0pt}
\setlength{\csromangathawakwidth}{0pt}
\csromangathameasuregroup{อนิสฺสิเตน กมฺมญฺจ, \\ โจทนา จ ธุตงฺคา จ, \\ อุพฺพาหิกาธิกรณํ, \\ อาวาสิกา กถินญฺจ,}{\csromangathabat{อนิสฺสิเตน กมฺมญฺจ,}{โวหาราวิกมฺเมน จ.} \\ \csromangathabat{โจทนา จ ธุตงฺคา จ,}{มุสา ภิกฺขุนิเมว จ.} \\ \csromangathabat{อุพฺพาหิกาธิกรณํ,}{เภทกา ปญฺจมา ปุเร.} \\ \csromangathabat{อาวาสิกา กถินญฺจ,}{จุทฺทสา สุปฺปกาสิตาติ.}}
\csromangathasetleft{อนิสฺสิเตน กมฺมญฺจ, \\ โจทนา จ ธุตงฺคา จ, \\ อุพฺพาหิกาธิกรณํ, \\ อาวาสิกา กถินญฺจ,}
\csromangathagroup{\csromangathastanza{\csromangathabat{อนิสฺสิเตน กมฺมญฺจ,}{โวหาราวิกมฺเมน จ.} \\ \csromangathabat{โจทนา จ ธุตงฺคา จ,}{มุสา ภิกฺขุนิเมว จ.}}\csromangathastanza{\csromangathabat{อุพฺพาหิกาธิกรณํ,}{เภทกา ปญฺจมา ปุเร.} \\ \csromangathabat{อาวาสิกา กถินญฺจ,}{จุทฺทสา สุปฺปกาสิตาติ.}}}

\csromanmatikaanchor{mtk.234}
\csromantocmarknopage{section}{อตฺถาปตฺติสมุฏฺฐาน}{mtk.234}
\bookmark[dest={mtk.234},level=1]{อตฺถาปตฺติสมุฏฺฐาน}
\csromanheader{1}{อตฺถา\-ปตฺติสมุฏฺฐาน}
\csromanmatikaanchor{mtk.235}
\csromantocmarkref{subsection}{1. ปาราชิก}{mtk.235}
\bookmark[dest={mtk.235},level=2]{1. ปาราชิก}
\csromanheader{2}{1. \textbf{ปาราชิก}}
\proseitem{470}{\csromanfolio{359}อตฺถาปตฺติ อจิตฺตโก อาปชฺชติ, สจิตฺตโก วุฏฺฐาติ, อตฺถาปตฺติ สจิตฺตโก อาปชฺชติ, อจิตฺตโก วุฏฺฐาติ, อตฺถาปตฺติ อจิตฺตโก อาปชฺชติ, อจิตฺตโก วุฏฺฐาติ, อตฺถาปตฺติ สจิตฺตโก อาปชฺชติ, สจิตฺตโก วุฏฺฐาติ, อตฺถาปตฺติ กุสลจิตฺโต อาปชฺชติ, กุสลจิตฺโต วุฏฺฐาติ, อตฺถาปตฺติ กุสลจิตฺโต อาปชฺชติ, อกุสลจิตฺโต วุฏฺฐาติ, อตฺถาปตฺติ กุสลจิตฺโต อาปชฺชติ, อพฺยากตจิตฺโต วุฏฺฐาติ, อตฺถาปตฺติ อกุสลจิตฺโต อาปชฺชติ, กุสลจิตฺโต วุฏฺฐาติ, อตฺถาปตฺติ อกุสลจิตฺโต อาปชฺชติ, อกุสลจิตฺโต วุฏฺฐาติ, อตฺถาปตฺติ อกุสลจิตฺโต อาปชฺชติ, อพฺยากตจิตฺโต วุฏฺฐาติ, อตฺถาปตฺติ อพฺยากตจิตฺโต อาปชฺชติ, กุสลจิตฺโต วุฏฺฐาติ, อตฺถาปตฺติ อพฺยากตจิตฺโต อาปชฺชติ, อกุสลจิตฺโต วุฏฺฐาติ, อตฺถาปตฺติ อพฺยากตจิตฺโต อาปชฺชติ, อพฺยากตจิตฺโต วุฏฺฐาติ.}

\prose{ปฐมํ ปาราชิกํ กติหิ สมุฏฺฐาเนหิ สมุฏฺฐาติ. ปฐมํ ปาราชิกํ เอเกน สมุฏฺฐาเนน สมุฏฺฐาติ, กายโต จ จิตฺตโต จ สมุฏฺฐาติ, น วาจโต.}

\prose{ทุติยํ ปาราชิกํ กติหิ สมุฏฺฐาเนหิ สมุฏฺฐาติ. ทุติยํ ปาราชิกํ ตีหิ สมุฏฺฐาเนหิ สมุฏฺฐาติ, สิยา กายโต จ จิตฺตโต จ สมุฏฺฐาติ, น วาจโต, สิยา วาจโต จ จิตฺตโต จ สมุฏฺฐาติ, น กายโต, สิยา กายโต จ วาจโต จ จิตฺตโต จ สมุฏฺฐาติ.}

\prose{ตติยํ ปาราชิกํ กติหิ สมุฏฺฐาเนหิ สมุฏฺฐาติ. ตติยํ ปาราชิกํ ตีหิ สมุฏฺฐาเนหิ สมุฏฺฐาติ, สิยา กายโต จ จิตฺตโต จ สมุฏฺฐาติ, น วาจโต, สิยา วาจโต จ จิตฺตโต จ สมุฏฺฐาติ, น กายโต, สิยา กายโต จ วาจโต จ จิตฺตโต จ สมุฏฺฐาติ.}

\prosewithcloserruled{จตุตฺถํ ปาราชิกํ กติหิ สมุฏฺฐาเนหิ สมุฏฺฐาติ. จตุตฺถํ ปาราชิกํ ตีหิ สมุฏฺฐาเนหิ สมุฏฺฐาติ, สิยา กายโต จ จิตฺตโต จ \csromanfolio{360}สมุฏฺฐาติ, น วาจโต, สิยา วาจโต จ จิตฺตโต จ สมุฏฺฐาติ, น กายโต, สิยา กายโต จ วาจโต จ จิตฺตโต จ สมุฏฺฐาติ.}{จตฺตาโร ปาราชิกา นิฏฺฐิตา.}
\csromanmatikaanchor{mtk.236}
\csromantocmarkref{subsection}{2. สํฆาทิเสส}{mtk.236}
\bookmark[dest={mtk.236},level=2]{2. สํฆาทิเสส}
\csromanheader{2}{2. \textbf{สํฆาทิเสส}}
\proseitem{471}{อุปกฺกมิตฺวา อสุจิํ โมเจนฺตสฺส สํฆาทิเสโส กติหิ สมุฏฺฐาเนหิ สมุฏฺฐาติ. อุปกฺกมิตฺวา อสุจิํ โมเจนฺตสฺส สํฆาทิเสโส เอเกน สมุฏฺฐาเนน สมุฏฺฐาติ, กายโต จ จิตฺตโต จ สมุฏฺฐาติ, น วาจโต.}

\prose{มาตุคาเมน สทฺธิํ กายสํสคฺคํ สมา\-ป\-ชฺช\-นฺตสฺส สํฆทิเสโส กติหิ สมุฏฺฐาเนหิ สมุฏฺฐาติ. มาตุคาเมน สทฺธิํ กายสํสคฺคํ สมา\-ป\-ชฺช\-นฺตสฺส สํฆาทิเสโส เอเกน สมุฏฺฐาเนน สมุฏฺฐาติ, กายโต จ จิตฺตโต จ สมุฏฺฐาติ, น วาจโต.}

\prose{มาตุคามํ ทุฏฺฐุลฺลาหิ วาจาหิ โอภาเสนฺตสฺส สํฆาทิเสโส กติหิ สมุฏฺฐาเนหิ สมุฏฺฐาติ. มาตุคามํ ทุฏฺฐุลฺลาหิ วาจาหิ โอภาเสนฺตสฺส สํฆาทิเสโส ตีหิ สมุฏฺฐาเนหิ สมุฏฺฐาติ, สิยา กายโต จ จิตฺตโต จ สมุฏฺฐาติ, น วาจโต, สิยา วาจโต จ จิตฺตโต จ สมุฏฺฐาติ, น กายโต, สิยา กายโต จ วาจโต จ จิตฺตโต จ สมุฏฺฐาติ.}

\prose{มาตุคามสฺส สนฺติเก อตฺต\-กาม\-ปาริจริยาย วณฺณํ ภาสนฺตสฺส สํฆาทิเสโส กติหิ สมุฏฺฐาเนหิ สมุฏฺฐาติ. มาตุคามสฺส สนฺติเก อตฺต\-กาม\-ปาริจริยาย วณฺณํ ภาสนฺตสฺส สํฆาทิเสโส ตีหิ สมุฏฺฐาเนหิ}

\prose{สญฺจริตฺตํ สมา\-ป\-ชฺช\-นฺตสฺส สํฆาทิเสโส กติหิ สมุฏฺฐาเนหิ สมุฏฺฐาติ. สญฺจริตฺตํ สมา\-ป\-ชฺช\-นฺตสฺส สํฆาทิเสโส ฉหิ สมุฏฺฐาเนหิ สมุฏฺฐาติ, สิยา กายโต สมุฏฺฐาติ, น วาจโต, น จิตฺตโต, สิยา วาจโต สมุฏฺฐาติ, น กายโต, น จิตฺตโต, สิยา กายโต จ วาจโต จ สมุฏฺฐาติ, น จิตฺตโต, สิยา กายโต จ จิตฺตโต จ สมุฏฺฐาติ, น วาจโต, สิยา วาจโต จ จิตฺตโต จ \csromanfolio{361}สมุฏฺฐาติ, น กายโต, สิยา กายโต จ วาจโต จ จิตฺตโต จ สมุฏฺฐาติ.}

\prose{สญฺญาจิกาย กุฏิํ การาเปนฺตสฺส สํฆาทิเสโส กติหิ สมุฏฺฐาเนหิ สมุฏฺฐาติ. สญฺญาจิกาย กุฏิํ การาเปนฺตสฺส สํฆาทิเสโส ฉหิ สมุฏฺฐาเนหิ สมุฏฺฐาติ ฯเปฯ.}

\prose{มหลฺลกํ วิหารํ การาเปนฺตสฺส สํฆาทิเสโส กติหิ สมุฏฺฐาเนหิ สมุฏฺฐาติ. มหลฺลกํ วิหารํ การาเปนฺตสฺส สํฆาทิเสโส ฉหิ สมุฏฺฐาเนหิ สมุฏฺฐาติ ฯเปฯ.}

\prose{ภิกฺขุํ อมูลเกน ปาราชิเกน ธมฺเมน อนุทฺธํเส\-นฺตสฺส สํฆาทิเสโส กติหิ สมุฏฺฐาเนหิ สมุฏฺฐาติ. ภิกฺขุํ อมูลเกน ปาราชิเกน ธมฺเมน อนุทฺธํเส\-นฺตสฺส สํฆาทิเสโส ตีหิ สมุฏฺฐาเนหิ}

\prose{ภิกฺขุํ อญฺญภาคิยสฺส อธิกรณสฺส กิญฺจิ เทสํ เลสมตฺตํ อุปาทาย ปาราชิเกน ธมฺเมน อนุทฺธํเส\-นฺตสฺส สํฆาทิเสโส กติหิ สมุฏฺฐาเนหิ สมุฏฺฐาติ. ภิกฺขุํ อญฺญภาคิยสฺส อธิกรณสฺส กิญฺจิ เทสํ เลสมตฺตํ อุปาทาย ปาราชิเกน ธมฺเมน อนุทฺธํเส\-นฺตสฺส สํฆาทิเสโส ตีหิ สมุฏฺฐาเนหิ สมุฏฺฐาติ ฯเปฯ.}

\prose{สํฆเภทกสฺส ภิกฺขุโน ยาวตติยํ สมนุภาสนาย น ปฏินิสฺสชฺชนฺตสฺส สํฆาทิเสโส กติหิ สมุฏฺฐาเนหิ สมุฏฺฐาติ. สํฆเภทกสฺส ภิกฺขุโน ยาวตติยํ สมนุภาสนาย น ปฏินิสฺสชฺชนฺตสฺส สํฆาทิเสโส เอเกน สมุฏฺฐาเนน สมุฏฺฐาติ, กายโต จ วาจโต จ จิตฺตโต จ สมุฏฺฐาติ.}

\prose{เภทกา\-นุวตฺตกานํ ภิกฺขูนํ ยาวตติยํ สมนุภาสนาย น ปฏินิสฺสชฺชนฺตานํ สํฆาทิเสโส กติหิ สมุฏฺฐาเนหิ สมุฏฺฐาติ. เภทกา\-นุวตฺตกานํ ภิกฺขูนํ ยาวตติยํ สมนุภาสนาย น ปฏินิสฺสชฺชนฺตานํ สํฆาเทเสโส เอเกน สมุฏฺฐาเนน สมุฏฺฐาติ, กายโต จ วาจโต จ จิตฺตโต จ สมุฏฺฐาติ.}

\prose{ทุพฺพจสฺส ภิกฺขุโน ยาวตติยํ สมนุภาสนาย น ปฏินิสฺสชฺชนฺตสฺส สํฆาทิเสโส กติหิ สมุฏฺฐาเนหิ สมุฏฺฐาติ. ทุพฺพจสฺส ภิกฺขุโน ยาวตติยํ สมนุภาสนาย ปฏินิสฺสชฺชนฺตสฺส สํฆาทิเสโส เอเกน สมุฏฺฐาเนน สมุฏฺฐาติ, กายโต จ วาจโต จ จิตฺตโต จ สมุฏฺฐาติ.}

\prosewithcloserruled{\csromanfolio{362}กุลทูสกสฺส ภิกฺขุโน ยาวตติยํ สมนุภาสนาย น ปฏินิสฺสชฺชนฺตสฺส สํฆาทิเสโส กติหิ สมุฏฺฐาเนหิ สมุฏฺฐาติ. กุลทูสกสฺส ภิกฺขุโน ยาวตติยํ สมนุภาสนาย น ปฏินิสฺสชฺชนฺตสฺส สํฆาทิเสโส เอเกน สมุฏฺฐาเนน สมุฏฺฐาติ, กายโต จ วาจโต จ จิตฺตโต จ สมุฏฺฐาติ.}{เตรส สํฆาทิเสสา นิฏฺฐิตา.}
\proseitemwithcloserruled{472}{ฯเปฯ อนาทริยํ ปฏิจฺจ อุทเก อุจฺจารํ วา ปสฺสาวํ วา เขฬํ วา กโรนฺตสฺส ทุกฺกฏํ กติหิ สมุฏฺฐาเนหิ สมุฏฺฐาติ. อนาทริยํ ปฏิจฺจ อุทเก อุจฺจรํ วา ปสฺสาวํ เขฬํ วา กโรนฺตสฺส ทุกฺกฏํ เอเกน สมุฏฺฐาเนน สมุฏฺฐาติ, กายโต จ จิตฺตโต จ สมุฏฺฐาติ, น วาจโต.}{เสขิยา นิฏฺฐิตา.}
\csromanmatikaanchor{mtk.237}
\csromantocmarkref{subsection}{3. ปาราชิกาทิ}{mtk.237}
\bookmark[dest={mtk.237},level=2]{3. ปาราชิกาทิ}
\csromanheader{2}{3. \textbf{ปาราชิกาทิ}}
\proseitem{473}{จตฺตาโร ปาราชิกา กติหิ สมุฏฺฐาเนหิ สมุฏฺฐนฺติ. จตฺตาโร ปาราชิกา ตีหิ สมุฏฺฐาเนหิ สมุฏฺฐนฺติ, สิยา กายโต จ จิตฺตโต จ สมุฏฺฐนฺติ, น วาจโต, สิยา วาจโต จ จิตฺตโต จ สมุฏฺฐนฺติ, น กายโต, สิยา กายโต จ วาจโต จ จิตฺตโต จ สมุฏฺฐนฺติ.}

\prose{เตรส สํฆาทิเสสา กติหิ สมุฏฺฐาเนหิ สมุฏฺฐนฺติ. เตรส สํฆาทิเสสา ฉหิ สมุฏฺฐาเนหิ สมุฏฺฐนฺติ, สิยา กายโต สมุฏฺฐนฺติ, น วาจโต, น จิตฺตโต, สิยา วาจโต สมุฏฺฐนฺติ, น กายโต, น จิตฺตโต, สิยา กายโต จ วาจโต จ สมุฏฺฐนฺติ, น จิตฺตโต, สิยา กายโต จ จิตฺตโต จ สมุฏฺฐนฺติ, น วาจโต, สิยา วาจโต จ จิตฺตโต จ สมุฏฺฐนฺติ, น กายโต, สิยา กายโต จ วาจโต จ จิตฺตโต จ สมุฏฺฐนฺติ.}

\prose{เทฺว อนิยตา กติหิ สมุฏฺฐาเนหิ สมุฏฺฐนฺติ. เทฺว อนิยตา ตีหิ สมุฏฺฐาเนหิ สมุฏฺฐนฺติ, สิยา กายโต จ จิตฺตโต จ สมุฏฺฐนฺติ, น \csromanfolio{363}วาจโต, สิยา วาจโต จ จิตฺตโต จ สมุฏฺฐนฺติ, น กายโต, สิยา กายโต จ วาจโต จ จิตฺตโต จ สมุฏฺฐนฺติ.}

\prose{ติํส นิสฺสคฺคิยา ปาจิตฺติยา กติหิ สมุฏฺฐาเนหิ สมุฏฺฐนฺติ. ติํส นิสฺสคฺคิยา ปาจิตฺติยา ฉหิ สมุฏฺฐาเนหิ สมุฏฺฐนฺติ, สิยา กายโต สมุฏฺฐนฺติ, น วาจโต, น จิตฺตโต, สิยา วาจโต สมุฏฺฐนฺติ, น กายโต, น จิตฺตโต, สิยา กายโต จ วาจโต จ สมุฏฺฐนฺติ, น จิตฺตโต, สิยา กายโต จ จิตฺตโต จ สมุฏฺฐนฺติ, น วาจโต, สิยา วาจโต จ จิตฺตโต จ สมุฏฺฐนฺติ, น กายโต, สิยา กายโต จ วาจโต จ จิตฺตโต จ สมุฏฺฐนฺติ.}

\prose{เทฺวนวุติ ปาจิตฺติยา กติหิ สมุฏฺฐาเนหิ สมุฏฺฐนฺติ. เทฺวนวุติ ปาจิตฺติยา ฉหิ สมุฏฺฐาเนหิ สมุฏฺฐนฺติ, สิยา กายโต สมุฏฺฐนฺติ, น วาจโต, น จิตฺตโต, สิยา วาจโต สมุฏฺฐนฺติ, น กายโต, น จิตฺตโต, สิยา กายโต จ วาจโต จ สมุฏฺฐนฺติ, น จิตฺตโต, สิยา กายโต จ จิตฺตโต จ สมุฏฺฐนฺติ, น วาจโต, สิยา วาจโต จ จิตฺตโต จ สมุฏฺฐนฺติ, น กายโต, สิยา กายโต จ วาจโต จ จิตฺตโต จ สมุฏฺฐนฺติ.}

\prose{จตฺตาโร ปาฏิเทสนียา กติหิ สมุฏฺฐาเนหิ สมุฏฺฐนฺติ. จตฺตาโร ปาฏิเทสนียา จตูหิ สมุฏฺฐาเนหิ สมุฏฺฐนฺติ, สิยา กายโต สมุฏฺฐนฺติ, น วาจโต, น จิตฺตโต, สิยา กายโต จ วาจโต จ สมุฏฺฐนฺติ, น จิตฺตโต, สิยา กายโต จ จิตฺตโต จ สมุฏฺฐนฺติ, น วาจโต, สิยา กายโต จ วาจโต จ จิตฺตโต จ สมุฏฺฐนฺติ.}

\prosewithcloser{ปญฺจสตฺตติ เสขิยา กติหิ สมุฏฺฐาเนหิ สมุฏฺฐนฺติ. ปญฺจสตฺตติ เสขิยา ตีหิ สมุฏฺฐาเนหิ สมุฏฺฐนฺติ, สิยา กายโต จ จิตฺตโต จ สมุฏฺฐนฺติ, น วาจโต, สิยา วาจโต จ จิตฺตโต จ สมุฏฺฐนฺติ, น กายโต, สิยา กายโต จ วาจโต จ จิตฺตโต จ สมุฏฺฐนฺติ.}{สมุฏฺฐานํ นิฏฺฐิตํ.}
\csromancenter{ตสฺสุทฺทานํ.}
\setlength{\csromangathapagewidth}{0pt}
\setlength{\csromangathawakwidth}{0pt}
\csromangathameasuregroup{อจิตฺตกุสลา เจว, \\ ยถาธมฺเมน ญาเยน,}{\csromangathabat{อจิตฺตกุสลา เจว,}{สมุฏฺฐานญฺจ สพฺพถา.} \\ \csromangathabat{ยถาธมฺเมน ญาเยน,}{สมุฏฺฐานํ วิชานถาติ.}}
\csromangathasetleft{อจิตฺตกุสลา เจว, \\ ยถาธมฺเมน ญาเยน,}
\csromangathagroup{\csromangathastanza{\csromangathabat{อจิตฺตกุสลา เจว,}{สมุฏฺฐานญฺจ สพฺพถา.} \\ \csromangathabat{ยถาธมฺเมน ญาเยน,}{สมุฏฺฐานํ วิชานถาติ.}}}

\csromanmatikaanchor{mtk.238}
\csromantocmarknopage{section}{ทุติยคาถาสงฺคณิ}{mtk.238}
\bookmark[dest={mtk.238},level=1]{ทุติยคาถาสงฺคณิ}
\csromanheader{1}{ทุติย\-คาถาสงฺคณิก}
\csromanmatikaanchor{mtk.239}
\csromantocmarkref{paragraph}{1. กายิกาทิอาปตฺติ}{mtk.239}
\bookmark[dest={mtk.239},level=4]{1. กายิกาทิอาปตฺติ}
\csromanheader{4}{1. กายิกาทิอาปตฺติ}
\setlength{\csromangathapagewidth}{0pt}
\setlength{\csromangathawakwidth}{0pt}
\csromangathameasuregroup{กติ อาปตฺติโย กายิกา, \\ ฉาเทนฺตสฺส กติ อาปตฺติโย, \\ ฉาปตฺติโย กายิกา, \\ ฉาเทนฺตสฺส ติสฺโส อาปตฺติโย, \\ อรุณุคฺเค กติ อาปตฺติโย, \\ กเตตฺถ อฏฺฐ วตฺถุกา, \\ อรุณุคฺเค ติสฺโส อาปตฺติโย, \\ เอเกตฺถ อฏฺฐ วตฺถุกา, \\ วินยสฺส กติ มูลานิ, \\ วินยครุกา กติ วุตฺตา, \\ วินยสฺส เทฺว มูลานิ, \\ วินยครุกา เทฺว วุตฺตา, \\ คามนฺตเร กติ อาปตฺติโย, \\ กติมํเสสุ ถุลฺลจฺจยํ, \\ คามนฺตเร จตสฺโส อาปตฺติโย, \\ เอกมํเส ถุลฺลจฺจยํ, \\ กติ วาจสิกา รตฺติํ, \\ ททมานสฺส กติ อาปตฺติโย, \\ เทฺว วาจสิกา รตฺติํ, \\ ททมานสฺส ติสฺโส อาปตฺติโย,}{\csromangathabat{กติ อาปตฺติโย กายิกา,}{กติ วาจสิกา กตา.} \\ \csromangathabat{ฉาเทนฺตสฺส กติ อาปตฺติโย,}{กติ สํสคฺคปจฺจยา.} \\ \csromangathabat{ฉาปตฺติโย กายิกา,}{ฉ วาจสิกา กตา.} \\ \csromangathabat{ฉาเทนฺตสฺส ติสฺโส อาปตฺติโย,}{ปญฺจ สํสคฺคปจฺจยา.} \\ \csromangathabat{อรุณุคฺเค กติ อาปตฺติโย,}{กติ ยาวตติยกา.} \\ \csromangathabat{กเตตฺถ อฏฺฐ วตฺถุกา,}{กติหิ สพฺพสงฺคโห.} \\ \csromangathabat{อรุณุคฺเค ติสฺโส อาปตฺติโย,}{เทฺว ยาวตติยกา.} \\ \csromangathabat{เอเกตฺถ อฏฺฐ วตฺถุกา,}{เอเกน สพฺพสงฺคโห.} \\ \csromangathabat{วินยสฺส กติ มูลานิ,}{ยานิ พุทฺเธน ปญฺญตฺตา.} \\ \csromangathabat{วินยครุกา กติ วุตฺตา,}{ทุฏฺฐุลฺลจฺฉาทนา กติ.} \\ \csromangathabat{วินยสฺส เทฺว มูลานิ,}{ยานิ พุทฺเธน ปญฺญตฺตา.} \\ \csromangathabat{วินยครุกา เทฺว วุตฺตา,}{เทฺว ทุฏฺฐุลฺลจฺฉาทนา.} \\ \csromangathabat{คามนฺตเร กติ อาปตฺติโย,}{กติ นทิปารปจฺจยา.} \\ \csromangathabat{กติมํเสสุ ถุลฺลจฺจยํ,}{กติมํเสสุ ทุกฺกฏํ.} \\ \csromangathabat{คามนฺตเร จตสฺโส อาปตฺติโย,}{จตสฺโส นทิปารปจฺจยา.} \\ \csromangathabat{เอกมํเส ถุลฺลจฺจยํ,}{นวมํเสสุ ทุกฺกฏํ.} \\ \csromangathabat{กติ วาจสิกา รตฺติํ,}{กติ วาจสิกา ทิวา.} \\ \csromangathabat{ททมานสฺส กติ อาปตฺติโย,}{ปฏิคฺคณฺหนฺตสฺส กิตฺตกา.} \\ \csromangathabat{เทฺว วาจสิกา รตฺติํ,}{เทฺว วาจสิกา ทิวา.} \\ \csromangathabat{ททมานสฺส ติสฺโส อาปตฺติโย,}{จตฺตาโร จ ปฏิคฺคเห.}}
\csromangathasetleft{กติ อาปตฺติโย กายิกา, \\ ฉาเทนฺตสฺส กติ อาปตฺติโย, \\ ฉาปตฺติโย กายิกา, \\ ฉาเทนฺตสฺส ติสฺโส อาปตฺติโย, \\ อรุณุคฺเค กติ อาปตฺติโย, \\ กเตตฺถ อฏฺฐ วตฺถุกา, \\ อรุณุคฺเค ติสฺโส อาปตฺติโย, \\ เอเกตฺถ อฏฺฐ วตฺถุกา, \\ วินยสฺส กติ มูลานิ, \\ วินยครุกา กติ วุตฺตา, \\ วินยสฺส เทฺว มูลานิ, \\ วินยครุกา เทฺว วุตฺตา, \\ คามนฺตเร กติ อาปตฺติโย, \\ กติมํเสสุ ถุลฺลจฺจยํ, \\ คามนฺตเร จตสฺโส อาปตฺติโย, \\ เอกมํเส ถุลฺลจฺจยํ, \\ กติ วาจสิกา รตฺติํ, \\ ททมานสฺส กติ อาปตฺติโย, \\ เทฺว วาจสิกา รตฺติํ, \\ ททมานสฺส ติสฺโส อาปตฺติโย,}
\csromangathagroup{\csromangathastanza{\csromanfolio{364}\csromangathaitembat{474}{กติ อาปตฺติโย กายิกา,}{กติ วาจสิกา กตา.} \\ \csromangathacontbat{ฉาเทนฺตสฺส กติ อาปตฺติโย,}{กติ สํสคฺคปจฺจยา.} \\ \csromangathacontbat{ฉาปตฺติโย กายิกา,}{ฉ วาจสิกา กตา.} \\ \csromangathacontbat{ฉาเทนฺตสฺส ติสฺโส อาปตฺติโย,}{ปญฺจ สํสคฺคปจฺจยา.} \\ \csromangathacontbat{อรุณุคฺเค กติ อาปตฺติโย,}{กติ ยาวตติยกา.} \\ \csromangathacontbat{กเตตฺถ อฏฺฐ วตฺถุกา,}{กติหิ สพฺพสงฺคโห.} \\ \csromangathacontbat{อรุณุคฺเค ติสฺโส อาปตฺติโย,}{เทฺว ยาวตติยกา.} \\ \csromangathacontbat{เอเกตฺถ อฏฺฐ วตฺถุกา,}{เอเกน สพฺพสงฺคโห.} \\ \csromangathacontbat{วินยสฺส กติ มูลานิ,}{ยานิ พุทฺเธน ปญฺญตฺตา.} \\ \csromangathacontbat{วินยครุกา กติ วุตฺตา,}{ทุฏฺฐุลฺลจฺฉาทนา กติ.} \\ \csromangathacontbat{วินยสฺส เทฺว มูลานิ,}{ยานิ พุทฺเธน ปญฺญตฺตา.} \\ \csromangathacontbat{วินยครุกา เทฺว วุตฺตา,}{เทฺว ทุฏฺฐุลฺลจฺฉาทนา.} \\ \csromangathacontbat{คามนฺตเร กติ อาปตฺติโย,}{กติ นทิปารปจฺจยา.} \\ \csromangathacontbat{กติมํเสสุ ถุลฺลจฺจยํ,}{กติมํเสสุ ทุกฺกฏํ.} \\ \csromangathacontbat{คามนฺตเร จตสฺโส อาปตฺติโย,}{จตสฺโส นทิปารปจฺจยา.} \\ \csromangathacontbat{เอกมํเส ถุลฺลจฺจยํ,}{นวมํเสสุ ทุกฺกฏํ.} \\ \csromangathacontbat{กติ วาจสิกา รตฺติํ,}{กติ วาจสิกา ทิวา.} \\ \csromangathacontbat{ททมานสฺส กติ อาปตฺติโย,}{ปฏิคฺคณฺหนฺตสฺส กิตฺตกา.} \\ \csromangathacontbat{เทฺว วาจสิกา รตฺติํ,}{เทฺว วาจสิกา ทิวา.} \\ \csromangathacontbat{ททมานสฺส ติสฺโส อาปตฺติโย,}{จตฺตาโร จ ปฏิคฺคเห.}}}

\csromanmatikaanchor{mtk.240}
\csromantocmarkref{paragraph}{2. เทสนาคามินิยาทิอาปตฺติ}{mtk.240}
\bookmark[dest={mtk.240},level=4]{2. เทสนาคามินิยาทิอาปตฺติ}
\csromanheader{4}{2. เทสนาคามินิยาทิอาปตฺติ}
\setlength{\csromangathapagewidth}{0pt}
\setlength{\csromangathawakwidth}{0pt}
\csromangathameasuregroup{กติ เทสนาคามินิโย, \\ กเตตฺถ อปฺปฏิกมฺมา วุตฺตา, \\ ปญฺจ เทสนาคามินิโย, \\ เอเกตฺถ อปฺปฏิกมฺมา วุตฺตา, \\ วินยครุกา กติ วุตฺตา, \\ กติ วิกาเล ธญฺญรโส, \\ วินยครุกา เทฺว วุตฺตา, \\ เอโก วิกาเล ธญฺญรโส, \\ ปาราชิกา กายิกา กติ, \\ กตินํ รตฺติจฺเฉโท, \\ ปาราชิกา กายิกา เทฺว, \\ ทฺวินฺนํ รตฺติจฺเฉโท, \\ กตตฺตานํ วธิตฺวาน, \\ กเตตฺถ ปฐมาปตฺติกา, \\ เทฺว อตฺตานํ วธิตฺวาน, \\ เทฺวตฺถ ปฏมาปตฺติกา, \\ ปาณาติปาเต กติ อาปตฺติโย, \\ โอภาสนา กติ วุตฺตา, \\ กติ ปุคฺคลา น อุปสมฺปาเทตพฺพา, \\ นาสิตกา กติ วุตฺตา, \\ ตโย ปุคฺคลา น อุปสมฺปาเทตพฺพา, \\ นาสิตกา ตโย วุตฺตา, \\ อทินฺนาทาเน กติ อาปตฺติโย, \\ ฉินฺทนฺตสฺส กติ อาปตฺติโย, \\ อทินฺนาทาเน ติสฺโส อาปตฺติโย, \\ ฉินฺทนฺตสฺส ติสฺโส อาปตฺติโย, \\ ภิกฺขุโนวาทกวคฺคสฺมิํ, \\ กเตตฺถ นวกา วุตฺตา, \\ ภิกฺขุโนวาทกวคฺคสฺมิํ, \\ จตุเรตฺถ นวกา วุตฺตา, \\ ภิกฺขุนีนญฺจ อกฺขาตา, \\ ภุญฺชนฺตามกธญฺเญน, \\ ภิกฺขุนีนญฺจ อกฺขาตา, \\ ภุญฺชนฺตามกธญฺเญน, \\ คจฺฉนฺตสฺส กติ อาปตฺติโย, \\ นิสินฺนสฺส กติ อาปตฺติโย, \\ คจฺฉนฺตสฺส จตสฺโส อาปตฺติโย, \\ นิสินฺนสฺส จตสฺโส อาปตฺติโย,}{\csromangathabat{กติ เทสนาคามินิโย,}{กติ สปฺปฏิกมฺมา กตา.} \\ \csromangathabat{กเตตฺถ อปฺปฏิกมฺมา วุตฺตา,}{พุทฺเธนาทิจฺจพนฺธุนา.} \\ \csromangathabat{ปญฺจ เทสนาคามินิโย,}{ฉ สปฺปฏิกมฺมา กตา.} \\ \csromangathabat{เอเกตฺถ อปฺปฏิกมฺมา วุตฺตา,}{พุทฺเธนาทิจฺจพนฺธุนา.} \\ \csromangathabat{วินยครุกา กติ วุตฺตา,}{กายวาจสิกานิ จ.} \\ \csromangathabat{กติ วิกาเล ธญฺญรโส,}{กติ ญตฺติจตุตฺเถน สมฺมุติ.} \\ \csromangathabat{วินยครุกา เทฺว วุตฺตา,}{กายวาจสิกานิ จ.} \\ \csromangathabat{เอโก วิกาเล ธญฺญรโส,}{เอกา ญตฺติจตุตฺเถน สมฺมุติ.} \\ \csromangathabat{ปาราชิกา กายิกา กติ,}{กติ สํวาสกภูมิโย.} \\ \csromangathabat{กตินํ รตฺติจฺเฉโท,}{ปญฺญตฺตา ทฺวงฺคุลา กติ.} \\ \csromangathabat{ปาราชิกา กายิกา เทฺว,}{เทฺว สํวาสกภูมิโย.} \\ \csromangathabat{ทฺวินฺนํ รตฺติจฺเฉโท,}{ปญฺญตฺตา ทฺวงฺคุลา ทุเว.} \\ \csromangathabat{กตตฺตานํ วธิตฺวาน,}{กติหิ สํโฆ ภิชฺชติ.} \\ \csromangathabat{กเตตฺถ ปฐมาปตฺติกา,}{ญตฺติยา กรณา กติ.} \\ \csromangathabat{เทฺว อตฺตานํ วธิตฺวาน,}{ทฺวีหิ สํโฆ ภิชฺชติ.} \\ \csromangathabat{เทฺวตฺถ ปฏมาปตฺติกา,}{ญตฺติยา กรณา ทุเว.} \\ \csromangathabat{ปาณาติปาเต กติ อาปตฺติโย,}{วาจา ปาราชิกา กติ.} \\ \csromangathabat{โอภาสนา กติ วุตฺตา,}{สญฺจริตฺเตน วา กติ.} \\ ปาณาติปาเต ติสฺโส อาปตฺติโย, \\ วาจา ปาราชิกา ตโย. \\ โอภาสนา ตโย วุตฺตา, \\ สญฺจริตฺเตน วา ตโย. \\ \csromangathabat{กติ ปุคฺคลา น อุปสมฺปาเทตพฺพา,}{กติ กมฺมานํ สงฺคหา.} \\ \csromangathabat{นาสิตกา กติ วุตฺตา,}{กตินํ เอกวาจิกา.} \\ \csromangathabat{ตโย ปุคฺคลา น อุปสมฺปาเทตพฺพา,}{ตโย กมฺมานํ สงฺคหา.} \\ \csromangathabat{นาสิตกา ตโย วุตฺตา,}{ติณฺณนฺนํ เอกวาจิกา.} \\ \csromangathabat{อทินฺนาทาเน กติ อาปตฺติโย,}{กติ เมถุนปจฺจยา.} \\ \csromangathabat{ฉินฺทนฺตสฺส กติ อาปตฺติโย,}{กติ ฉฑฺฑิตปจฺจยา.} \\ \csromangathabat{อทินฺนาทาเน ติสฺโส อาปตฺติโย,}{จตสฺโส เมถุนปจฺจยา.} \\ \csromangathabat{ฉินฺทนฺตสฺส ติสฺโส อาปตฺติโย,}{ปญฺจ ฉฑฺฑิตปจฺจยา.} \\ \csromangathabat{ภิกฺขุโนวาทกวคฺคสฺมิํ,}{ปาจิตฺติเยน ทุกฺกฏา.} \\ \csromangathabat{กเตตฺถ นวกา วุตฺตา,}{กตินํ จีวเรน จ.} \\ \csromangathabat{ภิกฺขุโนวาทกวคฺคสฺมิํ,}{ปาจิตฺติเยน ทุกฺกฏา กตา.} \\ \csromangathabat{จตุเรตฺถ นวกา วุตฺตา,}{ทฺวินฺนํ จีวเรน จ.} \\ \csromangathabat{ภิกฺขุนีนญฺจ อกฺขาตา,}{ปาฏิเทสนิยา กติ.} \\ \csromangathabat{ภุญฺชนฺตามกธญฺเญน,}{ปาจิตฺติเยน ทุกฺกฏา กติ.} \\ \csromangathabat{ภิกฺขุนีนญฺจ อกฺขาตา,}{อฏฺฐ ปาฏิเทสนียา กตา.} \\ \csromangathabat{ภุญฺชนฺตามกธญฺเญน,}{ปาจิตฺติเยน ทุกฺกฏา กตา.} \\ \csromangathabat{คจฺฉนฺตสฺส กติ อาปตฺติโย,}{ฐิตสฺส จาปิ กิตฺตกา.} \\ \csromangathabat{นิสินฺนสฺส กติ อาปตฺติโย,}{นิปนฺนสฺสาปิ กิตฺตกา.} \\ \csromangathabat{คจฺฉนฺตสฺส จตสฺโส อาปตฺติโย,}{ฐิตสฺส จาปิ ตตฺตกา.} \\ \csromangathabat{นิสินฺนสฺส จตสฺโส อาปตฺติโย,}{นิปนฺนสฺสาปิ ตตฺตกา.}}
\csromangathameasurewak{ปาณาติปาเต ติสฺโส อาปตฺติโย, \\ วาจา ปาราชิกา ตโย. \\ โอภาสนา ตโย วุตฺตา, \\ สญฺจริตฺเตน วา ตโย.}
\csromangathasetleft{กติ เทสนาคามินิโย, \\ กเตตฺถ อปฺปฏิกมฺมา วุตฺตา, \\ ปญฺจ เทสนาคามินิโย, \\ เอเกตฺถ อปฺปฏิกมฺมา วุตฺตา, \\ วินยครุกา กติ วุตฺตา, \\ กติ วิกาเล ธญฺญรโส, \\ วินยครุกา เทฺว วุตฺตา, \\ เอโก วิกาเล ธญฺญรโส, \\ ปาราชิกา กายิกา กติ, \\ กตินํ รตฺติจฺเฉโท, \\ ปาราชิกา กายิกา เทฺว, \\ ทฺวินฺนํ รตฺติจฺเฉโท, \\ กตตฺตานํ วธิตฺวาน, \\ กเตตฺถ ปฐมาปตฺติกา, \\ เทฺว อตฺตานํ วธิตฺวาน, \\ เทฺวตฺถ ปฏมาปตฺติกา, \\ ปาณาติปาเต กติ อาปตฺติโย, \\ โอภาสนา กติ วุตฺตา, \\ กติ ปุคฺคลา น อุปสมฺปาเทตพฺพา, \\ นาสิตกา กติ วุตฺตา, \\ ตโย ปุคฺคลา น อุปสมฺปาเทตพฺพา, \\ นาสิตกา ตโย วุตฺตา, \\ อทินฺนาทาเน กติ อาปตฺติโย, \\ ฉินฺทนฺตสฺส กติ อาปตฺติโย, \\ อทินฺนาทาเน ติสฺโส อาปตฺติโย, \\ ฉินฺทนฺตสฺส ติสฺโส อาปตฺติโย, \\ ภิกฺขุโนวาทกวคฺคสฺมิํ, \\ กเตตฺถ นวกา วุตฺตา, \\ ภิกฺขุโนวาทกวคฺคสฺมิํ, \\ จตุเรตฺถ นวกา วุตฺตา, \\ ภิกฺขุนีนญฺจ อกฺขาตา, \\ ภุญฺชนฺตามกธญฺเญน, \\ ภิกฺขุนีนญฺจ อกฺขาตา, \\ ภุญฺชนฺตามกธญฺเญน, \\ คจฺฉนฺตสฺส กติ อาปตฺติโย, \\ นิสินฺนสฺส กติ อาปตฺติโย, \\ คจฺฉนฺตสฺส จตสฺโส อาปตฺติโย, \\ นิสินฺนสฺส จตสฺโส อาปตฺติโย,}
\csromangathagroup{\csromangathastanza{\csromangathaitembat{475}{กติ เทสนาคามินิโย,}{กติ สปฺปฏิกมฺมา กตา.} \\ \csromangathacontbat{กเตตฺถ อปฺปฏิกมฺมา วุตฺตา,}{พุทฺเธนา\-ทิจฺจพนฺธุนา.}}\csromangathastanza{\csromanfolio{365}\csromangathabat{ปญฺจ เทสนาคามินิโย,}{ฉ สปฺปฏิกมฺมา กตา.} \\ \csromangathabat{เอเกตฺถ อปฺปฏิกมฺมา วุตฺตา,}{พุทฺเธนา\-ทิจฺจพนฺธุนา.}}\csromangathastanza{\csromangathabat{วินยครุกา กติ วุตฺตา,}{กายวาจสิกานิ จ.} \\ \csromangathabat{กติ วิกาเล ธญฺญรโส,}{กติ ญตฺติ\-จตุตฺเถน สมฺมุติ.}}\csromangathastanza{\csromangathabat{วินยครุกา เทฺว วุตฺตา,}{กายวาจสิกานิ จ.} \\ \csromangathabat{เอโก วิกาเล ธญฺญรโส,}{เอกา ญตฺติ\-จตุตฺเถน สมฺมุติ.}}\csromangathastanza{\csromangathabat{ปาราชิกา กายิกา กติ,}{กติ สํวาสกภูมิโย.} \\ \csromangathabat{กตินํ รตฺติจฺเฉโท,}{ปญฺญตฺตา ทฺวงฺคุลา กติ.}}\csromangathastanza{\csromangathabat{ปาราชิกา กายิกา เทฺว,}{เทฺว สํวาสกภูมิโย.} \\ \csromangathabat{ทฺวินฺนํ รตฺติจฺเฉโท,}{ปญฺญตฺตา ทฺวงฺคุลา ทุเว.}}\csromangathastanza{\csromangathabat{กตตฺตานํ วธิตฺวาน,}{กติหิ สํโฆ ภิชฺชติ.} \\ \csromangathabat{กเตตฺถ ปฐมาปตฺติกา,}{ญตฺติยา กรณา กติ.}}\csromangathastanza{\csromangathabat{เทฺว อตฺตานํ วธิตฺวาน,}{ทฺวีหิ สํโฆ ภิชฺชติ.} \\ \csromangathabat{เทฺวตฺถ ปฏมาปตฺติกา,}{ญตฺติยา กรณา ทุเว.}}\csromangathastanza{\csromangathabat{ปาณาติปาเต กติ อาปตฺติโย,}{วาจา ปาราชิกา กติ.} \\ \csromangathabat{โอภาสนา กติ วุตฺตา,}{สญฺจริตฺเตน วา กติ.}}\csromangathastanza{\csromangathawak{ปาณาติปาเต ติสฺโส อาปตฺติโย, \\ วาจา ปาราชิกา ตโย. \\ โอภาสนา ตโย วุตฺตา, \\ สญฺจริตฺเตน วา ตโย.}}\csromangathastanza{\csromangathabat{กติ ปุคฺคลา น อุปสมฺ\-ปาเท\-ตพฺพา,}{กติ กมฺมานํ สงฺคหา.} \\ \csromangathabat{นาสิตกา กติ วุตฺตา,}{กตินํ เอกวาจิกา.}}\csromangathastanza{\csromangathabat{ตโย ปุคฺคลา น อุปสมฺ\-ปาเท\-ตพฺพา,}{ตโย กมฺมานํ สงฺคหา.} \\ \csromangathabat{นาสิตกา ตโย วุตฺตา,}{ติณฺณนฺนํ เอกวาจิกา.}}\csromangathastanza{\csromangathabat{อทินฺนาทาเน กติ อาปตฺติโย,}{กติ เมถุนปจฺจยา.} \\ \csromangathabat{ฉินฺทนฺตสฺส กติ อาปตฺติโย,}{กติ ฉฑฺฑิตปจฺจยา.}}\csromangathastanza{\csromangathabat{อทินฺนาทาเน ติสฺโส อาปตฺติโย,}{จตสฺโส เมถุนปจฺจยา.} \\ \csromangathabat{ฉินฺทนฺตสฺส ติสฺโส อาปตฺติโย,}{ปญฺจ ฉฑฺฑิตปจฺจยา.}}\csromangathastanza{\csromanfolio{366}\csromangathabat{ภิกฺขุโนวาทกวคฺคสฺมิํ,}{ปาจิตฺติเยน ทุกฺกฏา.} \\ \csromangathabat{กเตตฺถ นวกา วุตฺตา,}{กตินํ จีวเรน จ.}}\csromangathastanza{\csromangathabat{ภิกฺขุโนวาทกวคฺคสฺมิํ,}{ปาจิตฺติเยน ทุกฺกฏา กตา.} \\ \csromangathabat{จตุเรตฺถ นวกา วุตฺตา,}{ทฺวินฺนํ จีวเรน จ.}}\csromangathastanza{\csromangathabat{ภิกฺขุนีนญฺจ อกฺขาตา,}{ปาฏิเทสนิยา กติ.} \\ \csromangathabat{ภุญฺชนฺตามกธญฺเญน,}{ปาจิตฺติเยน ทุกฺกฏา กติ.}}\csromangathastanza{\csromangathabat{ภิกฺขุนีนญฺจ อกฺขาตา,}{อฏฺฐ ปาฏิเทสนียา กตา.} \\ \csromangathabat{ภุญฺชนฺตามกธญฺเญน,}{ปาจิตฺติเยน ทุกฺกฏา กตา.}}\csromangathastanza{\csromangathabat{คจฺฉนฺตสฺส กติ อาปตฺติโย,}{ฐิตสฺส จาปิ กิตฺตกา.} \\ \csromangathabat{นิสินฺนสฺส กติ อาปตฺติโย,}{นิปนฺนสฺสาปิ กิตฺตกา.}}\csromangathastanza{\csromangathabat{คจฺฉนฺตสฺส จตสฺโส อาปตฺติโย,}{ฐิตสฺส จาปิ ตตฺตกา.} \\ \csromangathabat{นิสินฺนสฺส จตสฺโส อาปตฺติโย,}{นิปนฺนสฺสาปิ ตตฺตกา.}}}

\csromanmatikaanchor{mtk.241}
\csromantocmarkref{paragraph}{3. ปาจิตฺติย}{mtk.241}
\bookmark[dest={mtk.241},level=4]{3. ปาจิตฺติย}
\csromanheader{4}{3. \textbf{ปาจิตฺติย}}
\setlength{\csromangathapagewidth}{0pt}
\setlength{\csromangathawakwidth}{0pt}
\csromangathameasuregroup{กติ ปาจิตฺติยานิ, \\ อปุพฺพํ อจริมํ, \\ ปญฺจ ปาจิตฺติยานิ, \\ อปุพฺพํ อจริมํ, \\ กติ ปาจิตฺติยานิ, \\ อปุพฺพํ อจริมํ, \\ นว ปาจิตฺติยานิ, \\ อปุพฺพํ อจริมํ, \\ กติ ปาจิตฺติยานิ, \\ กติ วาจาย เทเสยฺย, \\ ปญฺจ ปาจิตฺติยานิ, \\ เอกวาจาย เทเสยฺย, \\ กติ ปาจิตฺติยานิ, \\ กติ วาจาย เทเสยฺย, \\ นว ปาจิตฺติยานิ, \\ เอกวาจาย เทเสยฺย, \\ กติ ปาจิตฺติยานิ, \\ กิญฺจ กิตฺเตตฺวา เทเสยฺย, \\ ปญฺจ ปาจิตฺติยานิ, \\ วตฺถุํ กิตฺเตตฺวา เทเสยฺย, \\ กติ ปาจิตฺติยานิ, \\ กิญฺจ กิตฺเตตฺวา เทเสยฺย, \\ นว ปาจิตฺติยานิ, \\ วตฺถุํ กิตฺเตตฺวา เทเสยฺย, \\ ยาวตติยเก กติ อาปตฺติโย, \\ ขาทนฺตสฺส กติ อาปตฺติโย, \\ ยาวตติยเก ติสฺโส อาปตฺติโย, \\ ขาทนฺตสฺส ติสฺโส อาปตฺติโย, \\ สพฺพา ยาวตติยกา, \\ กตินญฺเจว อาปตฺติ, \\ สพฺพา ยาวตติยกา, \\ ปญฺจนฺนญฺเจว อาปตฺติ, \\ กตินํ วินิจฺฉโย โหติ, \\ กตินญฺเจว อนาปตฺติ, \\ ปญฺจนฺนํ วินิจฺฉโย โหติ, \\ ปญฺจนฺนญฺเจว อนาปตฺติ, \\ กติ กายิกา รตฺติํ, \\ นิชฺฌายนฺตสฺส กติ อาปตฺติ, \\ เทฺว กายิกา รตฺติํ, \\ นิชฺฌายนฺตสฺส เอกา อาปตฺติ, \\ กตานิสํเส สมฺปสฺสํ, \\ อุกฺขิตฺตกา กติ วุตฺตา, \\ อฏฺฐานิสํเส สมฺปสฺสํ, \\ อุกฺขิตฺตกา ตโย วุตฺตา, \\ กติ ฐาเน มุสาวาโท, \\ กติ ปาฏิเทสนียา,}{\csromangathabat{กติ ปาจิตฺติยานิ,}{สพฺพานิ นานาวตฺถุกานิ.} \\ \csromangathabat{อปุพฺพํ อจริมํ,}{อาปชฺเชยฺย เอกโต.} \\ \csromangathabat{ปญฺจ ปาจิตฺติยานิ,}{สพฺพานิ นานาวตฺถุกานิ.} \\ \csromangathabat{อปุพฺพํ อจริมํ,}{อาปชฺเชยฺย เอกโต.} \\ \csromangathabat{กติ ปาจิตฺติยานิ,}{สพฺพานิ นานาวตฺถุกานิ.} \\ \csromangathabat{อปุพฺพํ อจริมํ,}{อาปชฺเชยฺย เอกโต.} \\ \csromangathabat{นว ปาจิตฺติยานิ,}{สพฺพานิ นานาวตฺถุกานิ.} \\ \csromangathabat{อปุพฺพํ อจริมํ,}{อาปชฺเชยฺย เอกโต.} \\ \csromangathabat{กติ ปาจิตฺติยานิ,}{สพฺพานิ นานาวตฺถุกานิ.} \\ \csromangathabat{กติ วาจาย เทเสยฺย,}{วุตฺตา อาทิจฺจพนฺธุนา.} \\ \csromangathabat{ปญฺจ ปาจิตฺติยานิ,}{สพฺพานิ นานาวตฺถุกานิ.} \\ \csromangathabat{เอกวาจาย เทเสยฺย,}{วุตฺตา อาทิจฺจพนฺธุนา.} \\ \csromangathabat{กติ ปาจิตฺติยานิ,}{สพฺพานิ นานาวตฺถุกานิ.} \\ \csromangathabat{กติ วาจาย เทเสยฺย,}{วุตฺตา อาทิจฺจพนฺธุนา.} \\ \csromangathabat{นว ปาจิตฺติยานิ,}{สพฺพานิ นานาวตฺถุกานิ.} \\ \csromangathabat{เอกวาจาย เทเสยฺย,}{วุตฺตา อาทิจฺจพนฺธุนา.} \\ \csromangathabat{กติ ปาจิตฺติยานิ,}{สพฺพานิ นานาวตฺถุกานิ.} \\ \csromangathabat{กิญฺจ กิตฺเตตฺวา เทเสยฺย,}{วุตฺตา อาทิจฺจพนฺธุนา.} \\ \csromangathabat{ปญฺจ ปาจิตฺติยานิ,}{สพฺพานิ นานาวตฺถุกานิ.} \\ \csromangathabat{วตฺถุํ กิตฺเตตฺวา เทเสยฺย,}{วุตฺตา อาทิจฺจพนฺธุนา.} \\ \csromangathabat{กติ ปาจิตฺติยานิ,}{สพฺพานิ นานาวตฺถุกานิ.} \\ \csromangathabat{กิญฺจ กิตฺเตตฺวา เทเสยฺย,}{วุตฺตา อาทิจฺจพนฺธุนา.} \\ \csromangathabat{นว ปาจิตฺติยานิ,}{สพฺพานิ นานาวตฺถุกานิ.} \\ \csromangathabat{วตฺถุํ กิตฺเตตฺวา เทเสยฺย,}{วุตฺตา อาทิจฺจพนฺธุนา.} \\ \csromangathabat{ยาวตติยเก กติ อาปตฺติโย,}{กติ โวหารปจฺจยา.} \\ \csromangathabat{ขาทนฺตสฺส กติ อาปตฺติโย,}{กติ โภชนปจฺจยา.} \\ \csromangathabat{ยาวตติยเก ติสฺโส อาปตฺติโย,}{ฉ โวหารปจฺจยา.} \\ \csromangathabat{ขาทนฺตสฺส ติสฺโส อาปตฺติโย,}{ปญฺจ โภชนปจฺจยา.} \\ \csromangathabat{สพฺพา ยาวตติยกา,}{กติ ฐานานิ คจฺฉนฺติ.} \\ \csromangathabat{กตินญฺเจว อาปตฺติ,}{กตินํ อธิกรเณน จ.} \\ \csromangathabat{สพฺพา ยาวตติยกา,}{ปญฺจ ฐานานิ คจฺฉนฺติ.} \\ \csromangathabat{ปญฺจนฺนญฺเจว อาปตฺติ,}{ปญฺจนฺนํ อธิกรเณน จ.} \\ \csromangathabat{กตินํ วินิจฺฉโย โหติ,}{กตินํ วูปสเมน จ.} \\ \csromangathabat{กตินญฺเจว อนาปตฺติ,}{กติหิ ฐาเนหิ โสภติ.} \\ \csromangathabat{ปญฺจนฺนํ วินิจฺฉโย โหติ,}{ปญฺจนฺนํ วูปสเมน จ.} \\ \csromangathabat{ปญฺจนฺนญฺเจว อนาปตฺติ,}{ตีหิ ฐาเนหิ โสภติ.} \\ \csromangathabat{กติ กายิกา รตฺติํ,}{กติ กายิกา ทิวา.} \\ \csromangathabat{นิชฺฌายนฺตสฺส กติ อาปตฺติ,}{กติ ปิณฺฑปาตปจฺจยา.} \\ \csromangathabat{เทฺว กายิกา รตฺติํ,}{เทฺว กายิกา ทิวา.} \\ \csromangathabat{นิชฺฌายนฺตสฺส เอกา อาปตฺติ,}{เอกา ปิณฺฑปาตปจฺจยา.} \\ \csromangathabat{กตานิสํเส สมฺปสฺสํ,}{ปเรสํ สทฺธาย เทสเย.} \\ \csromangathabat{อุกฺขิตฺตกา กติ วุตฺตา,}{กติ สมฺมาวตฺตนา.} \\ \csromangathabat{อฏฺฐานิสํเส สมฺปสฺสํ,}{ปเรสํ สทฺธาย เทสเย.} \\ \csromangathabat{อุกฺขิตฺตกา ตโย วุตฺตา,}{เตจตฺตาลีส สมฺมาวตฺตนา.} \\ \csromangathabat{กติ ฐาเน มุสาวาโท,}{กติ “ปรมนฺ”ติ วุจฺจติ.} \\ \csromangathabat{กติ ปาฏิเทสนียา,}{กตินํ เทสนาย จ.}}
\csromangathasetleft{กติ ปาจิตฺติยานิ, \\ อปุพฺพํ อจริมํ, \\ ปญฺจ ปาจิตฺติยานิ, \\ อปุพฺพํ อจริมํ, \\ กติ ปาจิตฺติยานิ, \\ อปุพฺพํ อจริมํ, \\ นว ปาจิตฺติยานิ, \\ อปุพฺพํ อจริมํ, \\ กติ ปาจิตฺติยานิ, \\ กติ วาจาย เทเสยฺย, \\ ปญฺจ ปาจิตฺติยานิ, \\ เอกวาจาย เทเสยฺย, \\ กติ ปาจิตฺติยานิ, \\ กติ วาจาย เทเสยฺย, \\ นว ปาจิตฺติยานิ, \\ เอกวาจาย เทเสยฺย, \\ กติ ปาจิตฺติยานิ, \\ กิญฺจ กิตฺเตตฺวา เทเสยฺย, \\ ปญฺจ ปาจิตฺติยานิ, \\ วตฺถุํ กิตฺเตตฺวา เทเสยฺย, \\ กติ ปาจิตฺติยานิ, \\ กิญฺจ กิตฺเตตฺวา เทเสยฺย, \\ นว ปาจิตฺติยานิ, \\ วตฺถุํ กิตฺเตตฺวา เทเสยฺย, \\ ยาวตติยเก กติ อาปตฺติโย, \\ ขาทนฺตสฺส กติ อาปตฺติโย, \\ ยาวตติยเก ติสฺโส อาปตฺติโย, \\ ขาทนฺตสฺส ติสฺโส อาปตฺติโย, \\ สพฺพา ยาวตติยกา, \\ กตินญฺเจว อาปตฺติ, \\ สพฺพา ยาวตติยกา, \\ ปญฺจนฺนญฺเจว อาปตฺติ, \\ กตินํ วินิจฺฉโย โหติ, \\ กตินญฺเจว อนาปตฺติ, \\ ปญฺจนฺนํ วินิจฺฉโย โหติ, \\ ปญฺจนฺนญฺเจว อนาปตฺติ, \\ กติ กายิกา รตฺติํ, \\ นิชฺฌายนฺตสฺส กติ อาปตฺติ, \\ เทฺว กายิกา รตฺติํ, \\ นิชฺฌายนฺตสฺส เอกา อาปตฺติ, \\ กตานิสํเส สมฺปสฺสํ, \\ อุกฺขิตฺตกา กติ วุตฺตา, \\ อฏฺฐานิสํเส สมฺปสฺสํ, \\ อุกฺขิตฺตกา ตโย วุตฺตา, \\ กติ ฐาเน มุสาวาโท, \\ กติ ปาฏิเทสนียา,}
\csromangathagroup{\csromangathastanza{\csromangathaitembat{476}{กติ ปาจิตฺติยานิ,}{สพฺพานิ นานาวตฺถุกานิ.} \\ \csromangathacontbat{อปุพฺพํ อจริมํ,}{อาปชฺเชยฺย เอกโต.} \\ \csromangathacontbat{ปญฺจ ปาจิตฺติยานิ,}{สพฺพานิ นานาวตฺถุกานิ.} \\ \csromangathacontbat{อปุพฺพํ อจริมํ,}{อาปชฺเชยฺย เอกโต.} \\ \csromangathacontbat{กติ ปาจิตฺติยานิ,}{สพฺพานิ นานาวตฺถุกานิ.} \\ \csromangathacontbat{อปุพฺพํ อจริมํ,}{อาปชฺเชยฺย เอกโต.} \\ \csromangathacontbat{นว ปาจิตฺติยานิ,}{สพฺพานิ นานาวตฺถุกานิ.} \\ \csromangathacontbat{อปุพฺพํ อจริมํ,}{อาปชฺเชยฺย เอกโต.} \\ \csromangathacontbat{กติ ปาจิตฺติยานิ,}{สพฺพานิ นานาวตฺถุกานิ.} \\ \csromangathacontbat{กติ วาจาย เทเสยฺย,}{วุตฺตา อาทิจฺจพนฺธุนา.} \\ \csromangathacontbat{ปญฺจ ปาจิตฺติยานิ,}{สพฺพานิ นานาวตฺถุกานิ.} \\ \csromangathacontbat{เอกวาจาย เทเสยฺย,}{วุตฺตา อาทิจฺจพนฺธุนา.} \\ \csromangathacontbat{กติ ปาจิตฺติยานิ,}{สพฺพานิ นานาวตฺถุกานิ.} \\ \csromangathacontbat{กติ วาจาย เทเสยฺย,}{วุตฺตา อาทิจฺจพนฺธุนา.}}\csromangathastanza{\csromanfolio{367}\csromangathabat{นว ปาจิตฺติยานิ,}{สพฺพานิ นานาวตฺถุกานิ.} \\ \csromangathabat{เอกวาจาย เทเสยฺย,}{วุตฺตา อาทิจฺจพนฺธุนา.}}\csromangathastanza{\csromangathabat{กติ ปาจิตฺติยานิ,}{สพฺพานิ นานาวตฺถุกานิ.} \\ \csromangathabat{กิญฺจ กิตฺเตตฺวา เทเสยฺย,}{วุตฺตา อาทิจฺจพนฺธุนา.}}\csromangathastanza{\csromangathabat{ปญฺจ ปาจิตฺติยานิ,}{สพฺพานิ นานาวตฺถุกานิ.} \\ \csromangathabat{วตฺถุํ กิตฺเตตฺวา เทเสยฺย,}{วุตฺตา อาทิจฺจพนฺธุนา.}}\csromangathastanza{\csromangathabat{กติ ปาจิตฺติยานิ,}{สพฺพานิ นานาวตฺถุกานิ.} \\ \csromangathabat{กิญฺจ กิตฺเตตฺวา เทเสยฺย,}{วุตฺตา อาทิจฺจพนฺธุนา.}}\csromangathastanza{\csromangathabat{นว ปาจิตฺติยานิ,}{สพฺพานิ นานาวตฺถุกานิ.} \\ \csromangathabat{วตฺถุํ กิตฺเตตฺวา เทเสยฺย,}{วุตฺตา อาทิจฺจพนฺธุนา.}}\csromangathastanza{\csromangathabat{ยาวตติยเก กติ อาปตฺติโย,}{กติ โวหารปจฺจยา.} \\ \csromangathabat{ขาทนฺตสฺส กติ อาปตฺติโย,}{กติ โภชนปจฺจยา.}}\csromangathastanza{\csromangathabat{ยาวตติยเก ติสฺโส อาปตฺติโย,}{ฉ โวหารปจฺจยา.} \\ \csromangathabat{ขาทนฺตสฺส ติสฺโส อาปตฺติโย,}{ปญฺจ โภชนปจฺจยา.}}\csromangathastanza{\csromangathabat{สพฺพา ยาวตติยกา,}{กติ ฐานานิ คจฺฉนฺติ.} \\ \csromangathabat{กตินญฺเจว อาปตฺติ,}{กตินํ อธิกรเณน จ.}}\csromangathastanza{\csromangathabat{สพฺพา ยาวตติยกา,}{ปญฺจ ฐานานิ คจฺฉนฺติ.} \\ \csromangathabat{ปญฺจนฺนญฺเจว อาปตฺติ,}{ปญฺจนฺนํ อธิกรเณน จ.}}\csromangathastanza{\csromangathabat{กตินํ วินิจฺฉโย โหติ,}{กตินํ วูปสเมน จ.} \\ \csromangathabat{กตินญฺเจว อนาปตฺติ,}{กติหิ ฐาเนหิ โสภติ.}}\csromangathastanza{\csromangathabat{ปญฺจนฺนํ วินิจฺฉโย โหติ,}{ปญฺจนฺนํ วูปสเมน จ.} \\ \csromangathabat{ปญฺจนฺนญฺเจว อนาปตฺติ,}{ตีหิ ฐาเนหิ โสภติ.}}\csromangathastanza{\csromangathabat{กติ กายิกา รตฺติํ,}{กติ กายิกา ทิวา.} \\ \csromangathabat{นิชฺฌายนฺตสฺส กติ อาปตฺติ,}{กติ ปิณฺฑปาตปจฺจยา.}}\csromangathastanza{\csromangathabat{เทฺว กายิกา รตฺติํ,}{เทฺว กายิกา ทิวา.} \\ \csromangathabat{นิชฺฌายนฺตสฺส เอกา อาปตฺติ,}{เอกา ปิณฺฑปาตปจฺจยา.}}\csromangathastanza{\csromanfolio{368}\csromangathabat{กตานิสํเส สมฺปสฺสํ,}{ปเรสํ สทฺธาย เทสเย.} \\ \csromangathabat{อุกฺขิตฺตกา กติ วุตฺตา,}{กติ สมฺมาวตฺตนา.}}\csromangathastanza{\csromangathabat{อฏฺฐานิสํเส สมฺปสฺสํ,}{ปเรสํ สทฺธาย เทสเย.} \\ \csromangathabat{อุกฺขิตฺตกา ตโย วุตฺตา,}{เตจตฺตาลีส สมฺมาวตฺตนา.}}\csromangathastanza{\csromangathaitembat{476}{กติ ฐาเน มุสาวาโท,}{กติ “ปรมนฺ”ติ วุจฺจติ.} \\ \csromangathacontbat{กติ ปาฏิเทสนียา,}{กตินํ เทสนาย จ.}}}

\prose{ปญฺจ ฐาเน มุสาวาโท, จุทฺทส “ปรมนฺ”ติ วุจฺจติ.}

\setlength{\csromangathapagewidth}{0pt}
\setlength{\csromangathawakwidth}{0pt}
\csromangathameasuregroup{ทฺวาทส ปาฏิเทสนียา, \\ กตงฺคิโก มุสาวาโท, \\ กติ ทูเตยฺยงฺคานิ, \\ อฏฺฐงฺคิโก มุสาวาโท, \\ อฏฺฐ ทูเตยฺยงฺคานิ, \\ กติวาจิกา อุปสมฺปทา, \\ กตินํ อาสนํ ทาตพฺพํ, \\ อฏฺฐวาจิกา อุปสมฺปทา, \\ อฏฺฐนฺนํ อาสนํ ทาตพฺพํ, \\ กตินํ เฉชฺชํ โหติ, \\ กตินญฺเจว อนาปตฺติ, \\ เอกสฺส เฉชฺชํ โหติ, \\ จตุนฺนญฺเจว อนาปตฺติ, \\ กติ อาฆาตวตฺถูนิ, \\ กเตตฺถ ปฐมาปตฺติกา, \\ นว อาฆาตวตฺถูนิ, \\ นเวตฺถ ปฐมาปตฺติกา,}{\csromangathabat{ทฺวาทส ปาฏิเทสนียา,}{จตุนฺนํ เทสนาย จ.} \\ \csromangathabat{กตงฺคิโก มุสาวาโท,}{กติ อุโปสถงฺคานิ.} \\ \csromangathabat{กติ ทูเตยฺยงฺคานิ,}{กติ ติตฺถิยวตฺตนา.} \\ \csromangathabat{อฏฺฐงฺคิโก มุสาวาโท,}{อฏฺฐ อุโปสถงฺคานิ.} \\ \csromangathabat{อฏฺฐ ทูเตยฺยงฺคานิ,}{อฏฺฐ ติตฺถิยวตฺตนา.} \\ \csromangathabat{กติวาจิกา อุปสมฺปทา,}{กตินํ ปจฺจุฏฺฐาตพฺพํ.} \\ \csromangathabat{กตินํ อาสนํ ทาตพฺพํ,}{ภิกฺขุโนวาทโก กติหิ.} \\ \csromangathabat{อฏฺฐวาจิกา อุปสมฺปทา,}{อฏฺฐนฺนํ ปจฺจุฏฺฐาตพฺพํ.} \\ \csromangathabat{อฏฺฐนฺนํ อาสนํ ทาตพฺพํ,}{ภิกฺขุโนวาทโก อฏฺฐหิ.} \\ \csromangathabat{กตินํ เฉชฺชํ โหติ,}{กตินํ ถุลฺลจฺจยํ.} \\ \csromangathabat{กตินญฺเจว อนาปตฺติ,}{สพฺเพสํ เอกวตฺถุกา.} \\ \csromangathabat{เอกสฺส เฉชฺชํ โหติ,}{จตุนฺนํ ถุลฺลจฺจยํ.} \\ \csromangathabat{จตุนฺนญฺเจว อนาปตฺติ,}{สพฺเพสํ เอกวตฺถุกา.} \\ \csromangathabat{กติ อาฆาตวตฺถูนิ,}{กติหิ สํโฆ ภิชฺชติ.} \\ \csromangathabat{กเตตฺถ ปฐมาปตฺติกา,}{ญตฺติยา กรณา กติ.} \\ \csromangathabat{นว อาฆาตวตฺถูนิ,}{นวหิ สํโฆ ภิชฺชติ.} \\ \csromangathabat{นเวตฺถ ปฐมาปตฺติกา,}{ญตฺติยา กรณา นว.}}
\csromangathasetleft{ทฺวาทส ปาฏิเทสนียา, \\ กตงฺคิโก มุสาวาโท, \\ กติ ทูเตยฺยงฺคานิ, \\ อฏฺฐงฺคิโก มุสาวาโท, \\ อฏฺฐ ทูเตยฺยงฺคานิ, \\ กติวาจิกา อุปสมฺปทา, \\ กตินํ อาสนํ ทาตพฺพํ, \\ อฏฺฐวาจิกา อุปสมฺปทา, \\ อฏฺฐนฺนํ อาสนํ ทาตพฺพํ, \\ กตินํ เฉชฺชํ โหติ, \\ กตินญฺเจว อนาปตฺติ, \\ เอกสฺส เฉชฺชํ โหติ, \\ จตุนฺนญฺเจว อนาปตฺติ, \\ กติ อาฆาตวตฺถูนิ, \\ กเตตฺถ ปฐมาปตฺติกา, \\ นว อาฆาตวตฺถูนิ, \\ นเวตฺถ ปฐมาปตฺติกา,}
\csromangathagroup{\csromangathastanza{\csromangathabat{ทฺวาทส ปาฏิเทสนียา,}{จตุนฺนํ เทสนาย จ.} \\ \csromangathabat{กตงฺคิโก มุสาวาโท,}{กติ อุโปสถงฺคานิ.}}\csromangathastanza{\csromangathabat{กติ ทูเตยฺยงฺคานิ,}{กติ ติตฺถิยวตฺตนา.} \\ \csromangathabat{อฏฺฐงฺคิโก มุสาวาโท,}{อฏฺฐ อุโปสถงฺคานิ.}}\csromangathastanza{\csromangathabat{อฏฺฐ ทูเตยฺยงฺคานิ,}{อฏฺฐ ติตฺถิยวตฺตนา.} \\ \csromangathabat{กติวาจิกา อุปสมฺปทา,}{กตินํ ปจฺจุฏฺฐาตพฺพํ.}}\csromangathastanza{\csromangathabat{กตินํ อาสนํ ทาตพฺพํ,}{ภิกฺขุโน\-วาทโก กติหิ.} \\ \csromangathabat{อฏฺฐวาจิกา อุปสมฺปทา,}{อฏฺฐนฺนํ ปจฺจุฏฺฐาตพฺพํ.}}\csromangathastanza{\csromangathabat{อฏฺฐนฺนํ อาสนํ ทาตพฺพํ,}{ภิกฺขุโน\-วาทโก อฏฺฐหิ.} \\ \csromangathabat{กตินํ เฉชฺชํ โหติ,}{กตินํ ถุลฺลจฺจยํ.}}\csromangathastanza{\csromangathabat{กตินญฺเจว อนาปตฺติ,}{สพฺเพสํ เอกวตฺถุกา.} \\ \csromangathabat{เอกสฺส เฉชฺชํ โหติ,}{จตุนฺนํ ถุลฺลจฺจยํ.}}\csromangathastanza{\csromangathabat{จตุนฺนญฺเจว อนาปตฺติ,}{สพฺเพสํ เอกวตฺถุกา.} \\ \csromangathabat{กติ อาฆาตวตฺถูนิ,}{กติหิ สํโฆ ภิชฺชติ.}}\csromangathastanza{\csromangathabat{กเตตฺถ ปฐมาปตฺติกา,}{ญตฺติยา กรณา กติ.} \\ \csromangathabat{นว อาฆาตวตฺถูนิ,}{นวหิ สํโฆ ภิชฺชติ.} \\ \csromangathabat{นเวตฺถ ปฐมาปตฺติกา,}{ญตฺติยา กรณา นว.}}}

\csromanmatikaanchor{mtk.242}
\csromantocmarkref{paragraph}{4. อวนฺทนียปุคฺคลาทิ}{mtk.242}
\bookmark[dest={mtk.242},level=4]{4. อวนฺทนียปุคฺคลาทิ}
\csromanheader{4}{4. อวนฺท\-นีย\-ปุคฺคลาทิ}
\setlength{\csromangathapagewidth}{0pt}
\setlength{\csromangathawakwidth}{0pt}
\csromangathameasuregroup{กติ ปุคฺคลา นาภิวาเทตพฺพา, \\ กตินํ ทุกฺกฏํ โหติ, \\ ทส ปุคฺคลา นาภิวาเทตพฺพา, \\ ทสนฺนํ ทุกฺกฏํ โหติ, \\ กตินํ วสฺสํวุฏฺฐานํ, \\ กตินํ สนฺเต ทาตพฺพํ, \\ ปญฺจนฺนํ วสฺสํวุฏฺฐานํ, \\ สตฺตนฺนํ สนฺเต ทาตพฺพํ, \\ กติสตํ รตฺติสตํ, \\ กติ รตฺติโย วสิตฺวาน, \\ ทสสตํ รตฺติสตํ, \\ ทส รตฺติโย วสิตฺวาน, \\ กติ กมฺมโทสา วุตฺตา, \\ จมฺปายํ วินยวตฺถุสฺมิํ, \\ ทฺวาทส กมฺมโทสา วุตฺตา, \\ จมฺปายํ วินยวตฺถุสฺมิํ, \\ กติ กมฺมสมฺปตฺติโย วุตฺตา, \\ จมฺปายํ วินยวตฺถุสฺมิํ, \\ จตสฺโส กมฺมสมฺปตฺติโย วุตฺตา, \\ จมฺปายํ วินยวตฺถุสฺมิํ, \\ กติ กมฺมานิ วุตฺตานิ, \\ จมฺปายํ วินยวตฺถุสฺมิํ, \\ ฉ กมฺมานิ วุตฺตานิ, \\ จมฺปายํ วินยวตฺถุสฺมิํ, \\ ปญฺจ อธมฺมิกา วุตฺตา, \\ กติ กมฺมานิ วุตฺตานิ, \\ จมฺปายํ วินยวตฺถุสฺมิํ,}{\csromangathabat{กติ ปุคฺคลา นาภิวาเทตพฺพา,}{อญฺชลิสามิเจน จ.} \\ \csromangathabat{กตินํ ทุกฺกฏํ โหติ,}{กติ จีวรธารณา.} \\ \csromangathabat{ทส ปุคฺคลา นาภิวาเทตพฺพา,}{อญฺชลิสามิเจน จ.} \\ \csromangathabat{ทสนฺนํ ทุกฺกฏํ โหติ,}{ทส จีวรธารณา.} \\ \csromangathabat{กตินํ วสฺสํวุฏฺฐานํ,}{ทาตพฺพํ อิธ จีวรํ.} \\ \csromangathabat{กตินํ สนฺเต ทาตพฺพํ,}{กตินญฺเจว น ทาตพฺพํ.} \\ \csromangathabat{ปญฺจนฺนํ วสฺสํวุฏฺฐานํ,}{ทาตพฺพํ อิธ จีวรํ.} \\ \csromangathabat{สตฺตนฺนํ สนฺเต ทาตพฺพํ,}{โสฬสนฺนํ น ทาตพฺพํ.} \\ \csromangathabat{กติสตํ รตฺติสตํ,}{อาปตฺติโย ฉาทยิตฺวาน.} \\ \csromangathabat{กติ รตฺติโย วสิตฺวาน,}{มุจฺเจยฺย ปาริวาสิโก.} \\ \csromangathabat{ทสสตํ รตฺติสตํ,}{อาปตฺติโย ฉาทยิตฺวาน.} \\ \csromangathabat{ทส รตฺติโย วสิตฺวาน,}{มุจฺเจยฺย ปาริวาสิโก.} \\ \csromangathabat{กติ กมฺมโทสา วุตฺตา,}{พุทฺเธนาทิจฺจพนฺธุนา.} \\ \csromangathabat{จมฺปายํ วินยวตฺถุสฺมิํ,}{สพฺเพว อธมฺมิกา1 กติ.} \\ \csromangathabat{ทฺวาทส กมฺมโทสา วุตฺตา,}{พุทฺเธนาทิจฺจพนฺธุนา.} \\ \csromangathabat{จมฺปายํ วินยวตฺถุสฺมิํ,}{สพฺเพว อธมฺมิกา2 กตา.} \\ \csromangathabat{กติ กมฺมสมฺปตฺติโย วุตฺตา,}{พุทฺเธนาทิจฺจพนฺธุนา.} \\ \csromangathabat{จมฺปายํ วินยวตฺถุสฺมิํ,}{สพฺเพว ธมฺมิกา กติ.} \\ \csromangathabat{จตสฺโส กมฺมสมฺปตฺติโย วุตฺตา,}{พุทฺเธนาทิจฺจพนฺธุนา.} \\ \csromangathabat{จมฺปายํ วินยวตฺถุสฺมิํ,}{สพฺเพว ธมฺมิกา กตา.} \\ \csromangathabat{กติ กมฺมานิ วุตฺตานิ,}{พุทฺเธนาทิจฺจพนฺธุนา.} \\ \csromangathabat{จมฺปายํ วินยวตฺถุสฺมิํ,}{สพฺเพว อธมฺมิกา กติ.} \\ \csromangathabat{ฉ กมฺมานิ วุตฺตานิ,}{พุทฺเธนาทิจฺจพนฺธุนา.} \\ \csromangathabat{จมฺปายํ วินยวตฺถุสฺมิํ,}{เอเกตฺถ ธมฺมิกา กตา.} \\ \csromangathabat{ปญฺจ อธมฺมิกา วุตฺตา,}{พุทฺเธนาทิจฺจพนฺธุนา.} \\ \csromangathabat{กติ กมฺมานิ วุตฺตานิ,}{พุทฺเธนาทิจฺจพนฺธุนา.} \\ \csromangathabat{จมฺปายํ วินยวตฺถุสฺมิํ,}{ธมฺมิกา อธมฺมิกา กติ.}}
\csromangathasetleft{กติ ปุคฺคลา นาภิวาเทตพฺพา, \\ กตินํ ทุกฺกฏํ โหติ, \\ ทส ปุคฺคลา นาภิวาเทตพฺพา, \\ ทสนฺนํ ทุกฺกฏํ โหติ, \\ กตินํ วสฺสํวุฏฺฐานํ, \\ กตินํ สนฺเต ทาตพฺพํ, \\ ปญฺจนฺนํ วสฺสํวุฏฺฐานํ, \\ สตฺตนฺนํ สนฺเต ทาตพฺพํ, \\ กติสตํ รตฺติสตํ, \\ กติ รตฺติโย วสิตฺวาน, \\ ทสสตํ รตฺติสตํ, \\ ทส รตฺติโย วสิตฺวาน, \\ กติ กมฺมโทสา วุตฺตา, \\ จมฺปายํ วินยวตฺถุสฺมิํ, \\ ทฺวาทส กมฺมโทสา วุตฺตา, \\ จมฺปายํ วินยวตฺถุสฺมิํ, \\ กติ กมฺมสมฺปตฺติโย วุตฺตา, \\ จมฺปายํ วินยวตฺถุสฺมิํ, \\ จตสฺโส กมฺมสมฺปตฺติโย วุตฺตา, \\ จมฺปายํ วินยวตฺถุสฺมิํ, \\ กติ กมฺมานิ วุตฺตานิ, \\ จมฺปายํ วินยวตฺถุสฺมิํ, \\ ฉ กมฺมานิ วุตฺตานิ, \\ จมฺปายํ วินยวตฺถุสฺมิํ, \\ ปญฺจ อธมฺมิกา วุตฺตา, \\ กติ กมฺมานิ วุตฺตานิ, \\ จมฺปายํ วินยวตฺถุสฺมิํ,}
\csromangathagroup{\csromangathastanza{\csromangathaitembat{477}{กติ ปุคฺคลา นาภิวาเทตพฺพา,}{อญฺชลิสามิเจน จ.} \\ \csromangathacontbat{กตินํ ทุกฺกฏํ โหติ,}{กติ จีวรธารณา.}}\csromangathastanza{\csromanfolio{369}\csromangathabat{ทส ปุคฺคลา นาภิวาเทตพฺพา,}{อญฺชลิสามิเจน จ.} \\ \csromangathabat{ทสนฺนํ ทุกฺกฏํ โหติ,}{ทส จีวรธารณา.}}\csromangathastanza{\csromangathabat{กตินํ วสฺสํ\-วุฏฺฐานํ,}{ทาตพฺพํ อิธ จีวรํ.} \\ \csromangathabat{กตินํ สนฺเต ทาตพฺพํ,}{กตินญฺเจว น ทาตพฺพํ.}}\csromangathastanza{\csromangathabat{ปญฺจนฺนํ วสฺสํ\-วุฏฺฐานํ,}{ทาตพฺพํ อิธ จีวรํ.} \\ \csromangathabat{สตฺตนฺนํ สนฺเต ทาตพฺพํ,}{โสฬสนฺนํ น ทาตพฺพํ.}}\csromangathastanza{\csromangathabat{กติสตํ รตฺติสตํ,}{อาปตฺติโย ฉาทยิตฺวาน.} \\ \csromangathabat{กติ รตฺติโย วสิตฺวาน,}{มุจฺเจยฺย ปาริวาสิโก.}}\csromangathastanza{\csromangathabat{ทสสตํ รตฺติสตํ,}{อาปตฺติโย ฉาทยิตฺวาน.} \\ \csromangathabat{ทส รตฺติโย วสิตฺวาน,}{มุจฺเจยฺย ปาริวาสิโก.}}\csromangathastanza{\csromangathabat{กติ กมฺมโทสา วุตฺตา,}{พุทฺเธนา\-ทิจฺจพนฺธุนา.} \\ \csromangathabat{จมฺปายํ วินยวตฺถุสฺมิํ,}{สพฺเพว อธมฺมิกา1 กติ.}}\csromangathastanza{\csromangathabat{ทฺวาทส กมฺมโทสา วุตฺตา,}{พุทฺเธนา\-ทิจฺจพนฺธุนา.} \\ \csromangathabat{จมฺปายํ วินยวตฺถุสฺมิํ,}{สพฺเพว อธมฺมิกา2 กตา.}}\csromangathastanza{\csromangathabat{กติ กมฺม\-สมฺปตฺติโย วุตฺตา,}{พุทฺเธนา\-ทิจฺจพนฺธุนา.} \\ \csromangathabat{จมฺปายํ วินยวตฺถุสฺมิํ,}{สพฺเพว ธมฺมิกา กติ.}}\csromangathastanza{\csromangathabat{จตสฺโส กมฺม\-สมฺปตฺติโย วุตฺตา,}{พุทฺเธนา\-ทิจฺจพนฺธุนา.} \\ \csromangathabat{จมฺปายํ วินยวตฺถุสฺมิํ,}{สพฺเพว ธมฺมิกา กตา.}}\csromangathastanza{\csromangathabat{กติ กมฺมานิ วุตฺตานิ,}{พุทฺเธนา\-ทิจฺจพนฺธุนา.} \\ \csromangathabat{จมฺปายํ วินยวตฺถุสฺมิํ,}{สพฺเพว อธมฺมิกา กติ.}}\csromangathastanza{\csromangathabat{ฉ กมฺมานิ วุตฺตานิ,}{พุทฺเธนา\-ทิจฺจพนฺธุนา.} \\ \csromangathabat{จมฺปายํ วินยวตฺถุสฺมิํ,}{เอเกตฺถ ธมฺมิกา กตา.}}\csromangathastanza{\csromangathabat{ปญฺจ อธมฺมิกา วุตฺตา,}{พุทฺเธนา\-ทิจฺจพนฺธุนา.} \\ \csromangathabat{กติ กมฺมานิ วุตฺตานิ,}{พุทฺเธนา\-ทิจฺจพนฺธุนา.} \\ \csromangathabat{จมฺปายํ วินยวตฺถุสฺมิํ,}{ธมฺมิกา อธมฺมิกา กติ.}}}

\proseitem{1}{สพฺเพ อธมฺมิกา (สี, สฺยา) 2. สพฺเพ\-วา\-ธมฺมิกา (สี), สพฺเพ อธมฺมิกา (สฺยา)}

\setlength{\csromangathapagewidth}{0pt}
\setlength{\csromangathawakwidth}{0pt}
\csromangathameasuregroup{จตฺตาริ กมฺมานิ วุตฺตานิ, \\ จมฺปายํ วินยวตฺถุสฺมิํ, \\ ตโย อธมฺมิกา วุตฺตา, \\ ยํ เทสิตํนนฺตชิเนน ตาทินา, \\ เอเกตฺถ สมฺมติ วินา สมเถหิ, \\ กติ อาปายิกา วุตฺตา, \\ วินยํ ปฏิชานนฺตสฺส, \\ ฉอูนทิยฑฺฒสตา วุตฺตา, \\ อาปายิกา เนรยิกา, \\ กติ อฏฺฐกา วุตฺตา, \\ วินยํ ปฏิชานนฺตสฺส, \\ อฏฺฐารส อฏฺฐกา วุตฺตา, \\ วินยํ ปฏิชานนฺตสฺส,}{\csromangathabat{จตฺตาริ กมฺมานิ วุตฺตานิ,}{พุทฺเธนาทิจฺจพนฺธุนา.} \\ \csromangathabat{จมฺปายํ วินยวตฺถุสฺมิํ,}{เอเกตฺถ ธมฺมิกา กตา.} \\ \csromangathabat{ตโย อธมฺมิกา วุตฺตา,}{พุทฺเธนาทิจฺจพนฺธุนา.} \\ \csromangathabat{ยํ เทสิตํนนฺตชิเนน ตาทินา,}{อาปตฺติกฺขนฺธานิ วิเวกทสฺสินา.} \\ กเตตฺถ สมฺมนฺติ วินา สมเถหิ, \\ ปุจฺฉามิ ตํ พฺรูหิ วิภงฺคโกวิท. \\ ยํ เทสิตํนนฺตชิเนน ตาทินา, \\ อาปตฺติกฺขนฺธานิ วิเวกทสฺสินา. \\ \csromangathabat{เอเกตฺถ สมฺมติ วินา สมเถหิ,}{เอตํ เต อกฺขามิ วิภงฺคโกวิท.} \\ \csromangathabat{กติ อาปายิกา วุตฺตา,}{พุทฺเธนาทิจฺจพนฺธุนา.} \\ \csromangathabat{วินยํ ปฏิชานนฺตสฺส,}{วินยานิ\textsuperscript{1} สุโณม เต.} \\ \csromangathabat{ฉอูนทิยฑฺฒสตา วุตฺตา,}{พุทฺเธนาทิจฺจพนฺธุนา.} \\ \csromangathabat{อาปายิกา เนรยิกา,}{กปฺปฏฺฐา สํฆเภทกา.} \\ วินยํ ปฏิชานนฺตสฺส,วินยานิ สุโณหิ เม. \\ กติ นาปายิกา วุตฺตา, \\ พุทฺเธนาทิจฺจพนฺธุนา. \\ วินยํ ปฏิชานนฺตสฺส, \\ วินยานิ สุโณม เต. \\ อฏฺฐารส นาปายิกา วุตฺตา, \\ พุทฺเธนาทิจฺจพนฺธุนา. \\ วินยํ ปฏิชานนฺตสฺส,วินยานิ สุโณหิ เม. \\ \csromangathabat{กติ อฏฺฐกา วุตฺตา,}{พุทฺเธนาทิจฺจพนฺธุนา.} \\ \csromangathabat{วินยํ ปฏิชานนฺตสฺส,}{วินยานิ สุโณม เต.} \\ \csromangathabat{อฏฺฐารส อฏฺฐกา วุตฺตา,}{พุทฺเธนาทิจฺจพนฺธุนา.} \\ \csromangathabat{วินยํ ปฏิชานนฺตสฺส,}{วินยานิ สุโณหิ เม.}}
\csromangathameasurewak{กเตตฺถ สมฺมนฺติ วินา สมเถหิ, \\ ปุจฺฉามิ ตํ พฺรูหิ วิภงฺคโกวิท. \\ ยํ เทสิตํนนฺตชิเนน ตาทินา, \\ อาปตฺติกฺขนฺธานิ วิเวกทสฺสินา. \\ พุทฺเธนาทิจฺจพนฺธุนา. \\ วินยํ ปฏิชานนฺตสฺส, \\ วินยานิ สุโณม เต. \\ อฏฺฐารส นาปายิกา วุตฺตา,}
\csromangathasetleft{จตฺตาริ กมฺมานิ วุตฺตานิ, \\ จมฺปายํ วินยวตฺถุสฺมิํ, \\ ตโย อธมฺมิกา วุตฺตา, \\ ยํ เทสิตํนนฺตชิเนน ตาทินา, \\ เอเกตฺถ สมฺมติ วินา สมเถหิ, \\ กติ อาปายิกา วุตฺตา, \\ วินยํ ปฏิชานนฺตสฺส, \\ ฉอูนทิยฑฺฒสตา วุตฺตา, \\ อาปายิกา เนรยิกา, \\ กติ อฏฺฐกา วุตฺตา, \\ วินยํ ปฏิชานนฺตสฺส, \\ อฏฺฐารส อฏฺฐกา วุตฺตา, \\ วินยํ ปฏิชานนฺตสฺส,}
\csromangathagroup{\csromangathastanza{\csromanfolio{370}\csromangathabat{จตฺตาริ กมฺมานิ วุตฺตานิ,}{พุทฺเธนา\-ทิจฺจพนฺธุนา.} \\ \csromangathabat{จมฺปายํ วินยวตฺถุสฺมิํ,}{เอเกตฺถ ธมฺมิกา กตา.}}\csromangathastanza{\csromangathabat{ตโย อธมฺมิกา วุตฺตา,}{พุทฺเธนา\-ทิจฺจพนฺธุนา.} \\ \csromangathabat{ยํ เทสิตํนนฺตชิเนน ตาทินา,}{อาปตฺติกฺขนฺธานิ วิเวกทสฺสินา.}}\csromangathastanza{\csromangathawak{กเตตฺถ สมฺมนฺติ วินา สมเถหิ, \\ ปุจฺฉามิ ตํ พฺรูหิ วิภงฺคโกวิท. \\ ยํ เทสิตํนนฺตชิเนน ตาทินา, \\ อาปตฺติกฺขนฺธานิ วิเวกทสฺสินา.}}\csromangathastanza{\csromangathabat{เอเกตฺถ สมฺมติ วินา สมเถหิ,}{เอตํ เต อกฺขามิ วิภงฺคโกวิท.} \\ \csromangathabat{กติ อาปายิกา วุตฺตา,}{พุทฺเธนา\-ทิจฺจพนฺธุนา.}}\csromangathastanza{\csromangathabat{วินยํ ปฏิชานนฺตสฺส,}{วินยานิ\csromansharedfootnote{cde764664dea9333}{วินยานิ (สี, สฺยา, เอวมุปริปิ)} สุโณม เต.} \\ \csromangathabat{ฉอูนทิยฑฺฒสตา วุตฺตา,}{พุทฺเธนา\-ทิจฺจพนฺธุนา.}}\csromangathastanza{\csromangathabat{อาปายิกา เนรยิกา,}{กปฺปฏฺฐา สํฆเภทกา.} \\ วินยํ ปฏิชานนฺตสฺส,วินยานิ สุโณหิ เม. \\ กติ นาปายิกา วุตฺตา,}\csromangathastanza{\csromangathawak{พุทฺเธนา\-ทิจฺจพนฺธุนา. \\ วินยํ ปฏิชานนฺตสฺส, \\ วินยานิ สุโณม เต. \\ อฏฺฐารส นาปายิกา วุตฺตา,}}\csromangathastanza{พุทฺเธนา\-ทิจฺจพนฺธุนา. \\ วินยํ ปฏิชานนฺตสฺส,วินยานิ สุโณหิ เม. \\ \csromangathabat{กติ อฏฺฐกา วุตฺตา,}{พุทฺเธนา\-ทิจฺจพนฺธุนา.}}\csromangathastanza{\csromangathabat{วินยํ ปฏิชานนฺตสฺส,}{วินยานิ สุโณม เต.} \\ \csromangathabat{อฏฺฐารส อฏฺฐกา วุตฺตา,}{พุทฺเธนา\-ทิจฺจพนฺธุนา.} \\ \csromangathabat{วินยํ ปฏิชานนฺตสฺส,}{วินยานิ สุโณหิ เม.}}}

\csromanmatikaanchor{mtk.243}
\csromantocmarkref{paragraph}{5. โสฬสกมฺมาทิ}{mtk.243}
\bookmark[dest={mtk.243},level=4]{5. โสฬสกมฺมาทิ}
\csromanheader{4}{5. \textbf{โสฬสกมฺมาทิ}}
\setlength{\csromangathapagewidth}{0pt}
\setlength{\csromangathawakwidth}{0pt}
\csromangathameasuregroup{โสฬส กมฺมานิ วุตฺตานิ,}{กติ กมฺมานิ วุตฺตานิ, พุทฺเธนาทิจฺจพนฺธุนา. \\ วินยํ ปฏิชานนฺตสฺส,วินยานิ สุโณม เต. \\ \csromangathabat{โสฬส กมฺมานิ วุตฺตานิ,}{พุทฺเธนาทิจฺจพนฺธุนา.}}
\csromangathameasurewak{กติ กมฺมานิ วุตฺตานิ, พุทฺเธนาทิจฺจพนฺธุนา. \\ วินยํ ปฏิชานนฺตสฺส,วินยานิ สุโณม เต.}
\csromangathasetleft{โสฬส กมฺมานิ วุตฺตานิ,}
\csromangathagroup{\csromangathastanza{\csromangathawak{\csromangathaitemline{478}{กติ กมฺมานิ วุตฺตานิ, พุทฺเธนา\-ทิจฺจพนฺธุนา.} \\ \csromangathacontline{วินยํ ปฏิชานนฺตสฺส,วินยานิ สุโณม เต.}}}\csromangathastanza{\csromanfolio{371}\csromangathabat{โสฬส กมฺมานิ วุตฺตานิ,}{พุทฺเธนา\-ทิจฺจพนฺธุนา.}}}

\prose{วินยํ ปฏิชานนฺตสฺส,วินยานิ สุโณหิ เม.}

\setlength{\csromangathapagewidth}{0pt}
\setlength{\csromangathawakwidth}{0pt}
\csromangathameasuregroup{กติ กมฺมโทสา วุตฺตา,}{\csromangathabat{กติ กมฺมโทสา วุตฺตา,}{พุทฺเธนาทิจฺจพนฺธุนา.}}
\csromangathasetleft{กติ กมฺมโทสา วุตฺตา,}
\csromangathagroup{\csromangathastanza{\csromangathabat{กติ กมฺมโทสา วุตฺตา,}{พุทฺเธนา\-ทิจฺจพนฺธุนา.}}}

\prose{วินยํ ปฏิชานนฺตสฺส,วินยานิ สุโณม เต.}

\setlength{\csromangathapagewidth}{0pt}
\setlength{\csromangathawakwidth}{0pt}
\csromangathameasuregroup{ทฺวาทส กมฺมโทสา วุตฺตา, \\ วินยํ ปฏิชานนฺตสฺส, \\ กติ กมฺมสมฺปตฺติโย วุตฺตา,}{\csromangathabat{ทฺวาทส กมฺมโทสา วุตฺตา,}{พุทฺเธนาทิจฺจพนฺธุนา.} \\ \csromangathabat{วินยํ ปฏิชานนฺตสฺส,}{วินยานิ สุโณหิ เม.} \\ \csromangathabat{กติ กมฺมสมฺปตฺติโย วุตฺตา,}{พุทฺเธนาทิจฺจพนฺธุนา.}}
\csromangathasetleft{ทฺวาทส กมฺมโทสา วุตฺตา, \\ วินยํ ปฏิชานนฺตสฺส, \\ กติ กมฺมสมฺปตฺติโย วุตฺตา,}
\csromangathagroup{\csromangathastanza{\csromangathabat{ทฺวาทส กมฺมโทสา วุตฺตา,}{พุทฺเธนา\-ทิจฺจพนฺธุนา.} \\ \csromangathabat{วินยํ ปฏิชานนฺตสฺส,}{วินยานิ สุโณหิ เม.} \\ \csromangathabat{กติ กมฺม\-สมฺปตฺติโย วุตฺตา,}{พุทฺเธนา\-ทิจฺจพนฺธุนา.}}}

\prose{วินยํ ปฏิชานนฺตสฺส,วินยานิ สุโณม เต.}

\setlength{\csromangathapagewidth}{0pt}
\setlength{\csromangathawakwidth}{0pt}
\csromangathameasuregroup{จตสฺโส กมฺมสมฺปตฺติโย วุตฺตา, \\ วินยํ ปฏิชานนฺตสฺส, \\ กติ กมฺมานิ วุตฺตานิ, \\ วินยํปฏิชานนฺตสฺส, \\ ฉ กมฺมานิ วุตฺตานิ, \\ วินยํ ปฏิชานนฺตสฺส, \\ กติ กมฺมานิ วุตฺตานิ, \\ วินยํ ปฏิชานนฺตสฺส, \\ จตฺตาริ กมฺมานิ วุตฺตานิ,}{\csromangathabat{จตสฺโส กมฺมสมฺปตฺติโย วุตฺตา,}{พุทฺเธนาทิจฺจพนฺธุนา.} \\ \csromangathabat{วินยํ ปฏิชานนฺตสฺส,}{วินยานิ สุโณหิ เม.} \\ \csromangathabat{กติ กมฺมานิ วุตฺตานิ,}{พุทฺเธนาทิจฺจพนฺธุนา.} \\ \csromangathabat{วินยํปฏิชานนฺตสฺส,}{วินยานิ สุโณม เต.} \\ \csromangathabat{ฉ กมฺมานิ วุตฺตานิ,}{พุทฺเธนาทิจฺจพนฺธุนา.} \\ \csromangathabat{วินยํ ปฏิชานนฺตสฺส,}{วินยานิ สุโณหิ เม.} \\ \csromangathabat{กติ กมฺมานิ วุตฺตานิ,}{พุทฺเธนาทิจฺจพนฺธุนา.} \\ \csromangathabat{วินยํ ปฏิชานนฺตสฺส,}{วินยานิ สุโณม เต.} \\ \csromangathabat{จตฺตาริ กมฺมานิ วุตฺตานิ,}{พุทฺเธนาทิจฺจพนฺธุนา.}}
\csromangathasetleft{จตสฺโส กมฺมสมฺปตฺติโย วุตฺตา, \\ วินยํ ปฏิชานนฺตสฺส, \\ กติ กมฺมานิ วุตฺตานิ, \\ วินยํปฏิชานนฺตสฺส, \\ ฉ กมฺมานิ วุตฺตานิ, \\ วินยํ ปฏิชานนฺตสฺส, \\ กติ กมฺมานิ วุตฺตานิ, \\ วินยํ ปฏิชานนฺตสฺส, \\ จตฺตาริ กมฺมานิ วุตฺตานิ,}
\csromangathagroup{\csromangathastanza{\csromangathabat{จตสฺโส กมฺม\-สมฺปตฺติโย วุตฺตา,}{พุทฺเธนา\-ทิจฺจพนฺธุนา.} \\ \csromangathabat{วินยํ ปฏิชานนฺตสฺส,}{วินยานิ สุโณหิ เม.}}\csromangathastanza{\csromangathabat{กติ กมฺมานิ วุตฺตานิ,}{พุทฺเธนา\-ทิจฺจพนฺธุนา.} \\ \csromangathabat{วินยํปฏิชานนฺตสฺส,}{วินยานิ สุโณม เต.}}\csromangathastanza{\csromangathabat{ฉ กมฺมานิ วุตฺตานิ,}{พุทฺเธนา\-ทิจฺจพนฺธุนา.} \\ \csromangathabat{วินยํ ปฏิชานนฺตสฺส,}{วินยานิ สุโณหิ เม.}}\csromangathastanza{\csromangathabat{กติ กมฺมานิ วุตฺตานิ,}{พุทฺเธนา\-ทิจฺจพนฺธุนา.} \\ \csromangathabat{วินยํ ปฏิชานนฺตสฺส,}{วินยานิ สุโณม เต.} \\ \csromangathabat{จตฺตาริ กมฺมานิ วุตฺตานิ,}{พุทฺเธนา\-ทิจฺจพนฺธุนา.}}}

\prose{วินยํ ปฏิชานนฺตสฺส,วินยานิ สุโณหิ เม.}

\setlength{\csromangathapagewidth}{0pt}
\setlength{\csromangathawakwidth}{0pt}
\csromangathameasuregroup{กติ ปาราชิกา วุตฺตา,}{\csromangathabat{กติ ปาราชิกา วุตฺตา,}{พุทฺเธนาทิจฺจพนฺธุนา.}}
\csromangathasetleft{กติ ปาราชิกา วุตฺตา,}
\csromangathagroup{\csromangathastanza{\csromangathabat{กติ ปาราชิกา วุตฺตา,}{พุทฺเธนา\-ทิจฺจพนฺธุนา.}}}

\prose{วินยํ ปฏิชานนฺตสฺส,วินยานิ สุโณม เต.}

\setlength{\csromangathapagewidth}{0pt}
\setlength{\csromangathawakwidth}{0pt}
\csromangathameasuregroup{อฏฺฐ ปาราชิกา วุตฺตา, \\ วินยํ ปฏิชานนฺตสฺส, \\ กติ สํฆาทิเสสา วุตฺตา,}{\csromangathabat{อฏฺฐ ปาราชิกา วุตฺตา,}{พุทฺเธนาทิจฺจพนฺธุนา.} \\ \csromangathabat{วินยํ ปฏิชานนฺตสฺส,}{วินยานิ สุโณหิ เม.} \\ \csromangathabat{กติ สํฆาทิเสสา วุตฺตา,}{พุทฺเธนาทิจฺจพนฺธุนา.}}
\csromangathasetleft{อฏฺฐ ปาราชิกา วุตฺตา, \\ วินยํ ปฏิชานนฺตสฺส, \\ กติ สํฆาทิเสสา วุตฺตา,}
\csromangathagroup{\csromangathastanza{\csromangathabat{อฏฺฐ ปาราชิกา วุตฺตา,}{พุทฺเธนา\-ทิจฺจพนฺธุนา.} \\ \csromangathabat{วินยํ ปฏิชานนฺตสฺส,}{วินยานิ สุโณหิ เม.} \\ \csromangathabat{กติ สํฆาทิเสสา วุตฺตา,}{พุทฺเธนา\-ทิจฺจพนฺธุนา.}}}

\prose{วินยํ ปฏิชานนฺตสฺส,วินยานิ สุโณม เต.}

\setlength{\csromangathapagewidth}{0pt}
\setlength{\csromangathawakwidth}{0pt}
\csromangathameasuregroup{เตวีส สํฆาทิเสสา วุตฺตา,}{\csromangathabat{เตวีส สํฆาทิเสสา วุตฺตา,}{พุทฺเธนาทิจฺจพนฺธุนา.}}
\csromangathasetleft{เตวีส สํฆาทิเสสา วุตฺตา,}
\csromangathagroup{\csromangathastanza{\csromangathabat{เตวีส สํฆาทิเสสา วุตฺตา,}{พุทฺเธนา\-ทิจฺจพนฺธุนา.}}}

\prose{วินยํ ปฏิชานนฺตสฺส,วินยานิ สุโณหิ เม.}

\setlength{\csromangathapagewidth}{0pt}
\setlength{\csromangathawakwidth}{0pt}
\csromangathameasuregroup{กติ อนิยตา วุตฺตา, \\ วินยํ ปฏิชานนฺตสฺส, \\ เทฺว อนิยตา วุตฺตา, \\ วินยํ ปฏิชานนฺตสฺส, \\ กติ นิสฺสคฺคิยา วุตฺตา, \\ วินยํ ปฏิชานนฺตสฺส, \\ เทฺวจตฺตาลีส นิสฺสคฺคิยา วุตฺตา, \\ วินยํ ปฏิชานนฺตสฺส, \\ กติ ปาจิตฺติยา วุตฺตา, \\ วินยํ ปฏิชานนฺตสฺส, \\ อฏฺฐาสีติสตํ ปาจิตฺติยา วุตฺตา, \\ วินยํ ปฏิชานนฺตสฺส, \\ กติ ปาฏิเทสนียา วุตฺตา, \\ วินยํ ปฏิชานนฺตสฺส, \\ ทฺวาทส ปาฏิเทสนียา วุตฺตา, \\ วินยํปฏิชานนฺตสฺส, \\ กติ เสขิยา วุตฺตา, \\ วินยํ ปฏิชานนฺตสฺส, \\ ปญฺจสตฺตติ เสขิยา วุตฺตา, \\ วินยํ ปฏิชานนฺตสฺส, \\ ยาว สุปุจฺฉิตํ ตยา, \\ ปุจฺฉาวิสฺสชฺชนาย วา,}{\csromangathabat{กติ อนิยตา วุตฺตา,}{พุทฺเธนาทิจฺจพนฺธุนา.} \\ \csromangathabat{วินยํ ปฏิชานนฺตสฺส,}{วินยานิ สุโณม เต.} \\ \csromangathabat{เทฺว อนิยตา วุตฺตา,}{พุทฺเธนาทิจฺจพนฺธุนา.} \\ \csromangathabat{วินยํ ปฏิชานนฺตสฺส,}{วินยานิ สุโณหิ เม.} \\ \csromangathabat{กติ นิสฺสคฺคิยา วุตฺตา,}{พุทฺเธนาทิจฺจพนฺธุนา.} \\ \csromangathabat{วินยํ ปฏิชานนฺตสฺส,}{วินยานิ สุโณม เต.} \\ \csromangathabat{เทฺวจตฺตาลีส นิสฺสคฺคิยา วุตฺตา,}{พุทฺเธนาทิจฺจพนฺธุนา.} \\ \csromangathabat{วินยํ ปฏิชานนฺตสฺส,}{วินยานิ สุโณหิ เม.} \\ \csromangathabat{กติ ปาจิตฺติยา วุตฺตา,}{พุทฺเธนาทิจฺจพนฺธุนา.} \\ \csromangathabat{วินยํ ปฏิชานนฺตสฺส,}{วินยานิ สุโณม เต.} \\ \csromangathabat{อฏฺฐาสีติสตํ ปาจิตฺติยา วุตฺตา,}{พุทฺเธนาทิจฺจพนฺธุนา.} \\ \csromangathabat{วินยํ ปฏิชานนฺตสฺส,}{วินยานิ สุโณหิ เม.} \\ \csromangathabat{กติ ปาฏิเทสนียา วุตฺตา,}{พุทฺเธนาทิจฺจพนฺธุนา.} \\ \csromangathabat{วินยํ ปฏิชานนฺตสฺส,}{วินยานิ สุโณม เต.} \\ \csromangathabat{ทฺวาทส ปาฏิเทสนียา วุตฺตา,}{พุทฺเธนาทิจฺจพนฺธุนา.} \\ \csromangathabat{วินยํปฏิชานนฺตสฺส,}{วินยานิ สุโณหิ เม.} \\ \csromangathabat{กติ เสขิยา วุตฺตา,}{พุทฺเธนาทิจฺจพนฺธุนา.} \\ \csromangathabat{วินยํ ปฏิชานนฺตสฺส,}{วินยานิ สุโณม เต.} \\ \csromangathabat{ปญฺจสตฺตติ เสขิยา วุตฺตา,}{พุทฺเธนาทิจฺจพนฺธุนา.} \\ \csromangathabat{วินยํ ปฏิชานนฺตสฺส,}{วินยานิ สุโณหิ เม.} \\ \csromangathabat{ยาว สุปุจฺฉิตํ ตยา,}{ยาว สุวิสฺสชฺชิตํ มยา.} \\ \csromangathabat{ปุจฺฉาวิสฺสชฺชนาย วา,}{นตฺถิ กิญฺจิ อสุตฺตกนฺติ.}}
\csromangathasetleft{กติ อนิยตา วุตฺตา, \\ วินยํ ปฏิชานนฺตสฺส, \\ เทฺว อนิยตา วุตฺตา, \\ วินยํ ปฏิชานนฺตสฺส, \\ กติ นิสฺสคฺคิยา วุตฺตา, \\ วินยํ ปฏิชานนฺตสฺส, \\ เทฺวจตฺตาลีส นิสฺสคฺคิยา วุตฺตา, \\ วินยํ ปฏิชานนฺตสฺส, \\ กติ ปาจิตฺติยา วุตฺตา, \\ วินยํ ปฏิชานนฺตสฺส, \\ อฏฺฐาสีติสตํ ปาจิตฺติยา วุตฺตา, \\ วินยํ ปฏิชานนฺตสฺส, \\ กติ ปาฏิเทสนียา วุตฺตา, \\ วินยํ ปฏิชานนฺตสฺส, \\ ทฺวาทส ปาฏิเทสนียา วุตฺตา, \\ วินยํปฏิชานนฺตสฺส, \\ กติ เสขิยา วุตฺตา, \\ วินยํ ปฏิชานนฺตสฺส, \\ ปญฺจสตฺตติ เสขิยา วุตฺตา, \\ วินยํ ปฏิชานนฺตสฺส, \\ ยาว สุปุจฺฉิตํ ตยา, \\ ปุจฺฉาวิสฺสชฺชนาย วา,}
\csromangathagroupwithcloser{\csromangathastanza{\csromanfolio{372}\csromangathabat{กติ อนิยตา วุตฺตา,}{พุทฺเธนา\-ทิจฺจพนฺธุนา.} \\ \csromangathabat{วินยํ ปฏิชานนฺตสฺส,}{วินยานิ สุโณม เต.}}\csromangathastanza{\csromangathabat{เทฺว อนิยตา วุตฺตา,}{พุทฺเธนา\-ทิจฺจพนฺธุนา.} \\ \csromangathabat{วินยํ ปฏิชานนฺตสฺส,}{วินยานิ สุโณหิ เม.}}\csromangathastanza{\csromangathabat{กติ นิสฺสคฺคิยา วุตฺตา,}{พุทฺเธนา\-ทิจฺจพนฺธุนา.} \\ \csromangathabat{วินยํ ปฏิชานนฺตสฺส,}{วินยานิ สุโณม เต.}}\csromangathastanza{\csromangathabat{เทฺวจตฺตาลีส นิสฺสคฺคิยา วุตฺตา,}{พุทฺเธนา\-ทิจฺจพนฺธุนา.} \\ \csromangathabat{วินยํ ปฏิชานนฺตสฺส,}{วินยานิ สุโณหิ เม.}}\csromangathastanza{\csromangathabat{กติ ปาจิตฺติยา วุตฺตา,}{พุทฺเธนา\-ทิจฺจพนฺธุนา.} \\ \csromangathabat{วินยํ ปฏิชานนฺตสฺส,}{วินยานิ สุโณม เต.}}\csromangathastanza{\csromangathabat{อฏฺฐาสีติสตํ ปาจิตฺติยา วุตฺตา,}{พุทฺเธนา\-ทิจฺจพนฺธุนา.} \\ \csromangathabat{วินยํ ปฏิชานนฺตสฺส,}{วินยานิ สุโณหิ เม.}}\csromangathastanza{\csromangathabat{กติ ปาฏิเทสนียา วุตฺตา,}{พุทฺเธนา\-ทิจฺจพนฺธุนา.} \\ \csromangathabat{วินยํ ปฏิชานนฺตสฺส,}{วินยานิ สุโณม เต.}}\csromangathastanza{\csromangathabat{ทฺวาทส ปาฏิเทสนียา วุตฺตา,}{พุทฺเธนา\-ทิจฺจพนฺธุนา.} \\ \csromangathabat{วินยํปฏิชานนฺตสฺส,}{วินยานิ สุโณหิ เม.}}\csromangathastanza{\csromangathabat{กติ เสขิยา วุตฺตา,}{พุทฺเธนา\-ทิจฺจพนฺธุนา.} \\ \csromangathabat{วินยํ ปฏิชานนฺตสฺส,}{วินยานิ สุโณม เต.}}\csromangathastanza{\csromangathabat{ปญฺจสตฺตติ เสขิยา วุตฺตา,}{พุทฺเธนา\-ทิจฺจพนฺธุนา.} \\ \csromangathabat{วินยํ ปฏิชานนฺตสฺส,}{วินยานิ สุโณหิ เม.}}\csromangathastanza{\csromangathabat{ยาว สุปุจฺฉิตํ ตยา,}{ยาว สุวิสฺสชฺชิตํ มยา.} \\ \csromangathabat{ปุจฺฉาวิสฺสชฺชนาย วา,}{นตฺถิ กิญฺจิ อสุตฺตกนฺติ.}}}{ทุติย\-คาถา\-สงฺคณิกํ นิฏฺฐิตํ.}
\csromanmatikaanchor{mtk.244}
\csromantocmarknopage{section}{เสทโมจนคาถา}{mtk.244}
\bookmark[dest={mtk.244},level=1]{เสทโมจนคาถา}
\csromanheader{1}{เสท\-โมจน\-คาถา}
\csromanmatikaanchor{mtk.245}
\csromantocmarkref{subparagraph}{1. อวิปฺปวาสปญฺหา}{mtk.245}
\bookmark[dest={mtk.245},level=5]{1. อวิปฺปวาสปญฺหา}
\csromanheader{5}{1. อวิปฺปวาส\-ปญฺหา}
\setlength{\csromangathapagewidth}{0pt}
\setlength{\csromangathawakwidth}{0pt}
\csromangathameasuregroup{อวิสฺสชฺชิยํ อเวภงฺคิยํ,}{อสํวาโส ภิกฺขูหิ จ ภิกฺขุนีหิ จ, \\ สมฺโภโค เอกจฺโจ ตหิํ น ลพฺภติ. \\ อวิปฺปวาเสน อนาปตฺติ, \\ ปญฺหา เมสา กุสเลหิ จินฺติตา. \\ \csromangathabat{อวิสฺสชฺชิยํ อเวภงฺคิยํ,}{ปญฺจ วุตฺตา มเหสินา.} \\ วิสฺสชฺชนฺตสฺส ปริภุญฺชนฺตสฺส อนาปตฺติ, \\ ปญฺหา เมสา กุสเลหิ จินฺติตา. \\ ทส ปุคฺคเล น วทามิ, เอกาทส วิวชฺชิย. \\ วุฑฺฒํ วนฺทนฺตสฺส อาปตฺติ, ปญฺหา เมสา กุสเลหิ จินฺติตา. \\ น อุกฺขิตฺตโก, น จ ปน ปาริวาสิโก, \\ น สํฆภินฺโน, น จ ปน ปกฺขสงฺกนฺโต. \\ สมานสํวาสกภูมิยา ฐิโต, \\ กถํ นุ สิกฺขาย อสาธารโณ สิยา. \\ ปญฺหา เมสา กุสเลหิ จินฺติตา. \\ อุเปติ ธมฺมํ ปริปุจฺฉมาโน, กุสลํ อตฺถูปสญฺหิตํ. \\ น ชีวติ, น มโต, น นิพฺพุโต, \\ ตํ ปุคฺคลํ กตมํ วทนฺติ พุทฺธา. \\ ปญฺหา เมสา กุสเลหิ จินฺติตา. \\ อุพฺภกฺขเก น วทามิ, อโธ นาภิํ วิวชฺชิย. \\ เมถุนธมฺมปจฺจยา, กถํ ปาราชิโก สิยา. \\ ปญฺหา เมสา กุสเลหิ จินฺติตา. \\ ภิกฺขุ สญฺญาจิกาย กุฏิํ กโรติ, \\ อเทสิตวตฺถุกํ ปมาณาติกฺกนฺตํ. \\ สารมฺภํ อปริกฺกมนํ อนาปตฺติ, \\ ปญฺหา เมสา กุสเลหิ จินฺติตา. \\ ภิกฺขุ สญฺญาจิกาย กุฏิํ กโรติ, \\ เทสิตวตฺถุกํ ปมาณิกํ. \\ อนารมฺภํ สปริกฺกมนํ อาปตฺติ, \\ ปญฺหา เมสา กุสเลหิ จินฺติ ตา. \\ น กายิกํ กิญฺจิ ปโยคมาจเร, \\ น จาปิ วาจาย ปเร ภเณยฺย. \\ อาปชฺเชยฺย ครุกํ เฉชฺชวตฺถุํ, \\ ปญฺหา เมสา กุสเลหิ จินฺติตา. \\ น กายิกํ วาจสิกญฺจ กิญฺจิ, \\ มนสาปิ สนฺโต น กเรยฺย ปาปํ. \\ โส นาสิโต กินฺติ สุนาสิโต ภเว, \\ ปญฺหา เมสา กุสเลหิ จินฺติตา. \\ อนาลปนฺโต มนุเชน เกนจิ, \\ วาจาคิรํ โน จ ปเร ภเณยฺย. \\ อาปชฺเชยฺย วาจสิกํ น กายิกํ, \\ ปญฺหา เมสา กุสเลหิ จินฺติตา. \\ สิกฺขาปทา พุทฺธวเรน วณฺณิตา, \\ สํฆาทิเสสา จตุโร ภเวยฺยุํ. \\ อาปชฺเชยฺย เอกปโยเคน สพฺเพ, \\ ปญฺหา เมสา กุสเลหิ จินฺติตา. \\ อุโภ เอกโต อุปสมฺปนฺนา, \\ อุภินฺนํ หตฺถโต จีวรํ ปฏิคฺคเณฺหยฺย. \\ สิยา อาปตฺติโย นานา, \\ ปญฺหา เมสา กุสลิหิ จินฺติตา. \\ จตุโร ชานา สํวิธาย, \\ ครุภณฺฑํ อวาหรุํ. \\ ตโย ปาราชิกา, เอโก น ปาราชิโก, \\ ปญฺหา เมสา กุสเลหิ จินฺติตา.}
\csromangathameasurewak{ภิกฺขุ สญฺญาจิกาย กุฏิํ กโรติ, \\ เทสิตวตฺถุกํ ปมาณิกํ. \\ อนารมฺภํ สปริกฺกมนํ อาปตฺติ, \\ ปญฺหา เมสา กุสเลหิ จินฺติ ตา. \\ น กายิกํ กิญฺจิ ปโยคมาจเร, \\ น จาปิ วาจาย ปเร ภเณยฺย. \\ อาปชฺเชยฺย ครุกํ เฉชฺชวตฺถุํ, \\ ปญฺหา เมสา กุสเลหิ จินฺติตา. \\ น กายิกํ วาจสิกญฺจ กิญฺจิ, \\ มนสาปิ สนฺโต น กเรยฺย ปาปํ. \\ โส นาสิโต กินฺติ สุนาสิโต ภเว, \\ ปญฺหา เมสา กุสเลหิ จินฺติตา. \\ อนาลปนฺโต มนุเชน เกนจิ, \\ วาจาคิรํ โน จ ปเร ภเณยฺย. \\ อาปชฺเชยฺย วาจสิกํ น กายิกํ, \\ ปญฺหา เมสา กุสเลหิ จินฺติตา. \\ สิกฺขาปทา พุทฺธวเรน วณฺณิตา, \\ สํฆาทิเสสา จตุโร ภเวยฺยุํ. \\ อาปชฺเชยฺย เอกปโยเคน สพฺเพ, \\ ปญฺหา เมสา กุสเลหิ จินฺติตา. \\ อุโภ เอกโต อุปสมฺปนฺนา, \\ อุภินฺนํ หตฺถโต จีวรํ ปฏิคฺคเณฺหยฺย. \\ สิยา อาปตฺติโย นานา, \\ ปญฺหา เมสา กุสลิหิ จินฺติตา. \\ จตุโร ชานา สํวิธาย, \\ ครุภณฺฑํ อวาหรุํ. \\ ตโย ปาราชิกา, เอโก น ปาราชิโก, \\ ปญฺหา เมสา กุสเลหิ จินฺติตา.}
\csromangathasetleft{อวิสฺสชฺชิยํ อเวภงฺคิยํ,}
\csromangathagroup{\csromangathastanza{\csromanfolio{373}\csromangathaitemline{479}{อสํวาโส ภิกฺขูหิ จ ภิกฺขุนีหิ จ,} \\ \csromangathacontline{สมฺโภโค เอกจฺโจ ตหิํ น ลพฺภติ.} \\ \csromangathacontline{อวิปฺปวาเสน อนาปตฺติ,} \\ \csromangathacontline{ปญฺหา เมสา กุสเลหิ จินฺติตา.} \\ \csromangathacontbat{อวิสฺสชฺชิยํ อเวภงฺคิยํ,}{ปญฺจ วุตฺตา มเหสินา.} \\ \csromangathacontline{วิสฺสชฺชนฺตสฺส ปริภุญฺช\-นฺตสฺส อนาปตฺติ,} \\ \csromangathacontline{ปญฺหา เมสา กุสเลหิ จินฺติตา.} \\ \csromangathacontline{ทส ปุคฺคเล น วทามิ, เอกาทส วิวชฺชิย.} \\ \csromangathacontline{วุฑฺฒํ วนฺทนฺตสฺส อาปตฺติ, ปญฺหา เมสา กุสเลหิ จินฺติตา.} \\ \csromangathacontline{น อุกฺขิตฺตโก, น จ ปน ปาริวาสิโก,} \\ \csromangathacontline{น สํฆภินฺโน, น จ ปน ปกฺขสงฺกนฺโต.} \\ \csromangathacontline{สมานสํวาสกภูมิยา ฐิโต,} \\ \csromangathacontline{กถํ นุ สิกฺขาย อสาธารโณ สิยา.} \\ \csromangathacontline{ปญฺหา เมสา กุสเลหิ จินฺติตา.} \\ \csromangathacontline{อุเปติ ธมฺมํ ปริปุจฺฉมาโน, กุสลํ อตฺถูปสญฺหิตํ.} \\ \csromangathacontline{น ชีวติ, น มโต, น นิพฺพุโต,} \\ \csromangathacontline{ตํ ปุคฺคลํ กตมํ วทนฺติ พุทฺธา.} \\ \csromangathacontline{ปญฺหา เมสา กุสเลหิ จินฺติตา.} \\ \csromangathacontline{อุพฺภกฺขเก น วทามิ, อโธ นาภิํ วิวชฺชิย.} \\ \csromangathacontline{เมถุนธมฺมปจฺจยา, กถํ ปาราชิโก สิยา.} \\ \csromangathacontline{ปญฺหา เมสา กุสเลหิ จินฺติตา.} \\ \csromangathacontline{ภิกฺขุ สญฺญาจิกาย กุฏิํ กโรติ,} \\ \csromangathacontline{อเทสิต\-วตฺถุกํ ปมาณา\-ติกฺกนฺตํ.} \\ \csromangathacontline{สารมฺภํ อปริกฺกมนํ อนาปตฺติ,} \\ \csromangathacontline{ปญฺหา เมสา กุสเลหิ จินฺติตา.}}\csromangathastanza{\csromangathawak{\csromanfolio{374}ภิกฺขุ สญฺญาจิกาย กุฏิํ กโรติ, \\ เทสิต\-วตฺถุกํ ปมาณิกํ. \\ อนารมฺภํ สปริกฺกมนํ อาปตฺติ, \\ ปญฺหา เมสา กุสเลหิ จินฺติ ตา.}}\csromangathastanza{\csromangathawak{น กายิกํ กิญฺจิ ปโยคมาจเร, \\ น จาปิ วาจาย ปเร ภเณยฺย. \\ อาปชฺเชยฺย ครุกํ เฉชฺชวตฺถุํ, \\ ปญฺหา เมสา กุสเลหิ จินฺติตา.}}\csromangathastanza{\csromangathawak{น กายิกํ วาจสิกญฺจ กิญฺจิ, \\ มนสาปิ สนฺโต น กเรยฺย ปาปํ. \\ โส นาสิโต กินฺติ สุนาสิโต ภเว, \\ ปญฺหา เมสา กุสเลหิ จินฺติตา.}}\csromangathastanza{\csromangathawak{อนาลปนฺโต มนุเชน เกนจิ, \\ วาจาคิรํ โน จ ปเร ภเณยฺย. \\ อาปชฺเชยฺย วาจสิกํ น กายิกํ, \\ ปญฺหา เมสา กุสเลหิ จินฺติตา.}}\csromangathastanza{\csromangathawak{สิกฺขาปทา พุทฺธวเรน วณฺณิตา, \\ สํฆาทิเสสา จตุโร ภเวยฺยุํ. \\ อาปชฺเชยฺย เอกปโยเคน สพฺเพ, \\ ปญฺหา เมสา กุสเลหิ จินฺติตา.}}\csromangathastanza{\csromangathawak{อุโภ เอกโต อุปสมฺปนฺนา, \\ อุภินฺนํ หตฺถโต จีวรํ ปฏิคฺคเณฺหยฺย. \\ สิยา อาปตฺติโย นานา, \\ ปญฺหา เมสา กุสลิหิ จินฺติตา.}}\csromangathastanza{\csromangathawak{จตุโร ชานา สํวิธาย, \\ ครุภณฺฑํ อวาหรุํ. \\ ตโย ปาราชิกา, เอโก น ปาราชิโก, \\ ปญฺหา เมสา กุสเลหิ จินฺติตา.}}}

\csromanmatikaanchor{mtk.246}
\csromantocmarkref{subparagraph}{2. ปาราชิกาทิปญฺหา}{mtk.246}
\bookmark[dest={mtk.246},level=5]{2. ปาราชิกาทิปญฺหา}
\csromanheader{5}{2. ปาราชิกา\-ทิ\-ปญฺหา}
\setlength{\csromangathapagewidth}{0pt}
\setlength{\csromangathawakwidth}{0pt}
\csromangathameasuregroup{อิตฺถี จ อพฺภนฺตเรสิยา, \\ ฉิทฺทํ ตสฺมิํ ฆเร นตฺถิ, \\ นิวตฺโถ อนฺตรวาสเกน,}{\csromangathabat{อิตฺถี จ อพฺภนฺตเรสิยา,}{ภิกฺขุ จ พหิทฺธา สิยา.} \\ \csromangathabat{ฉิทฺทํ ตสฺมิํ ฆเร นตฺถิ,}{เมถุนธมฺมปจฺจยา.} \\ กถํ ปาราชิโก สิยา, \\ ปญฺหา เมสา กุสเลหิ จินฺติตา. \\ เตลํ มธุํ ผาณิตญฺจาปิ สปฺปิํ, \\ สามํ คเหตฺวาน นิกฺขิเปยฺย. \\ อวีติวตฺเต สตฺตาเห, \\ สติ ปจฺจเย ปริภุญฺชนฺตสฺส อาปตฺติ. \\ ปญฺหา เมสา กุสเลหิ จินฺติตา. \\ นิสฺสคฺคิเยน อาปตฺติ, \\ สุทฺธเกน ปาจิตฺติยํ. \\ อาปชฺชนฺตสฺส เอกโต, \\ ปญฺหา เมสา กุสเลหิ จินฺติตา. \\ ภิกฺขู สิยา วีสติยา สมาคตา, \\ กมฺมํ กเรยฺยุํ สมคฺคสญฺญิโน. \\ ภิกฺขุ สิยา ทฺวาทสโยชเน ฐิโต, \\ กมฺมญฺจ ตํ กุปฺเปยฺย วคฺคปจฺจยา. \\ ปญฺหา เมสา กุสเลหิ จินฺติตา. \\ ปทวีติหารมตฺเตน วาจาย ภณิเตน จ, \\ สพฺพานิ ครุกานิ สปฺปฏิกมฺมานิ. \\ จตุสฏฺฐิ อาปตฺติโย อาปชฺเชยฺย เอกโต, \\ ปญฺหา เมสา กุสเลหิ จินฺติตา. \\ \csromangathabat{นิวตฺโถ อนฺตรวาสเกน,}{ทิคุณํ สํฆาฏิํ ปารุโต.} \\ สพฺพานิ ตานิ นิสฺสคฺคิยานิ โหนฺติ, \\ ปญฺหา เมสา กุสเลหิ จินฺติตา.}
\csromangathasetleft{อิตฺถี จ อพฺภนฺตเรสิยา, \\ ฉิทฺทํ ตสฺมิํ ฆเร นตฺถิ, \\ นิวตฺโถ อนฺตรวาสเกน,}
\csromangathagroup{\csromangathastanza{\csromanfolio{375}\csromangathaitembat{480}{อิตฺถี จ อพฺภนฺตเรสิยา,}{ภิกฺขุ จ พหิทฺธา สิยา.} \\ \csromangathacontbat{ฉิทฺทํ ตสฺมิํ ฆเร นตฺถิ,}{เมถุนธมฺมปจฺจยา.} \\ \csromangathacontline{กถํ ปาราชิโก สิยา,} \\ \csromangathacontline{ปญฺหา เมสา กุสเลหิ จินฺติตา.} \\ \csromangathacontline{เตลํ มธุํ ผาณิตญฺจาปิ สปฺปิํ,} \\ \csromangathacontline{สามํ คเหตฺวาน นิกฺขิเปยฺย.} \\ \csromangathacontline{อวีติวตฺเต สตฺตาเห,} \\ \csromangathacontline{สติ ปจฺจเย ปริภุญฺช\-นฺตสฺส อาปตฺติ.} \\ \csromangathacontline{ปญฺหา เมสา กุสเลหิ จินฺติตา.} \\ \csromangathacontline{นิสฺสคฺคิเยน อาปตฺติ,} \\ \csromangathacontline{สุทฺธเกน ปาจิตฺติยํ.} \\ \csromangathacontline{อาปชฺชนฺตสฺส เอกโต,} \\ \csromangathacontline{ปญฺหา เมสา กุสเลหิ จินฺติตา.} \\ \csromangathacontline{ภิกฺขู สิยา วีสติยา สมาคตา,} \\ \csromangathacontline{กมฺมํ กเรยฺยุํ สมคฺคสญฺญิโน.} \\ \csromangathacontline{ภิกฺขุ สิยา ทฺวาทสโยชเน ฐิโต,} \\ \csromangathacontline{กมฺมญฺจ ตํ กุปฺเปยฺย วคฺคปจฺจยา.} \\ \csromangathacontline{ปญฺหา เมสา กุสเลหิ จินฺติตา.} \\ \csromangathacontline{ปทวีติหารมตฺเตน วาจาย ภณิเตน จ,} \\ \csromangathacontline{สพฺพานิ ครุกานิ สปฺปฏิกมฺมานิ.} \\ \csromangathacontline{จตุสฏฺฐิ อาปตฺติโย อาปชฺเชยฺย เอกโต,} \\ \csromangathacontline{ปญฺหา เมสา กุสเลหิ จินฺติตา.} \\ \csromangathacontbat{นิวตฺโถ อนฺตรวาสเกน,}{ทิคุณํ สํฆาฏิํ ปารุโต.} \\ \csromangathacontline{สพฺพานิ ตานิ นิสฺสคฺคิยานิ โหนฺติ,} \\ \csromangathacontline{ปญฺหา เมสา กุสเลหิ จินฺติตา.}}}

\prose{\csromanfolio{376}น จาปิ ญตฺติ, น จ ปน กมฺมวาจา,}

\prose{น “เจหิ ภิกฺขู”ติ ชิโน อโวจ.}

\setlength{\csromangathapagewidth}{0pt}
\setlength{\csromangathawakwidth}{0pt}
\csromangathameasuregroup{สรณคมนมฺปิ น ตสฺส อตฺถิ, \\ อิตฺถิํ หเน จ มาตรํ, \\ มาตรํ ปิตรํ หนฺตฺวา, \\ โจทยิตฺวา สารยิตฺวา, \\ กตญฺจ กมฺมํ อกตํ ภเวยฺย, \\ สจฺจํ ภณนฺโต ครุกํ, \\ มุสา ภณนฺโต ครุกํ,}{\csromangathabat{สรณคมนมฺปิ น ตสฺส อตฺถิ,}{อุปสมฺปทา จสฺส อกุปฺปา.} \\ ปญฺหา เมสา กุสเลหิ จินฺติตา. \\ อิตฺถิํ หเน, \\ น มาตรํ, ปุริสญฺจ, น ปิตรํ หเน. \\ หเนยฺย อนริยํ มนฺโท, \\ เตน จานนฺตรํ ผุเส. \\ ปญฺหา เมสา กุสเลหิ จินฺติตา. \\ \csromangathabat{อิตฺถิํ หเน จ มาตรํ,}{ปุริสญฺจ ปิตรํ หเน.} \\ \csromangathabat{มาตรํ ปิตรํ หนฺตฺวา,}{น เตนานนฺตรํ ผุเส.} \\ ปญฺหา เมสา กุสเลหิ จินฺติตา. \\ อโจทยิตฺวา อสฺสารยิตฺวา, \\ อสมฺมุขีภูตสฺส กเรยฺย กมฺมํ. \\ กตญฺจ กมฺมํ สุกตํ ภเวยฺย, \\ การโก จ สํโฆ อนาปตฺติโก สิยา. \\ ปญฺหา เมสา กุสเลหิ จินฺติตา. \\ \csromangathabat{โจทยิตฺวา สารยิตฺวา,}{สมฺมุขีภูตสฺส กเรยฺย กมฺมํ.} \\ \csromangathabat{กตญฺจ กมฺมํ อกตํ ภเวยฺย,}{การโก จ สํโฆ สาปตฺติโก สิยา.} \\ ปญฺหา เมสา กุสเลหิ จินฺติตา. \\ ฉินฺทนฺตสฺส อาปตฺติ, \\ ฉินฺทนฺตสฺส อนาปตฺติ. \\ ฉาเทนฺตสฺส อาปตฺติ, \\ ฉาเทนฺตสฺส อนาปตฺติ. \\ ปญฺหา เมสา กุสเลหิ จินฺติตา. \\ \csromangathabat{สจฺจํ ภณนฺโต ครุกํ,}{มุสา จ ลหุ ภาสโต.} \\ \csromangathabat{มุสา ภณนฺโต ครุกํ,}{สจฺจญฺจ ลหุ ภาสโต.} \\ ปญฺหา เมสา กุสเลหิ จินฺติตา.}
\csromangathameasurewak{น มาตรํ, ปุริสญฺจ, น ปิตรํ หเน. \\ หเนยฺย อนริยํ มนฺโท, \\ เตน จานนฺตรํ ผุเส. \\ ปญฺหา เมสา กุสเลหิ จินฺติตา. \\ ปญฺหา เมสา กุสเลหิ จินฺติตา. \\ อโจทยิตฺวา อสฺสารยิตฺวา, \\ อสมฺมุขีภูตสฺส กเรยฺย กมฺมํ. \\ กตญฺจ กมฺมํ สุกตํ ภเวยฺย, \\ ฉินฺทนฺตสฺส อนาปตฺติ. \\ ฉาเทนฺตสฺส อาปตฺติ, \\ ฉาเทนฺตสฺส อนาปตฺติ. \\ ปญฺหา เมสา กุสเลหิ จินฺติตา. \\ ปญฺหา เมสา กุสเลหิ จินฺติตา.}
\csromangathasetleft{สรณคมนมฺปิ น ตสฺส อตฺถิ, \\ อิตฺถิํ หเน จ มาตรํ, \\ มาตรํ ปิตรํ หนฺตฺวา, \\ โจทยิตฺวา สารยิตฺวา, \\ กตญฺจ กมฺมํ อกตํ ภเวยฺย, \\ สจฺจํ ภณนฺโต ครุกํ, \\ มุสา ภณนฺโต ครุกํ,}
\csromangathagroup{\csromangathastanza{\csromangathabat{สรณคมนมฺปิ น ตสฺส อตฺถิ,}{อุปสมฺปทา จสฺส อกุปฺปา.} \\ ปญฺหา เมสา กุสเลหิ จินฺติตา. \\ อิตฺถิํ หเน,}\csromangathastanza{\csromangathawak{น มาตรํ, ปุริสญฺจ, น ปิตรํ หเน. \\ หเนยฺย อนริยํ มนฺโท, \\ เตน จานนฺตรํ ผุเส. \\ ปญฺหา เมสา กุสเลหิ จินฺติตา.}}\csromangathastanza{\csromangathabat{อิตฺถิํ หเน จ มาตรํ,}{ปุริสญฺจ ปิตรํ หเน.} \\ \csromangathabat{มาตรํ ปิตรํ หนฺตฺวา,}{น เตนานนฺตรํ ผุเส.}}\csromangathastanza{\csromangathawak{ปญฺหา เมสา กุสเลหิ จินฺติตา. \\ อโจทยิตฺวา อสฺสารยิตฺวา, \\ อสมฺมุขีภูตสฺส กเรยฺย กมฺมํ. \\ กตญฺจ กมฺมํ สุกตํ ภเวยฺย,}}\csromangathastanza{การโก จ สํโฆ อนาปตฺติโก สิยา. \\ ปญฺหา เมสา กุสเลหิ จินฺติตา. \\ \csromangathabat{โจทยิตฺวา สารยิตฺวา,}{สมฺมุขีภูตสฺส กเรยฺย กมฺมํ.}}\csromangathastanza{\csromangathabat{กตญฺจ กมฺมํ อกตํ ภเวยฺย,}{การโก จ สํโฆ สาปตฺติโก สิยา.} \\ ปญฺหา เมสา กุสเลหิ จินฺติตา. \\ ฉินฺทนฺตสฺส อาปตฺติ,}\csromangathastanza{\csromangathawak{ฉินฺทนฺตสฺส อนาปตฺติ. \\ ฉาเทนฺตสฺส อาปตฺติ, \\ ฉาเทนฺตสฺส อนาปตฺติ. \\ ปญฺหา เมสา กุสเลหิ จินฺติตา.}}\csromangathastanza{\csromangathabat{สจฺจํ ภณนฺโต ครุกํ,}{มุสา จ ลหุ ภาสโต.} \\ \csromangathabat{มุสา ภณนฺโต ครุกํ,}{สจฺจญฺจ ลหุ ภาสโต.}}\csromangathastanza{\csromangathawak{ปญฺหา เมสา กุสเลหิ จินฺติตา.}}}

\csromanmatikaanchor{mtk.247}
\csromantocmarkref{subparagraph}{3. ปาจิตฺติยาทิปญฺหา}{mtk.247}
\bookmark[dest={mtk.247},level=5]{3. ปาจิตฺติยาทิปญฺหา}
\csromanheader{5}{3. ปาจิตฺติยา\-ทิ\-ปญฺหา}
\setlength{\csromangathapagewidth}{0pt}
\setlength{\csromangathawakwidth}{0pt}
\csromangathameasuregroup{อธิฏฺฐิตํ รชนาย รตฺตํ, \\ น จาปิ โส เวทนาฏฺโฏ ภเวยฺย, \\ สลากํ เทนฺตสฺส โหติ เฉชฺชํ, \\ กุทฺโธ อาราธโก โหติ, \\ อถ โก นาม โส ธมฺโม, \\ อาปชฺชติ ครุกํ น ลหุกํ,}{\csromangathabat{อธิฏฺฐิตํ รชนาย รตฺตํ,}{กปฺปกตมฺปิ สนฺตํ.} \\ ปริภุญฺชนฺตสฺส อาปตฺติ, \\ ปญฺหา เมสา กุสเลหิ จินฺติตา. \\ อตฺถงฺคเต สูริเย ภิกฺขุ มํสานิ ขาทติ, \\ น อุมฺมตฺตโก, น จ ปน ขิตฺตจิตฺโต. \\ \csromangathabat{น จาปิ โส เวทนาฏฺโฏ ภเวยฺย,}{น จ สฺส โหติ อาปตฺติ.} \\ โส จ ธมฺโม สุคเตน เทสิโต, \\ ปญฺหา เมสา กุสเลหิ จินฺติตา. \\ น รตฺตจิตฺโต, น จ ปน เถยฺยจิตฺโต, \\ น จาปิ โส ปรํ มรณาย เจตยิ. \\ \csromangathabat{สลากํ เทนฺตสฺส โหติ เฉชฺชํ,}{ปฏิคฺคณฺหนฺตสฺส ถุลฺลจฺจยํ.} \\ ปญฺหา เมสา กุสเลหิ จินฺติตา. \\ น จาปิ อารญฺญกํ สาสงฺกสมฺมตํ, \\ น จาปิ สํเฆน สมฺมุติ ทินฺนา. \\ น จสฺส กถินํ อตฺถตํ ตตฺเถว, \\ จีวรํ นิกฺขิปิตฺวา คจฺเฉยฺย อฑฺฒโยชนํ. \\ ตตฺเถว อรุณํ อุคฺคจฺฉนฺตสฺส อนาปตฺติ, \\ ปญฺหา เมสา กุสเลหิ จินฺติตา. \\ กายิกานิ น วาจสิกานิ, \\ สพฺพานิ นานาวตฺถุกานิ. \\ อปุพฺพํ อจริมํ อาปชฺเชยฺย เอกโต, \\ ปญฺหา เมสา กุสเลหิ จินฺติตา. \\ วาจสิกานิ น กายิกานิ, \\ สพฺพานิ นานาวตฺถุกานิ. \\ อปุพฺพํ อจริมํ อาปชฺเชยฺย เอกโต, \\ ปญฺหา เมสา กุสเลหิ จินฺติตา. \\ ติสฺสิตฺถิโย เมถุนํ ตํ น เสเว, \\ ตโย ปุริเส ตโย อนริยปณฺฑเก. \\ น จาจเร เมถุนํ พฺยญฺชนสฺมิํ, \\ เฉชฺชํ สิยา เมถุนธมฺมปจฺจยา. \\ ปญฺหา เมสา กุสเลหิ จินฺติตา. \\ มาตรํ จีวรํ ยาเจ, \\ โน จ สํเฆ1 ปริณตํ. \\ เกนสฺส โหติ อาปตฺติ, \\ อนาปตฺติ จ ญาตเก. \\ ปญฺหา เมสา กุสเลหิ จินฺติตา. \\ \csromangathabat{กุทฺโธ อาราธโก โหติ,}{กุทฺโธ โหติ ครหิโย.} \\ \csromangathabat{อถ โก นาม โส ธมฺโม,}{เยน กุทฺโธ ปสํสิโย.} \\ ปญฺหา เมสา กุสเลหิ จินฺติตา. \\ ตุฏฺโฐ อาราธโก โหติ, \\ ตุฏฺโฐ โหติ ครหิโย. \\ อถ โก นาม โส ธมฺโม, \\ เยน ตุฏฺโฐ ครหิโย. \\ ปญฺหา เมสา กุสเลหิ จินฺติตา. \\ สํฆาทิเสสํ ถุลฺลจฺจยํ, \\ ปาจิตฺติยํ ปาฏิเทสนียํ. \\ ทุกฺกฏํ อาปชฺเชยฺย เอกโต, \\ ปญฺหา เมสา กุสเลหิ จินฺติตา. \\ อุโภ ปริปุณฺณวีสติวสฺสา, \\ อุภินฺนํ เอกุปชฺฌาโย. \\ เอกาจริโย เอกา กมฺมวาจา, \\ เอโก อุปสมฺปนฺโน เอโก อนุปสมฺปนฺโน. \\ ปญฺหา เมสา กุสเลหิ จินฺติตา. \\ อกปฺปกตํ นาปิ รชนาย รตฺตํ, \\ เตน นิวตฺโถ เยน กามํ วเชยฺย. \\ น จสฺส โหติ อาปตฺติ, \\ โส จ ธมฺโม สุคเตน เทสิโต. \\ ปญฺหา เมสา กุสเลหิ จินฺติตา. \\ น เทติ, \\ น ปฏิคฺคณฺหาติ, ปฏิคฺคโห เตน น วิชฺชติ. \\ \csromangathabat{อาปชฺชติ ครุกํ น ลหุกํ,}{ตญฺจ ปริโภคปจฺจยา.} \\ ปญฺหา เมสา กุสเลหิ จินฺติตา. \\ น เทติ, \\ น ปฏิคฺคณฺหาติ, ปฏิคฺคโห เตน น วิชฺชติ. \\ อาปชฺชติ ลหุกํ น ครุกํ, \\ ตญฺจ ปริโภคปจฺจยา. \\ ปญฺหา เมสา กุสเลหิ จินฺติตา. \\ อาปชฺชติ ครุกํ สาวเสสํ, \\ ฉาเทติ อนาทริยํ ปฏิจฺจ. \\ น ภิกฺขุนี, โน จ ผุเสยฺย วชฺชํ, \\ ปญฺหา เมสา กุสเลหิ จินฺติตา.}
\csromangathameasurewak{ติสฺสิตฺถิโย เมถุนํ ตํ น เสเว, \\ ตโย ปุริเส ตโย อนริยปณฺฑเก. \\ น จาจเร เมถุนํ พฺยญฺชนสฺมิํ, \\ เฉชฺชํ สิยา เมถุนธมฺมปจฺจยา. \\ ปญฺหา เมสา กุสเลหิ จินฺติตา. \\ มาตรํ จีวรํ ยาเจ, \\ โน จ สํเฆ1 ปริณตํ. \\ เกนสฺส โหติ อาปตฺติ, \\ ตุฏฺโฐ โหติ ครหิโย. \\ อถ โก นาม โส ธมฺโม, \\ เยน ตุฏฺโฐ ครหิโย. \\ ปญฺหา เมสา กุสเลหิ จินฺติตา. \\ สํฆาทิเสสํ ถุลฺลจฺจยํ, \\ ปาจิตฺติยํ ปาฏิเทสนียํ. \\ ทุกฺกฏํ อาปชฺเชยฺย เอกโต, \\ ปญฺหา เมสา กุสเลหิ จินฺติตา. \\ อุโภ ปริปุณฺณวีสติวสฺสา, \\ อุภินฺนํ เอกุปชฺฌาโย. \\ เอกาจริโย เอกา กมฺมวาจา, \\ เอโก อุปสมฺปนฺโน เอโก อนุปสมฺปนฺโน. \\ ปญฺหา เมสา กุสเลหิ จินฺติตา. \\ อกปฺปกตํ นาปิ รชนาย รตฺตํ, \\ เตน นิวตฺโถ เยน กามํ วเชยฺย. \\ น จสฺส โหติ อาปตฺติ, \\ น ปฏิคฺคณฺหาติ, ปฏิคฺคโห เตน น วิชฺชติ. \\ อาปชฺชติ ลหุกํ น ครุกํ, \\ ตญฺจ ปริโภคปจฺจยา. \\ ปญฺหา เมสา กุสเลหิ จินฺติตา. \\ อาปชฺชติ ครุกํ สาวเสสํ, \\ ฉาเทติ อนาทริยํ ปฏิจฺจ. \\ น ภิกฺขุนี, โน จ ผุเสยฺย วชฺชํ, \\ ปญฺหา เมสา กุสเลหิ จินฺติตา.}
\csromangathasetleft{อธิฏฺฐิตํ รชนาย รตฺตํ, \\ น จาปิ โส เวทนาฏฺโฏ ภเวยฺย, \\ สลากํ เทนฺตสฺส โหติ เฉชฺชํ, \\ กุทฺโธ อาราธโก โหติ, \\ อถ โก นาม โส ธมฺโม, \\ อาปชฺชติ ครุกํ น ลหุกํ,}
\csromangathagroupwithcloser{\csromangathastanza{\csromanfolio{377}\csromangathaitembat{481}{อธิฏฺฐิตํ รชนาย รตฺตํ,}{กปฺปกตมฺปิ สนฺตํ.} \\ \csromangathacontline{ปริภุญฺช\-นฺตสฺส อาปตฺติ,} \\ \csromangathacontline{ปญฺหา เมสา กุสเลหิ จินฺติตา.} \\ \csromangathacontline{อตฺถงฺคเต สูริเย ภิกฺขุ มํสานิ ขาทติ,} \\ \csromangathacontline{น อุมฺมตฺตโก, น จ ปน ขิตฺตจิตฺโต.} \\ \csromangathacontbat{น จาปิ โส เวทนาฏฺโฏ ภเวยฺย,}{น จ สฺส โหติ อาปตฺติ.} \\ \csromangathacontline{โส จ ธมฺโม สุคเตน เทสิโต,} \\ \csromangathacontline{ปญฺหา เมสา กุสเลหิ จินฺติตา.} \\ \csromangathacontline{น รตฺตจิตฺโต, น จ ปน เถยฺยจิตฺโต,} \\ \csromangathacontline{น จาปิ โส ปรํ มรณาย เจตยิ.} \\ \csromangathacontbat{สลากํ เทนฺตสฺส โหติ เฉชฺชํ,}{ปฏิคฺคณฺหนฺตสฺส ถุลฺลจฺจยํ.} \\ \csromangathacontline{ปญฺหา เมสา กุสเลหิ จินฺติตา.} \\ \csromangathacontline{น จาปิ อารญฺญกํ สาสงฺกสมฺมตํ,} \\ \csromangathacontline{น จาปิ สํเฆน สมฺมุติ ทินฺนา.} \\ \csromangathacontline{น จสฺส กถินํ อตฺถตํ ตตฺเถว,} \\ \csromangathacontline{จีวรํ นิกฺขิปิตฺวา คจฺเฉยฺย อฑฺฒโยชนํ.} \\ \csromangathacontline{ตตฺเถว อรุณํ อุคฺคจฺฉนฺตสฺส อนาปตฺติ,} \\ \csromangathacontline{ปญฺหา เมสา กุสเลหิ จินฺติตา.} \\ \csromangathacontline{กายิกานิ น วาจสิกานิ,} \\ \csromangathacontline{สพฺพานิ นานาวตฺถุกานิ.} \\ \csromangathacontline{อปุพฺพํ อจริมํ อาปชฺเชยฺย เอกโต,} \\ \csromangathacontline{ปญฺหา เมสา กุสเลหิ จินฺติตา.} \\ \csromangathacontline{วาจสิกานิ น กายิกานิ,} \\ \csromangathacontline{สพฺพานิ นานาวตฺถุกานิ.} \\ \csromangathacontline{อปุพฺพํ อจริมํ อาปชฺเชยฺย เอกโต,} \\ \csromangathacontline{ปญฺหา เมสา กุสเลหิ จินฺติตา.}}\csromangathastanza{\csromangathawak{\csromanfolio{378}ติสฺสิตฺถิโย เมถุนํ ตํ น เสเว, \\ ตโย ปุริเส ตโย อนริยปณฺฑเก. \\ น จาจเร เมถุนํ พฺยญฺชนสฺมิํ, \\ เฉชฺชํ สิยา เมถุนธมฺมปจฺจยา.}}\csromangathastanza{\csromangathawak{ปญฺหา เมสา กุสเลหิ จินฺติตา. \\ มาตรํ จีวรํ ยาเจ, \\ โน จ สํเฆ1 ปริณตํ. \\ เกนสฺส โหติ อาปตฺติ,}}\csromangathastanza{อนาปตฺติ จ ญาตเก. \\ ปญฺหา เมสา กุสเลหิ จินฺติตา. \\ \csromangathabat{กุทฺโธ อาราธโก โหติ,}{กุทฺโธ โหติ ครหิโย.}}\csromangathastanza{\csromangathabat{อถ โก นาม โส ธมฺโม,}{เยน กุทฺโธ ปสํสิโย.} \\ ปญฺหา เมสา กุสเลหิ จินฺติตา. \\ ตุฏฺโฐ อาราธโก โหติ,}\csromangathastanza{\csromangathawak{ตุฏฺโฐ โหติ ครหิโย. \\ อถ โก นาม โส ธมฺโม, \\ เยน ตุฏฺโฐ ครหิโย. \\ ปญฺหา เมสา กุสเลหิ จินฺติตา.}}\csromangathastanza{\csromangathawak{สํฆาทิเสสํ ถุลฺลจฺจยํ, \\ ปาจิตฺติยํ ปาฏิเทสนียํ. \\ ทุกฺกฏํ อาปชฺเชยฺย เอกโต, \\ ปญฺหา เมสา กุสเลหิ จินฺติตา.}}\csromangathastanza{\csromangathawak{อุโภ ปริปุณฺณ\-วีสติวสฺสา, \\ อุภินฺนํ เอกุปชฺฌาโย. \\ เอกาจริโย เอกา กมฺมวาจา, \\ เอโก อุปสมฺปนฺโน เอโก อนุปสมฺปนฺโน.}}\csromangathastanza{\csromangathawak{ปญฺหา เมสา กุสเลหิ จินฺติตา. \\ อกปฺปกตํ นาปิ รชนาย รตฺตํ, \\ เตน นิวตฺโถ เยน กามํ วเชยฺย. \\ น จสฺส โหติ อาปตฺติ,}}\csromangathastanza{โส จ ธมฺโม สุคเตน เทสิโต. \\ ปญฺหา เมสา กุสเลหิ จินฺติตา. \\ น เทติ, \\ น ปฏิคฺคณฺหาติ, ปฏิคฺคโห เตน น วิชฺชติ.}\csromangathastanza{\csromanfolio{379}\csromangathabat{อาปชฺชติ ครุกํ น ลหุกํ,}{ตญฺจ ปริโภคปจฺจยา.} \\ ปญฺหา เมสา กุสเลหิ จินฺติตา. \\ น เทติ,}\csromangathastanza{\csromangathawak{น ปฏิคฺคณฺหาติ, ปฏิคฺคโห เตน น วิชฺชติ. \\ อาปชฺชติ ลหุกํ น ครุกํ, \\ ตญฺจ ปริโภคปจฺจยา. \\ ปญฺหา เมสา กุสเลหิ จินฺติตา.}}\csromangathastanza{\csromangathawak{อาปชฺชติ ครุกํ สาวเสสํ, \\ ฉาเทติ อนาทริยํ ปฏิจฺจ. \\ น ภิกฺขุนี, โน จ ผุเสยฺย วชฺชํ, \\ ปญฺหา เมสา กุสเลหิ จินฺติตา.}}}{เสท\-โมจน\-คาถา นิฏฺฐิตา.}
\csromancenter{ตสฺสุทฺทานํ.}
\setlength{\csromangathapagewidth}{0pt}
\setlength{\csromangathawakwidth}{0pt}
\csromangathameasuregroup{อสํวาโส อวิสฺสชฺชิ, \\ อุเปติ ธมฺมํ อุพฺภกฺขกํ, \\ น กายิกญฺจ ครุกํ, \\ อนาลปนฺโต สิกฺขา จ, \\ อิตฺถี เตลญฺจ นิสฺสคฺคิ, \\ นิวตฺโถ จ น จ ญตฺติ, \\ อโจทยิตฺวา โจทยิตฺวา, \\ อธิฏฺฐิตญฺจตฺถงฺคเต, \\ กายิกา วาจสิกา จ, \\ กุทฺโธ อาราธโก ตุฏฺโฐ, \\ อกปฺปกตํ น เทติ,}{\csromangathabat{อสํวาโส อวิสฺสชฺชิ,}{ทส จ อนุกฺขิตฺตโก.} \\ \csromangathabat{อุเปติ ธมฺมํ อุพฺภกฺขกํ,}{ตโต สญฺญาจิกา จ เทฺว.} \\ \csromangathabat{น กายิกญฺจ ครุกํ,}{น กายิกํ น วาจสิกํ\textsuperscript{1}.} \\ \csromangathabat{อนาลปนฺโต สิกฺขา จ,}{อุโภ จ จตุโร ชนา.} \\ \csromangathabat{อิตฺถี เตลญฺจ นิสฺสคฺคิ,}{ภิกฺขุ จ ปทวีติโย.} \\ \csromangathabat{นิวตฺโถ จ น จ ญตฺติ,}{น มาตรํ ปิตรํ หเน.} \\ \csromangathabat{อโจทยิตฺวา โจทยิตฺวา,}{ฉินฺทนฺตํ สจฺจเมว จ.} \\ \csromangathabat{อธิฏฺฐิตญฺจตฺถงฺคเต,}{น รตฺตํ น จารญฺญกํ.} \\ \csromangathabat{กายิกา วาจสิกา จ,}{ติสฺสิตฺถี จาปิ มาตรํ.} \\ \csromangathabat{กุทฺโธ อาราธโก ตุฏฺโฐ,}{สํฆาทิเสสา จ อุโภ.} \\ \csromangathabat{อกปฺปกตํ น เทติ,}{น เทตาปชฺชตี ครุํ.}}
\csromangathasetleft{อสํวาโส อวิสฺสชฺชิ, \\ อุเปติ ธมฺมํ อุพฺภกฺขกํ, \\ น กายิกญฺจ ครุกํ, \\ อนาลปนฺโต สิกฺขา จ, \\ อิตฺถี เตลญฺจ นิสฺสคฺคิ, \\ นิวตฺโถ จ น จ ญตฺติ, \\ อโจทยิตฺวา โจทยิตฺวา, \\ อธิฏฺฐิตญฺจตฺถงฺคเต, \\ กายิกา วาจสิกา จ, \\ กุทฺโธ อาราธโก ตุฏฺโฐ, \\ อกปฺปกตํ น เทติ,}
\csromangathagroup{\csromangathastanza{\csromangathabat{อสํวาโส อวิสฺสชฺชิ,}{ทส จ อนุกฺขิตฺตโก.} \\ \csromangathabat{อุเปติ ธมฺมํ อุพฺภกฺขกํ,}{ตโต สญฺญาจิกา จ เทฺว.}}\csromangathastanza{\csromangathabat{น กายิกญฺจ ครุกํ,}{น กายิกํ น วาจสิกํ\csromansharedfootnote{826a70cb5f65511b}{น กายิกํ สุนาสิตํ (สฺยา)}.} \\ \csromangathabat{อนาลปนฺโต สิกฺขา จ,}{อุโภ จ จตุโร ชนา.}}\csromangathastanza{\csromangathabat{อิตฺถี เตลญฺจ นิสฺสคฺคิ,}{ภิกฺขุ จ ปทวีติโย.} \\ \csromangathabat{นิวตฺโถ จ น จ ญตฺติ,}{น มาตรํ ปิตรํ หเน.}}\csromangathastanza{\csromangathabat{อโจทยิตฺวา โจทยิตฺวา,}{ฉินฺทนฺตํ สจฺจเมว จ.} \\ \csromangathabat{อธิฏฺฐิตญฺจตฺถงฺคเต,}{น รตฺตํ น จารญฺญกํ.}}\csromangathastanza{\csromangathabat{กายิกา วาจสิกา จ,}{ติสฺสิตฺถี จาปิ มาตรํ.} \\ \csromangathabat{กุทฺโธ อาราธโก ตุฏฺโฐ,}{สํฆาทิเสสา จ อุโภ.} \\ \csromangathabat{อกปฺปกตํ น เทติ,}{น เทตาปชฺชตี ครุํ.}}}

\csromanheader{2}{เสทโมจนิกา คาถา, ปญฺหา วิญฺญูหิ วิภาวิตาติ\csromansharedfootnote{492e105b8e02d419}{วิญฺญูวิภาวิตา (สี, สฺยา)}.}
\csromanmatikaanchor{mtk.248}
\csromantocmarknopage{section}{ปญฺจวคฺค}{mtk.248}
\bookmark[dest={mtk.248},level=1]{ปญฺจวคฺค}
\csromanheader{1}{\textbf{ปญฺจวคฺค}}
\csromanmatikaanchor{mtk.249}
\csromantocmarkref{subparagraph}{1. กมฺมวคฺค}{mtk.249}
\bookmark[dest={mtk.249},level=5]{1. กมฺมวคฺค}
\csromanheader{6}{1. \textbf{กมฺมวคฺค}}
\proseitem{482}{\csromanfolio{380}จตฺตาริ กมฺมานิ อปโลกน\-กมฺมํ ญตฺติกมฺมํ ญตฺติ\-ทุติย\-กมฺมํ ญตฺติ\-จตุตฺถ\-กมฺมํ. อิมานิ จตฺตาริ กมฺมานิ กติหากาเรหิ วิปชฺชนฺติ. อิมานิ จตฺตาริ กมฺมานิ ปญฺจหากาเรหิ วิปชฺชนฺติ วตฺถุโต วา ญตฺติโต วา อนุสฺสาวนโต วา สีมโต วา ปริสโต วา.}

\proseitem{483}{กถํ วตฺถุโต กมฺมานิ วิปชฺชนฺติ. สมฺมุขา\-กรณียํ กมฺมํ อสมฺมุขา กโรติ, วตฺถุวิปนฺนํ อธมฺมกมฺมํ, ปฏิปุจฺฉา\-กรณียํ กมฺมํ อปฏิปุจฺฉา กโรติ, วตฺถุวิปนฺนํ อธมฺมกมฺมํ. ปฏิญฺญาย กรณียํ กมฺมํ อปฏิญฺญาย กโรติ, วตฺถุวิปนฺนํ อธมฺมกมฺมํ. สติวินยา\-รหสฺส อมูฬฺหวินยํ เทติ, วตฺถุวิปนฺนํ อธมฺมกมฺมํ. อมูฬฺหวินยา\-รหสฺส ตสฺสปาปิยสิกา\-กมฺมํ กโรติ, วตฺถุวิปนฺนํ อธมฺมกมฺมํ. ตสฺสปาปิยสิกา\-กมฺมา\-รหสฺส ตชฺชนีย\-กมฺมํ กโรติ, วตฺถุวิปนฺนํ อธมฺมกมฺมํ. ตชฺชนียกมฺมา\-รหสฺส นิยสฺสกมฺมํ กโรติ, วตฺถุวิปนฺนํ อธมฺมกมฺมํ. นิยสฺสกมฺมา\-รหสฺส ปพฺพาชนีย\-กมฺมํ กโรติ, วตฺถุวิปนฺนํ อธมฺมกมฺมํ. ปพฺพาชนียกมฺมา\-รหสฺส ปฏิสารณีย\-กมฺมํ กโรติ, วตฺถุวิปนํ อธมฺมกมฺมํ. ปฏิสารณียกมฺมา\-รหสฺส อุกฺเขปนีย\-กมฺมํ กโรติ, วตฺถุวิปนฺนํ อธมฺมกมฺมํ. อุกฺเขปนียกมฺมา\-รหสฺส ปริวาสํ เทติ, วตฺถุวิปนฺนํ อธมฺมกมฺมํ. ปริวาสารหํ มูลาย ปฏิกสฺสติ, วตฺถุวิปนฺนํ อธมฺมกมฺมํ. มูลาย\-ปฏิกสฺสนา\-รหสฺส มานตฺตํ เทติ, วตฺถุวิปนฺนํ อธมฺมกมฺมํ. มานตฺตารหํ อพฺเภติ, วตฺถุวิปนฺนํ อธมฺมกมฺมํ. อพฺภานารหํ อุปสมฺปาเทติ, วตฺถุวิปนฺนํ อธมฺมกมฺมํ. อนุโปสเถ อุโปสถํ กโรติ, วตฺถุวิปนฺนํ อธมฺมกมฺมํ. อปวารณาย ปวาเรติ, วตฺถุวิปนฺนํ อธมฺมกมฺมํ. เอวํ วตฺถุโต กมฺมานิ วิปชฺชนฺติ.}

\proseitem{484}{กถํ \textbf{ญตฺติโต กมฺมานิ วิปชฺชนฺติ.} ปญฺจหากาเรหิ ญตฺติโต กมฺมานิ วิปชฺชนฺติ. วตฺถุํ น ปรามสติ, สํฆํ น ปรามสติ, ปุคฺคลํ น ปรามสติ, ญตฺติํ น ปรามสติ, ปจฺฉา วา ญตฺติํ ฐเปติ. อิเมหิ ปญฺจหากาเรหิ ญตฺติโต กมฺมานิ วิปชฺชนฺติ.}

\proseitem{485}{\csromanfolio{381}กถํ \textbf{อนุสฺสาวนโต กมฺมานิ วิปชฺชนฺติ.} ปญฺจหากาเรหิ อนุสฺสาวนโต กมฺมานิ วิปชฺชนฺติ. วตฺถุํ น ปรามสติ, สํฆํ น ปรามสติ, ปุคฺคลํ น ปรามสติ, สาวนํ หาเปติ, อกาเล วา สาเวติ. อิเมหิ ปญฺจหากาเรหิ อนุสฺสาวนโต กมฺมานิ วิปชฺชนฺติ.}

\proseitem{486}{กถํ \textbf{สีมโต กมฺมานิ วิปชฺชนฺติ.} เอกาทสหิ อากาเรหิ สีมโต กมฺมานิ วิปชฺชนฺติ. อติขุทฺทกํ สีมํ สมฺมนฺนติ, อติมหติํ สีมํ สมฺมนฺนติ, ขณฺฑนิมิตฺตํ สีมํ สมฺมนฺนติ, ฉายานิมิตฺตํ สีมํ สมฺมนฺนติ, อนิมิตฺตํ สีมํ สมฺมนฺนติ, พหิสีเม ฐิโต สีมํ สมฺมนฺนติ, นทิยา สีมํ สมฺมนฺนติ, สมุทฺเท สีมํ สมฺมนฺนติ, ชาตสฺสเร สีมํ สมฺมนฺนติ, สีมาย สีมํ สมฺภินฺทติ, สีมาย สีมํ อชฺโฌตฺถรติ. อิเมหิ เอกาทสหิ อากาเรหิ สีมโต กมฺมานิ วิปชฺชนฺติ.}

\proseitem{487}{กถํ ปริสโต กมฺมานิ วิปชฺชนฺติ. ทฺวาทสหิ อากาเรหิ ปริสโต กมฺมานิ วิปชฺชนฺติ. จตุวคฺค\-กรเณ กมฺเม ยาวติกา ภิกฺขู กมฺมปตฺตา, เต อนาคตา โหนฺติ, ฉนฺทารหานํ ฉนฺโท อนาหโฏ โหติ, สมฺมุขีภูตา ปฏิกฺโกสนฺติ. จตุวคฺค\-กรเณ กมฺเม ยาวติกา ภิกฺขู กมฺมปตฺตา เต อาคตา โหนฺติ, ฉนฺทารหานํ ฉนฺโท อนาหโฏ โหติ, สมฺมุขีภูตา ปฏิกฺโกสนฺติ. จตุวคฺค\-กรเณ กมฺเม ยาวติกา ภิกฺขู กมฺมปตฺตา, เต อาคตา โหนฺติ, ฉนฺทารหานํ ฉนฺโท อาหโฏ โหติ, สมฺมุขีภูตา ปฏิกฺโกสนฺติ. ปญฺจวคฺค\-กรเณ กมฺเม ฯเปฯ. ทสวคฺค\-กรเณ กมฺเม ฯเปฯ. วีสติวคฺค\-กรเณ กมฺเม ยาวติกา ภิกฺขู กมฺมปตฺตา, เต อนาคตา โหนฺติ, ฉนฺทารหานํ ฉนฺโท อนาหโฏ โหติ, สมฺมุขีภูตา ปฏิกฺโกสนฺติ. วีสติวคฺค\-กรเณ กมฺเม ยาวติกา ภิกฺขู กมฺมปตฺตา, เต อาคตา โหนฺติ, ฉนฺทารหานํ ฉนฺโท อนาหโฏ โหติ, สมฺมุขีภูตา ปฏิกฺโกสนฺติ. วีสติวคฺค\-กรเณ กมฺเม ยาวติกา ภิกฺขู กมฺมปตฺตา, เต อาคตา โหนฺติ, ฉนฺทารหานํ ฉนฺโท อาหโฏ โหติ, สมฺมุขีภูตา ปฏิกฺโกสนฺติ. อิเมหิ ทฺวาทสหิ อากาเรหิ ปริสโต กมฺมานิ วิปชฺชนฺติ.}

\proseitem{488}{จตุวคฺค\-กรเณ กมฺเม จตฺตาโร ภิกฺขู ปกตตฺตา กมฺมปตฺตา, อวเสสา ปกตตฺตา ฉนฺทารหา. ยสฺส สํโฆ กมฺมํ กโรติ, โส เนว กมฺมปตฺโต, นาปิ ฉนฺทารโห, อปิ จ กมฺมารโห. \csromanfolio{382}ปญฺจวคฺค\-กรเณ กมฺเม ปญฺจ ภิกฺขู ปกตตฺตา กมฺมปตฺตา, อวเสสา ปกตตฺตา ฉนฺทารหา, ยสฺส สํโฆ กมฺมํ กโรติ, โส เนว กมฺมปตฺโต, นาปิ ฉนฺทารโห, อปิ จ กมฺมารโห. ทสวคฺค\-กรเณ กมฺเม ทส ภิกฺขู ปกตตฺตา กมฺมปตฺตา, อวเสสา ปกตตฺตา ฉนฺทารหา, ยสฺส สํโฆ กมฺมํ กโรติ, โส เนว กมฺมปตฺโต, นาปิ ฉนฺทารโห, อปิ จ กมฺมารโห. วีสติวคฺค\-กรเณ กมฺเม วีสติ ภิกฺขู ปกตตฺตา กมฺมปตฺตา, อวเสสา ปกตตฺตา ฉนฺทารหา, ยสฺส สํโฆ กมฺมํ กโรติ, โส เนว กมฺมปตฺโต, นาปิ ฉนฺทารโห, อปิ จ กมฺมารโห.}

\proseitem{489}{จตฺตาริ กมฺมานิ, อปโลกน\-กมฺมํ ญตฺติกมฺมํ ญตฺติ\-ทุติย\-กมฺมํ ญตฺติ\-จตุตฺถ\-กมฺมํ. อิมานิ จตฺตาริ กมฺมานิ กติหากาเรหิ วิปชฺชนฺติ. อิมานิ จตฺตาริ กมฺมานิ ปญฺจหากาเรหิ วิปชฺชนฺติ, วตฺถุโต วา ญตฺติโต วา อนุสฺสาวนโต วา สีมโต วา ปริสโต วา.}

\proseitem{490}{กถํ วตฺถุโต กมฺมานิ วิปชฺชนฺติ. ปณฺฑกํ อุปสมฺปาเทติ, วตฺถุวิปนฺนํ อธมฺมกมฺมํ. เถยฺย\-สํวาสกํ อุปสมฺปาเทติ, วตฺถุวิปนฺนํ อธมฺมกมฺมํ. ติตฺถิย\-ปกฺกนฺตกํ อุปสมฺปาเทติ, วตฺถุวิปนฺนํ อธมฺมกมฺมํ. ติรจฺฉานคตํ อุปสมฺปาเทติ, วตฺถุวิปนฺนํ อธมฺมกมฺมํ. มาตุฆาตกํ อุปสมฺปาเทติ, วตฺถุวิปนฺนํ อธมฺมกมฺมํ. ปิตุฆาตกํ อุปสมฺปาเทติ, วตฺถุวิปนฺนํ อธมฺมกมฺมํ. อรหนฺต\-ฆาตกํ อุปสมฺปาเทติ, วตฺถุวิปนฺนํ อธมฺมกมฺมํ. ภิกฺขุนิ\-ทูสกํ อุปสมฺปาเทติ, วตฺถุวิปนฺนํ อธมฺมกมฺมํ. สํฆเภทกํ อุปสมฺปาเทติ, วตฺถุวิปนฺนํ อธมฺมกมฺมํ. โลหิตุปฺปาทกํ อุปสมฺปาเทติ, วตฺถุวิปนฺนํ อธมฺมกมฺมํ. อุภโต\-พฺยญฺชนํ อุปสมฺปาเทติ, วตฺถุวิปนฺนํ อธมฺมกมฺมํ. อูนวีสติวสฺสํ ปุคฺคลํ อุปสมฺปาเทติ, วตฺถุวิปนฺนํ อธมฺมกมฺมํ. เอวํ วตฺถุโต กมฺมานิ วิปชฺชนฺติ.}

\proseitem{491}{กถํ \textbf{ญตฺติโต กมฺมานิ วิปชฺชนฺติ.} ปญฺจหากาเรหิ ญตฺติโต กมฺมานิ วิปชฺชนฺติ. วตฺถุํ น ปรามสติ, สํฆํ น ปรามสติ, ปุคฺคลํ น ปรามสติ, ญตฺติํ น ปรามสติ, ปจฺฉา วา ญตฺติํ ฐเปติ. อิเมหิ ปญฺจหากาเรหิ ญตฺติโต กมฺมานิ วิปชฺชนฺติ.}

\proseitem{492}{กถํ \textbf{อนุสฺสาวนโต กมฺมานิ วิปชฺชนฺติ.} ปญฺจหากาเรหิ อนุสฺสาวนโต กมฺมานิ วิปชฺชนฺติ. วตฺถุํ น ปรามสติ, สํฆํ น ปรามสติ, \csromanfolio{383}ปุคฺคลํ น ปรามสติ, สาวนํ หาเปติ, อกาเล วา สาเวติ. อิเมหิ ปญฺจหากาเรหิ อนุสฺสาวนโต กมฺมานิ วิปชฺชนฺติ.}

\proseitem{493}{กถํ \textbf{สีมโต กมฺมานิ วิปชฺชนฺติ.} เอกาทสหิ อากาเรหิ สีมโต กมฺมานิ วิปชฺชนฺติ. อติขุทฺทกํ สีมํ สมฺมนฺนติ, อติมหติํ สีมํ สมฺมนฺนติ, ขณฺฑนิมิตฺตํ สีมํ สมฺมนฺนติ, ฉายานิมิตฺตํ สีมํ สมฺมนฺนติ, อนิมิตฺตํ สีมํ สมฺมนฺนติ, พหิสีเม ฐิโต สีมํ สมฺมนฺนติ, นทิยา สีมํ สมฺมนฺนติ, สมุทฺเท สีมํ สมฺมนฺนติ, ชาตสฺสเร สีมํ สมฺมนฺนติ, สีมาย สีมํ สมฺภินฺทติ, สีมาย สีมํ อชฺโฌตฺถรติ. อิเมหิ เอกาทสหิ อากาเรหิ สีมโต กมฺมานิ วิปชฺชนฺติ.}

\proseitem{494}{กถํ ปริสโต กมฺมานิ วิปชฺชนฺติ. ทฺวาทสหิ อากาเรหิ ปริสโต กมฺมานิ วิปชฺชนฺติ. จตุวคฺค\-กรเณ กมฺเม ยาวติกา ภิกฺขู กมฺมปตฺตา, เต อนาคตา โหนฺติ, ฉนฺทารหานํ ฉนฺทา อนาหโฏ โหติ, สมฺมุขีภูตา ปฏิกฺโกสนฺติ. จตุวคฺค\-กรเณ กมฺเม ยาวติกา ภิกฺขู กมฺมปตฺตา, เต อาคตา โหนฺติ, ฉนฺทารหานํ ฉนฺโท อนาหโฏ โหติ, สมฺมุขีภูตา ปฏิกฺโกสนฺติ. จตุวคฺค\-กรเณ กมฺเม ยาวติกา ภิกฺขู กมฺมปตฺตา, เต อาคตา โหนฺติ, ฉนฺทารหานํ ฉนฺโท อาหโฏ โหติ, สมฺมุขีภูตา ปฏิกฺโกสนฺติ. ปญฺจวคฺค\-กรเณ กมฺเม ฯเปฯ. ทสวคฺค\-กรเณ กมฺเม ฯเปฯ. วีสติวคฺค\-กรเณ กมฺเม ยาวติกา ภิกฺขู กมฺมปตฺตา, เต อนาคตา โหนฺติ, ฉนฺทารหานํ ฉนฺโท อนาหโฏ โหติ, สมฺมุขีภูตา ปฏิกฺโกสนฺติ. วีสติวคฺค\-กรเณ กมฺเม ยาวติกา ภิกฺขู กมฺมปตฺตา, เต อาคตา โหนฺติ, ฉนฺทารหานํ ฉนฺโท อนาหโฏ โหติ, สมฺมุขีภูตา ปฏิกฺโกสนฺติ. วีสติวคฺค\-กรเณ กมฺเม ยาวติกา ภิกฺขู กมฺมปตฺตา, เต อาคตา โหนฺติ, ฉนฺทารหานํ ฉนฺโท อาหโฏ โหติ, สมฺมุขีภูตา ปฏิกฺโกสนฺติ. อิเมหิ ทฺวาทสหิ อากาเรหิ ปริสโต กมฺมานิ วิปชฺชนฺติ.}

\proseitem{495}{อปโลกน\-กมฺมํ กติ ฐานานิ คจฺฉติ. ญตฺติกมฺมํ กติ ฐานานิ คจฺฉติ. ญตฺติ\-ทุติย\-กมฺมํ กติ ฐานานิ คจฺฉติ. ญตฺติ\-จตุตฺถ\-กมฺมํ กติ ฐานานิ คจฺฉติ. อปโลกน\-กมฺมํ ปญฺจ ฐานานิ คจฺฉติ. ญตฺติกมฺมํ นว ฐานานิ คจฺฉติ. ญตฺติ\-ทุติย\-กมฺมํ สตฺต ฐานานิ คจฺฉติ. ญตฺติ\-จตุตฺถ\-กมฺมํ สตฺต ฐานานิ คจฺฉติ.}

\proseitem{496}{\csromanfolio{384}อปโลกน\-กมฺมํ กตมานิ ปญฺจ ฐานานิ คจฺฉติ. โอสารณํ นิสฺสารณํ ภณฺฑุกมฺมํ พฺรหฺมทณฺฑํ กมฺม\-ลกฺขณญฺเญว ปญฺจมํ. อปโลกน\-กมฺมํ อิมานิ ปญฺจ ฐานานิ คจฺฉติ. ญตฺติกมฺมํ กตมานิ นว ฐานานิ คจฺฉติ. โอสารณํ นิสฺสารณํ อุโปสถํ ปวารณํ สมฺมุติํ ทานํ ปฏิคฺคหํ ปจฺจุกฺกฑฺฒนํ กมฺม\-ลกฺขณญฺเญว นวมํ. ญตฺติกมฺมํ อิมานิ นว ฐานานิ คจฺฉติ. ญตฺติ\-ทุติย\-กมฺมํ กตมานิ สตฺต ฐานานิ คจฺฉติ. โอสารณํ นิสฺสารณํ สมฺมุติํ ทานํ อุทฺธรณํ เทสนํ กมฺม\-ลกฺขณญฺเญว สตฺตมํ. ญตฺติ\-ทุติย\-กมฺมํ อิมานิ สตฺต ฐานานิ คจฺฉติ. ญตฺติ\-จตุตฺถ\-กมฺมํ กตมานิ สตฺต ฐานานิ คจฺฉติ. โอสารณํ นิสฺสารณํ สมฺมุติํ ทานํ นิคฺคหํ สมนุภาสนํ กมฺม\-ลกฺขณญฺเญว สตฺตมํ. ญตฺติ\-จตุตฺถ\-กมฺมํ อิมานิ สตฺต ฐานานิ คจฺฉติ.}

\proseitemwithclosermidruled{497}{จตุวคฺค\-กรเณ กมฺเม จตฺตาโร ภิกฺขู ปกตตฺตา กมฺมปตฺตา, อวเสสา ปกตตฺตา ฉนฺทารหา, ยสฺส สํโฆ กมฺมํ กโรติ, โส เนว กมฺมปตฺโต, นาปิ ฉนฺทารโห, อปิ จ กมฺมารโห. ปญฺจวคฺค\-กรเณ กมฺเม ปญฺจ ภิกฺขู ปกตตฺตา กมฺมปตฺตา, อวเสสา ปกตตฺตา ฉนฺทารหา, ยสฺส สํโฆ กมฺมํ กโรติ, โส เนว กมฺมปตฺโต, นาปิ ฉนฺทารโห, อปิ จ กมฺมารโห. ทสวคฺค\-กรเณ กมฺเม ทส ภิกฺขู ปกตตฺตา กมฺมปตฺตา, อวเสสา ปกตตฺตา ฉนฺทารหา, ยสฺส สํโฆ กมฺมํ กโรติ, โส เนว กมฺมปตฺโต, นาปิ ฉนฺธารโห, อปิ จ กมฺมารโห. วีสติวคฺค\-กรเณ กมฺเม วีสติ ภิกฺขู ปกตตฺตา กมฺมปตฺตา, อวเสสา ปกตตฺตา ฉนฺทารหา, ยสฺส สํโฆ กมฺมํ กโรติ, โส เนว กมฺมปตฺโต, นาปิ ฉนฺทารโห, อปิ จ กมฺมารโห.}{กมฺมวคฺโค นิฏฺฐิโต ปฐโม.}
\csromanmatikaanchor{mtk.250}
\csromantocmarkref{subparagraph}{2. อตฺถวสวคฺค}{mtk.250}
\bookmark[dest={mtk.250},level=5]{2. อตฺถวสวคฺค}
\csromanheader{6}{2. \textbf{อตฺถวสวคฺค}}
\proseitemwithclosermidruled{498}{\csromansymbolfootnote{*}{อํ 1. 97 ปิฏฺเฐปิ.}เทฺว อตฺถวเส ปฏิจฺจ ตถาคเตน สาวกานํ สิกฺขาปทํ ปญฺญตฺตํ สํฆสุฏฺฐุตาย สํฆผาสุตาย, อิเม เทฺว อตฺถวเส ปฏิจฺจ ตถาคเตน สาวกานํ สิกฺขาปทํ ปญฺญตฺตํ. เทฺว อตฺถวเส ปฏิจฺจ ตถาคเตน สาวกานํ สิกฺขาปทํ ปญฺญตฺตํ ทุมฺมงฺกูนํ ปุคฺคลานํ นิคฺคหาย เปสลานํ ภิกฺขูนํ ผาสุวิหาราย, อิเม เทฺว อตฺถวเส ปฏิจฺจ \csromanfolio{385}ตถาคเตน สาวกานํ สิกฺขาปทํ ปญฺญตฺตํ. เทฺว อตฺถวเส ปฏิจฺจ ตถาคเตน สาวกานํ สิกฺขาปทํ ปญฺญตฺตํ ทิฏฺฐ\-ธมฺมิกานํ อาสวานํ สํวราย สมฺปรายิกานํ อาสวานํ ปฏิฆาตาย, อิเม เทฺว อตฺถวเส ปฏิจฺจ ตถาคเตน สาวกานํ สิกฺขาปทํ ปญฺญตฺตํ. เทฺว อตฺถวเส ปฏิจฺจ ตถาคเตน สาวกานํ สิกฺขาปทํ ปญฺญตฺตํ ทิฏฺฐ\-ธมฺมิกานํ เวรานํ สํวราย สมฺปรายิกานํ เวรานํ ปฏิฆาตาย, อิเม เทฺว อตฺถวเส ปฏิจฺจ ตถาคเตน สาวกานํ สิกฺขาปทํ ปญฺญตฺตํ. เทฺว อตฺถวเส ปฏิจฺจ ตถาคเตน สาวกานํ สิกฺขาปทํ ปญฺญตฺตํ ทิฏฺฐ\-ธมฺมิกานํ วชฺชานํ สํวราย สมฺปรายิกานํ วชฺชานํ ปฏิฆาตาย, อิเม เทฺว อตฺถวเส ปฏิจฺจ ตถาคเตน สาวกานํ สิกฺขาปทํ ปญฺญตฺตํ. เทฺว อตฺถวเส ปฏิจฺจ ตถาคเตน สาวกานํ สิกฺขาปทํ ปญฺญตฺตํ ทิฏฺฐ\-ธมฺมิกานํ ภยานํ สํวราย สมฺปรายิกานํ ภยานํ ปฏิฆาตาย, อิเม เทฺว อตฺถวเส ปฏิจฺจ ตถาคเตน สาวกานํ สิกฺขาปทํ ปญฺญตฺตํ. เทฺว อตฺถวเส ปฏิจฺจ ตถาคเตน สาวกานํ สิกฺขาปทํ ปญฺญตฺตํ ทิฏฺฐ\-ธมฺมิกานํ อกุสลานํ ธมฺมานํ สํวราย สมฺปรายิกานํ อกุสลานํ ธมฺมานํ ปฏิฆาตาย, อิเม เทฺว อตฺถวเส ปฏิจฺจ ตถาคเตน สาวกานํ สิกฺขาปทํ ปญฺญตฺตํ. เทฺว อตฺถวเส ปฏิจฺจ ตถาคเตน สาวกานํ สิกฺขาปทํ ปญฺญตฺตํ คิหีนํ อนุกมฺปาย ปาปิจฺฉานํ ปกฺขุ\-ปจฺเฉทาย, อิเม เทฺว อตฺถวเส ปฏิจฺจ ตถาคเตน สาวกานํ สิกฺขาปทํ ปญฺญตฺตํ. เทฺว อตฺถวเส ปฏิจฺจ ตถาคเตน สาวกานํ สิกฺขาปทํ ปญฺญตฺตํ อปฺปสนฺนานํ ปสาทาย ปสนฺนานํ ภิยฺโยภาวาย, อิเม เทฺว อตฺถวเส ปฏิจฺจ ตถาคเตน สาวกานํ สิกฺขาปทํ ปญฺญตฺตํ. เทฺว อตฺถวเส ปฏิจฺจ ตถาคเตน สาวกานํ สิกฺขาปทํ ปญฺญตฺตํ สทฺธมฺม\-ฏฺฐิติยา วินยา\-นุคฺคหาย, อิเม เทฺว อตฺถวเส ปฏิจฺจ ตถาคเตน สาวกานํ สิกฺขาปทํ ปญฺญตฺตํ.}{อตฺถวสวคฺโค นิฏฺฐิโต ทุติโย.}
\csromanmatikaanchor{mtk.251}
\csromantocmarkref{subparagraph}{3. ปญฺญตฺตวคฺค}{mtk.251}
\bookmark[dest={mtk.251},level=5]{3. ปญฺญตฺตวคฺค}
\csromanheader{6}{3. \textbf{ปญฺญตฺตวคฺค}}
\proseitemwithclosermidruled{499}{เทฺว อตฺถวเส ปฏิจฺจ ตถาคเตน สาวกานํ ปาติโมกฺขํ ปญฺญตฺตํ ฯเปฯ. ปาติโมกฺขุ\-ทฺเทโส ปญฺญตฺโต. ปาติโมกฺข\-ฏฺฐปนํ ปญฺญตฺตํ. ปวารณา ปญฺญตฺตา. ปวารณา\-ฐปนํ ปญฺญตฺตํ. ตชฺชนีย\-กมฺมํ ปญฺญตฺตํ. \csromanfolio{386}นิยสฺสกมฺมํ ปญฺญตฺตํ. ปพฺพาชนีย\-กมฺมํ ปญฺญตฺตํ. ปฏิสารณีย\-กมฺมํ ปญฺญตฺตํ. อุกฺเขปนีย\-กมฺมํ ปญฺญตฺตํ. ปริวาสทานํ ปญฺญตฺตํ. มูลาย\-ปฏิกสฺสนา ปญฺญตฺตา. มานตฺตทานํ ปญฺญตฺตํ. อพฺภานํ ปญฺญตฺตํ. โอสารณียํ ปญฺญตฺตํ. นิสฺสารณียํ ปญฺญตฺตํ. อุปสมฺปทํ ปญฺญตฺตํ. อปโลกน\-กมฺมํ ปญฺญตฺตํ. ญตฺติกมฺมํ ปญฺญตฺตํ. ญตฺติ\-ทุติย\-กมฺมํ ปญฺญตฺตํ. ญตฺติ\-จตุตฺถ\-กมฺมํ ปญฺญตฺตํ ฯเปฯ.}{ปญฺญตฺตวคฺโค นิฏฺฐิโต ตติโย.}
\csromanmatikaanchor{mtk.252}
\csromantocmarkref{subparagraph}{4. อปญฺญตฺเต ปญฺญตฺตวคฺค}{mtk.252}
\bookmark[dest={mtk.252},level=5]{4. อปญฺญตฺเต ปญฺญตฺตวคฺค}
\csromanheader{6}{4. \textbf{อปญฺญตฺเต ปญฺญตฺตวคฺค}}
\prose{500 ฯเปฯ. อปญฺญตฺเต ปญฺญตฺตํ ปญฺญตฺเต อนุปญฺญตฺตํ ฯเปฯ. สมฺมุขาวินโย ปญฺญตฺโต ฯเปฯ. สติวินโย ปญฺญตฺโต ฯเปฯ. อมูฬฺหวินโย ปญฺญตฺโต ฯเปฯ. ปฏิญฺญาต\-กรณํ ปญฺญตฺตํ ฯเปฯ. เยภุยฺยสิกา ปญฺญตฺตา ฯเปฯ. ตสฺสปาปิยสิกา ปญฺญตฺตา ฯเปฯ. ติณวตฺถารโก ปญฺญตฺโต สํฆสุฏฺฐุตาย สํฆผาสุตาย, อิเม เทฺว อตฺถวเส ปฏิจฺจ ตถาคเตน สาวกานํ ติณวตฺถารโก ปญฺญตฺโต. เทฺว อตฺถวเส ปฏิจฺจ ตถาคเตน สาวกานํ ติณวตฺถารโก ปญฺญตฺโต ทุมฺมงฺกูนํ ปุคฺคลานํ นิคฺคหาย เปสลานํ ภิกฺขูนํ ผาสุวิหาราย, อิเม เทฺว อตฺถวเส ปฏิจฺจ ตถาคเตน สาวกานํ ติณวตฺถารโก ปญฺญตฺโต. เทฺว อตฺถวเส ปฏิจฺจ ตถาคเตน สาวกานํ ติณวตฺถารโก ปญฺญตฺโต ทิฏฺฐ\-ธมฺมิกานํ อาสวานํ สํวราย สมฺปรายิกานํ อาสวานํ ปฏิฆาตาย, อิเม เทฺว อตฺถวเส ปฏิจฺจ ตถาคเตน สาวกานํ ติณวตฺถารโก ปญฺญตฺโต. เทฺว อตฺถวเส ปฏิจฺจ ตถาคเตน สาวกานํ ติณวตฺถารโก ปญฺญตฺโต ทิฏฺฐ\-ธมฺมิกานํ เวรานํ สํวราย สมฺปรายิกานํ เวรานํ ปฏิฆาตาย, อิเม เทฺว อตฺถวเส ปฏิจฺจ ตถาคเตน สาวกานํ ติณวตฺถารโก ปญฺญตฺโต. เทฺว อตฺถวเส ปฏิจฺจ ตถาคเตน สาวกานํ ติณวตฺถารโก ปญฺญตฺโต ทิฏฺฐ\-ธมฺมิกานํ วชฺชานํ สํวราย สมฺปรายิกานํ วชฺชานํ ปฏิฆาตาย, อิเม เทฺว อตฺถวเส ปฏิจฺจ ตถาคเตน สาวกานํ ติณวตฺถารโก ปญฺญตฺโต. เทฺว อตฺถวเส ปฏิจฺจ ตถาคเตน สาวกานํ ติณวตฺถารโก ปญฺญตฺโต ทิฏฺฐ\-ธมฺมิกานํ ภยานํ สํวราย สมฺปรายิกานํ ภยานํ ปฏิฆาตาย, อิเม เทฺว อตฺถวเส ปฏิจฺจ ตถาคเตน สาวกานํ ติณวตฺถารโก ปญฺญตฺโต.}

\prosewithclosermidruled{\csromanfolio{387}เทฺว อตฺถวเส ปฏิจฺจ ตถาคเตน สาวกานํ ติณวตฺถารโก ปญฺญตฺโต ทิฏฺฐ\-ธมฺมิกานํ อกุสลานํ ธมฺมานํ สํวราย สมฺปรายิกานํ อกุสลานํ ธมฺมานํ ปฏิฆาตาย, อิเม เทฺว อตฺถวเส ปฏิจฺจ ตถาคเตน สาวกานํ ติณวตฺถารโก ปญฺญตฺโต. เทฺว อตฺถวเส ปฏิจฺจ ตถาคเตน สาวกานํ ติณวตฺถารโก ปญฺญตฺโต คิหีนํ อนุกมฺปาย ปาปิจฺฉานํ ปกฺขุ\-ปจฺเฉทาย, อิเม เทฺว อตฺถวเส ปฏิจฺจ ตถาคเตน สาวกานํ ติณวตฺถารโก ปญฺญตฺโต. เทฺว อตฺถวเส ปฏิจฺจ ตถาคเตน สาวกานํ ติณวตฺถารโก ปญฺญตฺโต อปฺปสนฺนานํ ปสาทาย ปสนฺนานํ ภิยฺโยภาวาย, อิเม เทฺว อตฺถวเส ปฏิจฺจ ตถาคเตน สาวกานํ ติณวตฺถารโก ปญฺญตฺโต. เทฺว อตฺถวเส ปฏิจฺจ ตถาคเตน สาวกานํ ติณวตฺถารโก ปญฺญตฺโต สทฺธมฺม\-ฏฺฐิติยา วินยา\-นุคฺคหาย, อิเม เทฺว อตฺถวเส ปฏิจฺจ ตถาคเตน สาวกานํ ติณวตฺถารโก ปญฺญตฺโต.}{อปญฺญตฺเต ปญฺญตฺตวคฺโค นิฏฺฐิโต จตุตฺโถ.}
\csromanmatikaanchor{mtk.253}
\csromantocmarkref{subparagraph}{5. นวสงฺคหวคฺค}{mtk.253}
\bookmark[dest={mtk.253},level=5]{5. นวสงฺคหวคฺค}
\csromanheader{6}{5. นวสงฺคห\-วคฺค}
\proseitem{501}{นวสงฺคหา วตฺถุสงฺคโห วิปตฺติสงฺคโห อาปตฺติสงฺคโห นิทานสงฺคโห ปุคฺคลสงฺคโห ขนฺธสงฺคโห สมุฏฺฐาน\-สงฺคโห อธิกรณ\-สงฺคโห สมถ\-สงฺค\-โหติ.}

\prose{อธิกรเณ สมุปฺปนฺเน สเจ อุโภ อตฺถ\-ปจฺจตฺถิกา อาคจฺฉนฺติ, อุภินฺนมฺปิ วตฺถุ อาโรจาเปตพฺพํ, อุภินฺนมฺปิ วตฺถุํ อาโรจาเปตฺวา อุภินฺนมฺปิ ปฏิญฺญา โสตพฺพา, อุภินฺนมฺปิ ปฏิญฺญํ สุตฺวา อุโภปิ วตฺตพฺพา “อมฺหากํ อิมสฺมิํ อธิกรเณ วูปสมิเต\csromansharedfootnote{7c7cac40f06cda64}{วูปสเมปิ (ก)} อุโภปิ ตุฏฺฐา ภวิสฺสถา”ติ. สเจ อาหํสุ “อุโภปิ ตุฏฺฐา ภวิสฺสามา”ติ, สํเฆน ตํ อธิกรณํ สมฺปฏิจฺฉิตพฺพํ. สเจ อลชฺชุสฺสนฺนา โหติ ปริสา, อุพฺพาหิกาย วูปสเมตพฺพํ. สเจ พาลุสฺสนฺนา โหติ ปริสา, วินยธโร ปริเยสิตพฺโพ. เยน ธมฺเมน เยน วินเยน เยน สตฺถุสาสเนน ตํ อธิกรณํ วูปสมฺมติ, ตถา ตํ อธิกรฺณํ วูปสเมตพฺพํ.}

\prose{วตฺถุ ชานิตพฺพํ, โคตฺตํ ชานิตพฺพํ, นามํ ชานิตพฺพํ, อาปตฺติ ชานิตพฺพา.}

\prose{\csromanfolio{388}เมถุน\-ธมฺโมติ วตฺถุ เจว โคตฺตญฺจ. ปาราชิกนฺติ นามญฺเจว อาปตฺติ จ.}

\prose{อทินฺนาทานนฺติ วตฺถุ เจว โคตฺตญฺจ. ปาราชิกนฺติ นามญฺเจว อาปตฺติ จ.}

\prose{มนุสฺส\-วิคฺคโหติ วตฺถุ เจว โคตฺตญฺจ. ปาราชิกนฺติ นามญฺเจว อาปตฺติ จ.}

\prose{อุตฺตริ\-มนุสฺสธมฺโมติ วตฺถุ เจว โคตฺตญฺจ. ปาราชิกนฺติ นามญฺเจว อาปตฺติ จ.}

\prose{สุกฺกวิสฺสฏฺฐีติ วตฺถุ เจว โคตฺตญฺจ. สํฆาทิเสโสติ นามญฺเจว อาปตฺติ จ.}

\prose{กายสํสคฺโคติ วตฺถุ เจว โคตฺตญฺจ. สํฆาทิเสโสติ นามญฺเจว อาปตฺติ จ.}

\prose{ทุฏฺฐุลฺล\-วาจาติ วตฺถุ เจว โคตฺตญฺจ. สํฆาทิเสโสติ นามญฺเจว อาปตฺติ จ.}

\prose{อตฺตกามนฺติ วตฺถุ เจว โคตฺตญฺจ. สํฆาทิเสโสติ นามญฺเจว อาปตฺติ จ.}

\prose{สญฺจริตฺตนฺติ วตฺถุ เจว โคตฺตญฺจ. สํฆาทิเสโสติ นามญฺเจว อาปตฺติ จ.}

\prose{สญฺญาจิกาย กุฏิํ การาปนนฺติ วตฺถุ เจว โคตฺตญฺจ. สํฆาทิเสโสติ นามญฺเจว อาปตฺติ จ.}

\prose{มหลฺลกํ วิหารํ การาปนนฺติ วตฺถุ เจว โคตฺตญฺจ. สํฆาทิเสโสติ นามญฺเจว อาปตฺติ จ.}

\prose{ภิกฺขุํ อมูลเกน ปาราชิเกน ธมฺเมน อนุทฺธํส\-นนฺติ วตฺถุ เจว โคตฺตญฺจ. สํฆาทิเสโสติ นามญฺเจว อาปตฺติ จ.}

\prose{ภิกฺขุํ อญฺญภาคิยสฺส อธิกรณสฺส กิญฺจิ เทสํ เลสมตฺตํ อุปาทาย ปาราชิเกน ธมฺเมน อนุทฺธํส\-นนฺติ วตฺถุ เจว โคตฺตญฺจ. สํฆาทิเสโสติ นามญฺเจว อาปตฺติ จ.}

\prose{สํฆเภทกสฺส ภิกฺขุโน ยาวตติยํ สมนุภาสนาย น ปฏินิสฺสชฺชนนฺติ วตฺถุ เจว โคตฺตญฺจ. สํฆาทิเสโสติ นามญฺเจว อาปตฺติ จ.}

\prose{\csromanfolio{389}เภทกา\-นุวตฺตกานํ ภิกฺขูนํ ยาวตติยํ สมนุภาสนาย น ปฏินิสฺสชฺชนนฺติ วตฺถุ เจว โคตฺตญฺจ. สํฆาทิเสโสติ นามญฺเจว อาปตฺติ จ.}

\prose{ทุพฺพจสฺส ภิกฺขุโน ยาวตติยํ สมนุภาสนาย น ปฏินิสฺสชฺชนนฺติ วตฺถุ เจว โคตฺตญฺจ. สํฆาทิเสโสติ นามญฺเจว อาปตฺติ จ.}

\prose{กุลทูสกสฺส ภิกฺขุโน ยาวตติยํ สมนุภาสนาย น ปฏินิสฺสชฺชนนฺติ วตฺถุ เจว โคตฺตญฺจ. สํฆาทิเสโสติ นามญฺเจว อาปตฺติ จ ฯเปฯ.}

\prose{อนาทริยํ ปฏิจฺจ อุทเก อุจฺจารํ วา ปสฺสาวํ วา เขฬํ วา กรณนฺติ วตฺถุ เจว โคตฺตญฺจ. ทุกฺกฏนฺติ นามญฺเจว อาปตฺติ จาติ.}

\vspace{\tipitakaparskip}
\csromancenter{นวสงฺคห\-วคฺโค นิฏฺฐิโต ปญฺจโม.}
\csromancenter{ตสฺสุทฺทานํ.}
\setlength{\csromangathapagewidth}{0pt}
\setlength{\csromangathawakwidth}{0pt}
\csromangathameasuregroup{อปโลกนํ ญตฺติ จ, \\ วตฺถุ ญตฺติ อนุสฺสาวนํ, \\ สมฺมุขา ปฏิปุจฺฉา จ, \\ วตฺถุ สํฆปุคฺคลญฺจ, \\ วตฺถุํ สํฆปุคฺคลญฺจ, \\ อติขุทฺทกา มหนฺตา จ, \\ พหินที สมุทฺเท จ, \\ อชฺโฌตฺถรติ สีมาย, \\ ทส วีสติวคฺคา จ, \\ กมฺมปตฺตา ฉนฺทารหา, \\ อปโลกนํ ปญฺจฏฺฐานํ, \\ ญตฺติ ทุติยํ สตฺตฏฺฐานํ, \\ สุฏฺฐุ ผาสุ จ ทุมฺมงฺกุ, \\ เวรวชฺชภยญฺเจว, \\ ปาปิจฺฉา อปฺปสนฺนานํ, \\ วินยานุคฺคหา เจว, \\ ปาติโมกฺขญฺจ ฐปนา, \\ ตชฺชนีย นิยสฺสญฺจ, \\ อุกฺเขปน ปริวาสํ, \\ โอสารณํ นิสฺสารณํ, \\ อปโลกนญตฺติ จ, \\ อปญฺญตฺเตนุปญฺญตฺตํ, \\ อมูฬฺหปฏิเยภุยฺย, \\ วตฺถุ วิปตฺติ อาปตฺติ, \\ ขนฺธา เจว สมุฏฺฐานา, \\ สมถา สงฺคหา เจว,}{\csromangathabat{อปโลกนํ ญตฺติ จ,}{ทุติยํ จตุตฺเถน จ.} \\ \csromangathabat{วตฺถุ ญตฺติ อนุสฺสาวนํ,}{สีมา ปริสเมว จ.} \\ \csromangathabat{สมฺมุขา ปฏิปุจฺฉา จ,}{ปฏิญฺญา วินยารโห.} \\ \csromangathabat{วตฺถุ สํฆปุคฺคลญฺจ,}{ญตฺติํ น ปจฺฉา ญตฺติ จ.} \\ \csromangathabat{วตฺถุํ สํฆปุคฺคลญฺจ,}{สาวนํ อกาเลน จ.} \\ \csromangathabat{อติขุทฺทกา มหนฺตา จ,}{ขณฺฑจฺฉายา นิมิตฺตกา.} \\ \csromangathabat{พหินที สมุทฺเท จ,}{ชาตสฺสเร จ ภินฺทติ.} \\ \csromangathabat{อชฺโฌตฺถรติ สีมาย,}{จตุ ปญฺจ จ วคฺคิกา.} \\ \csromangathabat{ทส วีสติวคฺคา จ,}{อนาหฏา จ อาหฏา.} \\ \csromangathabat{กมฺมปตฺตา ฉนฺทารหา,}{กมฺมารหา จ ปุคฺคลา.} \\ \csromangathabat{อปโลกนํ ปญฺจฏฺฐานํ,}{ญตฺติ จ นวฐานิกา.} \\ \csromangathabat{ญตฺติ ทุติยํ สตฺตฏฺฐานํ,}{จตุตฺถา สตฺตฐานิกา.} \\ \csromangathabat{สุฏฺฐุ ผาสุ จ ทุมฺมงฺกุ,}{เปสลา จาปิ อาสวา.} \\ \csromangathabat{เวรวชฺชภยญฺเจว,}{อกุสลํ คิหีนญฺจ.} \\ \csromangathabat{ปาปิจฺฉา อปฺปสนฺนานํ,}{ปสนฺนา ธมฺมฏฺฐปนา\textbf{.}} \\ \csromangathabat{วินยานุคฺคหา เจว,}{ปาติโมกฺขุทฺเทเสน จ\textbf{.}} \\ \csromangathabat{ปาติโมกฺขญฺจ ฐปนา,}{ปวารณญฺจ ฐปนํ\textbf{.}} \\ \csromangathabat{ตชฺชนีย นิยสฺสญฺจ,}{ปพฺพาชนีย ปฏิสารณี\textbf{.}} \\ \csromangathabat{อุกฺเขปน ปริวาสํ,}{มูลมานตฺตอพฺภานํ\textbf{.}} \\ \csromangathabat{โอสารณํ นิสฺสารณํ,}{ตเถว อุปสมฺปทา\textbf{.}} \\ \csromangathabat{อปโลกนญตฺติ จ,}{ทุติยญฺจ จตุตฺถกํ\textbf{.}} \\ \csromangathabat{อปญฺญตฺเตนุปญฺญตฺตํ,}{สมฺมุขาวินโย สติ\textbf{.}} \\ \csromangathabat{อมูฬฺหปฏิเยภุยฺย,}{ปาปิย ติณวตฺถารกํ\textbf{.}} \\ \csromangathabat{วตฺถุ วิปตฺติ อาปตฺติ,}{นิทานํ ปุคฺคเลน จ\textbf{.}} \\ \csromangathabat{ขนฺธา เจว สมุฏฺฐานา,}{อธิกรณเมว จ\textbf{.}} \\ \csromangathabat{สมถา สงฺคหา เจว,}{นามอาปตฺติกา ตถาติ\textbf{.}}}
\csromangathasetleft{อปโลกนํ ญตฺติ จ, \\ วตฺถุ ญตฺติ อนุสฺสาวนํ, \\ สมฺมุขา ปฏิปุจฺฉา จ, \\ วตฺถุ สํฆปุคฺคลญฺจ, \\ วตฺถุํ สํฆปุคฺคลญฺจ, \\ อติขุทฺทกา มหนฺตา จ, \\ พหินที สมุทฺเท จ, \\ อชฺโฌตฺถรติ สีมาย, \\ ทส วีสติวคฺคา จ, \\ กมฺมปตฺตา ฉนฺทารหา, \\ อปโลกนํ ปญฺจฏฺฐานํ, \\ ญตฺติ ทุติยํ สตฺตฏฺฐานํ, \\ สุฏฺฐุ ผาสุ จ ทุมฺมงฺกุ, \\ เวรวชฺชภยญฺเจว, \\ ปาปิจฺฉา อปฺปสนฺนานํ, \\ วินยานุคฺคหา เจว, \\ ปาติโมกฺขญฺจ ฐปนา, \\ ตชฺชนีย นิยสฺสญฺจ, \\ อุกฺเขปน ปริวาสํ, \\ โอสารณํ นิสฺสารณํ, \\ อปโลกนญตฺติ จ, \\ อปญฺญตฺเตนุปญฺญตฺตํ, \\ อมูฬฺหปฏิเยภุยฺย, \\ วตฺถุ วิปตฺติ อาปตฺติ, \\ ขนฺธา เจว สมุฏฺฐานา, \\ สมถา สงฺคหา เจว,}
\csromangathagroupwithcloser{\csromangathastanza{\csromangathabat{อปโลกนํ ญตฺติ จ,}{ทุติยํ จตุตฺเถน จ.} \\ \csromangathabat{วตฺถุ ญตฺติ อนุสฺสาวนํ,}{สีมา ปริสเมว จ.}}\csromangathastanza{\csromangathabat{สมฺมุขา ปฏิปุจฺฉา จ,}{ปฏิญฺญา วินยารโห.} \\ \csromangathabat{วตฺถุ สํฆปุคฺคลญฺจ,}{ญตฺติํ น ปจฺฉา ญตฺติ จ.}}\csromangathastanza{\csromangathabat{วตฺถุํ สํฆปุคฺคลญฺจ,}{สาวนํ อกาเลน จ.} \\ \csromangathabat{อติขุทฺทกา มหนฺตา จ,}{ขณฺฑจฺฉายา นิมิตฺตกา.}}\csromangathastanza{\csromangathabat{พหินที สมุทฺเท จ,}{ชาตสฺสเร จ ภินฺทติ.} \\ \csromangathabat{อชฺโฌตฺถรติ สีมาย,}{จตุ ปญฺจ จ วคฺคิกา.}}\csromangathastanza{\csromangathabat{ทส วีสติวคฺคา จ,}{อนาหฏา จ อาหฏา.} \\ \csromangathabat{กมฺมปตฺตา ฉนฺทารหา,}{กมฺมารหา จ ปุคฺคลา.}}\csromangathastanza{\csromangathabat{อปโลกนํ ปญฺจฏฺฐานํ,}{ญตฺติ จ นวฐานิกา.} \\ \csromangathabat{ญตฺติ ทุติยํ สตฺตฏฺฐานํ,}{จตุตฺถา สตฺตฐานิกา.}}\csromangathastanza{\csromangathabat{สุฏฺฐุ ผาสุ จ ทุมฺมงฺกุ,}{เปสลา จาปิ อาสวา.} \\ \csromangathabat{เวรวชฺชภยญฺเจว,}{อกุสลํ คิหีนญฺจ.}}\csromangathastanza{\csromanfolio{390}\csromangathabat{ปาปิจฺฉา อปฺปสนฺนานํ,}{ปสนฺนา ธมฺมฏฺฐปนา\textbf{.}} \\ \csromangathabat{วินยานุคฺคหา เจว,}{ปาติโมกฺขุทฺเทเสน จ\textbf{.}}}\csromangathastanza{\csromangathabat{ปาติโมกฺขญฺจ ฐปนา,}{ปวารณญฺจ ฐปนํ\textbf{.}} \\ \csromangathabat{ตชฺชนีย นิยสฺสญฺจ,}{ปพฺพาชนีย ปฏิสารณี\textbf{.}}}\csromangathastanza{\csromangathabat{อุกฺเขปน ปริวาสํ,}{มูลมานตฺตอพฺภานํ\textbf{.}} \\ \csromangathabat{โอสารณํ นิสฺสารณํ,}{ตเถว อุปสมฺปทา\textbf{.}}}\csromangathastanza{\csromangathabat{อปโลกนญตฺติ จ,}{ทุติยญฺจ จตุตฺถกํ\textbf{.}} \\ \csromangathabat{อปญฺญตฺเตนุปญฺญตฺตํ,}{สมฺมุขาวินโย สติ\textbf{.}}}\csromangathastanza{\csromangathabat{อมูฬฺหปฏิเยภุยฺย,}{ปาปิย ติณวตฺถารกํ\textbf{.}} \\ \csromangathabat{วตฺถุ วิปตฺติ อาปตฺติ,}{นิทานํ ปุคฺคเลน จ\textbf{.}}}\csromangathastanza{\csromangathabat{ขนฺธา เจว สมุฏฺฐานา,}{อธิกรณเมว จ\textbf{.}} \\ \csromangathabat{สมถา สงฺคหา เจว,}{นามอาปตฺติกา ตถาติ\textbf{.}}}}{ปริวาโร นิฏฺฐิโต.}
\nitthitammajor{\textbf{ปริวารปาฬิ นิฏฺฐิตา.}}
